\IfFileExists{stacks-project.cls}{%
\documentclass{stacks-project}
}{%
\documentclass{amsart}
}

% For dealing with references we use the comment environment
\usepackage{verbatim}
\newenvironment{reference}{\comment}{\endcomment}
%\newenvironment{reference}{}{}
\newenvironment{slogan}{\comment}{\endcomment}
\newenvironment{history}{\comment}{\endcomment}

% For commutative diagrams we use Xy-pic
\usepackage[all]{xy}

% We use 2cell for 2-commutative diagrams.
\xyoption{2cell}
\UseAllTwocells

% We use multicol for the list of chapters between chapters
\usepackage{multicol}

% This is generally recommended for better output.  The normal Stacks build
% remains pdfTeX/T1; localized CJK editions opt into XeTeX by defining
% \GAGACJKMainFont before loading this preamble.
\usepackage{iftex}
\ifPDFTeX
  \usepackage{lmodern}
  \usepackage[T1]{fontenc}
\else
  \usepackage{fontspec}
  \usepackage{xeCJK}
  \setmainfont{Latin Modern Roman}
  \setsansfont{Latin Modern Sans}
  \setmonofont{Latin Modern Mono}
  \ifdefined\GAGACJKMainFont
    \setCJKmainfont{\GAGACJKMainFont}
  \else
    \setCJKmainfont{Yu Gothic}
  \fi
\fi

% For cross-file-references
\usepackage{xr-hyper}

% Package for hypertext links:
\usepackage{hyperref}

% For any local file, say "hello.tex" you want to link to please
% use \externaldocument[hello-]{hello}
\externaldocument[introduction-]{introduction}
\externaldocument[conventions-]{conventions}
\externaldocument[sets-]{sets}
\externaldocument[categories-]{categories}
\externaldocument[topology-]{topology}
\externaldocument[sheaves-]{sheaves}
\externaldocument[sites-]{sites}
\externaldocument[stacks-]{stacks}
\externaldocument[fields-]{fields}
\externaldocument[algebra-]{algebra}
\externaldocument[brauer-]{brauer}
\externaldocument[homology-]{homology}
\externaldocument[derived-]{derived}
\externaldocument[simplicial-]{simplicial}
\externaldocument[more-algebra-]{more-algebra}
\externaldocument[smoothing-]{smoothing}
\externaldocument[modules-]{modules}
\externaldocument[sites-modules-]{sites-modules}
\externaldocument[injectives-]{injectives}
\externaldocument[cohomology-]{cohomology}
\externaldocument[sites-cohomology-]{sites-cohomology}
\externaldocument[dga-]{dga}
\externaldocument[dpa-]{dpa}
\externaldocument[sdga-]{sdga}
\externaldocument[hypercovering-]{hypercovering}
\externaldocument[schemes-]{schemes}
\externaldocument[constructions-]{constructions}
\externaldocument[properties-]{properties}
\externaldocument[morphisms-]{morphisms}
\externaldocument[coherent-]{coherent}
\externaldocument[divisors-]{divisors}
\externaldocument[limits-]{limits}
\externaldocument[varieties-]{varieties}
\externaldocument[topologies-]{topologies}
\externaldocument[descent-]{descent}
\externaldocument[perfect-]{perfect}
\externaldocument[more-morphisms-]{more-morphisms}
\externaldocument[flat-]{flat}
\externaldocument[groupoids-]{groupoids}
\externaldocument[more-groupoids-]{more-groupoids}
\externaldocument[etale-]{etale}
\externaldocument[chow-]{chow}
\externaldocument[intersection-]{intersection}
\externaldocument[pic-]{pic}
\externaldocument[weil-]{weil}
\externaldocument[adequate-]{adequate}
\externaldocument[dualizing-]{dualizing}
\externaldocument[duality-]{duality}
\externaldocument[discriminant-]{discriminant}
\externaldocument[derham-]{derham}
\externaldocument[local-cohomology-]{local-cohomology}
\externaldocument[algebraization-]{algebraization}
\externaldocument[curves-]{curves}
\externaldocument[resolve-]{resolve}
\externaldocument[models-]{models}
\externaldocument[functors-]{functors}
\externaldocument[equiv-]{equiv}
\externaldocument[gaga-]{gaga}
\externaldocument[pione-]{pione}
\externaldocument[etale-cohomology-]{etale-cohomology}
\externaldocument[proetale-]{proetale}
\externaldocument[relative-cycles-]{relative-cycles}
\externaldocument[more-etale-]{more-etale}
\externaldocument[trace-]{trace}
\externaldocument[crystalline-]{crystalline}
\externaldocument[spaces-]{spaces}
\externaldocument[spaces-properties-]{spaces-properties}
\externaldocument[spaces-morphisms-]{spaces-morphisms}
\externaldocument[decent-spaces-]{decent-spaces}
\externaldocument[spaces-cohomology-]{spaces-cohomology}
\externaldocument[spaces-limits-]{spaces-limits}
\externaldocument[spaces-divisors-]{spaces-divisors}
\externaldocument[spaces-over-fields-]{spaces-over-fields}
\externaldocument[spaces-topologies-]{spaces-topologies}
\externaldocument[spaces-descent-]{spaces-descent}
\externaldocument[spaces-perfect-]{spaces-perfect}
\externaldocument[spaces-more-morphisms-]{spaces-more-morphisms}
\externaldocument[spaces-flat-]{spaces-flat}
\externaldocument[spaces-groupoids-]{spaces-groupoids}
\externaldocument[spaces-more-groupoids-]{spaces-more-groupoids}
\externaldocument[bootstrap-]{bootstrap}
\externaldocument[spaces-pushouts-]{spaces-pushouts}
\externaldocument[spaces-chow-]{spaces-chow}
\externaldocument[groupoids-quotients-]{groupoids-quotients}
\externaldocument[spaces-more-cohomology-]{spaces-more-cohomology}
\externaldocument[spaces-simplicial-]{spaces-simplicial}
\externaldocument[spaces-duality-]{spaces-duality}
\externaldocument[formal-spaces-]{formal-spaces}
\externaldocument[restricted-]{restricted}
\externaldocument[spaces-resolve-]{spaces-resolve}
\externaldocument[formal-defos-]{formal-defos}
\externaldocument[defos-]{defos}
\externaldocument[cotangent-]{cotangent}
\externaldocument[examples-defos-]{examples-defos}
\externaldocument[algebraic-]{algebraic}
\externaldocument[examples-stacks-]{examples-stacks}
\externaldocument[stacks-sheaves-]{stacks-sheaves}
\externaldocument[criteria-]{criteria}
\externaldocument[artin-]{artin}
\externaldocument[quot-]{quot}
\externaldocument[stacks-properties-]{stacks-properties}
\externaldocument[stacks-morphisms-]{stacks-morphisms}
\externaldocument[stacks-limits-]{stacks-limits}
\externaldocument[stacks-cohomology-]{stacks-cohomology}
\externaldocument[stacks-perfect-]{stacks-perfect}
\externaldocument[stacks-introduction-]{stacks-introduction}
\externaldocument[stacks-more-morphisms-]{stacks-more-morphisms}
\externaldocument[stacks-geometry-]{stacks-geometry}
\externaldocument[moduli-]{moduli}
\externaldocument[moduli-curves-]{moduli-curves}
\externaldocument[examples-]{examples}
\externaldocument[exercises-]{exercises}
\externaldocument[guide-]{guide}
\externaldocument[desirables-]{desirables}
\externaldocument[coding-]{coding}
\externaldocument[obsolete-]{obsolete}
\externaldocument[fdl-]{fdl}
\externaldocument[index-]{index}

% Theorem environments.
%
\theoremstyle{plain}
\newtheorem{theorem}[subsection]{Theorem}
\newtheorem{proposition}[subsection]{Proposition}
\newtheorem{lemma}[subsection]{Lemma}

\theoremstyle{definition}
\newtheorem{definition}[subsection]{Definition}
\newtheorem{example}[subsection]{Example}
\newtheorem{exercise}[subsection]{Exercise}
\newtheorem{situation}[subsection]{Situation}

\theoremstyle{remark}
\newtheorem{remark}[subsection]{Remark}
\newtheorem{remarks}[subsection]{Remarks}

\numberwithin{equation}{subsection}

% Macros
%
\def\lim{\mathop{\mathrm{lim}}\nolimits}
\def\colim{\mathop{\mathrm{colim}}\nolimits}
\def\Spec{\mathop{\mathrm{Spec}}}
\def\Hom{\mathop{\mathrm{Hom}}\nolimits}
\def\Ext{\mathop{\mathrm{Ext}}\nolimits}
\def\SheafHom{\mathop{\mathcal{H}\!\mathit{om}}\nolimits}
\def\SheafExt{\mathop{\mathcal{E}\!\mathit{xt}}\nolimits}
\def\Sch{\mathit{Sch}}
\def\Mor{\mathop{\mathrm{Mor}}\nolimits}
\def\Ob{\mathop{\mathrm{Ob}}\nolimits}
\def\Sh{\mathop{\mathit{Sh}}\nolimits}
\def\NL{\mathop{N\!L}\nolimits}
\def\CH{\mathop{\mathrm{CH}}\nolimits}
\def\proetale{{pro\text{-}\acute{e}tale}}
\def\etale{{\acute{e}tale}}
\def\QCoh{\mathit{QCoh}}
\def\Ker{\mathop{\mathrm{Ker}}}
\def\Im{\mathop{\mathrm{Im}}}
\def\Coker{\mathop{\mathrm{Coker}}}
\def\Coim{\mathop{\mathrm{Coim}}}

% Boxtimes
%
\DeclareMathSymbol{\boxtimes}{\mathbin}{AMSa}{"02}

%
% Macros for moduli stacks/spaces
%
\def\QCohstack{\mathcal{QC}\!\mathit{oh}}
\def\Cohstack{\mathcal{C}\!\mathit{oh}}
\def\Spacesstack{\mathcal{S}\!\mathit{paces}}
\def\Quotfunctor{\mathrm{Quot}}
\def\Hilbfunctor{\mathrm{Hilb}}
\def\Curvesstack{\mathcal{C}\!\mathit{urves}}
\def\Polarizedstack{\mathcal{P}\!\mathit{olarized}}
\def\Complexesstack{\mathcal{C}\!\mathit{omplexes}}
% \Pic is the operator that assigns to X its picard group, usage \Pic(X)
% \Picardstack_{X/B} denotes the Picard stack of X over B
% \Picardfunctor_{X/B} denotes the Picard functor of X over B
\def\Pic{\mathop{\mathrm{Pic}}\nolimits}
\def\Picardstack{\mathcal{P}\!\mathit{ic}}
\def\Picardfunctor{\mathrm{Pic}}
\def\Deformationcategory{\mathcal{D}\!\mathit{ef}}

% Standalone editorial-package reference adapter; cumulative source is unchanged.
\InputIfFileExists{external-reference-locators.tex}{}{}

% AMS symbols required by the reviewed strict-inclusion correction.
\usepackage{amssymb}
\usepackage{graphicx}
\usepackage{longtable,booktabs,array,calc}
\usepackage{unicode-math}
\setmainfont{Cambria}
\setmathfont{Cambria Math}

% Preserve the original S_\mathfrak q syntax with unicode-math.
% The wrapper passes exactly the original argument to the original alphabet.
\let\AlgebraOriginalMathfrak\mathfrak
\def\mathfrak#1{{\AlgebraOriginalMathfrak{#1}}}
\providecommand{\tightlist}{\setlength{\itemsep}{0pt}\setlength{\parskip}{0pt}}
\providecommand{\pandocbounded}[1]{#1}
\begin{document}
\title{Commutative Algebra: Conceptual Errata and Supplementary Proofs}
\maketitle
\noindent
This companion identifies editorial corrections, complete supplementary
proofs and stronger consequences alongside their original source passages.
The Stacks Project authors remain the authors of the cited source.
Editorial review and proof work: OpenAI Codex---GPT-6 Astra, Ultra effort.
No human review, novelty or official endorsement is claimed.
The original chapter and its source records accompany this document.
\tableofcontents
\section{Source passages and editorial destinations}
This index connects the original chapter to the complete editorial arguments. A review record can preserve a valid statement, correct an error, supply a missing argument, or establish a consequence; inclusion does not by itself classify the source as erroneous. Source-line links refer to the frozen original edition. The review ledger retains the individual decisions and their evidence.

\subsection{Colimits}
\hypertarget{algebra-source-section-colimits}{}
\noindent\href{https://github.com/stacks/stacks-project/blob/a04446e57ec1fbc252a871afcec7752fb2807b14/algebra.tex\#L702}{Original source section}. \href{algebra.pdf\#nameddest=section-colimits}{Accompanying chapter}.

\begin{description}
\item[ALGEBRA-RECON-005] Source lines 1000--1031: \hyperlink{algebra-context-001}{Context 1}.
\end{description}

\subsection{Localization}
\hypertarget{algebra-source-section-localization}{}
\noindent\href{https://github.com/stacks/stacks-project/blob/a04446e57ec1fbc252a871afcec7752fb2807b14/algebra.tex\#L1038}{Original source section}. \href{algebra.pdf\#nameddest=section-localization}{Accompanying chapter}.

\begin{description}
\item[ALGEBRA-RECON-017] Source lines 1420--1433: \hyperlink{algebra-context-003}{Context 3}.
\item[ALGEBRA-RECON-018] Source lines 1170--1211, 1307--1324: \hyperlink{algebra-context-002}{Context 2}.
\end{description}

\subsection{Characterizing finite and finitely presented modules}
\hypertarget{algebra-source-section-colim-and-hom}{}
\noindent\href{https://github.com/stacks/stacks-project/blob/a04446e57ec1fbc252a871afcec7752fb2807b14/algebra.tex\#L1580}{Original source section}. \href{algebra.pdf\#nameddest=section-colim-and-hom}{Accompanying chapter}.

\begin{description}
\item[ALGEBRA-RECON-020] Source lines 1608--1614: \hyperlink{algebra-context-004}{Context 4}.
\end{description}

\subsection{Tensor products}
\hypertarget{algebra-source-section-tensor-product}{}
\noindent\href{https://github.com/stacks/stacks-project/blob/a04446e57ec1fbc252a871afcec7752fb2807b14/algebra.tex\#L1702}{Original source section}. \href{algebra.pdf\#nameddest=section-tensor-product}{Accompanying chapter}.

\begin{description}
\item[ALGEBRA-RECON-021] Source lines 1743--1767: \hyperlink{algebra-context-005}{Context 5}.
\item[ALGEBRA-RECON-031] Source lines 1933--1950: \hyperlink{algebra-context-006}{Context 6}.
\item[ALGEBRA-RECON-032] Source lines 2098--2121: \hyperlink{algebra-context-007}{Context 7}.
\end{description}

\subsection{Tensor algebra}
\hypertarget{algebra-source-section-tensor-algebra}{}
\noindent\href{https://github.com/stacks/stacks-project/blob/a04446e57ec1fbc252a871afcec7752fb2807b14/algebra.tex\#L2164}{Original source section}. \href{algebra.pdf\#nameddest=section-tensor-algebra}{Accompanying chapter}.

\begin{description}
\item[ALGEBRA-RECON-033] Source lines 2240--2268: \hyperlink{algebra-context-008}{Context 8}.
\end{description}

\subsection{Cayley-Hamilton}
\hypertarget{algebra-source-section-cayley-hamilton}{}
\noindent\href{https://github.com/stacks/stacks-project/blob/a04446e57ec1fbc252a871afcec7752fb2807b14/algebra.tex\#L2806}{Original source section}. \href{algebra.pdf\#nameddest=section-cayley-hamilton}{Accompanying chapter}.

\begin{description}
\item[ALGEBRA-RECON-040] Source lines 2834--2887: \hyperlink{algebra-context-009}{Context 9}; \hyperlink{algebra-context-010}{Context 10}.
\item[ALGEBRA-RECON-041] Source lines 2892--2912: \hyperlink{algebra-context-011}{Context 11}.
\end{description}

\subsection{The spectrum of a ring}
\hypertarget{algebra-source-section-spectrum-ring}{}
\noindent\href{https://github.com/stacks/stacks-project/blob/a04446e57ec1fbc252a871afcec7752fb2807b14/algebra.tex\#L2964}{Original source section}. \href{algebra.pdf\#nameddest=section-spectrum-ring}{Accompanying chapter}.

\begin{description}
\item[ALGEBRA-RECON-048] Source lines 3029--3050: \hyperlink{algebra-context-012}{Context 12}.
\item[ALGEBRA-RECON-049] Source lines 3059--3068: \hyperlink{algebra-context-012}{Context 12}.
\item[ALGEBRA-RECON-050] Source lines 3014--3016, 3088--3091: \hyperlink{algebra-context-012}{Context 12}.
\item[ALGEBRA-RECON-051] Source lines 3177--3204: \hyperlink{algebra-context-013}{Context 13}.
\end{description}

\subsection{Local rings}
\hypertarget{algebra-source-section-local-rings}{}
\noindent\href{https://github.com/stacks/stacks-project/blob/a04446e57ec1fbc252a871afcec7752fb2807b14/algebra.tex\#L3305}{Original source section}. \href{algebra.pdf\#nameddest=section-local-rings}{Accompanying chapter}.

\begin{description}
\item[ALGEBRA-RECON-052] Source lines 3366--3397: \hyperlink{algebra-context-014}{Context 14}.
\item[ALGEBRA-RECON-057] Source lines 3500--3517: \hyperlink{algebra-context-015}{Context 15}.
\item[ALGEBRA-RECON-058] Source lines 3520--3538: \hyperlink{algebra-context-016}{Context 16}.
\end{description}

\subsection{The Jacobson radical of a ring}
\hypertarget{algebra-source-section-radical}{}
\noindent\href{https://github.com/stacks/stacks-project/blob/a04446e57ec1fbc252a871afcec7752fb2807b14/algebra.tex\#L3546}{Original source section}. \href{algebra.pdf\#nameddest=section-radical}{Accompanying chapter}.

\begin{description}
\item[ALGEBRA-RECON-059] Source lines 3586--3601: \hyperlink{algebra-context-017}{Context 17}.
\end{description}

\subsection{Nakayama's lemma}
\hypertarget{algebra-source-section-nakayama}{}
\noindent\href{https://github.com/stacks/stacks-project/blob/a04446e57ec1fbc252a871afcec7752fb2807b14/algebra.tex\#L3606}{Original source section}. \href{algebra.pdf\#nameddest=section-nakayama}{Accompanying chapter}.

\begin{description}
\item[ALGEBRA-RECON-060] Source lines 3714--3733: \hyperlink{algebra-context-018}{Context 18}.
\end{description}

\subsection{Connected components of spectra}
\hypertarget{algebra-source-section-connected-components}{}
\noindent\href{https://github.com/stacks/stacks-project/blob/a04446e57ec1fbc252a871afcec7752fb2807b14/algebra.tex\#L3958}{Original source section}. \href{algebra.pdf\#nameddest=section-connected-components}{Accompanying chapter}.

\begin{description}
\item[ALGEBRA-RECON-064] Source lines 4007--4031: \hyperlink{algebra-context-019}{Context 19}.
\end{description}

\subsection{Glueing properties}
\hypertarget{algebra-source-section-more-glueing}{}
\noindent\href{https://github.com/stacks/stacks-project/blob/a04446e57ec1fbc252a871afcec7752fb2807b14/algebra.tex\#L4041}{Original source section}. \href{algebra.pdf\#nameddest=section-more-glueing}{Accompanying chapter}.

\begin{description}
\item[ALGEBRA-RECON-065] Source lines 4167--4176, 4199--4209: \hyperlink{algebra-context-020}{Context 20}.
\end{description}

\subsection{Glueing functions}
\hypertarget{algebra-source-section-tilde-module-sheaf}{}
\noindent\href{https://github.com/stacks/stacks-project/blob/a04446e57ec1fbc252a871afcec7752fb2807b14/algebra.tex\#L4284}{Original source section}. \href{algebra.pdf\#nameddest=section-tilde-module-sheaf}{Accompanying chapter}.

\begin{description}
\item[ALGEBRA-RECON-066] Source lines 4306--4351: \hyperlink{algebra-context-021}{Context 21}.
\item[ALGEBRA-RECON-070] Source lines 4438--4523: \hyperlink{algebra-context-022}{Context 22}.
\end{description}

\subsection{Zerodivisors and total rings of fractions}
\hypertarget{algebra-source-section-total-quotient-ring}{}
\noindent\href{https://github.com/stacks/stacks-project/blob/a04446e57ec1fbc252a871afcec7752fb2807b14/algebra.tex\#L4536}{Original source section}. \href{algebra.pdf\#nameddest=section-total-quotient-ring}{Accompanying chapter}.

\begin{description}
\item[ALGEBRA-RECON-071] Source lines 4592--4604: \hyperlink{algebra-context-023}{Context 23}.
\end{description}

\subsection{Irreducible components of spectra}
\hypertarget{algebra-source-section-irreducible}{}
\noindent\href{https://github.com/stacks/stacks-project/blob/a04446e57ec1fbc252a871afcec7752fb2807b14/algebra.tex\#L4656}{Original source section}. \href{algebra.pdf\#nameddest=section-irreducible}{Accompanying chapter}.

\begin{description}
\item[ALGEBRA-RECON-072] Source lines 4688--4695: \hyperlink{algebra-context-024}{Context 24}.
\end{description}

\subsection{Examples of spectra of rings}
\hypertarget{algebra-source-section-examples-spectra}{}
\noindent\href{https://github.com/stacks/stacks-project/blob/a04446e57ec1fbc252a871afcec7752fb2807b14/algebra.tex\#L4816}{Original source section}. \href{algebra.pdf\#nameddest=section-examples-spectra}{Accompanying chapter}.

\begin{description}
\item[ALGEBRA-RECON-074] Source lines 4869--4876: \hyperlink{algebra-context-025}{Context 25}.
\item[ALGEBRA-RECON-077] Source lines 4877--4886: \hyperlink{algebra-context-025}{Context 25}.
\item[ALGEBRA-RECON-078] Source lines 4899--4928: \hyperlink{algebra-context-026}{Context 26}.
\item[ALGEBRA-RECON-079] Source lines 4941--4950: \hyperlink{algebra-context-027}{Context 27}.
\item[ALGEBRA-RECON-087] Source lines 5061--5070: \hyperlink{algebra-context-027}{Context 27}.
\item[ALGEBRA-RECON-088] Source lines 5045--5049: \hyperlink{algebra-context-027}{Context 27}.
\item[ALGEBRA-RECON-089] Source lines 4973--5055: \hyperlink{algebra-context-027}{Context 27}.
\item[ALGEBRA-RECON-090] Source lines 5073--5082: \hyperlink{algebra-context-027}{Context 27}.
\end{description}

\subsection{A meta-observation about prime ideals}
\hypertarget{algebra-source-section-oka-families}{}
\noindent\href{https://github.com/stacks/stacks-project/blob/a04446e57ec1fbc252a871afcec7752fb2807b14/algebra.tex\#L5089}{Original source section}. \href{algebra.pdf\#nameddest=section-oka-families}{Accompanying chapter}.

\begin{description}
\item[ALGEBRA-RECON-093] Source lines 5182--5190: \hyperlink{algebra-context-028}{Context 28}.
\item[ALGEBRA-RECON-094] Source lines 5254--5261, 5279--5285: \hyperlink{algebra-context-029}{Context 29}; \hyperlink{algebra-context-030}{Context 30}.
\item[ALGEBRA-RECON-095] Source lines 5299--5306: \hyperlink{algebra-context-031}{Context 31}.
\item[ALGEBRA-RECON-096] Source lines 5310--5327: \hyperlink{algebra-context-032}{Context 32}.
\item[ALGEBRA-RECON-097] Source lines 5334--5362: \hyperlink{algebra-context-033}{Context 33}.
\end{description}

\subsection{Images of ring maps of finite presentation}
\hypertarget{algebra-source-section-images-finite-presentation}{}
\noindent\href{https://github.com/stacks/stacks-project/blob/a04446e57ec1fbc252a871afcec7752fb2807b14/algebra.tex\#L5371}{Original source section}. \href{algebra.pdf\#nameddest=section-images-finite-presentation}{Accompanying chapter}.

\begin{description}
\item[ALGEBRA-RECON-100] Source lines 5444--5450: \hyperlink{algebra-context-034}{Context 34}.
\item[ALGEBRA-RECON-101] Source lines 5461--5472: \hyperlink{algebra-context-035}{Context 35}.
\item[ALGEBRA-RECON-103] Source lines 5554--5595: \hyperlink{algebra-context-036}{Context 36}.
\item[ALGEBRA-RECON-105] Source lines 5607--5612: \hyperlink{algebra-context-037}{Context 37}.
\item[ALGEBRA-RECON-107] Source lines 5668--5678: \hyperlink{algebra-context-038}{Context 38}.
\item[ALGEBRA-RECON-108] Source lines 5687--5708: \hyperlink{algebra-context-038}{Context 38}.
\end{description}

\subsection{More on images}
\hypertarget{algebra-source-section-more-images}{}
\noindent\href{https://github.com/stacks/stacks-project/blob/a04446e57ec1fbc252a871afcec7752fb2807b14/algebra.tex\#L5721}{Original source section}. \href{algebra.pdf\#nameddest=section-more-images}{Accompanying chapter}.

\begin{description}
\item[ALGEBRA-RECON-109] Source lines 5738--5766: \hyperlink{algebra-context-039}{Context 39}.
\item[ALGEBRA-RECON-111] Source lines 5796--5813: \hyperlink{algebra-context-040}{Context 40}.
\item[ALGEBRA-RECON-112] Source lines 5834--5868: \hyperlink{algebra-context-041}{Context 41}.
\item[ALGEBRA-RECON-114] Source lines 5933--5943: \hyperlink{algebra-context-042}{Context 42}.
\item[ALGEBRA-RECON-115] Source lines 5963--5974: \hyperlink{algebra-context-043}{Context 43}.
\end{description}

\subsection{Noetherian rings}
\hypertarget{algebra-source-section-Noetherian}{}
\noindent\href{https://github.com/stacks/stacks-project/blob/a04446e57ec1fbc252a871afcec7752fb2807b14/algebra.tex\#L5980}{Original source section}. \href{algebra.pdf\#nameddest=section-Noetherian}{Accompanying chapter}.

\begin{description}
\item[ALGEBRA-RECON-116] Source lines 5984--5988: \hyperlink{algebra-context-044}{Context 44}.
\item[ALGEBRA-RECON-117] Source lines 6008--6024: \hyperlink{algebra-context-044}{Context 44}.
\item[ALGEBRA-RECON-119] Source lines 6034--6066: \hyperlink{algebra-context-044}{Context 44}.
\item[ALGEBRA-RECON-121] Source lines 6188--6215: \hyperlink{algebra-context-044}{Context 44}.
\item[ALGEBRA-RECON-122] Source lines 6224--6229: \hyperlink{algebra-context-044}{Context 44}.
\end{description}

\subsection{Locally nilpotent ideals}
\hypertarget{algebra-source-section-locally-nilpotent}{}
\noindent\href{https://github.com/stacks/stacks-project/blob/a04446e57ec1fbc252a871afcec7752fb2807b14/algebra.tex\#L6238}{Original source section}. \href{algebra.pdf\#nameddest=section-locally-nilpotent}{Accompanying chapter}.

\begin{description}
\item[ALGEBRA-RECON-123] Source lines 6255--6264: \hyperlink{algebra-context-045}{Context 45}.
\item[ALGEBRA-RECON-124] Source lines 6274--6276: \hyperlink{algebra-context-046}{Context 46}.
\item[ALGEBRA-RECON-125] Source lines 6291--6296: \hyperlink{algebra-context-047}{Context 47}.
\item[ALGEBRA-RECON-126] Source lines 6305--6315: \hyperlink{algebra-context-048}{Context 48}.
\item[ALGEBRA-RECON-127] Source lines 6343--6360: \hyperlink{algebra-context-049}{Context 49}.
\item[ALGEBRA-RECON-128] Source lines 6395--6398: \hyperlink{algebra-context-050}{Context 50}.
\item[ALGEBRA-RECON-130] Source lines 6403--6437: \hyperlink{algebra-context-051}{Context 51}.
\end{description}

\subsection{Curiosity}
\hypertarget{algebra-source-section-curiosity}{}
\noindent\href{https://github.com/stacks/stacks-project/blob/a04446e57ec1fbc252a871afcec7752fb2807b14/algebra.tex\#L6445}{Original source section}. \href{algebra.pdf\#nameddest=section-curiosity}{Accompanying chapter}.

\begin{description}
\item[ALGEBRA-RECON-131] Source lines 6462--6476: \hyperlink{algebra-context-052}{Context 52}.
\item[ALGEBRA-RECON-132] Source lines 6481--6504: \hyperlink{algebra-context-053}{Context 53}.
\end{description}

\subsection{Hilbert Nullstellensatz}
\hypertarget{algebra-source-section-nullstellensatz}{}
\noindent\href{https://github.com/stacks/stacks-project/blob/a04446e57ec1fbc252a871afcec7752fb2807b14/algebra.tex\#L6520}{Original source section}. \href{algebra.pdf\#nameddest=section-nullstellensatz}{Accompanying chapter}.

\begin{description}
\item[ALGEBRA-RECON-133] Source lines 6550--6583: \hyperlink{algebra-context-054}{Context 54}.
\item[ALGEBRA-RECON-134] Source lines 6603--6606: \hyperlink{algebra-context-054}{Context 54}.
\item[ALGEBRA-RECON-135] Source lines 6623--6626: \hyperlink{algebra-context-055}{Context 55}.
\end{description}

\subsection{Jacobson rings}
\hypertarget{algebra-source-section-ring-jacobson}{}
\noindent\href{https://github.com/stacks/stacks-project/blob/a04446e57ec1fbc252a871afcec7752fb2807b14/algebra.tex\#L6649}{Original source section}. \href{algebra.pdf\#nameddest=section-ring-jacobson}{Accompanying chapter}.

\begin{description}
\item[ALGEBRA-RECON-136] Source lines 6729--6775: \hyperlink{algebra-context-056}{Context 56}.
\item[ALGEBRA-RECON-137] Source lines 6781--6806: \hyperlink{algebra-context-057}{Context 57}.
\item[ALGEBRA-RECON-138] Source lines 6814--6818: \hyperlink{algebra-context-058}{Context 58}.
\item[ALGEBRA-RECON-139] Source lines 6834--6841: \hyperlink{algebra-context-059}{Context 59}.
\item[ALGEBRA-RECON-140] Source lines 6882--6885: \hyperlink{algebra-context-060}{Context 60}.
\item[ALGEBRA-RECON-141] Source lines 6904--6917: \hyperlink{algebra-context-061}{Context 61}.
\item[ALGEBRA-RECON-142] Source lines 6936--6942: \hyperlink{algebra-context-061}{Context 61}.
\item[ALGEBRA-RECON-143] Source lines 6947--6962: \hyperlink{algebra-context-062}{Context 62}.
\item[ALGEBRA-RECON-144] Source lines 6970--6977: \hyperlink{algebra-context-063}{Context 63}.
\item[ALGEBRA-RECON-146] Source lines 7065--7086: \hyperlink{algebra-context-064}{Context 64}.
\item[ALGEBRA-RECON-148] Source lines 7149--7172: \hyperlink{algebra-context-065}{Context 65}.
\item[ALGEBRA-RECON-152] Source lines 7189--7214: \hyperlink{algebra-context-066}{Context 66}.
\item[ALGEBRA-RECON-153] Source lines 7249--7252: \hyperlink{algebra-context-067}{Context 67}.
\item[ALGEBRA-RECON-156] Source lines 7323--7326: \hyperlink{algebra-context-067}{Context 67}.
\item[ALGEBRA-RECON-157] Source lines 7277--7360: \hyperlink{algebra-context-067}{Context 67}.
\item[ALGEBRA-RECON-159] Source lines 7373--7379: \hyperlink{algebra-context-068}{Context 68}.
\item[ALGEBRA-RECON-161] Source lines 7392--7395: \hyperlink{algebra-context-068}{Context 68}.
\item[ALGEBRA-RECON-162] Source lines 7364--7396: \hyperlink{algebra-context-068}{Context 68}.
\end{description}

\subsection{Finite and integral ring extensions}
\hypertarget{algebra-source-section-finite-ring-extensions}{}
\noindent\href{https://github.com/stacks/stacks-project/blob/a04446e57ec1fbc252a871afcec7752fb2807b14/algebra.tex\#L7434}{Original source section}. \href{algebra.pdf\#nameddest=section-finite-ring-extensions}{Accompanying chapter}.

\begin{description}
\item[ALGEBRA-RECON-163] Source lines 7457--7466: \hyperlink{algebra-context-069}{Context 69}.
\item[ALGEBRA-RECON-164] Source lines 7491--7499: \hyperlink{algebra-context-070}{Context 70}.
\item[ALGEBRA-RECON-165] Source lines 7513--7515: \hyperlink{algebra-context-071}{Context 71}.
\item[ALGEBRA-RECON-166] Source lines 7552--7553: \hyperlink{algebra-context-072}{Context 72}.
\item[ALGEBRA-RECON-167] Source lines 7558--7565: \hyperlink{algebra-context-073}{Context 73}.
\item[ALGEBRA-RECON-168] Source lines 7568--7593: \hyperlink{algebra-context-074}{Context 74}.
\item[ALGEBRA-RECON-169] Source lines 7614--7626: \hyperlink{algebra-context-075}{Context 75}.
\item[ALGEBRA-RECON-170] Source lines 7638--7661: \hyperlink{algebra-context-076}{Context 76}.
\item[ALGEBRA-RECON-171] Source lines 7678--7685: \hyperlink{algebra-context-077}{Context 77}.
\item[ALGEBRA-RECON-172] Source lines 7699--7706: \hyperlink{algebra-context-078}{Context 78}.
\item[ALGEBRA-RECON-173] Source lines 7711--7719: \hyperlink{algebra-context-079}{Context 79}.
\item[ALGEBRA-RECON-174] Source lines 7724--7731: \hyperlink{algebra-context-080}{Context 80}.
\item[ALGEBRA-RECON-175] Source lines 7734--7762: \hyperlink{algebra-context-081}{Context 81}.
\item[ALGEBRA-RECON-176] Source lines 7764--7779: \hyperlink{algebra-context-082}{Context 82}.
\item[ALGEBRA-RECON-178] Source lines 7797--7809: \hyperlink{algebra-context-083}{Context 83}.
\item[ALGEBRA-RECON-179] Source lines 7812--7829: \hyperlink{algebra-context-084}{Context 84}.
\item[ALGEBRA-RECON-181] Source lines 7831--7850: \hyperlink{algebra-context-085}{Context 85}.
\item[ALGEBRA-RECON-182] Source lines 7852--7868: \hyperlink{algebra-context-086}{Context 86}.
\item[ALGEBRA-RECON-184] Source lines 7884--7924: \hyperlink{algebra-context-087}{Context 87}.
\item[ALGEBRA-RECON-185] Source lines 7927--7963: \hyperlink{algebra-context-088}{Context 88}.
\end{description}

\subsection{Normal rings}
\hypertarget{algebra-source-section-normal-rings}{}
\noindent\href{https://github.com/stacks/stacks-project/blob/a04446e57ec1fbc252a871afcec7752fb2807b14/algebra.tex\#L7971}{Original source section}. \href{algebra.pdf\#nameddest=section-normal-rings}{Accompanying chapter}.

\begin{description}
\item[ALGEBRA-RECON-186] Source lines 7984--7993: \hyperlink{algebra-context-089}{Context 89}.
\item[ALGEBRA-RECON-189] Source lines 8017--8045: \hyperlink{algebra-context-089}{Context 89}.
\item[ALGEBRA-RECON-190] Source lines 8047--8063: \hyperlink{algebra-context-089}{Context 89}.
\item[ALGEBRA-RECON-191] Source lines 8065--8079: \hyperlink{algebra-context-089}{Context 89}.
\item[ALGEBRA-RECON-193] Source lines 8081--8127: \hyperlink{algebra-context-089}{Context 89}.
\item[ALGEBRA-RECON-194] Source lines 8129--8147: \hyperlink{algebra-context-089}{Context 89}.
\item[ALGEBRA-RECON-195] Source lines 8149--8174: \hyperlink{algebra-context-089}{Context 89}.
\item[ALGEBRA-RECON-196] Source lines 8188--8198: \hyperlink{algebra-context-089}{Context 89}.
\item[ALGEBRA-RECON-197] Source lines 8212--8215: \hyperlink{algebra-context-089}{Context 89}.
\item[ALGEBRA-RECON-198] Source lines 8223--8234: \hyperlink{algebra-context-089}{Context 89}.
\item[ALGEBRA-RECON-199] Source lines 8237--8244: \hyperlink{algebra-context-089}{Context 89}.
\item[ALGEBRA-RECON-200] Source lines 8246--8257: \hyperlink{algebra-context-089}{Context 89}.
\item[ALGEBRA-RECON-202] Source lines 8259--8272: \hyperlink{algebra-context-089}{Context 89}.
\item[ALGEBRA-RECON-204] Source lines 8301--8308: \hyperlink{algebra-context-089}{Context 89}.
\item[ALGEBRA-RECON-205] Source lines 8313--8316: \hyperlink{algebra-context-089}{Context 89}.
\item[ALGEBRA-RECON-206] Source lines 8320--8330: \hyperlink{algebra-context-089}{Context 89}.
\end{description}

\subsection{Going down for integral over normal}
\hypertarget{algebra-source-section-going-down-integral-over-normal}{}
\noindent\href{https://github.com/stacks/stacks-project/blob/a04446e57ec1fbc252a871afcec7752fb2807b14/algebra.tex\#L8339}{Original source section}. \href{algebra.pdf\#nameddest=section-going-down-integral-over-normal}{Accompanying chapter}.

\begin{description}
\item[ALGEBRA-RECON-207] Source lines 8347--8396: \hyperlink{algebra-context-090}{Context 90}; \hyperlink{algebra-context-091}{Context 91}.
\item[ALGEBRA-RECON-209] Source lines 8409--8419: \hyperlink{algebra-context-092}{Context 92}.
\item[ALGEBRA-RECON-210] Source lines 8421--8430: \hyperlink{algebra-context-093}{Context 93}.
\item[ALGEBRA-RECON-211] Source lines 8430--8430: \hyperlink{algebra-context-093}{Context 93}.
\item[ALGEBRA-RECON-212] Source lines 8432--8496: \hyperlink{algebra-context-094}{Context 94}.
\item[ALGEBRA-RECON-214] Source lines 8496--8496: \hyperlink{algebra-context-094}{Context 94}.
\item[ALGEBRA-RECON-215] Source lines 8498--8515: \hyperlink{algebra-context-095}{Context 95}.
\item[ALGEBRA-RECON-216] Source lines 8532--8549: \hyperlink{algebra-context-096}{Context 96}.
\item[ALGEBRA-RECON-217] Source lines 8552--8560: \hyperlink{algebra-context-096}{Context 96}.
\item[ALGEBRA-RECON-218] Source lines 8563--8571: \hyperlink{algebra-context-096}{Context 96}.
\end{description}

\subsection{Flat modules and flat ring maps}
\hypertarget{algebra-source-section-flat}{}
\noindent\href{https://github.com/stacks/stacks-project/blob/a04446e57ec1fbc252a871afcec7752fb2807b14/algebra.tex\#L8587}{Original source section}. \href{algebra.pdf\#nameddest=section-flat}{Accompanying chapter}.

\begin{description}
\item[ALGEBRA-RECON-219] Source lines 8643--8657: \hyperlink{algebra-context-097}{Context 97}.
\item[ALGEBRA-RECON-220] Source lines 8660--8669: \hyperlink{algebra-context-097}{Context 97}.
\item[ALGEBRA-RECON-221] Source lines 8680--8705: \hyperlink{algebra-context-097}{Context 97}.
\item[ALGEBRA-RECON-222] Source lines 8729--8733: \hyperlink{algebra-context-097}{Context 97}.
\item[ALGEBRA-RECON-223] Source lines 8803--8818: \hyperlink{algebra-context-097}{Context 97}.
\item[ALGEBRA-RECON-224] Source lines 8821--8853: \hyperlink{algebra-context-097}{Context 97}.
\item[ALGEBRA-RECON-225] Source lines 8870--8886: \hyperlink{algebra-context-097}{Context 97}.
\item[ALGEBRA-RECON-226] Source lines 8894--8900: \hyperlink{algebra-context-097}{Context 97}.
\item[ALGEBRA-RECON-227] Source lines 8909--8927: \hyperlink{algebra-context-097}{Context 97}.
\item[ALGEBRA-RECON-228] Source lines 8940--8955: \hyperlink{algebra-context-097}{Context 97}.
\item[ALGEBRA-RECON-229] Source lines 8957--8978: \hyperlink{algebra-context-097}{Context 97}.
\item[ALGEBRA-RECON-230] Source lines 8980--9034: \hyperlink{algebra-context-097}{Context 97}.
\item[ALGEBRA-RECON-232] Source lines 9048--9067: \hyperlink{algebra-context-097}{Context 97}.
\item[ALGEBRA-RECON-233] Source lines 9078--9094: \hyperlink{algebra-context-097}{Context 97}.
\item[ALGEBRA-RECON-235] Source lines 9108--9121: \hyperlink{algebra-context-097}{Context 97}.
\item[ALGEBRA-RECON-237] Source lines 9142--9161: \hyperlink{algebra-context-097}{Context 97}.
\item[ALGEBRA-RECON-238] Source lines 9175--9190: \hyperlink{algebra-context-097}{Context 97}.
\item[ALGEBRA-RECON-239] Source lines 9195--9262: \hyperlink{algebra-context-097}{Context 97}.
\item[ALGEBRA-RECON-240] Source lines 9272--9279: \hyperlink{algebra-context-097}{Context 97}.
\item[ALGEBRA-RECON-241] Source lines 9292--9301: \hyperlink{algebra-context-097}{Context 97}.
\end{description}

\subsection{Supports and annihilators}
\hypertarget{algebra-source-section-supp-and-ann}{}
\noindent\href{https://github.com/stacks/stacks-project/blob/a04446e57ec1fbc252a871afcec7752fb2807b14/algebra.tex\#L9307}{Original source section}. \href{algebra.pdf\#nameddest=section-supp-and-ann}{Accompanying chapter}.

\begin{description}
\item[ALGEBRA-RECON-242] Source lines 9339--9344: \hyperlink{algebra-context-098}{Context 98}.
\item[ALGEBRA-RECON-243] Source lines 9363--9382: \hyperlink{algebra-context-099}{Context 99}.
\item[ALGEBRA-RECON-244] Source lines 9384--9410: \hyperlink{algebra-context-100}{Context 100}.
\item[ALGEBRA-RECON-245] Source lines 9419--9444: \hyperlink{algebra-context-101}{Context 101}.
\item[ALGEBRA-RECON-246] Source lines 9453--9458: \hyperlink{algebra-context-102}{Context 102}.
\item[ALGEBRA-RECON-247] Source lines 9467--9484: \hyperlink{algebra-context-103}{Context 103}.
\item[ALGEBRA-RECON-248] Source lines 9486--9508: \hyperlink{algebra-context-104}{Context 104}.
\end{description}

\subsection{Going up and going down}
\hypertarget{algebra-source-section-going-up}{}
\noindent\href{https://github.com/stacks/stacks-project/blob/a04446e57ec1fbc252a871afcec7752fb2807b14/algebra.tex\#L9514}{Original source section}. \href{algebra.pdf\#nameddest=section-going-up}{Accompanying chapter}.

\begin{description}
\item[ALGEBRA-RECON-253] Source lines 9573--9586: \hyperlink{algebra-context-105}{Context 105}.
\item[ALGEBRA-RECON-254] Source lines 9601--9603: \hyperlink{algebra-context-106}{Context 106}.
\item[ALGEBRA-RECON-255] Source lines 9611--9615: \hyperlink{algebra-context-107}{Context 107}.
\item[ALGEBRA-RECON-257] Source lines 9617--9649: \hyperlink{algebra-context-108}{Context 108}.
\item[ALGEBRA-RECON-258] Source lines 9660--9679: \hyperlink{algebra-context-109}{Context 109}.
\item[ALGEBRA-RECON-259] Source lines 9690--9705: \hyperlink{algebra-context-110}{Context 110}.
\item[ALGEBRA-RECON-260] Source lines 9715--9737: \hyperlink{algebra-context-111}{Context 111}.
\item[ALGEBRA-RECON-261] Source lines 9754--9772: \hyperlink{algebra-context-112}{Context 112}.
\item[ALGEBRA-RECON-262] Source lines 9782--9791: \hyperlink{algebra-context-113}{Context 113}.
\item[ALGEBRA-RECON-263] Source lines 9814--9845: \hyperlink{algebra-context-114}{Context 114}.
\item[ALGEBRA-RECON-264] Source lines 9858--9882: \hyperlink{algebra-context-115}{Context 115}.
\end{description}

\subsection{Separable extensions}
\hypertarget{algebra-source-section-separability}{}
\noindent\href{https://github.com/stacks/stacks-project/blob/a04446e57ec1fbc252a871afcec7752fb2807b14/algebra.tex\#L9906}{Original source section}. \href{algebra.pdf\#nameddest=section-separability}{Accompanying chapter}.

\begin{description}
\item[ALGEBRA-RECON-265] Source lines 9940--9942: \hyperlink{algebra-context-116}{Context 116}.
\item[ALGEBRA-RECON-266] Source lines 9944--9959: \hyperlink{algebra-context-117}{Context 117}.
\item[ALGEBRA-RECON-267] Source lines 9975--10032: \hyperlink{algebra-context-118}{Context 118}.
\end{description}

\subsection{Geometrically reduced algebras}
\hypertarget{algebra-source-section-geometrically-reduced}{}
\noindent\href{https://github.com/stacks/stacks-project/blob/a04446e57ec1fbc252a871afcec7752fb2807b14/algebra.tex\#L10037}{Original source section}. \href{algebra.pdf\#nameddest=section-geometrically-reduced}{Accompanying chapter}.

\begin{description}
\item[ALGEBRA-RECON-268] Source lines 10078--10081: \hyperlink{algebra-context-119}{Context 119}.
\item[ALGEBRA-RECON-269] Source lines 10093--10096: \hyperlink{algebra-context-119}{Context 119}.
\item[ALGEBRA-RECON-270] Source lines 10098--10103: \hyperlink{algebra-context-119}{Context 119}.
\item[ALGEBRA-RECON-271] Source lines 10121--10127: \hyperlink{algebra-context-119}{Context 119}.
\item[ALGEBRA-RECON-272] Source lines 10137--10149: \hyperlink{algebra-context-119}{Context 119}.
\item[ALGEBRA-RECON-273] Source lines 10160--10194: \hyperlink{algebra-context-119}{Context 119}.
\item[ALGEBRA-RECON-274] Source lines 10204--10218: \hyperlink{algebra-context-119}{Context 119}.
\item[ALGEBRA-RECON-275] Source lines 10228--10255: \hyperlink{algebra-context-119}{Context 119}.
\item[ALGEBRA-RECON-276] Source lines 10264--10282: \hyperlink{algebra-context-119}{Context 119}.
\end{description}

\subsection{Separable extensions, continued}
\hypertarget{algebra-source-section-separability-continued}{}
\noindent\href{https://github.com/stacks/stacks-project/blob/a04446e57ec1fbc252a871afcec7752fb2807b14/algebra.tex\#L10288}{Original source section}. \href{algebra.pdf\#nameddest=section-separability-continued}{Accompanying chapter}.

\begin{description}
\item[ALGEBRA-RECON-277] Source lines 10311--10349: \hyperlink{algebra-context-120}{Context 120}.
\item[ALGEBRA-RECON-278] Source lines 10374--10406: \hyperlink{algebra-context-121}{Context 121}.
\item[ALGEBRA-RECON-279] Source lines 10415--10416: \hyperlink{algebra-context-122}{Context 122}.
\item[ALGEBRA-RECON-280] Source lines 10424--10478: \hyperlink{algebra-context-123}{Context 123}.
\end{description}

\subsection{Perfect fields}
\hypertarget{algebra-source-section-perfect-fields}{}
\noindent\href{https://github.com/stacks/stacks-project/blob/a04446e57ec1fbc252a871afcec7752fb2807b14/algebra.tex\#L10482}{Original source section}. \href{algebra.pdf\#nameddest=section-perfect-fields}{Accompanying chapter}.

\begin{description}
\item[ALGEBRA-RECON-281] Source lines 10502--10509: \hyperlink{algebra-context-124}{Context 124}.
\item[ALGEBRA-RECON-282] Source lines 10529--10541: \hyperlink{algebra-context-124}{Context 124}.
\item[ALGEBRA-RECON-283] Source lines 10555--10578: \hyperlink{algebra-context-124}{Context 124}.
\item[ALGEBRA-RECON-284] Source lines 10580--10594: \hyperlink{algebra-context-124}{Context 124}.
\end{description}

\subsection{Universal homeomorphisms}
\hypertarget{algebra-source-section-universal-homeomorphism}{}
\noindent\href{https://github.com/stacks/stacks-project/blob/a04446e57ec1fbc252a871afcec7752fb2807b14/algebra.tex\#L10604}{Original source section}. \href{algebra.pdf\#nameddest=section-universal-homeomorphism}{Accompanying chapter}.

\begin{description}
\item[ALGEBRA-RECON-285] Source lines 10623--10629: \hyperlink{algebra-context-125}{Context 125}.
\item[ALGEBRA-RECON-286] Source lines 10647--10714: \hyperlink{algebra-context-126}{Context 126}.
\item[ALGEBRA-RECON-287] Source lines 10731--10766: \hyperlink{algebra-context-127}{Context 127}.
\item[ALGEBRA-RECON-288] Source lines 10783--10810: \hyperlink{algebra-context-128}{Context 128}.
\item[ALGEBRA-RECON-289] Source lines 10821--10844: \hyperlink{algebra-context-129}{Context 129}.
\item[ALGEBRA-RECON-291] Source lines 10861--10864: \hyperlink{algebra-context-130}{Context 130}.
\item[ALGEBRA-RECON-292] Source lines 10883--10909: \hyperlink{algebra-context-131}{Context 131}.
\item[ALGEBRA-RECON-293] Source lines 10924--10965: \hyperlink{algebra-context-132}{Context 132}.
\item[ALGEBRA-RECON-294] Source lines 10982--10985: \hyperlink{algebra-context-133}{Context 133}.
\item[ALGEBRA-RECON-295] Source lines 10993--11002: \hyperlink{algebra-context-134}{Context 134}.
\item[ALGEBRA-RECON-296] Source lines 11020--11060: \hyperlink{algebra-context-135}{Context 135}.
\end{description}

\subsection{Geometrically irreducible algebras}
\hypertarget{algebra-source-section-algebras-over-fields}{}
\noindent\href{https://github.com/stacks/stacks-project/blob/a04446e57ec1fbc252a871afcec7752fb2807b14/algebra.tex\#L11073}{Original source section}. \href{algebra.pdf\#nameddest=section-algebras-over-fields}{Accompanying chapter}.

\begin{description}
\item[ALGEBRA-RECON-297] Source lines 11098--11102: \hyperlink{algebra-context-136}{Context 136}.
\item[ALGEBRA-RECON-298] Source lines 11113--11142: \hyperlink{algebra-context-137}{Context 137}.
\item[ALGEBRA-RECON-299] Source lines 11163--11203: \hyperlink{algebra-context-138}{Context 138}.
\item[ALGEBRA-RECON-300] Source lines 11224--11231: \hyperlink{algebra-context-139}{Context 139}.
\item[ALGEBRA-RECON-301] Source lines 11248--11253: \hyperlink{algebra-context-140}{Context 140}.
\item[ALGEBRA-RECON-302] Source lines 11269--11289: \hyperlink{algebra-context-141}{Context 141}.
\item[ALGEBRA-RECON-303] Source lines 11303--11319: \hyperlink{algebra-context-142}{Context 142}.
\item[ALGEBRA-RECON-304] Source lines 11330--11337: \hyperlink{algebra-context-143}{Context 143}.
\item[ALGEBRA-RECON-305] Source lines 11351--11374: \hyperlink{algebra-context-144}{Context 144}.
\item[ALGEBRA-RECON-306] Source lines 11385--11388: \hyperlink{algebra-context-145}{Context 145}.
\item[ALGEBRA-RECON-307] Source lines 11402--11424: \hyperlink{algebra-context-146}{Context 146}.
\item[ALGEBRA-RECON-308] Source lines 11436--11442: \hyperlink{algebra-context-147}{Context 147}.
\item[ALGEBRA-RECON-309] Source lines 11454--11478: \hyperlink{algebra-context-148}{Context 148}.
\end{description}

\subsection{Geometrically connected algebras}
\hypertarget{algebra-source-section-geometrically-connected}{}
\noindent\href{https://github.com/stacks/stacks-project/blob/a04446e57ec1fbc252a871afcec7752fb2807b14/algebra.tex\#L11484}{Original source section}. \href{algebra.pdf\#nameddest=section-geometrically-connected}{Accompanying chapter}.

\begin{description}
\item[ALGEBRA-RECON-311] Source lines 11490--11527: \hyperlink{algebra-context-149}{Context 149}.
\item[ALGEBRA-RECON-312] Source lines 11544--11557: \hyperlink{algebra-context-150}{Context 150}.
\item[ALGEBRA-RECON-313] Source lines 11583--11584: \hyperlink{algebra-context-151}{Context 151}.
\item[ALGEBRA-RECON-314] Source lines 11601--11603: \hyperlink{algebra-context-152}{Context 152}.
\item[ALGEBRA-RECON-315] Source lines 11627--11648: \hyperlink{algebra-context-153}{Context 153}.
\end{description}

\subsection{Geometrically integral algebras}
\hypertarget{algebra-source-section-geometrically-integral}{}
\noindent\href{https://github.com/stacks/stacks-project/blob/a04446e57ec1fbc252a871afcec7752fb2807b14/algebra.tex\#L11654}{Original source section}. \href{algebra.pdf\#nameddest=section-geometrically-integral}{Accompanying chapter}.

\begin{description}
\item[ALGEBRA-RECON-318] Source lines 11682--11682: \hyperlink{algebra-context-154}{Context 154}.
\item[ALGEBRA-RECON-319] Source lines 11699--11701: \hyperlink{algebra-context-155}{Context 155}.
\item[ALGEBRA-RECON-320] Source lines 11712--11717: \hyperlink{algebra-context-156}{Context 156}.
\end{description}

\subsection{Valuation rings}
\hypertarget{algebra-source-section-valuation-rings}{}
\noindent\href{https://github.com/stacks/stacks-project/blob/a04446e57ec1fbc252a871afcec7752fb2807b14/algebra.tex\#L11723}{Original source section}. \href{algebra.pdf\#nameddest=section-valuation-rings}{Accompanying chapter}.

\begin{description}
\item[ALGEBRA-RECON-321] Source lines 11757--11779: \hyperlink{algebra-context-157}{Context 157}.
\item[ALGEBRA-RECON-322] Source lines 11789--11796: \hyperlink{algebra-context-157}{Context 157}.
\item[ALGEBRA-RECON-323] Source lines 11804--11819: \hyperlink{algebra-context-157}{Context 157}.
\item[ALGEBRA-RECON-324] Source lines 11825--11840: \hyperlink{algebra-context-157}{Context 157}.
\item[ALGEBRA-RECON-325] Source lines 11854--11859: \hyperlink{algebra-context-157}{Context 157}.
\item[ALGEBRA-RECON-326] Source lines 11869--11872: \hyperlink{algebra-context-157}{Context 157}.
\item[ALGEBRA-RECON-327] Source lines 11883--11887: \hyperlink{algebra-context-157}{Context 157}.
\item[ALGEBRA-RECON-328] Source lines 11894--11899: \hyperlink{algebra-context-157}{Context 157}.
\item[ALGEBRA-RECON-329] Source lines 11912--11920: \hyperlink{algebra-context-157}{Context 157}.
\item[ALGEBRA-RECON-330] Source lines 11937--11959: \hyperlink{algebra-context-157}{Context 157}.
\item[ALGEBRA-RECON-332] Source lines 11980--11983: \hyperlink{algebra-context-157}{Context 157}.
\item[ALGEBRA-RECON-333] Source lines 12001--12003: \hyperlink{algebra-context-157}{Context 157}.
\item[ALGEBRA-RECON-334] Source lines 12017--12017: \hyperlink{algebra-context-157}{Context 157}.
\item[ALGEBRA-RECON-335] Source lines 12031--12039: \hyperlink{algebra-context-157}{Context 157}.
\item[ALGEBRA-RECON-336] Source lines 12072--12072: \hyperlink{algebra-context-157}{Context 157}.
\item[ALGEBRA-RECON-337] Source lines 12093--12093: \hyperlink{algebra-context-157}{Context 157}.
\item[ALGEBRA-RECON-338] Source lines 12103--12124: \hyperlink{algebra-context-157}{Context 157}.
\end{description}

\subsection{More Noetherian rings}
\hypertarget{algebra-source-section-Noetherian-again}{}
\noindent\href{https://github.com/stacks/stacks-project/blob/a04446e57ec1fbc252a871afcec7752fb2807b14/algebra.tex\#L12149}{Original source section}. \href{algebra.pdf\#nameddest=section-Noetherian-again}{Accompanying chapter}.

\begin{description}
\item[ALGEBRA-RECON-339] Source lines 12163--12183: \hyperlink{algebra-context-158}{Context 158}.
\item[ALGEBRA-RECON-341] Source lines 12195--12217: \hyperlink{algebra-context-159}{Context 159}.
\item[ALGEBRA-RECON-342] Source lines 12235--12237: \hyperlink{algebra-context-160}{Context 160}.
\item[ALGEBRA-RECON-343] Source lines 12248--12253: \hyperlink{algebra-context-161}{Context 161}.
\item[ALGEBRA-RECON-344] Source lines 12261--12279: \hyperlink{algebra-context-162}{Context 162}.
\item[ALGEBRA-RECON-345] Source lines 12284--12293: \hyperlink{algebra-context-163}{Context 163}.
\item[ALGEBRA-RECON-346] Source lines 12305--12322: \hyperlink{algebra-context-164}{Context 164}.
\end{description}

\subsection{Length}
\hypertarget{algebra-source-section-length}{}
\noindent\href{https://github.com/stacks/stacks-project/blob/a04446e57ec1fbc252a871afcec7752fb2807b14/algebra.tex\#L12346}{Original source section}. \href{algebra.pdf\#nameddest=section-length}{Accompanying chapter}.

\begin{description}
\item[ALGEBRA-RECON-347] Source lines 12359--12368: \hyperlink{algebra-context-165}{Context 165}.
\item[ALGEBRA-RECON-348] Source lines 12390--12390: \hyperlink{algebra-context-165}{Context 165}.
\item[ALGEBRA-RECON-349] Source lines 12405--12417: \hyperlink{algebra-context-165}{Context 165}.
\item[ALGEBRA-RECON-350] Source lines 12429--12441: \hyperlink{algebra-context-165}{Context 165}.
\item[ALGEBRA-RECON-351] Source lines 12452--12456: \hyperlink{algebra-context-165}{Context 165}.
\item[ALGEBRA-RECON-352] Source lines 12464--12473: \hyperlink{algebra-context-165}{Context 165}.
\item[ALGEBRA-RECON-353] Source lines 12484--12486: \hyperlink{algebra-context-165}{Context 165}.
\item[ALGEBRA-RECON-354] Source lines 12492--12505: \hyperlink{algebra-context-165}{Context 165}.
\item[ALGEBRA-RECON-356] Source lines 12569--12599: \hyperlink{algebra-context-165}{Context 165}.
\item[ALGEBRA-RECON-357] Source lines 12623--12638: \hyperlink{algebra-context-165}{Context 165}.
\item[ALGEBRA-RECON-358] Source lines 12650--12670: \hyperlink{algebra-context-165}{Context 165}.
\item[ALGEBRA-RECON-359] Source lines 12685--12686: \hyperlink{algebra-context-165}{Context 165}.
\end{description}

\subsection{Artinian rings}
\hypertarget{algebra-source-section-artinian}{}
\noindent\href{https://github.com/stacks/stacks-project/blob/a04446e57ec1fbc252a871afcec7752fb2807b14/algebra.tex\#L12699}{Original source section}. \href{algebra.pdf\#nameddest=section-artinian}{Accompanying chapter}.

\begin{description}
\item[ALGEBRA-RECON-360] Source lines 12720--12720: \hyperlink{algebra-context-166}{Context 166}.
\item[ALGEBRA-RECON-361] Source lines 12729--12735: \hyperlink{algebra-context-167}{Context 167}.
\item[ALGEBRA-RECON-362] Source lines 12744--12754: \hyperlink{algebra-context-168}{Context 168}.
\item[ALGEBRA-RECON-363] Source lines 12765--12783: \hyperlink{algebra-context-169}{Context 169}.
\item[ALGEBRA-RECON-364] Source lines 12795--12818: \hyperlink{algebra-context-170}{Context 170}.
\end{description}

\subsection{Homomorphisms essentially of finite type}
\hypertarget{algebra-source-section-essentially-of-finite-type}{}
\noindent\href{https://github.com/stacks/stacks-project/blob/a04446e57ec1fbc252a871afcec7752fb2807b14/algebra.tex\#L12829}{Original source section}. \href{algebra.pdf\#nameddest=section-essentially-of-finite-type}{Accompanying chapter}.

\begin{description}
\item[ALGEBRA-RECON-365] Source lines 12854--12854: \hyperlink{algebra-context-171}{Context 171}.
\item[ALGEBRA-RECON-366] Source lines 12865--12865: \hyperlink{algebra-context-172}{Context 172}.
\item[ALGEBRA-RECON-367] Source lines 12883--12923: \hyperlink{algebra-context-173}{Context 173}.
\item[ALGEBRA-RECON-368] Source lines 12930--12972: \hyperlink{algebra_localization_kgroups_graded_notes_20260927}{Supplementary proof 1}, Valuation quantifier.
\end{description}

\subsection{K-groups}
\hypertarget{algebra-source-section-K-groups}{}
\noindent\href{https://github.com/stacks/stacks-project/blob/a04446e57ec1fbc252a871afcec7752fb2807b14/algebra.tex\#L12995}{Original source section}. \href{algebra.pdf\#nameddest=section-K-groups}{Accompanying chapter}.

\begin{description}
\item[ALGEBRA-RECON-369] Source lines 13013--13022: \hyperlink{algebra_localization_kgroups_graded_notes_20260927}{Supplementary proof 1}, Nonnegative presentation rank.
\item[ALGEBRA-RECON-370] Source lines 13068--13153: \hyperlink{algebra_localization_kgroups_graded_notes_20260927}{Supplementary proof 1}, Both split exact sequences.
\item[ALGEBRA-RECON-371] Source lines 13078--13090: \hyperlink{algebra_localization_kgroups_graded_notes_20260927}{Supplementary proof 1}, PID subtraction and zero factors.
\end{description}

\subsection{Graded rings}
\hypertarget{algebra-source-section-graded}{}
\noindent\href{https://github.com/stacks/stacks-project/blob/a04446e57ec1fbc252a871afcec7752fb2807b14/algebra.tex\#L13199}{Original source section}. \href{algebra.pdf\#nameddest=section-graded}{Accompanying chapter}.

\begin{description}
\item[ALGEBRA-RECON-372] Source lines 13224--13228: \hyperlink{algebra_localization_kgroups_graded_notes_20260927}{Supplementary proof 1}, Graded wording.
\item[ALGEBRA-RECON-373] Source lines 13231--13247: \hyperlink{algebra_localization_kgroups_graded_notes_20260927}{Supplementary proof 1}, Graded scalar action.
\item[ALGEBRA-RECON-374] Source lines 13249--13272: \hyperlink{algebra_graded_proj_hilbert_notes_20260927}{Supplementary proof 2}, Bounded-below graded Nakayama.
\item[ALGEBRA-RECON-375] Source lines 13307--13355: \hyperlink{algebra_graded_proj_hilbert_notes_20260927}{Supplementary proof 2}, Integral-closure wording.
\end{description}

\subsection{Proj of a graded ring}
\hypertarget{algebra-source-section-proj}{}
\noindent\href{https://github.com/stacks/stacks-project/blob/a04446e57ec1fbc252a871afcec7752fb2807b14/algebra.tex\#L13364}{Original source section}. \href{algebra.pdf\#nameddest=section-proj}{Accompanying chapter}.

\begin{description}
\item[ALGEBRA-RECON-376] Source lines 13369--13380: \hyperlink{algebra_graded_proj_hilbert_notes_20260927}{Supplementary proof 2}, Proper homogeneous prime test.
\item[ALGEBRA-RECON-377] Source lines 13402--13408: \hyperlink{algebra_graded_proj_hilbert_notes_20260927}{Supplementary proof 2}, Degree-zero localization wording.
\item[ALGEBRA-RECON-378] Source lines 13650--13667: \hyperlink{algebra_graded_proj_hilbert_notes_20260927}{Supplementary proof 2}, Homogeneous core of a prime.
\item[ALGEBRA-RECON-379] Source lines 13669--13683: \hyperlink{algebra_graded_proj_hilbert_notes_20260927}{Supplementary proof 2}, Homogeneous core of a prime.
\item[ALGEBRA-RECON-380] Source lines 13685--13727: \hyperlink{algebra_graded_proj_hilbert_notes_20260927}{Supplementary proof 2}, Negative degrees in the finite generator list.
\item[ALGEBRA-RECON-381] Source lines 13729--13769: \hyperlink{algebra_graded_proj_hilbert_notes_20260927}{Supplementary proof 2}, Positive-degree homogenization and the original base ring.
\end{description}

\subsection{Noetherian graded rings}
\hypertarget{algebra-source-section-noetherian-graded}{}
\noindent\href{https://github.com/stacks/stacks-project/blob/a04446e57ec1fbc252a871afcec7752fb2807b14/algebra.tex\#L13774}{Original source section}. \href{algebra.pdf\#nameddest=section-noetherian-graded}{Accompanying chapter}.

\begin{description}
\item[ALGEBRA-RECON-382] Source lines 13812--13830: \hyperlink{algebra_graded_proj_hilbert_notes_20260927}{Supplementary proof 2}, Noetherian and numerical wording.
\item[ALGEBRA-RECON-383] Source lines 13913--13943: \hyperlink{algebra_graded_proj_hilbert_notes_20260927}{Supplementary proof 2}, Nilpotence induction in the Hilbert proof.
\item[ALGEBRA-RECON-384] Source lines 13964--13982: \hyperlink{algebra_graded_proj_hilbert_notes_20260927}{Supplementary proof 2}, Hilbert bound notation.
\end{description}

\subsection{Noetherian local rings}
\hypertarget{algebra-source-section-Noetherian-local}{}
\noindent\href{https://github.com/stacks/stacks-project/blob/a04446e57ec1fbc252a871afcec7752fb2807b14/algebra.tex\#L13991}{Original source section}. \href{algebra.pdf\#nameddest=section-Noetherian-local}{Accompanying chapter}.

\begin{description}
\item[ALGEBRA-RECON-385] Source lines 14026--14042: \hyperlink{algebra_local_hilbert_dimension_notes_20260927}{Supplementary proof 3}, Finite-length wording.
\item[ALGEBRA-RECON-386] Source lines 14044--14085: \hyperlink{algebra_local_hilbert_dimension_notes_20260927}{Supplementary proof 3}, Finite-colength comparison and explicit constants.
\item[ALGEBRA-RECON-387] Source lines 14087--14132: \hyperlink{algebra_local_hilbert_dimension_notes_20260927}{Supplementary proof 3}, Artin–Rees endpoint.
\item[ALGEBRA-RECON-388] Source lines 14212--14223: \hyperlink{algebra_local_hilbert_dimension_notes_20260927}{Supplementary proof 3}, Degree comparison and highest coefficients.
\item[ALGEBRA-RECON-397] Source lines 14225--14269: \hyperlink{algebra_local_hilbert_dimension_notes_20260927}{Supplementary proof 3}, Degree comparison and highest coefficients.
\end{description}

\subsection{Dimension}
\hypertarget{algebra-source-section-dimension}{}
\noindent\href{https://github.com/stacks/stacks-project/blob/a04446e57ec1fbc252a871afcec7752fb2807b14/algebra.tex\#L14305}{Original source section}. \href{algebra.pdf\#nameddest=section-dimension}{Accompanying chapter}.

\begin{description}
\item[ALGEBRA-RECON-389] Source lines 14331--14343: \hyperlink{algebra_local_hilbert_dimension_notes_20260927}{Supplementary proof 3}, Punctuation in the dimension definition.
\item[ALGEBRA-RECON-390] Source lines 14369--14378: \hyperlink{algebra_local_hilbert_dimension_notes_20260927}{Supplementary proof 3}, Minimal prime wording.
\item[ALGEBRA-RECON-391] Source lines 14382--14390: \hyperlink{algebra_local_hilbert_dimension_notes_20260927}{Supplementary proof 3}, Length over the original ring is correct.
\item[ALGEBRA-RECON-392] Source lines 14363--14443: \hyperlink{algebra_local_hilbert_dimension_notes_20260927}{Supplementary proof 3}, Zero ring and dimension at most zero.
\item[ALGEBRA-RECON-393] Source lines 14482--14487: \hyperlink{algebra_local_hilbert_dimension_notes_20260927}{Supplementary proof 3}, One-dimensional wording.
\item[ALGEBRA-RECON-394] Source lines 14506--14592: \hyperlink{algebra_local_hilbert_dimension_notes_20260927}{Supplementary proof 3}, Multiplication maps in both dimension proofs.
\item[ALGEBRA-RECON-395] Source lines 14594--14596: \hyperlink{algebra_local_hilbert_dimension_notes_20260927}{Supplementary proof 3}, Closing sentence and induction scope.
\item[ALGEBRA-RECON-396] Source lines 14599--14608: \hyperlink{algebra_local_hilbert_dimension_notes_20260927}{Supplementary proof 3}, Elements generating the maximal ideal.
\end{description}

\subsection{Applications of dimension theory}
\hypertarget{algebra-source-section-applications-dimension-theory}{}
\noindent\href{https://github.com/stacks/stacks-project/blob/a04446e57ec1fbc252a871afcec7752fb2807b14/algebra.tex\#L14731}{Original source section}. \href{algebra.pdf\#nameddest=section-applications-dimension-theory}{Accompanying chapter}.

\begin{description}
\item[ALGEBRA-RECON-398] Source lines 14767--14770: \hyperlink{algebra_dimension_support_associated_notes_20260927}{Supplementary proof 4}, The p-adic example.
\item[ALGEBRA-RECON-399] Source lines 14790--14839: \hyperlink{algebra_dimension_support_associated_notes_20260927}{Supplementary proof 4}, Source editorial reminders.
\item[ALGEBRA-RECON-400] Source lines 14862--14869: \hyperlink{algebra_dimension_support_associated_notes_20260927}{Supplementary proof 4}, The nonmaximal prime and strict localization.
\item[ALGEBRA-RECON-401] Source lines 14870--14887: \hyperlink{algebra_dimension_support_associated_notes_20260927}{Supplementary proof 4}, The nonmaximal prime and strict localization.
\end{description}

\subsection{Support and dimension of modules}
\hypertarget{algebra-source-section-support}{}
\noindent\href{https://github.com/stacks/stacks-project/blob/a04446e57ec1fbc252a871afcec7752fb2807b14/algebra.tex\#L14909}{Original source section}. \href{algebra.pdf\#nameddest=section-support}{Accompanying chapter}.

\begin{description}
\item[ALGEBRA-RECON-402] Source lines 14949--14958: \hyperlink{algebra_dimension_support_associated_notes_20260927}{Supplementary proof 4}, Filtration wording and the original multiplicities.
\item[ALGEBRA-RECON-403] Source lines 14960--15051: \hyperlink{algebra_dimension_support_associated_notes_20260927}{Supplementary proof 4}, Filtration wording and the original multiplicities.
\item[ALGEBRA-RECON-404] Source lines 14996--15014: \hyperlink{algebra_dimension_support_associated_notes_20260927}{Supplementary proof 4}, Support and a finitely generated annihilating ideal.
\end{description}

\subsection{Associated primes}
\hypertarget{algebra-source-section-ass}{}
\noindent\href{https://github.com/stacks/stacks-project/blob/a04446e57ec1fbc252a871afcec7752fb2807b14/algebra.tex\#L15099}{Original source section}. \href{algebra.pdf\#nameddest=section-ass}{Accompanying chapter}.

\begin{description}
\item[ALGEBRA-RECON-405] Source lines 15138--15150: \hyperlink{algebra_dimension_support_associated_notes_20260927}{Supplementary proof 4}, Annihilator spelling in the exact sequence.
\item[ALGEBRA-RECON-406] Source lines 15296--15326: \hyperlink{algebra_dimension_support_associated_notes_20260927}{Supplementary proof 4}, One-equation module wording.
\item[ALGEBRA-RECON-407] Source lines 15436--15448: \hyperlink{algebra_dimension_support_associated_notes_20260927}{Supplementary proof 4}, Prime type and the two localization maps.
\item[ALGEBRA-RECON-408] Source lines 15470--15538: \hyperlink{algebra_dimension_support_associated_notes_20260927}{Supplementary proof 4}, Nonzerodivisors in arbitrary ideals.
\item[ALGEBRA-RECON-409] Source lines 15518--15538: \hyperlink{algebra_dimension_support_associated_notes_20260927}{Supplementary proof 4}, Source editorial reminders.
\end{description}

\subsection{Symbolic powers}
\hypertarget{algebra-source-section-symbolic-power}{}
\noindent\href{https://github.com/stacks/stacks-project/blob/a04446e57ec1fbc252a871afcec7752fb2807b14/algebra.tex\#L15542}{Original source section}. \href{algebra.pdf\#nameddest=section-symbolic-power}{Accompanying chapter}.

\begin{description}
\item[ALGEBRA-RECON-410] Source lines 15548--15575: \hyperlink{algebra_symbolic_weakass_notes_20260927}{Supplementary proof 5}, Symbolic powers and the associated-prime strengthening.
\item[ALGEBRA-RECON-411] Source lines 15577--15582: \hyperlink{algebra_symbolic_weakass_notes_20260927}{Supplementary proof 5}, Flat extension requires injection, not invertibility.
\item[ALGEBRA-RECON-412] Source lines 15584--15612: \hyperlink{algebra_symbolic_weakass_notes_20260927}{Supplementary proof 5}, Flat extension requires injection, not invertibility.
\end{description}

\subsection{Relative assassin}
\hypertarget{algebra-source-section-relative-assassin}{}
\noindent\href{https://github.com/stacks/stacks-project/blob/a04446e57ec1fbc252a871afcec7752fb2807b14/algebra.tex\#L15616}{Original source section}. \href{algebra.pdf\#nameddest=section-relative-assassin}{Accompanying chapter}.

\begin{description}
\item[ALGEBRA-RECON-413] Source lines 15624--15630: \hyperlink{algebra_symbolic_weakass_notes_20260927}{Supplementary proof 5}, Relative-assassin wording and exact local tensor maps.
\item[ALGEBRA-RECON-414] Source lines 15642--15657: \hyperlink{algebra_symbolic_weakass_notes_20260927}{Supplementary proof 5}, Relative-assassin wording and exact local tensor maps.
\item[ALGEBRA-RECON-415] Source lines 15666--15676: \hyperlink{algebra_symbolic_weakass_notes_20260927}{Supplementary proof 5}, Relative-assassin wording and exact local tensor maps.
\item[ALGEBRA-RECON-416] Source lines 15707--15719: \hyperlink{algebra_symbolic_weakass_notes_20260927}{Supplementary proof 5}, Relative-assassin wording and exact local tensor maps.
\item[ALGEBRA-RECON-417] Source lines 15773--15903: \hyperlink{algebra_symbolic_weakass_notes_20260927}{Supplementary proof 5}, Relative-assassin wording and exact local tensor maps.
\item[ALGEBRA-RECON-418] Source lines 15804--15829: \hyperlink{algebra_symbolic_weakass_notes_20260927}{Supplementary proof 5}, Relative-assassin wording and exact local tensor maps.
\item[ALGEBRA-RECON-419] Source lines 15812--15817: \hyperlink{algebra_symbolic_weakass_notes_20260927}{Supplementary proof 5}, Relative-assassin wording and exact local tensor maps.
\item[ALGEBRA-RECON-420] Source lines 15905--15918: \hyperlink{algebra_symbolic_weakass_notes_20260927}{Supplementary proof 5}, Relative-assassin wording and exact local tensor maps.
\item[ALGEBRA-RECON-433] Source lines 15853--16358: \hyperlink{algebra_symbolic_weakass_notes_20260927}{Supplementary proof 5}, Torsion-free generic-fibre comparison.
\end{description}

\subsection{Weakly associated primes}
\hypertarget{algebra-source-section-weakly-ass}{}
\noindent\href{https://github.com/stacks/stacks-project/blob/a04446e57ec1fbc252a871afcec7752fb2807b14/algebra.tex\#L15925}{Original source section}. \href{algebra.pdf\#nameddest=section-weakly-ass}{Accompanying chapter}.

\begin{description}
\item[ALGEBRA-RECON-421] Source lines 15925--15941: \hyperlink{algebra_symbolic_weakass_notes_20260927}{Supplementary proof 5}, Weakly associated primes and the original counterexample.
\item[ALGEBRA-RECON-422] Source lines 15985--16001: \hyperlink{algebra_symbolic_weakass_notes_20260927}{Supplementary proof 5}, Weakly associated primes and the original counterexample.
\item[ALGEBRA-RECON-423] Source lines 16039--16050: \hyperlink{algebra_symbolic_weakass_notes_20260927}{Supplementary proof 5}, Weakly associated primes and the original counterexample.
\item[ALGEBRA-RECON-424] Source lines 16039--16050: \hyperlink{algebra_symbolic_weakass_notes_20260927}{Supplementary proof 5}, Weakly associated primes and the original counterexample.
\item[ALGEBRA-RECON-425] Source lines 16114--16144: \hyperlink{algebra_symbolic_weakass_notes_20260927}{Supplementary proof 5}, Weakly associated primes and the original counterexample.
\item[ALGEBRA-RECON-426] Source lines 16188--16254: \hyperlink{algebra_symbolic_weakass_notes_20260927}{Supplementary proof 5}, Finite maps and the repeated localization type correction.
\item[ALGEBRA-RECON-427] Source lines 16221--16254: \hyperlink{algebra_symbolic_weakass_notes_20260927}{Supplementary proof 5}, Finite maps and the repeated localization type correction.
\item[ALGEBRA-RECON-428] Source lines 16270--16293: \hyperlink{algebra_symbolic_weakass_notes_20260927}{Supplementary proof 5}, Finite maps and the repeated localization type correction.
\item[ALGEBRA-RECON-429] Source lines 16360--16439: \hyperlink{algebra_symbolic_weakass_notes_20260927}{Supplementary proof 5}, Field-extension proof and the nonzero witness.
\item[ALGEBRA-RECON-430] Source lines 16406--16423: \hyperlink{algebra_symbolic_weakass_notes_20260927}{Supplementary proof 5}, Field-extension proof and the nonzero witness.
\item[ALGEBRA-RECON-431] Source lines 16419--16423: \hyperlink{algebra_symbolic_weakass_notes_20260927}{Supplementary proof 5}, Field-extension proof and the nonzero witness.
\item[ALGEBRA-RECON-432] Source lines 16426--16439: \hyperlink{algebra_symbolic_weakass_notes_20260927}{Supplementary proof 5}, Field-extension proof and the nonzero witness.
\item[ALGEBRA-RECON-433] Source lines 15853--16358: \hyperlink{algebra_symbolic_weakass_notes_20260927}{Supplementary proof 5}, Torsion-free generic-fibre comparison.
\end{description}

\subsection{Embedded primes}
\hypertarget{algebra-source-section-embedded-primes}{}
\noindent\href{https://github.com/stacks/stacks-project/blob/a04446e57ec1fbc252a871afcec7752fb2807b14/algebra.tex\#L16445}{Original source section}. \href{algebra.pdf\#nameddest=section-embedded-primes}{Accompanying chapter}.

\begin{description}
\item[ALGEBRA-RECON-434] Source lines 16467--16533: \hyperlink{algebra_embedded_regular_notes_20260927}{Supplementary proof 6}, The embedded-prime quotient and its exact kernel.
\item[ALGEBRA-RECON-435] Source lines 16529--16532: \hyperlink{algebra_embedded_regular_notes_20260927}{Supplementary proof 6}, The embedded-prime quotient and its exact kernel.
\item[ALGEBRA-RECON-436] Source lines 16467--16546: \hyperlink{algebra_embedded_regular_notes_20260927}{Supplementary proof 6}, Compatibility with arbitrary localization.
\end{description}

\subsection{Regular sequences}
\hypertarget{algebra-source-section-regular-sequences}{}
\noindent\href{https://github.com/stacks/stacks-project/blob/a04446e57ec1fbc252a871afcec7752fb2807b14/algebra.tex\#L16587}{Original source section}. \href{algebra.pdf\#nameddest=section-regular-sequences}{Accompanying chapter}.

\begin{description}
\item[ALGEBRA-RECON-437] Source lines 16638--16645: \hyperlink{algebra_embedded_regular_notes_20260927}{Supplementary proof 6}, Regular sequences retain their module and order.
\item[ALGEBRA-RECON-438] Source lines 16663--16672: \hyperlink{algebra_embedded_regular_notes_20260927}{Supplementary proof 6}, Regular sequences retain their module and order.
\item[ALGEBRA-RECON-439] Source lines 16675--17002: \hyperlink{algebra_embedded_regular_notes_20260927}{Supplementary proof 6}, Faithfully flat preservation and reflection.
\item[ALGEBRA-RECON-440] Source lines 16828--16844: \hyperlink{algebra_embedded_regular_notes_20260927}{Supplementary proof 6}, Regular sequences retain their module and order.
\end{description}

\subsection{Quasi-regular sequences}
\hypertarget{algebra-source-section-quasi-regular}{}
\noindent\href{https://github.com/stacks/stacks-project/blob/a04446e57ec1fbc252a871afcec7752fb2807b14/algebra.tex\#L16855}{Original source section}. \href{algebra.pdf\#nameddest=section-quasi-regular}{Accompanying chapter}.

\begin{description}
\item[ALGEBRA-RECON-439] Source lines 16675--17002: \hyperlink{algebra_embedded_regular_notes_20260927}{Supplementary proof 6}, Faithfully flat preservation and reflection.
\item[ALGEBRA-RECON-441] Source lines 16927--16970: \hyperlink{algebra_embedded_regular_notes_20260927}{Supplementary proof 6}, The two induction boundaries and module coefficients.
\item[ALGEBRA-RECON-442] Source lines 16972--16977: \hyperlink{algebra_embedded_regular_notes_20260927}{Supplementary proof 6}, The two induction boundaries and module coefficients.
\item[ALGEBRA-RECON-443] Source lines 17046--17056: \hyperlink{algebra_embedded_regular_notes_20260927}{Supplementary proof 6}, Truncation in degree zero and the original polynomial module.
\item[ALGEBRA-RECON-444] Source lines 17049--17056: \hyperlink{algebra_embedded_regular_notes_20260927}{Supplementary proof 6}, Truncation in degree zero and the original polynomial module.
\item[ALGEBRA-RECON-445] Source lines 17116--17130: \hyperlink{algebra_embedded_regular_notes_20260927}{Supplementary proof 6}, The separated quotient and its explicit ideal.
\end{description}

\subsection{Blow up algebras}
\hypertarget{algebra-source-section-blow-up}{}
\noindent\href{https://github.com/stacks/stacks-project/blob/a04446e57ec1fbc252a871afcec7752fb2807b14/algebra.tex\#L17136}{Original source section}. \href{algebra.pdf\#nameddest=section-blow-up}{Accompanying chapter}.

\begin{description}
\item[ALGEBRA-RECON-446] Source lines 17183--17207: \hyperlink{algebra_blowup_ext_depth_notes_20260927}{Supplementary proof 7}, Blowup fractions and base change.
\item[ALGEBRA-RECON-447] Source lines 17209--17226: \hyperlink{algebra_blowup_ext_depth_notes_20260927}{Supplementary proof 7}, Polynomial charts and the named base ring.
\item[ALGEBRA-RECON-448] Source lines 17249--17268: \hyperlink{algebra_blowup_ext_depth_notes_20260927}{Supplementary proof 7}, Polynomial charts and the named base ring.
\item[ALGEBRA-RECON-449] Source lines 17351--17391: \hyperlink{algebra_blowup_ext_depth_notes_20260927}{Supplementary proof 7}, The field boundary and the directed chart construction.
\item[ALGEBRA-RECON-450] Source lines 17384--17390: \hyperlink{algebra_blowup_ext_depth_notes_20260927}{Supplementary proof 7}, The field boundary and the directed chart construction.
\item[ALGEBRA-RECON-451] Source lines 17351--17391: \hyperlink{algebra_blowup_ext_depth_notes_20260927}{Supplementary proof 7}, The exact overring strengthening.
\end{description}

\subsection{Ext groups}
\hypertarget{algebra-source-section-ext}{}
\noindent\href{https://github.com/stacks/stacks-project/blob/a04446e57ec1fbc252a871afcec7752fb2807b14/algebra.tex\#L17402}{Original source section}. \href{algebra.pdf\#nameddest=section-ext}{Accompanying chapter}.

\begin{description}
\item[ALGEBRA-RECON-452] Source lines 17428--17435: \hyperlink{algebra_blowup_ext_depth_notes_20260927}{Supplementary proof 7}, Resolutions, chain maps and contravariance.
\item[ALGEBRA-RECON-453] Source lines 17453--17476: \hyperlink{algebra_blowup_ext_depth_notes_20260927}{Supplementary proof 7}, Resolutions, chain maps and contravariance.
\item[ALGEBRA-RECON-454] Source lines 17560--17563: \hyperlink{algebra_blowup_ext_depth_notes_20260927}{Supplementary proof 7}, Resolutions, chain maps and contravariance.
\item[ALGEBRA-RECON-455] Source lines 17574--17620: \hyperlink{algebra_blowup_ext_depth_notes_20260927}{Supplementary proof 7}, Resolutions, chain maps and contravariance.
\item[ALGEBRA-RECON-456] Source lines 17592--17595: \hyperlink{algebra_blowup_ext_depth_notes_20260927}{Supplementary proof 7}, Resolutions, chain maps and contravariance.
\item[ALGEBRA-RECON-457] Source lines 17607--17612: \hyperlink{algebra_blowup_ext_depth_notes_20260927}{Supplementary proof 7}, Resolutions, chain maps and contravariance.
\item[ALGEBRA-RECON-458] Source lines 17644--17657: \hyperlink{algebra_blowup_ext_depth_notes_20260927}{Supplementary proof 7}, Resolutions, chain maps and contravariance.
\item[ALGEBRA-RECON-459] Source lines 17681--17721: \hyperlink{algebra_blowup_ext_depth_notes_20260927}{Supplementary proof 7}, Resolutions, chain maps and contravariance.
\item[ALGEBRA-RECON-460] Source lines 17499--17741: \hyperlink{algebra_blowup_ext_depth_notes_20260927}{Supplementary proof 7}, Projective comparison with the same maps.
\end{description}

\subsection{Depth}
\hypertarget{algebra-source-section-depth}{}
\noindent\href{https://github.com/stacks/stacks-project/blob/a04446e57ec1fbc252a871afcec7752fb2807b14/algebra.tex\#L17759}{Original source section}. \href{algebra.pdf\#nameddest=section-depth}{Accompanying chapter}.

\begin{description}
\item[ALGEBRA-RECON-461] Source lines 17960--17982: \hyperlink{algebra_blowup_ext_depth_notes_20260927}{Supplementary proof 7}, Depth quotients and propagation to large powers.
\item[ALGEBRA-RECON-462] Source lines 17984--18020: \hyperlink{algebra_blowup_ext_depth_notes_20260927}{Supplementary proof 7}, Depth quotients and propagation to large powers.
\item[ALGEBRA-RECON-463] Source lines 17960--18020: \hyperlink{algebra_blowup_ext_depth_notes_20260927}{Supplementary proof 7}, Depth quotients and propagation to large powers.
\item[ALGEBRA-RECON-464] Source lines 18047--18082: \hyperlink{algebra_depth_tor_projective_notes_20260927}{Supplementary proof 8}, Finite-map depth and the infinite boundary.
\end{description}

\subsection{Functorialities for Ext}
\hypertarget{algebra-source-section-functoriality-ext}{}
\noindent\href{https://github.com/stacks/stacks-project/blob/a04446e57ec1fbc252a871afcec7752fb2807b14/algebra.tex\#L18087}{Original source section}. \href{algebra.pdf\#nameddest=section-functoriality-ext}{Accompanying chapter}.

\begin{description}
\item[ALGEBRA-RECON-465] Source lines 18103--18111: \hyperlink{algebra_depth_tor_projective_notes_20260927}{Supplementary proof 8}, Ext base change and the original splitting obstruction.
\item[ALGEBRA-RECON-466] Source lines 18120--18127: \hyperlink{algebra_depth_tor_projective_notes_20260927}{Supplementary proof 8}, Ext base change and the original splitting obstruction.
\end{description}

\subsection{An application of Ext groups}
\hypertarget{algebra-source-section-ext-application}{}
\noindent\href{https://github.com/stacks/stacks-project/blob/a04446e57ec1fbc252a871afcec7752fb2807b14/algebra.tex\#L18135}{Original source section}. \href{algebra.pdf\#nameddest=section-ext-application}{Accompanying chapter}.

\begin{description}
\item[ALGEBRA-RECON-467] Source lines 18152--18158: \hyperlink{algebra_depth_tor_projective_notes_20260927}{Supplementary proof 8}, Ext base change and the original splitting obstruction.
\item[ALGEBRA-RECON-468] Source lines 18141--18203: \hyperlink{algebra_depth_tor_projective_notes_20260927}{Supplementary proof 8}, Splitting on one principal neighbourhood.
\end{description}

\subsection{Tor groups and flatness}
\hypertarget{algebra-source-section-tor}{}
\noindent\href{https://github.com/stacks/stacks-project/blob/a04446e57ec1fbc252a871afcec7752fb2807b14/algebra.tex\#L18215}{Original source section}. \href{algebra.pdf\#nameddest=section-tor}{Accompanying chapter}.

\begin{description}
\item[ALGEBRA-RECON-469] Source lines 18416--18417: \hyperlink{algebra_depth_tor_projective_notes_20260927}{Supplementary proof 8}, The zig-zag quotient with its original signs.
\item[ALGEBRA-RECON-470] Source lines 18419--18449: \hyperlink{algebra_depth_tor_projective_notes_20260927}{Supplementary proof 8}, The zig-zag quotient with its original signs.
\item[ALGEBRA-RECON-471] Source lines 18465--18467: \hyperlink{algebra_depth_tor_projective_notes_20260927}{Supplementary proof 8}, The zig-zag quotient with its original signs.
\item[ALGEBRA-RECON-472] Source lines 18474--18476: \hyperlink{algebra_depth_tor_projective_notes_20260927}{Supplementary proof 8}, The zig-zag quotient with its original signs.
\item[ALGEBRA-RECON-473] Source lines 18487--18493: \hyperlink{algebra_depth_tor_projective_notes_20260927}{Supplementary proof 8}, The zig-zag quotient with its original signs.
\item[ALGEBRA-RECON-474] Source lines 18549--18551: \hyperlink{algebra_depth_tor_projective_notes_20260927}{Supplementary proof 8}, Tensor comparison, swap and finiteness.
\item[ALGEBRA-RECON-475] Source lines 18624--18631: \hyperlink{algebra_depth_tor_projective_notes_20260927}{Supplementary proof 8}, Tensor comparison, swap and finiteness.
\item[ALGEBRA-RECON-476] Source lines 18643--18646: \hyperlink{algebra_depth_tor_projective_notes_20260927}{Supplementary proof 8}, Tensor comparison, swap and finiteness.
\item[ALGEBRA-RECON-480] Source lines 18397--18620: \hyperlink{algebra_depth_tor_projective_notes_20260927}{Supplementary proof 8}, Flat resolutions compute the same Tor.
\end{description}

\subsection{Functorialities for Tor}
\hypertarget{algebra-source-section-functoriality-tor}{}
\noindent\href{https://github.com/stacks/stacks-project/blob/a04446e57ec1fbc252a871afcec7752fb2807b14/algebra.tex\#L18713}{Original source section}. \href{algebra.pdf\#nameddest=section-functoriality-tor}{Accompanying chapter}.

\begin{description}
\item[ALGEBRA-RECON-477] Source lines 18761--18764: \hyperlink{algebra_depth_tor_projective_notes_20260927}{Supplementary proof 8}, Tensor comparison, swap and finiteness.
\end{description}

\subsection{Projective modules}
\hypertarget{algebra-source-section-projective}{}
\noindent\href{https://github.com/stacks/stacks-project/blob/a04446e57ec1fbc252a871afcec7752fb2807b14/algebra.tex\#L18774}{Original source section}. \href{algebra.pdf\#nameddest=section-projective}{Accompanying chapter}.

\begin{description}
\item[ALGEBRA-RECON-478] Source lines 18821--18828: \hyperlink{algebra_depth_tor_projective_notes_20260927}{Supplementary proof 8}, Projectors and the existing projective definition.
\item[ALGEBRA-RECON-479] Source lines 18826--18827: \hyperlink{algebra_depth_tor_projective_notes_20260927}{Supplementary proof 8}, Projectors and the existing projective definition.
\item[ALGEBRA-RECON-481] Source lines 18974--18977: \hyperlink{algebra_projective_lifting_notes_20260927}{Supplementary proof 9}, Patching over the two original ideals.
\item[ALGEBRA-RECON-482] Source lines 18982--18986: \hyperlink{algebra_projective_lifting_notes_20260927}{Supplementary proof 9}, Patching over the two original ideals.
\item[ALGEBRA-RECON-491] Source lines 18897--19266: \hyperlink{algebra_projective_lifting_notes_20260927}{Supplementary proof 9}, Constant fibre rank and the exact nilradical kernel.
\end{description}

\subsection{Finite projective modules}
\hypertarget{algebra-source-section-finite-projective-modules}{}
\noindent\href{https://github.com/stacks/stacks-project/blob/a04446e57ec1fbc252a871afcec7752fb2807b14/algebra.tex\#L19000}{Original source section}. \href{algebra.pdf\#nameddest=section-finite-projective-modules}{Accompanying chapter}.

\begin{description}
\item[ALGEBRA-RECON-483] Source lines 19024--19024: \hyperlink{algebra_projective_lifting_notes_20260927}{Supplementary proof 9}, Finite-projective criteria and the localization element.
\item[ALGEBRA-RECON-484] Source lines 19083--19084: \hyperlink{algebra_projective_lifting_notes_20260927}{Supplementary proof 9}, Finite-projective criteria and the localization element.
\item[ALGEBRA-RECON-485] Source lines 19113--19118: \hyperlink{algebra_projective_lifting_notes_20260927}{Supplementary proof 9}, Finite-projective criteria and the localization element.
\item[ALGEBRA-RECON-486] Source lines 19161--19165: \hyperlink{algebra_projective_lifting_notes_20260927}{Supplementary proof 9}, Finite-projective criteria and the localization element.
\item[ALGEBRA-RECON-487] Source lines 19212--19220: \hyperlink{algebra_projective_lifting_notes_20260927}{Supplementary proof 9}, The smooth-function counterexample.
\item[ALGEBRA-RECON-488] Source lines 19279--19283: \hyperlink{algebra_projective_lifting_notes_20260927}{Supplementary proof 9}, Local flatness, descent and semilocal bases.
\item[ALGEBRA-RECON-489] Source lines 19299--19301: \hyperlink{algebra_projective_lifting_notes_20260927}{Supplementary proof 9}, Local flatness, descent and semilocal bases.
\item[ALGEBRA-RECON-490] Source lines 19349--19352: \hyperlink{algebra_projective_lifting_notes_20260927}{Supplementary proof 9}, Evaluation and the converse dual-basis criterion.
\item[ALGEBRA-RECON-491] Source lines 18897--19266: \hyperlink{algebra_projective_lifting_notes_20260927}{Supplementary proof 9}, Constant fibre rank and the exact nilradical kernel.
\item[ALGEBRA-RECON-492] Source lines 19349--19369: \hyperlink{algebra_projective_lifting_notes_20260927}{Supplementary proof 9}, Evaluation and the converse dual-basis criterion.
\end{description}

\subsection{Open loci defined by module maps}
\hypertarget{algebra-source-section-loci-maps}{}
\noindent\href{https://github.com/stacks/stacks-project/blob/a04446e57ec1fbc252a871afcec7752fb2807b14/algebra.tex\#L19375}{Original source section}. \href{algebra.pdf\#nameddest=section-loci-maps}{Accompanying chapter}.

\begin{description}
\item[ALGEBRA-RECON-493] Source lines 19424--19431: \hyperlink{algebra_open_loci_pure_descent_notes_20260927}{Supplementary proof 10}, The original isomorphism neighbourhood.
\item[ALGEBRA-RECON-494] Source lines 19430--19436: \hyperlink{algebra_open_loci_pure_descent_notes_20260927}{Supplementary proof 10}, The original isomorphism neighbourhood.
\item[ALGEBRA-RECON-495] Source lines 19523--19527: \hyperlink{algebra_open_loci_pure_descent_notes_20260927}{Supplementary proof 10}, Split kernels and the original minors.
\item[ALGEBRA-RECON-500] Source lines 19383--19406: \hyperlink{algebra_open_loci_pure_descent_notes_20260927}{Supplementary proof 10}, Surjectivity needs only a finite cokernel.
\end{description}

\subsection{Characterizing flatness}
\hypertarget{algebra-source-section-characterize-flatness}{}
\noindent\href{https://github.com/stacks/stacks-project/blob/a04446e57ec1fbc252a871afcec7752fb2807b14/algebra.tex\#L19583}{Original source section}. \href{algebra.pdf\#nameddest=section-characterize-flatness}{Accompanying chapter}.

\begin{description}
\item[ALGEBRA-RECON-501] Source lines 19635--20356: \hyperlink{algebra_open_loci_pure_descent_notes_20260927}{Supplementary proof 10}, Descent along the original universally injective ring map.
\end{description}

\subsection{Universally injective module maps}
\hypertarget{algebra-source-section-universally-injective}{}
\noindent\href{https://github.com/stacks/stacks-project/blob/a04446e57ec1fbc252a871afcec7752fb2807b14/algebra.tex\#L19742}{Original source section}. \href{algebra.pdf\#nameddest=section-universally-injective}{Accompanying chapter}.

\begin{description}
\item[ALGEBRA-RECON-496] Source lines 19995--20038: \hyperlink{algebra_open_loci_pure_descent_notes_20260927}{Supplementary proof 10}, The nonzero boundary in the examples.
\item[ALGEBRA-RECON-497] Source lines 20165--20194: \hyperlink{algebra_open_loci_pure_descent_notes_20260927}{Supplementary proof 10}, Localization and the rejected proof-gap report.
\item[ALGEBRA-RECON-501] Source lines 19635--20356: \hyperlink{algebra_open_loci_pure_descent_notes_20260927}{Supplementary proof 10}, Descent along the original universally injective ring map.
\end{description}

\subsection{Descent for finite projective modules}
\hypertarget{algebra-source-section-finite-projective}{}
\noindent\href{https://github.com/stacks/stacks-project/blob/a04446e57ec1fbc252a871afcec7752fb2807b14/algebra.tex\#L20274}{Original source section}. \href{algebra.pdf\#nameddest=section-finite-projective}{Accompanying chapter}.

\begin{description}
\item[ALGEBRA-RECON-498] Source lines 20306--20319: \hyperlink{algebra_open_loci_pure_descent_notes_20260927}{Supplementary proof 10}, The finite descent statement and coefficients.
\item[ALGEBRA-RECON-499] Source lines 20322--20324: \hyperlink{algebra_open_loci_pure_descent_notes_20260927}{Supplementary proof 10}, The finite descent statement and coefficients.
\item[ALGEBRA-RECON-501] Source lines 19635--20356: \hyperlink{algebra_open_loci_pure_descent_notes_20260927}{Supplementary proof 10}, Descent along the original universally injective ring map.
\end{description}

\subsection{Transfinite d\'evissage of modules}
\hypertarget{algebra-source-section-transfinite-devissage}{}
\noindent\href{https://github.com/stacks/stacks-project/blob/a04446e57ec1fbc252a871afcec7752fb2807b14/algebra.tex\#L20374}{Original source section}. \href{algebra.pdf\#nameddest=section-transfinite-devissage}{Accompanying chapter}.

\begin{description}
\item[ALGEBRA-RECON-502] Source lines 20409--20453: \hyperlink{algebra_transfinite_devissage_notes_20260927}{Supplementary proof 11}, Sections exist only at successor stages in the index.
\item[ALGEBRA-RECON-503] Source lines 20455--20467: \hyperlink{algebra_transfinite_devissage_notes_20260927}{Supplementary proof 11}, Include the terminal partial sum.
\item[ALGEBRA-RECON-504] Source lines 20469--20585: \hyperlink{algebra_transfinite_devissage_notes_20260927}{Supplementary proof 11}, The original projector closure and a bounded stopping stage.
\end{description}

\subsection{Mittag-Leffler systems}
\hypertarget{algebra-source-section-mittag-leffler}{}
\noindent\href{https://github.com/stacks/stacks-project/blob/a04446e57ec1fbc252a871afcec7752fb2807b14/algebra.tex\#L20697}{Original source section}. \href{algebra.pdf\#nameddest=section-mittag-leffler}{Accompanying chapter}.

\begin{description}
\item[ALGEBRA-RECON-510] Source lines 20697--20877: \hyperlink{algebra_ml_review_notes_20260927}{Supplementary proof 12}, Stable images and countable cofinality.
\end{description}

\subsection{Inverse systems}
\hypertarget{algebra-source-section-inverse-systems}{}
\noindent\href{https://github.com/stacks/stacks-project/blob/a04446e57ec1fbc252a871afcec7752fb2807b14/algebra.tex\#L20816}{Original source section}. \href{algebra.pdf\#nameddest=section-inverse-systems}{Accompanying chapter}.

\begin{description}
\item[ALGEBRA-RECON-510] Source lines 20697--20877: \hyperlink{algebra_ml_review_notes_20260927}{Supplementary proof 12}, Stable images and countable cofinality.
\end{description}

\subsection{Mittag-Leffler modules}
\hypertarget{algebra-source-section-mittag-leffler-modules}{}
\noindent\href{https://github.com/stacks/stacks-project/blob/a04446e57ec1fbc252a871afcec7752fb2807b14/algebra.tex\#L20881}{Original source section}. \href{algebra.pdf\#nameddest=section-mittag-leffler-modules}{Accompanying chapter}.

\begin{description}
\item[ALGEBRA-RECON-505] Source lines 20999--21057: \hyperlink{algebra_ml_review_notes_20260927}{Supplementary proof 12}, Equalize the lifted maps at a finite stage.
\item[ALGEBRA-RECON-506] Source lines 21078--21103: \hyperlink{algebra_ml_review_notes_20260927}{Supplementary proof 12}, Use a common upper bound for incomparable indices.
\item[ALGEBRA-RECON-507] Source lines 21132--21160: \hyperlink{algebra_ml_review_notes_20260927}{Supplementary proof 12}, Retain the two tensor corrections and the finite-free test.
\item[ALGEBRA-RECON-508] Source lines 21162--21197: \hyperlink{algebra_ml_review_notes_20260927}{Supplementary proof 12}, Retain the two tensor corrections and the finite-free test.
\item[ALGEBRA-RECON-511] Source lines 20999--21274: \hyperlink{algebra_ml_review_notes_20260927}{Supplementary proof 12}, Mittag-Leffler base change and descent along pure ring maps.
\end{description}

\subsection{Interchanging direct products with tensor}
\hypertarget{algebra-source-section-products-tensor}{}
\noindent\href{https://github.com/stacks/stacks-project/blob/a04446e57ec1fbc252a871afcec7752fb2807b14/algebra.tex\#L21277}{Original source section}. \href{algebra.pdf\#nameddest=section-products-tensor}{Accompanying chapter}.

\begin{description}
\item[ALGEBRA-RECON-509] Source lines 21315--21356: \hyperlink{algebra_ml_review_notes_20260927}{Supplementary proof 12}, The coefficient projection and the original product criteria.
\item[ALGEBRA-RECON-527] Source lines 21533--21952: \hyperlink{algebra_ml_coherent_notes_20260927}{Supplementary proof 13}, Arbitrary intersections and base change of the smallest submodule.
\end{description}

\subsection{Coherent rings}
\hypertarget{algebra-source-section-coherent}{}
\noindent\href{https://github.com/stacks/stacks-project/blob/a04446e57ec1fbc252a871afcec7752fb2807b14/algebra.tex\#L21690}{Original source section}. \href{algebra.pdf\#nameddest=section-coherent}{Accompanying chapter}.

\begin{description}
\item[ALGEBRA-RECON-512] Source lines 21748--21761: \hyperlink{algebra_ml_coherent_notes_20260927}{Supplementary proof 13}, Coherent-module extensions and the retained image correction.
\item[ALGEBRA-RECON-527] Source lines 21533--21952: \hyperlink{algebra_ml_coherent_notes_20260927}{Supplementary proof 13}, Arbitrary intersections and base change of the smallest submodule.
\end{description}

\subsection{Examples and non-examples of Mittag-Leffler modules}
\hypertarget{algebra-source-section-examples}{}
\noindent\href{https://github.com/stacks/stacks-project/blob/a04446e57ec1fbc252a871afcec7752fb2807b14/algebra.tex\#L21852}{Original source section}. \href{algebra.pdf\#nameddest=section-examples}{Accompanying chapter}.

\begin{description}
\item[ALGEBRA-RECON-513] Source lines 21859--21874: \hyperlink{algebra_ml_coherent_notes_20260927}{Supplementary proof 13}, The finite-free content criterion and products over the source ring.
\item[ALGEBRA-RECON-514] Source lines 21876--21878: \hyperlink{algebra_ml_coherent_notes_20260927}{Supplementary proof 13}, The finite-free content criterion and products over the source ring.
\item[ALGEBRA-RECON-515] Source lines 21880--21952: \hyperlink{algebra_ml_coherent_notes_20260927}{Supplementary proof 13}, The finite-free content criterion and products over the source ring.
\item[ALGEBRA-RECON-516] Source lines 21954--21981: \hyperlink{algebra_ml_coherent_notes_20260927}{Supplementary proof 13}, The finite-free content criterion and products over the source ring.
\item[ALGEBRA-RECON-517] Source lines 21995--22025: \hyperlink{algebra_ml_coherent_notes_20260927}{Supplementary proof 13}, The non-ML examples with their original elements.
\item[ALGEBRA-RECON-518] Source lines 22020--22025: \hyperlink{algebra_ml_coherent_notes_20260927}{Supplementary proof 13}, The non-ML examples with their original elements.
\item[ALGEBRA-RECON-527] Source lines 21533--21952: \hyperlink{algebra_ml_coherent_notes_20260927}{Supplementary proof 13}, Arbitrary intersections and base change of the smallest submodule.
\end{description}

\subsection{Countably generated Mittag-Leffler modules}
\hypertarget{algebra-source-section-ML-countable}{}
\noindent\href{https://github.com/stacks/stacks-project/blob/a04446e57ec1fbc252a871afcec7752fb2807b14/algebra.tex\#L22045}{Original source section}. \href{algebra.pdf\#nameddest=section-ML-countable}{Accompanying chapter}.

\begin{description}
\item[ALGEBRA-RECON-519] Source lines 22062--22071: \hyperlink{algebra_ml_coherent_notes_20260927}{Supplementary proof 13}, Countable subsystem, parity and cofinal sequences.
\item[ALGEBRA-RECON-520] Source lines 22062--22075: \hyperlink{algebra_ml_coherent_notes_20260927}{Supplementary proof 13}, Countable subsystem, parity and cofinal sequences.
\item[ALGEBRA-RECON-521] Source lines 22097--22114: \hyperlink{algebra_ml_coherent_notes_20260927}{Supplementary proof 13}, Countable subsystem, parity and cofinal sequences.
\end{description}

\subsection{Characterizing projective modules}
\hypertarget{algebra-source-section-characterize-projective}{}
\noindent\href{https://github.com/stacks/stacks-project/blob/a04446e57ec1fbc252a871afcec7752fb2807b14/algebra.tex\#L22184}{Original source section}. \href{algebra.pdf\#nameddest=section-characterize-projective}{Accompanying chapter}.

\begin{description}
\item[ALGEBRA-RECON-522] Source lines 22272--22285: \hyperlink{algebra_ml_coherent_notes_20260927}{Supplementary proof 13}, Projectivity and the original countability boundary.
\item[ALGEBRA-RECON-523] Source lines 22287--22299: \hyperlink{algebra_ml_coherent_notes_20260927}{Supplementary proof 13}, Projectivity and the original countability boundary.
\item[ALGEBRA-RECON-528] Source lines 22248--22417: \hyperlink{algebra_ml_coherent_notes_20260927}{Supplementary proof 13}, Pure descent of countable generation and the resulting projective cases.
\end{description}

\subsection{Ascending properties of modules}
\hypertarget{algebra-source-section-ascent}{}
\noindent\href{https://github.com/stacks/stacks-project/blob/a04446e57ec1fbc252a871afcec7752fb2807b14/algebra.tex\#L22303}{Original source section}. \href{algebra.pdf\#nameddest=section-ascent}{Accompanying chapter}.

\begin{description}
\item[ALGEBRA-RECON-528] Source lines 22248--22417: \hyperlink{algebra_ml_coherent_notes_20260927}{Supplementary proof 13}, Pure descent of countable generation and the resulting projective cases.
\end{description}

\subsection{Descending properties of modules}
\hypertarget{algebra-source-section-descent}{}
\noindent\href{https://github.com/stacks/stacks-project/blob/a04446e57ec1fbc252a871afcec7752fb2807b14/algebra.tex\#L22332}{Original source section}. \href{algebra.pdf\#nameddest=section-descent}{Accompanying chapter}.

\begin{description}
\item[ALGEBRA-RECON-524] Source lines 22341--22382: \hyperlink{algebra_ml_coherent_notes_20260927}{Supplementary proof 13}, Source ascent, descent and the countable image construction.
\item[ALGEBRA-RECON-525] Source lines 22368--22382: \hyperlink{algebra_ml_coherent_notes_20260927}{Supplementary proof 13}, Source ascent, descent and the countable image construction.
\item[ALGEBRA-RECON-526] Source lines 22423--22435: \hyperlink{algebra_ml_coherent_notes_20260927}{Supplementary proof 13}, Source ascent, descent and the countable image construction.
\item[ALGEBRA-RECON-528] Source lines 22248--22417: \hyperlink{algebra_ml_coherent_notes_20260927}{Supplementary proof 13}, Pure descent of countable generation and the resulting projective cases.
\item[ALGEBRA-RECON-529] Source lines 22482--22501: \hyperlink{algebra_projectivity_completion_notes_20260927}{Supplementary proof 14}, Faithfully flat projectivity descent and the retained corrections.
\item[ALGEBRA-RECON-530] Source lines 22461--22515: \hyperlink{algebra_projectivity_completion_notes_20260927}{Supplementary proof 14}, Faithfully flat projectivity descent and the retained corrections.
\item[ALGEBRA-RECON-531] Source lines 22503--22515: \hyperlink{algebra_projectivity_completion_notes_20260927}{Supplementary proof 14}, Faithfully flat projectivity descent and the retained corrections.
\item[ALGEBRA-RECON-540] Source lines 22461--23008: \hyperlink{algebra_projectivity_completion_notes_20260927}{Supplementary proof 14}, Faithful flatness criterion and projectivity after completion.
\end{description}

\subsection{Completion}
\hypertarget{algebra-source-section-completion}{}
\noindent\href{https://github.com/stacks/stacks-project/blob/a04446e57ec1fbc252a871afcec7752fb2807b14/algebra.tex\#L22528}{Original source section}. \href{algebra.pdf\#nameddest=section-completion}{Accompanying chapter}.

\begin{description}
\item[ALGEBRA-RECON-532] Source lines 22755--22773: \hyperlink{algebra_projectivity_completion_notes_20260927}{Supplementary proof 14}, Units and the principal-generator completion argument.
\item[ALGEBRA-RECON-533] Source lines 22775--22803: \hyperlink{algebra_projectivity_completion_notes_20260927}{Supplementary proof 14}, Units and the principal-generator completion argument.
\item[ALGEBRA-RECON-534] Source lines 22859--22883: \hyperlink{algebra_projectivity_completion_notes_20260927}{Supplementary proof 14}, Quotients and finite modules over a complete ring.
\item[ALGEBRA-RECON-535] Source lines 22871--22883: \hyperlink{algebra_projectivity_completion_notes_20260927}{Supplementary proof 14}, Quotients and finite modules over a complete ring.
\item[ALGEBRA-RECON-539] Source lines 22701--22731: \hyperlink{algebra_projectivity_completion_notes_20260927}{Supplementary proof 14}, The exact split defect of a second completion.
\item[ALGEBRA-RECON-540] Source lines 22461--23008: \hyperlink{algebra_projectivity_completion_notes_20260927}{Supplementary proof 14}, Faithful flatness criterion and projectivity after completion.
\end{description}

\subsection{Completion for Noetherian rings}
\hypertarget{algebra-source-section-completion-noetherian}{}
\noindent\href{https://github.com/stacks/stacks-project/blob/a04446e57ec1fbc252a871afcec7752fb2807b14/algebra.tex\#L22893}{Original source section}. \href{algebra.pdf\#nameddest=section-completion-noetherian}{Accompanying chapter}.

\begin{description}
\item[ALGEBRA-RECON-536] Source lines 22900--22953: \hyperlink{algebra_projectivity_completion_notes_20260927}{Supplementary proof 14}, Noetherian exactness and the tensor comparison.
\item[ALGEBRA-RECON-537] Source lines 22955--22974: \hyperlink{algebra_projectivity_completion_notes_20260927}{Supplementary proof 14}, Noetherian exactness and the tensor comparison.
\item[ALGEBRA-RECON-538] Source lines 22976--22992: \hyperlink{algebra_projectivity_completion_notes_20260927}{Supplementary proof 14}, Noetherian exactness and the tensor comparison.
\item[ALGEBRA-RECON-540] Source lines 22461--23008: \hyperlink{algebra_projectivity_completion_notes_20260927}{Supplementary proof 14}, Faithful flatness criterion and projectivity after completion.
\item[ALGEBRA-RECON-541] Source lines 23010--23062: \hyperlink{algebra_graded_completion_notes_20260927}{Supplementary proof 15}, Associated-graded generation and the completion map.
\item[ALGEBRA-RECON-542] Source lines 23099--23132: \hyperlink{algebra_graded_completion_notes_20260927}{Supplementary proof 15}, Local completion and the redundant maximal-ideal assumption.
\item[ALGEBRA-RECON-543] Source lines 23215--23237: \hyperlink{algebra_graded_completion_notes_20260927}{Supplementary proof 15}, Compatible splittings and successive completion.
\item[ALGEBRA-RECON-544] Source lines 23215--23237: \hyperlink{algebra_graded_completion_notes_20260927}{Supplementary proof 15}, Compatible splittings and successive completion.
\item[ALGEBRA-RECON-547] Source lines 23010--23132: \hyperlink{algebra_graded_completion_notes_20260927}{Supplementary proof 15}, Local completion and the redundant maximal-ideal assumption.
\end{description}

\subsection{Taking limits of modules}
\hypertarget{algebra-source-section-limits}{}
\noindent\href{https://github.com/stacks/stacks-project/blob/a04446e57ec1fbc252a871afcec7752fb2807b14/algebra.tex\#L23247}{Original source section}. \href{algebra.pdf\#nameddest=section-limits}{Accompanying chapter}.

\begin{description}
\item[ALGEBRA-RECON-545] Source lines 23299--23333: \hyperlink{algebra_graded_completion_notes_20260927}{Supplementary proof 15}, Correct stabilization of graded modules.
\item[ALGEBRA-RECON-546] Source lines 23335--23401: \hyperlink{algebra_graded_completion_notes_20260927}{Supplementary proof 15}, The shifted degree bound and separate injectivity and surjectivity.
\item[ALGEBRA-RECON-548] Source lines 23299--23333: \hyperlink{algebra_graded_completion_notes_20260927}{Supplementary proof 15}, Graded algebraization without Noetherianity.
\item[ALGEBRA-RECON-549] Source lines 23335--23401: \hyperlink{algebra_graded_completion_notes_20260927}{Supplementary proof 15}, The shifted degree bound and separate injectivity and surjectivity.
\end{description}

\subsection{Criteria for flatness}
\hypertarget{algebra-source-section-criteria-flatness}{}
\noindent\href{https://github.com/stacks/stacks-project/blob/a04446e57ec1fbc252a871afcec7752fb2807b14/algebra.tex\#L23415}{Original source section}. \href{algebra.pdf\#nameddest=section-criteria-flatness}{Accompanying chapter}.

\begin{description}
\item[ALGEBRA-RECON-550] Source lines 23423--23504: \hyperlink{algebra_flatness_criteria_notes_20260928}{Supplementary proof 16}, Injection modulo the maximal ideal and its original quotients.
\item[ALGEBRA-RECON-551] Source lines 23559--23580: \hyperlink{algebra_flatness_criteria_notes_20260928}{Supplementary proof 16}, Finite length, Artin--Rees and the local criterion.
\item[ALGEBRA-RECON-552] Source lines 23582--23641: \hyperlink{algebra_flatness_criteria_notes_20260928}{Supplementary proof 16}, Finite length, Artin--Rees and the local criterion.
\item[ALGEBRA-RECON-553] Source lines 23601--23640: \hyperlink{algebra_flatness_criteria_notes_20260928}{Supplementary proof 16}, Finite length, Artin--Rees and the local criterion.
\item[ALGEBRA-RECON-554] Source lines 23650--23716: \hyperlink{algebra_flatness_criteria_notes_20260928}{Supplementary proof 16}, The correct free index and the ideal-adic Tor criterion.
\item[ALGEBRA-RECON-555] Source lines 23664--23691: \hyperlink{algebra_flatness_criteria_notes_20260928}{Supplementary proof 16}, The correct free index and the ideal-adic Tor criterion.
\item[ALGEBRA-RECON-556] Source lines 23880--23918: \hyperlink{algebra_flatness_criteria_notes_20260928}{Supplementary proof 16}, Tor comparisons and the exact change-of-rings sequence.
\item[ALGEBRA-RECON-557] Source lines 23920--23968: \hyperlink{algebra_flatness_criteria_notes_20260928}{Supplementary proof 16}, The localized comparison criterion and the fibre criterion.
\item[ALGEBRA-RECON-558] Source lines 24025--24168: \hyperlink{algebra_flatness_criteria_notes_20260928}{Supplementary proof 16}, The localization gluing lemmas and faithful flatness.
\item[ALGEBRA-RECON-559] Source lines 23423--24023: \hyperlink{algebra_flatness_criteria_notes_20260928}{Supplementary proof 16}, Removing the unused base Noetherian hypothesis.
\item[ALGEBRA-RECON-560] Source lines 23650--23843: \hyperlink{algebra_flatness_criteria_notes_20260928}{Supplementary proof 16}, Elementwise ideal-power torsion and the associated graded map.
\item[ALGEBRA-RECON-561] Source lines 23845--23968: \hyperlink{algebra_flatness_criteria_notes_20260928}{Supplementary proof 16}, Tor comparisons and the exact change-of-rings sequence.
\end{description}

\subsection{Base change and flatness}
\hypertarget{algebra-source-section-base-change-flat}{}
\noindent\href{https://github.com/stacks/stacks-project/blob/a04446e57ec1fbc252a871afcec7752fb2807b14/algebra.tex\#L24177}{Original source section}. \href{algebra.pdf\#nameddest=section-base-change-flat}{Accompanying chapter}.

\begin{description}
\item[ALGEBRA-RECON-562] Source lines 24184--24237: \hyperlink{algebra_nilpotent_flatness_notes_20260928}{Supplementary proof 17}, The original base-change and local-ring diagrams.
\end{description}

\subsection{Flatness criteria over Artinian rings}
\hypertarget{algebra-source-section-flatness-artinian}{}
\noindent\href{https://github.com/stacks/stacks-project/blob/a04446e57ec1fbc252a871afcec7752fb2807b14/algebra.tex\#L24286}{Original source section}. \href{algebra.pdf\#nameddest=section-flatness-artinian}{Accompanying chapter}.

\begin{description}
\item[ALGEBRA-RECON-563] Source lines 24470--24472: \hyperlink{algebra_nilpotent_flatness_notes_20260928}{Supplementary proof 17}, The Artinian statement and the retained module corrections.
\item[ALGEBRA-RECON-564] Source lines 24494--24500: \hyperlink{algebra_nilpotent_flatness_notes_20260928}{Supplementary proof 17}, The Artinian statement and the retained module corrections.
\item[ALGEBRA-RECON-565] Source lines 24501--24517: \hyperlink{algebra_nilpotent_flatness_notes_20260928}{Supplementary proof 17}, The Artinian statement and the retained module corrections.
\item[ALGEBRA-RECON-566] Source lines 24528--24567: \hyperlink{algebra_nilpotent_flatness_notes_20260928}{Supplementary proof 17}, The nilpotent fibre criterion and its faithful support.
\item[ALGEBRA-RECON-567] Source lines 24383--24442: \hyperlink{algebra_nilpotent_flatness_notes_20260928}{Supplementary proof 17}, Exact contraction powers and the kernel quotient.
\item[ALGEBRA-RECON-568] Source lines 24418--24517: \hyperlink{algebra_nilpotent_flatness_notes_20260928}{Supplementary proof 17}, Injective descent beyond Artinian source rings.
\end{description}

\subsection{What makes a complex exact?}
\hypertarget{algebra-source-section-complex-exact}{}
\noindent\href{https://github.com/stacks/stacks-project/blob/a04446e57ec1fbc252a871afcec7752fb2807b14/algebra.tex\#L24576}{Original source section}. \href{algebra.pdf\#nameddest=section-complex-exact}{Accompanying chapter}.

\begin{description}
\item[ALGEBRA-RECON-569] Source lines 24607--24654: \hyperlink{algebra_complex_exactness_notes_20260928}{Supplementary proof 18}, The original pivot and its exact splitting maps.
\item[ALGEBRA-RECON-570] Source lines 24689--24789: \hyperlink{algebra_complex_exactness_notes_20260928}{Supplementary proof 18}, Depth-zero splitting and the notation corrections.
\item[ALGEBRA-RECON-571] Source lines 24702--24713: \hyperlink{algebra_complex_exactness_notes_20260928}{Supplementary proof 18}, Depth-zero splitting and the notation corrections.
\item[ALGEBRA-RECON-572] Source lines 24721--24763: \hyperlink{algebra_complex_exactness_notes_20260928}{Supplementary proof 18}, Determinantal ranks and the zero-ring boundary.
\item[ALGEBRA-RECON-573] Source lines 24738--24763: \hyperlink{algebra_complex_exactness_notes_20260928}{Supplementary proof 18}, Determinantal ranks and the zero-ring boundary.
\item[ALGEBRA-RECON-574] Source lines 24792--24831: \hyperlink{algebra_complex_exactness_notes_20260928}{Supplementary proof 18}, Reduction by a nonzerodivisor and the acyclicity proof.
\item[ALGEBRA-RECON-575] Source lines 24810--24830: \hyperlink{algebra_complex_exactness_notes_20260928}{Supplementary proof 18}, Reduction by a nonzerodivisor and the acyclicity proof.
\item[ALGEBRA-RECON-576] Source lines 24854--24924: \hyperlink{algebra_complex_exactness_notes_20260928}{Supplementary proof 18}, The exactness criterion with all rank-zero cases retained.
\item[ALGEBRA-RECON-577] Source lines 24887--24892: \hyperlink{algebra_complex_exactness_notes_20260928}{Supplementary proof 18}, The exactness criterion with all rank-zero cases retained.
\item[ALGEBRA-RECON-578] Source lines 24607--24713: \hyperlink{algebra_complex_exactness_notes_20260928}{Supplementary proof 18}, Splitting characterized beyond Noetherian depth zero.
\item[ALGEBRA-RECON-579] Source lines 24792--24924: \hyperlink{algebra_complex_exactness_notes_20260928}{Supplementary proof 18}, A stronger depth conclusion for the highest nonzero homology.
\end{description}

\subsection{Cohen-Macaulay modules}
\hypertarget{algebra-source-section-CM}{}
\noindent\href{https://github.com/stacks/stacks-project/blob/a04446e57ec1fbc252a871afcec7752fb2807b14/algebra.tex\#L24952}{Original source section}. \href{algebra.pdf\#nameddest=section-CM}{Accompanying chapter}.

\begin{description}
\item[ALGEBRA-RECON-580] Source lines 24959--25282: \hyperlink{algebra_cm_modules_notes_20260928}{Supplementary proof 19}, The zero module and the good-element boundary.
\item[ALGEBRA-RECON-581] Source lines 24977--25015: \hyperlink{algebra_cm_modules_notes_20260928}{Supplementary proof 19}, The zero module and the good-element boundary.
\item[ALGEBRA-RECON-582] Source lines 25027--25249: \hyperlink{algebra_cm_modules_notes_20260928}{Supplementary proof 19}, One-element cuts and complete parameter sequences.
\item[ALGEBRA-RECON-583] Source lines 25051--25075: \hyperlink{algebra_cm_modules_notes_20260928}{Supplementary proof 19}, One-element cuts and complete parameter sequences.
\item[ALGEBRA-RECON-584] Source lines 25201--25225: \hyperlink{algebra_cm_modules_notes_20260928}{Supplementary proof 19}, Quotients, associated primes, chains and localization.
\item[ALGEBRA-RECON-590] Source lines 25096--25225: \hyperlink{algebra_cm_modules_notes_20260928}{Supplementary proof 19}, Support-relative dimension and chain formulas.
\item[ALGEBRA-RECON-591] Source lines 25051--25132: \hyperlink{algebra_cm_modules_notes_20260928}{Supplementary proof 19}, Every positive parameter power and its associated primes.
\end{description}

\subsection{Cohen-Macaulay rings}
\hypertarget{algebra-source-section-CM-ring}{}
\noindent\href{https://github.com/stacks/stacks-project/blob/a04446e57ec1fbc252a871afcec7752fb2807b14/algebra.tex\#L25298}{Original source section}. \href{algebra.pdf\#nameddest=section-CM-ring}{Accompanying chapter}.

\begin{description}
\item[ALGEBRA-RECON-585] Source lines 25311--25316: \hyperlink{algebra_cm_modules_notes_20260928}{Supplementary proof 19}, The ring specializations and the exact syzygy depths.
\item[ALGEBRA-RECON-586] Source lines 25344--25355: \hyperlink{algebra_cm_modules_notes_20260928}{Supplementary proof 19}, The ring specializations and the exact syzygy depths.
\item[ALGEBRA-RECON-587] Source lines 25344--25355: \hyperlink{algebra_cm_modules_notes_20260928}{Supplementary proof 19}, The ring specializations and the exact syzygy depths.
\item[ALGEBRA-RECON-588] Source lines 25388--25396: \hyperlink{algebra_cm_modules_notes_20260928}{Supplementary proof 19}, The complete polynomial-module argument.
\item[ALGEBRA-RECON-589] Source lines 25398--25423: \hyperlink{algebra_cm_modules_notes_20260928}{Supplementary proof 19}, The ring specializations and the exact syzygy depths.
\item[ALGEBRA-RECON-592] Source lines 25425--25447: \hyperlink{algebra_cm_modules_notes_20260928}{Supplementary proof 19}, The original image-parameter construction and its module extension.
\end{description}

\subsection{Catenary rings}
\hypertarget{algebra-source-section-catenary}{}
\noindent\href{https://github.com/stacks/stacks-project/blob/a04446e57ec1fbc252a871afcec7752fb2807b14/algebra.tex\#L25456}{Original source section}. \href{algebra.pdf\#nameddest=section-catenary}{Accompanying chapter}.

\begin{description}
\item[ALGEBRA-RECON-593] Source lines 25534--25572: \hyperlink{algebra_catenary_regular_notes_20260928}{Supplementary proof 20}, Prime intervals, localization and the omitted universal argument.
\item[ALGEBRA-RECON-594] Source lines 25574--25604: \hyperlink{algebra_catenary_regular_notes_20260928}{Supplementary proof 20}, Arbitrary closed covers and nil-ideal invariance.
\item[ALGEBRA-RECON-595] Source lines 25640--25664: \hyperlink{algebra_catenary_regular_notes_20260928}{Supplementary proof 20}, The local dimension function without Noetherianity.
\item[ALGEBRA-RECON-596] Source lines 25606--25638: \hyperlink{algebra_catenary_regular_notes_20260928}{Supplementary proof 20}, Universal catenarity of module supports and covering families.
\end{description}

\subsection{Regular local rings}
\hypertarget{algebra-source-section-regular}{}
\noindent\href{https://github.com/stacks/stacks-project/blob/a04446e57ec1fbc252a871afcec7752fb2807b14/algebra.tex\#L25679}{Original source section}. \href{algebra.pdf\#nameddest=section-regular}{Accompanying chapter}.

\begin{description}
\item[ALGEBRA-RECON-597] Source lines 25714--25729: \hyperlink{algebra_catenary_regular_notes_20260928}{Supplementary proof 20}, Regular graded rings, domains and parameter quotients.
\item[ALGEBRA-RECON-598] Source lines 25758--25761: \hyperlink{algebra_catenary_regular_notes_20260928}{Supplementary proof 20}, Regular graded rings, domains and parameter quotients.
\item[ALGEBRA-RECON-599] Source lines 25759--25776: \hyperlink{algebra_catenary_regular_notes_20260928}{Supplementary proof 20}, Regular graded rings, domains and parameter quotients.
\item[ALGEBRA-RECON-600] Source lines 25789--25797: \hyperlink{algebra_catenary_regular_notes_20260928}{Supplementary proof 20}, Free lifting, maximal CM modules and regular lifting.
\item[ALGEBRA-RECON-601] Source lines 25840--25844: \hyperlink{algebra_catenary_regular_notes_20260928}{Supplementary proof 20}, Directed colimits with noninjective transition maps.
\item[ALGEBRA-RECON-602] Source lines 25840--25855: \hyperlink{algebra_catenary_regular_notes_20260928}{Supplementary proof 20}, Directed colimits with noninjective transition maps.
\item[ALGEBRA-RECON-603] Source lines 25845--25855: \hyperlink{algebra_catenary_regular_notes_20260928}{Supplementary proof 20}, Directed colimits with noninjective transition maps.
\item[ALGEBRA-RECON-604] Source lines 25779--25797: \hyperlink{algebra_catenary_regular_notes_20260928}{Supplementary proof 20}, Two weakenings of Noetherianity in free lifting.
\item[ALGEBRA-RECON-605] Source lines 25690--25729: \hyperlink{algebra_catenary_regular_notes_20260928}{Supplementary proof 20}, The valuation determined by the original maximal-ideal powers.
\end{description}

\subsection{Epimorphisms of rings}
\hypertarget{algebra-source-section-epimorphism}{}
\noindent\href{https://github.com/stacks/stacks-project/blob/a04446e57ec1fbc252a871afcec7752fb2807b14/algebra.tex\#L25862}{Original source section}. \href{algebra.pdf\#nameddest=section-epimorphism}{Accompanying chapter}.

\begin{description}
\item[ALGEBRA-RECON-606] Source lines 25914--25918: \hyperlink{algebra_epimorphism_notes_20260928}{Supplementary proof 21}, Tensor characterization and permanence.
\item[ALGEBRA-RECON-607] Source lines 25920--25941: \hyperlink{algebra_epimorphism_notes_20260928}{Supplementary proof 21}, Local comparison and finite or field targets.
\item[ALGEBRA-RECON-608] Source lines 25987--25998: \hyperlink{algebra_epimorphism_notes_20260928}{Supplementary proof 21}, Local comparison and finite or field targets.
\item[ALGEBRA-RECON-609] Source lines 26022--26058: \hyperlink{algebra_epimorphism_notes_20260928}{Supplementary proof 21}, Finite relations and the original matrix construction.
\item[ALGEBRA-RECON-610] Source lines 26060--26097: \hyperlink{algebra_epimorphism_notes_20260928}{Supplementary proof 21}, Finite relations and the original matrix construction.
\item[ALGEBRA-RECON-611] Source lines 26095--26096: \hyperlink{algebra_epimorphism_notes_20260928}{Supplementary proof 21}, Finite relations and the original matrix construction.
\item[ALGEBRA-RECON-612] Source lines 26099--26123: \hyperlink{algebra_epimorphism_notes_20260928}{Supplementary proof 21}, Typed triples and the cardinality argument.
\item[ALGEBRA-RECON-613] Source lines 26116--26155: \hyperlink{algebra_epimorphism_notes_20260928}{Supplementary proof 21}, Typed triples and the cardinality argument.
\item[ALGEBRA-RECON-614] Source lines 25920--25941: \hyperlink{algebra_epimorphism_notes_20260928}{Supplementary proof 21}, Maximal local tests and pure descent.
\item[ALGEBRA-RECON-615] Source lines 26125--26156: \hyperlink{algebra_epimorphism_notes_20260928}{Supplementary proof 21}, Surjectivity over nilpotent extensions of finite products of fields.
\end{description}

\subsection{Pure ideals}
\hypertarget{algebra-source-section-pure-ideals}{}
\noindent\href{https://github.com/stacks/stacks-project/blob/a04446e57ec1fbc252a871afcec7752fb2807b14/algebra.tex\#L26192}{Original source section}. \href{algebra.pdf\#nameddest=section-pure-ideals}{Accompanying chapter}.

\begin{description}
\item[ALGEBRA-RECON-616] Source lines 26229--26235: \hyperlink{algebra_pure_global_notes_20260928}{Supplementary proof 22}, The eleven pure-ideal conditions.
\item[ALGEBRA-RECON-617] Source lines 26256--26261: \hyperlink{algebra_pure_global_notes_20260928}{Supplementary proof 22}, The eleven pure-ideal conditions.
\item[ALGEBRA-RECON-618] Source lines 26405--26420: \hyperlink{algebra_pure_global_notes_20260928}{Supplementary proof 22}, Finite flat modules and their rank strata.
\item[ALGEBRA-RECON-619] Source lines 26426--26430: \hyperlink{algebra_pure_global_notes_20260928}{Supplementary proof 22}, Finite flat modules and their rank strata.
\item[ALGEBRA-RECON-624] Source lines 26199--26286: \hyperlink{algebra_pure_global_notes_20260928}{Supplementary proof 22}, Pure ideals under base change and pure descent.
\item[ALGEBRA-RECON-625] Source lines 26288--26469: \hyperlink{algebra_pure_global_notes_20260928}{Supplementary proof 22}, Nil-ideal invariance with the actual quotient maps.
\end{description}

\subsection{Rings of finite global dimension}
\hypertarget{algebra-source-section-ring-finite-gl-dim}{}
\noindent\href{https://github.com/stacks/stacks-project/blob/a04446e57ec1fbc252a871afcec7752fb2807b14/algebra.tex\#L26477}{Original source section}. \href{algebra.pdf\#nameddest=section-ring-finite-gl-dim}{Accompanying chapter}.

\begin{description}
\item[ALGEBRA-RECON-620] Source lines 26484--26566: \hyperlink{algebra_pure_global_notes_20260928}{Supplementary proof 22}, Schanuel maps with their original signs.
\item[ALGEBRA-RECON-621] Source lines 26677--26722: \hyperlink{algebra_pure_global_notes_20260928}{Supplementary proof 22}, Syzygies, augmented indices and the reported boundary cases.
\item[ALGEBRA-RECON-622] Source lines 26724--26759: \hyperlink{algebra_pure_global_notes_20260928}{Supplementary proof 22}, Syzygies, augmented indices and the reported boundary cases.
\item[ALGEBRA-RECON-623] Source lines 26677--26722: \hyperlink{algebra_pure_global_notes_20260928}{Supplementary proof 22}, Finite resolution data in place of Noetherianity.
\item[ALGEBRA-RECON-626] Source lines 26880--26903: \hyperlink{algebra_pure_global_notes_20260928}{Supplementary proof 22}, Global dimension along flat epimorphisms.
\item[ALGEBRA-RECON-627] Source lines 26724--26759: \hyperlink{algebra_pure_global_notes_20260928}{Supplementary proof 22}, Retained FAC additions and their degree-zero clarification.
\end{description}

\subsection{Regular rings and global dimension}
\hypertarget{algebra-source-section-regular-finite-gl-dim}{}
\noindent\href{https://github.com/stacks/stacks-project/blob/a04446e57ec1fbc252a871afcec7752fb2807b14/algebra.tex\#L26914}{Original source section}. \href{algebra.pdf\#nameddest=section-regular-finite-gl-dim}{Accompanying chapter}.

\begin{description}
\item[ALGEBRA-RECON-628] Source lines 26921--26936: \hyperlink{algebra_regular_ab_notes_20260928}{Supplementary proof 23}, Finite depth and the maximal-localization criterion.
\item[ALGEBRA-RECON-629] Source lines 26946--26950: \hyperlink{algebra_regular_ab_notes_20260928}{Supplementary proof 23}, Finite depth and the maximal-localization criterion.
\item[ALGEBRA-RECON-630] Source lines 26938--26963: \hyperlink{algebra_regular_ab_notes_20260928}{Supplementary proof 23}, Finite depth and the maximal-localization criterion.
\item[ALGEBRA-RECON-631] Source lines 27079--27093: \hyperlink{algebra_regular_ab_notes_20260928}{Supplementary proof 23}, The dimension bound, exact global dimension and zero rings.
\item[ALGEBRA-RECON-632] Source lines 27071--27093: \hyperlink{algebra_regular_ab_notes_20260928}{Supplementary proof 23}, The dimension bound, exact global dimension and zero rings.
\item[ALGEBRA-RECON-633] Source lines 27163--27186: \hyperlink{algebra_regular_ab_notes_20260928}{Supplementary proof 23}, The dimension bound, exact global dimension and zero rings.
\item[ALGEBRA-RECON-634] Source lines 27188--27213: \hyperlink{algebra_regular_ab_notes_20260928}{Supplementary proof 23}, Flat descent with the actual low-degree kernels.
\item[ALGEBRA-RECON-635] Source lines 27202--27208: \hyperlink{algebra_regular_ab_notes_20260928}{Supplementary proof 23}, Flat descent with the actual low-degree kernels.
\item[ALGEBRA-RECON-638] Source lines 26965--27069: \hyperlink{algebra_regular_ab_notes_20260928}{Supplementary proof 23}, The Koszul map, its signs and a characteristic-free left inverse.
\item[ALGEBRA-RECON-639] Source lines 27188--27213: \hyperlink{algebra_regular_ab_notes_20260928}{Supplementary proof 23}, Flat descent with the actual low-degree kernels.
\item[ALGEBRA-RECON-641] Source lines 27188--27213: \hyperlink{algebra_regular_ab_notes_20260928}{Supplementary proof 23}, Bounded literature reading and a separate source overclaim.
\end{description}

\subsection{Auslander-Buchsbaum}
\hypertarget{algebra-source-section-Auslander-Buchsbaum}{}
\noindent\href{https://github.com/stacks/stacks-project/blob/a04446e57ec1fbc252a871afcec7752fb2807b14/algebra.tex\#L27222}{Original source section}. \href{algebra.pdf\#nameddest=section-Auslander-Buchsbaum}{Accompanying chapter}.

\begin{description}
\item[ALGEBRA-RECON-636] Source lines 27238--27265: \hyperlink{algebra_regular_ab_notes_20260928}{Supplementary proof 23}, Auslander–Buchsbaum with every short case.
\item[ALGEBRA-RECON-637] Source lines 27268--27275: \hyperlink{algebra_regular_ab_notes_20260928}{Supplementary proof 23}, Auslander–Buchsbaum with every short case.
\item[ALGEBRA-RECON-640] Source lines 27228--27295: \hyperlink{algebra_regular_ab_notes_20260928}{Supplementary proof 23}, The retained FAC quotient result and all Betti numbers.
\end{description}

\subsection{Homomorphisms and dimension}
\hypertarget{algebra-source-section-homomorphism-dimension}{}
\noindent\href{https://github.com/stacks/stacks-project/blob/a04446e57ec1fbc252a871afcec7752fb2807b14/algebra.tex\#L27305}{Original source section}. \href{algebra.pdf\#nameddest=section-homomorphism-dimension}{Accompanying chapter}.

\begin{description}
\item[ALGEBRA-RECON-642] Source lines 27334--27345: \hyperlink{algebra_dimension_maps_notes_20260928}{Supplementary proof 24}, Chains, integral maps and the original fibre rings.
\item[ALGEBRA-RECON-643] Source lines 27436--27456: \hyperlink{algebra_dimension_maps_notes_20260928}{Supplementary proof 24}, The parameter bound and the zero-length going-down case.
\item[ALGEBRA-RECON-644] Source lines 27488--27503: \hyperlink{algebra_dimension_maps_notes_20260928}{Supplementary proof 24}, Cohen–Macaulay ascent with an arbitrary Cohen–Macaulay fibre.
\item[ALGEBRA-RECON-649] Source lines 27458--27480: \hyperlink{algebra_dimension_maps_notes_20260928}{Supplementary proof 24}, Regular ascent under going down and the resulting flatness.
\item[ALGEBRA-RECON-650] Source lines 27488--27503: \hyperlink{algebra_dimension_maps_notes_20260928}{Supplementary proof 24}, Cohen–Macaulay ascent with an arbitrary Cohen–Macaulay fibre.
\item[ALGEBRA-RECON-653] Source lines 27390--27456: \hyperlink{algebra_dimension_maps_notes_20260928}{Supplementary proof 24}, The parameter bound and the zero-length going-down case.
\end{description}

\subsection{The dimension formula}
\hypertarget{algebra-source-section-dimension-formula}{}
\noindent\href{https://github.com/stacks/stacks-project/blob/a04446e57ec1fbc252a871afcec7752fb2807b14/algebra.tex\#L27512}{Original source section}. \href{algebra.pdf\#nameddest=section-dimension-formula}{Accompanying chapter}.

\begin{description}
\item[ALGEBRA-RECON-645] Source lines 27570--27575: \hyperlink{algebra_dimension_maps_notes_20260928}{Supplementary proof 24}, The dimension formula with the original tower and height-one prime.
\item[ALGEBRA-RECON-646] Source lines 27585--27587: \hyperlink{algebra_dimension_maps_notes_20260928}{Supplementary proof 24}, The dimension formula with the original tower and height-one prime.
\item[ALGEBRA-RECON-647] Source lines 27588--27595: \hyperlink{algebra_dimension_maps_notes_20260928}{Supplementary proof 24}, The dimension formula with the original tower and height-one prime.
\item[ALGEBRA-RECON-648] Source lines 27624--27637: \hyperlink{algebra_dimension_maps_notes_20260928}{Supplementary proof 24}, The missing topological hypothesis and its exact repair.
\item[ALGEBRA-RECON-651] Source lines 27519--27638: \hyperlink{algebra_dimension_maps_notes_20260928}{Supplementary proof 24}, Essential finite type, local catenarity and the full codimension-one fibre.
\item[ALGEBRA-RECON-652] Source lines 27519--27596: \hyperlink{algebra_dimension_maps_notes_20260928}{Supplementary proof 24}, The original literature counterexample and propagation boundaries.
\end{description}

\subsection{Dimension of finite type algebras over fields}
\hypertarget{algebra-source-section-dimension-finite-type-algebras}{}
\noindent\href{https://github.com/stacks/stacks-project/blob/a04446e57ec1fbc252a871afcec7752fb2807b14/algebra.tex\#L27648}{Original source section}. \href{algebra.pdf\#nameddest=section-dimension-finite-type-algebras}{Accompanying chapter}.

\begin{description}
\item[ALGEBRA-RECON-654] Source lines 27669--27673: \hyperlink{algebra_affine_dimension_notes_20260928}{Supplementary proof 25}, The original triangular polynomials and the corrected induction.
\item[ALGEBRA-RECON-655] Source lines 27673--27675: \hyperlink{algebra_affine_dimension_notes_20260928}{Supplementary proof 25}, The original triangular polynomials and the corrected induction.
\item[ALGEBRA-RECON-656] Source lines 27688--27704: \hyperlink{algebra_affine_dimension_notes_20260928}{Supplementary proof 25}, The original triangular polynomials and the corrected induction.
\item[ALGEBRA-RECON-657] Source lines 27869--27873: \hyperlink{algebra_affine_dimension_notes_20260928}{Supplementary proof 25}, The decomposition with its weaker exact hypothesis.
\item[ALGEBRA-RECON-658] Source lines 27885--27891: \hyperlink{algebra_affine_dimension_notes_20260928}{Supplementary proof 25}, The decomposition with its weaker exact hypothesis.
\item[ALGEBRA-RECON-659] Source lines 27668--27720: \hyperlink{algebra_affine_dimension_notes_20260928}{Supplementary proof 25}, Explicit freeness over the original parameter algebra.
\item[ALGEBRA-RECON-660] Source lines 27786--27903: \hyperlink{algebra_affine_dimension_notes_20260928}{Supplementary proof 25}, The decomposition with its weaker exact hypothesis.
\item[ALGEBRA-RECON-661] Source lines 27869--27903: \hyperlink{algebra_affine_dimension_notes_20260928}{Supplementary proof 25}, Product maps, gaps in the dimensions and empty spectra.
\end{description}

\subsection{Noether normalization}
\hypertarget{algebra-source-section-Noether-normalization}{}
\noindent\href{https://github.com/stacks/stacks-project/blob/a04446e57ec1fbc252a871afcec7752fb2807b14/algebra.tex\#L27920}{Original source section}. \href{algebra.pdf\#nameddest=section-Noether-normalization}{Accompanying chapter}.

\begin{description}
\item[ALGEBRA-RECON-662] Source lines 27927--27960: \hyperlink{algebra_noether_dimension_notes_20260928}{Supplementary proof 26}, Finite supports and explicit positive weights.
\item[ALGEBRA-RECON-663] Source lines 27975--27979: \hyperlink{algebra_noether_dimension_notes_20260928}{Supplementary proof 26}, Finite supports and explicit positive weights.
\item[ALGEBRA-RECON-664] Source lines 28013--28021: \hyperlink{algebra_noether_dimension_notes_20260928}{Supplementary proof 26}, The original substitutions, full coefficients and one relation.
\item[ALGEBRA-RECON-665] Source lines 28024--28054: \hyperlink{algebra_noether_dimension_notes_20260928}{Supplementary proof 26}, The full induction and a neighborhood of the original point.
\item[ALGEBRA-RECON-666] Source lines 28056--28287: \hyperlink{algebra_noether_dimension_notes_20260928}{Supplementary proof 26}, Dimension, transcendence degree and the zero-length chains.
\item[ALGEBRA-RECON-667] Source lines 28086--28123: \hyperlink{algebra_noether_dimension_notes_20260928}{Supplementary proof 26}, The refined map, every unit and a finite free strengthening.
\item[ALGEBRA-RECON-668] Source lines 28137--28142: \hyperlink{algebra_noether_dimension_notes_20260928}{Supplementary proof 26}, Generic comparison over a domain without losing torsion.
\item[ALGEBRA-RECON-672] Source lines 27927--27994: \hyperlink{algebra_noether_dimension_notes_20260928}{Supplementary proof 26}, Finite supports and explicit positive weights.
\item[ALGEBRA-RECON-673] Source lines 28073--28123: \hyperlink{algebra_noether_dimension_notes_20260928}{Supplementary proof 26}, The refined map, every unit and a finite free strengthening.
\item[ALGEBRA-RECON-676] Source lines 28125--28161: \hyperlink{algebra_noether_dimension_notes_20260928}{Supplementary proof 26}, Generic comparison over a domain without losing torsion.
\end{description}

\subsection{Dimension of finite type algebras over fields, reprise}
\hypertarget{algebra-source-section-dimension-finite-type-algebras-reprise}{}
\noindent\href{https://github.com/stacks/stacks-project/blob/a04446e57ec1fbc252a871afcec7752fb2807b14/algebra.tex\#L28167}{Original source section}. \href{algebra.pdf\#nameddest=section-dimension-finite-type-algebras-reprise}{Accompanying chapter}.

\begin{description}
\item[ALGEBRA-RECON-666] Source lines 28056--28287: \hyperlink{algebra_noether_dimension_notes_20260928}{Supplementary proof 26}, Dimension, transcendence degree and the zero-length chains.
\item[ALGEBRA-RECON-669] Source lines 28343--28346: \hyperlink{algebra_noether_dimension_notes_20260928}{Supplementary proof 26}, Field extension and the common original fibre ring.
\item[ALGEBRA-RECON-670] Source lines 28373--28387: \hyperlink{algebra_noether_dimension_notes_20260928}{Supplementary proof 26}, Field extension and the common original fibre ring.
\item[ALGEBRA-RECON-671] Source lines 28221--28304: \hyperlink{algebra_noether_dimension_notes_20260928}{Supplementary proof 26}, Dimension, transcendence degree and the zero-length chains.
\item[ALGEBRA-RECON-674] Source lines 28358--28387: \hyperlink{algebra_noether_dimension_notes_20260928}{Supplementary proof 26}, Exact maximum fibre dimension and both height-jump extremes.
\end{description}

\subsection{Dimension of graded algebras over a field}
\hypertarget{algebra-source-section-dimension-graded}{}
\noindent\href{https://github.com/stacks/stacks-project/blob/a04446e57ec1fbc252a871afcec7752fb2807b14/algebra.tex\#L28395}{Original source section}. \href{algebra.pdf\#nameddest=section-dimension-graded}{Accompanying chapter}.

\begin{description}
\item[ALGEBRA-RECON-675] Source lines 28402--28434: \hyperlink{algebra_noether_dimension_notes_20260928}{Supplementary proof 26}, The graded vertex, its Hilbert function and its full completion.
\end{description}

\subsection{Generic flatness}
\hypertarget{algebra-source-section-generic-flatness}{}
\noindent\href{https://github.com/stacks/stacks-project/blob/a04446e57ec1fbc252a871afcec7752fb2807b14/algebra.tex\#L28448}{Original source section}. \href{algebra.pdf\#nameddest=section-generic-flatness}{Accompanying chapter}.

\begin{description}
\item[ALGEBRA-RECON-677] Source lines 28518--28523: \hyperlink{algebra_generic_freeness_notes_20260928}{Supplementary proof 27}, The Noetherian argument and the empty generic fibre.
\item[ALGEBRA-RECON-678] Source lines 28543--28561: \hyperlink{algebra_generic_freeness_notes_20260928}{Supplementary proof 27}, Finite coefficient models and the original relation kernel.
\item[ALGEBRA-RECON-679] Source lines 28587--28615: \hyperlink{algebra_generic_freeness_notes_20260928}{Supplementary proof 27}, Finite coefficient models and the original relation kernel.
\item[ALGEBRA-RECON-680] Source lines 28762--28773: \hyperlink{algebra_generic_freeness_notes_20260928}{Supplementary proof 27}, Reduced bases, coefficients and finite presentation across the substitution.
\item[ALGEBRA-RECON-681] Source lines 28858--28870: \hyperlink{algebra_generic_freeness_notes_20260928}{Supplementary proof 27}, Reduced bases, coefficients and finite presentation across the substitution.
\item[ALGEBRA-RECON-682] Source lines 28473--28495: \hyperlink{algebra_generic_freeness_notes_20260928}{Supplementary proof 27}, The Noetherian argument and the empty generic fibre.
\item[ALGEBRA-RECON-683] Source lines 28618--28867: \hyperlink{algebra_generic_freeness_notes_20260928}{Supplementary proof 27}, Reduced bases, coefficients and finite presentation across the substitution.
\item[ALGEBRA-RECON-684] Source lines 28526--28625: \hyperlink{algebra_generic_freeness_notes_20260928}{Supplementary proof 27}, Finite complexes and arbitrary base change after one localization.
\item[ALGEBRA-RECON-685] Source lines 28627--28870: \hyperlink{algebra_generic_freeness_notes_20260928}{Supplementary proof 27}, A dense good locus with no dense principal open.
\end{description}

\subsection{Around Krull-Akizuki}
\hypertarget{algebra-source-section-krull-akizuki}{}
\noindent\href{https://github.com/stacks/stacks-project/blob/a04446e57ec1fbc252a871afcec7752fb2807b14/algebra.tex\#L28877}{Original source section}. \href{algebra.pdf\#nameddest=section-krull-akizuki}{Accompanying chapter}.

\begin{description}
\item[ALGEBRA-RECON-686] Source lines 28903--28919: \hyperlink{algebra_krull_akizuki_notes_20260928}{Supplementary proof 28}, The original domination chart.
\item[ALGEBRA-RECON-687] Source lines 28904--28919: \hyperlink{algebra_krull_akizuki_notes_20260928}{Supplementary proof 28}, The original domination chart.
\item[ALGEBRA-RECON-688] Source lines 28928--29006: \hyperlink{algebra_krull_akizuki_notes_20260928}{Supplementary proof 28}, The four local cases and an explicit annihilator exponent.
\item[ALGEBRA-RECON-689] Source lines 28987--28994: \hyperlink{algebra_krull_akizuki_notes_20260928}{Supplementary proof 28}, The four local cases and an explicit annihilator exponent.
\item[ALGEBRA-RECON-690] Source lines 29072--29080: \hyperlink{algebra_krull_akizuki_notes_20260928}{Supplementary proof 28}, Globalizing the specified local modification.
\item[ALGEBRA-RECON-691] Source lines 29095--29100: \hyperlink{algebra_krull_akizuki_notes_20260928}{Supplementary proof 28}, Termination, a numerical bound and the completion obstruction.
\item[ALGEBRA-RECON-692] Source lines 29114--29116: \hyperlink{algebra_krull_akizuki_notes_20260928}{Supplementary proof 28}, Uniformizers and the original discrete valuation.
\item[ALGEBRA-RECON-693] Source lines 29129--29144: \hyperlink{algebra_krull_akizuki_notes_20260928}{Supplementary proof 28}, Uniformizers and the original discrete valuation.
\item[ALGEBRA-RECON-694] Source lines 29192--29218: \hyperlink{algebra_krull_akizuki_notes_20260928}{Supplementary proof 28}, Length inequalities with the original ambient vector space retained.
\item[ALGEBRA-RECON-695] Source lines 29221--29234: \hyperlink{algebra_krull_akizuki_notes_20260928}{Supplementary proof 28}, Length inequalities with the original ambient vector space retained.
\item[ALGEBRA-RECON-696] Source lines 29253--29264: \hyperlink{algebra_krull_akizuki_notes_20260928}{Supplementary proof 28}, Residue fields, ramification and Krull–Akizuki.
\item[ALGEBRA-RECON-697] Source lines 29311--29338: \hyperlink{algebra_krull_akizuki_notes_20260928}{Supplementary proof 28}, A DVR in the original finitely generated extension.
\item[ALGEBRA-RECON-698] Source lines 28960--28968: \hyperlink{algebra_krull_akizuki_notes_20260928}{Supplementary proof 28}, The four local cases and an explicit annihilator exponent.
\item[ALGEBRA-RECON-699] Source lines 29072--29100: \hyperlink{algebra_krull_akizuki_notes_20260928}{Supplementary proof 28}, Termination, a numerical bound and the completion obstruction.
\item[ALGEBRA-RECON-700] Source lines 29170--29309: \hyperlink{algebra_krull_akizuki_notes_20260928}{Supplementary proof 28}, Residue fields, ramification and Krull–Akizuki.
\item[ALGEBRA-RECON-701] Source lines 29031--29068: \hyperlink{algebra_krull_akizuki_notes_20260928}{Supplementary proof 28}, The two original examples and their exact completions.
\end{description}

\subsection{Factorization}
\hypertarget{algebra-source-section-factoring}{}
\noindent\href{https://github.com/stacks/stacks-project/blob/a04446e57ec1fbc252a871afcec7752fb2807b14/algebra.tex\#L29347}{Original source section}. \href{algebra.pdf\#nameddest=section-factoring}{Accompanying chapter}.

\begin{description}
\item[ALGEBRA-RECON-702] Source lines 29354--29365: \hyperlink{algebra_factorization_notes_20260928}{Supplementary proof 29}, Prime elements, associates and the rejected factorization reports.
\item[ALGEBRA-RECON-703] Source lines 29374--29379: \hyperlink{algebra_factorization_notes_20260928}{Supplementary proof 29}, Prime elements, associates and the rejected factorization reports.
\item[ALGEBRA-RECON-704] Source lines 29374--29379: \hyperlink{algebra_factorization_notes_20260928}{Supplementary proof 29}, Prime elements, associates and the rejected factorization reports.
\item[ALGEBRA-RECON-705] Source lines 29395--29409: \hyperlink{algebra_factorization_notes_20260928}{Supplementary proof 29}, Prime elements, associates and the rejected factorization reports.
\item[ALGEBRA-RECON-706] Source lines 29434--29443: \hyperlink{algebra_factorization_notes_20260928}{Supplementary proof 29}, The UFD criterion with every unit factor retained.
\item[ALGEBRA-RECON-707] Source lines 29434--29443: \hyperlink{algebra_factorization_notes_20260928}{Supplementary proof 29}, The UFD criterion with every unit factor retained.
\item[ALGEBRA-RECON-708] Source lines 29434--29443: \hyperlink{algebra_factorization_notes_20260928}{Supplementary proof 29}, The UFD criterion with every unit factor retained.
\item[ALGEBRA-RECON-709] Source lines 29448--29455: \hyperlink{algebra_factorization_notes_20260928}{Supplementary proof 29}, The UFD criterion with every unit factor retained.
\item[ALGEBRA-RECON-710] Source lines 29484--29540: \hyperlink{algebra_factorization_notes_20260928}{Supplementary proof 29}, Nagata cancellation with repeated factors.
\item[ALGEBRA-RECON-711] Source lines 29542--29554: \hyperlink{algebra_factorization_notes_20260928}{Supplementary proof 29}, ACC, polynomial rings and the full normality calculation.
\item[ALGEBRA-RECON-712] Source lines 29556--29575: \hyperlink{algebra_factorization_notes_20260928}{Supplementary proof 29}, ACC, polynomial rings and the full normality calculation.
\item[ALGEBRA-RECON-713] Source lines 29577--29598: \hyperlink{algebra_factorization_notes_20260928}{Supplementary proof 29}, ACC, polynomial rings and the full normality calculation.
\item[ALGEBRA-RECON-714] Source lines 29600--29619: \hyperlink{algebra_factorization_notes_20260928}{Supplementary proof 29}, ACC, polynomial rings and the full normality calculation.
\item[ALGEBRA-RECON-715] Source lines 29637--29644: \hyperlink{algebra_factorization_notes_20260928}{Supplementary proof 29}, Dedekind domains and the full uniqueness proof.
\item[ALGEBRA-RECON-716] Source lines 29647--29658: \hyperlink{algebra_factorization_notes_20260928}{Supplementary proof 29}, Dedekind domains and the full uniqueness proof.
\item[ALGEBRA-RECON-717] Source lines 29685--29694: \hyperlink{algebra_factorization_notes_20260928}{Supplementary proof 29}, Dedekind domains and the full uniqueness proof.
\item[ALGEBRA-RECON-718] Source lines 29696--29726: \hyperlink{algebra_factorization_notes_20260928}{Supplementary proof 29}, Dedekind domains and the full uniqueness proof.
\item[ALGEBRA-RECON-719] Source lines 29713--29717: \hyperlink{algebra_factorization_notes_20260928}{Supplementary proof 29}, Dedekind domains and the full uniqueness proof.
\item[ALGEBRA-RECON-720] Source lines 29719--29729: \hyperlink{algebra_factorization_notes_20260928}{Supplementary proof 29}, Dedekind domains and the full uniqueness proof.
\item[ALGEBRA-RECON-721] Source lines 29734--29747: \hyperlink{algebra_factorization_notes_20260928}{Supplementary proof 29}, Dedekind domains and the full uniqueness proof.
\item[ALGEBRA-RECON-722] Source lines 29738--29746: \hyperlink{algebra_factorization_notes_20260928}{Supplementary proof 29}, Dedekind domains and the full uniqueness proof.
\item[ALGEBRA-RECON-723] Source lines 29745--29747: \hyperlink{algebra_factorization_notes_20260928}{Supplementary proof 29}, Dedekind domains and the full uniqueness proof.
\item[ALGEBRA-RECON-724] Source lines 29753--29771: \hyperlink{algebra_factorization_notes_20260928}{Supplementary proof 29}, Dedekind domains and the full uniqueness proof.
\item[ALGEBRA-RECON-725] Source lines 29484--29540: \hyperlink{algebra_factorization_notes_20260928}{Supplementary proof 29}, Exact localization exponents and the atomicity boundary.
\item[ALGEBRA-RECON-726] Source lines 29660--29683: \hyperlink{algebra_factorization_notes_20260928}{Supplementary proof 29}, Explicit dual maps for a product of ideals.
\item[ALGEBRA-RECON-727] Source lines 29685--29771: \hyperlink{algebra_factorization_notes_20260928}{Supplementary proof 29}, Fractional ideal coordinates and the semilocal consequence.
\end{description}

\subsection{Orders of vanishing}
\hypertarget{algebra-source-section-orders-of-vanishing}{}
\noindent\href{https://github.com/stacks/stacks-project/blob/a04446e57ec1fbc252a871afcec7752fb2807b14/algebra.tex\#L29779}{Original source section}. \href{algebra.pdf\#nameddest=section-orders-of-vanishing}{Accompanying chapter}.

\begin{description}
\item[ALGEBRA-RECON-728] Source lines 29794--29800: \hyperlink{algebra_orders_notes_20260928}{Supplementary proof 30}, Nonzerodivisor lengths and the citation hypothesis.
\item[ALGEBRA-RECON-729] Source lines 29868--29942: \hyperlink{algebra_orders_notes_20260928}{Supplementary proof 30}, The original lattices and their signed distance.
\item[ALGEBRA-RECON-730] Source lines 29885--29894: \hyperlink{algebra_orders_notes_20260928}{Supplementary proof 30}, The original lattices and their signed distance.
\item[ALGEBRA-RECON-731] Source lines 29923--29931: \hyperlink{algebra_orders_notes_20260928}{Supplementary proof 30}, The original lattices and their signed distance.
\item[ALGEBRA-RECON-732] Source lines 29995--30000: \hyperlink{algebra_orders_notes_20260928}{Supplementary proof 30}, The original determinant identity and every generator.
\item[ALGEBRA-RECON-733] Source lines 30002--30034: \hyperlink{algebra_orders_notes_20260928}{Supplementary proof 30}, The original determinant identity and every generator.
\item[ALGEBRA-RECON-734] Source lines 29782--29800: \hyperlink{algebra_orders_notes_20260928}{Supplementary proof 30}, Nonzerodivisor lengths and the citation hypothesis.
\item[ALGEBRA-RECON-735] Source lines 29782--30035: \hyperlink{algebra_orders_notes_20260928}{Supplementary proof 30}, Finite modules, torsion terms and exact-sequence indices.
\item[ALGEBRA-RECON-736] Source lines 30037--30078: \hyperlink{algebra_orders_notes_20260928}{Supplementary proof 30}, The original finite-extension norm and its global version.
\item[ALGEBRA-RECON-737] Source lines 29802--30078: \hyperlink{algebra_orders_notes_20260928}{Supplementary proof 30}, Two exact boundaries: valuations and finite extensions.
\end{description}

\subsection{Quasi-finite maps}
\hypertarget{algebra-source-section-quasi-finite}{}
\noindent\href{https://github.com/stacks/stacks-project/blob/a04446e57ec1fbc252a871afcec7752fb2807b14/algebra.tex\#L30093}{Original source section}. \href{algebra.pdf\#nameddest=section-quasi-finite}{Accompanying chapter}.

\begin{description}
\item[ALGEBRA-RECON-738] Source lines 30190--30193: \hyperlink{algebra_quasifinite_notes_20260928}{Supplementary proof 31}, Isolated points and the exact fibre maps.
\item[ALGEBRA-RECON-739] Source lines 30200--30206: \hyperlink{algebra_quasifinite_notes_20260928}{Supplementary proof 31}, Isolated points and the exact fibre maps.
\item[ALGEBRA-RECON-740] Source lines 30205--30218: \hyperlink{algebra_quasifinite_notes_20260928}{Supplementary proof 31}, Isolated points and the exact fibre maps.
\item[ALGEBRA-RECON-741] Source lines 30292--30306: \hyperlink{algebra_quasifinite_notes_20260928}{Supplementary proof 31}, The original four-ring, composition and base-change proofs.
\item[ALGEBRA-RECON-742] Source lines 30329--30338: \hyperlink{algebra_quasifinite_notes_20260928}{Supplementary proof 31}, The original four-ring, composition and base-change proofs.
\item[ALGEBRA-RECON-743] Source lines 30346--30358: \hyperlink{algebra_quasifinite_notes_20260928}{Supplementary proof 31}, The original four-ring, composition and base-change proofs.
\item[ALGEBRA-RECON-744] Source lines 30376--30389: \hyperlink{algebra_quasifinite_notes_20260928}{Supplementary proof 31}, The original four-ring, composition and base-change proofs.
\item[ALGEBRA-RECON-745] Source lines 30428--30445: \hyperlink{algebra_quasifinite_notes_20260928}{Supplementary proof 31}, Generic finiteness with actual coefficient lifts and nilpotence powers.
\item[ALGEBRA-RECON-746] Source lines 30105--30390: \hyperlink{algebra_quasifinite_notes_20260928}{Supplementary proof 31}, The canonical finite-dimensional factor and field extension.
\item[ALGEBRA-RECON-747] Source lines 30292--30339: \hyperlink{algebra_quasifinite_notes_20260928}{Supplementary proof 31}, A finite four-ring comparison and its full multiplicity bound.
\item[ALGEBRA-RECON-748] Source lines 30417--30446: \hyperlink{algebra_quasifinite_notes_20260928}{Supplementary proof 31}, The finite locus and the boundary of the minimal-prime assumption.
\end{description}

\subsection{Zariski's Main Theorem}
\hypertarget{algebra-source-section-Zariski}{}
\noindent\href{https://github.com/stacks/stacks-project/blob/a04446e57ec1fbc252a871afcec7752fb2807b14/algebra.tex\#L30454}{Original source section}. \href{algebra.pdf\#nameddest=section-Zariski}{Accompanying chapter}.

\begin{description}
\item[ALGEBRA-RECON-749] Source lines 30483--30495: \hyperlink{algebra_zariski_main_notes_20260928}{Supplementary proof 32}, Integral equations with the original localization kernels.
\item[ALGEBRA-RECON-750] Source lines 30488--30495: \hyperlink{algebra_zariski_main_notes_20260928}{Supplementary proof 32}, Integral equations with the original localization kernels.
\item[ALGEBRA-RECON-751] Source lines 30570--30580: \hyperlink{algebra_zariski_main_notes_20260928}{Supplementary proof 32}, The conductor and the coefficient argument.
\item[ALGEBRA-RECON-752] Source lines 30582--30589: \hyperlink{algebra_zariski_main_notes_20260928}{Supplementary proof 32}, The conductor and the coefficient argument.
\item[ALGEBRA-RECON-753] Source lines 30713--30732: \hyperlink{algebra_zariski_main_notes_20260928}{Supplementary proof 32}, The original nowhere quasi-finite argument.
\item[ALGEBRA-RECON-754] Source lines 30745--30756: \hyperlink{algebra_zariski_main_notes_20260928}{Supplementary proof 32}, Monogenic localization and the explicit Horner witness.
\item[ALGEBRA-RECON-755] Source lines 30759--30767: \hyperlink{algebra_zariski_main_notes_20260928}{Supplementary proof 32}, Monogenic localization and the explicit Horner witness.
\item[ALGEBRA-RECON-756] Source lines 30766--30775: \hyperlink{algebra_zariski_main_notes_20260928}{Supplementary proof 32}, Monogenic localization and the explicit Horner witness.
\item[ALGEBRA-RECON-757] Source lines 30778--30786: \hyperlink{algebra_zariski_main_notes_20260928}{Supplementary proof 32}, The main induction with both final denominators.
\item[ALGEBRA-RECON-758] Source lines 30788--30863: \hyperlink{algebra_zariski_main_notes_20260928}{Supplementary proof 32}, The main induction with both final denominators.
\item[ALGEBRA-RECON-759] Source lines 30788--30793: \hyperlink{algebra_zariski_main_notes_20260928}{Supplementary proof 32}, The main induction with both final denominators.
\item[ALGEBRA-RECON-760] Source lines 30789--30793: \hyperlink{algebra_zariski_main_notes_20260928}{Supplementary proof 32}, The main induction with both final denominators.
\item[ALGEBRA-RECON-761] Source lines 30853--30863: \hyperlink{algebra_zariski_main_notes_20260928}{Supplementary proof 32}, The main induction with both final denominators.
\item[ALGEBRA-RECON-762] Source lines 30629--30732: \hyperlink{algebra_zariski_main_notes_20260928}{Supplementary proof 32}, Strong transcendence and the reducedness boundary.
\item[ALGEBRA-RECON-763] Source lines 30735--30776: \hyperlink{algebra_zariski_main_notes_20260928}{Supplementary proof 32}, Monogenic localization and the explicit Horner witness.
\item[ALGEBRA-RECON-764] Source lines 30779--30949: \hyperlink{algebra_zariski_main_notes_20260928}{Supplementary proof 32}, Finite envelopes and the exact open immersion.
\end{description}

\subsection{Applications of Zariski's Main Theorem}
\hypertarget{algebra-source-section-apply-main-theorem}{}
\noindent\href{https://github.com/stacks/stacks-project/blob/a04446e57ec1fbc252a871afcec7752fb2807b14/algebra.tex\#L30954}{Original source section}. \href{algebra.pdf\#nameddest=section-apply-main-theorem}{Accompanying chapter}.

\begin{description}
\item[ALGEBRA-RECON-765] Source lines 30965--31015: \hyperlink{algebra_zariski_applications_notes_20260928}{Supplementary proof 33}, The field exception and the original order formula.
\item[ALGEBRA-RECON-766] Source lines 31074--31084: \hyperlink{algebra_zariski_applications_notes_20260928}{Supplementary proof 33}, The completion factor and the citation scope.
\item[ALGEBRA-RECON-767] Source lines 31094--31102: \hyperlink{algebra_zariski_applications_notes_20260928}{Supplementary proof 33}, The completion factor and the citation scope.
\item[ALGEBRA-RECON-768] Source lines 31065--31069: \hyperlink{algebra_zariski_applications_notes_20260928}{Supplementary proof 33}, The completion factor and the citation scope.
\item[ALGEBRA-RECON-769] Source lines 30964--31016: \hyperlink{algebra_zariski_applications_notes_20260928}{Supplementary proof 33}, The exact deficit, including the field case.
\item[ALGEBRA-RECON-770] Source lines 31052--31103: \hyperlink{algebra_zariski_applications_notes_20260928}{Supplementary proof 33}, All isolated completed factors and their maximality.
\end{description}

\subsection{Dimension of fibres}
\hypertarget{algebra-source-section-dimension-fibres}{}
\noindent\href{https://github.com/stacks/stacks-project/blob/a04446e57ec1fbc252a871afcec7752fb2807b14/algebra.tex\#L31140}{Original source section}. \href{algebra.pdf\#nameddest=section-dimension-fibres}{Accompanying chapter}.

\begin{description}
\item[ALGEBRA-RECON-771] Source lines 31171--31208: \hyperlink{algebra_fibre_dimension_notes_20260928}{Supplementary proof 34}, The original polynomial chart and every localization.
\item[ALGEBRA-RECON-772] Source lines 31182--31187: \hyperlink{algebra_fibre_dimension_notes_20260928}{Supplementary proof 34}, The original polynomial chart and every localization.
\item[ALGEBRA-RECON-773] Source lines 31198--31204: \hyperlink{algebra_fibre_dimension_notes_20260928}{Supplementary proof 34}, The original polynomial chart and every localization.
\item[ALGEBRA-RECON-774] Source lines 31220--31227: \hyperlink{algebra_fibre_dimension_notes_20260928}{Supplementary proof 34}, The refined chart and its contracted prime.
\item[ALGEBRA-RECON-775] Source lines 31280--31287: \hyperlink{algebra_fibre_dimension_notes_20260928}{Supplementary proof 34}, Dimension bounds and exact base change.
\item[ALGEBRA-RECON-776] Source lines 31295--31304: \hyperlink{algebra_fibre_dimension_notes_20260928}{Supplementary proof 34}, Dimension bounds and exact base change.
\item[ALGEBRA-RECON-777] Source lines 31337--31338: \hyperlink{algebra_fibre_dimension_notes_20260928}{Supplementary proof 34}, Dimension bounds and exact base change.
\item[ALGEBRA-RECON-778] Source lines 31358--31376: \hyperlink{algebra_fibre_dimension_notes_20260928}{Supplementary proof 34}, Dimension bounds and exact base change.
\item[ALGEBRA-RECON-779] Source lines 31379--31388: \hyperlink{algebra_fibre_dimension_notes_20260928}{Supplementary proof 34}, The valuation-ring proof with the original coefficient factor.
\item[ALGEBRA-RECON-780] Source lines 31381--31397: \hyperlink{algebra_fibre_dimension_notes_20260928}{Supplementary proof 34}, The valuation-ring proof with the original coefficient factor.
\item[ALGEBRA-RECON-781] Source lines 31170--31377: \hyperlink{algebra_fibre_dimension_notes_20260928}{Supplementary proof 34}, Minimal chart dimension and finite charts.
\item[ALGEBRA-RECON-782] Source lines 31308--31377: \hyperlink{algebra_fibre_dimension_notes_20260928}{Supplementary proof 34}, Universal dimension ideals and the finite-presentation boundary.
\item[ALGEBRA-RECON-783] Source lines 31380--31409: \hyperlink{algebra_fibre_dimension_notes_20260928}{Supplementary proof 34}, Every nonempty valuation fibre and the generically finite consequence.
\end{description}

\subsection{Algebras and modules of finite presentation}
\hypertarget{algebra-source-section-finite-presentation}{}
\noindent\href{https://github.com/stacks/stacks-project/blob/a04446e57ec1fbc252a871afcec7752fb2807b14/algebra.tex\#L31434}{Original source section}. \href{algebra.pdf\#nameddest=section-finite-presentation}{Accompanying chapter}.

\begin{description}
\item[ALGEBRA-RECON-784] Source lines 31454--31460: \hyperlink{algebra_finite_presentation_notes_20260928}{Supplementary proof 35}, The original descent families.
\item[ALGEBRA-RECON-785] Source lines 31478--31486: \hyperlink{algebra_finite_presentation_notes_20260928}{Supplementary proof 35}, The original descent families.
\item[ALGEBRA-RECON-786] Source lines 31575--31578: \hyperlink{algebra_finite_presentation_notes_20260928}{Supplementary proof 35}, The original module lifts and stalk specialization.
\item[ALGEBRA-RECON-787] Source lines 31639--31654: \hyperlink{algebra_finite_presentation_notes_20260928}{Supplementary proof 35}, Local isomorphisms with all original denominators.
\item[ALGEBRA-RECON-788] Source lines 31655--31660: \hyperlink{algebra_finite_presentation_notes_20260928}{Supplementary proof 35}, Local isomorphisms with all original denominators.
\item[ALGEBRA-RECON-789] Source lines 31684--31694: \hyperlink{algebra_finite_presentation_notes_20260928}{Supplementary proof 35}, Nilpotent lifting and the source automorphism.
\item[ALGEBRA-RECON-790] Source lines 31703--31737: \hyperlink{algebra_finite_presentation_notes_20260928}{Supplementary proof 35}, The original isomorphism criteria.
\item[ALGEBRA-RECON-791] Source lines 31749--31752: \hyperlink{algebra_finite_presentation_notes_20260928}{Supplementary proof 35}, The original isomorphism criteria.
\item[ALGEBRA-RECON-792] Source lines 31442--31487: \hyperlink{algebra_finite_presentation_notes_20260928}{Supplementary proof 35}, Algebra descent along the original pure map.
\item[ALGEBRA-RECON-793] Source lines 31677--31695: \hyperlink{algebra_finite_presentation_notes_20260928}{Supplementary proof 35}, Finite inverse for endomorphisms over a nil ideal.
\item[ALGEBRA-RECON-794] Source lines 31697--31753: \hyperlink{algebra_finite_presentation_notes_20260928}{Supplementary proof 35}, The exact Tor obstruction and weaker sufficient hypotheses.
\item[ALGEBRA-RECON-795] Source lines 31580--31660: \hyperlink{algebra_finite_presentation_notes_20260928}{Supplementary proof 35}, Finite type does not replace finite presentation in the local splitting.
\end{description}

\subsection{Colimits and maps of finite presentation}
\hypertarget{algebra-source-section-colimits-flat}{}
\noindent\href{https://github.com/stacks/stacks-project/blob/a04446e57ec1fbc252a871afcec7752fb2807b14/algebra.tex\#L31759}{Original source section}. \href{algebra.pdf\#nameddest=section-colimits-flat}{Accompanying chapter}.

\begin{description}
\item[ALGEBRA-RECON-796] Source lines 31763--31765: \hyperlink{algebra_colimit_categories_notes_20260928}{Supplementary proof 36}, Filtered presentations with their actual structural maps.
\item[ALGEBRA-RECON-797] Source lines 31772--31782: \hyperlink{algebra_colimit_categories_notes_20260928}{Supplementary proof 36}, Filtered presentations with their actual structural maps.
\item[ALGEBRA-RECON-798] Source lines 31770--31811: \hyperlink{algebra_colimit_categories_notes_20260928}{Supplementary proof 36}, Filtered presentations with their actual structural maps.
\item[ALGEBRA-RECON-799] Source lines 31814--31817: \hyperlink{algebra_colimit_categories_notes_20260928}{Supplementary proof 36}, Filtered presentations with their actual structural maps.
\item[ALGEBRA-RECON-800] Source lines 31816--31827: \hyperlink{algebra_colimit_categories_notes_20260928}{Supplementary proof 36}, Filtered presentations with their actual structural maps.
\item[ALGEBRA-RECON-801] Source lines 31861--31867: \hyperlink{algebra_colimit_categories_notes_20260928}{Supplementary proof 36}, Filtered presentations with their actual structural maps.
\item[ALGEBRA-RECON-802] Source lines 31946--31949: \hyperlink{algebra_colimit_categories_notes_20260928}{Supplementary proof 36}, Filtered presentations with their actual structural maps.
\item[ALGEBRA-RECON-803] Source lines 31951--31960: \hyperlink{algebra_colimit_categories_notes_20260928}{Supplementary proof 36}, Filtered presentations with their actual structural maps.
\item[ALGEBRA-RECON-804] Source lines 31982--31986: \hyperlink{algebra_colimit_categories_notes_20260928}{Supplementary proof 36}, Module maps and the exact finite-stage witnesses.
\item[ALGEBRA-RECON-805] Source lines 32016--32022: \hyperlink{algebra_colimit_categories_notes_20260928}{Supplementary proof 36}, Module maps and the exact finite-stage witnesses.
\item[ALGEBRA-RECON-806] Source lines 32025--32147: \hyperlink{algebra_colimit_categories_notes_20260928}{Supplementary proof 36}, Module maps and the exact finite-stage witnesses.
\item[ALGEBRA-RECON-807] Source lines 32034--32037: \hyperlink{algebra_colimit_categories_notes_20260928}{Supplementary proof 36}, Module maps and the exact finite-stage witnesses.
\item[ALGEBRA-RECON-808] Source lines 32050--32057: \hyperlink{algebra_colimit_categories_notes_20260928}{Supplementary proof 36}, Module maps and the exact finite-stage witnesses.
\item[ALGEBRA-RECON-809] Source lines 32131--32138: \hyperlink{algebra_colimit_categories_notes_20260928}{Supplementary proof 36}, Algebra maps, rejected target changes and inverse identities.
\item[ALGEBRA-RECON-810] Source lines 32141--32147: \hyperlink{algebra_colimit_categories_notes_20260928}{Supplementary proof 36}, Algebra maps, rejected target changes and inverse identities.
\item[ALGEBRA-RECON-811] Source lines 32150--32154: \hyperlink{algebra_colimit_categories_notes_20260928}{Supplementary proof 36}, Algebra maps, rejected target changes and inverse identities.
\item[ALGEBRA-RECON-812] Source lines 32160--32164: \hyperlink{algebra_colimit_categories_notes_20260928}{Supplementary proof 36}, Algebra maps, rejected target changes and inverse identities.
\item[ALGEBRA-RECON-813] Source lines 32168--32172: \hyperlink{algebra_colimit_categories_notes_20260928}{Supplementary proof 36}, Algebra maps, rejected target changes and inverse identities.
\item[ALGEBRA-RECON-814] Source lines 32172--32174: \hyperlink{algebra_colimit_categories_notes_20260928}{Supplementary proof 36}, Algebra maps, rejected target changes and inverse identities.
\item[ALGEBRA-RECON-815] Source lines 32194--32199: \hyperlink{algebra_colimit_categories_notes_20260928}{Supplementary proof 36}, Algebra maps, rejected target changes and inverse identities.
\item[ALGEBRA-RECON-816] Source lines 32123--32127: \hyperlink{algebra_colimit_categories_notes_20260928}{Supplementary proof 36}, Algebra maps, rejected target changes and inverse identities.
\item[ALGEBRA-RECON-817] Source lines 31879--31944: \hyperlink{algebra_colimit_categories_notes_20260928}{Supplementary proof 36}, Finite type and the boundary between unions and arbitrary colimits.
\item[ALGEBRA-RECON-818] Source lines 31988--32216: \hyperlink{algebra_colimit_categories_notes_20260928}{Supplementary proof 36}, Descent of finite diagrams and finite equations.
\item[ALGEBRA-RECON-819] Source lines 31988--32216: \hyperlink{algebra_colimit_categories_notes_20260928}{Supplementary proof 36}, Split structures descend but unrestricted exactness need not.
\item[ALGEBRA-RECON-820] Source lines 32234--32294: \hyperlink{algebra_noetherian_approximation_notes_20260928}{Supplementary proof 37}, Subrings, localizations and the actual finite-type maps.
\item[ALGEBRA-RECON-821] Source lines 32345--32351: \hyperlink{algebra_noetherian_approximation_notes_20260928}{Supplementary proof 37}, Finite presentations require a compatible local base system.
\item[ALGEBRA-RECON-822] Source lines 32375--32377: \hyperlink{algebra_noetherian_approximation_notes_20260928}{Supplementary proof 37}, Finite presentations require a compatible local base system.
\item[ALGEBRA-RECON-823] Source lines 32380--32382: \hyperlink{algebra_noetherian_approximation_notes_20260928}{Supplementary proof 37}, Finite presentations require a compatible local base system.
\item[ALGEBRA-RECON-824] Source lines 32383--32390: \hyperlink{algebra_noetherian_approximation_notes_20260928}{Supplementary proof 37}, Finite presentations require a compatible local base system.
\item[ALGEBRA-RECON-825] Source lines 32416--32420: \hyperlink{algebra_noetherian_approximation_notes_20260928}{Supplementary proof 37}, The original wrong system and every later transition.
\item[ALGEBRA-RECON-826] Source lines 32521--32536: \hyperlink{algebra_noetherian_approximation_notes_20260928}{Supplementary proof 37}, Integral approximation with the original monic equations.
\item[ALGEBRA-RECON-827] Source lines 32345--32628: \hyperlink{algebra_noetherian_approximation_notes_20260928}{Supplementary proof 37}, Simultaneous Noetherian models for finite diagrams.
\item[ALGEBRA-RECON-828] Source lines 32394--32424: \hyperlink{algebra_noetherian_approximation_notes_20260928}{Supplementary proof 37}, The original wrong system and every later transition.
\item[ALGEBRA-RECON-829] Source lines 32505--32564: \hyperlink{algebra_noetherian_approximation_notes_20260928}{Supplementary proof 37}, Persistent defects and uniformly finite approximations.
\end{description}

\subsection{More flatness criteria}
\hypertarget{algebra-source-section-more-flatness-criteria}{}
\noindent\href{https://github.com/stacks/stacks-project/blob/a04446e57ec1fbc252a871afcec7752fb2807b14/algebra.tex\#L32643}{Original source section}. \href{algebra.pdf\#nameddest=section-more-flatness-criteria}{Accompanying chapter}.

\begin{description}
\item[ALGEBRA-RECON-830] Source lines 32658--32663: \hyperlink{algebra_more_flatness_notes_20260928}{Supplementary proof 38}, Regular bases and eventual flatness.
\item[ALGEBRA-RECON-831] Source lines 32673--32677: \hyperlink{algebra_more_flatness_notes_20260928}{Supplementary proof 38}, Regular bases and eventual flatness.
\item[ALGEBRA-RECON-832] Source lines 32771--32888: \hyperlink{algebra_more_flatness_notes_20260928}{Supplementary proof 38}, Fibre injectivity with actual fraction representatives.
\item[ALGEBRA-RECON-833] Source lines 32785--32791: \hyperlink{algebra_more_flatness_notes_20260928}{Supplementary proof 38}, Source dispositions and translation treatment.
\item[ALGEBRA-RECON-834] Source lines 32811--32814: \hyperlink{algebra_more_flatness_notes_20260928}{Supplementary proof 38}, Fibre injectivity with actual fraction representatives.
\item[ALGEBRA-RECON-835] Source lines 32833--32847: \hyperlink{algebra_more_flatness_notes_20260928}{Supplementary proof 38}, Fibre injectivity with actual fraction representatives.
\item[ALGEBRA-RECON-836] Source lines 32881--32903: \hyperlink{algebra_more_flatness_notes_20260928}{Supplementary proof 38}, Source dispositions and translation treatment.
\item[ALGEBRA-RECON-837] Source lines 32923--33011: \hyperlink{algebra_more_flatness_notes_20260928}{Supplementary proof 38}, The local Tor criterion with its quotient bases.
\item[ALGEBRA-RECON-838] Source lines 32925--32929: \hyperlink{algebra_more_flatness_notes_20260928}{Supplementary proof 38}, The local Tor criterion with its quotient bases.
\item[ALGEBRA-RECON-839] Source lines 32930--32934: \hyperlink{algebra_more_flatness_notes_20260928}{Supplementary proof 38}, The local Tor criterion with its quotient bases.
\item[ALGEBRA-RECON-840] Source lines 32964--32969: \hyperlink{algebra_more_flatness_notes_20260928}{Supplementary proof 38}, Source dispositions and translation treatment.
\item[ALGEBRA-RECON-841] Source lines 32970--32976: \hyperlink{algebra_more_flatness_notes_20260928}{Supplementary proof 38}, The local Tor criterion with its quotient bases.
\item[ALGEBRA-RECON-842] Source lines 33024--33041: \hyperlink{algebra_more_flatness_notes_20260928}{Supplementary proof 38}, The fibre criterion with all relations and localizations.
\item[ALGEBRA-RECON-843] Source lines 33042--33052: \hyperlink{algebra_more_flatness_notes_20260928}{Supplementary proof 38}, The fibre criterion with all relations and localizations.
\item[ALGEBRA-RECON-844] Source lines 33105--33110: \hyperlink{algebra_more_flatness_notes_20260928}{Supplementary proof 38}, The fibre criterion with all relations and localizations.
\item[ALGEBRA-RECON-845] Source lines 33124--33137: \hyperlink{algebra_more_flatness_notes_20260928}{Supplementary proof 38}, Finite presentation from the finite-type hypothesis.
\item[ALGEBRA-RECON-846] Source lines 33138--33140: \hyperlink{algebra_more_flatness_notes_20260928}{Supplementary proof 38}, Finite presentation from the finite-type hypothesis.
\item[ALGEBRA-RECON-847] Source lines 33142--33146: \hyperlink{algebra_more_flatness_notes_20260928}{Supplementary proof 38}, Source dispositions and translation treatment.
\item[ALGEBRA-RECON-848] Source lines 33157--33160: \hyperlink{algebra_more_flatness_notes_20260928}{Supplementary proof 38}, Finite presentation from the finite-type hypothesis.
\item[ALGEBRA-RECON-849] Source lines 33159--33162: \hyperlink{algebra_more_flatness_notes_20260928}{Supplementary proof 38}, Source dispositions and translation treatment.
\item[ALGEBRA-RECON-850] Source lines 33197--33210: \hyperlink{algebra_more_flatness_notes_20260928}{Supplementary proof 38}, The support prime required by the nil-ideal proof.
\item[ALGEBRA-RECON-851] Source lines 32769--32903: \hyperlink{algebra_more_flatness_notes_20260928}{Supplementary proof 38}, Universal exactness and regular sequences.
\item[ALGEBRA-RECON-852] Source lines 32991--33169: \hyperlink{algebra_more_flatness_notes_20260928}{Supplementary proof 38}, Faithfulness and a pure receiving map.
\item[ALGEBRA-RECON-853] Source lines 33171--33211: \hyperlink{algebra_more_flatness_notes_20260928}{Supplementary proof 38}, Nilness on module support suffices.
\end{description}

\subsection{Openness of the flat locus}
\hypertarget{algebra-source-section-open-flat}{}
\noindent\href{https://github.com/stacks/stacks-project/blob/a04446e57ec1fbc252a871afcec7752fb2807b14/algebra.tex\#L33227}{Original source section}. \href{algebra.pdf\#nameddest=section-open-flat}{Accompanying chapter}.

\begin{description}
\item[ALGEBRA-RECON-854] Source lines 33237--33244: \hyperlink{algebra_open_flat_notes_20260928}{Supplementary proof 39}, The nonempty zero locus and its exact equivalence.
\item[ALGEBRA-RECON-855] Source lines 33241--33244: \hyperlink{algebra_open_flat_notes_20260928}{Supplementary proof 39}, The nonempty zero locus and its exact equivalence.
\item[ALGEBRA-RECON-856] Source lines 33237--33258: \hyperlink{algebra_open_flat_notes_20260928}{Supplementary proof 39}, The nonempty zero locus and its exact equivalence.
\item[ALGEBRA-RECON-857] Source lines 33328--33335: \hyperlink{algebra_open_flat_notes_20260928}{Supplementary proof 39}, Exact fibre complexes under the equidimensional hypothesis.
\item[ALGEBRA-RECON-858] Source lines 33391--33402: \hyperlink{algebra_open_flat_notes_20260928}{Supplementary proof 39}, Exact fibre complexes under the equidimensional hypothesis.
\item[ALGEBRA-RECON-859] Source lines 33332--33402: \hyperlink{algebra_open_flat_notes_20260928}{Supplementary proof 39}, Exact fibre complexes under the equidimensional hypothesis.
\item[ALGEBRA-RECON-860] Source lines 33460--33466: \hyperlink{algebra_open_flat_notes_20260928}{Supplementary proof 39}, Openness of flatness with the zero and one variable cases.
\item[ALGEBRA-RECON-861] Source lines 33475--33490: \hyperlink{algebra_open_flat_notes_20260928}{Supplementary proof 39}, Openness of flatness with the zero and one variable cases.
\item[ALGEBRA-RECON-862] Source lines 33469--33505: \hyperlink{algebra_open_flat_notes_20260928}{Supplementary proof 39}, Openness of flatness with the zero and one variable cases.
\item[ALGEBRA-RECON-863] Source lines 33235--33259: \hyperlink{algebra_open_flat_notes_20260928}{Supplementary proof 39}, The nonempty zero locus and its exact equivalence.
\item[ALGEBRA-RECON-864] Source lines 33306--33403: \hyperlink{algebra_open_flat_notes_20260928}{Supplementary proof 39}, Exact-fibre openness without Cohen–Macaulay assumptions.
\item[ALGEBRA-RECON-865] Source lines 33406--33506: \hyperlink{algebra_open_flat_notes_20260928}{Supplementary proof 39}, The flat open and its radical ideal under base change.
\end{description}

\subsection{Openness of Cohen-Macaulay loci}
\hypertarget{algebra-source-section-CM-open}{}
\noindent\href{https://github.com/stacks/stacks-project/blob/a04446e57ec1fbc252a871afcec7752fb2807b14/algebra.tex\#L33515}{Original source section}. \href{algebra.pdf\#nameddest=section-CM-open}{Accompanying chapter}.

\begin{description}
\item[ALGEBRA-RECON-866] Source lines 33518--33522: \hyperlink{algebra_cm_open_notes_20260928}{Supplementary proof 40}, Report dispositions and the empty-fibre boundary.
\item[ALGEBRA-RECON-867] Source lines 33529--33535: \hyperlink{algebra_cm_open_notes_20260928}{Supplementary proof 40}, The original quasi-finite flat chart criterion.
\item[ALGEBRA-RECON-868] Source lines 33539--33542: \hyperlink{algebra_cm_open_notes_20260928}{Supplementary proof 40}, The original quasi-finite flat chart criterion.
\item[ALGEBRA-RECON-869] Source lines 33571--33576: \hyperlink{algebra_cm_open_notes_20260928}{Supplementary proof 40}, The original quasi-finite flat chart criterion.
\item[ALGEBRA-RECON-870] Source lines 33608--33612: \hyperlink{algebra_cm_open_notes_20260928}{Supplementary proof 40}, The original quasi-finite flat chart criterion.
\item[ALGEBRA-RECON-871] Source lines 33641--33650: \hyperlink{algebra_cm_open_notes_20260928}{Supplementary proof 40}, The original relative proof with fixed dimension retained.
\item[ALGEBRA-RECON-872] Source lines 33662--33664: \hyperlink{algebra_cm_open_notes_20260928}{Supplementary proof 40}, The original relative proof with fixed dimension retained.
\item[ALGEBRA-RECON-873] Source lines 33738--33770: \hyperlink{algebra_cm_open_notes_20260928}{Supplementary proof 40}, Field extension and the actual local tensor comparison.
\item[ALGEBRA-RECON-874] Source lines 33752--33754: \hyperlink{algebra_cm_open_notes_20260928}{Supplementary proof 40}, Report dispositions and the empty-fibre boundary.
\item[ALGEBRA-RECON-875] Source lines 33756--33762: \hyperlink{algebra_cm_open_notes_20260928}{Supplementary proof 40}, Report dispositions and the empty-fibre boundary.
\item[ALGEBRA-RECON-876] Source lines 33756--33763: \hyperlink{algebra_cm_open_notes_20260928}{Supplementary proof 40}, Report dispositions and the empty-fibre boundary.
\item[ALGEBRA-RECON-877] Source lines 33765--33771: \hyperlink{algebra_cm_open_notes_20260928}{Supplementary proof 40}, Canonical dimension idempotents with all nilpotents retained.
\item[ALGEBRA-RECON-878] Source lines 33622--33681: \hyperlink{algebra_cm_open_notes_20260928}{Supplementary proof 40}, A maximal open without assuming global flatness.
\item[ALGEBRA-RECON-879] Source lines 33622--33754: \hyperlink{algebra_cm_open_notes_20260928}{Supplementary proof 40}, Universal dimension strata and their radical ideals.
\item[ALGEBRA-RECON-880] Source lines 33756--33771: \hyperlink{algebra_cm_open_notes_20260928}{Supplementary proof 40}, Canonical dimension idempotents with all nilpotents retained.
\end{description}

\subsection{Differentials}
\hypertarget{algebra-source-section-differentials}{}
\noindent\href{https://github.com/stacks/stacks-project/blob/a04446e57ec1fbc252a871afcec7752fb2807b14/algebra.tex\#L33788}{Original source section}. \href{algebra.pdf\#nameddest=section-differentials}{Accompanying chapter}.

\begin{description}
\item[ALGEBRA-RECON-881] Source lines 33964--33969: \hyperlink{algebra_differentials_notes_20260928}{Supplementary proof 41}, Universal presentation, functoriality and directed colimits.
\item[ALGEBRA-RECON-882] Source lines 33972--33974: \hyperlink{algebra_differentials_notes_20260928}{Supplementary proof 41}, Report dispositions and preservation of the source.
\item[ALGEBRA-RECON-883] Source lines 33999--34003: \hyperlink{algebra_differentials_notes_20260928}{Supplementary proof 41}, The surjective square and both kernel arguments.
\item[ALGEBRA-RECON-884] Source lines 34004--34009: \hyperlink{algebra_differentials_notes_20260928}{Supplementary proof 41}, The surjective square and both kernel arguments.
\item[ALGEBRA-RECON-885] Source lines 34013--34014: \hyperlink{algebra_differentials_notes_20260928}{Supplementary proof 41}, The surjective square and both kernel arguments.
\item[ALGEBRA-RECON-886] Source lines 34024--34025: \hyperlink{algebra_differentials_notes_20260928}{Supplementary proof 41}, The surjective square and both kernel arguments.
\item[ALGEBRA-RECON-887] Source lines 34032--34037: \hyperlink{algebra_differentials_notes_20260928}{Supplementary proof 41}, The surjective square and both kernel arguments.
\item[ALGEBRA-RECON-888] Source lines 34093--34098: \hyperlink{algebra_differentials_notes_20260928}{Supplementary proof 41}, Localization with the complete denominators.
\item[ALGEBRA-RECON-889] Source lines 34093--34094: \hyperlink{algebra_differentials_notes_20260928}{Supplementary proof 41}, Localization with the complete denominators.
\item[ALGEBRA-RECON-890] Source lines 34107--34117: \hyperlink{algebra_differentials_notes_20260928}{Supplementary proof 41}, Fixed base, conormal maps and powers of the ideal.
\item[ALGEBRA-RECON-891] Source lines 34121--34124: \hyperlink{algebra_differentials_notes_20260928}{Supplementary proof 41}, Fixed base, conormal maps and powers of the ideal.
\item[ALGEBRA-RECON-892] Source lines 34127--34130: \hyperlink{algebra_differentials_notes_20260928}{Supplementary proof 41}, Fixed base, conormal maps and powers of the ideal.
\item[ALGEBRA-RECON-893] Source lines 34147--34154: \hyperlink{algebra_differentials_notes_20260928}{Supplementary proof 41}, Report dispositions and preservation of the source.
\item[ALGEBRA-RECON-894] Source lines 34178--34186: \hyperlink{algebra_differentials_notes_20260928}{Supplementary proof 41}, Fixed base, conormal maps and powers of the ideal.
\item[ALGEBRA-RECON-895] Source lines 34267--34272: \hyperlink{algebra_differentials_notes_20260928}{Supplementary proof 41}, Base change, the diagonal and finite presentations.
\item[ALGEBRA-RECON-896] Source lines 34275--34278: \hyperlink{algebra_differentials_notes_20260928}{Supplementary proof 41}, Base change, the diagonal and finite presentations.
\item[ALGEBRA-RECON-897] Source lines 34278--34280: \hyperlink{algebra_differentials_notes_20260928}{Supplementary proof 41}, Report dispositions and preservation of the source.
\item[ALGEBRA-RECON-898] Source lines 34302--34303: \hyperlink{algebra_differentials_notes_20260928}{Supplementary proof 41}, Report dispositions and preservation of the source.
\item[ALGEBRA-RECON-899] Source lines 34351--34362: \hyperlink{algebra_differentials_notes_20260928}{Supplementary proof 41}, Base change, the diagonal and finite presentations.
\item[ALGEBRA-RECON-900] Source lines 34368--34372: \hyperlink{algebra_differentials_notes_20260928}{Supplementary proof 41}, Base change, the diagonal and finite presentations.
\item[ALGEBRA-RECON-901] Source lines 34043--34048: \hyperlink{algebra_differentials_notes_20260928}{Supplementary proof 41}, The surjective square and both kernel arguments.
\item[ALGEBRA-RECON-902] Source lines 34102--34163: \hyperlink{algebra_differentials_notes_20260928}{Supplementary proof 41}, Fixed base, conormal maps and powers of the ideal.
\item[ALGEBRA-RECON-903] Source lines 34102--34210: \hyperlink{algebra_differentials_notes_20260928}{Supplementary proof 41}, First-order sections are exactly conormal retractions.
\item[ALGEBRA-RECON-904] Source lines 34157--34210: \hyperlink{algebra_differentials_notes_20260928}{Supplementary proof 41}, All sections and the exact change of splitting.
\item[ALGEBRA-RECON-905] Source lines 34212--34304: \hyperlink{algebra_differentials_notes_20260928}{Supplementary proof 41}, The diagonal quotient as an algebra, its flip and base change.
\end{description}

\subsection{The de Rham complex}
\hypertarget{algebra-source-section-de-rham-complex}{}
\noindent\href{https://github.com/stacks/stacks-project/blob/a04446e57ec1fbc252a871afcec7752fb2807b14/algebra.tex\#L34387}{Original source section}. \href{algebra.pdf\#nameddest=section-de-rham-complex}{Accompanying chapter}.

\begin{description}
\item[ALGEBRA-RECON-906] Source lines 34391--34393: \hyperlink{algebra_de_rham_notes_20260928}{Supplementary proof 42}, Construction in degree one.
\item[ALGEBRA-RECON-907] Source lines 34408--34416: \hyperlink{algebra_de_rham_notes_20260928}{Supplementary proof 42}, Higher degrees, alternating signs and the graded product rule.
\item[ALGEBRA-RECON-908] Source lines 34461--34467: \hyperlink{algebra_de_rham_notes_20260928}{Supplementary proof 42}, Higher degrees, alternating signs and the graded product rule.
\item[ALGEBRA-RECON-909] Source lines 34489--34500: \hyperlink{algebra_de_rham_notes_20260928}{Supplementary proof 42}, Higher degrees, alternating signs and the graded product rule.
\item[ALGEBRA-RECON-910] Source lines 34533--34537: \hyperlink{algebra_de_rham_notes_20260928}{Supplementary proof 42}, Dispositions and the additive convention.
\item[ALGEBRA-RECON-911] Source lines 34540--34547: \hyperlink{algebra_de_rham_notes_20260928}{Supplementary proof 42}, Functoriality, base change and the quotient construction.
\item[ALGEBRA-RECON-912] Source lines 34547--34551: \hyperlink{algebra_de_rham_notes_20260928}{Supplementary proof 42}, Functoriality, base change and the quotient construction.
\item[ALGEBRA-RECON-913] Source lines 34576--34579: \hyperlink{algebra_de_rham_notes_20260928}{Supplementary proof 42}, Functoriality, base change and the quotient construction.
\item[ALGEBRA-RECON-914] Source lines 34583--34593: \hyperlink{algebra_de_rham_notes_20260928}{Supplementary proof 42}, Functoriality, base change and the quotient construction.
\item[ALGEBRA-RECON-915] Source lines 34595--34610: \hyperlink{algebra_de_rham_notes_20260928}{Supplementary proof 42}, Functoriality, base change and the quotient construction.
\item[ALGEBRA-RECON-916] Source lines 34613--34614: \hyperlink{algebra_de_rham_notes_20260928}{Supplementary proof 42}, Functoriality, base change and the quotient construction.
\item[ALGEBRA-RECON-917] Source lines 34626--34635: \hyperlink{algebra_de_rham_notes_20260928}{Supplementary proof 42}, Functoriality, base change and the quotient construction.
\item[ALGEBRA-RECON-918] Source lines 34525--34537: \hyperlink{algebra_de_rham_notes_20260928}{Supplementary proof 42}, Dispositions and the additive convention.
\item[ALGEBRA-RECON-919] Source lines 34424--34523: \hyperlink{algebra_de_rham_notes_20260928}{Supplementary proof 42}, Higher degrees, alternating signs and the graded product rule.
\item[ALGEBRA-RECON-920] Source lines 34574--34638: \hyperlink{algebra_de_rham_notes_20260928}{Supplementary proof 42}, The exact obstruction to an exterior quotient complex.
\item[ALGEBRA-RECON-921] Source lines 34574--34638: \hyperlink{algebra_de_rham_notes_20260928}{Supplementary proof 42}, The smallest differential quotient when the obstruction is nonzero.
\item[ALGEBRA-RECON-922] Source lines 34559--34572: \hyperlink{algebra_de_rham_notes_20260928}{Supplementary proof 42}, Complex base change and the exact polynomial cohomology boundary.
\end{description}

\subsection{Finite order differential operators}
\hypertarget{algebra-source-section-differential-operators}{}
\noindent\href{https://github.com/stacks/stacks-project/blob/a04446e57ec1fbc252a871afcec7752fb2807b14/algebra.tex\#L34651}{Original source section}. \href{algebra.pdf\#nameddest=section-differential-operators}{Accompanying chapter}.

\begin{description}
\item[ALGEBRA-RECON-923] Source lines 34676--34681: \hyperlink{algebra_diff_operators_notes_20260928}{Supplementary proof 43}, Commutators, composition and the original universal presentation.
\item[ALGEBRA-RECON-924] Source lines 34712--34717: \hyperlink{algebra_diff_operators_notes_20260928}{Supplementary proof 43}, Commutators, composition and the original universal presentation.
\item[ALGEBRA-RECON-925] Source lines 34719--34722: \hyperlink{algebra_diff_operators_notes_20260928}{Supplementary proof 43}, Commutators, composition and the original universal presentation.
\item[ALGEBRA-RECON-926] Source lines 34740--34746: \hyperlink{algebra_diff_operators_notes_20260928}{Supplementary proof 43}, Commutators, composition and the original universal presentation.
\item[ALGEBRA-RECON-927] Source lines 34761--34771: \hyperlink{algebra_diff_operators_notes_20260928}{Supplementary proof 43}, Commutators, composition and the original universal presentation.
\item[ALGEBRA-RECON-928] Source lines 34807--34816: \hyperlink{algebra_diff_operators_notes_20260928}{Supplementary proof 43}, The first principal-parts sequence for every module.
\item[ALGEBRA-RECON-929] Source lines 34817--34819: \hyperlink{algebra_diff_operators_notes_20260928}{Supplementary proof 43}, The first principal-parts sequence for every module.
\item[ALGEBRA-RECON-930] Source lines 34872--34882: \hyperlink{algebra_diff_operators_notes_20260928}{Supplementary proof 43}, Functoriality, base change and order-one de Rham maps.
\item[ALGEBRA-RECON-931] Source lines 34882--34895: \hyperlink{algebra_diff_operators_notes_20260928}{Supplementary proof 43}, Functoriality, base change and order-one de Rham maps.
\item[ALGEBRA-RECON-932] Source lines 34896--34900: \hyperlink{algebra_diff_operators_notes_20260928}{Supplementary proof 43}, Functoriality, base change and order-one de Rham maps.
\item[ALGEBRA-RECON-933] Source lines 34929--34940: \hyperlink{algebra_diff_operators_notes_20260928}{Supplementary proof 43}, The diagonal comparison with both actions retained.
\item[ALGEBRA-RECON-934] Source lines 34942--34945: \hyperlink{algebra_diff_operators_notes_20260928}{Supplementary proof 43}, The diagonal comparison with both actions retained.
\item[ALGEBRA-RECON-935] Source lines 34979--34986: \hyperlink{algebra_diff_operators_notes_20260928}{Supplementary proof 43}, Functoriality, base change and order-one de Rham maps.
\item[ALGEBRA-RECON-936] Source lines 34998--35006: \hyperlink{algebra_diff_operators_notes_20260928}{Supplementary proof 43}, Functoriality, base change and order-one de Rham maps.
\item[ALGEBRA-RECON-937] Source lines 35017--35022: \hyperlink{algebra_diff_operators_notes_20260928}{Supplementary proof 43}, Functoriality, base change and order-one de Rham maps.
\item[ALGEBRA-RECON-938] Source lines 35039--35040: \hyperlink{algebra_diff_operators_notes_20260928}{Supplementary proof 43}, Localization and the complete fraction comparison.
\item[ALGEBRA-RECON-939] Source lines 35043--35047: \hyperlink{algebra_diff_operators_notes_20260928}{Supplementary proof 43}, Localization and the complete fraction comparison.
\item[ALGEBRA-RECON-940] Source lines 35050--35057: \hyperlink{algebra_diff_operators_notes_20260928}{Supplementary proof 43}, Localization and the complete fraction comparison.
\item[ALGEBRA-RECON-941] Source lines 35057--35059: \hyperlink{algebra_diff_operators_notes_20260928}{Supplementary proof 43}, Localization and the complete fraction comparison.
\item[ALGEBRA-RECON-942] Source lines 35065--35074: \hyperlink{algebra_diff_operators_notes_20260928}{Supplementary proof 43}, Localization and the complete fraction comparison.
\item[ALGEBRA-RECON-943] Source lines 35075--35084: \hyperlink{algebra_diff_operators_notes_20260928}{Supplementary proof 43}, Localization and the complete fraction comparison.
\item[ALGEBRA-RECON-944] Source lines 35099--35108: \hyperlink{algebra_diff_operators_notes_20260928}{Supplementary proof 43}, Functoriality, base change and order-one de Rham maps.
\item[ALGEBRA-RECON-945] Source lines 34786--34870: \hyperlink{algebra_diff_operators_notes_20260928}{Supplementary proof 43}, The first principal-parts sequence for every module.
\item[ALGEBRA-RECON-946] Source lines 35025--35086: \hyperlink{algebra_diff_operators_notes_20260928}{Supplementary proof 43}, Localization and the complete fraction comparison.
\item[ALGEBRA-RECON-947] Source lines 34786--34969: \hyperlink{algebra_diff_operators_notes_20260928}{Supplementary proof 43}, Symbols and the exact connection criterion.
\item[ALGEBRA-RECON-948] Source lines 34929--35086: \hyperlink{algebra_diff_operators_notes_20260928}{Supplementary proof 43}, Localized principal parts and finite denominator formulas.
\item[ALGEBRA-RECON-949] Source lines 34700--34969: \hyperlink{algebra_diff_operators_notes_20260928}{Supplementary proof 43}, Polynomial principal parts with all characteristic coefficients.
\end{description}

\subsection{The naive cotangent complex}
\hypertarget{algebra-source-section-netherlander}{}
\noindent\href{https://github.com/stacks/stacks-project/blob/a04446e57ec1fbc252a871afcec7752fb2807b14/algebra.tex\#L35116}{Original source section}. \href{algebra.pdf\#nameddest=section-netherlander}{Accompanying chapter}.

\begin{description}
\item[ALGEBRA-RECON-950] Source lines 35120--35121: \hyperlink{algebra_naive_cotangent_notes_20260928}{Supplementary proof 44}, Presentations and the exact homotopy maps.
\item[ALGEBRA-RECON-951] Source lines 35123--35124: \hyperlink{algebra_naive_cotangent_notes_20260928}{Supplementary proof 44}, Presentations and the exact homotopy maps.
\item[ALGEBRA-RECON-952] Source lines 35151--35153: \hyperlink{algebra_naive_cotangent_notes_20260928}{Supplementary proof 44}, Presentations and the exact homotopy maps.
\item[ALGEBRA-RECON-953] Source lines 35159--35165: \hyperlink{algebra_naive_cotangent_notes_20260928}{Supplementary proof 44}, The original truncation really is isomorphic.
\item[ALGEBRA-RECON-954] Source lines 35165--35170: \hyperlink{algebra_naive_cotangent_notes_20260928}{Supplementary proof 44}, The original truncation really is isomorphic.
\item[ALGEBRA-RECON-955] Source lines 35190--35191: \hyperlink{algebra_naive_cotangent_notes_20260928}{Supplementary proof 44}, Propagation and bounded follow-up.
\item[ALGEBRA-RECON-956] Source lines 35193--35197: \hyperlink{algebra_naive_cotangent_notes_20260928}{Supplementary proof 44}, Presentations and the exact homotopy maps.
\item[ALGEBRA-RECON-957] Source lines 35201--35204: \hyperlink{algebra_naive_cotangent_notes_20260928}{Supplementary proof 44}, Presentations and the exact homotopy maps.
\item[ALGEBRA-RECON-958] Source lines 35229--35231: \hyperlink{algebra_naive_cotangent_notes_20260928}{Supplementary proof 44}, Presentations and the exact homotopy maps.
\item[ALGEBRA-RECON-959] Source lines 35253--35259: \hyperlink{algebra_naive_cotangent_notes_20260928}{Supplementary proof 44}, Presentations and the exact homotopy maps.
\item[ALGEBRA-RECON-960] Source lines 35305--35306: \hyperlink{algebra_naive_cotangent_notes_20260928}{Supplementary proof 44}, Presentations and the exact homotopy maps.
\item[ALGEBRA-RECON-961] Source lines 35321--35327: \hyperlink{algebra_naive_cotangent_notes_20260928}{Supplementary proof 44}, Presentations and the exact homotopy maps.
\item[ALGEBRA-RECON-962] Source lines 35335--35339: \hyperlink{algebra_naive_cotangent_notes_20260928}{Supplementary proof 44}, Presentations and the exact homotopy maps.
\item[ALGEBRA-RECON-963] Source lines 35373--35376: \hyperlink{algebra_naive_cotangent_notes_20260928}{Supplementary proof 44}, Propagation and bounded follow-up.
\item[ALGEBRA-RECON-964] Source lines 35377--35380: \hyperlink{algebra_naive_cotangent_notes_20260928}{Supplementary proof 44}, Propagation and bounded follow-up.
\item[ALGEBRA-RECON-965] Source lines 35380--35384: \hyperlink{algebra_naive_cotangent_notes_20260928}{Supplementary proof 44}, Presentations and the exact homotopy maps.
\item[ALGEBRA-RECON-966] Source lines 35397--35400: \hyperlink{algebra_naive_cotangent_notes_20260928}{Supplementary proof 44}, Propagation and bounded follow-up.
\item[ALGEBRA-RECON-967] Source lines 35417--35432: \hyperlink{algebra_naive_cotangent_notes_20260928}{Supplementary proof 44}, Jacobi-Zariski and the tensor comparison sequence.
\item[ALGEBRA-RECON-968] Source lines 35432--35442: \hyperlink{algebra_naive_cotangent_notes_20260928}{Supplementary proof 44}, Propagation and bounded follow-up.
\item[ALGEBRA-RECON-969] Source lines 35493--35500: \hyperlink{algebra_naive_cotangent_notes_20260928}{Supplementary proof 44}, Jacobi-Zariski and the tensor comparison sequence.
\item[ALGEBRA-RECON-970] Source lines 35513--35515: \hyperlink{algebra_naive_cotangent_notes_20260928}{Supplementary proof 44}, Propagation and bounded follow-up.
\item[ALGEBRA-RECON-971] Source lines 35529--35533: \hyperlink{algebra_naive_cotangent_notes_20260928}{Supplementary proof 44}, Propagation and bounded follow-up.
\item[ALGEBRA-RECON-972] Source lines 35589--35593: \hyperlink{algebra_naive_cotangent_notes_20260928}{Supplementary proof 44}, Localization and why the reported splitting is valid.
\item[ALGEBRA-RECON-973] Source lines 35597--35601: \hyperlink{algebra_naive_cotangent_notes_20260928}{Supplementary proof 44}, Propagation and bounded follow-up.
\item[ALGEBRA-RECON-974] Source lines 35620--35627: \hyperlink{algebra_naive_cotangent_notes_20260928}{Supplementary proof 44}, Localization and why the reported splitting is valid.
\item[ALGEBRA-RECON-975] Source lines 35634--35654: \hyperlink{algebra_naive_cotangent_notes_20260928}{Supplementary proof 44}, Localization and why the reported splitting is valid.
\item[ALGEBRA-RECON-976] Source lines 35660--35661: \hyperlink{algebra_naive_cotangent_notes_20260928}{Supplementary proof 44}, Localization and why the reported splitting is valid.
\item[ALGEBRA-RECON-977] Source lines 35665--35669: \hyperlink{algebra_naive_cotangent_notes_20260928}{Supplementary proof 44}, Propagation and bounded follow-up.
\item[ALGEBRA-RECON-978] Source lines 35672--35674: \hyperlink{algebra_naive_cotangent_notes_20260928}{Supplementary proof 44}, Propagation and bounded follow-up.
\item[ALGEBRA-RECON-979] Source lines 35673--35679: \hyperlink{algebra_naive_cotangent_notes_20260928}{Supplementary proof 44}, Localization and why the reported splitting is valid.
\item[ALGEBRA-RECON-980] Source lines 35680--35685: \hyperlink{algebra_naive_cotangent_notes_20260928}{Supplementary proof 44}, Localization and why the reported splitting is valid.
\item[ALGEBRA-RECON-981] Source lines 35702--35704: \hyperlink{algebra_naive_cotangent_notes_20260928}{Supplementary proof 44}, Localization and why the reported splitting is valid.
\item[ALGEBRA-RECON-982] Source lines 35710--35712: \hyperlink{algebra_naive_cotangent_notes_20260928}{Supplementary proof 44}, Stable conormal isomorphisms with both products.
\item[ALGEBRA-RECON-983] Source lines 35734--35756: \hyperlink{algebra_naive_cotangent_notes_20260928}{Supplementary proof 44}, Stable conormal isomorphisms with both products.
\item[ALGEBRA-RECON-984] Source lines 35734--35757: \hyperlink{algebra_naive_cotangent_notes_20260928}{Supplementary proof 44}, Stable conormal isomorphisms with both products.
\item[ALGEBRA-RECON-985] Source lines 35760--35764: \hyperlink{algebra_naive_cotangent_notes_20260928}{Supplementary proof 44}, Stable conormal isomorphisms with both products.
\item[ALGEBRA-RECON-986] Source lines 35765--35768: \hyperlink{algebra_naive_cotangent_notes_20260928}{Supplementary proof 44}, Stable conormal isomorphisms with both products.
\item[ALGEBRA-RECON-987] Source lines 35775--35780: \hyperlink{algebra_naive_cotangent_notes_20260928}{Supplementary proof 44}, Stable conormal isomorphisms with both products.
\item[ALGEBRA-RECON-988] Source lines 35781--35785: \hyperlink{algebra_naive_cotangent_notes_20260928}{Supplementary proof 44}, Stable conormal isomorphisms with both products.
\item[ALGEBRA-RECON-989] Source lines 35159--35164: \hyperlink{algebra_naive_cotangent_notes_20260928}{Supplementary proof 44}, The original truncation really is isomorphic.
\item[ALGEBRA-RECON-990] Source lines 35580--35587: \hyperlink{algebra_naive_cotangent_notes_20260928}{Supplementary proof 44}, Presentations and the exact homotopy maps.
\item[ALGEBRA-RECON-991] Source lines 35734--35754: \hyperlink{algebra_naive_cotangent_notes_20260928}{Supplementary proof 44}, Stable conormal isomorphisms with both products.
\item[ALGEBRA-RECON-992] Source lines 35788--35797: \hyperlink{algebra_naive_cotangent_notes_20260928}{Supplementary proof 44}, Stable conormal isomorphisms with both products.
\item[ALGEBRA-RECON-993] Source lines 35157--35176: \hyperlink{algebra_naive_cotangent_notes_20260928}{Supplementary proof 44}, The original truncation really is isomorphic.
\item[ALGEBRA-RECON-994] Source lines 35411--35486: \hyperlink{algebra_naive_cotangent_notes_20260928}{Supplementary proof 44}, Jacobi-Zariski and the tensor comparison sequence.
\item[ALGEBRA-RECON-995] Source lines 35536--35554: \hyperlink{algebra_naive_cotangent_notes_20260928}{Supplementary proof 44}, Base change with its full obstruction module.
\item[ALGEBRA-RECON-996] Source lines 35691--35706: \hyperlink{algebra_naive_cotangent_notes_20260928}{Supplementary proof 44}, Two-term quasi-isomorphisms and stronger localization.
\item[ALGEBRA-RECON-997] Source lines 35404--35409: \hyperlink{algebra_naive_cotangent_notes_20260928}{Supplementary proof 44}, Propagation and bounded follow-up.
\end{description}

\subsection{Local complete intersections}
\hypertarget{algebra-source-section-lci}{}
\noindent\href{https://github.com/stacks/stacks-project/blob/a04446e57ec1fbc252a871afcec7752fb2807b14/algebra.tex\#L35805}{Original source section}. \href{algebra.pdf\#nameddest=section-lci}{Accompanying chapter}.

\begin{description}
\item[ALGEBRA-RECON-998] Source lines 35837--35844: \hyperlink{algebra_lci_notes_20260928}{Supplementary proof 45}, Original finite-type presentations and local criteria.
\item[ALGEBRA-RECON-999] Source lines 35863--35865: \hyperlink{algebra_lci_notes_20260928}{Supplementary proof 45}, Report dispositions and propagation.
\item[ALGEBRA-RECON-1000] Source lines 35882--35887: \hyperlink{algebra_lci_notes_20260928}{Supplementary proof 45}, Report dispositions and propagation.
\item[ALGEBRA-RECON-1001] Source lines 36000--36006: \hyperlink{algebra_lci_notes_20260928}{Supplementary proof 45}, Original finite-type presentations and local criteria.
\item[ALGEBRA-RECON-1002] Source lines 36039--36046: \hyperlink{algebra_lci_notes_20260928}{Supplementary proof 45}, Report dispositions and propagation.
\item[ALGEBRA-RECON-1003] Source lines 36049--36052: \hyperlink{algebra_lci_notes_20260928}{Supplementary proof 45}, Changing regular presentations and the quantitative defect.
\item[ALGEBRA-RECON-1004] Source lines 36065--36070: \hyperlink{algebra_lci_notes_20260928}{Supplementary proof 45}, Changing regular presentations and the quantitative defect.
\item[ALGEBRA-RECON-1005] Source lines 36171--36174: \hyperlink{algebra_lci_notes_20260928}{Supplementary proof 45}, Local presentations, neighbourhoods and field change.
\item[ALGEBRA-RECON-1006] Source lines 36234--36257: \hyperlink{algebra_lci_notes_20260928}{Supplementary proof 45}, Local presentations, neighbourhoods and field change.
\item[ALGEBRA-RECON-1007] Source lines 36282--36285: \hyperlink{algebra_lci_notes_20260928}{Supplementary proof 45}, Local presentations, neighbourhoods and field change.
\item[ALGEBRA-RECON-1008] Source lines 36286--36293: \hyperlink{algebra_lci_notes_20260928}{Supplementary proof 45}, Local presentations, neighbourhoods and field change.
\item[ALGEBRA-RECON-1009] Source lines 36307--36310: \hyperlink{algebra_lci_notes_20260928}{Supplementary proof 45}, Local presentations, neighbourhoods and field change.
\item[ALGEBRA-RECON-1010] Source lines 36315--36320: \hyperlink{algebra_lci_notes_20260928}{Supplementary proof 45}, Report dispositions and propagation.
\item[ALGEBRA-RECON-1011] Source lines 36332--36339: \hyperlink{algebra_lci_notes_20260928}{Supplementary proof 45}, Local presentations, neighbourhoods and field change.
\item[ALGEBRA-RECON-1012] Source lines 36362--36369: \hyperlink{algebra_lci_notes_20260928}{Supplementary proof 45}, Report dispositions and propagation.
\item[ALGEBRA-RECON-1013] Source lines 36369--36375: \hyperlink{algebra_lci_notes_20260928}{Supplementary proof 45}, Permanence without the unused regularity assumptions.
\item[ALGEBRA-RECON-1014] Source lines 36389--36396: \hyperlink{algebra_lci_notes_20260928}{Supplementary proof 45}, Permanence without the unused regularity assumptions.
\item[ALGEBRA-RECON-1015] Source lines 36414--36418: \hyperlink{algebra_lci_notes_20260928}{Supplementary proof 45}, Report dispositions and propagation.
\item[ALGEBRA-RECON-1016] Source lines 36360--36420: \hyperlink{algebra_lci_notes_20260928}{Supplementary proof 45}, The permanence proof has a genuine hypothesis gap.
\item[ALGEBRA-RECON-1017] Source lines 35908--36259: \hyperlink{algebra_lci_notes_20260928}{Supplementary proof 45}, Changing regular presentations and the quantitative defect.
\item[ALGEBRA-RECON-1018] Source lines 35908--36353: \hyperlink{algebra_lci_notes_20260928}{Supplementary proof 45}, The closed loci of each defect value.
\item[ALGEBRA-RECON-1019] Source lines 36360--36420: \hyperlink{algebra_lci_notes_20260928}{Supplementary proof 45}, Permanence without the unused regularity assumptions.
\end{description}

\subsection{Syntomic morphisms}
\hypertarget{algebra-source-section-syntomic}{}
\noindent\href{https://github.com/stacks/stacks-project/blob/a04446e57ec1fbc252a871afcec7752fb2807b14/algebra.tex\#L36427}{Original source section}. \href{algebra.pdf\#nameddest=section-syntomic}{Accompanying chapter}.

\begin{description}
\item[ALGEBRA-RECON-1020] Source lines 36465--36470: \hyperlink{algebra_syntomic_notes_20260928}{Supplementary proof 46}, Descent, base change and the original presentations.
\item[ALGEBRA-RECON-1021] Source lines 36491--36494: \hyperlink{algebra_syntomic_notes_20260928}{Supplementary proof 46}, Descent, base change and the original presentations.
\item[ALGEBRA-RECON-1022] Source lines 36542--36545: \hyperlink{algebra_syntomic_notes_20260928}{Supplementary proof 46}, Report dispositions and propagation.
\item[ALGEBRA-RECON-1023] Source lines 36577--36582: \hyperlink{algebra_syntomic_notes_20260928}{Supplementary proof 46}, Finite locally free factorization algebras and their exact rank.
\item[ALGEBRA-RECON-1024] Source lines 36589--36590: \hyperlink{algebra_syntomic_notes_20260928}{Supplementary proof 46}, Report dispositions and propagation.
\item[ALGEBRA-RECON-1025] Source lines 36623--36625: \hyperlink{algebra_syntomic_notes_20260928}{Supplementary proof 46}, Report dispositions and propagation.
\item[ALGEBRA-RECON-1026] Source lines 36640--36643: \hyperlink{algebra_syntomic_notes_20260928}{Supplementary proof 46}, Report dispositions and propagation.
\item[ALGEBRA-RECON-1027] Source lines 36651--36653: \hyperlink{algebra_syntomic_notes_20260928}{Supplementary proof 46}, Descent, base change and the original presentations.
\item[ALGEBRA-RECON-1028] Source lines 36731--36733: \hyperlink{algebra_syntomic_notes_20260928}{Supplementary proof 46}, Neighbourhoods, approximation and local regularity.
\item[ALGEBRA-RECON-1029] Source lines 36744--36746: \hyperlink{algebra_syntomic_notes_20260928}{Supplementary proof 46}, Neighbourhoods, approximation and local regularity.
\item[ALGEBRA-RECON-1030] Source lines 36773--36775: \hyperlink{algebra_syntomic_notes_20260928}{Supplementary proof 46}, Report dispositions and propagation.
\item[ALGEBRA-RECON-1031] Source lines 36805--36807: \hyperlink{algebra_syntomic_notes_20260928}{Supplementary proof 46}, Neighbourhoods, approximation and local regularity.
\item[ALGEBRA-RECON-1032] Source lines 36834--36835: \hyperlink{algebra_syntomic_notes_20260928}{Supplementary proof 46}, Neighbourhoods, approximation and local regularity.
\item[ALGEBRA-RECON-1033] Source lines 36836--36839: \hyperlink{algebra_syntomic_notes_20260928}{Supplementary proof 46}, Report dispositions and propagation.
\item[ALGEBRA-RECON-1034] Source lines 36842--36851: \hyperlink{algebra_syntomic_notes_20260928}{Supplementary proof 46}, Neighbourhoods, approximation and local regularity.
\item[ALGEBRA-RECON-1035] Source lines 36903--36906: \hyperlink{algebra_syntomic_notes_20260928}{Supplementary proof 46}, Local criteria, composition and lifting.
\item[ALGEBRA-RECON-1036] Source lines 36925--36929: \hyperlink{algebra_syntomic_notes_20260928}{Supplementary proof 46}, Local criteria, composition and lifting.
\item[ALGEBRA-RECON-1037] Source lines 36977--36979: \hyperlink{algebra_syntomic_notes_20260928}{Supplementary proof 46}, Report dispositions and propagation.
\item[ALGEBRA-RECON-1038] Source lines 36984--36988: \hyperlink{algebra_syntomic_notes_20260928}{Supplementary proof 46}, Local criteria, composition and lifting.
\item[ALGEBRA-RECON-1039] Source lines 37000--37010: \hyperlink{algebra_syntomic_notes_20260928}{Supplementary proof 46}, Local criteria, composition and lifting.
\item[ALGEBRA-RECON-1040] Source lines 37023--37027: \hyperlink{algebra_syntomic_notes_20260928}{Supplementary proof 46}, Local criteria, composition and lifting.
\item[ALGEBRA-RECON-1041] Source lines 37021--37037: \hyperlink{algebra_syntomic_notes_20260928}{Supplementary proof 46}, Local criteria, composition and lifting.
\item[ALGEBRA-RECON-1042] Source lines 36550--36874: \hyperlink{algebra_syntomic_notes_20260928}{Supplementary proof 46}, Finite locally free factorization algebras and their exact rank.
\item[ALGEBRA-RECON-1043] Source lines 36638--36854: \hyperlink{algebra_syntomic_notes_20260928}{Supplementary proof 46}, Universal Koszul resolution and every infinitesimal neighbourhood.
\end{description}

\subsection{Smooth ring maps}
\hypertarget{algebra-source-section-smooth}{}
\noindent\href{https://github.com/stacks/stacks-project/blob/a04446e57ec1fbc252a871afcec7752fb2807b14/algebra.tex\#L37055}{Original source section}. \href{algebra.pdf\#nameddest=section-smooth}{Accompanying chapter}.

\begin{description}
\item[ALGEBRA-RECON-1044] Source lines 37067--37070: \hyperlink{algebra_smooth_notes_20260928}{Supplementary proof 47}, The source definitions and the split conormal map.
\item[ALGEBRA-RECON-1045] Source lines 37121--37125: \hyperlink{algebra_smooth_notes_20260928}{Supplementary proof 47}, The source definitions and the split conormal map.
\item[ALGEBRA-RECON-1046] Source lines 37134--37136: \hyperlink{algebra_smooth_notes_20260928}{Supplementary proof 47}, Report decisions and propagation.
\item[ALGEBRA-RECON-1047] Source lines 37228--37231: \hyperlink{algebra_smooth_notes_20260928}{Supplementary proof 47}, Standard smooth charts with exact Jacobian matrices.
\item[ALGEBRA-RECON-1048] Source lines 37237--37242: \hyperlink{algebra_smooth_notes_20260928}{Supplementary proof 47}, Standard smooth charts with exact Jacobian matrices.
\item[ALGEBRA-RECON-1049] Source lines 37293--37299: \hyperlink{algebra_smooth_notes_20260928}{Supplementary proof 47}, Standard smooth charts with exact Jacobian matrices.
\item[ALGEBRA-RECON-1050] Source lines 37316--37319: \hyperlink{algebra_smooth_notes_20260928}{Supplementary proof 47}, Standard smooth charts with exact Jacobian matrices.
\item[ALGEBRA-RECON-1051] Source lines 37383--37386: \hyperlink{algebra_smooth_notes_20260928}{Supplementary proof 47}, Empty fibres and the composition lifts.
\item[ALGEBRA-RECON-1052] Source lines 37524--37526: \hyperlink{algebra_smooth_notes_20260928}{Supplementary proof 47}, Report decisions and propagation.
\item[ALGEBRA-RECON-1053] Source lines 37528--37533: \hyperlink{algebra_smooth_notes_20260928}{Supplementary proof 47}, The source definitions and the split conormal map.
\item[ALGEBRA-RECON-1054] Source lines 37558--37564: \hyperlink{algebra_smooth_notes_20260928}{Supplementary proof 47}, The local criterion and the correct splitting.
\item[ALGEBRA-RECON-1055] Source lines 37573--37575: \hyperlink{algebra_smooth_notes_20260928}{Supplementary proof 47}, Report decisions and propagation.
\item[ALGEBRA-RECON-1056] Source lines 37635--37639: \hyperlink{algebra_smooth_notes_20260928}{Supplementary proof 47}, Standard smooth charts with exact Jacobian matrices.
\item[ALGEBRA-RECON-1057] Source lines 37647--37650: \hyperlink{algebra_smooth_notes_20260928}{Supplementary proof 47}, Report decisions and propagation.
\item[ALGEBRA-RECON-1058] Source lines 37656--37661: \hyperlink{algebra_smooth_notes_20260928}{Supplementary proof 47}, Report decisions and propagation.
\item[ALGEBRA-RECON-1059] Source lines 37690--37694: \hyperlink{algebra_smooth_notes_20260928}{Supplementary proof 47}, The local criterion and the correct splitting.
\item[ALGEBRA-RECON-1060] Source lines 37714--37717: \hyperlink{algebra_smooth_notes_20260928}{Supplementary proof 47}, The local criterion and the correct splitting.
\item[ALGEBRA-RECON-1061] Source lines 37719--37720: \hyperlink{algebra_smooth_notes_20260928}{Supplementary proof 47}, The local criterion and the correct splitting.
\item[ALGEBRA-RECON-1062] Source lines 37726--37737: \hyperlink{algebra_smooth_notes_20260928}{Supplementary proof 47}, The source definitions and the split conormal map.
\item[ALGEBRA-RECON-1063] Source lines 37762--37765: \hyperlink{algebra_smooth_notes_20260928}{Supplementary proof 47}, The local criterion and the correct splitting.
\item[ALGEBRA-RECON-1064] Source lines 37764--37774: \hyperlink{algebra_smooth_notes_20260928}{Supplementary proof 47}, The local criterion and the correct splitting.
\item[ALGEBRA-RECON-1065] Source lines 37326--37336: \hyperlink{algebra_smooth_notes_20260928}{Supplementary proof 47}, Empty fibres and the composition lifts.
\item[ALGEBRA-RECON-1066] Source lines 37383--37434: \hyperlink{algebra_smooth_notes_20260928}{Supplementary proof 47}, Empty fibres and the composition lifts.
\item[ALGEBRA-RECON-1067] Source lines 37672--37758: \hyperlink{algebra_smooth_notes_20260928}{Supplementary proof 47}, Smooth loci under arbitrary base change of a flat map.
\item[ALGEBRA-RECON-1068] Source lines 37628--37758: \hyperlink{algebra_smooth_notes_20260928}{Supplementary proof 47}, Intrinsic Jacobian defect and its full Fitting ideals.
\item[ALGEBRA-RECON-1069] Source lines 37628--37670: \hyperlink{algebra_smooth_notes_20260928}{Supplementary proof 47}, Intrinsic Jacobian defect and its full Fitting ideals.
\end{description}

\subsection{Formally smooth maps}
\hypertarget{algebra-source-section-formally-smooth}{}
\noindent\href{https://github.com/stacks/stacks-project/blob/a04446e57ec1fbc252a871afcec7752fb2807b14/algebra.tex\#L37795}{Original source section}. \href{algebra.pdf\#nameddest=section-formally-smooth}{Accompanying chapter}.

\begin{description}
\item[ALGEBRA-RECON-1070] Source lines 37855--37859: \hyperlink{algebra_formal_smooth_notes_20260928}{Supplementary proof 48}, Square-zero lifts and the original conormal criterion.
\item[ALGEBRA-RECON-1071] Source lines 37863--37869: \hyperlink{algebra_formal_smooth_notes_20260928}{Supplementary proof 48}, Square-zero lifts and the original conormal criterion.
\item[ALGEBRA-RECON-1072] Source lines 37912--37920: \hyperlink{algebra_formal_smooth_notes_20260928}{Supplementary proof 48}, Square-zero lifts and the original conormal criterion.
\item[ALGEBRA-RECON-1073] Source lines 37934--38141: \hyperlink{algebra_formal_smooth_notes_20260928}{Supplementary proof 48}, Square-zero lifts and the original conormal criterion.
\item[ALGEBRA-RECON-1074] Source lines 37946--37962: \hyperlink{algebra_formal_smooth_notes_20260928}{Supplementary proof 48}, Square-zero lifts and the original conormal criterion.
\item[ALGEBRA-RECON-1075] Source lines 37995--37998: \hyperlink{algebra_formal_smooth_notes_20260928}{Supplementary proof 48}, Square-zero lifts and the original conormal criterion.
\item[ALGEBRA-RECON-1076] Source lines 38023--38028: \hyperlink{algebra_formal_smooth_notes_20260928}{Supplementary proof 48}, Square-zero lifts and the original conormal criterion.
\item[ALGEBRA-RECON-1077] Source lines 38031--38035: \hyperlink{algebra_formal_smooth_notes_20260928}{Supplementary proof 48}, Square-zero lifts and the original conormal criterion.
\item[ALGEBRA-RECON-1078] Source lines 38080--38083: \hyperlink{algebra_formal_smooth_notes_20260928}{Supplementary proof 48}, Projective differentials and the direction of the splittings.
\item[ALGEBRA-RECON-1079] Source lines 38154--38160: \hyperlink{algebra_formal_smooth_notes_20260928}{Supplementary proof 48}, Square-zero lifts and the original conormal criterion.
\item[ALGEBRA-RECON-1080] Source lines 38168--38174: \hyperlink{algebra_formal_smooth_notes_20260928}{Supplementary proof 48}, Square-zero lifts and the original conormal criterion.
\item[ALGEBRA-RECON-1081] Source lines 38313--38318: \hyperlink{algebra_formal_smooth_notes_20260928}{Supplementary proof 48}, Finite witnesses, descent and report accounting.
\item[ALGEBRA-RECON-1082] Source lines 37965--37980: \hyperlink{algebra_formal_smooth_notes_20260928}{Supplementary proof 48}, Square-zero lifts and the original conormal criterion.
\item[ALGEBRA-RECON-1083] Source lines 38078--38165: \hyperlink{algebra_formal_smooth_notes_20260928}{Supplementary proof 48}, Projective differentials and the direction of the splittings.
\item[ALGEBRA-RECON-1084] Source lines 38100--38165: \hyperlink{algebra_formal_smooth_notes_20260928}{Supplementary proof 48}, Nilpotent base ideals and the exact conormal obstruction.
\item[ALGEBRA-RECON-1085] Source lines 37804--38333: \hyperlink{algebra_formal_smooth_notes_20260928}{Supplementary proof 48}, All lifts, finite relation errors and the finite-presentation boundary.
\end{description}

\subsection{Smoothness and differentials}
\hypertarget{algebra-source-section-smooth-differential}{}
\noindent\href{https://github.com/stacks/stacks-project/blob/a04446e57ec1fbc252a871afcec7752fb2807b14/algebra.tex\#L38340}{Original source section}. \href{algebra.pdf\#nameddest=section-smooth-differential}{Accompanying chapter}.

\begin{description}
\item[ALGEBRA-RECON-1086] Source lines 38346--38352: \hyperlink{algebra_smooth_fields_notes_20260928}{Supplementary proof 49}, Exact sequences and the completion along the original section.
\item[ALGEBRA-RECON-1087] Source lines 38355--38363: \hyperlink{algebra_smooth_fields_notes_20260928}{Supplementary proof 49}, Generic smoothness and report dispositions.
\item[ALGEBRA-RECON-1088] Source lines 38385--38390: \hyperlink{algebra_smooth_fields_notes_20260928}{Supplementary proof 49}, Exact sequences and the completion along the original section.
\item[ALGEBRA-RECON-1089] Source lines 38406--38412: \hyperlink{algebra_smooth_fields_notes_20260928}{Supplementary proof 49}, Exact sequences and the completion along the original section.
\item[ALGEBRA-RECON-1090] Source lines 38408--38412: \hyperlink{algebra_smooth_fields_notes_20260928}{Supplementary proof 49}, Exact sequences and the completion along the original section.
\item[ALGEBRA-RECON-1091] Source lines 38430--38437: \hyperlink{algebra_smooth_fields_notes_20260928}{Supplementary proof 49}, Exact sequences and the completion along the original section.
\item[ALGEBRA-RECON-1092] Source lines 38438--38440: \hyperlink{algebra_smooth_fields_notes_20260928}{Supplementary proof 49}, Exact sequences and the completion along the original section.
\item[ALGEBRA-RECON-1103] Source lines 38346--38460: \hyperlink{algebra_smooth_fields_notes_20260928}{Supplementary proof 49}, Exact sequences and the completion along the original section.
\end{description}

\subsection{Smooth algebras over fields}
\hypertarget{algebra-source-section-smooth-over-field}{}
\noindent\href{https://github.com/stacks/stacks-project/blob/a04446e57ec1fbc252a871afcec7752fb2807b14/algebra.tex\#L38467}{Original source section}. \href{algebra.pdf\#nameddest=section-smooth-over-field}{Accompanying chapter}.

\begin{description}
\item[ALGEBRA-RECON-1093] Source lines 38540--38550: \hyperlink{algebra_smooth_fields_notes_20260928}{Supplementary proof 49}, The regular-local hypotheses and the field smoothness criterion.
\item[ALGEBRA-RECON-1094] Source lines 38551--38557: \hyperlink{algebra_smooth_fields_notes_20260928}{Supplementary proof 49}, The regular-local hypotheses and the field smoothness criterion.
\item[ALGEBRA-RECON-1095] Source lines 38564--38575: \hyperlink{algebra_smooth_fields_notes_20260928}{Supplementary proof 49}, The regular-local hypotheses and the field smoothness criterion.
\item[ALGEBRA-RECON-1096] Source lines 38577--38588: \hyperlink{algebra_smooth_fields_notes_20260928}{Supplementary proof 49}, The regular-local hypotheses and the field smoothness criterion.
\item[ALGEBRA-RECON-1097] Source lines 38752--38755: \hyperlink{algebra_smooth_fields_notes_20260928}{Supplementary proof 49}, Residue-field sections without the Noetherian restriction.
\item[ALGEBRA-RECON-1098] Source lines 38753--38757: \hyperlink{algebra_smooth_fields_notes_20260928}{Supplementary proof 49}, Residue-field sections without the Noetherian restriction.
\item[ALGEBRA-RECON-1099] Source lines 38767--38780: \hyperlink{algebra_smooth_fields_notes_20260928}{Supplementary proof 49}, Residue-field sections without the Noetherian restriction.
\item[ALGEBRA-RECON-1100] Source lines 38869--38884: \hyperlink{algebra_smooth_fields_notes_20260928}{Supplementary proof 49}, Characteristic-p examples with every coefficient retained.
\item[ALGEBRA-RECON-1101] Source lines 38897--38909: \hyperlink{algebra_smooth_fields_notes_20260928}{Supplementary proof 49}, Generic smoothness and report dispositions.
\item[ALGEBRA-RECON-1102] Source lines 38783--38815: \hyperlink{algebra_smooth_fields_notes_20260928}{Supplementary proof 49}, The characteristic-zero derivation and its exact boundary.
\item[ALGEBRA-RECON-1104] Source lines 38686--38780: \hyperlink{algebra_smooth_fields_notes_20260928}{Supplementary proof 49}, Residue-field sections without the Noetherian restriction.
\item[ALGEBRA-RECON-1105] Source lines 38783--38867: \hyperlink{algebra_smooth_fields_notes_20260928}{Supplementary proof 49}, The characteristic-zero derivation and its exact boundary.
\item[ALGEBRA-RECON-1106] Source lines 38869--38884: \hyperlink{algebra_smooth_fields_notes_20260928}{Supplementary proof 49}, Characteristic-p examples with every coefficient retained.
\end{description}

\subsection{Smooth ring maps in the Noetherian case}
\hypertarget{algebra-source-section-smooth-Noetherian}{}
\noindent\href{https://github.com/stacks/stacks-project/blob/a04446e57ec1fbc252a871afcec7752fb2807b14/algebra.tex\#L38933}{Original source section}. \href{algebra.pdf\#nameddest=section-smooth-Noetherian}{Accompanying chapter}.

\begin{description}
\item[ALGEBRA-RECON-1107] Source lines 38949--38954: \hyperlink{algebra_smooth_artin_notes_20260928}{Supplementary proof 50}, The original small-extension argument.
\item[ALGEBRA-RECON-1108] Source lines 38957--38968: \hyperlink{algebra_smooth_artin_notes_20260928}{Supplementary proof 50}, The original small-extension argument.
\item[ALGEBRA-RECON-1109] Source lines 38998--39013: \hyperlink{algebra_smooth_artin_notes_20260928}{Supplementary proof 50}, The original small-extension argument.
\item[ALGEBRA-RECON-1110] Source lines 39004--39008: \hyperlink{algebra_smooth_artin_notes_20260928}{Supplementary proof 50}, The original small-extension argument.
\item[ALGEBRA-RECON-1111] Source lines 39057--39065: \hyperlink{algebra_smooth_artin_notes_20260928}{Supplementary proof 50}, The original small-extension argument.
\item[ALGEBRA-RECON-1112] Source lines 39099--39110: \hyperlink{algebra_smooth_artin_notes_20260928}{Supplementary proof 50}, The original small-extension argument.
\item[ALGEBRA-RECON-1113] Source lines 39111--39119: \hyperlink{algebra_smooth_artin_notes_20260928}{Supplementary proof 50}, The original small-extension argument.
\item[ALGEBRA-RECON-1118] Source lines 38991--39138: \hyperlink{algebra_smooth_artin_notes_20260928}{Supplementary proof 50}, A residue conormal criterion and its full base-change defect.
\item[ALGEBRA-RECON-1119] Source lines 38949--39138: \hyperlink{algebra_smooth_artin_notes_20260928}{Supplementary proof 50}, A single Artinian order and the weaker local Noetherian hypothesis.
\end{description}

\subsection{Overview of results on smooth ring maps}
\hypertarget{algebra-source-section-smooth-overview}{}
\noindent\href{https://github.com/stacks/stacks-project/blob/a04446e57ec1fbc252a871afcec7752fb2807b14/algebra.tex\#L39148}{Original source section}. \href{algebra.pdf\#nameddest=section-smooth-overview}{Accompanying chapter}.

\begin{description}
\item[ALGEBRA-RECON-1114] Source lines 39186--39189: \hyperlink{algebra_smooth_artin_notes_20260928}{Supplementary proof 50}, Report accounting and propagation.
\item[ALGEBRA-RECON-1115] Source lines 39186--39195: \hyperlink{algebra_smooth_artin_notes_20260928}{Supplementary proof 50}, Report accounting and propagation.
\item[ALGEBRA-RECON-1116] Source lines 39196--39216: \hyperlink{algebra_smooth_artin_notes_20260928}{Supplementary proof 50}, Report accounting and propagation.
\item[ALGEBRA-RECON-1117] Source lines 39179--39185: \hyperlink{algebra_smooth_artin_notes_20260928}{Supplementary proof 50}, The overview's finite-presentation boundary.
\end{description}

\subsection{\'Etale ring maps}
\hypertarget{algebra-source-section-etale}{}
\noindent\href{https://github.com/stacks/stacks-project/blob/a04446e57ec1fbc252a871afcec7752fb2807b14/algebra.tex\#L39236}{Original source section}. \href{algebra.pdf\#nameddest=section-etale}{Accompanying chapter}.

\begin{description}
\item[ALGEBRA-RECON-1120] Source lines 39434--39444: \hyperlink{algebra_etale_notes_20260928}{Supplementary proof 51}, Original presentations and permanence arguments.
\item[ALGEBRA-RECON-1121] Source lines 39460--39466: \hyperlink{algebra_etale_notes_20260928}{Supplementary proof 51}, Original presentations and permanence arguments.
\item[ALGEBRA-RECON-1122] Source lines 39484--39501: \hyperlink{algebra_etale_notes_20260928}{Supplementary proof 51}, Original presentations and permanence arguments.
\item[ALGEBRA-RECON-1123] Source lines 39533--39538: \hyperlink{algebra_etale_notes_20260928}{Supplementary proof 51}, Original presentations and permanence arguments.
\item[ALGEBRA-RECON-1124] Source lines 39622--39625: \hyperlink{algebra_etale_notes_20260928}{Supplementary proof 51}, The full factorization matrix and its fibres.
\item[ALGEBRA-RECON-1125] Source lines 39684--39747: \hyperlink{algebra_etale_notes_20260928}{Supplementary proof 51}, Original factors and report accounting.
\item[ALGEBRA-RECON-1126] Source lines 39727--39735: \hyperlink{algebra_etale_notes_20260928}{Supplementary proof 51}, Original factors and report accounting.
\item[ALGEBRA-RECON-1127] Source lines 39564--39620: \hyperlink{algebra_etale_notes_20260928}{Supplementary proof 51}, Canonical infinitesimal map and nil ideals.
\item[ALGEBRA-RECON-1128] Source lines 39622--39748: \hyperlink{algebra_etale_notes_20260928}{Supplementary proof 51}, The full factorization matrix and its fibres.
\end{description}

\subsection{Local structure of \'etale ring maps}
\hypertarget{algebra-source-section-etale-local-structure}{}
\noindent\href{https://github.com/stacks/stacks-project/blob/a04446e57ec1fbc252a871afcec7752fb2807b14/algebra.tex\#L39754}{Original source section}. \href{algebra.pdf\#nameddest=section-etale-local-structure}{Accompanying chapter}.

\begin{description}
\item[ALGEBRA-RECON-1129] Source lines 39768--39774: \hyperlink{algebra_etale_local_notes_20260928}{Supplementary proof 52}, Report accounting and propagation.
\item[ALGEBRA-RECON-1130] Source lines 39809--39818: \hyperlink{algebra_etale_local_notes_20260928}{Supplementary proof 52}, Report accounting and propagation.
\item[ALGEBRA-RECON-1131] Source lines 39861--39868: \hyperlink{algebra_etale_local_notes_20260928}{Supplementary proof 52}, Original local-structure proof and the missing scalar.
\item[ALGEBRA-RECON-1132] Source lines 39863--39866: \hyperlink{algebra_etale_local_notes_20260928}{Supplementary proof 52}, Original local-structure proof and the missing scalar.
\item[ALGEBRA-RECON-1133] Source lines 39871--39881: \hyperlink{algebra_etale_local_notes_20260928}{Supplementary proof 52}, Original local-structure proof and the missing scalar.
\item[ALGEBRA-RECON-1134] Source lines 39884--39896: \hyperlink{algebra_etale_local_notes_20260928}{Supplementary proof 52}, Original local-structure proof and the missing scalar.
\item[ALGEBRA-RECON-1135] Source lines 39930--39934: \hyperlink{algebra_etale_local_notes_20260928}{Supplementary proof 52}, Original local-structure proof and the missing scalar.
\item[ALGEBRA-RECON-1136] Source lines 39934--39937: \hyperlink{algebra_etale_local_notes_20260928}{Supplementary proof 52}, Report accounting and propagation.
\item[ALGEBRA-RECON-1137] Source lines 39963--39969: \hyperlink{algebra_etale_local_notes_20260928}{Supplementary proof 52}, Original local-structure proof and the missing scalar.
\item[ALGEBRA-RECON-1138] Source lines 39981--39995: \hyperlink{algebra_etale_local_notes_20260928}{Supplementary proof 52}, Original local-structure proof and the missing scalar.
\item[ALGEBRA-RECON-1139] Source lines 39998--40001: \hyperlink{algebra_etale_local_notes_20260928}{Supplementary proof 52}, Original local-structure proof and the missing scalar.
\item[ALGEBRA-RECON-1140] Source lines 40069--40074: \hyperlink{algebra_etale_local_notes_20260928}{Supplementary proof 52}, Report accounting and propagation.
\item[ALGEBRA-RECON-1141] Source lines 40111--40118: \hyperlink{algebra_etale_local_notes_20260928}{Supplementary proof 52}, Explicit finite free covers and every chosen fibre point.
\item[ALGEBRA-RECON-1142] Source lines 40166--40170: \hyperlink{algebra_etale_local_notes_20260928}{Supplementary proof 52}, Explicit finite free covers and every chosen fibre point.
\item[ALGEBRA-RECON-1143] Source lines 40168--40172: \hyperlink{algebra_etale_local_notes_20260928}{Supplementary proof 52}, Explicit finite free covers and every chosen fibre point.
\item[ALGEBRA-RECON-1144] Source lines 40170--40178: \hyperlink{algebra_etale_local_notes_20260928}{Supplementary proof 52}, Explicit finite free covers and every chosen fibre point.
\item[ALGEBRA-RECON-1145] Source lines 39972--40044: \hyperlink{algebra_etale_local_notes_20260928}{Supplementary proof 52}, Original local-structure proof and the missing scalar.
\item[ALGEBRA-RECON-1146] Source lines 39781--39803: \hyperlink{algebra_etale_local_notes_20260928}{Supplementary proof 52}, Monogenic étale algebras over fields.
\item[ALGEBRA-RECON-1147] Source lines 40069--40180: \hyperlink{algebra_etale_local_notes_20260928}{Supplementary proof 52}, Explicit finite free covers and every chosen fibre point.
\item[ALGEBRA-RECON-1148] Source lines 39843--40062: \hyperlink{algebra_etale_local_notes_20260928}{Supplementary proof 52}, Direct finite-cokernel comparison over an arbitrary base.
\end{description}

\subsection{\'Etale local structure of quasi-finite ring maps}
\hypertarget{algebra-source-section-etale-local-quasi-finite}{}
\noindent\href{https://github.com/stacks/stacks-project/blob/a04446e57ec1fbc252a871afcec7752fb2807b14/algebra.tex\#L40188}{Original source section}. \href{algebra.pdf\#nameddest=section-etale-local-quasi-finite}{Accompanying chapter}.

\begin{description}
\item[ALGEBRA-RECON-1149] Source lines 40210--40210: \hyperlink{algebra_etale_quasifinite_notes_20260928}{Supplementary proof 53}, The original finite-component construction.
\item[ALGEBRA-RECON-1150] Source lines 40212--40213: \hyperlink{algebra_etale_quasifinite_notes_20260928}{Supplementary proof 53}, The original finite-component construction.
\item[ALGEBRA-RECON-1151] Source lines 40219--40220: \hyperlink{algebra_etale_quasifinite_notes_20260928}{Supplementary proof 53}, The original finite-component construction.
\item[ALGEBRA-RECON-1152] Source lines 40232--40424: \hyperlink{algebra_etale_quasifinite_notes_20260928}{Supplementary proof 53}, The original finite-component construction.
\item[ALGEBRA-RECON-1153] Source lines 40242--40243: \hyperlink{algebra_etale_quasifinite_notes_20260928}{Supplementary proof 53}, Report accounting and propagation.
\item[ALGEBRA-RECON-1154] Source lines 40258--40274: \hyperlink{algebra_etale_quasifinite_notes_20260928}{Supplementary proof 53}, The original finite-component construction.
\item[ALGEBRA-RECON-1155] Source lines 40259--40260: \hyperlink{algebra_etale_quasifinite_notes_20260928}{Supplementary proof 53}, The original finite-component construction.
\item[ALGEBRA-RECON-1156] Source lines 40294--40295: \hyperlink{algebra_etale_quasifinite_notes_20260928}{Supplementary proof 53}, The original finite-component construction.
\item[ALGEBRA-RECON-1157] Source lines 40330--40330: \hyperlink{algebra_etale_quasifinite_notes_20260928}{Supplementary proof 53}, Report accounting and propagation.
\item[ALGEBRA-RECON-1158] Source lines 40366--40367: \hyperlink{algebra_etale_quasifinite_notes_20260928}{Supplementary proof 53}, Report accounting and propagation.
\item[ALGEBRA-RECON-1159] Source lines 40375--40376: \hyperlink{algebra_etale_quasifinite_notes_20260928}{Supplementary proof 53}, Report accounting and propagation.
\item[ALGEBRA-RECON-1160] Source lines 40375--40443: \hyperlink{algebra_etale_quasifinite_notes_20260928}{Supplementary proof 53}, The original finite-component construction.
\item[ALGEBRA-RECON-1161] Source lines 40386--40387: \hyperlink{algebra_etale_quasifinite_notes_20260928}{Supplementary proof 53}, The original finite-component construction.
\item[ALGEBRA-RECON-1162] Source lines 40408--40411: \hyperlink{algebra_etale_quasifinite_notes_20260928}{Supplementary proof 53}, The original finite-component construction.
\item[ALGEBRA-RECON-1163] Source lines 40425--40426: \hyperlink{algebra_etale_quasifinite_notes_20260928}{Supplementary proof 53}, Report accounting and propagation.
\item[ALGEBRA-RECON-1164] Source lines 40434--40435: \hyperlink{algebra_etale_quasifinite_notes_20260928}{Supplementary proof 53}, Report accounting and propagation.
\item[ALGEBRA-RECON-1165] Source lines 40451--40457: \hyperlink{algebra_etale_quasifinite_notes_20260928}{Supplementary proof 53}, Separable splitting, inseparable degree and every fibre length.
\item[ALGEBRA-RECON-1166] Source lines 40451--40461: \hyperlink{algebra_etale_quasifinite_notes_20260928}{Supplementary proof 53}, Separable splitting, inseparable degree and every fibre length.
\item[ALGEBRA-RECON-1167] Source lines 40452--40453: \hyperlink{algebra_etale_quasifinite_notes_20260928}{Supplementary proof 53}, Separable splitting, inseparable degree and every fibre length.
\item[ALGEBRA-RECON-1168] Source lines 40200--40418: \hyperlink{algebra_etale_quasifinite_notes_20260928}{Supplementary proof 53}, Full Artinian fibres and a selected set of components.
\item[ALGEBRA-RECON-1169] Source lines 40420--40489: \hyperlink{algebra_etale_quasifinite_notes_20260928}{Supplementary proof 53}, Separable splitting, inseparable degree and every fibre length.
\end{description}

\subsection{Local homomorphisms}
\hypertarget{algebra-source-section-local-homomorphisms}{}
\noindent\href{https://github.com/stacks/stacks-project/blob/a04446e57ec1fbc252a871afcec7752fb2807b14/algebra.tex\#L40497}{Original source section}. \href{algebra.pdf\#nameddest=section-local-homomorphisms}{Accompanying chapter}.

\begin{description}
\item[ALGEBRA-RECON-1170] Source lines 40511--40516: \hyperlink{algebra_local_integral_notes_20260928}{Supplementary proof 54}, Local homomorphisms and the original parameter.
\item[ALGEBRA-RECON-1171] Source lines 40511--40516: \hyperlink{algebra_local_integral_notes_20260928}{Supplementary proof 54}, Local homomorphisms and the original parameter.
\item[ALGEBRA-RECON-1172] Source lines 40532--40539: \hyperlink{algebra_local_integral_notes_20260928}{Supplementary proof 54}, Local homomorphisms and the original parameter.
\item[ALGEBRA-RECON-1173] Source lines 40556--40567: \hyperlink{algebra_local_integral_notes_20260928}{Supplementary proof 54}, Local homomorphisms and the original parameter.
\item[ALGEBRA-RECON-1179] Source lines 40506--40580: \hyperlink{algebra_local_integral_notes_20260928}{Supplementary proof 54}, Torsion, completion and an exact gluing square.
\end{description}

\subsection{Integral closure and smooth base change}
\hypertarget{algebra-source-section-integral-closure-smooth-base-change}{}
\noindent\href{https://github.com/stacks/stacks-project/blob/a04446e57ec1fbc252a871afcec7752fb2807b14/algebra.tex\#L40586}{Original source section}. \href{algebra.pdf\#nameddest=section-integral-closure-smooth-base-change}{Accompanying chapter}.

\begin{description}
\item[ALGEBRA-RECON-1174] Source lines 40623--40627: \hyperlink{algebra_local_integral_notes_20260928}{Supplementary proof 54}, Report accounting and propagation.
\item[ALGEBRA-RECON-1175] Source lines 40659--40688: \hyperlink{algebra_local_integral_notes_20260928}{Supplementary proof 54}, Integral closure, the derivative and smooth base change.
\item[ALGEBRA-RECON-1176] Source lines 40721--40729: \hyperlink{algebra_local_integral_notes_20260928}{Supplementary proof 54}, Integral closure, the derivative and smooth base change.
\item[ALGEBRA-RECON-1177] Source lines 40734--40745: \hyperlink{algebra_local_integral_notes_20260928}{Supplementary proof 54}, Report accounting and propagation.
\item[ALGEBRA-RECON-1178] Source lines 40748--40761: \hyperlink{algebra_local_integral_notes_20260928}{Supplementary proof 54}, Integral closure, the derivative and smooth base change.
\item[ALGEBRA-RECON-1180] Source lines 40589--40762: \hyperlink{algebra_local_integral_notes_20260928}{Supplementary proof 54}, Integral closure, the derivative and smooth base change.
\item[ALGEBRA-RECON-1181] Source lines 40630--40762: \hyperlink{algebra_local_integral_notes_20260928}{Supplementary proof 54}, Finite conductors and an infinite boundary example.
\end{description}

\subsection{Formally unramified maps}
\hypertarget{algebra-source-section-formally-unramified}{}
\noindent\href{https://github.com/stacks/stacks-project/blob/a04446e57ec1fbc252a871afcec7752fb2807b14/algebra.tex\#L40769}{Original source section}. \href{algebra.pdf\#nameddest=section-formally-unramified}{Accompanying chapter}.

\begin{description}
\item[ALGEBRA-RECON-1182] Source lines 40884--40886: \hyperlink{algebra_conormal_thickening_notes_20260928}{Supplementary proof 55}, Report decisions and propagation.
\item[ALGEBRA-RECON-1183] Source lines 40893--40899: \hyperlink{algebra_conormal_thickening_notes_20260928}{Supplementary proof 55}, Uniqueness, colimits and the nilpotence boundary.
\item[ALGEBRA-RECON-1204] Source lines 40893--41239: \hyperlink{algebra_conormal_thickening_notes_20260928}{Supplementary proof 55}, Uniqueness, colimits and the nilpotence boundary.
\end{description}

\subsection{Conormal modules and universal thickenings}
\hypertarget{algebra-source-section-conormal}{}
\noindent\href{https://github.com/stacks/stacks-project/blob/a04446e57ec1fbc252a871afcec7752fb2807b14/algebra.tex\#L40911}{Original source section}. \href{algebra.pdf\#nameddest=section-conormal}{Accompanying chapter}.

\begin{description}
\item[ALGEBRA-RECON-1184] Source lines 40935--40966: \hyperlink{algebra_conormal_thickening_notes_20260928}{Supplementary proof 55}, The universal thickening and its exact obstruction.
\item[ALGEBRA-RECON-1185] Source lines 40949--40952: \hyperlink{algebra_conormal_thickening_notes_20260928}{Supplementary proof 55}, The universal thickening and its exact obstruction.
\item[ALGEBRA-RECON-1186] Source lines 40984--40995: \hyperlink{algebra_conormal_thickening_notes_20260928}{Supplementary proof 55}, The universal thickening and its exact obstruction.
\item[ALGEBRA-RECON-1187] Source lines 40985--40989: \hyperlink{algebra_conormal_thickening_notes_20260928}{Supplementary proof 55}, The universal thickening and its exact obstruction.
\item[ALGEBRA-RECON-1188] Source lines 40993--40995: \hyperlink{algebra_conormal_thickening_notes_20260928}{Supplementary proof 55}, Report decisions and propagation.
\item[ALGEBRA-RECON-1189] Source lines 41048--41069: \hyperlink{algebra_conormal_thickening_notes_20260928}{Supplementary proof 55}, Report decisions and propagation.
\item[ALGEBRA-RECON-1190] Source lines 41103--41123: \hyperlink{algebra_conormal_thickening_notes_20260928}{Supplementary proof 55}, The differential graph and the original base contribution.
\item[ALGEBRA-RECON-1191] Source lines 41139--41144: \hyperlink{algebra_conormal_thickening_notes_20260928}{Supplementary proof 55}, The differential graph and the original base contribution.
\item[ALGEBRA-RECON-1192] Source lines 41146--41150: \hyperlink{algebra_conormal_thickening_notes_20260928}{Supplementary proof 55}, The differential graph and the original base contribution.
\item[ALGEBRA-RECON-1202] Source lines 41144--41153: \hyperlink{algebra_conormal_thickening_notes_20260928}{Supplementary proof 55}, The differential graph and the original base contribution.
\item[ALGEBRA-RECON-1204] Source lines 40893--41239: \hyperlink{algebra_conormal_thickening_notes_20260928}{Supplementary proof 55}, Uniqueness, colimits and the nilpotence boundary.
\item[ALGEBRA-RECON-1205] Source lines 40919--41154: \hyperlink{algebra_conormal_thickening_notes_20260928}{Supplementary proof 55}, The universal thickening and its exact obstruction.
\end{description}

\subsection{Formally \'etale maps}
\hypertarget{algebra-source-section-formally-etale}{}
\noindent\href{https://github.com/stacks/stacks-project/blob/a04446e57ec1fbc252a871afcec7752fb2807b14/algebra.tex\#L41164}{Original source section}. \href{algebra.pdf\#nameddest=section-formally-etale}{Accompanying chapter}.

\begin{description}
\item[ALGEBRA-RECON-1193] Source lines 41224--41230: \hyperlink{algebra_conormal_thickening_notes_20260928}{Supplementary proof 55}, Uniqueness, colimits and the nilpotence boundary.
\item[ALGEBRA-RECON-1194] Source lines 41232--41238: \hyperlink{algebra_conormal_thickening_notes_20260928}{Supplementary proof 55}, Report decisions and propagation.
\item[ALGEBRA-RECON-1195] Source lines 41266--41284: \hyperlink{algebra_conormal_thickening_notes_20260928}{Supplementary proof 55}, Formal étale jets and differential operators.
\item[ALGEBRA-RECON-1196] Source lines 41290--41303: \hyperlink{algebra_conormal_thickening_notes_20260928}{Supplementary proof 55}, Formal étale jets and differential operators.
\item[ALGEBRA-RECON-1197] Source lines 41319--41331: \hyperlink{algebra_conormal_thickening_notes_20260928}{Supplementary proof 55}, Formal étale jets and differential operators.
\item[ALGEBRA-RECON-1198] Source lines 41323--41327: \hyperlink{algebra_conormal_thickening_notes_20260928}{Supplementary proof 55}, Formal étale jets and differential operators.
\item[ALGEBRA-RECON-1199] Source lines 41333--41352: \hyperlink{algebra_conormal_thickening_notes_20260928}{Supplementary proof 55}, Formal étale jets and differential operators.
\item[ALGEBRA-RECON-1200] Source lines 41333--41350: \hyperlink{algebra_conormal_thickening_notes_20260928}{Supplementary proof 55}, Formal étale jets and differential operators.
\item[ALGEBRA-RECON-1201] Source lines 41351--41358: \hyperlink{algebra_conormal_thickening_notes_20260928}{Supplementary proof 55}, Formal étale jets and differential operators.
\item[ALGEBRA-RECON-1203] Source lines 41360--41371: \hyperlink{algebra_conormal_thickening_notes_20260928}{Supplementary proof 55}, Formal étale jets and differential operators.
\item[ALGEBRA-RECON-1204] Source lines 40893--41239: \hyperlink{algebra_conormal_thickening_notes_20260928}{Supplementary proof 55}, Uniqueness, colimits and the nilpotence boundary.
\item[ALGEBRA-RECON-1206] Source lines 41256--41371: \hyperlink{algebra_conormal_thickening_notes_20260928}{Supplementary proof 55}, Formal étale jets and differential operators.
\end{description}

\subsection{Unramified ring maps}
\hypertarget{algebra-source-section-unramified}{}
\noindent\href{https://github.com/stacks/stacks-project/blob/a04446e57ec1fbc252a871afcec7752fb2807b14/algebra.tex\#L41377}{Original source section}. \href{algebra.pdf\#nameddest=section-unramified}{Accompanying chapter}.

\begin{description}
\item[ALGEBRA-RECON-1207] Source lines 41377--41383: \hyperlink{algebra_unramified_notes_20260928}{Supplementary proof 56}, Source statements and the two reference repairs.
\item[ALGEBRA-RECON-1208] Source lines 41440--41452: \hyperlink{algebra_unramified_notes_20260928}{Supplementary proof 56}, Source statements and the two reference repairs.
\item[ALGEBRA-RECON-1209] Source lines 41475--41477: \hyperlink{algebra_unramified_notes_20260928}{Supplementary proof 56}, Source statements and the two reference repairs.
\item[ALGEBRA-RECON-1210] Source lines 41518--41520: \hyperlink{algebra_unramified_notes_20260928}{Supplementary proof 56}, Source statements and the two reference repairs.
\item[ALGEBRA-RECON-1211] Source lines 41567--41579: \hyperlink{algebra_unramified_notes_20260928}{Supplementary proof 56}, Source statements and the two reference repairs.
\item[ALGEBRA-RECON-1212] Source lines 41653--41665: \hyperlink{algebra_unramified_notes_20260928}{Supplementary proof 56}, Source statements and the two reference repairs.
\item[ALGEBRA-RECON-1213] Source lines 41698--41706: \hyperlink{algebra_unramified_notes_20260928}{Supplementary proof 56}, Source statements and the two reference repairs.
\item[ALGEBRA-RECON-1214] Source lines 41549--41565: \hyperlink{algebra_unramified_notes_20260928}{Supplementary proof 56}, The diagonal, its determinant and equality of maps.
\end{description}

\subsection{Local structure of unramified ring maps}
\hypertarget{algebra-source-section-local-structure-unramified}{}
\noindent\href{https://github.com/stacks/stacks-project/blob/a04446e57ec1fbc252a871afcec7752fb2807b14/algebra.tex\#L41712}{Original source section}. \href{algebra.pdf\#nameddest=section-local-structure-unramified}{Accompanying chapter}.

\begin{description}
\item[ALGEBRA-RECON-1215] Source lines 41753--41765: \hyperlink{algebra_local_unramified_notes_20260929}{Supplementary proof 57}, The source reductions and all scalar factors.
\item[ALGEBRA-RECON-1216] Source lines 41847--41854: \hyperlink{algebra_local_unramified_notes_20260929}{Supplementary proof 57}, The source reductions and all scalar factors.
\item[ALGEBRA-RECON-1217] Source lines 41857--41865: \hyperlink{algebra_local_unramified_notes_20260929}{Supplementary proof 57}, The source reductions and all scalar factors.
\item[ALGEBRA-RECON-1218] Source lines 41861--41865: \hyperlink{algebra_local_unramified_notes_20260929}{Supplementary proof 57}, The source reductions and all scalar factors.
\item[ALGEBRA-RECON-1219] Source lines 41861--41909: \hyperlink{algebra_local_unramified_notes_20260929}{Supplementary proof 57}, The source reductions and all scalar factors.
\item[ALGEBRA-RECON-1220] Source lines 41861--41916: \hyperlink{algebra_local_unramified_notes_20260929}{Supplementary proof 57}, The source reductions and all scalar factors.
\item[ALGEBRA-RECON-1221] Source lines 41866--41880: \hyperlink{algebra_local_unramified_notes_20260929}{Supplementary proof 57}, The source reductions and all scalar factors.
\item[ALGEBRA-RECON-1222] Source lines 41720--42051: \hyperlink{algebra_local_unramified_notes_20260929}{Supplementary proof 57}, The exact conormal object on an etale quotient chart.
\end{description}

\subsection{Henselian local rings}
\hypertarget{algebra-source-section-henselian}{}
\noindent\href{https://github.com/stacks/stacks-project/blob/a04446e57ec1fbc252a871afcec7752fb2807b14/algebra.tex\#L42062}{Original source section}. \href{algebra.pdf\#nameddest=section-henselian}{Accompanying chapter}.

\begin{description}
\item[ALGEBRA-RECON-1223] Source lines 42224--42238: \hyperlink{algebra_henselian_notes_20260929}{Supplementary proof 58}, Taylor formulas and all thirteen characterizations.
\item[ALGEBRA-RECON-1224] Source lines 42241--42258: \hyperlink{algebra_henselian_notes_20260929}{Supplementary proof 58}, Taylor formulas and all thirteen characterizations.
\item[ALGEBRA-RECON-1225] Source lines 42309--42317: \hyperlink{algebra_henselian_notes_20260929}{Supplementary proof 58}, Taylor formulas and all thirteen characterizations.
\item[ALGEBRA-RECON-1226] Source lines 42140--42381: \hyperlink{algebra_henselian_notes_20260929}{Supplementary proof 58}, Taylor formulas and all thirteen characterizations.
\item[ALGEBRA-RECON-1227] Source lines 42460--42470: \hyperlink{algebra_henselian_notes_20260929}{Supplementary proof 58}, Finite algebras, residue functors and arbitrary local targets.
\item[ALGEBRA-RECON-1228] Source lines 42474--42486: \hyperlink{algebra_henselian_notes_20260929}{Supplementary proof 58}, Finite algebras, residue functors and arbitrary local targets.
\item[ALGEBRA-RECON-1229] Source lines 42493--42502: \hyperlink{algebra_henselian_notes_20260929}{Supplementary proof 58}, Finite algebras, residue functors and arbitrary local targets.
\item[ALGEBRA-RECON-1230] Source lines 42545--42736: \hyperlink{algebra_henselian_notes_20260929}{Supplementary proof 58}, Finite algebras, residue functors and arbitrary local targets.
\item[ALGEBRA-RECON-1231] Source lines 42600--42609: \hyperlink{algebra_henselian_notes_20260929}{Supplementary proof 58}, Finite algebras, residue functors and arbitrary local targets.
\item[ALGEBRA-RECON-1232] Source lines 42718--42750: \hyperlink{algebra_henselian_notes_20260929}{Supplementary proof 58}, The original endomorphism and the missing square.
\item[ALGEBRA-RECON-1233] Source lines 42737--42744: \hyperlink{algebra_henselian_notes_20260929}{Supplementary proof 58}, The original endomorphism and the missing square.
\item[ALGEBRA-RECON-1234] Source lines 42718--42760: \hyperlink{algebra_henselian_notes_20260929}{Supplementary proof 58}, The original endomorphism and the missing square.
\item[ALGEBRA-RECON-1235] Source lines 42062--42594: \hyperlink{algebra_henselian_notes_20260929}{Supplementary proof 58}, Complete and nil-ideal lifting with exact errors.
\item[ALGEBRA-RECON-1236] Source lines 42449--42690: \hyperlink{algebra_henselian_notes_20260929}{Supplementary proof 58}, Finite algebras, residue functors and arbitrary local targets.
\end{description}

\subsection{Filtered colimits of \'etale ring maps}
\hypertarget{algebra-source-section-ind-etale}{}
\noindent\href{https://github.com/stacks/stacks-project/blob/a04446e57ec1fbc252a871afcec7752fb2807b14/algebra.tex\#L42761}{Original source section}. \href{algebra.pdf\#nameddest=section-ind-etale}{Accompanying chapter}.

\begin{description}
\item[ALGEBRA-RECON-1237] Source lines 42770--42780: \hyperlink{algebra_indetale_notes_20260929}{Supplementary proof 59}, The source colimits and their exact comparison maps.
\item[ALGEBRA-RECON-1238] Source lines 42791--42811: \hyperlink{algebra_indetale_notes_20260929}{Supplementary proof 59}, The source colimits and their exact comparison maps.
\item[ALGEBRA-RECON-1239] Source lines 42834--42839: \hyperlink{algebra_indetale_notes_20260929}{Supplementary proof 59}, The source colimits and their exact comparison maps.
\item[ALGEBRA-RECON-1240] Source lines 42882--42899: \hyperlink{algebra_indetale_notes_20260929}{Supplementary proof 59}, Residue maps, local targets and filtered local rings.
\item[ALGEBRA-RECON-1241] Source lines 42882--42899: \hyperlink{algebra_indetale_notes_20260929}{Supplementary proof 59}, Residue maps, local targets and filtered local rings.
\item[ALGEBRA-RECON-1242] Source lines 42761--42878: \hyperlink{algebra_indetale_notes_20260929}{Supplementary proof 59}, Arbitrary small colimits and the full idempotent relations.
\item[ALGEBRA-RECON-1243] Source lines 42761--42878: \hyperlink{algebra_indetale_notes_20260929}{Supplementary proof 59}, Nil ideals and the exact reduction equivalence.
\item[ALGEBRA-RECON-1244] Source lines 42880--42954: \hyperlink{algebra_indetale_notes_20260929}{Supplementary proof 59}, Residue maps, local targets and filtered local rings.
\end{description}

\subsection{Henselization and strict henselization}
\hypertarget{algebra-source-section-henselization}{}
\noindent\href{https://github.com/stacks/stacks-project/blob/a04446e57ec1fbc252a871afcec7752fb2807b14/algebra.tex\#L42959}{Original source section}. \href{algebra.pdf\#nameddest=section-henselization}{Accompanying chapter}.

\begin{description}
\item[ALGEBRA-RECON-1245] Source lines 43084--43089: \hyperlink{algebra_henselization_notes_20260929}{Supplementary proof 60}, Local colimits and the source henselian property.
\item[ALGEBRA-RECON-1246] Source lines 43137--43140: \hyperlink{algebra_henselization_notes_20260929}{Supplementary proof 60}, Report dispositions and the deferred follow-on work.
\item[ALGEBRA-RECON-1247] Source lines 43229--43288: \hyperlink{algebra_henselization_notes_20260929}{Supplementary proof 60}, Pointed systems and the original coequalizer.
\item[ALGEBRA-RECON-1248] Source lines 43339--43361: \hyperlink{algebra_henselization_notes_20260929}{Supplementary proof 60}, Residue choices and functoriality.
\item[ALGEBRA-RECON-1249] Source lines 43375--43384: \hyperlink{algebra_henselization_notes_20260929}{Supplementary proof 60}, Residue choices and functoriality.
\item[ALGEBRA-RECON-1250] Source lines 43417--43427: \hyperlink{algebra_henselization_notes_20260929}{Supplementary proof 60}, Pointed systems and the original coequalizer.
\item[ALGEBRA-RECON-1251] Source lines 43434--43449: \hyperlink{algebra_henselization_notes_20260929}{Supplementary proof 60}, Pointed systems and the original coequalizer.
\item[ALGEBRA-RECON-1252] Source lines 43455--43463: \hyperlink{algebra_henselization_notes_20260929}{Supplementary proof 60}, Pointed systems and the original coequalizer.
\item[ALGEBRA-RECON-1253] Source lines 43517--43525: \hyperlink{algebra_henselization_notes_20260929}{Supplementary proof 60}, Residue choices and functoriality.
\end{description}

\subsection{Henselization and quasi-finite ring maps}
\hypertarget{algebra-source-section-henselization-quasi-finite}{}
\noindent\href{https://github.com/stacks/stacks-project/blob/a04446e57ec1fbc252a871afcec7752fb2807b14/algebra.tex\#L43573}{Original source section}. \href{algebra.pdf\#nameddest=section-henselization-quasi-finite}{Accompanying chapter}.

\begin{description}
\item[ALGEBRA-RECON-1254] Source lines 43650--43698: \hyperlink{algebra_quasifinite_henselization_notes_20260929}{Supplementary proof 61}, The original primes and finite local comparisons.
\item[ALGEBRA-RECON-1255] Source lines 43820--43832: \hyperlink{algebra_quasifinite_henselization_notes_20260929}{Supplementary proof 61}, The original filtered systems and the remaining cases.
\end{description}

\subsection{Serre's criterion for normality}
\hypertarget{algebra-source-section-serre-criterion}{}
\noindent\href{https://github.com/stacks/stacks-project/blob/a04446e57ec1fbc252a871afcec7752fb2807b14/algebra.tex\#L43865}{Original source section}. \href{algebra.pdf\#nameddest=section-serre-criterion}{Accompanying chapter}.

\begin{description}
\item[ALGEBRA-RECON-1256] Source lines 44031--44042: \hyperlink{algebra_serre_notes_20260929}{Supplementary proof 62}, The original Serre argument and its zero quotient.
\item[ALGEBRA-RECON-1257] Source lines 44063--44066: \hyperlink{algebra_serre_notes_20260929}{Supplementary proof 62}, Report accounting and propagation.
\item[ALGEBRA-RECON-1258] Source lines 44069--44110: \hyperlink{algebra_serre_notes_20260929}{Supplementary proof 62}, Height-one intersections with all original factors retained.
\item[ALGEBRA-RECON-1259] Source lines 44077--44079: \hyperlink{algebra_serre_notes_20260929}{Supplementary proof 62}, Height-one intersections with all original factors retained.
\end{description}

\subsection{Formal smoothness of fields}
\hypertarget{algebra-source-section-p-bases}{}
\noindent\href{https://github.com/stacks/stacks-project/blob/a04446e57ec1fbc252a871afcec7752fb2807b14/algebra.tex\#L44127}{Original source section}. \href{algebra.pdf\#nameddest=section-p-bases}{Accompanying chapter}.

\begin{description}
\item[ALGEBRA-RECON-1260] Source lines 44257--44260: \hyperlink{algebra_field_formal_smooth_notes_20260929}{Supplementary proof 63}, The corrected coefficient field and both differential maps.
\item[ALGEBRA-RECON-1261] Source lines 44282--44300: \hyperlink{algebra_field_formal_smooth_notes_20260929}{Supplementary proof 63}, The corrected coefficient field and both differential maps.
\item[ALGEBRA-RECON-1262] Source lines 44293--44318: \hyperlink{algebra_field_formal_smooth_notes_20260929}{Supplementary proof 63}, The corrected coefficient field and both differential maps.
\item[ALGEBRA-RECON-1263] Source lines 44308--44325: \hyperlink{algebra_field_formal_smooth_notes_20260929}{Supplementary proof 63}, The corrected coefficient field and both differential maps.
\item[ALGEBRA-RECON-1264] Source lines 44383--44393: \hyperlink{algebra_field_formal_smooth_notes_20260929}{Supplementary proof 63}, Formal lifting and the characteristic-zero assertions.
\item[ALGEBRA-RECON-1265] Source lines 44396--44406: \hyperlink{algebra_field_formal_smooth_notes_20260929}{Supplementary proof 63}, Formal lifting and the characteristic-zero assertions.
\item[ALGEBRA-RECON-1266] Source lines 44435--44440: \hyperlink{algebra_field_formal_smooth_notes_20260929}{Supplementary proof 63}, Formal lifting and the characteristic-zero assertions.
\item[ALGEBRA-RECON-1267] Source lines 44474--44486: \hyperlink{algebra_field_formal_smooth_notes_20260929}{Supplementary proof 63}, Formal lifting and the characteristic-zero assertions.
\item[ALGEBRA-RECON-1268] Source lines 44499--44504: \hyperlink{algebra_field_formal_smooth_notes_20260929}{Supplementary proof 63}, Formal lifting and the characteristic-zero assertions.
\item[ALGEBRA-RECON-1269] Source lines 44544--44548: \hyperlink{algebra_field_formal_smooth_notes_20260929}{Supplementary proof 63}, The original triangular complete-intersection algebras.
\item[ALGEBRA-RECON-1270] Source lines 44540--44551: \hyperlink{algebra_field_formal_smooth_notes_20260929}{Supplementary proof 63}, The original triangular complete-intersection algebras.
\end{description}

\subsection{Constructing flat ring maps}
\hypertarget{algebra-source-section-constructing-flat}{}
\noindent\href{https://github.com/stacks/stacks-project/blob/a04446e57ec1fbc252a871afcec7752fb2807b14/algebra.tex\#L44566}{Original source section}. \href{algebra.pdf\#nameddest=section-constructing-flat}{Accompanying chapter}.

\begin{description}
\item[ALGEBRA-RECON-1271] Source lines 44590--44617: \hyperlink{algebra_construct_flat_notes_20260929}{Supplementary proof 64}, The fixed base and the directed local colimit.
\item[ALGEBRA-RECON-1272] Source lines 44595--44603: \hyperlink{algebra_construct_flat_notes_20260929}{Supplementary proof 64}, The fixed base and the directed local colimit.
\item[ALGEBRA-RECON-1273] Source lines 44620--44622: \hyperlink{algebra_construct_flat_notes_20260929}{Supplementary proof 64}, Report accounting and propagation.
\item[ALGEBRA-RECON-1274] Source lines 44625--44629: \hyperlink{algebra_construct_flat_notes_20260929}{Supplementary proof 64}, The original transfinite construction and its endpoints.
\item[ALGEBRA-RECON-1275] Source lines 44625--44638: \hyperlink{algebra_construct_flat_notes_20260929}{Supplementary proof 64}, The original transfinite construction and its endpoints.
\item[ALGEBRA-RECON-1276] Source lines 44633--44638: \hyperlink{algebra_construct_flat_notes_20260929}{Supplementary proof 64}, The original transfinite construction and its endpoints.
\item[ALGEBRA-RECON-1277] Source lines 44624--44656: \hyperlink{algebra_construct_flat_notes_20260929}{Supplementary proof 64}, The original transfinite construction and its endpoints.
\item[ALGEBRA-RECON-1278] Source lines 44661--44671: \hyperlink{algebra_construct_flat_notes_20260929}{Supplementary proof 64}, Finite étale local stages and descent of locality.
\item[ALGEBRA-RECON-1279] Source lines 44667--44672: \hyperlink{algebra_construct_flat_notes_20260929}{Supplementary proof 64}, Finite étale local stages and descent of locality.
\item[ALGEBRA-RECON-1280] Source lines 44683--44686: \hyperlink{algebra_construct_flat_notes_20260929}{Supplementary proof 64}, Finite free realization with the original minimal polynomial.
\item[ALGEBRA-RECON-1281] Source lines 44687--44692: \hyperlink{algebra_construct_flat_notes_20260929}{Supplementary proof 64}, Finite free realization with the original minimal polynomial.
\item[ALGEBRA-RECON-1282] Source lines 44742--44749: \hyperlink{algebra_construct_flat_notes_20260929}{Supplementary proof 64}, The cardinal bound and the actual relation witnesses.
\item[ALGEBRA-RECON-1283] Source lines 44755--44758: \hyperlink{algebra_construct_flat_notes_20260929}{Supplementary proof 64}, The cardinal bound and the actual relation witnesses.
\end{description}

\subsection{The Cohen structure theorem}
\hypertarget{algebra-source-section-cohen-structure-theorem}{}
\noindent\href{https://github.com/stacks/stacks-project/blob/a04446e57ec1fbc252a871afcec7752fb2807b14/algebra.tex\#L44779}{Original source section}. \href{algebra.pdf\#nameddest=section-cohen-structure-theorem}{Accompanying chapter}.

\begin{description}
\item[ALGEBRA-RECON-1284] Source lines 44799--44802: \hyperlink{algebra_cohen_notes_20260929}{Supplementary proof 65}, Completeness, proper quotients and finite products.
\item[ALGEBRA-RECON-1285] Source lines 44813--44813: \hyperlink{algebra_cohen_notes_20260929}{Supplementary proof 65}, Completeness, proper quotients and finite products.
\item[ALGEBRA-RECON-1286] Source lines 44811--44815: \hyperlink{algebra_cohen_notes_20260929}{Supplementary proof 65}, Completeness, proper quotients and finite products.
\item[ALGEBRA-RECON-1287] Source lines 44873--44874: \hyperlink{algebra_cohen_notes_20260929}{Supplementary proof 65}, The original Cohen ring and all its finite quotients.
\item[ALGEBRA-RECON-1288] Source lines 44985--44995: \hyperlink{algebra_cohen_notes_20260929}{Supplementary proof 65}, The compatible coefficient-ring lifting maps.
\item[ALGEBRA-RECON-1289] Source lines 45016--45019: \hyperlink{algebra_cohen_notes_20260929}{Supplementary proof 65}, Power-series surjectivity with the respective ideals.
\item[ALGEBRA-RECON-1290] Source lines 45065--45068: \hyperlink{algebra_cohen_notes_20260929}{Supplementary proof 65}, Regular power-series rings and the two finite-extension cases.
\item[ALGEBRA-RECON-1291] Source lines 45098--45102: \hyperlink{algebra_cohen_notes_20260929}{Supplementary proof 65}, Regular power-series rings and the two finite-extension cases.
\end{description}

\subsection{Japanese rings}
\hypertarget{algebra-source-section-japanese}{}
\noindent\href{https://github.com/stacks/stacks-project/blob/a04446e57ec1fbc252a871afcec7752fb2807b14/algebra.tex\#L45128}{Original source section}. \href{algebra.pdf\#nameddest=section-japanese}{Accompanying chapter}.

\begin{description}
\item[ALGEBRA-RECON-1292] Source lines 45172--45181: \hyperlink{algebra_japanese_notes_20260929}{Supplementary proof 66}, Integral closures, localization and the zero-ring boundary.
\item[ALGEBRA-RECON-1293] Source lines 45310--45313: \hyperlink{algebra_japanese_notes_20260929}{Supplementary proof 66}, The trace lattice and the original involution example.
\item[ALGEBRA-RECON-1294] Source lines 45325--45326: \hyperlink{algebra_japanese_notes_20260929}{Supplementary proof 66}, The derivation and every original denominator factor.
\item[ALGEBRA-RECON-1295] Source lines 45364--45366: \hyperlink{algebra_japanese_notes_20260929}{Supplementary proof 66}, The derivation and every original denominator factor.
\item[ALGEBRA-RECON-1296] Source lines 45429--45433: \hyperlink{algebra_japanese_notes_20260929}{Supplementary proof 66}, Purely inseparable reduction and polynomial N-2.
\item[ALGEBRA-RECON-1297] Source lines 45434--45437: \hyperlink{algebra_japanese_notes_20260929}{Supplementary proof 66}, Purely inseparable reduction and polynomial N-2.
\item[ALGEBRA-RECON-1298] Source lines 45458--45464: \hyperlink{algebra_japanese_notes_20260929}{Supplementary proof 66}, Normal loci and the exact equality of local rings.
\item[ALGEBRA-RECON-1299] Source lines 45505--45510: \hyperlink{algebra_japanese_notes_20260929}{Supplementary proof 66}, Normal loci and the exact equality of local rings.
\item[ALGEBRA-RECON-1300] Source lines 45523--45529: \hyperlink{algebra_japanese_notes_20260929}{Supplementary proof 66}, Normal loci and the exact equality of local rings.
\item[ALGEBRA-RECON-1301] Source lines 45561--45563: \hyperlink{algebra_japanese_notes_20260929}{Supplementary proof 66}, Tate finiteness with the original extension retained.
\end{description}

\subsection{Nagata rings}
\hypertarget{algebra-source-section-nagata}{}
\noindent\href{https://github.com/stacks/stacks-project/blob/a04446e57ec1fbc252a871afcec7752fb2807b14/algebra.tex\#L45615}{Original source section}. \href{algebra.pdf\#nameddest=section-nagata}{Accompanying chapter}.

\begin{description}
\item[ALGEBRA-RECON-1302] Source lines 45739--45751: \hyperlink{algebra_nagata_notes_20260929}{Supplementary proof 67}, Reduced algebras and nonzero localization charts.
\item[ALGEBRA-RECON-1303] Source lines 45872--45872: \hyperlink{algebra_nagata_notes_20260929}{Supplementary proof 67}, Complete rings and integral closure in the total quotient ring.
\item[ALGEBRA-RECON-1304] Source lines 45877--45877: \hyperlink{algebra_nagata_notes_20260929}{Supplementary proof 67}, Complete rings and integral closure in the total quotient ring.
\item[ALGEBRA-RECON-1305] Source lines 46019--46021: \hyperlink{algebra_nagata_notes_20260929}{Supplementary proof 67}, Local Nagata domains and the finite local test.
\item[ALGEBRA-RECON-1306] Source lines 46025--46027: \hyperlink{algebra_nagata_notes_20260929}{Supplementary proof 67}, Local Nagata domains and the finite local test.
\item[ALGEBRA-RECON-1307] Source lines 46040--46043: \hyperlink{algebra_nagata_notes_20260929}{Supplementary proof 67}, Local Nagata domains and the finite local test.
\item[ALGEBRA-RECON-1308] Source lines 46042--46051: \hyperlink{algebra_nagata_notes_20260929}{Supplementary proof 67}, Local Nagata domains and the finite local test.
\item[ALGEBRA-RECON-1309] Source lines 46052--46052: \hyperlink{algebra_nagata_notes_20260929}{Supplementary proof 67}, Local Nagata domains and the finite local test.
\item[ALGEBRA-RECON-1310] Source lines 46053--46054: \hyperlink{algebra_nagata_notes_20260929}{Supplementary proof 67}, Local Nagata domains and the finite local test.
\item[ALGEBRA-RECON-1311] Source lines 46098--46107: \hyperlink{algebra_nagata_notes_20260929}{Supplementary proof 67}, Finite-type stability with every original ring retained.
\item[ALGEBRA-RECON-1312] Source lines 46110--46128: \hyperlink{algebra_nagata_notes_20260929}{Supplementary proof 67}, Finite-type stability with every original ring retained.
\item[ALGEBRA-RECON-1313] Source lines 46142--46144: \hyperlink{algebra_nagata_notes_20260929}{Supplementary proof 67}, Finite-type stability with every original ring retained.
\item[ALGEBRA-RECON-1314] Source lines 46143--46145: \hyperlink{algebra_nagata_notes_20260929}{Supplementary proof 67}, Finite-type stability with every original ring retained.
\item[ALGEBRA-RECON-1315] Source lines 46143--46161: \hyperlink{algebra_nagata_notes_20260929}{Supplementary proof 67}, Finite-type stability with every original ring retained.
\item[ALGEBRA-RECON-1316] Source lines 46152--46161: \hyperlink{algebra_nagata_notes_20260929}{Supplementary proof 67}, Finite-type stability with every original ring retained.
\item[ALGEBRA-RECON-1317] Source lines 46167--46181: \hyperlink{algebra_nagata_notes_20260929}{Supplementary proof 67}, Finite-type stability with every original ring retained.
\item[ALGEBRA-RECON-1318] Source lines 46185--46185: \hyperlink{algebra_nagata_notes_20260929}{Supplementary proof 67}, The linear kernel, fractional ideal and local DVR map.
\item[ALGEBRA-RECON-1319] Source lines 46280--46285: \hyperlink{algebra_nagata_notes_20260929}{Supplementary proof 67}, The original non-Japanese DVR and the completion root test.
\end{description}

\subsection{Ascending properties}
\hypertarget{algebra-source-section-ascending-properties}{}
\noindent\href{https://github.com/stacks/stacks-project/blob/a04446e57ec1fbc252a871afcec7752fb2807b14/algebra.tex\#L46293}{Original source section}. \href{algebra.pdf\#nameddest=section-ascending-properties}{Accompanying chapter}.

\begin{description}
\item[ALGEBRA-RECON-1320] Source lines 46321--46321: \hyperlink{algebra_ascent_notes_20260929}{Supplementary proof 68}, The module depth formula, including zero modules.
\item[ALGEBRA-RECON-1321] Source lines 46314--46371: \hyperlink{algebra_ascent_notes_20260929}{Supplementary proof 68}, The module depth formula, including zero modules.
\item[ALGEBRA-RECON-1322] Source lines 46388--46400: \hyperlink{algebra_ascent_notes_20260929}{Supplementary proof 68}, Cohen–Macaulay ascent and the original local fibre maps.
\item[ALGEBRA-RECON-1323] Source lines 46459--46508: \hyperlink{algebra_ascent_notes_20260929}{Supplementary proof 68}, Serre conditions, reducedness and normality in the Noetherian case.
\item[ALGEBRA-RECON-1324] Source lines 46528--46537: \hyperlink{algebra_ascent_notes_20260929}{Supplementary proof 68}, Smooth reduced ascent with the finite-stage witnesses retained.
\item[ALGEBRA-RECON-1325] Source lines 46609--46626: \hyperlink{algebra_ascent_notes_20260929}{Supplementary proof 68}, Regular smooth ascent and report accounting.
\end{description}

\subsection{Descending properties}
\hypertarget{algebra-source-section-descending-properties}{}
\noindent\href{https://github.com/stacks/stacks-project/blob/a04446e57ec1fbc252a871afcec7752fb2807b14/algebra.tex\#L46633}{Original source section}. \href{algebra.pdf\#nameddest=section-descending-properties}{Accompanying chapter}.

\begin{description}
\item[ALGEBRA-RECON-1326] Source lines 46637--46665: \hyperlink{algebra_descent_notes_20260929}{Supplementary proof 69}, Faithful contraction, Noetherianity and the regularity receiver.
\item[ALGEBRA-RECON-1327] Source lines 46683--46712: \hyperlink{algebra_descent_notes_20260929}{Supplementary proof 69}, Normal descent with the actual principal-ideal map.
\item[ALGEBRA-RECON-1328] Source lines 46803--46825: \hyperlink{algebra_descent_notes_20260929}{Supplementary proof 69}, Nagata descent and every total-quotient denominator.
\item[ALGEBRA-RECON-1329] Source lines 46826--46840: \hyperlink{algebra_descent_notes_20260929}{Supplementary proof 69}, Nagata descent and every total-quotient denominator.
\item[ALGEBRA-RECON-1330] Source lines 46839--46844: \hyperlink{algebra_descent_notes_20260929}{Supplementary proof 69}, Nagata descent and every total-quotient denominator.
\item[ALGEBRA-RECON-1331] Source lines 46847--46868: \hyperlink{algebra_descent_notes_20260929}{Supplementary proof 69}, The conductor ring and a finite étale counterexample.
\item[ALGEBRA-RECON-1332] Source lines 46860--46867: \hyperlink{algebra_descent_notes_20260929}{Supplementary proof 69}, The conductor ring and a finite étale counterexample.
\item[ALGEBRA-RECON-1333] Source lines 46851--46859: \hyperlink{algebra_descent_notes_20260929}{Supplementary proof 69}, The original series construction and its missing coefficient factor.
\end{description}

\subsection{Geometrically normal algebras}
\hypertarget{algebra-source-section-geometrically-normal}{}
\noindent\href{https://github.com/stacks/stacks-project/blob/a04446e57ec1fbc252a871afcec7752fb2807b14/algebra.tex\#L46878}{Original source section}. \href{algebra.pdf\#nameddest=section-geometrically-normal}{Accompanying chapter}.

\begin{description}
\item[ALGEBRA-RECON-1334] Source lines 47000--47008: \hyperlink{algebra_geonormal_notes_20260929}{Supplementary proof 70}, The original local reduction and normal finite subalgebras.
\item[ALGEBRA-RECON-1335] Source lines 47011--47024: \hyperlink{algebra_geonormal_notes_20260929}{Supplementary proof 70}, The original local reduction and normal finite subalgebras.
\item[ALGEBRA-RECON-1336] Source lines 47012--47020: \hyperlink{algebra_geonormal_notes_20260929}{Supplementary proof 70}, The original local reduction and normal finite subalgebras.
\item[ALGEBRA-RECON-1337] Source lines 47023--47024: \hyperlink{algebra_geonormal_notes_20260929}{Supplementary proof 70}, Report accounting and propagation.
\item[ALGEBRA-RECON-1338] Source lines 47037--47059: \hyperlink{algebra_geonormal_notes_20260929}{Supplementary proof 70}, Both tensor localizations inside the total quotient ring.
\end{description}

\subsection{Geometrically regular algebras}
\hypertarget{algebra-source-section-geometrically-regular}{}
\noindent\href{https://github.com/stacks/stacks-project/blob/a04446e57ec1fbc252a871afcec7752fb2807b14/algebra.tex\#L47092}{Original source section}. \href{algebra.pdf\#nameddest=section-geometrically-regular}{Accompanying chapter}.

\begin{description}
\item[ALGEBRA-RECON-1339] Source lines 47092--47318: \hyperlink{algebra_georeg_cm_notes_20260929}{Supplementary proof 71}, Finite purely inseparable descent to an actual subfield.
\end{description}

\subsection{Geometrically Cohen-Macaulay algebras}
\hypertarget{algebra-source-section-geometrically-CM}{}
\noindent\href{https://github.com/stacks/stacks-project/blob/a04446e57ec1fbc252a871afcec7752fb2807b14/algebra.tex\#L47255}{Original source section}. \href{algebra.pdf\#nameddest=section-geometrically-CM}{Accompanying chapter}.

\begin{description}
\item[ALGEBRA-RECON-1339] Source lines 47092--47318: \hyperlink{algebra_georeg_cm_notes_20260929}{Supplementary proof 71}, Finite purely inseparable descent to an actual subfield.
\end{description}

\subsection{Colimits and maps of finite presentation, II}
\hypertarget{algebra-source-section-colimits-finite-presentation}{}
\noindent\href{https://github.com/stacks/stacks-project/blob/a04446e57ec1fbc252a871afcec7752fb2807b14/algebra.tex\#L47319}{Original source section}. \href{algebra.pdf\#nameddest=section-colimits-finite-presentation}{Accompanying chapter}.

\begin{description}
\item[ALGEBRA-RECON-1340] Source lines 47330--47369: \hyperlink{algebra_fp_flat_notes_20260929}{Supplementary proof 72}, Original presentations and their localized maps.
\item[ALGEBRA-RECON-1341] Source lines 47394--47404: \hyperlink{algebra_fp_flat_notes_20260929}{Supplementary proof 72}, The entire flat locus and its direction under base change.
\item[ALGEBRA-RECON-1342] Source lines 47397--47404: \hyperlink{algebra_fp_flat_notes_20260929}{Supplementary proof 72}, The entire flat locus and its direction under base change.
\item[ALGEBRA-RECON-1343] Source lines 47407--47417: \hyperlink{algebra_fp_flat_notes_20260929}{Supplementary proof 72}, Every coefficient of the finite unit relation.
\item[ALGEBRA-RECON-1344] Source lines 47469--47485: \hyperlink{algebra_fp_flat_notes_20260929}{Supplementary proof 72}, Faithful flatness, spectral images and a finite relation algebra.
\item[ALGEBRA-RECON-1345] Source lines 47487--47519: \hyperlink{algebra_fp_flat_notes_20260929}{Supplementary proof 72}, Faithful flatness, spectral images and a finite relation algebra.
\item[ALGEBRA-RECON-1346] Source lines 47521--47769: \hyperlink{algebra_fp_end_notes_20260929}{Supplementary proof 73}, Original stage objects and the first four properties.
\item[ALGEBRA-RECON-1347] Source lines 47849--47851: \hyperlink{algebra_fp_end_notes_20260929}{Supplementary proof 73}, The redundant localization and the complete denominator.
\item[ALGEBRA-RECON-1348] Source lines 47836--47858: \hyperlink{algebra_fp_end_notes_20260929}{Supplementary proof 73}, The redundant localization and the complete denominator.
\item[ALGEBRA-RECON-1349] Source lines 47854--47874: \hyperlink{algebra_fp_end_notes_20260929}{Supplementary proof 73}, The redundant localization and the complete denominator.
\end{description}


\section{Contextual editorial arguments}

Each passage below is an explicitly identified editorial treatment. The original statement and its exact source location remain available in the source records. The retained context supplies the definitions, statement and surrounding proof needed to read the argument. Source corrections, completed proofs and stronger conclusions keep their individual decisions in the accompanying ledger. No novelty or official endorsement is claimed.

\subsection{Editorial treatment of Algebra lines 1000--1031}

\hypertarget{algebra-context-001}{}

\noindent\href{https://github.com/stacks/stacks-project/blob/a04446e57ec1fbc252a871afcec7752fb2807b14/algebra.tex\#L1000}{Original source, lines 1000--1031}. Complete original and reviewed TeX passages accompany this note. Every intervention is separately recorded; the following treatment is editorial.


\begin{lemma}
\label{editorial-source-lemma-almost-directed-colimit-exact}
Let $R$ be a ring. Let $\mathcal{I}$ be an index category satisfying the assumptions of
Categories, Lemma \ref{categories-lemma-split-into-directed}.
Then taking colimits of diagrams of $R$-modules over $\mathcal{I}$
is exact (i.e., the analogue of
Lemma \ref{algebra-lemma-directed-colimit-exact}
holds in this situation).
\end{lemma}

\begin{proof}
By
Categories, Lemma \ref{categories-lemma-split-into-directed}
we may write $\mathcal{I} = \coprod_{j \in J} \mathcal{I}_j$ with each
$\mathcal{I}_j$ a filtered category, and $J$ possibly empty. By
Categories, Lemma \ref{categories-lemma-directed-category-system}
taking colimits over the index categories $\mathcal{I}_j$ is
the same as taking the colimit over some directed set. Hence
Lemma \ref{algebra-lemma-directed-colimit-exact}
applies to these colimits. This reduces the problem to showing that
coproducts in the category of $R$-modules over the set $J$ are exact.
In other words, given exact sequences
$L_j \to M_j \to N_j$ of $R$-modules, we have to show that
$$
\bigoplus\nolimits_{j \in J} L_j
\longrightarrow
\bigoplus\nolimits_{j \in J} M_j
\longrightarrow
\bigoplus\nolimits_{j \in J} N_j
$$
is exact. Indeed, denote its component maps by
$\alpha_j : L_j \to M_j$ and $\beta_j : M_j \to N_j$.
An element $(m_j)_{j \in J}$ of $\bigoplus_{j \in J} M_j$
has finite support. It is in the kernel of $\bigoplus_j\beta_j$
if and only if every $m_j$ is in $\Ker(\beta_j)$.
For each index in its finite support, exactness gives
$l_j \in L_j$ such that $\alpha_j(l_j) = m_j$.
Set $l_j = 0$ at every other index. Then $(l_j)_{j \in J}$
has finite support and maps to $(m_j)_{j \in J}$.
Conversely, $\beta_j\alpha_j = 0$ at every index, so the image
of $\bigoplus_j\alpha_j$ is contained in this kernel.

\medskip\noindent
More generally, for complexes $L_j \xrightarrow{\alpha_j} M_j
\xrightarrow{\beta_j} N_j$, the canonical $R$-linear map
$$
\bigoplus_{j \in J}
\bigl(\Ker(\beta_j)/\Im(\alpha_j)\bigr)
\longrightarrow
\Ker\bigl(\bigoplus_j\beta_j\bigr)
/\Im\bigl(\bigoplus_j\alpha_j\bigr),
\qquad ([m_j])_j \longmapsto [(m_j)_j]
$$
is an isomorphism. Choose representatives at the finitely many
nonzero components and zero representatives elsewhere. Changing these
choices changes $(m_j)_j$ by the image of a finitely supported
family in $\bigoplus_j L_j$, so the map is well defined.
Every element of the target has a finitely supported representative
whose components lie in $\Ker(\beta_j)$; this proves surjectivity.
If a representative lies in $\Im(\bigoplus_j\alpha_j)$, each
$m_j$ lies in $\Im(\alpha_j)$, proving injectivity.
All maps used are $R$-linear. Applying this to the colimits over
the components $\mathcal{I}_j$ proves the homology assertion as well.
If $J$ is empty, all the displayed direct sums and homology modules
are zero, and the same assertions hold.
\end{proof}


\subsection{Editorial treatment of Algebra lines 1301--1324}

\hypertarget{algebra-context-002}{}

\noindent\href{https://github.com/stacks/stacks-project/blob/a04446e57ec1fbc252a871afcec7752fb2807b14/algebra.tex\#L1301}{Original source, lines 1301--1324}. Complete original and reviewed TeX passages accompany this note. Every intervention is separately recorded; the following treatment is editorial.


\begin{proposition}
\label{editorial-source-proposition-localize-twice-module}
View $S'^{-1}M$ as an $A$-module, then $S^{-1}(S'^{-1}M)$ is
isomorphic to $(SS')^{-1}M$.
\end{proposition}

\begin{proof}
Note that given an $A$-module $M$, the universal property of
Lemma \ref{algebra-lemma-universal-property-localization-module} is available.
We give an explicit construction of the isomorphism.

\medskip\noindent
We define the maps as follows
\begin{align*}
& f : S^{-1}(S'^{-1}M) \longrightarrow (SS')^{-1}M, \quad \frac{x/s'}{s}\mapsto
x/ss'\\
& g : (SS')^{-1}M \longrightarrow S^{-1}(S'^{-1}M), \quad x/t\mapsto
\frac{x/s'}{s}\ \text{for some }s\in S, s'\in S', \text{ and }
t = ss'
\end{align*}
We have to check that these homomorphisms are well-defined, that is,
independent of the choice of the fraction. The universal property proves that these maps are well defined:
every element of $S'$ and of $S$ acts invertibly on $(SS')^{-1}M$,
so the map $M \to (SS')^{-1}M$ extends first to $S'^{-1}M$
and then to $S^{-1}(S'^{-1}M)$, giving $f$ with the displayed formula.
On $S^{-1}(S'^{-1}M)$ every $s' \in S'$ acts invertibly already
on the inner localization, and every $s \in S$ acts invertibly
on the outer localization. Thus every $ss' \in SS'$ acts invertibly,
and $M \to S^{-1}(S'^{-1}M)$ extends uniquely to $g$.
The displayed formula for $g$ follows by inverting the actions of
$s'$ and $s$; it therefore does not depend on the chosen factorization
of the denominator. Both composites restrict to the original
localization map on $M$. Uniqueness in the universal property,
applied once for $(SS')^{-1}M$ and successively for the two
localizations, gives $f \circ g = \mathrm{id}_{(SS')^{-1}M}$ and
$g \circ f = \mathrm{id}_{S^{-1}(S'^{-1}M)}$.
\end{proof}


\subsection{Editorial treatment of Algebra lines 1420--1433}

\hypertarget{algebra-context-003}{}

\noindent\href{https://github.com/stacks/stacks-project/blob/a04446e57ec1fbc252a871afcec7752fb2807b14/algebra.tex\#L1420}{Original source, lines 1420--1433}. Complete original and reviewed TeX passages accompany this note. Every intervention is separately recorded; the following treatment is editorial.


\begin{lemma}
\label{editorial-source-lemma-submodule-localization}
Any submodule $N'$ of $S^{-1}M$ is of the form $S^{-1}N$ for some
$N\subset M$. Indeed one can take $N$ to be the inverse image of
$N'$ in $M$.
\end{lemma}

\begin{proof}
Let $N$ be the inverse image of $N'$ in $M$. Then $S^{-1}N\supset N'$:
if $x/s \in N'$, then $x/1 = s(x/s) \in N'$, so $x \in N$
and hence $x/s \in S^{-1}N$. To show they are equal, take $x/s$ in
$S^{-1}N$, where $s\in S$ and $x\in N$. This yields that $x/1\in
N'$. Since $N'$ is an $S^{-1}A$-submodule we have
$x/s = x/1\cdot 1/s\in N'$. This finishes the proof.
\end{proof}


\subsection{Editorial treatment of Algebra lines 1588--1614}

\hypertarget{algebra-context-004}{}

\noindent\href{https://github.com/stacks/stacks-project/blob/a04446e57ec1fbc252a871afcec7752fb2807b14/algebra.tex\#L1588}{Original source, lines 1588--1614}. Complete original and reviewed TeX passages accompany this note. Every intervention is separately recorded; the following treatment is editorial.


\begin{lemma}
\label{editorial-source-lemma-characterize-finite-module-hom}
Let $R$ be a ring. Let $N$ be an $R$-module. The following are equivalent
\begin{enumerate}
\item $N$ is a finite $R$-module,
\item for any filtered colimit $M = \colim M_i$ of $R$-modules the map
$\colim \Hom_R(N, M_i) \to \Hom_R(N, M)$ is injective.
\end{enumerate}
\end{lemma}

\begin{proof}
Assume (1) and choose generators $x_1, \ldots, x_m$ for $N$.
If $N \to M_i$ is a module map and the composition
$N \to M_i \to M$ is zero, then because $M = \colim_{i' \geq i} M_{i'}$
for each $j \in \{1, \ldots, m\}$ we can find an $i' \geq i$ such that
$x_j$ maps to zero in $M_{i'}$. Since there are finitely many
$x_j$ we can find a single $i'$ which works for all of them.
Then the composition $N \to M_i \to M_{i'}$ is zero and we conclude
the map is injective, i.e., part (2) holds.

\medskip\noindent
Assume (2). For a finite subset $E \subset N$ denote $N_E \subset N$
the $R$-submodule generated by the elements of $E$. Then
$0 = \colim_E N/N_E$ is a filtered colimit. At the stage
$E = \emptyset$ we have $N/N_E = N$, and the class of
$\mathrm{id}_N$ in $\colim_E\Hom_R(N, N/N_E)$ maps to zero in
$\Hom_R(N, 0)$. By (2) this class is zero. Thus, at some finite
stage $E$, its image is the zero map. That image is the quotient
map $q_E : N \to N/N_E$, so $q_E = 0$ and $N = N_E$.
Equivalently, $\mathrm{id}_N$ factors through the inclusion
$N_E \hookrightarrow N$. Hence $N$ is finitely generated.
\end{proof}


\subsection{Editorial treatment of Algebra lines 1715--1776}

\hypertarget{algebra-context-005}{}

\noindent\href{https://github.com/stacks/stacks-project/blob/a04446e57ec1fbc252a871afcec7752fb2807b14/algebra.tex\#L1715}{Original source, lines 1715--1776}. Complete original and reviewed TeX passages accompany this note. Every intervention is separately recorded; the following treatment is editorial.


\begin{lemma}
\label{editorial-source-lemma-tensor-product}
Let $M, N$ be $R$-modules. Then there exists a pair $(T, g)$
where $T$ is an $R$-module, and
$g : M \times N \to T$ an $R$-bilinear
mapping, with the following universal property:
For any $R$-module $P$ and any $R$-bilinear mapping
$f : M \times N \to P$, there
exists a unique $R$-linear
mapping $\tilde{f} : T \to P$ such that $f = \tilde{f} \circ g$.
In other words, the following diagram commutes:
$$
\xymatrix{
M \times N \ar[rr]^f \ar[dr]_g & & P\\
& T \ar[ur]_{\tilde f}
}
$$
Moreover, if $(T, g)$ and $(T', g')$
are two pairs with this property, then there
exists a unique isomorphism
$j : T \to T'$ such that $j\circ g = g'$.
\end{lemma}

\noindent
The $R$-module $T$ which satisfies the above universal property is called
the  {\it tensor product} of $R$-modules $M$ and $N$, denoted as
$M \otimes_R N$.

\begin{proof}
We first prove the existence of such $R$-module $T$.
Let $M, N$ be $R$-modules.
Let $T$ be the quotient module
$P/Q$, where $P$ is the free $R$-module $R^{(M \times N)}$ and $Q$ is the
$R$-module generated by all elements of
the following types: ($a\in R$, $x,x'\in M$, $y,y'\in N$)
\begin{align*}
(x + x', y) - (x, y) - (x', y), \\
(x, y + y') - (x, y) - (x, y'), \\
(ax, y) - a(x, y), \\
(x, ay) - a(x, y)
\end{align*}
Let $\pi : M \times N \to T$ denote the natural map.
This map is $R$-bilinear, as
implied by the above relations
when we check the bilinearity conditions. Denote the image by
$\pi(x, y) = x \otimes
y$. Then these elements generate
$T$. Now let $P'$ be any $R$-module and let
$f : M \times N \to P'$ be an $R$-bilinear map. The unique
$R$-linear map from the free module $P$ to $P'$ that sends the
basis element $(x,y)$ to $f(x,y)$ kills every displayed generator
of $Q$, by bilinearity. It therefore induces an $R$-linear map
$f' : T = P/Q \to P'$ with $f'(x\otimes y) = f(x,y)$.
Thus $f = f'\circ\pi$. The elements $x\otimes y$ generate $T$,
so these values determine $f'$ uniquely.
Suppose there is another pair $(T', g')$ satisfying the same properties.
Then there is a unique $j : T \to T'$ and
also $j' : T' \to T$ such that $g' = j\circ g$, $g = j'\circ g'$.
But then both the maps $(j'\circ j) \circ g$ and $g$
satisfy the universal properties, so by uniqueness they are equal,
and hence $j'\circ j$ is identity on $T$.
Similarly $(j\circ j') \circ g' = g'$ and $j\circ j'$ is identity on $T'$.
So $j$ is an isomorphism.
\end{proof}


\subsection{Editorial treatment of Algebra lines 1933--1991}

\hypertarget{algebra-context-006}{}

\noindent\href{https://github.com/stacks/stacks-project/blob/a04446e57ec1fbc252a871afcec7752fb2807b14/algebra.tex\#L1933}{Original source, lines 1933--1991}. Complete original and reviewed TeX passages accompany this note. Every intervention is separately recorded; the following treatment is editorial.


\begin{lemma}[Tensor products commute with colimits]
\label{editorial-source-lemma-tensor-products-commute-with-limits}
Let $(M_i, \mu_{ij})$ be a system over the preordered set $I$.
Let $N$ be an $R$-module. Then
$$
\colim (M_i \otimes N) \cong (\colim M_i)\otimes N.
$$
Moreover, the isomorphism is induced by the homomorphisms
$\mu_i \otimes 1: M_i \otimes N \to M \otimes N$
where $M = \colim_i M_i$ with natural maps $\mu_i : M_i \to M$.
More generally, for any small index category $\mathcal{I}$ and any
diagram $i\mapsto M_i$ of $R$-modules, the same canonical map
$$
\colim_{i\in\mathcal{I}}(M_i\otimes_R N)
\longrightarrow
\left(\colim_{i\in\mathcal{I}}M_i\right)\otimes_R N,
\qquad [m_i\otimes n]\longmapsto\mu_i(m_i)\otimes n,
$$
is an isomorphism.
\end{lemma}

\begin{proof}
First proof. The functor $M' \mapsto M' \otimes_R N$ is left adjoint
to the functor $N' \mapsto \Hom_R(N, N')$ by
Lemma \ref{algebra-lemma-hom-from-tensor-product}. Thus $M' \mapsto M' \otimes_R N$
commutes with all colimits, see
Categories, Lemma \ref{categories-lemma-adjoint-exact}.

\medskip\noindent
Second direct proof. Let $P = \colim (M_i \otimes N)$ with coprojections
$\lambda_i : M_i \otimes N \to P$. Let $M = \colim M_i$ with coprojections
$\mu_i : M_i \to M$. Then for all $i\leq j$, the following diagram commutes:
$$
\xymatrix{
M_i \otimes N \ar[r]_{\mu_i \otimes 1} \ar[d]_{\mu_{ij} \otimes 1} &
M \otimes N \ar[d]^{\text{id}} \\
M_j \otimes N \ar[r]^{\mu_j \otimes 1} &
M \otimes N
}
$$
By Lemma \ref{algebra-lemma-homomorphism-limit} these maps induce a unique homomorphism
$\psi : P \to M \otimes N$ such that
$\mu_i \otimes 1 = \psi \circ \lambda_i$.

\medskip\noindent
To construct the inverse map, for each $i\in I$, there is the canonical
$R$-bilinear mapping $g_i : M_i \times N \to
M_i \otimes N$. This induces a unique mapping
$\widehat{\phi} : M \times N \to P$
such that $\widehat{\phi} \circ (\mu_i \times 1) = \lambda_i \circ g_i$.
It is $R$-bilinear. Thus it induces an
$R$-linear mapping $\phi : M \otimes N \to P$.
From the commutative diagram below:
$$
\xymatrix{
M_i \times N \ar[r]^{g_i} \ar[d]^{\mu_i \times \text{id}} &
M_i \otimes N\ar[r]_{\text{id}} \ar[d]_{\lambda_i} &
M_i \otimes N \ar[d]_{\mu_i \otimes \text{id}} \ar[rd]^{\lambda_i} \\
M \times N \ar[r]^{\widehat{\phi}} &
P \ar[r]^{\psi} & M \otimes N \ar[r]^{\phi} & P
}
$$
we see that $\psi\circ\widehat{\phi} = g$, the canonical $R$-bilinear mapping
$g : M \times N \to M \otimes N$. So
$\psi\circ\phi$ is identity on $M \otimes N$. From the right-hand square and
triangle, $\phi\circ\psi$ is also
identity on $P$.
\end{proof}


\subsection{Editorial treatment of Algebra lines 2087--2121}

\hypertarget{algebra-context-007}{}

\noindent\href{https://github.com/stacks/stacks-project/blob/a04446e57ec1fbc252a871afcec7752fb2807b14/algebra.tex\#L2087}{Original source, lines 2087--2121}. Complete original and reviewed TeX passages accompany this note. Every intervention is separately recorded; the following treatment is editorial.


\begin{lemma}
\label{editorial-source-lemma-tensor-localization}
Let $M$ be an $R$-module. Then the $S^{-1}R$-modules $S^{-1}M$
and $S^{-1}R \otimes_R M$ are canonically isomorphic, and the
canonical isomorphism $f : S^{-1}R \otimes_R M \to S^{-1}M$
is given by
$$
f((a/s) \otimes m) = am/s, \forall a \in R, m \in M, s \in S
$$
\end{lemma}

\begin{proof}
Obviously, the map
$f' : S^{-1}R \times M \to S^{-1}M$ given by $f'(a/s, m) = am/s$ is
bilinear, and thus by the
universal property, this map induces a unique $S^{-1}R$-module homomorphism
$f : S^{-1}R \otimes_R M \to S^{-1}M$ as in the statement of the lemma.
Actually every element in $S^{-1}M$ is of the form $m/s$, $m\in M, s\in S$ and
every element in
$S^{-1}R \otimes_R M$ is of the form $1/s \otimes m$. To see the latter fact,
write an element in $S^{-1}R\otimes_R M$ as a finite sum
indexed by a finite set $K$. Put
$s = \prod_{k\in K}s_k$ and
$t_k = \prod_{\ell\in K\setminus\{k\}}s_\ell$ for $k\in K$.
Thus $s=s_kt_k$ and all these elements belong to $S$.
Every sum in the next display is over $k\in K$.
For the empty sum take $s=1$, so its tensor representative is
$1\otimes 0=0$. We obtain
$$
\sum_k \frac{a_k}{s_k} \otimes m_k =
\sum_k \frac{a_k t_k}{s} \otimes m_k =
\frac{1}{s} \otimes \sum_k {a_k t_k}m_k = \frac{1}{s} \otimes m
$$
Here $m = \sum_k {a_k t_k}m_k$. Then it is obvious that $f$ is surjective,
and if $f(\frac{1}{s} \otimes m) = m/s = 0$ then there exists $t \in S$ with
$tm = 0$ in $M$. Then we have
$$
\frac{1}{s} \otimes m = \frac{1}{st} \otimes tm = \frac{1}{st} \otimes 0 = 0
$$
Therefore $f$ is injective.
\end{proof}


\subsection{Editorial treatment of Algebra lines 2240--2268}

\hypertarget{algebra-context-008}{}

\noindent\href{https://github.com/stacks/stacks-project/blob/a04446e57ec1fbc252a871afcec7752fb2807b14/algebra.tex\#L2240}{Original source, lines 2240--2268}. Complete original and reviewed TeX passages accompany this note. Every intervention is separately recorded; the following treatment is editorial.


\begin{lemma}
\label{editorial-source-lemma-presentation-sym-exterior}
Let $R$ be a ring.
Let $M_2 \to M_1 \to M \to 0$ be an exact sequence of $R$-modules.
For every integer $n \geq 1$ there are exact sequences
$$
M_2 \otimes_R \text{Sym}^{n - 1}(M_1)
\to
\text{Sym}^n(M_1)
\to
\text{Sym}^n(M)
\to
0
$$
and similarly
$$
M_2 \otimes_R \wedge^{n - 1}(M_1)
\to
\wedge^n(M_1)
\to
\wedge^n(M)
\to
0
$$
In degree zero both induced maps are the identity map $R \to R$.
\end{lemma}

\begin{proof}
Denote the maps by $u : M_2 \to M_1$ and $v : M_1 \to M$.
We prove both assertions, writing $C(M_1)$ for the symmetric algebra
or the exterior algebra in the respective case. Let $J$ be the
homogeneous two sided ideal of $C(M_1)$ generated by all
$u(y)$ in degree one, for $y \in M_2$.
The map $v$ induces a graded $R$-algebra map
$C(M_1)/J \to C(M)$. Conversely, the assignment
$v(x) \mapsto x \bmod J$ defines an $R$-linear map
$M \to (C(M_1)/J)^1$: it is independent of the lift $x$
because $\Ker(v) = \Im(u)$. Sending a pure tensor in $M$
to the product of these images defines a map from its tensor algebra,
with the identity map on the degree-zero part $R$.
In the symmetric case it kills the commutator relations.
In the exterior case it kills $m\otimes m$, since a lift $x$ of $m$
satisfies $x\wedge x = 0$ in $C(M_1)$.
Thus it induces a graded $R$-algebra map
$C(M) \to C(M_1)/J$. The two maps are inverse on $R$ and on
the degree-one generators, and hence on the entire algebras.

\medskip\noindent
It follows that the kernel of $C^n(M_1) \to C^n(M)$ is $J^n$.
For $n \geq 1$ every homogeneous term generating $J^n$ has the form
$a\,u(y)\,b$, with $a$ of degree $p$, $b$ of degree $q$ and
$p+1+q=n$. In the symmetric algebra this equals $u(y)ab$.
In the exterior algebra it equals $(-1)^p u(y)ab$:
on a pure degree-$p$ tensor this is obtained by $p$ successive
interchanges of degree-one factors, and the assertion extends by
$R$-linearity. Therefore $J^n$ is exactly the image of the map
$$
M_2\otimes_R C^{n-1}(M_1) \longrightarrow C^n(M_1),
\qquad y\otimes z \longmapsto u(y)z.
$$
The reverse inclusion holds because every $u(y)$ belongs to $J$.
This proves the two displayed exact sequences, including their
surjective last maps. In degree zero $J^0=0$, and both maps
$C^0(M_1)=R \to C^0(M)=R$ are the identity.
\end{proof}


\subsection{Editorial treatment of Algebra lines 2834--2858}

\hypertarget{algebra-context-009}{}

\noindent\href{https://github.com/stacks/stacks-project/blob/a04446e57ec1fbc252a871afcec7752fb2807b14/algebra.tex\#L2834}{Original source, lines 2834--2858}. Complete original and reviewed TeX passages accompany this note. Every intervention is separately recorded; the following treatment is editorial.


\begin{lemma}
\label{editorial-source-lemma-charpoly-module}
Let $R$ be a ring.
Let $M$ be a finite $R$-module.
Let $\varphi : M \to M$ be an endomorphism.
Then there exists a monic polynomial $P \in R[T]$ such that
$P(\varphi) = 0$ as an endomorphism of $M$.
\end{lemma}

\begin{proof}
Choose a surjective $R$-module map $R^{\oplus n} \to M$, given by
$(a_1, \ldots, a_n) \mapsto \sum a_ix_i$ for some generators $x_i \in M$.
Choose $(a_{i1}, \ldots, a_{in}) \in R^{\oplus n}$ such that
$\varphi(x_i) = \sum a_{ij} x_j$. In other words the diagram
$$
\xymatrix{
R^{\oplus n} \ar[d]_{A^{\mathrm{t}}} \ar[r] & M \ar[d]^\varphi \\
R^{\oplus n} \ar[r] & M
}
$$
is commutative where $A = (a_{ij})$ and $A^{\mathrm{t}}$
denotes its transpose: the lift sends the $i$th basis vector to
$\sum_j a_{ij}e_j$, which maps to $\varphi(x_i)$.
By
Lemma \ref{algebra-lemma-charpoly}
there exists a monic polynomial $P$ such that $P(A) = 0$.
Taking transposes gives $P(A^{\mathrm{t}})=P(A)^{\mathrm{t}}=0$,
since the coefficient matrices are scalar and transpose preserves
powers of this one matrix.
Then it follows that $P(\varphi) = 0$.
\end{proof}


\subsection{Editorial treatment of Algebra lines 2860--2887}

\hypertarget{algebra-context-010}{}

\noindent\href{https://github.com/stacks/stacks-project/blob/a04446e57ec1fbc252a871afcec7752fb2807b14/algebra.tex\#L2860}{Original source, lines 2860--2887}. Complete original and reviewed TeX passages accompany this note. Every intervention is separately recorded; the following treatment is editorial.


\begin{lemma}
\label{editorial-source-lemma-charpoly-module-ideal}
Let $R$ be a ring. Let $I \subset R$ be an ideal.
Let $M$ be a finite $R$-module.
Let $\varphi : M \to M$ be an endomorphism such
that $\varphi(M) \subset IM$.
Then there exists a monic polynomial
$P = T^n + a_1 T^{n - 1} + \ldots + a_n \in R[T]$
such that $a_j \in I^j$ and $P(\varphi) = 0$ as an endomorphism of $M$.
\end{lemma}

\begin{proof}
Choose a surjective $R$-module map $R^{\oplus n} \to M$, given by
$(a_1, \ldots, a_n) \mapsto \sum a_ix_i$ for some generators $x_i \in M$.
Choose $(a_{i1}, \ldots, a_{in}) \in I^{\oplus n}$ such that
$\varphi(x_i) = \sum a_{ij} x_j$. In other words the diagram
$$
\xymatrix{
R^{\oplus n} \ar[d]_{A^{\mathrm{t}}} \ar[r] & M \ar[d]^\varphi \\
I^{\oplus n} \ar[r] & M
}
$$
is commutative where $A = (a_{ij})$: its transpose sends
$e_i$ to $\sum_j a_{ij}e_j \in I^{\oplus n}$, whose image in
$M$ is $\varphi(x_i)$. By
Lemma \ref{algebra-lemma-charpoly}
the polynomial
$P(t) = \det(t\text{id}_{n \times n} - A)$
is monic of degree $n$ and satisfies $P(A)=0$.
Transposing yields $P(A^{\mathrm{t}})=0$, so the commutative diagram
and the surjection $R^{\oplus n}\to M$ give $P(\varphi)=0$.
In the determinant expansion, a term contributing to the coefficient
of $t^{n-j}$ uses exactly $j$ entries of $-A$ and $n-j$ copies of $t$
from diagonal entries. Each such term is therefore a product of
$j$ elements of $I$, with its determinant sign, and lies in $I^j$.
Their sum is the coefficient $a_j$, so $a_j\in I^j$ as required.
When $n=0$, $M=0$ and the empty determinant is $P=1$;
its action is the zero endomorphism of the zero module.
\end{proof}


\subsection{Editorial treatment of Algebra lines 2892--2912}

\hypertarget{algebra-context-011}{}

\noindent\href{https://github.com/stacks/stacks-project/blob/a04446e57ec1fbc252a871afcec7752fb2807b14/algebra.tex\#L2892}{Original source, lines 2892--2912}. Complete original and reviewed TeX passages accompany this note. Every intervention is separately recorded; the following treatment is editorial.


\begin{lemma}
\label{editorial-source-lemma-fun}
Let $R$ be a ring.
Let $M$ be a finite $R$-module.
Let $\varphi : M \to M$ be a surjective $R$-module map.
Then $\varphi$ is an isomorphism. Moreover, its inverse belongs
to the subalgebra $R[\varphi]\subset\operatorname{End}_R(M)$:
there is a polynomial $Q\in R[T]$ with $\varphi^{-1}=Q(\varphi)$.
\end{lemma}

\begin{proof}[First proof]
Write $R' = R[x]$ and think of $M$ as a finite $R'$-module with
$x$ acting via $\varphi$. Set $I = (x) \subset R'$. By our assumption that
$\varphi$ is surjective we have $IM = M$. Hence we may apply
Lemma \ref{editorial-source-lemma-charpoly-module-ideal}
to $M$ as an $R'$-module, the ideal $I$ and the endomorphism $\text{id}_M$.
We conclude that
$(1 + a_1 + \ldots + a_n)\text{id}_M = 0$ with $a_j \in I$.
Write $a_j = b_j(x)x$ for some $b_j(x) \in R[x]$.
Translating back into $\varphi$ we see that
$\text{id}_M = -(\sum_{j = 1}^{n} b_j(\varphi)) \varphi$, and hence
$\varphi$ is invertible. More explicitly, put
$Q(T)=-\sum_{j=1}^{n}b_j(T)\in R[T]$, where $b_j(T)$ is
obtained from $b_j(x)$ by replacing its indeterminate $x$ by $T$.
The displayed identity gives $Q(\varphi)\varphi=\mathrm{id}_M$.
Every power of $\varphi$ and every scalar multiplication commute
with $\varphi$, so also $\varphi Q(\varphi)=\mathrm{id}_M$.
Thus the same polynomial gives the two-sided inverse.
If $M=0$, the zero polynomial gives its unique endomorphism,
which is also its identity and inverse.
\end{proof}


\subsection{Editorial treatment of Algebra lines 2986--3102}

\hypertarget{algebra-context-012}{}

\noindent\href{https://github.com/stacks/stacks-project/blob/a04446e57ec1fbc252a871afcec7752fb2807b14/algebra.tex\#L2986}{Original source, lines 2986--3102}. Complete original and reviewed TeX passages accompany this note. Every intervention is separately recorded; the following treatment is editorial.


\begin{lemma}
\label{editorial-source-lemma-Zariski-topology}
Let $R$ be a ring.
\begin{enumerate}
\item The spectrum of a ring $R$ is empty if and only if $R$
is the zero ring.
\item Every nonzero ring has a maximal ideal.
\item Every nonzero ring has a minimal prime ideal.
\item Given an ideal $I \subset R$ and a prime ideal
$I \subset \mathfrak p$ there exists a prime
$I \subset \mathfrak q \subset \mathfrak p$ such
that $\mathfrak q$ is minimal over $I$.
\item If $T \subset R$, and if $(T)$ is the ideal generated by
$T$ in $R$, then $V((T)) = V(T)$.
\item If $I$ is an ideal and $\sqrt{I}$ is its radical,
see basic notion (\ref{algebra-item-radical-ideal}), then $V(I) = V(\sqrt{I})$.
\item Given an ideal $I$ of $R$ we have $\sqrt{I} =
\bigcap_{I \subset \mathfrak p} \mathfrak p$.
\item If $I$ is an ideal then $V(I) = \emptyset$ if and only
if $I$ is the unit ideal.
\item If $I$, $J$ are ideals of $R$ then $V(I) \cup V(J) =
V(I \cap J)$.
\item If $(I_a)_{a\in A}$ is a set of ideals of $R$ then
$\bigcap_{a\in A} V(I_a) = V(\bigcup_{a\in A} I_a)$.
\item If $f \in R$, then $D(f) \amalg V(f) = \Spec(R)$.
\item If $f \in R$ then $D(f) = \emptyset$ if and only if $f$
is nilpotent.
\item If $f = u f'$ for some unit $u \in R$, then $D(f) = D(f')$.
\item If $I \subset R$ is an ideal, and $\mathfrak p$ is a prime of
$R$ with $\mathfrak p \not\in V(I)$, then there exists an $f \in I$
such that $\mathfrak p \in D(f)$, and $D(f) \cap V(I) = \emptyset$.
\item If $f, g \in R$, then $D(fg) = D(f) \cap D(g)$.
\item If $f_i \in R$ for $i \in I$, then
$\bigcup_{i\in I} D(f_i)$ is the complement of $V(\{f_i \}_{i\in I})$
in $\Spec(R)$.
\item If $f \in R$ and $D(f) = \Spec(R)$, then $f$ is a unit.
\end{enumerate}
\end{lemma}

\begin{proof}
We address each part in the corresponding item below.
\begin{enumerate}
\item This is a direct consequence of (2) or (3).
\item Let $\mathfrak{A}$ be the set of all proper ideals of $R$. This set is
ordered by inclusion and is non-empty, since $(0) \in \mathfrak{A}$ is a proper
ideal. An empty totally ordered subset has the upper bound $(0)$.
Let $A$ be a nonempty totally ordered subset of $\mathfrak A$.
Then $\bigcup_{I \in A} I$ is in
fact an ideal. Since $1 \notin I$ for all $I \in A$, the union does not contain
$1$ and thus is proper. Hence $\bigcup_{I \in A} I$ is in $\mathfrak{A}$ and is
an upper bound for the set $A$. Thus by Zorn's lemma $\mathfrak{A}$ has a
maximal element, which is the sought-after maximal ideal.
\item Since $R$ is nonzero, it contains a maximal ideal which is a prime ideal.
Thus the set $\mathfrak{A}$ of all prime ideals of $R$ is nonempty.
$\mathfrak{A}$ is ordered by reverse-inclusion. An empty totally
ordered subset has any one of these prime ideals as an upper bound.
Let $A$ be a nonempty totally ordered subset of $\mathfrak{A}$.
The intersection $J = \bigcap_{I \in A} I$ is an ideal, and it is proper
because it is contained in any chosen member of $A$.
To prove that it is prime, let $xy \in J$.
Then $xy \in I$ for all $I \in A$. Now let $B = \{I \in A | y \in I\}$. Let $K
= \bigcap_{I \in B} I$. Since $A$ is totally ordered, either $K = J$ (and we're
done, since then $y \in J$) or $K \supset J$ and for all $I \in A$ such that
$I$ is properly contained in $K$, we have $y \notin I$. But that means that for
all those $I$, we have $x \in I$, since they are prime.
Here is why this implies $x \in J$. Since $K\ne J$, choose
$a\in K\setminus J$ and $I_0\in A$ with $a\notin I_0$.
For each $H\in B$, total ordering forces $I_0\subset H$:
the other possibility $H\subset I_0$ would put $a\in K\subset H$
inside $I_0$. Consequently $I_0\subset K$, strictly because
$a\in K\setminus I_0$. This also holds when $B$ is empty,
with its intersection interpreted as $R$.
For any $L\in A$, either $L\subset I_0$, in which case
$L\subsetneq K$ and the preceding primality argument gives $x\in L$,
or $I_0\subset L$, in which case $x\in I_0\subset L$.
Thus $x$ belongs to every $L\in A$ and hence to $J$. In either case,
$J$ is prime as desired. Hence by Zorn's lemma we get a maximal element which
in this case is a minimal prime ideal.
\item This is the same exact argument as (3) except you only consider prime
ideals contained in $\mathfrak{p}$ and containing $I$.
This collection is nonempty because it contains $\mathfrak p$,
which also supplies an upper bound for the empty chain in reverse inclusion.
For a nonempty chain its intersection still contains $I$, is contained
in $\mathfrak p$, and is prime by the argument in (3).
The resulting minimal member $\mathfrak q$ is minimal over $I$
among all primes: any prime between $I$ and $\mathfrak q$ also
lies inside $\mathfrak p$ and therefore belongs to the same collection.
\item $(T)$ is the smallest ideal containing $T$. Hence if $T \subset I$, some
ideal, then $(T) \subset I$ as well. Hence if $I \in V(T)$, then $I \in V((T))$
as well. The other inclusion is obvious.
\item Since $I \subset \sqrt{I}, V(\sqrt{I}) \subset V(I)$. Now let
$\mathfrak{p} \in V(I)$. Let $x \in \sqrt{I}$. Then $x^n \in I$ for some $n$.
Hence $x^n \in \mathfrak{p}$. But since $\mathfrak{p}$ is prime, a boring
induction argument gets you that $x \in \mathfrak{p}$. Hence $\sqrt{I} \subset
\mathfrak{p}$ and $\mathfrak{p} \in V(\sqrt{I})$.
\item Let $f \in R \setminus \sqrt{I}$. Then $f^n \notin I$ for all $n$. Hence
$S = \{1, f, f^2, \ldots\}$ is a multiplicative subset, not containing $0$.
Moreover, $S\cap I=\emptyset$ and $S^{-1}I$ is proper.
Indeed, if $1=a/s$ with $a\in I$ and $s\in S$, equality of fractions
gives $u(s-a)=0$ for some $u\in S$. Then $us=ua\in I\cap S$,
a contradiction. Thus $(S^{-1}R)/(S^{-1}I)$ is a nonzero ring.
By (2) it has a maximal ideal; its inverse image supplies a
prime ideal $\bar{\mathfrak{p}} \subset S^{-1}R$ containing $S^{-1}I$. Then the
pull-back $\mathfrak{p}$ in $R$ of $\bar{\mathfrak{p}}$ is a prime ideal
containing $I$ that does not intersect $S$. This shows that $\bigcap_{I \subset
\mathfrak p} \mathfrak p \subset \sqrt{I}$. Now if $a \in \sqrt{I}$, then $a^n
\in I$ for some $n$. Hence if $I \subset \mathfrak{p}$, then $a^n \in
\mathfrak{p}$. But since $\mathfrak{p}$ is prime, we have $a \in \mathfrak{p}$.
Thus the equality is shown. If $I=R$, both sides are $R$,
with the intersection over the empty set of prime ideals interpreted as $R$.
\item $I$ is not the unit ideal if and only if $I$
is contained in some maximal ideal (to
see this, apply (2) to the ring $R/I$) which is therefore prime.
\item If $\mathfrak{p} \in V(I) \cup V(J)$, then $I \subset \mathfrak{p}$ or $J
\subset \mathfrak{p}$ which means that $I \cap J \subset \mathfrak{p}$. Now if
$I \cap J \subset \mathfrak{p}$, then $IJ \subset \mathfrak{p}$ and hence
either $I \subset \mathfrak{p}$ or $J \subset \mathfrak{p}$, since $\mathfrak{p}$ is prime.
\item $\mathfrak{p} \in \bigcap_{a \in A} V(I_a) \Leftrightarrow
I_a \subset \mathfrak{p}, \forall a \in A \Leftrightarrow
\mathfrak{p} \in V(\bigcup_{a\in A} I_a)$
\item If $\mathfrak{p}$ is a prime ideal and $f \in R$, then either $f \in
\mathfrak{p}$ or $f \notin \mathfrak{p}$ (strictly) which is what the disjoint
union says.
\item If $a \in R$ is nilpotent, then $a^n = 0$ for some $n$. Hence $a^n \in
\mathfrak{p}$ for any prime ideal. Thus $a \in \mathfrak{p}$ as can be shown by
induction and $D(a) = \emptyset$. Now, as shown in (7), if $a \in R$ is not
nilpotent, then there is a prime ideal that does not contain it.
\item $f \in \mathfrak{p} \Leftrightarrow uf \in \mathfrak{p}$, since $u$ is
invertible.
\item If $\mathfrak{p} \notin V(I)$, then $\exists f \in I \setminus
\mathfrak{p}$. Then $f \notin \mathfrak{p}$ so $\mathfrak{p} \in D(f)$. Also if
$\mathfrak{q} \in D(f)$, then $f \notin \mathfrak{q}$ and thus $I$ is not
contained in $\mathfrak{q}$. Thus $D(f) \cap V(I) = \emptyset$.
\item If $fg \in \mathfrak{p}$, then $f \in \mathfrak{p}$ or $g \in
\mathfrak{p}$. Hence if $f \notin \mathfrak{p}$ and $g \notin \mathfrak{p}$,
then $fg \notin \mathfrak{p}$. Since $\mathfrak{p}$ is an ideal, if $fg \notin
\mathfrak{p}$, then $f \notin \mathfrak{p}$ and $g \notin \mathfrak{p}$.
\item $\mathfrak{p} \in \bigcup_{i \in I} D(f_i) \Leftrightarrow \exists i \in
I, f_i \notin \mathfrak{p} \Leftrightarrow \mathfrak{p} \in \Spec(R)
\setminus V(\{f_i\}_{i \in I})$
\item If $D(f) = \Spec(R)$, then $V(f) = \emptyset$ and
hence $fR = R$, so $f$ is a unit.
\end{enumerate}
\end{proof}


\subsection{Editorial treatment of Algebra lines 3162--3204}

\hypertarget{algebra-context-013}{}

\noindent\href{https://github.com/stacks/stacks-project/blob/a04446e57ec1fbc252a871afcec7752fb2807b14/algebra.tex\#L3162}{Original source, lines 3162--3204}. Complete original and reviewed TeX passages accompany this note. Every intervention is separately recorded; the following treatment is editorial.


\begin{lemma}
\label{editorial-source-lemma-spec-localization}
Let $R$ be a ring. Let $S \subset R$ be a multiplicative subset.
The map $R \to S^{-1}R$ induces via the functoriality of $\Spec$
a homeomorphism
$$
\Spec(S^{-1}R)
\longrightarrow
\{\mathfrak p \in \Spec(R) \mid S \cap \mathfrak p = \emptyset \}
$$
where the topology on the right hand side is that induced from the
Zariski topology on $\Spec(R)$. The inverse map is given
by $\mathfrak p \mapsto S^{-1}\mathfrak p = \mathfrak p(S^{-1}R)$.
\end{lemma}

\begin{proof}
Denote the right hand side of the arrow of the lemma by $D$.
Choose a prime $\mathfrak p' \subset S^{-1}R$ and let $\mathfrak p$ be
the inverse image of $\mathfrak p'$ in $R$. Since $\mathfrak p'$
does not contain $1$ we see that $\mathfrak p$ does not contain
any element of $S$. Hence $\mathfrak p \in D$ and we see that
the image is contained in $D$. Let $\mathfrak p \in D$.
By assumption the image $\overline{S}$ does not contain $0$.
By basic notion (\ref{algebra-item-localization-zero})
$\overline{S}^{-1}(R/\mathfrak p)$ is not the zero ring.
By basic notion (\ref{algebra-item-localize-ideal}) we see
$S^{-1}R / S^{-1}\mathfrak p = \overline{S}^{-1}(R/\mathfrak p)$
is a domain, and hence $S^{-1}\mathfrak p$ is a prime.
The equality of rings also shows that the inverse image of
$S^{-1}\mathfrak p$ in $R$ is equal to $\mathfrak p$,
because $R/\mathfrak p \to \overline{S}^{-1}(R/\mathfrak p)$
is injective by basic notion (\ref{algebra-item-localize-nonzerodivisors}).
In the other direction, start with $\mathfrak p'\subset S^{-1}R$
and put $\mathfrak p$ equal to its inverse image.
For $a\in R$ and $s\in S$, the element $s/1$ is a unit, so
$a/s\in\mathfrak p'$ if and only if $a/1\in\mathfrak p'$,
if and only if $a\in\mathfrak p$.
It follows that $S^{-1}\mathfrak p=\mathfrak p'$.
Together with the preceding contraction identity, this proves that
the map $\Spec(S^{-1}R) \to \Spec(R)$
is bijective onto $D$ with inverse as given.
It is continuous by Lemma \ref{algebra-lemma-spec-functorial}.
Finally, let $D(g) \subset \Spec(S^{-1}R)$ be a standard
open. Write $g = h/s$ for some $h\in R$ and $s\in S$.
Since $g$ and $h/1$ differ by a unit we have $D(g) =
D(h/1)$ in $\Spec(S^{-1}R)$.
Hence by Lemma \ref{algebra-lemma-spec-functorial} and the bijectivity
above the image of $D(g) = D(h/1)$ is $D \cap D(h)$.
This proves the map is open as well.
\end{proof}


\subsection{Editorial treatment of Algebra lines 3366--3397}

\hypertarget{algebra-context-014}{}

\noindent\href{https://github.com/stacks/stacks-project/blob/a04446e57ec1fbc252a871afcec7752fb2807b14/algebra.tex\#L3366}{Original source, lines 3366--3397}. Complete original and reviewed TeX passages accompany this note. Every intervention is separately recorded; the following treatment is editorial.


\begin{lemma}
\label{editorial-source-lemma-characterize-local-ring}
Let $R$ be a ring. The following are equivalent:
\begin{enumerate}
\item $R$ is a local ring,
\item $\Spec(R)$ has exactly one closed point,
\item $R$ has a maximal ideal $\mathfrak m$
and every element of $R \setminus \mathfrak m$
is a unit, and
\item $R$ is not the zero ring and for every $x \in R$ either $x$
or $1 - x$ is invertible or both.
\end{enumerate}
\end{lemma}

\begin{proof}
Let $R$ be a ring, and $\mathfrak m$ a maximal ideal.
If $x \in R \setminus \mathfrak m$, and $x$ is not a unit
then there is a maximal ideal $\mathfrak m'$ containing $x$.
Hence $R$ has at least two maximal ideals. Conversely,
if $\mathfrak m'$ is another maximal ideal, then choose
$x \in \mathfrak m'$, $x \not \in \mathfrak m$. Clearly
$x$ is not a unit. This proves the equivalence of (1) and (3).
The equivalence of (1) and (2) follows because the closure of $\{\mathfrak p\}$ is
$V(\mathfrak p)$. Indeed, this is a closed set containing
$\mathfrak p$, and any closed set $V(T)$ containing $\mathfrak p$
has $T\subset\mathfrak p$, hence contains $V(\mathfrak p)$.
A maximal ideal therefore gives a closed point. Conversely, if
$\{\mathfrak p\}$ is closed, apply (2) of
Lemma \ref{editorial-source-lemma-Zariski-topology} to $R/\mathfrak p$.
The inverse image of a maximal ideal of that quotient is a maximal
ideal $\mathfrak m$ of $R$ containing $\mathfrak p$.
Then $\mathfrak m\in V(\mathfrak p)=\{\mathfrak p\}$,
so $\mathfrak m=\mathfrak p$.
If $R$ is local then (4) holds since $x$ is either in $\mathfrak m$
or not. If (4) holds, and $\mathfrak m$, $\mathfrak m'$ are distinct
maximal ideals then we may choose $x \in R$ such that
$x \bmod \mathfrak m' = 0$ and $x \bmod \mathfrak m = 1$
by the Chinese remainder theorem
(Lemma \ref{algebra-lemma-chinese-remainder}).
This element $x$ is not invertible and neither is $1 - x$ which is
a contradiction. Thus (4) and (1) are equivalent.
\end{proof}


\subsection{Editorial treatment of Algebra lines 3499--3517}

\hypertarget{algebra-context-015}{}

\noindent\href{https://github.com/stacks/stacks-project/blob/a04446e57ec1fbc252a871afcec7752fb2807b14/algebra.tex\#L3499}{Original source, lines 3499--3517}. Complete original and reviewed TeX passages accompany this note. Every intervention is separately recorded; the following treatment is editorial.


\begin{lemma}
\label{editorial-source-lemma-in-image}
Let $\varphi : R \to S$ be a ring map. Let $\mathfrak p$
be a prime of $R$. The following are equivalent
\begin{enumerate}
\item $\mathfrak p$ is in the image of
$\Spec(S) \to \Spec(R)$,
\item $S \otimes_R \kappa(\mathfrak p) \not = 0$,
\item $S_{\mathfrak p}/\mathfrak p S_{\mathfrak p} \not = 0$,
\item $(S/\mathfrak pS)_{\mathfrak p} \not = 0$, and
\item $\mathfrak p = \varphi^{-1}(\mathfrak pS)$.
\end{enumerate}
\end{lemma}

\begin{proof}
We have already seen the equivalence of the first two
in Remark \ref{algebra-remark-fundamental-diagram}.
Put $T=R\setminus\mathfrak p$ and $Q=S/\mathfrak pS$, with
$T$ acting on $Q$ through $\varphi$.
The identifications in that remark identify the three rings in
(2), (3), and (4) with $T^{-1}Q$. In particular the first map is
$$
S\otimes_R\kappa(\mathfrak p)\longrightarrow T^{-1}Q,\qquad
s\otimes\overline{a/u}\longmapsto
\frac{\varphi(a)s+\mathfrak pS}{u},
$$
where $a\in R$, $u\in T$, and the bar denotes the class in
$R_{\mathfrak p}/\mathfrak pR_{\mathfrak p}$.
The formula is $R$-balanced, respects multiplication, kills
$\mathfrak pR_{\mathfrak p}$, and inverts every member of $T$.
Its inverse sends $(s+\mathfrak pS)/u$ to
$s\otimes\overline{1/u}$: the relations defining the quotient
and localization hold because $\mathfrak p$ acts as zero and
each $u\in T$ as a unit in the second tensor factor.
Both composites are the identity on the displayed generators.
The corresponding quotient-localization map is
$s/\varphi(u)+\mathfrak pS_{\mathfrak p}\mapsto
(s+\mathfrak pS)/u$, with the reverse formula as inverse.
Finally, $T^{-1}Q$ is the zero ring if and only if there is
$u\in T$ with $\varphi(u)\in\mathfrak pS$, by the equality
criterion for $1/1=0/1$ in a localization.
The ideal $\varphi^{-1}(\mathfrak pS)$ always contains
$\mathfrak p$. Hence it misses $T$ if and only if it equals
$\mathfrak p$. This proves the equivalence with (5).
\end{proof}


\subsection{Editorial treatment of Algebra lines 3519--3538}

\hypertarget{algebra-context-016}{}

\noindent\href{https://github.com/stacks/stacks-project/blob/a04446e57ec1fbc252a871afcec7752fb2807b14/algebra.tex\#L3519}{Original source, lines 3519--3538}. Complete original and reviewed TeX passages accompany this note. Every intervention is separately recorded; the following treatment is editorial.


\begin{remark}
\label{editorial-source-remark-local-ring-fibre}
Let $R \to S$ be a ring map. Let $\mathfrak q$ be a prime
ideal of $S$ lying over the prime ideal $\mathfrak p$ of $R$.
According to Remark \ref{algebra-remark-fundamental-diagram}
the prime $\mathfrak q$ corresponds to a unique prime
$\overline{\mathfrak q}$ of the fibre ring
$F = S \otimes_R \kappa(\mathfrak p)$.
Then we have
$$
F_{\overline{\mathfrak q}}
\cong S_\mathfrak q \otimes_{R_\mathfrak p} \kappa(\mathfrak p)
\cong S_\mathfrak q/\mathfrak p S_\mathfrak q
$$
Namely, there is an obvious ring map
$F \to S_\mathfrak q \otimes_{R_\mathfrak p} \kappa(\mathfrak p)$
which, under the displayed isomorphism, identifies with $F \to F_{\overline{\mathfrak q}}$.
The second isomorphism follows from the fact that $\kappa(\mathfrak p)$
is the quotient of $R_\mathfrak p$ by $\mathfrak pR_\mathfrak p$.
Explicitly, write $T=R\setminus\mathfrak p$. The fibre-ring
description is $F=T^{-1}(S/\mathfrak pS)$, and the prime in it is
$\overline{\mathfrak q}=T^{-1}(\mathfrak q/\mathfrak pS)$.
For $a\in S$ and $u\in T$, the fraction
$(a+\mathfrak pS)/u$ belongs to $\overline{\mathfrak q}$
if and only if $a\in\mathfrak q$; here every $\varphi(u)$
is outside $\mathfrak q$.
Thus the mutually inverse maps between $F_{\overline{\mathfrak q}}$
and $S_{\mathfrak q}/\mathfrak pS_{\mathfrak q}$ are
$$
\frac{(a+\mathfrak pS)/u}{(b+\mathfrak pS)/v}
\longmapsto
\frac{a\varphi(v)}{\varphi(u)b}+\mathfrak pS_{\mathfrak q},
\qquad
\frac{a}{b}+\mathfrak pS_{\mathfrak q}
\longmapsto
\frac{(a+\mathfrak pS)/1}{(b+\mathfrak pS)/1},
$$
where $u,v\in T$ and $b\in S\setminus\mathfrak q$.
The first map is induced by $S\to S_{\mathfrak q}/\mathfrak pS_{\mathfrak q}$:
it kills $\mathfrak pS$, inverts $T$, and inverts exactly the
additional denominators displayed above. The second is induced by
$S\to F_{\overline{\mathfrak q}}$, which inverts
$S\setminus\mathfrak q$ and kills $\mathfrak pS$.
The formulas show that both composites fix every fraction.
The tensor comparison sends
$$
\frac{a}{b}\otimes\overline{r/t}
\longmapsto
\frac{a\varphi(r)}{b\varphi(t)}+\mathfrak pS_{\mathfrak q},
$$
for $r\in R$, $t\in T$, and has inverse
$a/b+\mathfrak pS_{\mathfrak q}\mapsto(a/b)\otimes\overline{1}$.
Balancing over $R_{\mathfrak p}$ and the vanishing of
$\mathfrak pR_{\mathfrak p}$ prove well-definedness and both
inverse identities. These maps identify the original map out of
$F$ with localization at $\overline{\mathfrak q}$.
\end{remark}


\subsection{Editorial treatment of Algebra lines 3586--3601}

\hypertarget{algebra-context-017}{}

\noindent\href{https://github.com/stacks/stacks-project/blob/a04446e57ec1fbc252a871afcec7752fb2807b14/algebra.tex\#L3586}{Original source, lines 3586--3601}. Complete original and reviewed TeX passages accompany this note. Every intervention is separately recorded; the following treatment is editorial.


\begin{lemma}
\label{editorial-source-lemma-surjective-on-spec-units}
Let $\varphi : R \to S$ be a ring map such that the induced map
$\Spec(S) \to \Spec(R)$ is surjective. Then an element $x \in R$
is a unit if and only if $\varphi(x) \in S$ is a unit.
The same conclusion holds under the weaker assumption that the image
of $\Spec(S)\to\Spec(R)$ contains every closed point of $\Spec(R)$.
\end{lemma}

\begin{proof}
If $x$ is a unit, then so is $\varphi(x)$. Conversely, if $\varphi(x)$
is a unit, then $\varphi(x) \not \in \mathfrak q$ for all
$\mathfrak q \in \Spec(S)$. Hence
$x \not \in \varphi^{-1}(\mathfrak q) = \Spec(\varphi)(\mathfrak q)$
for all $\mathfrak q \in \Spec(S)$. Since $\Spec(\varphi)$ is surjective
we conclude that $x$ is a unit by
part (17) of Lemma \ref{editorial-source-lemma-Zariski-topology}.
For the weaker assumption, suppose that $x$ is not a unit.
Then $xR$ is proper, and a maximal ideal of the nonzero ring $R/xR$
pulls back to a maximal ideal $\mathfrak m$ of $R$ containing $x$.
This is a closed point, so the assumption supplies
$\mathfrak q\in\Spec(S)$ with
$\varphi^{-1}(\mathfrak q)=\mathfrak m$.
Consequently $\varphi(x)\in\mathfrak q$, so $\varphi(x)$ is not a unit.
The forward implication still follows by applying $\varphi$ to an inverse.
The assumption is strictly weaker: for a prime integer $p$, the residue map
$\mathbf Z_{(p)}\to\mathbf F_p$ hits the unique closed point $(p)$
of $\Spec(\mathbf Z_{(p)})$, but misses its prime ideal $(0)$.
Indeed, localization of the integer prime ideals leaves exactly $(0)$
and $(p)$, and the only prime of the field $\mathbf F_p$ contracts to $(p)$.
\end{proof}


\subsection{Editorial treatment of Algebra lines 3698--3733}

\hypertarget{algebra-context-018}{}

\noindent\href{https://github.com/stacks/stacks-project/blob/a04446e57ec1fbc252a871afcec7752fb2807b14/algebra.tex\#L3698}{Original source, lines 3698--3733}. Complete original and reviewed TeX passages accompany this note. Every intervention is separately recorded; the following treatment is editorial.


\begin{lemma}
\label{editorial-source-lemma-NAK-localization}
Let $R$ be a ring, let $S \subset R$ be a multiplicative subset,
let $I \subset R$ be an ideal, and let $M$ be a finite $R$-module.
If $x_1, \ldots, x_r \in M$ generate $S^{-1}(M/IM)$
as an $S^{-1}(R/I)$-module, then there exists an $f \in S + I$
such that $x_1, \ldots, x_r$ generate $M_f$ as an
$R_f$-module.\footnote{Special cases: (I) $I = 0$. The lemma says
if $x_1, \ldots, x_r$ generate $S^{-1}M$, then $x_1, \ldots, x_r$
generate $M_f$ for some $f \in S$. (II) $I = \mathfrak p$ is
a prime ideal and $S = R \setminus \mathfrak p$. The lemma says if
$x_1, \ldots, x_r$ generate $M \otimes_R \kappa(\mathfrak p)$
then $x_1, \ldots, x_r$ generate $M_f$ for some
$f \in R$, $f \not \in \mathfrak p$.}
\end{lemma}

\begin{proof}
Special case $I = 0$. Let $y_1, \ldots, y_s$ be generators for $M$ over $R$.
Since $S^{-1}M$ is generated by $x_1, \ldots, x_r$, for each $i$
we can write $y_i = \sum (a_{ij}/s_{ij})x_j$ in $S^{-1}M$
for some $a_{ij} \in R$ and $s_{ij} \in S$. Multiplying by the product
$s=\prod_{k=1}^{s_0}\prod_{\ell=1}^r s_{k\ell}\in S$,
where $s_0$ denotes the number of chosen generators $y_i$, we see that
$sy_i = \sum_{j=1}^r a'_{ij}x_j$ in $S^{-1}M$, with
$$
a'_{ij}=a_{ij}
\prod_{\substack{1\leq k\leq s_0,\ 1\leq\ell\leq r\\(k,\ell)\ne(i,j)}}
s_{k\ell}\in R.
$$
Empty products are $1$ and empty sums are $0$.
This in turn means there exist $t_i \in S$ such that
$t_isy_i = \sum t_ia'_{ij}x_j$ in $M$. Thus if $t \in S$
is the product of the $t_i$, put
$t^{(i)}=\prod_{k\ne i}t_k$, so $t=t_it^{(i)}$.
Multiplying the equality in $M$ by $t^{(i)}$ gives
$sty_i=\sum_{j=1}^r t a'_{ij}x_j$.
In $M_{st}$ it follows that
$y_i=\sum_{j=1}^r(t a'_{ij}/(st))x_j$.
Thus $y_i$ is in the $R_{st}$-submodule generated by
$x_1,\ldots,x_r$. Hence $x_1, \ldots, x_r$ generate $M_{st}$.

\medskip\noindent
General case. By the special case, we can find an $s \in S$
such that $x_1, \ldots, x_r$ generate $(M/IM)_s$ over $(R/I)_s$.
By Lemma \ref{algebra-lemma-NAK} we can find a $g \in 1 + I_s \subset R_s$
such that $x_1, \ldots, x_r$ generate $(M_s)_g$ over $(R_s)_g$.
Write $g=1+i/s'$ with $i\in I$ and $s'=s^n$ for some
$n\geq0$. Put $f=ss'+is=s(s'+i)$. Since $ss'=s^{n+1}\in S$
and $is\in I$, we have $f\in S+I$.
In $(R_s)_g$, the image of $f$ equals $ss'g$ and is a unit.
Conversely, in $R_f$ the element $s$ has inverse $(s'+i)/f$,
and the image of $g$ has inverse $ss'/f$:
$$
s\frac{s'+i}{f}=1,\qquad
\left(1+\frac{i}{s'}\right)\frac{ss'}{f}=1.
$$
The two universal properties therefore give mutually inverse maps
$R_f\longleftrightarrow(R_s)_g$ which agree with the original
maps from $R$. More explicitly, the first sends $r/f^k$ to
$(r/1)/(ss'g)^k$. The reverse map sends
$r/s^a$ to $r(s'+i)^a/f^a$ and $g^{-1}$ to $ss'/f$.
Their composites fix the image of every $r\in R$ and each inverse
adjoined in the respective localizations, so both composites are identities.
The corresponding module maps are also inverse: they send
$m/f^k$ to $(m/1)/(ss'g)^k$ and send
$(m/s^a)/g^b$ to $m(s'+i)^a(ss')^b/f^{a+b}$.
They preserve the images of all the chosen $x_j$.
Consequently those images generate $M_f$ over $R_f$, as required.
These localization identities remain valid when a localization is
the zero ring and when the original lists of generators are empty.
\end{proof}


\subsection{Editorial treatment of Algebra lines 4007--4031}

\hypertarget{algebra-context-019}{}

\noindent\href{https://github.com/stacks/stacks-project/blob/a04446e57ec1fbc252a871afcec7752fb2807b14/algebra.tex\#L4007}{Original source, lines 4007--4031}. Complete original and reviewed TeX passages accompany this note. Every intervention is separately recorded; the following treatment is editorial.


\begin{lemma}
\label{editorial-source-lemma-connected-component}
Let $R$ be a ring.
A connected component of
$\Spec(R)$ is of the form $V(I)$,
where $I$ is an ideal generated by idempotents
such that every idempotent of $R$ either maps to
$0$ or $1$ in $R/I$.
Conversely, if $I$ is a proper ideal generated by idempotents and
every idempotent of $R$ maps to $0$ or $1$ in $R/I$, then $V(I)$
is a connected component of $\Spec(R)$.
\end{lemma}

\begin{proof}
Let $\mathfrak p$ be a prime of $R$. By
Lemma \ref{algebra-lemma-topology-spec}
we see that the hypotheses of
Topology, Lemma \ref{topology-lemma-connected-component-intersection}
are satisfied for the topological space $\Spec(R)$.
Hence the connected component of $\mathfrak p$ in $\Spec(R)$
is the intersection of open and closed subsets of $\Spec(R)$
containing $\mathfrak p$. Hence it equals $V(I)$ where
$I$ is generated by the idempotents $e \in R$ such that $e$ maps to $0$
in $\kappa(\mathfrak p)$, see
Lemma \ref{algebra-lemma-disjoint-decomposition}.
Any idempotent $e$ which is not in this collection maps to $1$
in $\kappa(\mathfrak p)$. Its complementary idempotent $1-e$
therefore belongs to $I$, so $e$ maps to $1$ in $R/I$.

\medskip\noindent
For the converse, we first show that idempotents lift modulo any ideal
$I$ generated by idempotents. Let $\overline a\in R/I$ be idempotent
and choose a lift $a\in R$. Write
$a^2-a=\sum_{j=1}^m r_je_j$, where $r_j\in R$ and the $e_j\in I$
are idempotent. Define
$$
c=\prod_{j=1}^m(1-e_j),\qquad
e=1-c=
\sum_{\emptyset\ne J\subset\{1,\ldots,m\}}
(-1)^{|J|+1}\prod_{j\in J}e_j.
$$
Every factor $1-e_j$ is idempotent, so $c^2=c$ and $e^2=e$.
The expanded sum shows $e\in I$. For each $j$ the product $ce_j$
contains the factor $(1-e_j)e_j=0$, hence
$c(a^2-a)=\sum_{j=1}^m r_jce_j=0$.
It follows that $b=ca$ satisfies
$b^2-b=c^2a^2-ca=c(a^2-a)=0$ and $b-a=-ea\in I$.
Thus $b$ is an idempotent lift of $\overline a$.
This also covers $m=0$, when $c=1$, $e=0$, and $a$ is already idempotent.
Under the hypothesis of the converse, every such $b$ maps to $0$ or $1$,
so $R/I$ has no nontrivial idempotents. Since $I$ is proper, $R/I$ is
nonzero. Lemmas \ref{algebra-lemma-characterize-spec-connected} and
\ref{algebra-lemma-spec-closed} show that $V(I)$ is nonempty and connected.
By Lemma \ref{algebra-lemma-closed-union-connected-components}, it is also
a union of connected components of $\Spec(R)$. A nonempty connected
subset lies in a single connected component; since $V(I)$ is a union
of such components, it is exactly that component.
\end{proof}


\subsection{Editorial treatment of Algebra lines 4133--4217}

\hypertarget{algebra-context-020}{}

\noindent\href{https://github.com/stacks/stacks-project/blob/a04446e57ec1fbc252a871afcec7752fb2807b14/algebra.tex\#L4133}{Original source, lines 4133--4217}. Complete original and reviewed TeX passages accompany this note. Every intervention is separately recorded; the following treatment is editorial.


\begin{lemma}
\label{editorial-source-lemma-cover}
\begin{slogan}
Zariski-local properties of modules and algebras
\end{slogan}
Let $R$ be a ring. Let $M$ be an $R$-module. Let $S$ be an $R$-algebra.
Suppose that $f_1, \ldots, f_n$ is a finite list of
elements of $R$ such that $\bigcup D(f_i) = \Spec(R)$,
in other words $(f_1, \ldots, f_n) = R$.
\begin{enumerate}
\item If each $M_{f_i} = 0$ then $M = 0$.
\item If each $M_{f_i}$ is a finite $R_{f_i}$-module,
then $M$ is a finite $R$-module.
\item If each $M_{f_i}$ is a finitely presented $R_{f_i}$-module,
then $M$ is a finitely presented $R$-module.
\item Let $M \to N$ be a map of $R$-modules. If $M_{f_i} \to N_{f_i}$
is an isomorphism for each $i$ then $M \to N$ is an isomorphism.
\item Let $0 \to M'' \to M \to M' \to 0$ be a complex of $R$-modules.
If $0 \to M''_{f_i} \to M_{f_i} \to M'_{f_i} \to 0$ is exact for each $i$,
then $0 \to M'' \to M \to M' \to 0$ is exact.
\item If each $R_{f_i}$ is Noetherian, then $R$ is Noetherian.
\item If each $S_{f_i}$ is a finite type $R_{f_i}$-algebra, then
$S$ is a finite type $R$-algebra.
\item If each $S_{f_i}$ is of finite presentation over $R_{f_i}$, then
$S$ is a finitely presented $R$-algebra.
\end{enumerate}
\end{lemma}

\begin{proof}
We prove each of the parts in turn.
\begin{enumerate}
\item By Proposition \ref{algebra-proposition-localize-twice}
this implies $M_\mathfrak p = 0$ for all $\mathfrak p \in \Spec(R)$,
so we conclude by Lemma \ref{algebra-lemma-characterize-zero-local}.
\item For each $i$ take a finite generating set $X_i$ of $M_{f_i}$.
For each $x\in X_i$, write $x=y_x/f_i^{a_x}$ with
$y_x\in M$ and $a_x\geq0$, and put $Y_i=\{y_x\mid x\in X_i\}$.
Keep the original elements $x$: the image of $y_x$ is $f_i^{a_x}x$,
and $x=f_i^{-a_x}(y_x/1)$ in $M_{f_i}$.
Since $f_i$ is a unit in $R_{f_i}$, the images of $Y_i$ and the
original set $X_i$ generate the same $R_{f_i}$-submodule. Let $Y$
be the union of these sets. This is still a finite set.
Consider the obvious $R$-linear map $R^Y \rightarrow M$ sending the basis
element $e_y$ to $y$. By assumption this map is surjective after localizing
at an arbitrary prime ideal $\mathfrak p$ of $R$, so it is surjective by
Lemma \ref{algebra-lemma-characterize-zero-local}
and $M$ is finitely generated.
\item By (2) we have a short exact sequence
$$
0 \rightarrow K \rightarrow R^m \rightarrow M \rightarrow 0
$$
Since localization is an exact functor and $M_{f_i}$ is finitely
presented we see that $K_{f_i}$ is finitely generated for all
$1 \leq i \leq n$ by Lemma \ref{algebra-lemma-extension}.
By (2) this implies that $K$ is a finite $R$-module and therefore
$M$ is finitely presented.
\item By Proposition \ref{algebra-proposition-localize-twice}
the assumption implies that the induced morphism
on localizations at all prime ideals is an isomorphism, so we conclude
by Lemma \ref{algebra-lemma-characterize-zero-local}.
\item By Proposition \ref{algebra-proposition-localize-twice} the assumption
implies that the induced
sequence of localizations at all prime ideals is short exact, so we
conclude by Lemma \ref{algebra-lemma-characterize-zero-local}.
\item We will show that every ideal of $R$ has a finite generating set:
For this, let $I \subset R$ be an arbitrary ideal. By
Proposition \ref{algebra-proposition-localization-exact}
each $I_{f_i} \subset R_{f_i}$ is an ideal. These are all
finitely generated by assumption, so we conclude by (2).
\item For each $i$ take a finite generating set $X_i$ of $S_{f_i}$.
For each $x\in X_i$, write $x=y_x/f_i^{a_x}$ with
$y_x\in S$ and $a_x\geq0$, and put $Y_i=\{y_x\mid x\in X_i\}$.
The equalities $y_x/1=f_i^{a_x}x$ and
$x=f_i^{-a_x}(y_x/1)$ show that the images of $Y_i$ generate
the same $R_{f_i}$-subalgebra as $X_i$, because both scalar factors
belong to the base ring $R_{f_i}$. Let $Y$ be the union of the
finite sets $Y_i$; it is finite.
Consider the algebra homomorphism $R[X_y]_{y \in Y} \rightarrow S$
induced by $Y$. Since it is an algebra homomorphism, the image $T$
is an $R$-submodule of the $R$-module $S$, so we can consider the
quotient module $S/T$. By assumption, this is zero if we localize
at the $f_i$, so it is zero by (1) and therefore $S$ is an
$R$-algebra of finite type.
\item By the previous item, there exists a surjective $R$-algebra
homomorphism $R[X_1, \ldots, X_n] \rightarrow S$. Let $K$ be the kernel
of this map. This is an ideal in $R[X_1, \ldots, X_n]$, finitely generated
in each localization at $f_i$. Since the $f_i$ generate the unit ideal
in $R$, they also generate the unit ideal in $R[X_1, \ldots, X_n]$, so an
application of (2) finishes the proof.
\end{enumerate}
\end{proof}


\subsection{Editorial treatment of Algebra lines 4306--4351}

\hypertarget{algebra-context-021}{}

\noindent\href{https://github.com/stacks/stacks-project/blob/a04446e57ec1fbc252a871afcec7752fb2807b14/algebra.tex\#L4306}{Original source, lines 4306--4351}. Complete original and reviewed TeX passages accompany this note. Every intervention is separately recorded; the following treatment is editorial.


\begin{lemma}
\label{editorial-source-lemma-cover-module}
Let $R$ be a ring. Let $f_1, \ldots, f_n$ be elements of $R$
generating the unit ideal. Let $M$ be an $R$-module.
The sequence
$$
0 \to
M \xrightarrow{\alpha}
\bigoplus\nolimits_{i = 1}^n M_{f_i} \xrightarrow{\beta}
\bigoplus\nolimits_{i, j = 1}^n M_{f_i f_j}
$$
is exact, where $\alpha(m) = (m/1, \ldots, m/1)$
and $\beta(m_1/f_1^{e_1}, \ldots, m_n/f_n^{e_n})
= (m_i/f_i^{e_i} - m_j/f_j^{e_j})_{(i, j)}$.
\end{lemma}

\begin{proof}
It suffices to prove exactness after localization at each maximal
ideal $\mathfrak m$, by Lemma \ref{algebra-lemma-characterize-zero-local}.
The finite direct sums commute with localization. For each original
index $i$, the comparison of the two localization orders is
$$
(M_{f_i})_{\mathfrak m}\longrightarrow
(M_{\mathfrak m})_{f_i/1},\qquad
\frac{x/f_i^a}{s}\longmapsto\frac{x/s}{(f_i/1)^a},
$$
where $s\in R\setminus\mathfrak m$. Its reverse formula is the
inverse, by Proposition \ref{editorial-source-proposition-localize-twice-module}.
The same formulas with $f_if_j$ give the comparisons on the last term.
They commute with the original maps $\alpha$ and $\beta$, including
the difference in each ordered pair $(i,j)$.

\medskip\noindent
Choose an original index $k$ with $f_k\notin\mathfrak m$, which exists
because $(f_1,\ldots,f_n)=R$. Keep all indices and all elements $f_i$.
The natural map
$\iota_k:M_{\mathfrak m}\to(M_{f_k})_{\mathfrak m}$
has the explicit inverse
$$
\theta_k\left(\frac{x/f_k^a}{s}\right)=\frac{x}{sf_k^a}
\quad\text{in }M_{\mathfrak m}.
$$
The denominator $sf_k^a$ lies outside $\mathfrak m$.
For each $i$, the natural map
$(M_{f_i})_{\mathfrak m}\to(M_{f_kf_i})_{\mathfrak m}$
is also an isomorphism, with inverse
$$
\theta_{k,i}\left(\frac{x/(f_kf_i)^a}{s}\right)
=\frac{x/f_i^a}{sf_k^a}.
$$
Its forward formula sends $(x/f_i^a)/s$ to
$(f_k^ax/(f_kf_i)^a)/s$. The two composites are identities by
the equality-of-fractions relations. These maps come from localization
at elements which already act invertibly, so the formulas are
well-defined and preserve the original module structures.

\medskip\noindent
Projection to the $k$th summand followed by $\theta_k$ is a left
inverse of $\alpha_{\mathfrak m}$, proving injectivity.
Now let $(z_i)_{i=1}^n$ be in the kernel of $\beta_{\mathfrak m}$
and put $z=\theta_k(z_k)\in M_{\mathfrak m}$.
The $(k,i)$ coordinate says that the images of $z_k$ and $z_i$
in $(M_{f_kf_i})_{\mathfrak m}$ agree. Applying $\theta_{k,i}$
shows that $z_i$ is the image of $z$ in $(M_{f_i})_{\mathfrak m}$.
For completeness, on a representative $z_k=(x/f_k^a)/s$ the image
in the pair localization is $(f_i^ax/(f_kf_i)^a)/s$; applying
$\theta_{k,i}$ gives $(f_i^ax/f_i^a)/(sf_k^a)$, which is precisely
the image of $x/(sf_k^a)=z$. Thus
$\alpha_{\mathfrak m}(z)=(z_i)_{i=1}^n$.
The reverse inclusion follows from the original identity
$\beta\alpha=0$, so the localized sequence is exact.
The local exactness criterion proves the original sequence exact.
If the covering list is empty, the unit-ideal hypothesis gives $R=0$,
and every unital $R$-module is zero, so the same assertion holds.
\end{proof}


\subsection{Editorial treatment of Algebra lines 4438--4523}

\hypertarget{algebra-context-022}{}

\noindent\href{https://github.com/stacks/stacks-project/blob/a04446e57ec1fbc252a871afcec7752fb2807b14/algebra.tex\#L4438}{Original source, lines 4438--4523}. Complete original and reviewed TeX passages accompany this note. Every intervention is separately recorded; the following treatment is editorial.


\begin{lemma}
\label{editorial-source-lemma-glue-modules}
Let $R$ be a ring. Let $f_1, \ldots, f_n \in R$. Suppose we are given
the following data:
\begin{enumerate}
\item For each $i$ an $R_{f_i}$-module $M_i$.
\item For each pair $i, j$ an $R_{f_if_j}$-module isomorphism
$\psi_{ij} : (M_i)_{f_j} \to (M_j)_{f_i}$.
\end{enumerate}
which satisfy the ``cocycle condition'' that all the diagrams
$$
\xymatrix{
(M_i)_{f_jf_k}
\ar[rd]_{\psi_{ij}}
\ar[rr]^{\psi_{ik}}
& &
(M_k)_{f_if_j} \\
&
(M_j)_{f_if_k} \ar[ru]_{\psi_{jk}}
}
$$
commute (for all triples $i, j, k$). Given this data define
$$
M = \Ker\left(
\bigoplus\nolimits_{1 \leq i \leq n} M_i
\longrightarrow
\bigoplus\nolimits_{1 \leq i, j \leq n} (M_i)_{f_j}
\right)
$$
where $(m_1, \ldots, m_n)$ maps to the element whose
$(i, j)$th entry is $m_i/1 - \psi_{ji}(m_j/1)$.
Then the natural map $M \to M_i$ induces an isomorphism
$M_{f_i} \to M_i$. Moreover $\psi_{ij}(m/1) = m/1$
for all $m \in M$ (with obvious notation).
\end{lemma}

\begin{proof}
Fix an original index $k\in\{1,\ldots,n\}$.
Localizing the defining kernel at $f_k$ and using exactness of
localization gives the isomorphism
$$
M_{f_k}\longrightarrow H_k:=
\Ker\left(
\bigoplus_{i=1}^n(M_i)_{f_k}\longrightarrow
\bigoplus_{i,j=1}^n((M_i)_{f_j})_{f_k}
\right),\qquad
\frac{(m_i)_{i=1}^n}{f_k^a}\longmapsto
\left(\frac{m_i}{f_k^a}\right)_{i=1}^n.
$$
The map defining $H_k$ keeps the original $(i,j)$ component
$z_i/1-\psi_{ji}(z_j/1)$, after the indicated localizations.
The comparison between localization orders sends
$(x/f_j^a)/f_k^b$ to $(x/f_k^b)/f_j^a$ and has the reverse
formula as inverse, by Proposition \ref{editorial-source-proposition-localize-twice-module}.

\medskip\noindent
Since $M_k$ is already an $R_{f_k}$-module, the canonical map
$M_k\to(M_k)_{f_k}$ has inverse
$x/f_k^a\mapsto f_k^{-a}x$.
The cocycle condition for $(k,k,k)$ gives $\psi_{kk}^2=\psi_{kk}$
under these comparisons; since $\psi_{kk}$ is invertible, it is the
identity. Likewise the conditions with repeated indices give
$\psi_{ik}\psi_{ki}=\mathrm{id}$ on each overlap.
Define
$$
\eta_k:M_k\longrightarrow\bigoplus_{i=1}^n(M_i)_{f_k},\qquad
a\longmapsto\bigl(\psi_{ki}(a/1)\bigr)_{i=1}^n,
$$
where $a/1$ is the image of $a$ in $(M_k)_{f_i}$ for the $i$th entry.
The cocycle for $(k,j,i)$ says
$\psi_{ji}\psi_{kj}=\psi_{ki}$ after localizing at $f_if_j$,
so every defining difference for $H_k$ vanishes on this tuple.
Thus $\eta_k$ takes values in $H_k$.
Projection onto the $k$th summand, followed by
$(M_k)_{f_k}\to M_k$, defines $\pi_k:H_k\to M_k$.
The identity $\psi_{kk}=\mathrm{id}$ proves $\pi_k\eta_k=\mathrm{id}$.
For $z=(z_i)\in H_k$, its $(i,k)$ equation identifies $z_i$
with $\psi_{ki}(\pi_k(z)/1)$; hence $\eta_k\pi_k=\mathrm{id}$.

\medskip\noindent
There is also an explicit fraction representing this inverse in $M_{f_k}$.
Write $\psi_{ki}(a/1)=b_i/f_k^{r_i}$ with $b_i\in M_i$ and $r_i\geq0$.
Choose $N\geq r_i$ for all $i$, and put $b'_i=f_k^{N-r_i}b_i$.
For each pair $(i,j)$ the element
$d_{ij}=b'_i/1-\psi_{ji}(b'_j/1)$ of $(M_i)_{f_j}$
vanishes after localization at $f_k$.
Choose $s_{ij}\geq0$ with $f_k^{s_{ij}}d_{ij}=0$, and choose
$L\geq s_{ij}$ for all pairs. The tuple
$m=(f_k^Lb'_i)_{i=1}^n$ belongs to $M$, and the fraction
$m/f_k^{N+L}$ maps to $\eta_k(a)$.
Independence of these choices follows from the injectivity of the
displayed kernel comparison, since their images are the same tuple.
Consequently the original projection induces $M_{f_k}\cong M_k$.
Finally, for $m=(m_i)\in M$, its $(j,i)$ equation reads
$m_j/1=\psi_{ij}(m_i/1)$, which is the asserted compatibility.
The proof uses no cover assumption on the $f_i$ and changes none of them.
If $n=0$, then $M=0$ and there are no indexed assertions to prove.
\end{proof}


\subsection{Editorial treatment of Algebra lines 4592--4604}

\hypertarget{algebra-context-023}{}

\noindent\href{https://github.com/stacks/stacks-project/blob/a04446e57ec1fbc252a871afcec7752fb2807b14/algebra.tex\#L4592}{Original source, lines 4592--4604}. Complete original and reviewed TeX passages accompany this note. Every intervention is separately recorded; the following treatment is editorial.


\begin{lemma}
\label{editorial-source-lemma-total-ring-fractions}
Let $R$ be a ring.
Let $S \subset R$ be a multiplicative subset consisting of nonzerodivisors.
Then $Q(R) \cong Q(S^{-1}R)$.
In particular $Q(R) \cong Q(Q(R))$.
\end{lemma}

\begin{proof}
Let $T\subset R$ be the set of nonzerodivisors.
Since $S\subset T$, the map $R\to S^{-1}R$ is injective.
An element $x=r/f$ of $S^{-1}R$, with $r\in R$ and $f\in S$,
is a nonzerodivisor if and only if $r\in T$.
Indeed, if $x$ is a nonzerodivisor and $ra=0$ in $R$, then
$x(a/1)=0$, hence $a/1=0$, and injectivity gives $a=0$.
Conversely, if $r\in T$ and $x(a/s)=0$, then
$ura=0$ for some $u\in S$. Both $u$ and $r$ are nonzerodivisors,
so $a=0$ and $a/s=0$.
The two ring maps are
$$
Q(R)\longrightarrow Q(S^{-1}R),\qquad
\frac{r}{t}\longmapsto\frac{r/1}{t/1},
$$
and
$$
Q(S^{-1}R)\longrightarrow Q(R),\qquad
\frac{r/s}{a/u}\longmapsto\frac{ru}{sa},
\quad s,u\in S,\quad a\in T.
$$
The first is induced by $R\to S^{-1}R$ because each $t\in T$
remains a nonzerodivisor. The second is induced by $S^{-1}R\to Q(R)$:
it sends every nonzerodivisor $a/u$ to a unit with inverse $u/a$.
Thus both formulas are well-defined by the localization universal
properties. Their composites fix each element of $R$ and every
adjoined inverse; explicitly, $r/t$ returns to $r/t$, and
$(r/s)/(a/u)$ returns to $(ru/1)/(sa/1)$, the original fraction.
Taking $S=T$ gives the final assertion, with the same maps.
\end{proof}


\subsection{Editorial treatment of Algebra lines 4664--4695}

\hypertarget{algebra-context-024}{}

\noindent\href{https://github.com/stacks/stacks-project/blob/a04446e57ec1fbc252a871afcec7752fb2807b14/algebra.tex\#L4664}{Original source, lines 4664--4695}. Complete original and reviewed TeX passages accompany this note. Every intervention is separately recorded; the following treatment is editorial.


\begin{lemma}
\label{editorial-source-lemma-irreducible}
Let $R$ be a ring.
\begin{enumerate}
\item For a prime $\mathfrak p \subset R$ the closure
of $\{\mathfrak p\}$ in the Zariski topology is $V(\mathfrak p)$.
In a formula $\overline{\{\mathfrak p\}} = V(\mathfrak p)$.
\item The irreducible closed subsets of $\Spec(R)$ are
exactly the subsets $V(\mathfrak p)$, with $\mathfrak p \subset R$
a prime.
\item The irreducible components (see Topology,
Definition \ref{topology-definition-irreducible-components})
of $\Spec(R)$ are  exactly the subsets $V(\mathfrak p)$,
with $\mathfrak p \subset R$ a minimal prime.
\end{enumerate}
\end{lemma}

\begin{proof}
Note that if $ \mathfrak p \in V(I)$, then
$I \subset \mathfrak p$. Hence,
clearly $\overline{\{\mathfrak p\}} = V(\mathfrak p)$.
In particular $V(\mathfrak p)$ is the closure of
a singleton and hence irreducible.
The second assertion implies the third.
To show the second, let
$V(I) \subset \Spec(R)$ with $I$ a radical ideal.
If $I=R$, then $V(I)=\emptyset$, which is not irreducible.
We may therefore suppose that $I$ is proper.
If $I$ is not prime, then choose $a, b\in R$, $a, b\not \in I$
with $ab\in I$. In this case $V(I, a) \cup V(I, b) = V(I)$,
but neither $V(I, b) = V(I)$ nor $V(I, a) = V(I)$, by
Lemma \ref{editorial-source-lemma-Zariski-topology}. Hence $V(I)$ is not
irreducible.
\end{proof}


\subsection{Editorial treatment of Algebra lines 4863--4887}

\hypertarget{algebra-context-025}{}

\noindent\href{https://github.com/stacks/stacks-project/blob/a04446e57ec1fbc252a871afcec7752fb2807b14/algebra.tex\#L4863}{Original source, lines 4863--4887}. Complete original and reviewed TeX passages accompany this note. Every intervention is separately recorded; the following treatment is editorial.


\begin{example}
\label{editorial-source-example-spec-Zx}
In this example we describe $X = \Spec(\mathbf{Z}[x])$.
Fix $\mathfrak p \in X$.
Let $\phi : \mathbf{Z} \to \mathbf{Z}[x]$ and notice
that $\phi^{-1}(\mathfrak p) \in \Spec(\mathbf{Z})$.
If $\phi^{-1}(\mathfrak p) = (q)$ for $q$ a prime number $q > 0$,
then $\mathfrak p$ corresponds to a prime in $(\mathbf{Z}/(q))[x]$.
The zero prime of that polynomial ring gives $\mathfrak p=(q)$.
Every nonzero prime is generated by an irreducible polynomial
$\overline{f_q}\in(\mathbf{Z}/(q))[x]$.
Choose any coefficient lift $f_q\in\mathbf{Z}[x]$ of the same degree.
The inverse image of $(\overline{f_q})$ is $(q,f_q)$: if the reduction
of $F\in\mathbf Z[x]$ is $\overline A\,\overline{f_q}$, choose a lift
$A\in\mathbf Z[x]$; then $F-Af_q\in q\mathbf Z[x]$.
The converse containment follows by reduction.
Thus the primes over $(q)$ are $(q)$ and the ideals $(q,f_q)$ with
$\overline{f_q}$ irreducible. The lift $f_q$ itself need not be
irreducible in $\mathbf Z[x]$: for $q=5$, the degree-one reduction
of $2x+2$ is irreducible, while $2x+2=2(x+1)$ in $\mathbf Z[x]$.
Now assume that $\phi^{-1}(\mathfrak p)=(0)$.
In this case, if $\mathfrak p = (0)$ there is nothing more to prove. Otherwise, $\mathfrak p$ contains nonconstant polynomials
and it contains no nonzero integer. Factoring any nonzero member
in the unique factorization domain $\mathbf Z[x]$ and using primality
gives an irreducible $f\in\mathfrak p$. It is nonconstant because
$\mathfrak p\cap\mathbf Z=(0)$, and it is primitive because a
nonunit integer content would be a nontrivial factor.
Every other nonconstant irreducible $g\in\mathfrak p$ is likewise primitive.
Gauss' lemma makes $f$ and $g$ irreducible in $\mathbf Q[x]$.
Moreover, if $g=(d/e)f$ with coprime nonzero integers $d,e$,
the equality $eg=df$ has contents $|e|$ and $|d|$, respectively.
Thus $|e|=|d|=1$, so association over $\mathbf Q[x]$ implies
association over $\mathbf Z[x]$ for these original polynomials.
If $f$ and $g$ were not associates, the Euclidean algorithm in
$\mathbf Q[x]$ would give $1=af+bg$ with $a,b\in\mathbf Q[x]$.
Choose a positive common denominator $m$ of every coefficient of
$a$ and $b$, and put $\bar a=ma$, $\bar b=mb$ in $\mathbf Z[x]$.
Then $m=\bar a f+\bar b g\in\mathfrak p\cap\mathbf Z$, a contradiction.
Therefore all irreducibles belonging to $\mathfrak p$ are associates of $f$.
For every nonzero $F\in\mathfrak p$, a factorization of $F$ has
an irreducible factor in $\mathfrak p$ by primality, and that factor
is a unit multiple of $f$. Thus $f$ divides every $F\in\mathfrak p$,
while $f\in\mathfrak p$, proving $\mathfrak p=(f)$.
Conversely, every nonconstant irreducible $f\in\mathbf Z[x]$
generates a prime ideal, because $\mathbf Z[x]$ is a unique
factorization domain. Its contraction to $\mathbf Z$ is zero:
a nonzero constant cannot be divisible by a positive-degree polynomial.
Together with the separately retained zero prime, this describes
all primes with zero contraction.
\end{example}


\subsection{Editorial treatment of Algebra lines 4889--4929}

\hypertarget{algebra-context-026}{}

\noindent\href{https://github.com/stacks/stacks-project/blob/a04446e57ec1fbc252a871afcec7752fb2807b14/algebra.tex\#L4889}{Original source, lines 4889--4929}. Complete original and reviewed TeX passages accompany this note. Every intervention is separately recorded; the following treatment is editorial.


\begin{example}
\label{editorial-source-example-spec-kxy}
In this example we describe $X = \Spec(k[x, y])$
when $k$ is an arbitrary field.
Clearly $(0)$ is prime, and any principal ideal generated by an
irreducible polynomial will also be a prime since $k[x, y]$ is a
unique factorization domain. Now assume $\mathfrak p$ is an
element of $X$ that is not principal. Since $k[x, y]$ is a
Noetherian UFD, the prime ideal $\mathfrak p$ can be generated
by a finite number of irreducible polynomials $(f_1, \ldots, f_n)$.
Now, I claim that if $f, g$ are irreducible polynomials in $k[x, y]$
that are not associates, then $(f, g) \cap k[x] \neq 0$. To do this,
if either $f$ or $g$ lies in $k[x]$, that nonzero polynomial
itself belongs to $(f,g)\cap k[x]$ and proves the assertion.
Otherwise both have positive degree in $y$. Each is primitive as
a polynomial in $y$ over $k[x]$, since a nonunit content would
factor the original irreducible polynomial in $k[x,y]$.
Gauss' lemma therefore makes both irreducible in $k(x)[y]$.
In this case it is enough to show that $f$ and $g$ are relatively prime when
viewed in $k(x)[y]$. In this case, $k(x)[y]$ is a Euclidean domain,
so by applying the Euclidean algorithm and clearing denominators, we
obtain $p = af + bg$ for $0 \ne p \in k[x]$ and $a, b \in k[x, y]$. Thus, assume this is not
the case, that is, that some nonunit $h \in k(x)[y]$ divides both
$f$ and $g$. Then, by Gauss's lemma, for some $a, b \in k(x)$ we
have $ah | f$ and $bh | g$ for $ah, bh \in k[x, y]$. By
irreducibility, $ah = f$ and
$bh = g$ (since $h \notin k(x)$). So, back in $k(x)[y]$, $f, g $
are associates, as $\frac{a}{b} g = f$. Since
$k(x)$ is the fraction field of $k[x]$, we can write $g = \frac{r}{s} f $
for elements $r , s \in k[x]$ sharing no common factors. This
implies that $sg = rf$ in $k[x, y]$ and so $s$ must divide $f$
since $k[x,y]$ is a UFD. If $s$ is a unit in $k[x]$, write
$s=c\in k^\times$. Then $g=(r/c)f$; irreducibility of $g$
forces $r/c$ to be a unit, contradicting that $f$ and $g$ are
not associates. If $s$ is not a unit, the divisibility $s\mid f$
and irreducibility of $f$ give $f=us$ for a unit $u\in k^\times$.
This would put $f$ in $k[x]$, contrary to its positive $y$-degree.
Both alternatives contradict the existence of the common divisor.
Thus $f$ and $g$ are relatively prime in $k(x)[y]$, and the
cleared Bezout identity gives a nonzero polynomial in $(f,g)\cap k[x]$.
To distinguish it from a selected irreducible factor below, denote
this resulting polynomial by $P$.
The irreducible generators of the nonprincipal prime $\mathfrak p$
include two which are not associates, since otherwise they generate
a principal ideal. Applying the argument to those two puts
$P\in\mathfrak p\cap k[x]$. It cannot be a nonzero constant,
because $\mathfrak p$ is proper. Write its full factorization as
$P=c\prod_{j=1}^d p_j$, with $c\in k^\times$ and irreducible
$p_j\in k[x]$, repeating factors according to their multiplicities.
Primality gives some $p_j\in\mathfrak p$; denote this chosen
factor by $p$. Thus $\mathfrak p$ contains the irreducible
polynomial $p\in k[x]\subset k[x,y]$.
By the above, $\mathfrak p$ corresponds to a prime in the ring
$k[x,y]/(p)\cong(k[x]/(p))[y]=k(\alpha)[y]$, where
$\alpha=x+(p)$ in the field $k[x]/(p)$.
The isomorphism sends $F(x,y)+(p)$ to $F(\alpha,y)$:
coefficientwise reduction is surjective and its kernel consists
exactly of the polynomials whose coefficients are multiples of $p$.
Every element of $k[x]/(p)$ is a polynomial in $\alpha$, and this
quotient is a field, so it is $k(\alpha)$. No leading coefficient
of the chosen $p$ is changed or omitted. This is a
PID, and so any prime ideal corresponds to $(0)$ or an
irreducible polynomial in $k(\alpha)[y]$. Thus, $\mathfrak p$
is of the form $(p)$ or $(p, f)$ where $f$ is a
polynomial in $k[x, y]$ that is irreducible in the quotient
$k[x, y]/(p)$.
\end{example}


\subsection{Editorial treatment of Algebra lines 4931--5083}

\hypertarget{algebra-context-027}{}

\noindent\href{https://github.com/stacks/stacks-project/blob/a04446e57ec1fbc252a871afcec7752fb2807b14/algebra.tex\#L4931}{Original source, lines 4931--5083}. Complete original and reviewed TeX passages accompany this note. Every intervention is separately recorded; the following treatment is editorial.


\begin{example}
\label{editorial-source-example-affine-open-not-standard}
Consider the ring
$$
R = \{ f \in \mathbf{Q}[z]\text{ with }f(0) = f(1) \}.
$$
Consider the map
$$
\varphi : \mathbf{Q}[A, B] \to R
$$
defined by $\varphi(A)=A(z)=z^2-z$ and
$\varphi(B)=B(z)=z^3-z^2=zA(z)$.
Put $F=A^3-B^2+AB$. The exact identity
$$
F(A(z),B(z))=A(z)^3-z^2A(z)^2+zA(z)^2
=A(z)^2\bigl(A(z)-z^2+z\bigr)=0
$$
shows $(F)\subset\Ker(\varphi)$.
Since the leading coefficient of $F$ as a polynomial in $B$ is $-1$,
division in $\mathbf Q[A][B]$ writes any polynomial as
$Q(A,B)F+C(A)+BD(A)$.
Suppose its remainder evaluates to zero, so
$C(A(z))+zA(z)D(A(z))=0$.
Substituting $1-z$ for $z$ fixes $A(z)$ and sends $B(z)$ to
$(1-z)A(z)$; thus also $C(A(z))+(1-z)A(z)D(A(z))=0$.
Subtracting gives $(2z-1)A(z)D(A(z))=0$ in the domain $\mathbf Q[z]$.
The first two factors are nonzero, hence $D(A(z))=0$.
Substitution $\mathbf Q[A]\to\mathbf Q[z]$, $A\mapsto z^2-z$,
is injective: a nonzero polynomial of degree $d$ has after substitution
degree $2d$ with the same nonzero leading coefficient.
Consequently $D=0$ and then $C=0$, proving $\Ker(\varphi)=(F)$.
This ideal is prime because the image is a subring of $\mathbf Q[z]$.
It follows that $F$ is irreducible: in a factorization $F=GH$, primality
puts one factor, say $G$, in $(F)$, and writing $G=FK$ and cancelling
the nonzero $F$ gives $KH=1$.
The surjectivity proved next therefore gives the induced isomorphism
$$
\overline\varphi:\mathbf Q[A,B]/(A^3-B^2+AB)\longrightarrow R.
$$

\medskip\noindent
To see that $\varphi$ is surjective, we must express any
$f\in R$ as a $\mathbf{Q}$-coefficient polynomial in $A(z) = z^2-z$
and $B(z) = z^3-z^2$. Note the relation $zA(z) = B(z)$. Let
$a = f(0) = f(1)$. Then $z(z-1)$ must divide $f(z)-a$, so we can write
$f(z) = z(z-1)g(z)+a = A(z)g(z)+a$.  If $\deg(g) < 2$, then
$g(z) = c_1z + c_0$ and $f(z) = A(z)(c_1z + c_0)+a = c_1B(z)+c_0A(z)+a$, so we
are done.  If $\deg(g)\geq 2$, then by the polynomial division
algorithm, we can write $g(z) = A(z)h(z)+b_1z + b_0$
($\deg(h)\leq\deg(g)-2$), so $f(z) = A(z)^2h(z)+b_1B(z)+b_0A(z)+a$.
Applying division to $h(z)$ and iterating, we obtain an expression
for $f(z)$ as a polynomial in $A(z)$ and $B(z)$; hence $\varphi$ is
surjective.

\medskip\noindent
Now let $a \in \mathbf{Q}$, $a \neq 0, \frac{1}{2}, 1$ and
consider
$$
R_a = \{ f \in \mathbf{Q}[z, \frac{1}{z-a}]\text{ with }f(0) = f(1)
\}.
$$
This is a finitely generated $\mathbf Q$-algebra. Here is a proof
with the stated generators. Put
$$
\delta=a^2-a,\qquad
w=z+\frac{\delta}{z-a},\qquad
D=\mathbf Q[z,1/(z-a)],\qquad
E=\mathbf Q[z^2-z,z^3-z,w].
$$
The values of the first two generators at $0$ and $1$ are $0$;
the values of $w$ are both $1-a$. Thus $E\subset R_a$.
Since $\delta\ne0$ and $(z-a)^{-1}=(w-z)/\delta$, we have
$D=E[z]$. The original relation $z^2=z+A(z)$, with $A(z)\in E$,
shows $D=E+Ez$ as an $E$-module.
Also $R_a=\mathbf Q+A(z)D$: if $f\in R_a$ and
$\lambda=f(0)=f(1)$, write $f-\lambda=P(z)/(z-a)^N$.
Its numerator vanishes at both $0$ and $1$, so $P(z)=z(z-1)H(z)$
for a polynomial $H$, giving $f-\lambda\in A(z)D$.
The reverse inclusion follows by evaluation.
Finally $A(z)\in E$ and
$zA(z)=z^3-z^2=(z^3-z)-(z^2-z)\in E$, hence
$A(z)D\subset E+E$ and $R_a\subset E$.
This proves $R_a=E$, with all three original generators retained.
We have the following inclusions:
$$
R\subset R_a\subset\mathbf{Q}[z, \frac{1}{z-a}], \quad
R\subset\mathbf{Q}[z]\subset\mathbf{Q}[z, \frac{1}{z-a}].
$$
Recall (Lemma \ref{editorial-source-lemma-spec-localization}) that for a ring $T$ and a
multiplicative subset $S\subset T$, the ring map $T \to S^{-1}T$
induces a map on spectra $\Spec(S^{-1}T) \to \Spec(T)$
which is a homeomorphism onto the subset
$$
\{\mathfrak p \in \Spec(T) \mid S \cap \mathfrak p = \emptyset\}
\subset \Spec(T).
$$
When $S = \{ 1, f, f^2, \ldots\}$ for some $f\in T$, this is
the open set $D(f)\subset \Spec(T)$.  We now verify a corresponding
property for the ring map $R \to R_a$: we will show that the map
$\theta : \Spec(R_a) \to \Spec(R)$ induced by inclusion
$R\subset R_a$ is a homeomorphism onto an open subset of
$\Spec(R)$ by verifying that $\theta$ is an injective local
homeomorphism.  We do so with respect to an open cover of
$\Spec(R_a)$ by two distinguished opens, as we now describe.
For any $r\in\mathbf{Q}$, let $\text{ev}_r : R \to \mathbf{Q}$ be the
homomorphism given by evaluation at $r$.  Note that for $r = 0$ and
$r = 1-a$, this can be extended to a homomorphism
$\text{ev}_r' : R_a \to \mathbf{Q}$ (the latter because $\frac{1}{z-a}$
is well-defined at $z = 1-a$, since $a\neq\frac{1}{2}$).  However,
$\text{ev}_a$ does not extend to $R_a$.  Write
$\mathfrak{m}_r = \Ker(\text{ev}_r)$. We have
$$
\mathfrak{m}_0 = (z^2-z, z^3-z),
$$
$$
\mathfrak{m}_a = ((z-1 + a)(z-a), (z^2-1 + a)(z-a)), \text{ and}
$$
$$
\mathfrak{m}_{1-a} = ((z-1 + a)(z-a), (z-1 + a)(z^2-a)).
$$
For every $r\in\mathbf Q$, the presentation already proved gives
$\mathfrak m_r=(A-r(r-1),B-r^2(r-1))$.
Indeed, evaluating $A$ and $B$ at those two constants gives a quotient
equal to $\mathbf Q$, so this ideal is exactly the evaluation kernel.
Write
$$
u=(z-1+a)(z-a)=A-\delta,\qquad
v=(z^2+z+2a-2)(z-a)=B+(2-a)A-2\delta.
$$
The displayed generators above are verified by the identities
$$
\begin{aligned}
z^3-z&=B+A,\\
(z^2-1+a)(z-a)&=(B-a\delta)+(1-a)(A-\delta),\\
(z-1+a)(z^2-a)&=(B-(1-a)\delta)+a(A-\delta).
\end{aligned}
$$
Each formula is reversible by subtracting the indicated multiple
of $A-\delta$, proving the asserted ideal equalities.
In particular $(u,v)=(A-\delta,B-a\delta)=\mathfrak m_a$,
since $v=(B-a\delta)+(2-a)(A-\delta)$.
Note that
$\mathfrak{m}_a$ is not in the image of $\theta$: we have
$$
(z^2 - z)^2(z - a)\left(\frac{a^2 - a}{z - a} + z\right) =
(z^2 - z)^2(a^2 - a) + (z^2 - z)^2(z - a)z
$$
The left hand side is in $\mathfrak m_a R_a$ because
$(z^2 - z)(z - a)$ is in $\mathfrak m_a$ and because
$(z^2 - z)(\frac{a^2 - a}{z - a} + z)$ is in $R_a$. Similarly
the element $(z^2 - z)^2(z - a)z$ is in $\mathfrak m_a R_a$
because $(z^2 - z)$ is in $R_a$ and $(z^2 - z)(z - a)$ is in $\mathfrak m_a$.
As $a \not \in \{0, 1\}$ we conclude that
$(z^2 - z)^2 \in \mathfrak m_a R_a$. Hence
no ideal $I$ of $R_a$ can satisfy $I \cap R = \mathfrak m_a$, as such
an $I$ would have to contain $(z^2 - z)^2$, which is in $R$ but not in
$\mathfrak m_a$. The distinguished open set
$D((z-1 + a)(z-a))\subset\Spec(R)$ is equal to the complement of
the closed set $\{\mathfrak{m}_a, \mathfrak{m}_{1-a}\}$.
Then check that $R_{(z-1 + a)(z-a)} = (R_a)_{(z-1 + a)(z-a)}$; calling
this localized ring $R'$, then, it follows that the map $R \to R'$
factors as $R \to R_a \to R'$.  By Lemma
\ref{editorial-source-lemma-spec-localization}, then, these maps express
$\Spec(R') \subset \Spec(R_a)$ and
$\Spec(R') \subset \Spec(R)$ as open subsets; hence
$\theta : \Spec(R_a) \to \Spec(R)$, when restricted to
$D((z-1 + a)(z-a))$, is a homeomorphism onto an open subset.
Similarly, $\theta$ restricted to
$D((z^2 + z + 2a-2)(z-a)) \subset \Spec(R_a)$ is a homeomorphism
onto the open subset $D((z^2 + z + 2a-2)(z-a)) \subset \Spec(R)$.
In $\Spec(R)$, the complement of this distinguished open consists
of $\mathfrak m_a$ and the one or two closed points associated to the
irreducible factors of $z^2+z+2a-2$; a repeated factor gives only one point.
In $\Spec(R_a)$ the point $\mathfrak m_a$ is absent.
The verification below identifies these zero sets as well.
Furthermore, the ideal in $R_a$ generated by the
elements $(z^2 + z + 2a-2)(z-a)$ and $(z-1 + a)(z-a)$ is all of $R_a$, so
these two distinguished open sets cover $\Spec(R_a)$. Hence in
order to show that $\theta$ is a homeomorphism onto
$\Spec(R)-\{\mathfrak{m}_a\}$, it suffices to show
that these one or two points can never equal $\mathfrak{m}_{1-a}$.
And this is indeed the case, since $1-a$ is a root of $z^2 + z + 2a-2$
if and only if $a = 0$ or $a = 1$, both of which do not occur.

\medskip\noindent
Here are the exact localization maps and the cover certificate.
As identities in the original rational-function field,
$$
uw=B+(a-1)\delta\in R,
\qquad
vw=zv+\delta(z^2+z+2a-2)\in R.
$$
The second numerator is polynomial and has value
$2\delta(a-1)$ at both $z=0$ and $z=1$, proving the stated membership.
Thus $w=(B+(a-1)\delta)/u$ belongs to $R_u$, and
$w=(zv+\delta(z^2+z+2a-2))/v$ belongs to $R_v$.
Since $R_a$ is generated by $A$, $B+A$ and $w$, inclusion induces
isomorphisms
$R_u\to(R_a)_u$ and $R_v\to(R_a)_v$.
Their inverse maps send the original $w$ to the respective fractions
just displayed, send each element of $R$ to itself, and send
$u^{-1}$ or $v^{-1}$ to that same inverse in the rational-function field.
These formulas prove both composites are identities; on the overlap
they agree as maps into $\mathbf Q(z)$.
The full unit-ideal identity is
$$
u(w-a+2)-v=\delta(2a-1).
$$
Its right hand side is a nonzero rational constant under the present
hypotheses. Division by this displayed constant gives
$1=((w-a+2)/(\delta(2a-1)))u-(1/(\delta(2a-1)))v$ in $R_a$.
Hence $D(u)$ and $D(v)$ cover $\Spec(R_a)$.
Their images are respectively $D_R(u)$ and $D_R(v)$, and those images
have union $\Spec(R)\setminus V(u,v)=\Spec(R)\setminus\{\mathfrak m_a\}$.
Contraction of prime ideals gives
$\theta^{-1}(D_R(u))=D_{R_a}(u)$ and the analogous equality for $v$.
If two source points have the same image, choose one of these target
opens containing it; both points belong to its inverse image and
are equal by that chart isomorphism. Thus $\theta$ is injective.
The two chart maps also prove surjectivity onto the stated union and
openness there, establishing the asserted homeomorphism.

\medskip\noindent
For the zero-set descriptions, the presentation of $R$ shows that the
only prime containing $A$ is $\mathfrak m_0$: modulo $A$ the relation
gives $B^2=0$. The values of $u$ and $v$ at this point are respectively
$-\delta$ and $-2\delta$, both nonzero.
Moreover $R_A=\mathbf Q[z,1/A(z)]$ inside $\mathbf Q(z)$, since
$z=B/A$ and the converse inclusion preserves $A(z)^{-1}$.
Thus all zeros under consideration can be read in this polynomial
localization. The two roots of $u$ are $a$ and $1-a$, distinct and
different from $0,1$. The other factors of $v$ come from
$q(z)=z^2+z+2a-2$, with $q(0)=2a-2\ne0$ and $q(1)=2a\ne0$.
Also $q(a)=a^2+3a-2\ne0$ for rational $a$: equality would give
$(2a+3)^2=17$, whereas $17$ is not a square of a rational number.
Consequently the irreducible factors of $q$ give one or two distinct
closed points in addition to $\mathfrak m_a$ in the target spectrum.
Finally $q(1-a)=a^2-a=\delta\ne0$, so none is $\mathfrak m_{1-a}$.
Deleting $\mathfrak m_a$ through the proved chart correspondence
gives exactly the claimed source zero set.

\medskip\noindent
Despite this homeomorphism which mimics the behavior of a
localization at an element of $R$, while
$\mathbf{Q}[z, \frac{1}{z-a}]$ is the localization of $\mathbf{Q}[z]$
at the element $z-a$, the ring $R_a$ is {\it not} a
localization of $R$: Any proper localization $S^{-1}R$ results in more
units than the original ring $R$.  The units of $R$ are
$\mathbf{Q}^\times$, the units of $\mathbf{Q}$.  In fact, it is
easy to see that the units of $R_a$ are $\mathbf{Q}^*$.
Namely, the units of $\mathbf{Q}[z, \frac{1}{z - a}]$ are
$c (z - a)^n$ for $c \in \mathbf{Q}^*$ and $n \in \mathbf{Z}$
and such a unit belongs to $R_a$ if and only if
$c(-a)^n=c(1-a)^n$. Its inverse then has equal values too, so
this is exactly the condition for it to be a unit of $R_a$.
For $n\ne0$ the equality implies
$\bigl((-a)/(1-a)\bigr)^n=1$.
A rational number with a nonzero integer power equal to $1$ is $1$ or $-1$:
writing it as a quotient of coprime integers shows that their absolute
values must agree. The value $1$ would give $-a=1-a$, which is impossible;
the value $-1$ gives $a=1/2$, which is excluded here.
Thus $n=0$ and $R_a^\times=\mathbf Q^\times=R^\times$.
The ambient description of units follows as well from polynomial
factorization: if $P(z)/(z-a)^r$ has inverse $Q(z)/(z-a)^s$, then
$P(z)Q(z)=(z-a)^{r+s}$, so $P(z)=c(z-a)^j$ and the original unit
is $c(z-a)^{j-r}$, with the coefficient and exponent retained.
Moreover $R_a\ne R$: the element
$w=z+(a^2-a)/(z-a)$ has a pole at $z=a$, since its numerator
$z(z-a)+(a^2-a)$ has nonzero value $a^2-a$ there, whereas all
elements of $R$ are polynomials.
If the inclusion $R\to R_a$ were a localization at a multiplicative
subset $S$, every $s\in S$ would become a unit of $R_a$ and hence
would already be a nonzero rational constant, a unit of $R$.
Then $S^{-1}R=R$ under the original map, contradicting $w\notin R$.
Therefore there is no $R$-algebra isomorphism $S^{-1}R\cong R_a$
compatible with the specified inclusion.

\medskip\noindent
We used the fact that $a\neq 0, 1$ to ensure that
$\frac{1}{z-a}$ makes sense at $z = 0, 1$.  We used the fact that
$a\neq 1/2$ in a few places: (1) In order to be able to talk about
the kernel of $\text{ev}_{1-a}$ on $R_a$, which ensures that
$\mathfrak{m}_{1-a}$ is a point of $R_a$ (i.e., that $R_a$ is
missing just one point of $R$). (2) At the end in order to conclude
that $(z-a)^n$ can belong to $R_a$ only for $n = 0$; indeed, if
$a = 1/2$, then it belongs to $R_a$ whenever $n$ is even. Hence
there are indeed more units in the ring defined by the same formula
at $a=1/2$. In this exceptional case we can decide the localization
question exactly. Set $h=(z-1/2)^2$. Its values at $0$ and $1$ are
both $1/4$, so $h\in R$ and $h^{-1}\in R_{1/2}$.
The original inclusion induces
$$
R_h\longrightarrow
R_{1/2}=\{f\in\mathbf Q[z,1/(z-1/2)]\mid f(0)=f(1)\}.
$$
It is injective because both rings are subrings of $\mathbf Q(z)$.
For surjectivity, write any $f\in R_{1/2}$ as
$P(z)/(z-1/2)^N$ and choose $m\geq0$ with $2m\geq N$.
Then $h^mf=P(z)(z-1/2)^{2m-N}$ is a polynomial, and its values at
$0$ and $1$ are equal because $h(0)=h(1)=1/4$ and $f(0)=f(1)$.
Thus $h^mf\in R$ and $f=(h^mf)/h^m\in R_h$.
This proves $R_{1/2}=R_{(z-1/2)^2}$ via the specified map.
The same endpoint calculation gives its exact unit group
$R_{1/2}^{\times}=\{c(z-1/2)^{2j}\mid c\in\mathbf Q^{\times},\ j\in\mathbf Z\}$,
including negative exponents. Its spectrum is the standard open
$D_R((z-1/2)^2)=\Spec(R)\setminus\{\mathfrak m_{1/2}\}$:
its only zero lies in $D_R(A)$ and corresponds there to the factor
$z-1/2$, while its value at $\mathfrak m_0$ is $1/4\ne0$.
\end{example}


\subsection{Editorial treatment of Algebra lines 5182--5190}

\hypertarget{algebra-context-028}{}

\noindent\href{https://github.com/stacks/stacks-project/blob/a04446e57ec1fbc252a871afcec7752fb2807b14/algebra.tex\#L5182}{Original source, lines 5182--5190}. Complete original and reviewed TeX passages accompany this note. Every intervention is separately recorded; the following treatment is editorial.


\begin{example}
\label{editorial-source-example-oka-family-bound-cardinality}
Let $R$ be a ring.
Let $\kappa$ be an infinite cardinal.
The family of ideals which can be generated by at most $\kappa$ elements
is an Oka family. To verify this, suppose $(I:a)$ has a
generating family $(u_\lambda)_{\lambda\in L}$ and $(I,a)$ has
a generating family $(v_\mu)_{\mu\in P}$, with $|L|,|P|\leq\kappa$.
Write $v_\mu=b_\mu+c_\mu a$, with $b_\mu\in I$ and $c_\mu\in R$.
Then $(I,a)=(a,b_\mu\mid\mu\in P)$.
For $x\in I$, an expression
$x=ra+\sum_{\mu\in P_0}s_\mu b_\mu$ with $P_0$ finite gives
$ra\in I$, hence $r\in(I:a)$. Writing
$r=\sum_{\lambda\in L_0}t_\lambda u_\lambda$ with $L_0$ finite yields
$$
x=\sum_{\lambda\in L_0}t_\lambda a u_\lambda+
\sum_{\mu\in P_0}s_\mu b_\mu.
$$
Thus $I=(a u_\lambda,b_\mu\mid\lambda\in L,\mu\in P)$, with
at most $\kappa+\kappa=\kappa$ generators. The ideal $R$ is
generated by $1$, so it belongs to this family as well.
\end{example}


\subsection{Editorial treatment of Algebra lines 5237--5262}

\hypertarget{algebra-context-029}{}

\noindent\href{https://github.com/stacks/stacks-project/blob/a04446e57ec1fbc252a871afcec7752fb2807b14/algebra.tex\#L5237}{Original source, lines 5237--5262}. Complete original and reviewed TeX passages accompany this note. Every intervention is separately recorded; the following treatment is editorial.


\begin{lemma}
\label{editorial-source-lemma-cohen}
Let $R$ be a ring.
\begin{enumerate}
\item An ideal $I \subset R$ maximal with respect to not being
finitely generated is prime.
\item If every prime ideal of $R$ is
finitely generated, then
every ideal of $R$ is finitely generated\footnote{Later we will say
that $R$ is Noetherian.}.
\end{enumerate}
\end{lemma}

\begin{proof}
The first assertion is an immediate consequence of
Example \ref{algebra-example-oka-family-finitely-generated} and
Proposition \ref{algebra-proposition-oka}. For the second,
suppose that there exists an ideal $I \subset R$ which is not finitely
generated. In the nonempty set of such ideals, the empty chain
has the chosen $I$ as an upper bound. The union of a nonempty totally ordered chain $\left\{I_\alpha\right\}$
of ideals that are not finitely generated is not finitely generated;
indeed, if $I = \bigcup I_\alpha$ were generated by
$a_1, \ldots, a_n$, then all the generators would belong to some
$I_\alpha $ and would consequently generate it.
By Zorn's lemma, there is an ideal maximal with respect to being not finitely
generated. By the first part this ideal is prime.
\end{proof}


\subsection{Editorial treatment of Algebra lines 5264--5286}

\hypertarget{algebra-context-030}{}

\noindent\href{https://github.com/stacks/stacks-project/blob/a04446e57ec1fbc252a871afcec7752fb2807b14/algebra.tex\#L5264}{Original source, lines 5264--5286}. Complete original and reviewed TeX passages accompany this note. Every intervention is separately recorded; the following treatment is editorial.


\begin{lemma}
\label{editorial-source-lemma-primes-principal}
Let $R$ be a ring.
\begin{enumerate}
\item An ideal $I \subset R$ maximal with respect to not being
principal is prime.
\item If every prime ideal of $R$ is principal, then
every ideal of $R$ is principal.
\end{enumerate}
\end{lemma}

\begin{proof}
The first part follows from
Example \ref{algebra-example-oka-family-principal} and
Proposition \ref{algebra-proposition-oka}.
For the second, suppose that there exists an ideal $I \subset R$
which is not principal. In the nonempty set of nonprincipal ideals,
the empty chain has the chosen $I$ as an upper bound.
The union of a nonempty totally ordered chain
$\left\{I_\alpha\right\}$ of ideals that are not principal is not principal;
indeed, if $I = \bigcup I_\alpha$ were generated by
$a$, then $a$ would belong to some $I_\alpha $ and $a$ would generate it.
By Zorn's lemma, there is an ideal maximal with respect to not being
principal. This ideal is necessarily prime by the first part.
\end{proof}


\subsection{Editorial treatment of Algebra lines 5288--5306}

\hypertarget{algebra-context-031}{}

\noindent\href{https://github.com/stacks/stacks-project/blob/a04446e57ec1fbc252a871afcec7752fb2807b14/algebra.tex\#L5288}{Original source, lines 5288--5306}. Complete original and reviewed TeX passages accompany this note. Every intervention is separately recorded; the following treatment is editorial.


\begin{lemma}
\label{editorial-source-lemma-characterize-domain}
Let $R$ be a ring.
\begin{enumerate}
\item An ideal maximal among the ideals which do not contain a
nonzerodivisor is prime.
\item If $R$ is nonzero and every nonzero prime ideal in $R$
contains a nonzerodivisor, then $R$ is a domain.
\end{enumerate}
\end{lemma}

\begin{proof}
Consider the set $S$ of nonzerodivisors. It is a multiplicative
subset of $R$. Hence any ideal maximal with respect to not intersecting
$S$ is prime, see
Lemma \ref{algebra-lemma-simple}.
For the second assertion, assume $R\ne0$. Then $0\notin S$,
so $(0)$ is an ideal disjoint from $S$. In the set of ideals disjoint
from $S$, every nonempty totally ordered chain has its union as an
upper bound: an element of that union belongs to some member, so
the union still misses $S$. The empty chain is bounded by $(0)$.
Zorn's lemma supplies a maximal member $\mathfrak p$, and the first
assertion shows it is prime. By hypothesis every nonzero prime
meets $S$, so this $\mathfrak p$ must equal $(0)$.
Thus $(0)$ is prime, and the nonzero ring $R$ is a domain.
\end{proof}


\subsection{Editorial treatment of Algebra lines 5308--5327}

\hypertarget{algebra-context-032}{}

\noindent\href{https://github.com/stacks/stacks-project/blob/a04446e57ec1fbc252a871afcec7752fb2807b14/algebra.tex\#L5308}{Original source, lines 5308--5327}. Complete original and reviewed TeX passages accompany this note. Every intervention is separately recorded; the following treatment is editorial.


\begin{remark}
\label{editorial-source-remark-cohen-bound-cardinality}
Let $R$ be a ring. Let $\kappa$ be an infinite cardinal.
By applying
Example \ref{editorial-source-example-oka-family-bound-cardinality} and
Proposition \ref{algebra-proposition-oka}
we see that any ideal maximal with respect to the property of not being
generated by $\kappa$ elements is prime. This result is not so
useful because there exists a ring for which every prime ideal
of $R$ can be generated by $\aleph_0$ elements, but some
ideal cannot. Namely, let $k$ be a field, let $T$ be a set whose
cardinality is greater than $\aleph_0$ and let
$$
R = k[\{x_n\}_{n \geq 1}, \{z_{t, n}\}_{t \in T, n \geq 0}]/
\bigl(x_n^2\ (n\geq1),\ z_{t,n}^2\ (t\in T,n\geq0),\
x_nz_{t,n}-z_{t,n-1}\ (t\in T,n\geq1)\bigr)
$$
This is a local ring with unique prime ideal
$\mathfrak m=(x_n\mid n\geq1)$. Indeed, every prime contains all
$x_n$ and $z_{t,n}$ because their squares vanish, and the quotient
by all these elements is the field $k$. The relations
$z_{t,n}=x_{n+1}z_{t,n+1}$ put every $z_{t,n}$ in $(x_n)$.
Every element of this ideal is nilpotent: it uses only finitely many
of the square-zero generators, and a product of more factors than
the number of those generators has a repeated factor and vanishes.
The augmentation to $k$ shows the ring is nonzero.
Thus $\mathfrak m$ is the unique prime and is countably generated.

\medskip\noindent
The ideal $J=(z_{t,n}\mid t\in T,n\geq0)$ in fact requires exactly
$|T|$ generators. To prove the lower bound we construct maps detecting
each original $T$-label. For $N\geq0$ let
$$
C_N=k[X_1,\ldots,X_N,W_N]/(X_1^2,\ldots,X_N^2,W_N^2).
$$
Define $C_N\to C_{N+1}$ by $X_i\mapsto X_i$ for $i\leq N$
and $W_N\mapsto X_{N+1}W_{N+1}$.
The square-free monomials
$X_E=\prod_{i\in E}X_i$ and $X_EW_N$, for $E\subset\{1,\ldots,N\}$,
form a $k$-basis of $C_N$.
Their images are respectively $X_E$ and $X_EX_{N+1}W_{N+1}$,
distinct members of the corresponding basis of $C_{N+1}$.
Thus all these maps are injective. Put $C=\colim_N C_N$.
Concretely, take classes of pairs $(N,c)$ with $c\in C_N$, identifying
two pairs when their images agree at a common later stage. Addition
and multiplication are computed at a common stage; compatibility of
the ring maps makes them independent of the stage and representatives.
The injectivity of every transition map makes each $C_N\to C$ injective.
In $C$ the original classes satisfy
$W_{n-1}=X_nW_n$ for $n\geq1$, and
$W_0=X_1\cdots X_NW_N\ne0$ at every stage and in the colimit.

\medskip\noindent
For each $t_0\in T$, the assignments
$$
\rho_{t_0}:R\longrightarrow C,\qquad
x_n\longmapsto X_n,\qquad
z_{t,n}\longmapsto
\begin{cases}W_n,&t=t_0,\\0,&t\ne t_0\end{cases}
$$
respect every displayed defining relation, including its index range,
and hence define a ring map. In particular
$\rho_{t_0}(z_{t_0,0})=W_0\ne0$.
Suppose a family $(j_\lambda)_{\lambda\in L}$ with $|L|<|T|$
generated $J$. Each $j_\lambda$ is a finite sum of multiples
of generators $z_{t,n}$, so the labels occurring in all these
finite expressions form a set $T'\subset T$ of cardinality at most
$\max(|L|,\aleph_0)<|T|$; here $T$ is uncountable.
Choose $t_0\in T\setminus T'$. Then $\rho_{t_0}$ kills every
$j_\lambda$, irrespective of its coefficients, but does not kill
$z_{t_0,0}\in J$, a contradiction.
This proves the lower bound. The displayed generating set for $J$
has cardinality $|T\times\mathbf Z_{\geq0}|=|T|$, proving equality.
In particular $J$ is not countably generated.
More generally, for any prescribed infinite cardinal $\kappa$, choosing
$|T|>\kappa$ gives a ring whose only prime ideal is countably
generated, hence generated by at most $\kappa$ elements, while $J$
cannot be generated by $\kappa$ elements.
\end{remark}


\subsection{Editorial treatment of Algebra lines 5329--5362}

\hypertarget{algebra-context-033}{}

\noindent\href{https://github.com/stacks/stacks-project/blob/a04446e57ec1fbc252a871afcec7752fb2807b14/algebra.tex\#L5329}{Original source, lines 5329--5362}. Complete original and reviewed TeX passages accompany this note. Every intervention is separately recorded; the following treatment is editorial.


\begin{example}
\label{editorial-source-example-noetherian-topology-on-spec}
\begin{reference}
Comment by Lukas Heger of November 12, 2020.
\end{reference}
Let $R$ be a ring and $X = \Spec(R)$. Since closed subsets of $X$ correspond to
radical ideals of $R$ (Lemma \ref{editorial-source-lemma-Zariski-topology}) we see that $X$ is a
Noetherian topological space if and only if we have ACC for radical ideals.
This holds if and only if every radical ideal is the radical of a finitely
generated ideal. For completeness, under ACC the set of radicals
of finitely generated subideals of a radical ideal $K$ has a maximal
member $\sqrt{F}$. Otherwise successive strict enlargements would
give an infinite ascending chain. For every $x\in K$ the ideal
$\sqrt{F+(x)}$ is another such member containing $\sqrt F$,
so maximality gives $x\in\sqrt F$ and hence $K=\sqrt F$.
Conversely, the union $K$ of an ascending sequence of radical ideals
$K_1\subset K_2\subset\cdots$ is radical: if a power of an element
lies in the union, it lies in one $K_i$, which then contains the element.
Write $K=\sqrt{(f_1,\ldots,f_n)}$. Each $f_j$ belongs to some
$K_i$, and a single $K_N$ contains all of them. Therefore
$K=\sqrt{(f_1,\ldots,f_n)}\subset K_N\subset K$, so the sequence
stabilizes. This proves ACC. Let
$$
\mathcal{F} = \{I \subset R \mid
\sqrt{I} = \sqrt{(f_1, \ldots, f_n)}\text{ for some }n
\text{ and }f_1, \ldots, f_n \in R\}.
$$
We now prove that $\mathcal F$ is an Oka family. First, $R\in\mathcal F$
using the generator $1$. For any $I$ and $a$, the product
$(I,a)(I:a)$ is contained in $I$. Conversely, for $x\in I$ we have
$x\in(I,a)$ and $x\in(I:a)$, since $xa\in I$; hence
$x^2\in(I,a)(I:a)$. These two containments prove the identity
$$
\sqrt{I} = \sqrt{(I, a)(I : a)}
$$
for every ideal $I\subset R$ and element $a\in R$.
If $\sqrt{(I,a)}=\sqrt{(f_1,\ldots,f_r)}$ and
$\sqrt{(I:a)}=\sqrt{(g_1,\ldots,g_s)}$, it follows that
$$
\sqrt I=\sqrt{(f_i g_j\mid 1\leq i\leq r,\ 1\leq j\leq s)}.
$$
Indeed, a prime contains a product of two ideals if and only if it
contains one of them, and replacing each ideal by another with the
same radical preserves precisely those primes. The radical is the
intersection of its containing primes by Lemma \ref{editorial-source-lemma-Zariski-topology}.
The displayed finite set of products therefore proves the Oka condition.
If the complement of $\mathcal F$ is nonempty, choose an ideal
in it as an upper bound for the empty chain. On the other hand,
if we have a nonempty totally ordered chain of ideals
$\{I_\alpha\}$ none of which are in $\mathcal{F}$, then the union
$I = \bigcup I_\alpha$ cannot be in $\mathcal{F}$ either. Otherwise
$\sqrt{I} = \sqrt{(f_1, \ldots, f_n)}$, then $f_i^e \in I$ for some $e$,
then $f_i^e \in I_\alpha$ for some $\alpha$ independent of $i$, then
$\sqrt{I_\alpha} = \sqrt{(f_1, \ldots, f_n)}$, contradiction.
Thus if the set of ideals not in $\mathcal{F}$ is nonempty, then it
has maximal elements and
exactly as in Lemma \ref{editorial-source-lemma-cohen} we conclude that $X$ is a
Noetherian topological space if and only if every prime ideal of $R$
is equal to $\sqrt{(f_1, \ldots, f_n)}$ for some $f_1, \ldots, f_n \in R$.
In detail, if every prime belongs to $\mathcal F$ but its complement
were nonempty, the chain argument and Zorn's lemma would give a maximal
excluded ideal, which is prime by Proposition \ref{algebra-proposition-oka},
a contradiction. Thus every ideal belongs to $\mathcal F$, giving
ACC for radical ideals by the equivalence just proved.
Conversely, ACC gives the finite radical presentation of every radical
ideal, in particular every prime. This proves the stated criterion
for the spectrum to be Noetherian.
\end{example}


\subsection{Editorial treatment of Algebra lines 5433--5451}

\hypertarget{algebra-context-034}{}

\noindent\href{https://github.com/stacks/stacks-project/blob/a04446e57ec1fbc252a871afcec7752fb2807b14/algebra.tex\#L5433}{Original source, lines 5433--5451}. Complete original and reviewed TeX passages accompany this note. Every intervention is separately recorded; the following treatment is editorial.


\begin{lemma}
\label{editorial-source-lemma-constructible}
Let $R$ be a ring. A subset of $\Spec(R)$ is constructible if and only
if it can be written as a finite union of subsets of the form
$D(f) \cap V(g_1, \ldots, g_m)$ for $f, g_1, \ldots, g_m \in R$.
\end{lemma}

\begin{proof}
By Lemma \ref{algebra-lemma-qc-open} the subset $D(f)$ and the complement of
$V(g_1, \ldots, g_m)$ are retro-compact open. Hence
$D(f) \cap V(g_1, \ldots, g_m)$ is a constructible subset and so is
any finite union of such. Conversely, by Topology, Definition
\ref{topology-definition-constructible}, a constructible subset has an
expression $T=\bigcup_{\alpha=1}^k(U_\alpha\cap V_\alpha^c)$ with
$U_\alpha,V_\alpha$ retrocompact open. By Lemma \ref{algebra-lemma-qc-open}, write
$U_\alpha=\bigcup_{i=1}^{n_\alpha}D(f_{\alpha i})$ and
$V_\alpha=\bigcup_{j=1}^{m_\alpha}D(g_{\alpha j})$. Then
$$
T=\bigcup_{\alpha=1}^k\bigcup_{i=1}^{n_\alpha}
\left(D(f_{\alpha i})\cap
V(g_{\alpha1},\ldots,g_{\alpha m_\alpha})\right).
$$
These are finite unions; an empty inner union is empty, and
$V()=\Spec(R)$ when $m_\alpha=0$.
\end{proof}


\subsection{Editorial treatment of Algebra lines 5453--5472}

\hypertarget{algebra-context-035}{}

\noindent\href{https://github.com/stacks/stacks-project/blob/a04446e57ec1fbc252a871afcec7752fb2807b14/algebra.tex\#L5453}{Original source, lines 5453--5472}. Complete original and reviewed TeX passages accompany this note. Every intervention is separately recorded; the following treatment is editorial.


\begin{lemma}
\label{editorial-source-lemma-constructible-is-image}
Let $R$ be a ring and let $T \subset \Spec(R)$
be constructible. Then there exists a ring map $R \to S$ of
finite presentation such that $T$ is the image of
$\Spec(S)$ in $\Spec(R)$.
\end{lemma}

\begin{proof}
Write
$T=\bigcup_{i=1}^k(D(f_i)\cap V(g_{i1},\ldots,g_{im_i}))$
as in Lemma \ref{editorial-source-lemma-constructible}, and set
$S_i=(R/(g_{i1},\ldots,g_{im_i}))_{f_i}$.
The canonical maps on spectra identify $\Spec(S_i)$ with the indicated
subset, by Lemmas \ref{algebra-lemma-standard-open} and \ref{algebra-lemma-spec-closed}.
For $k=0$, take $S=R/(1)$, whose spectrum is empty and which is finitely
presented over $R$. For $k>0$, take $S=\prod_{i=1}^kS_i$, with its
diagonal $R$-algebra structure. Lemma \ref{algebra-lemma-spec-product} identifies
its spectrum with the disjoint union of the component spectra; the
contraction map to $\Spec(R)$ therefore has image exactly $T$.

\medskip\noindent
Here is a finite presentation, including all component relations. Let
$C=R[E_1,\ldots,E_k,Y_1,\ldots,Y_k]/J$, where $J$ is generated by
$$
\begin{gathered}
E_i^2-E_i\quad(1\leq i\leq k),\qquad
E_iE_j\quad(1\leq i<j\leq k),\qquad
\sum_{i=1}^kE_i-1,\\
(1-E_i)Y_i,\quad f_iY_i-E_i\quad(1\leq i\leq k),\qquad
g_{ij}E_i\quad(1\leq i\leq k,\ 1\leq j\leq m_i).
\end{gathered}
$$
Map $C$ to $S$ by sending $E_i$ to the component idempotent and $Y_i$
to the element whose $i$th component is $1/f_i$ and whose other
components are zero. Every displayed relation vanishes. In $C$, the
orthogonal idempotents $E_i$ sum to $1$, so
$C\longrightarrow\prod_iE_iC$, $c\longmapsto(E_ic)_i$, has inverse
$(c_i)_i\longmapsto\sum_i c_i$. The ring $E_iC$, with identity $E_i$,
has the $R$-algebra map $r\mapsto rE_i$; all $g_{ij}$ vanish and
$f_iE_i$ has inverse $Y_i$. This defines a map $S_i\to E_iC$ sending
$\bar r/f_i^a$ to $rY_i^a$ for $a>0$, and $\bar r$ to $rE_i$.
The localization relations make the map well defined. The map in the
other direction sends $E_i$ to $1$ and $Y_i$ to $1/f_i$.
The composites fix the coefficient images and these generators; all
other $E_j,Y_j$ vanish in $E_iC$. Thus these maps are inverse, and
$C\cong S$ as $R$-algebras. This proves finite presentation without
any additional finiteness hypothesis on $R$.
\end{proof}


\subsection{Editorial treatment of Algebra lines 5554--5595}

\hypertarget{algebra-context-036}{}

\noindent\href{https://github.com/stacks/stacks-project/blob/a04446e57ec1fbc252a871afcec7752fb2807b14/algebra.tex\#L5554}{Original source, lines 5554--5595}. Complete original and reviewed TeX passages accompany this note. Every intervention is separately recorded; the following treatment is editorial.


\begin{lemma}
\label{editorial-source-lemma-characteristic-polynomial-prime}
Let $R \to A$ be a ring homomorphism.
Assume $A \cong R^{\oplus n}$ as an $R$-module.
Let $f \in A$. The multiplication map $m_f: A
\to A$ is $R$-linear and hence
has a characteristic polynomial
$P(T) = T^n + r_{n-1}T^{n-1} + \ldots + r_0 \in R[T]$.
For any prime
$\mathfrak{p} \in \Spec(R)$, $f$ acts nilpotently on $A
\otimes_R \kappa(\mathfrak{p})$ if and only if $\mathfrak p \in
V(r_0, \ldots, r_{n-1})$.
Moreover, the image of $\Spec(A_f)$ in $\Spec(R)$ is
$\bigcup_{i=0}^{n-1}D(r_i)$. This formula commutes with every ring
base change $R\to R'$, and $\Spec(A)\to\Spec(R)$ is open after every
such base change.
\end{lemma}

\begin{proof}
This follows quite easily once we prove that the characteristic
polynomial $\bar P(T) \in \kappa(\mathfrak p)[T]$ of the
multiplication map $m_{\bar f}: A \otimes_R \kappa(\mathfrak p) \to
A \otimes_R \kappa(\mathfrak p)$ which multiplies elements of $A
\otimes_R \kappa(\mathfrak p)$ by $\bar f$, the image of $f$ viewed in
$A\otimes_R\kappa(\mathfrak p)$, is just the image of $P(T)$ in
$\kappa(\mathfrak p)[T]$. Let $(a_{ij})$ be the matrix of the map
$m_f$ with entries in $R$, using a basis $e_1, \ldots, e_n$
of $A$ as an $R$-module.
Then, $A \otimes_R \kappa(\mathfrak p) \cong (R \otimes_R
\kappa(\mathfrak p))^{\oplus n} = \kappa(\mathfrak p)^n$, which is
an $n$-dimensional vector space over $\kappa(\mathfrak p)$ with
basis $e_1 \otimes 1, \ldots, e_n \otimes 1$. The image $\bar f = f
\otimes 1$, and so the multiplication map $m_{\bar f}$ has matrix
$(a_{ij} \otimes 1)$. Thus, the characteristic polynomial is
precisely the image of $P(T)$.

\medskip\noindent
For a linear endomorphism $N$ of a finite dimensional vector space,
nilpotence implies that the filtration
$0=\ker(N^0)\subset\ker(N)\subset\cdots\subset\ker(N^a)$
eventually reaches the whole space. A basis extending bases along this
filtration gives a strictly triangular matrix, since
$N(\ker(N^j))\subset\ker(N^{j-1})$; its characteristic polynomial
is $T^n$. Conversely, if the characteristic polynomial is $T^n$,
Cayley--Hamilton gives $N^n=0$ for $n>0$.
For $n=0$ the space is zero, its only endomorphism is zero and its
characteristic polynomial is $1=T^0$. Applying this to $m_{\bar f}$
gives nilpotence exactly when all the images of
$r_0,\ldots,r_{n-1}$ in $\kappa(\mathfrak p)$ vanish.
The kernel of $R\to\kappa(\mathfrak p)$ is $\mathfrak p$, proving
the first assertion, including $V()=\Spec(R)$ when $n=0$.

\medskip\noindent
Put $B=A\otimes_R\kappa(\mathfrak p)$. Multiplication by $\bar f$
is nilpotent if and only if $\bar f$ is nilpotent: one implication
follows by applying the power of the multiplication map to $1$, and
the other follows by multiplication by $\bar f^a=0$.
The localization $B_{\bar f}$ is zero if and only if
$\bar f^a=0$ for some $a\geq0$, by the equality criterion for
$1/1=0/1$. Also
$A_f\otimes_R\kappa(\mathfrak p)\cong B_{\bar f}$:
the map sends $(a/f^j)\otimes\lambda$ to
$(a\otimes\lambda)/\bar f^j$, and its inverse is induced by
$a\otimes\lambda\mapsto(a/1)\otimes\lambda$ and the inverse
$(1/f)\otimes1$ of $\bar f$. These maps fix the generators and
the inverse of $f$, so their composites are identities.
Lemma \ref{editorial-source-lemma-in-image} now gives
$$
\mathfrak p\in\operatorname{Im}(\Spec(A_f)\to\Spec(R))
\ \Longleftrightarrow\ B_{\bar f}\ne0
\ \Longleftrightarrow\ \mathfrak p\in\bigcup_{i=0}^{n-1}D(r_i).
$$
For $n=0$, $A=0$ and the union and image are both empty.

\medskip\noindent
For an arbitrary ring map $R\to R'$, the images of the chosen basis
give a basis of $A'=A\otimes_RR'$. Its multiplication matrix for
$f'=f\otimes1$ is the image $(a'_{ij})$ of $(a_{ij})$. The determinant
formula
$$
\det(TI-(a_{ij}))=
\sum_{\sigma\in S_n}\operatorname{sgn}(\sigma)
\prod_{i=1}^n(T\delta_{i,\sigma(i)}-a_{i,\sigma(i)})
$$
shows that every coefficient of its characteristic polynomial is the
image of the corresponding original coefficient, with all signs and
terms retained. Consequently the image of $\Spec(A'_{f'})$ is
$\bigcup_{i=0}^{n-1}D(r'_i)$, the inverse image of
$\bigcup_{i=0}^{n-1}D(r_i)$ in $\Spec(R')$.
Finally, the same image formula applies to every element of $A'$,
whether or not it comes from an element of $A$. Standard opens form
a basis of $\Spec(A')$, and arbitrary unions of their open images
are open. Hence $\Spec(A')\to\Spec(R')$ is open.
\end{proof}


\subsection{Editorial treatment of Algebra lines 5597--5643}

\hypertarget{algebra-context-037}{}

\noindent\href{https://github.com/stacks/stacks-project/blob/a04446e57ec1fbc252a871afcec7752fb2807b14/algebra.tex\#L5597}{Original source, lines 5597--5643}. Complete original and reviewed TeX passages accompany this note. Every intervention is separately recorded; the following treatment is editorial.


\begin{lemma}
\label{editorial-source-lemma-affineline-special}
Let $R$ be a ring. Let $f, g \in R[x]$ be polynomials.
Assume the leading coefficient of $g$ is a unit of $R$.
There exist elements $r_i\in R$, $i = 1\ldots, n$ such that
the image of $D(f) \cap V(g)$ in $\Spec(R)$ is
$\bigcup_{i = 1, \ldots, n} D(r_i)$.
\end{lemma}

\begin{proof}
Write $g = ux^d + a_{d-1}x^{d-1} + \ldots + a_0$, where
$d$ is the degree of $g$, and hence $u \in R^*$.
If $d=0$, then $g=u$ is a unit, so $V(g)$ and its image are empty;
the required union may be taken to have no terms. Suppose $d>0$ and
consider the original ring $A=R[x]/(g)$. It is finite free with basis
the images of $1,x,\ldots,x^{d-1}$. Indeed, if a polynomial has leading
term $bx^e$ with $e\geq d$, subtracting
$bu^{-1}x^{e-d}g$ cancels that term and lowers its degree. Iteration
gives a remainder of degree less than $d$, keeping $g$ and all its
coefficients unchanged. If a nonzero multiple $qg$ had degree less
than $d$, its leading coefficient would be the product of the nonzero
leading coefficient of $q$ with the unit $u$, hence nonzero, and its
degree would be $\deg(q)+d\geq d$, a contradiction. This proves both
spanning and independence. Consider multiplication by the image of
$f$ on $A$. This is an $R$-module map.
Hence we can let $P(T) \in R[T]$ be the characteristic polynomial
of this map. Write $P(T) = T^d + r_{d-1} T^{d-1} + \ldots + r_0$.
We claim that $r_0, \ldots, r_{d-1}$ have the desired property.
We will use below the property of characteristic polynomials
that
$$
\mathfrak p \in V(r_0, \ldots, r_{d-1})
\Leftrightarrow
\text{multiplication by }f\text{ is nilpotent on }
A \otimes_R \kappa(\mathfrak p).
$$
This was proved in Lemma \ref{editorial-source-lemma-characteristic-polynomial-prime}.

\medskip\noindent
Suppose $\mathfrak q\in D(f) \cap V(g)$, and let
$\mathfrak p = \mathfrak q \cap R$. Then there is a nonzero map
$A \otimes_R \kappa(\mathfrak p) \to \kappa(\mathfrak q)$ which
is compatible with multiplication by $f$.
And $f$ acts as a unit on $\kappa(\mathfrak q)$.
Thus we conclude $\mathfrak p \not \in  V(r_0, \ldots, r_{d-1})$.

\medskip\noindent
On the other hand, suppose that $r_i \not\in \mathfrak p$ for some
prime $\mathfrak p$ of $R$ and some $0 \leq i \leq d - 1$.
Then multiplication by $f$ is not nilpotent on the algebra
$A \otimes_R \kappa(\mathfrak p)$.
Hence there exists a prime ideal $\overline{\mathfrak q} \subset
A \otimes_R \kappa(\mathfrak p)$ not containing the image of $f$.
The inverse image of $\overline{\mathfrak q}$ in $R[x]$
is an element of $D(f) \cap V(g)$ mapping to $\mathfrak p$.
\end{proof}


\subsection{Editorial treatment of Algebra lines 5645--5709}

\hypertarget{algebra-context-038}{}

\noindent\href{https://github.com/stacks/stacks-project/blob/a04446e57ec1fbc252a871afcec7752fb2807b14/algebra.tex\#L5645}{Original source, lines 5645--5709}. Complete original and reviewed TeX passages accompany this note. Every intervention is separately recorded; the following treatment is editorial.


\begin{theorem}[Chevalley's Theorem]
\label{editorial-source-theorem-chevalley}
Suppose that $R \to S$ is of finite presentation.
The image of a constructible subset of
$\Spec(S)$ in $\Spec(R)$ is constructible.
\end{theorem}

\begin{proof}
Write $S = R[x_1, \ldots, x_n]/(f_1, \ldots, f_m)$.
We may factor $R \to S$ as $R \to R[x_1] \to R[x_1, x_2]
\to \ldots \to R[x_1, \ldots, x_{n-1}] \to S$. Hence
we may assume that $S = R[x]/(f_1, \ldots, f_m)$.
In this case we factor the map as $R \to R[x] \to S$,
and by Lemma \ref{algebra-lemma-closed-fp} we reduce to
the case $S = R[x]$. By Lemma \ref{editorial-source-lemma-constructible} it suffices
to show that if
$T = (\bigcup_{i = 1\ldots n} D(f_i)) \cap V(g_1, \ldots, g_m)$
for $f_i , g_j \in R[x]$ then the image in $\Spec(R)$ is
constructible. Since finite unions of constructible sets
are constructible, it suffices to deal with the case $n = 1$,
i.e., when $T = D(f) \cap V(g_1, \ldots, g_m)$.

\medskip\noindent
For $c \in R$, the following are disjoint decompositions of sets:
$$
\Spec(R) =
V(c) \amalg D(c) =
\Spec(R/(c)) \amalg \Spec(R_c),
$$
and correspondingly $\Spec(R[x]) =
V(c) \amalg D(c) = \Spec(R/(c)[x]) \amalg
\Spec(R_c[x])$. The intersection of $T = D(f) \cap V(g_1, \ldots, g_m)$
with each part still has the same shape, with $f$, $g_i$ replaced
by their images in $R/(c)[x]$, respectively $R_c[x]$.
Note that the image of $T$
in $\Spec(R)$ is the union of the image of
$T \cap V(c)$ and $T \cap D(c)$. Using Lemmas \ref{algebra-lemma-open-fp}
and \ref{algebra-lemma-closed-fp} it suffices to prove the images of both
parts are constructible in $\Spec(R/(c))$, respectively
$\Spec(R_c)$.

\medskip\noindent
We prove the assertion for
$T=D(f)\cap V(g_1,\ldots,g_m)$ over every ring by induction on
$$
N=\sum_{j=1}^m w(g_j),\qquad
w(g)=\begin{cases}0&g=0,\\ \deg(g)+1&g\ne0.\end{cases}
$$
All original polynomial labels are retained, including zero relations.
For $N=0$ all $g_j$ are zero, so $T=D(f)$ and
Lemma \ref{algebra-lemma-affineline-open} applies. If the coefficient ring is
zero, the spectrum is empty and the assertion also holds. For $N>0$
choose a nonzero $g_i$ of least degree among the nonzero relations and
write its full polynomial as
$g_i=c x^{d_i}+\sum_{a=0}^{d_i-1}c_{ia}x^a$, with $c\ne0$.

\medskip\noindent
Use the preceding decomposition into $R/(c)$ and $R_c$. Over $R/(c)$,
the weight of the image of $g_i$ strictly decreases, including when
$d_i=0$ and that image is zero. Every other weight is at most its
original value, so induction applies on this piece.
If $R_c=0$, its piece is empty. Otherwise the image of $c$ is a unit
and the image of $g_i$ retains degree $d_i$. If $g_i$ was the only
nonzero relation over $R$, Lemma \ref{editorial-source-lemma-affineline-special}
applies over $R_c$. If there was another nonzero relation, choose
one, say
$g_j=c_jx^{d_j}+\sum_{a=0}^{d_j-1}c_{ja}x^a$, where $d_j\geq d_i$
by the choice of $i$. In $R_c[x]$ replace only this relation by
$$
g'_j=g_j-(c_j/c)x^{d_j-d_i}g_i.
$$
Here every coefficient denotes its image in $R_c$; in particular
$c_j/c$ may be zero. The coefficient of $x^{d_j}$ cancels exactly,
so $g'_j=0$ or $\deg(g'_j)<d_j$. The other relations have weight
no larger than before base change, and the total weight is therefore
strictly less than the original $N$. Moreover,
$$
g'_j=g_j-(c_j/c)x^{d_j-d_i}g_i,\qquad
g_j=g'_j+(c_j/c)x^{d_j-d_i}g_i
$$
prove both inclusions between the original ideal and the ideal with
$g_j$ replaced by $g'_j$. Thus the two lists define exactly the same
subset $T$ over $R_c$, and induction proves its image constructible
there. In the notation $i=1,j=2$ this is precisely
$V(g_1,g'_2,g_3,\ldots,g_m)$, with every separator present.
The two induction base cases are a standard open and a standard open
intersected with the zero set of one polynomial having unit leading
coefficient, as in Lemmas \ref{algebra-lemma-affineline-open} and
\ref{editorial-source-lemma-affineline-special}. Finally,
Lemmas \ref{algebra-lemma-open-fp} and \ref{algebra-lemma-closed-fp} carry the two
constructible images back to $\Spec(R)$, and their finite union is
the desired image.
\end{proof}


\subsection{Editorial treatment of Algebra lines 5729--5767}

\hypertarget{algebra-context-039}{}

\noindent\href{https://github.com/stacks/stacks-project/blob/a04446e57ec1fbc252a871afcec7752fb2807b14/algebra.tex\#L5729}{Original source, lines 5729--5767}. Complete original and reviewed TeX passages accompany this note. Every intervention is separately recorded; the following treatment is editorial.


\begin{lemma}
\label{editorial-source-lemma-generic-finite-presentation}
Let $R \subset S$ be an inclusion of domains.
Assume that $R \to S$ is of finite type.
There exists a nonzero $f \in R$, and a nonzero $g \in S$
such that $R_f \to S_{fg}$ is of finite presentation.
\end{lemma}

\begin{proof}
We use induction on the number of generators of the original
$R$-algebra $S$. For no generators, $S=R$ and we take $f=g=1$.
For one generator write $S=R[x]/\mathfrak q$, where
$\mathfrak q\cap R=(0)$. If $\mathfrak q=(0)$, again take $f=g=1$.
Otherwise choose a nonzero polynomial $h\in\mathfrak q$ of least
degree and retain its full expression
$$
h=f x^d+a_{d-1}x^{d-1}+\cdots+a_0,
\qquad f\ne0,\quad f,a_0,\ldots,a_{d-1}\in R.
$$
Here $d\geq1$ since $R\to S$ is injective. Over $R_f$ its leading
coefficient $f$ is a unit, but the polynomial $h$ is unchanged.
Division by $h$, cancelling a leading term $b x^e$ by subtracting
$bf^{-1}x^{e-d}h$, gives a unique remainder of degree less than $d$.
Uniqueness holds because the leading coefficient of a nonzero
multiple of $h$ is a nonzero coefficient times the unit $f$.
The canonical map
$$
R_f[x]/(h)\longrightarrow S_f
$$
is surjective. If a remainder $v\in R_f[x]$ of degree less than $d$
maps to zero, choose $a\geq0$ with $f^av\in R[x]$.
Its image in $S_f$ is zero, so $f^b f^av$ maps to zero in $S$ for
some $b\geq0$. Since $S$ is a domain and $f\ne0$ in $S$, already
$f^av$ maps to zero in $S$. If $v\ne0$, the domain property of $R$
implies $f^av\ne0$ of degree less than $d$, contrary to the choice
of $h$. Thus $v=0$, proving injectivity and finite presentation.
This proves the one-generator case with the original $f$ and $g=1$.

\medskip\noindent
For $n>1$ let $x_1,\ldots,x_n$ be the original generators and
$S'=R[x_1,\ldots,x_{n-1}]\subset S$. By the induction hypothesis
there are nonzero $f\in R$ and $g\in S'$ for which
$R_f\to S'_{fg}$ is finitely presented. The extension
$S'_{fg}\subset S_{fg}$ is an inclusion of domains generated by
$x_n$. Apply the one-generator case to obtain nonzero
$f'\in S'_{fg}$ and $g'\in S_{fg}$ with
$$
(S'_{fg})_{f'}\longrightarrow (S_{fg})_{f'g'}
$$
finitely presented. Retain both elements and write their actual
fractions as $f'=a/(fg)^r$ and $g'=b/(fg)^s$, with
$a\in S'$, $b\in S$, $r,s\geq0$ and $a,b\ne0$.
Set $G=gab\in S$, which is nonzero. There are canonical inverse
$S$-algebra maps
$$
(S_{fg})_{f'g'}\ \longleftrightarrow\ S_{fG}.
$$
Indeed, in $S_{fG}$ the product $fgab$ is a unit, so each of
$f,g,a,b$ is a unit. Thus $fg,f',g'$ are units and the map from
the left ring exists. Conversely, in the left ring $fg$ and $f'g'$
are units, so $f,g,f',g'$ are units, and then
$a=f'(fg)^r$, $b=g'(fg)^s$ are units. Hence $fG=fgab$ is a unit
and the reverse map exists. The two composites fix $S$ and every
inverse adjoined in the indicated localizations, so they are identities.

\medskip\noindent
Finally $S'_{fg}\to(S'_{fg})_{f'}$ has the finite presentation
$S'_{fg}[z]/(f'z-1)$. Finite presentations compose: if
$P=A[X_1,\ldots,X_l]/(F_1,\ldots,F_v)$ and
$Q=P[Y_1,\ldots,Y_t]/(H_1,\ldots,H_w)$, lift each of the finitely
many coefficients in the $H_j$ to $A[X_1,\ldots,X_l]$ to obtain
$\widetilde H_j$. The canonical map identifies
$Q$ with $A[X_1,\ldots,X_l,Y_1,\ldots,Y_t]/(F_i,\widetilde H_j)$:
maps in both directions send every variable to its displayed class,
and their composites fix the coefficient ring and all variables.
Applying this twice proves that
$R_f\to(S_{fg})_{f'g'}\cong S_{fG}$ is finitely presented, with
the required elements $f\in R$ and $G\in S$.
\end{proof}


\subsection{Editorial treatment of Algebra lines 5769--5814}

\hypertarget{algebra-context-040}{}

\noindent\href{https://github.com/stacks/stacks-project/blob/a04446e57ec1fbc252a871afcec7752fb2807b14/algebra.tex\#L5769}{Original source, lines 5769--5814}. Complete original and reviewed TeX passages accompany this note. Every intervention is separately recorded; the following treatment is editorial.


\begin{lemma}
\label{editorial-source-lemma-characterize-image-finite-type}
Let $R \to S$ be a finite type ring map.
Denote $X = \Spec(R)$ and $Y = \Spec(S)$.
Write $f : Y \to X$ for the induced
map of spectra. Let $E \subset Y = \Spec(S)$ be a
constructible set.
If a point $\xi \in X$ is in $f(E)$, then
$\overline{\{\xi\}} \cap f(E)$ contains an open
dense subset of $\overline{\{\xi\}}$.
\end{lemma}

\begin{proof}
Let $\xi \in X$ be a point of $f(E)$. Choose a point $\eta \in E$
mapping to $\xi$. Let $\mathfrak p \subset R$ be the prime
corresponding to $\xi$ and let $\mathfrak q \subset S$ be the
prime corresponding to $\eta$. Consider the diagram
$$
\xymatrix{
\eta \ar[r] \ar@{|->}[d] & E \cap Y' \ar[r] \ar[d] &
Y' = \Spec(S/\mathfrak q) \ar[r] \ar[d] &
Y \ar[d] \\
\xi \ar[r] & f(E) \cap X' \ar[r] &
X' = \Spec(R/\mathfrak p) \ar[r] &
X
}
$$
By Lemma \ref{algebra-lemma-affine-map-quasi-compact} the set $E \cap Y'$
is constructible in $Y'$.
The original ring map induces an injection of domains
$R/\mathfrak p\hookrightarrow S/\mathfrak q$, still of finite type:
the original finite set of generators maps to a generating set of the
quotient. Let $\bar f:Y'\to X'$ be its map of spectra. The image
$\bar f(E\cap Y')$ is contained in $f(E)\cap X'$, so it suffices
to find a dense open subset of $X'$ contained in this smaller image.
The points $\xi,\eta$ are the generic points of $X',Y'$ respectively.
Lemma \ref{editorial-source-lemma-generic-finite-presentation} supplies nonzero
$a\in R/\mathfrak p$ and $b\in S/\mathfrak q$ with
$(R/\mathfrak p)_a\to(S/\mathfrak q)_{ab}$ finitely presented.
Thus $U=D(a)\subset X'$ and $V=D(ab)\subset Y'$ are dense opens,
contain $\xi$ and $\eta$, and $\bar f(V)\subset U$.
The set $E\cap Y'\cap V$ is constructible in $V$, by Topology,
Lemma \ref{topology-lemma-open-immersion-constructible-inverse-image}.
Chevalley's theorem (Theorem \ref{editorial-source-theorem-chevalley}) implies its
image $C$ is constructible in $U$ and contains $\xi$.
By Topology, Lemma \ref{topology-lemma-generic-point-in-constructible},
$C$ contains a dense open $U'\subset U$.
For completeness, a finite expression
$C=\bigcup_j(U_j\cap V_j^c)$ in the irreducible space $U$ has a
term containing $\xi$. Its open set $V_j$ cannot be nonempty:
every nonempty open of an irreducible space contains its generic
point. Hence $V_j=\emptyset$ and the nonempty open $U_j$ lies in
$C$. Such an open is dense in $U$. Since $U$ is open dense in $X'$,
$U'$ is open dense in $X'$ and is contained in $f(E)\cap X'$.
Under the original closed immersion
$X'=V(\mathfrak p)=\overline{\{\xi\}}\subset X$, this is exactly
the required assertion.
\end{proof}


\subsection{Editorial treatment of Algebra lines 5821--5869}

\hypertarget{algebra-context-041}{}

\noindent\href{https://github.com/stacks/stacks-project/blob/a04446e57ec1fbc252a871afcec7752fb2807b14/algebra.tex\#L5821}{Original source, lines 5821--5869}. Complete original and reviewed TeX passages accompany this note. Every intervention is separately recorded; the following treatment is editorial.


\begin{lemma}
\label{editorial-source-lemma-surjective-spec-radical-ideal}
Let $\varphi : R \to S$ be a ring map.
The following are equivalent:
\begin{enumerate}
\item The map $\Spec(S) \to \Spec(R)$ is surjective.
\item For any ideal $I \subset R$
the inverse image of $\sqrt{IS}$ in $R$ is equal to $\sqrt{I}$.
\item For any radical ideal $I \subset R$ the inverse image
of $IS$ in $R$ is equal to $I$.
\item For every prime $\mathfrak p$ of $R$ the inverse
image of $\mathfrak p S$ in $R$ is $\mathfrak p$.
\end{enumerate}
In this case the same is true after any base change: Given a ring map
$R \to R'$ the ring map $R' \to R' \otimes_R S$ has the equivalent
properties (1), (2), (3), and (4) as well.
\end{lemma}

\begin{proof}
For any ideal $J\subset S$, an element $x\in R$ belongs to
$\sqrt{\varphi^{-1}(J)}$ if and only if
$\varphi(x)^a\in J$ for some $a\geq1$. Thus
$\sqrt{\varphi^{-1}(J)}=\varphi^{-1}(\sqrt J)$.
Assume (1). For any ideal $I\subset R$, the inclusion
$\sqrt I\subset\varphi^{-1}(\sqrt{IS})$ follows by applying
$\varphi$ to a power lying in $I$. If $x\notin\sqrt I$, there is
a prime $\mathfrak p\supset I$ with $x\notin\mathfrak p$, by
Lemma \ref{editorial-source-lemma-Zariski-topology}. Choose a prime $\mathfrak q$
of $S$ over $\mathfrak p$. Then $IS\subset\mathfrak q$ but
$\varphi(x)\notin\mathfrak q$, so $\varphi(x)\notin\sqrt{IS}$.
This proves (2). If $I$ is radical, (2) gives
$I\subset\varphi^{-1}(IS)\subset\varphi^{-1}(\sqrt{IS})=I$,
proving (3), which immediately implies (4).
Finally assume (4). The quotient map $R/\mathfrak p\to S/\mathfrak pS$
is injective for every prime $\mathfrak p$. Inverting the nonzero
elements of the domain $R/\mathfrak p$ gives a nonzero ring:
otherwise some such element would map to zero in $S/\mathfrak pS$
by the equality criterion for $1/1=0/1$, contradicting injectivity.
This ring is $S\otimes_R\kappa(\mathfrak p)$, so
Lemma \ref{editorial-source-lemma-in-image} places $\mathfrak p$ in the image.
This proves (1), and hence the equivalence of all four properties.

\medskip\noindent
Assume (1) holds. Let $R \to R'$ be a ring map. Let
$\mathfrak p' \subset R'$ be a prime ideal lying over the prime
$\mathfrak p$ of $R$. To see that $\mathfrak p'$ is in the image
of $\Spec(R' \otimes_R S) \to \Spec(R')$ we have to show
that $(R' \otimes_R S) \otimes_{R'} \kappa(\mathfrak p')$ is not zero, see
Lemma \ref{editorial-source-lemma-in-image}.
Write $K=\kappa(\mathfrak p)$ and $K'=\kappa(\mathfrak p')$.
The given map induces an injection $\iota:K\hookrightarrow K'$:
its map $R/\mathfrak p\to R'/\mathfrak p'$ is injective and every
nonzero denominator remains nonzero in $K'$. There is an isomorphism
$$
(R'\otimes_RS)\otimes_{R'}K'
\longrightarrow (S\otimes_RK)\otimes_KK'
$$
sending $(r'\otimes s)\otimes\lambda$ to
$(s\otimes1)\otimes\overline{r'}\lambda$, where
$\overline{r'}$ denotes the image in $K'$.
Its inverse sends $(s\otimes a)\otimes\lambda$ to
$(1\otimes s)\otimes\iota(a)\lambda$.
The balancing relations over $R$, $R'$ and $K$ hold because the
coefficient maps into $K'$ agree; the maps preserve products and
units. Their composites are identities: in one direction move
$r'$ across the tensor product over $R'$, and in the other move
$a$ across the tensor product over $K$.
By (1) and Lemma \ref{editorial-source-lemma-in-image}, $S\otimes_RK$ is nonzero.
A $K$-basis containing its nonzero element $1$ gives a direct sum
of copies of $K$, whose tensor product with $K'$ is the corresponding
direct sum of copies of $K'$. It is nonzero. Thus every
$\mathfrak p'$ is in the image, proving (1) after base change and
therefore all four equivalent properties there.
\end{proof}


\subsection{Editorial treatment of Algebra lines 5920--5944}

\hypertarget{algebra-context-042}{}

\noindent\href{https://github.com/stacks/stacks-project/blob/a04446e57ec1fbc252a871afcec7752fb2807b14/algebra.tex\#L5920}{Original source, lines 5920--5944}. Complete original and reviewed TeX passages accompany this note. Every intervention is separately recorded; the following treatment is editorial.


\begin{lemma}
\label{editorial-source-lemma-image-dense-generic-points}
Let $R \to S$ be a ring map. The following are equivalent:
\begin{enumerate}
\item The kernel of $R \to S$ consists of nilpotent elements.
\item The minimal primes of $R$ are in the image of
$\Spec(S) \to \Spec(R)$.
\item The image of $\Spec(S) \to \Spec(R)$ is dense
in $\Spec(R)$.
\end{enumerate}
\end{lemma}

\begin{proof}
Write $\varphi:R\to S$ for the original map and $I=\ker\varphi$.
By Lemma \ref{editorial-source-lemma-Zariski-topology} and the power criterion for a
radical,
$$
\bigcap_{\mathfrak q\in\Spec(S)}\varphi^{-1}(\mathfrak q)
=\varphi^{-1}(\sqrt{(0)_S})=\sqrt I.
$$
Indeed $\varphi(x)$ is nilpotent exactly when $x^a\in I$ for some
$a\geq1$. A closed subset $V(J)\subset\Spec(R)$ contains the image
exactly when $J$ lies in this intersection. Hence its closure is
$V(\sqrt I)=V(I)$, proving (1) $\Leftrightarrow$ (3).
For (2) $\Rightarrow$ (3), every prime of $R$ contains a minimal
prime, and lies in the closure of that point, by
Lemma \ref{editorial-source-lemma-Zariski-topology}.
For (1) $\Rightarrow$ (2), the original map factors as
$$
R\longrightarrow R/I\lhook\joinrel\longrightarrow S.
$$
Every prime of $R$ contains $I$ under (1), and the quotient map
identifies $\Spec(R/I)$ homeomorphically with $\Spec(R)$ by
contraction, with inverse $\mathfrak p\mapsto\mathfrak p/I$.
These inverse order-preserving maps carry minimal primes to minimal
primes. The injection $R/I\hookrightarrow S$ hits all those points
by Lemma \ref{algebra-lemma-injective-minimal-primes-in-image}. Its contraction
composed with the quotient contraction is the contraction for
$\varphi$, proving (2) for the original map.
\end{proof}


\subsection{Editorial treatment of Algebra lines 5961--5975}

\hypertarget{algebra-context-043}{}

\noindent\href{https://github.com/stacks/stacks-project/blob/a04446e57ec1fbc252a871afcec7752fb2807b14/algebra.tex\#L5961}{Original source, lines 5961--5975}. Complete original and reviewed TeX passages accompany this note. Every intervention is separately recorded; the following treatment is editorial.


\begin{lemma}
\label{editorial-source-lemma-intersection-not-zero}
Let $A \subset B$ be an inclusion of domains inducing an algebraic extension
of fraction fields. If $J \subset B$ is a nonzero ideal, then $A \cap J$
is nonzero too. Thus the image of a proper closed subset of $\Spec(B)$
is not dense in $\Spec(A)$.
For a specified ideal $J$, the same conclusion $A\cap J\ne0$ already
follows if $J$ contains one nonzero element algebraic over
$\operatorname{Frac}(A)$; the entire fraction-field extension need
not be algebraic for this conclusion.
\end{lemma}

\begin{proof}
Let $x\in J$ be nonzero and algebraic over $K=\operatorname{Frac}(A)$.
Choose a nonzero polynomial $\sum_{i=0}^dc_iT^i\in K[T]$ of least
degree vanishing at $x$, with $c_d\ne0$. Then $d\geq1$.
Also $c_0\ne0$: otherwise dividing the equation by the nonzero
$x\in\operatorname{Frac}(B)$ gives a nonzero polynomial relation
of smaller degree. Write every coefficient as $c_i=b_i/t_i$ with
$b_i,t_i\in A$, $t_i\ne0$, keeping zero coefficients as well, and set
$$
D=\prod_{i=0}^dt_i,\qquad
a_i=b_i\prod_{\substack{0\leq j\leq d\\j\ne i}}t_j.
$$
Then $D\ne0$, $a_i=Dc_i$, $a_0,a_d\ne0$, and the complete equation is
$$
a_dx^d+a_{d-1}x^{d-1}+\cdots+a_0=0.
$$
Consequently $a_0=-\sum_{i=1}^d a_ix^i$ is a nonzero element of
$A\cap J$. This uses only the stated algebraicity of $x$.
For a proper closed subset $V(J)\subset\Spec(B)$, the domain property
ensures $J\ne0$, since $(0)$ is a prime. Its image is contained in
$V(a_0)\subset\Spec(A)$, a proper closed subset because
$a_0\ne0$ and $(0)$ is a prime of $A$. Thus that image is not dense.
\end{proof}


\subsection{Editorial treatment of Algebra lines 5980--6238}

\hypertarget{algebra-context-044}{}

\noindent\href{https://github.com/stacks/stacks-project/blob/a04446e57ec1fbc252a871afcec7752fb2807b14/algebra.tex\#L5980}{Original source, lines 5980--6238}. Complete original and reviewed TeX passages accompany this note. Every intervention is separately recorded; the following treatment is editorial.


\subsubsection{Noetherian rings}
\label{editorial-source-section-Noetherian}

\noindent
A ring $R$ is {\it Noetherian} if any ideal of $R$ is finitely generated.
This is equivalent to the ascending chain condition for ideals of $R$.
Indeed, the union of an ascending sequence of ideals is an ideal; if
it has a finite generating set, all those generators lie in one term,
and the sequence stabilizes there. Conversely, if an ideal is not
finitely generated, start with $(0)$ and successively choose an element
of the ideal outside the ideal generated by all the preceding choices.
This gives a strictly increasing sequence of ideals, contradicting
the ascending chain condition.
By
Lemma \ref{editorial-source-lemma-cohen}
it suffices to check that every prime ideal of $R$ is finitely generated.

\begin{lemma}
\label{editorial-source-lemma-Noetherian-permanence}
\begin{slogan}
Noetherian property is stable by passage to finite type extension
and localization.
\end{slogan}
Any finitely generated ring over a Noetherian ring
is Noetherian. Any localization of a Noetherian ring
is Noetherian.
\end{lemma}

\begin{proof}
The statement on localizations follows from the fact
that any ideal $J \subset S^{-1}R$ is of the form
$I \cdot S^{-1}R$. Any quotient $R/I$ of a Noetherian
ring $R$ is Noetherian because any ideal $\overline{J} \subset R/I$
is of the form $J/I$ for some ideal $I \subset J \subset R$.
Thus it suffices to show that if $R$ is Noetherian so
is $R[X]$. Suppose $J_1\subset J_2\subset\cdots$ is an ascending
chain of ideals in $R[X]$. For $i\geq1$ and $d\geq0$ define
$$
I_{i,d}=\{a\in R\mid aX^d+\sum_{e=0}^{d-1}b_eX^e\in J_i
\text{ for some }b_0,\ldots,b_{d-1}\in R\}.
$$
This contains zero, and addition and scalar multiplication of the
displayed polynomials show that it is an ideal. For a nonzero $a$,
the displayed polynomial has degree exactly $d$ and leading
coefficient $a$. Multiplication by $X^{d'-d}$ and the inclusion
$J_i\subset J_{i'}$ prove $I_{i,d}\subset I_{i',d'}$ when
$i\leq i'$ and $d\leq d'$.

\medskip\noindent
Put $K_d=\bigcup_{i\geq1}I_{i,d}$, an ideal because the union is
ascending in $i$. The sequence $K_0\subset K_1\subset\cdots$
stabilizes at some $d_0$. Each $K_d$ for $0\leq d\leq d_0$ has a
finite generating set. Every generator lies in some $I_{i,d}$;
taking the largest of these finitely many indices gives $i_d$ with
$I_{i_d,d}=K_d$ (take $i_d=1$ if the generating set is empty).
Let $i_0=\max_{0\leq d\leq d_0}i_d$. For $d\leq d_0$ and
$i\geq i_0$ we have $I_{i,d}=I_{i_0,d}=K_d$.
For $d>d_0$ the inclusions
$$
K_{d_0}=I_{i_0,d_0}\subset I_{i_0,d}\subset I_{i,d}
\subset K_d=K_{d_0}
$$
give the same equality. Thus the required $i_0$ works simultaneously
for every degree.

\medskip\noindent
For $i\geq i_0$ and $f\in J_i$, prove $f\in J_{i_0}$ by induction
on its degree, with $f=0$ treated separately. If $f\ne0$ has degree
$d$ and leading coefficient $a\ne0$, the equality
$I_{i,d}=I_{i_0,d}$ supplies $g\in J_{i_0}$ of degree $d$ with
the same coefficient $a$. Then $f-g\in J_i$ is zero or has degree
less than $d$. For $d=0$ it is zero; otherwise induction applies.
Hence $f-g\in J_{i_0}$ and $f\in J_{i_0}$. All $J_i$ with
$i\geq i_0$ are therefore equal, as required.
\end{proof}

\begin{lemma}
\label{editorial-source-lemma-Noetherian-power-series}
If $R$ is a Noetherian ring, then so is the formal power
series ring $R[[x_1, \ldots, x_n]]$.
\end{lemma}

\begin{proof}
The isomorphism
$R[[x_1,\ldots,x_{n+1}]]\cong R[[x_1,\ldots,x_n]][[x_{n+1}]]$
sends a coefficient family $(a_{\alpha,e})$ to
$\sum_{e\geq0}(\sum_{\alpha\in\mathbf Z_{\geq0}^n}a_{\alpha,e}
x_1^{\alpha_1}\cdots x_n^{\alpha_n})x_{n+1}^e$.
Its inverse reads off the same coefficients. For each fixed
multi-index, the Cauchy product has only finitely many contributing
pairs, so these mutually inverse maps preserve addition, multiplication
and the identity. It therefore suffices to prove that $R[[x]]$ is
Noetherian when $R$ is Noetherian. Let $I\subset R[[x]]$ be an ideal.
We have to show that $I$ is a finitely generated ideal.
For each nonnegative integer
$d$ denote $I_d = \{a \in R \mid ax^d + \text{h.o.t.} \in I\}$.
Then we see that $I_0 \subset I_1 \subset \ldots$ stabilizes as $R$
is Noetherian. Choose $d_0$ such that $I_{d_0} = I_{d_0 + 1} = \ldots$.
For each $d \leq d_0$ choose elements $f_{d, j} \in I \cap (x^d)$,
$j = 1, \ldots, n_d$ such that if we write
$f_{d, j} = a_{d, j}x^d + \text{h.o.t}$ then $I_d = (a_{d, j})$.
Denote $I' = (\{f_{d, j}\}_{d = 0, \ldots, d_0, j = 1, \ldots, n_d})$.
Then it is clear that $I' \subset I$. Pick $f \in I$.
Choose the coefficients $c_{d,i}$ successively for $d=0,\ldots,d_0$.
Before the step for $d$, the current remainder lies in $I\cap(x^d)$.
Its coefficient of $x^d$ belongs to $I_d$, so express it as
$\sum_{i=1}^{n_d}c_{d,i}a_{d,i}$ and subtract
$\sum_{i=1}^{n_d}c_{d,i}f_{d,i}$. The new remainder lies in
$I\cap(x^{d+1})$. Thus
$$
h_0=f-\sum_{d=0}^{d_0}\sum_{i=1}^{n_d}c_{d,i}f_{d,i}
\in I\cap(x^{d_0+1}).
$$
For every $e\geq1$, having constructed
$h_{e-1}\in I\cap(x^{d_0+e})$, its coefficient of $x^{d_0+e}$
lies in $I_{d_0+e}=I_{d_0}$. Choose $c_{i,e}\in R$ such that this
coefficient is $\sum_{i=1}^{n_{d_0}}c_{i,e}a_{d_0,i}$, and put
$$
h_e=h_{e-1}-\sum_{i=1}^{n_{d_0}}c_{i,e}x^e f_{d_0,i}
\in I\cap(x^{d_0+e+1}).
$$
These choices define actual formal power series
$C_i(x)=\sum_{e=1}^\infty c_{i,e}x^e\in R[[x]]$.
For any fixed degree $N$, the coefficient of $x^N$ in
$C_i(x)f_{d_0,i}$ receives contributions only from
$1\leq e\leq N-d_0$, a finite set. The remainder formula with
$d_0+e+1>N$ therefore proves, coefficient by coefficient, the identity
$$
f=\sum_{d=0}^{d_0}\sum_{i=1}^{n_d}c_{d,i}f_{d,i}
 +\sum_{i=1}^{n_{d_0}}\left(\sum_{e=1}^\infty c_{i,e}x^e\right)f_{d_0,i}.
$$
This is a finite $R[[x]]$-linear combination of the original chosen
generators, so $f$ belongs to $I'$. No assertion that $I'$ is closed
under limits is needed: the equality is an equality of all formal
coefficients with finitely many generator terms.
\end{proof}

\noindent
The following lemma, although easy, is useful because
finite type $\mathbf{Z}$-algebras come up quite often in
a technique called ``absolute Noetherian reduction''.

\begin{lemma}
\label{editorial-source-lemma-obvious-Noetherian}
Any finite type algebra over a field is Noetherian.
Any finite type algebra over $\mathbf{Z}$ is Noetherian.
\end{lemma}

\begin{proof}
This is immediate from Lemma \ref{editorial-source-lemma-Noetherian-permanence}
and the fact that fields are Noetherian rings and that
$\mathbf{Z}$ is a Noetherian ring (because it is a
principal ideal domain).
\end{proof}

\begin{lemma}
\label{editorial-source-lemma-Noetherian-finite-type-is-finite-presentation}
Let $R$ be a Noetherian ring.
\begin{enumerate}
\item Any finite $R$-module is of finite presentation.
\item Any submodule of a finite $R$-module is finite.
\item Any finite type $R$-algebra is of finite presentation over $R$.
\end{enumerate}
\end{lemma}

\begin{proof}
Let $M$ be a finite $R$-module. By
Lemma \ref{algebra-lemma-trivial-filter-finite-module}
we can find a finite filtration of $M$ whose successive quotients are
of the form $R/I$. Since any ideal is finitely generated, each of
the quotients $R/I$ is finitely presented. Hence $M$ is finitely
presented by
Lemma \ref{algebra-lemma-extension}.
This proves (1).

\medskip\noindent
Let $N \subset M$ be a submodule. As $M$ is finite, the quotient
$M/N$ is finite. Thus $M/N$ is of finite presentation by part (1).
Thus we see that $N$ is finite by Lemma \ref{algebra-lemma-extension} part (5).
This proves part (2).

\medskip\noindent
To see (3) note that any ideal of
$R[x_1, \ldots, x_n]$ is finitely generated by
Lemma \ref{editorial-source-lemma-Noetherian-permanence}.
\end{proof}

\begin{lemma}
\label{editorial-source-lemma-Noetherian-topology}
If $R$ is a Noetherian ring then $\Spec(R)$
is a Noetherian topological space, see Topology,
Definition \ref{topology-definition-noetherian}.
\end{lemma}

\begin{proof}
This is because any closed subset of $\Spec(R)$
is uniquely of the form $V(I)$ with $I$ a radical ideal,
see Lemma \ref{editorial-source-lemma-Zariski-topology}.
And this correspondence is inclusion reversing.
Thus the result follows from the definitions.
\end{proof}

\begin{lemma}
\label{editorial-source-lemma-Noetherian-irreducible-components}
\begin{slogan}
A Noetherian affine scheme has finitely many generic points.
\end{slogan}
If $R$ is a Noetherian ring then $\Spec(R)$
has finitely many irreducible components. In other words
$R$ has finitely many minimal primes.
\end{lemma}

\begin{proof}
By Lemma \ref{editorial-source-lemma-Noetherian-topology} and
Topology, Lemma \ref{topology-lemma-Noetherian}
we see there are finitely many irreducible components.
By Lemma \ref{editorial-source-lemma-irreducible} these correspond to
minimal primes of $R$.
\end{proof}

\begin{lemma}
\label{editorial-source-lemma-Noetherian-base-change-finite-type}
Let $R \to S$ be a ring map. Let $R \to R'$ be of finite type.
If $S$ is Noetherian, then the base change $S' = R' \otimes_R S$
is Noetherian.
\end{lemma}

\begin{proof}
By Lemma \ref{algebra-lemma-base-change-finiteness} finite type is stable under
base change. Thus $S \to S'$ is of finite type. Since $S$ is Noetherian we
can apply Lemma \ref{editorial-source-lemma-Noetherian-permanence}.
\end{proof}

\begin{lemma}
\label{editorial-source-lemma-Noetherian-field-extension}
Let $k$ be a field and let $R$ be a Noetherian $k$-algebra.
If $K/k$ is a finitely generated field extension then
$K \otimes_k R$ is Noetherian.
\end{lemma}

\begin{proof}
Since $K/k$ is a finitely generated field extension, there exists
a finitely generated $k$-algebra $B \subset K$ such that $K$ is
the fraction field of $B$. In other words, $K = S^{-1}B$
with $S = B \setminus \{0\}$. Then $K \otimes_k R = S^{-1}(B \otimes_k R)$.
Then $B \otimes_k R$ is Noetherian by
Lemma \ref{editorial-source-lemma-Noetherian-base-change-finite-type}.
Finally, $K \otimes_k R = S^{-1}(B \otimes_k R)$ is Noetherian by
Lemma \ref{editorial-source-lemma-Noetherian-permanence}.
\end{proof}

\noindent
Here are some fun lemmas that are sometimes useful.

\begin{lemma}
\label{editorial-source-lemma-subring-of-local-ring}
Let $R$ be a ring and $\mathfrak p \subset R$ be a prime.
There exists an $f \in R$, $f \not \in \mathfrak p$ such
that $R_f \to R_\mathfrak p$ is injective in each of the
following cases
\begin{enumerate}
\item $R$ is a domain,
\item $R$ is Noetherian, or
\item $R$ is reduced and has finitely many minimal primes.
\end{enumerate}
In (2), it suffices more generally that
$\ker(R\to R_{\mathfrak p})$ is finitely generated.
\end{lemma}

\begin{proof}
If $R$ is a domain, then $R \subset R_\mathfrak p$, hence $f = 1$ works.
More generally, suppose the kernel $I=\ker(R\to R_{\mathfrak p})$
is generated by $a_1,\ldots,a_t$. For each $a_j$, the localization
criterion gives $s_j\notin\mathfrak p$ with $s_ja_j=0$.
Set $f=\prod_{j=1}^ts_j$, using $f=1$ when $t=0$.
Then $f\notin\mathfrak p$ and $fI=0$.
If $r/f^e$ maps to zero in $R_{\mathfrak p}$, the image of $r$
is zero, so $r\in I$ and $fr=0$. Hence $r/f^e=0$ in $R_f$.
This proves injectivity, and applies in particular when $R$ is Noetherian.

\medskip\noindent
For (3), retain all the minimal primes
$\mathfrak p_1,\ldots,\mathfrak p_n$. For each $i$ with
$\mathfrak p_i\not\subset\mathfrak p$, choose
$f_i\in\mathfrak p_i\setminus\mathfrak p$, and put
$$
f=\prod_{\mathfrak p_i\not\subset\mathfrak p}f_i.
$$
The empty product is $1$, and primality of $\mathfrak p$ gives
$f\notin\mathfrak p$. If $r\in\ker(R\to R_{\mathfrak p})$,
some $s\notin\mathfrak p$ satisfies $sr=0$. Whenever
$\mathfrak p_i\subset\mathfrak p$, we have $s\notin\mathfrak p_i$,
so $r\in\mathfrak p_i$ by primality. For every other $i$ the
chosen factor $f_i$ gives $f\in\mathfrak p_i$.
It follows that $fr$ belongs to every minimal prime. The intersection
of the minimal primes is the nilradical, which is zero since $R$
is reduced. Thus $fr=0$. The same fraction argument as above proves
that the original canonical map $R_f\to R_{\mathfrak p}$ is injective.
\end{proof}

\begin{lemma}
\label{editorial-source-lemma-surjective-endo-noetherian-ring-is-iso}
Any surjective endomorphism of a Noetherian ring is an isomorphism.
\end{lemma}

\begin{proof}
Suppose $f:R\to R$ is surjective and not injective, and choose
$0\ne a\in\ker(f)$. For each $n\geq1$, surjectivity of $f^n$
provides $b_n\in R$ with $f^n(b_n)=a$. Then
$f^{n+1}(b_n)=f(a)=0$ but $f^n(b_n)\ne0$.
Thus every inclusion in
$$
\ker(f)\subset\ker(f^2)\subset\ker(f^3)\subset\cdots
$$
is strict, witnessed by $b_n$ at the $n$th inclusion. This contradicts
the ascending chain condition on ideals of the Noetherian ring $R$.
Consequently $f$ is injective as well as surjective, and is an isomorphism.
\end{proof}










\subsection{Editorial treatment of Algebra lines 6253--6265}

\hypertarget{algebra-context-045}{}

\noindent\href{https://github.com/stacks/stacks-project/blob/a04446e57ec1fbc252a871afcec7752fb2807b14/algebra.tex\#L6253}{Original source, lines 6253--6265}. Complete original and reviewed TeX passages accompany this note. Every intervention is separately recorded; the following treatment is editorial.


\begin{example}
\label{editorial-source-example-locally-nilpotent-not-nilpotent}
Let $R = k[x_n | n \in \mathbf{N}]$ be the polynomial ring in infinitely
many variables over a field $k$. Let $I$ be the ideal generated by
the elements $x_n^n$ for $n \in \mathbf{N}$ and $S = R/I$. Then the ideal
$J \subset S$ generated by the images of $x_n$, $n \in \mathbf{N}$
is locally nilpotent, but not nilpotent. The convention here is
$\mathbf N=\{1,2,3,\ldots\}$, so no relation with exponent zero occurs.
An element $y\in J$ has a finite expression
$y=\sum_{i=1}^t b_i\overline{x}_{n_i}$ with $b_i\in S$ and $n_i\geq1$.
Put $N=1+\sum_{i=1}^t(n_i-1)$. In the full multinomial expansion
of $y^N$, every term has an exponent at least $n_i$ on one factor
$\overline{x}_{n_i}$, and hence vanishes. The empty sum gives $y=0$.
Thus $J$ is locally nilpotent.
For each $n\geq1$ define a $k$-algebra map
$$
R\longrightarrow k[t]/(t^{n+1}),\qquad
x_{n+1}\longmapsto t,\qquad x_j\longmapsto0\quad(j\ne n+1).
$$
Every generator $x_j^j$ of $I$ maps to zero, so this map factors
through $S$. The element $\overline{x}_{n+1}^n\in J^n$ maps to
$t^n\ne0$, since $1,t,\ldots,t^n$ form a $k$-basis of the target.
Consequently $J^n\ne0$ for every $n\geq1$ and $J$ is not nilpotent.
\end{example}


\subsection{Editorial treatment of Algebra lines 6267--6277}

\hypertarget{algebra-context-046}{}

\noindent\href{https://github.com/stacks/stacks-project/blob/a04446e57ec1fbc252a871afcec7752fb2807b14/algebra.tex\#L6267}{Original source, lines 6267--6277}. Complete original and reviewed TeX passages accompany this note. Every intervention is separately recorded; the following treatment is editorial.


\begin{lemma}
\label{editorial-source-lemma-locally-nilpotent}
Let $R \to R'$ be a ring map and let $I \subset R$ be a locally nilpotent
ideal. Then $IR'$ is a locally nilpotent ideal of $R'$.
\end{lemma}

\begin{proof}
Write $\varphi:R\to R'$ for the map. Every element of $IR'$ has a
finite expression $y=\sum_{i=1}^t b_i\varphi(x_i)$ with $x_i\in I$
and $b_i\in R'$. Choose positive integers $d_i$ with $x_i^{d_i}=0$
and put $N=1+\sum_{i=1}^t(d_i-1)$. The complete expansion is
$$
y^N=\sum_{\substack{\alpha_1,\ldots,\alpha_t\geq0\\
\alpha_1+\cdots+\alpha_t=N}}
\frac{N!}{\alpha_1!\cdots\alpha_t!}
\prod_{i=1}^t b_i^{\alpha_i}\varphi(x_i)^{\alpha_i}.
$$
The displayed multinomial coefficients are integers, acting through
the canonical map $\mathbf Z\to R'$. Every summand has
$\alpha_i\geq d_i$ for some $i$, since otherwise its total degree
would be at most $N-1$. Thus every summand is zero. If $t=0$ then
$y=0$ and $N=1$ suffices. In particular the case of two nilpotents
$x^n=y^m=0$ gives $(x+y)^{n+m-1}=0$, with both original exponents
retained. This proves local nilpotence after extension.
\end{proof}


\subsection{Editorial treatment of Algebra lines 6279--6297}

\hypertarget{algebra-context-047}{}

\noindent\href{https://github.com/stacks/stacks-project/blob/a04446e57ec1fbc252a871afcec7752fb2807b14/algebra.tex\#L6279}{Original source, lines 6279--6297}. Complete original and reviewed TeX passages accompany this note. Every intervention is separately recorded; the following treatment is editorial.


\begin{lemma}
\label{editorial-source-lemma-locally-nilpotent-unit}
Let $R$ be a ring and let $I \subset R$ be a locally nilpotent
ideal.
An element $x$ of $R$ is a unit if and only if the image of $x$
in $R/I$ is a unit.
\end{lemma}

\begin{proof}
If $x$ is a unit in $R$, then its image is clearly a unit in $R/I$.
It remains to prove the converse.
Assume the image of $y \in R$ in $R/I$ is the inverse of the image of $x$.
Then $xy=1-z$ for some $z\in I$, and $z^N=0$ for a positive
integer $N$. Retain all terms in the finite geometric sum. They give
$$
x\left(y\sum_{j=0}^{N-1}z^j\right)
=(1-z)\sum_{j=0}^{N-1}z^j=1-z^N=1.
$$
Commutativity gives the reverse product as well, so the displayed
element $y\sum_{j=0}^{N-1}z^j$ is the inverse of the original $x$.
In particular, $1+I$ is a group under multiplication: it is closed
under products, and if $z\in I$ with $z^N=0$ then
$(1+z)^{-1}=\sum_{j=0}^{N-1}(-z)^j\in1+I$.
\end{proof}


\subsection{Editorial treatment of Algebra lines 6299--6316}

\hypertarget{algebra-context-048}{}

\noindent\href{https://github.com/stacks/stacks-project/blob/a04446e57ec1fbc252a871afcec7752fb2807b14/algebra.tex\#L6299}{Original source, lines 6299--6316}. Complete original and reviewed TeX passages accompany this note. Every intervention is separately recorded; the following treatment is editorial.


\begin{lemma}
\label{editorial-source-lemma-Noetherian-power}
\begin{slogan}
An ideal in a Noetherian ring is nilpotent if each element
of the ideal is nilpotent.
\end{slogan}
Let $R$ be a Noetherian ring. Let $I, J$ be ideals of $R$.
Suppose $J \subset \sqrt{I}$. Then $J^n \subset I$ for some $n$.
In particular, in a Noetherian ring the notions of
``locally nilpotent ideal''
and ``nilpotent ideal'' coincide.
More generally, the conclusion $J^n\subset I$ needs only finite
generation of $J$, not Noetherianity of $R$. If
$J=(f_1,\ldots,f_s)$ and $f_i^{d_i}\in I$ with $d_i\geq1$, one may
take $n=1+\sum_{i=1}^s(d_i-1)$.
\end{lemma}

\begin{proof}
Choose generators $J=(f_1,\ldots,f_s)$ and positive exponents
$d_i$ such that $f_i^{d_i}\in I$. Put $N=1+\sum_{i=1}^s(d_i-1)$.
The ideal $J^N$ is generated by the full list of monomials
$\prod_i f_i^{\alpha_i}$ with $\alpha_i\geq0$ and $\sum_i\alpha_i=N$.
Every such list has $\alpha_i\geq d_i$ for some $i$, so that monomial
belongs to $I$. Hence $J^N\subset I$. For $s=0$, $J=0$ and $N=1$
works. The original larger choice $d_1+\cdots+d_s+1$ also works,
since every higher power of $J$ is contained in $J^N$.
Only the existence of a finite generating set for $J$ was used.

\medskip\noindent
The bound cannot be decreased using only the displayed exponents.
In $B=\mathbf Z[T_1,\ldots,T_s]/(T_1^{d_1},\ldots,T_s^{d_s})$,
the residue classes of monomials with $0\leq\alpha_i<d_i$ form a
$\mathbf Z$-basis: the defining monomial ideal is the free subgroup
spanned by all the other monomials. Therefore
$\prod_i T_i^{d_i-1}$ is a nonzero element of
$(T_1,\ldots,T_s)^{N-1}$. Together with the upper bound this proves
sharpness, including the empty product in degree zero.
\end{proof}


\subsection{Editorial treatment of Algebra lines 6318--6383}

\hypertarget{algebra-context-049}{}

\noindent\href{https://github.com/stacks/stacks-project/blob/a04446e57ec1fbc252a871afcec7752fb2807b14/algebra.tex\#L6318}{Original source, lines 6318--6383}. Complete original and reviewed TeX passages accompany this note. Every intervention is separately recorded; the following treatment is editorial.


\begin{lemma}
\label{editorial-source-lemma-lift-idempotents}
Let $R$ be a ring. Let $I \subset R$ be a locally nilpotent ideal.
Then $R \to R/I$ induces a bijection on idempotents.
\end{lemma}

\begin{proof}[First proof of Lemma \ref{editorial-source-lemma-lift-idempotents}]
As $I$ is locally nilpotent it is contained in every prime ideal.
Hence $\Spec(R/I) = V(I) = \Spec(R)$. Hence the
lemma follows from Lemma \ref{algebra-lemma-disjoint-decomposition}.
\end{proof}

\begin{proof}[Second proof of Lemma \ref{editorial-source-lemma-lift-idempotents}]
Suppose $\overline{e} \in R/I$ is an idempotent.
We have to lift $\overline{e}$ to an idempotent of $R$.

\medskip\noindent
First, choose any lift $f \in R$ of $\overline{e}$, and set
$x = f^2 - f$. Then, $x \in I$, so $x$ is nilpotent (since $I$
is locally nilpotent). Let now $J$ be the ideal of $R$ generated
by $x$. Then, $J$ is nilpotent (not just locally nilpotent),
since it is generated by the nilpotent $x$.

\medskip\noindent
Now, assume that we have found a lift $e \in R$ of $\overline{e}$
such that $e^2 - e \in J^k$ for some $k \geq 1$.
Let $e' = e - (2e - 1)(e^2 - e) = 3e^2 - 2e^3$, which is another
lift of $\overline{e}$ (since the idempotency of $\overline{e}$
yields $e^2 - e \in I$). Then
$$
(e')^2 - e' = (4e^2 - 4e - 3)(e^2 - e)^2 \in J^{2k}
$$
because expansion of either side gives
$4e^6-12e^5+9e^4+2e^3-3e^2$.
In particular every coefficient and sign in the original update is retained.

\medskip\noindent
Choose $N\geq1$ with $x^N=0$, so $J^N=0$.
Starting with $e_0=f$, define for every $j\geq0$
$$
e_{j+1}=e_j-(2e_j-1)(e_j^2-e_j)=3e_j^2-2e_j^3.
$$
The identity just proved gives inductively
$e_j^2-e_j\in J^{2^j}$, and every $e_j$ still maps to
$\overline e$ since each update differs by an element of $J$.
Any integer $j$ with $2^j\geq N$ therefore gives
$e_j^2-e_j=0$. This constructs the required lift after finitely
many explicitly specified updates.

\medskip\noindent
We thus have seen that if $\overline{e} \in R/I$ is any
idempotent, then there exists a lift of $\overline{e}$
which is an idempotent of $R$.
It remains to prove that this lift is unique. Indeed, let
$e_1$ and $e_2$ be two such lifts. We
need to show that $e_1 = e_2$.

\medskip\noindent
By definition of $e_1$ and $e_2$, we have $e_1 \equiv e_2
\mod I$, and both $e_1$ and $e_2$ are idempotent. From
$e_1 \equiv e_2 \mod I$, we see that $e_1 - e_2 \in I$,
so that $e_1 - e_2$ is nilpotent (since $I$ is locally nilpotent).
A straightforward
computation (using the idempotency of $e_1$ and $e_2$)
reveals that $(e_1 - e_2)^3 = e_1 - e_2$. Using this and
induction, we obtain $(e_1 - e_2)^k = e_1 - e_2$ for any
positive odd integer $k$. Since all high enough $k$ satisfy
$(e_1 - e_2)^k = 0$ (since $e_1 - e_2$ is nilpotent),
this shows $e_1 - e_2 = 0$, so that $e_1 = e_2$, which
completes our proof.
\end{proof}


\subsection{Editorial treatment of Algebra lines 6385--6399}

\hypertarget{algebra-context-050}{}

\noindent\href{https://github.com/stacks/stacks-project/blob/a04446e57ec1fbc252a871afcec7752fb2807b14/algebra.tex\#L6385}{Original source, lines 6385--6399}. Complete original and reviewed TeX passages accompany this note. Every intervention is separately recorded; the following treatment is editorial.


\begin{lemma}
\label{editorial-source-lemma-lift-idempotents-noncommutative}
Let $A$ be a possibly noncommutative algebra.
Let $e \in A$ be an element such that $x = e^2 - e$ is nilpotent.
Then there exists an idempotent of the form
$e' = e + x(\sum a_{i, j}e^ix^j) \in A$
with $a_{i, j} \in \mathbf{Z}$.
\end{lemma}

\begin{proof}
Choose $n\geq1$ with $x^n=0$, and distinguish the indeterminate
$E$ from the original element $e\in A$. Evaluation gives a unital
ring map
$$
\mathbf Z[E]/((E^2-E)^n)\longrightarrow A,\qquad E\longmapsto e.
$$
This map exists even if $A$ is noncommutative: integer multiples of
$1_A$ are central and every polynomial in the single element $e$
commutes with every other such polynomial. The relation vanishes
because $(e^2-e)^n=x^n=0$.
In the commutative source ring let $I=(E^2-E)$ and apply
Lemma \ref{editorial-source-lemma-lift-idempotents} to the class of $E$ modulo $I$.
Its idempotent lift $P$ differs from $E$ by an element of $I$,
so $P=E+(E^2-E)Q(E)$ for a polynomial $Q\in\mathbf Z[E]$.
Its image is the idempotent $e'=e+xQ(e)$ in $A$.
Writing $Q(E)=\sum_i a_{i,0}E^i$ gives the required finite sum,
with $a_{i,j}=0$ for $j>0$. Alternatively the same explicit cubic
updates in the preceding proof construct $P$. No uniqueness of
arbitrary idempotent lifts in the whole noncommutative algebra is asserted.
\end{proof}


\subsection{Editorial treatment of Algebra lines 6401--6438}

\hypertarget{algebra-context-051}{}

\noindent\href{https://github.com/stacks/stacks-project/blob/a04446e57ec1fbc252a871afcec7752fb2807b14/algebra.tex\#L6401}{Original source, lines 6401--6438}. Complete original and reviewed TeX passages accompany this note. Every intervention is separately recorded; the following treatment is editorial.


\begin{lemma}
\label{editorial-source-lemma-lift-nth-roots}
Let $R$ be a ring. Let $I \subset R$ be a locally nilpotent ideal.
Let $n \geq 1$ be an integer which is invertible in $R/I$. Then
\begin{enumerate}
\item the $n$th power map $1 + I \to 1 + I$, $1 + x \mapsto (1 + x)^n$
is a bijection,
\item a unit of $R$ is an $n$th power if and only if its image in $R/I$
is an $n$th power.
\end{enumerate}
More precisely, for each unit $a\in R$, reduction induces a bijection
between the solutions $u\in R^*$ of $u^n=a$ and the solutions
$\overline u\in(R/I)^*$ of $\overline u^{,n}=\overline a$.
\end{lemma}

\begin{proof}
Let $a \in R$ be a unit whose image in $R/I$ is the same as the image
of $b^n$ with $b \in R$. Then $b$ is a unit
(Lemma \ref{editorial-source-lemma-locally-nilpotent-unit}) and
$ab^{-n} = 1 + x$ for some $x \in I$. Hence $ab^{-n} = c^n$ by
part (1). Thus (2) follows from (1).

\medskip\noindent
Proof of (1). This is true because there is an inverse to the
map $1 + x \mapsto (1 + x)^n$. Namely, we can consider the map
which sends $1 + x$ to
\begin{align*}
(1 + x)^{1/n}
& =
1 + {1/n \choose 1}x +
{1/n \choose 2}x^2 +
{1/n \choose 3}x^3 + \ldots \\
& =
1 + \frac{1}{n} x + \frac{1 - n}{2n^2}x^2 +
\frac{(1 - n)(1 - 2n)}{6n^3}x^3 + \ldots
\end{align*}
We justify the coefficients and both inverse properties over the original ring.
For $k\geq1$ retain the full coefficient formula
$$
c_k={1/n\choose k}
=\frac{\prod_{j=0}^{k-1}(1-jn)}{n^k k!},\qquad c_0=1.
$$
For each prime $p$ not dividing $n$ and each $a\geq1$, exactly one
residue class of $j$ modulo $p^a$ solves $1-jn\equiv0\pmod{p^a}$.
Among $k$ consecutive integers $j=0,\ldots,k-1$ there are at least
$\lfloor k/p^a\rfloor$ elements of this class. If the numerator is
nonzero, its $p$-adic valuation is therefore at least
$\sum_{a\geq1}\lfloor k/p^a\rfloor$, which equals the valuation
of $k!$: each multiple of $p^a$ among $1,\ldots,k$ contributes once
to that sum. The denominator $n^k$ contributes nothing at $p$.
If the numerator is zero, the coefficient is zero and the conclusion
is immediate. Thus the denominator of $c_k$ has no prime divisor
outside those dividing $n$, and $c_k\in\mathbf Z[1/n]$.
This statement does not require $k!$ to be invertible in $R$.
The integer $n$ is invertible in $R$ by
Lemma \ref{editorial-source-lemma-locally-nilpotent-unit}, so every $c_k$ has a
well-defined image in $R$ via $\mathbf Z[1/n]\to R$.

\medskip\noindent
Set $F(X)=\sum_{k\geq0}c_kX^k\in\mathbf Z[1/n][[X]]$.
The full product formula gives
$n(k+1)c_{k+1}=(1-kn)c_k$; hence in $\mathbf Q[[X]]$ it gives
$n(1+X)F'(X)=F(X)$ and $F(0)=1$.
For $H=F^n$, the product rule yields $(1+X)H'=H$, so the derivative
of $H/(1+X)$ is zero. A formal series over $\mathbf Q$ with zero
derivative is constant, and its constant term here is $1$.
Therefore $F(X)^n=1+X$. The coefficientwise injection
$\mathbf Z[1/n][[X]]\hookrightarrow\mathbf Q[[X]]$ proves the same
identity over $\mathbf Z[1/n]$.
For $x\in I$ with $x^N=0$, reduction modulo $X^N$ followed by
$X\mapsto x$ makes the finite expression
$c=\sum_{k=0}^{N-1}c_kx^k\in1+I$ satisfy $c^n=1+x$.
Every larger truncation gives the same element, since the additional
powers of $x$ vanish. This proves surjectivity in (1).

\medskip\noindent
For injectivity, if $u,v\in1+I$ and $u^n=v^n$, use the complete
factorization
$$
0=u^n-v^n=(u-v)\sum_{j=0}^{n-1}u^{n-1-j}v^j.
$$
The sum maps to the unit $n$ in $R/I$, hence is a unit of $R$ by
Lemma \ref{editorial-source-lemma-locally-nilpotent-unit}. Thus $u=v$.
The constructed map is consequently the inverse in (1).

\medskip\noindent
Finally fix a unit $a$ and a root $\overline b$ of $\overline a$.
Choose a lift $b\in R$; it is a unit, and $ab^{-n}\in1+I$.
The unique root $c\in1+I$ furnished by (1) gives $u=bc$ with
$u^n=a$ and $\overline u=\overline b$.
If another such lift is $v$, then $uv^{-1}\in1+I$ and
$(uv^{-1})^n=1$, so injectivity in (1) gives $u=v$.
This proves the asserted bijection of root sets and also (2).
\end{proof}


\subsection{Editorial treatment of Algebra lines 6455--6477}

\hypertarget{algebra-context-052}{}

\noindent\href{https://github.com/stacks/stacks-project/blob/a04446e57ec1fbc252a871afcec7752fb2807b14/algebra.tex\#L6455}{Original source, lines 6455--6477}. Complete original and reviewed TeX passages accompany this note. Every intervention is separately recorded; the following treatment is editorial.


\begin{lemma}
\label{editorial-source-lemma-invert-closed-quotient}
Let $R$ be a ring. Let $S \subset R$ be a multiplicative subset.
Assume the image of the map $\Spec(S^{-1}R) \to \Spec(R)$
is closed. Then $S^{-1}R \cong R/I$ for some ideal $I \subset R$.
\end{lemma}

\begin{proof}
Let $I = \Ker(R \to S^{-1}R)$ so that $V(I)$ contains the image.
Say the image is the closed subset $V(I') \subset \Spec(R)$ for
some ideal $I' \subset R$. So $V(I') \subset V(I)$.
For $f \in I'$ we see that $f/1 \in S^{-1}R$
belongs to every prime ideal. Hence $f^n$ maps to zero in $S^{-1}R$
for some $n \geq 1$ (Lemma \ref{editorial-source-lemma-Zariski-topology}).
Hence $V(I') = V(I)$.
Every $g\in S$ avoids all primes in the image, hence every prime
containing $I$. Its class $\overline g$ in $R/I$ belongs to no maximal
ideal and is therefore a unit. The quotient and localization properties
give ring maps
$$
\alpha:R/I\longrightarrow S^{-1}R,\quad \overline r\longmapsto r/1,
\qquad
\beta:S^{-1}R\longrightarrow R/I,\quad r/s\longmapsto
\overline r\,\overline s^{-1}.
$$
For every $\overline r$, $\beta\alpha(\overline r)=\overline r$.
For every $r/s$, $\alpha\beta(r/s)=(r/1)(s/1)^{-1}=r/s$.
Thus the maps are inverse, and their composite with $R\to R/I$
is the original localization map. In particular the localization map
is surjective. Conversely, if $R\to S^{-1}R$ is surjective, it is
the quotient by its kernel and its spectrum has closed image by
Lemma \ref{algebra-lemma-spec-closed}.
\end{proof}


\subsection{Editorial treatment of Algebra lines 6479--6505}

\hypertarget{algebra-context-053}{}

\noindent\href{https://github.com/stacks/stacks-project/blob/a04446e57ec1fbc252a871afcec7752fb2807b14/algebra.tex\#L6479}{Original source, lines 6479--6505}. Complete original and reviewed TeX passages accompany this note. Every intervention is separately recorded; the following treatment is editorial.


\begin{lemma}
\label{editorial-source-lemma-invert-closed-split}
Let $R$ be a ring. Let $S \subset R$ be a multiplicative subset.
Assume the image of the map $\Spec(S^{-1}R) \to \Spec(R)$
is closed. If $R$ is Noetherian, or $\Spec(R)$ is a
Noetherian topological space, or $S$ is finitely generated as a monoid,
then $R\cong S^{-1}R\times R'$ for some ring $R'$.
More generally, the same conclusion holds whenever the complement
of the closed image is quasi-compact.
\end{lemma}

\begin{proof}
By Lemma \ref{editorial-source-lemma-invert-closed-quotient} we have $S^{-1}R \cong R/I$
for some ideal $I \subset R$. By Lemma \ref{algebra-lemma-disjoint-implies-product}
it suffices to show that $V(I)$ is open.
If $R$ is Noetherian then $\Spec(R)$ is a Noetherian
topological space, see Lemma \ref{editorial-source-lemma-Noetherian-topology}.
If $\Spec(R)$ is Noetherian, the complement of $V(I)$ is
quasi-compact by Topology, Lemma
\ref{topology-lemma-Noetherian-quasi-compact}. More generally assume
only that this complement is quasi-compact. Its cover by the opens
$D(f)$, $f\in I$, has a finite subcover $D(f_1),\ldots,D(f_n)$,
so $V(I)=V(f_1,\ldots,f_n)$. For each $i$ choose $s_i\in S$ with
$s_if_i=0$, using the localization kernel criterion, and set
$g=\prod_{i=1}^n s_i\in S$ (with empty product $1$).
Then $gf_i=0$ for every $i$. If a prime does not contain $g$, it
contains every $f_i$, so $D(g)\subset V(I)$. Conversely, every prime
in $V(I)$ lies in the localization image and avoids all of $S$,
including $g$. Hence $V(I)=D(g)$, which is open.

\medskip\noindent
If instead $S$ is generated as a monoid by $g_1,\ldots,g_m$, put
$t=g_1\cdots g_m$, with $t=1$ when the list is empty.
Each $g_i$ is a unit in $R_t$, because a factor of a unit in a
commutative ring is a unit. Conversely $t$ is a unit in $S^{-1}R$.
The resulting canonical maps $S^{-1}R\rightleftarrows R_t$ are
inverse: their composites fix $R$ and the unique inverses adjoined.
Thus their spectrum image is $D(t)$, again proving $V(I)$ open.

\medskip\noindent
To identify the original localization as the first product factor,
choose an idempotent $e$ with $D(e)=V(I)$ by
Lemma \ref{algebra-lemma-disjoint-decomposition}. Each $s\in S$ avoids every
prime of $R_e$, because those primes correspond to $D(e)=V(I)$;
thus it is a unit in $R_e$. Likewise $e/1$ avoids every prime of
$S^{-1}R$, so it is a unit there. Being idempotent, it equals $1$.
The canonical maps $S^{-1}R\rightleftarrows R_e$ therefore exist and
are inverse by the same identity-on-$R$ and inverse-element argument.
Finally
$$
R\longrightarrow R_e\times R_{1-e},\qquad
r\longmapsto(r/1,r/1)
$$
has inverse $(a/e^u,b/(1-e)^v)\longmapsto ea+(1-e)b$.
This formula is independent of fraction representatives: equality
in $R_e$ implies $e(a-a')=0$, and equality in $R_{1-e}$ implies
$(1-e)(b-b')=0$, by the localization equality criterion and
idempotency. The two composites are identities, since $e$ maps to
$1$ in $R_e$ and $0$ in $R_{1-e}$, and $e+(1-e)=1$ in $R$.
The map is a ring map, and hence so is its inverse. Taking
$R'=R_{1-e}$ gives the claimed decomposition, whose first projection
is exactly the original map $R\to S^{-1}R$.
\end{proof}


\subsection{Editorial treatment of Algebra lines 6523--6608}

\hypertarget{algebra-context-054}{}

\noindent\href{https://github.com/stacks/stacks-project/blob/a04446e57ec1fbc252a871afcec7752fb2807b14/algebra.tex\#L6523}{Original source, lines 6523--6608}. Complete original and reviewed TeX passages accompany this note. Every intervention is separately recorded; the following treatment is editorial.


\begin{theorem}[Hilbert Nullstellensatz]
\label{editorial-source-theorem-nullstellensatz}
Let $k$ be a field.
\begin{enumerate}
\item
\label{editorial-source-item-finite-kappa}
For any maximal ideal $\mathfrak m \subset k[x_1, \ldots, x_n]$
the field extension $\kappa(\mathfrak m)/k$ is finite.
\item
\label{editorial-source-item-polynomial-ring-Jacobson}
Any radical ideal $I \subset k[x_1, \ldots, x_n]$
is the intersection of maximal ideals containing it.
\end{enumerate}
The same is true in any finite type $k$-algebra.
\end{theorem}

\begin{proof}
It is enough to prove part (\ref{editorial-source-item-finite-kappa}) of
the theorem for the case of a polynomial
algebra $k[x_1, \ldots, x_n]$, because any finitely generated
$k$-algebra is a quotient of such a polynomial algebra.
We prove this by induction on $n$. The case $n = 0$ is clear.
Suppose that $\mathfrak m$ is a maximal ideal in $k[x_1, \ldots, x_n]$.
Let $\mathfrak p \subset k[x_n]$ be the intersection
of $\mathfrak m$ with $k[x_n]$.

\medskip\noindent
If $\mathfrak p \not = (0)$,
then $\mathfrak p$ is maximal and generated by an irreducible
monic polynomial $P$ (because of the Euclidean algorithm
in $k[x_n]$). Then
$k' = k[x_n]/\mathfrak p$ is a finite field extension of $k$
and contained in $\kappa(\mathfrak m)$. In this case
we get a surjection
$$
k'[x_1, \ldots, x_{n-1}]
\to
k'[x_1, \ldots, x_n] =
k' \otimes_k k[x_1, \ldots, x_n]
\longrightarrow
\kappa(\mathfrak m)
$$
where the last map sends the coefficient field $k'$ into
$\kappa(\mathfrak m)$ and each variable to its original residue
class. The image of $x_n$ already belongs to this coefficient field:
it is the class of $x_n$ in $k'=k[x_n]/\mathfrak p$. Hence the
composite from $k'[x_1,\ldots,x_{n-1}]$ is surjective. Its kernel
is maximal, and the induction hypothesis makes
$\kappa(\mathfrak m)/k'$ finite. The extension $k'/k$ is finite,
with basis $1,\overline{x}_n,\ldots,\overline{x}_n^{d-1}$ where
$d=\deg(P)$. Products of this basis with a $k'$-basis of
$\kappa(\mathfrak m)$ span the latter over $k$, proving finiteness
over the original field.

\medskip\noindent
If $\mathfrak p = (0)$ we consider the ring
extension $k[x_n] \subset k[x_1, \ldots, x_n]/\mathfrak m$.
This is a finitely generated ring extension, hence
of finite presentation by
Lemmas \ref{editorial-source-lemma-obvious-Noetherian} and
\ref{editorial-source-lemma-Noetherian-finite-type-is-finite-presentation}.
Thus the image of $\Spec(k[x_1, \ldots, x_n]/\mathfrak m)$
in $\Spec(k[x_n])$ is constructible by
Theorem \ref{editorial-source-theorem-chevalley}. Since the image
contains $(0)$ we conclude that it contains a standard
open $D(f)$ for a nonzero $f\in k[x_n]$.
This open contains infinitely many closed points, over finite fields
as well as infinite fields. Indeed, given any finite list of
irreducible polynomials $p_1,\ldots,p_t$ defining points already
chosen in $D(f)$, form the original polynomial
$$
q=1+x_n f\prod_{i=1}^t p_i.
$$
It has positive degree and hence an irreducible divisor $h$.
Such a divisor exists by induction on degree, factoring any reducible
nonconstant polynomial until an irreducible factor is reached.
No irreducible divisor of $f$ or of any $p_i$ can divide $q$, since
$q$ is congruent to $1$ modulo each such divisor. Thus $(h)$ is a
new maximal ideal not containing $f$. The empty list is allowed,
and this construction proves infinitude. But the field
$k[x_1,\ldots,x_n]/\mathfrak m$ has just one point in its spectrum;
its image cannot contain $D(f)$. This is the required contradiction.

\medskip\noindent
Proof of (\ref{editorial-source-item-polynomial-ring-Jacobson}). Let
$I \subset R$ be a radical ideal, with $R$ of finite type over $k$.
Let $f \in R$, $f \not \in I$. We have to find a maximal ideal
$\mathfrak m \subset R$ with $I \subset \mathfrak m$ and
$f \not \in \mathfrak m$. The ring $(R/I)_f$ is nonzero, since
$1 = 0$ in this ring would mean $f^n \in I$ and since $I$ is
radical this would mean $f \in I$ contrary to our assumption on $f$.
Thus we may choose a maximal ideal $\mathfrak m'$
in $(R/I)_f$, see Lemma \ref{editorial-source-lemma-Zariski-topology}.
Let $\mathfrak m \subset R$
be the inverse image of $\mathfrak m'$ in $R$. We see that
$I \subset \mathfrak m$
and $f \not \in \mathfrak m$. If we show that $\mathfrak m$ is a maximal
ideal of $R$, then we are done. We clearly have
$$
k \subset R/\mathfrak m \subset \kappa(\mathfrak m').
$$
The algebra $(R/I)_f$ is of finite type over $k$: it is the quotient
of $R[t]$ by the ideal generated by $I$ and $ft-1$, with $t$ mapping
to the inverse of the class of $f$. The mutually inverse maps send
$r/f^j$ to $rt^j$ and $t$ to $1/f$.
Thus part (\ref{editorial-source-item-finite-kappa}) applies and
$\kappa(\mathfrak m')/k$ is finite. The original injection
$R/\mathfrak m\hookrightarrow\kappa(\mathfrak m')$ makes
$R/\mathfrak m$ a finite dimensional $k$-vector space and a domain.
For each nonzero $a\in R/\mathfrak m$, multiplication by $a$ is
an injective linear endomorphism of this finite dimensional space,
hence surjective. Its image contains $1$, giving the inverse of $a$.
Consequently $R/\mathfrak m$ is a field, as also stated in Fields,
Lemma \ref{fields-lemma-subalgebra-algebraic-extension-field}.
Thus $\mathfrak m$ is maximal and the proof is complete.
\end{proof}


\subsection{Editorial treatment of Algebra lines 6610--6627}

\hypertarget{algebra-context-055}{}

\noindent\href{https://github.com/stacks/stacks-project/blob/a04446e57ec1fbc252a871afcec7752fb2807b14/algebra.tex\#L6610}{Original source, lines 6610--6627}. Complete original and reviewed TeX passages accompany this note. Every intervention is separately recorded; the following treatment is editorial.


\begin{lemma}
\label{editorial-source-lemma-field-finite-type-over-domain}
Let $R$ be a ring. Let $K$ be a field.
If $R \subset K$ and $K$ is of finite type over $R$,
then there exists an $f \in R$ such that $R_f$ is a field,
and $K/R_f$ is a finite field extension.
\end{lemma}

\begin{proof}
By Lemma \ref{editorial-source-lemma-characterize-image-finite-type} there
exist a nonempty open $U \subset \Spec(R)$
contained in the image $\{(0)\}$ of $\Spec(K) \to \Spec(R)$.
Choose $f \in R$, $f \not = 0$ such that $D(f) \subset U$, i.e.,
$D(f)=\{(0)\}$. The localization $R_f$ is a domain whose only
prime ideal is $(0)$. Every nonunit belongs to a maximal ideal,
so its only nonunit is zero and $R_f$ is a field. The original
inclusion $R\subset K$ extends to $R_f\subset K$, and the original
finite set of $R$-algebra generators of $K$ is also a set of
$R_f$-algebra generators. The Hilbert Nullstellensatz
(Theorem \ref{editorial-source-theorem-nullstellensatz}) therefore makes $K/R_f$ finite.
In fact the canonical inclusion $R_f\subset\operatorname{Frac}(R)$
is an equality: for $r/s\in\operatorname{Frac}(R)$ with $s\ne0$,
the element $s/1$ is nonzero in the field $R_f$ and has an inverse
there, so $r/s$ belongs to $R_f$. Thus the same $f$ realizes the
original fraction field as a principal localization of $R$.
\end{proof}


\subsection{Editorial treatment of Algebra lines 6727--6776}

\hypertarget{algebra-context-056}{}

\noindent\href{https://github.com/stacks/stacks-project/blob/a04446e57ec1fbc252a871afcec7752fb2807b14/algebra.tex\#L6727}{Original source, lines 6727--6776}. Complete original and reviewed TeX passages accompany this note. Every intervention is separately recorded; the following treatment is editorial.


\begin{lemma}
\label{editorial-source-lemma-characterize-jacobson}
Let $R$ be a ring. If $R$ is not Jacobson there exist
a prime $\mathfrak p \subset R$, an element $f \in R$
such that the following hold
\begin{enumerate}
\item $\mathfrak p$ is not a maximal ideal,
\item $f \not \in \mathfrak p$,
\item $V(\mathfrak p) \cap D(f) = \{\mathfrak p\}$, and
\item $(R/\mathfrak p)_f$ is a field.
\end{enumerate}
On the other hand, if $R$ is Jacobson, then for any pair $(\mathfrak p, f)$
such that (1) and (2) hold, $V(\mathfrak p)\cap D(f)$ contains
infinitely many maximal ideals of $R$; in particular it is infinite.
\end{lemma}

\begin{proof}
Assume $R$ is not Jacobson.
By Lemma \ref{algebra-lemma-jacobson} this means there exists a
closed subset $T \subset \Spec(R)$
whose set $T_0 \subset T$ of closed points is not dense in $T$.
Choose an $f \in R$ such that $T_0 \subset V(f)$ but
$T \not \subset V(f)$. Note that $T \cap D(f)$
is homeomorphic to $\Spec((R/I)_f)$ if $T = V(I)$, see
Lemmas \ref{algebra-lemma-spec-closed} and \ref{algebra-lemma-standard-open}.
The set $T\cap D(f)$ is nonempty, so its displayed coordinate ring
is nonzero. Every nonzero ring has a maximal ideal
(Lemma \ref{editorial-source-lemma-Zariski-topology}), so choose a closed point $t$
of the space $T\cap D(f)$. Then $t$ corresponds to a prime ideal
$\mathfrak p \subset R$ which is not maximal (as $t \not \in T_0$).
Thus (1) holds. By construction $f \not \in \mathfrak p$, hence (2).
As $t$ is a closed point of $T \cap D(f)$ we see that
$V(\mathfrak p) \cap D(f) = \{\mathfrak p\}$, i.e., (3) holds. Hence we
conclude that $(R/\mathfrak p)_f$ is a domain whose
spectrum has one point, hence (4) holds
(for example combine Lemmas \ref{editorial-source-lemma-characterize-local-ring} and
\ref{algebra-lemma-minimal-prime-reduced-ring}).

\medskip\noindent
Conversely, suppose $R$ is Jacobson and $(\mathfrak p,f)$ satisfies
(1) and (2). Suppose only finitely many maximal ideals
$\mathfrak m_1,\ldots,\mathfrak m_t$ lie in
$V(\mathfrak p)\cap D(f)$. Because $\mathfrak p$ is not maximal,
choose $g_i\in\mathfrak m_i\setminus\mathfrak p$ and let
$g=\prod_{i=1}^t g_i$, with $g=1$ when $t=0$.
Primality gives $fg\notin\mathfrak p$, so the locally closed subset
$V(\mathfrak p)\cap D(fg)$ is nonempty. Since $\Spec(R)$ is Jacobson,
it contains a point closed in $\Spec(R)$. That point would be one of
the listed maximal ideals, since $D(fg)\subset D(f)$, but $g$ belongs
to every listed ideal, excluding each of them. This contradiction
proves that there are infinitely many such maximal ideals.
\end{proof}


\subsection{Editorial treatment of Algebra lines 6778--6807}

\hypertarget{algebra-context-057}{}

\noindent\href{https://github.com/stacks/stacks-project/blob/a04446e57ec1fbc252a871afcec7752fb2807b14/algebra.tex\#L6778}{Original source, lines 6778--6807}. Complete original and reviewed TeX passages accompany this note. Every intervention is separately recorded; the following treatment is editorial.


\begin{lemma}
\label{editorial-source-lemma-pid-jacobson}
The ring $\mathbf{Z}$ is a Jacobson ring.
More generally, let $R$ be a ring such that
\begin{enumerate}
\item $R$ is a domain,
\item $R$ is Noetherian,
\item any nonzero prime ideal is a maximal ideal, and
\item $R$ has infinitely many maximal ideals.
\end{enumerate}
Then $R$ is a Jacobson ring. In (2), it suffices more generally
that $\Spec(R)$ is a Noetherian topological space.
\end{lemma}

\begin{proof}
Assume (1), (3), (4) and that $\Spec(R)$ is Noetherian. The statement
means that $(0) = \bigcap_{\mathfrak m \subset R} \mathfrak m$.
Since $R$ has infinitely many maximal ideals it suffices to
show that any nonzero $x\in R$ belongs to only finitely many
maximal ideals, or equivalently that $V(x)$ is finite.
By Lemma \ref{algebra-lemma-spec-closed}
we see that $V(x)$ is homeomorphic to $\Spec(R/xR)$.
By assumption (3) every prime of $R/xR$ is minimal and hence
corresponds to an irreducible component of $\Spec(R/xR)$
(Lemma \ref{editorial-source-lemma-irreducible}).
If $R$ is Noetherian, then $\Spec(R)$ is Noetherian by
Lemma \ref{editorial-source-lemma-Noetherian-topology}. More generally assume only
this topological condition. Its closed subspace $V(x)$ is Noetherian,
and therefore has finitely many irreducible components by Topology,
Lemma \ref{topology-lemma-Noetherian}. Every prime in $V(x)$ is
nonzero and hence maximal by (3), so each irreducible component is
a singleton. Thus $V(x)$ is finite. The infinitely many maximal
ideals cannot all contain this $x$, proving their intersection is zero.
The nonzero primes are already maximal by (3), so every prime is an
intersection of maximal ideals and Lemma \ref{algebra-lemma-jacobson-prime}
proves that $R$ is Jacobson.
\end{proof}


\subsection{Editorial treatment of Algebra lines 6809--6827}

\hypertarget{algebra-context-058}{}

\noindent\href{https://github.com/stacks/stacks-project/blob/a04446e57ec1fbc252a871afcec7752fb2807b14/algebra.tex\#L6809}{Original source, lines 6809--6827}. Complete original and reviewed TeX passages accompany this note. Every intervention is separately recorded; the following treatment is editorial.


\begin{example}
\label{editorial-source-example-infinite-product-fields-jacobson}
Let $A$ be an infinite set.
For each $\alpha \in A$, let $k_\alpha$ be a field.
We claim that $R = \prod_{\alpha\in A} k_\alpha$ is Jacobson.
For the original element $f=(f_\alpha)_\alpha$, define
$$
e_\alpha=\begin{cases}1&f_\alpha\ne0,\\0&f_\alpha=0,\end{cases}
\qquad
u_\alpha=\begin{cases}f_\alpha&f_\alpha\ne0,\\1&f_\alpha=0.\end{cases}
$$
Then $e=(e_\alpha)$ is idempotent, $u=(u_\alpha)$ is a unit with
inverse $(u_\alpha^{-1})$, and $f=ue$ coordinate by coordinate.
Thus $D(f)=D(e)$ and the canonical maps identify
$R_f\cong R_e\cong R/(1-e)$. The first pair of inverse maps comes
from inverting $f=ue$ with $u$ already a unit. For the second, an
invertible idempotent is $1$, so $R\to R_e$ kills $1-e$; conversely
$e$ becomes $1$ in $R/(1-e)$. These universal maps are inverse
because their composites fix the image of $R$ and the adjoined inverse.

In fact every prime of this product ring is maximal. In the domain
$R/\mathfrak p$, the image of an idempotent is either $0$ or $1$.
For any $f\notin\mathfrak p$, its decomposition $f=ue$ forces the
image of $e$ to be $1$, so the image of $f$ is the unit given by
the image of $u$. Every nonzero element of $R/\mathfrak p$ is
therefore a unit. All primes are maximal, and each radical ideal
is the intersection of primes containing it, proving that $R$ is Jacobson.
This also covers finite index sets and the empty product, whose
spectrum is empty. More generally, the following original argument
proves the Jacobson property whenever each canonical localization
$R\to R_f$ is surjective.
Namely, say $\mathfrak p \subset R$ is a prime ideal
and $f \in R$, $f \not \in \mathfrak p$. We have to find
a maximal ideal $\mathfrak m$ of $R$ such that
$\mathfrak p \subset \mathfrak m$ and $f \not\in \mathfrak m$.
Because $R_f$ is a quotient of $R$ we see that any maximal
ideal of $R_f$ corresponds to a maximal ideal of $R$
not containing $f$. Hence the result follows
by choosing a maximal ideal of $R_f$ containing $\mathfrak p R_f$.
\end{example}


\subsection{Editorial treatment of Algebra lines 6829--6842}

\hypertarget{algebra-context-059}{}

\noindent\href{https://github.com/stacks/stacks-project/blob/a04446e57ec1fbc252a871afcec7752fb2807b14/algebra.tex\#L6829}{Original source, lines 6829--6842}. Complete original and reviewed TeX passages accompany this note. Every intervention is separately recorded; the following treatment is editorial.


\begin{example}
\label{editorial-source-example-not-jacobson}
A domain $R$ with finitely many maximal ideals
$\mathfrak m_i$, $i = 1, \ldots, n$ is not a
Jacobson ring, except when it is a field.
If this domain is not a field, every maximal ideal $\mathfrak m_i$
is nonzero. Choose $0\ne a_i\in\mathfrak m_i$ for $1\leq i\leq n$.
The nonzero domain has at least one maximal ideal, so $n\geq1$.
The product $a_1\cdots a_n$ is nonzero and belongs to
$$
\mathfrak m_1\cdots\mathfrak m_n
\subset\mathfrak m_1\cap\cdots\cap\mathfrak m_n.
$$
Hence the radical ideal $(0)$ is not the intersection of the maximal
ideals containing it, proving the assertion. In particular a discrete
valuation ring is not Jacobson. More generally, in any local ring
having a prime $\mathfrak p$ distinct from its unique maximal ideal
$\mathfrak m$, the intersection of maximal ideals containing the
radical ideal $\mathfrak p$ is $\mathfrak m\ne\mathfrak p$.
Thus any local ring with at least two prime ideals is not Jacobson,
without a domain hypothesis in this last assertion.
\end{example}


\subsection{Editorial treatment of Algebra lines 6872--6886}

\hypertarget{algebra-context-060}{}

\noindent\href{https://github.com/stacks/stacks-project/blob/a04446e57ec1fbc252a871afcec7752fb2807b14/algebra.tex\#L6872}{Original source, lines 6872--6886}. Complete original and reviewed TeX passages accompany this note. Every intervention is separately recorded; the following treatment is editorial.


\begin{lemma}
\label{editorial-source-lemma-dimension}
Suppose that $k$ is a field and suppose that $V$ is a nonzero vector
space over $k$. Assume the dimension of $V$ (which is a cardinal number)
is smaller than the cardinality of $k$. Then for any linear operator
$T : V \to V$ there exists some monic polynomial $P(t) \in k[t]$ such that
$P(T)$ is not invertible.
\end{lemma}

\begin{proof}
Suppose every monic polynomial $P$ gives an invertible $P(T)$.
For an arbitrary nonzero $Q\in k[t]$ with leading coefficient $c$,
retain the exact comparison
$Q(T)=c\,(c^{-1}Q)(T)$, where $c^{-1}Q$ is monic.
Its inverse is $((c^{-1}Q)(T))^{-1}c^{-1}$, so every such $Q(T)$
is invertible, with its scalar factor retained.
Define the action of the original rational-function field by
$$
\rho:k(t)\longrightarrow\operatorname{End}_k(V),\qquad
p/q\longmapsto p(T)q(T)^{-1}.
$$
All polynomial operators commute, and their inverses commute with
them and with each other. If $p/q=p'/q'$, then $pq'=p'q$;
evaluation at $T$ and multiplication by $q(T)^{-1}q'(T)^{-1}$
prove equality of the proposed operators. The common-denominator
formulas prove additivity and multiplicativity, and $\rho(1)$ is
the identity. This is a $k(t)$-vector-space structure extending
the original $k$-structure on $V$. For $0\ne v\in V$, the map
$k(t)\to V$, $a\mapsto\rho(a)v$, is injective: if $a\ne0$,
multiplication by $\rho(a^{-1})$ recovers $v$.

\medskip\noindent
The elements $1/(t-\lambda)$ for $\lambda\in k$ are linearly
independent over $k$. Indeed, multiplying a finite relation
$\sum_{i=1}^r c_i/(t-\lambda_i)=0$ for distinct $\lambda_i$ by
the full product $\prod_{i=1}^r(t-\lambda_i)$ gives
$\sum_{i=1}^r c_i\prod_{j\ne i}(t-\lambda_j)=0$.
Evaluation at $t=\lambda_i$ gives
$c_i\prod_{j\ne i}(\lambda_i-\lambda_j)=0$, and every factor
in that product is nonzero. Thus every $c_i=0$.
Their images in $V$ are linearly independent, so
$\dim_k(V)\geq\#k$, contradicting the hypothesis.
\end{proof}


\subsection{Editorial treatment of Algebra lines 6891--6943}

\hypertarget{algebra-context-061}{}

\noindent\href{https://github.com/stacks/stacks-project/blob/a04446e57ec1fbc252a871afcec7752fb2807b14/algebra.tex\#L6891}{Original source, lines 6891--6943}. Complete original and reviewed TeX passages accompany this note. Every intervention is separately recorded; the following treatment is editorial.


\begin{theorem}
\label{editorial-source-theorem-uncountable-nullstellensatz}
Let $k$ be a field. Let $S$ be a $k$-algebra generated over $k$
by the elements $\{x_i\}_{i \in I}$. Assume the cardinality of $I$
is smaller than the cardinality of $k$. Then
\begin{enumerate}
\item for all maximal ideals $\mathfrak m \subset S$ the field
extension $\kappa(\mathfrak m)/k$
is algebraic, and
\item $S$ is a Jacobson ring.
\end{enumerate}
If $I$ is finite, both conclusions hold for every field $k$, without
the cardinality restriction on the finite generating set.
\end{theorem}

\begin{proof}
For finite $I$, both assertions follow from the Hilbert
Nullstellensatz (Theorem \ref{editorial-source-theorem-nullstellensatz}) over every
field, regardless of the size of that finite set. Now assume $I$
is infinite and put $\kappa=\#I<\#k$.
Let $P=k[\{X_i\}_{i\in I}]$ be the polynomial ring mapping onto
the original algebra $S$ by $X_i\mapsto x_i$. Every monomial has
finite support and is specified by a finite list of pairs
$(i,e)$ with $i\in I$, $e\in\mathbf Z_{\geq0}$.
For each finite length $r$, the set of such lists has cardinality
at most $\kappa$, since $\kappa$ is infinite; the union over all
$r\geq0$ still has cardinality $\kappa$.
Distinct variables supply $\kappa$ distinct monomials, so the
monomial basis of $P$ has cardinality exactly $\kappa$.
The images of these monomials span $S$ and every quotient $S/\mathfrak m$.
In particular, for a maximal ideal $\mathfrak m\subset S$, its
residue field is exactly $S/\mathfrak m$ and has $k$-dimension
at most $\kappa$. These are maps from the same original polynomial
ring and algebra, without replacing $S$ by $P$.

\medskip\noindent
To arrive at a contradiction pick $T \in \kappa(\mathfrak m)$ transcendental
over $k$. Note that the $k$-linear map
$T : \kappa(\mathfrak m) \to \kappa(\mathfrak m)$
given by multiplication by $T$ has the property that $P(T)$ is invertible
for all monic polynomials $P(t) \in k[t]$.
Also, $\kappa(\mathfrak m)$ has dimension at most the cardinality of $I$
over $k$ since it is a quotient of the vector space
$k[\{x_i\}_{i \in I}]$ over $k$ (whose dimension is $\# I$ as we saw above).
This is impossible by Lemma \ref{editorial-source-lemma-dimension}.

\medskip\noindent
To show that $S$ is Jacobson we argue as follows. If not then
there exists a prime $\mathfrak q \subset S$ and an element $f \in S$,
$f \not \in \mathfrak q$ such that $\mathfrak q$ is not maximal
and $(S/\mathfrak q)_f$ is a field, see
Lemma \ref{editorial-source-lemma-characterize-jacobson}.
The field $L=(S/\mathfrak q)_f$ is generated over $k$ by the
images of the original $x_i$ and the inverse of the image of $f$.
Thus it has a generating set of cardinality at most
$\kappa+1=\kappa<\#k$. Apply the first assertion already proved
to the maximal ideal $(0)$ of the field $L$: the extension $L/k$
is algebraic. The canonical inclusion
$S/\mathfrak q\subset L\subset\operatorname{Frac}(S/\mathfrak q)$
identifies $L$ with the entire fraction field, since a field containing
$S/\mathfrak q$ contains the inverse of every nonzero element of
that domain. Hence $L=\kappa(\mathfrak q)$ canonically.
The prime $\mathfrak q$ lies over the maximal ideal $(0)$ of $k$,
so Lemma \ref{algebra-lemma-finite-residue-extension-closed} makes
$\mathfrak q$ maximal in $S$, a contradiction.
\end{proof}


\subsection{Editorial treatment of Algebra lines 6945--6963}

\hypertarget{algebra-context-062}{}

\noindent\href{https://github.com/stacks/stacks-project/blob/a04446e57ec1fbc252a871afcec7752fb2807b14/algebra.tex\#L6945}{Original source, lines 6945--6963}. Complete original and reviewed TeX passages accompany this note. Every intervention is separately recorded; the following treatment is editorial.


\begin{lemma}
\label{editorial-source-lemma-base-change-Jacobson}
Let $k$ be a field. Let $S$ be a $k$-algebra.
For any field extension $K/k$ whose cardinality is larger
than the cardinality of $S$ we have
\begin{enumerate}
\item for every maximal ideal $\mathfrak m$ of $S_K$ the field
$\kappa(\mathfrak m)$ is algebraic over $K$, and
\item $S_K$ is a Jacobson ring.
\end{enumerate}
It suffices more generally that $S$ have a $k$-algebra generating
set of cardinality less than $\#K$. If this set is finite, no
cardinality restriction on $K$ is needed.
\end{lemma}

\begin{proof}
Fix the given field extension $K/k$ and write
$S_K=K\otimes_kS$. If $(s_i)_{i\in I}$ generates $S$ as a
$k$-algebra, then $(1\otimes s_i)_{i\in I}$ generates $S_K$ as a
$K$-algebra. Indeed every tensor is a finite sum of terms
$\lambda\otimes s$, and an expression for $s$ as a polynomial
in finitely many $s_i$ gives the corresponding expression for
$\lambda\otimes s$, with all coefficients carried by $k\to K$.
Theorem \ref{editorial-source-theorem-uncountable-nullstellensatz} therefore applies
whenever $\#I<\#K$, and its finite-generator case applies without
that restriction. Taking the generating family indexed by every
element of $S$ gives the hypotheses and conclusion stated originally.
\end{proof}


\subsection{Editorial treatment of Algebra lines 6965--6978}

\hypertarget{algebra-context-063}{}

\noindent\href{https://github.com/stacks/stacks-project/blob/a04446e57ec1fbc252a871afcec7752fb2807b14/algebra.tex\#L6965}{Original source, lines 6965--6978}. Complete original and reviewed TeX passages accompany this note. Every intervention is separately recorded; the following treatment is editorial.


\begin{example}
\label{editorial-source-example-countable-trick-does-not-work}
The trick in the proof of
Theorem \ref{editorial-source-theorem-uncountable-nullstellensatz}
really does not work if $k$ is a countable field and $I$ is countable too.
Let $k$ be a countable field. Let $x$ be a variable,
and let $k(x)$ be the field of rational functions in $x$.
Consider the polynomial algebra $R = k[x, \{x_f\}_{f \in k[x]-\{0\}}]$.
Let $I = (\{fx_f - 1\}_{f\in k[x] - \{0\}})$. Note that
Evaluation defines a map
$$
R\longrightarrow k(x),\qquad x\longmapsto x,\qquad
x_f\longmapsto f^{-1}\quad(0\ne f\in k[x]),
$$
which kills every original relation $fx_f-1$ and factors through
$R/I$. Conversely the map $k[x]\to R/I$ sends every nonzero
$f$ to a unit, with inverse the class of $x_f$. The localization
property therefore gives
$$
k(x)\longrightarrow R/I,\qquad p/q\longmapsto
\overline p\,\overline{x}_q\quad(q\ne0).
$$
These maps are inverse. On $k(x)$ their composite fixes all fractions;
on $R/I$ it fixes $\overline x$ and every $\overline{x}_f$, since
each is the unique inverse of $\overline f$. Thus $R/I\cong k(x)$
as $k$-algebras, so $I$ itself is maximal and any maximal ideal
containing $I$ is equal to $I$. This extension is transcendental.

The same presentation and inverse maps work over every field $k$.
For countable $k$, the polynomial set $k[x]$ is countably infinite,
so the original algebra has a countable generating family.
For infinite $k$, that family has cardinality exactly $\#k$, since
$\#k[x]=\max(\#k,\aleph_0)=\#k$ and the nonzero constants already
give $\#k$ distinct indices. Thus equality at the cardinal threshold
can fail to force algebraic residue extensions over every infinite
field, with the same explicit quotient.
\end{example}


\subsection{Editorial treatment of Algebra lines 7050--7087}

\hypertarget{algebra-context-064}{}

\noindent\href{https://github.com/stacks/stacks-project/blob/a04446e57ec1fbc252a871afcec7752fb2807b14/algebra.tex\#L7050}{Original source, lines 7050--7087}. Complete original and reviewed TeX passages accompany this note. Every intervention is separately recorded; the following treatment is editorial.


\begin{proposition}
\label{editorial-source-proposition-Jacobson-permanence}
Let $R$ be a Jacobson ring. Let $R \to S$ be a
ring map of finite type. Then
\begin{enumerate}
\item The ring $S$ is Jacobson.
\item The map $\Spec(S) \to \Spec(R)$ transforms
closed points to closed points.
\item For $\mathfrak m' \subset S$ maximal lying over $\mathfrak m \subset R$
the field extension $\kappa(\mathfrak m')/\kappa(\mathfrak m)$
is finite.
\end{enumerate}
\end{proposition}

\begin{proof}
Let $\mathfrak m' \subset S$ be a maximal ideal and
let $\varphi:R\to S$ be the original map and set
$\mathfrak m=\varphi^{-1}(\mathfrak m')$.
Then $R/\mathfrak m \to S/\mathfrak m'$ satisfies
the conditions of Lemma \ref{algebra-lemma-silly-jacobson}
by Lemma \ref{algebra-lemma-Jacobson-mod-ideal}.
Hence $R/\mathfrak m$ is a field and
$\mathfrak m$ a maximal ideal and the induced
residue field extension is finite. This proves (2) and (3).  

\medskip\noindent
If $S$ is not Jacobson, then by Lemma \ref{editorial-source-lemma-characterize-jacobson} there
exists a non-maximal prime ideal $\mathfrak q$ of $S$ and an
$g \in S$, $g \not\in \mathfrak q$ such that $(S/\mathfrak q)_g$ is a field.
To arrive at a contradiction we show that $\mathfrak q$ is a maximal ideal.
Let $\mathfrak p=\varphi^{-1}(\mathfrak q)$.
The map $R/\mathfrak p\to S/\mathfrak q$ is injective and of finite
type, using the images of the original finite set of algebra generators.
Adjoining the inverse of the image of $g$ adds one generator, so
$R/\mathfrak p\to(S/\mathfrak q)_g$ is also of finite type.
The target is a field and embeds canonically in
$\operatorname{Frac}(S/\mathfrak q)$, so it equals that fraction
field: it already contains the domain and all inverses of its nonzero
elements. Thus
$R/\mathfrak p \to (S/\mathfrak q)_g$ satisfies the conditions of
Lemma \ref{algebra-lemma-silly-jacobson} by
Lemma \ref{algebra-lemma-Jacobson-mod-ideal}.
Hence $R/\mathfrak p$ is a field and the field extension
$\kappa(\mathfrak p) \to (S/\mathfrak q)_g = \kappa(\mathfrak q)$ is
finite, thus algebraic. Then $\mathfrak q$ is a maximal ideal of $S$ by
Lemma \ref{algebra-lemma-finite-residue-extension-closed}. Contradiction.
\end{proof}


\subsection{Editorial treatment of Algebra lines 7099--7173}

\hypertarget{algebra-context-065}{}

\noindent\href{https://github.com/stacks/stacks-project/blob/a04446e57ec1fbc252a871afcec7752fb2807b14/algebra.tex\#L7099}{Original source, lines 7099--7173}. Complete original and reviewed TeX passages accompany this note. Every intervention is separately recorded; the following treatment is editorial.


\begin{lemma}
\label{editorial-source-lemma-image-finite-type-map-Jacobson-rings}
Let $R \to S$ be a finite type ring map of Jacobson rings.
Denote $X = \Spec(R)$ and $Y = \Spec(S)$.
Write $f : Y \to X$ for the induced
map of spectra. Let $E \subset Y = \Spec(S)$ be a
constructible set. Denote with a subscript ${}_0$ the set
of closed points of a topological space.
\begin{enumerate}
\item We have $f(E)_0 = f(E_0) = X_0 \cap f(E)$.
\item A point $\xi \in X$ is in $f(E)$ if and only if
$\overline{\{\xi\}} \cap f(E_0)$ is dense in $\overline{\{\xi\}}$.
\end{enumerate}
\end{lemma}

\begin{proof}
We have a commutative diagram of continuous maps
$$
\xymatrix{
E \ar[r] \ar[d] & Y \ar[d] \\
f(E) \ar[r] & X
}
$$
Suppose $x \in f(E)$ is closed in $f(E)$. Then $f^{-1}(\{x\})\cap E$
is nonempty and closed in $E$. Applying
Topology, Lemma \ref{topology-lemma-jacobson-inherited}
to both inclusions
$$
f^{-1}(\{x\}) \cap E \subset E \subset Y
$$
we find there exists a point $y \in f^{-1}(\{x\}) \cap E$ which is
closed in $Y$. In other words, there exists $y \in Y_0$ and $y \in E_0$
mapping to $x$. Hence $x \in f(E_0)$.
This proves that $f(E)_0 \subset f(E_0)$.
Proposition \ref{editorial-source-proposition-Jacobson-permanence} implies that
$f(E_0) \subset X_0 \cap f(E)$. The inclusion
$X_0 \cap f(E) \subset f(E)_0$ is trivial. This proves the
first assertion.

\medskip\noindent
Suppose that $\xi \in f(E)$. According to
Lemma \ref{editorial-source-lemma-characterize-image-finite-type}
the set $f(E) \cap \overline{\{\xi\}}$ contains a dense
open subset of $\overline{\{\xi\}}$. Since $X$ is Jacobson
we conclude that $f(E) \cap \overline{\{\xi\}}$ contains a
dense set of closed points, see Topology,
Lemma \ref{topology-lemma-jacobson-inherited}.
We conclude by part (1) of the lemma.

\medskip\noindent
Conversely, suppose $\overline{\{\xi\}}\cap f(E_0)$ is dense
in $\overline{\{\xi\}}$. By Lemma \ref{editorial-source-lemma-constructible-is-image}
choose a finitely presented map $S\to S'$ whose spectrum map
$g:Y'=\Spec(S')\to Y$ has image exactly $E$.
The ring $S'$ is Jacobson by Proposition \ref{editorial-source-proposition-Jacobson-permanence}.
Part (1) of the present lemma, already proved and applied to $S\to S'$,
gives $g(Y'_0)=E_0$. Put $h=f\circ g$, so
$h(Y'_0)=f(E_0)$ and $h(Y')=f(E)$, retaining the original $S$ and $E$.
If $\xi$ corresponds to $\mathfrak p\subset R$, use the exact diagram
$$
\xymatrix{
S'\ar[r] & S'/\mathfrak pS'\\
R\ar[r]\ar[u] & R/\mathfrak p\ar[u]
}
$$
Every point of $Y'_0$ whose image lies in $V(\mathfrak p)$ corresponds
to a prime of $S'$ containing $\mathfrak pS'$, hence to a point of
$\Spec(S'/\mathfrak pS')$. The spectrum image of
$R/\mathfrak p\to S'/\mathfrak pS'$ therefore contains the given
dense subset $h(Y'_0)\cap V(\mathfrak p)$ of
$V(\mathfrak p)=\Spec(R/\mathfrak p)$.
By Lemma \ref{algebra-lemma-domain-image-dense-set-points-generic-point},
there is a prime $\overline{\mathfrak q}'$ of $S'/\mathfrak pS'$
contracting to $(0)$ in $R/\mathfrak p$. Its inverse image
$\mathfrak q'\subset S'$ contracts to $\mathfrak p$ in $R$.
The point $g(\mathfrak q')$ belongs to $E$ and its image under the
original $f$ is $\xi$. Consequently $\xi\in f(E)$, as required.
\end{proof}


\subsection{Editorial treatment of Algebra lines 7189--7215}

\hypertarget{algebra-context-066}{}

\noindent\href{https://github.com/stacks/stacks-project/blob/a04446e57ec1fbc252a871afcec7752fb2807b14/algebra.tex\#L7189}{Original source, lines 7189--7215}. Complete original and reviewed TeX passages accompany this note. Every intervention is separately recorded; the following treatment is editorial.


\begin{lemma}
\label{editorial-source-lemma-conclude-jacobson-Noetherian}
With notation as above, assume that $R$ is a Noetherian Jacobson ring.
Further assume $R \to S$ is of finite type.
There is a commutative diagram
$$
\xymatrix{
\text{Constr}(Y) \ar[r]^{E \mapsto E_0} \ar[d]^{E \mapsto f(E)} &
\text{Constr}(Y_0) \ar[d]^{E \mapsto f(E)} \\
\text{Constr}(X) \ar[r]^{E \mapsto E_0} &
\text{Constr}(X_0)
}
$$
where the horizontal arrows are the bijections from
Topology, Lemma \ref{topology-lemma-jacobson-equivalent-constructible}.
The same diagram and conclusions hold more generally when $R$ is
Jacobson and $R\to S$ is of finite presentation, without requiring
$R$ to be Noetherian.
\end{lemma}

\begin{proof}
Under the original hypotheses, the finite type map $R\to S$ is of
finite presentation by
Lemma \ref{editorial-source-lemma-Noetherian-finite-type-is-finite-presentation}.
For the more general assertion, finite presentation is assumed directly.
In both cases it implies finite type, so $S$ is Jacobson by
Proposition \ref{editorial-source-proposition-Jacobson-permanence}. Therefore $X$ and
$Y$ are Jacobson spaces and the two horizontal maps are bijections
by Topology, Lemma \ref{topology-lemma-jacobson-equivalent-constructible},
which requires no Noetherian hypothesis.
Chevalley's theorem (Theorem \ref{editorial-source-theorem-chevalley}) states that
the image of a constructible subset of the source $Y$ is constructible
in the target $X$. For each such $E\subset Y$,
Lemma \ref{editorial-source-lemma-image-finite-type-map-Jacobson-rings} gives
$f(E_0)=f(E)_0$, proving commutativity.
The right vertical arrow is well defined as a direct-image map:
given a constructible subset $C\subset Y_0$, the top horizontal
bijection supplies a unique constructible $E\subset Y$ with $E_0=C$,
and then $f(C)=f(E)_0$ is constructible in $X_0$ by the bottom
horizontal bijection. This proves every arrow and the square in
the stated greater generality.
\end{proof}


\subsection{Editorial treatment of Algebra lines 7220--7361}

\hypertarget{algebra-context-067}{}

\noindent\href{https://github.com/stacks/stacks-project/blob/a04446e57ec1fbc252a871afcec7752fb2807b14/algebra.tex\#L7220}{Original source, lines 7220--7361}. Complete original and reviewed TeX passages accompany this note. Every intervention is separately recorded; the following treatment is editorial.


\begin{example}
\label{editorial-source-example-product-matrices-zero}
Let $k$ be a field. The space $\Spec(k[x, y]/(xy))$
has two irreducible components: namely the $x$-axis and the
$y$-axis. As a generalization, let
$$
R = k[x_{11}, x_{12}, x_{21}, x_{22}, y_{11}, y_{12}, y_{21}, y_{22}]/
\mathfrak a,
$$
where $\mathfrak a$ is the ideal in
$k[x_{11}, x_{12}, x_{21}, x_{22}, y_{11}, y_{12}, y_{21}, y_{22}]$
generated by the entries of the $2 \times 2$ product matrix
$$
\left(
\begin{matrix}
x_{11} & x_{12}\\
x_{21} & x_{22}
\end{matrix}
\right)
 \left(
\begin{matrix}
y_{11} & y_{12}\\
y_{21} & y_{22}
\end{matrix}
\right).
$$
In this example we will describe $\Spec(R)$.

\medskip\noindent
To prove the statement about $\Spec(k[x, y]/(xy))$ we argue as follows.
If $\mathfrak p\subset k[x,y]$ is any prime ideal containing $xy$,
then either $x$ or $y$ belongs to $\mathfrak p$. Hence the minimal such
prime ideals are just $(x)$ and $(y)$. In case $k$ is
algebraically closed, the $\text{max-Spec}$ of these components
can then be visualized as the point sets of $y$- and $x$-axis.

\medskip\noindent
For the generalization, note that we may identify the closed
points of the spectrum of
$k[x_{11}, x_{12}, x_{21}, x_{22}, y_{11}, y_{12}, y_{21}, y_{22}]$
with the space of matrices
$$
\left\{ (X, Y) \in \text{Mat}(2, k)\times \text{Mat}(2, k) \mid
X = \left(
\begin{matrix}
x_{11} & x_{12}\\
x_{21} & x_{22}
\end{matrix}
\right),
Y= \left(
\begin{matrix}
y_{11} & y_{12}\\
y_{21} & y_{22}
\end{matrix}
\right)
\right\}
$$
if $k$ is algebraically closed. We assume this for the following
description in terms of closed $k$-points. The coordinate-ring proof
below establishes the component statement over every field $k$.
Now define a group action of
$\text{GL}(2, k)\times \text{GL}(2, k)\times \text{GL}(2, k)$
on the space of matrices $\{(X, Y)\}$ by
$$
(g_1, g_2, g_3) \times (X, Y) \mapsto ((g_1Xg_2^{-1}, g_2Yg_3^{-1})).
$$
Here, also observe that the algebraic set
$$
\text{GL}(2, k)\times \text{GL}(2, k)\times \text{GL}(2, k) \subset
\text{Mat}(2, k)\times \text{Mat}(2, k) \times \text{Mat}(2, k)
$$
is irreducible since it is the max spectrum of the domain
$$
k[x_{11}, x_{12}, \ldots, z_{21}, z_{22}, (x_{11}x_{22}-x_{12}x_{21})^{-1}
, (y_{11}y_{22}-y_{12}y_{21})^{-1}, (z_{11}z_{22}-z_{12}z_{21})^{-1}].
$$
Since the image of an irreducible algebraic set is still
irreducible, the orbit closures below are irreducible. There are
six orbits, indexed by the pairs of ranks $(r,s)$ with $r+s\leq2$:
$$
(2,0),\quad(1,1),\quad(0,2),\quad(1,0),\quad(0,1),\quad(0,0).
$$
Indeed $XY=0$ says $\operatorname{Im}(Y)\subset\ker(X)$.
Choose a basis of the middle space starting with a basis of
$\operatorname{Im}(Y)$, extend it to $\ker(X)$ and then to the
whole middle space. Lift a basis of $\operatorname{Im}(Y)$ through
$Y$ and extend those lifts by a basis of $\ker(Y)$ in the source.
In the target extend the images under $X$ of the remaining middle
basis vectors to a basis. These choices show that the ranks determine
the orbit. If their basis matrices in the original coordinates are
$C_1,C_2,C_3$, respectively, then
$g_1=C_1^{-1}$, $g_2=C_2^{-1}$, $g_3=C_3^{-1}$ give exactly
the displayed action, with transformed matrices
$C_1^{-1}XC_2$ and $C_2^{-1}YC_3$. It preserves the original
equations because the transformed product is $g_1XYg_3^{-1}$,
and the inverse transformation uses $(g_1^{-1},g_2^{-1},g_3^{-1})$.
The three cases according to the rank of $X$ are:
\begin{enumerate}
\item $\exists (g_1, g_2)$ such that $g_1Xg_2^{-1} = I_{2\times 2}$.
Then $Y$ is necessarily $0$, which as an algebraic set is
invariant under the group action. It follows that this orbit is
contained in the irreducible algebraic set defined by the prime
ideal $(y_{11}, y_{12}, y_{21}, y_{22})$. Taking the closure, we see
that $(y_{11}, y_{12}, y_{21}, y_{22})$ is actually a component.
\item $\exists (g_1, g_2)$ such that
$$
g_1Xg_2^{-1} = \left(
\begin{matrix}
1 & 0 \\
0 & 0
\end{matrix}
\right).
$$
This case occurs if and only if $X$ is a rank 1 matrix,
and furthermore, $Y$ is killed by such an $X$ if and only if
$$
x_{11}y_{11}+x_{12}y_{21} = 0; \quad x_{11}y_{12}+x_{12}y_{22} = 0;
$$
$$
x_{21}y_{11}+x_{22}y_{21} = 0; \quad x_{21}y_{12}+x_{22}y_{22} = 0.
$$
Fix a rank $1$ matrix $X$. The nonzero matrices $Y$ satisfying these
equations form an irreducible locally closed set, for the following
reason (the pair $(X,0)$ belongs to the closure of the first orbit):
$0 = g_1Xg_2^{-1}g_2Y$ implies that
$$
g_2Y = \left(
\begin{matrix}
0 & 0 \\
y_{21}' & y_{22}'
\end{matrix}
\right).
$$
With a further $\text{GL}(2, k)$-action on the right by $g_3$,
$g_2Y$ can be brought into
$$
g_2Yg_3^{-1} = \left(
\begin{matrix}
0 & 0 \\
0 & 1
\end{matrix}
\right),
$$
and thus such $Y$'s form an irreducible algebraic set
isomorphic to the image of $\text{GL}(2, k)$ under this action. Finally,
notice that the ``rank 1" condition for $X$'s forms an open dense
subset of the irreducible algebraic set
$\det X = x_{11}x_{22} - x_{12}x_{21} = 0$.
On the open set where both matrices are nonzero, the four displayed
product equations and $\det X=0$ describe the rank-$(1,1)$ orbit.
To identify its closure, retain the polynomial ring
$B=k[x_{11},x_{12},x_{21},x_{22},y_{11},y_{12},y_{21},y_{22}]$
and put
$$
\begin{aligned}
\mathfrak q_M=(
&x_{11}y_{11}+x_{12}y_{21},\ x_{11}y_{12}+x_{12}y_{22},\\
&x_{21}y_{11}+x_{22}y_{21},\ x_{21}y_{12}+x_{22}y_{22},\\
&x_{11}x_{22}-x_{12}x_{21},\ y_{11}y_{22}-y_{12}y_{21}).
\end{aligned}
$$
The proof below shows that this is prime and cuts out that closure.
The sixth equation is already a consequence of the first four on
the open set where $X$ is nonzero: writing $a_{ij}=(XY)_{ij}$,
the exact matrix identity $X\det(Y)=(XY)\operatorname{adj}(Y)$ gives
$$
\begin{aligned}
x_{11}\det Y&=a_{11}y_{22}-a_{12}y_{21},&
x_{12}\det Y&=-a_{11}y_{12}+a_{12}y_{11},\\
x_{21}\det Y&=a_{21}y_{22}-a_{22}y_{21},&
x_{22}\det Y&=-a_{21}y_{12}+a_{22}y_{11}.
\end{aligned}
$$
Thus on each standard open where an entry of $X$ is invertible, the
original five-equation ideal and $\mathfrak q_M$ give the same quotient.
No sixth equation is discarded in the closure.
\item $\exists (g_1, g_2)$ such that $g_1Xg_2^{-1} = 0$, or
equivalently, $X = 0$. Then $Y$ can be arbitrary and this component
is thus defined by $(x_{11}, x_{12}, x_{21}, x_{22})$.
\end{enumerate}

\medskip\noindent
Here is the coordinate-ring proof, valid over the original arbitrary
field $k$. Keep $B$ and the four-entry ideal $\mathfrak a$ of the
original product $XY$, so $R=B/\mathfrak a$. Put
$$
\mathfrak q_X=(x_{11},x_{12},x_{21},x_{22}),\qquad
\mathfrak q_Y=(y_{11},y_{12},y_{21},y_{22}),
$$
and retain $\mathfrak q_M=\mathfrak a+(\det X,\det Y)$ as displayed
above. The first two ideals are prime, since their quotient rings
are polynomial rings in the four remaining entries.

For the middle ideal use the exact rearrangement of the original entries
$$
M=\begin{pmatrix}
x_{11}&x_{21}&y_{21}&y_{22}\\
x_{12}&x_{22}&-y_{11}&-y_{12}
\end{pmatrix}.
$$
Its six $2\times2$ minors, in column order
$(12),(13),(14),(23),(24),(34)$, are
$$
\det X,\quad-a_{11},\quad-a_{12},\quad-a_{21},\quad-a_{22},\quad\det Y.
$$
The entries of $M$ are eight independent polynomial coordinates:
the inverse rearrangement recovers each original entry, including
$y_{11}=-M_{23}$ and $y_{12}=-M_{24}$. Define
$$
\Phi:B\longrightarrow k[\lambda,\mu,c_1,c_2,c_3,c_4]
$$
by
$$
\begin{aligned}
x_{11}&\mapsto\lambda c_1,&x_{12}&\mapsto\mu c_1,&
x_{21}&\mapsto\lambda c_2,&x_{22}&\mapsto\mu c_2,\\
y_{11}&\mapsto-\mu c_3,&y_{12}&\mapsto-\mu c_4,&
y_{21}&\mapsto\lambda c_3,&y_{22}&\mapsto\lambda c_4.
\end{aligned}
$$
Thus the two rows of $M$ map to $(\lambda c_i)_i$ and
$(\mu c_i)_i$. Every minor maps to zero, proving
$\mathfrak q_M\subset\ker\Phi$.

To prove equality, write $s_i=M_{1i}$ and $t_i=M_{2i}$ for
$1\leq i\leq4$. The image of a monomial
$\prod_i s_i^{u_i}t_i^{v_i}$ is
$\lambda^{\sum_i u_i}\mu^{\sum_i v_i}\prod_i c_i^{u_i+v_i}$.
Two monomials have the same image exactly when all their column
totals $u_i+v_i$ agree and their first-row totals $\sum_i u_i$ agree.
For two such monomials with first-row exponents $u_i,u'_i$, if
$u\ne u'$ choose $i,j$ with $u_i>u'_i$ and $u_j<u'_j$.
Then $s_i$ and $t_j$ divide the first monomial. Replacing the factor
$s_it_j$ by $t_is_j$ changes it by a multiple of the minor
$s_it_j-t_is_j$, preserves all row and column totals, and decreases
$\sum_i|u_i-u'_i|$ by $2$. Repetition reaches the second monomial,
proving their difference belongs to the minor ideal.
For an arbitrary polynomial in $\ker\Phi$, group its finitely many
terms according to their image monomial. Distinct image monomials
are linearly independent over $k$, so the sum of coefficients in
each group is zero. Choosing one monomial in each group expresses
that whole group as a $k$-linear combination of the differences
just proved to lie in the minor ideal. Consequently
$\ker\Phi=\mathfrak q_M$. The induced map
$B/\mathfrak q_M\hookrightarrow k[\lambda,\mu,c_1,c_2,c_3,c_4]$
is injective, so $\mathfrak q_M$ is prime over every field,
with all original signs retained also in characteristic $2$.

Now let $\mathfrak p\subset B$ be any prime containing
$\mathfrak a$. In the field $\operatorname{Frac}(B/\mathfrak p)$,
the original matrices satisfy $XY=0$. If $\det X\notin\mathfrak p$,
then $X$ is invertible over this field and $Y=0$, so
$\mathfrak q_Y\subset\mathfrak p$. If
$\det Y\notin\mathfrak p$, the same argument gives
$\mathfrak q_X\subset\mathfrak p$. If both determinants belong
to $\mathfrak p$, then $\mathfrak q_M\subset\mathfrak p$.
These three prime ideals are pairwise incomparable.
For example $x_{11}\in\mathfrak q_X\setminus\mathfrak q_M$
since $\Phi(x_{11})=\lambda c_1\ne0$, while
$\det Y\in\mathfrak q_M\setminus\mathfrak q_X$ since it remains
a nonzero polynomial in the quotient by $\mathfrak q_X$.
The analogous facts with $X,Y$ exchanged prove incomparability
with $\mathfrak q_Y$, and the two coordinate ideals are
incomparable by their distinct sets of variables.
Therefore these are exactly the three minimal primes over
$\mathfrak a$. Their images in $R$ define exactly the three
irreducible components of $\Spec(R)$.

Finally, over an algebraically closed field the rank-$(2,0)$ orbit
is the nonempty open $\det X\ne0$ in $V(\mathfrak q_Y)$ and is
dense there; the rank-$(0,2)$ orbit is dense in $V(\mathfrak q_X)$
by the identical argument. In the irreducible $V(\mathfrak q_M)$,
the open where both $X$ and $Y$ are nonzero is precisely the
rank-$(1,1)$ orbit. It is nonempty, as witnessed by
$$
X=\begin{pmatrix}1&0\\0&0\end{pmatrix},\qquad
Y=\begin{pmatrix}0&0\\0&1\end{pmatrix},
$$
and hence dense. Density of closed points in these affine finite
type spectra follows from the Jacobson results above, so these
are the same closures in the stated description by closed points.
The other three rank orbits lie in these components and give no
additional irreducible component.
\end{example}


\subsection{Editorial treatment of Algebra lines 7364--7396}

\hypertarget{algebra-context-068}{}

\noindent\href{https://github.com/stacks/stacks-project/blob/a04446e57ec1fbc252a871afcec7752fb2807b14/algebra.tex\#L7364}{Original source, lines 7364--7396}. Complete original and reviewed TeX passages accompany this note. Every intervention is separately recorded; the following treatment is editorial.


\begin{example}
\label{editorial-source-example-idempotent-matrices}
For another example, consider
$R = k[\{t_{ij}\}_{i, j = 1}^{n}]/\mathfrak a$, where $\mathfrak a$ is
the ideal generated by the entries of the product matrix $T^2-T$,
$T = (t_{ij})$. From linear algebra, we know that under the
$GL(n, k)$-action defined by $g, T \mapsto gTg^{-1}$, $T$ is
classified by its rank and each $T$ is conjugate to some
$\text{diag}(1, \ldots, 1, 0, \ldots, 0)$, which has $r$ 1's and $n-r$ 0's.
Every idempotent matrix with entries in $k$ belongs to exactly one
of these rank orbits. Indeed each vector has the decomposition
$v=Tv+(v-Tv)$ into $\operatorname{Im}(T)$ and $\ker(T)$;
their intersection is zero since $Tw=w$ on the image and $Tw=0$
on the kernel. Bases of these two subspaces give the stated
conjugacy, and conjugation preserves rank. To deduce the irreducible
components of $\operatorname{Spec}(R)$, including its nonclosed
points, we give a coordinate-ring proof below. It proves that the
rank loci are disjoint open and closed irreducible components over
every field $k$. To separate ranks in characteristic zero one may use
$f(t_{ij}) = trace(T) = \sum_{i = 1}^nt_{ii}$. In positive
characteristic cases, things are slightly more tricky since we
might have $trace(T) = 0$ even if $T \neq 0$. For instance, if $k$ has characteristic $3$,
$$
trace\left(
\begin{matrix}
1 & & \\
& 1 & \\
& & 1
\end{matrix}
\right) = 3 = 0
$$
For every $0\leq j\leq n$ put
$e_j(T)=\operatorname{tr}(\bigwedge^jT)$, with $e_0(T)=1$.
On a rank-$r$ idempotent this equals $\binom rj$ in $k$:
in a basis of the image followed by the kernel, the exterior map
has eigenvalue $1$ on the $j$-fold wedges of image basis vectors
and $0$ on all other basis wedges. The exterior basis change is
invertible, so it preserves the trace of the original exterior map.
Thus if $r<s$, $e_s$ is zero on rank $r$ and one on rank $s$.
The complete tuple $(e_1,\ldots,e_n)$ separates all ranks in every
characteristic. A single chosen entry need not do so: in
characteristic $3$, ranks $3$ and $4$ both give $e_3=1$, since
$\binom33=1$ and $\binom43=4=1$ in $k$.
The displayed function $e_3$ isolates rank $3$ when $n=3$;
it does not isolate that rank for arbitrary $n$.
\medskip\noindent
We now prove the component assertions in the original ring $R$,
without assuming that $k$ is algebraically closed or that $R$ is
reduced. Keep the original matrix $T$ and define its coefficients
$d_0,\ldots,d_n\in R$ by
$$
F(z)=\det(I_n-T+zT)=\sum_{r=0}^n d_rz^r.
$$
The identity $T^2=T$ gives, in matrices over $R[x,y]$,
$$
(I_n-T+xT)(I_n-T+yT)=I_n-T+xyT.
$$
Hence $F(x)F(y)=F(xy)$. Comparing the coefficient of each
$x^ry^s$ proves $d_r^2=d_r$ and $d_rd_s=0$ for $r\ne s$;
evaluation at $z=1$ gives $\sum_r d_r=1$.
Put $R_r=R/(1-d_r)$ and $X_r=\operatorname{Spec}(R_r)$.
Inverting the idempotent $d_r$ is the same as imposing $d_r=1$,
so $X_r=D(d_r)=V(1-d_r)$ in $\operatorname{Spec}(R)$.
In particular these loci are open and closed. The exact ring maps
$$
R\longrightarrow\prod_{r=0}^n R_r,
\qquad a\longmapsto(a\bmod(1-d_r))_r,
\qquad
(\overline{a_r})_r\longmapsto\sum_{r=0}^n d_ra_r
$$
are inverse: the second expression does not depend on lifts since
$d_r(1-d_r)=0$, and orthogonality and the sum identity verify both
composites on every element.

For any prime $\mathfrak p\subset R$, apply the image/kernel
decomposition above to the original matrix over $\kappa(\mathfrak p)$.
If its rank is $r$, its determinant polynomial is exactly $z^r$.
Thus $d_s$ has residue $1$ for $s=r$ and $0$ for $s\ne r$.
It follows that $X_r$ is precisely the residue-field rank-$r$
locus, including all nonclosed points. Each $X_r$ is nonempty:
evaluation at $P_r=\operatorname{diag}(1,\ldots,1,0,\ldots,0)$,
with exactly $r$ ones, is a $k$-point.

For irreducibility take
$A=k[g_{ij}:1\leq i,j\leq n][\det(G)^{-1}]$,
where $G=(g_{ij})$. This is an integral domain. The formula
$T\mapsto GP_rG^{-1}$ gives the coordinate-ring homomorphism
$$
R_r\longrightarrow A,\qquad
t_{ij}\longmapsto\sum_{a=1}^r g_{ia}(G^{-1})_{aj}.
$$
The inverse entries are in $A$ by the adjugate formula; the
substituted matrix is idempotent and has determinant polynomial
$z^r$, so the displayed homomorphism is well defined. Its spectrum
map $\operatorname{Spec}(A)\to X_r$ is surjective. In fact for
any $\mathfrak p\in X_r$, bases of the image and kernel over
$\kappa(\mathfrak p)$ provide an invertible matrix $G_{\mathfrak p}$
with $T_{\mathfrak p}=G_{\mathfrak p}P_rG_{\mathfrak p}^{-1}$.
Evaluation $A\to\kappa(\mathfrak p)$ at these entries has a
prime kernel whose contraction to $R_r$ is the given prime,
because its composite sends every original $t_{ij}$ to its
residue at $\mathfrak p$.
The continuous image of an irreducible space is irreducible:
a union of two relatively closed proper subsets would pull back
to two closed subsets covering the source, and surjectivity makes
both proper. Therefore each $X_r$ is irreducible. An irreducible
subset cannot meet two members of the disjoint open-and-closed
cover $(X_r)$, since that would separate it into two nonempty
relatively closed subsets. The same separation excludes a connected
subset from meeting two members. Each $X_r$ is itself irreducible
and connected. Hence the irreducible components and the connected
components of $\operatorname{Spec}(R)$ are exactly $X_0,\ldots,X_n$.

The same original coordinates give useful explicit charts.
For a subset $I\subseteq\{1,\ldots,n\}$ let $P_I$ be the diagonal
matrix with entry $1$ at indices in $I$ and $0$ elsewhere, and set
$$
Q_I=TP_I+(I_n-T)(I_n-P_I),\qquad q_I=\det Q_I.
$$
Multilinearity in the columns of $F(z)$, retaining their original
order, gives
$$
d_r=\sum_{|I|=r}q_I:
$$
the term indexed by $I$ chooses the column $Te_j$ for $j\in I$
and $(I_n-T)e_j$ for $j\notin I$.
Thus $X_r$ is covered by the opens $D(q_I)$ with $|I|=r$.
To check that each such open lies in $X_r$, compute
$TQ_I=Q_IP_I$. After inverting $q_I$ this gives
$T=Q_IP_IQ_I^{-1}$, so $F(z)=z^r$ there.

Write $K=Q_I^{-1}$ in this localization. From $KT=P_IK$,
the definition of $Q_I$ gives the exact identity
$$
I_n=KQ_I=P_IKP_I+(I_n-P_I)K(I_n-P_I).
$$
Consequently the two diagonal blocks of $K$, indexed by $I$ and
its complement without changing the original row and column labels,
are identity matrices. Introduce variables $u_{ab}$ for the ordered
pairs $(a,b)$ with exactly one of $a,b$ in $I$. Let $K_I(u)$ have
these off-diagonal entries and the two identity diagonal blocks.
Put $h_I(u)=\det K_I(u)$ and
$$
B_I=k[u_{ab}: \text{exactly one of }a,b\text{ belongs to }I]
[h_I(u)^{-1}].
$$
There are $2r(n-r)$ such variables and $h_I(0)=1$.
The matrix $T'=K_I(u)^{-1}P_IK_I(u)$ is idempotent, and
$$
\begin{aligned}
Q_I(T')
&=K_I(u)^{-1}\bigl(P_IK_I(u)P_I
 +(I_n-P_I)K_I(u)(I_n-P_I)\bigr)\\
&=K_I(u)^{-1}.
\end{aligned}
$$
Thus $q_I(T')=h_I(u)^{-1}$ is invertible. Sending each original
$t_{ab}$ to the corresponding entry of $T'$ defines
$R[q_I^{-1}]\to B_I$. Conversely, sending each $u_{ab}$ to the
corresponding entry of $Q_I^{-1}$ defines $B_I\to R[q_I^{-1}]$:
the diagonal blocks agree by the identity just proved, and
$h_I$ maps to $q_I^{-1}$. These homomorphisms are inverse.
On the original matrix the composite is $Q_IP_IQ_I^{-1}=T$;
on the matrix of chart variables it is $Q_I(T')^{-1}=K_I(u)$.
The inverted determinants have the corresponding inverse images.
We have therefore proved the exact chart isomorphism
$$
R[q_I^{-1}]\cong B_I.
$$

In particular every chart ring is an integral domain and is reduced.
These charts prove that $R$ itself is reduced. If $a\in R$ is
nilpotent, its image is zero in each chart, so for each $I$ some
positive power $q_I^{N_I}$ annihilates $a$. The finitely many
$q_I$ cover the spectrum, since their sum over all $I$ is
$\sum_r d_r=1$. The powers $q_I^{N_I}$ also generate the unit
ideal: a maximal ideal containing them would contain every $q_I$,
contrary to that sum identity. Taking a linear combination equal
to $1$ proves $a=0$.
Each factor $R_r$ is reduced and has irreducible spectrum, hence is
a domain: if $ab=0$ and neither factor were zero, reducedness would
make both $D(a)$ and $D(b)$ nonempty, whereas their intersection is
$D(ab)=\varnothing$, contradicting irreducibility.
All constructions commute with every field extension $L/k$:
$R\otimes_k L=L[t_{ij}]/((T^2-T)_{ij})$, with the same coefficient
polynomials $d_r$ and chart maps. Repeating the argument over $L$
shows that each $R_r\otimes_k L$ is a domain. Thus each component
is geometrically integral, in the sense of being integral after
every extension of the original base field.

We can also identify its defining ideal exactly. Let $J_\ell(T)$
be the ideal generated by the $\ell\times\ell$ minors, with
$J_0(T)=R$ and $J_{n+1}(T)=0$. The diagonal entries of
$\bigwedge^jT$ in its indexed exterior basis are the principal
$j\times j$ minors. Expanding the determinant therefore gives
$$
\det(I_n+uT)=\sum_{j=0}^n e_j(T)u^j,
\qquad
d_s=\sum_{j=s}^n(-1)^{j-s}\binom js e_j(T).
$$
For $s\geq\ell$, every minor appearing in the latter sum belongs
to $J_\ell(T)$, by Laplace expansion along any fixed $\ell$ of
its rows. Hence $(d_\ell,\ldots,d_n)\subseteq J_\ell(T)$.
Conversely, on a chart of $X_s$ with $s<\ell$, the identity
$T=Q_IP_IQ_I^{-1}$ gives
$$
\bigwedge^\ell T
=(\bigwedge^\ell Q_I)(\bigwedge^\ell P_I)
 (\bigwedge^\ell Q_I^{-1})=0,
$$
since $P_I$ has only $s$ nonzero diagonal entries.
All $\ell\times\ell$ minors therefore vanish on each chart of
$X_s$, hence in $R_s$ by the same finite-cover annihilator argument.
In the proved product decomposition of $R$, the ideal of elements
whose coordinates for $s<\ell$ are zero is exactly
$(d_\ell,\ldots,d_n)$. This proves
$$
J_\ell(T)=(d_\ell,\ldots,d_n)
\quad(0\leq\ell\leq n+1).
$$
For the complementary idempotent matrix the Laurent-polynomial
identity
$$
\det(T+z(I_n-T))
=z^n\det(I_n-T+z^{-1}T)
$$
shows that $d_s(I_n-T)=d_{n-s}(T)$. Applying the minor formula
to both matrices gives the exact component ideal
$$
(1-d_r)=J_{r+1}(T)+J_{n-r+1}(I_n-T).
$$
Indeed the right side is generated by all $d_s$ with $s\ne r$;
their sum is $1-d_r$ and each satisfies $d_s(1-d_r)=d_s$.
These are also precisely the minimal prime ideals of $R$, since
each quotient $R_r$ is a domain and the $X_r$ are all components.
Finally $F(1+u)=\det(I_n+uT)$ gives
$$
e_j(T)=\sum_{r=0}^n\binom rj d_r,
$$
an equality in the original ring, strengthening the residue-field
trace calculation. Every integer coefficient is mapped to $k$;
no characteristic restriction has been introduced. If $n=0$,
empty determinants are $1$, $R=k$, $d_0=1$, and the same statements
hold with their empty matrices and lists.
\end{example}


\subsection{Editorial treatment of Algebra lines 7455--7467}

\hypertarget{algebra-context-069}{}

\noindent\href{https://github.com/stacks/stacks-project/blob/a04446e57ec1fbc252a871afcec7752fb2807b14/algebra.tex\#L7455}{Original source, lines 7455--7467}. Complete original and reviewed TeX passages accompany this note. Every intervention is separately recorded; the following treatment is editorial.


\begin{lemma}
\label{editorial-source-lemma-characterize-integral-element}
Let $\varphi : R \to S$ be a ring map. Let $y \in S$. If there exists a
finite $R$-submodule $M$ of $S$ such that $1 \in M$ and $yM \subset M$,
then $y$ is integral over $R$.
Conversely, if $y$ is integral, the original subalgebra
$\varphi(R)[y]\subset S$ is a finite $R$-submodule containing $1$
and stable under multiplication by $y$.
\end{lemma}

\begin{proof}
Keep $\varphi:R\to S$ for the original coefficient map and write
$m_y:M\to M$, $x\mapsto yx$, for the $R$-linear endomorphism.
By Lemma \ref{editorial-source-lemma-charpoly-module} there is a monic
$P(T)=T^d+\sum_{i=1}^d a_iT^{d-i}\in R[T]$ with $P(m_y)=0$.
Evaluation at the original element $1\in M$ gives in $S$
$$
P^\varphi(y)=y^d+\sum_{i=1}^d\varphi(a_i)y^{d-i}
=P(m_y)(1)=0.
$$
Conversely, such a monic equation expresses $y^d$ as
$-\sum_{i=1}^d\varphi(a_i)y^{d-i}$. Multiplying this equation by
each needed power of $y$ and using induction expresses every power
of $y$ in the $R$-span of $1,y,\ldots,y^{d-1}$.
That span is exactly $\varphi(R)[y]$, contains $1$ and is stable
under $m_y$. If $S$ is the zero ring the same conclusion holds with
the zero submodule; a degree-one equation may be used.
\end{proof}


\subsection{Editorial treatment of Algebra lines 7480--7500}

\hypertarget{algebra-context-070}{}

\noindent\href{https://github.com/stacks/stacks-project/blob/a04446e57ec1fbc252a871afcec7752fb2807b14/algebra.tex\#L7480}{Original source, lines 7480--7500}. Complete original and reviewed TeX passages accompany this note. Every intervention is separately recorded; the following treatment is editorial.


\begin{lemma}
\label{editorial-source-lemma-characterize-integral}
Let $\varphi : R \to S$ be a ring map. Let $s_1, \ldots, s_n$
be a finite set of elements of $S$.
In this case $s_i$ is integral over $R$ for all $i = 1, \ldots, n$
if and only if
there exists an $R$-subalgebra $S' \subset S$ finite over $R$
containing all of the $s_i$.
\end{lemma}

\begin{proof}
For each original $s_i$ choose its monic equation
$s_i^{d_i}+\sum_{j=0}^{d_i-1}\varphi(a_{ij})s_i^j=0$ with
$d_i\geq1$. Consider the finite list of original monomials
$s_1^{e_1}\cdots s_n^{e_n}$ with $0\leq e_i<d_i$.
Their $R$-span contains $1$ and is stable under each $s_i$:
if multiplication reaches exponent $d_i$, substitute the full equation
$s_i^{d_i}=-\sum_{j=0}^{d_i-1}\varphi(a_{ij})s_i^j$, keeping every
coefficient and all the other exponents. Thus this span is precisely
$\varphi(R)[s_1,\ldots,s_n]$ and is generated by at most
$\prod_{i=1}^n d_i$ elements. For $n=0$ the list is the single
empty monomial $1$ and the subalgebra is $\varphi(R)$.
Conversely, suppose given a finite $R$-subalgebra
$S'$ containing all the $s_i$. Then all of the
$s_i$ are integral by Lemma \ref{algebra-lemma-finite-is-integral}.
\end{proof}


\subsection{Editorial treatment of Algebra lines 7502--7515}

\hypertarget{algebra-context-071}{}

\noindent\href{https://github.com/stacks/stacks-project/blob/a04446e57ec1fbc252a871afcec7752fb2807b14/algebra.tex\#L7502}{Original source, lines 7502--7515}. Complete original and reviewed TeX passages accompany this note. Every intervention is separately recorded; the following treatment is editorial.


\begin{lemma}
\label{editorial-source-lemma-characterize-finite-in-terms-of-integral}
Let $R \to S$ be a ring map. The following are equivalent
\begin{enumerate}
\item $R \to S$ is finite,
\item $R \to S$ is integral and of finite type, and
\item there exist $x_1, \ldots, x_n \in S$ which generate $S$ as an
algebra over $R$ such that each $x_i$ is integral over $R$.
\end{enumerate}
\end{lemma}

\begin{proof}
If $S$ is finite as an $R$-module, a finite module generating
family also generates it as an $R$-algebra: the latter subalgebra
already contains every linear combination of those elements.
Lemma \ref{algebra-lemma-finite-is-integral} proves integrality, giving
(1) $\Rightarrow$ (2). Under (2), choose the finite algebra
generators; integrality holds for each, proving (3).
Under (3), Lemma \ref{editorial-source-lemma-characterize-integral} shows that the
subalgebra generated by those same elements is finite over $R$.
It is the original $S$, so (1) follows.
\end{proof}


\subsection{Editorial treatment of Algebra lines 7541--7554}

\hypertarget{algebra-context-072}{}

\noindent\href{https://github.com/stacks/stacks-project/blob/a04446e57ec1fbc252a871afcec7752fb2807b14/algebra.tex\#L7541}{Original source, lines 7541--7554}. Complete original and reviewed TeX passages accompany this note. Every intervention is separately recorded; the following treatment is editorial.


\begin{lemma}
\label{editorial-source-lemma-integral-closure-is-ring}
Let $R \to S$ be a ring homomorphism.
The set
$$
S' = \{s \in S \mid s\text{ is integral over }R\}
$$
is an $R$-subalgebra of $S$.
\end{lemma}

\begin{proof}
The image of each $r\in R$ satisfies $T-r$, hence belongs to $S'$;
this includes $0$ and $1$. If $u,v\in S'$,
Lemma \ref{editorial-source-lemma-characterize-integral} places both in the finite
$R$-subalgebra generated by these two original elements. It contains
$u+v$, $uv$ and $-u$, each of which is integral over $R$ by
Lemma \ref{algebra-lemma-finite-is-integral}. Therefore $S'$ is a subring
containing the image of $R$, with its original coefficient map.
\end{proof}


\subsection{Editorial treatment of Algebra lines 7556--7566}

\hypertarget{algebra-context-073}{}

\noindent\href{https://github.com/stacks/stacks-project/blob/a04446e57ec1fbc252a871afcec7752fb2807b14/algebra.tex\#L7556}{Original source, lines 7556--7566}. Complete original and reviewed TeX passages accompany this note. Every intervention is separately recorded; the following treatment is editorial.


\begin{lemma}
\label{editorial-source-lemma-finite-product-integral}
Let $R_i\to S_i$ be ring maps $i = 1, \ldots, n$.
Let $R$ and $S$ denote the product of the $R_i$ and $S_i$ respectively.
Then an element $s = (s_1, \ldots, s_n) \in S$ is integral over $R$
if and only if each $s_i$ is integral over $R_i$.
For a product indexed by an arbitrary set, the exact corresponding
condition is the existence of monic equations for all $s_i$ whose
degrees have one common finite bound.
\end{lemma}

\begin{proof}
Let $\varphi_i:R_i\to S_i$ be the original maps. If the tuple $s$
satisfies a monic equation of degree $D$ over $R=\prod_iR_i$,
projection to the $i$th coordinate gives an equation of degree $D$
for $s_i$ over $R_i$, with each coefficient projected separately.
Conversely choose monic equations
$P_i(T)=T^{d_i}+\sum_{j=0}^{d_i-1}a_{ij}T^j$ and a finite bound
$D\geq d_i$ for every $i$. Keep these polynomials and form
$Q_i(T)=T^{D-d_i}P_i(T)$; its coefficient of $T^j$ is
$a_{i,j-(D-d_i)}$ for $D-d_i\leq j<D$ and is zero for
$0\leq j<D-d_i$. Let $b_j\in\prod_iR_i$ have precisely these
coordinates. Then
$$
Q(T)=T^D+\sum_{j=0}^{D-1}b_jT^j
$$
is monic and its evaluation at the original tuple has $i$th coordinate
$s_i^{D-d_i}P_i^{\varphi_i}(s_i)=0$. This proves the assertion for
arbitrary products with that uniform bound. For a nonempty finite
product take $D=\max_i d_i$; for an empty product the rings are zero
and the degree-one polynomial $T$ works.

The bound is necessary for an infinite product. Fix a field $k$ and
use $R_i=k[v_i]$, $S_i=k[u_i]$ with $v_i\mapsto u_i^i$ for $i\geq1$.
The original element $s_i=u_i$ satisfies $T^i-v_i$.
No nonzero polynomial of degree less than $i$ in $T$ with
coefficients in $k[v_i]$ vanishes at $u_i$: after substitution,
the terms $a_j(u_i^i)u_i^j$ for $0\leq j<i$ have powers of $u_i$
in distinct residue classes modulo $i$, so they cannot cancel.
Thus the least monic degree is $i$. Each $s_i$ is integral but
no uniform finite bound exists. The preceding equivalence shows
that $(s_i)_i$ is not integral over $\prod_{i\geq1}R_i$.
\end{proof}


\subsection{Editorial treatment of Algebra lines 7582--7594}

\hypertarget{algebra-context-074}{}

\noindent\href{https://github.com/stacks/stacks-project/blob/a04446e57ec1fbc252a871afcec7752fb2807b14/algebra.tex\#L7582}{Original source, lines 7582--7594}. Complete original and reviewed TeX passages accompany this note. Every intervention is separately recorded; the following treatment is editorial.


\begin{lemma}
\label{editorial-source-lemma-finite-product-integral-closure}
Let $R_i\to S_i$ be ring maps $i = 1, \ldots, n$.
Denote the integral closure of $R_i$ in $S_i$ by $S'_i$.
Further let $R$ and $S$ denote the product of the $R_i$ and $S_i$ respectively.
Then the integral closure of $R$ in $S$
is the product of the $S'_i$. In particular it equals the image
of $R\to S$ if and only if each $S'_i$ equals the image of
$R_i\to S_i$. When all the given maps are injective, this says
that $R$ is integrally closed in $S$ if and only if every $R_i$
is integrally closed in $S_i$, using Definition \ref{algebra-definition-integral-closure}.
\end{lemma}

\begin{proof}
By Lemma \ref{editorial-source-lemma-finite-product-integral}, the integral tuples
are exactly $\prod_iS'_i$, with coordinatewise operations and the
original product coefficient map. The image of that map is
$\prod_i\varphi_i(R_i)$: each tuple of image elements has a tuple
of preimages in the finite product. Equality of these two products
is equivalent to equality in each coordinate, by projection and
by filling the other coordinates with zero. If the maps are
injective, their product is injective and the image equality is
exactly the subring property defined above. The empty product
case gives the zero ring and the assertions hold vacuously.
\end{proof}


\subsection{Editorial treatment of Algebra lines 7596--7627}

\hypertarget{algebra-context-075}{}

\noindent\href{https://github.com/stacks/stacks-project/blob/a04446e57ec1fbc252a871afcec7752fb2807b14/algebra.tex\#L7596}{Original source, lines 7596--7627}. Complete original and reviewed TeX passages accompany this note. Every intervention is separately recorded; the following treatment is editorial.


\begin{lemma}
\label{editorial-source-lemma-integral-closure-localize}
Integral closure commutes with localization: If $A \to B$ is a ring
map, and $S \subset A$ is a multiplicative subset, then the integral
closure of $S^{-1}A$ in $S^{-1}B$ is $S^{-1}B'$, where $B' \subset B$
is the integral closure of $A$ in $B$.
\end{lemma}

\begin{proof}
Since localization is exact we see that $S^{-1}B' \subset S^{-1}B$.
Suppose $x \in B'$ and $f \in S$. Then
$x^d + \sum_{i = 1, \ldots, d} a_i x^{d - i} = 0$
in $B$ for some $a_i \in A$. Hence also
$$
(x/f)^d + \sum\nolimits_{i = 1, \ldots, d} a_i/f^i (x/f)^{d - i} = 0
$$
in $S^{-1}B$. In this way we see that $S^{-1}B'$ is contained in
the integral closure of $S^{-1}A$ in $S^{-1}B$. Conversely, suppose
that $x/f \in S^{-1}B$ is integral over $S^{-1}A$. Choosing a monic equation of positive degree, we have
$$
(x/f)^d + \sum\nolimits_{i = 1, \ldots, d} (a_i/f_i) (x/f)^{d - i} = 0
$$
in $S^{-1}B$ for the original $a_i\in A$ and $f_i\in S$.
Put $F=\prod_{j=1}^d f_j\in S$ and define the actual element of $B$
$$
H=Fx^d+\sum_{i=1}^d a_i f^i
\left(\prod_{j\ne i}f_j\right)x^{d-i},
$$
where all coefficients from $A$ act through the given map $A\to B$.
Multiplying the original equality by $f^dF$ shows that $H/1=0$
in $S^{-1}B$. Therefore there is $f'\in S$ with $f'H=0$ in $B$.
No injectivity of $B\to S^{-1}B$ is assumed. Multiplying this
last equality by $(f')^{d-1}F^{d-1}$ gives exactly
$$
(f'Fx)^d+
\sum_{i=1}^d f^i(f')^i
\left(f_i^{i-1}\prod_{j\ne i}f_j^i\right)a_i
(f'Fx)^{d-i}=0.
$$
Indeed every term equals its original term in $H$ multiplied by
$(f')^dF^{d-1}$, with no denominator remaining. This is a monic
equation for the original scaled element $y=f'Fx$, so $y\in B'$.
The equality of fractions
$$
x/f=y/(ff'F)
$$
holds in $S^{-1}B$ by multiplying numerator and denominator by
$f'F\in S$. The right side is the image of an element of
$S^{-1}B'$, proving the reverse inclusion with all original
denominators and the annihilating factor $f'$ retained.
\end{proof}


\subsection{Editorial treatment of Algebra lines 7629--7662}

\hypertarget{algebra-context-076}{}

\noindent\href{https://github.com/stacks/stacks-project/blob/a04446e57ec1fbc252a871afcec7752fb2807b14/algebra.tex\#L7629}{Original source, lines 7629--7662}. Complete original and reviewed TeX passages accompany this note. Every intervention is separately recorded; the following treatment is editorial.


\begin{lemma}
\label{editorial-source-lemma-integral-closure-stalks}
\begin{slogan}
An element of an algebra over a ring is integral over the ring
if and only if it is locally integral at every prime ideal of the ring.
\end{slogan}
Let $\varphi : R \to S$ be a ring map.
Let $x \in S$. The following are equivalent:
\begin{enumerate}
\item $x$ is integral over $R$,
\item for every prime ideal $\mathfrak p \subset R$ the element
$x/1\in S_{\mathfrak p}$ is integral over $R_{\mathfrak p}$, and
\item for every maximal ideal $\mathfrak m\subset R$, the element
$x/1\in S_{\mathfrak m}$ is integral over $R_{\mathfrak m}$.
\end{enumerate}
\end{lemma}

\begin{proof}
The monic equation for (1) remains a monic equation after every
localization, so (1) implies (2), and (2) implies (3).
For the reverse implication let $B\subset S$ be the integral
closure of $R$ in $S$, with its original inclusion. It is an
$R$-submodule because it is an $R$-subalgebra. Lemma
\ref{editorial-source-lemma-integral-closure-localize} identifies its localization
$B_{\mathfrak m}\subset S_{\mathfrak m}$ with the integral closure
of $R_{\mathfrak m}$ there. Under (3), the class $\overline x$
in the quotient $R$-module $S/B$ therefore vanishes after localization
at every maximal ideal. This uses the exact quotient comparison
$(S/B)_{\mathfrak m}\cong S_{\mathfrak m}/B_{\mathfrak m}$,
which sends $(s+B)/u$ to $s/u+B_{\mathfrak m}$; the inverse is
induced by $s/u\mapsto(s+B)/u$, and the fraction equality criterion
and quotient relation prove both maps well defined and inverse.
If $\overline x\ne0$, its annihilator in $R$ is a proper ideal
and lies in a maximal ideal $\mathfrak m$. Vanishing of
$\overline x/1$ would give $u\notin\mathfrak m$ with
$u\overline x=0$, contradicting that containment. Thus
$\overline x=0$ and $x\in B$, proving (1). If $R$ is zero,
the unital map forces $S$ zero and the conclusion holds directly.

For completeness, retain also the original finite-subalgebra
calculation. Let $S'\subset S$ be generated by $\varphi(R)$ and
the original $x$. At a prime $\mathfrak p$ under (2), write its
monic equation of positive degree with coefficients
$a_i=b_i/u_i\in R_{\mathfrak p}$,
where $b_i\in R$ and $u_i\notin\mathfrak p$ for $1\leq i\leq d$.
Take $U=\prod_{i=1}^d u_i$ and the element
$$
H=\varphi(U)x^d+
\sum_{i=1}^d\varphi\left(b_i\prod_{j\ne i}u_j\right)x^{d-i}
\quad\text{of }S.
$$
Its image in $S_{\mathfrak p}$ is zero. Hence some
$v\notin\mathfrak p$ satisfies $\varphi(v)H=0$ in $S$.
Set $f=vU\notin\mathfrak p$. In $R_f$, each $u_i$ and $v$
is a unit because each is a factor of $f$. Dividing the last
equality by $vU$ in $S_f$ recovers the complete original equation
$$
x^d+\sum_{i=1}^d\varphi_f(b_i/u_i)x^{d-i}=0,
$$
where $\varphi_f:R_f\to S_f$ is the original localized map.
Thus $S'_f$, still the localization of that specified subalgebra,
is spanned over $R_f$ by $1,x,\ldots,x^{d-1}$.
The opens $D(f)$ so obtained cover $\operatorname{Spec}(R)$;
quasi-compactness (Lemma \ref{algebra-lemma-quasi-compact}) gives finitely
many $f_1,\ldots,f_n$ whose opens cover. They generate the unit
ideal, since no prime contains all of them. Lemma \ref{editorial-source-lemma-cover}
then shows that the original $S'$ is finite over $R$. Finally
Lemma \ref{editorial-source-lemma-characterize-integral} proves $x$ integral.
The factors $u_i$, $v$, $U$ and $f$ explicitly account both for
coefficient denominators and for the kernel of localization.
\end{proof}


\subsection{Editorial treatment of Algebra lines 7664--7686}

\hypertarget{algebra-context-077}{}

\noindent\href{https://github.com/stacks/stacks-project/blob/a04446e57ec1fbc252a871afcec7752fb2807b14/algebra.tex\#L7664}{Original source, lines 7664--7686}. Complete original and reviewed TeX passages accompany this note. Every intervention is separately recorded; the following treatment is editorial.


\begin{lemma}
\label{editorial-source-lemma-base-change-integral}
\begin{slogan}
Integrality and finiteness are preserved under base change.
\end{slogan}
Let $R \to S$ and $R \to R'$ be ring maps.
Set $S' = R' \otimes_R S$.
\begin{enumerate}
\item If $R \to S$ is integral so is $R' \to S'$.
\item If $R \to S$ is finite so is $R' \to S'$.
\end{enumerate}
\end{lemma}

\begin{proof}
For (1), write an arbitrary element of the original tensor algebra
as a finite sum $z=\sum_{i=1}^m r'_i\otimes s_i$.
Each $s_i$ satisfies its monic equation
$s_i^{d_i}+\sum_{j=0}^{d_i-1}\varphi(a_{ij})s_i^j=0$.
Tensoring gives the same equation for $1\otimes s_i$, with the
coefficients $a_{ij}$ mapped into $R'$. The finite-subalgebra
argument of Lemma \ref{editorial-source-lemma-characterize-integral} places all
$1\otimes s_i$, and hence their $R'$-linear combination $z$,
in a finite $R'$-subalgebra of $S'$. Lemma
\ref{algebra-lemma-finite-is-integral} proves $z$ integral.
For (2), choose original $R$-module generators $s_1,\ldots,s_m$
of $S$. An expression $s=\sum_i a_i s_i$ gives
$r'\otimes s=\sum_i r'\varphi'(a_i)(1\otimes s_i)$, where
$\varphi':R\to R'$ is the given map. Every tensor is a finite
sum of such terms, so $1\otimes s_1,\ldots,1\otimes s_m$
generate $S'$ as an $R'$-module.
\end{proof}


\subsection{Editorial treatment of Algebra lines 7688--7707}

\hypertarget{algebra-context-078}{}

\noindent\href{https://github.com/stacks/stacks-project/blob/a04446e57ec1fbc252a871afcec7752fb2807b14/algebra.tex\#L7688}{Original source, lines 7688--7707}. Complete original and reviewed TeX passages accompany this note. Every intervention is separately recorded; the following treatment is editorial.


\begin{lemma}
\label{editorial-source-lemma-integral-local}
Let $R \to S$ be a ring map.
Let $f_1, \ldots, f_n \in R$ generate the unit ideal.
\begin{enumerate}
\item If each $R_{f_i} \to S_{f_i}$ is integral, so is $R \to S$.
\item If each $R_{f_i} \to S_{f_i}$ is finite, so is $R \to S$.
\end{enumerate}
\end{lemma}

\begin{proof}
For (1), fix the original $s\in S$ and keep
$I=\{P\in R[x]:P^\varphi(s)=0\}$, the kernel of evaluation.
Let $J\subset R$ consist of zero and the leading coefficients
of the nonzero polynomials in $I$. This is an ideal. To add two
nonzero leading coefficients, multiply the corresponding polynomials
by powers of $x$ to make their degrees equal. If the sum of the
leading coefficients is nonzero, it is the leading coefficient of
the sum polynomial; if it is zero, it belongs to $J$ by definition.
For multiplication by $r\in R$, a nonzero product of leading
coefficients remains the leading coefficient of $rP$; a zero
product again belongs to $J$. Thus all ideal axioms hold.

For each original $f_i$, choose the local monic equation
$$
s^{d_i}+\sum_{j=0}^{d_i-1}
\varphi_{f_i}(a_{ij}/f_i^{c_{ij}})s^j=0
\quad\text{in }S_{f_i},
$$
where $d_i\geq1$, $a_{ij}\in R$ and $c_{ij}\geq0$.
Choose $N_i\geq c_{ij}$ for all $j$. The element
$\varphi(f_i^{N_i})s^{d_i}+
\sum_j\varphi(a_{ij}f_i^{N_i-c_{ij}})s^j$ becomes zero in
$S_{f_i}$, so some $M_i\geq0$ makes it zero upon multiplication
by $\varphi(f_i^{M_i})$. Consequently the original polynomial
$$
f_i^{N_i+M_i}x^{d_i}
+\sum_{j=0}^{d_i-1}a_{ij}f_i^{N_i+M_i-c_{ij}}x^j
$$
belongs to $I$. If $f_i^{N_i+M_i}\ne0$ it is its leading
coefficient; if it is zero it belongs to $J$ by definition.
In either case $f_i^{N_i+M_i}\in J$. A proper $J$ would lie
in a maximal ideal containing every $f_i$ (or would already
contain $1$ when an exponent is zero), contrary to the given
unit-ideal hypothesis. Therefore $1\in J$, so $I$ contains a
monic polynomial and $s$ is integral over $R$.

The maximal-ideal criterion above gives the same conclusion:
every maximal ideal omits some $f_i$ and the corresponding local
map is a further localization of $R_{f_i}\to S_{f_i}$.

For (2), choose finite $R_{f_i}$-module generators of $S_{f_i}$.
Write each as $s_{ij}/f_i^{a_{ij}}$, keeping its numerator
$s_{ij}\in S$ and exponent $a_{ij}\geq0$.
Let $N\subset S$ be the $R$-submodule generated by all these
finitely many numerators. Since $f_i$ is a unit after localization,
$N_{f_i}=S_{f_i}$. Thus the quotient $Q=S/N$ satisfies $Q_{f_i}=0$
for each $i$. For $q\in Q$ choose $b_i\geq1$ with $f_i^{b_i}q=0$.
The ideal generated by the $f_i^{b_i}$ is the unit ideal: if it were
proper, a maximal ideal containing it would contain all $f_i$,
contrary to the original unit-ideal assumption. Choose $c_i\in R$
with $\sum_i c_if_i^{b_i}=1$. Then
$q=\sum_i c_if_i^{b_i}q=0$. Hence $Q=0$ and the original numerators
$s_{ij}$ generate $S$ over $R$. The empty-cover case forces $R=0$
and hence $S=0$, where the assertion also holds.
\end{proof}


\subsection{Editorial treatment of Algebra lines 7709--7720}

\hypertarget{algebra-context-079}{}

\noindent\href{https://github.com/stacks/stacks-project/blob/a04446e57ec1fbc252a871afcec7752fb2807b14/algebra.tex\#L7709}{Original source, lines 7709--7720}. Complete original and reviewed TeX passages accompany this note. Every intervention is separately recorded; the following treatment is editorial.


\begin{lemma}
\label{editorial-source-lemma-integral-permanence}
Let $A \to B \to C$ be ring maps.
\begin{enumerate}
\item If $A \to C$ is integral so is $B \to C$.
\item If $A \to C$ is finite so is $B \to C$.
\end{enumerate}
\end{lemma}

\begin{proof}
Write $\alpha:A\to B$ and $\beta:B\to C$ for the original maps.
If $c\in C$ satisfies
$c^d+\sum_{i=0}^{d-1}(\beta\alpha)(a_i)c^i=0$, the same monic
polynomial with coefficients $\alpha(a_i)\in B$ proves integrality
over $B$. This proves (1) for every original $c$.
If $c_1,\ldots,c_m$ generate $C$ as an $A$-module, every expression
$c=\sum_i(\beta\alpha)(a_i)c_i$ is also a $B$-linear expression,
with coefficients $\alpha(a_i)$. The same family generates $C$
as a $B$-module, proving (2).
\end{proof}


\subsection{Editorial treatment of Algebra lines 7722--7732}

\hypertarget{algebra-context-080}{}

\noindent\href{https://github.com/stacks/stacks-project/blob/a04446e57ec1fbc252a871afcec7752fb2807b14/algebra.tex\#L7722}{Original source, lines 7722--7732}. Complete original and reviewed TeX passages accompany this note. Every intervention is separately recorded; the following treatment is editorial.


\begin{lemma}
\label{editorial-source-lemma-integral-closure-transitive}
Let $A \to B \to C$ be ring maps.
Let $B'$ be the integral closure of $A$ in $B$,
let $C'$ be the integral closure of $B'$ in $C$. Then
$C'$ is the integral closure of $A$ in $C$.
\end{lemma}

\begin{proof}
The coefficient map $A\to B$ lands in $B'$ by the degree-one
equations for its image elements, and $A\to B'$ is integral by
definition. If $c\in C'$ is integral over $B'$, transitivity of
integrality (Lemma \ref{algebra-lemma-integral-transitive}) shows that $c$
is integral over $A$. Conversely, every monic equation for $c$
over $A$ becomes a monic equation over $B'$ under the original
map $A\to B'$, since its composite into $C$ is unchanged.
Thus the two sets of integral elements of the same ring $C$ agree,
and their subalgebra structures and coefficient maps agree as well.
\end{proof}


\subsection{Editorial treatment of Algebra lines 7734--7762}

\hypertarget{algebra-context-081}{}

\noindent\href{https://github.com/stacks/stacks-project/blob/a04446e57ec1fbc252a871afcec7752fb2807b14/algebra.tex\#L7734}{Original source, lines 7734--7762}. Complete original and reviewed TeX passages accompany this note. Every intervention is separately recorded; the following treatment is editorial.


\begin{lemma}
\label{editorial-source-lemma-integral-overring-surjective}
Suppose that $R \to S$ is an integral
ring extension with $R \subset S$.
Then $\varphi : \Spec(S) \to \Spec(R)$
is surjective. More generally, for an arbitrary integral ring map
$\varphi:R\to S$ its spectrum image is $V(\ker\varphi)$.
\end{lemma}

\begin{proof}
Let $\mathfrak p \subset R$ be a prime ideal.
We have to show $\mathfrak pS_{\mathfrak p} \not = S_{\mathfrak p}$, see
Lemma \ref{editorial-source-lemma-in-image}.
The localization $R_{\mathfrak p} \to S_{\mathfrak p}$ is injective
(as localization is exact) and integral by
Lemma \ref{editorial-source-lemma-integral-closure-localize} or
\ref{editorial-source-lemma-base-change-integral}.
Keep these localizations and put $\mathfrak m=\mathfrak pR_{\mathfrak p}$.
Suppose, for contradiction, that
$1=\sum_i f_is_i$ in $S_{\mathfrak p}$ with
$f_i\in\mathfrak m$ and $s_i\in S_{\mathfrak p}$, using the
original localized coefficient map. Lemma
\ref{editorial-source-lemma-characterize-integral} gives an
$R_{\mathfrak p}$-subalgebra $S'\subset S_{\mathfrak p}$
containing all these $s_i$ and finite over $R_{\mathfrak p}$.
The same equation for $1$ holds in $S'$. Multiplying it by every
element of $S'$ proves $S'=\mathfrak mS'$. Nakayama's
Lemma \ref{algebra-lemma-NAK} gives $S'=0$, which contradicts the injection
of the nonzero local ring $R_{\mathfrak p}$ into $S'$.
Thus the original fibre has a prime and the desired preimage exists.

For arbitrary $\varphi$, the induced map
$R/\ker\varphi\to S$, $r+\ker\varphi\mapsto\varphi(r)$,
is injective and integral: reduce each coefficient of a monic
equation modulo the kernel. The injective case just proved makes
its spectrum map surjective. Composing with the original quotient
map identifies its image in $\operatorname{Spec}(R)$ with
$V(\ker\varphi)$, because primes of the quotient are exactly
the primes of $R$ containing that kernel.
\end{proof}


\subsection{Editorial treatment of Algebra lines 7764--7779}

\hypertarget{algebra-context-082}{}

\noindent\href{https://github.com/stacks/stacks-project/blob/a04446e57ec1fbc252a871afcec7752fb2807b14/algebra.tex\#L7764}{Original source, lines 7764--7779}. Complete original and reviewed TeX passages accompany this note. Every intervention is separately recorded; the following treatment is editorial.


\begin{lemma}
\label{editorial-source-lemma-integral-under-field}
Let $R$ be a ring. Let $K$ be a field.
If $R \subset K$ and $K$ is integral over $R$,
then $R$ is a field and $K$ is an algebraic extension.
If $R \subset K$ and $K$ is finite over $R$,
then $R$ is a field and $K$ is a finite algebraic extension.
\end{lemma}

\begin{proof}
Assume that $R \subset K$ is integral.
By Lemma \ref{editorial-source-lemma-integral-overring-surjective} we see that
$\Spec(R)$ has $1$ point. Since clearly $R$ is a domain we see
that $R = R_{(0)}$ is a field (Lemma \ref{algebra-lemma-minimal-prime-reduced-ring}).
For an explicit inverse calculation, if $0\ne r\in R$, the
element $r^{-1}\in K$ satisfies a monic equation
$(r^{-1})^d+\sum_{i=1}^d a_i(r^{-1})^{d-i}=0$ with $a_i\in R$
and $d\geq1$. Multiplication by $r^{d-1}$ in the original field
gives the full formula
$$
r^{-1}=-\sum_{i=1}^d a_i r^{i-1}\in R.
$$
Every element of $K$ is then algebraic over the field $R$, since
its monic integral equation is a nonzero polynomial over $R$.
In the finite case, Lemma \ref{algebra-lemma-finite-is-integral} first
gives integrality, and the original finite module generators
become a finite vector-space spanning family over this same field
$R$. Hence the algebraic extension has finite degree.
\end{proof}


\subsection{Editorial treatment of Algebra lines 7781--7810}

\hypertarget{algebra-context-083}{}

\noindent\href{https://github.com/stacks/stacks-project/blob/a04446e57ec1fbc252a871afcec7752fb2807b14/algebra.tex\#L7781}{Original source, lines 7781--7810}. Complete original and reviewed TeX passages accompany this note. Every intervention is separately recorded; the following treatment is editorial.


\begin{lemma}
\label{editorial-source-lemma-integral-over-field}
Let $k$ be a field. Let $S$ be a $k$-algebra over $k$.
\begin{enumerate}
\item If $S$ is a domain and finite dimensional over $k$,
then $S$ is a field.
\item If $S$ is integral over $k$ and a domain,
then $S$ is a field.
\item If $S$ is integral over $k$ then every prime of
$S$ is a maximal ideal (see
Lemma \ref{algebra-lemma-ring-with-only-minimal-primes}
for more consequences).
\end{enumerate}
\end{lemma}

\begin{proof}
For (3), a quotient $S/\mathfrak q$ by any prime is an integral
$k$-algebra and a domain. Its coefficient map from $k$ is injective,
since the quotient is nonzero and $k$ is a field. Assertion (2)
makes this quotient a field, hence $\mathfrak q$ is maximal.
Let $S$ be integral over $k$ and assume $S$ is a domain.
Take a nonzero $s\in S$. By Lemma
\ref{editorial-source-lemma-characterize-integral} we may find a
finite dimensional $k$-subalgebra $k \subset S' \subset S$
containing $s$. This subalgebra is a domain since it is a subring
of $S$. Once (1) is proved, $S'$ is a field and the inverse of
the original nonzero $s$ belongs to $S'\subset S$, proving (2).
For (1), assume that the original $S$ is finite dimensional over
$k$ and is a domain. Pick a nonzero $s\in S$. The $k$-linear
map $S\to S$, $a\mapsto sa$, is injective because $S$ is a
domain and $s\ne0$. Equality of the finite source and target
dimensions makes it surjective. A preimage $s_1$ of $1$ satisfies
$ss_1=1$. Thus every nonzero element is invertible and $S$ is a field.
\end{proof}


\subsection{Editorial treatment of Algebra lines 7812--7829}

\hypertarget{algebra-context-084}{}

\noindent\href{https://github.com/stacks/stacks-project/blob/a04446e57ec1fbc252a871afcec7752fb2807b14/algebra.tex\#L7812}{Original source, lines 7812--7829}. Complete original and reviewed TeX passages accompany this note. Every intervention is separately recorded; the following treatment is editorial.


\begin{lemma}
\label{editorial-source-lemma-integral-no-inclusion}
Suppose $R \to S$ is integral.
Let $\mathfrak q, \mathfrak q' \in \Spec(S)$
be distinct primes
having the same image in $\Spec(R)$.
Then neither $\mathfrak q \subset \mathfrak q'$
nor $\mathfrak q' \subset \mathfrak q$.
\end{lemma}

\begin{proof}
Let $\mathfrak p=\varphi^{-1}(\mathfrak q)=
\varphi^{-1}(\mathfrak q')$ and set $U=R\setminus\mathfrak p$.
Retain the fibre algebra
$A=S\otimes_R\kappa(\mathfrak p)\cong U^{-1}S/\mathfrak pU^{-1}S$.
The two ideals are sent to
$U^{-1}\mathfrak q/\mathfrak pU^{-1}S$ and
$U^{-1}\mathfrak q'/\mathfrak pU^{-1}S$. Localization and quotient
preserve inclusion, and contraction recovers each original prime:
$s/1\in U^{-1}\mathfrak q$ means
$\varphi(u)s\in\mathfrak q$ for some $u\in U$, which by primality
and $\varphi(u)\notin\mathfrak q$ is equivalent to $s\in\mathfrak q$.
Thus the two resulting primes are distinct. Lemma
\ref{editorial-source-lemma-base-change-integral} makes $A$ integral over
$\kappa(\mathfrak p)$, so Lemma \ref{editorial-source-lemma-integral-over-field}
makes both primes maximal. An inclusion between them would force
equality, and contraction would contradict the original distinction.
\end{proof}


\subsection{Editorial treatment of Algebra lines 7831--7850}

\hypertarget{algebra-context-085}{}

\noindent\href{https://github.com/stacks/stacks-project/blob/a04446e57ec1fbc252a871afcec7752fb2807b14/algebra.tex\#L7831}{Original source, lines 7831--7850}. Complete original and reviewed TeX passages accompany this note. Every intervention is separately recorded; the following treatment is editorial.


\begin{lemma}
\label{editorial-source-lemma-finite-finite-fibres}
Suppose $R \to S$ is finite.
Then the fibres of $\Spec(S)\to\Spec(R)$ are finite.
More precisely, if $s_1,\ldots,s_m$ generate $S$ as an $R$-module,
then for every $\mathfrak p\in\operatorname{Spec}(R)$,
$$
\begin{aligned}
\#\{\mathfrak q:\varphi^{-1}(\mathfrak q)=\mathfrak p\}
&\leq\sum_{\varphi^{-1}(\mathfrak q)=\mathfrak p}
[\kappa(\mathfrak q):\kappa(\mathfrak p)]
\\
&\leq\dim_{\kappa(\mathfrak p)}(S\otimes_R\kappa(\mathfrak p))
\leq m.
\end{aligned}
$$
The middle difference is exactly the dimension of the nilradical
of this fibre algebra. In particular the residue-degree sum equals
the fibre dimension exactly when this fibre is reduced.
\end{lemma}

\begin{proof}
By the discussion in
Remark \ref{algebra-remark-fundamental-diagram}
the fibres are the spectra of the rings $S \otimes_R \kappa(\mathfrak p)$.
As $R \to S$ is finite, these fibre rings are finite over
$\kappa(\mathfrak p)$ hence Noetherian by
Lemma \ref{editorial-source-lemma-Noetherian-permanence}.
By
Lemma \ref{editorial-source-lemma-integral-no-inclusion}
every prime of $S \otimes_R \kappa(\mathfrak p)$ is a minimal
prime. Hence by
Lemma \ref{editorial-source-lemma-Noetherian-irreducible-components}
there are at most finitely many.
\medskip\noindent
For the quantitative assertions, keep the original map
$\varphi:R\to S$ and the given module generators $s_1,\ldots,s_m$.
Fix $\mathfrak p\in\operatorname{Spec}(R)$ and put
$K=\kappa(\mathfrak p)$, $U=R\setminus\mathfrak p$,
$A=S\otimes_R K$ and $d=\dim_K A$.
The original surjection $R^m\to S$ induces the surjection
$$
K^m\longrightarrow A,\qquad
(\lambda_j)_j\longmapsto\sum_{j=1}^m s_j\otimes\lambda_j.
$$
Indeed $s=\sum_j\varphi(a_j)s_j$ gives
$s\otimes\lambda=\sum_j s_j\otimes(\overline{a_j}\lambda)$,
and every tensor is a finite sum of pure tensors. Thus $d\leq m$.

The exact fibre-ring isomorphism and its inverse are
$$
\begin{aligned}
A&\longrightarrow U^{-1}S/\mathfrak pU^{-1}S,&
s\otimes\frac{r+\mathfrak p}{u+\mathfrak p}
&\longmapsto\frac{s\varphi(r)}{\varphi(u)}+\mathfrak pU^{-1}S,\\
\frac{s}{\varphi(u)}+\mathfrak pU^{-1}S
&\longmapsto s\otimes\frac1{u+\mathfrak p}.&&
\end{aligned}
$$
Here $U^{-1}S$ means localization at the image $\varphi(U)$.
Tensor balancing respects the original coefficient map; a change
of fraction representatives is killed by a denominator from $U$
on both sides, and an element of $\mathfrak p$ acts as zero on
$K$. These facts verify the defining relations of both maps.
Their composites on the displayed fractions and pure tensors
are identities, proving the isomorphism.
Write $\beta:S\to A$, $s\mapsto s\otimes1$.
A prime $\mathfrak q$ of $S$ with
$\varphi^{-1}(\mathfrak q)=\mathfrak p$ corresponds to
$$
\mathfrak n_{\mathfrak q}
=U^{-1}\mathfrak q/\mathfrak pU^{-1}S\subset A.
$$
It is proper because $\mathfrak q$ avoids $\varphi(U)$, and
contraction recovers $\mathfrak q$: the condition
$s/1\in U^{-1}\mathfrak q$ is equivalent to
$\varphi(u)s\in\mathfrak q$ for some $u\in U$, hence to
$s\in\mathfrak q$ by primality. Conversely, a prime
$\mathfrak n\subset A$ contracts to a prime $\mathfrak q\subset S$
whose contraction in $R$ is $\mathfrak p$: elements of
$\mathfrak p$ act as zero and elements of $U$ act as units.
Every element of $A$ is represented by a fraction $s/\varphi(u)$
modulo $\mathfrak p$, and it belongs to $\mathfrak n$ precisely
when $\beta(s)$ does. This verifies the remaining composite of
the two prime correspondences. In particular the correspondence
is a bijection for the original, possibly noninjective map.

For every prime $\mathfrak n$ of $A$, the quotient $A/\mathfrak n$
is a finite dimensional domain over $K$, hence a field by
Lemma \ref{editorial-source-lemma-integral-over-field}. Therefore all these primes
are maximal. If $\mathfrak n=\mathfrak n_{\mathfrak q}$,
its field is canonically the original residue field:
$$
A/\mathfrak n_{\mathfrak q}
\cong U^{-1}(S/\mathfrak q)
\xrightarrow{\ \sim\ }\kappa(\mathfrak q),\qquad
\frac{s+\mathfrak q}{\varphi(u)+\mathfrak q}
\longmapsto\frac{s+\mathfrak q}{\varphi(u)+\mathfrak q}.
$$
The last map is injective since $S/\mathfrak q$ is a domain.
Its source is already a field containing $S/\mathfrak q$,
so it contains the inverse of each nonzero element of that domain;
it therefore contains every fraction and the map is surjective.
This also identifies the coefficient field $K$ in both targets.

Here is the Chinese remainder map with its actual preimages.
For any finite family of distinct maximal ideals
$\mathfrak n_1,\ldots,\mathfrak n_t$ of $A$, comaximality gives
$b_{ij}\in\mathfrak n_j$ with $1-b_{ij}\in\mathfrak n_i$
when $i\ne j$. Put $e_i=\prod_{j\ne i}b_{ij}$.
It is $1$ modulo $\mathfrak n_i$ and $0$ modulo each other
$\mathfrak n_j$. Hence
$$
\rho:A\longrightarrow\prod_{i=1}^t A/\mathfrak n_i,
\qquad a\longmapsto(a+\mathfrak n_i)_i
$$
is surjective: representatives $a_i$ of a prescribed tuple have
preimage $\sum_i e_ia_i$. Its kernel is exactly
$\bigcap_i\mathfrak n_i$. Comparing $K$-dimensions gives
$\sum_i\dim_K(A/\mathfrak n_i)\leq d$. Every field factor has
positive dimension, so there cannot be more than $d$ maximal
ideals: any $d+1$ distinct ones would contradict the inequality.
If $A\ne0$ it has a maximal ideal, for example a proper ideal of
greatest $K$-dimension. Thus its primes are a finite family and
the preceding construction applies to all of them at once.
Let $L_i=A/\mathfrak n_i\cong\kappa(\mathfrak q_i)$ and
$f_i=[L_i:K]$. With $J=\bigcap_i\mathfrak n_i$, it gives
$$
0\longrightarrow J\longrightarrow A\xrightarrow{\rho}
\prod_iL_i\longrightarrow0,
\qquad
d=\sum_i f_i+\dim_KJ.
$$

The ideal $J$ is precisely the nilradical, and is nilpotent.
To prove the latter assertion directly, the decreasing chain
$J,J^2,J^3,\ldots$ stabilizes since these are subspaces of the
finite dimensional $K$-space $A$. Choose $N\geq1$ with
$J^N=J^{N+1}$. A finite $K$-basis of $J^N$ generates it as an
$A$-module. Write each generator as a $J$-linear combination
of all generators, giving a matrix $C$ with entries in $J$.
Multiplication by the adjugate of $I-C$ gives
$\det(I-C)J^N=0$. This determinant belongs to $1+J$ and is a
unit: if it lay in a maximal ideal, that ideal would also contain
$J$, and hence would contain $1$. Therefore $J^N=0$.
Every nilpotent element lies in all prime ideals, and conversely
every element of $J$ is now nilpotent, proving the assertion.
Consequently the middle inequality is an equality exactly when
$A$ is reduced. Together with $f_i\geq1$ and $d\leq m$, this
proves all the displayed inequalities. For $A=0$ there are no
prime ideals, its dimension is zero, and the formulas use empty sums.

There is a more precise local formula, retaining every multiplicity.
For $A\ne0$ and the above $N$, the ideals $\mathfrak n_i^N$ are
pairwise comaximal. In fact, if $a+b=1$ with
$a\in\mathfrak n_i$, $b\in\mathfrak n_j$, every term of
$(a+b)^{2N-1}=1$ belongs to $\mathfrak n_i^N$ or $\mathfrak n_j^N$.
For two comaximal ideals their intersection is their product:
if $x$ is in both and $a+b=1$ with one element in each ideal,
then $x=xa+xb$ belongs to the product. Induction gives the finite
version; the product of the first ideals stays comaximal with
the next since multiplying equations congruent to $1$ modulo
the next ideal gives another such equation. It follows that
$$
\bigcap_i\mathfrak n_i^N
=\prod_i\mathfrak n_i^N
=\left(\prod_i\mathfrak n_i\right)^N
=J^N=0.
$$
The same Chinese remainder construction therefore gives
$$
A\xrightarrow{\ \sim\ }\prod_i C_i,
\qquad C_i=A/\mathfrak n_i^N.
$$
Each $C_i$ has unique maximal ideal
$\mathfrak m_i=\mathfrak n_i/\mathfrak n_i^N$: a prime
containing $\mathfrak n_i^N$ contains every element of
$\mathfrak n_i$, by taking its $N$th power, and thus equals
that maximal ideal. The residue field is $L_i$.

The identification $C_i\cong A_{\mathfrak n_i}$ has explicit maps.
Choose the Chinese remainder element $e_i$ that is $1$ modulo
$\mathfrak n_i^N$ and $0$ modulo the other powers. For
$a\in\mathfrak n_i^N$, the product $e_ia$ is in the intersection
of all powers and is zero. Since $e_i\notin\mathfrak n_i$,
this gives $a/1=0$ in $A_{\mathfrak n_i}$. Thus
$\overline a\mapsto a/1$ defines $C_i\to A_{\mathfrak n_i}$.
Its inverse sends $a/u$ to $\overline a\,\overline u^{-1}$;
the inverse exists since $u\notin\mathfrak n_i$ and $C_i$ is
local. The two composites on elements and fractions are identities.
There is also the exact original-ring comparison
$$
A_{\mathfrak n_i}\cong S_{\mathfrak q_i}/\mathfrak pS_{\mathfrak q_i}.
$$
One direction sends $s/v$ modulo $\mathfrak p$ to
$\beta(s)/\beta(v)$ for $v\notin\mathfrak q_i$.
The other is induced by
$s\otimes((r+\mathfrak p)/(u+\mathfrak p))
\mapsto s\varphi(r)/\varphi(u)$ modulo $\mathfrak p$.
It extends to localization at $\mathfrak n_i$, since an element
outside that maximal ideal maps to an element with nonzero
residue in the identified field $\kappa(\mathfrak q_i)$ and
hence to a unit of the target local ring. Tensor, fraction and
quotient relations are respected as in the preceding fibre map,
and the composites fix every such fraction, proving the comparison.

Use the finite filtration
$C_i\supset\mathfrak m_i\supset\mathfrak m_i^2\supset\cdots\supset0$.
Each quotient $\mathfrak m_i^j/\mathfrak m_i^{j+1}$ is annihilated
by $\mathfrak m_i$ and is a finite dimensional $L_i$-space.
Put
$$
\mu_i=\sum_{j\geq0}
\dim_{L_i}(\mathfrak m_i^j/\mathfrak m_i^{j+1}).
$$
This positive integer is the length of $C_i$ over itself:
refine each quotient by a vector-space flag, and note that every
simple module over the local ring $C_i$ is $L_i$, because the
annihilator of a nonzero simple-module generator is its unique
maximal ideal. Every composition series has the same number of
factors, since each such factor has $K$-dimension $f_i$.
Taking dimensions of the filtration and of the product yields
the exact formulas
$$
\begin{aligned}
\dim_K(S\otimes_RK)&=\sum_i\mu_i[\kappa(\mathfrak q_i):K],\\
\dim_KJ&=\sum_i(\mu_i-1)[\kappa(\mathfrak q_i):K],\\
\mu_i&=\operatorname{length}_{S_{\mathfrak q_i}/\mathfrak pS_{\mathfrak q_i}}
 (S_{\mathfrak q_i}/\mathfrak pS_{\mathfrak q_i}).
\end{aligned}
$$
Thus no local multiplicity has been absorbed into a residue degree.

The empty fibre also has an exact criterion:
$A=0$ if and only if $\ker\varphi\not\subseteq\mathfrak p$.
If some $u\in U$ is in the kernel, its zero image is inverted
in $U^{-1}S$, giving the zero ring. Conversely $A=0$ implies
$U^{-1}S=\mathfrak pR_{\mathfrak p}(U^{-1}S)$.
This is a finite $R_{\mathfrak p}$-module generated by the images
of the specified $s_j$. Nakayama's Lemma \ref{algebra-lemma-NAK} gives
$U^{-1}S=0$, so $1/1=0$ implies $\varphi(u)=0$ for some
$u\in U$. This proves the criterion without assuming injectivity.

For clarity, all equality cases follow from these formulas.
The number of primes equals the residue-degree sum precisely
when every $f_i=1$. That sum equals $d$ precisely when $J=0$,
equivalently all $\mu_i=1$. The upper equality $d=m$ holds
precisely when the original images $s_j\otimes1$ are a $K$-basis.
In particular the number of primes equals $d$ exactly when
$A\cong K^t$, where $t$ is the number of primes: the equalities
give $J=0$, $f_i=1$ and the proved Chinese remainder isomorphism;
the converse follows from that product's coordinate maximal ideals.
The number of primes equals $m$ exactly when $A\cong K^m$.
These statements include the empty product, interpreted as the zero ring.
In terms of the original relation module
$H=\ker(R^m\to S)$, the kernel of $K^m\to A$ is exactly
$$
\frac{H_{\mathfrak p}+\mathfrak pR_{\mathfrak p}^m}
{\mathfrak pR_{\mathfrak p}^m},
$$
because localization gives $S_{\mathfrak p}=R_{\mathfrak p}^m/H_{\mathfrak p}$
and reduction gives that quotient. Thus $d=m$ exactly when
$H_{\mathfrak p}\subseteq\mathfrak pR_{\mathfrak p}^m$, equivalently
$H\subseteq\mathfrak pR^m$. The last equivalence is coordinatewise:
$a/1\in\mathfrak pR_{\mathfrak p}$ if and only if $a\in\mathfrak p$,
by the fraction equality criterion and primality.

The bounds and their limitations can be seen in explicit examples.
For $R=k$, $S=k^m$ and the coordinate idempotent generators,
the fibre has $m$ primes and all inequalities are equalities.
For $R=k[t]$, $S=k[x]$ and $\varphi(t)=x^m$ with $m\geq1$,
the original generators $1,x,\ldots,x^{m-1}$ form a free
$R$-basis: collecting powers by their residue modulo $m$ gives
spanning and uniqueness of the coefficients in $k[t]$.
At $\mathfrak p=(t)$ the fibre is $k[x]/(x^m)$, so it has
one prime, residue degree $1$, multiplicity $m$ and nilradical
dimension $m-1$, although both original rings are domains.
For the noninjective map $k[t]\to k$, $t\mapsto0$, with generator
$1$, the fibre at $(t-1)$ is zero, while the fibre at $(t)$ is $k$.
Repeating the generator $1$ for $R=S=k$ gives $d=1<m$ whenever
the supplied family has $m>1$ entries.

Equality in the residue-degree bound does not imply separability.
In characteristic $\ell>0$, let $K=\mathbb F_\ell(t)$ and
$L=\mathbb F_\ell(u)$ with $t\mapsto u^\ell$.
The elements $1,u,\ldots,u^{\ell-1}$ are linearly independent
over $K$: clearing denominators in any relation gives
$\sum_{j=0}^{\ell-1}u^jP_j(u^\ell)=0$, whose powers in distinct
residue classes modulo $\ell$ cannot cancel. Their span is closed
under multiplication since $u^\ell=t$, and is a finite dimensional
domain over $K$, hence a field. It contains $u$ and therefore is
the whole $L$, proving $[L:K]=\ell$ and
$L\cong K[X]/(X^\ell-t)$. Its unique residue field has degree
equal to its dimension and it is reduced. But
$$
L\otimes_KL\cong L[X]/(X^\ell-u^\ell)
=L[X]/((X-u)^\ell)
$$
is not reduced: $X-u$ is a nonzero element of nilpotence order
$\ell$, as its degree is smaller than that of the defining monic
polynomial. Thus the equality asserts reducedness of the specified
fibre, not reducedness after every field extension.

Finally every finite choice of positive integers $f_i,\mu_i$
is realized, showing the numerical formula is sharp. Take
$K=F(t_0)$ for a field $F$ and
$L_i=F(u_i)$ with $t_0\mapsto u_i^{f_i}$.
The same distinct-residue-class argument proves
$[L_i:K]=f_i$, with basis $1,u_i,\ldots,u_i^{f_i-1}$;
the spanning ring is a finite dimensional domain and hence a field
containing $u_i$, so the argument applies to the full rational
function field. Set
$$
S=\prod_i L_i[\varepsilon_i]/(\varepsilon_i^{\mu_i}).
$$
Its $K$-basis consists of the elements supported in factor $i$
with entries $u_i^a\varepsilon_i^b$ for
$0\leq a<f_i$, $0\leq b<\mu_i$. The nilpotent ideal in each
factor is $(\varepsilon_i)$, its quotient is the field $L_i$,
and its successive powers have one dimensional quotients over
$L_i$ until exponent $\mu_i$. Thus its unique prime has residue
degree $f_i$ and local length $\mu_i$. A prime of the product
selects exactly one factor because the coordinate idempotents
are orthogonal and sum to $1$. Consequently the fibre over the
zero prime of $K$ has exactly these degrees and multiplicities.
The displayed basis has size $\sum_i\mu_if_i$; appending zeros
gives a specified generating family of any larger size $m$.
Thus positivity and $\sum_i\mu_if_i\leq m$ are the complete
universal numerical restrictions on these data.
\end{proof}


\subsection{Editorial treatment of Algebra lines 7852--7868}

\hypertarget{algebra-context-086}{}

\noindent\href{https://github.com/stacks/stacks-project/blob/a04446e57ec1fbc252a871afcec7752fb2807b14/algebra.tex\#L7852}{Original source, lines 7852--7868}. Complete original and reviewed TeX passages accompany this note. Every intervention is separately recorded; the following treatment is editorial.


\begin{lemma}
\label{editorial-source-lemma-integral-going-up}
Let $R \to S$ be a ring map such that
$S$ is integral over $R$.
Let $\mathfrak p \subset \mathfrak p' \subset R$
be primes. Let $\mathfrak q$ be a prime of $S$ mapping
to $\mathfrak p$. Then there exists a prime $\mathfrak q'$
with $\mathfrak q \subset \mathfrak q'$
mapping to $\mathfrak p'$.
\end{lemma}

\begin{proof}
Keep the original map $\varphi:R\to S$ and form
$$
\overline\varphi:R/\mathfrak p\longrightarrow S/\mathfrak q,
\qquad r+\mathfrak p\longmapsto\varphi(r)+\mathfrak q.
$$
It is well defined and injective because
$\varphi^{-1}(\mathfrak q)=\mathfrak p$; both rings are domains.
It is integral because the monic equation for each original
$s\in S$ descends coefficient by coefficient to its class.
Lemma \ref{editorial-source-lemma-integral-overring-surjective} therefore gives
a prime $\overline{\mathfrak q'}$ of $S/\mathfrak q$ contracting
to $\mathfrak p'/\mathfrak p$. Let $\mathfrak q'$ be its inverse
image in $S$. This is prime, contains the original $\mathfrak q$,
and its contraction is exactly $\mathfrak p'$: an element $r\in R$
maps into it if and only if $r+\mathfrak p\in\mathfrak p'/\mathfrak p$,
which is equivalent to $r\in\mathfrak p'$.
\end{proof}


\subsection{Editorial treatment of Algebra lines 7875--7925}

\hypertarget{algebra-context-087}{}

\noindent\href{https://github.com/stacks/stacks-project/blob/a04446e57ec1fbc252a871afcec7752fb2807b14/algebra.tex\#L7875}{Original source, lines 7875--7925}. Complete original and reviewed TeX passages accompany this note. Every intervention is separately recorded; the following treatment is editorial.


\begin{lemma}
\label{editorial-source-lemma-finite-finitely-presented-extension}
Let $R \to S$ be a finite and finitely presented ring map.
Let $M$ be an $S$-module.
Then $M$ is finitely presented as an $R$-module if and only if
$M$ is finitely presented as an $S$-module.
\end{lemma}

\begin{proof}
If $M$ is finitely presented over $R$, Lemma
\ref{algebra-lemma-finitely-presented-over-subring} applies: the given finite
map is of finite type, because its module generators also generate
its algebra. Hence $M$ is finitely presented over $S$.
To see the other assume that $M$ is finitely presented as an $S$-module.
Pick a presentation
$$
S^{\oplus m} \longrightarrow
S^{\oplus n} \longrightarrow
M \longrightarrow 0
$$
The kernel of $S^{\oplus n}\to M$ is the image of the displayed
$S^{\oplus m}$. A finite $R$-generating family of $S$ in each of
these $m$ summands maps to a finite $R$-generating family of that
kernel. Thus from
Lemma \ref{algebra-lemma-extension}
we see that it suffices to prove that $S$ is finitely presented as an
$R$-module.

\medskip\noindent
Pick $y_1, \ldots, y_n \in S$ such that $y_1, \ldots, y_n$ generate $S$
as an $R$-module. By Lemma \ref{editorial-source-lemma-characterize-integral-element}
each $y_i$ is integral over $R$. Choose monic polynomials
$P_i(x) \in R[x]$ with $P_i(y_i) = 0$. Consider the ring
$$
S' = R[x_1, \ldots, x_n]/(P_1(x_1), \ldots, P_n(x_n))
$$
Then we see that $S$ is of finite presentation as an $S'$-algebra
by Lemma \ref{algebra-lemma-compose-finite-type}. Since $S' \to S$ is surjective,
the kernel $J = \Ker(S' \to S)$ is finitely generated as an ideal by
Lemma \ref{algebra-lemma-finite-presentation-independent}. Hence $J$ is a finite
$S'$-module (immediate from the definitions).
Thus $S = \Coker(J \to S')$ is  of finite presentation as an $S'$-module
by Lemma \ref{algebra-lemma-extension}.
Hence, arguing as in the first paragraph, it suffices to show that $S'$ is
of finite presentation as an $R$-module. Actually, $S'$ is free as an
$R$-module with basis the monomials $x_1^{e_1} \ldots x_n^{e_n}$
for $0 \leq e_i < \deg(P_i)$. Namely, write $R \to S'$ as the composition
$$
R \to R[x_1]/(P_1(x_1)) \to R[x_1, x_2]/(P_1(x_1), P_2(x_2)) \to
\ldots \to S'
$$
This shows that the $i$th ring in this sequence is free as a module over the
$(i - 1)$st one with basis $1, x_i, \ldots, x_i^{\deg(P_i) - 1}$. The result
follows by induction, with the following division argument proving
both spanning and independence over every intermediate coefficient
ring. If $P_i(x_i)=x_i^{d_i}+\sum_{j<d_i}a_{ij}x_i^j$ is monic,
subtract the leading coefficient times $x_i^{e-d_i}P_i$ from
a polynomial of degree $e\geq d_i$. Repeating decreases the degree
and leaves a remainder of degree less than $d_i$, retaining each
coefficient $a_{ij}$. Uniqueness holds because for nonzero $Q$,
the product $P_iQ$ has leading coefficient equal to that of $Q$
and degree $d_i+\deg Q$, even if the coefficient ring has
zero divisors. Thus no nonzero remainder of degree less than
$d_i$ belongs to $(P_i)$. Choosing $d_i\geq1$ is always possible;
a monic equation of degree zero can be multiplied by $x_i$.
Applying the unique remainder successively in $x_1,\ldots,x_n$
gives exactly the displayed monomial basis of $S'$ over $R$,
of size $D=\prod_i d_i$ (the empty product is $1$).

To spell out the final presentations, let $h_1,\ldots,h_a$
generate the original ideal $J\subset S'$. The $aD$ elements
$h_j x_1^{e_1}\cdots x_n^{e_n}$, with the displayed exponent
bounds, generate $J$ as an $R$-module. Indeed expand each
coefficient in an ideal expression $\sum_j c_jh_j$ uniquely
in that monomial basis. Their coordinate vectors give a map
$R^{\oplus aD}\to R^{\oplus D}$ with cokernel the original
$S=S'/J$. This is an explicit finite $R$-presentation of $S$.
For the original $S^{\oplus n}$ of the first paragraph, take
one copy of this presentation per summand. That number of summands
is independent of the number of chosen $y_i$ in this paragraph.
Lift the finite $R$-generating family
of the kernel there to this finite free $R$-module. The lifted
vectors together with its existing finitely many relation vectors
generate the full kernel onto $M$: subtracting the lifts from
any vector mapping to zero leaves an old relation. This supplies
a finite $R$-presentation of the same $M$ and finishes the proof.
\end{proof}


\subsection{Editorial treatment of Algebra lines 7927--7963}

\hypertarget{algebra-context-088}{}

\noindent\href{https://github.com/stacks/stacks-project/blob/a04446e57ec1fbc252a871afcec7752fb2807b14/algebra.tex\#L7927}{Original source, lines 7927--7963}. Complete original and reviewed TeX passages accompany this note. Every intervention is separately recorded; the following treatment is editorial.


\begin{lemma}
\label{editorial-source-lemma-silly-normal}
Let $R$ be a ring. Let $x, y \in R$ be nonzerodivisors.
Let $R[x/y] \subset R_{xy}$ be the $R$-subalgebra generated
by $x/y$, and similarly for the subalgebras $R[y/x]$ and $R[x/y, y/x]$.
If $R$ is integrally closed in $R_x$ or $R_y$, then the sequence
$$
0 \to R \xrightarrow{(-1, 1)} R[x/y] \oplus R[y/x] \xrightarrow{(1, 1)}
R[x/y, y/x] \to 0
$$
is a short exact sequence of $R$-modules.
\end{lemma}

\begin{proof}
Because $x$ and $y$ are nonzerodivisors, their product is a
nonzerodivisor. Thus $R\to R_{xy}$ is injective, and so are the
maps $R_x\to R_{xy}$ and $R_y\to R_{xy}$: a fraction that becomes
zero has its numerator killed by a power of $xy$, hence has zero
numerator. All subrings in the statement therefore have their
specified inclusions in $R_{xy}$. The product of the original
elements $x/y$ and $y/x$ is $1$. Consequently each monomial
$(x/y)^a(y/x)^b$ with $a,b\geq0$ equals $(x/y)^{a-b}$ when
$a\geq b$, and equals $(y/x)^{b-a}$ otherwise. Splitting a
polynomial into these terms expresses it as a sum of elements
of $R[x/y]$ and $R[y/x]$, proving surjectivity of the map
$(u,v)\mapsto u+v$. Its kernel consists exactly of pairs
$(-\alpha,\alpha)$ with $\alpha\in R[x/y]\cap R[y/x]$.
The first map in the statement is injective and sends
$a\mapsto(-a,a)$, with the original signs.
Let $\alpha \in R[x/y] \cap R[y/x]$. To show exactness in the middle
we have to prove that $\alpha \in R$. By assumption we may write
$$
\alpha = a_0 + a_1 x/y + \ldots + a_n (x/y)^n =
b_0 + b_1 y/x + \ldots + b_m(y/x)^m
$$
for some $n, m \geq 0$ and $a_i, b_j \in R$.
Pick some $N > \max(n, m)$.
Consider the finite $R$-submodule $M$ of $R_{xy}$ generated by the elements
$$
(x/y)^N, (x/y)^{N - 1}, \ldots, x/y, 1, y/x, \ldots, (y/x)^{N - 1}, (y/x)^N
$$
We check $\alpha M\subset M$ using the full original expressions.
For $0\leq i\leq N$, multiply $(x/y)^i$ by
$\sum_{j=0}^m b_j(y/x)^j$. If $i\geq j$, the resulting
monomial is $(x/y)^{i-j}$ with $0\leq i-j\leq N$; if $i<j$,
it is $(y/x)^{j-i}$ with $0<j-i\leq m<N$.
Thus every summand, with its original coefficient $b_j$, belongs
to $M$. Multiplying $(y/x)^i$ instead by
$\sum_{j=0}^n a_j(x/y)^j$ gives $(y/x)^{i-j}$ for $i\geq j$
and $(x/y)^{j-i}$ for $i<j$, with the same bounds using $n<N$.
This proves stability for every listed generator. Since $1\in M$,
Lemma \ref{editorial-source-lemma-characterize-integral-element} proves $\alpha$
integral over $R$. The expression with coefficients $b_j$ places
$\alpha$ in the original $R_x$, and the expression with $a_j$
places it in $R_y$. Either of the two integral-closedness
hypotheses therefore gives $\alpha\in R$, proving middle exactness.

The calculation also identifies precisely what is needed without
either integral-closedness assumption. With the same original
nonzerodivisors, let $C=R[x/y]\cap R[y/x]\subset R_{xy}$.
It is a subring containing $R$, and the preceding finite-submodule
argument proves every element of $C$ integral over $R$.
There is always an exact sequence with its actual intersection
$$
0\longrightarrow C\xrightarrow{\alpha\mapsto(-\alpha,\alpha)}
R[x/y]\oplus R[y/x]\xrightarrow{(u,v)\mapsto u+v}
R[x/y,y/x]\longrightarrow0.
$$
Surjectivity and the kernel were proved above, and injection uses
the original inclusions. In the sequence with $R$ at the left,
the middle homology is exactly $C/R$: the map sends
$\alpha+R$ to the class of $(-\alpha,\alpha)$, is well defined
by the original map $R\to R[x/y]\oplus R[y/x]$, and has inverse
given by the second coordinate of a kernel pair modulo $R$.
Both composites are identities. Thus the original sequence is
exact precisely when this specified integral extension $C$ equals
$R$; the two stated hypotheses are sufficient conditions for that
equality, and the formula retains the precise defect otherwise.
\end{proof}


\subsection{Editorial treatment of Algebra lines 7971--8339}

\hypertarget{algebra-context-089}{}

\noindent\href{https://github.com/stacks/stacks-project/blob/a04446e57ec1fbc252a871afcec7752fb2807b14/algebra.tex\#L7971}{Original source, lines 7971--8339}. Complete original and reviewed TeX passages accompany this note. Every intervention is separately recorded; the following treatment is editorial.


\subsubsection{Normal rings}
\label{editorial-source-section-normal-rings}

\noindent
We first introduce the notion of a normal domain, and then we
introduce the (very general) notion of a normal ring.

\begin{definition}
\label{editorial-source-definition-domain-normal}
A domain $R$ is called {\it normal} if it is integrally
closed in its field of fractions.
\end{definition}

\begin{lemma}
\label{editorial-source-lemma-integral-closure-in-normal}
Let $R \to S$ be a ring map.
If $S$ is a normal domain, then the integral closure of $R$
in $S$ is a normal domain. It also equals the integral closure
of $R$ in the fraction field of the original $S$.
\end{lemma}

\begin{proof}
Write $B\subset S$ for that integral closure, retaining the
original map $R\to S$. The ring $B$ is a domain, since it is
a subring of the domain $S$. The inclusion induces the injective
map $\operatorname{Frac}(B)\to\operatorname{Frac}(S)$,
$a/b\mapsto a/b$; denominators remain nonzero and equality of
fractions is the same equality $ab'=a'b$ in either domain.
If $z\in\operatorname{Frac}(B)$ is integral over $B$, its monic
equation is also a monic equation over $S$. Normality of $S$
therefore places the image of $z$ in $S$. The subalgebra $B[z]$
is integral over $B$, and $R\to B$ is integral by the definition
of $B$. Transitivity (Lemma \ref{algebra-lemma-integral-transitive}) makes
$z$ integral over $R$, so $z\in B$. This proves normality of $B$.
If instead $z\in\operatorname{Frac}(S)$ is integral over $R$,
its same equation, with coefficients mapped into $S$, makes it
integral over $S$. Thus $z\in S$ and then $z\in B$.
The reverse inclusion follows from the original inclusion
$B\subset S\subset\operatorname{Frac}(S)$ and its monic equations,
proving the asserted equality of integral closures.
\end{proof}

\noindent
The following notion is occasionally useful when
studying normality.

\begin{definition}
\label{editorial-source-definition-almost-integral}
Let $R$ be a domain.
\begin{enumerate}
\item An element $g$ of the fraction
field of $R$ is called {\it almost integral over $R$}
if there exists an element $r \in R$, $r\not = 0$
such that $rg^n \in R$ for all $n \geq 0$.
\item The domain $R$ is called {\it completely normal} if every
almost integral element of the fraction field of $R$ is
contained in $R$.
\end{enumerate}
\end{definition}

\noindent
The following lemma shows that a Noetherian domain is normal
if and only if it is completely normal.

\begin{lemma}
\label{editorial-source-lemma-almost-integral}
Let $R$ be a domain with fraction field $K$.
If $u, v \in K$ are almost integral over $R$, then so are
$u + v$ and $uv$. Any element $g \in K$ which is integral over $R$
is almost integral over $R$. If $R$ is Noetherian
then the converse holds as well. More precisely, for an almost
integral $g$ and any original witness $0\ne r\in R$, integrality
of $g$ is equivalent to finite generation of the specified ideal
$I_g=\sum_{n\geq0}R\,rg^n\subset R$.
\end{lemma}

\begin{proof}
Keep witnesses $0\ne r,r'\in R$ with $ru^n\in R$ and
$r'v^n\in R$ for all $n\geq0$. Their product is nonzero, and
the full identities
$$
rr'(uv)^n=(ru^n)(r'v^n),\qquad
rr'(u+v)^n=\sum_{j=0}^n\binom nj(ru^j)(r'v^{n-j})
$$
place both expressions in $R$ for every $n\geq0$. The binomial
coefficients are integers acting through $\mathbf Z\to R$;
no division in $R$ is used. Also $r(-u)^n=(-1)^nru^n\in R$,
and every element of $R$ has witness $1$. Consequently the
almost integral elements form an $R$-subalgebra of the original $K$.
Suppose $g \in K$ is integral over $R$.
In this case there exists a $d > 0$ such that
the ring $R[g]$ is generated by $1, g, \ldots, g^d$ as an $R$-module.
Write each original $g^i=a_i/b_i$ for $0\leq i\leq d$,
with $a_i,b_i\in R$ and $b_i\ne0$, and choose the actual common
denominator $r=\prod_{i=0}^d b_i\ne0$. Then
$rg^i=a_i\prod_{j\ne i}b_j\in R$ for each listed power.
Every element of $R[g]$ is an $R$-linear combination of these
powers, so multiplying that combination by $r$ gives an element
of $R$. In particular $rg^n\in R$ for all $n\geq0$,
which proves that the original $g$ is almost integral.

\medskip\noindent
Let $g\in K$ be almost integral with its original witness
$0\ne r\in R$. The set
$I_g=\sum_{n\geq0}R\,rg^n=rR[g]\subset R$ is an ideal.
Multiplication by $r$ gives the exact $R$-module isomorphism
$$
R[g]\longrightarrow I_g,\quad z\longmapsto rz,
\qquad I_g\longrightarrow R[g],\quad w\longmapsto w/r.
$$
The second map lands in $R[g]$ by the displayed description of
$I_g$, and both composites are identities in the original field $K$.
If $I_g$ is finitely generated, so is $R[g]$, and Lemma
\ref{algebra-lemma-finite-is-integral} proves $g$ integral. Conversely,
if $g$ is integral, Lemma \ref{editorial-source-lemma-characterize-integral}
makes $R[g]$ finite, so the same isomorphism makes $I_g$ finite
for every choice of witness. This proves the exact criterion.
When $R$ is Noetherian the ideal $I_g$ is finitely generated,
giving the claimed converse. Equivalently, the original inclusion
$R[g]\subset(1/r)R$ identifies it with a submodule of the finite
$R$-module $(1/r)R$; multiplication by $r$ identifies this
submodule with the very same ideal $I_g$.
\end{proof}

\begin{lemma}
\label{editorial-source-lemma-localize-normal-domain}
If $R$ is a normal domain and $S\subset R$ is multiplicative
with $0\notin S$, then $S^{-1}R$ is a normal domain.
If $0\in S$, the localization is the zero ring, which is a normal
ring in the sense of Definition \ref{editorial-source-definition-ring-normal}.
\end{lemma}

\begin{proof}
If $0\in S$, every fraction is zero and the spectrum is empty,
so the later definition of a normal ring applies vacuously.
Suppose $0\notin S$ and keep $K=\operatorname{Frac}(R)$.
The original localization injects into $K$ by $a/s\mapsto a/s$.
Its fraction field is identified with $K$ by the inverse maps
$$
\frac{a/s}{b/t}\longmapsto\frac{at}{sb},\qquad
\frac ab\longmapsto\frac{a/1}{b/1}.
$$
All denominators are nonzero; fraction equality in a domain
verifies well-definedness, and multiplication of the original
numerators and denominators verifies both composites.
Let $g\in K$ be integral over $S^{-1}R$, with its monic equation
$P(g)=0$ for $P(x)=x^d+\sum_{0\leq j<d}a_jx^j$, $d\geq1$.
Write $a_j=c_j/s_j$ with $c_j\in R$, $s_j\in S$, and keep
$s=\prod_{j=0}^{d-1}s_j\in S$.
Then $sa_j=c_j\prod_{h\ne j}s_h\in R$, and
$s^{d-j}a_j=s^{d-j-1}(sa_j)\in R$ for every $0\leq j<d$.
The full scaled equation is
$$
(sg)^d+\sum_{j=0}^{d-1}s^{d-j}a_j(sg)^j
=s^d\left(g^d+\sum_{j=0}^{d-1}a_jg^j\right)=0.
$$
It is monic with coefficients in $R$, so normality gives
$sg\in R$. Thus the original $g$ is the image of $(sg)/s$
in $S^{-1}R$, proving the assertion.
\end{proof}

\begin{lemma}
\label{editorial-source-lemma-PID-normal}
A principal ideal domain is normal.
\end{lemma}

\begin{proof}
Let $R$ be a principal ideal domain and retain the original
$g=a/b\in\operatorname{Frac}(R)$, where $b\ne0$, and the full
monic equation
$$
(a/b)^n+\sum_{i=0}^{n-1}r_i(a/b)^i=0,\qquad r_i\in R.
$$
Multiplication by the original $b^n$ gives
$a^n+\sum_{i=0}^{n-1}r_i a^i b^{n-i}=0$ in $R$.
Choose $c$ generating the original ideal $(a,b)$, and write
$a=ca'$, $b=cb'$, $c=ua+vb$ for $a',b',u,v\in R$.
Since $b\ne0$, both $c$ and $b'$ are nonzero. Cancellation in
$c=c(ua'+vb')$ gives $ua'+vb'=1$.
The full original equation, after these substitutions, is
$$
c^n\left((a')^n+\sum_{i=0}^{n-1}r_i(a')^i(b')^{n-i}\right)=0.
$$
The nonzero factor $c^n$ in this domain makes the parenthesized
expression zero. Expanding the original identity
$1=(ua'+vb')^n$ and substituting that expression gives $1=b'w$,
where the entire inverse is
$$
w=-u^n\sum_{i=0}^{n-1}r_i(a')^i(b')^{n-i-1}
+\sum_{j=0}^{n-1}\binom nj(ua')^jv^{n-j}(b')^{n-j-1}\in R.
$$
Indeed the omitted top term of the latter binomial sum is
$u^n(a')^n$, exactly the term replaced using the monic equation.
Returning to the original numerator and denominator gives
$b(a'w)=cb'a'w=ca'=a$. Hence $g=a/b=a'w\in R$.
Every factor from the original fraction and every coefficient
of its monic equation is retained in this comparison.
\end{proof}

\begin{lemma}
\label{editorial-source-lemma-prepare-polynomial-ring-normal}
Let $R$ be a domain with fraction field $K$.
Suppose $f = \sum \alpha_i x^i$ is an
element of $K[x]$.
\begin{enumerate}
\item $f$ is integral over $R[x]$ if and only if
all $\alpha_i$ are integral over $R$, and
\item $f$ is almost integral over $R[x]$ if and only if
all $\alpha_i$ are almost integral over $R$.
\end{enumerate}
\end{lemma}

\begin{proof}
If $f=0$, all its coefficients are zero and the assertions hold
directly. Suppose $f\ne0$. We first prove the forward implication
in the second statement, keeping all coefficients, including zeros.
Write $f = \alpha_0 + \alpha_1 x + \ldots + \alpha_r x^r$
with $\alpha_r \not = 0$.
By assumption there exists $h = b_0 + b_1 x + \ldots + b_s x^s \in R[x]$,
$b_s \not = 0$ such that $f^n h \in R[x]$ for all
$n \geq 0$. This implies that $b_s \alpha_r^n \in R$
for all $n \geq 0$. Hence $\alpha_r$ is almost
integral over $R$. Since the set of almost integral
elements forms a subring (Lemma \ref{editorial-source-lemma-almost-integral}) we deduce that
$f - \alpha_r x^r = \alpha_0 + \alpha_1 x + \ldots + \alpha_{r - 1} x^{r - 1}$
is almost integral over $R[x]$. Here $\alpha_rx^r$ has the
actual witness $b_s$, since
$b_s(\alpha_rx^r)^n=(b_s\alpha_r^n)x^{rn}\in R[x]$.
The product $b_sh$ is consequently a witness for their difference
by the binomial calculation in Lemma \ref{editorial-source-lemma-almost-integral}.
The induction can be made fully explicit. For $0\leq t\leq r$
keep $f_t=\sum_{i=0}^t\alpha_ix^i$ and set
$$
h_t=b_s^{2^{r-t}-1}h,\qquad c_t=b_s^{2^{r-t}}.
$$
The original hypothesis says $h_rf_r^n\in R[x]$ for all $n\geq0$.
If $h_tf_t^n\in R[x]$, its coefficient of $x^{tn+s}$ is
$c_t\alpha_t^n$, also when $\alpha_t=0$. Hence this coefficient
belongs to $R$. The full binomial formula for
$f_{t-1}=f_t-\alpha_tx^t$ shows
$$
h_tc_t f_{t-1}^n
=\sum_{j=0}^n\binom nj(h_tf_t^j)
 \bigl(c_t(-\alpha_tx^t)^{n-j}\bigr)\in R[x].
$$
Since $h_tc_t=h_{t-1}$, this proves the descending induction.
In particular each original $\alpha_t$ has nonzero witness $c_t$,
and the single nonzero element $b_s^{2^r}$ is a witness for
all coefficients, by multiplying their equations by the remaining
power of $b_s$.

Conversely, if $\alpha_i$ has witness $0\ne w_i\in R$ for
$0\leq i\leq r$, put $w=\prod_iw_i\ne0$. For every $n\geq0$
the complete multinomial expansion gives
$$
w f^n=\sum_{k_0+\cdots+k_r=n}
\binom{n}{k_0,\ldots,k_r}
\left(\prod_{i=0}^r w_i\alpha_i^{k_i}\right)
x^{\sum_{i=0}^r i k_i}\in R[x].
$$
The indices $k_i$ are nonnegative integers and the multinomial
coefficients are integers, not inverses of factorials in $R$.
This proves the reverse almost-integral implication.

\medskip\noindent
In order to prove the first statement we will use absolute Noetherian
reduction. Namely, write $\alpha_i = a_i / b_i$ and
let $P(t) = t^d + \sum_{j < d} f_j t^j$ be a polynomial
with coefficients $f_j \in R[x]$ such that $P(f) = 0$.
Let $f_j = \sum f_{ji}x^i$. Consider the subring
$R_0 \subset R$ generated by the finite list of elements
$a_i, b_i, f_{ji}$ of $R$. It is a domain; let
$K_0$ be its field of fractions. Since $R_0$ is a finite type
$\mathbf{Z}$-algebra it is Noetherian, see
Lemma \ref{editorial-source-lemma-obvious-Noetherian}. It is still
the case that $f \in K_0[x]$ is integral over $R_0[x]$,
because the inclusion $R_0\subset R$ induces the injective maps
$K_0\to K$, $a/b\mapsto a/b$, and $K_0[x]\to K[x]$.
Each original denominator $b_i$ is nonzero in $R_0$, all the
original coefficients $f_{ji}$ belong to it, and the same monic
polynomial $P$ and same $f$ have $P(f)\in K_0[x]$.
Its image is zero in $K[x]$, so injectivity gives $P(f)=0$
already in $K_0[x]$.
By Lemma \ref{editorial-source-lemma-almost-integral} the element
$f$ is almost integral over $R_0[x]$. By the second statement of
the lemma, the elements $\alpha_i$ are almost integral
over $R_0$. And since $R_0$ is Noetherian, they are
integral over $R_0$, see Lemma \ref{editorial-source-lemma-almost-integral}.
Each resulting monic equation over $R_0$ maps coefficient by
coefficient to a monic equation for the same $\alpha_i$ over $R$.
For the reverse implication, if all the finitely many $\alpha_i$
are integral over $R$, the original algebra
$B=R[\alpha_0,\ldots,\alpha_r]\subset K$ is finite over $R$
by Lemma \ref{editorial-source-lemma-characterize-integral}. A finite module
generating family of $B$ also generates $B[x]$ over $R[x]$:
express each coefficient of a polynomial in that family.
Thus the original $f\in B[x]$ is integral over $R[x]$ by
Lemma \ref{algebra-lemma-finite-is-integral}, proving the other implication.
\end{proof}

\begin{lemma}
\label{editorial-source-lemma-polynomial-domain-normal}
Let $R$ be a normal domain.
Then $R[x]$ is a normal domain. For any domain $R$, the polynomial
ring $R[x]$ is completely normal if and only if $R$ is completely normal.
\end{lemma}

\begin{proof}
The result is true if $R$ is a field $K$ because
$K[x]$ is a euclidean domain and hence a principal ideal
domain and hence normal by Lemma \ref{editorial-source-lemma-PID-normal}.
The inclusions $R[x]\subset K[x]$ identify their fraction fields
with $K(x)$. More explicitly a fraction of polynomials over $R$
maps to the same fraction over $K$. Conversely, for $A/B$ with
$A,B\in K[x]$, choose a common nonzero denominator $c\in R$
of every coefficient of both polynomials. Its inverse image is
$(cA)/(cB)$, with both original scaled polynomials in $R[x]$.
The field equality criterion verifies independence of the chosen
common denominator and both composites.
Let $g$ be an element of the fraction field of
$R[x]$ which is integral over $R[x]$. Because $g$
is integral over $K[x]$ where $K$ is the fraction
field of $R$ we may write $g = \alpha_d x^d + \alpha_{d-1}x^{d-1} +
\ldots + \alpha_0$ with $\alpha_i \in K$.
By Lemma \ref{editorial-source-lemma-prepare-polynomial-ring-normal}
the elements $\alpha_i$ are integral over $R$ and
hence are in $R$. Thus the original $g$ belongs to $R[x]$.

For the second assertion, suppose $R$ is completely normal and
$g\in K(x)$ is almost integral over $R[x]$. Its same nonzero
polynomial witness also makes it almost integral over $K[x]$.
The ring $K[x]$ is Noetherian, since it is a principal ideal
domain, so Lemma \ref{editorial-source-lemma-almost-integral} makes $g$ integral
over $K[x]$. Normality of $K[x]$ then gives $g\in K[x]$.
The proved coefficientwise almost-integral equivalence shows
that every coefficient of this polynomial is almost integral
over $R$, hence lies in $R$. Therefore $g\in R[x]$.
Conversely, an almost integral element $g\in K$ and its witness
$0\ne r\in R$ remain the same constant fraction and witness
over $R[x]$. Complete normality of $R[x]$ places $g$ in
$R[x]\cap K=R$, where equality follows from uniqueness of
polynomial coefficients. This proves complete normality of $R$.
\end{proof}

\begin{lemma}
\label{editorial-source-lemma-power-series-over-Noetherian-normal-domain}
Let $R$ be a Noetherian normal domain. Then $R[[x]]$ is
a Noetherian normal domain. More generally, for any domain $R$,
$R[[x]]$ is completely normal if and only if $R$ is completely normal.
\end{lemma}

\begin{proof}
First prove the second assertion in the forward direction.
Assume that $R$ is completely normal, and retain nonzero
$f,g\in R[[x]]$ and the original $w=f/g$ in its fraction field,
now assumed almost integral. The power series ring is a domain:
the first nonzero coefficients of two nonzero series multiply
to the first nonzero coefficient of their product. Coefficient
inclusion gives an injection $R[[x]]\to K[[x]]$, where
$K=\operatorname{Frac}(R)$, hence an injection of its fraction
field into $K((x))$. A nonzero series over $K$ is $x^a$ times
a series with nonzero constant coefficient. For a series
$x^a\sum_{j\geq0}c_jx^j$ with $c_0\ne0$, its inverse is
$x^{-a}\sum_{j\geq0}d_jx^j$, where
$d_0=c_0^{-1}$ and
$d_j=-c_0^{-1}\sum_{i=1}^j c_id_{j-i}$ for $j\geq1$.
The coefficient of degree zero in the product is $1$ and every
higher coefficient is zero by this recurrence, proving the
inverse formula in $K((x))$. Thus the
original $w$ has the unique Laurent expansion
$$
w=a_nx^n+a_{n+1}x^{n+1}+\cdots,
\qquad n\in\mathbf Z,\quad a_n\ne0.
$$
Keep its original nonzero almost-integrality witness
$$
h=b_mx^m+b_{m+1}x^{m+1}+\cdots\in R[[x]],
\qquad b_m\ne0,
$$
so $hw^e\in R[[x]]$ for every $e\geq0$.
The first nonzero term of $hw^e$ has degree $m+en$ and
coefficient $b_ma_n^e$. If $n<0$, choose $e$ with $m+en<0$;
that nonzero negative-degree term contradicts membership in
$R[[x]]$. Hence $n\geq0$. For every $e\geq0$ the displayed
coefficient belongs to $R$, so complete normality gives $a_n\in R$.

Suppose all original coefficients $a_n,\ldots,a_{N-1}$ have
been proved to belong to $R$, retaining zeros in this list, and
put $P_N=\sum_{i=n}^{N-1}a_ix^i\in R[x]$.
The same original $h$ is a witness for $w-P_N$, since the full
identity
$$
h(w-P_N)^e=\sum_{j=0}^e\binom ej(hw^j)(-P_N)^{e-j}
$$
has every term in $R[[x]]$. The coefficient of $x^{m+eN}$
on the left is $b_ma_N^e$, including the case $a_N=0$:
$w-P_N$ has no terms of degree below $N$, so only these leading
terms can contribute at that degree. Thus $b_ma_N^e\in R$
for every $e\geq0$. The original nonzero $b_m$ witnesses
almost integrality of $a_N$ over $R$, giving $a_N\in R$.
Induction proves every coefficient belongs to $R$, hence the
original fraction $w$ belongs to $R[[x]]$. The zero fraction
belongs to the ring directly. This proves complete normality
of $R[[x]]$ without any Noetherian assumption.

Conversely, if $R[[x]]$ is completely normal and $z\in K$ is
almost integral over $R$, its original constant numerator,
denominator and witness define the same almost integral constant
fraction in the fraction field of $R[[x]]$. Thus it lies in
$R[[x]]\cap K=R$: uniqueness of the Laurent expansion forces
all coefficients but the constant one to be zero. Hence $R$
is completely normal, proving the equivalence.

Finally if the original $R$ is Noetherian and normal, Lemma
\ref{editorial-source-lemma-almost-integral} says that it is completely normal.
The argument just given makes $R[[x]]$ completely normal, hence
normal since every integral element is almost integral.
Its Noetherian property follows from
Lemma \ref{editorial-source-lemma-Noetherian-power-series}, proving the first assertion.
\end{proof}

\begin{lemma}
\label{editorial-source-lemma-normality-is-local}
Let $R$ be a domain. The following are equivalent:
\begin{enumerate}
\item The domain $R$ is a normal domain,
\item for every prime $\mathfrak p \subset R$ the local ring
$R_{\mathfrak p}$ is a normal domain, and
\item for every maximal ideal $\mathfrak m$ the ring $R_{\mathfrak m}$
is a normal domain.
\end{enumerate}
\end{lemma}

\begin{proof}
We deduce (1) $\Rightarrow$ (2) from Lemma \ref{editorial-source-lemma-localize-normal-domain}.
The implication (2) $\Rightarrow$ (3) is immediate. The implication
(3) $\Rightarrow$ (1) follows from the fact that for any domain $R$ we have
$$
R = \bigcap\nolimits_{\mathfrak m} R_{\mathfrak m}
$$
inside the original fraction field of $R$. To prove this equality,
let $g$ belong to every $R_{\mathfrak m}$ and retain
$I=\{x\in R:xg\in R\}$. It is an ideal: multiplication by
$g$ respects sums and multiplication by elements of $R$.
For each maximal $\mathfrak m$, the expression $g=a/s$ in
$R_{\mathfrak m}\subset\operatorname{Frac}(R)$ has
$s\notin\mathfrak m$ and gives $sg=a\in R$. Thus $s\in I$
and $I\not\subseteq\mathfrak m$. If $I$ were proper it would
lie in a maximal ideal, a contradiction. Therefore $1\in I$
and $g\in R$. Finally if $g$ is integral over $R$, its original
monic equation remains monic over each $R_{\mathfrak m}$.
Under (3), each of those rings contains $g$, and the proved
intersection identity gives $g\in R$, establishing (1).
\end{proof}

\noindent
Lemma \ref{editorial-source-lemma-normality-is-local} shows that the following definition
is compatible with Definition \ref{editorial-source-definition-domain-normal}. (It is the
definition from EGA -- see \cite[IV, 5.13.5 and 0, 4.1.4]{EGA}.)

\begin{definition}
\label{editorial-source-definition-ring-normal}
A ring $R$ is called {\it normal} if for every prime
$\mathfrak p \subset R$ the localization $R_{\mathfrak p}$ is
a normal domain (see Definition \ref{editorial-source-definition-domain-normal}).
\end{definition}

\noindent
It suffices in this definition to test maximal ideals. Indeed,
if each $R_{\mathfrak m}$ is a normal domain and
$\mathfrak p\subseteq\mathfrak m$, there is the exact isomorphism
$$
(R_{\mathfrak m})_{\mathfrak pR_{\mathfrak m}}
\longrightarrow R_{\mathfrak p},\qquad
\frac{a/s}{b/t}\longmapsto\frac{at}{sb},
$$
with inverse $a/u\mapsto(a/1)/(u/1)$.
Here $s,t\notin\mathfrak m$ and $b\notin\mathfrak p$, so
$sb\notin\mathfrak p$. The prime correspondence for localization
identifies $\mathfrak pR_{\mathfrak m}$ as a proper prime with
contraction $\mathfrak p$. Fraction equality witnesses map to
fraction equality witnesses in these localizations; the displayed
formulas fix every original fraction in both composites.
Lemma \ref{editorial-source-lemma-localize-normal-domain} then makes
$R_{\mathfrak p}$ a normal domain. The converse is part of the
definition, and every prime lies in a maximal ideal.

A normal ring is reduced. More explicitly the natural map
$R\to\prod_{\mathfrak m}R_{\mathfrak m}$ is injective: if
$a\ne0$, its annihilator is a proper ideal and is contained in
a maximal ideal $\mathfrak m$; then $a/1$ cannot be zero there,
since a denominator outside $\mathfrak m$ killing $a$ would
belong to that annihilator. A nilpotent element maps to zero
in every domain $R_{\mathfrak m}$ and hence is zero in $R$.
This also proves the stated subring assertion using all primes
(see Lemma \ref{algebra-lemma-characterize-zero-local}).
For the zero ring the products are empty and all assertions
hold directly.

\begin{lemma}
\label{editorial-source-lemma-normal-ring-integrally-closed}
A normal ring is integrally closed in its total ring of fractions.
\end{lemma}

\begin{proof}
Let $R$ be normal, let $\Sigma\subset R$ be its original
multiplicative set of nonzerodivisors, and keep
$Q(R)=\Sigma^{-1}R$. The map $R\to Q(R)$ is injective,
since any element in its kernel is killed by a member of $\Sigma$.
Let $x\in Q(R)$ be integral over $R$, and retain the ideal
$I=\{f\in R\mid fx\in R\}$.
Fix a prime $\mathfrak p$ and put $A=R_{\mathfrak p}$, a
normal domain. Every $s\in\Sigma$ has nonzero image in $A$:
an equality $s/1=0$ would give $u\notin\mathfrak p$ with
$us=0$, forcing $u=0$, a contradiction.
There is the exact comparison
$$
Q(R)\otimes_RA\cong\Sigma^{-1}A
\lhook\joinrel\longrightarrow\operatorname{Frac}(A).
$$
The first map sends $(r/s)\otimes c$ to $(rc)/(s/1)$;
its inverse sends $c/(s/1)$ to $(1/s)\otimes c$.
The original coefficient maps agree on $R$, so tensor balancing
gives the first map. The second is defined because each $s/1$
maps to the unit $(s/1)\otimes1$ with inverse $(1/s)\otimes1$.
The formulas verify both composites. The last map is injective
since $A$ is a domain and every image of $\Sigma$ is nonzero.
For the original fractions its full formula is
$$
(r/s)\otimes(a/u)\longmapsto
\frac{(ra)/1}{(su)/1},\qquad
s\in\Sigma,\quad u\notin\mathfrak p.
$$
This embeds the indicated localization in the fraction field;
it does not require that every nonzero element of $A$ come from
an original nonzerodivisor of $R$.

The image of $x\otimes1$ satisfies the same monic equation
with coefficients mapped into $A$. Normality of $A$ places
it in $A$. The injective comparison just proved therefore gives
$x\otimes1=a\otimes(1/f)$ in
$Q(R)\otimes_RR_{\mathfrak p}\cong Q(R)_{\mathfrak p}$ for
some original $a,f\in R$, $f\notin\mathfrak p$.
This last isomorphism sends $z\otimes(a/u)$ to $az/u$
and has inverse $z/u\mapsto z\otimes(1/u)$; tensor balancing
and the fraction equality relation verify both maps and composites.
Hence $fx-a$ becomes zero in that localization, so there is
$f'\notin\mathfrak p$ with $f'(fx-a)=0$ in $Q(R)$.
The equality $f'fx=f'a$ has its right side in the embedded $R$,
so $ff'\in I$ and $ff'\notin\mathfrak p$.
As this holds for every prime, $I$ cannot be proper; otherwise
a maximal ideal containing it would contradict this property.
Thus $1\in I$ and the original $x$ lies in $R$.
If $R=0$, its total quotient ring is zero and the assertion
holds without any choice of a prime.
\end{proof}

\begin{lemma}
\label{editorial-source-lemma-localization-normal-ring}
A localization of a normal ring is a normal ring.
\end{lemma}

\begin{proof}
Let $S\subset R$ be the original multiplicative subset and
let $\mathfrak q$ be a prime of $S^{-1}R$. Its contraction
$\mathfrak p$ to $R$ is prime and disjoint from $S$, and
$\mathfrak q=S^{-1}\mathfrak p$. The maps
$$
(S^{-1}R)_{\mathfrak q}\longrightarrow R_{\mathfrak p},
\quad\frac{a/s}{b/t}\longmapsto\frac{at}{sb},
\qquad
R_{\mathfrak p}\longrightarrow(S^{-1}R)_{\mathfrak q},
\quad a/u\longmapsto\frac{a/1}{u/1}
$$
are defined because $s,t\in S$ and $b\notin\mathfrak p$,
while $u\notin\mathfrak p$ in the second formula. The original
maps from $R$ send each required denominator to a unit, hence
extend to these localizations; on fractions the displayed formulas
prove that the two extensions are inverse. Normality of $R$
makes $R_{\mathfrak p}$ a normal domain, proving the assertion
at every prime of $S^{-1}R$. If this localization is zero,
there are no primes to test and it is normal by definition.
\end{proof}

\begin{lemma}
\label{editorial-source-lemma-polynomial-ring-normal}
Let $R$ be a normal ring. Then $R[x]$ is a normal ring.
More generally, the polynomial ring over $R$ in any set of
indeterminates is normal.
\end{lemma}

\begin{proof}
Let $\mathfrak q\subset R[x]$ be prime and let
$\mathfrak p$ be its contraction to the original $R$.
The ring $R_{\mathfrak p}[x]$ is a normal domain by
Lemma \ref{editorial-source-lemma-polynomial-domain-normal}. It is canonically
$(R\setminus\mathfrak p)^{-1}R[x]$: a finite polynomial with
fraction coefficients is represented by multiplying all their
original denominators to obtain one common denominator, and
the inverse sends $F/s$ to the polynomial whose coefficients
are the corresponding fractions. These maps preserve coefficient
equalities and are inverse. The extended ideal
$\widetilde{\mathfrak q}=(R\setminus\mathfrak p)^{-1}\mathfrak q$
is prime, with contraction $\mathfrak q$. The exact local maps are
$$
(R[x])_{\mathfrak q}\longrightarrow
(R_{\mathfrak p}[x])_{\widetilde{\mathfrak q}},
\quad F/G\longmapsto(F/1)/(G/1),
$$
with inverse $(F/s)/(G/t)\mapsto tF/(sG)$.
Here $s,t\in R\setminus\mathfrak p$ and $G\notin\mathfrak q$,
so every displayed denominator is valid. The original coefficient
maps invert the indicated sets, proving well-definedness by the
localization relations, and the fraction formulas verify both
composites. Lemma \ref{editorial-source-lemma-localize-normal-domain} makes this
local ring a normal domain, proving the one-variable assertion.

Induction proves the assertion for any finite set of variables.
For an arbitrary set $J$, retain all its original labeled variables.
The polynomial ring $R[x_j:j\in J]$ is the directed colimit of
$R[x_j:j\in F]$ over the finite subsets $F\subset J$:
every polynomial has finite support, and coefficient equality
already holds in a finite set containing the variables of both
polynomials. The transition maps keep each original variable and
coefficient. Lemma \ref{editorial-source-lemma-colimit-normal-ring} applies to
these normal rings and gives the full assertion. Its proof below
does not use this polynomial-ring assertion.
\end{proof}

\begin{lemma}
\label{editorial-source-lemma-finite-product-normal}
A finite product of normal rings is normal. More generally,
for any set $I$, the product $\prod_{i\in I}R_i$ is normal
if and only if every original factor $R_i$ is normal.
Every ultraproduct of normal rings is normal as well.
\end{lemma}

\begin{proof}
It suffices to show that the product of two normal rings, say $R$ and $S$, is
normal. By Lemma \ref{algebra-lemma-disjoint-decomposition} the prime ideals of
$R\times S$ are of the form $\mathfrak{p}\times S$ and $R\times
\mathfrak{q}$, where $\mathfrak{p}$ and $\mathfrak{q}$ are primes of $R$
and $S$ respectively. Localization yields 
$(R\times S)_{\mathfrak{p}\times S}=R_{\mathfrak{p}}$ which is a normal domain
by assumption. Explicitly the map sends
$(a,b)/(s,t)$ to $a/s$, and its inverse sends $a/s$ to
$(a,0)/(s,1)$. Both are defined because $s\notin\mathfrak p$.
The idempotent $(1,0)$ is a unit in this localization, so $(0,1)$
is zero there, verifying the inverse identities. The corresponding
maps at $R\times\mathfrak q$ use the other coordinate, proving
the finite assertion without any assumption on the factor maps.
\medskip\noindent
For the arbitrary-product assertion we first prove an algebraic
characterization. A ring $B$ is normal if and only if, for every
$n\geq1$ and all $a,b,c_0,\ldots,c_{n-1}\in B$, it satisfies
$$
a^n+\sum_{j=0}^{n-1}c_ja^jb^{n-j}=0
\quad\Longrightarrow\quad a=bd\text{ for some }d\in B.
\qquad(\mathcal D_n)
$$
Here $b$ may be zero or a zero divisor.
Suppose first that $B$ is normal and the displayed equation holds.
At a prime $\mathfrak p$, write $a_{\mathfrak p},b_{\mathfrak p}$
for the original images in $B_{\mathfrak p}$. If
$b_{\mathfrak p}=0$, the equation gives $a_{\mathfrak p}^n=0$,
hence $a_{\mathfrak p}=0$ since this localization is a domain.
If $b_{\mathfrak p}\ne0$, division by $b_{\mathfrak p}^n$
gives the monic equation
$$
(a_{\mathfrak p}/b_{\mathfrak p})^n+
\sum_{j=0}^{n-1}c_{j,\mathfrak p}
 (a_{\mathfrak p}/b_{\mathfrak p})^j=0.
$$
Integral closedness then gives
$a_{\mathfrak p}\in b_{\mathfrak p}B_{\mathfrak p}$ in both cases.
This implies $a\in bB$: otherwise the ideal
$J=\{r\in B:ra\in bB\}$ is proper and lies in a maximal ideal
$\mathfrak p$. Membership in the localized principal ideal gives
$s\notin\mathfrak p$ with $sa\in bB$ (clear the denominator and
retain the further multiplier witnessing equality in the localization).
This contradicts $s\in J\subset\mathfrak p$ and proves
$(\mathcal D_n)$.

Conversely suppose all $(\mathcal D_n)$ hold. Condition
$(\mathcal D_2)$ with $b=0$ and both coefficients zero gives
$z^2=0\Rightarrow z=0$. If $ab=0$, apply $(\mathcal D_2)$
to the exact equation $a^2-a(a+b)=0$ to get
$a=(a+b)d$ for some $d\in B$. Put
$z=(1-d)a=db$. Then
$z^2=((1-d)a)(db)=d(1-d)ab=0$, so $z=0$ and
$$
(1-d)a=0,\qquad db=0,\qquad (1-d)+d=1.
$$
These identities show that every $B_{\mathfrak p}$ is a domain.
Indeed a zero product $(a/u)(b/v)=0$ has a witness
$w\notin\mathfrak p$ with $wab=0$ in $B$. Apply the identities
to $wa$ and $b$. At least one of $d,1-d$ is outside
$\mathfrak p$. If $d$ is outside, $b/v=0$ in the localization;
if $1-d$ is outside, $a/u=0$ since $w$ is also a unit there.
The localization is nonzero because $\mathfrak p$ is proper.

To prove this domain integrally closed, retain a fraction
$y=(a_0/s)/(b_0/t)$ with $s,t\notin\mathfrak p$ and
$b_0/t\ne0$. Set $a=a_0t$ and $b=b_0s$ in the original $B$,
so $b/1\ne0$ and $y=(a/1)/(b/1)$ with these factors recorded.
Suppose its original monic equation is
$$
y^n+\sum_{j=0}^{n-1}(r_j/t_j)y^j=0,
\qquad r_j\in B,\quad t_j\notin\mathfrak p.
$$
Put $u=\prod_{j=0}^{n-1}t_j$ and
$c_j=r_j\prod_{k\ne j}t_k$, retaining every denominator factor.
Then $r_j/t_j=c_j/u$ in the localization. Multiplication by
$(b/1)^n(u/1)^n$ says that the element
$$
F=(ua)^n+\sum_{j=0}^{n-1}c_j u^{n-1-j}(ua)^jb^{n-j}
$$
has zero image in $B_{\mathfrak p}$. Therefore some
$v\notin\mathfrak p$ satisfies $vF=0$ in the original ring.
The full equation
$$
(vua)^n+\sum_{j=0}^{n-1}c_j u^{n-1-j}(vua)^j(vb)^{n-j}
=v^nF=0
$$
holds there, because $n\geq1$. Condition $(\mathcal D_n)$ gives
$vua=vbd$ for some $d\in B$. Since $v/1$ and $u/1$ are units
and $b/1\ne0$ in the localized domain, the original fraction
satisfies $y=d/u\in B_{\mathfrak p}$. This proves the characterization,
including the zero ring, whose prime-local condition is vacuous.

Now let $R=\prod_{i\in I}R_i$ with every $R_i$ normal.
An equation as in $(\mathcal D_n)$ in $R$ means precisely
$$
a_i^n+\sum_{j=0}^{n-1}c_{j,i}a_i^jb_i^{n-j}=0
\quad\text{in every original }R_i.
$$
Choose $d_i\in R_i$ with $a_i=b_id_i$ using the proved
characterization, and put $d=(d_i)_{i\in I}$. It gives $a=bd$
in the original product. Thus $R$ satisfies every $(\mathcal D_n)$
and is normal. Only the equation degree and its finitely many
coefficients are common; no characteristic or finiteness condition
on $I$ or on the factors is used.

We describe the exact prime and localization maps as well,
including the primes of an infinite product that are not obtained
from a single coordinate. For $E\subseteq I$, let $e_E\in R$
have coordinate $1$ on $E$ and $0$ on its complement. For a
prime $\mathfrak p\subset R$, the family
$$
\mathcal U=\{E\subseteq I:e_E\notin\mathfrak p\}
$$
is an ultrafilter, that is, a family closed under finite intersections
and enlargement, containing $I$ but not the empty set, and containing
exactly one of each subset and its complement. These facts follow
respectively from
$$
e_Ee_F=e_{E\cap F},\quad
e_E=e_Ee_F\ (E\subseteq F),\quad
e_I=1,\quad e_\varnothing=0,
$$
and from $e_Ee_{I\setminus E}=0$,
$e_E+e_{I\setminus E}=1$, using primality.
Set
$$
K_{\mathcal U}=\{a=(a_i):\{i:a_i=0\}\in\mathcal U\},
\qquad A_{\mathcal U}=R/K_{\mathcal U},
\qquad\pi:R\to A_{\mathcal U}.
$$
The set $K_{\mathcal U}$ is an ideal: intersections of zero sets
control sums, and multiplication by any tuple enlarges a zero set.
The quotient $A_{\mathcal U}$ is the ultraproduct, here defined by
this exact ideal. Also $K_{\mathcal U}\subseteq\mathfrak p$:
if $E=\{i:a_i=0\}\in\mathcal U$, then $e_Ea=0$ and
$e_E\notin\mathfrak p$, so primality gives $a\in\mathfrak p$.
Therefore $\mathfrak q=\mathfrak p/K_{\mathcal U}$ is a prime
of $A_{\mathcal U}$, and the map
$$
R_{\mathfrak p}\longrightarrow (A_{\mathcal U})_{\mathfrak q},
\qquad a/s\longmapsto\pi(a)/\pi(s)
$$
is defined, because $s\notin\mathfrak p$ if and only if
$\pi(s)\notin\mathfrak q$. It is surjective by lifting each
numerator and denominator through the quotient. It is injective:
if $\pi(a)/\pi(s)=0$, choose $t\notin\mathfrak p$ with
$ta\in K_{\mathcal U}$. For
$E=\{i:t_ia_i=0\}\in\mathcal U$, one has
$e_Eta=0$ and $e_Et\notin\mathfrak p$. This is the original
localization witness that $a/s=0$. The inverse map is consequently
$\pi(a)/\pi(s)\mapsto a/s$, independent of the chosen lifts by
the same argument applied to their difference. Both composites fix
the original fractions.

Conversely a prime of an ultraproduct contracts to a prime of
the product and recovers its ultrafilter: the class of $e_E$ is
$1$ when $E\in\mathcal U$ and $0$ when its complement is in
$\mathcal U$. Thus, as sets, all product primes are uniquely the
pairs of an ultrafilter and a prime of its ultraproduct, with the
displayed localization isomorphism for each pair.

The ultraproduct of normal rings is itself normal. If a displayed
$(\mathcal D_n)$ equation holds in $A_{\mathcal U}$, it holds
coordinatewise on a set $E\in\mathcal U$ by the definition of
the quotient ideal. For $i\in E$ choose $d_i$ with $a_i=b_id_i$
in $R_i$; set $d_i=0$ elsewhere. Their classes satisfy
$\pi(a)=\pi(b)\pi(d)$. The characterization proves normality
of $A_{\mathcal U}$, so the localization isomorphism also gives
the product result through its actual primes.

When every factor is a nonzero normal domain, these maps have
a stronger description. The ultraproduct is then a domain:
the zero set of a product is the union of the zero sets of its
two factors, so membership in an ultrafilter forces one of those
zero sets to belong to it. Its identity is nonzero because the
coordinate identities are nonzero. If $\pi(a)/\pi(b)$ is integral
in its fraction field with $\pi(b)\ne0$, then the sets
$$
T=\{i:b_i\ne0\},\qquad
E=\{i:a_i^n+\sum_{j=0}^{n-1}c_{j,i}a_i^jb_i^{n-j}=0\}
$$
belong to $\mathcal U$. On $E\cap T$, the original fraction
$a_i/b_i$ satisfies its monic equation over $R_i$ and belongs
to $R_i$. Choose that value for $d_i$ there and zero elsewhere.
Then $\pi(a)=\pi(b)\pi(d)$, proving integral closedness directly.
Moreover in this domain-factor case
$\ker(R\to R_{\mathfrak p})=K_{\mathcal U}$.
One inclusion was proved by the idempotent witnesses above.
For the other, if $ta=0$ with $t\notin\mathfrak p$, then
$t\notin K_{\mathcal U}$, so $\{i:t_i\ne0\}\in\mathcal U$;
on that set the domains force $a_i=0$, giving $a\in K_{\mathcal U}$.

Finally normality of the product implies normality of each factor.
For $\mathfrak r\in\operatorname{Spec}(R_i)$, take
$\mathfrak p=\operatorname{pr}_i^{-1}(\mathfrak r)$.
The localization map
$$
R_{\mathfrak p}\longrightarrow(R_i)_{\mathfrak r},
\qquad a/s\longmapsto a_i/s_i
$$
has inverse obtained by using the numerator tuple supported at $i$
and the denominator tuple equal to $s_i$ at $i$ and to $1$ elsewhere.
Multiplication by $e_{\{i\}}\notin\mathfrak p$ kills any
differences outside that coordinate, which verifies both composites
and independence of fraction representatives. Hence these local
rings are isomorphic, proving normality of the factor.
Zero-ring factors cause no exception: if $J=\{i:R_i\ne0\}$,
then $e_J=1_R$, so any ultrafilter obtained from a product prime
contains $J$. An ultrafilter concentrated on zero factors gives
the zero ultraproduct and contributes no primes. The empty product
is the zero ring and satisfies the same conclusion.
\end{proof}

\begin{lemma}
\label{editorial-source-lemma-characterize-reduced-ring-normal}
Let $R$ be a ring. Assume $R$ is reduced and has finitely many
minimal primes. Then the following are equivalent:
\begin{enumerate}
\item $R$ is a normal ring,
\item $R$ is integrally closed in its total ring of fractions, and
\item $R$ is a finite product of normal domains.
\end{enumerate}
\end{lemma}

\begin{proof}
The implications (1) $\Rightarrow$ (2) and
(3) $\Rightarrow$ (1) hold in general,
see Lemmas \ref{editorial-source-lemma-normal-ring-integrally-closed} and
\ref{editorial-source-lemma-finite-product-normal}.

\medskip\noindent
Let $\mathfrak p_1, \ldots, \mathfrak p_n$ be the minimal primes of $R$.
By Lemmas \ref{algebra-lemma-reduced-ring-sub-product-fields} and
\ref{algebra-lemma-total-ring-fractions-no-embedded-points} we have
$Q(R) = R_{\mathfrak p_1} \times \ldots \times R_{\mathfrak p_n}$, and
by Lemma \ref{algebra-lemma-minimal-prime-reduced-ring} each factor is a field.
Let $e_i = (0, \ldots, 0, 1, 0, \ldots, 0)$ be the $i$th idempotent
of $Q(R)$.

\medskip\noindent
Suppose $R$ is integrally closed in its original $Q(R)$.
Each $e_i$ satisfies $e_i^2-e_i=0$, hence lies in $R$.
The orthogonality identities and $\sum_i e_i=1$ give inverse maps
$$
R\longrightarrow\prod_{i=1}^n e_iR,\quad
r\longmapsto(e_ir)_i,
\qquad (r_i)_i\longmapsto\sum_{i=1}^n r_i.
$$
The factor $e_iR$ has identity $e_i$ and embeds in the original
field $R_{\mathfrak p_i}$ by the $i$th coordinate of $Q(R)$.
The kernel of $R\to e_iR$ is exactly $\mathfrak p_i$.
Indeed the kernel of $R\to R_{\mathfrak p_i}$ is
$\mathfrak p_i$: the localization is a field, whose maximal ideal
$\mathfrak p_iR_{\mathfrak p_i}$ is zero; conversely an equation
$ur=0$ with $u\notin\mathfrak p_i$ forces $r\in\mathfrak p_i$
by primality. Multiplication by $e_i$ keeps precisely this
coordinate, and injectivity of $R\to Q(R)$ identifies its
vanishing with $e_ir=0$ in $R$. Therefore
$R/\mathfrak p_i\cong e_iR$ by $r+\mathfrak p_i\mapsto e_ir$,
with inverse $e_ir\mapsto r+\mathfrak p_i$.
The maps
$$
\operatorname{Frac}(R/\mathfrak p_i)\longrightarrow R_{\mathfrak p_i},
\quad\frac{a+\mathfrak p_i}{b+\mathfrak p_i}
\longmapsto\frac{a/1}{b/1},
$$
and $a/s\mapsto(a+\mathfrak p_i)/(s+\mathfrak p_i)$ are inverse:
the denominators avoid $\mathfrak p_i$, and the original fraction
equality witnesses give equality in the domain $R/\mathfrak p_i$.
This proves the asserted field-of-fractions comparison.

Each subring $R/\mathfrak p_i\subset R_{\mathfrak p_i}$ is
integrally closed, as also follows from
Lemma \ref{editorial-source-lemma-finite-product-integral-closure}.
Directly, if $z_i$ in that field satisfies a monic equation of
degree $d\geq1$ over $e_iR$, form $z\in Q(R)$ with $i$th
coordinate $z_i$ and all others zero. Regard each coefficient
as the element of $R$ supported by $e_i$. The same monic
polynomial vanishes on $z$: at the other coordinates it is
evaluated at zero with all coefficients zero, including its
constant coefficient. Thus $z$ is integral over $R$ and belongs
to $R$ by hypothesis. Its $i$th coordinate lies in $e_iR$,
so $z_i$ belongs to the stated subring. This proves that each
$R/\mathfrak p_i$ is a normal domain and gives (3).
If $R=0$, there are no minimal primes and (3) is the empty
product, so the same conclusion holds directly.
\end{proof}

\begin{lemma}
\label{editorial-source-lemma-colimit-normal-ring}
Let $(R_i,\varphi_{ii'})$ be a system of rings indexed by a nonempty
directed preordered set $I$
(Categories, Definition \ref{categories-definition-directed-system}).
If each $R_i$ is normal, so is
$R = \colim_i R_i$.
\end{lemma}

\begin{proof}
Write $\lambda_i:R_i\to R$ for the original colimit maps.
If $R=0$, the conclusion is vacuous. Otherwise fix a prime
$\mathfrak p\subset R$ and set
$\mathfrak p_i=\lambda_i^{-1}(\mathfrak p)$.
These are proper primes, and
$\varphi_{ij}^{-1}(\mathfrak p_j)=\mathfrak p_i$ because
$\lambda_j\varphi_{ij}=\lambda_i$.
Thus the localized transition maps are defined by
$$
\psi_{ij}:(R_i)_{\mathfrak p_i}\longrightarrow(R_j)_{\mathfrak p_j},
\qquad a_i/s_i\longmapsto
\varphi_{ij}(a_i)/\varphi_{ij}(s_i).
$$
Indeed denominators outside $\mathfrak p_i$ stay outside
$\mathfrak p_j$, and an equality witness
$u_i(a_it_i-b_is_i)=0$ with $u_i\notin\mathfrak p_i$
maps to the corresponding equality witness in $R_j$.
Put $A_i=(R_i)_{\mathfrak p_i}$ and $L=\colim_i A_i$,
with maps $\rho_i:A_i\to L$.

The element rule for a directed colimit says that every element
has a representative at one stage, and equality of two representatives
holds at a common later stage. To recall why, construct the underlying
set from pairs $(i,a_i)$, identifying two pairs when their images
agree at some common stage. Directedness makes this relation
transitive; addition and multiplication are computed at a common
stage and are independent of that stage by the same rule.
This construction satisfies the universal property of the ring
colimit, since every compatible family of maps sends equivalent
pairs to the same value. In particular a zero equality may need
a later stage; no transition map is assumed injective.

We prove the original localization comparison with both maps:
$$
\alpha:L\longrightarrow R_{\mathfrak p},\qquad
\rho_i(a_i/s_i)\longmapsto
\lambda_i(a_i)/\lambda_i(s_i),
$$
and
$$
\beta:R_{\mathfrak p}\longrightarrow L,\qquad
a/s\longmapsto\rho_i(a_i/s_i),
$$
where numerator and denominator are represented at one common
stage $i$. Since $s\notin\mathfrak p$, every such representative
$s_i$ is outside $\mathfrak p_i$. For well-definedness of the
second map, an equality $a/s=b/t$ has a witness
$u\notin\mathfrak p$ with $u(at-bs)=0$ in $R$.
Represent all five elements at one common stage. The equality
rule gives a later stage $j$ with
$u_j(a_jt_j-b_js_j)=0$ in $R_j$.
All three denominators $u_j,s_j,t_j$ are outside
$\mathfrak p_j$, so $a_j/s_j=b_j/t_j$ in $A_j$.
The same argument handles any change of representative.
Addition and multiplication use the ordinary fraction formulas
at a common stage. The first map is defined by the compatible
localized maps to $R_{\mathfrak p}$, whose fraction relations
are respected in exactly the same way.
The composites fix every $a/s$ and every generator
$\rho_i(a_i/s_i)$, respectively. Hence
$$
R_{\mathfrak p}\cong\colim_i(R_i)_{\mathfrak p_i}.
$$

Each $A_i$ is a nonzero normal domain by the hypothesis.
The colimit $L$ is nonzero: otherwise $1=0$ would hold at
some later stage, contradicting this property of $A_i$.
It is a domain even though the transition maps may have kernels.
For $xy=0$ in $L$, represent $x,y$ at one stage and then
pass to a later stage where their product is zero.
The domain at that stage forces one representative to be zero,
so one of the original $x,y$ is zero in $L$.
Thus $A=R_{\mathfrak p}$ is a nonzero domain.

Retain $a,b\in A$, $b\ne0$, and an original monic equation
$$
z^n+\sum_{k=0}^{n-1}c_kz^k=0,
\qquad z=a/b,\quad c_k\in A,\quad n\geq1.
$$
Multiplication by the full denominator power $b^n$ gives
$a^n+\sum_{k=0}^{n-1}c_ka^kb^{n-k}=0$ in $A$.
Choose a common stage $i$ representing each of the finitely
many localized elements as
$$
a=\frac{\lambda_i(a_i)}{\lambda_i(u_i)},\qquad
b=\frac{\lambda_i(b_i)}{\lambda_i(v_i)},\qquad
c_k=\frac{\lambda_i(c_{k,i})}{\lambda_i(w_{k,i})},
$$
with $u_i,v_i,w_{k,i}\notin\mathfrak p_i$.
Keep the products
$$
W_i=\prod_{k=0}^{n-1}w_{k,i},\qquad
W_{k,i}=\prod_{\ell\ne k}w_{\ell,i},\qquad
D_i=u_i^nv_i^nW_i
$$
and the full numerator
$$
H_i=a_i^nv_i^nW_i+
\sum_{k=0}^{n-1}c_{k,i}a_i^kb_i^{n-k}
 u_i^{n-k}v_i^kW_{k,i}.
$$
In $A_i$, without omitting any denominator factor, one has
$$
\frac{H_i}{D_i}
=\left(\frac{a_i}{u_i}\right)^n+
\sum_{k=0}^{n-1}\frac{c_{k,i}}{w_{k,i}}
 \left(\frac{a_i}{u_i}\right)^k
 \left(\frac{b_i}{v_i}\right)^{n-k}.
$$
The image of this fraction in $A$ is zero. Therefore some
$t\in R\setminus\mathfrak p$ satisfies
$t\lambda_i(H_i)=0$ in the original $R$.
Represent $t$ at a stage above $i$, then apply the equality rule
again to obtain a stage $q$ where $t_qH_q=0$ in $R_q$.
Here every entry of $H_q$ is the image of the corresponding
original entry of $H_i$ under the indicated transition map:
$$
H_q=a_q^nv_q^n\prod_{\ell=0}^{n-1}w_{\ell,q}
+\sum_{k=0}^{n-1}c_{k,q}a_q^kb_q^{n-k}u_q^{n-k}v_q^k
 \prod_{\ell\ne k}w_{\ell,q}.
$$
The witness $t_q$ and all the denominator factors are outside
$\mathfrak p_q$. Thus $t_q$ and $D_q$ are units in $A_q$,
and the complete displayed equation for $H_q/D_q$ is zero there.
The element $b_q/v_q$ is nonzero in $A_q$, since its image is
the original nonzero $b$ in $A$. Consequently
$$
z_q=\frac{a_q/u_q}{b_q/v_q}\in\operatorname{Frac}(A_q)
$$
is defined, and division by $(b_q/v_q)^n$ gives its monic equation
$$
z_q^n+\sum_{k=0}^{n-1}(c_{k,q}/w_{k,q})z_q^k=0.
$$
Normality gives $z_q=h_q\in A_q$. Transport the equality
$a_q/u_q=(b_q/v_q)h_q$ in $A_q$ through its actual map to $A$.
It becomes $a=bh$ for the image $h\in A$ of $h_q$.
Since $A$ is a domain and $b\ne0$, the original $z=a/b$ equals
$h$ and belongs to $A$. This proves integral closedness of $A$
and hence normality of $R$.
No map from $\operatorname{Frac}(A_q)$ to
$\operatorname{Frac}(A)$ has been assumed: it need not exist
when the stage map has a kernel. The equality transported is
the displayed equality inside $A_q$ itself.

The divisibility characterization in the arbitrary-product proof
gives a second direct verification. Represent the finitely many
elements of a $(\mathcal D_n)$ equation in $R$ at one stage,
then pass to a later stage where that whole equation holds.
The normal ring at that stage supplies $a_i=b_id_i$;
its image gives the required $a=bd$ in $R$. Thus all
$(\mathcal D_n)$ hold, in agreement with the localization proof.
Neither proof requires antisymmetry of the directed preorder
or injectivity of the transition maps.
\end{proof}










\subsection{Editorial treatment of Algebra lines 8347--8357}

\hypertarget{algebra-context-090}{}

\noindent\href{https://github.com/stacks/stacks-project/blob/a04446e57ec1fbc252a871afcec7752fb2807b14/algebra.tex\#L8347}{Original source, lines 8347--8357}. Complete original and reviewed TeX passages accompany this note. Every intervention is separately recorded; the following treatment is editorial.


\begin{definition}
\label{editorial-source-definition-integral-over-ideal}
Let $\varphi : R \to S$ be a ring map.
Let $I \subset R$ be an ideal.
We say an element $g \in S$ is
{\it integral over $I$} if
there exists a monic
polynomial $P = x^d + \sum_{j=0}^{d-1} a_j x^j$, with $d\geq1$,
with coefficients $a_j \in I^{d-j}$ such
that $P^\varphi(g) = 0$ in $S$.
\end{definition}


\subsection{Editorial treatment of Algebra lines 8365--8396}

\hypertarget{algebra-context-091}{}

\noindent\href{https://github.com/stacks/stacks-project/blob/a04446e57ec1fbc252a871afcec7752fb2807b14/algebra.tex\#L8365}{Original source, lines 8365--8396}. Complete original and reviewed TeX passages accompany this note. Every intervention is separately recorded; the following treatment is editorial.


\begin{lemma}
\label{editorial-source-lemma-characterize-integral-ideal}
Let $\varphi : R \to S$ be a ring map.
Let $I \subset R$ be an ideal.
Let $A = \bigoplus_{h\geq0} I^ht^h \subset R[t]$ be the
subring of the polynomial ring
generated by $R \oplus It \subset R[t]$.
An element $s \in S$ is integral over $I$ if
and only if the element $st \in S[t]$
is integral over $A$, with the coefficient map induced by $\varphi$.
More generally, for every integer $k\geq1$, the element $s$ is
integral over $I^k$ if and only if $st^k$ is integral over $A$.
\end{lemma}

\begin{proof}
Let $\widetilde\varphi:R[t]\to S[t]$ send
$\sum_h r_ht^h$ to $\sum_h\varphi(r_h)t^h$, and use its
restriction to $A$. Its additivity and multiplicativity follow
coefficient by coefficient from those of $\varphi$. The displayed
description of $A$ is the subring generated by $R\oplus It$:
products of $h$ elements of $It$ generate $I^ht^h$ additively,
and products of its homogeneous summands have the required powers
$I^hI^\ell=I^{h+\ell}$. Here $I^0=R$.

We prove the assertion for every $k\geq1$, retaining $k=1$ as
the first assertion. Suppose $st^k$ is integral over $A$, and
keep an original monic equation
$$
P=x^d+\sum_{j=0}^{d-1}a_jx^j,\qquad
a_j=\sum_{h\geq0}a_{jh}t^h,\quad a_{jh}\in I^h,
\qquad P^{\widetilde\varphi}(st^k)=0.
$$
Each coefficient sum has finite support. The full equation is
$$
s^dt^{kd}+\sum_{j=0}^{d-1}\sum_{h\geq0}
  \varphi(a_{jh})s^jt^{h+kj}=0.
$$
For every $\ell\geq0$ its coefficient of $t^\ell$ is
$$
\delta_{\ell,kd}s^d+
\sum_{\substack{0\leq j<d,\ h\geq0\\h+kj=\ell}}
  \varphi(a_{jh})s^j=0,
$$
where $\delta_{\ell,kd}$ equals $1$ for $\ell=kd$ and $0$
otherwise. Thus none of the other coefficient equations is discarded.
Taking $\ell=kd$ gives
$$
s^d+\sum_{j=0}^{d-1}\varphi(a_{j,k(d-j)})s^j=0,
\qquad a_{j,k(d-j)}\in I^{k(d-j)}=(I^k)^{d-j}.
$$
This proves integrality over $I^k$. In the original $k=1$
notation, $a_j''=a_{j,d-j}$ and $a_j'=a_j''t^{d-j}$,
so the degree-$d$ equation is exactly
$(st)^d+\sum_{j=0}^{d-1}\varphi(a_j'')t^{d-j}(st)^j=0$.

Conversely retain a monic $P=x^d+\sum_{j=0}^{d-1}a_jx^j$
with $a_j\in I^{k(d-j)}$ and $P^\varphi(s)=0$.
Every coefficient of
$$
Q=x^d+\sum_{j=0}^{d-1}(a_jt^{k(d-j)})x^j
$$
belongs to $A$, and the exact equation
$$
Q^{\widetilde\varphi}(st^k)
=t^{kd}\left(s^d+\sum_{j=0}^{d-1}\varphi(a_j)s^j\right)=0
$$
proves the converse, without cancellation of $t$ or any
injectivity assumption on $\varphi$.
\end{proof}


\subsection{Editorial treatment of Algebra lines 8398--8419}

\hypertarget{algebra-context-092}{}

\noindent\href{https://github.com/stacks/stacks-project/blob/a04446e57ec1fbc252a871afcec7752fb2807b14/algebra.tex\#L8398}{Original source, lines 8398--8419}. Complete original and reviewed TeX passages accompany this note. Every intervention is separately recorded; the following treatment is editorial.


\begin{lemma}
\label{editorial-source-lemma-integral-over-ideal-is-submodule}
Let $\varphi : R \to S$ be a ring map.
Let $I \subset R$ be an ideal.
The set of elements of $S$ which are integral
over $I$ forms an $R$-submodule of $S$.
Furthermore, if $s \in S$ is integral over
$R$, and $s'$ is integral over $I$, then
$ss'$ is integral over $I$.
\end{lemma}

\begin{proof}
Use $A$ and its map to $S[t]$ from
Lemma \ref{editorial-source-lemma-characterize-integral-ideal}. The elements of
$S[t]$ integral over $A$ form a subring by
Lemma \ref{editorial-source-lemma-integral-closure-is-ring}. If $u,v\in S$
are integral over $I$, the elements $ut$ and $vt$ belong to
that subring, so do $(u+v)t=ut+vt$ and $(-u)t=-ut$.
The same characterization gives closure under addition and negatives;
$0$ is integral over $I$ by the monic equation $x=0$.

If $s\in S$ is integral over $R$, its original monic equation
over $R$ is also a monic equation for the constant polynomial
$s\in S[t]$ over $A$, since the degree-zero part of $A$ is
$R$. For $s'$ integral over $I$, the element $s't$ is integral
over $A$. Their product $s(s't)=(ss')t$ is therefore integral
over $A$, and the characterization gives $ss'$ integral over
$I$, which proves the second assertion. Finally every
$\varphi(r)$ for $r\in R$ is integral over $R$ by the equation
$x-r$, so this product assertion supplies multiplication by
every original $R$-scalar and proves the submodule assertion.
\end{proof}


\subsection{Editorial treatment of Algebra lines 8421--8430}

\hypertarget{algebra-context-093}{}

\noindent\href{https://github.com/stacks/stacks-project/blob/a04446e57ec1fbc252a871afcec7752fb2807b14/algebra.tex\#L8421}{Original source, lines 8421--8430}. Complete original and reviewed TeX passages accompany this note. Every intervention is separately recorded; the following treatment is editorial.


\begin{lemma}
\label{editorial-source-lemma-integral-integral-over-ideal}
Suppose $\varphi : R \to S$ is integral.
Suppose $I \subset R$ is an ideal.
Then every element of $IS$ is integral over $I$.
\end{lemma}

\begin{proof}
An element of $IS$ is a finite sum
$\sum_{i=1}^r\varphi(a_i)s_i$ with $a_i\in I$ and $s_i\in S$.
Each $\varphi(a_i)$ is integral over $I$ by the degree-one
equation $x-a_i$. Every $s_i$ is integral over $R$ by the
hypothesis on $\varphi$. The product assertion and then addition
in Lemma \ref{editorial-source-lemma-integral-over-ideal-is-submodule} give
integrality of the entire sum over $I$. The empty sum is $0$,
already covered by its degree-one equation.
\end{proof}

\begin{lemma}
\label{editorial-source-lemma-normal-principal-ideals-integrally-closed}
A ring $R$ is normal if and only if every principal ideal
$bR$, for $b\in R$, is integrally closed in $R$ in the sense
of Definition \ref{editorial-source-definition-integral-over-ideal}.
This includes $b=0$ and zero divisors $b$.
\end{lemma}

\begin{proof}
For every integer $e\geq1$, the equality $(bR)^e=b^eR$
follows in both directions by multiplying generators and by
writing $b^er$ as a product of $e$ elements of $bR$.
Consequently an element $a\in R$ is integral over $bR$ exactly
when, for some $n\geq1$ and $c_0,\ldots,c_{n-1}\in R$,
$$
a^n+\sum_{j=0}^{n-1}c_ja^jb^{n-j}=0.
$$
Indeed each original coefficient of a defining equation lies in
$(bR)^{n-j}=b^{n-j}R$ and can be written $c_jb^{n-j}$;
conversely these displayed coefficients have the required ideal
powers. This uses no cancellation or uniqueness of $c_j$.
Also every $a\in bR$ is integral over $bR$ by $x-a$.
Thus every principal ideal is integrally closed exactly when
every displayed equation implies $a=bd$ for some $d\in R$.
This is precisely the family of conditions $(\mathcal D_n)$,
for all $n\geq1$, proved equivalent to normality in the proof of
Lemma \ref{editorial-source-lemma-finite-product-normal}. Both implications of
that characterization hold for arbitrary rings, including the
zero ring and the cases where $b$ is a zero divisor; applying
them proves the assertion.
\end{proof}


\subsection{Editorial treatment of Algebra lines 8432--8496}

\hypertarget{algebra-context-094}{}

\noindent\href{https://github.com/stacks/stacks-project/blob/a04446e57ec1fbc252a871afcec7752fb2807b14/algebra.tex\#L8432}{Original source, lines 8432--8496}. Complete original and reviewed TeX passages accompany this note. Every intervention is separately recorded; the following treatment is editorial.


\begin{lemma}
\label{editorial-source-lemma-polynomials-divide}
Let $K$ be a field. Let $n, m \in \mathbf{N}$ and
$a_0, \ldots, a_{n - 1}, b_0, \ldots, b_{m - 1} \in K$.
If the polynomial $x^n + a_{n - 1}x^{n - 1} + \ldots + a_0$
divides the polynomial $x^m + b_{m - 1} x^{m - 1} + \ldots + b_0$
in $K[x]$ then
\begin{enumerate}
\item $a_0, \ldots, a_{n - 1}$ are integral over any subring
$R_0$ of $K$ containing the elements $b_0, \ldots, b_{m - 1}$, and
\item each $a_i$ lies in $\sqrt{(b_0, \ldots, b_{m-1})R}$
for any subring $R \subset K$ containing the elements
$a_0, \ldots, a_{n - 1}, b_0, \ldots, b_{m - 1}$.
\end{enumerate}
\end{lemma}

\begin{proof}
Write
$$
Q=x^n+\sum_{j=0}^{n-1}a_jx^j,\qquad
P=x^m+\sum_{i=0}^{m-1}b_ix^i.
$$
Since both are monic and $Q$ divides $P$ over the field $K$,
we have $m\geq n$. Choose a splitting field $L/K$ and retain
the roots with their multiplicities:
$$
P=\prod_{i=1}^m(x-\beta_i),\qquad \beta_i\in L.
$$
See Fields, Section \ref{fields-section-splitting-fieds}.
Unique factorization in $L[x]$ gives an indexed subset
$J\subseteq\{1,\ldots,m\}$ of cardinality $n$ with
$Q=\prod_{i\in J}(x-\beta_i)$. Equal roots remain distinct
indices in this expression. Every $\beta_i$ is integral over
$R_0$ by the original $P$, and every original coefficient is
$$
a_j=(-1)^{n-j}
\sum_{\substack{E\subseteq J\\ |E|=n-j}}
\prod_{i\in E}\beta_i\qquad(0\leq j<n).
$$
The subring property in Lemma \ref{editorial-source-lemma-integral-closure-is-ring}
proves (1), retaining every sign and repeated-root contribution.

For (2), first consider $m=n$. The monic quotient has degree
zero and is $1$, so $a_i=b_i$ for every $i$, which gives the
assertion directly. Now let $e=m-n>0$ and write the original
monic quotient as
$$
C=x^e+\sum_{k=0}^{e-1}c_kx^k,\qquad P=QC.
$$
Put $a_n=c_e=1$. Comparing the coefficient of $x^{n+k}$
for $k=e-1,e-2,\ldots,0$ gives the exact recurrence
$$
c_k=b_{n+k}-
\sum_{i=\max(0,n+k-e)}^{n-1}a_i c_{n+k-i}.
$$
Every index $n+k-i$ in this sum is between $k+1$ and $e$.
Thus descending induction shows $c_k\in R$ for all $k$,
without negative coefficient indices or assumptions on $n>1$.
Set $H=\sqrt{(b_0,\ldots,b_{m-1})R}$ and
$\overline R=R/H$. This quotient is reduced. If it is the
zero ring, the required membership in $H=R$ already holds.
Otherwise the complete coefficient image of $P=QC$ is
$$
x^m+\sum_{i=0}^{m-1}\overline b_i x^i=x^m
=\left(x^n+\sum_{j=0}^{n-1}\overline a_jx^j\right)
 \left(x^e+\sum_{k=0}^{e-1}\overline c_kx^k\right).
$$
At a minimal prime $\mathfrak p$ of $\overline R$, the ring
$\overline R_{\mathfrak p}$ is a field by
Lemma \ref{algebra-lemma-minimal-prime-reduced-ring}. In its polynomial
ring unique factorization makes the first monic factor
$f=x^n+\sum_{j=0}^{n-1}\overline a_jx^j$ equal to $x^n$:
every irreducible factor divides $x$, its degree is $n$, and
its leading coefficient fixes the unit factor to be $1$.
For each $j$ choose $s_j\notin\mathfrak p$ with
$s_j\overline a_j=0$, using the equality in this localization,
and take $s=\prod_{j=0}^{n-1}s_j$.
Then there exists an $s\in\overline R$, $s\notin\mathfrak p$,
such that $s(f-x^n)=0$, with this explicit common choice.
Primality gives $\overline a_j\in\mathfrak p$ for every $j$.
As this holds at every minimal prime and their intersection
is zero in the reduced ring, every $\overline a_j$ is zero
(Lemma \ref{algebra-lemma-reduced-ring-sub-product-fields}).
Hence every $a_j$ belongs to $H$, proving (2).
\end{proof}

\begin{lemma}
\label{editorial-source-lemma-weighted-monic-factor-coefficients}
Let $R\subseteq K$ be a subring of a field and let $I\subseteq R$
be an ideal. Suppose
$$
P=x^m+\sum_{i=0}^{m-1}b_ix^i,\qquad b_i\in I^{m-i},
\qquad Q=x^n+\sum_{j=0}^{n-1}a_jx^j\in K[x]
$$
are monic, $1\leq n\leq m$, and $Q$ divides $P$.
For each $0\leq j<n$, the coefficient $a_j$ is integral over
$I^{n-j}$. More precisely, with $N=\binom{m}{n}$, it satisfies
an equation
$$
a_j^N+\sum_{\ell=0}^{N-1}C_{j,\ell}a_j^\ell=0,
\qquad C_{j,\ell}\in I^{(n-j)(N-\ell)}.
$$
\end{lemma}

\begin{proof}
In a splitting field $L/K$ retain the original indexed roots,
including multiplicities, so
$P=\prod_{r=1}^m(x-\beta_r)$. As in the proof of
Lemma \ref{editorial-source-lemma-polynomials-divide}, there is an indexed subset
$J\subseteq\{1,\ldots,m\}$ of cardinality $n$ with
$Q=\prod_{r\in J}(x-\beta_r)$. Fix $j$ and put $w=n-j$.
Let $\mathcal T$ be the set of all $n$-element subsets of
$\{1,\ldots,m\}$, so $|\mathcal T|=N$. For each $T\in\mathcal T$
define the full coefficient
$$
r_{j,T}=(-1)^w
\sum_{\substack{U\subseteq T\\|U|=w}}\prod_{r\in U}\beta_r.
$$
In particular $a_j=r_{j,J}$. Form
$$
H_j(Y)=\prod_{T\in\mathcal T}(Y-r_{j,T})
=Y^N+\sum_{\ell=0}^{N-1}C_{j,\ell}Y^\ell,
$$
where every coefficient, sign and repeated contribution is
$$
C_{j,\ell}=(-1)^{N-\ell}
\sum_{\substack{\mathcal A\subseteq\mathcal T\\|\mathcal A|=N-\ell}}
\prod_{T\in\mathcal A}r_{j,T}.
$$
Then $H_j(a_j)=0$.

Replace the $\beta_r$ temporarily by independent variables
$Z_r$. The same expression for $C_{j,\ell}$ is an integer
polynomial, symmetric in the $m$ variables and homogeneous of
total degree $D=w(N-\ell)$. It is an integer polynomial in
the elementary symmetric polynomials $e_1,\ldots,e_m$, each
monomial having weighted degree $D$ when $e_s$ has weight $s$.
Here is a proof retaining the weights. In lexicographic order
$Z_1>\cdots>Z_m$, the leading exponent vector
$(\lambda_1,\ldots,\lambda_m)$ of a nonzero symmetric
homogeneous polynomial satisfies
$\lambda_1\geq\cdots\geq\lambda_m$: an inversion would give
a larger monomial after interchanging the corresponding
variables, contradicting maximality. With $\lambda_{m+1}=0$,
the polynomial
$$
\prod_{s=1}^m e_s^{\lambda_s-\lambda_{s+1}}
$$
has that same leading monomial with coefficient $1$ and
weighted degree
$\sum_{s=1}^m s(\lambda_s-\lambda_{s+1})
=\sum_{s=1}^m\lambda_s=D$.
Subtract the integer multiple matching the leading coefficient.
The result remains symmetric and homogeneous of degree $D$,
with a smaller leading monomial. Only finitely many monomials
have total degree $D$, so this procedure terminates. It proves
an expression with integers $\kappa_{j,\ell,\nu}$:
$$
C_{j,\ell}
=\sum_{\substack{\nu\in\mathbf Z_{\geq0}^m\\
                  \sum_{s=1}^m s\nu_s=D}}
\kappa_{j,\ell,\nu}
\prod_{s=1}^m e_s(\beta_1,\ldots,\beta_m)^{\nu_s}.
$$
The original coefficients of $P$ give
$e_s(\beta_1,\ldots,\beta_m)=(-1)^s b_{m-s}$, so the
complete substituted expression is
$$
C_{j,\ell}
=\sum_{\substack{\nu\in\mathbf Z_{\geq0}^m\\
                  \sum_{s=1}^m s\nu_s=D}}
\kappa_{j,\ell,\nu}
\prod_{s=1}^m\bigl((-1)^s b_{m-s}\bigr)^{\nu_s}.
$$
Each product belongs to
$\prod_s(I^s)^{\nu_s}=I^D=I^{w(N-\ell)}$.
This proves the required equation over $R$. All identities
were proved over the integers before evaluation, so no
separability or characteristic restriction is present.

The exact Rees comparison gives a second proof of the
integrality conclusion. Put
$A=\bigoplus_{h\geq0}I^ht^h\subseteq R[t]$ and map it to
$L[t]$ by the given inclusion. Every $t\beta_r$ satisfies
$$
(t\beta_r)^m+
\sum_{i=0}^{m-1}(b_it^{m-i})(t\beta_r)^i=t^mP(\beta_r)=0.
$$
Its coefficients belong to $A$, so all these elements are
integral over $A$. The original signed identity
$$
a_jt^{n-j}=(-1)^{n-j}
\sum_{\substack{U\subseteq J\\|U|=n-j}}
\prod_{r\in U}(t\beta_r)
$$
and Lemma \ref{editorial-source-lemma-integral-closure-is-ring} show that this
element is integral over $A$. The weight-$n-j$ assertion of
Lemma \ref{editorial-source-lemma-characterize-integral-ideal} gives integrality
of the original $a_j$ over $I^{n-j}$.
\end{proof}

\begin{lemma}
\label{editorial-source-lemma-integral-ideal-minimal-polynomial-criterion}
Let $R$ be a domain, $K=\operatorname{Frac}(R)$, $I\subseteq R$
an ideal, and $g$ an algebraic element in a field extension
$L/K$. Write its monic minimal polynomial as
$Q=x^n+\sum_{j=0}^{n-1}a_jx^j\in K[x]$.
Then $g$ is integral over $I$ if and only if every $a_j$ is
integral over $I^{n-j}$. If $R$ is normal, this is equivalent
to $a_j\in\overline{I^{n-j}}\subseteq R$ for every $j$,
where the bar denotes integral closure of the ideal inside $R$.
\end{lemma}

\begin{proof}
If $g$ is integral over $I$, retain a witnessing equation
$P(g)=0$ with $P=x^m+\sum_{i=0}^{m-1}b_ix^i$ and
$b_i\in I^{m-i}$. Division by $Q$ in $K[x]$ gives a
remainder of degree less than $n$ that vanishes at $g$, hence
is zero by minimality. Thus $Q\mid P$, and
Lemma \ref{editorial-source-lemma-weighted-monic-factor-coefficients} gives
the asserted coefficient integrality.

Conversely suppose every coefficient has that integrality.
Use $A=\bigoplus_{h\geq0}I^ht^h\subseteq L[t]$ with the
original coefficient inclusion. By the weighted Rees
characterization every $h_j=a_jt^{n-j}$ is integral over $A$.
The algebra $B=A[h_0,\ldots,h_{n-1}]$ is finite over $A$:
if the chosen monic equation for $h_j$ has degree $d_j$,
the products $\prod_jh_j^{e_j}$ with $0\leq e_j<d_j$
generate it as an $A$-module, using those original monic
equations to reduce each higher power. The element $gt$
satisfies the full monic equation
$$
(gt)^n+\sum_{j=0}^{n-1}h_j(gt)^j=t^nQ(g)=0.
$$
Thus $B[gt]$ is generated over $A$ by all
$\bigl(\prod_jh_j^{e_j}\bigr)(gt)^r$ with $0\leq r<n$.
To see directly that this gives integrality over $A$, list
these generators as $v_1,\ldots,v_M$, including $1$.
Multiplication by $gt$ has relations
$(gt)v_i=\sum_k d_{ik}v_k$ with $d_{ik}\in A$.
Multiplying the vector equation by the adjugate of
$XI-(d_{ik})$ at $X=gt$ shows that the monic polynomial
$\det(XI-(d_{ik}))$ annihilates every $v_i$, including $1$.
It is therefore an integral equation for $gt$ over $A$.
The weight-one Rees characterization makes $g$ integral over
$I$, proving the converse. This argument does not require
the displayed module generators to be linearly independent.

If $R$ is normal, each integral equation for $a_j$ over an
ideal is in particular an integral equation over $R$.
Since $a_j\in K$, integral closedness of the normal domain
gives $a_j\in R$. The definition of integral closure of
$I^{n-j}$ inside $R$ now gives exactly the stated equivalence.
\end{proof}

\begin{example}
\label{editorial-source-example-weighted-coefficients-need-integral-closure}
Normality does not place these coefficients in the ideal
powers themselves. Let $k$ be any field,
$R=k[u,v]$, $K=k(u,v)$ and $I=(u^2,v^2)$.
The polynomial ring $R$ is a normal domain by
Lemma \ref{editorial-source-lemma-polynomial-domain-normal}, applied twice to $k$.
Choose a root $g$ of $x^2-u^3v$ in a field extension of $K$.
This quadratic is irreducible: in $k(u)[v]$ the exponent of
the factor $v$ defines an integer order on nonzero rational
functions by numerator order minus denominator order.
Multiplication adds this order, as follows by removing the
largest powers of $v$ and observing that the remaining
polynomials have nonzero values at $v=0$. Thus every nonzero
square in $K$ has even order, whereas $u^3v$ has order $1$.
It cannot be a square, so the quadratic has no root in $K$
and is irreducible, including in characteristic two.
Consequently $Q=x^2-u^3v$ is the minimal polynomial of $g$,
and $S=R[g]$ in the chosen extension field is a domain.
Yet the original equation
$$
g^4-u^6v^2=0,\qquad u^6v^2=(u^2)^3v^2\in I^4
$$
makes $g$ integral over $I$; its other three coefficients
are zero and belong to the required powers. The constant
coefficient $a_0=-u^3v$ of $Q$ is not in
$I^2=(u^4,u^2v^2,v^4)$: the monomial $u^3v$ is divisible
by none of these generators, and every polynomial multiple
of a generator is a sum of divisible monomials. Its distinct
monomial coefficient therefore cannot be obtained in their
sum. However
$$
a_0^2-u^6v^2=0,\qquad -u^6v^2\in I^4=(I^2)^2,
$$
so $a_0\in\overline{I^2}\setminus I^2$.
Even weight one need not give actual membership: $h=uv\in R$
satisfies $h^2-u^2v^2=0$ with $u^2v^2\in I^2$, but the
constant coefficient $-uv$ of its minimal polynomial $x-uv$
is outside $I$ by the same monomial argument.
\end{example}

\begin{lemma}
\label{editorial-source-lemma-integrally-closed-powers-monic-factors}
For a normal domain $R$, $K=\operatorname{Frac}(R)$ and an
ideal $I\subseteq R$, the following are equivalent:
\begin{enumerate}
\item Every positive power $I^w$ is integrally closed in $R$.
\item For every monic $P=x^m+\sum_{i=0}^{m-1}b_ix^i$ with
$b_i\in I^{m-i}$, every monic divisor
$Q=x^n+\sum_{j=0}^{n-1}a_jx^j\in K[x]$ of positive degree
has $a_j\in I^{n-j}$ for every $j$.
\item Every element integral over $I$ in a field extension of
$K$ has minimal-polynomial coefficients $a_j\in I^{n-j}$,
where $n$ is its minimal-polynomial degree.
\end{enumerate}
\end{lemma}

\begin{proof}
The weighted coefficient lemma gives (1)$\Rightarrow$(2):
its equations first give $a_j\in R$ by normality and then
$a_j\in I^{n-j}$ by integral closedness of that power.
Minimal-polynomial division as above gives (2)$\Rightarrow$(3).
For (3)$\Rightarrow$(2), factor the given $Q$ in $K[x]$
into monic irreducible factors, repeated with their original
multiplicities. A root of each factor in a splitting field
is a root of $P$, hence is integral over $I$. Condition (3)
gives the required powers for each irreducible factor.
Write the factors as
$Q_r=\sum_{j=0}^{d_r}c_{r,j}x^j$ for $1\leq r\leq q$,
with $c_{r,d_r}=1\in I^0$ and
$c_{r,j}\in I^{d_r-j}$ for $j<d_r$.
The full coefficient of $x^j$ in their product is
$$
\sum_{\substack{0\leq j_r\leq d_r\ (1\leq r\leq q)\\
                 j_1+\cdots+j_q=j}}
\prod_{r=1}^q c_{r,j_r}
\ \in I^{\sum_{r=1}^qd_r-j}=I^{n-j}.
$$
This retains the multiplicities and proves (2).

Finally assume (2), fix $w\geq1$, and take
$c\in\overline{I^w}\subseteq R$. Keep its defining equation
$$
c^d+\sum_{\ell=0}^{d-1}r_\ell c^\ell=0,
\qquad r_\ell\in I^{w(d-\ell)}.
$$
Set $F(Y)=Y^d+\sum_{\ell=0}^{d-1}r_\ell Y^\ell$.
The complete difference identity is
$$
F(Y)-F(c)=(Y-c)\left(
\sum_{k=0}^{d-1}Y^{d-1-k}c^k+
\sum_{\ell=1}^{d-1}r_\ell
   \sum_{k=0}^{\ell-1}Y^{\ell-1-k}c^k\right).
$$
Since $F(c)=0$, substituting $Y=x^w$ proves that $x^w-c$
divides
$P=x^{wd}+\sum_{\ell=0}^{d-1}r_\ell x^{w\ell}$.
The coefficient at $x^{w\ell}$ belongs to $I^{wd-w\ell}$;
every omitted coefficient is zero and is retained in its
required power. Condition (2) applied to this divisor gives
$-c\in I^w$, hence $c\in I^w$. The reverse inclusion
$I^w\subseteq\overline{I^w}$ follows from the degree-one
equation for each of its elements. This proves (1).
\end{proof}

\begin{remark}
\label{editorial-source-remark-weighted-coefficient-consequences}
Under the normal-domain hypotheses, the original coefficient
$a_j\in\overline{I^{n-j}}$ lies in $\sqrt I$: reduce its
defining equation modulo $\sqrt I$, where all nonleading
coefficients vanish; its positive leading power is zero, so
it belongs to that radical. If $I$ is integrally closed,
the same defining coefficients belong to $I^{N-\ell}$
because $(n-j)(N-\ell)\geq N-\ell$, and therefore $a_j$
is integral over $I$ and belongs to $I$. If $I$ is radical,
the radical conclusion already gives membership in $I$.
These claims do not assert membership in $I^{n-j}$ without
the corresponding integral-closedness condition.

The exponent $n-j$ cannot be increased in general.
For $R=k[u]$, $I=(u)$ and any $w\geq1$, take
$P=Q=x^w-u^w$. Its constant coefficient $a_0=-u^w$
belongs to $I^w$. If it were integral over $I^{w+1}$,
an equation of degree $d\geq1$ would have leading term
$(-u^w)^d$ of $u$-order $wd$, whereas each other term
$c_\ell(-u^w)^\ell$ would have order at least
$(w+1)(d-\ell)+w\ell=wd+d-\ell>wd$.
The coefficient of $u^{wd}$ in their sum would remain
$(-1)^d\ne0$, a contradiction in every characteristic.
\end{remark}


\subsection{Editorial treatment of Algebra lines 8498--8515}

\hypertarget{algebra-context-095}{}

\noindent\href{https://github.com/stacks/stacks-project/blob/a04446e57ec1fbc252a871afcec7752fb2807b14/algebra.tex\#L8498}{Original source, lines 8498--8515}. Complete original and reviewed TeX passages accompany this note. Every intervention is separately recorded; the following treatment is editorial.


\begin{lemma}
\label{editorial-source-lemma-minimal-polynomial-normal-domain}
Let $R \subset S$ be an inclusion of domains.
Assume $R$ is normal. Let $g \in S$ be integral
over $R$. Then the minimal polynomial of $g$
has coefficients in $R$.
\end{lemma}

\begin{proof}
The original inclusion $R\subset S$ of domains gives an
injective map $K=\operatorname{Frac}(R)\to\operatorname{Frac}(S)$,
sending $a/b$ to the same fraction; denominators stay nonzero,
and equality of fractions is preserved in both directions by
cross multiplication. Regard $g$ in this fixed extension field.
Choose its original monic equation
$P=x^m+\sum_{i=0}^{m-1}b_ix^i\in R[x]$, with $P(g)=0$,
and its monic minimal polynomial
$Q=x^n+\sum_{j=0}^{n-1}a_jx^j\in K[x]$.
Division in $K[x]$ gives $P=QC+D$ with $\deg D<n$ or $D=0$.
Evaluation at $g$ gives $D(g)=0$, so minimality forces $D=0$.
Part (1) of Lemma \ref{editorial-source-lemma-polynomials-divide} makes every
$a_j$ integral over $R$. Each already belongs to the original
$K$, so normality gives $a_j\in R$.
\end{proof}


\subsection{Editorial treatment of Algebra lines 8517--8572}

\hypertarget{algebra-context-096}{}

\noindent\href{https://github.com/stacks/stacks-project/blob/a04446e57ec1fbc252a871afcec7752fb2807b14/algebra.tex\#L8517}{Original source, lines 8517--8572}. Complete original and reviewed TeX passages accompany this note. Every intervention is separately recorded; the following treatment is editorial.


\begin{proposition}
\label{editorial-source-proposition-going-down-normal-integral}
Let $R \subset S$ be an inclusion of domains.
Assume $R$ is normal and $S$ integral over $R$.
Let $\mathfrak p \subset \mathfrak p' \subset R$
be primes. Let $\mathfrak q'$ be a prime of $S$
with $\mathfrak p' = R \cap \mathfrak q'$.
Then there exists a prime $\mathfrak q$
with $\mathfrak q \subset \mathfrak q'$
such that $\mathfrak p = R \cap \mathfrak q$. In other words:
the going down property holds for $R \to S$, see
Definition \ref{algebra-definition-going-up-down}.
\end{proposition}

\begin{proof}
Keep the stated primes and the original inclusion, and put
$B=S_{\mathfrak q'}$. It suffices to prove
$\mathfrak pB\cap R=\mathfrak p$, as in Lemma \ref{editorial-source-lemma-in-image}.
Here is the precise prime construction supplied by that equality.
The ring
$$
F=(R\setminus\mathfrak p)^{-1}(B/\mathfrak pB)
$$
is nonzero: if $1=0$ in $F$, some $r\in R\setminus\mathfrak p$
has zero image in $B/\mathfrak pB$, contradicting the equality.
A prime of $F$ pulls back through the quotient and localization
maps to a prime of $B$ containing $\mathfrak pB$ and avoiding
the image of $R\setminus\mathfrak p$. Its inverse image
$\mathfrak q$ under $S\to B$ has contraction $\mathfrak p$
in $R$. Every element of $S\setminus\mathfrak q'$ is a unit
in $B$, so none belongs to $\mathfrak q$ and
$\mathfrak q\subseteq\mathfrak q'$. This proves the desired
conclusion once the equality is established.

The inclusion $\mathfrak p\subseteq\mathfrak pB\cap R$ is
immediate. For the reverse inclusion choose an original
$z\in R$ whose image belongs to $\mathfrak pB$. Write that
image as a finite sum $\sum_{i=1}^r r_i h_i/g_i$ with
$r_i\in\mathfrak p$, $h_i\in S$ and $g_i\notin\mathfrak q'$.
Take
$$
g=\prod_{i=1}^r g_i,\qquad
y=\sum_{i=1}^r r_i h_i\prod_{k\ne i}g_k\in\mathfrak pS.
$$
Then $g\notin\mathfrak q'$ and $z=y/g$ in $B$. Since $S$
is a domain, the map $S\to B$ is injective, and the exact
identity $zg=y$ holds in $S$ itself. For the empty sum use
$g=1$ and $y=0$.

\medskip\noindent
By Lemma \ref{editorial-source-lemma-integral-integral-over-ideal}
there exists a monic polynomial
$P = x^m + b_{m-1} x^{m-1} + \ldots + b_0$
with $b_i\in\mathfrak p^{m-i}\subseteq\mathfrak p$
for every $0\leq i<m$, such that $P(y)=0$.
These are the full coefficient powers in the original definition
of integrality over $\mathfrak p$.

\medskip\noindent
By Lemma \ref{editorial-source-lemma-minimal-polynomial-normal-domain}, the
minimal polynomial of $g$ over $K=\operatorname{Frac}(R)$ has
coefficients in $R$. Write it as
$Q=x^n+\sum_{j=0}^{n-1}a_jx^j$. If every $a_j$ belonged to
$\mathfrak p$, its original equation would give
$$
g^n=-\sum_{j=0}^{n-1}a_jg^j\in\mathfrak pS
\subseteq\mathfrak p'S\subseteq\mathfrak q'.
$$
Since $n\geq1$ and $\mathfrak q'$ is prime, this would give
$g\in\mathfrak q'$, a contradiction. Therefore at least one
of the original $a_j$ is outside $\mathfrak p$.

\medskip\noindent
If $z=0$, the desired membership $z\in\mathfrak p$ already
holds. Hence assume $z\ne0$, so $z$ is a unit in the original
$K$, and retain $y=zg$. The coefficient-preserving substitution
map $K[x]\to K[x]$, $H(x)\mapsto H(zx)$, has inverse
$H(x)\mapsto H(x/z)$. Set
$$
Q'(x)=z^nQ(x/z)
=x^n+\sum_{j=0}^{n-1}z^{n-j}a_jx^j.
$$
This is monic of degree $n$ and satisfies
$Q'(y)=z^nQ(g)=0$. If $H(y)=0$ for $H\in K[x]$, then
$H(zg)=0$, so minimality of $Q$ gives $H(zx)=Q(x)D(x)$
for some $D\in K[x]$. The inverse substitution yields
$$
H(x)=Q(x/z)D(x/z)=Q'(x)\bigl(z^{-n}D(x/z)\bigr).
$$
Thus $Q'$ divides every annihilating polynomial of $y$ and is
its monic minimal polynomial; in particular $Q'$ divides the
original $P$. All coefficients of $P$ and $Q'$ belong to $R$.
Part (2) of Lemma \ref{editorial-source-lemma-polynomials-divide} gives
$$
z^ja_{n-j}\in\sqrt{(b_0,\ldots,b_{m-1})R}
\subseteq\mathfrak p,\qquad 1\leq j\leq n.
$$
The full powers $b_i\in\mathfrak p^{m-i}$ also yield the
stronger statement that each $z^ja_{n-j}$ is integral over
$\mathfrak p^j$, by
Lemma \ref{editorial-source-lemma-weighted-monic-factor-coefficients}.
Choose $a_{n-j}\notin\mathfrak p$ from the previous paragraph.
Primality first gives $z^j\in\mathfrak p$ and then
$z\in\mathfrak p$, since $j\geq1$. The required contraction
equality, and therefore the prime constructed above, follow.
\end{proof}


\subsection{Editorial treatment of Algebra lines 8587--9307}

\hypertarget{algebra-context-097}{}

\noindent\href{https://github.com/stacks/stacks-project/blob/a04446e57ec1fbc252a871afcec7752fb2807b14/algebra.tex\#L8587}{Original source, lines 8587--9307}. Complete original and reviewed TeX passages accompany this note. Every intervention is separately recorded; the following treatment is editorial.


\subsubsection{Flat modules and flat ring maps}
\label{editorial-source-section-flat}

\noindent
One often used result is that if $M = \colim_{i\in \mathcal{I}} M_i$
is a colimit of $R$-modules and if $N$ is an $R$-module then
$$
M \otimes N
=
\colim_{i\in \mathcal{I}} M_i \otimes_R N,
$$
see Lemma \ref{editorial-source-lemma-tensor-products-commute-with-limits}.
This property is usually expressed by saying
that {\it $\otimes$ commutes with colimits}.
Another often used result is that if $0 \to N_1 \to N_2 \to N_3 \to 0$
is an exact sequence and if $M$ is any $R$-module, then
$$
M \otimes_R N_1
\to
M \otimes_R N_2
\to
M \otimes_R N_3
\to
0
$$
is still exact, see Lemma \ref{algebra-lemma-tensor-product-exact}.
Both of these properties tell us that the functor
$N \mapsto M \otimes_R N$ {\it is right exact}.
See Categories, Section \ref{categories-section-exact-functor}
and Homology, Section \ref{homology-section-functors}.
An $R$-module $M$ is flat if $N \mapsto N \otimes_R M$ is also left exact,
i.e., if it is exact. Here is the precise definition.

\begin{definition}
\label{editorial-source-definition-flat}
Let $R$ be a ring.
\begin{enumerate}
\item An $R$-module $M$ is called {\it flat} if whenever
$N_1 \to N_2 \to N_3$ is an exact sequence of $R$-modules
the sequence $M \otimes_R N_1 \to M \otimes_R N_2 \to M \otimes_R N_3$
is exact as well.
\item An $R$-module $M$ is called {\it faithfully flat} if the
complex of $R$-modules
$N_1 \to N_2 \to N_3$ is exact if and only if
the sequence $M \otimes_R N_1 \to M \otimes_R N_2 \to M \otimes_R N_3$
is exact.
\item A ring map $R \to S$ is called {\it flat} if
$S$ is flat as an $R$-module.
\item A ring map $R \to S$ is called {\it faithfully flat} if
$S$ is faithfully flat as an $R$-module.
\end{enumerate}
\end{definition}

\noindent
Here is an example of how you can use the flatness condition.

\begin{lemma}
\label{editorial-source-lemma-flat-intersect-ideals}
Let $R$ be a ring. Let $I, J \subset R$ be ideals. Let $M$ be a flat
$R$-module. Then $IM\cap JM=(I\cap J)M$.
More generally, if $A_1,\ldots,A_r$ are submodules of an
$R$-module $N$, their tensors with $M$ identify with
submodules of $N\otimes_R M$, and
$$
\left(\bigcap_{i=1}^r A_i\right)\otimes_R M
\ \cong\ \bigcap_{i=1}^r(A_i\otimes_R M)
$$
under these original inclusions, for every finite $r\geq1$.
\end{lemma}

\begin{proof}
In the original ideal case use the exact sequence
$$
0\longrightarrow I\cap J\longrightarrow R
\longrightarrow R/I\oplus R/J,
$$
whose last map is $r\mapsto(r+I,r+J)$. Tensoring with $M$
preserves exactness. The maps
$$
R\otimes_R M\longrightarrow M,\quad r\otimes m\longmapsto rm,
\qquad
(R/I)\otimes_R M\longrightarrow M/IM,\quad
(r+I)\otimes m\longmapsto rm+IM
$$
are isomorphisms with inverses $m\mapsto1\otimes m$ and
$m+IM\mapsto(1+I)\otimes m$ respectively. The latter is
well-defined since each $a m$ with $a\in I$ maps to
$(a+I)\otimes m=0$; their composites fix the original tensor
generators and residue classes. The same formula applies to
$J$. Tensor product respects the indicated finite direct sum
by the component projections and inclusions. Thus the exact
sequence becomes
$$
0\longrightarrow (I\cap J)\otimes_R M\longrightarrow M
\longrightarrow M/IM\oplus M/JM.
$$
The middle kernel is $IM\cap JM$. The first map sends
$a\otimes m$ to $am$, so its image is exactly $(I\cap J)M$.
This proves the original equality.

For the general statement, tensoring $0\to A_i\to N$ with
$M$ proves that each original map $A_i\otimes_R M\to N\otimes_R M$
is injective. For two submodules $A,B\subseteq N$, use
$$
0\longrightarrow A\cap B\longrightarrow N
\longrightarrow (N/A)\oplus(N/B),\qquad
n\longmapsto(n+A,n+B).
$$
After tensoring, exactness identifies
$(A\cap B)\otimes_R M$ with the kernel of the two quotient
maps. Right exactness for $A\to N\to N/A\to0$ identifies
the kernel of the first quotient map with the image of
$A\otimes_R M$; the sequence for $B$ identifies the other
kernel. The kernel of the pair is their intersection, proving
the displayed isomorphism for two submodules. Its map is
induced by their original inclusions, so induction, intersecting
the first $r-1$ submodules with $A_r$, proves the finite case.
For $r=1$ it is the identity under the given inclusion.
\end{proof}

\begin{lemma}
\label{editorial-source-lemma-colimit-flat}
Let $R$ be a ring. Let $\{M_i, \varphi_{ii'}\}$ be a directed system of
flat $R$-modules. Then $\colim_i M_i$ is a flat $R$-module.
\end{lemma}

\begin{proof}
Write $M=\colim_i M_i$ with its original structure maps
$\lambda_i:M_i\to M$. For each module $N$, the natural
isomorphism
$$
\colim_i(N\otimes_R M_i)\longrightarrow N\otimes_R M,
\qquad [n\otimes m_i]\longmapsto n\otimes\lambda_i(m_i)
$$
is Lemma \ref{editorial-source-lemma-tensor-products-commute-with-limits};
it commutes with every map of $N$ because both composites
send the displayed generator to the tensor of its original
image. Given an exact complex $N_1\to N_2\to N_3$, flatness
of each $M_i$ makes its tensor complex exact at every stage.
Lemma \ref{algebra-lemma-directed-colimit-exact} gives exactness of
their directed colimit, and the displayed natural isomorphisms
identify it with the tensor complex for $M$. This proves
flatness without assuming that any transition map is injective.
\end{proof}

\begin{example}
\label{editorial-source-example-flat-infinite-intersections}
The finite-intersection assertion above cannot be extended to
arbitrary families of ideals for all flat modules. Take
$R=\mathbf Z$, $M=\mathbf Q$ and $I_n=2^n\mathbf Z$ for
$n\geq1$. The module $\mathbf Q$ is flat by
Lemma \ref{editorial-source-lemma-colimit-flat}: it is the colimit of copies
of $\mathbf Z$ indexed by $n\geq1$, with transition from
$n$ to $n+1$ multiplication by $n+1$ and structure map
$a\mapsto a/n!$. These maps are compatible. Every rational
$a/b$, $b\ne0$, has such a representation since $|b|$ divides
$n!$ for $n\geq|b|$. Equality $a/n!=a'/m!$ is equivalent
after a common stage $k\geq n,m$ to
$a(k!/n!)=a'(k!/m!)$, which is exactly the colimit relation.
Each copy of $\mathbf Z$ is flat because its tensor functor
is the identity via $a\otimes n\mapsto an$ and inverse
$n\mapsto1\otimes n$.

An integer divisible by every $2^n$ must be zero: for a
nonzero integer $a$, choose $n$ with $2^n>|a|$; it cannot
be a nonzero multiple of $2^n$. Thus $\bigcap_{n\geq1}I_n=0$.
On the other hand $I_nM=\mathbf Q$, since every original
$q\in\mathbf Q$ equals $2^n(q/2^n)$. Therefore
$$
\left(\bigcap_{n\geq1}I_n\right)M=0
\quad\hbox{whereas}\quad
\bigcap_{n\geq1}(I_nM)=\mathbf Q.
$$
\end{example}

\begin{lemma}
\label{editorial-source-lemma-composition-flat}
A composition of (faithfully) flat ring maps is
(faithfully) flat.
If $R \to R'$ is (faithfully) flat, and $M'$ is a
(faithfully) flat $R'$-module, then $M'$ is a
(faithfully) flat $R$-module.
\end{lemma}

\begin{proof}
For every original $R$-module $N$ there are mutually inverse
maps
$$
\begin{aligned}
\theta_N:M'\otimes_{R'}(R'\otimes_R N)&\longrightarrow M'\otimes_R N,
&m'\otimes(r'\otimes n)&\longmapsto r'm'\otimes n,\\
\eta_N:M'\otimes_R N&\longrightarrow M'\otimes_{R'}(R'\otimes_R N),
&m'\otimes n&\longmapsto m'\otimes(1\otimes n).
\end{aligned}
$$
The first respects $R'$-balancing because multiplication by
an $R'$-scalar can move between $m'$ and $r'$. It respects
$R$-balancing in the inner tensor by the restricted $R$-action
on $M'$. For the inverse, the equality
$(rm')\otimes(1\otimes n)=m'\otimes(\varphi(r)\otimes n)
=m'\otimes(1\otimes rn)$ proves $R$-balancing.
The composite $\theta_N\eta_N$ fixes $m'\otimes n$, and
$\eta_N\theta_N$ sends a generator to
$r'm'\otimes(1\otimes n)=m'\otimes(r'\otimes n)$.
Both maps commute with each homomorphism of $N$ on those
generators, so they are natural. Here $\varphi:R\to R'$
is the original ring map.

If $R'$ is flat over $R$ and $M'$ flat over $R'$, apply
their tensor functors successively to an exact complex
$N_1\to N_2\to N_3$. Both steps preserve exactness, and
$\theta_{N_i}$ identifies the resulting complex with
$M'\otimes_R N_1\to M'\otimes_R N_2\to M'\otimes_R N_3$.
This proves flatness over $R$. If both hypotheses are faithful,
an arbitrary original complex is exact if and only if the
first tensor complex is exact, and this in turn holds if
and only if the second tensor complex is exact. The same
natural maps prove faithful flatness over $R$. Taking $M'$
to be the target ring of a second ring map gives the first
assertion, for flat and faithfully flat maps alike.
\end{proof}

\begin{lemma}
\label{editorial-source-lemma-flat}
Let $M$ be an $R$-module. The following are equivalent:
\begin{enumerate}
\item
\label{editorial-source-item-flat}
$M$ is flat over $R$.
\item
\label{editorial-source-item-injective}
for every injection of $R$-modules $N \subset N'$
the map $N \otimes_R M \to N'\otimes_R M$ is injective.
\item
\label{editorial-source-item-f-ideal}
for every ideal $I \subset R$ the map
$I \otimes_R M \to R \otimes_R M = M$ is injective.
\item
\label{editorial-source-item-ffg-ideal}
for every finitely generated ideal $I \subset R$
the map $I \otimes_R M \to R \otimes_R M = M$ is injective.
\end{enumerate}
\end{lemma}

\begin{proof}
The successive implications
$(\ref{editorial-source-item-flat})\Rightarrow(\ref{editorial-source-item-injective})
\Rightarrow(\ref{editorial-source-item-f-ideal})\Rightarrow(\ref{editorial-source-item-ffg-ideal})$
follow directly: apply flatness to $0\to N\to N'$, then
specialize to inclusions $I\subseteq R$, and then to finitely
generated ideals. We prove
$(\ref{editorial-source-item-ffg-ideal})\Rightarrow(\ref{editorial-source-item-flat})$.
Suppose that $N_1 \to N_2 \to N_3$ is exact.
Let $K = \Ker(N_2 \to N_3)$ and $Q = \Im(N_2 \to N_3)$.
Then we get maps
$$
N_1 \otimes_R M \to
K \otimes_R M \to
N_2 \otimes_R M \to
Q \otimes_R M \to
N_3 \otimes_R M
$$
Observe that the first and third arrows are surjective. Thus if we show
that the second and fourth arrows are injective, then we are
done\footnote{Here is the argument in more detail:
Assume that we know that the second and fourth arrows are
injective. Lemma \ref{algebra-lemma-tensor-product-exact} (applied
to the exact sequence $K \to N_2 \to Q \to 0$) yields that
the sequence $K \otimes_R M \to N_2 \otimes_R M \to
Q \otimes_R M \to 0$ is exact. Hence,
$\Ker \left(N_2 \otimes_R M \to Q \otimes_R M\right)
= \Im \left(K \otimes_R M \to N_2 \otimes_R M\right)$.
Since
$\Im \left(K \otimes_R M \to N_2 \otimes_R M\right)
= \Im \left(N_1 \otimes_R M \to N_2 \otimes_R M\right)$
(due to the surjectivity of $N_1 \otimes_R M \to
K \otimes_R M$) and
$\Ker \left(N_2 \otimes_R M \to Q \otimes_R M\right)
= \Ker \left(N_2 \otimes_R M \to N_3 \otimes_R M\right)$
(due to the injectivity of $Q \otimes_R M \to
N_3 \otimes_R M$), this becomes
$\Ker \left(N_2 \otimes_R M \to N_3 \otimes_R M\right)
= \Im \left(N_1 \otimes_R M \to N_2 \otimes_R M\right)$,
which shows that the functor $- \otimes_R M$ is exact,
whence $M$ is flat.}.
Hence it suffices to show that $- \otimes_R M$ transforms
injective $R$-module maps into injective $R$-module maps.

\medskip\noindent
Assume $K \to N$ is an injective $R$-module map and
let $x \in \Ker(K \otimes_R M \to N \otimes_R M)$.
We have to show that $x$ is zero.
The $R$-module $K$ is the union of its finite
$R$-submodules; hence, $K \otimes_R M$ is
the colimit of $R$-modules of the form
$K_i \otimes_R M$ where $K_i$ runs over all finite
$R$-submodules of $K$
(because tensor product commutes with colimits).
Thus, for some $i$ our $x$ comes from an element
$x_i \in K_i \otimes_R M$. Thus we may assume that $K$
is a finite $R$-module. Assume this. We regard the
injection $K \to N$ as an inclusion, so that
$K \subset N$.

\medskip\noindent
The $R$-module $N$ is the union of its finite
$R$-submodules that contain $K$. Hence, $N \otimes_R M$
is the colimit of $R$-modules of the form
$N_i \otimes_R M$ where $N_i$ runs over all finite
$R$-submodules of $N$ that contain $K$
(again since tensor product commutes with colimits).
Notice that this is a colimit over a directed system
(since the sum of two finite submodules of $N$ is
again finite).
Hence, (by Lemma \ref{algebra-lemma-zero-directed-limit})
the element $x \in K \otimes_R M$ maps to
zero in at least one of these $R$-modules
$N_i \otimes_R M$ (since $x$ maps to zero
in $N \otimes_R M$).
Thus we may assume $N$ is a finite $R$-module.

\medskip\noindent
Assume $N$ is a finite $R$-module. Choose a surjection
$R^{\oplus n}\to N$, put $L$ equal to its kernel, and let
$L'$ be the inverse image of the original $K\subseteq N$.
Then $L\subseteq L'\subseteq R^{\oplus n}$,
$N=R^{\oplus n}/L$ and $K=L'/L$, under the quotient maps
just specified. For any submodule $G\subseteq R^{\oplus n}$,
the map $G\otimes_R M\to M^{\oplus n}$ sends
$(g_1,\ldots,g_n)\otimes m$ to $(g_1m,\ldots,g_nm)$.
Suppose these maps are injective for $G=L,L'$, and denote
them by $\beta$ and $\beta'$ respectively. Tensoring the
original inclusion gives the map in the direction
$$
\alpha:L\otimes_R M\longrightarrow L'\otimes_R M,
\qquad \ell\otimes m\longmapsto\ell\otimes m,
\qquad \beta'\alpha=\beta.
$$
It is injective since $\beta$ is injective. Right exactness
identifies the original $K\otimes_R M$ with
$(L'\otimes_R M)/\alpha(L\otimes_R M)$ and $N\otimes_R M$
with $M^{\oplus n}/\beta(L\otimes_R M)$.
The induced quotient map is injective: if
$v\in L'\otimes_R M$ has $\beta'(v)=\beta(u)$ for some
$u\in L\otimes_R M$, then
$\beta'(v-\alpha(u))=0$, hence $v=\alpha(u)$ by injectivity
of $\beta'$. This proves the required injection using the
original inclusions, without a reverse map from $L'$ to $L$.

\medskip\noindent
Thus it suffices to show that $L \otimes_R M \to M^{\oplus n}$
is injective when $L \subset R^{\oplus n}$ is an $R$-submodule.
We do this by induction on $n$. If $n=0$, both $L$ and
$M^{\oplus n}$ are zero, so the assertion holds. The base
case $n=1$ is proved below.
For the induction step assume $n > 1$ and set
$L'=L\cap(R\oplus0^{\oplus n-1})$. It is isomorphic to
the ideal $\{a\in R:(a,0,\ldots,0)\in L\}$ by the first
coordinate map, with inverse $a\mapsto(a,0,\ldots,0)$.
The last $n-1$ coordinate projection on $L$ has kernel $L'$,
so it induces the injection
$L''=L/L'\to R^{\oplus n-1}$ sending the class of
$(a_1,\ldots,a_n)$ to $(a_2,\ldots,a_n)$.
It is well-defined since elements of $L'$ have all these
coordinates zero, and injective by the kernel equality.
We obtain a diagram
$$
\xymatrix{
&
L' \otimes_R M \ar[r] \ar[d] &
L \otimes_R M \ar[r] \ar[d] &
L'' \otimes_R M \ar[r] \ar[d] &
0 \\
0 \ar[r] &
M \ar[r] &
M^{\oplus n} \ar[r] &
M^{\oplus n - 1} \ar[r] & 0
}
$$
The left vertical map is induced by the first coordinate,
the right by the last $n-1$ coordinates, and the lower maps
are the corresponding inclusion and projection. The upper
row is right exact, the lower row is exact, and the squares
commute on every tensor generator. The left and right
vertical maps are injective by the base case and induction.
If $v\in L\otimes_R M$ maps to zero in $M^{\oplus n}$,
its image in $L''\otimes_R M$ maps to zero on the right and
is therefore zero. Lift $v$ to $u\in L'\otimes_R M$ using
upper-row exactness. Its image in $M$ maps to zero in
$M^{\oplus n}$ and hence is zero, since the lower left map
is injective. Left vertical injectivity now gives $u=0$ and
then $v=0$, proving the middle injectivity.

\medskip\noindent
The base case of the induction above is when $L \subset R$ is an ideal.
In other words, we have to show that $I \otimes_R M \to M$ is injective
for any ideal $I$ of $R$. We know this is true when $I$ is finitely
generated. However, $I = \bigcup I_\alpha$ is the union of the
finitely generated ideals $I_\alpha$ contained in it. In other words,
$I=\colim_\alpha I_\alpha$, over the directed set of finitely
generated subideals; the sum of two such subideals contains
both. Tensor-colimit compatibility gives
$I\otimes_R M=\colim_\alpha(I_\alpha\otimes_R M)$.
If an element of this colimit maps to zero in $M$, choose
a representative $v_\alpha\in I_\alpha\otimes_R M$.
Its image in $M$ is already zero. The assumed injectivity
at that stage gives $v_\alpha=0$, so the original colimit
element is zero. This proves the base case for every ideal,
and completes the stated criterion.
\end{proof}

\begin{lemma}
\label{editorial-source-lemma-colimit-rings-flat}
Let $\{R_i, \varphi_{ii'}\}$ be a system of rings over the directed set $I$.
Let $R = \colim_i R_i$.
\begin{enumerate}
\item If $M$ is an $R$-module such that $M$ is flat as an $R_i$-module
for all $i$, then $M$ is flat as an $R$-module.
\item For $i \in I$ let $M_i$ be a flat $R_i$-module and
for $i' \geq i$ let $f_{ii'} : M_i \to M_{i'}$ be a $\varphi_{ii'}$-linear
map such that $f_{i' i''} \circ f_{i i'} = f_{i i''}$. Then
$M = \colim_{i \in I} M_i$ is a flat $R$-module.
\end{enumerate}
\end{lemma}

\begin{proof}
Write $\lambda_i:R_i\to R$ for the original ring maps and
$\mu_i:M_i\to M$ for the module structure maps. Part (1)
is the instance of part (2) with every $M_i=M$, acted on
through $\lambda_i$, and all underlying transition maps the
identity. They are semilinear because
$\lambda_{i'}\varphi_{ii'}=\lambda_i$.

We first specify the colimit constructions in part (2).
Only a nonempty directed indexing set is needed; none of
the arguments uses antisymmetry of its order. For a system
of abelian groups $B_i$ with transition maps $b_{ij}$,
the colimit is the direct sum of the $B_i$ modulo the
subgroup generated by the original relations
$\iota_i(x)-\iota_j(b_{ij}(x))$ for $i\leq j$.
Every element has a representative at a single stage:
move its finitely many summands to a common upper stage
and add their images there. Two representatives $x_i,x_j$
are equal in the colimit exactly when their images agree
at some common upper stage. To prove the nontrivial
direction, express their difference as a finite sum of
the generating relations and move every stage occurring
in that finite expression to one upper stage; compatibility
kills each relation there. The converse follows from the
defining transition relations. In particular a stage
element represents zero exactly when its image is zero
at a later stage. This description permits noninjective
transition maps throughout.

For $r=\lambda_i(r_i)$ and $m=\mu_j(m_j)$ choose $k\geq i,j$
and define the original scalar action by
$$
rm=\mu_k\bigl(\varphi_{ik}(r_i)f_{jk}(m_j)\bigr).
$$
Changing either representative or the chosen upper stage
does not change this element: move the relevant equalities
to one further stage and use semilinearity and compatibility.
The distributive laws and associativity of the action follow
by moving all finitely many entries of each identity to one
stage and applying its module axiom. The unit acts as the
identity because $[f_{ik}(m_i)]=[m_i]$.

Here is the varying-base tensor comparison needed below.
Let $A_i$ be $R_i$-modules with compatible semilinear maps
$h_{ij}:A_i\to A_j$, and put $A=\colim_i A_i$ with the
same scalar construction and structure maps $\alpha_i$.
There are transition maps
$$
\tau_{ij}:A_i\otimes_{R_i}M_i\longrightarrow A_j\otimes_{R_j}M_j,
\qquad a_i\otimes m_i\longmapsto h_{ij}(a_i)\otimes f_{ij}(m_i).
$$
They are well-defined: the two images of the balancing
relation $r_ia_i\otimes m_i=a_i\otimes r_im_i$ agree by
$$
\varphi_{ij}(r_i)h_{ij}(a_i)\otimes f_{ij}(m_i)
=h_{ij}(a_i)\otimes\varphi_{ij}(r_i)f_{ij}(m_i).
$$
Let $T=\colim_i(A_i\otimes_{R_i}M_i)$ with maps $\theta_i$.
The first comparison is
$$
F:T\longrightarrow A\otimes_R M,\qquad
\theta_i\left(\sum_{\ell=1}^q a_{i,\ell}\otimes m_{i,\ell}\right)
\longmapsto\sum_{\ell=1}^q
\alpha_i(a_{i,\ell})\otimes\mu_i(m_{i,\ell}).
$$
Each stage formula respects balancing and every transition
relation, so it defines a map of abelian groups. The scalar
formula above makes it $R$-linear by moving a scalar and
the tensor representative to one common stage.
For the inverse, define an additive balanced pairing on
representatives by
$$
B\bigl(\alpha_i(a_i),\mu_j(m_j)\bigr)
=\theta_k\bigl(h_{ik}(a_i)\otimes f_{jk}(m_j)\bigr),
\qquad k\geq i,j.
$$
Its value is independent of representatives and $k$:
move any two eventual equalities of representatives and
both choices of $k$ to one further stage, where the tensors
are identical. Additivity is checked at a common stage for
the two summands. For $r=\lambda_\ell(r_\ell)$, move
$r,a,m$ to one stage $k$. Then
$$
B(ra,m)=\theta_k\bigl(\varphi_{\ell k}(r_\ell)a_k\otimes m_k\bigr)
=\theta_k\bigl(a_k\otimes\varphi_{\ell k}(r_\ell)m_k\bigr)
=B(a,rm).
$$
Thus $B$ induces an $R$-linear map $G:A\otimes_R M\to T$.
On a pure tensor $FG$ gives
$\alpha_k(h_{ik}(a_i))\otimes\mu_k(f_{jk}(m_j))
=\alpha_i(a_i)\otimes\mu_j(m_j)$.
On a stage tensor $GF$ gives its image at a common later
stage, which is the same class $\theta_i(a_i\otimes m_i)$.
Finite sums of these generators prove both composites are
the identity. Hence the exact comparison, with both maps
specified, is
$$
\colim_i(A_i\otimes_{R_i}M_i)
\ \cong\ (\colim_i A_i)\otimes_R(\colim_i M_i).
$$

For completeness, stagewise injections induce an injection
on these directed colimits. Indeed, for a compatible family
of injections $u_i:X_i\to Y_i$, if $[x_i]$ maps to zero,
then at a later stage $j$ compatibility gives $u_j(x_j)=0$.
Stagewise injectivity gives $x_j=0$, hence $[x_i]=0$.
This uses no injectivity of the transition maps.

Now retain the original finitely generated ideal
$\mathfrak a=(a_1,\ldots,a_n)\subseteq R$. Choose one stage
$i$ containing representatives $a_{\ell,i}$ of all its
generators, and put
$$
\mathfrak a'=(a_{1,i},\ldots,a_{n,i})\subseteq R_i,
\qquad
\mathfrak a_j=\mathfrak a'R_j
=(\varphi_{ij}(a_{1,i}),\ldots,\varphi_{ij}(a_{n,i}))
\quad(j\geq i).
$$
The tail $j\geq i$ is cofinal: every original representative,
and every eventual equality, can be moved to a stage above
both its original stage and $i$. It has the same ring and
module colimits. The map
$$
\delta:\colim_{j\geq i}\mathfrak a_j\longrightarrow\mathfrak a,
\qquad [x_j]\longmapsto\lambda_j(x_j)
$$
is surjective. For $x=\sum_{\ell=1}^n r_\ell a_\ell$,
choose a stage $j\geq i$ representing all $r_\ell$ and use
$x_j=\sum_{\ell=1}^n r_{\ell,j}\varphi_{ij}(a_{\ell,i})$.
It is injective: if $\lambda_j(x_j)=0$ in $R$, some
$k\geq j$ has $\varphi_{jk}(x_j)=0$ in $R_k$; because
$\mathfrak a_k$ is an actual ideal in that ring, the image
is zero in $\mathfrak a_k$ itself. This proves that the
inverse formula $x\mapsto[x_j]$ is independent of its
expression and representatives; equal expressions differ
by an element eventually zero in the original ring system.
Both maps respect the scalar action and fix the represented
elements in their composites.

Apply the proved tensor comparison to $A_j=\mathfrak a_j$.
It identifies
$$
\colim_{j\geq i}(\mathfrak a_j\otimes_{R_j}M_j)
\ \cong\ \mathfrak a\otimes_R M.
$$
For every $j$, flatness of $M_j$ makes
$\mathfrak a_j\otimes_{R_j}M_j\to M_j$ injective. Its
formula is $x_j\otimes m_j\mapsto x_jm_j$, and compatibility
of these maps is exactly semilinearity of $f_{jk}$.
Their colimit is injective by the element argument above.
Under the displayed comparison it is the original map
$\mathfrak a\otimes_R M\to M$, $a\otimes m\mapsto am$.
Lemma \ref{editorial-source-lemma-flat} now proves flatness of $M$.
The zero ideal, including a presentation with no generators,
has zero tensor product and requires no exception.

The comparison also proves the full injection criterion
directly. For an injection $N'\to N$ of $R$-modules, regard
both as constant underlying systems with restricted
$R_i$-actions and identity transition maps. Each
$N'\otimes_{R_i}M_i\to N\otimes_{R_i}M_i$ is injective.
Their colimit comparison is exactly
$N'\otimes_R M\to N\otimes_R M$, hence it is injective.
This gives a second proof with all original modules retained.

The explicitly displayed composition law for $f_{ij}$ also
suffices if identities on the diagonal were not stipulated.
Indeed $e_i=f_{ii}$ is an idempotent $R_i$-linear map and
$$
M_i=\operatorname{Im}(e_i)\oplus\operatorname{Ker}(e_i),
\qquad m_i=e_i(m_i)+(m_i-e_i(m_i)).
$$
The intersection of the two summands is zero since $e_i$
is the identity on its image. Each summand of a flat module
is flat: its tensor map on any injection is a direct summand
of the injective tensor map for $M_i$. Compatibility gives
$f_{ij}=e_jf_{ij}=f_{ij}e_i$, so the restrictions to
$\operatorname{Im}(e_i)$ form a genuine directed system.
The colimit by all the original relations is isomorphic to
the colimit of these images. The maps send $[m_i]$ to
$[e_i(m_i)]$ and an image representative to $[m_i]$ in the
original colimit; they respect transitions by the displayed
identities and are inverse because $[e_i(m_i)]=[m_i]$ is
one of the original relations. Thus no extra identity
hypothesis is needed for the flatness conclusion.

\medskip\noindent
We record exactly what is additionally needed for faithful
flatness. For a maximal ideal $\mathfrak m\subset R$ set
$\mathfrak m_i=\lambda_i^{-1}(\mathfrak m)$. Each is a
proper prime, though it need not be maximal in $R_i$, and
$\varphi_{ij}^{-1}(\mathfrak m_j)=\mathfrak m_i$.
Consequently the induced maps $R_i/\mathfrak m_i\to R_j/\mathfrak m_j$
are injective. The comparison
$$
\colim_i(R_i/\mathfrak m_i)\longrightarrow R/\mathfrak m,
\qquad [r_i+\mathfrak m_i]\longmapsto\lambda_i(r_i)+\mathfrak m
$$
is surjective by representatives and injective because an
element mapping to zero already has $r_i\in\mathfrak m_i$.
Its inverse sends any representative $\lambda_i(r_i)+\mathfrak m$
to $[r_i+\mathfrak m_i]$; the same zero test proves independence
of the choice. Applying the varying-base tensor comparison
gives
$$
M/\mathfrak mM\ \cong\
\colim_i(M_i/\mathfrak m_iM_i).
$$
Here the stage quotient-tensor isomorphism sends
$(r_i+\mathfrak m_i)\otimes u_i$ to
$r_iu_i+\mathfrak m_iM_i$ and has inverse
$u_i+\mathfrak m_iM_i\mapsto(1+\mathfrak m_i)\otimes u_i$.
The same formulas identify the receiving quotient over $R$.
Balancing makes both maps well-defined, and both composites
fix the indicated generators. Thus the original transition
on the right is
$u_i+\mathfrak m_iM_i\mapsto f_{ij}(u_i)+\mathfrak m_jM_j$,
and the receiving map is
$u_i+\mathfrak m_iM_i\mapsto\mu_i(u_i)+\mathfrak mM$.
The eventual-zero rule now proves
$$
\mu_i(u_i)+\mathfrak mM=0
\quad\Longleftrightarrow\quad
f_{ij}(u_i)\in\mathfrak m_jM_j\text{ for some }j\geq i.
$$

A flat module is faithfully flat exactly when all these
maximal-ideal quotients are nonzero. Here is the full
argument, using the definition by exactness reflection.
Faithful flatness detects nonzero modules by applying that
definition to $0\to N\to0$, and in particular detects
$N=R/\mathfrak m$. Conversely suppose all the quotients
are nonzero. Given $0\ne n\in N$, let
$J=\operatorname{Ann}_R(n)$, choose a maximal ideal
$\mathfrak m\supseteq J$, and retain the injection
$R/J\to N$, $r+J\mapsto rn$. The module $M/JM$ surjects
onto the nonzero $M/\mathfrak mM$, so it is nonzero.
Flatness makes $(R/J)\otimes_R M=M/JM\to N\otimes_R M$
injective, proving $N\otimes_R M\ne0$ for every nonzero $N$.

To deduce exactness reflection, take an original complex
$N_1\xrightarrow{d_1}N_2\xrightarrow{d_2}N_3$ and write
$J=\operatorname{Im}(d_1)\subseteq K=\operatorname{Ker}(d_2)$.
Flatness embeds $J\otimes_R M$ and $K\otimes_R M$ in
$N_2\otimes_R M$. The exact sequences defining image and
kernel identify them with the tensor image and tensor
kernel, respectively: the map $N_1\to J$ is surjective,
and $0\to K\to N_2\to N_3$ stays exact. Right exactness
for $J\to K\to K/J\to0$ then identifies
$$
(K/J)\otimes_R M\ \cong\
\operatorname{Ker}(d_2\otimes\mathrm{id}_M)\big/
\operatorname{Im}(d_1\otimes\mathrm{id}_M).
$$
These are the maps induced by the original inclusions and
quotient $K\to K/J$, so the identification respects those
original submodules. If the tensor complex is exact, this
quotient is zero. Nonzero-module detection gives $K/J=0$,
hence the original complex is exact. This proves the criterion.

Combining it with the exact quotient comparison proves:
$M$ is faithfully flat over $R$ if and only if, for every
maximal $\mathfrak m\subset R$, there are an index $i$ and
$u_i\in M_i$ such that
$$
f_{ij}(u_i)\notin\mathfrak m_jM_j
\quad\hbox{for every }j\geq i.
$$
This quantifier keeps the same original element at every
later stage. It cannot be replaced by merely having a
nonzero quotient at each stage.

For example take constant $R_i=k$ and $M_i=k$ over a
nonzero field, indexed by $i\in\mathbf Z_{\geq0}$, with
identity ring maps and $f_{ij}=0$ for $i<j$, $f_{ii}$ the
identity. Each stage module is faithfully flat, but every
element dies at the next stage, so $R=k$ and $M=0$.
Even injective module transitions do not suffice. Take
$R_i=M_i=\mathbf Z$ with identity ring maps and
$f_{ij}(a)=2^{j-i}a$. The map $[a\text{ at }i]\mapsto a/2^i$
identifies $M$ with $\mathbf Z[1/2]$: it is surjective by
the description of its elements, and equality of two
fractions is exactly equality after multiplication to a
common power-of-two denominator. Every stage is faithfully
flat and every transition injective, but $M/2M=0$ since
$a/2^i=2(a/2^{i+1})$, so $M$ is not faithfully flat.

A useful sufficient condition follows with the actual
residue maps retained. Suppose every $M_i$ is faithfully
flat and, for each maximal $\mathfrak m\subset R$, there
is a tail $i\geq i_0$ on which all maps
$M_i/\mathfrak m_iM_i\to M_j/\mathfrak m_jM_j$ are injective.
Faithful flatness at $i_0$ makes its quotient nonzero because
$R_{i_0}/\mathfrak m_{i_0}\ne0$. A nonzero element there
remains nonzero at every later stage, which is precisely
the necessary and sufficient condition just proved.

One condition on the original transitions that guarantees
this is purity of the scalar-extended maps
$$
g_{ij}:R_j\otimes_{R_i}M_i\longrightarrow M_j,
\qquad r_j\otimes u_i\longmapsto r_jf_{ij}(u_i).
$$
Here purity means that tensoring this map with every
$R_j$-module gives an injection. The formula is balanced
by semilinearity, so $g_{ij}$ is an $R_j$-linear map.
The contraction equality gives an injection of
$R_i$-modules $R_i/\mathfrak m_i\to R_j/\mathfrak m_j$.
Flatness of $M_i$ and then purity give injections
$$
(R_i/\mathfrak m_i)\otimes_{R_i}M_i
\longrightarrow (R_j/\mathfrak m_j)\otimes_{R_i}M_i
\longrightarrow (R_j/\mathfrak m_j)\otimes_{R_j}M_j.
$$
For the second map the associativity comparison sends
$(s_j+\mathfrak m_j)\otimes(r_j\otimes u_i)$ to
$(s_jr_j+\mathfrak m_j)\otimes u_i$ and has inverse
$(s_j+\mathfrak m_j)\otimes u_i\mapsto
(s_j+\mathfrak m_j)\otimes(1\otimes u_i)$; balancing and
these formulas verify both composites. Under the quotient
identifications, the composite injection is exactly
$u_i+\mathfrak m_iM_i\mapsto f_{ij}(u_i)+\mathfrak m_jM_j$.
Thus faithful flatness of the stages together with these
pure transition maps suffices.

If $R=0$, its unital action forces $M=0$. The only unital
module over that ring is zero, so its tensor functor
preserves and reflects exactness; it is faithfully flat in
the stated sense. The maximal-ideal criterion gives the
same result because there are no maximal ideals.
\end{proof}

\begin{lemma}
\label{editorial-source-lemma-flat-base-change}
Suppose that $M$ is (faithfully) flat over $R$, and that $R \to R'$
is a ring map. Then $M \otimes_R R'$ is (faithfully) flat over $R'$.
\end{lemma}

\begin{proof}
For every original $R'$-module $N$, define
$$
\begin{aligned}
\theta_N:N\otimes_{R'}(R'\otimes_R M)&\longrightarrow N\otimes_R M,
&n\otimes(r'\otimes m)&\longmapsto r'n\otimes m,\\
\eta_N:N\otimes_R M&\longrightarrow N\otimes_{R'}(R'\otimes_R M),
&n\otimes m&\longmapsto n\otimes(1\otimes m).
\end{aligned}
$$
The first formula respects the outer $R'$-balancing and the
inner $R$-balancing by the original $R'$-action on $N$.
The inverse respects $R$-balancing because
$(rn)\otimes(1\otimes m)=n\otimes(\varphi(r)\otimes m)
=n\otimes(1\otimes rm)$. Both composites fix the tensor
generators, using $r'n\otimes(1\otimes m)=n\otimes(r'\otimes m)$.
The formulas also commute with every $R'$-linear map of $N$.
The flip $M\otimes_R R'\to R'\otimes_R M$, $m\otimes r'\mapsto r'\otimes m$,
identifies the module in the statement with this one and is
its own inverse after reversing the factors. Exactness of a
complex of $R'$-modules is equivalent to exactness of its
underlying $R$-module complex, since kernels and images are
the same subsets. Apply flatness, or preservation and
reflection of exactness, for the original $M$ over $R$ and
then use the displayed natural maps. This proves both claims.
\end{proof}

\begin{lemma}
\label{editorial-source-lemma-flatness-descends}
Let $R \to R'$ be a faithfully flat ring map.
Let $M$ be a module over $R$, and set $M' = R' \otimes_R M$.
Then $M$ is flat over $R$ if and only if $M'$ is flat over $R'$.
\end{lemma}

\begin{proof}
By Lemma \ref{editorial-source-lemma-flat-base-change} we see that if $M$ is flat
then $M'$ is flat. For the converse, suppose that $M'$ is flat.
Let $N_1 \to N_2 \to N_3$ be an exact sequence of $R$-modules.
We want to show that $N_1 \otimes_R M \to N_2 \otimes_R M \to N_3 \otimes_R M$
is exact. We know that
$N_1 \otimes_R R' \to N_2 \otimes_R R' \to N_3 \otimes_R R'$ is
exact, because $R \to R'$ is flat. Flatness of $M'$ implies that
$N_1 \otimes_R R' \otimes_{R'} M'
\to N_2 \otimes_R R' \otimes_{R'} M'
\to N_3 \otimes_R R' \otimes_{R'} M'$ is exact.
The exact natural comparison for every original $N$ is
\[
\begin{split}
(N\otimes_R R')\otimes_{R'}(R'\otimes_R M)
&\longrightarrow (N\otimes_R M)\otimes_R R',\\
(n\otimes a)\otimes(b\otimes m)&\longmapsto(n\otimes m)\otimes ab.
\end{split}
\]
Its inverse sends $(n\otimes m)\otimes c$ to
$(n\otimes c)\otimes(1\otimes m)$. All $R$-balancing
relations follow from the original coefficient map $R\to R'$;
the $R'$-balancing relation moves a scalar between $a$ and
$b$ and preserves their product $ab$. These observations
make both maps well-defined. One composite fixes
$(n\otimes m)\otimes c$. The other gives
$(n\otimes ab)\otimes(1\otimes m)
=(n\otimes a)\otimes(b\otimes m)$ by $R'$-balancing.
Their formulas commute with every map of $N$. Thus the
exact complex above is the tensor with $R'$ of the original
complex $N_1\otimes_R M\to N_2\otimes_R M\to N_3\otimes_R M$.
Faithful flatness of $R'$ over $R$ reflects exactness, so
this original complex is exact, as required.
\end{proof}

\begin{lemma}
\label{editorial-source-lemma-flatness-descends-more-general}
Let $R$ be a ring. Let $S \to S'$ be a flat map of $R$-algebras.
Let $M$ be a module over $S$, and set $M' = S' \otimes_S M$.
\begin{enumerate}
\item If $M$ is flat over $R$, then $M'$ is flat over $R$.
\item If $S \to S'$ is faithfully flat, then $M$ is flat
over $R$ if and only if $M'$ is flat over $R$.
\end{enumerate}
\end{lemma}

\begin{proof}
Let $f:N\to N'$ be an injection of $R$-modules and put
$K=\Ker(f\otimes\mathrm{id}_M)$. The map
$f\otimes\mathrm{id}_M:N\otimes_R M\to N'\otimes_R M$
is $S$-linear, with the $S$-action on its $M$ factor, so
$K$ is an $S$-module. Tensoring its defining exact sequence
with the flat $S$-module $S'$ identifies $K\otimes_S S'$
with the kernel after that tensor operation. For every $N$
the precise comparison is
\[
\begin{split}
(N\otimes_R M)\otimes_S S'&\longrightarrow
N\otimes_R(S'\otimes_S M),\\
(n\otimes m)\otimes s'&\longmapsto n\otimes(s'\otimes m),
\end{split}
\]
with inverse $n\otimes(s'\otimes m)\mapsto(n\otimes m)\otimes s'$.
The $S$-balancing relations agree on the two sides by moving
the original $S$-scalar between $m$ and $s'$; the $R$-balancing
relations agree through $R\to S\to S'$. The composites fix
the generators, and the formulas commute with $f$. Consequently
the kernel identification is exactly
$$
K\otimes_S S'\ \cong\
\Ker(N\otimes_R M'\longrightarrow N'\otimes_R M').
$$
If $M$ is flat over $R$, then $K=0$, so the right kernel
vanishes for every injection $f$ and Lemma \ref{editorial-source-lemma-flat}
gives flatness of $M'$ over $R$. Conversely, if $M'$ is
flat, then $K\otimes_S S'=0$. When $S'$ is faithfully flat
over $S$, apply reflection of exactness to the complex
$0\to K\to0$: its tensor complex is exact, hence $K=0$.
The same criterion now gives flatness of $M$ over $R$.
\end{proof}

\begin{lemma}
\label{editorial-source-lemma-flat-permanence}
Let $R \to S$ be a ring map. Let $M$ be an $S$-module.
If $M$ is flat as an $R$-module and faithfully flat as an $S$-module,
then $R\to S$ is flat.
More generally, whenever $M$ is faithfully flat over $S$,
the map $R\to S$ is flat (respectively faithfully flat)
if and only if $M$ is flat (respectively faithfully flat)
as an $R$-module.
\end{lemma}

\begin{proof}
For every original $R$-module $N$, the comparison
$$
(N\otimes_R S)\otimes_S M\longrightarrow N\otimes_R M,
\qquad (n\otimes s)\otimes m\longmapsto n\otimes sm
$$
has inverse $n\otimes m\mapsto(n\otimes1)\otimes m$.
Balancing over $S$ preserves the product $sm$, and balancing
over $R$ is compatible with the restricted action on $M$.
Both composites fix the generators; for one of them use
$(n\otimes1)\otimes sm=(n\otimes s)\otimes m$.
These maps are natural in $N$.

If $M$ is flat over $R$, an exact complex
$N_1\to N_2\to N_3$ stays exact after tensoring with $M$.
The comparison identifies this with the result of first
tensoring with $S$ and then with $M$ over $S$. Faithful
flatness over $S$ reflects exactness, so the first tensor
complex is exact. This proves the original conclusion that
$R\to S$ is flat. The converse under the same faithful
$S$-module hypothesis follows from
Lemma \ref{editorial-source-lemma-composition-flat}.

If $M$ is faithfully flat over $R$, take any original complex
$N_1\to N_2\to N_3$ whose tensor with $S$ is exact. Tensoring
that complex with the flat $S$-module $M$ and using the
comparison shows its tensor with $M$ over $R$ is exact.
Faithful flatness over $R$ reflects this to exactness of the
original complex. Together with flatness just proved, this
makes $R\to S$ faithfully flat. Conversely, if $R\to S$
is faithfully flat, composition with the faithfully flat
$S$-module $M$ gives faithful flatness over $R$, again by
Lemma \ref{editorial-source-lemma-composition-flat}.
\end{proof}

\noindent
Let $R$ be a ring.
Let $M$ be an $R$-module.
Let $\sum_{i=1}^n f_ix_i=0$ be a finite relation in $M$,
with $f_i\in R$, $x_i\in M$, and $n\geq0$.
We say the relation $\sum f_i x_i$
is {\it trivial} if there exist an integer $m \geq 0$,
elements $y_j \in M$, $j = 1, \ldots, m$, and elements $a_{ij} \in R$,
$i = 1, \ldots, n$, $j = 1, \ldots, m$ such that
$$
x_i = \sum\nolimits_j a_{ij} y_j, \forall i,
\quad\text{and}\quad
0 = \sum\nolimits_i f_ia_{ij}, \forall j.
$$

\begin{lemma}[Equational criterion of flatness]
\label{editorial-source-lemma-flat-eq}
A module $M$ over $R$ is flat if and only if
every relation in $M$ is trivial.
\end{lemma}

\begin{proof}
Assume $M$ is flat and let $\sum f_i x_i = 0$ be a relation in $M$.
Let $I = (f_1, \ldots, f_n)$, and let
$K = \Ker(R^n \to I, (a_1, \ldots, a_n) \mapsto \sum_i a_i f_i)$.
So we have the short exact sequence
$0 \to K \to R^n \to I \to 0$. Then $\sum f_i \otimes x_i$
is an element of $I \otimes_R M$ which maps
to zero in $R \otimes_R M = M$. By flatness
$\sum f_i \otimes x_i$ is zero in $I \otimes_R M$.
Thus there exists an element of $K \otimes_R M$ mapping
to $\sum e_i \otimes x_i \in R^n \otimes_R M$ where $e_i$
is the $i$th basis element of $R^n$.
Write this element as a finite sum
$\sum_{j=1}^m k_j\otimes y_j$, and write each original
$k_j\in K\subseteq R^n$ as $\sum_{i=1}^n a_{ij}e_i$.
Membership in $K$ gives $\sum_{i=1}^n f_i a_{ij}=0$
for every $j$. Under the coordinate isomorphism
$R^n\otimes_R M\to M^n$, $e_i\otimes y$ maps to the
vector having $y$ in coordinate $i$ and zero elsewhere;
its inverse sends $(x_1,\ldots,x_n)$ to
$\sum_i e_i\otimes x_i$. Comparing all coordinates of
the original lifted tensor gives
$x_i=\sum_{j=1}^m a_{ij}y_j$ for every $i$, exactly the
two defining conditions. The sum may be empty; when $n=0$
there are no coordinates and the empty relation is handled
with $m=0$.

\medskip\noindent
Assume every relation is trivial, let $I$
be a finitely generated ideal, and let $x = \sum f_i \otimes x_i$
be an element of $I \otimes_R M$ mapping to zero in $R \otimes_R M = M$.
This just means exactly that $\sum f_i x_i$ is a relation in
$M$. And the fact that it is trivial implies easily that
$x$ is zero, because
$$
x
=
\sum f_i \otimes x_i
=
\sum f_i \otimes \left(\sum a_{ij}y_j\right)
=
\sum \left(\sum f_i a_{ij}\right) \otimes y_j
=
0
$$
The sums are finite, so the middle equality follows by
additivity and the original balancing relation
$f_i\otimes a_{ij}y_j=(f_i a_{ij})\otimes y_j$ in
$I\otimes_R M$. Thus the map $I\otimes_R M\to M$ is
injective for every finitely generated ideal $I$, and
Lemma \ref{editorial-source-lemma-flat} proves that $M$ is flat.
\end{proof}

\begin{lemma}
\label{editorial-source-lemma-flat-finite-simultaneous-relations}
An $R$-module $M$ is flat if and only if it has the following
property. For every matrix $(f_{\ell i})\in\operatorname{Mat}_{r\times n}(R)$
and elements $x_1,\ldots,x_n\in M$ satisfying
$\sum_{i=1}^n f_{\ell i}x_i=0$ for every $1\leq\ell\leq r$,
there are an integer $m\geq0$, elements $y_1,\ldots,y_m\in M$
and coefficients $a_{ij}\in R$ such that
$$
x_i=\sum_{j=1}^m a_{ij}y_j\quad(1\leq i\leq n),\qquad
\sum_{i=1}^n f_{\ell i}a_{ij}=0
\quad(1\leq\ell\leq r,\ 1\leq j\leq m).
$$
Here $r,n$ may be zero; all matrices and sums are finite.
\end{lemma}

\begin{proof}
Keep the original matrix and let
$\psi:R^n\to R^r$ send the $i$th basis vector to the
column $(f_{1i},\ldots,f_{ri})$. Put $K=\Ker\psi$.
Flatness makes
$K\otimes_R M\to R^n\otimes_R M\to R^r\otimes_R M$
exact. The coordinate tensor isomorphisms in the preceding
proof identify its last map with the original matrix
acting on $M^n$. Thus the given vector is the image of
a finite sum $\sum_{j=1}^m k_j\otimes y_j$.
Write $k_j=\sum_i a_{ij}e_i$. All $r$ original equations
for $k_j\in K$ are precisely the second displayed family,
and coordinate equality gives the first family. No finite
generation assumption on $K$ is used. Conversely, the case
$r=1$ gives the equational criterion
of Lemma \ref{editorial-source-lemma-flat-eq}, hence flatness. When $r=0$
one may take $m=n$, $a_{ij}=\delta_{ij}$ and $y_j=x_j$;
when $n=0$ take $m=0$. These choices verify the boundary
cases with the original empty matrices retained.
\end{proof}

\begin{lemma}
\label{editorial-source-lemma-flat-tor-zero}
Suppose that $R$ is a ring, that $0\to M''\to M'\to M\to0$
is a short exact sequence, and that $N$ is an $R$-module. If $M$ is flat
then $N \otimes_R M'' \to N \otimes_R M'$ is injective, i.e., the
sequence
$$
0 \to N \otimes_R M'' \to N \otimes_R M' \to N \otimes_R M \to 0
$$
is a short exact sequence.
\end{lemma}

\begin{proof}
Retain a surjection $q:F=R^{(I)}\to N$ with kernel
$K$ and inclusion $j:K\to F$. Write $u:M''\to M'$ and
$v:M'\to M$ for the original maps. The coordinate
isomorphism $L\otimes_R F\to L^{(I)}$ for any $L$ sends
$\ell\otimes\sum_i r_ie_i$ to $(r_i\ell)_i$ and has
inverse $(\ell_i)_i\mapsto\sum_i\ell_i\otimes e_i$;
all supports are finite. The rows and columns of the
following diagram are induced by $u,v,j,q$:
$$
\begin{matrix}
 & & 0 & & 0 & & 0 & & \\
 & & \uparrow & & \uparrow & & \uparrow & & \\
 & & M''\otimes_R N & \to & M' \otimes_R N & \to & M \otimes_R N & \to & 0 \\
 & & \uparrow & & \uparrow & & \uparrow & & \\
0 & \to & (M'')^{(I)} & \to & (M')^{(I)} & \to & M^{(I)} & \to & 0 \\
 & & \uparrow & & \uparrow & & \uparrow & & \\
 & & M''\otimes_R K & \to & M' \otimes_R K & \to & M \otimes_R K & \to & 0 \\
 & & & & & & \uparrow & & \\
 & & & & & & 0 & &
\end{matrix}
$$
Each column ending in zero is right exact by tensoring
$K\to F\to N\to0$. The bottom and top rows are right
exact by tensoring $M''\to M'\to M\to0$. The middle row
is short exact coordinate by coordinate, including
surjectivity since a finitely supported tuple has finitely
many lifts. The bottom right vertical map is injective
because $M$ is flat and $K\to F$ is injective.
Every square commutes on its tensor generators.

Here is the full chase proving the remaining injection.
Let $x\in M''\otimes_R N$ map to zero in $M'\otimes_R N$.
Lift it to $a\in(M'')^{(I)}$. The image $u^{(I)}(a)$
has zero image in $M'\otimes_R N$, so it is the image
of some $b\in M'\otimes_R K$. The image of $b$ in
$M\otimes_R K$ maps to zero in $M^{(I)}$, because
$v^{(I)}u^{(I)}(a)=0$. Bottom right injectivity forces
that image to be zero. Bottom-row exactness supplies
$c\in M''\otimes_R K$ mapping to $b$.
Its image $d\in(M'')^{(I)}$ satisfies
$u^{(I)}(d)=u^{(I)}(a)$ by commutativity. Middle-row
injectivity gives $d=a$. As $d$ comes from
$M''\otimes_R K$, its image in $M''\otimes_R N$ is zero,
so $x=0$. This proves the asserted injection.
Finally the flips $L\otimes_R N\to N\otimes_R L$,
$\ell\otimes n\mapsto n\otimes\ell$, are mutual inverses
and commute with $u,v$, so the diagram conclusion is
exactly the sequence in the statement.
\end{proof}

\begin{lemma}
\label{editorial-source-lemma-flat-quotient-cyclic-tensor-criterion}
In the original short exact sequence
$0\to M''\xrightarrow{u}M'\xrightarrow{v}M\to0$,
assume $M'$ is flat. The following are equivalent:
\begin{enumerate}
\item $M$ is flat.
\item For every $R$-module $N$, the map
$N\otimes_R M''\to N\otimes_R M'$ is injective.
\item For every finitely generated ideal $I\subseteq R$,
the map $M''/IM''\to M'/IM'$ induced by $u$ is injective.
\end{enumerate}
\end{lemma}

\begin{proof}
The first implication is Lemma \ref{editorial-source-lemma-flat-tor-zero}.
The second follows by taking $N=R/I$ and using the
quotient-tensor maps $(r+I)\otimes m\mapsto rm+IM$
and $m+IM\mapsto(1+I)\otimes m$, on each original module.

Assume the third condition and retain any finite original
relation $\sum_{i=1}^n f_ix_i=0$ in $M$.
Choose lifts $\widetilde x_i\in M'$ with
$v(\widetilde x_i)=x_i$, and set $I=(f_1,\ldots,f_n)$.
The element $\sum_i f_i\widetilde x_i$ has zero image
under $v$, so it equals $u(w)$ for a unique $w\in M''$.
Its class in $M'/IM'$ is zero. Injectivity of the specified
quotient map gives $w\in IM''$.
Write $w=\sum_{\ell=1}^q c_\ell z_\ell$ with
$c_\ell\in I$, $z_\ell\in M''$, and express each original
$c_\ell=\sum_{i=1}^n f_i d_{i\ell}$.
Putting $w_i=\sum_{\ell=1}^q d_{i\ell}z_\ell$ gives the
exact equality $w=\sum_i f_iw_i$.
The corrected lifts
$\widehat x_i=\widetilde x_i-u(w_i)$ therefore satisfy
$$
\sum_i f_i\widehat x_i
=u(w)-u\left(\sum_i f_iw_i\right)=0,
\qquad v(\widehat x_i)=x_i.
$$
By flatness of the original $M'$ and
Lemma \ref{editorial-source-lemma-flat-eq}, there are $y_j\in M'$ and
$a_{ij}\in R$ with
$\widehat x_i=\sum_j a_{ij}y_j$ and
$\sum_i f_i a_{ij}=0$ for every $j$.
Applying $v$ gives $x_i=\sum_j a_{ij}v(y_j)$ with the
same coefficients and the same vanishing equations.
Thus every original relation in $M$ is trivial and
$M$ is flat by that criterion. The empty relation
requires no lifts and is already trivial.
\end{proof}


\begin{lemma}
\label{editorial-source-lemma-flat-ses}
Suppose that $0 \to M' \to M \to M'' \to 0$ is
a short exact sequence of $R$-modules.
If $M'$ and $M''$ are flat so is $M$.
If $M$ and $M''$ are flat so is $M'$.
\end{lemma}

\begin{proof}
We will use the criterion that a module $N$ is flat if for
every ideal $I \subset R$ the map $N \otimes_R I \to N$ is injective,
see Lemma \ref{editorial-source-lemma-flat}.
Consider an ideal $I \subset R$.
Consider the diagram
$$
\begin{matrix}
0 & \to & M' & \to & M & \to & M'' & \to & 0 \\
& & \uparrow & & \uparrow & & \uparrow & & \\
& & M'\otimes_R I & \to & M \otimes_R I & \to & M''\otimes_R I & \to & 0
\end{matrix}
$$
The vertical maps are $m\otimes a\mapsto am$ for the
original ideal $I$, so both squares commute on generators.
The upper row is the given short exact sequence and the
lower row is right exact. For the first assertion suppose
$M'$ and $M''$ are flat, and let $x\in M\otimes_R I$
map to zero in $M$. Its image in $M''\otimes_R I$ maps
to zero in $M''$ and hence is zero by flatness of $M''$.
Lower-row exactness lifts $x$ to $y\in M'\otimes_R I$.
The image of $y$ in $M'$ maps to zero in $M$, so it is
zero by upper-row injectivity. Flatness of $M'$ now gives
$y=0$, hence $x=0$. Lemma \ref{editorial-source-lemma-flat} proves that
$M$ is flat.

For the second assertion suppose $M$ and $M''$ are flat.
Lemma \ref{editorial-source-lemma-flat-tor-zero}, applied with quotient
$M''$ and tensor module $I$, makes the lower left map
injective. If $y\in M'\otimes_R I$ maps to zero in $M'$,
its image $x\in M\otimes_R I$ maps to zero in $M$.
Flatness of $M$ gives $x=0$, and lower left injectivity
gives $y=0$. The same ideal criterion proves flatness
of the original $M'$.
\end{proof}

\begin{lemma}
\label{editorial-source-lemma-easy-ff}
Let $R$ be a ring.
Let $M$ be an $R$-module.
The following are equivalent:
\begin{enumerate}
\item $M$ is faithfully flat, and
\item $M$ is flat and for all $R$-module homomorphisms $\alpha : N \to N'$
we have $\alpha = 0$ if and only if $\alpha \otimes \text{id}_M = 0$.
\end{enumerate}
\end{lemma}

\begin{proof}
Faithful flatness includes flatness. It also implies that
$L\otimes_R M=0$ forces $L=0$, by reflecting exactness
of $0\to L\to0$. For an original map $\alpha:N\to N'$,
retain $K=\Ker\alpha$. Flatness makes
$0\to K\otimes_R M\to N\otimes_R M\to N'\otimes_R M$
exact. If $\alpha\otimes\mathrm{id}_M=0$, the first
injection is also surjective. Right exactness of tensoring
$K\to N\to N/K\to0$ then gives $(N/K)\otimes_R M=0$.
The vanishing reflection just proved gives $N/K=0$, so
$K=N$ and $\alpha=0$. Conversely a zero map tensors
to a zero map by its formula on generators. This proves
the first implication with the original kernel retained.

Assume the second condition. Flatness already preserves
exactness. To prove reflection, take an original complex
$N_1\xrightarrow{d_1}N_2\xrightarrow{d_2}N_3$ whose
tensor complex is exact, and put $J=\operatorname{Im}d_1$.
For each $x\in\Ker d_2$, keep the original map
$\alpha:R\to N_2/J$, $r\mapsto rx+J$.
For every $m\in M$, the element $x\otimes m$ belongs to
$\Ker(d_2\otimes\mathrm{id}_M)$, hence to the image of
$d_1\otimes\mathrm{id}_M$. The quotient map $N_2\to N_2/J$
tensored with $M$ kills that image, so
$(x+J)\otimes m=0$. Therefore on every tensor generator
$$
(\alpha\otimes\mathrm{id}_M)(r\otimes m)
=(rx+J)\otimes m=(x+J)\otimes rm=0.
$$
The assumed reflection of zero morphisms gives $\alpha=0$,
and evaluation at $1$ gives $x\in J$. Thus
$\Ker d_2\subseteq\operatorname{Im}d_1$; the reverse
inclusion holds because the original sequence is a complex.
It is exact, proving faithful flatness.
\end{proof}

\begin{lemma}
\label{editorial-source-lemma-ff}
\begin{slogan}
A flat module is faithfully flat if and only if it has nonzero fibers.
\end{slogan}
Let $M$ be a flat $R$-module.
The following are equivalent:
\begin{enumerate}
\item $M$ is faithfully flat,
\item for every nonzero $R$-module $N$, the tensor product $M \otimes_R N$
is nonzero,
\item for all $\mathfrak p \in \Spec(R)$
the tensor product $M \otimes_R \kappa(\mathfrak p)$ is nonzero, and
\item for all maximal ideals $\mathfrak m$ of $R$
the tensor product $M \otimes_R \kappa(\mathfrak m) = M/{\mathfrak m}M$
is nonzero.
\end{enumerate}
\end{lemma}

\begin{proof}
Assume $M$ faithfully flat and $N \not = 0$. By Lemma \ref{editorial-source-lemma-easy-ff}
the nonzero map $1 : N \to N$ induces a nonzero map
$M \otimes_R N \to M \otimes_R N$, so $M \otimes_R N \not = 0$.
Thus (1) implies (2). The implications (2) $\Rightarrow$ (3) $\Rightarrow$ (4)
are immediate.

\medskip\noindent
Assume (4), and retain a complex
$N_1\xrightarrow{d_1}N_2\xrightarrow{d_2}N_3$ whose
tensor complex is exact. Write
$K=\Ker d_2$, $J=\operatorname{Im}d_1\subseteq K$,
and keep its original cohomology $H=K/J$.
Flatness embeds $J\otimes_R M$ and $K\otimes_R M$ in
$N_2\otimes_R M$. Tensoring the surjection $N_1\to J$
identifies the first image with
$\operatorname{Im}(d_1\otimes\mathrm{id}_M)$; tensoring
$0\to K\to N_2\to N_3$ identifies the second image with
$\operatorname{Ker}(d_2\otimes\mathrm{id}_M)$.
The tensor complex is exact, so these submodules coincide.
Right exactness applied to $J\to K\to H\to0$ gives
$H\otimes_R M=0$ with these exact original maps.

Take any $x\in H$ and set $I=\{f\in R:fx=0\}$.
The map $R/I\to H$, $f+I\mapsto fx$, is well-defined
and injective by the definition of $I$.
Flatness makes $(R/I)\otimes_R M\to H\otimes_R M=0$
injective. The quotient-tensor isomorphism
$(f+I)\otimes m\mapsto fm+IM$ with inverse
$m+IM\mapsto(1+I)\otimes m$ therefore gives $M/IM=0$.
If $I$ were proper, choose a maximal ideal
$\mathfrak m\supseteq I$. The quotient map
$M/IM\to M/\mathfrak mM$, $m+IM\mapsto m+\mathfrak mM$,
is well-defined and surjective. Its source is zero, forcing
$M/\mathfrak mM=0$, contrary to (4). Thus $I=R$,
$x=1x=0$, and $H=0$. This proves exactness reflection.
For the zero ring all modules are zero and all prime and
maximal-ideal conditions are vacuous, so the same equivalence
also covers that case.
\end{proof}

\begin{lemma}
\label{editorial-source-lemma-ff-rings}
Let $R \to S$ be a flat ring map.
The following are equivalent:
\begin{enumerate}
\item $R \to S$ is faithfully flat,
\item the induced map on $\Spec$ is surjective, and
\item any closed point $x \in \Spec(R)$ is
in the image of the map $\Spec(S) \to \Spec(R)$.
\end{enumerate}
\end{lemma}

\begin{proof}
For the original prime $\mathfrak p\subset R$, the fibre
ring $S\otimes_R\kappa(\mathfrak p)$ is nonzero exactly
when there is a prime of $S$ contracting to $\mathfrak p$,
by Remark \ref{algebra-remark-fundamental-diagram} and
Lemma \ref{editorial-source-lemma-in-image}. The prime correspondence is
through the quotient by $\mathfrak p$ followed by inversion
of the original $R\setminus\mathfrak p$; its primes
contract to primes of $S$ with precisely that contraction.
Thus condition (3) of Lemma \ref{editorial-source-lemma-ff} says exactly
that every prime is in the spectrum image. Its condition
(4) says exactly that every maximal ideal is in the image,
and maximal ideals are the closed points of $\Spec(R)$.
Since $S$ is flat over $R$, that lemma proves all three
equivalences asserted here.
\end{proof}

\begin{lemma}
\label{editorial-source-lemma-local-flat-ff}
A flat local ring homomorphism of local rings is faithfully flat.
\end{lemma}

\begin{proof}
Write $\varphi:(R,\mathfrak m)\to(S,\mathfrak n)$ for
the original local homomorphism. Its defining local
condition is $\varphi^{-1}(\mathfrak n)=\mathfrak m$.
Thus the unique closed point of $\Spec(R)$ is the image
of the prime $\mathfrak n$ of $S$. Flatness and
Lemma \ref{editorial-source-lemma-ff-rings} give faithful flatness.
\end{proof}

\noindent
Flatness meshes well with localization.

\begin{lemma}
\label{editorial-source-lemma-flat-localization}
Let $R$ be a ring. Let $S \subset R$ be a multiplicative subset.
\begin{enumerate}
\item The localization $S^{-1}R$ is a flat $R$-algebra.
\item If $M$ is an $S^{-1}R$-module, then $M$ is a flat $R$-module
if and only if $M$ is a flat $S^{-1}R$-module.
\item Suppose $M$ is an $R$-module. Then
$M$ is a flat $R$-module if and only if $M_{\mathfrak p}$ is a flat
$R_{\mathfrak p}$-module for all primes $\mathfrak p$ of $R$.
\item Suppose $M$ is an $R$-module. Then $M$ is a flat $R$-module if
and only if $M_{\mathfrak m}$ is a flat
$R_{\mathfrak m}$-module for all maximal ideals $\mathfrak m$ of $R$.
\item Suppose $R \to A$ is a ring map, $M$ is an $A$-module,
and $g_1, \ldots, g_m \in A$ are elements generating the unit
ideal of $A$. Then $M$ is flat over $R$ if and only if each localization
$M_{g_i}$ is flat over $R$.
\item Suppose $R \to A$ is a ring map, and $M$ is an $A$-module.
Then $M$ is a flat $R$-module if and only if the localization
$M_{\mathfrak q}$ is a flat $R_{\mathfrak p}$-module
(with $\mathfrak p$ the prime of $R$ lying under $\mathfrak q$)
for all primes $\mathfrak q$ of $A$.
\item Suppose $R \to A$ is a ring map, and $M$ is an $A$-module.
Then $M$ is a flat $R$-module if and only if the localization
$M_{\mathfrak m}$ is a flat $R_{\mathfrak p}$-module
(with $\mathfrak p$ the inverse image of $\mathfrak m$ under $R\to A$)
for all maximal ideals $\mathfrak m$ of $A$.
\end{enumerate}
\end{lemma}

\begin{proof}
Retain every original coefficient map; in (5)--(7) write it
as $\varphi:R\to A$. All contractions below are inverse
images under this map, which need not be injective.

We first recall the exact fraction facts used in the proof.
For a ring $B$, a multiplicative set $W\subseteq B$ and a
$B$-module $L$, the equality $\ell/w=\ell'/w'$ means
that $v(w'\ell-w\ell')=0$ for some original $v\in W$.
Thus $\ell/w=0$ means that some $v\in W$ kills $\ell$.
If $L'\to L$ is injective, this equality for an element
of $L'$ holds in $L'$ whenever it holds in $L$, so
localization preserves injections. For a map $u:L\to N$,
an element $\ell/w$ in the localized kernel has
$v u(\ell)=0$ for some $v\in W$. Then $v\ell\in\Ker u$
and $\ell/w=(v\ell)/(vw)$, proving that localization
preserves kernels. Its image consists exactly of fractions
$u(\ell)/w$, and a surjection remains surjective by lifting
the numerator. These facts prove exactness with the
original fraction witnesses retained.

For later zero detection, if $0\ne\ell\in L$ its
annihilator is a proper ideal of $B$ and is contained in
a maximal ideal $\mathfrak m$. Then $\ell/1\ne0$ in
$L_{\mathfrak m}$, since an element outside $\mathfrak m$
killing $\ell$ would belong to that annihilator. Hence
$L=0$ if all its maximal localizations vanish, as in
Lemma \ref{algebra-lemma-characterize-zero-local}.

For an original $R$-module $N$, the $A$-action on
$N\otimes_R M$ is $a(n\otimes m)=n\otimes am$.
For every multiplicative $W\subseteq A$, there are
mutually inverse maps
\[
\begin{split}
W^{-1}(N\otimes_R M)&\longrightarrow N\otimes_R W^{-1}M,\\
\frac{\sum_{\ell=1}^d n_\ell\otimes m_\ell}{w}
&\longmapsto\sum_{\ell=1}^d n_\ell\otimes\frac{m_\ell}{w},
\end{split}
\]
and $n\otimes(m/w)\mapsto(n\otimes m)/w$.
For the forward map start with $n\otimes m\mapsto
n\otimes(m/1)$, which is $A$-linear and $R$-balanced;
every element of $W$ acts invertibly on its target, so
it extends to the localization with the stated formula.
For the inverse, if $v(w'm-wm')=0$, then
$v(w'(n\otimes m)-w(n\otimes m'))=0$ by the original
$A$-action. This proves independence of the fraction.
The $R$-balancing relation is
$rn\otimes m=n\otimes\varphi(r)m$, and is preserved by
the formula. Additivity follows by common denominators.
Both composites fix the displayed generators and their
finite sums. The maps are natural in $N$ and $M$, since
they apply their homomorphisms only to the numerators.
Consequently, if $M$ is flat over $R$, then $W^{-1}M$
is flat over $R$: localize the injection obtained by
tensoring any injection of $R$-modules with $M$ and use
these natural isomorphisms.

For a prime $\mathfrak q\subset A$, put
$\mathfrak p=\varphi^{-1}(\mathfrak q)$. The map and action are
$$
R_{\mathfrak p}\longrightarrow A_{\mathfrak q},\quad
\frac rs\longmapsto\frac{\varphi(r)}{\varphi(s)},
\qquad
\frac rs\cdot\frac ma=\frac{\varphi(r)m}{\varphi(s)a}.
$$
All denominators are valid because $s\notin\mathfrak p$
is equivalent to $\varphi(s)\notin\mathfrak q$.
The exact two-sided comparison for every original $N$ is
$$
(N\otimes_R M)_{\mathfrak q}
\ \cong\ N_{\mathfrak p}\otimes_{R_{\mathfrak p}}M_{\mathfrak q},
\qquad
\frac{n\otimes m}{a}\longmapsto\frac n1\otimes\frac ma,
$$
with inverse
$$
\frac ns\otimes\frac ma\longmapsto
\frac{n\otimes m}{\varphi(s)a}.
$$
If $n/s=n'/s'$, choose $u\notin\mathfrak p$ with
$u(s'n-sn')=0$. The two proposed inverse images have
common denominator $\varphi(s)\varphi(s')a$ and numerator
$(s'n-sn')\otimes m$. Multiplication by
$\varphi(u)\notin\mathfrak q$ kills it, by $R$-balancing.
If $m/a=m'/a'$, choose $b\notin\mathfrak q$ with
$b(a'm-am')=0$. Using common denominator $\varphi(s)aa'$
gives numerator $n\otimes(a'm-am')$, killed by $b$.
For $r/t\in R_{\mathfrak p}$ the two balancing images are
$$
\frac{rn\otimes m}{\varphi(t)\varphi(s)a}
=\frac{n\otimes\varphi(r)m}{\varphi(s)\varphi(t)a},
$$
so the inverse is well-defined and additive.
The forward map is the localization of
$n\otimes m\mapsto(n/1)\otimes(m/1)$; elements outside
$\mathfrak q$ act invertibly on its target.
One composite fixes $(n\otimes m)/a$, while the other
sends $(n/s)\otimes(m/a)$ to
$(n/1)\otimes(m/(\varphi(s)a))$, which equals the
original tensor by balancing with $1/s$.
The formulas prove naturality for the original maps.

We also retain the exact ideal descent used in the source.
For any ideal $I\subseteq R_{\mathfrak p}$, put
$J=\{r\in R:r/1\in I\}$. Then $J_{\mathfrak p}=I$:
if $j\in J$, its fraction $j/s=(1/s)(j/1)$ is in $I$;
if $a/s\in I$, then $a/1=(s/1)(a/s)\in I$, so $a\in J$.
The localized map $J_{\mathfrak p}\to R_{\mathfrak p}$ is
injective by localization of $J\hookrightarrow R$.
For specified finite generators $a_\ell/s_\ell$ of $I$,
the smaller specified ideal $J_0=(a_1,\ldots,a_n)$ also
has $(J_0)_{\mathfrak p}=I$, because each $s_\ell/1$ is
a unit and both identities above apply to that generator.
This does not assert finite generation of the contraction $J$.

We now prove (1). Keep the original multiplicative set
$S\subseteq R$ and put $B=S^{-1}R$.
For every $N$ the maps
$$
N\otimes_R B\longrightarrow S^{-1}N,\quad
n\otimes(r/s)\longmapsto rn/s,
\qquad
n/s\longmapsto n\otimes(1/s)
$$
are inverse. To check the inverse on fraction equality,
if $u(s'n-sn')=0$ then its difference of images is
$$
(s'n-sn')\otimes\frac1{ss'}
=u(s'n-sn')\otimes\frac1{uss'}=0.
$$
The forward map respects equal fractions for the same
reason and respects the original $R$-balancing relation.
The two composites fix the generators, since
$rn\otimes(1/s)=n\otimes(r/s)$.
Localization preserves injections, so these natural maps
and Lemma \ref{editorial-source-lemma-flat} prove that $B$ is flat over $R$.

For (2), if $M$ is flat over $B$, then (1) and
Lemma \ref{editorial-source-lemma-composition-flat} imply flatness over $R$.
The receiving tensor comparison can also be written as
$$
(S^{-1}N)\otimes_B M\longrightarrow N\otimes_R M,
\qquad (n/s)\otimes m\longmapsto n\otimes(1/s)m,
$$
with inverse $n\otimes m\mapsto(n/1)\otimes m$.
For example if $u(s'n-sn')=0$, the difference between
the first images is
$(s'n-sn')\otimes(1/(ss'))m
=u(s'n-sn')\otimes(1/(uss'))m=0$.
Balancing by $r/t\in B$ agrees on both sides after
moving $r$ and the inverse action of $t$ to the $M$ factor.
Both composites are identities by these same relations.

Conversely let $M$ be a $B$-module flat over $R$.
For any $B$-module $P$, the quotient map
$P\otimes_R M\to P\otimes_B M$ has inverse
$p\otimes_B m\mapsto p\otimes_R m$.
To verify the additional balancing needed for this inverse,
for each original $s\in S$ use
$$
(1/s)p\otimes_R m
=(1/s)p\otimes_R s((1/s)m)
=p\otimes_R(1/s)m.
$$
Together with $R$-balancing this gives balancing by every
$r/s\in B$. The composites fix all generators.
An injection of $B$-modules is an injection of underlying
$R$-modules, so these natural maps and flatness over $R$
give the injection criterion over $B$. This proves (2).

Next prove (6) and (7), retaining the original ideal maps.
Suppose $M_{\mathfrak m}$ is flat over $R_{\mathfrak p}$
for each maximal $\mathfrak m\subset A$, where
$\mathfrak p=\varphi^{-1}(\mathfrak m)$.
For an ideal $I\subseteq R$, consider the original
$A$-linear map
$$
\eta:I\otimes_R M\longrightarrow M,
\qquad a\otimes m\longmapsto\varphi(a)m,
\qquad K=\Ker\eta.
$$
At $\mathfrak m$ the comparison already proved gives
$$
(I\otimes_R M)_{\mathfrak m}
\cong I_{\mathfrak p}\otimes_{R_{\mathfrak p}}M_{\mathfrak m}.
$$
It identifies $\eta_{\mathfrak m}$ with the multiplication
map $(a/s)\otimes(m/t)\mapsto\varphi(a)m/(\varphi(s)t)$,
so that map is injective by the asserted local flatness
and the inclusion $I_{\mathfrak p}\hookrightarrow R_{\mathfrak p}$.
Exactness of localization gives $K_{\mathfrak m}=0$.
Zero detection at all maximal ideals of $A$ gives $K=0$,
and Lemma \ref{editorial-source-lemma-flat} proves flatness over $R$.

Conversely suppose $M$ is flat over $R$ and retain any
prime $\mathfrak q\subset A$ with
$\mathfrak p=\varphi^{-1}(\mathfrak q)$.
For an ideal $I\subseteq R_{\mathfrak p}$ take its original
contraction $J\subseteq R$ above, so $J_{\mathfrak p}=I$.
Flatness makes $J\otimes_R M\to M$ injective. Localizing
at $\mathfrak q$ preserves its injectivity, and the exact
comparison identifies it with
$I\otimes_{R_{\mathfrak p}}M_{\mathfrak q}\to M_{\mathfrak q}$.
The ideal criterion proves flatness over $R_{\mathfrak p}$.
This proves necessity at every prime and sufficiency at
all maximal ideals. Since the former condition includes
the latter, (6) and (7) both follow. A maximal ideal of
$A$ may contract to a nonmaximal prime of $R$; no step
replaces the original contracted prime by a maximal one.

Taking $A=R$ and $\varphi=\mathrm{id}_R$ in (6) and (7)
gives (3) and (4), respectively, with the same module $M$
and the same prime and maximal localizations.

For (5) we prove the stronger arbitrary-family statement:
if $(g_\lambda)_{\lambda\in\Lambda}$ is any family in $A$
such that for every maximal $\mathfrak m$ with
$M_{\mathfrak m}\ne0$ some $g_\lambda\notin\mathfrak m$,
then $M$ is flat over $R$ if and only if every
$M_{g_\lambda}$ is flat over $R$.

First, for each original $g\in A$, multiplication by $g$
gives a directed system of the original $R$-module $M$,
indexed by $n\in\mathbf Z_{\geq0}$. Its colimit is
$$
\colim(M\xrightarrow{g}M\xrightarrow{g}\cdots)
\longrightarrow M_g,\qquad [m\text{ at }n]\longmapsto m/g^n.
$$
Compatibility is $gm/g^{n+1}=m/g^n$.
Every fraction occurs. A fraction $m/g^n$ is zero exactly
when some $g^k m=0$, which is exactly eventual vanishing
of its original representative. This proves injectivity
as well as surjectivity; the inverse sends a fraction to
that class, independently of its representative by the
same equality test. Lemma \ref{editorial-source-lemma-colimit-flat} then
proves that $M_g$ is flat whenever $M$ is flat.

For the converse retain any injection of $R$-modules
$N'\hookrightarrow N$ and its original $A$-module kernel
$K=\Ker(N'\otimes_R M\to N\otimes_R M)$.
If every $M_{g_\lambda}$ is flat over $R$, localization
and the first tensor comparison identify the localized
map with
$N'\otimes_R M_{g_\lambda}\to N\otimes_R M_{g_\lambda}$,
which is injective. Thus $K_{g_\lambda}=0$ for each $\lambda$.
At a maximal ideal $\mathfrak m$ with $M_{\mathfrak m}=0$,
the same comparison gives
$(N'\otimes_R M)_{\mathfrak m}=N'\otimes_R M_{\mathfrak m}=0$,
so $K_{\mathfrak m}=0$. At any other maximal ideal choose
$g_\lambda\notin\mathfrak m$ by the hypothesis.
For every $k\in K$, the equality $K_{g_\lambda}=0$ gives
$g_\lambda^e k=0$ for some $e\geq0$, which may depend
on $k$. This power is outside $\mathfrak m$, so $k/1=0$
in $K_{\mathfrak m}$. Thus $K_{\mathfrak m}=0$ at every
maximal ideal, and zero detection gives $K=0$.
The injection criterion proves flatness over $R$.

The stated finite unit-ideal hypothesis implies this cover:
if every $g_i$ belonged to a maximal ideal, so would their
generated ideal, contrary to $(g_1,\ldots,g_m)=A$.
This proves (5), including both directions.

The stronger covering hypothesis can equivalently be
stated on all primes $\mathfrak q$ with $M_{\mathfrak q}\ne0$.
Indeed choose $x\in M$ with $x/1\ne0$ at such a prime.
Then $\operatorname{Ann}_A(x)\subseteq\mathfrak q$ by
the exact fraction-zero test. At every maximal
$\mathfrak m\supseteq\mathfrak q$, the same inclusion
shows $x/1\ne0$ in $M_{\mathfrak m}$. A cover element
outside $\mathfrak m$ is also outside $\mathfrak q$.
The reverse implication specializes the prime condition
to maximal ideals. No finite generation of $M$ or of the
family is required.

This really weakens the unit-ideal requirement. For a
nonzero ring $R$, take $A=R\times R$ with the diagonal
map from $R$, and an $R$-module $M$ viewed as an $A$-module
through the first projection. The element $g=(1,0)$ acts
as the identity on $M$, so $M_g=M$ by the maps
$m/g^n\mapsto m$ and $m\mapsto m/1$.
The element $1-g=(0,1)$ acts as zero. If a prime contains
$g$, it cannot contain $1-g$, which becomes an invertible
annihilator of the localized module; hence that localization
is zero. Thus $D(g)$ covers every nonzero stalk of $M$,
although $gA=R\times0$ is a proper ideal of $A$.

Finally all assertions include zero cases. If $0$ lies
in a multiplicative set, its localized ring and modules
are zero by the fraction rule; a unital module over that
ring is zero. In particular $g=0$ or a nilpotent $g$ gives
$M_g=0$ and an empty principal open. If $A=0$, then $M=0$
and its prime and maximal conditions are empty. If the
cover family is empty, its hypothesis says all maximal
stalks of $M$ are zero, so $M=0$ by zero detection and the
equivalence is valid. The zero module is flat by the
definition in every case.
\end{proof}

\begin{lemma}
\label{editorial-source-lemma-flat-going-down}
Let $R \to S$ be flat. Let $\mathfrak p \subset \mathfrak p'$
be primes of $R$. Let $\mathfrak q' \subset S$ be a prime of $S$
mapping to $\mathfrak p'$. Then there exists a prime
$\mathfrak q \subset \mathfrak q'$ mapping to $\mathfrak p$.
\end{lemma}

\begin{proof}
Let $\varphi:R\to S$ be the original map and retain the
stated primes. The localization map
$$
\varphi':R_{\mathfrak p'}\longrightarrow S_{\mathfrak q'},
\qquad r/t\longmapsto\varphi(r)/\varphi(t)
$$
is defined because $t\notin\mathfrak p'$ implies
$\varphi(t)\notin\mathfrak q'$. Its maximal-ideal
contraction is $\mathfrak p'R_{\mathfrak p'}$, so it is
local, and Lemma \ref{editorial-source-lemma-flat-localization} makes it
flat. It is faithfully flat by Lemma \ref{editorial-source-lemma-local-flat-ff}.
The ideal $\mathfrak pR_{\mathfrak p'}$ is prime and has
contraction $\mathfrak p$ in $R$ by prime localization.
Lemma \ref{editorial-source-lemma-ff-rings} supplies a prime
$\widetilde{\mathfrak q}\subset S_{\mathfrak q'}$ whose
contraction under $\varphi'$ is that ideal. Define
$$
\mathfrak q=\{s\in S:s/1\in\widetilde{\mathfrak q}\}.
$$
It is prime as an inverse image of a prime. An element of
$S\setminus\mathfrak q'$ becomes a unit in $S_{\mathfrak q'}$
and cannot belong to this inverse image, so
$\mathfrak q\subseteq\mathfrak q'$. Finally
$$
r\in\varphi^{-1}(\mathfrak q)
\Longleftrightarrow
r/1\in\mathfrak pR_{\mathfrak p'}
\Longleftrightarrow r\in\mathfrak p.
$$
This proves the exact required contraction, without
assuming the original flat map is injective.
\end{proof}

\noindent
The property of $R \to S$ described in the lemma is called the
``going down property''. See Definition \ref{algebra-definition-going-up-down}.

\begin{lemma}
\label{editorial-source-lemma-colimit-faithfully-flat}
Let $R$ be a ring. Let $\{S_i, \varphi_{ii'}\}$ be a directed system of
faithfully flat $R$-algebras. Then $S = \colim_i S_i$ is a faithfully flat
$R$-algebra.
\end{lemma}

\begin{proof}
The underlying $R$-module colimit is flat by
Lemma \ref{editorial-source-lemma-colimit-flat}. For every maximal ideal
$\mathfrak m\subset R$, faithful flatness at each stage
makes $S_i/\mathfrak mS_i$ nonzero, by Lemma \ref{editorial-source-lemma-ff}.
The canonical ring comparison is
$$
\colim_i(S_i/\mathfrak mS_i)\longrightarrow S/\mathfrak mS,
\qquad [s_i+\mathfrak mS_i]\longmapsto[s_i]+\mathfrak mS.
$$
It is surjective by representatives. If $[s_i]\in\mathfrak mS$,
write its original expression as a finite sum
$[s_i]=\sum_{\ell=1}^q a_\ell[t_{j_\ell}]$ with
$a_\ell\in\mathfrak m$. Move all terms to a common stage
and then to a further stage where this equality holds.
There the image of $s_i$ belongs to the actual ideal
$\mathfrak mS_k$, so its quotient class is eventually zero.
This proves injectivity and gives the inverse by choosing
any original representative; equality of two choices is
handled by the same kernel argument.

If this colimit were the zero ring, the class of $1$ at
one stage would become zero at some later stage. Every
transition is a unital $R$-algebra map, so the image is
the original $1$ of that later quotient, contradicting
its nonzero ring structure. Hence $S/\mathfrak mS\ne0$
for every maximal ideal, and Lemma \ref{editorial-source-lemma-ff} gives
faithful flatness. This is the specific reason algebra
colimits preserve faithful flatness although arbitrary
module colimits need not: the same element $1$ survives
in every receiving residue quotient. For $R=0$ all
unital $R$-algebras and their modules are zero, and the
faithfully-flat conclusion holds by the definition.
\end{proof}








\subsection{Editorial treatment of Algebra lines 9328--9344}

\hypertarget{algebra-context-098}{}

\noindent\href{https://github.com/stacks/stacks-project/blob/a04446e57ec1fbc252a871afcec7752fb2807b14/algebra.tex\#L9328}{Original source, lines 9328--9344}. Complete original and reviewed TeX passages accompany this note. Every intervention is separately recorded; the following treatment is editorial.


\begin{lemma}
\label{editorial-source-lemma-support-zero}
\begin{slogan}
A module over a ring has empty support if and only if it is the trivial module.
\end{slogan}
Let $R$ be a ring. Let $M$ be an $R$-module. Then
$$
M = (0) \Leftrightarrow \text{Supp}(M) = \emptyset.
$$
\end{lemma}

\begin{proof}
If $M=0$, all its localizations are zero, so its support
is empty. Conversely, for any original nonzero $m\in M$,
the ideal $\{r\in R:rm=0\}$ is proper and lies in a
maximal ideal $\mathfrak p$. The fraction $m/1$ is nonzero
in $M_{\mathfrak p}$: if it were zero, a witness
$s\notin\mathfrak p$ would satisfy $sm=0$, contradicting
membership of every such annihilator in $\mathfrak p$.
Thus the support contains a maximal ideal, as also proved
in Lemma \ref{algebra-lemma-characterize-zero-local}.
\end{proof}


\subsection{Editorial treatment of Algebra lines 9363--9382}

\hypertarget{algebra-context-099}{}

\noindent\href{https://github.com/stacks/stacks-project/blob/a04446e57ec1fbc252a871afcec7752fb2807b14/algebra.tex\#L9363}{Original source, lines 9363--9382}. Complete original and reviewed TeX passages accompany this note. Every intervention is separately recorded; the following treatment is editorial.


\begin{lemma}
\label{editorial-source-lemma-annihilator-flat-base-change}
Let $R \to S$ be a flat ring map. Let $M$ be an $R$-module and
$m \in M$. Then $\text{Ann}_R(m) S = \text{Ann}_S(m \otimes 1)$.
If $M$ is a finite $R$-module, then
$\text{Ann}_R(M)S=\text{Ann}_S(M\otimes_R S)$.
The latter conclusion also holds if finitely many elements
$m_1,\ldots,m_n\in M$ have
$\operatorname{Ann}_R(M)=\bigcap_{i=1}^n\operatorname{Ann}_R(m_i)$,
without requiring that they generate $M$.
\end{lemma}

\begin{proof}
Write the original ring map as $\varphi:R\to S$ and retain
$I=\operatorname{Ann}_R(m)$. The exact sequence
$$
0\longrightarrow I\longrightarrow R\longrightarrow M,
\qquad r\longmapsto rm,
$$
remains exact after tensoring with the flat $R$-module $S$.
Under $R\otimes_RS\to S$, $r\otimes s\mapsto\varphi(r)s$,
the last map becomes $s\mapsto m\otimes s=(m\otimes1)s$.
The image of $I\otimes_RS\to S$, $a\otimes s\mapsto\varphi(a)s$,
is exactly the extended ideal $IS$. Its equality with the
kernel therefore proves
$$
\operatorname{Ann}_R(m)S=\operatorname{Ann}_S(m\otimes1).
$$

For the original finite generating family $m_1,\ldots,m_n$,
retain $I_i=\operatorname{Ann}_R(m_i)$. An element kills $M$
exactly when it kills every generator. The elements
$m_i\otimes1$ generate $M\otimes_RS$, so the full comparison is
\[
\begin{split}
\operatorname{Ann}_S(M\otimes_RS)
&=\bigcap_{i=1}^n\operatorname{Ann}_S(m_i\otimes1)\\
&=\bigcap_{i=1}^n(I_iS)
=\left(\bigcap_{i=1}^n I_i\right)S
=\operatorname{Ann}_R(M)S.
\end{split}
\]
The finite-intersection equality is
Lemma \ref{editorial-source-lemma-flat-intersect-ideals}.

For the stronger assertion, suppose only that the original
finite list detects the annihilator:
$$
J=\operatorname{Ann}_R(M)=\bigcap_{i=1}^n I_i.
$$
These elements need not generate $M$. The sequence
$$
0\longrightarrow J\longrightarrow R\longrightarrow M^n,
\qquad r\longmapsto(rm_1,\ldots,rm_n),
$$
is exact. Tensoring and using the coordinate maps
$M^n\otimes_RS\simeq(M\otimes_RS)^n$ gives
$JS=\bigcap_i\operatorname{Ann}_S(m_i\otimes1)$.
Consequently every element annihilating the whole tensor
module belongs to $JS$. Conversely each finite sum
$\sum_j\varphi(a_j)s_j$ with $a_j\in J$ kills every tensor:
$$
\left(\sum_j\varphi(a_j)s_j\right)(m\otimes s)
=\sum_j(a_jm)\otimes s_js=0.
$$
This proves equality for the weaker hypothesis. For $n=0$
the empty intersection is $R$, forcing $M=0$; both module
annihilators are then the whole rings, as required.
For example, any nonempty free module, even with an infinite
basis, has its annihilator detected by one original basis
vector: $re=0$ implies $r=0$ by that coordinate.

\medskip\noindent
For an arbitrary module the same element calculation gives
the exact formula
$$
\operatorname{Ann}_S(M\otimes_RS)
=\bigcap_{m\in M}\bigl(\operatorname{Ann}_R(m)S\bigr),
$$
since all the elements $m\otimes1$ generate the tensor module.
Extension of ideals need not commute with this infinite
intersection, even for a faithfully flat ring map.

Retain the explicit example
$$
R=\mathbf Z,\qquad
M=\bigoplus_{n\geq1}\mathbf Z/2^n\mathbf Z,\qquad
S=\mathbf Z\times\mathbf Q,
\qquad \varphi(a)=(a,a).
$$
The rational numbers are the localization at the nonzero
integers. For every module $N$ the tensor-localization map is
$N\otimes_{\mathbf Z}\mathbf Q\to
(\mathbf Z\setminus\{0\})^{-1}N$,
$n\otimes a/b\mapsto an/b$, with inverse $n/b\mapsto n\otimes1/b$.
Localization preserves injections: a fraction mapping to zero
has a nonzero integer witness killing its numerator, and
injectivity gives the same relation in the original module.
Thus $\mathbf Q$ is flat, as also established in
Example \ref{editorial-source-example-flat-infinite-intersections}.
The exact tensor comparison
$$
N\otimes_{\mathbf Z}(\mathbf Z\times\mathbf Q)
\longrightarrow N\oplus(N\otimes_{\mathbf Z}\mathbf Q),
\qquad n\otimes(a,q)\longmapsto(an,n\otimes q)
$$
has inverse
$$
\left(n,\sum_j n_j\otimes q_j\right)
\longmapsto n\otimes(1,0)+\sum_j n_j\otimes(0,q_j).
$$
Both are balanced and their composites are identities on
the displayed generators. This proves flatness of $S$ and
faithfulness: the tensor of a homomorphism $h$ has first
component $h$, so it is zero only if $h=0$.

The original annihilator is
$$
\operatorname{Ann}_{\mathbf Z}(M)
=\bigcap_{n\geq1}2^n\mathbf Z=0,
$$
because no nonzero integer is divisible by every $2^n$.
Each original $m\in M$ has finite support, hence is killed
by some $2^N$. Therefore
$m\otimes q=(2^Nm)\otimes(q/2^N)=0$ in
$M\otimes_{\mathbf Z}\mathbf Q$. The preceding comparison gives
$M\otimes_{\mathbf Z}S\simeq M\oplus0$, with the action
$(a,q)(m,0)=(am,0)$. It follows exactly that
$$
\operatorname{Ann}_{\mathbf Z}(M)S=0,
\qquad
\operatorname{Ann}_S(M\otimes_{\mathbf Z}S)=0\times\mathbf Q\ne0.
$$
Every finite family in $M$ is killed by a common nonzero
power of $2$, so no such family detects its zero annihilator.
This same module has support $\{(2)\}$: its first summand
has nonzero localization at $(2)$, while at every other
prime $2$ is a unit and every element is killed by a power
of $2$. Thus its support is closed but is strictly smaller
than $V(\operatorname{Ann}_{\mathbf Z}(M))=\Spec(\mathbf Z)$.

The counterexample even has a finitely presented faithfully
flat algebra version. Put $S'=\mathbf Z\times\mathbf Z[1/2]$.
The identical localization and direct-sum arguments prove
faithful flatness and give annihilator $0\times\mathbf Z[1/2]$.
More explicitly,
$$
\mathbf Z[T]/(2T^2-T)\longrightarrow S',
\qquad f(T)\longmapsto(f(0),f(1/2)),
$$
is an isomorphism. The ideals $(T)$ and $(2T-1)$ are
comaximal since $2T-(2T-1)=1$; hence their intersection
is their product $(2T^2-T)$. The Chinese remainder map
is surjective with this kernel. Its factors are $\mathbf Z$
and $\mathbf Z[1/2]$: in the second quotient $2T=1$, and
the inverse to evaluation at $1/2$ sends $a/2^j$ to $aT^j$.
This proves the asserted presentation with the original
integer factors retained.

Finally, faithful flatness does always imply the contraction
identity
$$
\varphi^{-1}\bigl(\operatorname{Ann}_S(M\otimes_RS)\bigr)
=\operatorname{Ann}_R(M).
$$
Multiplication by $r$ on $M$ becomes multiplication by
$\varphi(r)$ on its tensor module, and a faithfully flat
tensor functor detects whether this homomorphism is zero.
Thus the counterexample fails extension, while contraction
still recovers the original annihilator. For $M=0$ or the
zero ring, all the annihilator statements retain these meanings.
\end{proof}


\subsection{Editorial treatment of Algebra lines 9384--9410}

\hypertarget{algebra-context-100}{}

\noindent\href{https://github.com/stacks/stacks-project/blob/a04446e57ec1fbc252a871afcec7752fb2807b14/algebra.tex\#L9384}{Original source, lines 9384--9410}. Complete original and reviewed TeX passages accompany this note. Every intervention is separately recorded; the following treatment is editorial.


\begin{lemma}
\label{editorial-source-lemma-support-closed}
Let $R$ be a ring and let $M$ be an $R$-module.
If $M$ is finite, then $\text{Supp}(M)$ is closed.
More precisely, if $I = \text{Ann}(M)$ is the annihilator of $M$, then
$V(I) = \text{Supp}(M)$.
\end{lemma}

\begin{proof}
We will show that $V(I) = \text{Supp}(M)$.

\medskip\noindent
Suppose $\mathfrak p \in \text{Supp}(M)$. Then $M_{\mathfrak p} \not = 0$.
Choose an element $m \in M$ whose image in $M_\mathfrak p$ is nonzero.
Then the annihilator of $m$ is contained in $\mathfrak p$ by construction
of the localization $M_\mathfrak p$. Hence
a fortiori $I = \text{Ann}(M)$ must be contained in $\mathfrak p$.

\medskip\noindent
If $M=0$, then $I=R$ and both sets are empty.
Otherwise retain generators $x_1,\ldots,x_r$ with $r\geq1$.
Suppose $\mathfrak p\notin\operatorname{Supp}(M)$.
Each $x_i/1$ is zero in $M_{\mathfrak p}$, so
Lemma \ref{algebra-lemma-localization-colimit} gives a common
$f\notin\mathfrak p$ for which all are zero in $M_f$.
For each original generator choose $n_i\geq1$ with
$f^{n_i}x_i=0$, and keep $n=\max\{n_1,\ldots,n_r\}$.
The complete equality
$f^nx_i=f^{n-n_i}(f^{n_i}x_i)=0$ shows that $f^n$ kills
every generator and therefore $f^nM=0$.
Thus $f^n\in I$ while $f^n\notin\mathfrak p$, proving
$\mathfrak p\notin V(I)$ and the required equality.
\end{proof}

\begin{lemma}
\label{editorial-source-lemma-support-finite-annihilator-detection}
Let $M$ be an arbitrary $R$-module. If finitely many
$m_1,\ldots,m_n\in M$ satisfy
$\operatorname{Ann}_R(M)=\bigcap_{i=1}^n\operatorname{Ann}_R(m_i)$,
then $\operatorname{Supp}(M)=V(\operatorname{Ann}_R(M))$.
\end{lemma}

\begin{proof}
Retain the original submodule $N=\sum_{i=1}^n Rm_i\subseteq M$.
An element kills $N$ exactly when it kills each $m_i$,
so $\operatorname{Ann}_R(N)=\operatorname{Ann}_R(M)$ by
the stated hypothesis. Since $N$ is finite, the preceding
lemma gives $\operatorname{Supp}(N)=V(\operatorname{Ann}_R(M))$.
Localizing the actual injection $N\hookrightarrow M$
is injective, so $\operatorname{Supp}(N)\subseteq\operatorname{Supp}(M)$.
For every $M$, the first half of the preceding proof
gives $\operatorname{Supp}(M)\subseteq V(\operatorname{Ann}_R(M))$.
The inclusions prove equality without replacing the
original $M$ by $N$. If the list is empty its intersection
of annihilators is $R$, hence $M=0$, which also satisfies
the conclusion. For example an arbitrary nonempty free
module is covered by this hypothesis: one original basis
vector already has annihilator zero, although the basis
may be infinite.
\end{proof}


\subsection{Editorial treatment of Algebra lines 9412--9444}

\hypertarget{algebra-context-101}{}

\noindent\href{https://github.com/stacks/stacks-project/blob/a04446e57ec1fbc252a871afcec7752fb2807b14/algebra.tex\#L9412}{Original source, lines 9412--9444}. Complete original and reviewed TeX passages accompany this note. Every intervention is separately recorded; the following treatment is editorial.


\begin{lemma}
\label{editorial-source-lemma-support-base-change}
Let $R \to R'$ be a ring map and let $M$ be a finite $R$-module.
Then $\text{Supp}(M \otimes_R R')$ is the inverse image of
$\text{Supp}(M)$.
\end{lemma}

\begin{proof}
Write the original map as $\varphi:R\to R'$ and the induced
spectrum map as $f:\Spec(R')\to\Spec(R)$.
Fix an original prime $\mathfrak p'\subset R'$ and its
contraction $\mathfrak p=\varphi^{-1}(\mathfrak p')$.
The local map is
$$
R_{\mathfrak p}\longrightarrow R'_{\mathfrak p'},
\qquad r/t\longmapsto\varphi(r)/\varphi(t).
$$
All denominators are valid, because $t\notin\mathfrak p$
implies $\varphi(t)\notin\mathfrak p'$. Its inverse image
of $\mathfrak p'R'_{\mathfrak p'}$ is
$\mathfrak pR_{\mathfrak p}$: a fraction in the former
ideal has numerator in $\mathfrak p'$, and contraction
of that numerator is the stated condition in $R$.

For any original module $M$, not yet assumed finite, the
exact localization-tensor comparison is
$$
(M\otimes_RR')_{\mathfrak p'}
\longrightarrow M_{\mathfrak p}\otimes_{R_{\mathfrak p}}R'_{\mathfrak p'},
\qquad \frac{m\otimes s}{u}\longmapsto\frac m1\otimes\frac su.
$$
Its inverse is
$$
\frac mr\otimes\frac su
\longmapsto\frac{m\otimes s}{\varphi(r)u},
\qquad r\notin\mathfrak p,\quad u\notin\mathfrak p'.
$$
The forward map is induced by the balanced map to the
right-hand module and then by its invertible action of
every $u\notin\mathfrak p'$. For the inverse, if
$v(r'm-rm')=0$ with $v\notin\mathfrak p$, multiplying
the corresponding numerator difference by
$\varphi(v)\notin\mathfrak p'$ gives zero. This checks
the equivalence relation in $M_{\mathfrak p}$. If
$w(u's-us')=0$ with $w\notin\mathfrak p'$, the same
$w$ kills the tensor numerator difference, checking the
other localization relation. Additivity uses common
denominators and tensor additivity. For $a/t\in R_{\mathfrak p}$,
the two balanced expressions map respectively to
$$
\frac{am\otimes s}{\varphi(t)\varphi(r)u}
=\frac{m\otimes\varphi(a)s}{\varphi(r)\varphi(t)u}.
$$
This proves balancing with every original factor. Composing
in one direction gives the original fraction immediately;
in the other, $m/1\otimes s/(\varphi(r)u)=m/r\otimes s/u$.
Thus the maps are inverse. In particular, for every $M$
and every $R\to R'$,
$$
\operatorname{Supp}_{R'}(M\otimes_RR')
\subseteq f^{-1}\operatorname{Supp}_R(M).
$$

The original residue-field map is the injective map
$$
\iota:\kappa(\mathfrak p)\longrightarrow\kappa(\mathfrak p'),
\qquad
\frac{r+\mathfrak p}{t+\mathfrak p}
\longmapsto
\frac{\varphi(r)+\mathfrak p'}{\varphi(t)+\mathfrak p'}.
$$
It exists and is injective because $R/\mathfrak p\to
R'/\mathfrak p'$ is injective and every displayed
nonzero denominator stays nonzero. There are exact maps
\[
\begin{split}
(M\otimes_RR')\otimes_{R'}\kappa(\mathfrak p')
&\longrightarrow M\otimes_R\kappa(\mathfrak p'),\\
(m\otimes s)\otimes\lambda
&\longmapsto m\otimes\overline s\lambda,
\end{split}
\]
with inverse $m\otimes\lambda\mapsto(m\otimes1)\otimes\lambda$,
and
\[
\begin{split}
(M\otimes_R\kappa(\mathfrak p))
\otimes_{\kappa(\mathfrak p)}\kappa(\mathfrak p')
&\longrightarrow M\otimes_R\kappa(\mathfrak p'),\\
(m\otimes\alpha)\otimes\beta
&\longmapsto m\otimes\iota(\alpha)\beta,
\end{split}
\]
with inverse $m\otimes\beta\mapsto(m\otimes1)\otimes\beta$.
The tensor balancing relations show independence and give
both composites on the displayed generators. Extension
along the field inclusion detects a nonzero vector space:
if $B$ is a basis, the map from the direct sum of copies
of $\kappa(\mathfrak p')$ indexed by $B$, sending its
$b$-coordinate $\lambda$ to $b\otimes\lambda$, is an
isomorphism, with inverse given by the finite basis
expansion of each vector. Thus it is zero exactly when
$B$ is empty.

For completeness, for any local ring $(A,\mathfrak m)$
and module $N$, the map
$N\otimes_A A/\mathfrak m\to N/\mathfrak mN$,
$n\otimes\overline a\mapsto an+\mathfrak mN$, has inverse
$n+\mathfrak mN\mapsto n\otimes1$; balancing and the
relations $\mathfrak mN=0$ in both targets prove the claim.
Applying this after the localization comparisons gives the
original fibre quotient
$$
M\otimes_R\kappa(\mathfrak p)
\simeq M_{\mathfrak p}/\mathfrak pM_{\mathfrak p},
$$
and the corresponding quotient for $M\otimes_RR'$ at
$\mathfrak p'$. Explicitly,
$m\otimes(a/t+\mathfrak pR_{\mathfrak p})$ maps to
$am/t+\mathfrak pM_{\mathfrak p}$, with inverse
$m/t+\mathfrak pM_{\mathfrak p}\mapsto
m\otimes(1/t+\mathfrak pR_{\mathfrak p})$.

Now restore the original finiteness hypothesis. The module
$M_{\mathfrak p}$ is finite over the local ring, so
Nakayama's Lemma \ref{algebra-lemma-NAK} says that it is nonzero
exactly when this fibre quotient is nonzero. If
$\mathfrak p\in\operatorname{Supp}_R(M)$, the field
extension and the displayed fibre maps show that the
original quotient
$$
(M\otimes_RR')_{\mathfrak p'}/
\mathfrak p'(M\otimes_RR')_{\mathfrak p'}
$$
is nonzero. Therefore its module is nonzero. Together
with the already proved reverse inclusion, this proves
the stated lemma, keeping both original primes and all maps.

\medskip\noindent
There are two further exact forms. First, for an arbitrary
$M$ put
$$
F_R(M)=\{\mathfrak p\in\Spec(R):
M\otimes_R\kappa(\mathfrak p)\ne0\}.
$$
The field-extension argument above did not use finiteness
or flatness. Hence for every ring map it proves
$$
F_{R'}(M\otimes_RR')=f^{-1}F_R(M).
$$
The localization comparison always gives
$F_R(M)\subseteq\operatorname{Supp}_R(M)$.
Consequently localization support commutes with every
base change if and only if
$$
F_R(M)=\operatorname{Supp}_R(M).
$$
For sufficiency, a prime in the inverse image of support
is in the inverse image of $F_R(M)$, hence in the new
fibre set and therefore in the new support. The general
reverse inclusion was proved above. For necessity, apply
the asserted support equality to the particular original
map $R\to\kappa(\mathfrak p)$. The field's unique prime
contracts to $\mathfrak p$, and its module is nonzero
exactly when $M\otimes_R\kappa(\mathfrak p)\ne0$.
This proves both directions of the criterion.

This criterion strictly extends finite generation. An
arbitrary direct sum of finite modules satisfies it.
Localization commutes with direct sums by the maps
$(m_i)_i/s\mapsto(m_i/s)_i$ and their inverse obtained
from the finite sum of the localized inclusions.
For coordinates $m_i/s_i$ indexed by a finite set $E$,
the inverse has denominator $\prod_{i\in E}s_i$ and
$i$-th numerator $m_i\prod_{j\in E,\,j\ne i}s_j$.
Tensor products commute with these direct sums by
$(m_i)_i\otimes c\mapsto(m_i\otimes c)_i$, with inverse
given by the sum of the tensor products of the original
inclusions. Balancing verifies these maps, and both
composites are identities on elements of finite support.
A direct sum is nonzero exactly when one summand is;
Nakayama for each finite summand therefore proves the
criterion for the sum. For a nonzero ring, a free module
on an infinite set is not finite: any finite collection
of vectors involves only finitely many basis coordinates,
so cannot generate a basis vector outside their union.
It supplies the strictness example.

\medskip\noindent
Second, if the original map $R\to R'$ is flat, support
commutes with this base change for every $M$. The local
map $R_{\mathfrak p}\to R'_{\mathfrak p'}$ is flat:
tensor with $R'$ preserves injections, and subsequent
localization at $\mathfrak p'$ preserves them. For any
$R_{\mathfrak p}$-module $N$, the comparison
$N\otimes_RR'_{\mathfrak p'}\simeq
N\otimes_{R_{\mathfrak p}}R'_{\mathfrak p'}$
sends each original elementary tensor to itself.
Its inverse is balanced over $R_{\mathfrak p}$ since
every $t\notin\mathfrak p$ acts invertibly in both factors;
thus the comparison follows with no injectivity assumption
on the ring map.

A flat local ring map $(A,\mathfrak m)\to(B,\mathfrak n)$
detects nonzero arbitrary modules. Choose $x\ne0$ in such
a module $N$ and put $J=\operatorname{Ann}_A(x)$.
The map $A/J\hookrightarrow N$, $a+J\mapsto ax$, is
injective. Flatness gives $B/JB\hookrightarrow N\otimes_AB$.
Here $J\subseteq\mathfrak m$ and
$JB\subseteq\mathfrak mB\subseteq\mathfrak n$, so
$B/JB\ne0$. This proves faithful flatness of the local
map, and applies to the original localization map above.
Its tensor comparison proves
$$
M_{\mathfrak p}\ne0
\quad\Longleftrightarrow\quad
(M\otimes_RR')_{\mathfrak p'}\ne0
$$
for every $M$, proving the asserted flat-base-change extension.

Flatness of the module alone does not imply the universal
support criterion. Retain the example
$R=\mathbf Z$, $M=\mathbf Q$, $R'=\mathbf F_p$ for a
prime integer $p$, with the quotient map. The module
$\mathbf Q$ is flat by its exact localization comparison,
proved in Example \ref{editorial-source-example-flat-infinite-intersections}.
At every prime of $\mathbf Z$, its localization is still
$\mathbf Q$: the map sending an original fraction $q/s$
to the same rational number is inverse to $q\mapsto q/1$.
Thus its localization support is all of $\Spec(\mathbf Z)$.
But
$$
q\otimes\overline a=(q/p)\otimes p\overline a=0
$$
in $\mathbf Q\otimes_{\mathbf Z}\mathbf F_p$, so the
new support is empty, whereas the inverse image of the
old support is the nonempty spectrum of $\mathbf F_p$.
The fibre set is exactly $\{(0)\}$: all nonzero prime
fibres vanish as just proved, while multiplication gives
$\mathbf Q\otimes_{\mathbf Z}\mathbf Q\simeq\mathbf Q$,
with inverse $q\mapsto q\otimes1$. For a fraction $a/b$,
$q\otimes a/b=(qa/b)\otimes1$ because $b$ is invertible
on both sides of the tensor, proving the inverse formula.
Moreover the annihilator of $\mathbf Q$ is detected by
its element $1$. Thus finite annihilator detection does
not imply the universal support criterion. Empty spectra
and zero modules satisfy all the stated identities as well.
\end{proof}


\subsection{Editorial treatment of Algebra lines 9446--9458}

\hypertarget{algebra-context-102}{}

\noindent\href{https://github.com/stacks/stacks-project/blob/a04446e57ec1fbc252a871afcec7752fb2807b14/algebra.tex\#L9446}{Original source, lines 9446--9458}. Complete original and reviewed TeX passages accompany this note. Every intervention is separately recorded; the following treatment is editorial.


\begin{lemma}
\label{editorial-source-lemma-support-element}
Let $R$ be a ring, let $M$ be an $R$-module, and let $m \in M$.
Then $\mathfrak p \in V(\text{Ann}(m))$ if and only if
$m$ does not map to zero in $M_\mathfrak p$.
\end{lemma}

\begin{proof}
If $\operatorname{Ann}_R(m)\subseteq\mathfrak p$ and
$m/1=0$ in the original $M_{\mathfrak p}$, there would
be $s\notin\mathfrak p$ with $sm=0$. This contradicts
$s\in\operatorname{Ann}_R(m)\subseteq\mathfrak p$.
Conversely, if $\operatorname{Ann}_R(m)$ is not contained
in $\mathfrak p$, take an original annihilator
$s\notin\mathfrak p$. It becomes a unit at $\mathfrak p$
and kills $m/1$, so $m/1=0$.
The exact cyclic comparison underlying this argument is
the injection
$$
R/\operatorname{Ann}_R(m)\longrightarrow M,\qquad
r+\operatorname{Ann}_R(m)\longmapsto rm,
$$
whose image is the actual submodule $Rm$ and whose
inverse on that image sends $rm$ to the indicated class.
Equality of two such expressions is exactly an
annihilator relation, proving independence. Localization
gives the injection $(r+\operatorname{Ann}_R(m))/s\mapsto rm/s$
into the original $M_{\mathfrak p}$. Thus the cyclic
module comparison preserves, rather than replaces, $M$.
\end{proof}

\begin{lemma}
\label{editorial-source-lemma-support-arbitrary-cyclic-union}
For every original $R$-module $M$,
$$
\operatorname{Supp}(M)=\bigcup_{m\in M}V(\operatorname{Ann}_R(m)).
$$
Consequently supports of arbitrary modules are stable
under specialization, and every subset of $\Spec(R)$
stable under specialization is the support of an $R$-module.
\end{lemma}

\begin{proof}
A localization $M_{\mathfrak p}$ is nonzero exactly when
some original $m\in M$ has $m/1\ne0$: a nonzero fraction
$m/s$ has nonzero numerator image since $s/1$ is a unit.
Lemma \ref{editorial-source-lemma-support-element} gives the displayed union.
If $\mathfrak p\subseteq\mathfrak q$ and
$\operatorname{Ann}_R(m)\subseteq\mathfrak p$, the same
annihilator is contained in $\mathfrak q$. Hence support
is stable under specialization with the original element
as its witness.

Localization commutes with arbitrary direct sums by the
maps
$$
(\bigoplus_{\lambda\in\Lambda}M_\lambda)_{\mathfrak p}
\longrightarrow\bigoplus_{\lambda\in\Lambda}(M_\lambda)_{\mathfrak p},
\qquad (m_\lambda)_\lambda/s\longmapsto(m_\lambda/s)_\lambda.
$$
For finitely many nonzero coordinates indexed by a set
$E\subseteq\Lambda$, its inverse sends
$(m_\lambda/s_\lambda)_{\lambda\in E}$ to the vector
with numerator coordinate
$m_\lambda\prod_{\mu\in E,\mu\ne\lambda}s_\mu$
and common denominator $\prod_{\mu\in E}s_\mu$.
All factors are outside $\mathfrak p$.
Equivalently the inverse is the sum of the localized
original inclusions $M_\lambda\to\bigoplus M_\lambda$,
which proves independence of representations; these
formulas verify both composites. Thus the support of
a direct sum is the union of the supports of its summands.

For a specialization-stable subset $T\subseteq\Spec(R)$,
take the specified module
$M_T=\bigoplus_{\mathfrak p\in T}R/\mathfrak p$.
Each summand has support $V(\mathfrak p)$ by
Lemma \ref{editorial-source-lemma-support-closed}, since its annihilator
is exactly $\mathfrak p$. Their union is $T$: it contains
each original $\mathfrak p\in T$, and every prime
containing one of them belongs to $T$ by specialization
stability. The empty subset gives the zero direct sum.
\end{proof}


\subsection{Editorial treatment of Algebra lines 9460--9484}

\hypertarget{algebra-context-103}{}

\noindent\href{https://github.com/stacks/stacks-project/blob/a04446e57ec1fbc252a871afcec7752fb2807b14/algebra.tex\#L9460}{Original source, lines 9460--9484}. Complete original and reviewed TeX passages accompany this note. Every intervention is separately recorded; the following treatment is editorial.


\begin{lemma}
\label{editorial-source-lemma-support-finite-presentation-constructible}
Let $R$ be a ring and let $M$ be an $R$-module.
If $M$ is a finitely presented $R$-module, then $\text{Supp}(M)$ is a
closed subset of $\Spec(R)$ whose complement is quasi-compact.
\end{lemma}

\begin{proof}
Retain the original finite presentation
$$
R^{\oplus m}\xrightarrow{A}R^{\oplus n}
\xrightarrow{\pi}M\longrightarrow0,
$$
where $A=(a_{ri})$ is the original $n\times m$ matrix.
Write $e_1,\ldots,e_n$ for the original standard basis
of $R^{\oplus n}$ and $a_i$ for column $i$ of $A$.
For each sorted subset $S=\{i_1<\cdots<i_n\}$ of
$\{1,\ldots,m\}$, set
$$
A_S=(a_{i_1},\ldots,a_{i_n}),\qquad
\Delta_S=\det(A_S),\qquad
I=(\Delta_S:|S|=n).
$$
Thus $I$ is exactly the original ideal of maximal minors.

Fix $\mathfrak p\in\Spec(R)$ and retain
$K=\kappa(\mathfrak p)$.
Applying $R\to R_{\mathfrak p}\to K$ to every matrix
entry gives $A_{\mathfrak p}$ and $A(\mathfrak p)$.
The localized presentation is exact. In detail, the last
map sends $v/s$ to $\pi(v)/s$ and is surjective by lifting
the original numerator. If $\pi(v)/s=0$, there is
$u\notin\mathfrak p$ with $u\pi(v)=0$. Then $uv=Aw$
for some $w\in R^{\oplus m}$, and the full equality
$$
\frac vs=\frac{uv}{us}
=A_{\mathfrak p}\left(\frac w{us}\right)
$$
proves its kernel statement. No injectivity of
$R\to R_{\mathfrak p}$ is assumed.
The fibre quotient comparison is
$$
K^n/\operatorname{im}A(\mathfrak p)
\longrightarrow M\otimes_RK,
\qquad
[(\lambda_1,\ldots,\lambda_n)]
\longmapsto\sum_{r=1}^n\pi(e_r)\otimes\lambda_r.
$$
Its inverse sends $m\otimes\lambda$ to
$[\lambda\overline v]$ for a lift $\pi(v)=m$.
Changing $v$ by $Aw$ changes this vector by
$\lambda A(\mathfrak p)\overline w$, which is in the
matrix image. Additivity and $R$-balancing follow from
the original linear map $\pi$. Both composites are
identities on the displayed generators. Hence
$$
\dim_K(M\otimes_RK)=n-\operatorname{rank}_K A(\mathfrak p).
$$

For $n\geq1$, a minor $\Delta_S\notin\mathfrak p$
gives the explicit right inverse
$$
A_{\mathfrak p}
\left(E_S\frac{\operatorname{adj}(A_S)_{\mathfrak p}}
{\Delta_S/1}\right)=I_n,
$$
where $E_S:R^n\to R^m$ sends its $\ell$-th basis
vector to the $i_\ell$-th original basis vector.
Indeed $AE_S=A_S$ and the adjugate identity gives the
displayed product. Thus $M_{\mathfrak p}=0$.
If $M_{\mathfrak p}=0$, the localized map is surjective;
preimages of the original basis vectors reduce to a
right inverse for $A(\mathfrak p)$, so the fibre is zero.
Conversely, if the fibre is zero, choose
$C\in\operatorname{Mat}_{m\times n}(K)$ with
$A(\mathfrak p)C=I_n$. Cauchy--Binet gives
$$
1=\sum_{\substack{S\subseteq\{1,\ldots,m\}\\|S|=n}}
\overline{\Delta_S}\det(C_{S,\{1,\ldots,n\}}).
$$
At least one original minor therefore lies outside
$\mathfrak p$. This proves, with explicit maps, all
the original equivalences
$$
M_{\mathfrak p}\ne0
\quad\Longleftrightarrow\quad
M\otimes_RK\ne0
\quad\Longleftrightarrow\quad
\operatorname{rank}_K A(\mathfrak p)<n.
$$
In particular $\operatorname{Supp}(M)=V(I)$.

If $n=0$, the target is zero and $M=0$. There is one
$0\times0$ minor indexed by the empty subset, with
determinant $1$, so $I=R$ and both sets are empty.
If $m<n$, there are no maximal minors, so $I=0$.
Every fibre then has dimension at least $n-m>0$, giving
support all of $\Spec(R)$. For the zero ring there are
no primes; every module in the presentation is zero,
and the same support identities hold.

There are only finitely many original maximal minors,
and the complement is the exact finite union
$$
\Spec(R)\setminus\operatorname{Supp}(M)
=\bigcup_{\substack{S\subseteq\{1,\ldots,m\}\\|S|=n}}D(\Delta_S).
$$
Each $D(f)$ is quasi-compact. To recall the proof,
localization identifies $\Spec(R_f)$ with $D(f)$:
contraction is one map and
$\mathfrak p\mapsto\{r/f^a:r\in\mathfrak p,\ a\geq0\}$
is its inverse. The prime localization correspondence
gives these inverse maps, and $D(r/f^a)$ corresponds
to $D(r)\cap D(f)$, proving that they are homeomorphisms.
For any ring $B$, refine a cover of $\Spec(B)$ by
basic opens $D(b_\lambda)$. If their ideal were proper,
a maximal ideal containing it would be missed. Hence
$1=\sum_{j=1}^t c_jb_{\lambda_j}$ for finitely many
original cover elements, and those basic opens cover.
This proves quasi-compactness of every spectrum and
of $D(f)$. A finite union of quasi-compact subspaces
is quasi-compact, by restricting any cover to each
of the finitely many subspaces and joining their finite
subcovers. This proves the original assertion. The
empty union is allowed; for $n=0$ the union is $D(1)$.

\medskip\noindent
The presentation also proves the more precise ideal
bounds. Put $H=\operatorname{Ann}_R(M)$. For $n\geq1$,
$$
I\subseteq H,\qquad H^n\subseteq I.
$$
To prove the first, retain the original column indices
and every cofactor sign in the adjugate identity:
$$
\Delta_Se_r
=\sum_{\ell=1}^n(-1)^{r+\ell}
\det A_{\{1,\ldots,n\}\setminus\{r\},\,
S\setminus\{i_\ell\}}\ a_{i_\ell}.
$$
The remaining rows and columns have inherited increasing
order. In coordinate $k$ the right side is the cofactor
expansion along row $r$ of the matrix with that row
replaced by row $k$. It is $\Delta_S$ if $k=r$ and
zero if $k\ne r$, since then two rows coincide.
Applying $\pi$ proves $\Delta_S\pi(e_r)=0$ for all
$r$. These original elements generate $M$, so
$\Delta_S\in H$ and $I\subseteq H$.

For the second bound take arbitrary original
$t_1,\ldots,t_n\in H$. Since $\pi(t_je_j)=0$,
choose $b_j\in R^m$ with $Ab_j=t_je_j$.
Let $B=(b_{ij})$ be the $m\times n$ matrix with these
columns. Then $AB=\operatorname{diag}(t_1,\ldots,t_n)$.
The full ordered-column determinant expansion is
$$
t_1\cdots t_n=\det(AB)
=\sum_{u_1=1}^m\cdots\sum_{u_n=1}^m
\left(\prod_{j=1}^n b_{u_j,j}\right)
\det(a_{u_1},\ldots,a_{u_n}).
$$
Terms with repeated indices vanish because the determinant
has repeated columns. For distinct indices let
$S=\{i_1<\cdots<i_n\}$ be their sorted set and let
$u_j=i_{\sigma(j)}$. The determinant is
$\operatorname{sgn}(\sigma)\Delta_S$. Thus the full sum is
$$
\sum_{\substack{S=\{i_1<\cdots<i_n\}\\S\subseteq\{1,\ldots,m\}}}
\Delta_S
\left(\sum_{\sigma\in\mathfrak S_n}
\operatorname{sgn}(\sigma)\prod_{j=1}^n b_{i_{\sigma(j)},j}\right).
$$
For $\tau=\sigma^{-1}$ the inner product is
$\prod_{\ell=1}^n b_{i_\ell,\tau(\ell)}$, with the same
sign, so the exact Cauchy--Binet identity is
$$
t_1\cdots t_n
=\sum_{\substack{S\subseteq\{1,\ldots,m\}\\|S|=n}}
\Delta_S\det(B_{S,\{1,\ldots,n\}}).
$$
This proves that every product of $n$ possibly different
elements of $H$ is in $I$. Those products generate $H^n$,
so it proves the ideal-power statement, not only the
statement about individual $n$-th powers. If $m<n$
the sum is empty, giving $H^n=0$. If $n=0$, both $H$
and $I$ are $R$ and the convention $H^0=R$ retains the
two containments. They also hold over the zero ring.
For $n\geq1$ they imply $\sqrt I=\sqrt H$: one inclusion
uses $I\subseteq H$; for $h\in H$, the other uses
$h^n\in H^n\subseteq I$, followed by radicals.
Thus the exact support formula is also
$V(I)=V(\operatorname{Ann}_R(M))$.

The exponent $n$ is sharp uniformly. For a field $k$
take the original presentation $A=xI_n$ over $R=k[x]$,
with $m=n\geq1$. It gives
$$
M=(R/(x))^n,\qquad H=(x),\qquad I=(x^n).
$$
Here an element kills all basis quotient classes exactly
when it belongs to $(x)$, and the single maximal minor
is $x^n$. Thus $H^n=I$; for $n\geq2$ the polynomial
$x^{n-1}$ belongs to $H^{n-1}$ and not to $(x^n)$.
Also $I\subsetneq H$ for these $n$.

\medskip\noindent
There is an exact weaker condition for the topological
conclusion. For any original module $N$ and specified
ideal $L$ with $\operatorname{Supp}(N)=V(L)$, its
complement is quasi-compact if and only if
$$
\sqrt L=\sqrt{(f_1,\ldots,f_t)}
\quad\hbox{for finitely many }f_1,\ldots,f_t\in L.
$$
Indeed $D(L)=\bigcup_{f\in L}D(f)$, so quasi-compactness
gives a finite subcover of this specific cover. Taking
complements gives $V(L)=V((f_1,\ldots,f_t))$ and hence
the radical equality. Conversely that equality gives the
same finite union of quasi-compact basic opens.
The empty list is allowed, with generated ideal zero.
The radical characterization used here follows directly:
every prime containing $J$ contains $\sqrt J$. If
$a\notin\sqrt J$, the ring $(R/J)_a$ is nonzero, since
otherwise some original $a^N$ would be in $J$.
A maximal ideal of this localization contracts to a prime
containing $J$ and avoiding $a$. Consequently $\sqrt J$
is exactly the intersection of the primes containing $J$,
which proves the asserted equivalence of radicals and
closed sets without a finiteness assumption on $J$.

Finite presentation is not necessary for this topology.
Take the specified example
$$
R=k[x,y_1,y_2,\ldots],\qquad
I=(x^2,xy_i:i\geq1),\qquad M=R/I,
$$
where $k$ is a field. Here $I\subseteq(x)$ and $x^2\in I$,
while $R/(x)=k[y_1,y_2,\ldots]$ is a domain: two
nonzero polynomials belong to a polynomial ring in
finitely many variables over $k$, where their product
is nonzero. Thus $\sqrt I=(x)$.
The original cyclic quotient has support $V(I)$.
If $a\in I\setminus\mathfrak p$, the unit $a/1$
makes $(R/I)_{\mathfrak p}$ zero; if $I\subseteq\mathfrak p$,
this quotient maps onto the nonzero residue field.
Therefore the support is $V(x)$ and its complement
is $D(x)=D(x^2)$, a quasi-compact basic open. In the
weaker criterion one can take the actual original
element $f_1=x^2\in I$; there is no replacement of
the original ideal by its radical.

To prove that $I$ is not finitely generated, suppose
$I=(h_1,\ldots,h_r)$ and choose $N$ containing all
indices of variables $y_i$ occurring in any $h_j$.
Define $\rho:R\to k[x,z]$ by fixing $x$, sending
$y_{N+1}$ to $z$, and sending every other $y_i$ to
zero. Define also $\rho_0:R\to k[x]$ by fixing $x$
and sending all the $y_i$ to zero. Since $y_{N+1}$
does not occur in $h_j$, one has
$\rho(h_j)=\rho_0(h_j)\in(x^2)$: the last membership
follows by applying $\rho_0$ to the original generators
$x^2,xy_i$ of $I$. Every combination of the proposed
finite generators therefore maps into $(x^2)$.
But $xy_{N+1}\in I$ maps to $xz\notin(x^2)$.
Indeed its coefficient of $x^1$ over $k[z]$ is the
nonzero polynomial $z$, whereas every multiple of
$x^2$ has that coefficient zero. This is a contradiction.

Finally, for every ring and ideal, the cyclic module
$R/I$ is finitely presented exactly when $I$ is
finitely generated. A finite generating list gives
the presentation $R^t\to R\to R/I\to0$ with that
list as the first matrix. Conversely choose a finite
presentation with surjection $\pi:R^b\to R/I$ and
finitely generated kernel $K=(k_1,\ldots,k_t)$.
Let $q:R\to R/I$ be the original quotient. Choose
$v\in R^b$ with $\pi(v)=1+I$, and an $R$-linear
lift $f:R^b\to R$ of $\pi$ by choosing representatives
for the images of the original basis vectors.
Thus $qf=\pi$, $f(v)-1\in I$, and $f(k_j)\in I$.
For each original $a\in I$, the vector $av$ lies
in $K$, so $av=\sum_j c_jk_j$ for some coefficients.
Retaining both original terms gives the exact identity
$$
a=f(av)+a(1-f(v))
=\sum_{j=1}^t c_jf(k_j)+a(1-f(v)).
$$
Every listed generator is in $I$, and the identity
proves the reverse inclusion in
$$
I=(f(k_1),\ldots,f(k_t),1-f(v)).
$$
This covers $R/I=0$ as well, including $b=0$, where
$1-f(v)=1$ supplies the unit. Applying it to the
specified infinitely generated ideal proves that the
example's original $M$ is not finitely presented,
despite its quasi-compact support complement.
\end{proof}


\subsection{Editorial treatment of Algebra lines 9486--9508}

\hypertarget{algebra-context-104}{}

\noindent\href{https://github.com/stacks/stacks-project/blob/a04446e57ec1fbc252a871afcec7752fb2807b14/algebra.tex\#L9486}{Original source, lines 9486--9508}. Complete original and reviewed TeX passages accompany this note. Every intervention is separately recorded; the following treatment is editorial.


\begin{lemma}
\label{editorial-source-lemma-support-quotient}
Let $R$ be a ring, let $M$ be an $R$-module, and let $I\subseteq R$ be an ideal.
\begin{enumerate}
\item If $M$ is finite then the support
of $M/IM$ is $\text{Supp}(M) \cap V(I)$.
\item If $N \subset M$, then $\text{Supp}(N) \subset
\text{Supp}(M)$.
\item If $Q$ is a quotient module of $M$ then $\text{Supp}(Q) \subset
\text{Supp}(M)$.
\item If $0 \to N \to M \to Q \to 0$ is a short exact sequence
then $\text{Supp}(M) = \text{Supp}(Q) \cup \text{Supp}(N)$.
\end{enumerate}
\end{lemma}

\begin{proof}
Localization preserves the original injections, surjections
and short exact sequences. For (2), the map
$N_{\mathfrak p}\to M_{\mathfrak p}$ is injective, so
the first module can be nonzero only if the second is.
For (3), the map $M_{\mathfrak p}\to Q_{\mathfrak p}$
is surjective, so the same conclusion holds for a nonzero
quotient. In (4), the exact localized sequence shows
$M_{\mathfrak p}=0$ if and only if both
$N_{\mathfrak p}=0$ and $Q_{\mathfrak p}=0$, which is
precisely the stated union formula.

For (1), retain the exact quotient comparison
$$
(M/IM)_{\mathfrak p}\longrightarrow M_{\mathfrak p}/IM_{\mathfrak p},
\qquad (m+IM)/s\longmapsto m/s+IM_{\mathfrak p}.
$$
Its inverse sends $m/s+IM_{\mathfrak p}$ to $(m+IM)/s$.
The original quotient map localizes to a surjection with
kernel $(IM)_{\mathfrak p}=IM_{\mathfrak p}$: the equality
follows by writing finite sums of products over a common
denominator. This proves well-definedness and both
composites of the displayed maps. If $M_{\mathfrak p}=0$,
the quotient is zero. If some $a\in I$ lies outside
$\mathfrak p$, it is a unit in $R_{\mathfrak p}$ and
$IM_{\mathfrak p}=M_{\mathfrak p}$, so the quotient is
again zero. Conversely, if $M_{\mathfrak p}\ne0$ and
$I\subseteq\mathfrak p$, then $IR_{\mathfrak p}$ lies
in the maximal ideal of the local ring. The module
$M_{\mathfrak p}$ is finite, so Nakayama's
Lemma \ref{algebra-lemma-NAK} prevents
$M_{\mathfrak p}/IM_{\mathfrak p}$ from being zero.
These cases prove the exact intersection formula.
\end{proof}


\subsection{Editorial treatment of Algebra lines 9566--9586}

\hypertarget{algebra-context-105}{}

\noindent\href{https://github.com/stacks/stacks-project/blob/a04446e57ec1fbc252a871afcec7752fb2807b14/algebra.tex\#L9566}{Original source, lines 9566--9586}. Complete original and reviewed TeX passages accompany this note. Every intervention is separately recorded; the following treatment is editorial.


\begin{lemma}
\label{editorial-source-lemma-open-going-down}
Let $R \to S$ be a ring map. If the induced map
$\varphi : \Spec(S) \to \Spec(R)$ is open, then
$R \to S$ satisfies going down.
\end{lemma}

\begin{proof}
Retain $\mathfrak p\subseteq\mathfrak p'$ and the original
$\mathfrak q'$ lying over $\mathfrak p'$.
For $g\notin\mathfrak q'$, the open set $D(g)$ contains
$\mathfrak q'$, so its image is an open neighbourhood of
$\mathfrak p'$. This image contains a standard open
$D(f)$ with $f\notin\mathfrak p'$. Then
$f\notin\mathfrak p$, so it also contains $\mathfrak p$.
The fibre criterion in Lemma \ref{editorial-source-lemma-in-image} gives
$S_g\otimes_R\kappa(\mathfrak p)\ne0$ for every such $g$.

Here is the precise directed localization argument. Index
finite subsets $E\subseteq S\setminus\mathfrak q'$ by
inclusion and set $g_E=\prod_{g\in E}g$, with $g_\varnothing=1$.
For $E\subseteq F$ the map is
$$
S_{g_E}\longrightarrow S_{g_F},\qquad
\frac{s}{g_E^n}\longmapsto
\frac{s\,g_{F\setminus E}^{\,n}}{g_F^n}.
$$
It is the original localization map, since all factors
in the product $g_F$ become units. The map from their
colimit to $S_{\mathfrak q'}$ sends each fraction to the
same original fraction. It is surjective: $s/t$ comes
from $E=\{t\}$. If $s/g_E^n$ maps to zero, there is
$u\notin\mathfrak q'$ with $us=0$. At the stage
$F=E\cup\{u\}$, $u$ is a unit, so the image of $s$
and then that fraction is zero. This proves injectivity
and the representative formula for the inverse. It is
the comparison of Lemma \ref{algebra-lemma-localization-colimit}.

Tensor-colimit compatibility gives the unital ring
comparison
$$
S_{\mathfrak q'}\otimes_R\kappa(\mathfrak p)
\cong\colim_E(S_{g_E}\otimes_R\kappa(\mathfrak p)),
$$
sending a stage tensor $s/g_E^n\otimes c$ to the tensor
of the same fraction and $c$; the inverse follows by
representing the finitely many fractions at a common
stage, as in Lemma \ref{editorial-source-lemma-tensor-products-commute-with-limits}.
Every stage ring is nonzero by the previous paragraph,
because $g_E\notin\mathfrak q'$. Their colimit is nonzero:
if its unit were zero, it would become zero at a later
stage, but every transition is unital and that later
ring is nonzero.

Choose a prime of this nonzero fibre ring and pull it
back along $S\to S_{\mathfrak q'}\to
S_{\mathfrak q'}\otimes_R\kappa(\mathfrak p)$ to a prime
$\mathfrak q$ of $S$. Every element outside
$\mathfrak q'$ maps to a unit, so
$\mathfrak q\subseteq\mathfrak q'$.
The contraction to the field $\kappa(\mathfrak p)$ is
its zero prime, and the original map $R\to\kappa(\mathfrak p)$
has kernel $\mathfrak p$. Thus the contraction of
$\mathfrak q$ to $R$ is exactly $\mathfrak p$.

In fact this argument proves the exact weaker criterion:
$R\to S$ satisfies going down if and only if the image
of every principal open $D(g)\subseteq\Spec(S)$ is stable
under generalization. For sufficiency that stability
supplies $\mathfrak p$ in each image containing
$\mathfrak p'$, and the remaining argument just given
applies without openness. Conversely, if going down
holds and $\mathfrak q'\in D(g)$ lies over
$\mathfrak p'$, every $\mathfrak p\subseteq\mathfrak p'$
has a lift $\mathfrak q\subseteq\mathfrak q'$.
Since $g\notin\mathfrak q'$, also $g\notin\mathfrak q$,
so the lift still belongs to the original $D(g)$.
This proves necessity and both directions of the criterion.

Openness is stronger than this criterion. The flat map
$\mathbf Z\to\mathbf Q$ satisfies going down by
Lemma \ref{editorial-source-lemma-flat-going-down}, but the image of its
whole spectrum is $\{(0)\}$, which is not open in
$\Spec(\mathbf Z)$. Indeed any open neighbourhood of
$(0)$ contains $D(n)$ for some nonzero integer $n$.
A prime divisor $p$ of $|n|+1>1$ does not divide $n$,
so $(p)\in D(n)$ and $(p)\ne(0)$. This proves the
failure of openness. Flatness of $\mathbf Q$ was proved
in Example \ref{editorial-source-example-flat-infinite-intersections}.
\end{proof}


\subsection{Editorial treatment of Algebra lines 9588--9603}

\hypertarget{algebra-context-106}{}

\noindent\href{https://github.com/stacks/stacks-project/blob/a04446e57ec1fbc252a871afcec7752fb2807b14/algebra.tex\#L9588}{Original source, lines 9588--9603}. Complete original and reviewed TeX passages accompany this note. Every intervention is separately recorded; the following treatment is editorial.


\begin{lemma}
\label{editorial-source-lemma-going-up-down-specialization}
Let $R \to S$ be a ring map.
\begin{enumerate}
\item $R \to S$ satisfies going down if and only if
generalizations lift along the map $\Spec(S) \to \Spec(R)$,
see Topology, Definition \ref{topology-definition-lift-specializations}.
\item $R \to S$ satisfies going up if and only if
specializations lift along the map $\Spec(S) \to \Spec(R)$,
see Topology, Definition \ref{topology-definition-lift-specializations}.
\end{enumerate}
\end{lemma}

\begin{proof}
For the original map $\varphi:R\to S$, the spectrum map
sends a prime $\mathfrak q$ to $\varphi^{-1}(\mathfrak q)$.
The point of $\mathfrak p$ specializes to the point of
$\mathfrak p'$ exactly when $\mathfrak p\subseteq\mathfrak p'$,
because the closure of the point of $\mathfrak p$ is
$V(\mathfrak p)$. The same statement holds for primes of $S$.
Thus, given the original $\mathfrak p\subseteq\mathfrak p'$
and a prime $\mathfrak q'$ over $\mathfrak p'$, lifting
the generalization is precisely finding
$\mathfrak q\subseteq\mathfrak q'$ with contraction
$\mathfrak p$. This is the stated going-down condition,
in both directions. Given instead a prime $\mathfrak q$
over $\mathfrak p$, lifting the specialization is precisely
finding $\mathfrak q'\supseteq\mathfrak q$ with contraction
$\mathfrak p'$, the going-up condition, again in both
directions. These identifications agree exactly with
Topology, Definition \ref{topology-definition-lift-specializations}.
\end{proof}


\subsection{Editorial treatment of Algebra lines 9605--9615}

\hypertarget{algebra-context-107}{}

\noindent\href{https://github.com/stacks/stacks-project/blob/a04446e57ec1fbc252a871afcec7752fb2807b14/algebra.tex\#L9605}{Original source, lines 9605--9615}. Complete original and reviewed TeX passages accompany this note. Every intervention is separately recorded; the following treatment is editorial.


\begin{lemma}
\label{editorial-source-lemma-going-up-down-composition}
Suppose $R \to S$ and $S \to T$ are ring maps satisfying
going down. Then so does $R \to T$. Similarly for going up.
\end{lemma}

\begin{proof}
Write the original maps as $\varphi:R\to S$ and
$\psi:S\to T$. For going down, take
$\mathfrak p\subseteq\mathfrak p'$ in $R$ and a prime
$\mathfrak r'$ of $T$ contracting to $\mathfrak p'$.
Put $\mathfrak q'=\psi^{-1}(\mathfrak r')$.
Going down for $\varphi$ gives
$\mathfrak q\subseteq\mathfrak q'$ contracting to
$\mathfrak p$. Going down for $\psi$ then gives
$\mathfrak r\subseteq\mathfrak r'$ with
$\psi^{-1}(\mathfrak r)=\mathfrak q$. Its original
composite contraction is
$(\psi\circ\varphi)^{-1}(\mathfrak r)
=\varphi^{-1}(\mathfrak q)=\mathfrak p$.

For going up, take the same chain in $R$ and a prime
$\mathfrak r$ of $T$ contracting to $\mathfrak p$.
Set $\mathfrak q=\psi^{-1}(\mathfrak r)$.
Going up for $\varphi$ supplies
$\mathfrak q'\supseteq\mathfrak q$ contracting to
$\mathfrak p'$, and going up for $\psi$ supplies
$\mathfrak r'\supseteq\mathfrak r$ contracting to
$\mathfrak q'$. The same inverse-image composition gives
$(\psi\circ\varphi)^{-1}(\mathfrak r')=\mathfrak p'$.
These are the precise prime maps underlying
Lemma \ref{editorial-source-lemma-going-up-down-specialization} and
Topology, Lemma \ref{topology-lemma-lift-specialization-composition}.
\end{proof}


\subsection{Editorial treatment of Algebra lines 9617--9649}

\hypertarget{algebra-context-108}{}

\noindent\href{https://github.com/stacks/stacks-project/blob/a04446e57ec1fbc252a871afcec7752fb2807b14/algebra.tex\#L9617}{Original source, lines 9617--9649}. Complete original and reviewed TeX passages accompany this note. Every intervention is separately recorded; the following treatment is editorial.


\begin{lemma}
\label{editorial-source-lemma-image-stable-specialization-closed}
Let $R \to S$ be a ring map. Let $T \subset \Spec(R)$
be the image of $\Spec(S)$. If $T$ is stable under specialization,
then $T$ is closed. More precisely, for every ring map
$\varphi:R\to S$, without a specialization hypothesis,
$$
\overline T=V(\ker\varphi)
=\{\mathfrak p:\text{there is }\mathfrak p'\in T
\text{ with }\mathfrak p'\subseteq\mathfrak p\}.
$$
\end{lemma}

\begin{proof}
We retain both original proofs and all original rings.
Write the map as $\varphi:R\to S$ and put $I=\ker\varphi$.

\medskip\noindent
First proof. Let $\mathfrak p$ be a point of $\overline T$.
Every $f\notin\mathfrak p$ gives $D(f)\cap T\ne\varnothing$.
The prime localization correspondence identifies this
intersection with the image of $\Spec(S_{\varphi(f)})$,
so $S_{\varphi(f)}\ne0$.
The directed localization comparison with $S_{\mathfrak p}$
can be written without suppressing the original coefficient map.
Index finite subsets $E\subseteq R\setminus\mathfrak p$
by inclusion and put $f_E=\prod_{f\in E}f$.
For $E\subseteq F$, the transition is
$$
S_{\varphi(f_E)}\longrightarrow S_{\varphi(f_F)},\qquad
\frac{s}{\varphi(f_E)^n}\longmapsto
\frac{s\varphi(f_{F\setminus E})^n}{\varphi(f_F)^n}.
$$
The colimit map to $S_{\mathfrak p}=(R\setminus\mathfrak p)^{-1}S$
sends a fraction to the same fraction. It is surjective
because $s/\varphi(t)$ comes from $E=\{t\}$.
If an original numerator $s$ maps to zero, some
$u\notin\mathfrak p$ has $\varphi(u)s=0$; the numerator
then vanishes at the stage $E\cup\{u\}$ where
$\varphi(u)$ is a unit. Thus the map is also injective.
Every stage is nonzero, including $E=\varnothing$,
because $f_E\notin\mathfrak p$; their common unit
cannot die in the colimit, since that would make it
zero at a nonzero later stage. Hence $S_{\mathfrak p}\ne0$.

Choose a prime of this nonzero ring and contract it
to a prime $\mathfrak q\subset S$. All
$\varphi(t)$, $t\notin\mathfrak p$, are units in
$S_{\mathfrak p}$, so
$\mathfrak p'=\varphi^{-1}(\mathfrak q)\subseteq\mathfrak p$.
This is a point of the original image $T$.
Thus every point of $\overline T$ is a specialization
of a point of $T$. The reverse containment follows
because $\overline T$ is closed and contains $T$.

The same localization argument identifies the closed set
exactly. Every prime in $T$ contains $I$, so
$\overline T\subseteq V(I)$. Conversely, for
$I\subseteq\mathfrak p$, the ring $S_{\mathfrak p}$
cannot be zero: the equality $1/1=0$ would give
$t\notin\mathfrak p$ with $\varphi(t)1=0$,
contradicting $t\in I\subseteq\mathfrak p$.
Its prime again contracts to a point of $T$ contained
in $\mathfrak p$. This proves both displayed equalities.
If $T$ is stable under specialization, they imply
$T=V(I)$ and hence the asserted closedness.

\medskip\noindent
Second proof. Retain the original quotient map
$q:R\to R/I$ and the induced injective map
$\overline\varphi:R/I\to S$, $r+I\mapsto\varphi(r)$.
The spectrum map for $q$ is the homeomorphism onto
$V(I)$ given by contraction, with inverse
$\mathfrak p\mapsto\mathfrak p/I$.
For each $\mathfrak p\in V(I)$ choose a minimal prime
$\overline{\mathfrak p}_0\subseteq\mathfrak p/I$
of $R/I$. Existence follows from Zorn's lemma: the
intersection of a chain of primes contained in
$\mathfrak p/I$ is prime and is a lower bound.
Indeed if $ab$ belongs to that intersection and $a$
is absent from one chain member, primality forces
$b$ into that member and all smaller members, and
then into all larger members as well.
Lemma \ref{algebra-lemma-injective-minimal-primes-in-image}
gives a prime $\mathfrak q\subset S$ whose contraction
along $\overline\varphi$ is $\overline{\mathfrak p}_0$.
Along the original map its contraction is
$$
\varphi^{-1}(\mathfrak q)
=q^{-1}(\overline{\mathfrak p}_0)\subseteq\mathfrak p.
$$
Thus every point of $V(I)$ is again a specialization
of an original image point; no original ring has been
replaced in the comparison. If $S=0$, then $I=R$,
and all three sets in the statement are empty.
\end{proof}


\subsection{Editorial treatment of Algebra lines 9651--9679}

\hypertarget{algebra-context-109}{}

\noindent\href{https://github.com/stacks/stacks-project/blob/a04446e57ec1fbc252a871afcec7752fb2807b14/algebra.tex\#L9651}{Original source, lines 9651--9679}. Complete original and reviewed TeX passages accompany this note. Every intervention is separately recorded; the following treatment is editorial.


\begin{lemma}
\label{editorial-source-lemma-going-up-closed}
Let $R \to S$ be a ring map. The following are equivalent:
\begin{enumerate}
\item Going up holds for $R \to S$, and
\item the map $\Spec(S) \to \Spec(R)$ is closed.
\end{enumerate}
\end{lemma}

\begin{proof}
Write the original map as $\varphi:R\to S$ and its
spectrum map as $f$. Suppose first that $f$ is closed.
Given $\mathfrak p\subseteq\mathfrak p'$ and a prime
$\mathfrak q$ of $S$ with $\varphi^{-1}(\mathfrak q)=\mathfrak p$,
the closed set $V(\mathfrak q)$ contains $\mathfrak q$.
Its image is closed and contains $\mathfrak p$, so it
contains $\mathfrak p'$. Thus there exists
$\mathfrak q'\in V(\mathfrak q)$ with
$\varphi^{-1}(\mathfrak q')=\mathfrak p'$.
This says $\mathfrak q\subseteq\mathfrak q'$ and proves
going up. These are the exact prime maps in
Topology, Lemma \ref{topology-lemma-closed-open-map-specialization}.

Conversely suppose going up holds. For an original
ideal $I\subseteq S$, the image of $V(I)$ is stable
under specialization. Given a prime $\mathfrak q$
containing $I$ and a specialization of its contraction,
going up supplies a containing prime $\mathfrak q'$;
it still contains the original $I$ and maps to that
specialization. The quotient map $S\to S/I$ identifies
$\Spec(S/I)$ with $V(I)$ by
$\overline{\mathfrak q}\mapsto\{s:s+I\in\overline{\mathfrak q}\}$,
with inverse $\mathfrak q\mapsto\mathfrak q/I$.
The composite $R\to S/I$ has kernel $\varphi^{-1}(I)$.
Applying Lemma \ref{editorial-source-lemma-image-stable-specialization-closed}
to this composite proves not only that the image is
closed but also the exact formula
$$
f(V(I))=V(\varphi^{-1}(I)).
$$
This proves closedness for every original closed set.
The quotient map itself satisfies going up: a larger
prime of $S$ above one containing $I$ also contains
$I$, so its quotient supplies the lift. Together
with Lemma \ref{editorial-source-lemma-going-up-down-composition}, this
also verifies the composite-map argument in the original proof.

More generally, for any ring map and any ideal $I$
the same quotient comparison proves
$$
\overline{f(V(I))}=V(\varphi^{-1}(I)).
$$
Consequently going up is equivalent to the unclosed
image equality for every $I$. Necessity was proved
above; sufficiency implies every closed-set image is
closed, so the first paragraph supplies the prime lift.
\end{proof}


\subsection{Editorial treatment of Algebra lines 9681--9705}

\hypertarget{algebra-context-110}{}

\noindent\href{https://github.com/stacks/stacks-project/blob/a04446e57ec1fbc252a871afcec7752fb2807b14/algebra.tex\#L9681}{Original source, lines 9681--9705}. Complete original and reviewed TeX passages accompany this note. Every intervention is separately recorded; the following treatment is editorial.


\begin{lemma}
\label{editorial-source-lemma-constructible-stable-specialization-closed}
Let $R$ be a ring. Let $E \subset \Spec(R)$ be a constructible subset.
\begin{enumerate}
\item If $E$ is stable under specialization, then $E$ is closed.
\item If $E$ is stable under generalization, then $E$ is open.
\end{enumerate}
\end{lemma}

\begin{proof}
First proof. By Lemma \ref{editorial-source-lemma-constructible}, retain
a finite expression
$$
E=\bigcup_{i=1}^a\left(D(f_i)\cap
V(g_{i1},\ldots,g_{ib_i})\right).
$$
For $J_i=(g_{i1},\ldots,g_{ib_i})$, the original
quotient and localization maps give
$\Spec((R/J_i)_{f_i})$ the image $D(f_i)\cap V(J_i)$.
Contraction is the map; the inverse on this subset
sends $\mathfrak p$ to $(\mathfrak p/J_i)_{f_i}$.
The finite product $A=\prod_i(R/J_i)_{f_i}$ therefore
has spectrum image $E$ under the diagonal original
map $R\to A$, as in Lemma \ref{editorial-source-lemma-constructible-is-image}.
Indeed the coordinate idempotents sum to $1$ and
have pairwise zero products, so every prime selects
exactly one factor; this identifies its contraction
with that factor's contraction. The empty product
is the zero ring and gives the empty set.
If $E$ is stable under specialization, the preceding
image lemma proves it closed.

The complement $E^c$ is constructible by
Topology, Lemma \ref{topology-lemma-constructible}.
If $E$ is stable under generalization, $E^c$ is
stable under specialization: a specialization from
a point outside $E$ cannot enter $E$, since the
reverse generalization would then enter $E$ too.
Applying the first assertion to $E^c$ proves it
closed, hence $E$ open.
Each original basic term $D(f_i)\cap V(J_i)$ is
quasi-compact, being the image of the spectrum of
$(R/J_i)_{f_i}$. Therefore every constructible subset
of $\Spec(R)$ is quasi-compact by this finite union.
The complement is constructible as well. Thus in (1)
the closed set has quasi-compact complement, and in
(2) the open set is quasi-compact.

\medskip\noindent
Second proof, including its stronger scope. By
Lemma \ref{algebra-lemma-spec-spectral}, $X=\Spec(R)$ is
spectral. Its constructible topology has the
quasi-compact opens and their complements as a
subbasis; these sets are therefore both open and
closed in that topology. The topology is
quasi-compact Hausdorff by Topology, Lemma
\ref{topology-lemma-constructible-hausdorff-quasi-compact}.
Let $C\subseteq X$ be closed in the constructible
topology, which is weaker as a hypothesis than
being constructible, and let $x\in\overline C$
for the original topology. For every quasi-compact
open neighbourhood $U$ of $x$, the set $U\cap C$
is nonempty and is closed in the constructible
topology. Finite intersections of these $U$ are
again quasi-compact open neighbourhoods, so this
family has the finite intersection property.
Compactness gives $y\in C$ lying in every such $U$.
Since these sets form a neighbourhood basis of $x$,
$x$ is a specialization of $y$. Thus
$$
\overline C=\{x:\text{some }y\in C\text{ specializes to }x\}.
$$
If $C$ is stable under specialization it is closed
in the original topology. By taking complements,
any subset open in the constructible topology
and stable under generalization is open in the
original topology. This gives the full stronger
form of Topology, Lemma
\ref{topology-lemma-constructible-stable-specialization-closed},
as well as both original assertions. The stronger
scope does not assert that these sets or their
complements are constructible or quasi-compact in
all the cases just stated.
\end{proof}


\subsection{Editorial treatment of Algebra lines 9707--9737}

\hypertarget{algebra-context-111}{}

\noindent\href{https://github.com/stacks/stacks-project/blob/a04446e57ec1fbc252a871afcec7752fb2807b14/algebra.tex\#L9707}{Original source, lines 9707--9737}. Complete original and reviewed TeX passages accompany this note. Every intervention is separately recorded; the following treatment is editorial.


\begin{proposition}
\label{editorial-source-proposition-fppf-open}
Let $R \to S$ be flat and of finite presentation.
Then $\Spec(S) \to \Spec(R)$ is open.
More generally this holds for any ring map $R \to S$ of
finite presentation which satisfies going down.
\end{proposition}

\begin{proof}
If the original map $\varphi:R\to S$ is flat, it
satisfies going down by Lemma \ref{editorial-source-lemma-flat-going-down}.
Thus to prove the proposition it suffices to retain
finite presentation and going down.
For an original $g\in S$, write a finite presentation
$S=R[x_1,\ldots,x_n]/(h_1,\ldots,h_m)$ and choose
an original polynomial $G(x_1,\ldots,x_n)$ representing
$g$. Then the localization has the finite presentation
$$
S_g\simeq
R[x_1,\ldots,x_n,t]/(h_1,\ldots,h_m,tG-1).
$$
The map sends $x_i$ to its original class and $t$
to $1/g$. The inverse sends $s/g^d$ to the class
of $P(x)t^d$ for a polynomial lift $P$ of $s$.
Independence follows from the original relations
$h_j$ and the relation $tG=1$; these maps are
inverse on the algebra generators and fractions.
Hence Chevalley's Theorem \ref{editorial-source-theorem-chevalley}
applies to this particular $R\to S_g$ and shows
that the image $E_g$ of $D(g)$ is constructible.

Going down shows that $E_g$ is stable under
generalization. Given $\mathfrak q'\in D(g)$ over
$\mathfrak p'$ and $\mathfrak p\subseteq\mathfrak p'$
in $R$, take $\mathfrak q\subseteq\mathfrak q'$
over $\mathfrak p$. Since $g\notin\mathfrak q'$,
also $g\notin\mathfrak q$, and the lift remains
in the original $D(g)$. This is also the exact
going-down comparison for $S\to S_g$ and the
composition lemma. The preceding constructible-set
lemma now makes $E_g$ open. Every original open
of $\Spec(S)$ is a union of such $D(g)$ and images
commute with unions, proving openness of the map.

The proof only needs constructibility of the
images of all principal opens. Under that weaker
assumption, going down is equivalent to openness:
the proof just given gives one direction, and
Lemma \ref{editorial-source-lemma-open-going-down} gives the other.
In particular this is an equivalence for maps of
finite presentation. Finite presentation is not
necessary: for a field $k$, the map
$k\to k[x_i:i\geq1]$ has open spectrum map,
because the target spectrum $\Spec(k)$ has one
point. The algebra is not even of finite type:
finitely many polynomials use finitely many
variables and cannot generate an omitted original
indeterminate, as evaluation fixing those variables
and changing that indeterminate verifies.

For the flat finite-presentation case the map is
universally open. Indeed for every original
$R$-algebra $A$, the base change has presentation
$$
S\otimes_RA\simeq
A[x_1,\ldots,x_n]/(h_1^A,\ldots,h_m^A),
$$
where $h_j^A$ retains every coefficient through
the original map $R\to A$. The map sends
$s\otimes a$ to $a$ times a polynomial representative
of $s$; the inverse sends the original variables
to $x_i\otimes1$ and each $a$ to $1\otimes a$.
The displayed relations and tensor balancing make
these maps inverse. Flatness is preserved because
for an injection of $A$-modules $L\hookrightarrow L'$,
the original comparison
$$
L\otimes_A(S\otimes_RA)\simeq L\otimes_RS,
\qquad l\otimes(s\otimes a)\longmapsto al\otimes s,
$$
with inverse $l\otimes s\mapsto l\otimes(s\otimes1)$,
identifies its tensor map with an injection by
flatness of $S$ over $R$. Thus the proved proposition
applies after every base change. This last assertion
uses flatness, not merely the original going-down
hypothesis.
\end{proof}


\subsection{Editorial treatment of Algebra lines 9739--9772}

\hypertarget{algebra-context-112}{}

\noindent\href{https://github.com/stacks/stacks-project/blob/a04446e57ec1fbc252a871afcec7752fb2807b14/algebra.tex\#L9739}{Original source, lines 9739--9772}. Complete original and reviewed TeX passages accompany this note. Every intervention is separately recorded; the following treatment is editorial.


\begin{lemma}
\label{editorial-source-lemma-same-image}
Let $k$ be a field, and let $R$, $S$ be $k$-algebras.
Let $S' \subset S$ be a sub $k$-algebra, and let $f \in S' \otimes_k R$.
In the commutative diagram
$$
\xymatrix{
\Spec((S \otimes_k R)_f) \ar[rd] \ar[rr] & &
\Spec((S' \otimes_k R)_f) \ar[ld] \\
& \Spec(R) &
}
$$
the images of the diagonal arrows are the same.
\end{lemma}

\begin{proof}
Retain the original inclusion $i:S'\hookrightarrow S$,
the original $k$-algebra $R$, and the specified
$f\in S'\otimes_kR$. For a prime $\mathfrak p\subset R$
put $K=\kappa(\mathfrak p)$. The coefficient maps
$R\to K$ give elements $f_{S'}\in S'\otimes_kK$
and $f_S\in S\otimes_kK$, related by $i\otimes1$.
For either $B=S'$ or $B=S$, the exact comparison is
$$
((B\otimes_kR)_f)\otimes_RK
\longrightarrow (B\otimes_kK)_{f_B},
\qquad
\frac{\sum_j b_j\otimes r_j}{f^d}\otimes\lambda
\longmapsto
\frac{\sum_j b_j\otimes\overline{r_j}\lambda}{f_B^d}.
$$
Here the $f$ in the $B=S$ expression denotes its
image under the original inclusion, not a different
choice of element. The inverse sends
$$
\frac{b\otimes\lambda}{f_B^d}
\longmapsto \frac{b\otimes1}{f^d}\otimes\lambda
$$
and extends additively. To check these maps, first use
the balanced isomorphism
$(B\otimes_kR)\otimes_RK\simeq B\otimes_kK$,
$(b\otimes r)\otimes\lambda\mapsto b\otimes\overline r\lambda$,
whose inverse is $b\otimes\lambda\mapsto(b\otimes1)\otimes\lambda$.
On localizing, $f$ maps to $f_B$ and is inverted in
both targets. The universal property of localization
therefore extends the two inverse algebra maps to
the displayed fraction maps. In particular it checks
all zero-divisor witnesses as well as tensor balancing;
their composites are identities on the original
algebra generators and on $1/f$.

For a ring $C$ and $h\in C$, one has $C_h=0$
exactly when $h$ is nilpotent. Indeed $1/1=0$
holds exactly when $h^d\cdot1=0$ for some
$d\geq0$; in the case $d=0$ the ring is zero
and $h^1=0$ also holds. Conversely any such positive
power makes $h$ an invertible nilpotent, forcing the
localized unit to vanish. Thus the fibre criterion
of Lemma \ref{editorial-source-lemma-in-image} identifies the image
of each diagonal arrow with the primes for which
its original $f_B$ is not nilpotent.

Since $k$ is a field, tensoring the injection
$S'\hookrightarrow S$ with the $k$-vector space $K$
is injective. Explicitly, choose a $k$-basis of $K$.
Under the direct-sum identifications of both tensor
modules, the map is the direct sum of copies of the
original injection; hence its kernel is zero.
Consequently
$$
f_{S'}^d=0\quad\Longleftrightarrow\quad f_S^d=0
\qquad(d\geq1).
$$
The forward implication is functoriality; the reverse
uses this injection. Nonnilpotence is therefore the
same in both original fibres. The fibre criterion
and the fraction comparisons prove equality of the
two spectrum images. Independently, one inclusion
is also the factorization through the horizontal
arrow: every prime of the larger tensor algebra
contracts to a prime of the smaller one, with the
same contraction to $R$. The reverse inclusion is
the nonnilpotence argument just supplied.

\medskip\noindent
This proof gives a more general assertion with its
exact sufficient hypothesis. Let $A$ be any base
ring, let $B\to C$ be a map of $A$-algebras, and
let $R$ be an $A$-algebra. Suppose that for every
original $\mathfrak p\in\Spec(R)$ the map
$$
B\otimes_A\kappa(\mathfrak p)
\longrightarrow C\otimes_A\kappa(\mathfrak p)
$$
is injective. Then for every $f\in B\otimes_AR$,
the two original principal-localization spectrum
images in $\Spec(R)$ agree. All the preceding
fraction maps work with $A$ in place of $k$,
and injectivity of these particular maps supplies
exactly the nilpotence reflection used in the proof.
In particular the hypothesis holds when the
$A$-module map $B\to C$ is pure, meaning that
$B\otimes_AL\to C\otimes_AL$ is injective for
every $A$-module $L$. It suffices, more weakly,
that each of the displayed fibre maps reflect
nilpotence of the particular image of $f$.

An arbitrary injection over a general base ring
does not suffice. Take $A=R=B=\mathbf Z$,
$C=\mathbf Q$, and $f=1$. The spectrum image
for $B$ is all of $\Spec(\mathbf Z)$, whereas
the image for $C$ is $\{(0)\}$. At a nonzero
prime $(p)$ the receiving map of fibres is
$\mathbf F_p\to0$, since
$q\otimes\overline a=(q/p)\otimes p\overline a=0$.
This identifies the failed residue injection and
the failed image equality for the same original
maps and element. Finally, nonzero elements cannot
replace nonnilpotent elements in the proof: over
$k[\epsilon]/(\epsilon^2)$ the nonzero element
$\epsilon$ has empty principal open. Its square
is zero and its localization is the zero ring.
\end{proof}


\subsection{Editorial treatment of Algebra lines 9774--9791}

\hypertarget{algebra-context-113}{}

\noindent\href{https://github.com/stacks/stacks-project/blob/a04446e57ec1fbc252a871afcec7752fb2807b14/algebra.tex\#L9774}{Original source, lines 9774--9791}. Complete original and reviewed TeX passages accompany this note. Every intervention is separately recorded; the following treatment is editorial.


\begin{lemma}
\label{editorial-source-lemma-map-into-tensor-algebra-open}
Let $k$ be a field.
Let $R$ and $S$ be $k$-algebras.
The map $\Spec(S \otimes_k R) \to \Spec(R)$
is open.
\end{lemma}

\begin{proof}
Retain the original $k$-algebras $S,R$ and write the
specified tensor as a finite sum
$$
f=\sum_{i=1}^d s_i\otimes r_i\in S\otimes_kR.
$$
Set $S'=k[s_1,\ldots,s_d]\subseteq S$, keeping every
original coefficient $s_i$. Tensoring $S'\hookrightarrow S$
with $R$ is injective because $R$ is a $k$-vector space,
so the same displayed tensor is an element of
$S'\otimes_kR$ with its specified image $f$.
The $k$-algebra $S'$ is finitely presented: it is
a quotient of $k[x_1,\ldots,x_d]$, and this polynomial
ring is Noetherian by Lemma \ref{editorial-source-lemma-Noetherian-permanence}.
Hence the kernel of the original map $x_i\mapsto s_i$
has finitely many generators $h_1,\ldots,h_m$.
This also follows from
Lemma \ref{editorial-source-lemma-Noetherian-finite-type-is-finite-presentation}.

The base change retains the presentation
$$
S'\otimes_kR\simeq R[x_1,\ldots,x_d]/(h_1^R,\ldots,h_m^R),
$$
where every coefficient of $h_j^R$ is the image of
the corresponding coefficient of $h_j$ in $R$.
The map sends $s_i\otimes1$ to the class of $x_i$
and $1\otimes r$ to $r$, and its inverse sends the
original polynomial generators and coefficients back
to these tensors. The relations and balancing prove
both maps are well defined and inverse.
The algebra is flat over $R$. To see this directly,
choose a $k$-basis $(b_\lambda)$ of $S'$.
The map
$$
\bigoplus_\lambda R\longrightarrow S'\otimes_kR,
\qquad (a_\lambda)_\lambda\longmapsto
\sum_\lambda b_\lambda\otimes a_\lambda
$$
is an $R$-module isomorphism; its inverse is the
finite basis expansion in the first tensor factor.
Tensoring a free module preserves injections by
the corresponding direct-sum maps. Thus
Proposition \ref{editorial-source-proposition-fppf-open} applies to
$R\to S'\otimes_kR$ and makes the image of the
original principal open $D(f)$ in this spectrum
open. Lemma \ref{editorial-source-lemma-same-image}, with the exact
same $f$, identifies it with the image of $D(f)$
in $\Spec(S\otimes_kR)$. Principal opens form a
basis and images commute with unions, so the
original spectrum map is open, as claimed.

This also proves universal openness of the
original map. For an arbitrary original $R$-algebra
$A$, its base change is identified with
$\Spec(S\otimes_kA)\to\Spec(A)$ through the map
$$
(S\otimes_kR)\otimes_RA\longrightarrow S\otimes_kA,
\qquad (s\otimes r)\otimes a\longmapsto s\otimes\overline r a,
$$
whose inverse is $s\otimes a\mapsto(s\otimes1)\otimes a$.
Both maps respect the original algebra products
and balancing, and their composites fix all
elementary tensors. The result already proved
applies to this $k$-algebra $A$, giving openness
after every base change, without finite generation
of the full algebra $S$.

\medskip\noindent
One can make the principal-open image more explicit
without confusing nonzero and nonnilpotent elements.
For the fixed original list $s_1,\ldots,s_d$, there
exists an integer $N\geq1$, independent of $R$
and $r_1,\ldots,r_d$, such that the image of $D(f)$
is the union of the coefficient opens of the full
power $f^N$ in a $k$-basis of $S'$.
Here is a proof with all coefficients retained.

Put $A_0=k[T_1,\ldots,T_d]$ and
$F=\sum_{i=1}^d s_i\otimes T_i\in S'\otimes_kA_0$.
Choose a $k$-basis $(b_\lambda)_{\lambda\in\Lambda}$
of $S'$. For each ordered tuple $(i_1,\ldots,i_n)$,
write the complete finite expansion
$$
s_{i_1}\cdots s_{i_n}
=\sum_{\lambda\in\Lambda}
\mu_{i_1,\ldots,i_n,\lambda}b_\lambda,
\qquad \mu_{i_1,\ldots,i_n,\lambda}\in k.
$$
Then, with all ordered tuples and their repeated
terms retained, the coefficients of $F^n$ are
\[
\begin{split}
F^n&=\sum_{\lambda\in\Lambda}b_\lambda\otimes C_{n,\lambda}(T),\\
C_{n,\lambda}(T)
&=\sum_{i_1=1}^d\cdots\sum_{i_n=1}^d
\mu_{i_1,\ldots,i_n,\lambda}T_{i_1}\cdots T_{i_n}.
\end{split}
\]
For fixed $n$ only finitely many of these polynomials
are nonzero. Put $J_n=(C_{n,\lambda}:\lambda\in\Lambda)
\subseteq A_0$ and $Z_n=V(J_n)$. Extension of the
basis to every residue field shows that
$\mathfrak p\in Z_n$ exactly when the original
$F$ has $n$-th power zero in
$S'\otimes_k\kappa(\mathfrak p)$. Hence
$Z_n\subseteq Z_{n+1}$. The fibre criterion and
the tensor-localization maps of the preceding
lemma give
$$
\Spec(A_0)\setminus
\operatorname{image}\bigl(D(F)\to\Spec(A_0)\bigr)
=\bigcup_{n\geq1}Z_n.
$$
Call this complement $Z$. The proven openness
result, applied to the original algebra $S'$
and base $A_0$, says that $Z$ is closed.

The ring $A_0$ is Noetherian, so its closed set
$Z$ has finitely many irreducible components
$V(\mathfrak p_1),\ldots,V(\mathfrak p_e)$ with
their indicated generic points. This follows
by applying Lemma \ref{editorial-source-lemma-Noetherian-irreducible-components}
to the quotient defining $Z$. Each generic point
lies in some $Z_{n_j}$, because their union is
$Z$. Since $Z_{n_j}$ is closed, it contains
the full component $V(\mathfrak p_j)$.
Thus for $N=\max(1,n_1,\ldots,n_e)$,
$$
Z=Z_N.
$$
If $Z$ is empty, use $N=1$: every $Z_n$ is
already empty. This is a finite-component argument;
it does not infer stabilization of increasing
closed sets merely from Noetherianity.

For any original $R$ and $r_i$, use the original
map $A_0\to R$, $T_i\mapsto r_i$.
At a prime $\mathfrak p\subset R$, its contraction
$\mathfrak p_0\subset A_0$ induces an injective
field map $\kappa(\mathfrak p_0)\to\kappa(\mathfrak p)$.
The map
$S'\otimes_k\kappa(\mathfrak p_0)\to
S'\otimes_k\kappa(\mathfrak p)$ is injective by
extension of vector spaces, and carries the full
original $F$ to $f$. Therefore nilpotence is
equivalent in these two fibres, and by $Z=Z_N$
it is equivalent to the vanishing of the $N$-th
power. Its coefficients are exactly
$C_{N,\lambda}(r_1,\ldots,r_d)$, with the same
ordered-tuple formula. It follows that
$$
\operatorname{image}\bigl(D(f)\subseteq
\Spec(S\otimes_kR)\to\Spec(R)\bigr)
=\bigcup_{\lambda\in\Lambda}
D\bigl(C_{N,\lambda}(r_1,\ldots,r_d)\bigr).
$$
Only finitely many original coefficient polynomials
contribute, giving a finite union of basic opens
and hence a quasi-compact open image. The equality
for $S$ uses the same-image lemma after the equality
for its specified subalgebra $S'$. None of these
comparisons replaces $f$ by its linear nonzero
coefficient test: the required test uses the
proved power $N$. For instance in
$S'=k[\epsilon]/(\epsilon^2)$ with $f=\epsilon\otimes1$,
the first power is nonzero and has a unit coefficient,
while the actual principal open is empty; the
second power vanishes, giving the correct empty
union. The case $d=0$ has $F=f=0$ and can use
$N=1$, and zero tensor algebras give empty images
under the same formulas.
\end{proof}


\subsection{Editorial treatment of Algebra lines 9796--9845}

\hypertarget{algebra-context-114}{}

\noindent\href{https://github.com/stacks/stacks-project/blob/a04446e57ec1fbc252a871afcec7752fb2807b14/algebra.tex\#L9796}{Original source, lines 9796--9845}. Complete original and reviewed TeX passages accompany this note. Every intervention is separately recorded; the following treatment is editorial.


\begin{lemma}
\label{editorial-source-lemma-unique-prime-over-localize-below}
Let $R \to S$ be a ring map.
Let $\mathfrak p \subset R$ be a prime.
Assume that
\begin{enumerate}
\item there exists a unique prime $\mathfrak q \subset S$ lying over
$\mathfrak p$, and
\item either
\begin{enumerate}
\item going up holds for $R \to S$, or
\item going down holds for $R \to S$ and there is at most one prime
of $S$ above every prime of $R$.
\end{enumerate}
\end{enumerate}
Then $S_{\mathfrak p} = S_{\mathfrak q}$.
\end{lemma}

\begin{proof}
Write the original ring map as $\varphi:R\to S$.
The rings in the conclusion are the specified
localizations
$$
S_{\mathfrak p}=\varphi(R\setminus\mathfrak p)^{-1}S,
\qquad S_{\mathfrak q}=(S\setminus\mathfrak q)^{-1}S.
$$
Since $\varphi^{-1}(\mathfrak q)=\mathfrak p$,
the first multiplicative set is contained in the
second. Thus there is an original localization map
$j:S_{\mathfrak p}\to S_{\mathfrak q}$, sending
$s/\varphi(r)$ to the same fraction.
No injectivity of $\varphi$ or of these localization
maps is assumed.

A prime of $S_{\mathfrak p}$ corresponds exactly
to a prime $\mathfrak q'\subset S$ whose contraction
$\mathfrak p'=\varphi^{-1}(\mathfrak q')$ is contained
in $\mathfrak p$. A prime of $S_{\mathfrak q}$
corresponds exactly to a prime
$\mathfrak q'\subseteq\mathfrak q$. Both statements
follow from contraction and extension under
localization: the corresponding prime consists of
the fractions with numerator in $\mathfrak q'$,
and these maps are inverse because the denominator
set is disjoint from that prime.

Assume first (1) and (2)(a). For such a
$\mathfrak q'$ over $\mathfrak p'\subseteq\mathfrak p$,
going up supplies $\mathfrak q''\supseteq\mathfrak q'$
whose contraction is $\mathfrak p$. Uniqueness over
$\mathfrak p$ makes $\mathfrak q''=\mathfrak q$,
so $\mathfrak q'\subseteq\mathfrak q$.
Assume instead (1) and (2)(b). Going down applied
to $\mathfrak q$ gives
$\mathfrak q''\subseteq\mathfrak q$ over
$\mathfrak p'$. Uniqueness over $\mathfrak p'$
makes $\mathfrak q''=\mathfrak q'$, with the same
conclusion. This argument uses uniqueness only
over the primes $\mathfrak p'\subseteq\mathfrak p$;
the global uniqueness assumption in (2)(b) can
therefore be weakened to precisely those primes.

It follows in either case that for every
$u\in S\setminus\mathfrak q$, the element $u/1$
belongs to no prime of $S_{\mathfrak p}$: every
such prime contracts to a prime contained in
$\mathfrak q$. Hence $u/1$ is a unit. Indeed a
nonunit generates a proper ideal lying in a
maximal ideal, which is prime, a contradiction.
Write its inverse as $v/\varphi(r)$ with
$r\notin\mathfrak p$. The equality
$(u/1)(v/\varphi(r))=1$ means that for some
$t\notin\mathfrak p$,
$$
\varphi(t)(uv-\varphi(r))=0\quad\hbox{in }S.
$$
Thus the inverse to $j$ has the full fraction
formula
$$
k:S_{\mathfrak q}\longrightarrow S_{\mathfrak p},
\qquad
\frac su\longmapsto\frac{sv}{\varphi(r)},
$$
where $v,r$ satisfy that original inverse identity.
Independence of this choice follows from uniqueness
of the inverse of $u/1$. Independence of fraction
representatives can also be checked explicitly:
if $w(u's-us')=0$ for $w\notin\mathfrak q$,
then $w/1,u/1,u'/1$ are all units in
$S_{\mathfrak p}$, so multiplication by their
inverses gives $s/u=s'/u'$ there. Addition and
multiplication follow by the same common-denominator
identities. Consequently $k$ is the localization
extension of the original map $S\to S_{\mathfrak p}$.
Both composites fix $S$ and all the specified inverse
denominators, hence are identities. This proves
$S_{\mathfrak p}\simeq S_{\mathfrak q}$ through
the original canonical map, giving the equality
notation in the statement its exact meaning.

In fact for any original $\varphi$, $\mathfrak p$
and $\mathfrak q$ with
$\varphi^{-1}(\mathfrak q)=\mathfrak p$, the
following conditions are equivalent:
$$
\begin{gathered}
j:S_{\mathfrak p}\longrightarrow S_{\mathfrak q}
\text{ is an isomorphism};\\
u/1\text{ is a unit in }S_{\mathfrak p}
\text{ for every }u\notin\mathfrak q;\\
\varphi^{-1}(\mathfrak q')\subseteq\mathfrak p
\Longrightarrow\mathfrak q'\subseteq\mathfrak q
\text{ for every prime }\mathfrak q'\subset S.
\end{gathered}
$$
The first implies the second because these elements
are units in $S_{\mathfrak q}$. The second implies
the third since a prime of $S_{\mathfrak p}$ cannot
contain a unit; contraction then avoids every
element outside $\mathfrak q$. The third implies
the second by the maximal-ideal argument, and the
second supplies the proved inverse fraction map.
This is the exact criterion used by both cases
of the original lemma. It is specific to the
specified further localization; a bijective
spectrum map alone has not been substituted for
an isomorphism of rings.
\end{proof}


\subsection{Editorial treatment of Algebra lines 9851--9882}

\hypertarget{algebra-context-115}{}

\noindent\href{https://github.com/stacks/stacks-project/blob/a04446e57ec1fbc252a871afcec7752fb2807b14/algebra.tex\#L9851}{Original source, lines 9851--9882}. Complete original and reviewed TeX passages accompany this note. Every intervention is separately recorded; the following treatment is editorial.


\begin{lemma}
\label{editorial-source-lemma-going-down-flat-module}
Let $R \to S$ be a ring map. Let $N$ be a finite $S$-module flat over $R$.
Endow $\text{Supp}(N) \subset \Spec(S)$ with the induced topology.
Then generalizations lift along $\text{Supp}(N) \to \Spec(R)$.
\end{lemma}

\begin{proof}
Write the original map as $\varphi:R\to S$. The
statement means the following. Given original primes
$\mathfrak p\subseteq\mathfrak p'\subset R$ and
$\mathfrak q'\in\operatorname{Supp}_S(N)$ with
$\varphi^{-1}(\mathfrak q')=\mathfrak p'$, there
must be a prime $\mathfrak q\subseteq\mathfrak q'$
such that $\varphi^{-1}(\mathfrak q)=\mathfrak p$
and $\mathfrak q\in\operatorname{Supp}_S(N)$.
The contraction condition on $\mathfrak q'$ is
part of the lifting statement and is retained
explicitly here. For the induced topology on
support, specialization between two of its points
is the original prime inclusion: closure in a
subspace is the intersection with the original
closure. Thus these prime conditions are exactly
the required generalization lift.

Put
$$
A=R_{\mathfrak p'},\qquad C=S_{\mathfrak q'},
\qquad E=N_{\mathfrak q'},\qquad
\mathfrak m=\mathfrak p'R_{\mathfrak p'}.
$$
Keep the local ring map
$A\to C$, $r/t\mapsto\varphi(r)/\varphi(t)$.
Its denominators are valid because
$t\notin\mathfrak p'$ implies
$\varphi(t)\notin\mathfrak q'$.
The module $E$ is flat over $A$, by
Lemma \ref{editorial-source-lemma-flat-localization}. One may also
see the original comparison directly: tensor an
injection of $A$-modules with the $R$-flat $N$
and then localize the resulting $S$-modules at
$\mathfrak q'$. Exactness of localization preserves
the injection, and the maps
$$
(L\otimes_RN)_{\mathfrak q'}\longrightarrow L\otimes_RE,
\qquad (l\otimes n)/s\longmapsto l\otimes n/s
$$
and their inverses identify the result with
tensoring by $E$. Here every original element of
$R\setminus\mathfrak p'$ is invertible on both
$L$ and $E$, so balancing already factors uniquely
through $A$. This gives the same $A$-flatness.

The module $E$ is nonzero by the chosen support
point, and is finite over the local ring $C$.
Nakayama's Lemma \ref{algebra-lemma-NAK} gives
$E/\mathfrak q' E\ne0$. Since the image of
$\mathfrak m$ lies in $\mathfrak q'C$, the quotient
map
$$
E/\mathfrak mE\longrightarrow E/\mathfrak q'E,
\qquad e+\mathfrak mE\longmapsto e+\mathfrak q'E
$$
is surjective. Consequently
$E\otimes_A\kappa(\mathfrak p')=E/\mathfrak mE\ne0$.
The quotient-tensor identification sends
$e\otimes\overline a$ to $ae+\mathfrak mE$,
with inverse $e+\mathfrak mE\mapsto e\otimes1$.

The flat $A$-module $E$ is therefore faithfully
flat. For a direct verification, if an arbitrary
$A$-module $L$ has an element $l\ne0$, set
$H=\operatorname{Ann}_A(l)\subseteq\mathfrak m$.
The injection $A/H\hookrightarrow L$,
$a+H\mapsto al$, becomes the injection
$E/HE\hookrightarrow L\otimes_AE$ by flatness.
The source surjects onto $E/\mathfrak mE\ne0$,
so $L\otimes_AE\ne0$. This proves nonzero-module
detection and, together with flatness, faithful
flatness, in agreement with Lemma \ref{editorial-source-lemma-ff}.
In particular for the nonzero $A$-module
$K=\kappa(\mathfrak p)$, defined using the original
chain $\mathfrak p\subseteq\mathfrak p'$, one has
$$
P=E\otimes_AK\ne0,
\qquad B=C\otimes_AK.
$$

Here are the full localization and quotient maps
in this fibre. Let $W=R\setminus\mathfrak p$,
$D=W^{-1}C$ and $Q=W^{-1}E$, using the images of
$W$ under the original map. Then
$$
B\simeq D/\mathfrak pD,\qquad
P\simeq Q/\mathfrak pQ.
$$
The ring comparison sends
$$
c\otimes\frac{a+\mathfrak p}{w+\mathfrak p}
\longmapsto\frac{\varphi(a)c}{\varphi(w)}+\mathfrak pD,
$$
where coefficients in $C$ are interpreted via its
original map from $R$. Its inverse sends
$c/\varphi(w)+\mathfrak pD$ to
$c\otimes(1/(w+\mathfrak p))$.
The module comparison uses the same formulas
with $e\in E$ in place of $c$. The original
relations of localization are respected because
each image of $w\in W$ becomes a unit in $K$.
The original ideal $\mathfrak p$ is killed in
both targets. Thus the universal properties of
localization and quotient, followed by tensor
balancing over $A$, give these maps. Their
composites fix every original coefficient,
element of $E$, and inverse denominator, proving
the assertions. These expressions retain the
possibly noninjective maps and all denominators.

Since $N$ is finite over $S$, the localized
$E$ is finite over $C$, and its original finite
generators tensor to generators of $P$ over $B$.
Put $J=\operatorname{Ann}_B(P)$, the original
annihilator ideal in the source proof.
Because $P\ne0$, $J$ is proper. Choose a prime
$\mathfrak b$ of $B$ containing $J$.
The finite-module support equality of
Lemma \ref{editorial-source-lemma-support-closed} gives
$P_{\mathfrak b}\ne0$.
Equivalently this prime is obtained from a prime
of the original quotient $B/J$ by inverse image.
Write
$$
\theta:S\longrightarrow C\longrightarrow B,
\qquad s\longmapsto (s/1)\otimes1,
\qquad \mathfrak q=\theta^{-1}(\mathfrak b).
$$
It is prime. Every $s\notin\mathfrak q'$ is a
unit in $C$ and hence in $B$, so
$\mathfrak q\subseteq\mathfrak q'$.
Every $r\in\mathfrak p$ maps to zero, while
every $r\notin\mathfrak p$ maps to the nonzero
scalar $r+\mathfrak p$ in the field $K$, which
is a unit in the nonzero ring $B$. The map
$K\to B$ is unital and injective, since a
nonzero field has no proper nonzero kernel.
Thus
$\varphi^{-1}(\mathfrak q)=\mathfrak p$ exactly.

To prove support membership through the original
module, use the balanced isomorphism
$$
N\otimes_SB\longrightarrow P,
\qquad n\otimes(c\otimes a)\longmapsto(cn/1)\otimes a,
$$
with inverse $(n/s)\otimes a\mapsto
n\otimes((1/s)\otimes a)$ for $s\notin\mathfrak q'$.
All such $s$ are inverted in $C$, so localization
and tensor balancing make these maps independent
of representatives; both composites fix the
original elements and scalar tensors. Localizing
at $\mathfrak b$ gives the further comparison
$$
P_{\mathfrak b}\simeq
N_{\mathfrak q}\otimes_{S_{\mathfrak q}}B_{\mathfrak b}.
$$
The map on an elementary fraction sends
$(n\otimes b)/u$ to $n/1\otimes b/u$.
Its inverse sends $n/s\otimes b/u$ to
$(n\otimes b)/(\theta(s)u)$, where
$s\notin\mathfrak q$ implies
$\theta(s)\notin\mathfrak b$. These formulas
respect both localization witness relations and
the scalar balancing, and are inverse on the
generators. Since the left side is nonzero,
$N_{\mathfrak q}$ is nonzero. This proves the
required lift in the original support.

\medskip\noindent
The exact places where finiteness was used give
a weaker sufficient hypothesis. Suppose now that
$N$ is any $S$-module flat over $R$ and that for
every original $\mathfrak q'\in\operatorname{Supp}_S(N)$,
with $\mathfrak p'=\varphi^{-1}(\mathfrak q')$, one has
$$
N_{\mathfrak q'}\otimes_{R_{\mathfrak p'}}
\kappa(\mathfrak p')\ne0.
$$
The argument above still proves $A$-flatness of
$E$; the displayed hypothesis replaces the
Nakayama step, and gives its faithful flatness.
Hence the same original fibre $P$ is nonzero
for every $\mathfrak p\subseteq\mathfrak p'$.
It need not be finite. Choose an original
$z\ne0$ in $P$, and a maximal ideal
$\mathfrak b\supseteq\operatorname{Ann}_B(z)$.
Then $z/1\ne0$ in $P_{\mathfrak b}$: a witness
$u\notin\mathfrak b$ with $uz=0$ would contradict
the containment of its annihilator. The same
prime contractions and comparison maps now give
the lift. This replaces the finite-module use
of $\operatorname{Ann}_B(P)$ by the exact cyclic
support of the chosen original element.

In particular it suffices that the localization
support equal the residue-fibre set over $S$:
$$
\operatorname{Supp}_S(N)
=\{\mathfrak q:N\otimes_S\kappa(\mathfrak q)\ne0\}.
$$
At a support point this supplies
$E/\mathfrak q'E\ne0$, and the same quotient
surjection supplies the required nonzero fibre
over $R_{\mathfrak p'}$. Every arbitrary direct
sum of finite $S$-modules satisfies this equality,
as proved in the support base-change comparison;
when that sum is also $R$-flat, the lifting
conclusion therefore follows without finite
generation of $N$. This is a sufficient hypothesis,
not a claim that such fibre nonvanishing is
necessary for every possible support lift.

\medskip\noindent
There is also an openness consequence. Suppose the
original algebra $S$ is finitely presented over $R$,
and the original $S$-module $N$ is finitely presented
and $R$-flat. Then the map on its support is open,
and the corresponding support map remains open
after every base change $R\to A_1$.
To prove this with the specified objects, retain
a finite presentation $S^b\xrightarrow{H}S^a\to N\to0$.
Let $J_H$ be the ideal generated by the original
$a\times a$ minors of $H$, using determinant $1$
in size zero. The full presentation calculation in
Lemma \ref{editorial-source-lemma-support-finite-presentation-constructible}
gives $\operatorname{Supp}_S(N)=V(J_H)$, with the
quotient-spectrum homeomorphism
$\Spec(S/J_H)\to V(J_H)$ given by contraction.
The ideal $J_H$ is finitely generated. Appending
polynomial representatives of its finitely many
specified generators to a finite presentation of
$S$ over $R$ proves that $S/J_H$ is finitely
presented over $R$. The going-down result just
proved lifts every prime generalization inside
$V(J_H)$; under the quotient homeomorphism this
is exactly going down for $R\to S/J_H$.
Proposition \ref{editorial-source-proposition-fppf-open} therefore
proves openness of the original support map.

For the arbitrary original base change $R\to A_1$,
put $S_1=S\otimes_RA_1$ and $N_1=N\otimes_RA_1$.
Tensoring the displayed presentation and its
coefficient maps gives
$S_1^b\xrightarrow{H_1}S_1^a\to N_1\to0$,
where every entry of $H_1$ is the original entry
of $H$ tensored with $1$. The full determinant
polynomial, with all its signs, shows that its
minor ideal is $J_HS_1$. Thus
$$
\operatorname{Supp}_{S_1}(N_1)=V(J_HS_1),
\qquad
(S/J_H)\otimes_RA_1\simeq S_1/J_HS_1.
$$
The last map sends $(s+J_H)\otimes a$ to
$s\otimes a+J_HS_1$; its inverse is the
original quotient map followed by tensoring.
The ideal generators are killed on both sides,
so these are inverse maps with the indicated
formulas. The algebra $S_1$ is finitely presented
over $A_1$ by the base-change presentation.
The module $N_1$ is $A_1$-flat: for an injection
of $A_1$-modules $L\hookrightarrow L'$, the maps
$$
L\otimes_{A_1}N_1\simeq L\otimes_RN,
\qquad l\otimes(n\otimes a)\longmapsto al\otimes n,
$$
with inverse $l\otimes n\mapsto l\otimes(n\otimes1)$,
identify the tensor map with an injection by
the original $R$-flatness of $N$.
Applying the proven support argument to these
original base-change objects proves openness
for every $A_1$. Equivalently, the quotient map
$\Spec(S/J_H)\to\Spec(R)$ is universally open,
even though the argument did not require
$S/J_H$ itself to be flat over $R$.

Flatness alone, and even flatness together with
finite annihilator detection, does not suffice.
Retain the original ring and module
$$
R=S=k[x,y]/(xy),\qquad \varphi=\operatorname{id}_R,
\qquad N=R_x,
$$
where $k$ is a field and $x,y$ denote their
specified quotient classes. The module is flat
by exactness of localization. Its annihilator
is exactly $(y)$, and is detected by $1\in R_x$.
Indeed $y$ kills the entire localization since
$xy=0$ and $x$ is inverted. Conversely, if
$r\cdot1=0$ in $R_x$, then $x^nr=0$ for some
$n\geq0$. Modulo $(y)$ this says
$x^n\overline r=0$ in the domain $k[x]$,
so $r\in(y)$. The same condition is equivalent
to $r$ killing every element of $R_x$ because
all its elements are fractions with powers
of $x$ as denominators.

The exact comparison
$$
R_x\longrightarrow k[x,x^{-1}],\qquad
r/x^n\longmapsto r(x,0)/x^n,
$$
is an isomorphism, with inverse given by the
original class of $x$ and its inverse; $y$ is
zero in $R_x$ by $xy=0$. Thus as an $R$-module,
$N$ has support $V(y)$. To verify both directions,
if $y\notin\mathfrak a$, its unit image in
$R_{\mathfrak a}$ kills $N_{\mathfrak a}$,
so that localization is zero. If
$y\in\mathfrak a$, the prime
$\overline{\mathfrak a}=\mathfrak a/(y)$
in $k[x]$ identifies $N_{\mathfrak a}$ with
the localization of $k[x,x^{-1}]$ at the
images of $k[x]\setminus\overline{\mathfrak a}$.
These images are all nonzero in a domain,
so the localized ring is nonzero. The comparison
sends each original fraction to the same
fraction after $y=0$; the inverse follows from
the two specified localization maps.
Its fibre set, however, is exactly $D(x)$:
$$
N\otimes_R\kappa(\mathfrak a)
\simeq\kappa(\mathfrak a)_{\overline x}.
$$
This is zero if $x\in\mathfrak a$ and is the
nonzero field otherwise, with fraction maps
$r/x^n\otimes c\mapsto\overline r c/\overline x^n$
and the tensor inverse. For primes avoiding
$x$, the relation $xy=0$ forces $y$ into the
prime, consistently with $D(x)\subseteq V(y)$.

Take the original primes
$\mathfrak p=(x)\subset\mathfrak p'=(x,y)$
and $\mathfrak q'=\mathfrak p'$. The latter
is in $V(y)=\operatorname{Supp}(N)$, but the
former is not, because $y\notin(x)$.
Since the spectrum map for $\varphi$ is the
identity, the only possible lift over
$\mathfrak p$ is $\mathfrak p$ itself, which
is outside support. This proves the failure
of generalization lifting despite flatness
and finite annihilator detection. At
$\mathfrak q'$ the required base residue fibre
is zero, exactly identifying the failed
hypothesis in the extended proof.
\end{proof}


\subsection{Editorial treatment of Algebra lines 9933--9942}

\hypertarget{algebra-context-116}{}

\noindent\href{https://github.com/stacks/stacks-project/blob/a04446e57ec1fbc252a871afcec7752fb2807b14/algebra.tex\#L9933}{Original source, lines 9933--9942}. Complete original and reviewed TeX passages accompany this note. Every intervention is separately recorded; the following treatment is editorial.


\begin{lemma}
\label{editorial-source-lemma-subextensions-are-separable}
Let $K/k$ be a separable field extension.
For any subextension $K/K'/k$ the field
extension $K'/k$ is separable.
\end{lemma}

\begin{proof}
Let $k\subseteq L\subseteq K'$ be any original
intermediate field finitely generated over $k$.
Then $L$ is also an intermediate field of the
original $K/k$. Separability of $K/k$, with
exactly the given definition, supplies a
transcendence basis of $L/k$ over which $L$
is separable algebraic. Thus every such $L$
is separably generated over $k$, which is
precisely separability of $K'/k$. No finite
generation of $K'$ or $K$ is required.
\end{proof}


\subsection{Editorial treatment of Algebra lines 9944--9959}

\hypertarget{algebra-context-117}{}

\noindent\href{https://github.com/stacks/stacks-project/blob/a04446e57ec1fbc252a871afcec7752fb2807b14/algebra.tex\#L9944}{Original source, lines 9944--9959}. Complete original and reviewed TeX passages accompany this note. Every intervention is separately recorded; the following treatment is editorial.


\begin{lemma}
\label{editorial-source-lemma-generating-finitely-generated-separable-field-extensions}
Let $K/k$ be a separably generated and finitely generated
field extension.
Set $r = \text{trdeg}_k(K)$. Then there exist elements
$x_1, \ldots, x_{r + 1}$ of $K$ such that
\begin{enumerate}
\item $x_1, \ldots, x_r$ is a transcendence basis of $K$ over $k$,
\item $K = k(x_1, \ldots, x_{r + 1})$, and
\item $x_{r + 1}$ is separable over $k(x_1, \ldots, x_r)$.
\end{enumerate}
\end{lemma}

\begin{proof}
Choose a separating transcendence basis of the
original extension. Since $K/k$ is finitely
generated, it has finite size $r$; enumerate
the chosen elements as $x_1,\ldots,x_r$.
Put $F=k(x_1,\ldots,x_r)$, retaining this
original subfield of $K$. If
$K=k(z_1,\ldots,z_m)$ is a finite generating
list, each $z_j$ is algebraic over $F$,
and $K=F(z_1,\ldots,z_m)$. The successive
finite algebraic degrees therefore prove
that $K/F$ is finite. It is separable by
the chosen separating basis. Fields,
Lemma \ref{fields-lemma-primitive-element}
supplies an original $x_{r+1}\in K$ with
$K=F(x_{r+1})$. Its minimal polynomial over
$F$ is separable because the extension is.
This proves all three assertions with the
same original elements. If $K=F$, take
$x_{r+1}=0$, whose minimal polynomial $T$
is separable. If $r=0$, the basis is empty,
$F=k$, and the same finite-extension
argument applies without any transcendental
generators.
\end{proof}


\subsection{Editorial treatment of Algebra lines 9961--10032}

\hypertarget{algebra-context-118}{}

\noindent\href{https://github.com/stacks/stacks-project/blob/a04446e57ec1fbc252a871afcec7752fb2807b14/algebra.tex\#L9961}{Original source, lines 9961--10032}. Complete original and reviewed TeX passages accompany this note. Every intervention is separately recorded; the following treatment is editorial.


\begin{lemma}
\label{editorial-source-lemma-make-separably-generated}
Let $K/k$ be a finitely generated field extension.
There exists a diagram
$$
\xymatrix{
K \ar[r] & K' \\
k \ar[u] \ar[r] & k' \ar[u]
}
$$
where $k'/k$, $K'/K$ are finite purely inseparable field
extensions such that $K'/k'$ is a separably generated field extension.
\end{lemma}

\begin{proof}
Choose the original transcendence basis
$x_1,\ldots,x_r$ of $K$ over $k$, and put
$F=k(x_1,\ldots,x_r)$. Its size $r$ is finite
because the original field extension is finitely
generated. If $K=k(z_1,\ldots,z_m)$ is a specified
finite generating family, then every $z_j$ is
algebraic over $F$ and $K=F(z_1,\ldots,z_m)$.
Successively adjoining these original elements
therefore makes $K/F$ finite. In characteristic
zero every irreducible algebraic polynomial has
nonzero derivative and is separable, so this
basis already separates $K/k$ and one can take
$k'=k$, $K'=K$.

Suppose now that the characteristic is $p>0$.
Let $K/K_{sep}/F$ be the original separable-first
tower from Fields, Lemma \ref{fields-lemma-separable-first}.
Thus $K_{sep}/F$ is finite separable and
$K/K_{sep}$ is finite purely inseparable. Retain
$d=[K:K_{sep}]$. If $d=1$, the original basis
is separating and again no extension is needed.
We construct the strict induction step for $d>1$.

Choose an original $\eta\in K\setminus K_{sep}$.
Pure inseparability gives a least $h\geq1$
such that $\eta^{p^h}\in K_{sep}$. Set
$$
\beta=\eta^{p^{h-1}},\qquad
\alpha=\beta^p=\eta^{p^h}.
$$
Then $\beta\notin K_{sep}$ and
$\alpha\in K_{sep}$. The polynomial $T^p-\alpha$
is irreducible over $K_{sep}$, since a $p$-th
root of $\alpha$ there would equal $\beta$
in the ambient field, contradicting this choice;
see Fields, Lemma \ref{fields-lemma-take-pth-root}.
Hence $[K_{sep}(\beta):K_{sep}]=p$.
Keep the original monic minimal polynomial
$$
P(T)=T^n+a_1T^{n-1}+\cdots+a_n\in F[T]
$$
of $\alpha$. It is separable since
$\alpha\in K_{sep}$; in particular $P'(\alpha)\ne0$.

Here is the omitted finite coefficient construction.
For every original $a_i$, choose complete polynomial
numerator and denominator expressions
$$
a_i=\frac{A_i(x_1,\ldots,x_r)}{B_i(x_1,\ldots,x_r)},
\qquad A_i,B_i\in k[X_1,\ldots,X_r],\quad B_i\ne0.
$$
No common factors are canceled and all coefficients
and summands of these chosen polynomials are kept.
Let $\mathcal C\subset k$ be the finite set of
all their coefficients. Work in one fixed
algebraic closure $\Omega$ of the original $K$,
and choose the unique $p$-th roots there of
the elements of $\mathcal C$ and of the $x_j$.
Set
$$
k'=k(c^{1/p}:c\in\mathcal C),\qquad
y_j=x_j^{1/p},\qquad
F'=k'(y_1,\ldots,y_r).
$$
Each adjoined constant satisfies $U^p-c=0$.
Thus $k'/k$ is finite purely inseparable,
every element of $k'$ has $p$-th power in $k$,
and $[k':k]\leq p^{|\mathcal C|}$.
The last assertion about powers follows for
polynomials and fractions from the characteristic
$p$ identity $(a+b)^p=a^p+b^p$.
The elements $y_1,\ldots,y_r$ are algebraically
independent over $k'$: if
$\sum_\nu d_\nu y^\nu=0$ is a nonzero polynomial
relation, raising it to the $p$-th power gives
$\sum_\nu d_\nu^p x^\nu=0$ over $k$.
At least one coefficient remains nonzero because
Frobenius is injective on a field, contradicting
the original algebraic independence of the $x_j$.

Write the full original expansions as
$A_i(X)=\sum_\nu a_{i,\nu}X^\nu$ and
$B_i(X)=\sum_\nu b_{i,\nu}X^\nu$, with finite
support and nonnegative integer multi-indices.
Their coefficient-root polynomials are
$$
\widetilde A_i(Y)=\sum_\nu a_{i,\nu}^{1/p}Y^\nu,
\qquad
\widetilde B_i(Y)=\sum_\nu b_{i,\nu}^{1/p}Y^\nu.
$$
The denominator is nonzero at the algebraically
independent original $y_j$, and the full identities
are
$$
\widetilde A_i(y)^p=A_i(x),\qquad
\widetilde B_i(y)^p=B_i(x),\qquad
\left(\frac{\widetilde A_i(y)}{\widetilde B_i(y)}\right)^p=a_i.
$$
Put $b_i=\widetilde A_i(y)/\widetilde B_i(y)\in F'$.
This retains each original numerator, denominator,
monomial and coefficient and proves the required
$p$-th power property for every $a_i$, including zero.
The field $F'$ contains $F$ by the original
inclusions and the identities $x_j=y_j^p$.
Moreover $(F')^p\subseteq F$ and
$[F':F]\leq p^{|\mathcal C|+r}$.

Let $L=KF'\subseteq\Omega$ be the actual compositum.
It gives the original commutative diagram
$$
\xymatrix{
K\ar[r]&L\\
k(x_1,\ldots,x_r)\ar[u]\ar[r]&
k'(x_1^{1/p},\ldots,x_r^{1/p})\ar[u].
}
$$
Every arrow is the indicated inclusion in $\Omega$.
The extension $L/K$ is finite purely inseparable
and $L^p\subseteq K$, because the only new
generators are the displayed constant roots
and coordinate roots. In particular
$[L:K]\leq p^{|\mathcal C|+r}$.
The original field $L$ is also finite over $F'$,
as $K$ is finite over $F$.

Put $C=K_{sep}F'$. A finite collection of separable
generators of $K_{sep}/F$ remains separable after
extension to $F'$: their minimal polynomials over
the larger base divide their original separable
polynomials. Thus $C/F'$ is finite separable.
The extension $L/C$ is purely inseparable. To
check this on every element, take finite
generators of $K/K_{sep}$, choose a common
$p$-power placing all of them in $K_{sep}$,
and apply that power to any rational expression
in these generators with coefficients in $C$.
All resulting powers lie in $C$.
It follows that the original separable-first
field $L_{sep}$ over $F'$ is exactly $C$.
Indeed all of $C$ is separable over $F'$, while
an element of $L$ separable over $F'$ is also
separable over $C$, and is purely inseparable
over $C$ by what was just proved. Its minimal
polynomial over $C$ is then both separable and
purely inseparable, hence linear. This proves
the identification $L_{sep}=K_{sep}F'$ without
an assumption about linear disjointness.

The chosen $\beta$ already belongs to $L_{sep}$.
To verify this with the original polynomial, put
$$
Q(T)=T^n+b_1T^{n-1}+\cdots+b_n\in F'[T].
$$
Then $Q(\beta)^p=P(\alpha)=0$, so
$Q(\beta)=0$. Keeping all derivative factors gives
\[
\begin{split}
Q'(\beta)&=n\beta^{n-1}
+\sum_{i=1}^{n-1}(n-i)b_i\beta^{n-i-1},\\
Q'(\beta)^p&=n\alpha^{n-1}
+\sum_{i=1}^{n-1}(n-i)a_i\alpha^{n-i-1}
=P'(\alpha)\ne0.
\end{split}
\]
Here each integer coefficient is its element of
the prime field, whose $p$-th power is itself;
terms divisible by $p$ are retained with their
actual zero coefficients. Thus the minimal
polynomial of $\beta$ over $F'$ is separable:
it divides $Q$, and differentiating that
factorization at $\beta$ shows that its
derivative is nonzero. Consequently
$\beta\in L_{sep}$. Equivalently, Fields,
Lemma \ref{fields-lemma-pth-root} applies to
the original $\alpha$ and $P$ inside the separable
extension $L_{sep}/F'$; its root $\alpha'$
satisfies $(\alpha')^p=\alpha=\beta^p$ and
therefore $\alpha'=\beta$ by injectivity of
Frobenius in $\Omega$. This proves the original
root identification, including its coefficient map.

All entries of the original induction tower are
therefore the actual indicated subfields:
$$
\xymatrix{
K\ar[r]&L\\
K_{sep}(\beta)\ar[r]\ar[u]&L_{sep}\ar[u]\\
K_{sep}\ar[r]\ar[u]&L_{sep}\ar@{=}[u]\\
k(x_1,\ldots,x_r)\ar[r]\ar[u]&
k'(x_1^{1/p},\ldots,x_r^{1/p})\ar[u]\\
k\ar[r]\ar[u]&k'\ar[u].
}
$$
The decrease can be bounded exactly. Since
$K_{sep}(\beta)\subseteq L_{sep}$ and
$L=KL_{sep}$, a basis of $K$ over
$K_{sep}(\beta)$ spans $L$ over $L_{sep}$.
For completeness, its span is the image of
$L_{sep}\otimes_{K_{sep}(\beta)}K\to L$.
This image is a finite-dimensional algebra over
$L_{sep}$, is a domain as a subring of $L$,
and contains both original fields. Any nonzero
element has invertible multiplication on this
finite-dimensional vector space, so the image
is itself a field. It is therefore the
compositum $L$, proving the asserted spanning.
The original tower law now gives
$$
[L:L_{sep}]
\leq[K:K_{sep}(\beta)]
=\frac{[K:K_{sep}]}{[K_{sep}(\beta):K_{sep}]}
=\frac d p<d.
$$
This verifies the strict induction step without
replacing it by an unproved degree comparison.
Applying the same construction to $L/k'$ with
the retained basis $y_1,\ldots,y_r$ completes
the induction. Compositions of the finite
purely inseparable extensions remain finite
and purely inseparable: choose common powers
for finite generating families at successive
stages and compose those powers. Thus the
final inclusions $k\subseteq k''$ and
$K\subseteq L'$ give the required original
diagram, with $k''$ and $L'$ used as the
final $k'$ and $K'$ in the statement.

\medskip\noindent
The same proof provides finite quantitative
information. The original purely inseparable
degree is $d=p^e$ for some integer $e\geq0$.
Indeed adjoining a purely inseparable algebraic
element has degree a power of $p$: its irreducible
polynomial has only one root in an algebraic
closure, and repeated removal of zero derivatives
writes it as a separable irreducible polynomial
in $T^{p^a}$; the separable polynomial must be
linear because there is only one root. A finite
generating tower and the degree law give $d=p^e$.
Every nonterminal step divides the bound for
the remaining inseparable degree by at least $p$.
Consequently the construction terminates after
some $s\leq e$ steps.

Write the actual fields at these steps as
$k_h\subseteq K_h$ for $0\leq h\leq s$,
with $k_0=k$ and $K_0=K$, and retain the
separating-coordinate candidates
$x_j^{1/p^h}$ in the fixed $\Omega$.
At step $h$, let $c_h$ be the number of
coefficients in the finite set used for its
chosen complete rational numerator and
denominator expressions. The construction gives
$$
[k_{h+1}:k_h]\leq p^{c_h},\qquad
[K_{h+1}:K_h]\leq p^{c_h+r},\qquad
k_{h+1}^{,p}\subseteq k_h,\qquad
K_{h+1}^{,p}\subseteq K_h.
$$
Iterating these actual tower maps yields
$$
[k_s:k]\leq p^{\sum_{h=0}^{s-1}c_h},\qquad
[K_s:K]\leq p^{\sum_{h=0}^{s-1}(c_h+r)},\qquad
k_s^{,p^s}\subseteq k,\qquad K_s^{,p^s}\subseteq K.
$$
The terminal extension
$K_s/k_s(x_1^{1/p^s},\ldots,x_r^{1/p^s})$
is separable by the termination criterion.
These are bounds for the original finite
construction, with all constants and coordinate
root adjunctions retained. They do not assert
that the coefficient counts $c_h$ depend only
on $d$, nor that the original transcendence
basis separates $K/k$ before these extensions.
When $r=0$ the coordinate list is empty; when
$s=0$ the sums are empty, the bounds are $1$,
and the original fields already give the result.
\end{proof}


\subsection{Editorial treatment of Algebra lines 10037--10288}

\hypertarget{algebra-context-119}{}

\noindent\href{https://github.com/stacks/stacks-project/blob/a04446e57ec1fbc252a871afcec7752fb2807b14/algebra.tex\#L10037}{Original source, lines 10037--10288}. Complete original and reviewed TeX passages accompany this note. Every intervention is separately recorded; the following treatment is editorial.


\subsubsection{Geometrically reduced algebras}
\label{editorial-source-section-geometrically-reduced}

\noindent
The main result on geometrically reduced algebras is
Lemma \ref{editorial-source-lemma-geometrically-reduced-finite-purely-inseparable-extension}.
We suggest the reader skip to the lemma after reading the definition.

\begin{definition}
\label{editorial-source-definition-geometrically-reduced}
Let $k$ be a field. Let $S$ be a $k$-algebra.
We say $S$ is {\it geometrically reduced over $k$}
if for every field extension $K/k$ the
$K$-algebra $K \otimes_k S$ is reduced.
\end{definition}

\noindent
Let $k$ be a field and let $S$ be a reduced $k$-algebra.
To check that $S$ is geometrically reduced it will suffice
to check that $\overline{k} \otimes_k S$ is reduced (where
$\overline{k}$ denotes the algebraic closure of $k$).
In fact it is enough to check this for finite purely inseparable
field extensions $k'/k$. See
Lemma \ref{editorial-source-lemma-geometrically-reduced-finite-purely-inseparable-extension}.

\begin{lemma}
\label{editorial-source-lemma-subalgebra-separable}
Elementary properties of geometrically reduced algebras.
Let $k$ be a field. Let $S$ be a $k$-algebra.
\begin{enumerate}
\item If $S$ is geometrically reduced over $k$ so is every
$k$-subalgebra.
\item If all finitely generated $k$-subalgebras of $S$ are
geometrically reduced, then $S$ is geometrically reduced.
\item A directed colimit of geometrically reduced $k$-algebras
is geometrically reduced.
\item If $S$ is geometrically reduced over $k$, then any localization
of $S$ is geometrically reduced over $k$.
\end{enumerate}
\end{lemma}

\begin{proof}
Fix an arbitrary original field extension $K/k$.
For (1), if $S'\subseteq S$ is the specified
$k$-subalgebra, the map
$K\otimes_kS'\to K\otimes_kS$ is injective.
Indeed choosing a $k$-basis of $K$ identifies it
with the direct sum of copies of the original
injection. If an element of the first tensor
ring is nilpotent, its image in the reduced
second ring is zero, and injectivity makes the
original element zero. This proves (1).

For (2), retain an original nilpotent element
$z=\sum_{i=1}^n\lambda_i\otimes s_i$ of
$K\otimes_kS$, with $z^d=0$ for some $d\geq1$.
The same sum belongs to $K\otimes_kS_0$ for
the original finite subalgebra
$S_0=k[s_1,\ldots,s_n]\subseteq S$.
The map into $K\otimes_kS$ is injective by
the same basis argument, so the identical
$d$-th power is zero already in this subalgebra.
It is reduced by the hypothesis on $S_0$;
thus the sum is zero and so is its original
image $z$. This proves (2), including the
empty sum, which is zero from the outset.

For (3), write the original system as
$(S_i,\varphi_{ij})$ and its directed colimit
as $S$. The comparison
$$
\mathop{\rm colim}_i(K\otimes_kS_i)
\longrightarrow K\otimes_kS,
\qquad [\lambda\otimes s_i]\longmapsto
\lambda\otimes[s_i]
$$
is an isomorphism. Its inverse sends a finite
sum of tensors to representatives at a common
stage. Independence follows because any finite
list of equalities of colimit elements holds
at a later common stage; the tensor balancing
and additivity relations are finite such lists.
These formulas give both composites on the
original elementary tensors. If the class of
$z_i\in K\otimes_kS_i$ has $d$-th power zero,
then at some later stage $j$ the image of
$z_i^d$ is zero. The image of $z_i$ has
$d$-th power zero in the reduced ring
$K\otimes_kS_j$, and is therefore zero there.
Thus its original colimit class is zero.
No injectivity of the transition maps is used.
This proves the second and third properties
with the actual tensor-colimit comparison.

For (4), let $U\subseteq S$ be the original
multiplicative set and let $U_K$ be its image
$\{1\otimes u:u\in U\}$ in $K\otimes_kS$.
There are inverse algebra maps
$$
K\otimes_kU^{-1}S\longrightarrow
U_K^{-1}(K\otimes_kS),\qquad
\lambda\otimes s/u\longmapsto
\frac{\lambda\otimes s}{1\otimes u},
$$
and
$$
\frac{\sum_i\lambda_i\otimes s_i}{1\otimes u}
\longmapsto\sum_i\lambda_i\otimes s_i/u.
$$
Products of denominators still have the displayed
form. The original balancing maps send every
denominator to a unit, so localization gives
well-defined homomorphisms; their composites
fix both tensor factors and the inverses of
the original denominators. A localization of
a reduced ring $B$ is reduced: if
$(b/u)^d=0$, there is $v$ in the multiplicative
set with $vb^d=0$. Then
$(vb)^d=v^{d-1}(vb^d)=0$, so $vb=0$ in $B$.
This is a zero witness for the original fraction
$b/u$. If zero is inverted, the localization
is the zero ring and the same conclusion holds.
Apply this to the reduced $K\otimes_kS$.
Since $K$ was arbitrary, all four assertions
hold for the original algebras.
\end{proof}

\begin{lemma}
\label{editorial-source-lemma-geometrically-reduced-permanence}
Let $k$ be a field.
If $R$ is geometrically reduced over $k$,
and $S \subset R$ is a multiplicative subset, then the localization
$S^{-1}R$ is geometrically reduced over $k$.
If $R$ is geometrically reduced over $k$, then $R[x]$ is geometrically
reduced over $k$.
\end{lemma}

\begin{proof}
For the original multiplicative set $S\subseteq R$,
the preceding proof gives the exact maps
$$
K\otimes_kS^{-1}R\simeq
S_K^{-1}(K\otimes_kR),\qquad
\lambda\otimes r/s\longmapsto
\frac{\lambda\otimes r}{1\otimes s},
$$
where $S_K=\{1\otimes s:s\in S\}$.
Its inverse sends a fraction with numerator
$\sum_i\lambda_i\otimes r_i$ and denominator
$1\otimes s$ to $\sum_i\lambda_i\otimes r_i/s$.
That proof verifies the original balancing,
fraction witnesses and reducedness of this
localization, including the zero localization.

For the original polynomial ring the comparison is
$$
K\otimes_kR[x]\longrightarrow(K\otimes_kR)[x],
\qquad \lambda\otimes\sum_i r_i x^i
\longmapsto\sum_i(\lambda\otimes r_i)x^i.
$$
Its inverse sends each coefficient tensor
times $x^i$ to $\lambda\otimes r_i x^i$
and extends additively, retaining every
original coefficient and exponent. The
monomial bases verify both composites,
balancing and multiplication.
If a nonzero polynomial over a reduced
ring $B$ has highest nonzero coefficient
$b_d$, its $n$-th power has coefficient
$b_d^n$ at the original degree $nd$.
That coefficient is nonzero, since $b_d$
is nonzero and $B$ reduced. Thus a nonzero
polynomial cannot be nilpotent. Applying
this to $B=K\otimes_kR$ proves reducedness
for every original $K/k$, as required.

The same comparison and induction apply
to any finite list of original variables.
A polynomial in an arbitrary variable
set uses only a finite sublist; its
nilpotence relation therefore holds in
that finite polynomial subring by the
injective monomial inclusion. Thus the
polynomial assertion holds for arbitrary
sets of indeterminates. It is an equivalence:
the constant inclusion
$K\otimes_kR\hookrightarrow(K\otimes_kR)[x_i]$
is injective, with evaluation at all
$x_i=0$ a left inverse. Hence geometric
reducedness of the polynomial algebra
implies that of the original coefficient
algebra. Further localization of these
original polynomial rings is covered
by the proved localization maps.
\end{proof}

\noindent
In the proofs of the following lemmas we will repeatedly use
the following observation: Suppose that $R' \subset R$ and
$S' \subset S$ are inclusions of $k$-algebras.
Then the map $R' \otimes_k S' \to R \otimes_k S$
is injective. More explicitly it factors as
$$
R'\otimes_kS'\longrightarrow R\otimes_kS'
\longrightarrow R\otimes_kS.
$$
The first map is injective by choosing a
$k$-basis of $S'$ and applying the original
injection $R'\to R$ in each coordinate.
The second is injective by choosing a
$k$-basis of $R$ and applying $S'\to S$
in each coordinate. Both maps send every
original elementary tensor to its indicated
image and are algebra homomorphisms, so
their composite is the stated injection.

\begin{lemma}
\label{editorial-source-lemma-limit-argument}
Let $k$ be a field. Let $R$, $S$ be $k$-algebras.
\begin{enumerate}
\item If $R \otimes_k S$ is nonreduced, then there exist
finitely generated subalgebras $R' \subset R$,
$S' \subset S$ such that $R' \otimes_k S'$ is not reduced.
\item If $R \otimes_k S$ contains a nonzero zerodivisor, then there exist
finitely generated subalgebras $R' \subset R$,
$S' \subset S$ such that $R' \otimes_k S'$ contains a nonzero zerodivisor.
\item If $R \otimes_k S$ contains a nontrivial idempotent, then there exist
finitely generated subalgebras $R' \subset R$,
$S' \subset S$ such that $R' \otimes_k S'$ contains a nontrivial idempotent.
\end{enumerate}
\end{lemma}

\begin{proof}
For (1), retain a nonzero nilpotent original
$z=\sum_{i=1}^n x_i\otimes y_i$ with
$z^d=0$, $d\geq1$. Let $R'$ be the
original subalgebra generated by all $x_i$
and $S'$ that generated by all $y_i$.
The same full sum defines $z'\in R'\otimes_kS'$.
Its image is $z$, and the preceding
injection gives $(z')^d=0$. It is nonzero,
since its image is nonzero. Thus this
original tensor subalgebra is not reduced.

For (2), retain both original witnesses
$z\ne0$ and $w\ne0$ with $zw=0$.
Write each as a finite tensor sum, and
generate $R'$ and $S'$ using all the
coefficients from both sums. Their
preimages $z',w'$ are nonzero, and
$z'w'=0$ by the same tensor injection.
Thus $z'$ is a nonzero zerodivisor
in the stated finite subalgebra.
Keeping only the coefficients of $z$
would not by itself supply its required
nonzero annihilating witness.

For (3), retain an original idempotent
$e$ with $e^2=e$, $e\ne0$ and $e\ne1$.
Use all the coefficients of a finite
tensor expression of $e$. The tensor
injection is unital, so its preimage
$e'$ satisfies $(e')^2=e'$ and has
$e'\ne0,1$. This proves (3).

More generally any fixed finite collection
of original tensor elements and polynomial
equations and inequations over $k$ among
them can be witnessed in such finitely
generated $R',S'$. Include every coefficient
of the finite tensor expressions for all
the specified elements. Each original
equation descends because the algebra
map is injective, and each inequation
is preserved because its evaluated image
is nonzero. This assertion concerns the
complete specified finite list and keeps
all witnesses; no infinite collection
is replaced by a finite one.
\end{proof}

\begin{lemma}
\label{editorial-source-lemma-geometrically-reduced-any-reduced-base-change}
Let $k$ be a field.
Let $S$ be a geometrically reduced $k$-algebra.
Let $R$ be any reduced $k$-algebra.
Then $R \otimes_k S$ is reduced.
\end{lemma}

\begin{proof}
We first retain the original finite-subalgebra
argument and its maps. Let
$z=\sum_{i=1}^n r_i\otimes s_i$ be nilpotent
in the original $R\otimes_kS$, and let
$R_0=k[r_1,\ldots,r_n]\subseteq R$.
It is reduced as a subring of the reduced
original $R$ and Noetherian by
Lemma \ref{editorial-source-lemma-Noetherian-permanence}.
The map $R_0\otimes_kS\to R\otimes_kS$ is
injective by the vector-space tensor comparison,
so the same original sum has its power zero
already in $R_0\otimes_kS$.

The original minimal-prime localizations of
$R_0$ are fields $F_j=(R_0)_{\mathfrak p_j}$,
and there are finitely many of them. The maps
$$
R_0\longrightarrow\prod_j F_j,\qquad
r\longmapsto(r/1)_j
$$
are injective, by
Lemmas \ref{editorial-source-lemma-Noetherian-irreducible-components}
and \ref{algebra-lemma-minimal-prime-reduced-ring}.
More explicitly each localization has just its
zero prime, since it is reduced and has a
unique prime. Its kernel on $R_0$ is
$\mathfrak p_j$: a numerator outside that
prime becomes a unit, and a numerator inside
it maps into the zero maximal ideal. The
intersection of the minimal primes is the
nilradical, which is zero. This proves the
injection without replacing the original
$R_0$ by its product of fields.
It agrees with the original total-fraction
comparison from
Lemma \ref{algebra-lemma-total-ring-fractions-no-embedded-points}:
a fraction $r/t$ with $t$ a nonzerodivisor
maps to $(r/1)/(t/1)$ in every factor.
That lemma identifies the total fraction ring
with this finite product; the argument here
uses its indicated injection and maps.

Tensoring the injection with the $k$-vector
space $S$ preserves injectivity. For this
finite product the exact comparison is
$$
\left(\prod_j F_j\right)\otimes_kS
\longrightarrow\prod_j(F_j\otimes_kS),\qquad
(a_j)_j\otimes s\longmapsto(a_j\otimes s)_j.
$$
Its inverse sends $(w_j)_j$ to
$\sum_j(\iota_j\otimes1)(w_j)$, where
$\iota_j:F_j\to\prod_jF_j$ is the original
coordinate inclusion as a $k$-linear map.
The finite sum proves both inverse identities
and preserves multiplication, since distinct
coordinate factors multiply to zero.
Each $F_j\otimes_kS$ is reduced by geometric
reducedness of $S$. A product of reduced rings
is reduced, because nilpotence is checked
coordinate by coordinate. Thus the image of
the original $z$ in this product is zero.
Both injections show $z=0$ in the original
$R\otimes_kS$. If $R_0=0$, the original sum
is already zero; the empty product causes
no exception. This proves the lemma.

\medskip\noindent
The same result yields an exact formula with
an arbitrary original $k$-algebra $A$, which
need not be reduced. Write
$\mathcal N=\sqrt{(0)}\subseteq A$ for its
original nilradical. Then
$$
\sqrt{(0)}_{\,A\otimes_kS}
=\operatorname{image}\bigl(\mathcal N\otimes_kS
\longrightarrow A\otimes_kS\bigr)
=\mathcal N(A\otimes_kS).
$$
The tensor map is injective, by flatness over
the field. Its image is the displayed ideal:
multiplying a tensor $n\otimes s$ by
$a\otimes t$ gives $an\otimes st$, still
with $an\in\mathcal N$, and every generator
$(n\otimes1)(1\otimes s)$ is such a tensor.
Every original finite element
$w=\sum_{i=1}^m n_i\otimes s_i$ of this image
is nilpotent. Choose $d_i\geq1$ with
$n_i^{d_i}=0$ and retain
$D=1+\sum_{i=1}^m(d_i-1)$.
For $m\geq1$ its full multinomial expansion is
$$
w^D=\sum_{\substack{a_1+\cdots+a_m=D\\a_i\geq0}}
\frac{D!}{a_1!\cdots a_m!}
\left(\prod_{i=1}^m n_i^{a_i}\right)
\otimes\left(\prod_{i=1}^m s_i^{a_i}\right)=0.
$$
For each tuple some $a_i\geq d_i$; otherwise
their sum would be at most $D-1$. Thus every
displayed term is zero with its full integer
multinomial coefficient retained. These are
integer coefficients acting in the ring,
not divisions by possibly noninvertible
factorials. The empty sum is zero.

Conversely, exactness of tensoring gives the
quotient comparison
$$
(A\otimes_kS)/\operatorname{image}(\mathcal N\otimes_kS)
\longrightarrow(A/\mathcal N)\otimes_kS,
\qquad [a\otimes s]\longmapsto(a+\mathcal N)\otimes s.
$$
The inverse is induced by the same elementary
formula in the reverse direction, independent
of the representative because the difference
has numerator in $\mathcal N$. The original
quotient $A/\mathcal N$ is reduced, so the
proved lemma makes this tensor quotient reduced.
Every nilpotent in $A\otimes_kS$ consequently
maps to zero there, proving the reverse
inclusion and the exact nilradical identity.

If $S\ne0$, the map $A\to A\otimes_kS$,
$a\mapsto a\otimes1$, is injective: extend
$1\in S$ to a $k$-basis and take the linear
functional $\ell:S\to k$ with $\ell(1)=1$
and zero on the other basis elements.
Then $1\otimes\ell$ is a left inverse.
It follows that for this geometrically
reduced nonzero $S$, the original $A$ is
reduced if and only if $A\otimes_kS$ is
reduced. For $S=0$ the tensor is always
zero, so this last converse is not asserted;
the nilradical formula itself still holds.
\end{proof}

\begin{lemma}
\label{editorial-source-lemma-separable-extension-preserves-reducedness}
Let $k$ be a field.
Let $S$ be a reduced $k$-algebra.
Let $K/k$ be either a separable field extension,
or a separably generated field extension.
Then $K \otimes_k S$ is reduced.
\end{lemma}

\begin{proof}
We first prove the assertion for an original
finitely generated and separably generated
extension $K/k$ and an original domain $D$
over $k$. Choose the original elements
$x_1,\ldots,x_r,\alpha=x_{r+1}$ from
Lemma \ref{editorial-source-lemma-generating-finitely-generated-separable-field-extensions}.
Put $E=k(x_1,\ldots,x_r)$, so that
$K=E(\alpha)$, and retain the monic separable
minimal polynomial $P(T)\in E[T]$ of $\alpha$.
If $K=E$, the choice $\alpha=0$, $P(T)=T$
gives the same proof.

The tensor map
$$
k[x_1,\ldots,x_r]\otimes_kD
\longrightarrow D[x_1,\ldots,x_r],\qquad
\left(\sum_\nu c_\nu x^\nu\right)\otimes d
\longmapsto\sum_\nu c_\nu d x^\nu
$$
is an isomorphism. Its inverse sends
$\sum_\nu d_\nu x^\nu$ to
$\sum_\nu x^\nu\otimes d_\nu$; the original
monomial bases prove balancing and both
composites. With
$W=k[x_1,\ldots,x_r]\setminus\{0\}$,
the original fraction map gives
$$
A=E\otimes_kD\simeq W^{-1}D[x_1,\ldots,x_r],
\qquad (g/h)\otimes d\longmapsto dg/h.
$$
The inverse sends $\sum_\nu d_\nu x^\nu/h$
to $\sum_\nu(x^\nu/h)\otimes d_\nu$.
The universal property of inverting the
same original $h\in W$ verifies these maps,
including equality witnesses for fractions.
Every element of $W$ remains nonzero in the
polynomial ring over the domain $D$, because
the original coefficient map $k\to D$ is
injective. Thus $A$ is a domain. Retain its
fraction field $F=\operatorname{Frac}(A)$.

There is an exact original algebra comparison
$$
A[T]/(P)\longrightarrow K\otimes_kD,
\qquad T\longmapsto\alpha\otimes1,
$$
with coefficients $e\otimes d$ sent to the
same original tensor. The powers
$1,\alpha,\ldots,\alpha^{n-1}$, where
$n=\deg P$, form an $E$-basis of $K$.
Consequently both sides are free $A$-modules
with these corresponding bases, proving the
map an isomorphism. Its inverse sends
$\gamma\otimes d$, for the unique expansion
$\gamma=\sum_{i=0}^{n-1}e_i\alpha^i$, to
$\sum_{i=0}^{n-1}(e_i\otimes d)T^i$ modulo
the original $P$. Multiplication is preserved
because both sides impose exactly $P(\alpha)=0$.
The map $A[T]/(P)\to F[T]/(P)$ is injective:
monic division leaves a unique remainder
of degree less than $n$, and its coefficients
cannot become zero in the fraction field
unless they were zero in $A$.

Separability supplies polynomials $U,V\in E[T]$
with the full identity $UP+VP'=1$.
The original coefficient maps carry this
identity into $F[T]$. Thus $P$ has no repeated
irreducible factor there. Write its actual
factorization $P=\prod_{j=1}^t P_j$ into
distinct monic irreducible polynomials over
$F$. The Chinese remainder comparison sends
$$
F[T]/(P)\longrightarrow\prod_{j=1}^t F[T]/(P_j),
\qquad [H]\longmapsto([H])_j.
$$
For the inverse, put $H_j=P/P_j$ and choose
$U_j$ with $U_jH_j\equiv1\pmod{P_j}$.
A tuple $([G_j])_j$ maps to
$[\sum_jG_jU_jH_j]$ modulo $P$. It is
independent of lifts and composes to the
identity, since each term is $G_j$ in its
own factor and zero in the others; a
polynomial vanishing modulo all the coprime
$P_j$ is divisible by their original
product $P$. Each quotient is a field.
Hence $F[T]/(P)$ is reduced, and the
proved injections show $K\otimes_kD$
is reduced for the original $K$ and $D$.

Now suppose $K/k$ is separable according to
the definition, with no finiteness hypothesis,
and retain an original nilpotent
$z=\sum_{i=1}^m\lambda_i\otimes s_i$
in $K\otimes_kS$. Let
$K_0=k(\lambda_1,\ldots,\lambda_m)\subseteq K$
and $S_0=k[s_1,\ldots,s_m]\subseteq S$.
By definition $K_0/k$ is separably generated
and finitely generated. The original $S_0$
is reduced and Noetherian. The tensor map
$K_0\otimes_kS_0\to K\otimes_kS$ is injective,
so the same original sum is nilpotent already
in the smaller tensor ring.
As in the preceding reduced-base-change
proof, retain the injection
$S_0\to\prod_j (S_0)_{\mathfrak p_j}$
through its finitely many minimal-prime
fields, with every map $s\mapsto s/1$.
This is the original total-fraction comparison
of Lemmas \ref{algebra-lemma-minimal-prime-reduced-ring}
and \ref{algebra-lemma-total-ring-fractions-no-embedded-points}.
Tensoring with $K_0$ is injective, and the
finite-product tensor map is an isomorphism
with inverse given by the finite sum of
coordinate inclusions. Each factor
$K_0\otimes_k(S_0)_{\mathfrak p_j}$ is
reduced by the domain case just proved.
Thus the original sum is zero in their
product and, by injectivity, zero in
$K\otimes_kS$. This proves the separable case.
If $S_0=0$ the tensor was already zero.

Finally suppose $K/k$ is separably generated,
possibly by an infinite original separating
basis $(x_i)_{i\in I}$. We do not assume
the later theorem that this implies separability.
Put $E=k(x_i:i\in I)$ and take any finite
list $\lambda_1,\ldots,\lambda_m\in K$.
Each $\lambda_j$ has a monic separable
minimal polynomial $P_j\in E[T]$, with a
Bezout identity $U_jP_j+V_jP_j'=1$ in $E[T]$.
Every coefficient in this finite collection
of complete polynomials and identities is
a rational function in finitely many of
the original $x_i$. Choose a finite
$J\subseteq I$ containing all those indices.
The same polynomials and identities then
belong to $E_J[T]$, where
$E_J=k(x_i:i\in J)$, since its inclusion
into $E$ is injective. Each $\lambda_j$
is algebraic and separable over $E_J$:
its minimal polynomial divides the same
monic $P_j$, which has no repeated roots
by the retained Bezout identity.
Therefore
$$
K_{J,\lambda}=k(x_i:i\in J)(\lambda_1,\ldots,\lambda_m)
$$
is finite separable over $E_J$, and is a
finitely generated, separably generated
subextension of the original $K/k$.
The original algebraic independence of the
$x_i$ shows that the selected sublist is
indeed a transcendence basis for this field.

These subfields form a directed family:
for two finite lists and two finite index
sets, include their unions and the finitely
many coefficient indices required for the
combined list. Enlarging the rational base
retains separability because each minimal
polynomial still divides its specified
separable polynomial. Their union is $K$,
since every original element is included
by taking it in a one-element list.
Every nilpotent tensor in $K\otimes_kS$
therefore belongs to one of these original
finite subextensions tensored with $S$.
Its power relation holds there by injectivity,
and that tensor algebra is reduced by the
finite separably generated argument above
and the finite-subalgebra argument for $S$.
Hence the original tensor is zero. This
also gives explicitly the directed-colimit
description in the original proof.

As a consequence, either hypothesis on
$K/k$ makes $K$ geometrically reduced over
$k$: apply the result with an arbitrary
receiving field as the reduced algebra $S$
and use the tensor flip
$\lambda\otimes s\mapsto s\otimes\lambda$,
whose inverse is the same flip. Moreover,
for every original $k$-algebra $B$, the
nilradical formula of the preceding lemma
now gives the exact equality
$$
\sqrt{(0)}_{\,K\otimes_kB}
=\operatorname{image}\bigl(K\otimes_k\sqrt{(0)}_B
\longrightarrow K\otimes_kB\bigr).
$$
Thus this result preserves the entire
nilradical under the original scalar
extension, not only reducedness when
that ideal happens to vanish.
\end{proof}

\begin{lemma}
\label{editorial-source-lemma-generic-points-geometrically-reduced}
Let $k$ be a field and let $S$ be a $k$-algebra. Assume that
$S$ is reduced and that $S_{\mathfrak p}$ is geometrically
reduced for every minimal prime $\mathfrak p$ of $S$.
Then $S$ is geometrically reduced.
\end{lemma}

\begin{proof}
Retain the original localizations $S_{\mathfrak p}$
at every minimal prime. Since $S$ is reduced,
so is each localization. Its sole prime is
$\mathfrak pS_{\mathfrak p}$, and the intersection
of its primes is its zero nilradical; hence
$\mathfrak pS_{\mathfrak p}=0$ and this local
ring is a field. The original map
$S\to S_{\mathfrak p}$ has kernel exactly
$\mathfrak p$: elements in it map to zero,
whereas an element outside it maps to a unit
in this nonzero field. The intersection of
the minimal primes of $S$ is zero, so
$$
S\longrightarrow\prod_{\mathfrak p\text{ minimal}}S_{\mathfrak p},
\qquad s\longmapsto(s/1)_{\mathfrak p}
$$
is injective, as in
Lemma \ref{algebra-lemma-reduced-ring-sub-product-fields}.

For any original field extension $K/k$, tensoring
this map with $K$ preserves injectivity. The
second map in the source proof is
$$
\Theta:\left(\prod_{\mathfrak p}S_{\mathfrak p}\right)\otimes_kK
\longrightarrow\prod_{\mathfrak p}(S_{\mathfrak p}\otimes_kK),
\qquad (a_{\mathfrak p})_{\mathfrak p}\otimes\lambda
\longmapsto(a_{\mathfrak p}\otimes\lambda)_{\mathfrak p}.
$$
To prove its injectivity, write any original
tensor as a finite sum
$w=\sum_{j=1}^r a_j\otimes\lambda_j$ with
$k$-linearly independent $\lambda_j\in K$.
Such an expression is obtained by expanding
the finitely many original scalar coefficients
in a basis of their span, retaining the
resulting sums in each $a_j$.
Extend these $\lambda_j$ to a basis of $K$.
If $\Theta(w)=0$, then for every original
$\mathfrak p$ the finite basis expansion
$\sum_j(a_j)_{\mathfrak p}\otimes\lambda_j=0$
forces $(a_j)_{\mathfrak p}=0$ for each $j$.
Thus every original $a_j$ is zero and $w=0$.
The map is an algebra map since it is
coordinatewise on elementary products.

Each factor $S_{\mathfrak p}\otimes_kK$ is
reduced by its assumed geometric reducedness.
The product is reduced, so the two original
injections imply that $S\otimes_kK$ is reduced.
The tensor flip identifies this with the
ordering in the definition, and proves the
lemma for every $K$. If $S=0$ the index set
is empty and all the displayed rings are
zero, which is reduced.

The converse also holds: if $S$ is geometrically
reduced, taking $K=k$ shows that $S$ is reduced,
and Lemma \ref{editorial-source-lemma-subalgebra-separable} shows
that each of its original localizations is
geometrically reduced. Thus the original
conditions are an equivalence.

The same proof applies more generally whenever
there is a specified injective $k$-algebra map
$S\hookrightarrow\prod_{i\in I}B_i$ and each
$B_i$ is geometrically reduced. The tensor
map into the product of the base changes
is injective by the identical finite-basis
argument, so $S$ is geometrically reduced.
In particular, an arbitrary product of
geometrically reduced $k$-algebras is
geometrically reduced. Conversely if such
a product is geometrically reduced, fix
an original index $i$. Decomposing it as
the two factors $B_i\times\prod_{j\ne i}B_j$
and tensoring uses the finite-product
isomorphism. The first factor $B_i\otimes_kK$
is reduced for every $K$, because a nilpotent
there would give a nilpotent with zero
second coordinate in the reduced product.
Hence each $B_i$ is geometrically reduced.
The empty product is the zero algebra and
satisfies the assertion as well.

Injectivity of $\Theta$ must not be strengthened
to surjectivity without further hypotheses.
For example take every $B_i=k$, indexed by
$i\geq1$, and $K=k(t)$. The sequence
$(t^i)_{i\geq1}\in\prod_iK$ is outside the
image of $(\prod_i k)\otimes_kK$.
Indeed a finite tensor expression uses a
single finite list $\lambda_1,\ldots,\lambda_r$
of elements of $K$, so every coordinate
of its image belongs to their fixed finite
$k$-linear span. The powers $t^i$ are
$k$-linearly independent, since a nonzero
polynomial in the transcendental $t$ cannot
vanish. They cannot all lie in that span.
The original proof and its extensions use
only the established injection.
\end{proof}

\begin{lemma}
\label{editorial-source-lemma-separable-algebraic-diagonal}
Let $k'/k$ be a separable algebraic extension.
Then there exists a multiplicative subset $S \subset k' \otimes_k k'$
such that the multiplication map $k' \otimes_k k' \to k'$
is identified with $k' \otimes_k k' \to S^{-1}(k' \otimes_k k')$.
\end{lemma}

\begin{proof}
First let the original $k'/k$ be finite
separable, and choose its original primitive
element $\alpha$ using Fields,
Lemma \ref{fields-lemma-primitive-element}.
Retain its monic separable minimal polynomial
$P(x)\in k[x]$, of degree $n\geq1$.
The original tensor map, with both sums
included, is
$$
\psi:k'[x]/(P)\longrightarrow k'\otimes_kk',
\qquad
\left[\sum_i\alpha_i x^i\right]
\longmapsto\sum_i\alpha_i\otimes\alpha^i.
$$
The relation $P(1\otimes\alpha)=0$ makes this
well defined. Both rings are free over the
first $k'$ with bases corresponding to
$1,x,\ldots,x^{n-1}$ and
$1\otimes1,1\otimes\alpha,\ldots,1\otimes\alpha^{n-1}$.
The latter basis is obtained by tensoring
the original $k$-basis of the second $k'$.
Thus $\psi$ is an isomorphism. Its inverse
sends $a\otimes b$, with the unique original
expansion $b=\sum_{i=0}^{n-1}c_i\alpha^i$,
$c_i\in k$, to $[\sum_{i=0}^{n-1}ac_i x^i]$.
This also verifies its balancing and its
product law through the same polynomial $P$.
Under this map multiplication
$\mu:k'\otimes_kk'\to k'$ corresponds to
evaluation at $x=\alpha$.

Retain the factorization $P(x)=(x-\alpha)Q(x)$
in the original $k'[x]$. Differentiating gives
$P'(\alpha)=Q(\alpha)\ne0$ by separability.
For $B=k'[x]/(P)$ and the multiplicative set
$S'$ generated by the class of $Q$, the
evaluation map extends to
$$
B_Q\longrightarrow k',\qquad
\frac{[H(x)]}{[Q(x)]^d}\longmapsto
\frac{H(\alpha)}{Q(\alpha)^d}.
$$
Its inverse sends $a\in k'$ to $a/1$.
Indeed $(x-\alpha)Q=0$ in $B$; once the
original $Q$ is a unit, $x=\alpha$.
Consequently every displayed fraction is
the image of its stated scalar value,
and both composites are identities. The
fraction map is independent of lifts and
zero-divisor witnesses by the localization
universal property. Transporting this exact
set $S'$ under $\psi$ gives the required
original multiplicative set $S$.

There is also an explicit original direct
factor. Put $q=Q(\alpha)\in k'\setminus\{0\}$
and retain the polynomial $H$ defined by
$Q(x)-q=(x-\alpha)H(x)$. In $B$, set
$e=[Q(x)]/q$, where the scalar denominator
is the original invertible $q$.
Then the full identities are
$$
e^2-e=\frac{Q(x)(Q(x)-q)}{q^2}
=\frac{P(x)H(x)}{q^2}=0,
\qquad (x-\alpha)e=\frac{P(x)}q=0,
\qquad e(\alpha)=1.
$$
Thus $e$ is an idempotent and
$\ker(B\to k')=(x-\alpha)B=(1-e)B$.
For the equality of ideals, one inclusion
uses $1-e=-(x-\alpha)H/q$; the other uses
$(x-\alpha)=(x-\alpha)(1-e)$.
The map $eB\to k'$ is evaluation, with
inverse $a\mapsto ae$. Also $B_Q=B_e$,
because $Q=qe$ with $q$ already a unit.
These formulas retain the original $Q$
and its full scalar factor. If $n=1$,
then $Q=1$, $H=0$, $e=1$, and the kernel
is zero, so the same formulas apply.

Now retain an arbitrary separable algebraic
extension $k'/k$. Its finite intermediate
fields $k_i$ form a directed family: a
compositum of two finite subextensions
inside $k'$ is again finite and separable.
Every finite set of original coefficients
of $k'$ lies in some $k_i$. The injections
$$
k_i\otimes_kk_i\longrightarrow k'\otimes_kk'
$$
are injective by tensoring the two original
field inclusions over the field $k$.
Their images have union $k'\otimes_kk'$,
since every tensor is a finite sum of
elementary tensors. For each $i$, use the
actual multiplicative set $S_i$ and
idempotent $e_i$ constructed above for
$k_i/k$, and let $S$ be the multiplicative
closure of their images $\bigcup_iS_i$.
The original multiplication map $\mu$
sends every element of $S$ to a nonzero
element of $k'$, because this holds for
each original generator $Q_i$ and their
finite products. It therefore extends to
$S^{-1}(k'\otimes_kk')\to k'$.

If an original $z\in\ker\mu$ is represented
at a stage $k_i\otimes_kk_i$, its multiplication
there is zero: its image in the containing
field $k'$ is zero and $k_i\hookrightarrow k'$
is injective. The finite-stage proof makes
$z$ zero on inverting that stage's $Q_i$;
indeed $Q_i$ annihilates its diagonal kernel.
It is therefore zero in the full original
localization. Multiplication is surjective
already before localization, since
$\mu(a\otimes1)=a$ for every $a\in k'$.
The localized map is consequently bijective:
a fraction mapping to zero has numerator
in $\ker\mu$, which is killed as just proved.
Its inverse is explicitly
$a\mapsto(a\otimes1)/1$.
For a general original fraction $z/s$,
the difference from the image of
$\mu(z)/\mu(s)$ has numerator
$(\mu(s)\otimes1)z-(\mu(z)\otimes1)s$,
whose multiplication is zero; its denominator
is the original product $(\mu(s)\otimes1)s$.
This verifies the second composite and
keeps all scalar and localization factors.

\medskip\noindent
The diagonal kernel also gives a sharp finite
generation consequence. For this original
separable algebraic extension, the ideal
$I=\ker(k'\otimes_kk'\to k')$ is finitely
generated if and only if $k'/k$ is finite.
The finite case follows from the explicit
idempotent generator $1-e$ above.
For the converse, suppose
$I=(z_1,\ldots,z_m)$, and choose a finite
original subfield $k_i$ containing every
coefficient of all these tensors.
Every $z_j$ is in the finite diagonal
kernel, so $e_i z_j=0$ by its proved
idempotent formula. Therefore $e_iI=0$.
Since $\mu(e_i)=1$, also $1-e_i\in I$,
and the exact equality is
$I=(1-e_i)(k'\otimes_kk')$.

The quotient by this particular ideal is
canonically $k'\otimes_{k_i}k'$. In detail,
$\ker(k_i\otimes_kk_i\to k_i)$ is generated
by $a\otimes1-1\otimes a$, $a\in k_i$:
modulo these relations the two original
copies of every $a$ coincide, and multiplication
and $a\mapsto a\otimes1$ are inverse maps.
Extending these relations to $k'\otimes_kk'$
imposes precisely the balancing over $k_i$.
The map sends the original $a\otimes_k b$
to $a\otimes_{k_i}b$, and its inverse is
the same elementary formula into the
quotient; the just-stated relations prove
well-definedness and both composites.
Since the extended finite kernel is exactly
$(1-e_i)$, the full kernel assumption would
make multiplication
$k'\otimes_{k_i}k'\to k'$ an isomorphism.
If $k'\ne k_i$, choose an original
$b\in k'\setminus k_i$ and extend the
$k_i$-independent list $1,b$ to a basis
of $k'$. In the induced basis of the
second tensor factor, the element
$1\otimes_{k_i}b-b\otimes_{k_i}1$ is nonzero,
but its multiplication is zero. This
contradicts that isomorphism. Hence
$k'=k_i$ and the original extension is
finite. The argument includes an empty
generating list for the zero ideal.
\end{proof}

\begin{lemma}
\label{editorial-source-lemma-geometrically-reduced-over-separable-algebraic}
Let $k'/k$ be a separable algebraic field extension.
Let $A$ be an algebra over $k'$. Then $A$ is geometrically
reduced over $k$ if and only if it is geometrically reduced over $k'$.
\end{lemma}

\begin{proof}
First assume that the original $A$ is geometrically
reduced over $k'$, and let $K/k$ be any original
field extension. The ring $K\otimes_kk'$ is
reduced by
Lemma \ref{editorial-source-lemma-separable-extension-preserves-reducedness},
applied to the separable field extension $k'/k$
and the reduced $k$-algebra $K$, followed by
the tensor flip. Treat it as a $k'$-algebra
through its original second factor. Then
Lemma \ref{editorial-source-lemma-geometrically-reduced-any-reduced-base-change}
over $k'$ shows that
$(K\otimes_kk')\otimes_{k'}A$ is reduced.
The exact comparison is
$$
(K\otimes_kk')\otimes_{k'}A\longrightarrow K\otimes_kA,
\qquad (\lambda\otimes b)\otimes a\longmapsto\lambda\otimes ba,
$$
with inverse $\lambda\otimes a\mapsto
(\lambda\otimes1)\otimes a$.
Balancing over the original $k'$ moves $b$
into $a$, and balancing over $k$ is the
original tensor relation. These formulas
verify both composites and multiplication.
Hence $K\otimes_kA$ is reduced for every
$K/k$, proving this direction.

Conversely assume that the original $A$ is
geometrically reduced over $k$, and let $K/k'$
be an arbitrary original field extension.
Retain the rings
$$
H=k'\otimes_kk',\qquad B=K\otimes_kA.
$$
The original $H$-algebra structure on $B$
is the map
$$
\chi:H\longrightarrow B,\qquad
b\otimes c\longmapsto b\otimes c,
$$
where $b$ enters through $k'\to K$ and
$c$ through $k'\to A$. Give $k'$ the
$H$-algebra structure of multiplication
$\mu(b\otimes c)=bc$.
The source's exact comparison, now including
its original actions and punctuation, is
$$
B\otimes_Hk'\longrightarrow K\otimes_{k'}A,
\qquad (\lambda\otimes a)\otimes c
\longmapsto c\lambda\otimes a.
$$
It is balanced over $H$: multiplying the
first tensor by $b\otimes d$ gives
$cb\lambda\otimes da= cbd\lambda\otimes a$
over $k'$, equal to multiplying the last
scalar by $bd$. Its inverse sends
$\lambda\otimes_{k'}a$ to
$(\lambda\otimes_ka)\otimes1$.
For its $k'$-balancing, the difference
$c\lambda\otimes a-\lambda\otimes ca$
is the action of $c\otimes1-1\otimes c$
on $\lambda\otimes a$, and that element
of $H$ is killed by $\mu$. The inverse
is therefore well defined. Both composites
are identities, since the scalar $c$ can
be moved from the last factor by the
original element $c\otimes1\in H$.

Lemma \ref{editorial-source-lemma-separable-algebraic-diagonal}
gives an actual multiplicative set $U\subseteq H$
and the original identification $k'\simeq U^{-1}H$
through $\mu$. Thus the preceding tensor
ring is the localization $\chi(U)^{-1}B$.
The full comparison sends
$$
b\otimes\frac hu\longmapsto
\frac{b\chi(h)}{\chi(u)},
$$
with inverse $b/\chi(u)\mapsto b\otimes1/u$.
Balancing and the same original inverse
denominators prove both maps and composites.
Since $K$ is a field extension of $k$,
the hypothesis makes $B=K\otimes_kA$
reduced. Its localization is reduced by
the zero-witness proof already given.
Therefore the original $K\otimes_{k'}A$
is reduced for every $K/k'$, completing
the converse without identifying a general
quotient of a reduced ring with a reduced ring.

\medskip\noindent
The forward implication has a wider scope:
it holds if the original $k'/k$ is any
separable or separably generated extension,
with no algebraicity assumption. The first
paragraph uses only that
$K\otimes_kk'$ is reduced for every $K/k$,
and the preceding separable-extension lemma
proves exactly this in either case. Its
two displayed original tensor maps remain
unchanged. The reverse implication does
not hold for arbitrary separably generated
extensions, as the following example shows.

Let $k$ be a field of characteristic $p>0$,
take $k'=k(t)$ and $A=k(u)$ with the
original embedding $k'\to A$ given by
$t\mapsto u^p$, where $u$ is transcendental
over $k$. The field $k'/k$ is separably
generated using its basis $t$, and $A$
is geometrically reduced over $k$.
Indeed for any field $L/k$,
$$
L\otimes_kk(u)\simeq
(k[u]\setminus\{0\})^{-1}L[u]
$$
through $\lambda\otimes f(u)/g(u)\mapsto
\lambda f(u)/g(u)$, with inverse obtained
by expanding each numerator in its
original powers of $u$. The specified
nonzero denominators stay nonzero over
$L$, so this is a domain.

The polynomial $T^p-t$ is irreducible
over the original $k(t)$. There is no
$p$-th root of $t$ in $k(t)$: for nonzero
polynomials $f,g\in k[t]$, the equation
$f(t)^p=tg(t)^p$ would imply
$p\operatorname{ord}_t(f)=1+p\operatorname{ord}_t(g)$,
which is impossible. These orders are the
actual smallest powers of $t$ with nonzero
coefficients; the product law follows by
multiplying the corresponding first
nonzero coefficients. No common factor
of the original $f,g$ is removed.
Fields, Lemma \ref{fields-lemma-take-pth-root}
therefore gives irreducibility and the
original degree $[k(u):k(u^p)]=p$.
For the receiving field $K=A$ over $k'$
the original tensor map gives
$$
A\otimes_{k'}A\simeq A[T]/(T^p-u^p),
\qquad \sum_{i=0}^{p-1}a_i\otimes u^i
\longleftrightarrow\left[\sum_{i=0}^{p-1}a_iT^i\right].
$$
More precisely the first $A$ supplies
the coefficients and $1\otimes u$
corresponds to $T$; the original basis
$1,u,\ldots,u^{p-1}$ of the second
factor proves the isomorphism and its
inverse on all tensors. The element
$$
z=1\otimes u-u\otimes1
$$
is nonzero: its polynomial representative
$T-u$ has degree $1<p$ and nonzero leading
coefficient, so it is not a multiple of
the original monic $T^p-u^p$.
But its full binomial power is
$$
(T-u)^p=\sum_{j=0}^p
\binom pj T^j(-u)^{p-j}=T^p-u^p=0
$$
in this quotient, since each middle
integer coefficient is divisible by $p$.
Thus $A$ is not geometrically reduced
over the original $k'$, despite being
geometrically reduced over $k$ and
despite $k'/k$ being separably generated.
This locates the role of algebraicity
in the converse diagonal-localization proof.

Separability cannot simply be dropped
from the original algebraic assertion
either. Take $k=\mathbf F_p(t)$,
$k'=\mathbf F_p(u)$ with the same
$t\mapsto u^p$, and $A=k'$.
The identity algebra $A$ is geometrically
reduced over $k'$, because every field
base change is that field. Its tensor
with the receiving field $k'$ over $k$
is the same nonreduced ring
$k'[T]/(T^p-u^p)$ with the same original
nonzero nilpotent. Hence it is not
geometrically reduced over $k$.
\end{proof}








\subsection{Editorial treatment of Algebra lines 10295--10350}

\hypertarget{algebra-context-120}{}

\noindent\href{https://github.com/stacks/stacks-project/blob/a04446e57ec1fbc252a871afcec7752fb2807b14/algebra.tex\#L10295}{Original source, lines 10295--10350}. Complete original and reviewed TeX passages accompany this note. Every intervention is separately recorded; the following treatment is editorial.


\begin{lemma}
\label{editorial-source-lemma-mini-separability}
Let $k$ be a field of characteristic $p > 1$. Let $K/k$
be a field extension generated by $x_1, \ldots, x_{n + 1} \in K$ such that
\begin{enumerate}
\item $\{x_1, \ldots, x_n\}$ is a transcendence base of $K/k$,
\item for every $k$-linearly independent subset
$\{a_1, \ldots, a_m\}$ of $K$ the set
$\{a^p_1, \ldots, a_m^p\}$ is $k$-linearly independent.
\end{enumerate}
Then there is $1 \leq j \leq n+1$ such that
$\{ x_1, \ldots, \widehat{x}_j, \ldots, x_{n+1}\}$
is a separating transcendence base for $K / k$.
\end{lemma}

\begin{proof}
Retain the original generators $x_1,\ldots,x_{n+1}$.
Since $x_{n+1}$ is algebraic over $k(x_1,\ldots,x_n)$,
clearing the nonzero polynomial denominators in one of its
algebraic equations gives a nonzero
$F\in k[X_1,\ldots,X_{n+1}]$ with $F(x_1,\ldots,x_{n+1})=0$.
Choose such an $F$ of least total degree $D$. A nonzero
constant cannot vanish, so $D\geq1$. If $F=GH$ with both
factors of positive degree, evaluation in the field $K$
makes at least one factor zero. That factor has degree
less than $D$, a contradiction. Thus $F$ is irreducible.

Write the full polynomial as
$$
F=\sum_{\alpha\in A}\lambda_\alpha
X_1^{\alpha_1}\cdots X_{n+1}^{\alpha_{n+1}},
\qquad 0\neq\lambda_\alpha\in k,
$$
where $A$ is its finite set of distinct exponent tuples,
all of whose coordinates are integers at least zero.
Suppose every coordinate of every $\alpha\in A$ were
divisible by $p$. Put
$b_\alpha=x_1^{\alpha_1/p}\cdots x_{n+1}^{\alpha_{n+1}/p}$.
The family $(b_\alpha^p)_{\alpha\in A}$ is dependent,
since $\sum_{\alpha\in A}\lambda_\alpha b_\alpha^p=0$.
Hypothesis (2) also preserves independence of indexed
finite families: an independent family has distinct
entries, and Frobenius is injective on $K$, so its
$p$th powers still have distinct entries and hypothesis
(2) applies to the corresponding sets. Its contrapositive
therefore makes $(b_\alpha)_{\alpha\in A}$ dependent,
including the case where two evaluated monomials agree.
Choose $c_\alpha\in k$, not all zero, such that
$\sum_\alpha c_\alpha b_\alpha=0$. The polynomial
$$
G=\sum_{\alpha\in A}c_\alpha
X_1^{\alpha_1/p}\cdots X_{n+1}^{\alpha_{n+1}/p}
$$
is nonzero because its formal monomials are distinct.
It vanishes at the same original generators and has
degree at most $D/p<D$, a contradiction. Consequently
some exponent of some $X_i$ in $F$ is not divisible
by $p$, equivalently $\partial F/\partial X_i\neq0$.

Fix any such $i$, and retain the original coefficient
polynomials in
$$
F=\sum_{r=0}^h A_r(X_1,\ldots,\widehat X_i,\ldots,X_{n+1})X_i^r.
$$
There is $r\geq1$ with $p\nmid r$ and $A_r\neq0$.
Every nonzero $A_r$ with $r\geq1$ has total degree
at most $D-r<D$. Hence its value at the other original
generators is nonzero: otherwise that same coefficient
polynomial, viewed in all $n+1$ variables, would be a
smaller nonzero relation. In particular the polynomial
$$
P(T)=F(x_1,\ldots,x_{i-1},T,x_{i+1},\ldots,x_{n+1})
\in L[T],\qquad
L=k(x_1,\ldots,\widehat x_i,\ldots,x_{n+1}),
$$
is nonconstant and has $P(x_i)=0$. Thus $K=L(x_i)$
is finite algebraic over $L$. A transcendence basis of
$L/k$, chosen among its $n$ displayed generators by
Fields, Lemma \ref{fields-lemma-transcendence-degree},
is also a transcendence basis of $K/k$: algebraicity
in a tower is transitive. Its size must therefore be
$n$. All the displayed generators of $L$ are consequently
algebraically independent over $k$.

Evaluation thus identifies the polynomial ring
$B=k[X_1,\ldots,\widehat X_i,\ldots,X_{n+1}]$
with $k[x_1,\ldots,\widehat x_i,\ldots,x_{n+1}]$,
and identifies its fraction field with the original $L$.
The irreducible polynomial $F\in B[X_i]$ is primitive:
a nonunit common divisor of its coefficients would
factor $F$, with the other factor still of positive
$X_i$-degree. Gauss' lemma now makes the unchanged
polynomial $P(T)$ irreducible in $L[T]$.
Its derivative is the full sum
$$
P'(T)=\sum_{r=1}^h r A_r(x_1,\ldots,\widehat x_i,\ldots,x_{n+1})T^{r-1}.
$$
The previously chosen term is nonzero, so $P'\neq0$.
Irreducibility and $\deg P'<\deg P$ give
$\gcd(P,P')=1$. Hence $x_i$ is separable over $L$,
and $K/L$ is separable. This proves the assertion,
also when $n=0$ and the displayed basis of $L=k$ is empty.

The proof gives an exact criterion for every omitted
coordinate, not merely existence of one. For this same
minimal-degree $F$, omission of $x_j$ gives a separating
transcendence basis if and only if
$\partial F/\partial X_j\neq0$. Sufficiency was proved
for every such index. For necessity, suppose the other
$n$ elements form a separating basis and put
$L_j=k(x_1,\ldots,\widehat x_j,\ldots,x_{n+1})$.
The polynomial $F$ must depend on $X_j$, since otherwise
it would be a relation among that basis. Its coefficient
evaluation is injective, and the same primitive Gauss
argument makes its image $P_j(T)$ irreducible over $L_j$.
If $c_j$ is its actual leading coefficient and $M_j(T)$
is the monic minimal polynomial of $x_j$ over $L_j$,
then $P_j(T)=c_jM_j(T)$ with $c_j\neq0$. Separability
gives $M_j'\neq0$, so $P_j'=c_jM_j'\neq0$.
Injectivity of coefficient evaluation yields
$\partial F/\partial X_j\neq0$, as required.

There is also an exact relation-ring comparison for
these original generators. Evaluation induces
$$
k[X_1,\ldots,X_{n+1}]/(F)
\longrightarrow k[x_1,\ldots,x_{n+1}].
$$
It is surjective. If an original polynomial $H$ is in
its kernel, view $H$ in $B[X_i]$ for the index already
chosen. Over $L$, its image is divisible by the
minimal polynomial of $x_i$, hence by the original
$P=c_iM_i$ as $c_i$ is a unit of $L$. Primitive
divisibility in Gauss' lemma gives $F\mid H$ already
in $B[X_i]$. The kernel is therefore exactly $(F)$,
and the displayed map is an isomorphism with the
original evaluation map as its defining morphism.
\end{proof}


\subsection{Editorial treatment of Algebra lines 10359--10407}

\hypertarget{algebra-context-121}{}

\noindent\href{https://github.com/stacks/stacks-project/blob/a04446e57ec1fbc252a871afcec7752fb2807b14/algebra.tex\#L10359}{Original source, lines 10359--10407}. Complete original and reviewed TeX passages accompany this note. Every intervention is separately recorded; the following treatment is editorial.


\begin{lemma}
\label{editorial-source-lemma-characterize-separable-field-extensions}
Let $k$ be a field of characteristic $p > 0$.
Let $K/k$ be a field extension.
The following are equivalent:
\begin{enumerate}
\item $K$ is separable over $k$,
\item for every $k$-linearly independent subset $\{a_1, \ldots, a_m\}$
of $K$ the set $\{a^p_1, \ldots, a_m^p\}$ is $k$-linearly independent,
\item the ring $K \otimes_k k^{1/p}$ is reduced, and
\item $K$ is geometrically reduced over $k$.
\end{enumerate}
\end{lemma}

\begin{proof}
The implication (1) $\Rightarrow$ (4) is Lemma
\ref{editorial-source-lemma-separable-extension-preserves-reducedness}
applied to every field extension of $k$ as a reduced
$k$-algebra. The symmetry map $a\otimes b\mapsto b\otimes a$
has itself as inverse. The implication (4) $\Rightarrow$ (3)
then follows by using the specified field $k^{1/p}$.

For (3) $\Rightarrow$ (2), retain the original tensor
ring $B=K\otimes_k k^{1/p}$ and the original map
$$
m:B\longrightarrow K,\qquad
\sum_i\lambda_i\otimes\mu_i\longmapsto
\sum_i\lambda_i^p\mu_i^p.
$$
This is a ring homomorphism. Additivity and multiplicativity
follow from Frobenius in characteristic $p$, and balancing
follows from
$(c\lambda)^p\mu^p=\lambda^p(c\mu)^p$ for $c\in k$.
Every $\mu_i^p$ belongs to the original $k$, so the
displayed sum belongs to $K$. For every original
$z=\sum_i\lambda_i\otimes\mu_i$ one has
$$
z^p=\sum_i\lambda_i^p\otimes\mu_i^p
=\left(\sum_i\lambda_i^p\mu_i^p\right)\otimes1
=m(z)\otimes1.
$$
Here all mixed multinomial terms have coefficients
divisible by the prime $p$. The map
$K\to B$, $a\mapsto a\otimes1$, is injective, as
extension of scalars from one field to another is
faithful on vector spaces. If $B$ is reduced, $m(z)=0$
therefore implies $z^p=0$ and then $z=0$.

Let $a_1,\ldots,a_m$ be $k$-linearly independent,
and suppose $\sum_{i=1}^m c_i a_i^p=0$ with $c_i\in k$.
Let $\mu_i\in k^{1/p}$ be the unique root with
$\mu_i^p=c_i$. Then
$m(\sum_i a_i\otimes\mu_i)=0$, so the tensor is zero.
The elements $a_i\otimes1$ are independent over
$k^{1/p}$: extend the original finite family to a
$k$-basis of $K$ and extend that basis by scalars.
Consequently every $\mu_i=0$ and every $c_i=0$.
This proves (2). The argument uses the displayed
ring homomorphism, not an assertion that $m$ is
$k$-linear: in fact $m(cz)=c^p m(z)$ for the
original scalar action of $c\in k$.

For (2) $\Rightarrow$ (1), property (2) is inherited
by every intermediate field of the original $K/k$.
It suffices, by the definition of separability, to
prove that each finitely generated such field is
separably generated. Work with one of these fields,
denoted by $K$ in the remainder of this argument.
Choose a transcendence basis $x_1,\ldots,x_d$ and
retain $K'=k(x_1,\ldots,x_d)\subseteq K$. Since
$K/k$ is finitely generated and $K/K'$ is algebraic,
it is finite by Fields, Lemma
\ref{fields-lemma-algebraic-finitely-generated}.
Among these original bases choose one for which
the positive integer $[K:K']_i$ is least. Such a
least value exists because the set of bases is
nonempty and these values are positive integers.

If $K/K'$ is separable, this basis is separating.
Otherwise choose $x_{d+1}\in K$ not separable over
$K'$, and retain the intermediate field
$E=K'(x_{d+1})$. Then $[E:K']_i>1$.
The original $E/k$ is generated by the displayed
$d+1$ elements, has $x_1,\ldots,x_d$ as a
transcendence basis, and inherits (2) from $K/k$.
Apply Lemma \ref{editorial-source-lemma-mini-separability} to $E/k$,
with its integer $n$ equal to this $d$. It supplies
$1\leq j\leq d+1$ such that
$$
K''=k(x_1,\ldots,\widehat x_j,\ldots,x_{d+1})
\subseteq E
$$
has the indicated generators as a transcendence
basis and $E/K''$ is finite separable. They are
also a transcendence basis of $K/k$, since
$K/E$ is finite algebraic. The two finite towers
are $K/E/K'$ and $K/E/K''$. Fields, Lemma
\ref{fields-lemma-multiplicativity-all-degrees}
gives the full comparison
$$
[K:K']_i=[K:E]_i[E:K']_i,
\qquad
[K:K'']_i=[K:E]_i[E:K'']_i=[K:E]_i.
$$
Thus $[K:K'']_i<[K:K']_i$, contradicting the choice
of basis. This proves separable generation of
every finitely generated intermediate field, and
hence (1) for the original extension, with no
finite-generation assumption on that extension.
When $d=0$, the same argument uses an empty basis
and $K''=k$; no omitted index is out of range.

The same original tensor maps give a stronger
description even when the equivalent conditions
fail. Choose an algebraic closure $\Omega$ of $K$
and embed $k^{1/p}$ there by its unique $p$th roots.
Put $C=Kk^{1/p}\subseteq\Omega$ and
$D=K^p k\subseteq K$, retaining both composita
with their indicated embeddings. Multiplication
defines a ring map
$$
\nu:B\longrightarrow C,\qquad
\sum_i\lambda_i\otimes\mu_i\longmapsto
\sum_i\lambda_i\mu_i.
$$
Its image is the subring of finite such sums.
For a nonzero element $w$ of this subring,
$w^p=\sum_i\lambda_i^p\mu_i^p\in K\setminus\{0\}$,
so $w^{-1}=w^{p-1}/w^p$ lies in the same subring.
It is consequently a field containing $K$ and
$k^{1/p}$, and is exactly $C$. Hence $\nu$ is
surjective. Frobenius on $C$ maps into $K$ and
sends these original sums to $\sum_i\lambda_i^p\mu_i^p$.
Its image is exactly $D$: it is a field containing
$K^p$ and $k$, and each image lies in the field
generated by them. Frobenius on a field is injective,
so it is a ring isomorphism $C\longrightarrow D$.
The original $m$ is precisely its composite with
$\nu$, followed by the inclusion $D\subseteq K$.

Since the target of $\nu$ is a field, every
nilpotent of $B$ belongs to $\ker\nu$. Conversely,
if $\nu(z)=0$, then $m(z)=\nu(z)^p=0$, and the
displayed tensor identity gives $z^p=0$. Therefore
$$
\operatorname{Nil}(B)=\ker\nu=\ker m
=\{z\in B:z^p=0\}.
$$
The exact quotient maps identify
$B/\operatorname{Nil}(B)$ with $C$ by $[z]\mapsto\nu(z)$,
and with $D$ by $[z]\mapsto m(z)$. The first is a
$K$-algebra isomorphism for the original action;
the second sends the scalar $a\in K$ to $a^p\in D$.
Neither map replaces the original tensor ring or
changes its scalar structure silently. In particular
$B$ has a unique prime ideal, namely its nilradical:
primes contain all nilpotents, and its reduced
quotient is the nonzero field $C$. The equivalent
conditions of the lemma can also be expressed as
injectivity of $\nu$, or equivalently as the original
map $K\otimes_k k^{1/p}\to Kk^{1/p}$ being an
isomorphism of $K$-algebras.

Elementwise annihilation by the $p$th power in this
description does not assert that the $p$th ideal
power of the nilradical is zero. For an exact
example, take algebraically independent $u,v$
over $\mathbf F_p$, and retain
$$
s=u^p,\quad t=v^p,\quad
k=\mathbf F_p(s,t),\quad K=\mathbf F_p(u,v).
$$
Then $k^{1/p}=K$ inside an algebraic closure of $K$.
Indeed every rational function in $s,t$ has the
unique root obtained by using the same numerator
and denominator in $u,v$, since Frobenius fixes
each coefficient in $\mathbf F_p$. The elements
$u^a v^b$ for $0\leq a,b<p$ form a $k$-basis of
$K$. Independence follows on clearing the original
denominators in any proposed relation: the exponent
pairs in its polynomial terms have distinct residues
modulo $p$. They span because powers can be reduced
using $u^p=s$, $v^p=t$, and every nonzero polynomial
denominator $q(u,v)$ has inverse
$q(u,v)^{p-1}/q(u,v)^p$, with $q(u,v)^p\in k$.

The original tensor comparison is consequently
$$
B=K\otimes_k K\simeq
K[U,V]/(U^p-s,V^p-t),\qquad
1\otimes u\longleftrightarrow U,\quad
1\otimes v\longleftrightarrow V.
$$
It fixes $K$ through the left tensor factor.
Monic division gives the basis $U^aV^b$ for
$0\leq a,b<p$ on the polynomial side; the stated
basis of the right tensor factor proves the map
and its inverse on every basis vector. Put
$\epsilon=U-u$, $\eta=V-v$. The full binomial
expansions give $\epsilon^p=U^p-u^p=0$ and
$\eta^p=V^p-v^p=0$, because every intermediate
binomial coefficient is divisible by $p$.
Evaluation at $U=u,V=v$ identifies the quotient
by $(\epsilon,\eta)$ with $K$, and all elements
of this ideal have $p$th power zero. Thus it is
the nilradical $N$. Its power $N^{2p-1}$ is zero:
each generating monomial has either epsilon or
eta exponent at least $p$. But
$\epsilon^{p-1}\eta^{p-1}\neq0$, since in the
original basis its $U^{p-1}V^{p-1}$ coefficient
is $1$ and all other terms have smaller exponent
in at least one coordinate. Hence the ideal
nilpotence index is exactly $2p-1$, and in
particular $N^p\neq0$ for every prime $p$.
\end{proof}


\subsection{Editorial treatment of Algebra lines 10409--10417}

\hypertarget{algebra-context-122}{}

\noindent\href{https://github.com/stacks/stacks-project/blob/a04446e57ec1fbc252a871afcec7752fb2807b14/algebra.tex\#L10409}{Original source, lines 10409--10417}. Complete original and reviewed TeX passages accompany this note. Every intervention is separately recorded; the following treatment is editorial.


\begin{lemma}
\label{editorial-source-lemma-separably-generated-separable}
A separably generated field extension is separable.
\end{lemma}

\begin{proof}
In characteristic $p>0$, Lemma
\ref{editorial-source-lemma-separable-extension-preserves-reducedness}
applies directly to the original separably generated
extension $K/k$ and to every reduced $k$-algebra.
In particular $K\otimes_k L$ is reduced for every
field extension $L/k$. Thus $K/k$ is geometrically
reduced, and Lemma
\ref{editorial-source-lemma-characterize-separable-field-extensions}
makes this original extension separable.

In characteristic zero, every finitely generated
intermediate field $E/k$ has a finite transcendence
basis $y_1,\ldots,y_d$, and $E$ is finite algebraic
over $k(y_1,\ldots,y_d)$. Each irreducible polynomial
of positive degree over this field has nonzero
derivative: its leading derivative coefficient is
the original nonzero leading coefficient multiplied
by its positive integer degree. The derivative has
smaller degree, so irreducibility makes their greatest
common divisor equal to a unit. Every algebraic
element of $E$ is consequently separable over the
chosen rational function field. Thus $E/k$ is
separably generated. This holds for every such $E$,
which is precisely the definition of separability
of $K/k$. In fact this argument proves that every
field extension in characteristic zero is separable,
without requiring that the entire extension be
separably generated.
\end{proof}


\subsection{Editorial treatment of Algebra lines 10424--10478}

\hypertarget{algebra-context-123}{}

\noindent\href{https://github.com/stacks/stacks-project/blob/a04446e57ec1fbc252a871afcec7752fb2807b14/algebra.tex\#L10424}{Original source, lines 10424--10478}. Complete original and reviewed TeX passages accompany this note. Every intervention is separately recorded; the following treatment is editorial.


\begin{lemma}
\label{editorial-source-lemma-geometrically-reduced-finite-purely-inseparable-extension}
Let $k$ be a field. Let $S$ be a $k$-algebra.
The following are equivalent:
\begin{enumerate}
\item $k' \otimes_k S$ is reduced for every finite
purely inseparable extension $k'$ of $k$,
\item if $\operatorname{char}(k)=p>0$, then
$k^{1/p}\otimes_k S$ is reduced; if
$\operatorname{char}(k)=0$, then $S$ is reduced,
\item $k^{perf} \otimes_k S$ is reduced, where $k^{perf}$ is the
perfect closure of $k$,
\item $\overline{k} \otimes_k S$ is reduced, where $\overline{k}$ is the
algebraic closure of $k$, and
\item $S$ is geometrically reduced over $k$.
\end{enumerate}
\end{lemma}

\begin{proof}
We treat characteristic zero explicitly. Every field
extension in that characteristic is separable, as
proved in Lemma \ref{editorial-source-lemma-separably-generated-separable}.
Every finite purely inseparable extension is the
identity extension and $k^{perf}=k$. Thus (1), (2)
and (3) each say that the original $S$ is reduced.
If it is reduced, Lemma
\ref{editorial-source-lemma-separable-extension-preserves-reducedness}
makes $L\otimes_k S$ reduced for every field
extension $L/k$. This proves (5), and hence (4).
Conversely (4) implies that $S$ is reduced through
the injection $s\mapsto1\otimes s$ into
$\overline k\otimes_k S$. This proves all equivalences
in characteristic zero without using undefined
$p$th-root notation.

Now suppose $\operatorname{char}(k)=p>0$. Use an
algebraic closure $\overline k$ to contain the
specified perfect closure and $k^{1/p}$. A finite
purely inseparable extension $k'/k$ has a unique
$k$-embedding into $\overline k$, and its image lies
in $k^{perf}$, since each element has a $p$-power
in $k$. Extension of scalars by the $k$-vector
space $S$ preserves these field injections.
Consequently their original maps
$a\otimes s\mapsto a\otimes s$ prove
(5) $\Rightarrow$ (4) $\Rightarrow$ (3) $\Rightarrow$ (2)
and (3) $\Rightarrow$ (1). A subring of a reduced
ring is reduced, which is the exact assertion
used at each injection.

For (1) $\Rightarrow$ (5), let $K/k$ be any original
field extension and let $z=\sum_{i=1}^m a_i\otimes s_i$
be a nilpotent of $K\otimes_k S$. Retain the
finitely generated intermediate field
$K_0=k(a_1,\ldots,a_m)\subseteq K$. The map
$K_0\otimes_k S\to K\otimes_k S$ is injective,
so the same displayed tensor is already nilpotent
in $K_0\otimes_k S$. Apply Lemma
\ref{editorial-source-lemma-make-separably-generated} to this $K_0/k$.
It gives the original commutative field diagram
$$
\xymatrix{
K_0\ar[r]&K_0'\\
k\ar[u]\ar[r]&k'\ar[u]
}
$$
where $k'/k$ and $K_0'/K_0$ are finite purely
inseparable and $K_0'/k'$ is separably generated.
By (1), $S'=k'\otimes_k S$ is reduced. The original
comparison map is
$$
K_0'\otimes_{k'}S'\longrightarrow K_0'\otimes_k S,
\qquad b\otimes(c\otimes s)\longmapsto bc\otimes s.
$$
Its inverse sends $b\otimes s$ to
$b\otimes(1\otimes s)$. Balancing follows from
the displayed commutative field diagram, and
both composites are the identity on generating
pure tensors. Lemma
\ref{editorial-source-lemma-separable-extension-preserves-reducedness}
makes the left side reduced, so the right side
is reduced. The injection
$K_0\otimes_k S\to K_0'\otimes_k S$ now makes
the original $z$ zero. Its image in $K\otimes_k S$
is the tensor we started with. Thus every
nilpotent in that original tensor ring is zero,
proving (5).

For (2) $\Rightarrow$ (5), inject $S$ into
$k^{1/p}\otimes_k S$ by $s\mapsto1\otimes s$;
hence $S$ is reduced. Let $\mathfrak p$ be any
minimal prime of $S$ and put $W=S\setminus\mathfrak p$.
There are inverse original localization maps
$$
k^{1/p}\otimes_k S_{\mathfrak p}
\longrightarrow (1\otimes W)^{-1}(k^{1/p}\otimes_k S),
\qquad a\otimes(s/w)\longmapsto (a\otimes s)/(1\otimes w),
$$
and
$$
\frac{\sum_i a_i\otimes s_i}{1\otimes w}
\longmapsto\sum_i a_i\otimes(s_i/w).
$$
The first comes from inverting exactly $W$ in
the $S$-factor; the second comes from inverting
exactly $1\otimes W$ in the tensor ring. These
universal properties prove well-definedness,
and the formulas show both composites are
identities, with the same original denominators.
The localized tensor ring is reduced by (2).
Since $S_{\mathfrak p}$ is a field by Lemma
\ref{algebra-lemma-minimal-prime-reduced-ring}, Lemma
\ref{editorial-source-lemma-characterize-separable-field-extensions}
applied to this original field says it is
geometrically reduced over $k$. Lemma
\ref{editorial-source-lemma-generic-points-geometrically-reduced}
then proves that $S$ is geometrically reduced.
This also covers the zero ring: it has no
minimal primes and all its scalar extensions
are zero.

There are two useful exact forms of the test in
positive characteristic. First, it suffices in
(1) to test only finite purely inseparable
extensions $l/k$ with $l^p\subseteq k$.
Necessity follows from (5). Conversely every
finite list of coefficients in $k^{1/p}$ lies
in a field $l=k(b_1,\ldots,b_r)$ with $b_i^p\in k$.
This extension is finite of degree at most $p^r$,
and every element of $l$ has its $p$th power in
$k$, by applying Frobenius to its full rational
expression. If a tensor in $k^{1/p}\otimes_k S$
is nilpotent, capture its finite coefficient
list in this original $l$. Injectivity of
$l\otimes_k S\to k^{1/p}\otimes_k S$ places the
same nilpotence equation in the smaller tensor
ring. The restricted test makes it zero there
and therefore zero in the original ring.
Thus the restricted tests imply (2) and all
five conditions.

Second, the field's linear-independence criterion
extends to the original arbitrary $k$-algebra $S$:
geometric reducedness is equivalent to the
following condition. Every finite $k$-linearly
independent family $s_1,\ldots,s_r$ in $S$ has
$k$-linearly independent $p$th powers. To prove
this, retain $A=S\otimes_k k^{1/p}$ and the ring
homomorphism
$$
m_S:A\longrightarrow S,\qquad
\sum_i s_i\otimes b_i\longmapsto\sum_i s_i^p b_i^p.
$$
Balancing uses $(cs)^p b^p=s^p(cb)^p$, and the
full Frobenius identity gives
$z^p=m_S(z)\otimes1$. The original map
$S\to A$ is injective, so $\ker m_S$ is exactly
the kernel of the $p$th-power map on $A$.
That kernel is zero exactly when $A$ is reduced:
for the nontrivial direction, if $z$ is nilpotent,
choose $p^e$ at least a nilpotence exponent and
apply injectivity of Frobenius $e$ times to
$z^{p^e}=0$.

If $m_S$ is injective and
$\sum_i c_i s_i^p=0$, take the unique roots
$b_i^p=c_i$ in $k^{1/p}$. Injectivity makes
$\sum_i s_i\otimes b_i=0$, and extension of the
independent family to a vector-space basis
forces every $b_i=0$, hence every $c_i=0$.
Conversely any original $z\in A$ can be written
$z=\sum_i s_i\otimes b_i$ with $s_i$ independent:
choose a basis of the finite-dimensional span
of its original first-factor coefficients and
rewrite those coefficients in that basis,
using tensor balancing. If the asserted
independence of $p$th powers holds and $m_S(z)=0$,
then every $b_i^p=0$ and hence every $b_i=0$.
Thus $m_S$ is injective. We have proved the
equivalence with reducedness of $A$, which is
(2) after the explicit tensor-factor symmetry.
No domain, finiteness, or Noetherian hypothesis
on the original $S$ is used.
\end{proof}


\subsection{Editorial treatment of Algebra lines 10482--10604}

\hypertarget{algebra-context-124}{}

\noindent\href{https://github.com/stacks/stacks-project/blob/a04446e57ec1fbc252a871afcec7752fb2807b14/algebra.tex\#L10482}{Original source, lines 10482--10604}. Complete original and reviewed TeX passages accompany this note. Every intervention is separately recorded; the following treatment is editorial.


\subsubsection{Perfect fields}
\label{editorial-source-section-perfect-fields}

\noindent
Here is the definition.

\begin{definition}
\label{editorial-source-definition-perfect}
Let $k$ be a field. We say $k$ is {\it perfect}
if every field extension of $k$ is separable over $k$.
\end{definition}

\begin{lemma}
\label{editorial-source-lemma-perfect}
A field $k$ is perfect if and only if it is a field of characteristic $0$
or a field of characteristic $p > 0$ such that every element has a $p$th
root.
\end{lemma}

\begin{proof}
In characteristic zero every field extension is
separable, by the characteristic-zero argument in
Lemma \ref{editorial-source-lemma-separably-generated-separable}.
This is precisely the stated definition of perfect.

Suppose now that $\operatorname{char}(k)=p>0$.
If $k$ is perfect, take any $a\in k$ and its unique
$p$th root $b$ in an algebraic closure. The original
extension $k(b)/k$ is purely inseparable, since
$b^p=a$, and is separable by perfectness. A minimal
polynomial of $b$ divides $T^p-a=(T-b)^p$ over the
algebraic closure. Thus it has only one distinct
root, whereas separability makes all its roots
distinct. Its degree is therefore one, so $b\in k$.
This includes $a=0$, whose root is zero.

Conversely, suppose every $a\in k$ has a $p$th root
in $k$. Frobenius is injective on a field, so those
roots are unique and the specified extension
$k^{1/p}/k$ is the identity. For every original
field extension $K/k$ the tensor map
$$
K\otimes_k k\longrightarrow K,\qquad
\lambda\otimes a\longmapsto a\lambda
$$
has inverse $\lambda\mapsto\lambda\otimes1$.
It identifies $K\otimes_k k^{1/p}$ with the reduced
field $K$. Lemma
\ref{editorial-source-lemma-characterize-separable-field-extensions}
then makes every original $K/k$ separable, so $k$
is perfect.

In particular, in positive characteristic it
suffices to test the separability of the individual
extensions $k(a^{1/p})/k$. If $a$ has no root in $k$,
then $T^p-a$ is irreducible: a proper monic factor
over $k$ would be $(T-b)^d$ over the algebraic
closure for some $1\leq d<p$, whose coefficient
of $T^{d-1}$ is $-db$. Since $d\neq0$ in $k$,
that coefficient would put $b$ in $k$, a contradiction.
Thus every nontrivial one of these tests has degree
exactly $p$ and is inseparable. This proves the
test and its exact degree without assuming a
general extension is finitely generated.
\end{proof}

\begin{lemma}
\label{editorial-source-lemma-make-separable}
Let $K/k$ be a finitely generated field extension.
There exists a diagram
$$
\xymatrix{
K \ar[r] & K' \\
k \ar[u] \ar[r] & k' \ar[u]
}
$$
where $k'/k$, $K'/K$ are finite purely inseparable field
extensions such that $K'/k'$ is a separable field extension.
In this situation we can assume that $K' = k'K$ is the compositum,
and also that $K' = (k' \otimes_k K)_{red}$.
\end{lemma}

\begin{proof}
Lemma \ref{editorial-source-lemma-make-separably-generated} supplies
the original diagram, with its upper-right field
temporarily denoted by $K'_0$, in which $k'/k$
and $K'_0/K$ are finite purely inseparable and
$K'_0/k'$ is separably generated. By Lemma
\ref{editorial-source-lemma-separably-generated-separable}, the last
extension is separable. Retain both original
embeddings into $K'_0$ and let $L=k'K\subseteq K'_0$
be their compositum. Since $k'\subseteq L\subseteq K'_0$,
Lemma \ref{editorial-source-lemma-subextensions-are-separable} makes
$L/k'$ separable. Also $L/K$ is finite purely
inseparable, being an intermediate extension of
$K'_0/K$. Thus choosing the new upper-right field
to be $K'=L$, with its actual inclusion into
$K'_0$, proves the compositum assertion.

For the tensor assertion, first treat characteristic
zero. A purely inseparable extension is then the
identity, so $k'=k$ and $K'_0=K$. The multiplication
map $k\otimes_k K\to K$, $a\otimes b\mapsto ab$,
is an isomorphism with inverse $b\mapsto1\otimes b$.
Its nilradical is zero, and it is the required map.

In characteristic $p>0$, put $B=k'\otimes_k K$.
Because the original extension $k'/k$ is finite
purely inseparable, there is a single $q=p^e$
with $(k')^q\subseteq k$: choose finitely many
field generators, take the largest of their
inseparability exponents, and apply Frobenius
to their full rational expressions. Retain
the multiplication map
$$
\mu:B\longrightarrow K'_0,\qquad
\sum_i a_i\otimes b_i\longmapsto\sum_i a_i b_i.
$$
It is balanced because the original field diagram
commutes. For every original $z=\sum_i a_i\otimes b_i$
one has the full identity
$$
z^q=\sum_i a_i^q\otimes b_i^q
=1\otimes\left(\sum_i a_i^q b_i^q\right),
\qquad
\mu(z)^q=\sum_i a_i^q b_i^q\in K.
$$
The equality for $q$ follows by iterating the
$p$th-power binomial identity; each intermediate
mixed coefficient is divisible by $p$. The map
$K\to B$, $b\mapsto1\otimes b$, is injective.
If $\mu(z)=0$, the coefficient sum in $K$ is
zero, and hence $z^q=0$. Conversely every
nilpotent maps to zero in the field $K'_0$.
Therefore $\ker\mu=\operatorname{Nil}(B)$.

The image of $\mu$ is the subring of finite sums
of products $a_i b_i$. It is a field: for a
nonzero such sum $w$, its $q$th power is a
nonzero element of $K$, and
$w^{-1}=w^{q-1}/w^q$ belongs to the same subring.
The image contains $k'$ and $K$, so it is exactly
the original $L=k'K$. Consequently the original
map gives the canonical $K$-algebra isomorphism
$$
(k'\otimes_k K)_{red}\longrightarrow L,
\qquad [\sum_i a_i\otimes b_i]\longmapsto\sum_i a_i b_i.
$$
Its inverse sends a finite sum in $L$ to the
class of the corresponding tensor; if two such
sums are equal, their difference lies in the
kernel just calculated, proving independence
of the representation. Thus both asserted
choices of $K'$ agree through this exact map.

This last comparison needs neither finite
generation of $K/k$ nor finiteness of the purely
inseparable extension in the first tensor factor.
Indeed let $l/k$ be any algebraic purely
inseparable extension and embed it, by its unique
roots, in an algebraic closure of the original
$K$. Every individual tensor in $l\otimes_k K$
uses finitely many coefficients $a_i\in l$,
so there is a $q=p^e$ for that finite list with
all $a_i^q\in k$. The same two displayed power
identities prove that the multiplication kernel
is exactly the nilradical, and the same inverse
formula for each nonzero image proves that the
image is the original compositum $lK$.
Thus $(l\otimes_k K)_{red}\simeq lK$ with the
same explicit map in this full generality.
Only the exponent may now depend on the tensor.

Finite generation cannot, however, be discarded
from the earlier assertion that one finite
purely inseparable base extension makes the
whole receiving extension separable. Take
$$
k=\mathbf F_p(t),\qquad
K=\bigcup_{r\geq0}\mathbf F_p(t^{1/p^r})
$$
inside a fixed algebraic closure, with $t$
transcendental and all root inclusions retained.
Every finite purely inseparable $k'/k$ has a
unique image in $K$: a root of an original
rational function $A(t)/B(t)$ of order $p^r$
is $A(t^{1/p^r})/B(t^{1/p^r})$, since its
coefficients belong to $\mathbf F_p$.
The elements of $k'$ have
$p^e$th powers in $k$ for some finite $e$.
The element $t^{1/p^{e+1}}$ is not in $k'$:
otherwise its $p^e$th power $t^{1/p}$ would
belong to $k$. The latter is impossible since
a rational function in $t$ cannot have $p$th
power $t$. To see this directly, writing such
a function as $A(t)/B(t)$ would give
$A(t)^p=tB(t)^p$; the exponents on the left
are divisible by $p$ and those on the right
are congruent to $1$ modulo $p$, so a nonzero
polynomial equality is impossible.
Thus $K/k'$ is a nontrivial purely inseparable
extension and cannot be separable. Its
compositum with the original $k'$ is still
$K$. Moreover any purported larger receiving
field separable over $k'$ would have the
subextension $K/k'$ separable, again impossible.
This proves the necessity of some restriction
on the original $K/k$ for the finite-base-change
existence assertion, while retaining the
unrestricted tensor-quotient comparison.
\end{proof}

\begin{lemma}
\label{editorial-source-lemma-perfection}
\begin{slogan}
Every field has a unique perfect closure.
\end{slogan}
For every field $k$ there exists a purely inseparable extension
$k'/k$ such that $k'$ is perfect. The field extension
$k'/k$ is unique up to unique isomorphism.
\end{lemma}

\begin{proof}
In characteristic zero $k$ itself is perfect by
Lemma \ref{editorial-source-lemma-perfect}, and a purely inseparable
extension must be the identity. This gives
existence and uniqueness, including the specified
embedding of $k$.

Let $\operatorname{char}(k)=p>0$, and fix an
algebraic closure $\Omega$ of the original $k$.
For every integer $n\geq0$ put
$$
E_n=\{b\in\Omega:b^{p^n}\in k\}.
$$
The Frobenius identities show that this is a
field: sums and products have $p^n$th powers
in $k$, as does the inverse of a nonzero
element. One has $E_0=k$ and $E_n\subseteq E_{n+1}$.
Every $a\in k$ has a unique $p^n$th root in
$\Omega$, because $T^{p^n}-a$ has a root there
and the difference of any two roots has zero
$p^n$th power in a field. Thus $E_n$ has exactly
properties (a) and (b) of the stated construction.

To compare this with the original abstract
construction, let $k_n$ denote a copy of the
underlying field $k$, with structural map
$i_n:k\to k_n$, $a\mapsto a^{p^n}$.
The transition map $f_n:k_n\to k_{n+1}$ is
$r\mapsto r^p$, and
$f_n(i_n(a))=a^{p^{n+1}}=i_{n+1}(a)$.
The map
$$
\theta_n:k_n\longrightarrow E_n,\qquad
r\longmapsto r^{1/p^n}
$$
is a field isomorphism. Uniqueness of roots
proves additivity and multiplicativity;
surjectivity follows from the definition of
$E_n$, and injectivity follows from Frobenius.
It respects the original $k$-structure since
$\theta_n(i_n(a))=a$, and it respects the
transitions since
$\theta_{n+1}(r^p)=\theta_n(r)$.
Accordingly the union of the actual subfields
$E_n$ realizes the directed colimit of the
original copies with their Frobenius transition
maps. No literal inclusion of a copy with its
different structural map is being assumed.

Set $k'=\bigcup_{n\geq0}E_n$. This is a field,
and every one of its elements has a $p$-power
in the original $k$, so $k'/k$ is purely
inseparable. If $b\in E_n$, its unique $p$th
root in $\Omega$ lies in $E_{n+1}$, since its
$p^{n+1}$th power equals $b^{p^n}\in k$.
Thus Frobenius on $k'$ is surjective, and
Lemma \ref{editorial-source-lemma-perfect} makes $k'$ perfect.

For uniqueness, let $L/k$ be any perfect purely
inseparable extension, with its original
structural embedding of $k$. Every $a\in k$
has a unique $p^n$th root in $L$. Define
$\phi_n:E_n\to L$ by sending $b$ to that
root of $b^{p^n}\in k$. Frobenius and root
uniqueness prove that $\phi_n$ preserves
addition, multiplication, $1$, and the given
embedding of $k$. It is injective as a
nonzero field homomorphism. These maps agree
on the inclusions $E_n\subseteq E_{n+1}$,
again by root uniqueness, and give
$\phi:k'\to L$. Every $c\in L$ has some
$c^{p^n}\in k$, so it is the image of the
unique corresponding root in $E_n$.
Hence $\phi$ is surjective. Any $k$-homomorphism
from $k'$ to $L$ must take every such root
to that same unique root, so it equals $\phi$.
This proves uniqueness up to unique
$k$-isomorphism.

The construction proves the stronger universal
property with its maps: for every perfect
field $P$ and every field homomorphism
$\iota:k\to P$, there is exactly one field
homomorphism $\iota^{perf}:k'\to P$ extending
$\iota$. On $b\in E_n$ it is the unique
$p^n$th root in $P$ of $\iota(b^{p^n})$.
The preceding addition, multiplication and
compatibility proofs apply verbatim; they
did not use pure inseparability of the
receiving field for existence, only for
surjectivity. This also proves functoriality:
two composites that agree on $k$ send every
root to the same unique root. In
characteristic zero the assertion is just
the original map $\iota$ on $k'=k$.
\end{proof}

\begin{definition}
\label{editorial-source-definition-perfection}
Let $k$ be a field. The field extension $k'/k$ of Lemma \ref{editorial-source-lemma-perfection}
is called the {\it perfect closure} of $k$. Notation $k^{perf}/k$.
\end{definition}

\noindent
Let $l/k$ be any algebraic purely inseparable
extension. In positive characteristic, send an
original $a\in l$, with $a^{p^n}\in k$, to
the unique root of $T^{p^n}-a^{p^n}$ in
$k^{perf}$. The value is independent of the
chosen exponent, and a common exponent for
two elements proves additivity and
multiplicativity by Frobenius. This defines
the unique $k$-embedding $l\to k^{perf}$.
In characteristic zero $l=k$ and this is
the identity. Retain this actual embedding.
The extension $k^{perf}/l$ is purely
inseparable because every element has a
$p$-power in $k\subseteq l$, and its upper
field is perfect. It is therefore a
perfect closure of $l$. More explicitly,
the universal property just proved gives
maps $k^{perf}\to l^{perf}$ extending
$k\to l^{perf}$ and $l^{perf}\to k^{perf}$
extending the displayed embedding of $l$.
The first map fixes the original $l$
under these embeddings, since roots of
its elements' powers in $k$ are unique.
Both composites consequently fix the
base fields and every adjoining root,
so are identities. These are the
specified inverse isomorphisms of
perfect closures, not merely an
abstract equality of field types.

\begin{lemma}
\label{editorial-source-lemma-perfect-reduced}
Let $k$ be a perfect field.
Any reduced $k$-algebra is geometrically reduced over $k$.
Let $R$, $S$ be $k$-algebras.
Assume both $R$ and $S$ are reduced.
Then the $k$-algebra $R \otimes_k S$ is reduced.
\end{lemma}

\begin{proof}
By Lemma \ref{editorial-source-lemma-perfect}, either $k$ has
characteristic zero or its Frobenius map is
surjective. In characteristic zero the test
in Lemma
\ref{editorial-source-lemma-geometrically-reduced-finite-purely-inseparable-extension}
is precisely reducedness of the original $S$.
In positive characteristic $k^{1/p}=k$ through
the actual root embedding, and the same test
is reducedness of $k\otimes_k S\simeq S$,
with maps $a\otimes s\mapsto as$ and
$s\mapsto1\otimes s$. Thus every reduced
$k$-algebra is geometrically reduced. Applying
Lemma \ref{editorial-source-lemma-geometrically-reduced-any-reduced-base-change}
to the original reduced $R$ and geometrically
reduced $S$ proves that $R\otimes_k S$ is
reduced, including either zero-ring case.

The hypothesis on $k$ is also necessary for
this conclusion to hold for every pair of
reduced $k$-algebras, or even for every pair
of fields over $k$. If $k$ is not perfect,
Lemma \ref{editorial-source-lemma-perfect} gives characteristic
$p>0$ and an $a\in k$ with no $p$th root in
$k$. Retain $L=k(b)$, $b^p=a$. As proved in
Lemma \ref{editorial-source-lemma-perfect}, $T^p-a$ is
irreducible, so $1,b,\ldots,b^{p-1}$ is a
$k$-basis of $L$. The original map
$$
L[T]/(T^p-a)\longrightarrow L\otimes_k L,
\qquad T\longmapsto1\otimes b
$$
fixes $L$ through the left factor and is an
isomorphism by the corresponding bases.
Its inverse sends $c\otimes\sum_i d_i b^i$
to $\sum_i cd_i T^i$, for $d_i\in k$.
The element $\delta=1\otimes b-b\otimes1$
has $\delta^p=1\otimes a-a\otimes1=0$,
but is nonzero because its inverse image
$T-b$ has degree $1<p$ in the monic
quotient. Consequently the tensor of the
two original reduced fields is not reduced.
The same example shows that $L$ is not
geometrically reduced. Thus perfectness is
equivalent both to every reduced $k$-algebra
being geometrically reduced and to every
tensor product of two reduced $k$-algebras
being reduced.

For arbitrary original $k$-algebras $R,S$
over this perfect field, their nilradicals
$N_R=\operatorname{Nil}(R)$ and
$N_S=\operatorname{Nil}(S)$ give the exact
formula in $T=R\otimes_k S$
$$
\operatorname{Nil}(T)=J,
\qquad
J=\operatorname{im}(N_R\otimes_k S\to T)
+\operatorname{im}(R\otimes_k N_S\to T).
$$
This is a sum of ideals, with no assertion
that the sum is direct. The tensor quotient
map $T\to (R/N_R)\otimes_k(S/N_S)$ is
surjective and kills $J$. Its inverse
after passage to $T/J$ sends
$[r]\otimes[s]$ to $[r\otimes s]$.
Changing either representative changes
the tensor by an element of the indicated
ideal, so this is well-defined; balancing
and both composites follow on pure tensors.
Thus the kernel is exactly $J$.

Every element of $J$ is a finite sum
$z=z_1+\cdots+z_m$ of tensors having at
least one nilpotent factor. Choose positive
integers $d_i$ with $z_i^{d_i}=0$, and put
$D=1+\sum_{i=1}^m(d_i-1)$. Its full
multinomial expansion is
$$
z^D=
\sum_{\substack{a_1+\cdots+a_m=D\\a_i\geq0}}
\frac{D!}{a_1!\cdots a_m!}
z_1^{a_1}\cdots z_m^{a_m}=0.
$$
Indeed at least one $a_i\geq d_i$ in each
summand. The coefficients are the original
integers; no factorial is inverted in the
receiving ring. The empty sum is zero.
Hence $J\subseteq\operatorname{Nil}(T)$.
The proved quotient is reduced by the
first part, applied to $R/N_R$ and $S/N_S$.
Every nilpotent of $T$ therefore maps to
zero in that quotient and belongs to $J$,
proving the reverse inclusion and the
claimed formula with its exact quotient
isomorphism.
\end{proof}












\subsection{Editorial treatment of Algebra lines 10614--10630}

\hypertarget{algebra-context-125}{}

\noindent\href{https://github.com/stacks/stacks-project/blob/a04446e57ec1fbc252a871afcec7752fb2807b14/algebra.tex\#L10614}{Original source, lines 10614--10630}. Complete original and reviewed TeX passages accompany this note. Every intervention is separately recorded; the following treatment is editorial.


\begin{lemma}
\label{editorial-source-lemma-surjective-locally-nilpotent-kernel}
Let $\varphi : R \to S$ be a surjective map with locally nilpotent kernel.
Then $\varphi$ induces a homeomorphism of spectra and isomorphisms
on residue fields. For any ring map $R \to R'$ the ring map
$R' \to R' \otimes_R S$ is surjective with locally nilpotent kernel.
\end{lemma}

\begin{proof}
Put $I=\ker\varphi$ and retain the isomorphism
$R/I\to S$, $[r]\mapsto\varphi(r)$, supplied
by surjectivity. Every prime of $R$ contains
$I$, since each element of $I$ is nilpotent.
The inverse maps on primes are therefore
$$
\mathfrak q\longmapsto\varphi^{-1}(\mathfrak q),
\qquad \mathfrak p\longmapsto\varphi(\mathfrak p).
$$
The second is a prime ideal because its
quotient is the original domain $R/\mathfrak p$.
Surjectivity and $I\subseteq\mathfrak p$
show both composites are identities.
Moreover $D(\varphi(r))$ corresponds to
$D(r)$ for every original $r\in R$;
these sets form bases, so the bijection
and its inverse are continuous.

For corresponding primes $\mathfrak p$
and $\mathfrak q$, the residue-field map is
the explicit isomorphism
$$
\operatorname{Frac}(R/\mathfrak p)
\longrightarrow\operatorname{Frac}(S/\mathfrak q),
\qquad
\frac{r+\mathfrak p}{t+\mathfrak p}
\longmapsto
\frac{\varphi(r)+\mathfrak q}{\varphi(t)+\mathfrak q},
\quad t\notin\mathfrak p.
$$
The inverse lifts numerator and denominator
along $\varphi$; different lifts differ by
$I\subseteq\mathfrak p$, so the fractions
are unchanged. A nonzero denominator on
either side remains nonzero on the other.

For an arbitrary original map $f:R\to R'$,
there are inverse algebra maps
$$
R'\otimes_R S\longrightarrow R'/IR',\qquad
a\otimes\varphi(r)\longmapsto[af(r)],
\qquad
R'/IR'\longrightarrow R'\otimes_R S,\quad
[a]\longmapsto a\otimes1.
$$
Changing a lift $r$ of an element of $S$
adds an element of $I$, and balancing
uses the same $f(r)$, so the first map
is well-defined. The second kills $IR'$
because $f(i)\otimes1=1\otimes\varphi(i)=0$
for $i\in I$. The formulas prove both
composites are identities. Thus the
original map $R'\to R'\otimes_R S$ is
surjective with kernel exactly $IR'$.

Each element of that kernel is a finite
sum $\sum_{j=1}^m a_j f(i_j)$ with $i_j\in I$.
Choose $d_j\geq1$ with $i_j^{d_j}=0$.
In the full multinomial expansion of its
$D$th power, for $D=1+\sum_j(d_j-1)$,
every term contains some $f(i_j)^{d_j}$,
and is zero. Thus $IR'$ is locally
nilpotent, without a uniform exponent
being assumed for the whole ideal.
Applying the first part to this actual
quotient proves the asserted base-change
conclusions, with the same prime and
residue-field maps.

The kernel hypothesis is also necessary
for a surjective ring map to induce a
homeomorphism onto the entire spectrum:
its image is $V(I)$, so surjectivity on
primes requires $I$ to lie in every
prime, equivalently in the nilradical.
In addition, the given quotient induces
an exact isomorphism of reduced rings
$R_{red}\to S_{red}$. Indeed if
$\varphi(r)$ is nilpotent, then $r^n\in I$
for some $n$, and $(r^n)^m=0$ for some
$m$, so $r$ is nilpotent. The converse
is immediate under a ring homomorphism.
The preimage of $\operatorname{Nil}(S)$
is consequently exactly
$\operatorname{Nil}(R)$, proving the
stated reduced-quotient isomorphism and
the same assertion after each base change.
\end{proof}


\subsection{Editorial treatment of Algebra lines 10632--10715}

\hypertarget{algebra-context-126}{}

\noindent\href{https://github.com/stacks/stacks-project/blob/a04446e57ec1fbc252a871afcec7752fb2807b14/algebra.tex\#L10632}{Original source, lines 10632--10715}. Complete original and reviewed TeX passages accompany this note. Every intervention is separately recorded; the following treatment is editorial.


\begin{lemma}
\label{editorial-source-lemma-powers-field}
\begin{reference}
\cite[Lemma 3.1.6]{Alper-adequate}
\end{reference}
Let $k'/k$ be a field extension. The following are equivalent
\begin{enumerate}
\item for each $x \in k'$ there exists an $n > 0$ such that $x^n \in k$, and
\item $k' = k$ or $k$ and $k'$ have characteristic $p > 0$ and
either $k'/k$ is a purely inseparable extension or
$k$ and $k'$ are algebraic extensions of $\mathbf{F}_p$.
\end{enumerate}
\end{lemma}

\begin{proof}
Each case of (2) implies (1). For $k'=k$
take $n=1$. For a purely inseparable
extension in characteristic $p>0$, each
original element has a $p$-power in $k$.
If both fields are algebraic over
$\mathbf F_p$, a nonzero $x\in k'$
lies in the finite field $\mathbf F_p(x)$,
whose multiplicative group is finite;
some positive power of $x$ is $1\in k$.
For $x=0$ take $n=1$.

Assume (1) and retain the original extension
$k'/k$. Each $x\in k'$ is algebraic over
$k$, being a root of $T^n-x^n\in k[T]$.
Let $E$ be its maximal separable
subextension from Fields, Lemma
\ref{fields-lemma-separable-first}.
In positive characteristic, if $E=k$
then $k'/k$ is purely inseparable and
we have finished. In characteristic
zero $E=k'$, so $E=k$ is precisely the
identity case. Otherwise choose
$x\in E\setminus k$ and keep the finite
separable intermediate field $L=k(x)$.
Property (1) is inherited by $L/k$.
It is enough to show that $k$ has
positive characteristic and is algebraic
over its prime field: the original
$k'/k$ is already algebraic, and then
$k'$ is algebraic over the same prime
field by transitivity. This proves the
reduction for infinite $k'/k$ as well.

First let $z$ be any element of $L\setminus k$.
There are $k$-embeddings $\sigma,\tau$ of
$k(z)$ into an algebraic closure $\overline k$
with $A=\sigma(z)\neq B=\tau(z)$, because
its separable minimal polynomial has
degree at least two. This is the
embedding count of Fields, Lemma
\ref{fields-lemma-count-embeddings}.
Both $z$ and $z+1$ are nonzero, since
neither belongs to $k$. Choose the
original positive integers $n,m$ with
$z^n,(z+1)^m\in k$. Their images give
$A^n=B^n$ and $(A+1)^m=(B+1)^m$.
Thus, with all denominators nonzero,
$$
\zeta=A/B,\qquad \zeta'=(A+1)/(B+1),
\qquad \zeta^n=1,\quad(\zeta')^m=1.
$$
One has $\zeta'\neq1$ because $A\neq B$.
The full identity
$$
\zeta'(B+1)=A+1=\zeta B+1
$$
gives $B(\zeta'-\zeta)=1-\zeta'$.
Since the right side is nonzero,
$\zeta'-\zeta\neq0$, and hence
$$
B=\frac{1-\zeta'}{\zeta'-\zeta}.
$$
Both roots of unity are algebraic over
the prime field, so this formula makes
$B$ algebraic over that field. The
embedding $\tau$ fixes the prime field;
pulling back an annihilating polynomial
through the injective $\tau$ proves that
the original $z$ is algebraic over it.
This applies to every $z\in L\setminus k$.
For any $a\in k$, both $x$ and $x+a$
lie outside $k$ and are algebraic over
the prime field, so their difference
$a=(x+a)-x$ is algebraic. It follows
that $k$, and then the original $k'$,
are algebraic over the prime field.

It remains to exclude characteristic
zero. Use the fixed original $x\in L\setminus k$
and two embeddings as above, with
$A=\sigma(x)$, $B=\tau(x)$, and the
same $\zeta,\zeta'$. Both $A=\zeta B$
and $B=(1-\zeta')/(\zeta'-\zeta)$
belong to the number field
$F=\mathbf Q(\zeta,\zeta')$.
This field contains only finitely many
roots of unity, for the following exact
reason. Put $d=[F:\mathbf Q]$. The monic
minimal polynomial of a root of unity
in $F$ has degree $e\leq d$ and integer
coefficients: it divides some $T^N-1$
in $\mathbf Q[T]$, and the primitive
Gauss lemma gives a factor in $\mathbf Z[T]$
whose leading coefficient divides $1$.
Every complex root $\xi_i$ of that
minimal polynomial is also an $N$th
root of unity and has absolute value
$1$. Its coefficient of $T^{e-j}$ is
the full signed sum
$$
(-1)^j\sum_{1\leq i_1<\cdots<i_j\leq e}
\xi_{i_1}\cdots\xi_{i_j},
$$
so its absolute value is at most
$\binom ej$. There are finitely many
integer coefficient lists with these
bounds and $e\leq d$, and each resulting
nonzero polynomial has finitely many
roots. This proves the required
finiteness inside $F$.

For every original rational $y$, one
has $x+y\notin k$ and $B+y\neq0$.
Property (1) gives a positive exponent
for $x+y$, so
$$
\zeta_y=\frac{A+y}{B+y}
=\frac{\zeta B+y}{B+y}\in F
$$
is a root of unity, and $\zeta_y\neq1$.
Infinitely many rational $y$ and only
finitely many roots of unity in $F$
give distinct $y,y'$ with
$\zeta_y=\zeta_{y'}$. Subtracting the
two original identities gives
$$
y-y'=(A+y)-(A+y')
=\zeta_y\big((B+y)-(B+y')\big)
=\zeta_y(y-y').
$$
Since $y-y'\neq0$, this forces
$\zeta_y=1$, a contradiction. Thus the
nontrivial separable case has positive
characteristic, and the already proved
algebraicity gives exactly (2).

The argument also determines when a
single positive exponent can work for
all elements. Such an exponent exists
if and only if $k'=k$, or the fields
have characteristic $p>0$ and either
$k'/k$ is purely inseparable of bounded
exponent, or $k'$ is a finite field.
The converse is immediate with exponent
$1$, a common $p$-power, or $|k'|-1$,
respectively. For necessity retain a
single original exponent $N$ for all
elements. If a nontrivial separable
intermediate field exists, choose
$A=\sigma(x)\neq B=\tau(x)$ as above.
For every $a\in k$, the quotient
$(A+a)/(B+a)$ is an $N$th root of unity.
These quotients are distinct for
distinct $a$: cross multiplication
gives the full identity
$$
(A+a)(B+b)-(A+b)(B+a)=(A-B)(b-a).
$$
Their denominators are nonzero because
$x\notin k$. Hence $|k|\leq N$, as a
nonzero degree-$N$ polynomial has at
most $N$ roots. The field $k$ is finite.
Every element of $k'$ is then among
the roots of the finitely many original
polynomials $T^N-c$, $c\in k$. For
$c=0$ the only root is zero; for
each of the $|k|-1$ nonzero values
there are at most $N$ roots. Thus
$|k'|\leq1+N(|k|-1)$, and $k'$ is finite.

In the remaining nonidentity case
$k'/k$ is purely inseparable in
characteristic $p$. Write the original
$N=p^e m$ with $p\nmid m$. For every
nonzero $x\in k'$, put $u=x^{p^e}$.
It is purely inseparable over $k$ and
satisfies $u^m=x^N\in k$. The polynomial
$T^m-x^N$ has nonzero derivative at
each root, since $m$ and each root
are nonzero. Hence $u$ is also
separable over $k$, and belongs to
$k$. The zero element satisfies the
same conclusion directly. Therefore
$(k')^{p^e}\subseteq k$, giving a
uniform inseparability exponent from
the full original integer $N$.

For finite original fields the fixed-exponent
condition has an exact form. Put $q=|k|$
and $q'=|k'|$. The multiplicative group
of a finite field is cyclic, as follows
directly from the root bound. For a
finite subgroup $G$ of a field's
multiplicative group, let $h$ be the
least common multiple of its element
orders. For every prime $\ell$ dividing
$h$, choose an element whose order
has the largest $\ell$-power $\ell^{a_\ell}$
occurring in $h$, and raise that element
to the prime-to-$\ell$ part of its
order. The result has order exactly
$\ell^{a_\ell}$. The product of these
elements has order $h$: if its $m$th
power is $1$, raising to
$h/\ell^{a_\ell}$ shows
$\ell^{a_\ell}\mid m$ for each $\ell$,
since $h/\ell^{a_\ell}$ is coprime
to $\ell$. Conversely its $h$th power
is $1$. All elements of $G$ are
roots of $T^h-1$, so $|G|\leq h$;
the element of order $h$ gives
$|G|\geq h$. Thus it generates $G$.
For the trivial group the empty
product has order $1$, giving the
same conclusion.

Choose a generator $g$ of the original
$(k')^\times$, of order $q'-1$.
Inside $k'$, the nonzero roots of
$T^{q-1}-1$ are exactly $k^\times$:
all its $q-1$ elements are roots,
and a degree-$q-1$ polynomial has
no additional ones. Therefore
\[
\begin{split}
(\forall x\in k')\ x^N\in k
&\quad\Longleftrightarrow\quad g^N\in k^\times\\
&\quad\Longleftrightarrow\quad g^{N(q-1)}=1\\
&\quad\Longleftrightarrow\quad (q'-1)\mid N(q-1).
\end{split}
\]
The zero element satisfies the condition
for every positive $N$. If $d=[k':k]$,
then $q'=q^d$, by counting the coordinate
vectors of a $k$-basis. Retaining the
full factorization
$$
q'-1=q^d-1=(q-1)(1+q+\cdots+q^{d-1}),
$$
the displayed divisibility is equivalently
$(1+q+\cdots+q^{d-1})\mid N$. The least
positive common exponent is consequently
the full sum $1+q+\cdots+q^{d-1}$.
This also proves the sharp cardinality
bound above: equality holds for this
least exponent, with every contribution
from the original zero and nonzero
elements retained.
\end{proof}


\subsection{Editorial treatment of Algebra lines 10717--10767}

\hypertarget{algebra-context-127}{}

\noindent\href{https://github.com/stacks/stacks-project/blob/a04446e57ec1fbc252a871afcec7752fb2807b14/algebra.tex\#L10717}{Original source, lines 10717--10767}. Complete original and reviewed TeX passages accompany this note. Every intervention is separately recorded; the following treatment is editorial.


\begin{lemma}
\label{editorial-source-lemma-powers}
Let $\varphi : R \to S$ be a ring map. If
\begin{enumerate}
\item for any $x \in S$ there exists $n > 0$ such that
$x^n$ is in the image of $\varphi$, and
\item $\Ker(\varphi)$ is locally nilpotent,
\end{enumerate}
then $\varphi$ induces a homeomorphism on spectra and induces residue
field extensions satisfying the equivalent conditions of
Lemma \ref{editorial-source-lemma-powers-field}.
\end{lemma}

\begin{proof}
Write $I=\ker\varphi$. First use only
condition (1). Every original $x\in S$
has an equation $x^n-\varphi(y)=0$
for some $n>0$ and $y\in R$. This is
a monic equation over the actual image
$R/I$, so $R/I\to S$ is an integral
injection. Lemma
\ref{editorial-source-lemma-integral-overring-surjective}
therefore makes $\Spec(S)\to\Spec(R/I)$
surjective. Under the original quotient
map, $\Spec(R/I)$ identifies with
the closed subset $V(I)\subseteq\Spec(R)$.

To prove injectivity, suppose primes
$\mathfrak q,\mathfrak q'$ have the
same contraction $\mathfrak p$ and
$x\in\mathfrak q\setminus\mathfrak q'$.
Choose the original equation
$x^n=\varphi(y)$ from (1). It gives
$y\in\mathfrak p$, and therefore
$x^n\in\mathfrak q'$, contradicting
primality and $x\notin\mathfrak q'$.
Interchanging the two primes proves
equality. For this same equation,
the resulting bijection sends $D(x)$
exactly onto $D(y)\cap V(I)$: a prime
contains $x$ if and only if its
contraction contains $y$. Hence it
is open for the subspace topology
on $V(I)$, and is a homeomorphism
onto that closed subset, since the
original map on spectra is continuous.
Condition (2) says precisely that
$I\subseteq\operatorname{Nil}(R)$,
so $V(I)=\Spec(R)$ and gives the
homeomorphism in the statement.

Retain corresponding original primes
$\mathfrak q$ and $\mathfrak p$.
The map $R/\mathfrak p\to S/\mathfrak q$
is injective and thus extends to the
specified residue-field embedding.
Let $x$ in the receiving residue
field be represented by $y/z$ with
$y,z\in S$, $z\notin\mathfrak q$.
Choose original $a,b\in R$ and
positive $n,m$ with
$y^n=\varphi(a)$, $z^m=\varphi(b)$.
Then $b\notin\mathfrak p$, because
its image is $z^m\notin\mathfrak q$.
The full identity in $S_{\mathfrak q}$ is
$$
(y/z)^{nm}
=\frac{(y^n)^m}{(z^m)^n}
=\frac{\varphi(a)^m}{\varphi(b)^n}
=\varphi_{\mathfrak p,\mathfrak q}(a^m/b^n).
$$
Passing to residue fields proves
that the original $x^{nm}$ belongs
to the image of $\kappa(\mathfrak p)$.
Lemma \ref{editorial-source-lemma-powers-field} now
applies to these exact residue fields.
This part did not require (2).

Thus (1) alone gives a homeomorphism
onto $V(\ker\varphi)$ and the same
residue-field conclusion; under (1),
(2) is equivalent to surjectivity
on all primes of $R$. No universal
homeomorphism follows from these two
conditions alone. For an exact
counterexample take the original
inclusion
$$
R=\mathbf F_2\longrightarrow
S=\mathbf F_2[\alpha]/(\alpha^2+\alpha+1).
$$
The polynomial has no root in
$\mathbf F_2$, so $S$ is a field
with elements $0,1,\alpha,\alpha^2$,
and $\alpha^3=1$. Hence every
element of $S$ has cube in $R$,
and the kernel is zero. After the
original base change $R\to S$,
the tensor ring has the comparison
$$
S\otimes_R S\simeq S[T]/(T^2+T+1),
\qquad1\otimes\alpha\longleftrightarrow T,
$$
fixing $S$ through the first factor.
The bases $1,\alpha$ and $1,T$
prove this map and its inverse.
Retain the two distinct roots
$\alpha,\alpha^2$ and the full
factorization
$T^2+T+1=(T-\alpha)(T-\alpha^2)$.
Evaluation gives
$$
S[T]/(T^2+T+1)\longrightarrow S\times S,
\qquad[H]\longmapsto(H(\alpha),H(\alpha^2)).
$$
Its inverse sends $(a,b)$ to
$$
\left[
a\frac{T-\alpha^2}{\alpha-\alpha^2}
+b\frac{T-\alpha}{\alpha^2-\alpha}
\right].
$$
The original denominators are nonzero
in $S$; evaluating verifies both
coordinates, and a polynomial zero
at both roots is divisible by the
original product of the two distinct
linear factors. These are inverse
ring maps. The base-changed spectrum
therefore has two points mapping
to the single point of $\Spec(S)$,
and cannot be a homeomorphism.
\end{proof}


\subsection{Editorial treatment of Algebra lines 10769--10811}

\hypertarget{algebra-context-128}{}

\noindent\href{https://github.com/stacks/stacks-project/blob/a04446e57ec1fbc252a871afcec7752fb2807b14/algebra.tex\#L10769}{Original source, lines 10769--10811}. Complete original and reviewed TeX passages accompany this note. Every intervention is separately recorded; the following treatment is editorial.


\begin{lemma}
\label{editorial-source-lemma-2-3-ring-map}
Let $\varphi : R \to S$ be a ring map. Assume
\begin{enumerate}
\item[(a)] $S$ is generated as an $R$-algebra by elements $x$ such
that $x^2, x^3 \in \varphi(R)$, and
\item[(b)] $\Ker(\varphi)$ is locally nilpotent,
\end{enumerate}
Then $\varphi$ induces isomorphisms on residue fields and
a homeomorphism of spectra. For any ring map $R \to R'$
the ring map $R' \to R' \otimes_R S$ also satisfies (a) and (b).
\end{lemma}

\begin{proof}
For every original generator $x$, choose $a_x,b_x\in R$
with $\varphi(a_x)=x^2$ and $\varphi(b_x)=x^3$.
The monic equation $T^2-a_x$ proves integrality of
$x$ over $R$. Every element of the generated algebra
involves finitely many such generators, so the
integral-closure subring property proves that
$R\to S$ is integral. Put $I=\ker\varphi$.
The integral injection $R/I\to S$ is surjective
on spectra by Lemma \ref{editorial-source-lemma-integral-overring-surjective}.
Since $I$ is locally nilpotent, $V(I)=\Spec(R)$.
Thus the original map on spectra is surjective,
and is closed by Lemmas \ref{editorial-source-lemma-integral-going-up}
and \ref{editorial-source-lemma-going-up-closed}.

Let $\mathfrak q$ lie over $\mathfrak p$, and use
the original injection
$\kappa(\mathfrak p)\to\kappa(\mathfrak q)$.
If $\alpha$ is the image of $x$ in the receiving
field, then $\alpha^2=\overline a_x$ and
$\alpha^3=\overline b_x$, where the bars are the
images of these original coefficients in
$\kappa(\mathfrak p)$. If $\overline a_x=0$,
then $\alpha=0$. Otherwise
$\alpha=\overline b_x/\overline a_x$.
Consequently every generator belongs to the
image of the original base residue field.
Since $\kappa(\mathfrak q)=\operatorname{Frac}(S/\mathfrak q)$
is generated as a field by those images and
$\kappa(\mathfrak p)$, this residue-field map
is an isomorphism.

If $\mathfrak q'$ also lies over $\mathfrak p$,
compose the maps $S\to\kappa(\mathfrak q)$ and
$S\to\kappa(\mathfrak q')$ with the respective
inverses of the base residue-field isomorphisms.
Both resulting maps to $\kappa(\mathfrak p)$
agree on $R$ and send each original $x$ to
the same displayed value, zero or
$\overline b_x/\overline a_x$.
They agree on every polynomial in the
generators, hence on $S$. Their kernels
are respectively $\mathfrak q$ and
$\mathfrak q'$, so these primes are equal.
The original spectrum map is a continuous
closed bijection, hence a homeomorphism.

Retain any original $f:R\to R'$ and put
$S'=R'\otimes_R S$. It is generated over
$R'$ by $1\otimes x$, with
$$
(1\otimes x)^2=f(a_x)\otimes1,
\qquad (1\otimes x)^3=f(b_x)\otimes1.
$$
To verify the new kernel without assuming
$f$ flat, let $\mathfrak p'\subset R'$
contract to $\mathfrak p\subset R$.
Choose an original $\mathfrak q$ above
$\mathfrak p$, as is possible by surjectivity.
The tensor ring
$$
C=\kappa(\mathfrak p')\otimes_{\kappa(\mathfrak p)}
\kappa(\mathfrak q)
$$
is nonzero: a vector-space basis of either
nonzero field remains a basis after extending
scalars. The original maps give
$S'\to C$, $r'\otimes s\mapsto\overline r'\otimes\overline s$.
Take a prime of $C$. Its quotient is a
nonzero domain, and the induced unital
map from the field $\kappa(\mathfrak p')$
is injective. The inverse image of that
prime in $S'$ therefore contracts to
exactly $\mathfrak p'$ in $R'$.
Thus $\Spec(S')\to\Spec(R')$ is surjective,
and its kernel lies in every prime of
$R'$, hence consists of nilpotents.
This proves both base-change hypotheses;
the already proved argument gives the
same homeomorphism and residue-field
isomorphisms after every such base change.

The exponents $2,3$ can more generally
be replaced, separately for each original
generator $x$, by a finite nonempty list
of positive integers $d_1,\ldots,d_r$
with greatest common divisor $1$, provided
all $x^{d_i}$ belong to $\varphi(R)$.
One positive exponent supplies the monic
integral equation. Choose integers $u_i$
with $\sum_i u_i d_i=1$. For a nonzero
residue $\alpha$ of $x$, every $\alpha^{d_i}$
is a nonzero element of the base residue
field, and the full identity
$$
\alpha=\prod_{i=1}^r(\alpha^{d_i})^{u_i}
$$
puts it in that field, with negative
powers interpreted through its actual
field inverses. If any power is zero,
then $\alpha=0$. This formula determines
the same residue at every prime over
the given base prime. The preceding
integrality, kernel, and generator-map
arguments therefore prove the same
conclusions and their preservation by
arbitrary base change for these lists.

The greatest-common-divisor condition
cannot be weakened to an arbitrary
positive common divisor. For a list
with greatest common divisor $d>1$,
take the original field inclusion
$\mathbf Q(u^d)\to\mathbf Q(u)$, with
$u$ transcendental. The powers in the
list all belong to the smaller field.
The elements $1,u,\ldots,u^{d-1}$ are
independent over $\mathbf Q(u^d)$:
clearing the original denominators
leaves polynomial monomials with
distinct exponent residues modulo $d$.
They span the algebra generated by $u$
using $u^d\in\mathbf Q(u^d)$; that
finite-dimensional domain is a field,
since multiplication by a nonzero
element is injective and hence
surjective on its finite-dimensional
vector space. Thus it is the original
$\mathbf Q(u)$ and the residue-field
extension has degree exactly $d$,
not degree one.

Even under the original $2,3$ hypotheses,
an injective map can acquire a nonzero
nilpotent kernel after base change.
For any field $k$ let
$$
R=k[a,b]/(b^2-a^3)\longrightarrow S=k[t],
\qquad a\longmapsto t^2,\quad b\longmapsto t^3.
$$
This is injective: monic division in
$b$ writes an element uniquely as
$f(a)+b g(a)$, and its image
$f(t^2)+t^3g(t^2)$ has disjoint even
and odd exponent supports. Also $S$
is generated by the original $t$ with
$t^2,t^3$ in the image. Base change
by $R\to R/(a)=k[b]/(b^2)$ gives
$$
k[b]/(b^2)\longrightarrow k[t]/(t^2),
\qquad b\longmapsto t^3=0.
$$
Its kernel is exactly the nonzero
ideal $(b)$, whose square is zero.
The quotient-tensor comparison
$(R/(a))\otimes_R S\simeq S/(t^2)$
sends $[r]\otimes s$ to $[\varphi(r)s]$
and has inverse $[s]\mapsto1\otimes s$.
Thus this is the actual base-changed
map, with its nilpotent kernel fully
identified.
\end{proof}


\subsection{Editorial treatment of Algebra lines 10813--10845}

\hypertarget{algebra-context-129}{}

\noindent\href{https://github.com/stacks/stacks-project/blob/a04446e57ec1fbc252a871afcec7752fb2807b14/algebra.tex\#L10813}{Original source, lines 10813--10845}. Complete original and reviewed TeX passages accompany this note. Every intervention is separately recorded; the following treatment is editorial.


\begin{lemma}
\label{editorial-source-lemma-help-with-powers}
Let $p$ be a prime number. Let $n, m > 0$ be two integers. There exists
an integer $a$ such that
$(x + y)^{p^a}, p^a(x + y) \in \mathbf{Z}[x^{p^n}, p^nx, y^{p^m}, p^my]$.
\end{lemma}

\begin{proof}
Retain the original subring
$$
A=\mathbf Z[x^{p^n},p^n x,y^{p^m},p^m y]
\subseteq\mathbf Z[x,y].
$$
We calculate exactly which integer
coefficients can occur at each original
monomial $x^i y^j$. Write uniquely
$$
i=q p^n+r,\quad 0\leq r<p^n,
\qquad j=q' p^m+r',\quad 0\leq r'<p^m.
$$
In any product of the four displayed
generators contributing to this
monomial, the exponent on $p^n x$
is at least $r$, and the exponent
on $p^m y$ is at least $r'$.
Thus its coefficient is divisible
by $p^{nr+mr'}$. Conversely the
product
$$
(x^{p^n})^q(p^n x)^r
(y^{p^m})^{q'}(p^m y)^{r'}
=p^{nr+mr'}x^i y^j
$$
belongs to $A$. Since distinct formal
monomials are independent over
$\mathbf Z$, a polynomial belongs
to $A$ exactly when each of its
coefficients is divisible by this
corresponding integer. This proves
necessity as well as sufficiency
of the source divisibility test.

For a nonnegative integer $a$,
$p^a(x+y)$ belongs to $A$ exactly
when $a\geq b=\max(n,m)$: the
coefficient tests for $x$ and $y$
are respectively divisibility by
$p^n$ and $p^m$. In that case its
full expression in the generators is
$$
p^a(x+y)=p^{a-n}(p^n x)+p^{a-m}(p^m y).
$$
Now keep $a\geq b$ and $N=p^a$.
For $0<i<N$, the original binomial
coefficient has the exact valuation
$$
v_p\binom Ni=a-v_p(i).
$$
Indeed
$\binom Ni=(N/i)\binom{N-1}{i-1}$,
and the second factor is not divisible
by $p$. To verify the latter assertion,
reduce the full polynomial identity
$(1+T)^N=1+T^N$ modulo $p$ and divide
by $1+T$ in $\mathbf F_p[T]$.
It gives
$$
(1+T)^{N-1}=\sum_{h=0}^{N-1}(-T)^h
\quad\hbox{in }\mathbf F_p[T].
$$
Multiplication by $1+T$ verifies this
identity: the last sign is positive
for odd $p$, and signs coincide for
$p=2$. Every coefficient is nonzero,
giving the valuation assertion with
the original integer coefficient
still retained.

The two endpoint terms $i=0,N$
already belong to $A$. For an
intermediate $i$, put $t=i\bmod p^b$
with $0\leq t<p^b$. If $t=0$, both
remainders $r,r'$ are zero. If $t>0$,
then $v_p(i)=v_p(t)$ and
$$
r=t\bmod p^n,\qquad
r'=(-t)\bmod p^m,
$$
where both residues are chosen in
their original nonnegative ranges.
Every $1\leq t<p^b$ occurs among
$1\leq i<N$. Consequently the least
possible $a$ for both assertions
is exactly
$$
\max\left(b,
\max_{1\leq t<p^b}
\left(v_p(t)+n(t\bmod p^n)+m((-t)\bmod p^m)\right)\right).
$$
This is an exact test on the
original integer coefficients,
not merely a sufficient estimate.

It has the following explicit value:
$$
a_{min}=
\begin{cases}
n p^n+n-1,&n=m,\\
b(p^b-1)+c,&n\neq m,\quad
b=\max(n,m),\ c=\min(n,m).
\end{cases}
$$
For $n=m$, each nonzero residue $t$
has $r+r'=p^n$, and its valuation
is at most $n-1$, with equality at
$t=p^{n-1}$. This gives the first
value, which is at least $b$.
For $n<m$, write $t=q p^n+r$.
If $r>0$, its full weight is
$$
m p^m-m q p^n-(m-n)r+v_p(r).
$$
For every positive $r$ one has
$v_p(r)\leq r-1$: if $v=v_p(r)$,
then $r\geq p^v\geq v+1$, the
last inequality following by
induction on $v$. Since $m-n\geq1$,
this gives
$v_p(r)\leq(m-n)(r-1)$.
The weight is therefore at most
$m p^m-(m-n)$, and equality is
attained at the original $t=1$.
If $r=0$ and $t>0$, its weight is
$m(p^m-t)+v_p(t)$, at most
$m p^m-(m-1)t-1$. Since $t\geq1$,
this is at most $m p^m-m$, and
hence at most $m p^m-(m-n)$.
Thus the maximum is exactly
$m(p^m-1)+n$. For $m<n$, the
automorphism of $\mathbf Z[x,y]$
exchanging the original $x$ and
$y$ takes $A$ to the corresponding
subring with $n,m$ exchanged, and
fixes both $(x+y)^{p^a}$ and
$p^a(x+y)$. Applying this actual
inverse-to-itself map proves the
other unequal case. These values
are all at least $b$.

For every $a\geq a_{min}$, the full
desired expression is therefore
$$
(x+y)^{p^a}=
\sum_{\substack{i,j\geq0\\i+j=p^a}}
\frac{\binom{p^a}{i,j}}{p^{nr+mr'}}
(x^{p^n})^q(p^n x)^r
(y^{p^m})^{q'}(p^m y)^{r'}.
$$
Each displayed quotient is an
integer by the proved divisibility;
no power of $p$ is inverted in the
ring. All original binomial
coefficients and all four generator
factors are retained. If $b\leq a<a_{min}$,
the maximizing residue just described
gives an actual monomial whose
coefficient fails the necessary
divisibility. If $a<b$, the linear
expression already fails. This
proves minimality.

In particular the source choice
$a_0=n p^n+m p^m+n+m$ works.
For $n=m$, its difference from
$a_{min}$ is $n p^n+n+1$.
For $n\neq m$, with the same
$b,c$, the difference is $c p^c+2b$.
Both are positive. Thus the original
sufficient choice is preserved and
verified, alongside the sharp bound.
\end{proof}


\subsection{Editorial treatment of Algebra lines 10847--10865}

\hypertarget{algebra-context-130}{}

\noindent\href{https://github.com/stacks/stacks-project/blob/a04446e57ec1fbc252a871afcec7752fb2807b14/algebra.tex\#L10847}{Original source, lines 10847--10865}. Complete original and reviewed TeX passages accompany this note. Every intervention is separately recorded; the following treatment is editorial.


\begin{lemma}
\label{editorial-source-lemma-p-ring-map-field}
Let $k'/k$ be a field extension. Let $p$ be a prime number.
The following are equivalent
\begin{enumerate}
\item $k'$ is generated as a field extension of $k$ by elements
$x$ such that there exists an $n > 0$ with $x^{p^n} \in k$ and
$p^nx \in k$, and
\item $k = k'$ or the characteristic of $k$
and $k'$ is $p$ and $k'/k$ is purely inseparable.
\end{enumerate}
\end{lemma}

\begin{proof}
The original field embedding preserves $1$, hence
the characteristic of $k$ and $k'$ is the
same. If that characteristic differs from
the specified prime $p$, the image of
$p^n$ is a nonzero field element and
therefore a unit of $k$. Thus
$x=(p^n x)/p^n\in k$ for every generator
in (1). Those generators generate $k'/k$,
so $k'=k$.

If the common characteristic is $p$,
each such generator has a $p$-power
in the original $k$. Every element
of the generated field is a rational
expression in finitely many original
generators. Choose an exponent at
least all their witnessing exponents.
Applying Frobenius to that full
rational expression raises its
coefficients and each generator to
elements of $k$; its denominator
remains nonzero in the field.
Thus every element of $k'$ has a
$p$-power in $k$, proving pure
inseparability without a finiteness
assumption on the whole extension.

Conversely, if $k'=k$, each element
itself satisfies (1) with $n=1$.
If the extension is purely
inseparable in characteristic $p$,
each original element has a
$p^e$th power in $k$. Choose
$n=\max(1,e)$; then its $p^n$th
power is still in $k$ and
$p^n x=0\in k$. The set of all
original elements can serve as
the generating set, proving (1).
The proof also shows exactly
which condition is needed for
fields: outside characteristic
$p$ the integer-multiple condition
alone forces the identity extension;
in characteristic $p$ it is automatic,
and the power condition is sufficient.
\end{proof}


\subsection{Editorial treatment of Algebra lines 10867--10910}

\hypertarget{algebra-context-131}{}

\noindent\href{https://github.com/stacks/stacks-project/blob/a04446e57ec1fbc252a871afcec7752fb2807b14/algebra.tex\#L10867}{Original source, lines 10867--10910}. Complete original and reviewed TeX passages accompany this note. Every intervention is separately recorded; the following treatment is editorial.


\begin{lemma}
\label{editorial-source-lemma-p-ring-map}
Let $\varphi : R \to S$ be a ring map. Let $p$ be a prime number. Assume
\begin{enumerate}
\item[(a)] $S$ is generated as an $R$-algebra by elements $x$ such
that there exists an $n > 0$ with $x^{p^n} \in \varphi(R)$ and
$p^nx \in \varphi(R)$, and
\item[(b)] $\Ker(\varphi)$ is locally nilpotent,
\end{enumerate}
Then $\varphi$ induces a homeomorphism of spectra and induces
residue field extensions satisfying the equivalent conditions
of Lemma \ref{editorial-source-lemma-p-ring-map-field}. For any ring map $R \to R'$
the ring map $R' \to R' \otimes_R S$ also satisfies (a) and (b).
\end{lemma}

\begin{proof}
Retain the actual image $A=\varphi(R)\subseteq S$.
Let $T$ consist of the original $x\in S$ for
which some positive $n$ satisfies
$x^{p^n}\in A$ and $p^n x\in A$.
Every element of $A$ lies in $T$, by taking
$n=1$. If $x\in T$, then $-x\in T$ with
the same $n$, since both required expressions
change only by the respective integer signs.
For $r\in R$,
$$
(\varphi(r)x)^{p^n}=\varphi(r)^{p^n}x^{p^n}\in A,
\qquad
p^n\varphi(r)x=\varphi(r)(p^n x)\in A.
$$
Thus this verification never assumes that
the original $\varphi$ is injective.
If $x,y\in T$ have the original witnesses
$n,m$, then
$$
(xy)^{p^{n+m}}
=(x^{p^n})^{p^m}(y^{p^m})^{p^n}\in A,
\qquad
p^{n+m}xy=(p^n x)(p^m y)\in A.
$$
For addition, apply the actual evaluation
homomorphism $\mathbf Z[X,Y]\to S$,
$X\mapsto x$, $Y\mapsto y$, to Lemma
\ref{editorial-source-lemma-help-with-powers}. Its four
original generators map to
$x^{p^n},p^n x,y^{p^m},p^m y\in A$.
The displayed full coefficient formulas
in that lemma give both required expressions
for $x+y$ in $A$, with its specified
integer exponent. Therefore $T$ is an
$R$-subalgebra of $S$. Hypothesis (a)
places every generating element in $T$,
so $T=S$.

In particular every original element
of $S$ has a positive power in $A$.
Together with the nil-kernel hypothesis,
Lemma \ref{editorial-source-lemma-powers} proves that
the original map on spectra is a
homeomorphism. Let $\mathfrak q$ lie
over $\mathfrak p$. The field
$\kappa(\mathfrak q)$ is generated over
$\kappa(\mathfrak p)$ by the residues
of the original algebra generators:
the images generate $S/\mathfrak q$
as an $R/\mathfrak p$-algebra, and
fractions give its residue field.
Each such residue still satisfies
both original power and integer-multiple
conditions. Lemma \ref{editorial-source-lemma-p-ring-map-field}
therefore proves precisely the asserted
residue-field classification.

For any original $R\to R'$, the
algebra $S'=R'\otimes_R S$ has the
generators $1\otimes x$. If the
original preimages of $x^{p^n}$ and
$p^n x$ are $a,b\in R$, then
$$
(1\otimes x)^{p^n}=f(a)\otimes1,
\qquad p^n(1\otimes x)=f(b)\otimes1.
$$
Thus (a) is preserved with exactly
the same $n$ for each generator.
Surjectivity on spectra is preserved
as well: for $\mathfrak p'\subset R'$
contracting to $\mathfrak p$, choose
$\mathfrak q\subset S$ above
$\mathfrak p$ and use the nonzero
field tensor
$\kappa(\mathfrak p')\otimes_{\kappa(\mathfrak p)}\kappa(\mathfrak q)$.
The map $r'\otimes s\mapsto
\overline r'\otimes\overline s$
and any prime of that tensor have
inverse image in $S'$ contracting
to exactly $\mathfrak p'$, since
a field injects into every nonzero
unital algebra over it. This is
the same explicit fibre construction
as in Lemma \ref{editorial-source-lemma-2-3-ring-map}.
Every element of the new kernel
therefore lies in every prime of
$R'$ and is nilpotent, proving (b).
The theorem applies to this actual
base-changed map, giving a
homeomorphism and the stated
residue-field extensions after
every base change.

The integer-multiple hypothesis
also gives an exact global conclusion
after inverting the original prime
integer $p$. The map
$$
R[1/p]\longrightarrow S[1/p]
$$
is surjective: each original
generator with $p^n x=\varphi(b)$
is the image of the actual fraction
$b/p^n$. Its kernel is exactly
$I[1/p]$, for $I=\ker\varphi$.
Indeed if $r/p^j$ maps to zero,
there is an integer $h\geq0$
with $p^h\varphi(r)=0$, and then
$p^h r\in I$ and
$r/p^j=(p^h r)/p^{j+h}$.
The reverse containment is immediate.
Thus
$S[1/p]\simeq R[1/p]/I[1/p]$
through this original fraction map,
and the kernel remains locally
nilpotent. In particular the
localized map is an isomorphism
if the original $\varphi$ is
injective. This statement also
covers the case that inverting
$p$ makes both rings zero.
\end{proof}


\subsection{Editorial treatment of Algebra lines 10912--10966}

\hypertarget{algebra-context-132}{}

\noindent\href{https://github.com/stacks/stacks-project/blob/a04446e57ec1fbc252a871afcec7752fb2807b14/algebra.tex\#L10912}{Original source, lines 10912--10966}. Complete original and reviewed TeX passages accompany this note. Every intervention is separately recorded; the following treatment is editorial.


\begin{lemma}
\label{editorial-source-lemma-radicial}
Let $\varphi : R \to S$ be a ring map. Assume
\begin{enumerate}
\item $\varphi$ induces an injective map of spectra,
\item $\varphi$ induces purely inseparable residue field extensions.
\end{enumerate}
Then for any ring map $R \to R'$ properties (1) and (2) are true for
$R' \to R' \otimes_R S$.
\end{lemma}

\begin{proof}
Retain the original base-change algebra
$S'=R'\otimes_R S$ and its two structural
maps. Let $\mathfrak p'\subset R'$ contract
to $\mathfrak p\subset R$, and set
$K=\kappa(\mathfrak p)$ and
$K'=\kappa(\mathfrak p')$ with the original
field embedding $K\to K'$. The original
fibre algebra is $F=K\otimes_R S$.
By Remark \ref{algebra-remark-fundamental-diagram},
its primes correspond to the primes
of $S$ over $\mathfrak p$, with their
original residue-field maps.

If there is no such prime, $F$ has
no prime and is zero, since every
nonzero unital ring has a maximal
ideal. The new fibre is zero as
well by the explicit comparison
below. Otherwise there is exactly
one original prime $\mathfrak q$
above $\mathfrak p$. Let $\mathfrak r$
be its corresponding prime in $F$.
It is the only prime of $F$, hence
is both its nilradical and its
unique maximal ideal. Consequently
$F/\mathfrak r$ is a field, and
the canonical map $F\to\kappa(\mathfrak r)$
is surjective with kernel $\mathfrak r$.
The original fibre correspondence
identifies this field with
$E=\kappa(\mathfrak q)$ over $K$.
By hypothesis $E/K$ is purely
inseparable. This proves the
needed nil-kernel quotient without
assuming that $F$ itself is reduced.

There is an original fibre isomorphism
$$
K'\otimes_{R'} S'\longrightarrow
K'\otimes_K F,\qquad
c\otimes(r'\otimes s)\longmapsto
c\overline r'\otimes(1\otimes s).
$$
Its inverse sends
$c\otimes(a\otimes s)$ to
$ca\otimes(1\otimes s)$, using the
specified embedding $K\to K'$.
Balancing follows from the original
commutative scalar diagram, and
these formulas compose to the
identity on every pure tensor.
Write $H=K'\otimes_K F$.
Lemma \ref{editorial-source-lemma-surjective-locally-nilpotent-kernel}
applied to $F\to E$ gives a
surjection
$H\to K'\otimes_K E$ with locally
nilpotent kernel.

Embed the purely inseparable $E$
in an algebraic closure of $K'$
by the unique roots compatible
with $K$. The full tensor comparison
proved in Lemma \ref{editorial-source-lemma-make-separable}
gives
$$
(K'\otimes_K E)_{red}\simeq C=K'E,
\qquad [\sum_i c_i\otimes e_i]
\longmapsto\sum_i c_i e_i.
$$
It applies also when $E/K$ is
infinite; each individual tensor
uses only finitely many root
exponents. The field $C/K'$ is
purely inseparable, because every
finite sum or quotient in the
compositum has a sufficiently high
$p$-power in $K'$ in positive
characteristic; in characteristic
zero $E=K$ and $C=K'$.
Both kernels in the composite
$H\to K'\otimes_K E\to C$ are nil.
Its full kernel is nil as well:
if an element maps to zero, some
power lies in the first nil kernel,
and a further power is zero.
Thus $H_{red}\simeq C$. The fibre
$H$ has exactly one prime and its
residue field is the original $C$.
The fibre correspondence proves
both required base-change properties.

In particular we obtain the exact
prime and residue field, rather
than only their uniqueness. Whenever
$\mathfrak q$ exists, the unique
$\mathfrak q'\subset S'$ over
$\mathfrak p'$ is the kernel of
$$
S'\longrightarrow C,\qquad
r'\otimes s\longmapsto
\overline r'\,\overline s,
$$
with the two factors embedded in
$K'$ and $E$ as above. Its contraction
to $S$ is $\mathfrak q$ and its
contraction to $R'$ is $\mathfrak p'$.
The fraction field of its image
is exactly $C$: it contains the
fractions giving $K'$ and $E$,
and is contained in their compositum.
This identifies $\kappa(\mathfrak q')$
with $C$ through the original maps.
It also proves that the image on
spectra after base change is exactly
the inverse image of the original
image on spectra.

Conversely, these two original
hypotheses are necessary for
injectivity on spectra after every
base change. Injectivity before
base change follows by choosing
$R'=R$. For an original
$\mathfrak q$ over $\mathfrak p$,
use the same $K,E,F$ above; injectivity
makes $F\to E$ a surjective nil-kernel
map. If $E/K$ were not purely
inseparable, the ring $E\otimes_K E$
would have at least two primes, as
we now prove with the original maps.
Its multiplication map onto $E$
has a prime kernel containing every
$z\otimes1-1\otimes z$.
If some $z\in E$ is transcendental
over $K$, no positive power of
this difference is zero: in its
full binomial expansion, the
coefficient of $1\otimes z^n$
is $(-1)^n$, and the powers
$1,z,\ldots,z^n$ remain independent
after extension from $K$ to $E$.

Otherwise $E/K$ is algebraic but
not purely inseparable. The
separable-first theorem supplies
$z\in E\setminus K$ separable over
$K$. Let $P(T)$ be its original
monic minimal polynomial, of degree
at least two. The ring
$E\otimes_K K(z)\simeq E[T]/(P)$
is reduced because $P$ is separable,
by the original polynomial and
Bezout argument in Lemma
\ref{editorial-source-lemma-separable-extension-preserves-reducedness}.
The element corresponding to $T-z$
is nonzero, by monic division,
and so is not nilpotent. The map
$E\otimes_K K(z)\to E\otimes_K E$
is injective by extension of a
vector-space basis. Thus the same
diagonal difference is not nilpotent
in the larger tensor ring either.
In both cases some prime avoids
that difference, while the original
multiplication kernel contains it.
They are two distinct primes.

Now take the actual base change
$R\to E$ through $R\to K\to E$.
Its algebra is
$E\otimes_R S\simeq E\otimes_K F$.
The surjection from this algebra
to $E\otimes_K E$ pulls the two
distinct primes back to distinct
primes, both over the unique prime
of the field $E$. This contradicts
universal injectivity. Hence every
original residue-field extension
must be purely inseparable, proving
the converse and the exact
characterization.
\end{proof}


\subsection{Editorial treatment of Algebra lines 10968--10986}

\hypertarget{algebra-context-133}{}

\noindent\href{https://github.com/stacks/stacks-project/blob/a04446e57ec1fbc252a871afcec7752fb2807b14/algebra.tex\#L10968}{Original source, lines 10968--10986}. Complete original and reviewed TeX passages accompany this note. Every intervention is separately recorded; the following treatment is editorial.


\begin{lemma}
\label{editorial-source-lemma-radicial-integral}
Let $\varphi : R \to S$ be a ring map. Assume
\begin{enumerate}
\item $\varphi$ is integral,
\item $\varphi$ induces an injective map of spectra,
\item $\varphi$ induces purely inseparable residue field extensions.
\end{enumerate}
Then $\varphi$ induces a homeomorphism from $\Spec(S)$ onto a closed
subset of $\Spec(R)$ and for any ring map
$R \to R'$ properties (1), (2), (3) are true for $R' \to R' \otimes_R S$.
\end{lemma}

\begin{proof}
Put $I=\ker\varphi$. The original integral
injection $R/I\to S$ gives image $V(I)$
on spectra by lying over. More generally,
for every ideal $J\subseteq S$, the
map $R/\varphi^{-1}(J)\to S/J$ is an
integral injection, since its monic
equations are the images of the
original ones. Its spectrum is
surjective, so the image of $V(J)$
is exactly $V(\varphi^{-1}(J))$.
Thus the original spectrum map is
closed. Its assumed injectivity
makes it a homeomorphism onto
the closed subset $V(I)$: closed
subsets of the domain have closed
images in this subspace, which
proves continuity of the inverse.

For any original $R\to R'$, write
$S'=R'\otimes_R S$. Each original
monic equation for $s\in S$ maps
to the same monic equation for
$1\otimes s$ with coefficients
mapped into $R'$. Those elements
generate $S'$ over $R'$, so
integrality is preserved, including
infinite generating families by
the finite-support subring argument.
Lemma \ref{editorial-source-lemma-radicial} supplies
injectivity on spectra and pure
inseparability of residue fields
for this exact base-changed map.
The preceding closed-image proof
therefore applies to it as well.

There is also an exact radical
comparison for its kernel $I'$:
$$
\sqrt{I'}=\sqrt{IR'}.
$$
Indeed the original image is
$V(I)$, and the explicit fibre
construction in Lemma
\ref{editorial-source-lemma-radicial} identifies the
new image with its inverse image
in $\Spec(R')$, namely $V(IR')$.
Integrality also identifies the
new image with $V(I')$. Equality
of these closed sets gives the
stated equality of radicals.
The same formula in fact holds
for every integral ring map:
for a prime of $R'$ containing
$IR'$, its contraction contains
$I$, so lying over supplies a
prime of $S$ above that contraction;
the nonzero tensor of the two
original residue fields supplies
a prime after base change. This
proves equality of the images
without requiring uniqueness.
Equality of the kernels themselves
need not hold: the explicit cusp
base change in Lemma
\ref{editorial-source-lemma-2-3-ring-map} has $I=0$
and $I'=(b)\neq0$ in $k[b]/(b^2)$,
while both radicals are $(b)$.
All comparisons concern the
original extended ideal $IR'$.
\end{proof}


\subsection{Editorial treatment of Algebra lines 10988--11003}

\hypertarget{algebra-context-134}{}

\noindent\href{https://github.com/stacks/stacks-project/blob/a04446e57ec1fbc252a871afcec7752fb2807b14/algebra.tex\#L10988}{Original source, lines 10988--11003}. Complete original and reviewed TeX passages accompany this note. Every intervention is separately recorded; the following treatment is editorial.


\begin{lemma}
\label{editorial-source-lemma-radicial-integral-bijective}
Let $\varphi : R \to S$ be a ring map. Assume
\begin{enumerate}
\item $\varphi$ is integral,
\item $\varphi$ induces a bijective map of spectra,
\item $\varphi$ induces purely inseparable residue field extensions.
\end{enumerate}
Then $\varphi$ induces a homeomorphism on spectra and for any ring map
$R \to R'$ properties (1), (2), (3) are true for $R' \to R' \otimes_R S$.
\end{lemma}

\begin{proof}
By Lemma \ref{editorial-source-lemma-radicial-integral}, the
original integral injective map on spectra
is a homeomorphism onto its image.
Its assumed bijectivity makes that image
the entire $\Spec(R)$. For any original
$R\to R'$, the same lemma preserves
integrality, injectivity on spectra,
and purely inseparable residue fields.
Surjectivity on spectra is preserved
by the original nonzero residue-field
tensor construction of Lemma
\ref{editorial-source-lemma-surjective-spec-radical-ideal}.
Thus the receiving map is bijective
on spectra and still integral, hence
a homeomorphism by the closed-image
argument, with exactly the asserted
residue fields. Equivalently, for
an integral map already satisfying
the injectivity and residue-field
hypotheses, full universal surjectivity
is the condition $\ker\varphi$
locally nilpotent: its image is
the exact $V(\ker\varphi)$, and
the same holds after every base
change by the radical formula.
\end{proof}


\subsection{Editorial treatment of Algebra lines 11005--11061}

\hypertarget{algebra-context-135}{}

\noindent\href{https://github.com/stacks/stacks-project/blob/a04446e57ec1fbc252a871afcec7752fb2807b14/algebra.tex\#L11005}{Original source, lines 11005--11061}. Complete original and reviewed TeX passages accompany this note. Every intervention is separately recorded; the following treatment is editorial.


\begin{lemma}
\label{editorial-source-lemma-universally-bijective}
Let $\varphi : R \to S$ be a ring map such that
\begin{enumerate}
\item the kernel of $\varphi$ is locally nilpotent, and
\item $S$ is generated as an $R$-algebra by elements $x$
such that there exist $n > 0$ and a polynomial $P(T) \in R[T]$
whose image in $S[T]$ is $(T - x)^n$.
\end{enumerate}
Then $\Spec(S) \to \Spec(R)$ is a homeomorphism and $R \to S$
induces purely inseparable extensions of residue fields.
Moreover, conditions (1) and (2) remain true on arbitrary base change.
\end{lemma}

\begin{proof}
If $S=0$, the nil-kernel hypothesis
makes $1\in R$ nilpotent, so $R=0$.
Every unital base change of the zero
ring is again zero, and all assertions
follow. Hence assume $S\neq0$.
Retain the original map $\varphi:R\to S$
and its kernel $I$. For every original
generator $x$ retain the given positive
integer $n$ and the entire original
polynomial $P(T)=\sum_{j=0}^d c_jT^j$,
whose image is $(T-x)^n$. Coefficient
comparison gives $d\geq n$,
$c_n-1\in I$, and $c_j\in I$ for
$j>n$. The polynomial
$Q(T)=T^n+\sum_{j=0}^{n-1}c_jT^j$
is a monic annihilating polynomial
for $x$. Indeed the full difference
$$
P(T)-Q(T)=(c_n-1)T^n+
\sum_{j=n+1}^d c_jT^j
$$
has all coefficients in the original
kernel, so $Q$ has the same image
$(T-x)^n$. This retains every
coefficient of $P$ and proves
integrality without presuming that
the given lift itself is monic.
All generators, hence $S$, are
integral over $R$. The integral
injection $R/I\to S$ is surjective
on spectra. Since $I$ is nil,
$V(I)=\Spec(R)$, so the original
spectrum map is surjective; it is
also closed by integrality.

We first prove the global scalar
identity used in the fibre argument.
Write
$$
e_j=\binom nj x^j
=\varphi((-1)^j c_{n-j})\in\varphi(R),
\qquad 1\leq j\leq n.
$$
All signs here are the actual signs
in the given $(T-x)^n$. For
$1\leq h\leq n$, the full Newton
recurrence is
$$
n x^h=
\sum_{j=1}^{h-1}(-1)^{j-1}e_j\,n x^{h-j}
+(-1)^{h-1}h e_h.
$$
It follows directly on substituting
the original $e_j$: Pascal's identity
telescopes to
$\sum_{j=0}^{h-1}(-1)^j\binom nj
=(-1)^{h-1}\binom{n-1}{h-1}$,
and $n\binom{n-1}{h-1}=h\binom nh$.
These two integer identities give
exactly the coefficient $n$ on the
right, with every displayed summand
retained. Induction on $h$, starting
with $n x=e_1$, proves $n x^h\in\varphi(R)$.
Explicit preimages are supplied by
the same recurrence in $R$, taking
$E_j=(-1)^j c_{n-j}$ and
$$
U_h=\sum_{j=1}^{h-1}(-1)^{j-1}E_j U_{h-j}
+(-1)^{h-1}h E_h.
$$
The empty sum for $h=1$ gives
$U_1=E_1$, and $\varphi(U_h)=n x^h$.

Fix a prime integer $p$ and write
the original $n=p^e m$, with
$p\nmid m$, retaining $q=p^e$.
One has the exact integer identity
$$
\gcd\left(n,\binom nq\right)=m.
$$
For a prime $\ell\neq p$ dividing
$n$, the identity
$\binom nq=(n/q)\binom{n-1}{q-1}$
shows its $\ell$-adic valuation is
at least that of $n$, since
$q$ is prime to $\ell$.
For $p$, the factor $n/q=m$
is prime to $p$, and
$\binom{n-1}{q-1}$ is nonzero
modulo $p$. In fact
$(1+T)^{qm}=(1+T^q)^m$ modulo
$p$, so dividing by $1+T$ shows
that its degree-$q-1$ coefficient
is $(-1)^{q-1}=1$ in
$\mathbf F_p$. This also treats
$q=1$. These valuations prove
the displayed greatest common
divisor, including every prime
factor of the original $n$.
Choose integers $A,B$ with
$An+B\binom nq=m$. Then
$$
\varphi(AU_q+BE_q)
=A n x^q+B\binom nq x^q=m x^q.
$$
Thus the source assertion
$m x^{p^e}\in\varphi(R)$ is
indeed true globally, not merely
after taking a residue field.
The recurrence and this integer
identity supply its missing proof.

Now let $\mathfrak p\subset R$
and retain $K=\kappa(\mathfrak p)$
and the original fibre
$C=K\otimes_R S$. It is nonzero
by surjectivity on spectra. Thus
the map from the field $K$ to $C$
is injective. If $K$ has
characteristic zero, each generator
$y=1\otimes x$ satisfies
$n y=-\overline c_{n-1}$.
The original positive $n$ is
invertible in $K$, so
$y=-n^{-1}\overline c_{n-1}$
belongs to the image of $K$.
All generators do, and hence
$K\to C$ is an isomorphism.

If $K$ has characteristic $p>0$,
use the original decomposition
$n=q m$, $q=p^e$, for each
generator separately. The full
fibre polynomial is
$$
(T-y)^n=(T^q-y^q)^m
=\sum_{j=0}^m\binom mj(-1)^j
T^{q(m-j)}y^{qj}.
$$
Its coefficient of $T^{n-q}$
is $-m y^q=\overline c_{n-q}$.
Since $m$ is a unit of $K$,
$y^q=-m^{-1}\overline c_{n-q}$
is in the original image of $K$.
This is the direct fibre formula;
the global preimage proved above
gives the same conclusion before
reducing coefficients.
If $e=0$, use the exponent $1$
in Lemma \ref{editorial-source-lemma-p-ring-map};
otherwise use $e$. In either
case a positive $p$-power of $y$
is in $K$, and the corresponding
integer multiple is zero in
characteristic $p$. Thus that
lemma applies to $K\to C$, whose
kernel is zero. The fibre has
one prime and its residue field
is purely inseparable over $K$.
The same conclusion holds in
characteristic zero by the
already proved isomorphism.
Remark \ref{algebra-remark-fundamental-diagram}
now gives bijectivity on spectra
and the asserted original residue
extensions. A continuous closed
bijection is a homeomorphism.

Under any original $f:R\to R'$,
the image polynomial
$P_f(T)=\sum_j f(c_j)T^j$ has
image $(T-1\otimes x)^n$ in
$(R'\otimes_R S)[T]$, with all
coefficients and the same $n$
retained. These elements generate
the base-changed algebra, so
condition (2) persists. For
condition (1), the original
surjectivity on spectra persists
under base change by the explicit
nonzero residue-field tensor
construction in Lemma
\ref{editorial-source-lemma-2-3-ring-map}. Hence
the new kernel lies in every
prime and is locally nilpotent.
Both hypotheses, and therefore
all conclusions, hold for the
actual base-changed map.

The scalar identities give more
global information for a finite
original generating list
$x_1,\ldots,x_r$ with integers
$n_1,\ldots,n_r$. Put
$N=\prod_{i=1}^r n_i$, taking
$N=1$ for the empty list.
Since $n_i x_i$ has the
specified preimage $U_{i,1}$,
the map $R[1/N]\to S[1/N]$
is surjective, sending the
original fraction
$U_{i,1}/n_i=((N/n_i)U_{i,1})/N$
to $x_i$. Its kernel is exactly
$I[1/N]$: the fraction-witness
proof is the same as for
inverting $p$ in Lemma
\ref{editorial-source-lemma-p-ring-map}, replacing
each power of $p$ there by the
same power of this original $N$.
Thus it is an explicit nil-kernel
quotient, and is an isomorphism
if $\varphi$ was injective.

For a fixed prime $p$, retain
instead $n_i=p^{e_i}m_i$ and
$M=\prod_i m_i$. Over $R[1/M]$,
the global preimages above give
both
$$
x_i^{p^{e_i}}=
\frac{\varphi(A_iU_{i,p^{e_i}}+B_iE_{i,p^{e_i}})}{m_i},
\qquad
p^{e_i}x_i=\frac{\varphi(U_{i,1})}{m_i}.
$$
When $e_i=0$ these put $x_i$
itself in the image; one may
then choose exponent $1$.
For $e_i>0$ these are the
original two $p$-power conditions.
The kernel remains $I[1/M]$
and is nil. Thus the whole
localized map satisfies Lemma
\ref{editorial-source-lemma-p-ring-map}, with the
exact root exponents and all
prime-to-$p$ denominators retained.
\end{proof}


\subsection{Editorial treatment of Algebra lines 11082--11103}

\hypertarget{algebra-context-136}{}

\noindent\href{https://github.com/stacks/stacks-project/blob/a04446e57ec1fbc252a871afcec7752fb2807b14/algebra.tex\#L11082}{Original source, lines 11082--11103}. Complete original and reviewed TeX passages accompany this note. Every intervention is separately recorded; the following treatment is editorial.


\begin{lemma}
\label{editorial-source-lemma-flat-fibres-irreducible}
Let $R \to S$ be a ring map. Assume
\begin{enumerate}
\item[(a)] $\Spec(R)$ is irreducible,
\item[(b)] $R \to S$ is flat,
\item[(c)] $R \to S$ is of finite presentation,
\item[(d)] the fibre rings $S \otimes_R \kappa(\mathfrak p)$
have irreducible spectra for a dense collection of primes $\mathfrak p$ of $R$.
\end{enumerate}
Then $\Spec(S)$ is irreducible.
This is true more generally with (b) $+$ (c)
replaced by ``the map $\Spec(S) \to \Spec(R)$ is open''.
\end{lemma}

\begin{proof}
Write $\varphi:R\to S$ for the original
ring map. Put $X=\Spec(S)$, $Y=\Spec(R)$
and retain its spectrum map $f:X\to Y$. Proposition
\ref{editorial-source-proposition-fppf-open} makes $f$ open
under (b) and (c). We prove the stated
more general open-map assertion directly.
Let $D\subseteq Y$ be the given dense
set of points with irreducible fibres.
The irreducible space $Y$ is nonempty,
so $D$ is nonempty. Each fibre at a
point of $D$ is nonempty; consequently
$X$ is nonempty as well.

Let $U,V$ be nonempty open subsets of
$X$. The original images $f(U),f(V)$
are nonempty open subsets of the
irreducible $Y$, so their intersection
is nonempty and open. Density supplies
$y\in D\cap f(U)\cap f(V)$.
The subsets $U\cap f^{-1}(y)$ and
$V\cap f^{-1}(y)$ are nonempty open
subsets of the irreducible fibre.
Their intersection is nonempty,
so $U\cap V\neq\emptyset$.
This is the nonempty-space criterion
for irreducibility, proving the lemma
and supplying the nonemptiness step
in the cited Topology, Lemma
\ref{topology-lemma-irreducible-on-top}.

The same proof needs less than
openness: it is enough that the
image of every nonempty open
$U\subseteq X$ contain a nonempty
open subset of $Y$. Choose such
subsets inside $f(U)$ and $f(V)$
and perform the same intersection
argument with $D$. The whole proof
also works with the actual subspace
$f(X)\subseteq Y$ as codomain,
provided it is irreducible, the
specified interior condition is
relative to that subspace, and
the good fibres form a dense
subset of it. These are statements
about the original map with its
codomain restricted, not an
assumption that the image equals $Y$.

For flat ring maps there is a
different exact criterion without
finite presentation. Under (a),
let $\mathfrak p=\operatorname{Nil}(R)$
be the unique minimal prime of
the original $R$. Put
$C=S\otimes_R\kappa(\mathfrak p)$.
The original fibre comparison is
$$
C\simeq(R\setminus\mathfrak p)^{-1}(S/\mathfrak pS),
\qquad
s\otimes(\overline a/\overline b)
\longmapsto\varphi(a)s/\varphi(b),
\quad b\notin\mathfrak p.
$$
The inverse sends $\overline s/\varphi(b)$
to $s\otimes(1/\overline b)$;
quotient and localization balancing
prove well-definedness and both
composites. Primes of $C$ therefore
correspond exactly to original
primes of $S$ contracting to
$\mathfrak p$.

If $\mathfrak q$ is a minimal
prime of $S$, going down for
flat maps (Lemma \ref{editorial-source-lemma-flat-going-down}) lifts the inclusion
$\mathfrak p\subseteq\varphi^{-1}(\mathfrak q)$
to a prime contained in
$\mathfrak q$. Minimality forces
equality, so its contraction is
exactly $\mathfrak p$. Its fibre
prime is minimal, since a smaller
fibre prime would contract to
a smaller prime of $S$.
Conversely, if a minimal fibre
prime corresponds to $\mathfrak q$,
choose a minimal prime
$\mathfrak q_0\subseteq\mathfrak q$
of $S$. Its contraction is
$\mathfrak p$ by the previous
argument. Its fibre prime is
contained in the given minimal
one, so both are equal, and
$\mathfrak q_0=\mathfrak q$.
Thus these original fibre maps
give a bijection between the
minimal primes of $S$ and $C$.
If either ring is zero, both
sets are empty; if $S$ is nonzero,
it has a minimal prime and hence
$C$ is nonzero. Consequently,
under (a) and flatness alone,
$\Spec(S)$ is irreducible if and
only if this original generic
fibre has irreducible spectrum.

Dense irreducible fibres cannot
replace the generic-fibre condition
after dropping the openness
assumption. Retain the original
diagonal map
$$
\mathbf Z\longrightarrow\mathbf Q\times\mathbf Z,
\qquad n\longmapsto(n,n).
$$
It is flat: its underlying module
is the finite direct sum of the
flat localization $\mathbf Q$
and the free module $\mathbf Z$.
For every prime integer $\ell$,
its fibre is $\mathbf F_\ell$.
Indeed $\mathbf Q\otimes_{\mathbf Z}\mathbf F_\ell=0$,
since $1\otimes1=\ell^{-1}\otimes\ell=0$,
and the second factor gives
$\mathbf Z\otimes\mathbf F_\ell\simeq\mathbf F_\ell$.
These closed primes are dense:
every nonempty $D(n)$, $n\neq0$,
contains a prime dividing $|n|+1$
and hence not dividing $n$.
But the generic fibre is
$\mathbf Q\times\mathbf Q$,
through the multiplication map
$q\otimes(a,z)\mapsto(qa,qz)$
and its inverse
$(a,b)\mapsto a\otimes(1,0)+b\otimes(0,1)$.
The equality on the first factor
uses the original localization
identity $\mathbf Q\otimes_{\mathbf Z}\mathbf Q\simeq\mathbf Q$.
Both the total spectrum and its
generic fibre are reducible.
The nonempty open component
$\Spec(\mathbf Q)$ has image
$\{(0)\}\subseteq\Spec(\mathbf Z)$,
which has empty interior by the
same $D(n)$ argument. This
identifies exactly the failed
topological hypothesis.
\end{proof}


\subsection{Editorial treatment of Algebra lines 11105--11143}

\hypertarget{algebra-context-137}{}

\noindent\href{https://github.com/stacks/stacks-project/blob/a04446e57ec1fbc252a871afcec7752fb2807b14/algebra.tex\#L11105}{Original source, lines 11105--11143}. Complete original and reviewed TeX passages accompany this note. Every intervention is separately recorded; the following treatment is editorial.


\begin{lemma}
\label{editorial-source-lemma-separably-closed-irreducible}
Let $k$ be a separably closed field.
Let $R$, $S$ be $k$-algebras. If $R$, $S$ have a unique
minimal prime, so does $R \otimes_k S$.
\end{lemma}

\begin{proof}
Choose the original perfect closure
$\overline k/k$ of Definition
\ref{editorial-source-definition-perfection}. It is
algebraically closed in this situation.
Indeed take an algebraic closure
$\Omega/k$. Its maximal separable
subextension is $k$, because $k$
is separably closed. Fields, Lemma
\ref{fields-lemma-separable-first}
therefore makes $\Omega/k$ purely
inseparable. The algebraically
closed field $\Omega$ is perfect,
so uniqueness of perfect closure
gives the unique $k$-isomorphism
$\overline k\simeq\Omega$.
In characteristic zero both
extensions are already the identity.

For each original $k$-algebra $A$,
the map $A\to A\otimes_k\overline k$
is injective by extension of a
vector-space basis. In positive
characteristic the receiving algebra
is generated by $1\otimes\lambda$,
each with a $p$-power in the image
of $k$, and the corresponding
integer multiple is zero.
Lemma \ref{editorial-source-lemma-p-ring-map} gives
a homeomorphism on spectra; in
characteristic zero the map is
the identity under multiplication.
Apply this to $R$, $S$ and the
original $R\otimes_k S$.
The scalar comparison is the
isomorphism
$$
(R\otimes_k S)\otimes_k\overline k
\longrightarrow
(R\otimes_k\overline k)\otimes_{\overline k}
(S\otimes_k\overline k),
\qquad
(r\otimes s)\otimes c\longmapsto
(r\otimes c)\otimes(s\otimes1).
$$
Its inverse sends
$(r\otimes a)\otimes(s\otimes b)$
to $(r\otimes s)\otimes ab$.
The original scalar actions prove
balancing and both composites.
The two changed factors still
have a unique minimal prime.
It therefore suffices to prove
the assertion over the actual
algebraically closed field
$\overline k$, and then descend
through these homeomorphisms.

Work now over an algebraically
closed field $k$, keeping the
algebras then under consideration.
For arbitrary $R,S$, the tensor
quotient
$$
R\otimes_k S\longrightarrow R_{red}\otimes_k S_{red}
$$
is surjective with kernel
the sum of the two ideals
generated by the original factor
nilradicals. The inverse of its
induced quotient isomorphism
sends $[r]\otimes[s]$ to the
class of $r\otimes s$.
Each element of this kernel is
a finite sum of nilpotent pure
tensors; if their exponents are
$d_1,\ldots,d_m$, the full
multinomial expansion at exponent
$1+\sum_i(d_i-1)$ is zero.
Thus the kernel is locally
nilpotent and Lemma
\ref{editorial-source-lemma-surjective-locally-nilpotent-kernel}
identifies the spectra.
Since the original factors have
unique minimal primes, their
reductions are nonzero domains.
We may consequently prove the
remaining assertion for these
two domains, without losing the
nilpotent information in the
original quotient map.

For domains $R,S$ over this
algebraically closed $k$, their
tensor ring is nonzero, as is
seen by extending the vectors
$1\in R$ and $1\in S$ to bases.
It is reduced by Lemma
\ref{editorial-source-lemma-perfect-reduced}.
If it were not a domain, there
would be nonzero original tensors
$z,w$ with $zw=0$. Each uses
finitely many coefficients from
$R$ and $S$. Choose the subalgebras
$R_0,S_0$ generated by all those
coefficients. They are domains
of finite type over the
algebraically closed field $k$.
The map $R_0\otimes_k S_0\to R\otimes_k S$
is injective by applying the
two vector-space injections
successively. Hence the same
tensors $z,w$ are nonzero already
in $R_0\otimes_k S_0$, and
their product is zero there.

The original $R_0$-algebra
$R_0\otimes_k S_0$ is flat:
a $k$-basis of $S_0$ gives a
free $R_0$-module basis.
It is finitely presented since
$S_0$ has a finite presentation
over $k$ by the Hilbert basis
theorem, and the same generators
and relations give its scalar
extension over $R_0$.
For each maximal ideal
$\mathfrak m\subset R_0$, the
Nullstellensatz (Theorem \ref{editorial-source-theorem-nullstellensatz}) gives the
original scalar isomorphism
$k\to R_0/\mathfrak m$.
The fibre map
$$
(R_0\otimes_k S_0)\otimes_{R_0}R_0/\mathfrak m
\longrightarrow S_0,
\qquad(r\otimes s)\otimes\overline a
\longmapsto\overline r\,\overline a\,s
$$
uses that isomorphism; its
inverse is $s\mapsto(1\otimes s)\otimes1$.
These fibres are irreducible.
The maximal ideals are dense
because $R_0$ is a finite type
algebra over a field and is
Jacobson, by Lemma
\ref{algebra-lemma-finite-type-field-Jacobson}.
Lemma \ref{editorial-source-lemma-flat-fibres-irreducible}
therefore makes the tensor's
spectrum irreducible. It is
also reduced by perfectness,
so its nilradical is the zero
prime and it is a domain.
This contradicts the original
nonzero $z,w$ with product zero.
Thus the arbitrary domain tensor
is a domain, and the preceding
quotient and field-extension
maps give the result for the
original $R,S$ over the original
separably closed field.
\end{proof}


\subsection{Editorial treatment of Algebra lines 11145--11204}

\hypertarget{algebra-context-138}{}

\noindent\href{https://github.com/stacks/stacks-project/blob/a04446e57ec1fbc252a871afcec7752fb2807b14/algebra.tex\#L11145}{Original source, lines 11145--11204}. Complete original and reviewed TeX passages accompany this note. Every intervention is separately recorded; the following treatment is editorial.


\begin{lemma}
\label{editorial-source-lemma-geometrically-irreducible}
Let $k$ be a field.
Let $R$ be a $k$-algebra.
The following are equivalent
\begin{enumerate}
\item for every field extension $k'/k$ the
spectrum of $R \otimes_k k'$ is irreducible,
\item for every finite separable field extension $k'/k$ the
spectrum of $R \otimes_k k'$ is irreducible,
\item the spectrum of $R \otimes_k \overline{k}$ is irreducible
where $\overline{k}$ is the separable algebraic closure of $k$, and
\item the spectrum of $R \otimes_k \overline{k}$ is irreducible
where $\overline{k}$ is the algebraic closure of $k$.
\end{enumerate}
\end{lemma}

\begin{proof}
Condition (1) implies (2) directly.
Under (2), taking the identity field
extension shows that the original
$R$ has nonempty spectrum, so $R\neq0$.
Let $k_s$ be a separable algebraic
closure of $k$. Its finite
subextensions $E/k$ are separable,
are directed under compositum,
and have union $k_s$.
Every original finite list of
coefficients in $k_s$ belongs to
one such $E$, so
$R\otimes_k k_s=\bigcup_E R\otimes_k E$
through the actual injective
coefficient maps.

This tensor is nonzero by a
vector-space basis argument.
Suppose $\mathfrak q_1,\mathfrak q_2$
are minimal primes in it.
For each original finite $E$,
the map $R\otimes_k E\to R\otimes_k k_s$
is flat: its module comparison
is scalar extension by the field
$k_s/E$. Going down makes the
contractions of both primes
minimal. Indeed a smaller prime
under a contraction would lift
to a smaller prime under the
given minimal one. By (2), the
ring $R\otimes_k E$ has only
one minimal prime, so these
contractions agree. Each element
of $R\otimes_k k_s$ lies in some
one of the displayed finite
tensors; membership in either
$\mathfrak q_i$ is therefore
determined by equal contractions.
Thus $\mathfrak q_1=\mathfrak q_2$.
A nonzero ring has a minimal
prime, so this proves (3),
including the required existence.

For (3) $\Rightarrow$ (4), choose
an algebraic closure $\overline k$
containing the specified $k_s$.
The extension $\overline k/k_s$
is purely inseparable, or is the
identity in characteristic zero,
by the separable-first theorem.
The map
$R\otimes_k k_s\to(R\otimes_k k_s)\otimes_{k_s}\overline k$
is a homeomorphism on spectra
by Lemma \ref{editorial-source-lemma-p-ring-map}
in positive characteristic, and
is the identity in characteristic
zero. The latter tensor is
identified with $R\otimes_k\overline k$
by
$(r\otimes a)\otimes b\mapsto r\otimes ab$,
whose inverse is
$r\otimes b\mapsto(r\otimes1)\otimes b$.
This proves (4) for these original
field embeddings.

For (4) $\Rightarrow$ (1), keep
an arbitrary original extension
$L/k$. A common receiving field
must contain compatible copies
over $k$; construct it as follows.
The tensor $L\otimes_k\overline k$
is nonzero by extension of a
vector-space basis. Choose a
prime $\mathfrak a$ in it and put
$$
F=\operatorname{Frac}((L\otimes_k\overline k)/\mathfrak a).
$$
The maps $a\mapsto[a\otimes1]$
and $b\mapsto[1\otimes b]$ embed
$L$ and $\overline k$ into this
nonzero field and agree on the
original $k$ by tensor balancing.
Thus they provide the needed
commutative field diagram, not
merely unrelated abstract copies.
The scalar comparison
$$
(R\otimes_k\overline k)\otimes_{\overline k}F
\longrightarrow R\otimes_k F,
\qquad(r\otimes a)\otimes b\longmapsto r\otimes ab
$$
has inverse $r\otimes b\mapsto(r\otimes1)\otimes b$.
Its source has irreducible
spectrum by Lemma
\ref{editorial-source-lemma-separably-closed-irreducible},
applied over the algebraically
closed $\overline k$ to the
original $R\otimes_k\overline k$
and the field $F$.

The map $B=R\otimes_k L\to R\otimes_k F$
is scalar extension by $F/L$.
It is faithfully flat: for any
nonzero $B$-module $M$, the
comparison
$M\otimes_B(B\otimes_L F)\simeq M\otimes_L F$
sends $m\otimes(b\otimes a)$ to
$bm\otimes a$, with inverse
$m\otimes a\mapsto m\otimes(1\otimes a)$;
extension of a vector-space
basis preserves nonzeroness
and exact sequences. More
explicitly, over each prime
$\mathfrak b$ of $B$ the
nonzero ring $\kappa(\mathfrak b)\otimes_L F$
has a prime. Its inverse image
is a prime of $B\otimes_L F$
contracting to $\mathfrak b$.
Hence the spectrum map is
surjective. The continuous
surjective image of the
nonempty irreducible spectrum
just proved is irreducible,
giving (1) for the original
$R\otimes_k L$.

Two further forms follow with
the same maps. First, in (2)
it suffices to test finite
Galois extensions. Every finite
separable $E/k$ embeds in a
finite separable normal splitting
field $N/k$: take a primitive
element for $E/k$ and adjoin
all roots of its separable
minimal polynomial in an
algebraic closure. These are
finitely many algebraic roots,
all separable, so the resulting
extension is finite and
separable; each embedding
permutes the roots, so it is
normal. The field extension
$N/E$ makes
$R\otimes_k E\to R\otimes_k N$
faithfully flat by the same
module comparison above.
Irreducibility of the latter
spectrum therefore descends
to the former, proving (2).

Second, for every fixed field
extension $L/k$, the original
$R$ is geometrically irreducible
over $k$ if and only if
$R\otimes_k L$ is geometrically
irreducible over $L$.
The forward direction uses
$(R\otimes_k L)\otimes_L F\simeq R\otimes_k F$
for every field $F/L$.
For the reverse direction and
any original field $E/k$, use
the same prime-of-a-tensor
construction to place $L$ and
$E$ compatibly in a field $F$.
The hypothesis over $L$ makes
$R\otimes_k F$ irreducible,
and the faithfully flat map
$R\otimes_k E\to R\otimes_k F$
descends irreducibility as
above. Every scalar comparison
is the displayed tensor map
and its displayed inverse.
\end{proof}


\subsection{Editorial treatment of Algebra lines 11220--11232}

\hypertarget{algebra-context-139}{}

\noindent\href{https://github.com/stacks/stacks-project/blob/a04446e57ec1fbc252a871afcec7752fb2807b14/algebra.tex\#L11220}{Original source, lines 11220--11232}. Complete original and reviewed TeX passages accompany this note. Every intervention is separately recorded; the following treatment is editorial.


\begin{lemma}
\label{editorial-source-lemma-separably-closed-irreducible-implies-geometric}
Let $k$ be a field.
Let $R$ be a $k$-algebra.
If $k$ is separably closed then $R$ is
geometrically irreducible over $k$ if and only if the
spectrum of $R$ is irreducible.
\end{lemma}

\begin{proof}
If $R$ is geometrically irreducible,
take the identity field extension in
Definition \ref{algebra-definition-geometrically-irreducible}.
It gives irreducibility of $\Spec(R)$.
Conversely suppose $\Spec(R)$ is
irreducible and $k$ is separably closed.
For every finite separable extension
$E/k$, each original element of $E$
has a separable irreducible minimal
polynomial over $k$. Separably closedness
forces that polynomial to be linear,
so the given embedding $k\to E$ is
surjective and hence an isomorphism.
The map $R\otimes_k E\to R$ sends
$r\otimes a$ to $r\iota^{-1}(a)$,
where $\iota:k\to E$ is that given
embedding; its inverse is $r\mapsto r\otimes1$.
Thus every finite separable test has
irreducible spectrum. Lemma
\ref{editorial-source-lemma-geometrically-irreducible}
gives geometric irreducibility over
the original $k$. The proof imposes
neither perfectness nor algebraic
closedness on that separably closed
field. If $R=0$, both irreducibility
assertions are false, in accordance
with the nonemptiness convention.
\end{proof}


\subsection{Editorial treatment of Algebra lines 11234--11254}

\hypertarget{algebra-context-140}{}

\noindent\href{https://github.com/stacks/stacks-project/blob/a04446e57ec1fbc252a871afcec7752fb2807b14/algebra.tex\#L11234}{Original source, lines 11234--11254}. Complete original and reviewed TeX passages accompany this note. Every intervention is separately recorded; the following treatment is editorial.


\begin{lemma}
\label{editorial-source-lemma-subalgebra-geometrically-irreducible}
Let $k$ be a field. Let $S$ be a $k$-algebra.
\begin{enumerate}
\item If $S$ is geometrically irreducible over $k$ so is every
$k$-subalgebra.
\item If all finitely generated $k$-subalgebras of $S$ are
geometrically irreducible, then $S$ is geometrically irreducible.
\item A directed colimit of geometrically irreducible $k$-algebras
is geometrically irreducible.
\end{enumerate}
\end{lemma}

\begin{proof}
For a nonzero ring $A$, its spectrum
is irreducible exactly when its
nilradical is prime, by Lemma
\ref{editorial-source-lemma-irreducible}. Equivalently,
whenever $ab$ is nilpotent, at
least one of $a,b$ is nilpotent.
We keep nonzeroness in this
criterion; the zero ring has
empty spectrum and is not
irreducible.

For (1), retain the original
subalgebra $S'\subseteq S$ and
any original field extension
$L/k$. The map
$S'\otimes_k L\to S\otimes_k L$
is injective by extension of
a vector-space basis. For any
injective ring map $A\to B$,
the nilradical of $A$ is exactly
the inverse image of that of
$B$: a power becomes zero in
$B$ exactly when it was zero
in $A$. Here the receiving
nilradical is a proper prime,
so its inverse image is also
a proper prime. The source is
nonzero because its $1$ remains
the nonzero $1$ of the target.
The criterion proves that
$S'\otimes_k L$ is irreducible,
for every such $L$, proving (1).
This argument works for any
injective $k$-algebra map, with
its actual image subalgebra.

For (3), let $(S_i)$ be a
nonempty directed system of
geometrically irreducible
$k$-algebras with their original
unital transition maps; they
need not be injective. Put
$S=\mathop{\rm colim}_i S_i$.
For each original $L/k$ there
are inverse maps
$$
S\otimes_k L\longrightarrow
\mathop{\rm colim}_i(S_i\otimes_k L),
\qquad [i,s]\otimes a\longmapsto[i,s\otimes a],
$$
and
$$
[i,\sum_j s_j\otimes a_j]
\longmapsto\sum_j[i,s_j]\otimes a_j.
$$
All expressions use finitely
many elements, so a common
later index verifies balancing,
independence of representatives,
and both composites. Let $A_i$
denote the original $S_i\otimes_k L$.
Each is nonzero with prime
nilradical. Their unital colimit
$A$ is nonzero: an equality
$1=0$ there would hold at a
later stage, contradicting
$1\neq0$ in that $A_i$.

Suppose original elements
$a,b\in A$ satisfy $(ab)^n=0$.
Choose representatives at a
common stage $i$. The zero
equality holds at some later
stage $j$. Its nilradical is
prime, so the image of $a$
or of $b$ is nilpotent at
that stage. The corresponding
element in the original
colimit is then nilpotent.
This proves that the nilradical
of $A$ is prime, and proves
(3) after the explicit tensor
comparison. In particular no
injection assumption on the
transition maps is used.
If an empty colimit is allowed
in the category of $k$-algebras,
it is its initial object $k$,
which also satisfies the
assertion.

For (2), the original finitely
generated $k$-subalgebras of $S$
form a directed system under
inclusion: adjoining the union
of two finite lists gives a
common member. Every original
element of $S$ occurs in such
a subalgebra, so their colimit
is exactly $S$. Apply (3).
The hypothesis excludes $S=0$,
since its finitely generated
subalgebra generated by the
image of $k$ would itself be
zero and would not satisfy
the stated hypothesis.

The same nilradical criterion
also proves localization
stability with its necessary
nonemptiness qualification.
Assume $S$ is geometrically irreducible.
Let $W\subseteq S$ be a
multiplicative subset and
assume $W^{-1}S\neq0$.
Since $S$ is geometrically
irreducible, its nilradical
$N$ is prime; $W$ avoids $N$,
because a nilpotent becoming
invertible would make the
localization zero. For every
original $L/k$, the injection
$S\to S\otimes_k L$ shows
that no image of an element
of $W$ is nilpotent. The
prime nilradical of that
tensor therefore survives
localization and gives a
nonzero ring with one minimal
prime. Explicitly the nilradical
of a localization $B_W$ is
$W^{-1}\operatorname{Nil}(B)$:
if $(b/w)^n=0$, some $v\in W$
has $v b^n=0$, and then
$(vb)^n=v^{n-1}(v b^n)=0$,
while $b/w=(vb)/(vw)$.
The converse follows by taking
powers of a nilpotent numerator.
The original tensor-localization
maps are
$$
(W^{-1}S)\otimes_k L
\longrightarrow (W\otimes1)^{-1}(S\otimes_k L),
\qquad(s/w)\otimes a\longmapsto(s\otimes a)/(w\otimes1),
$$
and the inverse sends
$\sum_i(s_i\otimes a_i)/(w\otimes1)$
to $\sum_i(s_i/w)\otimes a_i$.
The two localization universal
properties prove these maps
well-defined and inverse.
Thus every nonzero localization
of the original $S$ remains
geometrically irreducible.
\end{proof}


\subsection{Editorial treatment of Algebra lines 11256--11290}

\hypertarget{algebra-context-141}{}

\noindent\href{https://github.com/stacks/stacks-project/blob/a04446e57ec1fbc252a871afcec7752fb2807b14/algebra.tex\#L11256}{Original source, lines 11256--11290}. Complete original and reviewed TeX passages accompany this note. Every intervention is separately recorded; the following treatment is editorial.


\begin{lemma}
\label{editorial-source-lemma-geometrically-irreducible-any-base-change}
Let $k$ be a field.
Let $S$ be a geometrically irreducible $k$-algebra.
Let $R$ be any $k$-algebra.
The map
$$
\Spec(R \otimes_k S) \longrightarrow \Spec(R)
$$
induces a bijection on irreducible components.
\end{lemma}

\begin{proof}
Write $A=R\otimes_k S$ and retain
$f:\Spec(A)\to\Spec(R)$. The
original $S$ is nonzero, since
geometric irreducibility includes
nonemptiness. If $R=0$, then
$A=0$ and both sets of components
are empty, so the assertion
holds. Otherwise, for every
$\mathfrak p\in\Spec(R)$ the
original quotient and fibre maps
give
$$
A/\mathfrak pA\simeq(R/\mathfrak p)\otimes_k S
\longrightarrow\kappa(\mathfrak p)\otimes_k S.
$$
The quotient isomorphism sends
$[r\otimes s]$ to $\overline r\otimes s$
and has inverse
$\overline r\otimes s\mapsto[r\otimes s]$.
The arrow is injective, since
$R/\mathfrak p\to\kappa(\mathfrak p)$
is an injection of $k$-vector
spaces and tensoring by $S$
preserves it. Its target has
prime nilradical by the given
geometric irreducibility of $S$.
The inverse image of this
nilradical is precisely the
nilradical of the source, by
injectivity, and is therefore
a proper prime. In particular
$A/\mathfrak pA$ has irreducible
nonempty spectrum.

Consequently the original ideal
$$
\sigma(\mathfrak p)=\sqrt{\mathfrak pA}
$$
is prime, and $f^{-1}(V(\mathfrak p))$
is the irreducible closed subset
$V(\sigma(\mathfrak p))$.
Its contraction is exactly
$\mathfrak p$. To verify this,
the injection
$R/\mathfrak p\to(R/\mathfrak p)\otimes_k S$
can be split as a linear map
by choosing a $k$-linear
functional $S\to k$ taking
$1$ to $1$. Hence an element
of $R/\mathfrak p$ becomes
nilpotent in that tensor
exactly when it was nilpotent
in the original domain, that
is, exactly when it is zero.
This proves the contraction
claim with the given scalar
map.

The map $R\to A$ is flat: a
$k$-basis of $S$ gives a free
$R$-module basis of $A$.
Going down therefore sends
every minimal prime of $A$
to a minimal prime of $R$.
If $\mathfrak p$ is minimal
in $R$, choose a minimal
$\mathfrak q_0\subseteq\sigma(\mathfrak p)$
of $A$. Its contraction is
a minimal prime contained
in $\mathfrak p$, hence is
exactly $\mathfrak p$.
It contains $\mathfrak pA$,
so it contains
$\sqrt{\mathfrak pA}=\sigma(\mathfrak p)$.
Thus equality holds, proving
that $\sigma(\mathfrak p)$
is minimal. Conversely, for
any minimal $\mathfrak q$
of $A$, put
$\mathfrak p=\mathfrak q\cap R$.
Then $\sigma(\mathfrak p)\subseteq\mathfrak q$,
so minimality makes them equal.
These actual inverse maps on
minimal primes prove the
bijection on components via
Lemma \ref{editorial-source-lemma-irreducible}.

The construction gives more:
$\sigma$ is a continuous section
of $f$ on the entire spectra,
selecting the generic point
of each original fibre.
For every $\mathfrak p$, its
prime is contained in every
prime of $A$ over $\mathfrak p$,
and its contraction equals
$\mathfrak p$, so it is the
unique generic fibre prime.
For an original $a\in A$ one
has the exact equality
$$
\sigma^{-1}(D(a))=f(D(a)).
$$
If a prime above $\mathfrak p$
does not contain $a$, the
smaller $\sigma(\mathfrak p)$
does not contain it. Conversely
if $a\notin\sigma(\mathfrak p)$,
this prime itself is a point
of $D(a)$ above $\mathfrak p$.
Lemma \ref{editorial-source-lemma-map-into-tensor-algebra-open}
makes the right side open,
proving continuity. Since
$f\sigma$ is the identity,
the section is a topological
embedding, with inverse the
restriction of $f$ to its
image. That image is dense:
every nonempty basic open
$D(a)$ has nonempty image,
and the displayed equality
puts a section point in it.
Also
$$
\overline{\{\sigma(\mathfrak p)\}}
=V(\sqrt{\mathfrak pA})
=f^{-1}(V(\mathfrak p)).
$$
Restricting the two continuous
inverse maps to the minimal
primes proves a homeomorphism
of their subspaces in the
respective spectra, strengthening
the stated component bijection.

Finally this section is
compatible with arbitrary
original $k$-algebra base change
$\theta:R\to R'$. If $\mathfrak p'$
contracts to $\mathfrak p$,
the induced map of fields
$\kappa(\mathfrak p)\to\kappa(\mathfrak p')$
gives an injection
$$
\kappa(\mathfrak p)\otimes_k S
\longrightarrow\kappa(\mathfrak p')\otimes_k S.
$$
The inverse image of the
nilradical on the right is
the nilradical on the left,
since injective ring maps
reflect nilpotence. The two
quotient-fibre descriptions
above therefore show that
$\sigma_{R'}(\mathfrak p')$
contracts to
$\sigma_R(\mathfrak p)$ in
$R\otimes_k S$. This proves
the stated compatibility
through the original tensor
map $r\otimes s\mapsto\theta(r)\otimes s$.
The section is a continuous
map of these spectra; no
extra residue-field
isomorphisms are asserted.
\end{proof}


\subsection{Editorial treatment of Algebra lines 11296--11320}

\hypertarget{algebra-context-142}{}

\noindent\href{https://github.com/stacks/stacks-project/blob/a04446e57ec1fbc252a871afcec7752fb2807b14/algebra.tex\#L11296}{Original source, lines 11296--11320}. Complete original and reviewed TeX passages accompany this note. Every intervention is separately recorded; the following treatment is editorial.


\begin{lemma}
\label{editorial-source-lemma-field-extension-geometrically-irreducible}
Let $K/k$ be a field extension. If $k$ is algebraically closed in $K$, then
$K$ is geometrically irreducible over $k$.
\end{lemma}

\begin{proof}
We prove the assertion under the weaker
hypothesis that every element of the
original $K$ separably algebraic over
$k$ belongs to $k$. The stated relative
algebraic closedness implies this.
By Lemma \ref{editorial-source-lemma-geometrically-irreducible}
it suffices to treat each original
finite separable extension $k'/k$.
Choose its primitive element $\alpha$
using Fields, Lemma
\ref{fields-lemma-primitive-element},
and retain its full minimal polynomial
$$
P=T^d+a_1T^{d-1}+\cdots+a_d\in k[T].
$$
The map $K[T]/(P)\to K\otimes_k k'$
sends $\sum_{j=0}^{d-1}b_jT^j$ to
$\sum_{j=0}^{d-1}b_j\otimes\alpha^j$.
Its inverse sends $b\otimes Q(\alpha)$
to the class of $bQ(T)$.
The original basis $1,\alpha,\ldots,\alpha^{d-1}$
and monic division by $P$ show that
these maps are well-defined inverses.

Suppose the original polynomial has
a nontrivial factorization $P=P_1P_2$
in $K[T]$. Set $r=\deg(P_1)$ and
$c=\operatorname{lc}(P_1)\in K^\times$.
Then $1\leq r\leq d-1$ and
$\operatorname{lc}(P_2)=c^{-1}$.
In an algebraic closure $\Omega_K$ of
$K$, the original separable $P$ has
distinct roots $\beta_1,\ldots,\beta_d$.
There is a subset $J$ of cardinality
$r$ such that the complete original
factors are
$$
P_1=c\prod_{i\in J}(T-\beta_i),\qquad
P_2=c^{-1}\prod_{i\notin J}(T-\beta_i).
$$
Define the auxiliary polynomials
$Q_1=c^{-1}P_1$ and $Q_2=cP_2$,
retaining the exact identity
$P=(cQ_1)(c^{-1}Q_2)=Q_1Q_2$.
For $0\leq j\leq r$, the coefficient
of $T^{r-j}$ in the original $P_1$ is
$$
c(-1)^j\sum_{\substack{H\subseteq J\\|H|=j}}
\prod_{i\in H}\beta_i.
$$
For $P_2$ the corresponding coefficient
has factor $c^{-1}$ and the same signed
sum over subsets of the complementary
root set. Empty products give the
original leading coefficients $c,c^{-1}$.

Each $\beta_i$ is separably algebraic
over $k$ because $P$ is separable.
The finite compositum $k(\beta_1,\ldots,\beta_d)$
is separable over $k$, by Fields,
Lemma \ref{fields-lemma-separable-permanence}
and permanence under enlarging the
base field. Consequently every displayed
coefficient ratio, namely each coefficient
of $Q_1$ and $Q_2$, is separably algebraic
over $k$. These ratios also lie in
the original $K$ by their definitions,
so the weakened hypothesis puts them
in $k$. Both $Q_1,Q_2$ thus belong to
$k[T]$ and have positive degree.
Their product is the original $P$,
contradicting its irreducibility over $k$.
It follows that $K[T]/(P)$ is a field;
in particular its spectrum is irreducible.
All finite separable tests now give
geometric irreducibility as asserted.

The leading factors cannot be discarded
from the coefficient argument in Lemma
\ref{editorial-source-lemma-polynomials-divide}, whose
statement concerns monic polynomials.
For example, over the original
$\mathbf Q(u)$ one has
$$
T^2-1=\bigl(u(T-1)\bigr)\bigl(u^{-1}(T+1)\bigr).
$$
The coefficients $u$ and $u^{-1}$
are not algebraic over $\mathbf Q$;
a nonzero polynomial equation for
$u^{-1}$ would, after multiplication
by its full power of $u$, give a
nonzero polynomial equation for $u$.
This example concerns the coefficient
inference for an arbitrary monic
product, not the irreducible-polynomial
hypothesis of the lemma. The proof
above keeps $c,c^{-1}$ and proves
exactly the required claim about
their coefficient ratios. It also
proves the stronger conclusion that
$K\otimes_k k'$ is a field for every
finite separable original $k'/k$.
\end{proof}


\subsection{Editorial treatment of Algebra lines 11322--11338}

\hypertarget{algebra-context-143}{}

\noindent\href{https://github.com/stacks/stacks-project/blob/a04446e57ec1fbc252a871afcec7752fb2807b14/algebra.tex\#L11322}{Original source, lines 11322--11338}. Complete original and reviewed TeX passages accompany this note. Every intervention is separately recorded; the following treatment is editorial.


\begin{lemma}
\label{editorial-source-lemma-geometrically-irreducible-transitive}
Let $K/k$ be a geometrically irreducible field extension.
Let $S$ be a geometrically irreducible $K$-algebra.
Then $S$ is geometrically irreducible over $k$.
\end{lemma}

\begin{proof}
Let $k'/k$ be any original finite
separable extension. By Fields,
Lemma \ref{fields-lemma-primitive-element},
write $k'=k(\alpha)$ with monic
separable minimal polynomial $P$.
The inverse tensor-polynomial maps
from Lemma
\ref{editorial-source-lemma-field-extension-geometrically-irreducible}
give $K'=K\otimes_k k'\simeq K[T]/(P)$.
Since $P$ is separable, there are
$U,V\in k[T]$ with $UP+VP'=1$.
The same full identity holds in
$K[T]$. Thus $P$ has no repeated
irreducible factor there: a repeated
factor would divide $P$ and $P'$
and hence divide $1$.

Factor $P=\prod_{i=1}^m Q_i$ into
its distinct monic irreducible
factors in $K[T]$. The leading
coefficient of $P$ is $1$, so this
product has exactly the original
leading coefficient. The Chinese
remainder map
$$
K[T]/(P)\longrightarrow\prod_{i=1}^m K[T]/(Q_i)
$$
sends the class of $F$ to its residue
classes. For each $i$, Bezout gives
$A_i\prod_{j\ne i}Q_j+B_iQ_i=1$;
the inverse sends $([F_i])_i$ to
$[\sum_i F_i A_i\prod_{j\ne i}Q_j]$.
These formulas give both composites
by reduction modulo every $Q_i$.
Each factor is a nonzero field,
so the spectrum has exactly $m$
points and is irreducible exactly
when $m=1$. The hypothesis on
$K/k$ therefore makes $K'$ a field.
It is finite over $K$ with basis
$1,T,\ldots,T^{\deg(P)-1}$, and
separable since its irreducible
polynomial $P$ has the displayed
Bezout identity with its derivative.

There are inverse original maps
$$
S\otimes_K(K\otimes_k k')\longrightarrow S\otimes_k k',
\qquad s\otimes(a\otimes b)\longmapsto as\otimes b,
$$
and $s\otimes b\mapsto s\otimes(1\otimes b)$.
The $K$-balancing relation proves
both composites are the identity.
Geometric irreducibility of $S/K$
applies to the actual field $K'/K$,
so this tensor has irreducible
spectrum. Every finite separable
$k'/k$ passes the test, and Lemma
\ref{editorial-source-lemma-geometrically-irreducible}
proves the result over the original $k$.
\end{proof}


\subsection{Editorial treatment of Algebra lines 11340--11375}

\hypertarget{algebra-context-144}{}

\noindent\href{https://github.com/stacks/stacks-project/blob/a04446e57ec1fbc252a871afcec7752fb2807b14/algebra.tex\#L11340}{Original source, lines 11340--11375}. Complete original and reviewed TeX passages accompany this note. Every intervention is separately recorded; the following treatment is editorial.


\begin{lemma}
\label{editorial-source-lemma-geometrically-irreducible-base-change-transcendental}
Let $K/k$ be a field extension. The following are equivalent
\begin{enumerate}
\item $K$ is geometrically irreducible over $k$, and
\item the induced extension $K(t)/k(t)$ of purely transcendental extensions
is geometrically irreducible.
\end{enumerate}
\end{lemma}

\begin{proof}
Assume (1). Let $\Omega$ be an
algebraic closure of the original
$k(t)$. The algebra $K\otimes_k\Omega$
has irreducible spectrum by (1).
Put $U=k[t]\setminus\{0\}$ and
$V=K[t]\setminus\{0\}$, keeping
the same indeterminate $t$.
There is an isomorphism
$$
K\otimes_k k(t)\longrightarrow U^{-1}K[t],
\qquad a\otimes(P(t)/Q(t))\longmapsto aP(t)/Q(t).
$$
The inverse sends
$(\sum_i a_it^i)/Q(t)$ to
$\sum_i a_i\otimes(t^i/Q(t))$.
The quotient relations for the
original denominators verify
independence of representatives,
and the formulas give both composites.
Localizing further at the images
of all of $V$ gives precisely the
original fraction field $K(t)$.

The two original scalar maps
$$
(K\otimes_k k(t))\otimes_{k(t)}\Omega
\longrightarrow K\otimes_k\Omega,
\qquad(a\otimes b)\otimes c\longmapsto a\otimes bc
$$
and $a\otimes c\mapsto(a\otimes1)\otimes c$
are inverse. Localization commutes
with these tensors: the fraction
$b/v\otimes c$ goes to
$(b\otimes c)/(v\otimes1)$, with
inverse determined by inverting
the same original $v$. Hence
$K(t)\otimes_{k(t)}\Omega$ is a
localization of $K\otimes_k\Omega$
at exactly those images of $V$.
It is nonzero because it is a
tensor product of two nonzero
fields over $k(t)$; extending
$1$ to vector-space bases proves
this directly. A nonzero
localization of a ring with
prime nilradical again has prime
nilradical, as proved in Lemma
\ref{editorial-source-lemma-subalgebra-geometrically-irreducible}.
Thus its spectrum is irreducible.
The algebraic-closure test of Lemma
\ref{editorial-source-lemma-geometrically-irreducible}
proves (2).

Conversely, assume (2) and retain
any original extension $k'/k$.
Put $A=K\otimes_k k'$ and let $W$
be the multiplicative subset of
$A[t]$ generated by the images of
all nonzero polynomials of $K[t]$
and of $k'[t]$. Both coefficient
maps into $A$ are injective:
extension of a vector-space basis
over $k$ proves this. More strongly,
the leading coefficient of each
such polynomial is a unit of $A$,
with inverse given by the inverse
of the original nonzero field
coefficient in the same tensor
factor. Multiplication by a
polynomial with unit leading
coefficient is injective on $A[t]$:
if $H$ is nonzero of degree $n$
and leading coefficient $h_n$,
the leading coefficient of its
product with a degree-$d$
polynomial of leading coefficient
$c$ is exactly $c h_n\neq0$.
Products of these injective
multiplication maps remain
injective. Thus every element
of $W$ is a nonzerodivisor and
$A\to A[t]\to W^{-1}A[t]$ is
injective, without assuming $A$
reduced or a domain.

The full receiving algebra is
exactly
$$
K(t)\otimes_{k(t)}k'(t)\simeq W^{-1}A[t].
$$
To construct the map from left
to right, use the original
coefficient embeddings and send
$$
\frac{F(t)}{f(t)}\otimes\frac{G(t)}{g(t)}
\longmapsto\frac{F(t)G(t)}{f(t)g(t)},
\qquad
f\in K[t]\setminus\{0\},\quad
g\in k'[t]\setminus\{0\}.
$$
The field-fraction relations hold
after precisely these denominators
are inverted. The two images of
every rational function in $k(t)$
agree, so the map is balanced over
$k(t)$. For the inverse, send
$a\otimes b\in A$ to $a\otimes b$
in the left-hand tensor and send
$t$ to $t\otimes1=1\otimes t$.
The image of each original
$f\in K[t]\setminus\{0\}$ is
the unit $f(t)\otimes1$, and
each original nonzero $g\in k'[t]$
gives the unit $1\otimes g(t)$.
The universal property therefore
extends this map to $W^{-1}A[t]$.
The composites fix $A$, $t$ and
each specified inverse denominator,
so they are identities on their
generating sets. In particular,
the injection just proved is
the original tensor-constant map
used in the source argument.

By (2), this receiving tensor
has nonempty irreducible spectrum,
so its nilradical is a proper
prime. Its inverse image under
the injection from $A$ is exactly
$\operatorname{Nil}(A)$, because
an injection reflects the vanishing
of every power. Hence $A$ is
nonzero with prime nilradical,
and its spectrum is irreducible.
This works for every original
$k'/k$, proving (1).

Iteration gives the same equivalence
after adjoining any specified finite
list of common independent
indeterminates. At each step the
proved one-variable equivalence
applies to the actual preceding
field extension and retains all
its polynomial denominators.
The empty list is the identity
extension, so it also satisfies
the assertion.
\end{proof}


\subsection{Editorial treatment of Algebra lines 11377--11389}

\hypertarget{algebra-context-145}{}

\noindent\href{https://github.com/stacks/stacks-project/blob/a04446e57ec1fbc252a871afcec7752fb2807b14/algebra.tex\#L11377}{Original source, lines 11377--11389}. Complete original and reviewed TeX passages accompany this note. Every intervention is separately recorded; the following treatment is editorial.


\begin{lemma}
\label{editorial-source-lemma-geometrically-irreducible-add-transcendental}
Let $K/L/M$ be a tower of fields with $L/M$ geometrically irreducible.
Let $x \in K$ be transcendental over $L$. Then $L(x)/M(x)$ is geometrically
irreducible.
\end{lemma}

\begin{proof}
Keep the original inclusions
$M\subseteq L\subseteq K$ and
the original $x\in K$.
Transcendence over $L$ makes
evaluation at $x$ injective on
$L[t]$, and hence on its subring
$M[t]$. In particular no original
nonzero denominator vanishes.
The maps
$$
\lambda:L(t)\longrightarrow L(x),\quad
\frac{P(t)}{Q(t)}\longmapsto\frac{P(x)}{Q(x)},
\qquad
\mu:M(t)\longrightarrow M(x)
$$
with the same evaluation formula
are injective field maps. They
are surjective because their
images are fields containing the
original coefficient fields and
$x$, respectively. Their inverses
replace $x$ by $t$ in a rational
expression; injectivity proves
independence of that expression.
The square with the original
inclusions commutes because the
formulas agree on every $M$
coefficient, on $t$, and on
every inverse denominator.

For any original field $E/M(x)$,
use its given embedding composed
with $\mu$ to view it over $M(t)$.
The map
$$
L(t)\otimes_{M(t)}E\longrightarrow L(x)\otimes_{M(x)}E,
\qquad a\otimes b\longmapsto\lambda(a)\otimes b
$$
is well-defined by that commuting
square, with inverse
$c\otimes b\mapsto\lambda^{-1}(c)\otimes b$.
Lemma
\ref{editorial-source-lemma-geometrically-irreducible-base-change-transcendental}
makes the left-hand spectrum
irreducible from the hypothesis
on $L/M$. This proves the assertion
for every given $E/M(x)$.
Conversely the same tensor maps,
using $\mu^{-1}$ to transport a
field over $M(t)$, show that
geometric irreducibility of
$L(x)/M(x)$ implies that of
$L(t)/M(t)$. The converse in the
preceding lemma then gives that
of $L/M$. Thus, under the stated
transcendence hypothesis, the
original conclusion is in fact
equivalent to geometric
irreducibility of $L/M$.
\end{proof}


\subsection{Editorial treatment of Algebra lines 11391--11425}

\hypertarget{algebra-context-146}{}

\noindent\href{https://github.com/stacks/stacks-project/blob/a04446e57ec1fbc252a871afcec7752fb2807b14/algebra.tex\#L11391}{Original source, lines 11391--11425}. Complete original and reviewed TeX passages accompany this note. Every intervention is separately recorded; the following treatment is editorial.


\begin{lemma}
\label{editorial-source-lemma-geometrically-irreducible-separable-elements}
Let $K/k$ be a field extension. The following are equivalent
\begin{enumerate}
\item $K/k$ is geometrically irreducible, and
\item every element $\alpha \in K$ separably algebraic over $k$ is
in $k$.
\end{enumerate}
\end{lemma}

\begin{proof}
Assume (1) and let the original
$\alpha\in K$ be separably algebraic
over $k$. Put $k'=k(\alpha)$ and
let $P\in k[T]$ be its monic
minimal polynomial of degree $d$.
Choose an algebraic closure
$\overline k$ of $k$ and retain
the distinct roots
$\alpha_1,\ldots,\alpha_d$ of $P$
in it. The original inclusion
$k'\subseteq K$ and Lemma
\ref{editorial-source-lemma-subalgebra-geometrically-irreducible}
make $k'/k$ geometrically irreducible.
The tensor-polynomial comparison
and the following explicit Chinese
remainder maps give
$$
k'\otimes_k\overline k\simeq
\overline k[T]/(P)\simeq\prod_{i=1}^d\overline k.
$$
The first sends $Q(\alpha)\otimes b$
to $[bQ(T)]$; its inverse sends
$[\sum_j b_jT^j]$ to
$\sum_j\alpha^j\otimes b_j$.
The second sends $[F]$ to
$(F(\alpha_i))_i$; its inverse is
$$
(b_1,\ldots,b_d)\longmapsto
\left[\sum_{i=1}^d b_i
\frac{\prod_{j\ne i}(T-\alpha_j)}
{\prod_{j\ne i}(\alpha_i-\alpha_j)}\right].
$$
All displayed denominators are
nonzero because the roots are
distinct. The interpolation
polynomials evaluate to the
coordinate vectors. Conversely,
a polynomial vanishing at all
these roots is divisible by
their full product $P$, proving
the other composite. These are
the embedding maps of Fields,
Lemma
\ref{fields-lemma-finite-separable-tensor-alg-closed},
with their inverse now explicit.
The product has $d$ distinct
maximal and minimal primes,
the kernels of its projections.
It has irreducible spectrum
only if $d=1$. Therefore the
original $\alpha$ belongs to $k$,
proving (2).

Assume (2) and retain the original
subfield $k'\subseteq K$ consisting
of all elements algebraic over $k$.
It is a field: a finite list of
such elements generates a finite
field extension of $k$, containing
their sums, products and inverses
of nonzero elements. Furthermore
$k'$ is algebraically closed in
$K$. Indeed if $\beta\in K$ is
algebraic over $k'$, the finitely
many coefficients of a polynomial
equation for $\beta$ belong to a
finite extension $E/k$ in $k'$.
Then $E(\beta)/E$ is finite, so
$E(\beta)/k$ is finite and
$\beta\in k'$ by the definition.
Lemma
\ref{editorial-source-lemma-field-extension-geometrically-irreducible}
therefore makes $K/k'$ geometrically
irreducible, using the full relative
algebraic closedness just proved.

By Fields, Lemma
\ref{fields-lemma-separable-first},
the algebraic extension $k'/k$
is purely inseparable over its
maximal separable subextension.
Condition (2) says precisely
that this subextension is $k$.
In characteristic zero this
forces $k'=k$. In characteristic
$p>0$, for every original $L/k$
the map $L\to L\otimes_k k'$
is injective by extension of
a vector-space basis. Its
receiving algebra is generated
over $L$ by $1\otimes a$ with
$a\in k'$, and each such
generator has some $p$-power
in the image of $L$, while
its corresponding integer
$p$-power multiple is zero.
Lemma \ref{editorial-source-lemma-p-ring-map}
thus identifies its spectrum
with the one-point $\Spec(L)$.
This proves geometric
irreducibility of $k'/k$,
including infinite extensions
with no common exponent assumed.
Lemma
\ref{editorial-source-lemma-geometrically-irreducible-transitive}
now proves (1) for the original
tower $K/k'/k$.

The criterion has an exact
tensor-field form: (1) and (2)
are also equivalent to
$K\otimes_k E$ being a field
for every separable algebraic
extension $E/k$. To prove it
from (2), each finite separable
subextension $E_i/k$ gives a
field $K\otimes_k E_i$ by the
coefficient proof of Lemma
\ref{editorial-source-lemma-field-extension-geometrically-irreducible}.
These fields form a directed
system of injective original
maps whose union is
$K\otimes_k E$: every tensor
has finitely many coefficients.
A nonzero element belongs to
one of those fields and has
its inverse there. Their union
is therefore a field. The
converse follows at once by
applying the finite separable
tests of Lemma
\ref{editorial-source-lemma-geometrically-irreducible}.

For any compatible embeddings
of the original $K$ and $E$
into a common field $F$, the
multiplication map
$$
K\otimes_k E\longrightarrow KE\subseteq F,
\qquad a\otimes b\longmapsto ab
$$
is consequently an isomorphism:
it is a unital map from a field,
so is injective; its image is
a field containing both original
embedded fields and consists of
their finite tensor sums, hence
equals their compositum.
For finite $E/k$, extension of
its original $k$-basis gives
the exact degree equality
$[KE:K]=[E:k]$. More generally,
every original finite linearly
independent family over $k$
in $E$ remains independent
over $K$: its tensor extension
is independent, and the proved
multiplication map is injective.
This identifies the precise
linear-disjointness consequence
without replacing the original
field embeddings.
\end{proof}


\subsection{Editorial treatment of Algebra lines 11427--11443}

\hypertarget{algebra-context-147}{}

\noindent\href{https://github.com/stacks/stacks-project/blob/a04446e57ec1fbc252a871afcec7752fb2807b14/algebra.tex\#L11427}{Original source, lines 11427--11443}. Complete original and reviewed TeX passages accompany this note. Every intervention is separately recorded; the following treatment is editorial.


\begin{lemma}
\label{editorial-source-lemma-make-geometrically-irreducible}
Let $K/k$ be a field extension. Consider the subextension $K/k'/k$ consisting
of elements separably algebraic over $k$. Then $K$ is geometrically irreducible
over $k'$. If $K/k$ is a finitely generated field extension, then
$[k' : k] < \infty$.
\end{lemma}

\begin{proof}
The original set $k'$ is a subfield
by Fields, Lemma \ref{fields-lemma-separable-elements}.
If $a\in K$ is separably algebraic
over $k'$, transitivity of separable
algebraic extensions (Fields, Lemma
\ref{fields-lemma-separable-permanence})
makes $a$ separably algebraic over
$k$. Hence $a\in k'$ by its
definition. Lemma
\ref{editorial-source-lemma-geometrically-irreducible-separable-elements}
now proves geometric irreducibility
of the original $K/k'$.

For completeness, retain the full
degree argument of Fields, Lemma
\ref{fields-lemma-algebraic-closure-in-finitely-generated}.
Choose a finite transcendence basis
$x_1,\ldots,x_r$ of the finitely
generated extension $K/k$ and put
$$
F=k(x_1,\ldots,x_r),\qquad n=[K:F]<\infty.
$$
Here $r=0$ is allowed, with $F=k$.
The remaining members of a finite
generating set for $K/k$ are
algebraic over $F$, so adjoining
them successively proves the
displayed finiteness. Let $E/k$
be any original finite algebraic
subextension of $K/k$, and choose
a basis $e_1=1,e_2,\ldots,e_d$ of
$E/k$, where $d=[E:k]$.

The original $x_i$ remain
algebraically independent over $E$.
Here is a direct proof retaining
every coefficient. For nonzero
$P\in E[T_1,\ldots,T_r]$, let $M_P$
be the matrix of multiplication
by $P$ on the free $k[T_1,\ldots,T_r]$-module
$E[T_1,\ldots,T_r]$ in the basis
$e_1,\ldots,e_d$. This multiplication
is injective because the latter
polynomial ring is a domain.
After localization to $k(T_1,\ldots,T_r)$,
the matrix is an injective square
linear map on a finite-dimensional
vector space, hence
$N_P=\det(M_P)\neq0$ in $k[T_1,\ldots,T_r]$.
Write $H_P\in E[T_1,\ldots,T_r]$
for the element whose coordinate
column is $\operatorname{adj}(M_P)(1,0,\ldots,0)^t$.
The full adjugate identity gives
$$
P H_P=N_P.
$$
If $P(x_1,\ldots,x_r)=0$, this
identity would imply
$N_P(x_1,\ldots,x_r)=0$, contradicting
the original independence over $k$.
This proves the asserted independence
over $E$, also for $r=0$.

Consequently the original evaluation
map induces inverse isomorphisms
$$
E\otimes_k F\longrightarrow E(x_1,\ldots,x_r),
\qquad a\otimes h(x)\longmapsto a h(x).
$$
To specify the inverse, write a
numerator in the basis $e_i$ and
retain each original denominator
$Q(x)$ with $Q\in E[T_1,\ldots,T_r]\setminus\{0\}$.
Use exactly the identity
$Q^{-1}=H_Q/N_Q$ just proved,
then expand the numerator $P H_Q$
in the same basis. This gives
its unique expression
$\sum_i e_i f_i(x)$ with $f_i\in F$,
and the inverse sends it to
$\sum_i e_i\otimes f_i(x)$.
Uniqueness follows by multiplying
all denominators in $k[T_1,\ldots,T_r]$
and comparing polynomial coefficients
in the basis $e_i$; no denominator
vanishes at the original $x$.
Thus these are well-defined inverse
maps and
$$
[E:k]=[E(x_1,\ldots,x_r):F]\leq[K:F]=n.
$$

In fact this bounds the full
algebraic constant field $k_a$
inside $K$, not only $k'$.
The degrees of its finite
subextensions lie among the
integers $1,\ldots,n$, so choose
one, $E_0$, of maximal degree.
For any $a\in k_a$, the original
$E_0(a)/k$ is finite and has
degree at least $[E_0:k]$.
Maximality and degree multiplicativity
give $[E_0(a):E_0]=1$. Hence every
$a$ lies in $E_0$, proving
$k_a=E_0$ and the precise bound
$$
[k':k]\leq[k_a:k]\leq[K:k(x_1,\ldots,x_r)].
$$

There is also an exact description
of the algebraic intermediate
fields $k\subseteq L\subseteq K$
over which $K$ is geometrically
irreducible: they are exactly
those containing $k'$.
For necessity, every $a\in k'$
is separably algebraic over $L$,
because its minimal polynomial
over $L$ divides its original
separable polynomial over $k$.
Geometric irreducibility of $K/L$
and the preceding criterion force
$a\in L$. For sufficiency, suppose
$k'\subseteq L$ and $L/k$ algebraic.
The full algebraic constant field
$k_a/k'$ is purely inseparable
by Fields, Lemma
\ref{fields-lemma-separable-first}.
If $a\in K$ is separably algebraic
over $L$, it is algebraic over
$k$, hence belongs to $k_a$.
Thus it is also purely inseparable
over $L$: some original $p$-power
lies in $k'\subseteq L$ in positive
characteristic, and in characteristic
zero $k_a=k'$. An element both
separable and purely inseparable
has degree one, since its minimal
polynomial is separable and divides
a polynomial with only one distinct
root. Therefore $a\in L$, and the
criterion proves sufficiency.
The restriction $L/k$ algebraic
is part of this characterization;
it is not removed.
\end{proof}


\subsection{Editorial treatment of Algebra lines 11445--11479}

\hypertarget{algebra-context-148}{}

\noindent\href{https://github.com/stacks/stacks-project/blob/a04446e57ec1fbc252a871afcec7752fb2807b14/algebra.tex\#L11445}{Original source, lines 11445--11479}. Complete original and reviewed TeX passages accompany this note. Every intervention is separately recorded; the following treatment is editorial.


\begin{lemma}
\label{editorial-source-lemma-Galois-orbit}
Let $K/k$ be an extension of fields.
Let $\overline{k}/k$ be a separable algebraic closure.
Then $\text{Gal}(\overline{k}/k)$ acts transitively on the
primes of $\overline{k} \otimes_k K$.
\end{lemma}

\begin{proof}
Retain the original separable closure
$\overline k$, put $G=\operatorname{Gal}(\overline k/k)$,
and let $K/k'/k$ be the original
subextension of Lemma
\ref{editorial-source-lemma-make-geometrically-irreducible}.
Write
$$
B=\overline k\otimes_k k',\qquad
A=\overline k\otimes_k K.
$$
Both tensors are nonzero by extension
of vector-space bases over $k$.
As $\overline k/k$ is algebraic,
these are integral algebras over
the respective fields $k'$ and
$K$: each generator $a\otimes1$
satisfies its original monic
polynomial over $k$. Thus all
their primes are maximal by
Lemma \ref{editorial-source-lemma-integral-no-inclusion},
and are also minimal since
there are no strict inclusions
between primes over a field.
The two original comparison maps
$$
B\otimes_{k'}K\longrightarrow A,
\qquad(a\otimes b)\otimes c\longmapsto a\otimes bc,
$$
and $a\otimes c\mapsto(a\otimes1)\otimes c$
are inverse by $k'$-balancing.
Since $K/k'$ is geometrically
irreducible, Lemma
\ref{editorial-source-lemma-geometrically-irreducible-any-base-change}
therefore identifies the spectra
of $A$ and $B$, including their
topologies by its minimal-prime
homeomorphism. The map is
contraction, and commutes with
the original automorphisms
$g\otimes1$ for $g\in G$.

For a prime $\mathfrak p$ of $B$,
its residue field is separably
algebraic over the original
$\overline k$. Indeed $B$ is
generated over $\overline k$ by
$1\otimes b$ for $b\in k'$,
and each of their images satisfies
the original separable polynomial
of $b$ over $k$. Minimal
polynomials over $\overline k$
divide these separable polynomials.
Every element of the field
$B/\mathfrak p$ belongs to a
finite subextension generated
by finitely many such images,
so that extension is separable.
As $\overline k$ is separably
closed, its structural embedding
into $B/\mathfrak p$ is an
isomorphism. Its inverse composed
with $b\mapsto[1\otimes b]$
defines a unique embedding
$\sigma:k'\to\overline k$ over $k$.
Consequently
$$
\mathfrak p=\mathfrak p_\sigma
=\ker(\mu_\sigma),\qquad
\mu_\sigma:B\to\overline k,\quad
\sum_i a_i\otimes b_i\longmapsto\sum_i a_i\sigma(b_i).
$$
Conversely this map is surjective
because $a=\mu_\sigma(a\otimes1)$,
so its kernel is a prime. The
structural field isomorphism
above proves uniqueness of
$\sigma$ from its kernel.
This gives the stated bijection
with $\operatorname{Hom}_k(k',\overline k)$.
Separability, rather than
maximality alone, is what
identifies these residue fields
with the separably closed field.

Given embeddings $\sigma,\tau$,
extend the isomorphism
$\tau\sigma^{-1}:\sigma(k')\to\tau(k')$
to a $k$-embedding of $\overline k$
in itself. Here are the needed
extension details, including
infinite $k'/k$. Order partial
extensions by inclusion and
take unions on chains. If a
maximal partial extension has
domain $L\neq\overline k$, choose
$a\in\overline k\setminus L$.
Its minimal polynomial over
$L$ divides a separable
polynomial over $k$, hence is
separable. Transporting its
coefficients gives a separable
polynomial over the image field,
which has a root in the
separably closed $\overline k$.
Sending $a$ to that root extends
the embedding to $L(a)$, a
contradiction. Thus the domain
is all of $\overline k$.
Any $k$-embedding of this field
in itself is surjective: each
$b\in\overline k$ lies in a
finite Galois subextension $N/k$,
obtained by adjoining all roots
of its separable minimal polynomial.
The embedding permutes those
roots, sends $N$ into itself,
and its injective $k$-linear
map on the finite-dimensional
$N$ is surjective. In particular
$b$ is in its image. We have
therefore constructed $g\in G$
with $g\sigma=\tau$.

Use the action on ideals
$\mathfrak p\mapsto(g\otimes1)(\mathfrak p)$.
For every original $z\in B$,
$$
\mu_{g\sigma}((g\otimes1)z)=g(\mu_\sigma(z)).
$$
It follows that
$(g\otimes1)(\mathfrak p_\sigma)=\mathfrak p_{g\sigma}$.
This proves transitivity on
$\Spec(B)$ and, by equivariant
contraction, on the original
$\Spec(A)$.

The correspondence also identifies
the actual prime and residue
field of $A$. Give $\overline k$
its $k'$-structure through
$\sigma$. The extension
$\overline k/\sigma(k')$ is
separable algebraic. The tensor-field
form of Lemma
\ref{editorial-source-lemma-geometrically-irreducible-separable-elements}
applied to $K/k'$ therefore makes
$$
C_\sigma=\overline k\otimes_{k',\sigma}K
$$
a field. The original quotient
comparison sends the class of
$a\otimes x$ to $a\otimes x$
and gives an isomorphism
$A/\mathfrak p_\sigma A\simeq C_\sigma$;
its inverse sends the same pure
tensor to its class. The additional
$k'$-balancing relations are
exactly the extended kernel
of $\mu_\sigma$, so both maps
are well-defined. Hence the
corresponding prime is exactly
$\mathfrak q_\sigma=\mathfrak p_\sigma A$,
and its residue field is
$C_\sigma$. No further radical
or residue-field identification
is needed or assumed.

For a fixed $\sigma$, its
stabilizer in $G$ is exactly
the subgroup fixing $\sigma(k')$
pointwise, by the injectivity
of the embedding-to-prime
correspondence. Thus the orbit
is the corresponding coset
set. If $K/k$ is finitely
generated, the preceding lemma
makes $k'/k$ finite separable;
Fields, Lemma
\ref{fields-lemma-separable-equality}
then gives exactly $[k':k]$
primes of $A$. Since this
finite spectrum has every
point closed, each point is
also open: its complement is
a finite union of closed
points. It is therefore a
finite discrete orbit. No
finite-degree count is asserted
for an infinite $k'/k$.
\end{proof}


\subsection{Editorial treatment of Algebra lines 11487--11528}

\hypertarget{algebra-context-149}{}

\noindent\href{https://github.com/stacks/stacks-project/blob/a04446e57ec1fbc252a871afcec7752fb2807b14/algebra.tex\#L11487}{Original source, lines 11487--11528}. Complete original and reviewed TeX passages accompany this note. Every intervention is separately recorded; the following treatment is editorial.


\begin{lemma}
\label{editorial-source-lemma-separably-closed-connected}
Let $k$ be a separably closed field.
Let $R$, $S$ be $k$-algebras. If $\Spec(R)$ and
$\Spec(S)$ are connected, then so is
$\Spec(R \otimes_k S)$.
\end{lemma}

\begin{proof}
The source convention for connectedness
includes nonemptiness (Topology,
Definition \ref{topology-definition-connected-components}).
Thus the original $R$ and $S$ are
nonzero, and Lemma
\ref{algebra-lemma-characterize-spec-connected}
applies to them. Their tensor is
also nonzero, since extending $1$
in each factor to a $k$-basis
shows $1\otimes1\neq0$.

Suppose this tensor had a nontrivial
idempotent $e$. A finite tensor
expression for this original $e$
uses coefficients in finite type
$k$-subalgebras $R_0\subseteq R$
and $S_0\subseteq S$. These are
nonzero and have no nontrivial
idempotents, since their inclusions
in the original factors preserve
idempotents and reflect $0$ and $1$.
They therefore have connected spectra.
The map $R_0\otimes_k S_0\to R\otimes_k S$
is injective by applying the two
vector-space injections successively.
It follows that the same $e$ is
a nontrivial idempotent in this
finite type tensor: its equation
$e^2-e=0$ is reflected by that
injection. Thus it suffices to
prove connectedness for the actual
finite type subalgebras, with
the original $e$ and all its
coefficients retained.

For the finite type assertion,
use $R,S$ for the two algebras
under consideration. They are
Noetherian and have finite
nonzero numbers $n,m$ of
irreducible components by Lemma
\ref{editorial-source-lemma-Noetherian-irreducible-components}.
We induct on $n+m$. For $n=m=1$,
Lemma \ref{editorial-source-lemma-separably-closed-irreducible}
gives irreducibility, hence
connectedness, of the tensor.
If $n=1$ and $m>1$, interchange
the factors through the inverse
maps $r\otimes s\mapsto s\otimes r$
and $s\otimes r\mapsto r\otimes s$.
It remains to handle $n>1$.

Here are the component-selection
details of Topology, Lemma
\ref{topology-lemma-remove-irreducible-connected}.
Form the finite graph whose
vertices are the original
irreducible components and whose
edges record nonempty intersection.
It is connected: otherwise the
unions over two parts with no
edges between them would be
disjoint nonempty closed sets
covering the spectrum, a separation.
Choose a spanning tree, obtained
by adding adjacent new vertices
until every vertex is reached.
A leaf exists by taking an
endpoint of a longest simple
path in this finite tree.
Remove that leaf. The remaining
tree is connected, so the union
of its connected vertex spaces
is connected: add them in tree
order, each meeting the existing
union. This proves the required
choice without changing any
irreducible component.

Let $T=V(\mathfrak p)$ be the
chosen component and let $T'$
be the union of the other
$n-1$ components. Take the
original ideal $I$ to be the
intersection of their minimal
primes. Then $T'=V(I)$, by
finite union of closed sets,
and its irreducible components
are those same $n-1$ components:
their defining primes are
pairwise incomparable. Both
$T$ and $T'$ are connected
and nonempty. They meet,
since otherwise these two
closed sets would separate
the original connected spectrum.
Their intersection is exactly
$V(\mathfrak p+I)$, so
$R/(\mathfrak p+I)\neq0$.

Under the original spectrum map,
the inverse images of $T$ and
$T'$ are respectively
$$
\Spec((R/\mathfrak p)\otimes_k S),\qquad
\Spec((R/I)\otimes_k S).
$$
For each of the ideals $J$
involved, the quotient comparison
sends $[r\otimes s]$ to
$[r]\otimes s$ and has inverse
$[r]\otimes s\mapsto[r\otimes s]$,
identifying $(R\otimes_k S)/J(R\otimes_k S)$
with $(R/J)\otimes_k S$.
The induction parameters for
these two inverse images are
$1+m$ and $(n-1)+m$, both
smaller than $n+m$; hence
both are connected. Their
intersection is the inverse
image of $T\cap T'$, namely
$$
\Spec((R/(\mathfrak p+I))\otimes_k S).
$$
It is nonempty because both
original tensor factors are
nonzero vector spaces and
their unit tensor is nonzero.
The two inverse images cover
the spectrum of $R\otimes_k S$.
A union of two connected sets
with nonempty intersection is
connected: a separation would
put each in one part, and
the common point forces the
same part for both.
This finishes the induction
and excludes the original
nontrivial idempotent $e$.
\end{proof}


\subsection{Editorial treatment of Algebra lines 11530--11558}

\hypertarget{algebra-context-150}{}

\noindent\href{https://github.com/stacks/stacks-project/blob/a04446e57ec1fbc252a871afcec7752fb2807b14/algebra.tex\#L11530}{Original source, lines 11530--11558}. Complete original and reviewed TeX passages accompany this note. Every intervention is separately recorded; the following treatment is editorial.


\begin{lemma}
\label{editorial-source-lemma-geometrically-connected}
Let $k$ be a field.
Let $R$ be a $k$-algebra.
The following are equivalent
\begin{enumerate}
\item for every field extension $k'/k$ the
spectrum of $R \otimes_k k'$ is connected, and
\item for every finite separable field extension $k'/k$ the
spectrum of $R \otimes_k k'$ is connected.
\end{enumerate}
\end{lemma}

\begin{proof}
Condition (1) implies (2).
Assume (2); its identity field
extension says that $\Spec(R)$
is connected and in particular
nonempty, so $R\neq0$.
Every scalar extension considered
below is nonzero by extension
of a vector-space basis.
For these nonzero rings, Lemma
\ref{algebra-lemma-characterize-spec-connected}
identifies connectedness with
absence of nontrivial idempotents.
The nonzero condition cannot be
omitted from that criterion:
the zero ring has no nontrivial
idempotents but has empty,
nonconnected spectrum under
Topology, Definition
\ref{topology-definition-connected-components}.

Let $k_s/k$ be the original
separable algebraic closure.
Every finite list of its
elements lies in a finite
separable subextension $E/k$.
The original coefficient maps
give the injective directed union
$$
R\otimes_k k_s=\bigcup_E R\otimes_k E.
$$
An idempotent in this union
lies in one such tensor, and
its equation is already true
there by injectivity. By (2)
it equals $0$ or $1$ there,
hence in the union. This
proves connectedness over $k_s$.

For any original field $L/k$,
choose a prime $\mathfrak a$
of the nonzero ring $L\otimes_k k_s$
and put
$$
F=\operatorname{Frac}((L\otimes_k k_s)/\mathfrak a).
$$
Its original field embeddings
$a\mapsto[a\otimes1]$ and
$b\mapsto[1\otimes b]$ are
injective and agree on $k$.
The scalar comparison maps
$$
(R\otimes_k k_s)\otimes_{k_s}F\longrightarrow R\otimes_k F,
\qquad(r\otimes a)\otimes b\longmapsto r\otimes ab,
$$
and $r\otimes b\mapsto(r\otimes1)\otimes b$
are inverse by balancing.
The source has connected
spectrum by Lemma
\ref{editorial-source-lemma-separably-closed-connected},
since both $R\otimes_k k_s$
and the field $F$ do.
The injection
$R\otimes_k L\to(R\otimes_k L)\otimes_L F\simeq R\otimes_k F$
reflects nontrivial idempotents,
again by a vector-space basis
and the displayed scalar maps.
Thus $R\otimes_k L$ is nonzero
and has no nontrivial idempotents.
It is connected, proving (1).

The same argument proves that
connectedness after just the
separable closure is an equivalent
test. One may instead use the
algebraic closure $\overline k$:
the extension $\overline k/k_s$
is purely inseparable, or the
identity in characteristic zero.
The original injection
$R\otimes_k k_s\to R\otimes_k\overline k$
is a homeomorphism on spectra
by Lemma \ref{editorial-source-lemma-p-ring-map}
and the same scalar comparison,
so it preserves and reflects
connectedness. Also finite Galois
extensions suffice in (2): every
finite separable $E/k$ embeds
in a finite Galois $N/k$, and
the injection $R\otimes_k E\to R\otimes_k N$
reflects nontrivial idempotents
and nonzeroness.

Finally, for every original $L/k$,
geometric connectedness of $R/k$
is equivalent to geometric
connectedness of $(R\otimes_k L)/L$.
The forward direction uses
the displayed scalar comparison
for any $F/L$. For the reverse
direction and any $E/k$, construct
a common receiving field for
$E$ and $L$ by the same
prime-quotient construction.
Geometric connectedness over
$L$ gives connectedness of
$R\otimes_k F$, and its
injected subalgebra $R\otimes_k E$
is nonzero and has no nontrivial
idempotents. This proves the
reverse implication with the
original scalar actions.
\end{proof}


\subsection{Editorial treatment of Algebra lines 11573--11585}

\hypertarget{algebra-context-151}{}

\noindent\href{https://github.com/stacks/stacks-project/blob/a04446e57ec1fbc252a871afcec7752fb2807b14/algebra.tex\#L11573}{Original source, lines 11573--11585}. Complete original and reviewed TeX passages accompany this note. Every intervention is separately recorded; the following treatment is editorial.


\begin{lemma}
\label{editorial-source-lemma-separably-closed-connected-implies-geometric}
Let $k$ be a field.
Let $R$ be a $k$-algebra.
If $k$ is separably closed then $R$ is
geometrically connected over $k$ if and only if the
spectrum of $R$ is connected.
\end{lemma}

\begin{proof}
Geometric connectedness implies
connectedness by the identity field
extension in Definition
\ref{algebra-definition-geometrically-connected}.
Conversely, if $\Spec(R)$ is
connected, then $R\neq0$ under
the source nonemptiness convention.
For any finite separable extension
$E/k$, separably closedness of $k$
makes the given embedding $\iota:k\to E$
an isomorphism: each element has
a separable minimal polynomial,
which must be linear over $k$.
The inverse maps
$R\otimes_k E\to R$, $r\otimes a\mapsto r\iota^{-1}(a)$,
and $r\mapsto r\otimes1$ show
that all the finite separable
tests have connected spectrum.
Lemma \ref{editorial-source-lemma-geometrically-connected}
proves geometric connectedness.
For $R=0$, neither side is
connected, so the same equivalence
holds without changing the definition.
\end{proof}


\subsection{Editorial treatment of Algebra lines 11587--11604}

\hypertarget{algebra-context-152}{}

\noindent\href{https://github.com/stacks/stacks-project/blob/a04446e57ec1fbc252a871afcec7752fb2807b14/algebra.tex\#L11587}{Original source, lines 11587--11604}. Complete original and reviewed TeX passages accompany this note. Every intervention is separately recorded; the following treatment is editorial.


\begin{lemma}
\label{editorial-source-lemma-subalgebra-geometrically-connected}
Let $k$ be a field. Let $S$ be a $k$-algebra.
\begin{enumerate}
\item If $S$ is geometrically connected over $k$ so is every
$k$-subalgebra.
\item If all finitely generated $k$-subalgebras of $S$ are
geometrically connected, then $S$ is geometrically connected.
\item A directed colimit of geometrically connected $k$-algebras
is geometrically connected.
\end{enumerate}
\end{lemma}

\begin{proof}
Use throughout the nonzero-ring
criterion of Lemma
\ref{algebra-lemma-characterize-spec-connected}.
For (1), the original inclusion
$S'\subseteq S$ gives an injection
$S'\otimes_k L\to S\otimes_k L$
for every field $L/k$ by extension
of a vector-space basis. Its
target is nonzero, and the
unital source is therefore
nonzero as well. Any nontrivial
idempotent of the source would
remain nontrivial in the target,
which is impossible. This proves (1).

For (3), retain the original
nonempty directed system $(S_i)$
and its unital transition maps;
no injectivity is required.
For every $L/k$, the maps
$$
(\mathop{\rm colim}_i S_i)\otimes_k L
\longrightarrow\mathop{\rm colim}_i(S_i\otimes_k L),
\qquad[i,s]\otimes a\longmapsto[i,s\otimes a]
$$
and $[i,\sum_j s_j\otimes a_j]\mapsto\sum_j[i,s_j]\otimes a_j$
are well-defined inverses:
all representatives and balancing
equations involve finitely many
terms and can be compared at
one common later stage.
Each original $S_i\otimes_k L$
is nonzero. The colimit is
nonzero, since $1=0$ there
would hold at a later stage.
If $e$ is an idempotent in
the colimit, represent it at
a stage $i$. Its equation
$e^2-e=0$ holds at some later
stage $j$, making that later
representative idempotent.
It is $0$ or $1$ there by
connectedness, and hence is
$0$ or $1$ in the colimit.
This proves connectedness after
each $L/k$, and gives (3).

For (2), the finite type
$k$-subalgebras form a directed
system under adjoining the union
of two finite lists of generators,
and their union is the original
$S$. Apply (3). The hypotheses
exclude $S=0$, since its
subalgebra generated by the
image of $k$ would already
be zero and would fail
geometric connectedness.

Unlike geometric irreducibility,
geometric connectedness need
not survive a nonzero localization.
For the exact example
$$
S=k[x,y]/(xy),\qquad s=x+y,
$$
each field extension gives
two affine lines meeting at
their origin. More explicitly,
$V(x)$ and $V(y)$ cover its
spectrum, are irreducible
because their quotients are
$L[y]$ and $L[x]$, and meet
in the nonempty $\Spec(L)$.
Their union is therefore
connected. But in the original
$S_s$ the element
$$
e=\frac{x}{x+y},\qquad
e^2-e=\frac{-xy}{(x+y)^2}=0
$$
is a nontrivial idempotent.
The map
$$
S_s\longrightarrow k[x,x^{-1}]\times k[y,y^{-1}],
\qquad\frac{h(x,y)}{(x+y)^n}
\longmapsto\left(\frac{h(x,0)}{x^n},\frac{h(0,y)}{y^n}\right)
$$
sends it to $(1,0)$, so it
is neither $0$ nor $1$.
This map is in fact an
isomorphism with the following
explicit inverse. In $eS_s$,
whose unit is $e$, send $x$
to the original $x=ex$ and
$x^{-1}$ to $e/s$; their
product is $e$ because
$x(e/s)=x^2/s^2=e$.
In $(1-e)S_s$, send $y$ to
$y=(1-e)y$ and $y^{-1}$ to
$(1-e)/s$; their product is
$1-e$. Add the images of
the two Laurent polynomials.
Orthogonality of $e,1-e$
makes this a ring map.
The composites fix both
component units and Laurent
generators. On the original
side they fix $x,y$ and
$s^{-1}$, since
$e/s+(1-e)/s=s^{-1}$.
Thus both composites are
identities, proving the full
product description and the
failure of localization stability.
\end{proof}


\subsection{Editorial treatment of Algebra lines 11610--11649}

\hypertarget{algebra-context-153}{}

\noindent\href{https://github.com/stacks/stacks-project/blob/a04446e57ec1fbc252a871afcec7752fb2807b14/algebra.tex\#L11610}{Original source, lines 11610--11649}. Complete original and reviewed TeX passages accompany this note. Every intervention is separately recorded; the following treatment is editorial.


\begin{lemma}
\label{editorial-source-lemma-geometrically-connected-any-base-change}
Let $k$ be a field.
Let $S$ be a geometrically connected $k$-algebra.
Let $R$ be any $k$-algebra.
The map
$$
R \longrightarrow R \otimes_k S
$$
induces a bijection on idempotents, and the map
$$
\Spec(R \otimes_k S) \longrightarrow \Spec(R)
$$
induces a bijection on connected components.
\end{lemma}

\begin{proof}
Put $A=R\otimes_k S$ and retain
$f:X=\Spec(A)\to Y=\Spec(R)$.
Geometric connectedness of the
original $S$ includes nonemptiness,
so $S\neq0$. If $R=0$, both
spectra are empty and both
rings have one idempotent,
so the two asserted bijections
hold. Assume $R\neq0$.
For every $\mathfrak p\in Y$,
the original fibre comparison
$$
A\otimes_R\kappa(\mathfrak p)\longrightarrow
\kappa(\mathfrak p)\otimes_k S,
\qquad(r\otimes s)\otimes a\longmapsto ra\otimes s
$$
has inverse $a\otimes s\mapsto(1\otimes s)\otimes a$.
The receiving spectrum is
connected by the hypothesis
on $S$ and is nonempty.
Hence $f$ is surjective with
connected fibres. It is open
without finite type assumptions,
by Lemma
\ref{editorial-source-lemma-map-into-tensor-algebra-open}.

Let $C\subseteq X$ be open
and closed. Connectedness
of each fibre implies that
the fibre is contained either
in $C$ or in its complement:
otherwise their intersections
would separate the fibre.
The images $f(C)$ and $f(X\setminus C)$
are therefore disjoint open
sets covering $Y$. In
particular $f(C)$ is open
and closed, and
$$
C=f^{-1}(f(C)).
$$
Conversely, inverse images of
open and closed subsets are
open and closed, and surjectivity
gives $f(f^{-1}(D))=D$ for
every such $D\subseteq Y$.
This proves a bijection of
the original open-and-closed
subsets in both directions.

By Lemma \ref{algebra-lemma-disjoint-decomposition},
the map $e\mapsto D(e)$ is
a bijection between idempotents
of each ring and its open-and-closed
subsets. Under the original
ring map $R\to A$, the
identity $f^{-1}(D(e))=D(e\otimes1)$
therefore proves exactly the
claimed idempotent bijection.
The inverse takes an idempotent
$a\in A$ to the unique
idempotent $e\in R$ such
that $D(e)=f(D(a))$.
These maps also preserve
complements and intersections,
represented by $1-e$ and
$ee'$, so give an isomorphism
of the corresponding Boolean
algebras. They commute with
arbitrary original $k$-algebra
base change: the image of
$e\otimes1$ is $\theta(e)\otimes1$,
and uniqueness of descent
gives the inverse compatibility.

We prove the component assertion
directly as well. For any
connected component $D$ of
$Y$, the restriction
$f^{-1}(D)\to D$ is open,
surjective and has the same
connected fibres. Openness
holds because for every open
$U\subseteq X$,
$$
f(U\cap f^{-1}(D))=f(U)\cap D.
$$
If $f^{-1}(D)$ were separated
into disjoint nonempty open
and closed parts, the same
fibre argument would give
a separation of $D$ by
their open images. Thus
$f^{-1}(D)$ is connected
and nonempty. It is maximal
connected: any larger connected
subset would have connected
image contained in a component
of $Y$, which must be $D$
because it meets $D$.
Conversely each component
of $X$ has connected image
in some component $D$ of
$Y$ and hence equals this
connected inverse image.
This proves the claimed
bijection and the exact
inverse-image description,
as in Topology, Lemma
\ref{topology-lemma-connected-fibres-connected-components}.

With the quotient topologies
on the sets of components,
the bijection is a homeomorphism.
Write $q_X,q_Y$ for the two
quotient maps and $h$ for
the induced bijection.
The relation $h q_X=q_Y f$
proves continuity by the
quotient property. If $V$
is open in the component
space of $X$, then
$$
q_Y^{-1}(h(V))=f(q_X^{-1}(V))
$$
by the exact component
correspondence. The right
side is open because $f$
is open, proving that $h$
is open and hence a
homeomorphism. None of this
requires individual components
to be open in their spectra.
\end{proof}


\subsection{Editorial treatment of Algebra lines 11673--11683}

\hypertarget{algebra-context-154}{}

\noindent\href{https://github.com/stacks/stacks-project/blob/a04446e57ec1fbc252a871afcec7752fb2807b14/algebra.tex\#L11673}{Original source, lines 11673--11683}. Complete original and reviewed TeX passages accompany this note. Every intervention is separately recorded; the following treatment is editorial.


\begin{lemma}
\label{editorial-source-lemma-geometrically-integral}
Let $k$ be a field.
Let $S$ be a $k$-algebra.
In this case $S$ is geometrically integral over $k$ if and only if
$S$ is geometrically irreducible as well as geometrically reduced over $k$.
\end{lemma}

\begin{proof}
For any original field extension $L/k$,
put $A=S\otimes_k L$. If $A$ is
a domain, it is nonzero and reduced,
and its zero ideal is prime. Its
spectrum is therefore irreducible.
Conversely, if $A$ is reduced and
its spectrum is irreducible, then
it is nonzero and its nilradical
is prime by Lemma \ref{editorial-source-lemma-irreducible}.
Reducedness makes that nilradical
the zero ideal, so the zero ideal
is prime and $A$ is a domain.
Apply this exact ring equivalence
for every original $L/k$ to obtain
both directions of the assertion.
Nonemptiness comes from irreducibility;
reducedness alone would allow the
zero ring and would not suffice.
\end{proof}


\subsection{Editorial treatment of Algebra lines 11685--11702}

\hypertarget{algebra-context-155}{}

\noindent\href{https://github.com/stacks/stacks-project/blob/a04446e57ec1fbc252a871afcec7752fb2807b14/algebra.tex\#L11685}{Original source, lines 11685--11702}. Complete original and reviewed TeX passages accompany this note. Every intervention is separately recorded; the following treatment is editorial.


\begin{lemma}
\label{editorial-source-lemma-characterize-geometrically-integral}
Let $k$ be a field. Let $S$ be a $k$-algebra.
The following are equivalent
\begin{enumerate}
\item $S$ is geometrically integral over $k$,
\item for every finite extension $k'/k$ of fields
the ring $S \otimes_k k'$ is a domain,
\item $S \otimes_k \overline{k}$ is a domain
where $\overline{k}$ is the algebraic closure of $k$.
\end{enumerate}
\end{lemma}

\begin{proof}
Condition (1) implies (2) by
restricting the field extensions.
Assume (2). In particular its
identity extension makes $S$
a nonzero domain. Choose the
original algebraic closure
$\overline k$ and consider its
finite subextensions $E/k$.
They are directed under compositum
and contain every finite list
of coefficients of a tensor.
The original injective coefficient
maps therefore give
$$
S\otimes_k\overline k=\bigcup_E S\otimes_k E.
$$
Each ring in this union is a
domain by (2). The union is
nonzero, since the maps preserve
the nonzero unit. For two
nonzero elements, choose one
stage containing both. Their
product there is nonzero and
remains so by injectivity.
This proves (3), including
nonemptiness and all zero-divisor
conditions.

Assume (3), and let $L/k$
be any original extension.
Choose a prime $\mathfrak a$
of the nonzero field tensor
$L\otimes_k\overline k$ and
form the compatible receiving
field
$$
F=\operatorname{Frac}((L\otimes_k\overline k)/\mathfrak a).
$$
The maps $a\mapsto[a\otimes1]$
and $b\mapsto[1\otimes b]$
embed the two original fields
and agree on $k$. Write
$B=S\otimes_k\overline k$.
This is a nonzero domain
by (3). Since $\overline k$
is perfect, Lemma
\ref{editorial-source-lemma-perfect-reduced}
makes $B\otimes_{\overline k}F$
reduced. Since $\overline k$
is separably closed and both
$B$ and $F$ have irreducible
spectra, Lemma
\ref{editorial-source-lemma-separably-closed-irreducible}
makes its spectrum irreducible.
Thus it is a domain by Lemma
\ref{editorial-source-lemma-geometrically-integral}'s
underlying ring criterion.

For either original subfield
$H=L$ or $H=\overline k$ of
the receiving $F$, the maps
$$
(S\otimes_k H)\otimes_H F\longrightarrow S\otimes_k F,
\qquad(s\otimes a)\otimes b\longmapsto s\otimes ab
$$
and $s\otimes b\mapsto(s\otimes1)\otimes b$
are inverse by balancing.
They identify the domain just
proved with $S\otimes_k F$.
The original scalar map
$S\otimes_k L\to(S\otimes_k L)\otimes_L F$
is injective by extension of
an $L$-vector-space basis.
Its source is therefore a
nonzero subring of that domain,
so is a domain. This proves
(1) for every original $L/k$.

The same receiving-field maps
prove geometric integrality
is preserved and reflected
by every field extension
$L/k$. Preservation follows
from the displayed comparison
for all $F/L$. For reflection,
suppose $S\otimes_k L$ is
geometrically integral over $L$.
For any $E/k$, construct a
compatible common field for
$E$ and $L$ as above. The
hypothesis makes $S\otimes_k F$
a domain, and the injection
$S\otimes_k E\to S\otimes_k F$
then makes the original
$S\otimes_k E$ a domain.
This is precisely geometric
integrality of $S/k$.
The zero algebra satisfies
none of (1)--(3), as each
condition includes a nonzero
domain after scalar extension.
\end{proof}


\subsection{Editorial treatment of Algebra lines 11704--11718}

\hypertarget{algebra-context-156}{}

\noindent\href{https://github.com/stacks/stacks-project/blob/a04446e57ec1fbc252a871afcec7752fb2807b14/algebra.tex\#L11704}{Original source, lines 11704--11718}. Complete original and reviewed TeX passages accompany this note. Every intervention is separately recorded; the following treatment is editorial.


\begin{lemma}
\label{editorial-source-lemma-geometrically-integral-any-integral-base-change}
Let $k$ be a field. Let $S$ be a geometrically integral $k$-algebra.
Let $R$ be a $k$-algebra and an integral domain. Then $R \otimes_k S$
is an integral domain.
\end{lemma}

\begin{proof}
Let $F=\operatorname{Frac}(R)$.
The original injection $R\to F$
gives an injection
$$
R\otimes_k S\longrightarrow F\otimes_k S
$$
because tensoring $k$-vector spaces
preserves injections. The receiving
ring is a domain by geometric
integrality of $S/k$. The source
is nonzero: extending $1$ to
bases of the two original
nonzero factors shows
$1\otimes1\neq0$. Hence it
is a domain, as asserted.
The corresponding localization
maps retain the original
denominators $W=R\setminus\{0\}$:
$$
F\otimes_k S\longrightarrow W^{-1}(R\otimes_k S),
\qquad(r/w)\otimes s\longmapsto(r\otimes s)/(w\otimes1),
$$
and the inverse sends
$(\sum_i r_i\otimes s_i)/(w\otimes1)$
to $\sum_i(r_i/w)\otimes s_i$.
The fraction equivalence relations
and tensor balancing make these
well-defined inverses. The
displayed injection is exactly
the original constant-fraction
map under this comparison.

The argument gives precise
ideal consequences for any
original $k$-algebra $R$, now
without a domain hypothesis.
Write $A=R\otimes_k S$ and
let $I\subseteq R$ be an
arbitrary ideal. The maps
$[r\otimes s]\mapsto[r]\otimes s$
and $[r]\otimes s\mapsto[r\otimes s]$
identify
$$
A/IA\simeq(R/I)\otimes_k S.
$$
Choose a $k$-linear functional
$\ell:S\to k$ with $\ell(1)=1$.
The map $\operatorname{id}\otimes\ell$
splits $R/I\to(R/I)\otimes_k S$
as a linear map, proving
$IA\cap R=I$ for every $I$.
If $I$ is prime, $R/I$ is
a domain and the argument
above makes $A/IA$ a domain,
so $IA$ is prime. Conversely,
if $IA$ is prime, its
contraction $I$ is prime.
This proves the exact equivalence
for the original extended ideal,
including the exclusion of
the unit ideal from prime ideals.

There is also the full radical
comparison
$$
\sqrt{IA}=(\sqrt I)A.
$$
To prove it, keep the quotient
map
$$
(R/I)\otimes_k S\longrightarrow(R/\sqrt I)\otimes_k S.
$$
Its kernel is the ideal
generated by the original
$\sqrt I/I$ factor, by the
same explicit quotient maps.
Each element of this kernel
is a finite sum of tensors
$r_j\otimes s_j$ with $r_j$
nilpotent. If $r_j^{d_j}=0$,
its power of degree
$N=1+\sum_j(d_j-1)$ vanishes:
in the full multinomial sum
$$
\sum_{i_1+\cdots+i_m=N}
\frac{N!}{i_1!\cdots i_m!}
\prod_{j=1}^m(r_j\otimes s_j)^{i_j},
$$
some $i_j\geq d_j$ in every
summand. The displayed coefficient
is an integer; no factorial
is inverted in the receiving ring.
The target is reduced by Lemma
\ref{editorial-source-lemma-geometrically-reduced-any-reduced-base-change},
since $S$ is geometrically reduced
and $R/\sqrt I$ is reduced.
Thus this kernel is exactly
the nilradical of $A/IA$.
Taking its inverse image in
$A$ proves the radical formula.

In particular for every prime
$\mathfrak p$ of $R$ the ideal
$\mathfrak pA$ is already prime,
so the continuous section of
Lemma
\ref{editorial-source-lemma-geometrically-irreducible-any-base-change}
has the exact value
$\sigma(\mathfrak p)=\mathfrak pA$.
Its contraction, topology,
base-change compatibility and
minimal-prime correspondence
are those already proved for
the original tensor map;
the stronger integrality
hypothesis removes no terms
from the original ideal.
\end{proof}


\subsection{Editorial treatment of Algebra lines 11723--12149}

\hypertarget{algebra-context-157}{}

\noindent\href{https://github.com/stacks/stacks-project/blob/a04446e57ec1fbc252a871afcec7752fb2807b14/algebra.tex\#L11723}{Original source, lines 11723--12149}. Complete original and reviewed TeX passages accompany this note. Every intervention is separately recorded; the following treatment is editorial.


\subsubsection{Valuation rings}
\label{editorial-source-section-valuation-rings}

\noindent
Here are some definitions.

\begin{definition}
\label{editorial-source-definition-valuation-ring}
Valuation rings.
\begin{enumerate}
\item Let $K$ be a field. Let $A$, $B$ be local rings contained
in $K$. We say that $B$ {\it dominates} $A$ if $A \subset B$
and $\mathfrak m_A = A \cap \mathfrak m_B$.
\item Let $A$ be a ring. We say $A$ is a {\it valuation ring}
if $A$ is a local domain and if $A$ is maximal
for the relation of domination among local rings contained in
the fraction field of $A$.
\item Let $A$ be a valuation ring with fraction field $K$.
If $R \subset K$ is a subring of $K$, then we say $A$
is {\it centered} on $R$ if $R \subset A$.
\end{enumerate}
\end{definition}

\noindent
With this definition a field is a valuation ring.

\begin{lemma}
\label{editorial-source-lemma-dominate}
Let $K$ be a field. Let $A \subset K$ be a local subring.
Then there exists a valuation ring with fraction field $K$
dominating $A$.
\end{lemma}

\begin{proof}
Consider the set of local subrings
$B\subseteq K$ dominating the
original $A$, ordered by domination.
It is nonempty since it contains
$A$. The relation is transitive:
if $B$ dominates $A$ and $C$
dominates $B$, then
$A\cap\mathfrak m_C=A\cap\mathfrak m_B=\mathfrak m_A$.
Antisymmetry follows from the
inclusions of the actual subrings.

For a nonempty chain $(B_i)$,
put $B=\bigcup_i B_i$ and
$\mathfrak n=\bigcup_i\mathfrak m_{B_i}$.
The maximal ideals are compatible
under contraction and inclusion.
A common chain member containing
any finite list of elements proves
that $B$ is a ring and
$\mathfrak n$ is an ideal.
It is proper because no member
contains $1$. If $b\in B\setminus\mathfrak n$,
it is a unit in any member
containing it, hence in $B$.
If $b\in\mathfrak n$, an inverse
in $B$ would put both elements
in a common later member and
contradict properness of its
maximal ideal. Thus $B$ is
local with maximal ideal
$\mathfrak n$, and
$B_i\cap\mathfrak n=\mathfrak m_{B_i}$
for every $i$. It dominates
all members and the original
$A$. The empty chain has
upper bound $A$. Zorn's lemma
therefore gives a maximal
member $V$.

We prove that its actual
fraction field $F=\operatorname{Frac}(V)$
equals $K$. Otherwise choose
$t\in K\setminus F$.
If $t$ is transcendental over
$F$, the evaluation map embeds
the original polynomial ring
$V[T]$ in $K$, with $T\mapsto t$.
The ideal $(\mathfrak m_V,t)$
of $V[t]$ is maximal, since
its quotient is $V/\mathfrak m_V$.
Its localization embeds in
$K$ by the original fraction
map and dominates $V$:
the maximal ideal contracts
to $\mathfrak m_V$. It contains
$t\notin V$, contradicting
maximality of $V$.

If $t$ is algebraic over
$F$, retain its monic equation
$$
t^d+c_1t^{d-1}+\cdots+c_d=0,
\qquad c_i\in F.
$$
Choose a nonzero $a\in V$
that is a product of denominators
for all original $c_i$, and
write $b_i=ac_i\in V$.
The original element $u=at$
satisfies the full equation
$$
u^d+b_1u^{d-1}+a b_2u^{d-2}
+\cdots+a^{d-1}b_d=0.
$$
This is exactly $a^d$ times
the displayed original equation,
using $a^{i-1}b_i=a^ic_i$;
every coefficient is retained.
Hence $V[u]$ is finite integral
over $V$, generated as a module
by $1,u,\ldots,u^{d-1}$.
By Lemma
\ref{editorial-source-lemma-integral-overring-surjective},
there is a prime $\mathfrak q$
of $V[u]$ contracting to
$\mathfrak m_V$. The original
fraction map embeds $V[u]_{\mathfrak q}$
in $K$, with maximal-ideal
contraction $\mathfrak m_V$.
It dominates $V$ and contains
$u\notin V$, since $u\in V$
would imply $t=u/a\in F$.
This is again a contradiction.
Thus $F=K$.

Every local subring of $K$
dominating $V$ also dominates
$A$. Maximality in the chosen
set therefore makes $V$ maximal
for domination in its own
fraction field $K$, exactly
the definition of a valuation
ring. This proves the lemma.

In particular, for any original
subring $R\subseteq K$ and
prime $\mathfrak p\subset R$,
apply the result to the
actual local subring $R_{\mathfrak p}$
of $K$. The resulting $V$
has fraction field $K$ and
$$
\mathfrak m_V\cap R
=(\mathfrak m_V\cap R_{\mathfrak p})\cap R
=\mathfrak pR_{\mathfrak p}\cap R
=\mathfrak p.
$$
Thus the center can be
prescribed as any given
prime of the original subring.
\end{proof}

\begin{lemma}
\label{editorial-source-lemma-valuation-ring-normal}
Let $A$ be a valuation ring.
Then $A$ is a normal domain.
\end{lemma}

\begin{proof}
Let $K=\operatorname{Frac}(A)$, retaining
the original fraction embedding. Suppose
$x\in K$ is integral over $A$.
The original $A'=A[x]\subseteq K$
is finite integral over $A$: a
monic equation of degree $d$ for
$x$ expresses all its powers in
the span of $1,x,\ldots,x^{d-1}$.
Lemma \ref{editorial-source-lemma-integral-overring-surjective}
gives a prime $\mathfrak m'$ of
$A'$ contracting to $\mathfrak m_A$.
The actual fraction map
$A'_{\mathfrak m'}\to K$, $a/s\mapsto a/s$,
is injective because $A'$ is a
domain and $s\notin\mathfrak m'$
is nonzero. The maximal ideal
of this localization contracts to
$\mathfrak m'$ in $A'$ and then
to $\mathfrak m_A$ in $A$.
Thus it dominates the original
$A$. The valuation-ring maximality
in this same $K$ gives
$A=A'_{\mathfrak m'}$, so $x\in A$.
Every element of the fraction field
integral over $A$ therefore belongs
to $A$, proving normality.
In particular any integral overring
$A\subseteq D\subseteq K$ equals
$A$, by applying this proof to
each of its original elements.
\end{proof}


\begin{lemma}
\label{editorial-source-lemma-valuation-ring-x-or-x-inverse}
Let $A$ be a valuation ring with maximal ideal $\mathfrak m$ and
fraction field $K$.
Let $x \in K^*$. Then either $x \in A$ or $x^{-1} \in A$ or both.
\end{lemma}

\begin{proof}
It suffices to treat the case
$x\notin A$; then $x\neq0$,
since $0\in A$. Retain the
original subring $A'=A[x]\subseteq K$.
If a prime $\mathfrak q$ of
$A'$ contracted to $\mathfrak m$,
the original fraction map would
embed $A'_{\mathfrak q}$ in $K$
as a local ring dominating
$A$. Maximality of the valuation
ring would give $A=A'_{\mathfrak q}$,
contrary to $x\notin A$.
There is therefore no such
prime. Every prime containing
$\mathfrak m A'$ would contract
to the maximal ideal $\mathfrak m$,
so $\mathfrak m A'$ is the
unit ideal. A finite expression
for its unit, expanded in
the original generator $x$,
gives
$$
1=\sum_{i=0}^d t_i x^i,\qquad t_i\in\mathfrak m.
$$
Here $d\geq1$, since $d=0$
would imply $1=t_0\in\mathfrak m$.
With $y=x^{-1}$, multiplication
by the full $x^{-d}$ gives
$$
(1-t_0)y^d-\sum_{i=1}^d t_i y^{d-i}=0.
$$
The index starts at $1$:
the original $t_0$ contribution
is already in $1-t_0$.
That coefficient is a unit
of the local $A$, so the
monic polynomial
$$
Q(T)=T^d-\sum_{i=1}^d(1-t_0)^{-1}t_iT^{d-i}
$$
annihilates $y$. The complete
original polynomial is
$(1-t_0)Q(T)$; its unit
factor and all coefficients
are kept in the comparison.
Thus $y$ is integral over
$A$ and belongs to $A$ by
Lemma \ref{editorial-source-lemma-valuation-ring-normal}.
This proves the assertion.

The inverse is in fact in
$\mathfrak m$: if $y$ were
a unit of $A$, then
$x=y^{-1}$ would belong to
$A$, contrary to the case
under consideration. Conversely,
the inverse of a nonzero
element of $\mathfrak m$
cannot belong to $A$, since
their product would put $1$
in $\mathfrak m$. Therefore
$$
K\setminus A=\{a^{-1}:0\neq a\in\mathfrak m\},
$$
and a nonzero $x$ and its
inverse both belong to $A$
exactly when $x\in A^*$.
The zero element belongs to
$A$ independently and is
never assigned an inverse.
\end{proof}

\begin{lemma}
\label{editorial-source-lemma-x-or-x-inverse-valuation-ring}
Let $A \subset K$ be a subring of a field $K$ such that for all
$x \in K^*$ either $x \in A$ or $x^{-1} \in A$ or both.
Then $A$ is a valuation ring with fraction field $K$.
\end{lemma}

\begin{proof}
The original $A$ is a nonzero
domain, as a unital subring
of the field $K$. The given
alternative proves that its
fraction field is all of $K$:
for nonzero $x\in K$, either
$x=x/1$ with $x\in A$ or
$x=1/x^{-1}$ with $x^{-1}\in A$.
These expressions are the
inverse of the canonical
injection $\operatorname{Frac}(A)\to K$.
The zero element is $0/1$.

If $A=K$, it is a local
ring with zero maximal ideal.
Otherwise it is not a field,
since a subfield equals its
own fraction field. Choose
a maximal ideal $\mathfrak m$
of $A$. If another maximal
ideal $\mathfrak m'$ were
distinct, their incomparability
would give
$$
x\in\mathfrak m\setminus\mathfrak m',\qquad
y\in\mathfrak m'\setminus\mathfrak m.
$$
Both original elements are
nonzero. If $x/y\in A$,
then $x=y(x/y)\in\mathfrak m'$,
a contradiction. If $y/x\in A$,
then $y=x(y/x)\in\mathfrak m$,
also a contradiction. The
given alternative excludes
two maximal ideals. Thus
$A$ is local, including
the field case.

Let the original local subring
$A'\subseteq K$ dominate $A$.
If $x\in A'\setminus A$,
then $x\neq0$ and the
hypothesis gives $x^{-1}\in A$.
It is not a unit of $A$,
since otherwise $x\in A$.
Hence $x^{-1}\in\mathfrak m_A$
and domination puts it in
$\mathfrak m_{A'}$. Multiplying
by the original $x\in A'$
would put $1$ in that ideal.
This contradiction gives
$A'=A$, proving the valuation
ring assertion in the exact
fraction field $K$.

This criterion also proves
that all ideals of a valuation
ring are linearly ordered
by inclusion. For nonzero
$a,b\in A$, applying the
alternative to $a/b$ gives
either $a=bc$ or $b=ac$
with $c\in A$. The zero
cases are immediate divisibility
equalities. Thus principal
ideals are comparable. If
ideals $I,J$ satisfy
$I\not\subseteq J$, choose
$a\in I\setminus J$.
For every $b\in J$, the
alternative $a=bc$ is
impossible because it puts
$a$ in $J$. The other
alternative gives $b\in(a)\subseteq I$;
this includes $b=0$.
Consequently $J\subseteq I$.
In particular all original
prime ideals are in one
inclusion chain. Conversely,
comparability of the original
principal ideals in a nonzero
domain gives the fraction
alternative in its own
fraction field, and the proof
above makes it a valuation ring.
\end{proof}

\begin{lemma}
\label{editorial-source-lemma-colimit-valuation-rings}
\begin{slogan}
Valuation rings are stable under filtered direct limits
\end{slogan}
Let $I$ be a directed set. Let $(A_i, \varphi_{ij})$
be a system of valuation rings over $I$.
Then $A = \colim A_i$ is a valuation ring.
\end{lemma}

\begin{proof}
The directed set is nonempty,
and all transition maps are
the original unital ring maps;
they need not be injective
or local. The colimit $A$
is nonzero: an equality
$1=0$ in it would hold
at a later stage, impossible
in a valuation domain $A_j$.
For $a,b\in A$ with $ab=0$,
choose representatives $a_i,b_i$
at one stage. At a later
stage their product is zero.
That stage is a domain,
so one of the two images
is zero and the corresponding
original $a$ or $b$ is zero.
Thus $A$ is a domain.

For arbitrary $a,b\in A$,
choose such representatives.
If either representative is
zero, one divides the other.
Otherwise Lemma
\ref{editorial-source-lemma-valuation-ring-x-or-x-inverse}
in $\operatorname{Frac}(A_i)$
gives either $a_i=b_i c_i$
or $b_i=a_i c_i$ with
$c_i\in A_i$. The same
original equality holds in
$A$ under the structure map.
For a fraction $a/b$ with
$b\neq0$, this proves that
$a/b\in A$ or, if $a\neq0$,
its inverse belongs to $A$.
Lemma
\ref{editorial-source-lemma-x-or-x-inverse-valuation-ring}
now proves the assertion.

There is an exact fraction-field
description even for the original
noninjective system. Put
$I_i=\ker(A_i\to A)$ and
$B_i=A_i/I_i$. Since $B_i$
injects in the domain $A$,
$I_i$ is prime and $B_i$
is a nonzero domain. The
divisibility equalities in
$A_i$ descend to $B_i$,
so the same criterion makes
$B_i$ a valuation ring.
The induced maps $B_i\to B_j$
are injective: their kernels
are zero because an element
vanishing in $B_j$ has zero
image in $A$, hence was
already in $I_i$.
These are the actual image
subrings of $A$, with
$A=\bigcup_i B_i$.

The field maps
$\operatorname{Frac}(B_i)\to\operatorname{Frac}(B_j)$
send each original $a_i/b_i$
to the ratio of its images;
denominators remain nonzero
by injectivity. They give
the isomorphism
$$
\mathop{\rm colim}_i\operatorname{Frac}(B_i)
\longrightarrow\operatorname{Frac}(A),
\qquad[i,a_i/b_i]\longmapsto a_i/b_i.
$$
Every fraction on the right
has both coefficients in
one common $B_i$, proving
surjectivity. Equality of two
fractions is the equality
of their cross products;
it holds in a common stage
because those stages inject
in $A$. This proves injectivity
and gives the inverse by
choosing a stage for the
two original coefficients.
No fraction-field map is
asserted for an original
transition map before its
kernel has been accounted for.

For a representative $a_i$,
its image in $A$ is a unit
exactly when its image is
a unit at some later stage.
One direction preserves the
original inverse equation.
For the other, an inverse
in $A$ has a representative
at a common stage, and
the equation giving product
$1$ holds at a later stage.
Hence the exact maximal-ideal
criterion is that every
later image of $a_i$ belong
to the corresponding maximal
ideal. A union of the
original maximal ideals need
not work: for the two-stage
system $V\to\operatorname{Frac}(V)$
with a nonfield valuation
ring $V$, each nonzero
element of $\mathfrak m_V$
becomes a unit at the
second stage. The colimit
is the field, with zero
maximal ideal, as the exact
criterion predicts.
\end{proof}

\begin{lemma}
\label{editorial-source-lemma-valuation-ring-cap-field}
Let $L/K$ be an extension of fields. If $B \subset L$
is a valuation ring, then $A = K \cap B$ is a valuation ring.
\end{lemma}

\begin{proof}
Retain the original fields $K\subseteq L$
and the original $B\subseteq L$.
Put $F=\operatorname{Frac}(B)\subseteq L$
and $E=K\cap F$. The ring in
the statement is exactly
$A=K\cap B=E\cap B$; these
are equal subrings of the
original $L$. For every
nonzero $x\in E$, the
valuation alternative in $B$
gives $x\in B$ or $x^{-1}\in B$.
Both $x,x^{-1}$ lie in $E$,
so the alternative holds in
$A$. Lemma
\ref{editorial-source-lemma-x-or-x-inverse-valuation-ring}
makes $A$ a valuation ring
whose exact fraction field
is $E=K\cap F$.
The canonical fraction map
$\operatorname{Frac}(A)\to E$
sends $a/b$ to the same
fraction in $L$. Its inverse
sends a nonzero $x$ to
$x/1$ if $x\in A$, or
to $1/x^{-1}$ otherwise;
the two possible expressions
agree when both apply.
This proves the fraction-field
identification without replacing
either original ambient field.

An original nonzero $a\in A$
is a unit of $A$ precisely
when it is a unit of $B$.
Indeed, if $a^{-1}\in B$,
then $a^{-1}\in K$ as well,
so it belongs to $A$.
The converse follows from
the inclusion. Hence
$$
\mathfrak m_A=A\cap\mathfrak m_B.
$$
The resulting residue map
$A/\mathfrak m_A\to B/\mathfrak m_B$
sends $[a]$ to $[a]$ and
is injective by this exact
kernel calculation. In particular
$B$ dominates the original
intersection $A$.

This also proves extension
of any given valuation ring
$A$ of a field $K$ to
any prescribed extension $L/K$.
Lemma \ref{editorial-source-lemma-dominate}
provides a valuation ring
$B\subseteq L$ of fraction
field $L$ dominating $A$.
Its intersection $A'=K\cap B$
has fraction field $K$ by
the proved formula and is
local with maximal ideal
$A'\cap\mathfrak m_B$.
It contains $A$ and its
maximal ideal contracts to
$\mathfrak m_A$, so it
dominates $A$. Maximality
of the original valuation
ring $A$ in $K$ forces
$A'=A$. Thus the extension
restricts to the original
ring and residue embedding
exactly, not merely to some
valuation ring containing it.
\end{proof}

\begin{lemma}
\label{editorial-source-lemma-valuation-ring-cap-field-finite}
Let $L/K$ be an algebraic extension of fields. If $B \subset L$
is a valuation ring with fraction field $L$ and not a field, then
$A = K \cap B$ is a valuation ring and not a field.
\end{lemma}

\begin{proof}
By Lemma \ref{editorial-source-lemma-valuation-ring-cap-field},
the original $A=K\cap B$ is
a valuation ring, and its
fraction field is exactly $K$
because $\operatorname{Frac}(B)=L$.
If $A$ were a field, it
would equal its fraction field,
so $A=K\subseteq B$.
Every $b\in B$ is algebraic
over $K$, hence integral over
this field by its monic
minimal polynomial. Thus
$K\subseteq B$ is integral.
Lemma \ref{editorial-source-lemma-integral-no-inclusion}
forbids a proper inclusion
of primes in $B$ over
the sole prime of $K$.
But the local domain $B$
has $(0)\subsetneq\mathfrak m_B$
because it is not a field,
a contradiction.

Equivalently the contradiction
has a direct inverse formula.
For nonzero $b\in B$, its
original monic minimal polynomial
$$
b^d+c_1b^{d-1}+\cdots+c_d=0
$$
has $c_d\neq0$ in $K$,
and yields
$$
b^{-1}=-c_d^{-1}
\left(b^{d-1}+\sum_{i=1}^{d-1}c_i b^{d-1-i}\right)\in B.
$$
For $d=1$ the parenthesis
is $1$. Thus every nonzero
$b$ would be a unit,
again making $B$ a field.

The algebraicity hypothesis
cannot be dropped. For any
field $k$, retain
$K=k$, $L=k(t)$ and
$B=k[t]_{(t)}$. Its fraction
field is the original $L$.
For a nonzero original
fraction $P(t)/Q(t)$, write
$P=t^rP_0$, $Q=t^sQ_0$
with $P_0(0),Q_0(0)\neq0$,
retaining both powers and
both residual factors.
If $r\geq s$, the fraction
$t^{r-s}P_0/Q_0$ lies in
$B$; if $r<s$, its inverse
$t^{s-r}Q_0/P_0$ lies in
$B$. The fraction criterion
therefore makes $B$ a
valuation ring. It is not
a field: evaluation at $t=0$
on this localization has
proper kernel $(t)B$, which
contains its nonzero $t$.
Nevertheless $K\cap B=k$
is a field. This identifies
the exact failure without
the original algebraicity
assumption.
\end{proof}

\begin{lemma}
\label{editorial-source-lemma-make-valuation-rings}
Let $A$ be a valuation ring. For any prime ideal $\mathfrak p \subset A$ the
quotient $A/\mathfrak p$ is a valuation ring. The same is true for the
localization $A_\mathfrak p$ and in fact any nonzero localization of $A$.
\end{lemma}

\begin{proof}
For the original prime
$\mathfrak p\subset A$, the
quotient $A/\mathfrak p$ is
a nonzero domain. Any two
nonzero residue classes have
original representatives $a,b$
outside $\mathfrak p$.
The valuation alternative in
$\operatorname{Frac}(A)$ gives
$a=bc$ or $b=ac$ for
some original $c\in A$.
The same equality in the
quotient gives its fraction
alternative, so Lemma
\ref{editorial-source-lemma-x-or-x-inverse-valuation-ring}
makes $A/\mathfrak p$ a
valuation ring. Its maximal
ideal is $\mathfrak m_A/\mathfrak p$:
a class outside that ideal
has an inverse inherited
from $A$. Its fraction field
is the original residue
field $\kappa(\mathfrak p)$,
through the inverse maps
$\overline a/\overline b\mapsto[a/b]$
and $[a/b]\mapsto\overline a/\overline b$
between $\operatorname{Frac}(A/\mathfrak p)$
and $A_{\mathfrak p}/\mathfrak pA_{\mathfrak p}$,
where $b\notin\mathfrak p$.

Let $W\subseteq A$ be an
original multiplicative subset.
If $0\in W$, Lemma
\ref{algebra-lemma-localization-zero}
gives $W^{-1}A=0$, which
is not a domain and hence
not a valuation ring.
This is an allowed localization
under Definition
\ref{algebra-definition-multiplicative-subset}.
If $0\notin W$, the
fraction map embeds $W^{-1}A$
in $K=\operatorname{Frac}(A)$:
equality of fractions is
exactly the original cross-product
equality, since nonzero elements
of $W$ do not annihilate
anything in the domain.
It contains $A$ and has
fraction field $K$. The
alternative $x\in A$ or
$x^{-1}\in A$ for every
nonzero $x\in K$ therefore
holds in $W^{-1}A$, making
it a valuation ring. This
includes $A_{\mathfrak p}$,
whose denominator set avoids zero.

More generally every original
overring $A\subseteq B\subseteq K$
is exactly $A_{\mathfrak q}$
for one prime $\mathfrak q$
of $A$. The same fraction
alternative first proves that
$B$ is a valuation ring
with fraction field $K$.
Put $\mathfrak q=A\cap\mathfrak m_B$.
The original fraction map
$A_{\mathfrak q}\to B$ is
defined because every denominator
outside $\mathfrak q$ is
a unit of $B$, and is
injective inside $K$.
For $b\in B$, if $b\in A$
it is already in $A_{\mathfrak q}$.
Otherwise $b\neq0$ and
$b^{-1}\in A$ by the
original valuation criterion.
It is a unit of $B$,
so $b^{-1}\notin\mathfrak q$
and $b=1/b^{-1}\in A_{\mathfrak q}$.
This proves equality of the
original subrings. The inverse
prime assignment is contraction
of the maximal ideal, since
$\mathfrak qA_{\mathfrak q}\cap A=\mathfrak q$.
Thus primes and overrings
correspond bijectively, reversing
inclusion: larger primes invert
fewer original elements.
Conversely an inclusion of overrings
preserves their units, so contraction
of their maximal ideals reverses
that inclusion as well.

For the given $W$, this
prime is explicitly
$$
\mathfrak q=\{a\in A:(a)\cap W=\emptyset\}.
$$
Indeed $a/1$ is a unit
in $W^{-1}A$ precisely when
$ab=w$ for some $b\in A$
and $w\in W$. An inverse
$b/w$ gives this exact
equality by cancelling the
nonzero localization witness
in the domain, and the
equality conversely gives
the inverse. This proves
the formula from the maximal
ideal of the valuation ring.
The prime is disjoint from
$W$ and contains every prime
disjoint from $W$: if
$a$ is in such a prime,
no multiple $ab$ can be
in $W$. It is therefore
the largest such prime,
and the canonical fraction
map identifies $W^{-1}A$
with exactly $A_{\mathfrak q}$.
\end{proof}

\begin{lemma}
\label{editorial-source-lemma-stack-valuation-rings}
Let $A'$ be a valuation ring with residue field $K$.
Let $A$ be a valuation ring with fraction field $K$.
Then
$C = \{\lambda \in A' \mid \lambda \bmod \mathfrak m_{A'} \in A\}$
is a valuation ring.
\end{lemma}

\begin{proof}
Keep the original residue map
$\pi:A'\to K$ and write
$\mathfrak p=\mathfrak m_{A'}$.
The original $C=\pi^{-1}(A)$
is a subring of $A'$,
and $\mathfrak p\subseteq C$.
The map $C/\mathfrak p\to A$
sending $[c]$ to $\pi(c)$
is an isomorphism: surjectivity
follows by lifting each $a\in A$
through $\pi$, and its kernel
is exactly $\mathfrak p$.
The inverse sends $a$ to
the class of any such lift;
different lifts differ in
$\mathfrak p$. In particular
$\mathfrak p$ is a prime
of the original $C$.

There is a more precise
comparison than just equality
of fraction fields:
$$
C_{\mathfrak p}\longrightarrow A',\qquad c/t\longmapsto c t^{-1}
$$
is an isomorphism inside
$F=\operatorname{Frac}(A')$.
It is defined and injective
because $t\in C\setminus\mathfrak p$
is a nonzero unit of $A'$.
For surjectivity, take the
original $b\in A'$ and
write its residue in
$K=\operatorname{Frac}(A)$
as $\pi(b)=a/d$ with
$a,d\in A$ and $d\neq0$;
for zero residue take $a=0,d=1$.
Choose $t\in A'$ lifting
this original $d$. Then
$t\in C\setminus\mathfrak p$
and $tb\in C$ because
$\pi(tb)=a$. Thus the
inverse sends $b$ to
$(tb)/t$. Independence of
the lift and denominator
follows from the equality
of cross products in the
original field $F$. Both
composites are identities.
It follows that
$\operatorname{Frac}(C)=F$,
including the case $A'$ is
a field and $\mathfrak p=0$.

For nonzero $x\in F$, if
$x\notin A'$, Lemma
\ref{editorial-source-lemma-valuation-ring-x-or-x-inverse}
gives $x^{-1}\in\mathfrak m_{A'}\subseteq C$.
If $x\in\mathfrak m_{A'}$,
then $x\in C$. In the
remaining case $x$ is a
unit of $A'$, so its
nonzero residue satisfies
$\pi(x)\in A$ or
$\pi(x)^{-1}\in A$ by the
valuation alternative in the
original residue field $K$.
Because $\pi(x^{-1})=\pi(x)^{-1}$,
this says $x\in C$ or
$x^{-1}\in C$. The criterion
of Lemma
\ref{editorial-source-lemma-x-or-x-inverse-valuation-ring}
therefore proves the assertion.

The maximal ideal and residue
field are also exact:
$$
\mathfrak m_C=\pi^{-1}(\mathfrak m_A),\qquad
C/\mathfrak m_C\simeq A/\mathfrak m_A.
$$
Indeed an element $c$ outside
that inverse image has
$\pi(c)\in A^*$, so it
is a unit of $A'$ and
$\pi(c^{-1})=\pi(c)^{-1}\in A$,
whence $c^{-1}\in C$.
An element inside cannot be
a unit, by applying $\pi$
to an inverse equation.
The residue isomorphism is
the map $[c]\mapsto[\pi(c)]$,
with inverse given by a
lift as above.

Finally the original prime
chain of $C$ is described
by these two exact maps.
Primes contained in $\mathfrak p$
correspond to primes of
$C_{\mathfrak p}=A'$ by
extension and contraction.
Primes containing $\mathfrak p$
correspond to primes of
$C/\mathfrak p=A$ by
quotient and inverse image.
All primes of the valuation
ring $C$ are comparable,
as proved in Lemma
\ref{editorial-source-lemma-x-or-x-inverse-valuation-ring},
so these two portions cover
its spectrum. They meet
at the single original prime
$\mathfrak p$, corresponding
to $\mathfrak m_{A'}$ on
the localization side and
to $(0)$ on the quotient
side. This proves the full
prime-chain composition.
\end{proof}

\begin{lemma}
\label{editorial-source-lemma-find-valuation-rings}
Let $A$ be a normal domain with fraction field $K$.
\begin{enumerate}
\item For every $x \in K$, $x \not \in A$ there exists a valuation ring
$A \subset V \subset K$ with fraction field $K$ such that $x \not \in V$.
\item If $A$ is local, we can moreover choose $V$ which dominates $A$.
\end{enumerate}
In other words, $A$ is the intersection of all valuation rings in $K$
containing $A$ and if $A$ is local, then $A$ is the intersection of
all valuation rings in $K$ dominating $A$.
\end{lemma}

\begin{proof}
For the given normal domain,
$x\in K\setminus A$ is
not integral over $A$.
We prove the construction
using exactly this nonintegrality
condition; it will also give
the integral-closure statement
for arbitrary domains below.
Since $0\in A$, the original
$x$ is nonzero. Put
$B=A[x^{-1}]\subseteq K$.
If $x\in B$, a finite
expression
$x=\sum_{i=0}^d a_i x^{-i}$
would give the full monic
equation
$$
x^{d+1}-\sum_{i=0}^d a_i x^{d-i}=0,
$$
contradicting nonintegrality.
Thus $x^{-1}$ is not a
unit of $B$, and
$x^{-1}B$ is a proper
ideal. Choose a prime
$\mathfrak p$ containing it
and apply Lemma
\ref{editorial-source-lemma-dominate} to the
original local subring
$B_{\mathfrak p}\subseteq K$.
It gives a valuation ring
$V$ of fraction field $K$
dominating $B_{\mathfrak p}$.
The element $x^{-1}$ belongs
to $\mathfrak m_V$, so
$x\notin V$: otherwise
their product would put
$1$ in $\mathfrak m_V$.
Since $A\subseteq B\subseteq V$,
this proves (1).

Suppose now the original
$A$ is local. Retain
$$
J=x^{-1}B+\mathfrak m_A B.
$$
If $J=B$, finite expressions
in the original $x^{-1}$,
with zero padding to a
common degree $d\geq0$,
give
$$
1=\left(\sum_{i=0}^d a_i x^{-i}\right)x^{-1}
+\sum_{i=0}^d a'_i x^{-i},
\qquad a_i\in A,\quad a'_i\in\mathfrak m_A.
$$
Multiplying by the full
$x^{d+1}$ and keeping the
two original coefficient lists
gives
$$
(1-a'_0)x^{d+1}
-\sum_{i=0}^d a_i x^{d-i}
-\sum_{i=1}^d a'_i x^{d+1-i}=0.
$$
For $d=0$ the last sum
is empty and the equation
is $(1-a'_0)x-a_0=0$.
In every case $u=1-a'_0$
is a unit of $A$.
The monic polynomial
$$
Q(T)=T^{d+1}
-\sum_{i=0}^d u^{-1}a_i T^{d-i}
-\sum_{i=1}^d u^{-1}a'_i T^{d+1-i}
$$
annihilates $x$, and the
complete polynomial originally
obtained is exactly $uQ(T)$.
Thus $x$ would be integral,
a contradiction. It follows
that the original $J$ is
proper. Choose a prime
$\mathfrak p\supseteq J$ and
then a valuation ring $V$
of fraction field $K$
dominating $B_{\mathfrak p}$,
as in the first construction.

The contraction $\mathfrak p\cap A$
contains $\mathfrak m_A$ and
is proper, so it equals
$\mathfrak m_A$. Localization
gives
$\mathfrak m_{B_{\mathfrak p}}\cap A=\mathfrak m_A$,
and domination gives
$\mathfrak m_V\cap A=\mathfrak m_A$.
Hence this original $V$
dominates $A$ and still
excludes $x$ because
$x^{-1}\in\mathfrak m_V$.
This proves (2). Each
asserted intersection equals
$A$: inclusion of $A$
is built into all its
members, and the construction
excludes every original element
of $K\setminus A$ from
at least one member.

More generally let $A$
be any domain with original
fraction field $K$, and
let $\widetilde A$ be its
integral closure in $K$.
Every $x\in\widetilde A$
belongs to every valuation
ring $A\subseteq V\subseteq K$:
its original monic equation
over $A$ is the same
monic equation over $V$,
and $V$ is normal by
Lemma \ref{editorial-source-lemma-valuation-ring-normal}.
If $x\notin\widetilde A$,
the construction above applies
because $x$ is not integral
over $A$, and gives a
member excluding it. Therefore
$$
\widetilde A=\bigcap_{A\subseteq V\subseteq K\text{ valuation}}V.
$$
For local $A$, exactly the
second construction gives
the same equality with the
intersection restricted to
valuation rings dominating $A$.

For a prescribed prime
$\mathfrak p\subset A$, the
valuation rings containing
$A$ with center
$\mathfrak m_V\cap A=\mathfrak p$
are precisely those dominating
the original $A_{\mathfrak p}$.
Indeed every $a\in A\setminus\mathfrak p$
is a unit of $V$, giving
the actual inclusion of
fractions $A_{\mathfrak p}\to V$,
and its maximal ideal
contracts to $\mathfrak pA_{\mathfrak p}$.
The converse follows by
contracting once more to
$A$. Their intersection
is thus exactly the
integral closure of
$A_{\mathfrak p}$ in $K$,
by the proved local formula.
\end{proof}

\noindent
A {\it totally ordered abelian group} is a pair $(\Gamma, \geq)$
consisting of an abelian group $\Gamma$ endowed with a total
ordering $\geq$ such that $\gamma \geq \gamma' \Rightarrow
\gamma + \gamma'' \geq \gamma' + \gamma''$ for all
$\gamma, \gamma', \gamma'' \in \Gamma$.

\begin{lemma}
\label{editorial-source-lemma-valuation-group}
Let $A$ be a valuation ring with field of fractions $K$.
Set $\Gamma = K^*/A^*$ (with group law written additively).
For $\gamma, \gamma' \in \Gamma$
define $\gamma \geq \gamma'$ if and only if
$\gamma - \gamma'$ is in the image of $A - \{0\} \to \Gamma$.
Then $(\Gamma, \geq)$ is a totally ordered abelian group.
\end{lemma}

\begin{proof}
Retain the original quotient
$\Gamma=K^*/A^*$ and write
$v(x)=[x]$ additively, so
$v(xy)=v(x)+v(y)$ by the
quotient group law. For original
$x,y\in K^*$, the proposed
order is exactly
$$
v(x)\geq v(y)\quad\Longleftrightarrow\quad x/y\in A.
$$
Indeed $v(x/y)$ is in the
image of $A\setminus\{0\}$
precisely when $x/y$ differs
from such an element by
a unit of $A$. Multiplication
by that unit and its inverse
proves the equivalence.
If the representatives are
changed to $xu,yw$ with
$u,w\in A^*$, the original
ratio becomes $(x/y)(u/w)$.
Its membership in $A$ is
unchanged, so the order is
well-defined on the quotient.

Reflexivity follows from $x/x=1$.
If $x/y,y/x\in A$, their
product is $1$, so $x/y$
is a unit and $v(x)=v(y)$;
this is antisymmetry. If
$x/y$ and $y/z$ belong
to $A$, their exact product
$x/z$ does too, proving
transitivity. For every pair,
the valuation alternative for
$x/y$ gives $x/y\in A$
or $y/x\in A$, proving
totality. Finally multiplying
both original representatives
by the same $z\in K^*$
does not change their ratio:
$(xz)/(yz)=x/y$. This
proves translation invariance
of the order and the
claimed ordered group structure.

Every totally ordered abelian
group is torsion-free. If
$\gamma>0$, repeated strict
translation invariance gives
$n\gamma>0$ for every integer
$n\geq1$; if $\gamma<0$,
apply the same argument to
$-\gamma$. Hence $n\gamma=0$
with $n\geq1$ forces
$\gamma=0$. In particular
this holds for the original
$\Gamma$.

The order recovers the
original ring and maximal
ideal exactly:
$$
A=\{0\}\cup\{x\in K^*:v(x)\geq0\},\qquad
\mathfrak m_A=\{0\}\cup\{x\in K^*:v(x)>0\}.
$$
The first equality is the
ratio criterion with $y=1$.
For $x\in A\setminus\{0\}$,
the quotient definition says
$v(x)=0$ exactly when
$x\in A^*$, giving the
second equality by localness.
Also $v(a/b)=v(a)-v(b)$
for all original nonzero
$a,b\in K$. The value
of $0$ is not introduced
into this group. If $A=K$,
the quotient is the zero
group, with its unique order.
Conversely, if this quotient
is zero, every nonzero
element of $K$ is a unit
of $A$, so $A=K$.
\end{proof}

\begin{definition}
\label{editorial-source-definition-value-group}
Let $A$ be a valuation ring.
\begin{enumerate}
\item The totally ordered abelian group $(\Gamma, \geq)$ of
Lemma \ref{editorial-source-lemma-valuation-group} is called the
{\it value group} of the valuation ring $A$.
\item The map $v : A - \{0\} \to \Gamma$ and also $v : K^* \to \Gamma$ is
called the {\it valuation} associated to $A$.
\item The valuation ring $A$ is called a {\it discrete valuation ring}
if $\Gamma \cong \mathbf{Z}$.
\end{enumerate}
\end{definition}

\noindent
Suppose the original ordered group
$\Gamma$ is isomorphic as a group
to $\mathbf Z$, and choose an
initial group isomorphism
$\psi:\Gamma\to\mathbf Z$.
Put $\gamma=\psi^{-1}(1)\neq0$.
Let $\varepsilon=1$ if
$\gamma>0$, and $\varepsilon=-1$
if $\gamma<0$. Then
$$
\phi=\varepsilon\psi:\Gamma\longrightarrow\mathbf Z
$$
is an order isomorphism with
the usual order on the
receiving integers. Indeed
$\varepsilon\gamma>0$, and
every original element is
$n\gamma$ for a unique
$n\in\mathbf Z$. Repeated
addition and translation
invariance give
$n\gamma\geq0$ precisely when
$\varepsilon n\geq0$, which
is precisely $\phi(n\gamma)\geq0$.
The same applies to differences,
proving both order implications.
Any other group isomorphism
is $\delta\psi$ for
$\delta\in\{1,-1\}$, since
the only generators of
$\mathbf Z$ are $1,-1$.
Requiring the inverse image
of $1$ to be positive
forces $\delta=\varepsilon$.
This is the precise unique
order-compatible isomorphism.
The original $\Gamma$ and
its value map $v$ are
retained; their integer comparison
is $\phi\circ v$, with
$v=\phi^{-1}\circ(\phi\circ v)$.

Write $\gamma_0=\phi^{-1}(1)>0$.
Surjectivity of the original
quotient $v:K^*\to\Gamma$
gives an element $\pi\in K^*$
with $v(\pi)=\gamma_0$.
The positive-cone description
puts $\pi$ in the maximal
ideal of $A$. For every
original nonzero $a\in A$,
put $n=\phi(v(a))\geq0$.
The exact original element
$u=a/\pi^n$ has value
$v(a)-n\gamma_0=0$, so
$u\in A^*$ and
$$
a=u\pi^n.
$$
Both the unit $u$ and
the original exponent are
part of this equality;
neither is discarded. For
the chosen $\pi$, the
integer $n$ is unique by
applying $\phi v$, and
then $u=a/\pi^n$ is
unique as well. Another
element $\pi'$ of the
same least positive value
has $\pi'/\pi\in A^*$,
which is the exact comparison
between such choices.

Every nonzero ideal $I$
has the form $(\pi^n)$
for exactly one $n\geq0$.
Choose the least integer
among $\phi(v(a))$ for
the original nonzero $a\in I$,
and choose such an $a$.
The full equality $a=u\pi^n$
with $u$ a unit puts
$\pi^n=u^{-1}a$ in $I$.
For every nonzero $b\in I$,
the inequality $v(b)\geq n\gamma_0$
puts $b/\pi^n$ in $A$,
and a zero member contributes
the equality $0=\pi^n0$.
Thus all of $I$ is
contained in $(\pi^n)$,
giving equality. Distinct
$n$ give distinct ideals
by their least original
values. The zero ideal
is retained as a separate
case and is not assigned
an infinite value in $\Gamma$.
In particular the maximal
ideal is $(\pi)$, every
ideal is principal, and
the only prime ideals are
$(0)$ and $(\pi)$: for
$n>1$, the product
$\pi\pi^{n-1}$ lies in
$(\pi^n)$ although neither
factor does. This proves
the ideal and prime descriptions
from the original ordered
quotient without suppressing
any unit factors.

\begin{lemma}
\label{editorial-source-lemma-properties-valuation}
Let $A$ be a valuation ring. The valuation $v : A -\{0\} \to \Gamma_{\geq 0}$
has the following properties:
\begin{enumerate}
\item $v(a) = 0 \Leftrightarrow a \in A^*$,
\item $v(ab) = v(a) + v(b)$,
\item $v(a + b) \geq \min(v(a), v(b))$ provided $a + b \not = 0$.
\end{enumerate}
\end{lemma}

\begin{proof}
The quotient map to the
original $\Gamma=K^*/A^*$
has kernel $A^*$, giving
(1), and its group law
gives (2). For (3), retain
the original nonzero $a,b$
and assume $a+b\neq0$.
Interchanging their roles
if necessary, suppose
$v(a)\leq v(b)$. The
order proved in Lemma
\ref{editorial-source-lemma-valuation-group}
then gives $b/a\in A$.
The exact equality
$$
\frac{a+b}{a}=1+\frac ba
$$
puts the nonzero left-hand
element in $A$, so
$v(a+b)-v(a)\geq0$.
This proves (3), retaining
both summands and their
original ratio.

If $v(a)<v(b)$, that
ratio $b/a$ belongs to
$\mathfrak m_A$. Thus
$1+b/a$ is a unit:
if it were in the maximal
ideal, subtracting $b/a$
would put $1$ there.
Consequently $a+b\neq0$
and $v(a+b)=v(a)$.
The reversed strict inequality
gives the symmetric equality.
This entire proof also
applies to nonzero $a,b\in K$
outside $A$: the comparison
still puts the same ratio
in $A$. The original value
map on $K^*$ therefore
satisfies the same inequality,
with equality whenever the
two original values differ.

For a finite original list
$a_1,\ldots,a_n\in K$, keep
all its terms, including
zeros. If at least one
term is nonzero, choose
$a_j$ of least value among
the nonzero terms. Every
ratio $a_i/a_j$ belongs
to $A$, with a zero
term giving the ratio $0$.
Hence if the original sum
$a=\sum_i a_i$ is nonzero,
then $a/a_j=\sum_i a_i/a_j$
is a nonzero element of
$A$, proving $v(a)\geq v(a_j)$.
If exactly one term has
the least value, every
other ratio is in
$\mathfrak m_A$, including
the zeros. The full sum
of ratios is $1$ plus
an element of that ideal,
hence a unit. In that
case the sum is nonzero
and has exactly the least
original value. For an
all-zero or empty list,
the sum is $0$ and no
value in $\Gamma$ is
assigned to it.
\end{proof}

\begin{lemma}
\label{editorial-source-lemma-characterize-valuation-ring}
Let $A$ be a ring. The following are equivalent
\begin{enumerate}
\item $A$ is a valuation ring,
\item $A$ is a local domain and every finitely generated
ideal of $A$ is principal.
\end{enumerate}
\end{lemma}

\begin{proof}
Suppose first that $A$ is
a valuation ring. Retain
the complete finite generator
list $f_1,\ldots,f_n$.
If the list is empty
or all generators are zero,
its ideal is $(0)$ and
is principal. Otherwise choose
an original nonzero $f_i$
whose value is least among
the nonzero generators, and
put $h=f_i$. For each
nonzero $f_j$, the comparison
$v(f_j)\geq v(h)$ gives
the original ratio $c_j=f_j/h\in A$.
For $f_j=0$, set $c_j=0$.
Then the exact equalities
$f_j=hc_j$ for every $j$
give $(f_1,\ldots,f_n)\subseteq(h)$,
while $h$ itself is in
the original list, giving
the reverse inclusion.
No value is taken at
a zero generator.

Conversely, assume the original
$A$ is a local domain
with every finitely generated
ideal principal. An original
fraction $f/g$ with $g\neq0$
belongs to $A$ if $f=0$.
For $f,g\neq0$, write
$(f,g)=(h)$. Then $h\neq0$,
and there are original
$a,b,c,d\in A$ with
$$
f=ah,\qquad g=bh,\qquad h=cf+dg.
$$
Substitution gives
$h=(ca+db)h$. The domain
property and $h\neq0$ permit
cancellation, giving $ca+db=1$.
Both $a,b$ cannot belong
to the maximal ideal,
so at least one is a
unit. If $a$ is a
unit, the original fraction
$g/f=b/a$ belongs to $A$;
if $b$ is a unit,
$f/g=a/b$ belongs to $A$.
This proves the nonzero
fraction alternative in the
actual $\operatorname{Frac}(A)$.
Lemma
\ref{editorial-source-lemma-x-or-x-inverse-valuation-ring}
therefore makes $A$ a
valuation ring.

The converse proof uses
only two-generated ideals.
Thus in (2) it suffices
to require that every
two-generated ideal be principal;
the forward proof then
recovers the assertion for
every original finite list.
The local hypothesis cannot
be discarded: every ideal
of $\mathbf Z$ is principal,
since for a nonzero ideal
its least positive element
divides all its elements
by division with remainder,
and the zero ideal is
$(0)$. But neither $2/3$
nor $3/2$ is an integer,
so the exact fraction
criterion shows that $\mathbf Z$
is not a valuation ring.
\end{proof}

\begin{lemma}
\label{editorial-source-lemma-valuation-valuation-ring}
Let $(\Gamma, \geq)$ be a totally ordered abelian group.
Let $K$ be a field. Let $v : K^* \to \Gamma$ be a homomorphism
of abelian groups such that $v(a + b) \geq \min(v(a), v(b))$ for
$a, b \in K$ with $a, b, a + b$ not zero. Then
$$
A =
\{
x \in K \mid x = 0 \text{ or } v(x) \geq 0
\}
$$
is a valuation ring with value group $\Im(v) \subset \Gamma$,
with maximal ideal
$$
\mathfrak m =
\{
x \in K \mid x = 0 \text{ or } v(x) > 0
\}
$$
and with group of units
$$
A^* =
\{
x \in K^* \mid v(x) = 0
\}.
$$
\end{lemma}

\begin{proof}
Keep the original ordered group
$\Gamma$ and homomorphism
$v:K^*\to\Gamma$, without
assuming it is surjective.
The homomorphism identities give
$v(1)=0$ and $v(x^{-1})=-v(x)$.
Also $2v(-1)=v(1)=0$.
A totally ordered abelian group
has no nonzero element killed
by a positive integer, by
the strict-order argument in
Lemma \ref{editorial-source-lemma-valuation-group}.
Hence $v(-1)=0$, and
$v(-x)=v(x)$ for nonzero $x$.

The original set $A$ contains
$0,1$ and is closed under
additive inverses. It is
closed under multiplication:
products involving zero are
zero, and otherwise the
sum of two nonnegative
original values is nonnegative.
For addition, a zero summand
causes no problem. If both
summands and their sum are
nonzero, the given inequality
puts the sum in $A$;
if their sum is zero,
it belongs to $A$ by
definition. Thus $A$ is
a unital subring of $K$.
For every nonzero $x\in K$,
totality gives $v(x)\geq0$
or $v(x)<0$. In the
second case $v(x^{-1})>0$,
so $x^{-1}\in A$.
The criterion of Lemma
\ref{editorial-source-lemma-x-or-x-inverse-valuation-ring}
therefore proves that $A$
is a valuation ring whose
fraction field is exactly $K$.

A nonzero original $x$
is a unit of $A$ exactly
when both $v(x)$ and
$-v(x)$ are nonnegative,
which by antisymmetry is
equivalent to $v(x)=0$.
Thus its nonunits are
precisely $0$ and the
elements of positive value,
giving the asserted maximal
ideal. Directly, that set
is a proper ideal: addition
uses the minimum inequality
unless the sum is zero,
and multiplying a positive
value by an element of
$A$ adds a nonnegative
value. It excludes $1$,
and every element outside
it has the inverse just
described. This also proves
maximality explicitly.

Let $q:K^*\to\Delta=K^*/A^*$
be the original associated
value-group quotient. The map
$$
\overline v:\Delta\longrightarrow\operatorname{Im}(v),
\qquad[x]\longmapsto v(x)
$$
is well-defined because
$A^*=\ker(v)$, injective
by the same kernel equality,
and surjective by the
definition of the image.
Its inverse sends a value
$\gamma=v(x)$ to $[x]$;
different original choices
$x,y$ have $v(x/y)=0$
and differ by a unit of
$A$, proving independence.
Moreover
$$
[x]\geq[y]\quad\Longleftrightarrow\quad
x/y\in A\quad\Longleftrightarrow\quad
v(x)-v(y)\geq0.
$$
Thus it is an order
isomorphism with exactly
the asserted subgroup of
the original $\Gamma$.
If $i:\operatorname{Im}(v)\hookrightarrow\Gamma$
is its original inclusion,
then $i\overline vq=v$;
the ambient codomain and
the original value map
have both been retained.

Finally two such original
maps $v:K^*\to\Gamma$ and
$w:K^*\to\Lambda$ define
the same valuation ring
exactly when there is an
order isomorphism
$\phi:\operatorname{Im}(v)\to\operatorname{Im}(w)$
with $\phi(v(x))=w(x)$
for every original $x$.
For a common ring, equality
of two $v$-values means
their ratio is a unit
of that ring, and hence
also has $w$-value zero.
This proves the formula
well-defined; the reverse
argument gives its inverse.
The same ratio criterion
proves that it preserves
and reflects order. Conversely
such an isomorphism preserves
the nonnegative cone, so
the two original sets
$\{0\}\cup v^{-1}(\Gamma_{\geq0})$
and $\{0\}\cup w^{-1}(\Lambda_{\geq0})$
are equal. Uniqueness follows
because every element of
the image is an original
value $v(x)$. No isomorphism
between unused parts of the
ambient groups is asserted.
\end{proof}

\noindent
Let $(\Gamma, \geq)$ be a totally ordered abelian group.
An {\it ideal of $\Gamma$} is a subset $I \subset \Gamma$ such
that all elements of $I$ are $\geq 0$ and $\gamma \in I$,
$\gamma' \geq \gamma$ implies $\gamma' \in I$. We say that such
an ideal is {\it prime} if $0 \not \in I$ and if
$\gamma + \gamma' \in I, \gamma, \gamma' \geq 0
\Rightarrow \gamma \in I \text{ or } \gamma' \in I$.

\begin{lemma}
\label{editorial-source-lemma-ideals-valuation-ring}
Let $A$ be a valuation ring.
Ideals in $A$ correspond $1 - 1$ with ideals of $\Gamma$.
This bijection is inclusion preserving, and maps prime
ideals to prime ideals.
\end{lemma}

\begin{proof}
Retain $K=\operatorname{Frac}(A)$, the original ordered quotient
$\Gamma=K^*/A^*$, and its quotient map $v:K^*\to\Gamma$.
For an ideal $J\subset A$ and an ideal $U\subset\Gamma$, define
$$
V(J)=\{v(x):x\in J,\ x\neq0\},\qquad
J(U)=\{0\}\cup\{x\in K^*:v(x)\in U\}.
$$
Every element of $V(J)$ is nonnegative. If $v(x)=\gamma\in V(J)$
and $\delta\geq\gamma$, choose $y\in K^*$ with $v(y)=\delta$.
Then $v(y/x)=\delta-\gamma\geq0$, so $y/x\in A$ and
$y=(y/x)x\in J$. Thus $V(J)$ is upward closed.
The same calculation shows that every representative of a value in
$V(J)$ lies in $J$; when the values are equal, $y/x$ is an actual unit.

Since $U\subset\Gamma_{\geq0}$, the set $J(U)$ lies in $A$.
It contains zero. If $x\neq0$ belongs to it, then
$v(-x)=v(x)\in U$. For nonzero $x,y\in J(U)$ with $x+y\neq0$,
the valuation inequality gives
$v(x+y)\geq\min\{v(x),v(y)\}\in U$, hence $x+y\in J(U)$.
If a summand or the sum is zero the closure follows directly from
the definition. For $a\in A$ and $x\in J(U)$, a zero factor gives
$ax=0$; otherwise $v(ax)=v(a)+v(x)\geq v(x)$, so $ax\in J(U)$.
This proves that $J(U)$ is an ideal.
Surjectivity of $v$ and the representative calculation give
$V(J(U))=U$ and $J(V(J))=J$. Both maps preserve inclusion.
In particular $J(\varnothing)=(0)$ and
$J(\Gamma_{\geq0})=A$; no value is assigned to the zero element.

The ideal $J$ is proper exactly when $0\notin V(J)$, since an
element of value zero is a unit. If $J$ is prime and
$\gamma,\delta\geq0$ with $\gamma+\delta\in V(J)$, choose
$x,y\in A\setminus\{0\}$ with these values. The representative
calculation puts $xy$ in $J$, so $x\in J$ or $y\in J$.
Consequently $\gamma\in V(J)$ or $\delta\in V(J)$.
Conversely, suppose $U$ is prime and $xy\in J(U)$ for $x,y\in A$.
If $xy=0$, the domain property gives a zero factor, which belongs
to $J(U)$. Otherwise the defining condition on
$v(x)+v(y)\in U$ puts one factor in $J(U)$.
Thus the correspondence preserves and reflects primality, including
the empty group ideal and the zero ring ideal.

The correspondence also describes the ideal operations exactly.
For every family $(J_\lambda)_{\lambda\in\Lambda}$, put
$U_\lambda=V(J_\lambda)$. Then
$$
V\left(\bigcap_{\lambda\in\Lambda}J_\lambda\right)
=\bigcap_{\lambda\in\Lambda}U_\lambda,
\qquad
V\left(\sum_{\lambda\in\Lambda}J_\lambda\right)
=\bigcup_{\lambda\in\Lambda}U_\lambda.
$$
The first equality follows by membership of each nonzero element.
For the second, a nonzero member of the sum is a finite sum of
members of the original ideals. Discarding no term from that
equality, list its nonzero terms. Their list is nonempty; its least
value belongs to one $U_\lambda$. The finite-sum inequality and
upward closure put the value of the sum in that same $U_\lambda$.
The reverse inclusion follows from each $J_\lambda$ being contained
in the sum. For an empty family the intersection is $A$ and its
value set is $\Gamma_{\geq0}$, while the sum is $(0)$ and its
value set is empty; the intersections on the right are taken
inside $\Gamma_{\geq0}$.
For two ideals with value sets $U,W$, one has
$$
V(J(U)J(W))=U+W
=\{\gamma+\delta:\gamma\in U,\ \delta\in W\},
$$
and
$$
V(\sqrt{J(U)})
=\{\gamma\in\Gamma_{\geq0}:n\gamma\in U
\text{ for some integer }n\geq1\}.
$$
Indeed $U+W$ is upward closed: if
$\eta\geq\gamma+\delta$, write
$\eta=\gamma+(\delta+\eta-\gamma-\delta)$ and use
$\delta+\eta-\gamma-\delta\geq\delta$.
Every value in $U+W$ is the value of a product of representatives
from the original ideals. A nonzero element of their product ideal
is a finite sum of such products, and the least-value argument
just used gives the reverse inclusion. If either ideal is zero,
both sides are empty. Finally, for nonzero $x\in A$, the equality
$v(x^n)=n v(x)$ proves the radical formula in both directions;
the original zero element lies in the radical separately.
\end{proof}

\begin{lemma}
\label{editorial-source-lemma-valuation-ring-Noetherian-discrete}
A valuation ring is Noetherian if and only if it is
a discrete valuation ring or a field.
\end{lemma}

\begin{proof}
Let $K=\operatorname{Frac}(A)$ and retain the original value map
$v:K^*\to\Gamma=K^*/A^*$. If $A$ is a field, its two ideals
are $(0)$ and $A$, so it is Noetherian.
If $A$ is a discrete valuation ring, let $\gamma_0>0$ generate
its original ordered group, and let
$\phi:\Gamma\to\mathbf Z$ be the order isomorphism
$\phi(n\gamma_0)=n$. Choose $\pi\in A$ with
$v(\pi)=\gamma_0$. The nonzero ideals, and only the nonzero
ideals, are
$$
I_n=\{0\}\cup\{x\in K^*:v(x)\geq n\gamma_0\}
=(\pi^n),\qquad n\geq0.
$$
To verify this directly, the values of a nonzero ideal under
$\phi$ form a nonempty upward closed subset of
$\mathbf Z_{\geq0}$, with a least element $n$.
For an original element $a$ of that value,
$u=a/\pi^n$ has value zero, so $u\in A^*$ and
$a=u\pi^n$. Every original nonzero $b$ of the ideal satisfies
$b/\pi^n\in A$, while $0=\pi^n0$. These equalities prove
that the ideal is $(\pi^n)$ and retain its original unit factor.
The remaining ideal $(0)$ is generated by zero. Thus every ideal
is finitely generated and $A$ is Noetherian. The map $v$ still
takes values in $\Gamma$; the integer-valued comparison is
$\phi\circ v$, with $v=\phi^{-1}\circ(\phi\circ v)$.

Conversely, suppose $A$ is Noetherian and is not a field.
Ascending chains of its ideals stabilize: their union is an ideal,
and finitely many generators of that union all occur in a single
member of the chain. By Lemma \ref{editorial-source-lemma-ideals-valuation-ring},
the same holds for the upward closed subsets of
$\Gamma_{\geq0}$. There is a positive element of $\Gamma$.
If there were no least positive element, starting from one such
element one could choose
$\gamma_1>\gamma_2>\gamma_3>\cdots>0$.
The corresponding ideals
$\gamma_i+\Gamma_{\geq0}$ would form a strictly ascending
chain: $\gamma_{i+1}$ belongs to its own set but not to the
preceding one. This is impossible. Denote the least positive
element by $\gamma_0$.

For any original $\gamma>0$, suppose that
$\gamma-n\gamma_0\geq0$ for every integer $n\geq0$.
These elements strictly decrease, and their upward closed sets
again give a strictly ascending chain of ideals, a contradiction.
Consequently there is a least integer $N\geq1$ such that
$\gamma-N\gamma_0<0$. Put
$r=\gamma-(N-1)\gamma_0$. Then
$0\leq r<\gamma_0$, so the least-positive property forces
$r=0$. Thus $\gamma=(N-1)\gamma_0$.
Zero is $0\gamma_0$, and applying the argument to $-\gamma$
expresses every negative element as a negative multiple of
$\gamma_0$. Distinct integer multiples are distinct, because
a positive multiple of $\gamma_0$ is positive. Therefore
$n\mapsto n\gamma_0$ is an order isomorphism
$\mathbf Z\to\Gamma$, and $A$ is a discrete valuation ring.
The argument uses the ascending chain condition twice; merely
having a least positive group element is not the assertion used
to conclude that all positive elements are its multiples.
\end{proof}


























\subsection{Editorial treatment of Algebra lines 12153--12184}

\hypertarget{algebra-context-158}{}

\noindent\href{https://github.com/stacks/stacks-project/blob/a04446e57ec1fbc252a871afcec7752fb2807b14/algebra.tex\#L12153}{Original source, lines 12153--12184}. Complete original and reviewed TeX passages accompany this note. Every intervention is separately recorded; the following treatment is editorial.


\begin{lemma}
\label{editorial-source-lemma-Noetherian-basic}
Let $R$ be a Noetherian ring.
Any finite $R$-module is of finite presentation.
Any submodule of a finite $R$-module is finite.
The ascending chain condition holds for $R$-submodules
of a finite $R$-module.
\end{lemma}

\begin{proof}
We first prove that every submodule $N$ of a finite $R$-module
$M$ is finite, by induction on a chosen number $n$ of generators
$x_1,\ldots,x_n$. If $n=0$, then $M=N=0$ and the empty list
generates $N$. If $n=1$, the surjection $R\to M$, $a\mapsto ax_1$,
has an ideal $I$ as kernel. The inverse image $J$ of $N$ is an
ideal containing $I$, and the induced map $J/I\to N$ is an
isomorphism. A finite generating list of $J$, which exists because
$R$ is Noetherian, maps to a generating list of $N$.

For $n>1$, let $M'$ be the submodule generated by
$x_1,\ldots,x_{n-1}$ and let $M''=M/M'$ with quotient map $q$.
The module $M''$ is generated by $q(x_n)$.
Set $N'=N\cap M'$ and $N''=q(N)$. The inclusion and restriction
of $q$ give the exact sequence
$$
0\longrightarrow N'\longrightarrow N\longrightarrow N''
\longrightarrow0.
$$
Both end modules are finite by induction. Choose generators
$u_1,\ldots,u_a$ of $N'$ and generators
$w_1,\ldots,w_b$ of $N''$, with lifts
$z_1,\ldots,z_b\in N$. If $z\in N$, write
$q(z)=\sum_{j=1}^b c_jw_j$. Then
$z-\sum_{j=1}^b c_jz_j\in N'$, so
$z=\sum_{i=1}^a d_i u_i+\sum_{j=1}^b c_jz_j$ for suitable
$d_i,c_j\in R$. This proves finite generation, including an
empty generating list at either end.
For a finite module $M$, choose a surjection $R^n\to M$.
Its kernel is a finite submodule of the finite module $R^n$.
Mapping a basis of $R^t$ to a finite generating list of this
kernel gives an exact presentation $R^t\to R^n\to M\to0$.

The last equivalence holds for every module $M$ over every ring
$R$, without a Noetherian or finite-generation hypothesis on
$R$ or $M$. If every submodule is finite and
$N_1\subset N_2\subset\cdots$ is an ascending chain,
its union $N$ is a submodule. Finitely many generators of $N$
all belong to one $N_j$, so $N=N_j=N_{j+1}=\cdots$.
Conversely, if a submodule $N$ is not finite, start with
$N_0=0$. Once $x_1,\ldots,x_j\in N$ have been chosen, their
span $N_j$ is not $N$, so choose $x_{j+1}\in N\setminus N_j$.
The strictly ascending chain $N_0\subsetneq N_1\subsetneq\cdots$
contradicts the ascending chain condition for submodules of $M$.
Thus that condition implies that every submodule is finite.
Applying the equivalence to the first part proves the final claim.
\end{proof}


\subsection{Editorial treatment of Algebra lines 12186--12218}

\hypertarget{algebra-context-159}{}

\noindent\href{https://github.com/stacks/stacks-project/blob/a04446e57ec1fbc252a871afcec7752fb2807b14/algebra.tex\#L12186}{Original source, lines 12186--12218}. Complete original and reviewed TeX passages accompany this note. Every intervention is separately recorded; the following treatment is editorial.


\begin{lemma}[Artin-Rees]
\label{editorial-source-lemma-Artin-Rees}
Suppose that $R$ is Noetherian. Let $I \subset R$ be an ideal.
Let $N \subset M$ be finite $R$-modules.
There exists a constant $c > 0$ such that
$I^n M \cap N  =  I^{n-c}(I^cM \cap N)$ for all $n \geq c$.
\end{lemma}

\begin{proof}
Keep the graded ring
$$
S=\bigoplus_{n\geq0}I^n,\qquad I^0=R,
$$
with its actual degree inclusions $\iota_n:I^n\to S$.
Thus $\iota_n(a)\iota_m(b)=\iota_{n+m}(ab)$.
Since $R$ is Noetherian, write $I=(f_1,\ldots,f_t)$.
The $R$-algebra map $R[X_1,\ldots,X_t]\to S$ sends $r\in R$
to $\iota_0(r)$ and sends a monomial with its coefficient to
$$
rX_1^{e_1}\cdots X_t^{e_t}
\longmapsto
\iota_{e_1+\cdots+e_t}
 (rf_1^{e_1}\cdots f_t^{e_t}).
$$
Every element of $I^n$ is a finite sum of these products of total
degree $n$, so the map is surjective. The polynomial ring, and
therefore $S$, is Noetherian by Lemma
\ref{editorial-source-lemma-Noetherian-permanence}. If $I=0$, the empty generator
list gives $S=R$ in degree zero and the same argument applies.

The graded $S$-module
$E=\bigoplus_{n\geq0}I^nM$ has action
$\iota_d(a)\jmath_n(x)=\jmath_{d+n}(ax)$, where $\jmath_n$
denotes its degree inclusion. If $x_1,\ldots,x_r$ generate
$M$, their degree-zero images generate $E$: each member of
$I^nM$ is a sum $\sum_{i=1}^r a_ix_i$ with $a_i\in I^n$.
The graded subgroup
$$
L=\bigoplus_{n\geq0}(I^nM\cap N)\subset E
$$
is an $S$-submodule, since
$I^d(I^nM\cap N)\subset I^{d+n}M\cap N$.
By Lemma \ref{editorial-source-lemma-Noetherian-basic}, $L$ is finite over $S$.
Take a finite generating list and replace it by the collection of
all its homogeneous components. Each component is in $L$, each
old generator is the sum of its components, and the generated
submodule is consequently still exactly $L$. Denote this finite
list by $\jmath_{d_j}(\xi_j)$ for $1\leq j\leq s$, where
$\xi_j\in I^{d_j}M\cap N$. The list may be empty when $L=0$.
Set
$$
c=\max\bigl(\{1\}\cup\{d_j:1\leq j\leq s\}\bigr)>0.
$$
For $n\geq c$, taking the degree-$n$ component of an expression
in these generators shows that every $x\in I^nM\cap N$ is
$$
x=\sum_{j=1}^s h_j\xi_j,\qquad h_j\in I^{n-d_j}.
$$
The degree inclusions above specify the graded expression from
which this equality in the original $M$ is obtained.
For each $j$, the ideal identity
$I^{n-d_j}=I^{n-c}I^{c-d_j}$ means that one can write
$h_j=\sum_{\ell=1}^{q_j}a_{j\ell}b_{j\ell}$ with
$a_{j\ell}\in I^{n-c}$ and $b_{j\ell}\in I^{c-d_j}$.
Each $b_{j\ell}\xi_j$ lies in $I^cM\cap N$, so the full
equality
$x=\sum_{j=1}^s\sum_{\ell=1}^{q_j}
a_{j\ell}(b_{j\ell}\xi_j)$ proves
$I^nM\cap N\subset I^{n-c}(I^cM\cap N)$.
The opposite inclusion holds because a product with a member of
$I^cM\cap N$ remains in $N$ and lies in $I^nM$.
If the generating list is empty then $L=0$, both modules in the
desired equality are zero, and the same choice $c=1$ is valid.

Every larger integer $d\geq c$ is also valid. Indeed, for $n\geq d$,
the equalities already proved at $n$ and at $d$ give
$$
I^nM\cap N=I^{n-c}(I^cM\cap N)
=I^{n-d}\bigl(I^{d-c}(I^cM\cap N)\bigr)
=I^{n-d}(I^dM\cap N).
$$
In particular, the two original filtrations on $N$ satisfy
$$
I^nN\subset I^nM\cap N\quad(n\geq0),\qquad
I^{k+c}M\cap N\subset I^kN\quad(k\geq0).
$$
The first containment follows from $N\subset M$; the second
uses the equality at $n=k+c$ and $I^cM\cap N\subset N$.
Thus the subspace $I$-adic topology from $M$ and the $I$-adic
topology on $N$ have cofinal neighbourhood systems, with the
displayed explicit index bound.
\end{proof}


\subsection{Editorial treatment of Algebra lines 12220--12238}

\hypertarget{algebra-context-160}{}

\noindent\href{https://github.com/stacks/stacks-project/blob/a04446e57ec1fbc252a871afcec7752fb2807b14/algebra.tex\#L12220}{Original source, lines 12220--12238}. Complete original and reviewed TeX passages accompany this note. Every intervention is separately recorded; the following treatment is editorial.


\begin{lemma}
\label{editorial-source-lemma-map-AR}
Suppose that $0 \to K \to M \xrightarrow{f} N$ is an
exact sequence of finitely generated modules
over a Noetherian ring $R$. Let $I \subset R$ be an ideal.
Then there exists a $c$ such that
$$
f^{-1}(I^nN) = K + I^{n-c}f^{-1}(I^cN)
\quad\text{and}\quad
f(M) \cap I^nN \subset f(I^{n - c}M)
$$
for all $n \geq c$.
\end{lemma}

\begin{proof}
Identify $K$ with its image in $M$, which is $\ker(f)$ by
exactness. The module $f(M)$ is finite over $R$.
Apply Lemma \ref{editorial-source-lemma-Artin-Rees} to $f(M)\subset N$, giving
an integer $c>0$ with
$$
f(M)\cap I^nN=I^{n-c}(f(M)\cap I^cN)\quad(n\geq c).
$$
Put $P=f^{-1}(I^cN)$. The restriction $f:P\to f(M)\cap I^cN$
is surjective: an element of the intersection already has an
inverse image in $M$, and every such inverse image lies in $P$.
For every integer $d\geq0$, linearity and lifting each member
of a finite sum give the exact equality
$f(I^dP)=I^d f(P)$.
Consequently, for $n\geq c$,
$$
f(M)\cap I^nN
=f\bigl(I^{n-c}P\bigr)
\subset f(I^{n-c}M).
$$

If $x\in f^{-1}(I^nN)$, the preceding equality expresses
$f(x)=\sum_{j=1}^t a_jf(p_j)$ for original coefficients
$a_j\in I^{n-c}$ and original elements $p_j\in P$.
Thus $x-\sum_{j=1}^t a_jp_j\in K$ and
$x\in K+I^{n-c}P$. Conversely, $f(K)=0$ and
$f(I^{n-c}P)\subset I^nN$, so every member of
$K+I^{n-c}P$ belongs to $f^{-1}(I^nN)$. This proves the
first equality as well as the second claimed containment.
The same maps give the more precise exact sequence
$$
0\longrightarrow K\cap I^{n-c}P
\longrightarrow I^{n-c}P
\xrightarrow{\ f\ }f(M)\cap I^nN
\longrightarrow0.
$$
Here the first arrow is the inclusion, its image is exactly the
kernel of the displayed restriction of the original $f$, and
surjectivity is the lifting calculation above. The kernel $K$
has not been assumed to vanish or omitted from the preimage formula.
\end{proof}


\subsection{Editorial treatment of Algebra lines 12240--12254}

\hypertarget{algebra-context-161}{}

\noindent\href{https://github.com/stacks/stacks-project/blob/a04446e57ec1fbc252a871afcec7752fb2807b14/algebra.tex\#L12240}{Original source, lines 12240--12254}. Complete original and reviewed TeX passages accompany this note. Every intervention is separately recorded; the following treatment is editorial.


\begin{lemma}[Krull's intersection theorem]
\label{editorial-source-lemma-intersect-powers-ideal-module-zero}
Let $R$ be a Noetherian local ring. Let $I \subset R$ be
a proper ideal. Let $M$ be a finite $R$-module.
Then $\bigcap_{n \geq 0} I^nM = 0$.
\end{lemma}

\begin{proof}
Let $N=\bigcap_{n\geq0}I^nM$. It is a submodule of the finite
module $M$, hence is finite by Lemma \ref{editorial-source-lemma-Noetherian-basic}.
For every $n\geq0$ one has $I^nM\cap N=N$.
Apply Lemma \ref{editorial-source-lemma-Artin-Rees} to $N\subset M$ and take
$n=c+1$, where $c>0$ is the integer it supplies. This gives
$$
N=I^{c+1}M\cap N=I(I^cM\cap N)=IN.
$$
The proper ideal $I$ lies in the maximal ideal $\mathfrak m$
of the local ring $R$. Hence
$N=IN\subset\mathfrak mN\subset N$, so
$N=\mathfrak mN$. Nakayama's Lemma \ref{algebra-lemma-NAK}, applied
to the finite module $N$, gives $N=0$.
\end{proof}


\subsection{Editorial treatment of Algebra lines 12256--12280}

\hypertarget{algebra-context-162}{}

\noindent\href{https://github.com/stacks/stacks-project/blob/a04446e57ec1fbc252a871afcec7752fb2807b14/algebra.tex\#L12256}{Original source, lines 12256--12280}. Complete original and reviewed TeX passages accompany this note. Every intervention is separately recorded; the following treatment is editorial.


\begin{lemma}
\label{editorial-source-lemma-intersection-powers-ideal-module}
Let $R$ be a Noetherian ring. Let $I \subset R$ be an ideal.
Let $M$ be a finite $R$-module. Let $N = \bigcap_n I^n M$.
\begin{enumerate}
\item For every prime $\mathfrak p$ containing $I$, there
exists $f \in R$, $f \not \in \mathfrak p$ such that $N_f = 0$.
\item If $I$ is contained in the Jacobson radical
of $R$, then $N = 0$.
\end{enumerate}
\end{lemma}

\begin{proof}
In fact there is one element $f\in1+I$ such that $fN=0$.
We prove this without a local hypothesis. The module $N$ is finite
by Lemma \ref{editorial-source-lemma-Noetherian-basic}. Applying
Lemma \ref{editorial-source-lemma-Artin-Rees} to $N\subset M$ at $n=c+1$
gives $N=IN$, since $I^nM\cap N=N$ for every $n\geq0$.
Choose generators $x_1,\ldots,x_t$ of $N$. Because $N=IN$,
there are actual coefficients $a_{ij}\in I$ with
$x_i=\sum_{j=1}^t a_{ij}x_j$ for $1\leq i\leq t$.
Put $A=(a_{ij})$ and $B=(\delta_{ij}-a_{ij})$.
The adjugate identity $\operatorname{adj}(B)B=\det(B)\,1_t$
applied to the column $(x_1,\ldots,x_t)^{\mathrm t}$ shows
that $f=\det(B)$ annihilates every $x_i$, hence annihilates $N$.
Retaining every coefficient and sign, this element is
$$
f=\sum_{\sigma\in\mathfrak S_t}\operatorname{sgn}(\sigma)
       \prod_{i=1}^t(\delta_{i,\sigma(i)}-a_{i,\sigma(i)}).
$$
Modulo $I$ the matrix $B$ is the identity matrix, so $f-1\in I$.
Equivalently $f=1-a$ for the specific element $a=1-f\in I$;
the determinant expression above records all contributions to $a$.
For an empty generating list, $N=0$ and the determinant of the
empty matrix is $f=1$, which has all the same properties.

For every prime $\mathfrak p$ containing $I$, the equality
$f\equiv1\pmod I$ implies $f\notin\mathfrak p$.
Since $fN=0$, inverting $f$ gives $N_f=0$. This proves (1)
with the same $f$ for all such primes. If $I$ is contained in the
Jacobson radical, $f=1-a$ is a unit: otherwise the proper ideal
$(f)$ would lie in a maximal ideal $\mathfrak m$, while
$a\in I\subset\mathfrak m$, contradicting $f+a=1$.
Thus $fN=0$ implies $N=0$, proving (2). If $R$ is the zero
ring then $M=N=0$ and both assertions hold directly.

The argument gives two exact descriptions of the same original
intersection. Let $W=1+I=\{1+b:b\in I\}$, a multiplicative
set because $(1+b)(1+d)=1+(b+d+bd)$ and $b+d+bd\in I$.
Then
$$
N=\{x\in M:(1-a)x=0\text{ for some }a\in I\}
=\ker\bigl(M\longrightarrow W^{-1}M\bigr).
$$
The inclusion from $N$ to the first set follows from the single
annihilator just constructed. Conversely, if $(1-a)x=0$,
then $x=ax$, and induction gives $x=a^nx\in I^nM$ for
every $n\geq0$. The localization kernel is exactly the set of
$x$ killed by some $w\in W$; writing $w=1-a$ gives the
second equality. This reasoning includes $I=R$, when $0\in W$,
$W^{-1}M=0$ and $N=M$, as well as $I=0$, when $W=\{1\}$
and $N=0$.

Moreover, $N$ is the largest submodule $L\subset M$ satisfying
$IL=L$. Indeed this equality implies $I^nL=L\subset I^nM$
for every $n$, so $L\subset N$, while $IN=N$ was proved above.
The quotient $M/N$ is $I$-adically separated:
$$
\bigcap_{n\geq0} I^n(M/N)=0.
$$
To prove this using the original quotient map $q:M\to M/N$,
note that $I^n(M/N)=q(I^nM)$ and $N\subset I^nM$.
Thus $q^{-1}(I^n(M/N))=I^nM+N=I^nM$.
A class belonging to every such image has any representative in
$\bigcap_n I^nM=N$ and is therefore zero. Finally, for every
$R$-linear map $h:M\to Q$ into an $I$-adically separated
module $Q$, one has $h(N)\subset\bigcap_n I^nQ=0$.
The original map $h$ therefore factors uniquely through $q$,
by $x+N\mapsto h(x)$. This proves the universal property of
this separated quotient with no finite-generation hypothesis on $Q$.
\end{proof}


\subsection{Editorial treatment of Algebra lines 12282--12294}

\hypertarget{algebra-context-163}{}

\noindent\href{https://github.com/stacks/stacks-project/blob/a04446e57ec1fbc252a871afcec7752fb2807b14/algebra.tex\#L12282}{Original source, lines 12282--12294}. Complete original and reviewed TeX passages accompany this note. Every intervention is separately recorded; the following treatment is editorial.


\begin{remark}
\label{editorial-source-remark-intersection-powers-ideal}
Lemma \ref{editorial-source-lemma-intersect-powers-ideal-module-zero} in particular
gives $\bigcap_{n\geq0}I^n=(0)$ for a proper ideal $I$ in a
Noetherian local ring. More generally let $R$ be Noetherian,
let $I\subset R$ be an ideal, and retain an original element
$f\in\bigcap_{n\geq0}I^n$. Apply the proof of
Lemma \ref{editorial-source-lemma-intersection-powers-ideal-module} with $M=R$.
It gives a single $g\in1+I$ annihilating the whole ideal
$\bigcap_{n\geq0}I^n$, so in particular $gf=0$.
Every prime containing $I$ avoids $g$, hence
$V(I)\subset D(g)$. In $R_g$, the equality $gf=0$ implies
$f/1=0$. Thus the original $f$, and indeed every member of the
intersection, is zero on the same principal open neighbourhood
$D(g)$ of $V(I)$, with the coefficient condition $g\equiv1\pmod I$.
When $I=R$, the closed set $V(I)$ is empty and $g=0$ is
allowed; no nonempty neighbourhood or nonzero denominator is
asserted in that case.
\end{remark}


\subsection{Editorial treatment of Algebra lines 12296--12323}

\hypertarget{algebra-context-164}{}

\noindent\href{https://github.com/stacks/stacks-project/blob/a04446e57ec1fbc252a871afcec7752fb2807b14/algebra.tex\#L12296}{Original source, lines 12296--12323}. Complete original and reviewed TeX passages accompany this note. Every intervention is separately recorded; the following treatment is editorial.


\begin{lemma}[Artin-Tate]
\label{editorial-source-lemma-Artin-Tate}
Let $R$ be a Noetherian ring. Let $S$ be a finitely
generated $R$-algebra. If $T \subset S$ is an $R$-subalgebra such
that $S$ is finitely generated as a $T$-module, then $T$ is of
finite type over $R$.
\end{lemma}

\begin{proof}
If $S=0$, then its subalgebra $T=0$ is the image of $R$ and
is of finite type over $R$. We may therefore assume $S\neq0$.
Choose original algebra generators $x_1,\ldots,x_n$ of $S$
over $R$ and original module generators $y_1,\ldots,y_m$ of
$S$ over $T$. Here $n$ may be zero, while $m\geq1$.
Choose coefficients $a_{ij},b_{ijk}\in T$ giving the full equations
$$
x_i=\sum_{j=1}^m a_{ij}y_j\quad(1\leq i\leq n),\qquad
y_i y_j=\sum_{k=1}^m b_{ijk}y_k\quad(1\leq i,j\leq m).
$$
Let $T'\subset T$ be the $R$-subalgebra generated by all these
coefficients. It is a quotient of a polynomial ring in finitely
many variables over the given $R$, and is Noetherian by
Lemma \ref{editorial-source-lemma-Noetherian-permanence}. This does not require
the structure map $R\to T'$ to be injective.
Form the algebra with every indicated relation
$$
S'=T'[Y_1,\ldots,Y_m]\big/
\left(Y_iY_j-\sum_{k=1}^m b_{ijk}Y_k:
1\leq i,j\leq m\right).
$$
Its underlying $T'$-module is generated by the classes of
$1,Y_1,\ldots,Y_m$. Indeed a monomial of total degree at least
two has a factor $Y_iY_j$. Replacing that pair by
$\sum_{k=1}^m b_{ijk}Y_k$ expresses it as a sum of monomials
of total degree one less, with their actual coefficients
$b_{ijk}$. Induction on total degree expresses every monomial,
and then every polynomial, in the asserted span. This argument
retains the constant generator $1$ and does not assume that one
of the chosen $y_j$ equals $1$.

The coefficient inclusion $T'\subset S$ and the assignments
$Y_j\mapsto y_j$ define a $T'$-algebra map $\rho:S'\to S$:
every displayed relation maps to its given zero expression in $S$.
For each $i$, the element $\sum_{j=1}^m a_{ij}Y_j$ maps to
the original $x_i$. The image also contains the image of $R$,
so it contains all of $S$, even if the algebra-generator list is
empty. Therefore $\rho$ is surjective and $S$ is finite over
$T'$. No assertion that $\rho$ is injective is needed or made.
Now $T\subset S$ is a $T'$-submodule. By
Lemma \ref{editorial-source-lemma-Noetherian-basic}, choose a finite generating
list $z_1,\ldots,z_q$ for this module. The exact equality
$$
T=R[a_{ij},b_{ijk},z_1,\ldots,z_q]
$$
uses all $a_{ij}$ for $1\leq i\leq n$, $1\leq j\leq m$,
all $b_{ijk}$ for $1\leq i,j,k\leq m$, and every $z_\ell$
for $1\leq\ell\leq q$, together with the image of $R$.
To verify it, every $t\in T$ is a sum
$\sum_{\ell=1}^q c_\ell z_\ell$ with $c_\ell\in T'$;
each $c_\ell$ is a polynomial in the original $a_{ij},b_{ijk}$
with coefficients from $R$. This gives containment in the
right-hand subalgebra. The reverse containment follows because
all its generators already lie in $T$. Hence $T$ is of finite
type over $R$, as required.
\end{proof}


\subsection{Editorial treatment of Algebra lines 12346--12699}

\hypertarget{algebra-context-165}{}

\noindent\href{https://github.com/stacks/stacks-project/blob/a04446e57ec1fbc252a871afcec7752fb2807b14/algebra.tex\#L12346}{Original source, lines 12346--12699}. Complete original and reviewed TeX passages accompany this note. Every intervention is separately recorded; the following treatment is editorial.


\subsubsection{Length}
\label{editorial-source-section-length}

\begin{definition}
\label{editorial-source-definition-length}
Let $R$ be a ring. For any $R$-module $M$
we define the {\it length} of $M$ over $R$ by the
formula
$$
\text{length}_R(M)
=
\sup
\{
n \in \mathbf Z_{\geq0}
\mid
\exists\ 0 = M_0 \subset M_1 \subset \ldots \subset M_n = M,
\text{ }M_i \not = M_{i + 1}\quad(0\leq i<n)
\}.
$$
\end{definition}

\noindent
The supremum is taken in $\mathbf Z_{\geq0}\cup\{\infty\}$.
For $M=0$ the length-zero chain is allowed, so its length is zero.
In other words it is the supremum of the lengths of chains
of submodules. There is an obvious notion of when a chain
of submodules is a refinement of another. This gives a
partial ordering on the collection of all chains of submodules,
with the smallest chain having the shape $0 = M_0 \subset M_1 = M$
if $M$ is not zero.
We note the obvious fact that if the length of
$M$ is finite, then every chain can be refined to a
maximal chain. But it is not as obvious that all maximal
chains have the same length (as we will see later).

\begin{lemma}
\label{editorial-source-lemma-finite-length-finite}
\begin{slogan}
Modules of finite length are finite.
\end{slogan}
Let $R$ be a ring.
Let $M$ be an $R$-module.
If $\text{length}_R(M) < \infty$ then $M$ is a finite $R$-module.
\end{lemma}

\begin{proof}
Write $\ell=\operatorname{length}_R(M)<\infty$.
If $M=0$, its empty generating list suffices. Otherwise start with
$N_0=0$. Whenever $N_j\neq M$, choose
$x_{j+1}\in M\setminus N_j$ and put
$N_{j+1}=N_j+Rx_{j+1}$. Each such step is strict.
There cannot be $\ell+1$ such steps: adjoining the final endpoint
$M$ if it is not already present would give a chain from zero
to $M$ with at least $\ell+1$ strict steps, contrary to the
definition of length. Therefore $N_j=M$ for some $j\leq\ell$,
and $x_1,\ldots,x_j$ generate the original $M$.

The same bound applies to every subquotient $N/P$ of $M$.
Indeed a strict chain in $N/P$ pulls back under the original
quotient map $N\to N/P$ to a strict chain from $P$ to $N$.
Adjoin zero if $P\neq0$ and $M$ if $N\neq M$. This is a
chain of submodules of $M$, so the original chain has at most
$\ell$ strict steps. Thus $N/P$ has length at most $\ell$
and is generated by at most its own length many elements by the
same construction. Moreover, no ascending or descending chain
of submodules of $M$ can have more than $\ell$ strict changes:
any finite selection of its strict changes, ordered increasingly
and completed with the endpoints, gives such a chain in $M$.
In particular both chain conditions hold for $M$.
\end{proof}

\begin{lemma}
\label{editorial-source-lemma-length-additive}
\begin{slogan}
Length is additive in short exact sequences.
\end{slogan}
If $0 \to M' \to M \to M'' \to 0$
is a short exact sequence of modules over $R$ then
the length of $M$ is the sum of the
lengths of $M'$ and $M''$.
\end{lemma}

\begin{proof}
Write the original sequence as
$$
0\longrightarrow M'\xrightarrow{u}M\xrightarrow{q}M''
\longrightarrow0.
$$
Lengths take values in $\mathbf Z_{\geq0}\cup\{\infty\}$,
with $a+\infty=\infty+a=\infty$ for every such $a$.
Given strict chains from zero to $M'$ and $M''$, of lengths
$a$ and $b$, respectively, first apply $u$ to the chain in
$M'$. Its last member is $u(M')=\ker(q)$.
Continue with the inverse images under $q$ of the nonzero steps
of the chain in $M''$. Injectivity of $u$ and surjectivity of
$q$ show that every indicated step is strict. This yields a
chain in $M$ of length $a+b$. If either length on the right is
infinite, finite chains of arbitrarily large length give
$\operatorname{length}_R(M)=\infty$. If both are finite, each
is an attained integer supremum, so choosing chains with those
lengths proves
$\operatorname{length}_R(M)\geq
\operatorname{length}_R(M')+\operatorname{length}_R(M'')$.
This construction includes a zero module with its length-zero chain.

Conversely, let $0=M_0\subsetneq\cdots\subsetneq M_n=M$
be a finite strict chain. Put
$M_i'=u^{-1}(M_i)$ and $M_i''=q(M_i)$, using the original
maps rather than silently identifying $M'$ with a different module.
Let $a$ and $b$ count only the strict changes in these two
nondecreasing sequences. Removing their repeated terms gives
chains of lengths $a$ and $b$ in $M'$ and $M''$.
At least one changes at each of the $n$ original steps.
For if $M_i'=M_{i+1}'$ and $M_i''=M_{i+1}''$, take
$x\in M_{i+1}$. There is $y\in M_i$ with $q(x)=q(y)$.
Then $x-y=u(z)$ for some $z\in M'$ and
$u(z)\in M_{i+1}$, so $z\in M_{i+1}'=M_i'$.
Consequently $u(z)\in M_i$ and $x=y+u(z)\in M_i$,
contradicting strictness. It follows that
$$
n\leq a+b\leq
\operatorname{length}_R(M')+\operatorname{length}_R(M'').
$$
Taking the supremum over the original chains proves the reverse
inequality, including infinite lengths. This proves additivity.
\end{proof}

\begin{lemma}
\label{editorial-source-lemma-length-infinite}
Let $R$ be a local ring with maximal ideal $\mathfrak m$.
If $M$ is an $R$-module and $\mathfrak m^n M \not = 0$ for all $n \geq 0$,
then $\text{length}_R(M) = \infty$. In other words, if $M$
has finite length then $\mathfrak m^nM = 0$ for some $n$.
\end{lemma}

\begin{proof}
Fix an integer $n\geq0$ such that $\mathfrak m^nM\neq0$.
A nonzero member of this module is a finite sum of products
$f_1\cdots f_nx$, with every $f_i\in\mathfrak m$ and $x\in M$.
At least one of these products is nonzero; choose it and put
$y_j=f_1\cdots f_jx$ for $0\leq j\leq n$, where $y_0=x$.
All $y_j$ are nonzero, since multiplying a zero earlier product
by the remaining original factors would give $y_n=0$.
The inclusions
$$
0\subsetneq Ry_n\subsetneq Ry_{n-1}\subsetneq\cdots
\subsetneq Ry_0\subset M
$$
are strict up to $Ry_0$. To verify any step with $1\leq j\leq n$,
suppose $Ry_{j-1}=Ry_j$. Then
$y_{j-1}=g y_j=g f_j y_{j-1}$ for some $g\in R$, hence
$(1-gf_j)y_{j-1}=0$. Since $gf_j\in\mathfrak m$, the original
factor $1-gf_j$ is a unit in the local ring. This forces
$y_{j-1}=0$, a contradiction. The first inclusion is strict
because $y_n\neq0$. Append $M$ if necessary; if $Ry_0=M$,
there is no extra final step. Thus
$\operatorname{length}_R(M)\geq n+1$.
For $n=0$ this uses the nonzero element $x$ and the empty product,
and gives the same bound.

If every $\mathfrak m^nM$ is nonzero, these bounds for all $n$
give infinite length. They also give the sharper quantitative
conclusion
$$
\mathfrak m^\ell M=0
\quad\text{if}\quad
\ell=\operatorname{length}_R(M)<\infty.
$$
Indeed a nonzero left-hand module would imply
$\ell\geq\ell+1$. When $\ell=0$, the definition of length
already gives $M=0$, so the displayed assertion remains valid.
\end{proof}

\begin{lemma}
\label{editorial-source-lemma-length-independent}
Let $R \to S$ be a ring map. Let $M$ be an $S$-module.
We always have $\text{length}_R(M) \geq \text{length}_S(M)$.
If $R \to S$ is surjective then equality holds.
\end{lemma}

\begin{proof}
Every $S$-submodule of the original $M$ is an $R$-submodule by
restriction of scalars. Hence every strict chain for the former
action gives the identical strict chain for the latter, proving
the inequality even for infinite length.
If $R\to S$ is surjective, then for any $R$-submodule $N\subset M$,
$s\in S$ and $x\in N$, choose $r\in R$ mapping to $s$.
The equality $sx=rx$ puts $sx$ in $N$. Thus the two collections
of submodules, and all their chains, agree exactly.

The proof gives a weaker sufficient hypothesis for equality:
surjectivity of the original composite
$R\to S/\operatorname{Ann}_S(M)$ suffices.
For then any $s\in S$ has an $r\in R$ such that their images
in this quotient agree; their difference annihilates every
original $x\in M$, so once again $sx=rx$.
The original actions are retained and their equality on each
element proves equality of the submodule lattices. When $M=0$,
both lattices have a single member and both lengths are zero.
\end{proof}

\begin{lemma}
\label{editorial-source-lemma-dimension-is-length}
Let $R$ be a ring with maximal ideal $\mathfrak m$.
Suppose that $M$ is an $R$-module with
$\mathfrak m M  =  0$. Then the length of $M$ as
an $R$-module agrees with its dimension as a vector space over
$R/\mathfrak m$, interpreting every infinite dimension as $\infty$.
The length is finite if and only if $M$ is a finite $R$-module.
\end{lemma}

\begin{proof}
Set $k=R/\mathfrak m$. The original action on $M$ factors
through $k$, since $\mathfrak mM=0$; explicitly
$(r+\mathfrak m)x=rx$, which is independent of the representative.
An $R$-submodule is exactly a $k$-vector subspace by
Lemma \ref{editorial-source-lemma-length-independent}.
If $\dim_k M=d<\infty$, a strict inclusion of vector subspaces
increases dimension by at least one. Thus every strict chain has
at most $d$ steps. Given a basis $x_1,\ldots,x_d$, the spans
of its first $i$ vectors, for $0\leq i\leq d$, give a chain
of exactly $d$ steps. Therefore the length is $d$, including $d=0$.
If the dimension is infinite, for each finite $n$ choose
$n$ independent vectors. Their initial spans form $n$ strict
steps, and adjoining $M$ as endpoint if needed gives length at
least $n$. The length is therefore $\infty$; it is this extended
integer value, not the particular infinite basis cardinal, that
is asserted here.
Finally, a finite $R$-generating list is exactly a finite
$k$-spanning list because the two actions just described have the
same spans. Such a list contains a basis after removing dependent
vectors, and a finite basis is itself a finite $R$-generating
list. This proves the final equivalence.
\end{proof}

\begin{lemma}
\label{editorial-source-lemma-length-localize}
Let $R$ be a ring. Let $M$ be an $R$-module. Let $S \subset R$ be
a multiplicative subset. Then
$\text{length}_R(M) \geq \text{length}_{S^{-1}R}(S^{-1}M)$.
\end{lemma}

\begin{proof}
Let $j:M\to S^{-1}M$ be the original map $x\mapsto x/1$.
For an $S^{-1}R$-submodule $L\subset S^{-1}M$, put
$j^{-1}(L)=\{x\in M:x/1\in L\}$. Its localization has image
exactly $L$. Indeed if $x/s\in L$, then $x/1=s(x/s)\in L$,
so $x\in j^{-1}(L)$; conversely $x/1\in L$ implies
$x/s=(1/s)(x/1)\in L$. Localization of the inclusion
$j^{-1}(L)\subset M$ is injective, so this identifies the actual
localized submodule with $L$.
Consequently a strict chain $0=L_0\subsetneq\cdots\subsetneq
L_n=S^{-1}M$ pulls back to a strict chain
$j^{-1}(L_0)\subsetneq\cdots\subsetneq j^{-1}(L_n)=M$:
equality of two adjacent inverse images would give equality of
their localized images. The first inverse image is $\ker(j)$,
which need not be zero. Adjoin zero when it is nonzero. This
produces at least $n$ strict steps from zero to $M$, proving
the length inequality. If $S^{-1}M=0$, its length is zero and
the inequality holds directly; no nonzero localization is assumed.
For finite length, the proof of Lemma \ref{editorial-source-lemma-simple-pieces}
below computes the exact loss: it is the number of original
composition factors whose maximal ideals meet $S$, counted
with their multiplicities.
\end{proof}

\begin{lemma}
\label{editorial-source-lemma-length-finite}
Let $R$ be a ring with finitely generated
maximal ideal $\mathfrak m$. (For example, this holds if $R$ is Noetherian.)
Suppose that $M$ is a finite $R$-module with
$\mathfrak m^n M  =  0$ for some integer $n\geq0$.
Then $\text{length}_R(M) < \infty$.
\end{lemma}

\begin{proof}
If $n=0$, then $M=\mathfrak m^0M=0$ and the claim follows.
Assume $n\geq1$ and write $k=R/\mathfrak m$.
Choose original generators $x_1,\ldots,x_r$ of $M$ and
$f_1,\ldots,f_t$ of $\mathfrak m$. For $i\geq0$, put
$$
E_i=\{(e_1,\ldots,e_t)\in\mathbf Z_{\geq0}^{,t}:
e_1+\cdots+e_t=i\}.
$$
The quotient $\mathfrak m^iM/\mathfrak m^{i+1}M$ is killed
by $\mathfrak m$ and has the following explicit surjection
of $k$-vector spaces:
$$
\begin{aligned}
k^{\{1,\ldots,r\}\times E_i}&\longrightarrow
\mathfrak m^iM/\mathfrak m^{i+1}M,\\
(\bar a_{j,e})_{j,e}&\longmapsto
\sum_{j=1}^r\sum_{e\in E_i}
a_{j,e} f_1^{e_1}\cdots f_t^{e_t}x_j
\pmod{\mathfrak m^{i+1}M}.
\end{aligned}
$$
Changing any representative $a_{j,e}$ by an element of
$\mathfrak m$ changes the indicated summand by a member of
$\mathfrak m^{i+1}M$, so the map is well defined. The products
in the specified ideal generators generate $\mathfrak m^i$;
expanding a member of $\mathfrak m^iM$ in these products and
the original $x_j$ proves surjectivity. Thus each layer has
finite dimension over $k$, and its length is that dimension by
Lemma \ref{editorial-source-lemma-dimension-is-length}. Applying additivity
successively to the original filtration by powers gives the exact formula
$$
\operatorname{length}_R(M)
=\sum_{i=0}^{n-1}\dim_k
 (\mathfrak m^iM/\mathfrak m^{i+1}M)
\leq r\sum_{i=0}^{n-1}\#E_i<\infty.
$$
Repeated adjacent powers contribute their actual zero-dimensional
quotients. If $t\geq1$, distributing $i$ positions among $t$
nonnegative exponents gives $\#E_i=\binom{t+i-1}{i}$: encode
the tuple by $i$ identical positions separated into $t$ blocks
by $t-1$ separators, choosing the $i$ positions among the
$t+i-1$ total positions. If $t=0$, there is one empty tuple
in $E_0$ and no tuple in $E_i$ for $i>0$.

In fact finite generation of $\mathfrak m$ can be replaced by
finite dimension of the original cotangent space
$\mathfrak m/\mathfrak m^2$ over $k$.
Choose lifts $f_1,\ldots,f_t\in\mathfrak m$ of a finite
spanning list and put $J=(f_1,\ldots,f_t)$.
Then $\mathfrak m=J+\mathfrak m^2$.
For $i\geq1$, expanding this actual ideal equality retains all terms:
$$
\mathfrak m^i
=(J+\mathfrak m^2)^i
=\sum_{h=0}^i J^{i-h}\mathfrak m^{2h}.
$$
Every term with $h\geq1$ lies in $\mathfrak m^{i+h}$,
hence in $\mathfrak m^{i+1}$. Conversely both $J^i$ and
$\mathfrak m^{i+1}$ lie in $\mathfrak m^i$. Therefore
$\mathfrak m^i=J^i+\mathfrak m^{i+1}$ and
$\mathfrak m^iM=J^iM+\mathfrak m^{i+1}M$.
All the terms with $h\geq1$ have zero image under this particular
quotient map; no equality $\mathfrak m=J$ has been assumed.
The same displayed coefficient map is consequently surjective
for every layer, with these original lifts $f_j$.
It proves the same finite length conclusion and the same explicit
bound whenever $M$ is finite and $\mathfrak m^nM=0$.
This argument covers $t=0$ as well: then
$\mathfrak m^iM=\mathfrak m^{i+1}M$ for $i\geq1$, so
nilpotence on $M$ forces all those layers and $\mathfrak mM$ to vanish.
\end{proof}

\begin{definition}
\label{editorial-source-definition-simple-module}
Let $R$ be a ring. Let $M$ be an $R$-module.
We say $M$ is {\it simple} if $M \not = 0$ and
every submodule of $M$ is either equal to $M$ or
to $0$.
\end{definition}

\begin{lemma}
\label{editorial-source-lemma-characterize-length-1}
Let $R$ be a ring. Let $M$ be an $R$-module.
The following are equivalent:
\begin{enumerate}
\item $M$ is simple,
\item $\text{length}_R(M) = 1$, and
\item $M \cong R/\mathfrak m$ for some maximal ideal
$\mathfrak m \subset R$.
\end{enumerate}
\end{lemma}

\begin{proof}
Let $\mathfrak m$ be a maximal ideal of $R$.
By Lemma \ref{editorial-source-lemma-dimension-is-length} the module
$R/\mathfrak m$ has length $1$. The equivalence of
the first two assertions is tautological.
Suppose that $M$ is simple. Choose $x \in M$, $x \not = 0$.
As $M$ is simple we have $M = R \cdot x$.
Let $I \subset R$ be the annihilator of $x$, i.e.,
$I = \{f \in R \mid fx = 0\}$. The map $R/I \to M$,
$f \bmod I \mapsto fx$ is an isomorphism, hence
$R/I$ is a simple $R$-module. Since $R/I \not = 0$ we see $I \not = R$.
Let $\mathfrak m$ be a maximal ideal containing $I$.
If $I \not = \mathfrak m$, then $\mathfrak m /I \subset R/I$
is a nontrivial submodule contradicting the simplicity
of $R/I$. Hence we see $I = \mathfrak m$ as desired.
\end{proof}

\begin{lemma}
\label{editorial-source-lemma-simple-pieces}
Let $R$ be a ring. Let $M$ be a finite length $R$-module.
Choose any maximal chain of submodules
$$
0 = M_0 \subset M_1 \subset M_2 \subset \ldots \subset M_n = M
$$
with $M_i \not = M_{i-1}$, $i = 1, \ldots, n$. Then
\begin{enumerate}
\item $n = \text{length}_R(M)$,
\item each $M_i/M_{i-1}$ is simple,
\item each $M_i/M_{i-1}$ is of the form
$R/\mathfrak m_i$ for some maximal ideal $\mathfrak m_i$,
\item given a maximal ideal $\mathfrak m \subset R$
we have
$$
\# \{i \mid \mathfrak m_i = \mathfrak m\}
=
\text{length}_{R_{\mathfrak m}} (M_{\mathfrak m}).
$$
\end{enumerate}
\end{lemma}

\begin{proof}
For each $i$, let $q_i:M_i\to M_i/M_{i-1}$ be the original
quotient map. If its nonzero target were not simple, a nonzero
proper submodule would have inverse image strictly between
$M_{i-1}$ and $M_i$. Inserting it would contradict maximality
of the given chain. Thus every quotient is simple, and
Lemma \ref{editorial-source-lemma-characterize-length-1} gives an isomorphism
$M_i/M_{i-1}\cong R/\mathfrak m_i$ and length one.
Applying Lemma \ref{editorial-source-lemma-length-additive} to the successive
short exact sequences gives $\operatorname{length}_R(M)=n$.
For $M=0$ the chain has no strict steps, the lists of quotients
and maximal ideals are empty, and this equality reads $0=0$.

For a maximal ideal $\mathfrak m$, localization gives the exact sequences
$$
0\longrightarrow(M_{i-1})_{\mathfrak m}
\longrightarrow(M_i)_{\mathfrak m}
\longrightarrow(M_i/M_{i-1})_{\mathfrak m}
\longrightarrow0.
$$
If $\mathfrak m_i\neq\mathfrak m$, choose
$a\in\mathfrak m_i\setminus\mathfrak m$; such an element exists
because distinct maximal ideals cannot contain one another.
It annihilates $R/\mathfrak m_i$ and is inverted at
$\mathfrak m$, so this quotient localizes to zero.
If $\mathfrak m_i=\mathfrak m$, the isomorphism with the
residue field of the original local ring is explicitly
$$
(R/\mathfrak m)_{\mathfrak m}
\longrightarrow R_{\mathfrak m}/\mathfrak mR_{\mathfrak m},
\qquad
(r+\mathfrak m)/s\longmapsto r/s+\mathfrak mR_{\mathfrak m}.
$$
Its inverse sends $r/s+\mathfrak mR_{\mathfrak m}$ to
$(r+\mathfrak m)/s$. The localization fraction relations hold
on both sides, and coefficients in $\mathfrak m$ act as zero,
so both maps are well defined; their composites are identities.
The resulting nonzero quotient has length one over $R_{\mathfrak m}$.
Additivity proves exactly the asserted count.

More generally let $S\subset R$ be any multiplicative subset.
For each factor $R/\mathfrak m_i$, its localization is zero
if $S\cap\mathfrak m_i\neq\varnothing$. Otherwise every
$s\in S$ has nonzero, hence invertible, image in the field
$R/\mathfrak m_i$, and
$$
S^{-1}(R/\mathfrak m_i)\longrightarrow R/\mathfrak m_i,
\qquad (r+\mathfrak m_i)/s\longmapsto
(r+\mathfrak m_i)(s+\mathfrak m_i)^{-1}
$$
is an isomorphism, with inverse $z\mapsto z/1$.
It identifies this localized module with the simple quotient
$S^{-1}R/S^{-1}\mathfrak m_i$.
Localizing every short exact sequence and adding the lengths gives
$$
\operatorname{length}_{S^{-1}R}(S^{-1}M)
=\#\{i: S\cap\mathfrak m_i=\varnothing\},
$$
and therefore the exact loss of length is
$\#\{i:S\cap\mathfrak m_i\neq\varnothing\}$.
This refines Lemma \ref{editorial-source-lemma-length-localize} for finite length.
For an arbitrary prime $\mathfrak p$, the same calculation with
$S=R\setminus\mathfrak p$ retains a factor precisely when
$\mathfrak m_i\subset\mathfrak p$, hence precisely when
$\mathfrak p=\mathfrak m_i$.
Thus the original support of $M$ is the finite set of distinct
$\mathfrak m_i$, and
$$
\operatorname{length}_R(M)
=\sum_{\mathfrak m\in\operatorname{Supp}_R(M)}
  \operatorname{length}_{R_{\mathfrak m}}(M_{\mathfrak m}).
$$
In particular the multiplicity of each simple module
$R/\mathfrak m$ is independent of the chosen maximal chain.
The ideal $\mathfrak m$ is determined by that module, since
its annihilator is exactly $\mathfrak m$; hence these counts
determine the multiset of simple factors, not only the number
of steps.

The two chain conditions also characterize finite length for
an arbitrary $R$-module. Necessity was proved in the proof of
Lemma \ref{editorial-source-lemma-finite-length-finite}.
Conversely assume both conditions. In every nonzero submodule
$L$, the ascending chain condition supplies a maximal proper
submodule: if none existed, starting from zero one could choose
an infinite strictly increasing sequence of proper submodules
of $L$. Successively choose such a maximal proper submodule
inside each nonzero module, starting from $M$.
The descending chain condition forces this sequence to reach
zero after finitely many strict steps. Reversing it gives a
finite chain whose quotients are simple. Each quotient has
length one by Lemma \ref{editorial-source-lemma-characterize-length-1}, so
additivity gives finite length for the original $M$.
\end{proof}

\begin{lemma}
\label{editorial-source-lemma-pushdown-module}
Let $A$ be a local ring with maximal ideal $\mathfrak m$.
Let $B$ be a semi-local ring with maximal ideals $\mathfrak m_i$,
$i = 1, \ldots, n$.
Suppose that $A \to B$ is a homomorphism such that each $\mathfrak m_i$
lies over $\mathfrak m$ and such that
$$
[\kappa(\mathfrak m_i) : \kappa(\mathfrak m)] < \infty.
$$
Let $M$ be a $B$-module of finite length.
Then
$$
\text{length}_A(M) = \sum\nolimits_{i = 1, \ldots, n}
[\kappa(\mathfrak m_i) : \kappa(\mathfrak m)]
\text{length}_{B_{\mathfrak m_i}}(M_{\mathfrak m_i}),
$$
in particular $\text{length}_A(M) < \infty$.
\end{lemma}

\begin{proof}
Choose a maximal chain of $B$-submodules of the original $M$
as in Lemma \ref{editorial-source-lemma-simple-pieces}. Each of its quotients is
isomorphic to $B/\mathfrak m_{i(j)}=\kappa(\mathfrak m_{i(j)})$.
The original map $A\to B$ induces an injective field map
$\kappa(\mathfrak m)=A/\mathfrak m\to B/\mathfrak m_i$:
its kernel before quotienting is exactly the contraction
$\mathfrak m_i\cap A=\mathfrak m$.
The $A$-action on this quotient is the action through that map.
By Lemma \ref{editorial-source-lemma-dimension-is-length}, its $A$-length is
exactly $[\kappa(\mathfrak m_i):\kappa(\mathfrak m)]$.
Additivity along the original chain, now regarded as a chain
of $A$-submodules, gives the sum of these degrees over its factors.
Lemma \ref{editorial-source-lemma-simple-pieces} identifies the number of factors
at $\mathfrak m_i$ with
$\operatorname{length}_{B_{\mathfrak m_i}}(M_{\mathfrak m_i})$.
Grouping the finite sum by those maximal ideals proves the
displayed formula and its finiteness. When $M=0$, the chain
is empty and all terms are zero, as required.

The assumptions need only concern the maximal ideals in the
support of $M$. More precisely, let $\alpha:A\to B$ be any
ring homomorphism and let $M$ have finite length over $B$.
The preceding composition-factor proof shows that
$Q=\operatorname{Supp}_B(M)$ is a finite set of maximal ideals.
For each $\mathfrak q\in Q$, put
$\mathfrak p_{\mathfrak q}=\alpha^{-1}(\mathfrak q)$.
Then $M$ has finite length over $A$ if and only if every
$\mathfrak p_{\mathfrak q}$ is maximal in $A$ and every
extension $\kappa(\mathfrak q)/\kappa(\mathfrak p_{\mathfrak q})$
has finite degree. Under these conditions the exact formula is
$$
\operatorname{length}_A(M)
=\sum_{\mathfrak q\in Q}
[\kappa(\mathfrak q):\kappa(\mathfrak p_{\mathfrak q})]
\operatorname{length}_{B_{\mathfrak q}}(M_{\mathfrak q}).
$$
To prove the criterion in both directions, fix one composition
factor $k=B/\mathfrak q$. The original map identifies
$D=A/\mathfrak p_{\mathfrak q}$ with a subring of the field $k$.
If $\mathfrak p_{\mathfrak q}$ is maximal, $D$ is its residue
field and the $A$-length of $k$ is its dimension over $D$ by
Lemma \ref{editorial-source-lemma-dimension-is-length}, finite exactly in the
stated finite-degree case.
If $\mathfrak p_{\mathfrak q}$ is not maximal, the domain $D$
is not a field. Choose a nonzero nonunit $a\in D$.
Its image is nonzero in $k$, and the original $A$-submodules
$$
D\subsetneq a^{-1}D\subsetneq a^{-2}D\subsetneq\cdots\subset k
$$
are strictly increasing. Indeed equality of two consecutive
members would give $a^{-(j+1)}=a^{-j}d$ for some $d\in D$,
and multiplying in $k$ by $a^{j+1}$ would yield $1=ad$ in
the embedded $D$, contradicting the choice of $a$.
Adjoining zero and the endpoint $k$ gives finite chains of
arbitrarily large length, so $\operatorname{length}_A(k)=\infty$.
Finally, the original finite composition chain in $M$ remains
an exact filtration of $A$-modules. Additivity shows that its
$A$-length is the sum of its factor lengths, and is finite
exactly when each occurring factor has finite $A$-length.
Every $\mathfrak q\in Q$ occurs with the strictly positive
multiplicity specified above. This proves necessity, sufficiency
and the formula without any local or semilocal hypothesis on
the whole rings. If $Q$ is empty, $M=0$ and the criterion is
vacuous with an empty sum equal to zero.
\end{proof}

\begin{lemma}
\label{editorial-source-lemma-pullback-module}
Let $A \to B$ be a flat local homomorphism of local rings.
Then for any $A$-module $M$ we have
$$
\text{length}_A(M) \text{length}_B(B/\mathfrak m_AB)
=
\text{length}_B(M \otimes_A B).
$$
Here products of lengths use $0\cdot\infty=\infty\cdot0=0$.
In particular, if $\text{length}_B(B/\mathfrak m_AB) < \infty$
then $M$ has finite length if and only if $M \otimes_A B$ has finite length.
\end{lemma}

\begin{proof}
Retain the original flat local homomorphism $\alpha:A\to B$.
It is faithfully flat by Lemma \ref{editorial-source-lemma-local-flat-ff}.
The original closed fibre
$F=B/\mathfrak m_AB$ is nonzero, because locality implies
$\mathfrak m_AB\subset\mathfrak m_B\neq B$.
Thus $d=\operatorname{length}_B(F)$ lies in
$\mathbf Z_{\geq1}\cup\{\infty\}$.
When $M=0$, also $M\otimes_A B=0$, and the formula reads
$0=0$ using $0\cdot\infty=0$ if needed.

For every strict chain $0=M_0\subsetneq\cdots\subsetneq M_n=M$,
flatness gives injections of the actual modules $M_i\otimes_A B$
into $M\otimes_A B$ and exact sequences
$$
0\longrightarrow M_{i-1}\otimes_A B
\longrightarrow M_i\otimes_A B
\longrightarrow (M_i/M_{i-1})\otimes_A B
\longrightarrow0.
$$
Every quotient on the right is nonzero by faithful flatness,
so every step remains strict. Therefore infinite length of $M$
implies infinite length after base change, as required by
$\infty\cdot d=\infty$ for the positive value $d$.

If $0<\ell=\operatorname{length}_A(M)<\infty$, choose a
composition chain of length $\ell$. Since $A$ is local, each
simple factor is isomorphic to the original $A/\mathfrak m_A$.
The exact sequences above identify each new factor with $F$.
In the tensor order being used, define
$\Phi:(A/\mathfrak m_A)\otimes_A B\to F$ and
$\Psi:F\to(A/\mathfrak m_A)\otimes_A B$ by
$$
\begin{aligned}
\Phi((a+\mathfrak m_A)\otimes b)&=\alpha(a)b+\mathfrak m_AB,\\
\Psi(b+\mathfrak m_AB)&=(1+\mathfrak m_A)\otimes b.
\end{aligned}
$$
The first map respects tensor balancing and annihilates the
indicated representatives from $\mathfrak m_A$.
For the second map, every summand $\alpha(a)b$ with
$a\in\mathfrak m_A$ maps to
$(a+\mathfrak m_A)\otimes b=0$. Thus both maps are well
defined, and tensor balancing verifies that their composites
are identities on each pure tensor and quotient class.
Additivity along the filtration obtained by tensoring gives
$\operatorname{length}_B(M\otimes_A B)=\ell d$,
also when $d=\infty$, since a positive finite number of
infinite summands has infinite sum.
Together with the zero and infinite cases this proves the formula.

In particular, without any advance finiteness assumption on $d$,
the exact finite-length criterion is
$$
\begin{aligned}
&\operatorname{length}_B(M\otimes_A B)<\infty\\
&\quad\Longleftrightarrow\quad
M=0\ \text{or}\quad
\bigl(\operatorname{length}_A(M)<\infty\ \text{and}\ d<\infty\bigr).
\end{aligned}
$$
The positivity of $d$, faithful flatness and the proved equality
give both directions. When $d<\infty$, this specializes to
the stated equivalence.
\end{proof}

\begin{lemma}
\label{editorial-source-lemma-pullback-transitive}
Let $A \to B \to C$ be flat local homomorphisms of local rings. Then
$$
\text{length}_B(B/\mathfrak m_A B)
\text{length}_C(C/\mathfrak m_B C)
=
\text{length}_C(C/\mathfrak m_A C)
$$
\end{lemma}

\begin{proof}
Write the original maps as $\alpha:A\to B$ and $\beta:B\to C$.
Apply Lemma \ref{editorial-source-lemma-pullback-module} to $B\to C$ and the
original module $B/\mathfrak m_AB$. The receiving module in
that formula is identified with $C/\mathfrak m_AC$ by the maps
$\Phi:(B/\mathfrak m_AB)\otimes_B C\to C/\mathfrak m_AC$
and $\Psi:C/\mathfrak m_AC\to(B/\mathfrak m_AB)\otimes_B C$:
$$
\begin{aligned}
\Phi((b+\mathfrak m_AB)\otimes c)&=\beta(b)c+\mathfrak m_AC,\\
\Psi(c+\mathfrak m_AC)&=(1+\mathfrak m_AB)\otimes c.
\end{aligned}
$$
Here $\mathfrak m_AC$ is generated by the actual elements
$\beta(\alpha(a))$ for $a\in\mathfrak m_A$.
The first map respects balancing and sends any member of
$\mathfrak m_AB$ to that ideal in $C$.
For the second map, a generator times $c$ maps to
$(\alpha(a)+\mathfrak m_AB)\otimes c=0$.
Thus both are well defined, and their composites are identities
by the balancing relation. The base-change length formula now
gives precisely the claimed equality. Each of the three original
closed fibres is nonzero because the maps are local. Their
lengths are therefore positive or infinite, so their product
has no zero-times-infinity ambiguity.
\end{proof}














\subsection{Editorial treatment of Algebra lines 12713--12721}

\hypertarget{algebra-context-166}{}

\noindent\href{https://github.com/stacks/stacks-project/blob/a04446e57ec1fbc252a871afcec7752fb2807b14/algebra.tex\#L12713}{Original source, lines 12713--12721}. Complete original and reviewed TeX passages accompany this note. Every intervention is separately recorded; the following treatment is editorial.


\begin{lemma}
\label{editorial-source-lemma-finite-dimensional-algebra}
Suppose $R$ is a finite dimensional algebra over a field.
Then $R$ is Artinian.
\end{lemma}

\begin{proof}
Let $k$ be the given field and let $d=\dim_k R<\infty$.
Every ideal is a $k$-vector subspace of the original $R$,
because scalar multiplication is multiplication by the image
of an element of $k$ in $R$. In any strict descending chain
of ideals, the dimensions strictly decrease among the integers
$0,\ldots,d$. Thus there are at most $d$ strict changes and
every descending chain stabilizes, proving that $R$ is Artinian.
The same argument bounds every strict increasing chain.
In particular, directly from the definition of length,
$\operatorname{length}_R(R)\leq d$.
If $R=0$, then $d=0$, its only ideal is zero, and these
assertions include its length-zero chain.
\end{proof}


\subsection{Editorial treatment of Algebra lines 12723--12736}

\hypertarget{algebra-context-167}{}

\noindent\href{https://github.com/stacks/stacks-project/blob/a04446e57ec1fbc252a871afcec7752fb2807b14/algebra.tex\#L12723}{Original source, lines 12723--12736}. Complete original and reviewed TeX passages accompany this note. Every intervention is separately recorded; the following treatment is editorial.


\begin{lemma}
\label{editorial-source-lemma-artinian-finite-nr-max}
If $R$ is Artinian then $R$ has only finitely many maximal ideals.
\end{lemma}

\begin{proof}
If there were infinitely many maximal ideals, choose distinct
$\mathfrak m_1,\mathfrak m_2,\ldots$ and put
$J_n=\bigcap_{i=1}^n\mathfrak m_i$ for $n\geq1$.
For every $i\leq n$, distinctness and maximality imply
$\mathfrak m_i+\mathfrak m_{n+1}=R$. Choose actual elements
$a_i\in\mathfrak m_i$ and $b_i\in\mathfrak m_{n+1}$ with
$a_i+b_i=1$. Their product
$a_1\cdots a_n$ belongs to every $\mathfrak m_i$ for
$i\leq n$, hence to $J_n$, and is congruent to $1$ modulo
$\mathfrak m_{n+1}$. It therefore does not belong to
$J_{n+1}=J_n\cap\mathfrak m_{n+1}$.
Consequently $J_1\supsetneq J_2\supsetneq\cdots$ is an
infinite strict descending chain of ideals, contradicting the
Artinian hypothesis. Thus the set of maximal ideals is finite.
For the zero ring the set is empty, which also satisfies the claim.
\end{proof}


\subsection{Editorial treatment of Algebra lines 12738--12755}

\hypertarget{algebra-context-168}{}

\noindent\href{https://github.com/stacks/stacks-project/blob/a04446e57ec1fbc252a871afcec7752fb2807b14/algebra.tex\#L12738}{Original source, lines 12738--12755}. Complete original and reviewed TeX passages accompany this note. Every intervention is separately recorded; the following treatment is editorial.


\begin{lemma}
\label{editorial-source-lemma-artinian-radical-nilpotent}
Let $R$ be Artinian. The Jacobson radical of $R$ is a nilpotent ideal.
\end{lemma}

\begin{proof}
Let $I$ be the original Jacobson radical. By the descending
chain condition applied to $I\supset I^2\supset\cdots$,
choose $n\geq1$ such that $I^n=I^{n+1}$.
Put $J=\operatorname{Ann}_R(I^n)=\{x\in R:xI^n=0\}$.
If $J=R$, then $1\in J$ gives $I^n=0$, as desired.
Otherwise the nonempty collection of ideals strictly containing
$J$ has a minimal member $J'$: absence of a minimal member
would let one choose an infinite strict descending chain in that
collection, contradicting the hypothesis.
The original quotient $J'/J$ is nonzero and has no nonzero
proper submodule. Indeed the inverse image of such a submodule
under $J'\to J'/J$ would be an ideal strictly between $J$ and
$J'$. Hence $J'/J$ is simple.
By Lemma \ref{editorial-source-lemma-characterize-length-1}, it is isomorphic
to $R/\mathfrak m$ for some maximal ideal $\mathfrak m$.
It follows that $\mathfrak mJ'\subset J$.

Since $I\subset\mathfrak m$, for every $a\in I$ and $y\in J'$
we have $ay\in J$. For every $b\in I^n$ this gives
$b(ay)=0$. An element of $I^{n+1}=II^n$ is a finite sum
$\sum_{j=1}^t a_jb_j$ with $a_j\in I$ and $b_j\in I^n$;
its product with $y$ is the full sum
$\sum_{j=1}^t b_j(a_jy)=0$.
Thus $I^{n+1}J'=0$. The chosen equality $I^{n+1}=I^n$
then gives $I^nJ'=0$, so every $y\in J'$ lies in
$\operatorname{Ann}_R(I^n)=J$. This contradicts $J'\supsetneq J$.
Therefore $J=R$ and $I^n=0$. No finite generation of $I^n$
or Noetherian property has been assumed in this proof.
\end{proof}


\subsection{Editorial treatment of Algebra lines 12757--12784}

\hypertarget{algebra-context-169}{}

\noindent\href{https://github.com/stacks/stacks-project/blob/a04446e57ec1fbc252a871afcec7752fb2807b14/algebra.tex\#L12757}{Original source, lines 12757--12784}. Complete original and reviewed TeX passages accompany this note. Every intervention is separately recorded; the following treatment is editorial.


\begin{lemma}
\label{editorial-source-lemma-product-local}
Any ring with finitely many maximal ideals and
locally nilpotent Jacobson radical is the product of its localizations
at its maximal ideals. Also, all primes are maximal.
\end{lemma}

\begin{proof}
Let the maximal ideals be the distinct ideals
$\mathfrak m_1,\ldots,\mathfrak m_n$, and retain
$I=\bigcap_{i=1}^n\mathfrak m_i$.
If $n=0$, the ring is zero, since every nonzero ring has a
maximal ideal. Its spectrum is empty and its map to the empty
product of rings is the isomorphism of zero rings. Assume $n\geq1$.
Every element of $I$ is nilpotent, so it belongs to every prime.
Since $\mathfrak m_1\cdots\mathfrak m_n\subset I$, a prime
$\mathfrak p$ must contain some $\mathfrak m_i$.
Indeed, otherwise one could choose $a_i\in\mathfrak m_i\setminus
\mathfrak p$ for each $i$, whose product would belong to
$\mathfrak p$ while all its factors avoid it, contradicting
primality. Maximality then gives $\mathfrak p=\mathfrak m_i$.
This proves the assertion about prime ideals.

We construct the precise product isomorphism. For each ordered
pair $i\neq j$, choose $a_{ij}\in\mathfrak m_j$ and
$b_{ij}\in\mathfrak m_i$ with $a_{ij}+b_{ij}=1$, and set
$c_i=\prod_{j\neq i}a_{ij}$. For $n=1$ this product is $1$.
The element $c_i$ has residue $1$ at $\mathfrak m_i$ and
residue $0$ at every other maximal ideal. Thus
$c_i(1-c_i)\in I$. Choose $N_i\geq1$ for which
$[c_i(1-c_i)]^{N_i}=0$. Define the actual element
$$
e_i=\sum_{k=N_i}^{2N_i-1}
\binom{2N_i-1}{k}c_i^k(1-c_i)^{2N_i-1-k}.
$$
All its coefficients and powers are part of this construction.
The binomial identity $(c_i+(1-c_i))^{2N_i-1}=1$ gives
$$
1-e_i=\sum_{k=0}^{N_i-1}
\binom{2N_i-1}{k}c_i^k(1-c_i)^{2N_i-1-k}.
$$
Every term of $e_i$ contains $c_i^{N_i}$, and every term of
$1-e_i$ contains $(1-c_i)^{N_i}$. Their product is zero by
the chosen nilpotence equation. Hence $e_i^2=e_i$.
The residues of $e_i$ at all maximal ideals are the same as
those of $c_i$: substituting $0$ gives zero, and substituting
$1$ leaves just the term with $k=2N_i-1$, equal to one.
For $i\neq j$, the product $e_ie_j$ is an idempotent in $I$.
An idempotent that is nilpotent is zero, since all its positive
powers equal itself. Thus $e_ie_j=0$.
The sum $\sum_i e_i$ is an idempotent, and
$1-\sum_i e_i$ is an idempotent in $I$, so it is zero too.
We have obtained $\sum_i e_i=1$ with pairwise orthogonal
idempotents, without assuming a uniform nilpotence exponent for $I$.

Put $R_i=e_iR$, viewed as a ring with identity $e_i$.
The maps
$$
R\longrightarrow\prod_{i=1}^n R_i,\quad r\longmapsto(e_ir)_i,
\qquad
\prod_{i=1}^n R_i\longrightarrow R,\quad (x_i)_i\longmapsto\sum_i x_i
$$
are inverse ring homomorphisms, by the orthogonality and sum
equations just proved. The first component map is surjective
and has kernel $(1-e_i)R$: $e_ir=0$ implies
$r=(1-e_i)r$, and the reverse implication follows from
$e_i(1-e_i)=0$.
Maximal ideals of $R_i$ correspond to maximal ideals of $R$
containing $(1-e_i)R$. By the computed residues, only
$\mathfrak m_i$ contains $1-e_i$.
Therefore $R_i$ is a nonzero local ring with unique maximal
ideal $e_i\mathfrak m_i$.

The comparison with the original localization is
$$
R_i\longrightarrow R_{\mathfrak m_i},\quad x\longmapsto x/1,
\qquad
R_{\mathfrak m_i}\longrightarrow R_i,\quad
r/s\longmapsto (e_ir)(e_is)^{-1}.
$$
Here $s\notin\mathfrak m_i$ implies
$e_is\notin e_i\mathfrak m_i$: reducing modulo
$\mathfrak m_i$ would otherwise give $s=0$ in its residue
field. Hence $e_is$ is a unit in the original ring $R_i$.
Also $e_i/1$ is an idempotent unit in $R_{\mathfrak m_i}$,
so it equals one there. The first map is consequently unital.
For the second map, a fraction relation
$t(s'r-sr')=0$ with $t\notin\mathfrak m_i$ becomes the
same relation multiplied by $e_i$ in $R_i$; the images of
$t,s,s'$ are units, so the two prescribed images agree.
The maps preserve sums and products by their fraction formulas.
Their composites fix $x\in R_i$ and each original $r/s$,
respectively. Thus both are isomorphisms.
Combining them proves that the canonical map in the assertion is
$$
R\longrightarrow\prod_i R_{\mathfrak m_i},\qquad r\longmapsto(r/1)_i,
$$
with inverse
$(r_i/s_i)_i\mapsto\sum_i(e_ir_i)(e_is_i)^{-1}$.
All terms of this sum are elements of their specified subrings
$e_iR\subset R$; the original denominators and their inverses
in those subrings are retained.

The same construction works for every original $R$-module $M$,
without a finite-generation hypothesis. The maps
$M\to\bigoplus_i e_iM$, $x\mapsto(e_ix)_i$, and
$(x_i)_i\mapsto\sum_i x_i$ are inverse module maps.
The maps $e_iM\to M_{\mathfrak m_i}$, $x\mapsto x/1$,
have inverses
$$
x/s\longmapsto(e_is)^{-1}e_ix.
$$
The same fraction relation and the unit $e_it$ verify that this
inverse is well defined; the composites are identities because
$e_i/1=1$ in the localization.
Consequently $M\to\prod_i M_{\mathfrak m_i}$ is the canonical
isomorphism, with the displayed component maps and inverse sum.
Using finite additivity of length and the surjections
$R\to e_iR$, Lemma \ref{editorial-source-lemma-length-independent} gives
$$
\operatorname{length}_R(M)
=\sum_{i=1}^n\operatorname{length}_{R_{\mathfrak m_i}}
 (M_{\mathfrak m_i}),
$$
also when one or more terms are infinite. In the zero-ring case
every unital module is zero and this is the empty sum.
\end{proof}


\subsection{Editorial treatment of Algebra lines 12786--12819}

\hypertarget{algebra-context-170}{}

\noindent\href{https://github.com/stacks/stacks-project/blob/a04446e57ec1fbc252a871afcec7752fb2807b14/algebra.tex\#L12786}{Original source, lines 12786--12819}. Complete original and reviewed TeX passages accompany this note. Every intervention is separately recorded; the following treatment is editorial.


\begin{lemma}
\label{editorial-source-lemma-artinian-finite-length}
A ring $R$ is Artinian if and only if it has finite length
as a module over itself. Any such ring $R$ is both Artinian and
Noetherian, any prime ideal of $R$ is a maximal ideal, and $R$ is equal
to the (finite) product of its localizations at its maximal ideals.
\end{lemma}

\begin{proof}
If the original module $R$ has finite length, the proof of
Lemma \ref{editorial-source-lemma-finite-length-finite} gives both chain conditions
for its submodules, which are exactly its ideals. The descending
condition says that $R$ is Artinian. The ascending condition says
that every ideal is finite by the arbitrary-module equivalence
proved in Lemma \ref{editorial-source-lemma-Noetherian-basic}, so $R$ is Noetherian.

Conversely, suppose $R$ is Artinian. Lemmas
\ref{editorial-source-lemma-artinian-finite-nr-max},
\ref{editorial-source-lemma-artinian-radical-nilpotent}, and
\ref{editorial-source-lemma-product-local} identify $R$ with the finite product
of its original localizations $R_{\mathfrak m_i}$ and prove
that every prime is maximal. The product comparison is the
canonical map $r\mapsto(r/1)_i$, with the inverse proved there.
Its component maps are surjective: in the preceding proof each
localization is identified with the quotient $e_iR$ of $R$.
Thus each $R_{\mathfrak m_i}$ is Artinian, since a descending
chain of its ideals pulls back to such a chain in $R$.
When $R=0$, the product is empty and the length is zero.
For nonzero $R$, the module decomposition and length formula
in that same proof show that it suffices to prove finite length
for each original local factor.

Consider one such local Artinian ring $C=R_{\mathfrak m_i}$,
with maximal ideal $\mathfrak n=\mathfrak m_i C$ and residue
field $k=C/\mathfrak n$. By the nilpotence lemma,
$\mathfrak n^N=0$ for some $N\geq1$.
Each layer $V_j=\mathfrak n^j/\mathfrak n^{j+1}$,
$0\leq j<N$, is a $k$-vector space via the original quotient
map, because $\mathfrak n V_j=0$.
This space must have finite dimension. If not, choose a basis
containing distinct elements $v_1,v_2,\ldots$.
For each $a\geq0$ let $W_a$ be the span of all basis elements
except $v_1,\ldots,v_a$. Then
$W_0\supsetneq W_1\supsetneq\cdots$ is a strict descending
chain of vector subspaces. Their inverse images under
$\mathfrak n^j\to V_j$ are ideals of $C$ containing
$\mathfrak n^{j+1}$, and strictness is preserved by the
surjectivity of that quotient map. This contradicts the
Artinian condition on $C$.
By Lemma \ref{editorial-source-lemma-dimension-is-length}, the length of each
layer is its finite $k$-dimension. Additivity along the original
filtration gives
$$
\operatorname{length}_C(C)
=\sum_{j=0}^{N-1}\dim_k(\mathfrak n^j/\mathfrak n^{j+1})<\infty.
$$
Summing these finite lengths over the original local factors
proves finite length of $R$ and completes all assertions.

There are quantitative consequences for the same ring.
Write $\ell_i=\operatorname{length}_{R_{\mathfrak m_i}}
(R_{\mathfrak m_i})$. Each is positive and finite, and
$\operatorname{length}_R(R)=\sum_i\ell_i$.
The bound proved in Lemma \ref{editorial-source-lemma-length-infinite} gives
$(\mathfrak m_iR_{\mathfrak m_i})^{\ell_i}=0$.
Under the canonical product map, the original Jacobson radical
$I$ maps onto $\prod_i\mathfrak m_iR_{\mathfrak m_i}$:
membership in each component maximal ideal is exactly membership
in the corresponding original maximal ideal, and the product
map is an isomorphism. Powers are componentwise, so for
$R\neq0$,
$$
I^{\max_i\ell_i}=0.
$$
For $R=0$ the radical is already zero and no maximum of an empty
set is used.
Finally, every finite $R$-module $M$ has finite length. If its
given generating list has $r$ members, the associated surjection
$R^r\to M$ and additivity of length give the precise bound
$$
\operatorname{length}_R(M)\leq
r\operatorname{length}_R(R).
$$
The reverse implication, from finite length to finite generation,
is Lemma \ref{editorial-source-lemma-finite-length-finite}. Thus over this
Artinian ring the two conditions on $M$ are equivalent, with
the displayed bound retaining the chosen number of generators.
\end{proof}


\subsection{Editorial treatment of Algebra lines 12846--12855}

\hypertarget{algebra-context-171}{}

\noindent\href{https://github.com/stacks/stacks-project/blob/a04446e57ec1fbc252a871afcec7752fb2807b14/algebra.tex\#L12846}{Original source, lines 12846--12855}. Complete original and reviewed TeX passages accompany this note. Every intervention is separately recorded; the following treatment is editorial.


\begin{lemma}
\label{editorial-source-lemma-composition-essentially-of-finite-type}
The class of ring maps which are essentially of finite type is
preserved under composition. Similarly for essentially of finite
presentation.
\end{lemma}

\begin{proof}
Let the original maps be $R\to S\to T$. Choose an $R$-algebra
$A$ of finite type and a multiplicative subset $U\subset A$,
together with the given $R$-algebra isomorphism $U^{-1}A\cong S$.
Choose a finite type $S$-algebra $B$, a multiplicative subset
$V\subset B$, and the given $S$-algebra isomorphism
$V^{-1}B\cong T$.
Write $B=S[X_1,\ldots,X_t]/H$ for an ideal $H$.
Let $H_0\subset A[X_1,\ldots,X_t]$ be its inverse image under
the coefficient map to $S[X_1,\ldots,X_t]$, and put
$C=A[X_1,\ldots,X_t]/H_0$. This is a finite type $R$-algebra.
The natural map $U^{-1}C\to B$ is an isomorphism.
To see the kernel assertion explicitly, every polynomial over
$U^{-1}A$ is $G/u$ with $G\in A[X_1,\ldots,X_t]$ and
$u\in U$, by taking a product of its finitely many coefficient
denominators. If $G/u\in H$, then $G/1=u(G/u)\in H$,
so $G\in H_0$. Conversely every polynomial in $H_0$ maps
into $H$. This proves that the localized ideal is exactly
$H$, and proves the asserted quotient isomorphism with its
original polynomial and fraction maps.

In either this construction or the finite-presentation construction
below, set
$$
W=\{c\in C:\text{the image of }c\text{ in the original }T
\text{ is a unit}\}.
$$
This is a multiplicative subset. The canonical map
$W^{-1}C\to T$ is an isomorphism as follows.
Every image of $u\in U$ in $C$ belongs to $W$, so
$C\to W^{-1}C$ extends to $B=U^{-1}C$.
For $v\in V$, write its original representative in this
localization as $d/u$ with $d\in C$ and $u\in U$.
Its image is a unit in $T$, as is the image of $u$, hence
the image of $d$ is a unit in $T$ and $d\in W$.
Thus $v$ maps to a unit of $W^{-1}C$, and there is a map
$T=V^{-1}B\to W^{-1}C$.
It sends a fraction represented as $(c/u)/(d/w)$, with
$c,d\in C$, $u,w\in U$ and $d/w\in V$, to
$cw/(ud)$. The preceding argument applied to this actual
$d/w$ gives $d\in W$, and also $u\in W$, so its
displayed denominator belongs to $W$.
Both maps are defined by the localization relations and their
composites fix $C$, the inverted images of $U$ and $V$,
and the inverted elements of $W$. They are therefore mutual
inverses. This proves the assertion for finite type.

For finite presentation choose $A$ of finite presentation over
$R$ and choose a finite presentation
$B=S[X_1,\ldots,X_t]/(F_1,\ldots,F_m)$.
The contraction $H_0$ above need not be finitely generated,
so instead perform the following finite construction, keeping
all actual coefficients and denominators.
For each $j$, write
$$
F_j=\sum_{\nu\in E_j}\frac{a_{j\nu}}{u_{j\nu}}X^\nu,
\qquad a_{j\nu}\in A,\quad u_{j\nu}\in U,
$$
where $E_j$ is the finite set of exponent tuples occurring in
the polynomial. Set
$$
d_j=\prod_{\nu\in E_j}u_{j\nu},\qquad
G_j=\sum_{\nu\in E_j}
a_{j\nu}\left(\prod_{\mu\in E_j\setminus\{\nu\}}
u_{j\mu}\right)X^\nu\in A[X_1,\ldots,X_t].
$$
The coefficient image of $G_j$ in $S[X_1,\ldots,X_t]$ is
exactly $(d_j/1)F_j$. The factor $d_j/1$ is retained and is
a unit; thus the ideals generated by the images of the $G_j$
and by the $F_j$ are equal. If $F_j=0$, use $E_j=\varnothing$,
$d_j=1$ and $G_j=0$.
Put $C=A[X_1,\ldots,X_t]/(G_1,\ldots,G_m)$.
Then $U^{-1}C\cong B$ by precisely these coefficient maps.
Moreover $C$ is finitely presented over $R$: if
$A=R[Y_1,\ldots,Y_a]/(P_1,\ldots,P_b)$, choose polynomial
lifts $\widetilde G_j\in R[Y_1,\ldots,Y_a,X_1,\ldots,X_t]$
of all the original coefficients in $G_j$.
The resulting presentation is
$$
C\cong R[Y_1,\ldots,Y_a,X_1,\ldots,X_t]
/(P_1,\ldots,P_b,\widetilde G_1,\ldots,\widetilde G_m).
$$
Indeed quotienting by the $P_i$ first gives the polynomial
ring over $A$, after which the images of the remaining listed
relations are exactly the $G_j$.
Using the same $W$ and the same inverse localization maps proves
$T\cong W^{-1}C$, as required. These arguments allow empty
variable and relation lists and multiplicative subsets containing
zero; in the latter case the corresponding localization is the
zero ring and the same universal maps apply.
\end{proof}


\subsection{Editorial treatment of Algebra lines 12857--12866}

\hypertarget{algebra-context-172}{}

\noindent\href{https://github.com/stacks/stacks-project/blob/a04446e57ec1fbc252a871afcec7752fb2807b14/algebra.tex\#L12857}{Original source, lines 12857--12866}. Complete original and reviewed TeX passages accompany this note. Every intervention is separately recorded; the following treatment is editorial.


\begin{lemma}
\label{editorial-source-lemma-base-change-essentially-of-finite-type}
The class of ring maps which are essentially of finite type is
preserved by base change. Similarly for essentially of finite
presentation.
\end{lemma}

\begin{proof}
Retain the original map $R\to S$ and a map $\rho:R\to R'$.
Write $S=U^{-1}A$ through a specified $R$-algebra isomorphism,
where $A$ has the required finite type or finite presentation
property. Put $A'=A\otimes_R R'$ and let $U'$ be the image
of $U$ under $a\mapsto a\otimes1$.
There is an isomorphism of the actual base-changed rings
$$
\Phi:S\otimes_R R'\longrightarrow (U')^{-1}A',
\qquad
(a/u)\otimes r'\longmapsto
\frac{a\otimes r'}{u\otimes1}.
$$
To construct it, the ring map $A\to(U')^{-1}A'$ sends
$a$ to $(a\otimes1)/1$ and inverts $U$, so extends to $S$.
Together with the map $R'\to(U')^{-1}A'$,
$r'\mapsto(1\otimes r')/1$, it induces the displayed map;
the two maps agree on the original $R$ by tensor balancing.
Conversely $A'\to S\otimes_R R'$ sends
$a\otimes r'$ to $(a/1)\otimes r'$ and sends $u\otimes1$
to a unit with inverse $(1/u)\otimes1$.
It therefore induces a map $\Psi:(U')^{-1}A'\to S\otimes_R R'$.
Every numerator is a finite sum of pure tensors and every
denominator in $U'$ is the image of an original $u\in U$;
on those fractions its formula is
$$
\Psi\left(\frac{\sum_{i=1}^q a_i\otimes r_i'}{u\otimes1}\right)
=\sum_{i=1}^q(a_i/u)\otimes r_i'.
$$
The construction from the localization maps proves independence
of both tensor expressions and fraction representatives.
The composites fix each pure tensor $(a/u)\otimes r'$ and
each displayed fraction, respectively, so they are identities.

If $a_1,\ldots,a_t$ generate $A$ as an $R$-algebra, their
images $a_i\otimes1$ generate $A'$ as an $R'$-algebra:
apply the original polynomial expression for each $a$ in any
finite sum of pure tensors, retaining its coefficients through
$\rho$. Thus finite type is preserved.
For finite presentation, write
$A=R[X_1,\ldots,X_t]/(P_1,\ldots,P_m)$.
The coefficient map gives
$$
A'\cong R'[X_1,\ldots,X_t]/(\rho(P_1),\ldots,\rho(P_m)),
$$
where for the full original polynomial
$P_j=\sum_\nu r_{j\nu}X^\nu$ one has
$\rho(P_j)=\sum_\nu\rho(r_{j\nu})X^\nu$.
Indeed $R[X_1,\ldots,X_t]\otimes_R R'$ maps isomorphically
to $R'[X_1,\ldots,X_t]$ by
$rX^\nu\otimes r'\mapsto\rho(r)r'X^\nu$, with inverse
$r'X^\nu\mapsto X^\nu\otimes r'$.
Tensoring the quotient presentation is right exact, so the
kernel of the resulting surjection is generated by the images
of the original ideal $(P_1,\ldots,P_m)$.
A finite sum of its polynomial multiples maps to the same
finite sum of multiples of $\rho(P_j)$, and each listed
$\rho(P_j)$ occurs, proving the asserted ideal equality.
This is a finite presentation of $A'$.
In either case the proved isomorphism $S\otimes_R R'\cong
(U')^{-1}A'$ establishes the required essential property.
It remains valid when a denominator maps to zero, since both
localization constructions and their inverse maps then yield
the zero ring.
\end{proof}


\subsection{Editorial treatment of Algebra lines 12868--12924}

\hypertarget{algebra-context-173}{}

\noindent\href{https://github.com/stacks/stacks-project/blob/a04446e57ec1fbc252a871afcec7752fb2807b14/algebra.tex\#L12868}{Original source, lines 12868--12924}. Complete original and reviewed TeX passages accompany this note. Every intervention is separately recorded; the following treatment is editorial.


\begin{lemma}
\label{editorial-source-lemma-essentially-of-finite-type-into-artinian-local}
Let $R \to S$ be a ring map. Assume $S$ is an Artinian local ring with
maximal ideal $\mathfrak m$. Then
\begin{enumerate}
\item $R \to S$ is finite if and only if $R \to S/\mathfrak m$ is finite,
\item $R \to S$ is of finite type if and only if $R \to S/\mathfrak m$
is of finite type.
\item $R \to S$ is essentially of finite type if and
only if the composition $R \to S/\mathfrak m$ is essentially
of finite type.
\end{enumerate}
\end{lemma}

\begin{proof}
Write $k=S/\mathfrak m$ and retain its quotient map from $S$.
If $S$ is finite as an $R$-module, its quotient $k$ is finite.
Conversely, suppose $k$ is generated over $R$ by
$\bar z_1,\ldots,\bar z_r$. By
Lemma \ref{editorial-source-lemma-artinian-finite-length}, the original $S$
has a composition chain
$$
0=I_0\subsetneq I_1\subsetneq\cdots\subsetneq I_\ell=S,
\qquad \ell=\operatorname{length}_S(S).
$$
Each $I_j$ is an ideal and each quotient $I_j/I_{j-1}$ is
isomorphic, as an $S$-module, to the original $k$, since $S$
is local. For every $j$, choose lifts $z_{j1},\ldots,z_{jr}\in I_j$
of the images of the specified $\bar z_a$ under such an
isomorphism $k\to I_j/I_{j-1}$.
These $r\ell$ elements generate $S$ over $R$.
In fact, for $x\in I_j$ its class is
$\sum_{a=1}^r r_a(z_{ja}+I_{j-1})$ for some original
$r_a\in R$, so $x-\sum_a r_a z_{ja}\in I_{j-1}$.
Induction on $j$ supplies the claimed expression with all the
chosen lifted generators. This proves (1), and gives the
quantitative bound $r\operatorname{length}_S(S)$ on a number
of $R$-generators, with the full original filtration retained.

For (2), a quotient of a finite type $R$-algebra is of finite
type: the images of the given algebra generators still generate.
Conversely, lift a finite algebra generating list of $k$ to
$f_1,\ldots,f_n\in S$ and form the original coefficient map
$R[X_1,\ldots,X_n]\to S$, $X_i\mapsto f_i$.
Its composite to $k$ is surjective, so $k$ is a cyclic module
over this polynomial ring. Part (1), applied with this ring as
the base, shows that $S$ is a finite module over it.
Choose its finite module generators $s_1,\ldots,s_q$.
Every $s\in S$ is a sum of these $s_j$ with coefficients
that are polynomials in the original $f_i$ and coefficients
from $R$. Therefore $f_1,\ldots,f_n,s_1,\ldots,s_q$ generate
$S$ as an $R$-algebra, proving (2).

For (3), a quotient is a finite type map, hence essentially of
finite type; the forward implication follows from
Lemma \ref{editorial-source-lemma-composition-essentially-of-finite-type}.
Conversely suppose a specified localization of
$D=R[X_1,\ldots,X_n]/H$ is isomorphic to the original field
$k$. Choose $f_i\in S$ lifting the corresponding images of
$X_i$ in $k$. Let $J$ be the kernel of
$R[X_1,\ldots,X_n]\to S$, $X_i\mapsto f_i$, and retain
the actual subalgebra
$A=R[X_1,\ldots,X_n]/J\subset S$ and its prime
$\mathfrak p=A\cap\mathfrak m$.
The subring $A/\mathfrak p\subset k$ is exactly the image
of $D$ in $k$, since both are generated by the same original
coefficient images and the same $f_i$ modulo $\mathfrak m$.
Every element of $k$ is a fraction of elements of this image,
with a nonzero denominator coming from the specified
multiplicative subset of $D$. Thus $k\subset
\operatorname{Frac}(A/\mathfrak p)$ in the given ambient field.
Conversely the field $k$ contains that domain and the inverse
of each of its nonzero elements, so
$\operatorname{Frac}(A/\mathfrak p)\subset k$.
These two original embeddings identify
$\kappa(\mathfrak p)=\operatorname{Frac}(A/\mathfrak p)$
with the original $k$.

Every $a\in A\setminus\mathfrak p$ has image outside
$\mathfrak m$ and hence is a unit in $S$. The inclusion
$A\subset S$ therefore extends to
$A_{\mathfrak p}\to S$, $a/u\mapsto a\,u^{-1}$,
using the actual inverse in $S$.
Its composite to $k$ is the residue map just identified and
is surjective. Part (1) shows that $S$ is finite over
$A_{\mathfrak p}$. Hence it is finite type over that ring.
The map $R\to A$ is finite type and $A\to A_{\mathfrak p}$
is a localization. Applying the proved composition lemma to
these original maps gives that $R\to S$ is essentially of
finite type, as required.
\end{proof}


\section{Complete supplementary proofs}
\hypertarget{algebra_localization_kgroups_graded_notes_20260927}{}
\subsection{Editorial evidence for Algebra 12926--13248}\label{algebra-proof-001-editorial-evidence-for-algebra-1292613248}

Source: the original \texttt{algebra.tex} at authority commit \texttt{a04446e57ec1fbc252a871afcec7752fb2807b14}, SHA-256 \texttt{FA8BB92E58A4F78A2BD01B3B6A4A87DE0A0D279F5DD90641B574DD5FBFFFA4F3}. The source interval and all eight received reports were read. The exact old and proposed readings are retained in \texttt{ALGEBRA\_INTAKE\_REVIEW\_20260925.json}, groups 368--373.

These are editorial arguments supporting small, source-linked corrections. This file is not replacement text for the translated chapter. Both current and diplomatic translations must retain the original exposition and identify mathematical source corrections visibly. The source's omitted proofs and its stated degree of generality do not create mandatory rewriting assignments. No novelty is claimed.

\subsubsection{Valuation quantifier}\label{algebra-proof-001-valuation-quantifier}

Source: \texttt{lemma-localization-at-closed-point-special-fibre}, lines 12930--12972. Report \texttt{OCC-00370} concerns lines 12951--12954: the homomorphism from the entire polynomial ring to the valuation ring requires \textbf{every} variable image to belong to that valuation ring.

Write \(P=R[x_1,\ldots,x_n]\), let \(\mathfrak q\subset P\) be the original prime, and let \(\lambda_j\) be the image of \(x_j\) in \(\kappa(\mathfrak q)\). The image \(D\) of \(R\) in this field is a quotient of the local ring \(R\), hence is a local domain. The valuation ring \(A\subset\kappa(\mathfrak q)\) chosen in the source dominates \(D\), so the inverse image of \(\mathfrak m_A\) in \(R\) is \(\mathfrak m_R\).

If all \(\lambda_j\) belong to \(A\), the original evaluation map factors as \(P\to A\subset\kappa(\mathfrak q)\). Its inverse image \(\mathfrak p\) of \(\mathfrak m_A\) is proper and contains both \(\mathfrak q\) and \(\mathfrak m_RP\). A maximal ideal \(\mathfrak m\) of \(P\) containing \(\mathfrak p\) contracts to \(\mathfrak m_R\), since its contraction is proper and contains that maximal ideal. Localizing \(P_{\mathfrak m}\) at the images of \(P\setminus\mathfrak q\) gives \(P_{\mathfrak q}\), with both comparison maps given by the original fractions.

Otherwise at least one nonzero \(\lambda_j\) lies outside \(A\). Use the additive order convention \(a\in A\setminus\{0\}\) if and only if \(v(a)\geq0\). Among the finitely many \textbf{nonzero} \(\lambda_j\), choose an index \(i\) of least value. Its value is negative. Thus \(v(\lambda_i^{-1})>0\) and \(v(\lambda_j/\lambda_i)\geq0\) whenever \(\lambda_j\ne0\). For a zero \(\lambda_j\), the ratio is zero and belongs to \(A\) directly; no value of zero is used.

For the original chart \(Q=R[y_0,y_1,\ldots,\widehat{y_i},\ldots,y_n]\), the maps

\[
\alpha:Q_{y_0}\longrightarrow P_{x_i},\qquad
y_0\longmapsto x_i^{-1},\quad y_j\longmapsto x_jx_i^{-1},
\] \[
\beta:P_{x_i}\longrightarrow Q_{y_0},\qquad
x_i\longmapsto y_0^{-1},\quad x_j\longmapsto y_jy_0^{-1}
\]

fix \(R\) and are inverse. Indeed, \(\beta\alpha(y_0)=y_0\), \(\beta\alpha(y_j)=(y_jy_0^{-1})y_0=y_j\), \(\alpha\beta(x_i)=x_i\), and \(\alpha\beta(x_j)=(x_jx_i^{-1})x_i=x_j\); these identities also give the inverses on denominators. Since \(\lambda_i\ne0\), \(x_i\notin\mathfrak q\). The evaluation maps to \(\kappa(\mathfrak q)\) agree under these two isomorphisms. Their kernel primes therefore correspond, giving the claimed isomorphism \(Q_{\mathfrak q'}\cong P_{\mathfrak q}\), where \(\mathfrak q'\) is the inverse image of \(\mathfrak qP_{\mathfrak q}\). Every chart variable, including \(y_0\), now has image in \(A\). The preceding proper-ideal argument applies to \(Q\). When \(n=0\), the all-images condition is vacuous and the first case applies; no index is chosen.

The initial quotient reduction also preserves the original ring. If the given local ring is \(S=W^{-1}(P/I)\), let \(\mathfrak q\) be the inverse image of its maximal ideal in \(P\), and let \(U\subset P\) be the inverse image of \(W\). Every element outside \(\mathfrak q\) maps to a unit in \(S\), and \(U\cap\mathfrak q=\varnothing\). Consequently \(P_{\mathfrak q}\to S\) is defined and surjective. If \(p/t\) maps to zero, some \(u\in U\) has \(up\in I\), so \(p/t\in IP_{\mathfrak q}\). The reverse inclusion is immediate, proving \(S=P_{\mathfrak q}/IP_{\mathfrak q}\). A quotient of a localization of \(P_{\mathfrak m}\) is a localization of its quotient by the contracted ideal: both rings have fractions with the same denominators and the same zero criterion. Thus the source reduction is valid. Only the all-images quantifier needs a textual clarification.

\subsubsection{Nonnegative presentation rank}\label{algebra-proof-001-nonnegative-presentation-rank}

Source: lines 13013--13022; report \texttt{OCC-00371}. The ordinary finite free module \(R^n\) in this construction has \(n\geq0\). A finite generating list of a module \(M\), of length \(n\), defines a surjection \(R^n\to M\); its kernel \(K\) yields \(R^n/K\cong M\). The union over nonnegative integers of the sets of submodules of \(R^n\) is a set. Zero is included: \(R^0=0\) represents the zero module. Negative ranks have no meaning in the displayed construction. The correction changes only ``all integers'' to ``all nonnegative integers''.

\subsubsection{Both split exact sequences}\label{algebra-proof-001-both-split-exact-sequences}

Source: lines 13068--13074 and 13147--13153. Reports \texttt{OCC-00372}, \texttt{OCC-00373}, and the overlapping report \texttt{OCC-12299} refer to the same two wrong summands. The second actual correction is at line \textbf{13152}; the overlapping report's proposed line 13157 and range 13154--13158 are stale.

In a split exact sequence \(0\to M'\xrightarrow{u}M\xrightarrow{q}M''\to0\), choose its section \(s:M''\to M\), with \(qs=1\). The maps

\[
\Phi:M'\oplus M''\to M,\quad (a,b)\mapsto u(a)+s(b),
\qquad
\Psi:M\to M'\oplus M'',\quad
m\mapsto\bigl(u^{-1}(m-sq(m)),q(m)\bigr)
\]

are defined: \(q(m-sq(m))=0\), and \(u\) is an isomorphism onto \(\ker q\). They are linear. Their composites satisfy \(\Phi\Psi(m)=m-sq(m)+sq(m)=m\) and \(\Psi\Phi(a,b)=(a,b)\). Hence the second summand is the actual quotient \(M''\). For finite free terms, adjoining their bases gives rank additivity. Over a PID or a local ring, uniqueness of finite free rank gives the source's stated relation. The wrong formula is false even for \(M'=0\), \(M=M''=R\ne0\), with the quotient map the identity. The three reports are accounted for once each.

Both minimal corrections already exist in the cumulative source as \texttt{MC-STK-ERR-1597-OP1} and \texttt{MC-STK-ERR-1597-OP2}, with the exact original authority lines 13073 and 13152. The manifest-bound unit and both current readings were inspected. This batch retains \texttt{MC-STK-ERR-1597} and records semantic agreement; it proposes no duplicate operation or new correction identity for either occurrence. An initial draft check wrongly required the whole interval to remain identical to the authority. It stopped before append. Inspection identified these two prior corrections and a separate retained FAC insertion after line 13247. The check now binds each actual edit and the retained canonical unit separately. The FAC insertion is preserved and is not newly certified by this batch.

This change does not alter the local finite-projective argument or its use at lines 13174--13180: a free module of rank \(n\) still has length \(n\operatorname{length}_R(R)\) over an Artinian local ring, by the original length additivity.

\subsubsection{PID subtraction and zero factors}\label{algebra-proof-001-pid-subtraction-and-zero-factors}

Source: \texttt{example-K0-PID}, lines 13078--13090; report \texttt{OCC-01373}. The exact sequence \(0\to(d_i)\to R\to R/(d_i)\to0\) gives

\[
[R/(d_i)]=[R]-[(d_i)].
\]

The source prints the subtraction in the opposite order. For \(d_i\ne0\), the map \(R\to(d_i)\), \(a\mapsto ad_i\), is surjective by definition of the ideal and injective because the PID is a domain. Its inverse sends \(ad_i\) to the uniquely determined coefficient \(a\). Thus \([(d_i)]=[R]\) and the correctly signed difference is zero. Nonzero unit factors are allowed and have zero quotient.

If a displayed factor has \(d_i=0\), then \((d_i)=0\), \(R/(d_i)=R\), and its class is \([R]\), not zero. For the original displayed decomposition with arbitrary values of the \(d_i\), the exact formula is

\[
[M]=r[R]+\sum_{\{i:d_i=0\}}[R]
       +\sum_{\{i:d_i\ne0\}}\bigl([R]-[R]\bigr)
     =\bigl(r+\#\{i:d_i=0\}\bigr)[R].
\]

Every contribution is retained in the first expression; the second equality follows from the group law. The structure theorem permits choosing the torsion cyclic summands with all \(d_i\ne0\); any zero factors contribute free summands. The small proposed clarification states that choice explicitly next to the decomposition. This makes the source's subsequent isomorphism \((d_i)\cong R\) legitimate without changing the theorem or its exposition.

For completeness of the receiving identification, put \(K=\operatorname{Frac}(R)\). Localizing an exact sequence at \(R\setminus\{0\}\) remains exact: an element mapping to zero is killed by a nonzero denominator, which multiplies its numerator into the original kernel; injectivity is checked by the same zero criterion, and lifts of numerators give surjectivity. Taking finite \(K\)-dimensions therefore defines an additive rank map on \(K'_0(R)\). A nonzero \(d_i\) gives zero after localization, whereas a zero \(d_i\) gives one copy of \(K\). The map sends \([R]\) to 1, so the displayed generator formula proves its inverse is \(n\mapsto n[R]\). The original PID conclusion and its \(k[x]\) consequence at lines 13093--13098 remain valid. The separate node example at lines 13100--13106 is not proved or certified by this correction.

\subsubsection{Graded wording}\label{algebra-proof-001-graded-wording}

Source: lines 13224--13228; report \texttt{OCC-01374}. Removing the comma between ``definitions'' and ``and lemmas'' changes punctuation only. The mathematical definitions, including the allowance of negative module degrees and exclusion of negative ring degrees, remain unchanged.

\subsubsection{Graded scalar action}\label{algebra-proof-001-graded-scalar-action}

Source: lines 13231--13247; report \texttt{OCC-00374}. With the source's exact convention \(N(n)_d=N_{n+d}\), an element \(f\in\operatorname{GrHom}_n(M,N)\) is an \(S\)-linear map with \(f(M_d)\subset N_{d+n}\). For homogeneous \(s\in S_e\), define \((sf)(m)=s f(m)\). The original ring is commutative, so \((sf)(am)=sa f(m)=a(sf)(m)\), proving \(S\)-linearity. The image of \(M_d\) lies in \(N_{d+n+e}\), so \(sf\) has degree \(n+e\).

For finite homogeneous decompositions \(s=\sum_e s_e\) and \(f=\sum_n f_n\), define the component of degree \(k\) of their action to be \(\sum_{e+n=k}s_ef_n\). Only finitely many components occur. Additivity, \((st)f=s(tf)\), and \(1f=f\) follow by evaluating each map on every element of \(M\) and using the module axioms of \(N\). Thus this is exactly the asserted graded \(S\)-module structure. Composable maps \(M\to N\to P\) have a graded composition, but for unrelated \(M,N\), the given data do not supply composition as an internal operation on \(\operatorname{GrHom}(M,N)\).

``Multiplication'' can also mean scalar multiplication, so the original sentence is not evidence that the source falsely asserts an internal algebra structure. Classify the proposed ``\(S\)-action'' wording as clarification. Keep the construction here as editorial evidence; do not replace the source's deliberate omission with this expanded argument.

\hypertarget{algebra_graded_proj_hilbert_notes_20260927}{}
\subsection{Editorial evidence for Algebra 13249--13990}\label{algebra-proof-002-editorial-evidence-for-algebra-1324913990}

Authority: original \texttt{algebra.tex}, commit \texttt{a04446e57ec1fbc252a871afcec7752fb2807b14}, SHA-256 \texttt{FA8BB92E58A4F78A2BD01B3B6A4A87DE0A0D279F5DD90641B574DD5FBFFFA4F3}. This interval and all twelve reports whose first loci lie in it were read. The proposed source operations are small corrections or clarifications. The arguments and strengthening below are separately identified editorial material, not replacement proof bodies for a translation. No novelty is claimed.

\subsubsection{Bounded-below graded Nakayama}\label{algebra-proof-002-bounded-below-graded-nakayama}

Source: \texttt{lemma-graded-NAK}, lines 13249--13272. Its four statements are valid as written. A finite graded module over the original nonnegatively graded ring is bounded below: take homogeneous parts of a finite generating list; multiplication by a ring element never lowers degree below the least generator degree. The original induction therefore works, including a module with no generators.

There is a strengthening, recorded here without changing the source. In statements (1), (3), and (4), replace finite generation of \(M\) by the assumption that \(M_d=0\) for all \(d<b\), for some integer \(b\). In (2), the corresponding weaker assumption is that \(N'\) is bounded below. The original nonnegative ring grading, possible negative module degrees, and all original maps remain unchanged.

For (1), suppose \(M=S_+M\). Induct on integers \(d\geq b\). For a homogeneous \(m\in M_d\), an expression in \(S_+M\), followed by taking its degree-\(d\) component, writes \(m\) as a finite sum of products \(s_e m_{d-e}\) with \(e>0\). Every module component has strictly smaller degree and is zero by induction. Thus \(M_d=0\) for every \(d\), so \(M=0\). This uses no finite generating set or Noetherian assumption.

For (2), let \(Q=M/N\), and let \(L\) be the image of \(N'\) in \(Q\). The original equation gives \(Q=S_+L\). Since \(L\) is a submodule, \(S_+L\subset L\subset Q\), hence \(Q=L=S_+Q\). The surjection \(N'\to Q\) is graded, so boundedness below of \(N'\) passes to \(Q\). Part (1) gives \(Q=0\), or \(M=N\).

For (3), if \(\varphi:N\to M\) has surjective reduction modulo \(S_+\), then \(Q=M/\operatorname{Im}\varphi\) satisfies \(Q=S_+Q\). It is a graded quotient of bounded-below \(M\), so (1) gives \(Q=0\). For (4), assign \(d_i=\deg(x_i)\) to every nonzero generator and any specified degree to a zero generator. The original degree-zero map \(\bigoplus_i S(-d_i)\to M\), \((s_i)\mapsto\sum_i s_i x_i\), has surjective reduction, so (3) applies. If the generating list is empty, its reduction condition is \(M/S_+M=0\), already covered by (1).

Boundedness cannot simply be discarded. For \(S=k[t]\), \(\deg(t)=1\), and the original graded module \(M=k[t,t^{-1}]\), one has \(S_+M=tM=M\ne0\), with nonzero components in arbitrarily negative degrees. The map \(0\to M\) becomes surjective modulo \(S_+\) but is not surjective. Thus the stated weaker hypothesis has a necessary role in the proof. No broader classification is asserted.

\subsubsection{Integral-closure wording}\label{algebra-proof-002-integral-closure-wording}

Source: \texttt{lemma-integral-closure-graded}, lines 13307--13355; report \texttt{OCC-01375}. The ring \(C\) consists of elements of \(S[t,t^{-1}]\) integral over the image of \(R[t,t^{-1}]\). The conventional expression is therefore ``the integral closure of \(R[t,t^{-1}]\) in \(S[t,t^{-1}]\)''. This corrects the description of the base and ambient ring; it does not assert that the original author's intended construction was different.

With the source's degree-zero variable \(t\), the displayed ring automorphism sends \(\sum_j s_j\) to \(\sum_j t^j s_j\) on homogeneous components and fixes \(t\). Its inverse sends each \(s_j\) to \(t^{-j}s_j\) and again fixes \(t\). Applying the two maps to each generator gives the identity. The analogous maps on \(R[t,t^{-1}]\) commute with the given structural homomorphism, so applying the automorphism to any monic integral equation gives another monic equation over the same base. Thus it preserves exactly the described \(C\).

The later descent in the source is also valid: for \(m>i>0\), the quotient by \(t^m-t^i-1\) is free with basis \(1,t,\ldots,t^{m-1}\) over its coefficient ring, by division by this monic polynomial. In that quotient \(t(t^{m-1}-t^{i-1})=1\), so adjoining \(t^{-1}\) does not change the quotient. The inclusion of the original coefficient ring is injective, since a nonzero multiple of the monic degree-\(m\) polynomial cannot be a nonzero constant. The finite base extension and integral transitivity used at lines 13348--13352 therefore return an actual monic equation for the original \(s_i\) in \(S\). No factor or coefficient-ring map is discarded.

\subsubsection{Proper homogeneous prime test}\label{algebra-proof-002-proper-homogeneous-prime-test}

Source: lines 13369--13380. The printed test needs the ideal to be \textbf{proper}. The whole ring satisfies the implication about products of homogeneous elements but is not a prime ideal.

Here is the exact test, also valid for a commutative \(\mathbf Z\)-graded ring. If a proper homogeneous ideal \(J\) has the property that a product of homogeneous elements belongs to \(J\) only when one factor belongs to \(J\), then \(S/J\) has no zero divisors. Indeed, two nonzero elements of \(S/J\) have finite homogeneous supports. Let their highest nonzero components have degrees \(a\) and \(b\). The degree-\(a+b\) component of the product is precisely the product of those two components: any term of lower degree in one factor cannot be paired with a degree exceeding the other factor's maximum. That component is nonzero by the homogeneous-product property. Thus the quotient is a nonzero domain and \(J\) is prime. The converse follows directly from primality. The minimal proposed wording inserts ``proper''.

This also supplies the precise input to the source's \(\mathbf Z\)-graded localization argument at lines 13410--13437. For a prime \(\mathfrak p_0\subset S_0\), \(I=\mathfrak p_0S\) is homogeneous, and taking degree zero gives \(I\cap S_0=\mathfrak p_0\). Its radical is homogeneous: in the quotient by any homogeneous ideal, if \(z=\sum z_j\) is nilpotent, its highest component is nilpotent by taking the highest degree of a power. Subtract that component and continue through its finite support. The subtraction remains nilpotent because a sum of two commuting nilpotents is nilpotent, as the binomial expansion shows. Hence every component belongs to the radical.

Write \(f\in S_d\), \(d>0\), for the actual homogeneous unit. For homogeneous \(a,b\) with \(ab\in\sqrt I\), choose \(N\geq1\) such that \((ab)^N\in I\). Then

\[
\left(a^{Nd}/f^{N\deg(a)}\right)
\left(b^{Nd}/f^{N\deg(b)}\right)\in I\cap S_0=\mathfrak p_0.
\]

One factor belongs to \(\mathfrak p_0\), so \(a\in\sqrt I\) or \(b\in\sqrt I\). Also \(\sqrt I\cap S_0=\mathfrak p_0\), making \(\sqrt I\) proper. The corrected prime test proves it is prime. Conversely, for a homogeneous prime \(\mathfrak p\) contracting to \(\mathfrak p_0\), any homogeneous \(a\in\mathfrak p\) has \(a^d/f^{\deg(a)}\in\mathfrak p_0\), so \(a\in\sqrt{\mathfrak p_0S}\). The reverse inclusion follows from \(\mathfrak p_0S\subset\mathfrak p\); thus the two prime correspondences are inverse. Negative degrees cause no problem because the original \(f\) is a unit. The source's open-image formula follows from \(a\notin\mathfrak p\) if and only if \(a^d/f^{\deg(a)}\notin\mathfrak p_0\), and homogeneity handles the finite sum of components. This explains why the properness correction is compatible with the receiving homeomorphism, without replacing that proof in a translation.

\subsubsection{Degree-zero localization wording}\label{algebra-proof-002-degree-zero-localization-wording}

Source: lines 13402--13408; report \texttt{OCC-01376}. ``Submodule'' is a single word. Its scalar ring remains the explicitly specified \(S_{(f)}\). This copyedit does not assert that the degree-zero part is an \(S_f\)-submodule, which need not be true.

\subsubsection{Homogeneous core of a prime}\label{algebra-proof-002-homogeneous-core-of-a-prime}

Source: \texttt{lemma-smear-out}, lines 13650--13667; reports \texttt{OCC-00375} and \texttt{OCC-01377}. Let \(\mathfrak q\) be the ideal generated by the homogeneous elements of the original prime \(\mathfrak p\). Every generator is in \(\mathfrak p\), so \(\mathfrak q\subset\mathfrak p\), making \(\mathfrak q\) proper. If homogeneous \(f,g\) satisfy \(fg\in\mathfrak q\), then \(fg\in\mathfrak p\), so one factor is in \(\mathfrak p\). That homogeneous factor is a generator of \(\mathfrak q\). The proper homogeneous prime test just proved shows that \(\mathfrak q\) is prime. The two small corrections reverse the printed containment and supply the missing verb ``is''.

In the receiving minimal-prime lemma, a minimal prime \(\mathfrak p\) contains the prime \(\mathfrak q\), so minimality gives equality and hence homogeneity. For a homogeneous ideal \(I\), the same argument in the graded quotient \(S/I\) applies to the prime \(\mathfrak p/I\) minimal over zero. Its inverse image \(\mathfrak p\) is homogeneous. Report \texttt{OCC-01378} corrects only the sentence fragment ``The second because'' to ``The second follows because''; the source already uses the correct containment in this receiving argument.

\subsubsection{Negative degrees in the finite generator list}\label{algebra-proof-002-negative-degrees-in-the-finite-generator-list}

Source: \texttt{lemma-dehomogenize-finite-type}, lines 13685--13727. Keep the original ring generators \(f_1,\ldots,f_n\), their degrees \(a_i\geq0\), the original homogeneous denominator \(f\) of degree \(d>0\), and module generators \(x_j\) of arbitrary integer degrees \(b_j\). The final finite generator list needs a convention when its resulting exponent \(e\) is negative. For \(S=k[t]\), \(f=t\), and \(M=S(1)\) with generator \(x\) of degree \(-1\), that list forces \(e=-1\), although \(M_{(f)}\ne0\) and is generated by \(tx\). Reading the list as requiring only nonnegative \(e\) would incorrectly give an empty list.

For an original monomial fraction \(f_1^{E_1}\cdots f_n^{E_n}x_j/f^e\) of degree zero, initially with \(e\geq0\), write \(E_i=q_i d+r_i\), \(0\leq r_i<d\), and retain \(e_0=e-\sum_i q_i a_i\). The exact identity in the original localization is

\[
\frac{f_1^{E_1}\cdots f_n^{E_n}x_j}{f^e}
=\left(\prod_i\left(\frac{f_i^d}{f^{a_i}}\right)^{q_i}\right)
\frac{f_1^{r_1}\cdots f_n^{r_n}x_j}{f^{e_0}},
\qquad de_0=\sum_i r_i a_i+b_j.
\]

Each factor in parentheses belongs to the original \(S_{(f)}\). If \(e_0<0\), the last term means the actual element \(f^{-e_0}f_1^{r_1}\cdots f_n^{r_n}x_j\) with denominator 1; its degree is zero by the displayed equality. All factors are retained. There are at most \(m d^n\) choices of \((j,r_1,\ldots,r_n)\), and for each choice the equation determines at most one integral \(e_0\). Expanding a homogeneous numerator in the original finite generators gives a finite \(R\)-linear combination of these fractions, so the list generates \(M_{(f)}\). This proves the asserted finiteness and the explicit bound, including an empty module generating list. A zero localization gives the zero module and satisfies the same bound. The minimal proposed clarification states the negative-exponent convention beside the list; it does not change the module or its grading.

\subsubsection{Positive-degree homogenization and the original base ring}\label{algebra-proof-002-positive-degree-homogenization-and-the-original-base-ring}

Source: \texttt{lemma-homogenize}, lines 13729--13769; reports \texttt{OCC-01309}, \texttt{OCC-00376}, \texttt{OCC-01379}, and \texttt{OCC-01380}. Keep the original presentation \(R'=P/I\), \(P=R[x_1,\ldots,x_n]\), and \(H=R[X_0,\ldots,X_n]\) with every \(X_i\) of degree 1. The printed minimal-degree homogenization fails to preserve \(S_0=R\): for \(R=\mathbf Z\), \(n=0\), and \(I=(2)\), it gives \(S=\mathbf Z[X_0]/(2)\), whose degree-zero ring is \(\mathbf Z/(2)\).

For every polynomial \(g=\sum_\alpha c_\alpha x^\alpha\), choose an integer \(D(g)\geq1\) at least as large as every total degree \(|\alpha|\) in its finite support. For the zero polynomial the empty support allows any specified positive integer. Define

\[
\widetilde g=\sum_\alpha c_\alpha
 X_0^{D(g)-|\alpha|}X_1^{\alpha_1}\cdots X_n^{\alpha_n}.
\]

This retains every coefficient and exponent, is homogeneous of its assigned degree \(D(g)\), and satisfies \(\widetilde g(1,x_1,\ldots,x_n)=g\). The zero polynomial is homogeneous of the assigned degree without assigning an intrinsic polynomial degree to zero. Set \(\widetilde I=(\widetilde g:g\in I)\), \(S=H/\widetilde I\), and let \(f\) be the image of \(X_0\). All generators of \(\widetilde I\) have positive degree, so every homogeneous part of any product used to generate this ideal has positive degree. Thus \(\widetilde I\cap H_0=0\), and the original structural map gives \(S_0=R\), even when \(R\to R'\) is not injective.

There are exact inverse graded ring maps

\[
H_{X_0}\longrightarrow P[t,t^{-1}],\quad X_0\mapsto t,\quad X_i\mapsto t x_i,
\] \[
P[t,t^{-1}]\longrightarrow H_{X_0},\quad t\mapsto X_0,\quad x_i\mapsto X_i/X_0,
\]

where \(P\) has degree zero and \(t\) has degree 1. Each composite fixes the original coefficients, variables, and the inverse variable. Under these maps \(\widetilde g\) becomes \(t^{D(g)}g\). Since that displayed power is a unit after localization, the localized ideal generated by all \(\widetilde g\) is exactly \(I P[t,t^{-1}]\). Passing to quotients gives \(S_f\cong R'[t,t^{-1}]\), and taking degree zero gives the original dehomogenization map \(S_{(f)}\cong R'\), \(X_i/X_0\mapsto x_i\). This includes \(R'=0\) and a zero image of \(f\). The ring \(S\) remains generated by the original finite list of degree-one variable images over \(R\).

For the module, choose a presentation with \(r\geq1\),

\[
M=(R')^{\oplus r}/\sum_{j\in J}R'k_j,\qquad
k_j=(k_{1j},\ldots,k_{rj}).
\]

Every finite module admits this choice; for the zero module one may use \(r=1\). Choose actual polynomial lifts \(h_{ij}\in P\) mapping to the original \(k_{ij}\in R'\). Assign positive homogenization degrees \(d_{ij}\) as above, including zero entries. Let \(d_j=\max_{1\leq i\leq r}d_{ij}\), and define the original presentation column with its corrected, now defined coefficients

\[
K_{ij}=X_0^{d_j-d_{ij}}\widetilde h_{ij},\qquad
N=\operatorname{Coker}\left(\bigoplus_{j\in J}S(-d_j)
 \xrightarrow{(K_{ij})} S^{\oplus r}\right).
\]

Every exponent is nonnegative and every column has its specified degree \(d_j\), so the map has degree zero with the original twist convention. The module \(N\) is finite because it is a quotient of \(S^{\oplus r}\); the set \(J\) need not be finite. Under \(S_f\cong R'[t,t^{-1}]\), the coefficient \(K_{ij}\) becomes \(t^{d_j}k_{ij}\). The degree-zero generator of \((S(-d_j))_f\) is \(t^{-d_j}\) times its original basis element, which has degree \(d_j\). The induced degree-zero matrix therefore has entries exactly \(k_{ij}\), with both displayed powers retained and multiplied to 1.

Localization preserves cokernels by the numerator-and-denominator definition. Taking degree zero preserves cokernels of degree-zero maps: a homogeneous degree-zero element in the image is the image of the degree-zero component of any preimage. Taking degree zero also commutes with the direct sum, since every element has finite support. Consequently the localized degree-zero cokernel is exactly \((R')^{\oplus r}/\sum_j R'k_j=M\). This proves the specified \(R'\)-linear isomorphism \(N_{(f)}\cong M\), together with every asserted property of the corrected construction. The other two reports insert the missing ``be'' and correct ``denote by''. These changes and the polynomial-lift choice are linked to the original passage, not attributed silently to the source author.

\subsubsection{Noetherian and numerical wording}\label{algebra-proof-002-noetherian-and-numerical-wording}

Source: lines 13812--13830. Report \texttt{OCC-00377} inserts the missing comma in \(S_0[X_1,\ldots,X_n]\). A nearby copyedit changes ``sufficient large'' to ``sufficiently large''. Neither changes a hypothesis, coefficient group, or the integer-valued binomial basis.

\subsubsection{Nilpotence induction in the Hilbert proof}\label{algebra-proof-002-nilpotence-induction-in-the-hilbert-proof}

Source: \texttt{proposition-graded-hilbert-polynomial}, lines 13889--13944; report \texttt{OCC-00378}. In the inner induction with least positive \(r>1\) satisfying \(x^rM=0\), take the actual graded submodule \(M'=xM\) and quotient \(M''=M/xM\). The maps are its inclusion and quotient map, giving the short exact sequence used by the source. The module \(M'\) is finite, generated by the products of \(x\) with any finite generating list of \(M\); \(M''\) is a finite quotient. The degrees are inherited, so taking each degree is exact. One has \(x^{r-1}M'=x^rM=0\) and \(xM''=0\). Their least positive annihilating exponents therefore satisfy \(r'\leq r-1<r\) and \(r''=1<r\), including zero modules under the least-positive convention. This supplies the previously unnamed objects and exponents without changing the proof's induction.

The outer reduction retains its original Noetherian hypothesis. The ascending chain \(\ker(x:M\to M)\subset\ker(x^2)\subset\cdots\) stabilizes because \(M\) is Noetherian. Its union is the graded submodule \(T\) of all \(x\)-power torsion and equals \(\ker(x^N)\) for some \(N\). Thus \(x^NT=0\). If \(xm\in T\), then \(x^{N+1}m=0\), so stability gives \(m\in T\). Hence multiplication by \(x\) on \(M/T\) is injective, as required by the receiving exact degree sequence. No new assumption or replacement theorem is needed. In degree \(d\), its maps give exactly \([M_{d+1}]-[M_d]=[(M/xM)_{d+1}]\) after that reduction; the original binomial antidifference step applies with its original sign and shift.

\subsubsection{Hilbert bound notation}\label{algebra-proof-002-hilbert-bound-notation}

Source: \texttt{lemma-quotient-smaller-d}, lines 13964--13982; report \texttt{OCC-00379}. The proposed change replaces the literal \texttt{\textgreater{}\textgreater{}} with \texttt{\textbackslash{}gg}, preserving ``sufficiently large''. Multiplication by the original nonzero homogeneous \(f\) in the polynomial domain is injective, so its degree-\(n\) image has the original dimension \(\binom{n-e-1+d}{d-1}\). Subtracting this lower bound for \(\dim I_n\) from \(\binom{n-1+d}{d-1}\) gives the displayed upper bound. For \(d\geq2\), the two degree-\(d-1\) polynomials have the same leading coefficient \(1/(d-1)!\), so their difference has smaller degree; a nonnegative eventual dimension polynomial bounded by it also has smaller degree. For \(d=1\) the difference is zero. If the allowed nonzero ideal contains a degree-zero unit, the quotient is zero directly. For zero variables, a nonzero ideal of the field is the whole field and the quotient is again zero, without using a binomial coefficient with lower entry \(-1\). These boundary explanations remain editorial material.

\subsubsection{Reading and propagation limits}\label{algebra-proof-002-reading-and-propagation-limits}

The other complete source units in this interval were read for their surrounding argument. In particular the Proj chart/localization formulas, the finite-type ring generator argument, and the receiving minimal-prime and Hilbert calculations were compared with the findings above. This is not a new certification of every omitted detail or of previously inserted FAC proofs. The original Veronese argument uses a positive number of positive-degree algebra generators; when that number is zero, \(S=S_0\) and its conclusion holds directly for every positive \(d\). That boundary explanation does not require expanding its translated proof. The periodic-polynomial remark is read in the preceding finite-module context. Broader combined-results synthesis and the short mathematical paper remain after the core correction work. Source from 13991 onward is not adjudicated by this batch.

\hypertarget{algebra_local_hilbert_dimension_notes_20260927}{}
\subsection{Editorial evidence for Algebra 13991--14710}\label{algebra-proof-003-editorial-evidence-for-algebra-1399114710}

Authority: \texttt{algebra.tex} at commit \texttt{a04446e57ec1fbc252a871afcec7752fb2807b14}, SHA-256 \texttt{FA8BB92E58A4F78A2BD01B3B6A4A87DE0A0D279F5DD90641B574DD5FBFFFA4F3}. The complete interval and its twelve reports were read. These arguments support separately recorded minimal corrections, retained source statements, and editorial consequences. They are not replacement proof bodies for translations, and no novelty is claimed.

The earlier files \texttt{algebra\_length\_restriction\_action\_20260927.tex} and \texttt{algebra\_semilocal\_nil\_localization\_product\_20260927.tex} were reread in full for the scalar-restriction and local-factor issue below. They are the complete proof files for review groups 351 and 363. The original and current \texttt{topology.tex} passages labelled \texttt{definition-Krull} were also read: authority lines 1391--1418, current lines 1450--1477. Both explicitly assign dimension \(-\infty\) to the empty space.

\subsubsection{Finite-length wording}\label{algebra-proof-003-finite-length-wording}

Source: lines 14026--14042; report \texttt{OCC-00380}. The proposed ``has finite length'' wording is a small copyedit. The finite-length assertion is mathematically unchanged. If \(\mathfrak m^r\subset I\) and \(I^aM=0\), then \(\mathfrak m^{ra}M\subset I^aM=0\). Every successive \(\mathfrak m\)-adic layer of the finite module is a finite-dimensional vector space over the original residue field. There are only finitely many nonzero layers, so their original \(R\)-lengths sum to a finite length. No scalar field, ideal, or exponent is replaced.

\subsubsection{Finite-colength comparison and explicit constants}\label{algebra-proof-003-finite-colength-comparison-and-explicit-constants}

Source: \texttt{lemma-differ-finite}, lines 14044--14085; reports \texttt{OCC-00381} and \texttt{OCC-12300} identify the same ``there exists a constants'' wording. One correction accounts for both reports.

Let \(Q=M/M'\) and \(l=\operatorname{length}_R(Q)\). In addition to correcting the wording, one can make the source's constants explicit: \textbf{\(c_1=c_2=l\) works}, retaining the original \(I\), \(M'\), and \(M\). Here is the proof. Every simple factor of a finite-length module over the original local ring is its residue field and is annihilated by \(\mathfrak m\). Along a composition series \(0=Q_0\subset Q_1\subset\cdots\subset Q_l=Q\), multiplication by \(\mathfrak m\) sends \(Q_j\) into \(Q_{j-1}\). Hence \(\mathfrak m^lQ=0\). Since an ideal of definition is contained in \(\mathfrak m\), one has \(I^lQ=0\), or \(I^lM\subset M'\). For \(l=0\), the quotient is zero and this inclusion is \(M=M'\).

For \(n\geq l\), the exact inclusion and quotient maps give

\[
0\longrightarrow M'/I^{n+1}M\longrightarrow M/I^{n+1}M
\longrightarrow Q\longrightarrow0.
\]

Thus \(\chi_{I,M}(n)=l+\operatorname{length}_R(M'/I^{n+1}M)\). The original submodules satisfy

\[
I^{n+1}M'\subset I^{n+1}M
 =I^{n+1-l}(I^lM)\subset I^{n+1-l}M'\subset M'.
\]

Their quotient maps are surjective in the reverse order. Additivity of length yields the exact bounds

\[
l+\chi_{I,M'}(n-l)\leq\chi_{I,M}(n)
\leq l+\chi_{I,M'}(n)\qquad(n\geq l).
\]

This explicit choice is an editorial consequence; the source's existence statement is retained.

\subsubsection{Artin--Rees endpoint}\label{algebra-proof-003-artinrees-endpoint}

Source: \texttt{lemma-hilbert-ses}, lines 14087--14132. The statement, including the endpoint \(n=c\), is valid. The displayed intermediate proof at lines 14119--14127 uses \(\chi_{I,N}(-1)\) when its original index equals \(c\), although the function was introduced only for nonnegative arguments. State the convention \(\chi_{I,Q}(-1)=0\) for every finite module \(Q\). It agrees with the original quotient formula because \(I^0Q=Q\), so \(\operatorname{length}_R(Q/I^0Q)=0\). It does not change any previously defined value.

Write the original exact sequence as \(0\to M'\xrightarrow{u}M\xrightarrow{v}M''\to0\), identifying \(M'\) with its image, and set \(N=M'\cap I^cM\). Artin--Rees gives \(M'\cap I^rM=I^{r-c}N\) for \(r\geq c\). Since \(I^cM'\subset N\), the quotient \(M'/N\) has finite length \(l\). For \(n\geq c\), the quotient sequence and the intermediate quotient by \(N\) give

\[
\chi_{I,M}(n)=\chi_{I,M''}(n)
 +\operatorname{length}_R(M'/I^{n+1-c}N)
=\chi_{I,M''}(n)+l+\chi_{I,N}(n-c).
\]

There is also a direct proof of the \(\varphi\)-formula at every stated index, including the endpoint. The exact degree sequence is

\[
0\longrightarrow I^{n-c}N/I^{n+1-c}N
\longrightarrow I^nM/I^{n+1}M
\longrightarrow I^nM''/I^{n+1}M''\longrightarrow0.
\]

The first map is the original inclusion; its kernel is zero because \(I^{n-c}N=M'\cap I^nM\) and \(M'\cap I^{n+1}M=I^{n+1-c}N\). The last map is induced by \(v\) and is surjective. If \(m\in I^nM\) maps into \(I^{n+1}M''\), lift that image to an element of \(I^{n+1}M\); subtracting that lift puts \(m\) in \(M'\cap I^nM\) modulo \(I^{n+1}M\). This proves exactness and, by length additivity, the source's \(\varphi\)-formula. At \(n=c\), the left term is exactly \(N/IN\), as required.

\subsubsection{Degree comparison and highest coefficients}\label{algebra-proof-003-degree-comparison-and-highest-coefficients}

Source: \texttt{lemma-differ-finite-chi} and \texttt{lemma-hilbert-ses-chi}, lines 14212--14269; report \texttt{OCC-00382} supplies the missing article before ``degree''. The source identifies the two polynomials by the immediately preceding \(M\) and \(M'\); no new mathematical hypothesis is needed for this copyedit.

Here are the precise polynomial consequences used in the receiving argument. For any nonzero finite module \(Q\), \(\chi_{I,Q}(n)\) is positive for all \(n\geq0\): its quotient maps onto the nonzero \(Q/IQ\), which is nonzero by Nakayama because \(I\subset\mathfrak m\). Consequently its eventual polynomial has a positive leading coefficient. Its degree is zero exactly when \(Q\) has finite length. One direction is immediate from a power of \(I\) killing a finite-length module. For the reverse direction, eventual constancy makes \(\operatorname{length}_R(I^nQ/I^{n+1}Q)=0\) for large \(n\), so \(I^nQ=I I^nQ\). The finite module \(I^nQ\) is zero by Nakayama. A module killed by that power of the ideal of definition has finite length, by the earlier layer argument. Thus a module without finite length has a positive-degree eventual polynomial.

In the finite-colength comparison, if \(M\) lacks finite length, then so does \(M'\); otherwise additivity in \(0\to M'\to M\to Q\to0\) would give finite length to \(M\). Let \(P(t)\) and \(P'(t)\) be their original eventual polynomials. The proved bounds give

\[
l+P'(n-l)\leq P(n)\leq l+P'(n)
\]

for all sufficiently large integers. Since \(P'\) has positive leading coefficient, these two bounding polynomials have the same positive degree and leading coefficient. If \(P\) had a larger or smaller degree, or a different leading coefficient at that degree, the corresponding polynomial difference would eventually have a sign contradicting one of the inequalities. Hence the two original polynomials have the same degree and leading coefficient, and \(P-P'\) has strictly smaller degree. This proves the receiving comparison without replacing either original polynomial by a rescaled one.

For the original exact sequence, retain its \(N,c,l\) from Artin--Rees and let \(P_Q\) denote each original eventual polynomial. The full error identity is

\[
P_M(t)-P_{M'}(t)-P_{M''}(t)
=\bigl(P_N(t-c)-P_N(t)\bigr)
 +\bigl(P_N(t)-P_{M'}(t)\bigr)+l.
\]

If \(M'\) has infinite length, its degree \(d'\) is positive; the finite-colength comparison applies to \(N\subset M'\). Each parenthesized term has degree strictly less than \(d'\). For the first, retain the original expansion \(P_N(t)=\sum_j a_jt^j\); the term is \(\sum_j a_j((t-c)^j-t^j)\), whose degree-\(j\) coefficients cancel for each \(j\). The constant \(l\) also has smaller degree. If \(M'\) has finite length, choose \(b\) with \(I^bM'=0\). Artin--Rees gives \(M'\cap I^rM=I^{r-c}N=0\) for \(r\geq c+b\), so the quotient exact sequence gives \(P_M-P_{M''}=\operatorname{length}_R(M')=P_{M'}\) exactly.

Thus one further consequence is proved without changing the source: for any such exact sequence with a nonzero term, put \(D=\max\{d(M'),d(M'')\}\). Then

\[
[t^D]P_M=[t^D]P_{M'}+[t^D]P_{M''}.
\]

For \(D\geq1\), this follows from the full error formula and its degree bound, including a finite-length \(M'\). For \(D=0\), both end terms have finite length, so all three polynomials are their actual lengths and the equality is the original exact-sequence additivity. If both end terms are zero, the middle term is zero and all three polynomials vanish; no coefficient with index \(-\infty\) is used. This is a highest-coefficient identity of the original polynomials, not a replacement of them by a normalized invariant. Zero polynomials have degree \(-\infty\) in these comparisons.

\subsubsection{Punctuation in the dimension definition}\label{algebra-proof-003-punctuation-in-the-dimension-definition}

Source: lines 14331--14343; report \texttt{OCC-12301}. Remove the period at the end of the displayed prime chain so that ``of length \(n\)'' continues its sentence. The strict inclusions and chain length do not change.

\subsubsection{Minimal prime wording}\label{algebra-proof-003-minimal-prime-wording}

Source: lines 14369--14378; report \texttt{OCC-00383}. An irreducible component \(V(\mathfrak p_i)\) corresponds to a minimal \textbf{prime} ideal. Minimality is within the prime ideals, because a smaller prime gives a larger irreducible closed subset and contradicts maximality of that component. ``Minimal ideal'' would ordinarily describe a different condition among all ideals; insert ``prime''.

\subsubsection{Length over the original ring is correct}\label{algebra-proof-003-length-over-the-original-ring-is-correct}

Source: lines 14382--14390; report \texttt{OCC-00384}. Retain the source's assertion that each \(A_i=R_{\mathfrak p_i}\) has finite length \textbf{over \(R\)}. It is true, and the preceding product decomposition already supplies the required reason. The report correctly observes that the cited local finite-length lemma initially gives finiteness over \(A_i\); its suggested objection does not account for the immediately preceding isomorphism.

The original canonical isomorphism \(\theta:R\to\prod_i A_i\) sends \(r\) to \((r/1)_i\). Its inverse, fully proved in \texttt{algebra\_semilocal\_nil\_localization\_product\_20260927.tex}, sends \((r_i/s_i)_i\) to \(\sum_i(e_i r_i)(e_i s_i)^{-1}\), using the original orthogonal idempotents and the inverses in \(e_iR\). In particular each original localization map \(R\to A_i\) is surjective: apply \(\theta^{-1}\) to the tuple with the requested element in slot \(i\) and zero elsewhere.

For any \(R\)-submodule \(L\subset A_i\), \(a\in A_i\), and \(x\in L\), choose \(r\in R\) with image \(a\). Then \(ax=rx\in L\). Hence every such \(L\) is an \(A_i\)-submodule, and the converse is immediate by restriction of scalars. The full submodule lattices and all strict chains coincide. Therefore

\[
\operatorname{length}_R(A_i)=\operatorname{length}_{A_i}(A_i)<\infty.
\]

The canonical product decomposition is also an \(R\)-module isomorphism; the product is finite, hence is its direct sum. Length additivity gives \(\operatorname{length}_R(R)=\sum_i\operatorname{length}_{A_i}(A_i)<\infty\), exactly the receiving conclusion needed by the source. The empty product is the zero ring with zero length. The same argument applies to every module over each factor, not just the factor itself; that stronger scalar-restriction statement is already proved in group 351 and is reused here. No source replacement or duplicate generalization is needed.

\subsubsection{Zero ring and dimension at most zero}\label{algebra-proof-003-zero-ring-and-dimension-at-most-zero}

Source: \texttt{lemma-Noetherian-dimension-0} at lines 14363--14397 and \texttt{proposition-dimension-zero-ring} at lines 14414--14443. The original topological definition explicitly gives \(\dim\varnothing=-\infty\); it is unchanged in the current chapter. Thus the zero ring is Artinian and Noetherian, with length zero, but its Krull dimension is \(-\infty\), not zero. This is a real exception to the printed converse and to item (2) of the ten-condition proposition.

The minimal corrections are ``dimension at most zero'' in the converse and \(\dim(R)\leq0\) in item (2). They preserve the original dimension convention. To verify the corrected scope, a nonzero ring has a maximal ideal, so its spectrum is nonempty and dimension is at least zero. Consequently the new bound is exactly dimension zero on every nonzero ring, where the source argument applies. For the zero ring, the bound holds, both chain conditions are automatic, and its sole unital module is zero. Every finite-product condition in the proposition holds using the empty product; the finite discrete spectrum and absence of strict prime inclusions hold for the empty space; and the Jacobson radical is the zero ideal, hence nilpotent. Thus all ten conditions remain equivalent, now including the original zero ring.

The subsequent local lemmas are unchanged: a local ring in those statements has its specified maximal ideal and residue field and is nonzero. Likewise every localization at a prime is nonzero. Hence the local assertions \(\dim(R)=0\iff d(R)=0\) and the height statements still use equality with zero. Do not propagate the new weak inequality into those different, already correct statements.

\subsubsection{One-dimensional wording}\label{algebra-proof-003-one-dimensional-wording}

Source: lines 14482--14487; reports \texttt{OCC-00385} and \texttt{OCC-12302} identify the same misplaced ``directly''. One operation changes ``from directly'' to ``directly from''. The equivalence of the two ideal-of-definition descriptions is unchanged.

\subsubsection{Multiplication maps in both dimension proofs}\label{algebra-proof-003-multiplication-maps-in-both-dimension-proofs}

Source: lines 14506--14517 and 14579--14592; report \texttt{OCC-00386} names the first occurrence. The identical later sequence requires the same visible map label, so the correction is propagated to both.

Since the original \(\mathfrak p\) is prime and \(x\notin\mathfrak p\), the element \(\bar x\) is nonzero in the original domain \(R/\mathfrak p\). Define

\[
\mu_x:R/\mathfrak p\longrightarrow R/\mathfrak p,
\qquad a+\mathfrak p\longmapsto xa+\mathfrak p.
\]

Changing a representative by an element of \(\mathfrak p\) changes its product by an element of \(\mathfrak p\), so the map is defined and \(R\)-linear. Its kernel is zero by primality and \(x\notin\mathfrak p\); its image is \((xR+\mathfrak p)/\mathfrak p\). The quotient map onto \(R/(xR+\mathfrak p)\) has exactly this kernel and is surjective, proving the displayed short exact sequence. It is generally not multiplication by 1: \(x\in\mathfrak q\subset\mathfrak m\) is retained throughout. Both source operations label the first nonzero arrow by \(x\).

The two polynomial contributions from \(R/\mathfrak p\) are those of the same original module with the same original ideal action, so their difference is zero and the full signed expression remains \(-\chi_{I,R/(xR+\mathfrak p)}\), as printed. The quotient is nonzero since \(xR+\mathfrak p\subset\mathfrak q\) is proper. In the first proof it has degree zero and hence finite length; the two displayed distinct primes contradict dimension zero. In the general proof its polynomial degree is strictly smaller than that of \(R/\mathfrak p\), hence smaller than \(d(R)\); the induction therefore applies to that exact quotient.

There is no mismatch between degrees calculated over \(R\) and over a quotient ring \(B=R/J\): the structural map is surjective, its action on each module is unchanged, and the submodule lattices coincide by the same scalar argument above. Also \(I^nB=(IB)^n\) as submodules of \(B\), and \(IB\) is an ideal of definition whenever \(J\subset\mathfrak m\). Thus every corresponding \(\chi\)-value is exactly equal, with no suppressed length factor.

\subsubsection{Closing sentence and induction scope}\label{algebra-proof-003-closing-sentence-and-induction-scope}

Source: lines 14594--14596; report \texttt{OCC-12303}. Retain ``Reading back the reader will see we proved the circular inequalities as desired.'' The participial phrase has the stated reader as its subject; ``reading back'' means reviewing the preceding argument, and the complement after ``see'' may omit ``that''. A comma or the proposed longer sentence is an optional style change, not a demonstrated grammatical or mathematical defect. The term ``circular inequalities'' describes the cycle of inequalities printed at lines 14544--14546, not a claim that a premise is assumed from its conclusion.

For the mathematical scope, the source's induction can be read exactly as follows. Assume the equivalence for every smaller integer. If \(\dim R=d\), its prime-avoidance quotient has dimension at most \(d-1\), so induction yields \(d'(R)\leq d\). Independently, an ideal of definition with \(d\) generators gives the full monomial length bound and hence \(d(R)\leq d\). If \(d(R)=d\), the labelled injection and the nonzero quotient in the preceding section give \(\dim R\leq d\) by induction on the quotient's smaller numerical degree; arbitrary finite chains are tested, so an infinite Krull dimension has not been assumed away. For any of the three invariants equal to \(d\), a strict drop of another to a smaller integer would, by that integer's induction equivalence, force all three to equal the smaller integer and contradict the premise. This proves equality. In addition \(d(R)\) is a finite nonnegative integer for a nonzero Noetherian local ring, by its already proved numerical polynomial, so the argument covers every such ring. This explanation stays in editorial material.

\subsubsection{Elements generating the maximal ideal}\label{algebra-proof-003-elements-generating-the-maximal-ideal}

Source: lines 14599--14608; report \texttt{OCC-00387}. The original maximal ideal is generated by elements of its original ring, which need not be specified polynomial variables. Replace ``variables'' by ``elements''. Nakayama identifies the minimal number with the dimension of the original residue vector space \(\mathfrak m/\mathfrak m^2\); the printed inequality and its original residue field remain unchanged.

\subsubsection{Receiving height results and limits}\label{algebra-proof-003-receiving-height-results-and-limits}

The complete proofs through \texttt{lemma-one-equation} at line 14710 were read. For a prime minimal over a given generated ideal, the original localization maps identify primes with the primes below it and show that the localized ideal has radical equal to the localized maximal ideal. The dimension bound uses that exact localized ring; the zero-ring correction does not change it. For primes above an original prime \(\mathfrak p\), quotient and localization give the same strict-chain correspondence used in the two height statements. The one-equation comparison likewise retains the original quotient and parameter lifts. No source changes or new certification of omitted details are asserted for these receiving statements. Source 14711--14720 is only the next lemma's statement and proof opening, and remains unadjudicated. Broader synthesis of the collected findings and its short paper remain after the core task.

\hypertarget{algebra_dimension_support_associated_notes_20260927}{}
\subsection{Editorial evidence for Algebra 14711--15541}\label{algebra-proof-004-editorial-evidence-for-algebra-1471115541}

Authority: original \texttt{algebra.tex}, commit \texttt{a04446e57ec1fbc252a871afcec7752fb2807b14}, SHA-256 \texttt{FA8BB92E58A4F78A2BD01B3B6A4A87DE0A0D279F5DD90641B574DD5FBFFFA4F3}. The complete interval and all seventeen reports whose first loci fall there were read. Minimal proposed corrections are recorded separately from these supporting arguments and extensions. The original source's editorial reminders remain identifiable. These are standard source-review consequences, with no novelty claim.

\subsubsection{Parameter quotients and the infinite-open argument}\label{algebra-proof-004-parameter-quotients-and-the-infinite-open-argument}

Source: \texttt{lemma-elements-generate-ideal-definition}, lines 14712--14723, and \texttt{lemma-Noetherian-local-domain-dim-2-infinite-opens}, lines 14738--14765. These statements are retained. For the first, let \(J_i=(x_1,\ldots,x_i)\). Iterating the one-equation inequality on the original quotients gives \(\dim(R/J_i)\geq d-i\). The remaining original parameters generate an ideal of definition in \(R/J_i\); the already proved dimension criterion gives \(\dim(R/J_i)\leq d-i\). Thus equality holds, also after any permutation of the original parameter list. All these quotients are nonzero because every parameter lies in the original maximal ideal. For \(d=0\) the asserted list of equalities is empty.

In the second proof, any nonempty open subset of the spectrum of a domain contains its generic point \((0)\). If such an open is finite, removing the finitely many closures of its other points leaves the generic point as an open singleton. A standard open \(D(x)\) inside that singleton has \(x\ne0\). It has \(x\in\mathfrak m\), since a unit would give the entire spectrum and the ring has dimension at least two. The domain has only the minimal prime \((0)\), so the original one-equation result gives \(\dim(R/xR)=d-1\). Lifting a parameter list and permuting it with \(y_2\) first and \(x\) second gives the two dimensions \(d-1\) and \(d-2\) printed in the source. In particular \(y_2\ne0\). If every prime containing \(y_2\) also contained \(x\), then \(V(y_2)=V(y_2,x)\), contradicting those dimensions. Hence the constructed prime contains nonzero \(y_2\) and avoids \(x\), contradicting \(D(x)=\{(0)\}\). This verifies the later Jacobson proof's required input without replacing either source proof.

\subsubsection{The p-adic example}\label{algebra-proof-004-the-p-adic-example}

Source: lines 14767--14770; report \texttt{OCC-00388}. The intended one-dimensional example is the ring of \(p\)-adic \textbf{integers} \(\mathbf Z_p\). The \(p\)-adic numbers \(\mathbf Q_p\) form a field and have dimension zero.

For the stated correction, retain the usual \(p\)-adic digit description: a nonzero \(z\in\mathbf Z_p\) has a least index \(n\) with nonzero digit and factors as \(p^n u\), where \(u\) has nonzero residue modulo \(p\). Such a \(u\) is a unit: choose an integer \(c\) with \(uc=1+pw\); completeness makes \(c\sum_{j\geq0}(-pw)^j\) an inverse. Products of nonzero elements have the sum of their least digit indices, so the ring is a domain. In a nonzero ideal choose an element of least index \(n\); its unit factor shows \(p^n\) belongs to that ideal, and every other element is divisible by \(p^n\). Thus every nonzero ideal is \((p^n)\). Its only prime ideals are \((0)\) and \((p)\): for \(n>1\), the product \(p\cdot p^{n-1}\) lies in \((p^n)\) while neither factor does. Ascending chains of these ideals stabilize because their nonnegative exponents decrease, so \(\mathbf Z_p\) is Noetherian and has dimension one. Its fraction field \(\mathbf Q_p\) has only the prime \((0)\). The proposed word change therefore repairs the example without changing the surrounding dimension theorem.

\subsubsection{Source editorial reminders}\label{algebra-proof-004-source-editorial-reminders}

Source: the final list item at line 14803 and the placement note at line 15519; reports \texttt{OCC-00389}, \texttt{OCC-12304}, \texttt{OCC-00396}, and \texttt{OCC-12310}.

The proof of \texttt{lemma-finite-type-algebra-finite-nr-primes} establishes exactly its seven mathematical conditions. ``Add more here'' is the source's editorial reminder, not an eighth proposition. The proposed treatment retains those words as a labelled, unnumbered source reminder in the list. It neither invents another theorem nor attributes an editorial deletion to the original author. The nonzero algebra hypothesis already excludes the zero-ring dimension exception found in the preceding batch; it remains unchanged.

The later sentence ``This lemma should probably be put somewhere else'' is also a source editorial note. Keep it in its original location with a visible label identifying that status. The reports do not establish an intended new location, and their suggestions do not require a source reorganization. In particular the reference to ``this lemma'' is not used as authority to move the preceding or following lemma. These choices apply the user's translation-fidelity instruction while making both notes distinguishable from mathematical assertions.

\subsubsection{The nonmaximal prime and strict localization}\label{algebra-proof-004-the-nonmaximal-prime-and-strict-localization}

Source: \texttt{lemma-noetherian-dim-1-Jacobson}, lines 14841--14888; reports \texttt{OCC-00390} and \texttt{OCC-00391}. The point \(\mathfrak p\in\operatorname{Spec}(R)\) supplied by the topological argument is a nonmaximal \textbf{prime} ideal. Insert the missing word.

Write \(A=R/\mathfrak p\), a nonzero domain. The locally closed singleton of its generic point is open: the closure of that point is all of \(\operatorname{Spec}(A)\), and any closed subset occurring in a representation of the singleton as a locally closed set must therefore be the whole space. For a prime \(\mathfrak q\supsetneq\mathfrak p\), let \(\mathfrak r=\mathfrak q/\mathfrak p\). This is nonzero. The original localization \(A_{\mathfrak r}\) is a local domain with a nonzero maximal ideal: a nonzero element of \(\mathfrak r\) remains nonzero in the domain localization and belongs to that maximal ideal. Thus its dimension is at least one. The open generic singleton pulls back to an open generic singleton in its spectrum. The infinite-open lemma excludes dimension at least two, so the dimension is exactly one.

At \(\mathfrak q=\mathfrak p\), by contrast, the localization is the original fraction field \(\operatorname{Frac}(A)\), whose dimension is zero. Both quantified occurrences in the proof must therefore specify strict containment. Since \(\mathfrak p\) is nonmaximal, maximal ideals above it exist and are strictly larger. Every localization of \(A\) at a maximal ideal has dimension one, so the supremum-of-heights formula gives \(\dim A=1\). The source hypothesis gives infinitely many primes of \(A\), all but its zero prime maximal.

For completeness of the receiving Jacobson step, in a one-dimensional Noetherian domain every prime containing a nonzero \(a\) is maximal and minimal over \((a)\); there are finitely many of these by Noetherian finiteness of minimal primes. Infinitely many maximal ideals therefore cannot all contain any one nonzero \(a\), so their intersection is zero. The other primes are maximal and are their own intersections of maximal ideals. The ring is Jacobson. Equivalently its closed points are dense. Its open generic singleton, whose point is nonmaximal, cannot meet those closed points, a contradiction. The original conclusion and each original quotient and localization are retained.

\subsubsection{Filtration wording and the original multiplicities}\label{algebra-proof-004-filtration-wording-and-the-original-multiplicities}

Source: lines 14949--15051; reports \texttt{OCC-12306}, \texttt{OCC-00392}, \texttt{OCC-12305}, and \texttt{OCC-01381}. The two reports about ``counter example'' share one copyedit. The two reports about the repeated ``Let \ldots{} as in'' opening share the two needed insertions of ``be''. The prime filtration and its original primes are unchanged.

The surrounding multiplicity argument is consistent with these corrections. Exact localization gives the filtration quotients \((R/\mathfrak p_i)_{\mathfrak p}=R_{\mathfrak p}/\mathfrak p_iR_{\mathfrak p}\). If \(\mathfrak p_i\not\subset\mathfrak p\), an element of \(\mathfrak p_i\setminus\mathfrak p\) acts as both zero and a unit, so this quotient is zero. If \(\mathfrak p\) is minimal in the support and \(\mathfrak p_i\subset\mathfrak p\), then \(\mathfrak p_i=\mathfrak p\), since all filtration primes are in that support. The quotient is then the original residue field \(\kappa(\mathfrak p)\), of length one over \(R_{\mathfrak p}\). Adding these lengths proves exactly the printed multiplicity formula. No extra residue-field factor occurs here: both the simple factors and their lengths use the same original local ring.

\subsubsection{Support and a finitely generated annihilating ideal}\label{algebra-proof-004-support-and-a-finitely-generated-annihilating-ideal}

Source: \texttt{lemma-Noetherian-power-ideal-kills-module}, lines 14996--15014. Keep the source statement unchanged. Its Noetherian-ring hypothesis can be weakened in a precise editorial extension: \textbf{for any commutative ring \(R\), any finite \(R\)-module \(M\), and any finitely generated ideal \(I\), one has \(\operatorname{Supp}(M)\subset V(I)\) if and only if \(I^NM=0\) for some \(N\geq0\)}.

First, the original finite module has \(\operatorname{Supp}(M)=V(\operatorname{Ann}_R M)\) over any ring. If \(M_{\mathfrak p}=0\), choose a finite generating list \(m_1,\ldots,m_t\) and, for each generator, an original \(s_j\notin\mathfrak p\) with \(s_jm_j=0\). The full product \(s=\prod_{j=1}^t s_j\) lies outside \(\mathfrak p\) and annihilates all of \(M\). Conversely, an annihilator outside \(\mathfrak p\) becomes a unit and kills the localization. For the empty generating list \(M=0\), the empty product is 1 and gives the same statement.

Let \(I=(a_1,\ldots,a_r)\). The support containment is therefore equivalent, by the original radical/intersection-of-primes identity, to \(a_j\in\sqrt{\operatorname{Ann}_R M}\) for every \(j\). Choose positive integers \(n_j\) such that \(a_j^{n_j}M=0\), and retain the explicit integer

\[
N=1+\sum_{j=1}^r(n_j-1).
\]

Every monomial generator \(a_1^{e_1}\cdots a_r^{e_r}\) of \(I^N\) has \(\sum e_j=N\), so at least one \(e_j\geq n_j\); otherwise the sum would be at most \(N-1\). It annihilates \(M\). Thus every term of every finite \(R\)-linear combination generating \(I^N\) annihilates \(M\), proving \(I^NM=0\). If the ideal has an empty generating list, it is zero and the formula gives \(N=1\), which works. Conversely, if \(I^NM=0\), at any prime outside \(V(I)\) an element of \(I\) becomes a unit and its \(N\)-th power kills the localization. If \(N=0\), then \(M=0\) directly. This proves both directions, retaining all original annihilators, generators, and factors.

The finite-generation condition on the ideal cannot be dropped from this general assertion. Use exactly the algebra already appearing in the source at lines 15347--15349:

\[
A=k[x_1,x_2,\ldots]/(x_i^2:i\geq1),\qquad
J=(x_1,x_2,\ldots),\qquad M=A.
\]

The squarefree monomials are a \(k\)-basis: the monomial ideal being divided out has as basis precisely the monomials divisible by some \(x_i^2\). Every element of \(J\) involves finitely many variables and has zero constant term, so a power greater than their number is zero. Therefore \(J\) is nil, and \(A/J=k\) shows that \(J\) is the unique prime. Hence \(\operatorname{Supp}_A(A)=V(J)\). Yet \(x_1\cdots x_N\ne0\) belongs to \(J^N\) for every \(N\geq1\). Thus no power of \(J\) kills the finite, cyclic module \(A\). The ideal is not finitely generated: a finite collection of its elements involves only finitely many variables, and the homomorphism setting those variables to zero sends that generated ideal to zero but leaves any fresh variable nonzero.

This same basis calculation verifies the source's associated-prime counterexample. For any nonzero \(a\in A\), choose a variable absent from its finite monomial support. Multiplication by that variable sends the distinct basis monomials of \(a\) to distinct nonzero squarefree monomials, so it does not annihilate \(a\). Thus \(\operatorname{Ann}_A(a)\ne J\); as \(J\) is the only prime, \(\operatorname{Ass}_A(A)=\varnothing\). Over the original coefficient field \(k\), every nonzero vector has annihilator zero, so \(\operatorname{Ass}_k(A)=\{(0)\}\). The original example and the new finiteness boundary use the very same algebra; neither is replaced by a different presentation.

\subsubsection{Annihilator spelling in the exact sequence}\label{algebra-proof-004-annihilator-spelling-in-the-exact-sequence}

Source: \texttt{lemma-ass}, lines 15130--15150; reports \texttt{OCC-12307} and \texttt{OCC-00393}. Replace ``annilator'' by ``annihilator'' once. The original dichotomy is valid for arbitrary modules and rings. If \(\operatorname{Ann}_R(m)=\mathfrak p\) and \(g\notin\mathfrak p\), then \(a(gm)=0\) is equivalent to \(ag\in\mathfrak p\), hence to \(a\in\mathfrak p\). If no such \(g\) sends \(m\) into the submodule, every annihilator of its quotient image lies in \(\mathfrak p\), while every element of \(\mathfrak p\) already kills \(m\). Thus the image has exactly that annihilator. The source maps and hypotheses are unchanged.

\subsubsection{One-equation module wording}\label{algebra-proof-004-one-equation-module-wording}

Source: lines 15296--15326; reports \texttt{OCC-12308} and \texttt{OCC-00394}. Both name the same ``Let \(R\) is'' error. FAC integration has already repaired it to ``Let \(R\) be'': the current lemma's history retains the earlier reading, and \texttt{fac/qa.json} records the correction under \texttt{algebra-lemma-one-equation-module}, kind \texttt{ungrammatical\_ring\_hypothesis}. The original finite-module and local-ring hypotheses and complete proof are unchanged. This is an already-composed correction, not a new edit. It lies outside the existing correction-only model, so its minimal original/replacement pair is retained separately in the review for the reusable fixes-only export at release. No second cumulative application is needed.

The support equality used in its proof is exact: outside \(V(f)\), the scalar \(f\) is invertible and the quotient localizes to zero. At \(\mathfrak p\) in the original support with \(f\in\mathfrak p\), the finite nonzero module \(M_{\mathfrak p}\) over the local ring \(R_{\mathfrak p}\) has nonzero quotient modulo \(f\), by Nakayama. Therefore \(\operatorname{Supp}(M/fM)=\operatorname{Supp}(M)\cap V(f)\). The original finite prime filtration expresses this closed support as a finite union of \(V(\mathfrak p_i)\). Applying the already proved original quotient dimension inequalities to each \(R/\mathfrak p_i\) gives the two inequalities, with the same prime-chain interpretation for equality when \(f\) avoids the minimal support primes. For \(M=0\), both supports are empty and both dimensions are \(-\infty\), with the extended-integer convention \(-\infty+1=-\infty\). No extra nonzero-module assumption is inserted.

\subsubsection{Prime type and the two localization maps}\label{algebra-proof-004-prime-type-and-the-two-localization-maps}

Source: \texttt{lemma-localize-ass}, lines 15422--15448; reports \texttt{OCC-12309} and \texttt{OCC-00395}. The object being quantified is \(\mathfrak p\in\operatorname{Spec}(R)\), not a ring element. Both reports are resolved by one operation.

For the original multiplicative subset \(S\subset R\) disjoint from \(\mathfrak p\), the exact comparison used by the proof is

\[
\alpha:M_{\mathfrak p}\longrightarrow(S^{-1}M)_{S^{-1}\mathfrak p},
\quad m/r\longmapsto(m/1)/(r/1),
\] \[
\beta:(S^{-1}M)_{S^{-1}\mathfrak p}\longrightarrow M_{\mathfrak p},
\quad (m/s)/(r/t)\longmapsto tm/(sr),
\]

where \(s,t\in S\) and \(r\notin\mathfrak p\). The latter condition follows from the original prime contraction correspondence, since \(r/t\notin S^{-1}\mathfrak p\). The maps are induced by the original module map to \(M_{\mathfrak p}\): each \(s\in S\) becomes a unit, and then each \(r/t\notin S^{-1}\mathfrak p\) becomes a unit. This proves their existence and well-definedness by the two localization universal properties. Conversely the map to the iterated localization inverts every \(r\notin\mathfrak p\), so it induces \(\alpha\). The composite \(\beta\alpha\) fixes every fraction \(m/r\). For the other composite the full expression \((tm/1)/(sr/1)\) equals \((m/s)/(r/t)\), by multiplying numerator and denominator using the original units \(s/1,t/1\); equivalently both maps agree on \(M\) and its two required inversions. Thus both composites are identities. The Noetherian condition still ensures finite generation of the prime used in the preceding associated-localization converse; the grammar correction does not remove it.

\subsubsection{Nonzerodivisors in arbitrary ideals}\label{algebra-proof-004-nonzerodivisors-in-arbitrary-ideals}

Source: \texttt{lemma-ideal-nonzerodivisor}, lines 15470--15493. The original lemma remains a translation of its stated local result. The same proof establishes a separate strengthening: \textbf{for every Noetherian ring \(R\), every finite \(R\)-module \(M\), and every ideal \(I\subset R\), the ideal contains a nonzerodivisor on \(M\) if and only if it is contained in none of the associated primes of \(M\)}. Locality and the condition \(I\subset\mathfrak m\) are unnecessary for this conclusion.

If \(x\in I\) acts injectively, it lies in no annihilator prime of a nonzero element, so \(I\) is contained in none of them. Conversely, the original Noetherian finite-module result gives a finite set \(\operatorname{Ass}_R(M)\). If \(I\) is contained in none, finite prime avoidance produces \(x\in I\) outside their union. The already proved zero-divisor characterization then makes multiplication by this exact \(x\) injective. These arguments use no maximal ideal or local ring. If \(M=0\), the associated set is empty and \(x=0\in I\) acts injectively on the zero module. If \(I=R\), the element 1 provides the required injection. Thus the previously excluded boundaries are retained and verified.

The proof also applies directly to the receiving lemma at lines 15521--15538. In fact \textbf{every Noetherian ring with infinitely many maximal ideals contains a nonzero nonunit nonzerodivisor}. Its associated set as a module over itself is finite, so choose a maximal ideal \(\mathfrak m\) absent from that set. It is contained in none of those proper prime ideals: containment would imply equality by maximality. The strengthened criterion gives an injective multiplier \(f\in\mathfrak m\). The ring is nonzero because it has maximal ideals, so injectivity makes \(f\ne0\); membership in \(\mathfrak m\) makes it a nonunit. The source's positive-dimensional finite-type algebra has infinitely many maximal ideals by its preceding equivalence, so this gives precisely its intended receiving application. The source's field, finite-type, and dimension hypotheses are kept in its translation; the broader result stays here as editorial material.

\subsubsection{Reading and propagation limits}\label{algebra-proof-004-reading-and-propagation-limits}

The source proofs concerning finite filtrations, associated primes, minimal support primes, scalar restriction, associated-prime localization, and detection by associated localizations were read through line 15538. The named corrections do not newly certify every omitted proof as a separate reconstruction. Support equalities in the original finite-module statements retain their hypotheses; no arbitrary-module dimension assertion is inferred from them. The earlier zero-ring correction was checked against the explicitly nonzero finite-type-algebra lemma and the positive-dimensional receiving lemma; both retain their original statements. The symbolic-power section through 15565 is lookahead, ending before the first lemma's proof, and is unadjudicated. Broader combined-results synthesis and the short mathematical paper remain after the core work.

\hypertarget{algebra_symbolic_weakass_notes_20260927}{}
\subsection{Editorial evidence for Algebra 15542--16444}\label{algebra-proof-005-editorial-evidence-for-algebra-1554216444}

Primary source: Stacks Project authors, original \texttt{algebra.tex} at commit \texttt{a04446e57ec1fbc252a871afcec7752fb2807b14}, SHA-256 \texttt{FA8BB92E58A4F78A2BD01B3B6A4A87DE0A0D279F5DD90641B574DD5FBFFFA4F3}. The complete interval and its thirty received reports were read. Source 16445--16480 is lookahead only, ending within the opening of the next lemma. The operations in the review contain minimal corrections; these arguments and strengthenings are separate editorial material. No novelty is claimed.

\subsubsection{Symbolic powers and the associated-prime strengthening}\label{algebra-proof-005-symbolic-powers-and-the-associated-prime-strengthening}

Source: definition-symbolic-power and lemma-symbolic-power-associated, 15548--15575. Retain the exact ideal

\[
J=\mathfrak p^{(n)}=\ker\bigl(R\longrightarrow R_{\mathfrak p}/\mathfrak p^nR_{\mathfrak p}\bigr).
\]

For every commutative ring, every prime \(\mathfrak p\), and every positive integer \(n\), one has

\[
\operatorname{Ass}_R(R/J)=\operatorname{WeakAss}_R(R/J)=\{\mathfrak p\}.
\]

The original Noetherian hypothesis is unnecessary for this particular statement. Here is the complete argument. The ideal \(\mathfrak p^nR_{\mathfrak p}\) lies in the proper maximal ideal, so its quotient is nonzero and \(J\) is proper. Directly from the defining map, \(\mathfrak p^n\subset J\subset\mathfrak p\): an element outside \(\mathfrak p\) becomes a unit and cannot lie in that proper localized ideal. These inclusions give \(\sqrt J=\mathfrak p\). If \(a\notin\mathfrak p\) and \(ab\in J\), the image of \(a\) is invertible in \(R_{\mathfrak p}\), so \(b/1\in\mathfrak p^nR_{\mathfrak p}\) and \(b\in J\). Thus every such \(a\) acts injectively on \(R/J\).

Choose the least integer \(r\in\{1,\ldots,n\}\) with \(\mathfrak p^r\subset J\). Because \(\mathfrak p^{r-1}\not\subset J\), there is a \(c\in\mathfrak p^{r-1}\setminus J\). Every element of \(\mathfrak p\) kills \(c+J\), while the preceding injection property says no element outside \(\mathfrak p\) kills it. Hence \(\operatorname{Ann}_R(c+J)=\mathfrak p\), proving existence of the associated prime without a Noetherian existence theorem. This includes \(r=1\), where one may choose \(c=1\).

For any nonzero \(b+J\), its annihilator contains \(J\) and is contained in \(\mathfrak p\), so its radical is exactly \(\mathfrak p\). Thus its unique minimal annihilator prime is \(\mathfrak p\); if its annihilator itself is prime, that prime is also \(\mathfrak p\). This proves both asserted equalities. At \(n=0\), the original definition gives \(J=R\), the quotient is zero, and both associated sets are empty; that boundary must not be included in the positive-power assertion. The original source statement remains unchanged, with this strengthening identified as editorial material.

\subsubsection{Flat extension requires injection, not invertibility}\label{algebra-proof-005-flat-extension-requires-injection-not-invertibility}

Source: lemma-symbolic-power-flat-extension, 15577--15612; reports OCC-00397, OCC-12311 and OCC-00398. Insert the article in ``a flat ring map'' and replace ``acts invertibly'' by ``acts injectively''. The theorem and its flatness hypothesis remain unchanged.

First verify the original base-ring maps used in its proof. Write \(\mathfrak q=\mathfrak pS\), assumed prime. The nonzero domain \(S/\mathfrak q\) is flat over \(R/\mathfrak p\). Multiplication by a nonzero element of \(R/\mathfrak p\) is injective there and remains injective after tensoring with \(S/\mathfrak q\). Such an element therefore cannot kill its nonzero identity element. Hence \(R/\mathfrak p\to S/\mathfrak q\) is injective and \(\mathfrak q\cap R=\mathfrak p\). In particular every original denominator in \(R\setminus\mathfrak p\) also avoids \(\mathfrak q\).

Set \(A=R_{\mathfrak p}\) and \(B=S_{\mathfrak p}\), retaining the source's localization at the image of \(R\setminus\mathfrak p\). Flatness gives

\[
\ker\bigl(S\to B/\mathfrak p^nB\bigr)=\mathfrak p^{(n)}S.
\]

The comparison \(B/\mathfrak p^nB\to S_{\mathfrak q}/\mathfrak p^nS_{\mathfrak q}\) localizes further at the elements of \(S\setminus\mathfrak q\). For \(n>0\), filter its source by \(\mathfrak p^iB/\mathfrak p^nB\), for \(0\leq i\leq n\). Exact tensoring with the flat \(A\)-algebra \(B\) identifies the successive quotients with

\[
\bigl(\mathfrak p^iA/\mathfrak p^{i+1}A\bigr)\otimes_A B
=\bigl(\mathfrak p^iA/\mathfrak p^{i+1}A\bigr)
 \otimes_{\kappa(\mathfrak p)}(S/\mathfrak q)_{\mathfrak p}.
\]

The first factor is the original residue-field vector space, possibly of infinite dimension. Choosing a basis expresses this tensor product as a direct sum of copies of the domain \((S/\mathfrak q)_{\mathfrak p}\). For \(f\in S\setminus\mathfrak q\), its image is nonzero in that domain, so multiplication by \(f\) is injective on every copy and hence on the direct sum. If \(fz=0\) in the filtered module, injection on its first quotient puts \(z\) in the next filtration term; repeating through the finite filtration gives \(z=0\). Thus every additional denominator is a nonzerodivisor, and the comparison map is injective. Its composite with the original map from \(S\) has kernel \(\mathfrak q^{(n)}\), so the two kernels coincide. At \(n=0\), both symbolic powers are unit ideals and the equality is immediate.

Surjectivity of multiplication is false even for \(R=k\), \(\mathfrak p=(0)\), \(S=k[t]\), \(\mathfrak q=(0)\), \(n=1\), \(f=t\). The filtration quotient is \(k[t]\); multiplication by \(t\) is injective but misses 1. The original further localization \(k[t]\to k(t)\) is nevertheless injective, exactly as required. No Noetherian hypothesis or finite-dimensionality of the vector-space quotients is added.

\subsubsection{Relative-assassin wording and exact local tensor maps}\label{algebra-proof-005-relative-assassin-wording-and-exact-local-tensor-maps}

Source: 15616--15918. The definition of \(A'\) requires the base \(R\) in its fibre ring, and the comparison lemma's opening list should include the already-defined sixth set \(A'_{fin}\). The remaining corrections restore ``proves'', the module \(N/\mathfrak p'N\), two occurrences of ``an \(R\)-module'', and finite verbs in the explanatory remark. These changes preserve every original set and hypothesis.

The identifications of associated sets in this interval are through the canonical maps on spectra, not identifications of arbitrary ideals in different rings. For an \(S/\mathfrak pS\)-module, annihilators over \(S\) are inverse images of annihilators over the quotient, with the inverse correspondence given by quotienting by \(\mathfrak pS\). For a localized module the same statement uses extension and contraction of its annihilator. This is precisely the quotient/localization comparison already proved in the preceding source lemmas and used in the definition of the fibre sets.

In lemma-bourbaki, the membership at line 15827 must use \(\operatorname{Ass}_{S_{\mathfrak q}}\), since the prime written there is \(\mathfrak qS_{\mathfrak q}\). The map at the preceding display is exact as printed, with an added terminal period. Its first isomorphism sends \((m\otimes n)/s\) to \(m\otimes(n/s)\), with inverse \(m\otimes(n/s)\mapsto(m\otimes n)/s\). Writing \(\varphi:R\to S\) and \(\mathfrak p=\varphi^{-1}(\mathfrak q)\), the next isomorphism and inverse are

\[
M\otimes_RN_{\mathfrak q}\longrightarrow
 M_{\mathfrak p}\otimes_{R_{\mathfrak p}}N_{\mathfrak q},
\quad m\otimes(n/s)\longmapsto(m/1)\otimes(n/s),
\] \[
(m/a)\otimes(n/s)\longmapsto m\otimes n/(\varphi(a)s),
\qquad a\notin\mathfrak p,\ s\notin\mathfrak q.
\]

Every displayed denominator is invertible in the indicated target; tensor balancing and the localization universal property make these maps well-defined. Their two composites fix each displayed pure tensor, which generates its module. Consequently the nonzerodivisor chosen in \(\mathfrak pR_{\mathfrak p}\), together with flatness of \(N_{\mathfrak q}\), acts injectively on the original localized tensor product. An associated element whose annihilator is \(\mathfrak qS_{\mathfrak q}\) would be a nonzero element killed by that same multiplier, giving the claimed contradiction over the correctly named local ring. The field and residue-field statements receiving this argument keep their original hypotheses.

\subsubsection{Weakly associated primes and the original counterexample}\label{algebra-proof-005-weakly-associated-primes-and-the-original-counterexample}

Source: 15925--16209. The received plural, sentence-fragment, apposition and infinitive corrections do not change the characterization of weak association. In the reduced local argument, an annihilator with radical the maximal ideal is proper, so its witness is nonzero; reducedness then contradicts the resulting nilpotence unless the maximal ideal is zero. In the arbitrary-module existence argument, a nonzero element has proper annihilator and the nonzero quotient has a minimal prime. The correction ``But as \ldots'' removes a doubled finite predicate while preserving that argument.

For the finite-generation lemma at 16114--16144, the positive exponents \(e_i\) apply to the original generators of \(\mathfrak p\). Replacing the nonzero localized witness by \(g_i^{e_i-1}m\) keeps it nonzero, sets that exponent to 1, and cannot increase the least annihilating exponents of the other generators. The finite sum decreases until every generator kills the witness. Its proper annihilator is then the maximal ideal, and the earlier finite-generator localization argument descends it. The source already states the weaker hypothesis that this particular prime is finitely generated; no global Noetherian condition is inserted.

The example at 16146--16166 retains exactly

\[
R=k[x_1,x_2,\ldots],\qquad
S=k[x_1,x_2,\ldots,y_1,y_2,\ldots]/(x_iy_i:i\geq1),\qquad M=S.
\]

Its prime \(\mathfrak q=\sum_i x_iS\) is minimal: the quotient is the domain \(k[y_1,y_2,\ldots]\); if a prime is contained in \(\mathfrak q\), none of the \(y_i\) belongs to it, so \(x_iy_i=0\) forces all \(x_i\) into it. Thus that prime equals \(\mathfrak q\).

The monomial relations express \(S\), as an \(R\)-module, as the direct sum over the finite-support exponent vectors \(\beta\) of \((R/I_\beta)y^\beta\), where \(I_\beta=(x_i:\beta_i>0)\). A nonzero element has finitely many nonzero components. Each component has annihilator \(I_\beta\), since \(R/I_\beta\) is a domain. Its total annihilator is therefore the intersection of these finitely many coordinate ideals. If a component has \(\beta=0\), the intersection is zero. Otherwise let \(F\) be the finite union of their supports. The intersection is generated by the monomials \(\prod_{i\in E}x_i\) as \(E\) ranges over the subsets of \(F\) meeting every component support: a monomial belongs to every coordinate ideal exactly when its support meets each of these finite sets, and the monomial basis gives the same criterion for polynomials. This proves finite generation. The annihilator is contained in the prime \((x_i:i\in F)\), strictly below \((x_1,x_2,\ldots)\). Hence the latter is not minimal over that annihilator. This proves the exact failure of weak-association functoriality asserted in the source; unequal radicals alone would not suffice. The source's omitted details remain editorial here, not a replacement passage.

\subsubsection{Finite maps and the repeated localization type correction}\label{algebra-proof-005-finite-maps-and-the-repeated-localization-type-correction}

Source: 16211--16314; reports OCC-00407, OCC-00408 and OCC-00409. Insert ``by'' at both occurrences of ``Denote'', and correct the target of \(ym\) from \(S_{\mathfrak q}\) to \(M_{\mathfrak q}\). The actual witness is nonzero there: its annihilator localizes to a proper ideal, and \(y\notin\mathfrak q\) is invertible. The witness is zero at the other maximal ideals because \(x\) is inverted there and \(yx^tm=0\). For each \(f\in\mathfrak pR_{\mathfrak p}\), a power kills its localization at \(\mathfrak q\), and it already vanishes at the other maximal ideals. Detection of zero by those localizations shows that power kills \(ym\) in \(M_{\mathfrak p}\). The witness remains nonzero, so its proper annihilator there has radical the maximal ideal. The earlier local characterization proves precisely the asserted contraction of weak association.

The source at 16291 repeats the earlier malformed quantifier from 15446. Correct it to \(\mathfrak p\in\operatorname{Spec}(R)\) with the given disjointness condition. This resolves the pending propagation entry in ALGEBRA-RECON-407. The exact two-stage module maps are the same maps already recorded in the preceding editorial note: \(m/r\mapsto(m/1)/(r/1)\) and \((m/s)/(r/t)\mapsto tm/(sr)\), with every original denominator retained. The proof uses the weakly-associated local characterization, so it needs no Noetherian hypothesis. Its original statement remains unchanged.

\subsubsection{Field-extension proof and the nonzero witness}\label{algebra-proof-005-field-extension-proof-and-the-nonzero-witness}

Source: lemma-weakly-ass-change-fields, 16360--16439. The polynomial ring in the purely transcendental stage uses \(x_1,\ldots,x_r\); the earlier finite-extension degree \(n\) does not name its variables. The annihilator membership is \(gf^n\in J\), equivalent to \(gf^nz=0\); the latter module element is not itself an element of the ideal \(J\). These are the two received mathematical type/index repairs.

The nonzerodivisor argument at 16410 should explicitly choose a nonzero \(z\). Injection is proved by showing a nonzero input has nonzero image; the source's subsequent contradiction would be false for \(z=0\). Write this original \(z\) as a finite sum \(\sum_\alpha m_\alpha\otimes t_\alpha\) in a fixed \(k\)-basis of \(K\), and let \(M'\) be generated by its actual coefficients. Coordinate extraction shows that every submodule containing \(z\) after tensoring contains all these coefficients. If \(z\) belonged to \(\mathfrak pM'\otimes_kK\), each coefficient would lie in \(\mathfrak pM'\), giving \(M'=\mathfrak pM'\). Nakayama would give \(M'=0\), contrary to \(z\ne0\). Thus its residue is nonzero.

In this stage \(K=k(x_1,\ldots,x_r)\). The ring \(\kappa(\mathfrak p)\otimes_kK\) is the localization of \(\kappa(\mathfrak p)[x_1,\ldots,x_r]\) at the nonzero polynomials from \(k[x_1,\ldots,x_r]\). Those denominators remain nonzero, so the localization is a domain. The original residue module is a direct sum of copies of this domain, and the chosen \(f\) has nonzero residue. Multiplication cannot kill the nonzero residue of \(z\). The comma after the Nakayama reference and the period after the residue-module sentence complete the received wording repairs.

For the final annihilator argument, \(\mathfrak p(R\otimes_kK)\subset\mathfrak q\), so \(g\notin\mathfrak q\) implies \(g\notin\mathfrak p(R\otimes_kK)\). The proved injection permits cancellation of this actual multiplier from \(gf^nz=0\). Every \(f\in\mathfrak p\) thus has a power annihilating \(z\), while any annihilator in \(R\) lies in \(\mathfrak q\cap R=\mathfrak p\). Its radical is exactly \(\mathfrak p\). The finitely many nonzero coefficients of \(z\) place it in a finite direct sum of copies of \(M\). Applying the source's weak-association short-exact-sequence inclusions successively yields weak association to one copy of \(M\), with the same original prime. This completes the receiving argument without altering its field-extension statement.

\subsubsection{Torsion-free generic-fibre comparison}\label{algebra-proof-005-torsion-free-generic-fibre-comparison}

Sources: lemma-post-bourbaki, 15853--15877, and lemma-weak-post-bourbaki, 16338--16358. A separate strengthening replaces their flatness hypothesis by the exact property used in both proofs: every nonzero element of the domain \(R\) acts injectively on the \(S\)-module \(N\). Retain its fraction field \(K\) and the original coefficient map \(R\to S\).

Under this hypothesis, localization gives an injection \(N\to N\otimes_RK\). For any numerator \(u\in N\), denominator \(s\in R\setminus\{0\}\) and \(a\in S\), the relation \(a(u/s)=0\) means there is an original nonzero \(t\in R\) with \(tau=0\); injectivity of multiplication by \(t\) makes this equivalent to \(au=0\). Thus the annihilator in \(S\) of every fraction equals the annihilator of its numerator. Conversely every numerator occurs with denominator 1. Their associated and weakly associated sets over \(S\) are consequently identical. Passing to the localized coefficient ring by the source localization correspondences gives

\[
\operatorname{Ass}_S(N)=\operatorname{Ass}_{S\otimes_RK}(N\otimes_RK),\qquad
\operatorname{WeakAss}_S(N)=\operatorname{WeakAss}_{S\otimes_RK}(N\otimes_RK),
\]

where both equalities use the canonical injection of localized spectra. The maps retain the original numerators and denominators. The zero module is included, with empty associated sets. Flatness implies the stated injectivity by tensoring \(0\to R\xrightarrow{t}R\), but is stronger than necessary.

For a strict example let \(R=S=k[x,y]\), \(N=(x,y)\), and \(\mathfrak m=(x,y)\). This module is torsion-free because it is a submodule of the domain. It is not flat: tensor the injection \(N\hookrightarrow R\) with \(N\). The element \(x\otimes y-y\otimes x\) maps to zero in \(R\otimes_RN\), but maps to the nonzero vector \(\bar x\otimes\bar y-\bar y\otimes\bar x\) in \((N/\mathfrak mN)\otimes_k(N/\mathfrak mN)\). The degree-one monomials \(\bar x,\bar y\) form a basis, and the two displayed ordered tensor basis elements are distinct in every characteristic. Thus tensoring fails to preserve that injection. Both generic-fibre conclusions still apply to this original torsion-free module. The source's flat statements stay unchanged.

\subsubsection{Bounded literature comparison and limits}\label{algebra-proof-005-bounded-literature-comparison-and-limits}

The machine-readable topic route and corpus-reading record preserve the four actual bounded queries and source identities. M. Richard Sayanagi, \emph{Power Series Rings over Zero-Dimensional Rings}, source \texttt{20251008\_Power\_Series\_Rings\_Sayanagi.tex}, lines 589--630, defines weak-Bourbaki primes of an ideal by minimal primes over \((I:a)\). The exact comparison with the source module convention is \(\operatorname{Ann}_R(a+I)=(I:a)\), since \(r(a+I)=0\) exactly when \(ra\in I\); minimal primes therefore agree for the same witness. This dictionary applies directly to the symbolic quotient above. The later power-series lemma's proof was only opened, not reviewed in full.

Roman Bezrukavnikov, Michael Finkelberg and Ivan Mirkovic, \emph{Equivariant (K-)homology of affine Grassmannian and Toda lattice}, \texttt{equiv\_K.tex}, lines 1695--1742, uses symbolic powers in the construction of an affine blow-up and discusses intersections over components of a reduced subscheme. This contextual passage does not replace the original prime-ideal definition under review. Its Proposition \texttt{blow} was only partially read, and no assertion of that proposition is adopted here. The exact TeX hashes and section checks are retained in the route. These bounded readings are neither a novelty survey nor a claim to have read all corpus hits. Broader combined-results synthesis remains after completion of the core integration work.

\hypertarget{algebra_embedded_regular_notes_20260927}{}
\subsection{Editorial evidence for Algebra 16445--17135}\label{algebra-proof-006-editorial-evidence-for-algebra-1644517135}

Authority: Stacks Project authors, \texttt{algebra.tex}, commit \texttt{a04446e57ec1fbc252a871afcec7752fb2807b14}, SHA-256 \texttt{FA8BB92E58A4F78A2BD01B3B6A4A87DE0A0D279F5DD90641B574DD5FBFFFA4F3}. The complete interval and twelve due reports were read. Source 17136--17170 is lookahead through the blowup definition and first lemma opening, not adjudicated. Corrections remain minimal proposed operations; the arguments and stronger statements below remain separate editorial material.

\subsubsection{The embedded-prime quotient and its exact kernel}\label{algebra-proof-006-the-embedded-prime-quotient-and-its-exact-kernel}

Source: lemma-remove-embedded-primes, 16467--16533; reports OCC-12152, OCC-00412, OCC-00413 and OCC-12069. The localization at 16508 must name the newly quantified submodule \(K'\). The empty-set correction at 16531 is already composed as MC-STK-ERR-0412. Its current text has \(D(f)\cap\operatorname{supp}(K)=\varnothing\). Lowercase \texttt{supp} identifies the same support; changing its capitalization is optional typography, not an additional mathematical correction. Retain the prior correction without allocating it again.

Let \(Q=\{\mathfrak q_1,\ldots,\mathfrak q_t\}\) be the original minimal support primes and \(E=\{\mathfrak p_1,\ldots,\mathfrak p_s\}\) the original embedded associated primes. The submodule constructed in the source is exactly

\[
K=\ker\left(M\longrightarrow\prod_{\mathfrak q\in Q}M_{\mathfrak q}\right).
\]

Indeed, a cyclic submodule supported in \(\bigcup_{\mathfrak p\in E}V(\mathfrak p)\) has zero localization at every \(\mathfrak q\in Q\): no embedded prime can be contained in a minimal support prime. Conversely, if \(m\) vanishes at all those localizations, every associated prime of \(Rm\) belongs to \(\operatorname{Ass}(M)\), and none is in \(Q\). Every minimal support prime of the finite module \(Rm\) is associated. Its closed support is therefore contained in the union of the closed sets of the remaining associated primes, all in \(E\). This puts \(m\) in the original constructed submodule.

Every submodule of the finite module \(M\) is finite over the Noetherian ring. Its support is closed; such a support is nowhere dense in \(\operatorname{Supp}(M)\) precisely when it omits all the generic points \(Q\). One direction follows because containing a generic point would contain that entire irreducible component. For the converse, a closed set containing a nonempty open set must contain the generic point of a component meeting that open set, since opens are stable under generization. Thus the displayed kernel is the \textbf{unique greatest} submodule with nowhere-dense support. The source's maximality assertion is consequently stronger than mere existence of a maximal choice, and \(M/K\) is determined canonically.

The quotient map is an isomorphism at every prime in \(Q\). Since both modules have closed support and every component of \(\operatorname{Supp}(M)\) is retained, their supports agree. If \(\bar m\in M/K\) has prime annihilator, its lift does not belong to \(K\), so its image is nonzero at some \(\mathfrak q\in Q\). The annihilator is then contained in \(\mathfrak q\). It belongs to the common support, so minimality of \(\mathfrak q\) forces equality. Thus the quotient has no embedded associated primes. If \(f\) belongs to every prime in \(E\), then the closed union containing \(\operatorname{Supp}(K)\) misses \(D(f)\), giving the already-recorded empty intersection and \(K_f=0\).

\subsubsection{Compatibility with arbitrary localization}\label{algebra-proof-006-compatibility-with-arbitrary-localization}

Source: lemma-remove-embedded-primes-localize, 16535--16546. Its principal-localization conclusion extends to any multiplicative subset \(T\subset R\), with no change to its original Noetherian and finite-module hypotheses. Let \(K_T^{new}\) denote the greatest submodule with nowhere-dense support constructed inside \(T^{-1}M\). Then

\[
K_T^{new}=T^{-1}K,\qquad
T^{-1}(M/K)\xrightarrow{\;\sim\;}(T^{-1}M)/(T^{-1}K),
\quad (m+K)/t\longmapsto(m/t)+T^{-1}K.
\]

The inverse sends the displayed class of \(m/t\) to \((m+K)/t\); exact localization makes both maps well-defined and their composites the identity.

The minimal primes of the localized support correspond exactly to the \(\mathfrak q\in Q\) disjoint from \(T\). In fact, any support prime disjoint from \(T\) contains a member of \(Q\), which is then disjoint from \(T\) too; minimality on either side gives the assertion. For such a surviving prime the original comparison of fractions identifies the iterated localization with \(M_{\mathfrak q}\). If \(\mathfrak q\cap T\ne\varnothing\), choose an actual \(t\in\mathfrak q\cap T\). The finite module \(M_{\mathfrak q}\) has support consisting of the maximal ideal alone. The support/power criterion, in the finite-principal-ideal form proved in ALGEBRA-RECON-404, gives \(t^N M_{\mathfrak q}=0\) for some \(N\). Inverting \(t\) therefore kills that factor. Localization is exact and commutes with the finite product indexed by \(Q\), so localizing the kernel formula above leaves exactly the kernel formula for the localized module. This proves the asserted equality.

If \(M=0\), the set \(Q\) is empty, its product in the module category is zero, and \(K=0\). If no support prime survives, \(T^{-1}M=0\) and the same formula applies. Multiplicative sets containing zero are included. The stronger result is editorial; the source's principal-localization lemma and its omitted proof remain identifiable.

\subsubsection{Regular sequences retain their module and order}\label{algebra-proof-006-regular-sequences-retain-their-module-and-order}

Source: 16593--16844; reports OCC-00414, OCC-00415 and OCC-00416. The source's definition requires a nonzero final quotient and preserves the order of the elements. These conditions are retained.

The permutation lemma must conclude \textbf{\(M\)-regular}, not unqualified ring regularity. For a concrete distinction, take

\[
R=\bigl(k[t,\epsilon]/(t\epsilon,\epsilon^2)\bigr)_{(t,\epsilon)},
\qquad M=R/(\epsilon).
\]

Multiplication by \(t\) is injective on \(M=k[t]_{(t)}\), and \(M/tM=k\ne0\). But \(t\epsilon=0\) in \(R\), and \(\epsilon\ne0\): the map to \(k[\epsilon]/(\epsilon^2)\) setting \(t=0\) extends through this localization, since its denominators have nonzero constant term, and sends \(\epsilon\) to a nonzero element. Thus even a one-term \(M\)-regular sequence need not be ring regular. The source's adjacent-transposition proof works on \(M/(x_1,\ldots,x_{i-2})M\), so the omitted final \(M\) in that quotient is restored. All the \(x_i\) are nonunits in the local ring, since the final quotient is nonzero; Nakayama in its finite kernel is therefore applicable as printed.

In the polynomial criterion the next multiplier is \(f_{i+1}x_{i+1}\), not \(f_{i+1}x_i\). Retain the source decomposition into monomial components

\[
R[x_1,\ldots,x_r]/(f_1x_1,\ldots,f_ix_i)
=\bigoplus_{E\in\mathbf Z_{\geq0}^r}(R/I_E)x^E,
\qquad I_E=(f_j:1\leq j\leq i,\ e_j>0).
\]

Multiplication by \(x_{i+1}\) sends the component \(E\) to \(E+\mathbf e_{i+1}\), and these two components have exactly the same ideal \(I_E\), because only indices at most \(i\) enter that ideal. The map on coefficients is multiplication by \(f_{i+1}\), and different components remain different. Thus injection holds precisely when it holds on every \(R/I_E\). Taking supports of \(E\) to be arbitrary subsets of \(\{1,\ldots,i\}\), rather than just the all-positive case mentioned in the source, verifies every ordered subsequence condition. The final quotients are nonzero because the ideal generated by the original \(f_j\) is proper; the polynomial quotient also has its unchanged degree-zero component \(R\). The case of no variables has the same conclusion under that proper-ideal hypothesis. No permutation or index is absorbed into different notation.

\subsubsection{Faithfully flat preservation and reflection}\label{algebra-proof-006-faithfully-flat-preservation-and-reflection}

Sources: lemma-flat-increases-depth, 16675--16693, and lemma-flat-base-change-quasi-regular, 16980--17002. Both comparisons have a separate common strengthening: for \textbf{any faithfully flat} map \(R\to S\), with no locality assumption, an ordered sequence is \(M\)-regular, respectively \(M\)-quasi-regular, if and only if its image is so on \(M\otimes_RS\).

For regularity, write \(Q_i=M/(f_1,\ldots,f_i)M\), including \(Q_0=M\), and \(L_i=\ker(f_i:Q_{i-1}\to Q_{i-1})\). The quotient comparison sends \((m+(f_1,\ldots,f_i)M)\otimes s\) to \(m\otimes s\) modulo the extended generated submodule. Its inverse is induced by the tensor of the original quotient map. Right exactness gives the inverse identifications. Flatness then identifies \(L_i\otimes_RS\) with the kernel of the same multiplier on the extended quotient. Faithfulness says these kernels vanish exactly when \(L_i=0\), and says \(Q_r\otimes_RS\ne0\) exactly when \(Q_r\ne0\). This proves every part of the original definition, including the empty sequence and zero-module boundaries, without changing the order of any element.

For quasi-regularity let \(J=(f_1,\ldots,f_c)\) and let \(\mu\) be the original canonical graded map. Exact tensoring sends \(J^nM\subset M\) to \((JS)^n(M\otimes_RS)\subset M\otimes_RS\), and sends each graded quotient to its corresponding extended quotient. On the polynomial side the comparison sends each \(((m+JM)\otimes X^E)\otimes s\) to \((m\otimes s+(JS)(M\otimes_RS))\otimes X^E\). The reverse tensor/quotient comparison defines the inverse; balancing over the original coefficient quotient proves well-definedness. These maps preserve every monomial and identify the target canonical map with \(\mu\otimes_RS\). Since \(\mu\) is always surjective, flatness gives \(\ker(\mu\otimes_RS)=\ker(\mu)\otimes_RS\), and faithfulness detects whether that kernel vanishes. No nonzero-quotient condition is inserted into the source's different definition of quasi-regularity. In particular the zero module remains included there.

Local flat maps in the source are faithfully flat, so its regularity statement is a receiving special case. The flat quasi-regularity statement retains its original forward implication. The stronger reflection assertion and removal of locality stay in this editorial note.

\subsubsection{The two induction boundaries and module coefficients}\label{algebra-proof-006-the-two-induction-boundaries-and-module-coefficients}

Source: lemma-regular-quasi-regular, 16895--16978; reports OCC-00418 and OCC-00417. The outer induction starts at the empty sequence, where the canonical map is the identity in degree zero and has zero source and target in positive degree. For \(c>0\), keep the original \(J'=(f_1,\ldots,f_{c-1})\), total degree \(n\), and last exponent \(l\). At \(l=0\), the relation involves only the first \(c-1\) elements, so the outer induction gives every coefficient in \(J'\). There is no term with exponent \(-1\). The rewrite involving \(l-1\) must begin only at \(l\geq1\); when \(l=1\), the sum from 0 through \(l-2\) is empty.

For the step, all terms with last exponent less than \(l\) have primed degree at least \(n-l+1\). Thus the leading primed sum, multiplied by \(f_c^l\), is in \((J')^{n-l+1}\). The outer induction gives \(f_c^l a_{I',l}\in J'\), and regularity on \(R/J'\) gives \(a_{I',l}\in J'\). Write each such coefficient as \(\sum_{j<c}f_jd_{I',j}\). Regrouping the actual terms of primed degree \(n-l+1\) gives coefficients \(b_{B,l-1}=\sum_{j<c,\ B_j>0}d_{B-\mathbf e_j,j}\). The leading term becomes precisely \((\sum_B f_cb_{B,l-1}f^B)f_c^{l-1}\). The inner induction applies, and subtracting \(f_cb_{B,l-1}\in J\) recovers the original coefficients in \(J\). If there are no primed variables, the corresponding sums are empty; the same argument uses \(J'=0\). All coefficients and summands are retained.

For the original module assertion, make this same argument with coefficients in \(M\), the graded map for the first \(c-1\) elements, and injective action by \(f_c\) on \(M/J'M\). A relation belonging to \(J^{n+1}M\) can first be expressed with degree-\(n\) coefficients in \(JM\) and subtracted, exactly as for the ring. The conclusion is \(m_I\in JM\), not membership in the ring ideal \(J\). The proposed changes state the missing inner base case and repair this type; they do not replace the full proof body.

\subsubsection{Truncation in degree zero and the original polynomial module}\label{algebra-proof-006-truncation-in-degree-zero-and-the-original-polynomial-module}

Source: lemma-truncate-quasi-regular, 17026--17057; reports OCC-00419 and OCC-00420. In degree zero the source quotient is \(M/JM\), also the degree-zero term of the desired polynomial quotient. Separate that case. For every \(n\geq1\) the required intersection is

\[
J^nM\cap f_1M=f_1J^{n-1}M.
\]

The right side is visibly contained in the left. For the reverse, suppose \(f_1m\in J^nM\). At \(n=1\), the claim is immediate because \(m\in M=J^0M\). At \(n\geq2\), if \(m\notin J^{n-1}M\), there is a greatest \(k\in\{0,\ldots,n-2\}\) for which \(m\in J^kM\). Its class in \(J^kM/J^{k+1}M\) is nonzero. The exact original graded identification is with the polynomial \textbf{module}

\[
(M/JM)\otimes_{R/J}(R/J)[X_1,\ldots,X_c].
\]

Multiplication by \(X_1\) is injective on this module: it shifts distinct monomial components to distinct components, leaving their coefficients unchanged. Thus \(f_1m\notin J^{k+2}M\), contrary to \(J^nM\subset J^{k+2}M\). This proves the intersection without assuming separatedness, finiteness, or Noetherianity. Taking the associated graded quotient by \(X_1\) gives exactly the tail-variable polynomial module for \(\overline M\), as required. There is no \(J^{-1}\). A coefficient module must not be called a polynomial algebra, and the source's omitted denominator \(M\) must be restored.

\subsubsection{The separated quotient and its explicit ideal}\label{algebra-proof-006-the-separated-quotient-and-its-explicit-ideal}

Source: lemma-quasi-regular-on-quotient, 17116--17130; report OCC-00421. Define the missing ideal \(\overline J=J\overline R=(\overline f_1,\ldots,\overline f_r)\). Write \(C=\bigcap_{n\geq0}J^n\) and \(D=\bigcap_{n\geq0}J^nM\), with the original \(\overline R=R/C\) and \(\overline M=M/D\). Since \(CM\subset D\), the module structure is defined. For every \(n\), the quotient map identifies \(\overline J^n\overline M\) with \(J^nM/D\), and \(D\subset J^{n+1}M\). The displayed comparison sends

\[
m+J^{n+1}M\longmapsto (m+D)+\overline J^{n+1}\overline M
\]

for \(m\in J^nM\); its inverse uses any representative in \(J^nM\). Changing representatives by \(D\) changes nothing in the source quotient. This proves both composites are identities in every degree. The degree-zero coefficient comparison and the quotient \(R/J\cong\overline R/\overline J\) identify the two polynomial modules, and the canonical maps commute on every original monomial. Consequently the original quasi-regularity equivalence is retained, including the unit-ideal and zero-module cases.

\subsubsection{Reading limits and propagation}\label{algebra-proof-006-reading-limits-and-propagation}

The original source was read through 17135, without turning every omitted source proof into a rewrite. The source-defined order and nonzero-final-quotient condition for regularity are not imposed on its distinct quasi-regularity definition. The earlier support/power result in ALGEBRA-RECON-404 is used at the precise localization step above. The original local regular/quasi-regular equivalence at 17059--17081 retains its Noetherian, finite, nonzero and maximal-ideal hypotheses; no unproved removal of separatedness is inferred from the graded calculation.

The bounded corpus route retains the previous actual queries and adds one exact-phrase query for quasi-regularity. Its top hits were not used as mathematical evidence. Additional reading of M. Richard Sayanagi's original TeX, lines 692--729, covered the stated grade/height result and its proof, plus commented material; its dependency Theorem \texttt{unmixed} was not reviewed here, so that grade/height result is not adopted. The source Stacks proofs above supply the current dependencies. No novelty or exhaustive literature claim is made, and broader synthesis remains after the core work.

\hypertarget{algebra_blowup_ext_depth_notes_20260927}{}
\subsection{Editorial evidence for Algebra 17136--18021}\label{algebra-proof-007-editorial-evidence-for-algebra-1713618021}

Authority: Stacks Project authors, algebra.tex, commit a04446e57ec1fbc252a871afcec7752fb2807b14, SHA-256 FA8BB92E58A4F78A2BD01B3B6A4A87DE0A0D279F5DD90641B574DD5FBFFFA4F3. This interval and its fifteen received reports have been read. Lines 18022--18030 are lookahead through the next statement and proof opening only. These arguments support the correction review; stronger results remain identifiable editorial material. No translated or cumulative chapter has been changed.

\subsubsection{Blowup fractions and base change}\label{algebra-proof-007-blowup-fractions-and-base-change}

Sources: definition-blow-up and lemma-affine-blowup, 17142--17181; lemma-blowup-base-change, 17183--17207; OCC-00424.

Retain the original graded algebra \(\bigoplus_{n\geq0}I^n\), including \(I^0=R\), and its degree-zero localization at \(a^{(1)}\). Its fractions \(x/a^n\), \(x\in I^n\), map to the same fractions in \(R_a\). The original equality relation \[
a^k(a^mx-a^ny)=0
\] is exactly the equality relation in \(R_a\), so this map is injective. This includes zero divisors in \(R\). The image is generated by all \(x/a\), \(x\in I\); products generate each \(I^n\). Multiplication by \(a\) on this subring of \(R_a\) is injective, \(IR'=aR'\) since \(x=a(x/a)\), and inverting \(a\) gives \(R_a\). For \(a=0\) the localization is the zero ring, where injectivity of its unique endomorphism is still valid.

For the base change retain \(b\), the image of \(a\), \(J=IS\), \(B=S\otimes_R R[I/a]\), and \(T_b(B)=\{u:\exists k\geq0,\ b^ku=0\}\). This is an ideal: a common larger exponent annihilates a sum, and multiplication by any element preserves annihilation. The map \[
\theta:B\longrightarrow S[J/b],\qquad
s\otimes(x/a^n)\longmapsto s\,\overline{x}/b^n
\] is well-defined by the original fraction equality, balancing, and the multiplication formula. It is surjective: every numerator in \(J^n=I^nS\) is a finite sum \(\sum_i\overline{x_i}s_i\), \(x_i\in I^n\). It kills \(T_b(B)\), because \(b\) is injective on the target.

There is also an exact direct proof of the claimed annihilation in \(S\otimes_R I^n\). If \(w=\sum_i s_i\otimes x_i\) maps to zero in \(J^n\), then \[
b^nw=\sum_i s_i\otimes a^nx_i
=\sum_i s_i\overline{x_i}\otimes a^n=0.
\] Here \(a^n\) lies in the original \(I^n\), so each balancing step has its specified domain. The source claim is valid in that tensor module, not merely after mapping into \(B\).

The inverse from \(S[J/b]\) to \(B/T_b(B)\) sends \(y/b^n\), with \(y=\sum_i\overline{x_i}s_i\), to \(\sum_i s_i\otimes x_i/a^n\) modulo torsion. A different decomposition at the same degree has difference killed by \(b^n\). For different representatives \(y/b^n=y'/b^m\), choose \(k\) with \(b^k(b^my-b^ny')=0\). If \(u,u'\) are the proposed images in \(B\), then \(b^nu=y\otimes1\) and \(b^mu'=y'\otimes1\), so \(b^{k+n+m}(u-u')=0\). The class is independent of both choices. Addition, multiplication and the identity follow by common denominators and these same balancing equations. On generators the two composites are the identity. Thus the original quotient statement holds for arbitrary ring maps, with every torsion term retained.

The report is a notation clarification: use \(b\)-power torsion and \(b^n\) on the \(S\)-side consistently. Identifying \(a\) with its image would also be legitimate if declared; it is not a false mathematical torsion assertion.

\subsubsection{Polynomial charts and the named base ring}\label{algebra-proof-007-polynomial-charts-and-the-named-base-ring}

Sources: examples at 17209--17247 and lemma-affine-blowup-quotient-description, 17249--17268; OCC-00422 and OCC-00423.

The source monomial is \(t_1^{e_1}\cdots t_n^{e_n}\), so the unrelated \(x_n\) is corrected to \(t_n\). The base-change map in the quotient-description proof is the declared \(\mathbf Z[t_1,\ldots,t_n]\to R\), not a new ring \(A\).

For the polynomial affine chart the displayed quotient identifies with \(R[t_1,x_2,\ldots,x_n]\) by the maps \(t_i\mapsto t_1x_i\) for \(i\geq2\), and the reverse map sending each retained variable to its class. The displayed relations verify both composites. Multiplication by \(t_1\) is injective even when \(R\) has zero divisors, by the monomial coefficients. After inverting \(t_1\), the maps \(x_i\leftrightarrow t_i/t_1\) are inverse. The source chart injects into its localization, as does the blowup fraction ring, so the surjection between them is injective. Base change along the stated map to \(R\), with the exact torsion quotient just proved, gives the claimed general presentation. Its \(n=1\) case retains the quotient of \(R\) by all \(a_1\)-power torsion; it need not equal \(R\).

\subsubsection{The field boundary and the directed chart construction}\label{algebra-proof-007-the-field-boundary-and-the-directed-chart-construction}

Source: lemma-valuation-ring-colimit-affine-blowups, 17351--17391; OCC-00425 and OCC-00426.

If \(R\) is a field, then \(\mathfrak m=0\) and \(R=A=K\). The original condition \(a\in I\subset\mathfrak m\) forces \(a=0\) and \(I=0\). That chart is zero and has zero fibre, so none of the prescribed objects exists. Preserve the theorem's field case by admitting the identity chart \((I,a)=(R,1)\), whose algebra is \(R\) and whose closed fibre is \(R/\mathfrak m\ne0\). Also state explicitly that the indexed charts are contained in the given \(A\); the properties of a chart alone do not select this particular valuation ring. These are corrections of the stated diagram, not a replacement of \(A\).

Suppose now \(R\) is not a field. For a finite set \(E=\{e_1,\ldots,e_n\}\subset A\), choose an actual common denominator \(0\ne d\in R\) with \(de_i\in R\). Choose \(0\ne t\in\mathfrak m\), and retain \(a=td\) and \(f_i=ae_i=t(de_i)\). Thus \(a,f_i\in\mathfrak m\) directly in \(R\), and \(a\ne0\). For \(I=(f_1,\ldots,f_n,a)\) the chart is \[
R[I/a]=R[e_1,\ldots,e_n]\subset A.
\] Indeed each fraction with numerator in \(I^k\) expands as a polynomial in the \(f_i/a=e_i\) and \(a/a=1\), and those generators are already chart fractions. If \(\mathfrak n\) is the maximal ideal of the dominating valuation ring, then \(\mathfrak n\cap R=\mathfrak m\); its contraction to the chart is a prime over \(\mathfrak m\). The chart's quotient by \(\mathfrak m\) is consequently nonzero, proving the required fibre condition.

For two such charts with pairs \((I,a),(J,b)\), their containing chart is \(R[IJ/(ab)]\). The first inclusion sends \(x/a^n\) to \(b^nx/(ab)^n\), retaining \(b^nx\in I^nJ^n\); the second inclusion is analogous. A numerator \(z\in(IJ)^n=I^nJ^n\) is a finite sum \(\sum_i x_i y_i\) and \[
z/(ab)^n=\sum_i(x_i/a^n)(y_i/b^n)\in A.
\] The product ideal remains finitely generated and contained in \(\mathfrak m\), the denominator remains nonzero, and contraction of \(\mathfrak n\) again gives a nonzero fibre. The identity chart combines with either chart to give that chart. Thus this is a nonempty directed system of subrings of the original \(K\), every finite subset of \(A\) lies in an object, and its colimit under the displayed inclusions is exactly \(A\). The empty finite subset is covered by the identity chart.

\subsubsection{The exact overring strengthening}\label{algebra-proof-007-the-exact-overring-strengthening}

The preceding argument does not require that the receiving overring be a valuation ring. For any local domain \((R,\mathfrak m)\) with fraction field \(K\), and any intermediate ring \(R\subset A\subset K\), the following are equivalent:

\begin{enumerate}
\def\labelenumi{\arabic{enumi}.}
\tightlist
\item
  \(A/\mathfrak mA\ne0\).
\item
  \(A\) is the directed union of its affine charts \(R[I/a]\subset A\) with \(I\) finitely generated, nonzero \(a\in I\subset\mathfrak m\), and nonzero closed fibre, together with the identity chart.
\end{enumerate}

For (1) choose a maximal ideal of \(A\) containing the proper ideal \(\mathfrak mA\). Its contraction to \(R\) is exactly \(\mathfrak m\). The finite-set common-denominator construction, product chart and fraction maps in the previous section apply without change; contraction of this ideal supplies the fibre condition for every chart inside \(A\). If \(R\) is a field, the intermediate ring is \(R\) and the identity chart suffices.

Conversely, if \(1=\sum_{j=1}^s r_j a_j\) with \(r_j\in\mathfrak m\), \(a_j\in A\), directedness puts the finite list \(a_j\) in one chart \(B\). Then \(1\in\mathfrak mB\), contrary to its nonzero closed fibre. Thus (2) implies (1). A finite-type \(A\) with nonzero closed fibre is itself one of the constructed charts, using a finite list of its algebra generators. This is a proved editorial strengthening; the translation retains the author's valuation-ring hypothesis and the visible field-case correction.

\subsubsection{Resolutions, chain maps and contravariance}\label{algebra-proof-007-resolutions-chain-maps-and-contravariance}

Sources: 17410--17754; OCC-00430, OCC-00427, OCC-01388, OCC-00428, OCC-00429, OCC-01389, OCC-01391 and OCC-01390.

In the finite-free construction the first kernel is \(\ker(R^{n_0}\to M)\). It is finite because \(R\) is Noetherian, so choose a finite free module surjecting onto it before using the printed recursion at \(e\geq1\). There is no \(R^{n_{-1}}\).

The chain identity is \[
\alpha_{i-1}d_{F,i}=d_{G,i}\alpha_i:F_i\longrightarrow G_{i-1}.
\] The original \(d_{G,i-1}\alpha_i\) is ill-typed. This correction is already present as MC-STK-ERR-1602.

For a chain map \(\alpha:F_\bullet\to G_\bullet\), the induced cochain map is \[
\alpha^*:\operatorname{Hom}_R(G_\bullet,N)
\longrightarrow\operatorname{Hom}_R(F_\bullet,N),\quad
u\longmapsto u\alpha_i\text{ in degree }i.
\] The original cochain differential is precomposition with \(d_{F,i+1}\); the chain identity verifies that this commutes with differentials. If \(\alpha-\beta=d_Gh+hd_F\), set \[
k^i(u)=u h_{i-1}.
\] Then \(\delta_F k^i(u)+k^{i+1}\delta_G(u)
=u h_{i-1}d_{F,i}+u d_{G,i+1}h_i=u(\alpha_i-\beta_i)\). Thus the cochain maps are homotopic and induce the same \(H^i\). In particular \(\alpha:F\to G\), \(\beta:G\to F\) give \[
H^i(\alpha^*)H^i(\beta^*)=H^i((\beta\alpha)^*)
\] on the cohomology of \(\operatorname{Hom}_R(F_\bullet,N)\), with identity \(\operatorname{id}_F\); the other composite is on \(G\). This verifies all parts of the retained MC-STK-ERR-1594 and MC-STK-ERR-1596 corrections. The \(H_i\) typo is already corrected by MC-STK-ERR-1595. The two grammar fixes are already present as MC-STK-ERR-1603 and MC-STK-ERR-1604. These six canonical units comprise nine operations and must not be applied a second time.

The remaining grammar changes at 17646 and 17698 leave both long exact sequences intact. Termwise free modules make Hom exact in the second variable; the first-variable construction is termwise split, so applying Hom reverses the arrows and remains exact. These are precisely the variance directions in the later annihilation and depth arguments.

\subsubsection{Projective comparison with the same maps}\label{algebra-proof-007-projective-comparison-with-the-same-maps}

The comparison lemma at 17499--17546 holds when each \(F_i\) is projective, rather than free, with no change to the arbitrary target resolution or the source complex condition. Here projective means that a map from \(F_i\) lifts through every surjection.

Let \(f:M\to N\) be the fixed map. Lift \(F_0\to M\to N\) through the surjection \(N_0\to N\) to obtain \(\alpha_0\). If \(\alpha_0,\ldots,\alpha_{i-1}\) are chosen, then \(\alpha_{i-1}d_{F,i}\) lands in \(\ker d_{N,i-1}\); for \(i=1\) this means the kernel of the augmentation. The assertion follows from the chain identities and \(d_F^2=0\), or from the augmentation condition when \(i=1\). Exactness makes \(N_i\to\ker d_{N,i-1}\) surjective, so projectivity gives \(\alpha_i\).

For two lifts put \(\gamma=\alpha-\beta\). Lift \(\gamma_0\) to \(h_0:F_0\to N_1\), as its augmentation is zero. Having chosen the lower \(h_j\), put \(u_i=\gamma_i-h_{i-1}d_{F,i}\). Then \[
d_{N,i}u_i
=\gamma_{i-1}d_{F,i}
-(\gamma_{i-1}-h_{i-2}d_{F,i-1})d_{F,i}=0.
\] Exactness and projectivity lift \(u_i\) to \(h_i:F_i\to N_{i+1}\). Hence \(\gamma_i=d_{N,i+1}h_i+h_{i-1}d_{F,i}\) at every degree. The augmentation supplies the initial case; no undefined negative module is needed.

Applying these constructions to projective resolutions in both directions yields inverse chain maps up to homotopy. The exact precomposition calculation above therefore identifies their Hom cohomology with the original Ext definition, naturally and contravariantly in the first module. Projective modules also make Hom exact in the second variable by the same lifting property. The original source chooses free resolutions, which already exist, and stays unchanged; this is supplementary generality, not a proposed replacement definition.

\subsubsection{Depth quotients and propagation to large powers}\label{algebra-proof-007-depth-quotients-and-propagation-to-large-powers}

Sources: 17759--18021; OCC-00431 and OCC-00432.

The source \(N\cong R/\mathfrak p\) gives the map \[
R/(\mathfrak p+(x))\longrightarrow N/xN,\qquad
r+(\mathfrak p+(x))\longmapsto r n_0+xN
\] for the chosen generator \(n_0\) with annihilator \(\mathfrak p\). It is surjective, and \(rn_0\in xN\) means \(r-xs\in\mathfrak p\) for some \(s\), proving its kernel is zero and giving the inverse through any representative. Thus both displayed occurrences require parentheses around \(\mathfrak p+(x)\). The original support, associated prime and dimension calculations keep the same quotient.

There is a uniform strengthening of lemma-inherit-minimal-primes. Let \(R\) be any Noetherian ring, \(M\) a finite module, \(\mathfrak p\in\operatorname{Ass}(M)\), and \(x\in R\). Choose the same \(N\cong R/\mathfrak p\subset M\). Artin--Rees gives \(c\geq0\) such that for every \(n\geq c\), \[
N\cap x^nM=x^{n-c}(N\cap x^cM).
\] For every \(n\geq c+1\), this is contained in \(xN\). Therefore \[
\overline N_n=N/(N\cap x^nM)\subset M/x^nM
\] surjects onto \(N/xN\), while it is annihilated by both \(\mathfrak p\) and \(x^n\). Consequently \[
\operatorname{Supp}(\overline N_n)=V(\mathfrak p+(x))
\] for all such \(n\): the surjection gives inclusion of the right-hand support, and the annihilators give the reverse inclusion. Every minimal prime of this closed set is an associated prime of the finite module \(\overline N_n\), hence of \(M/x^nM\). The single bound \(c+1\) works simultaneously for all primes minimal over \(\mathfrak p+(x)\). If that ideal is the unit ideal there are no such primes, and the assertion is empty. If \(x=0\), the formula remains valid and the assertion gives the original \(\mathfrak p\).

This proves the source conclusion for all sufficiently large \(n\), over an arbitrary Noetherian ring. It retains the actual submodule and its entire power-intersection formula. In the receiving depth proof at 17984--18020, choose any such \(n\); because its original \(x\) is a nonzerodivisor, so is \(x^n\), and the already established depth drop is one. No hypothesis of that local depth proof is removed.

The use of an \(x\) with depth drop exactly one in the earlier depth/Ext proof is not circular: finite depth attains a maximum length \(d\), choose the first element of a length-\(d\) regular sequence, and its tail shows depth at least \(d-1\). Any longer sequence on the quotient would concatenate with \(x\) to exceed \(d\), giving the reverse inequality. The Ext maps labelled by multiplication by this \(x\) vanish because \(x\) annihilates the original residue field. Induction gives finite first nonvanishing Ext and the stated equality. The later theorem then establishes this drop for every permitted nonzerodivisor. These source steps require no new correction.

\subsubsection{Reading boundaries and retained material}\label{algebra-proof-007-reading-boundaries-and-retained-material}

The entire current interval differs from the authority by nine already recorded canonical operations and eight reference/history environments attached to four Ext lemmas. Comparison after accounting for precisely those additions is exact. Those environments and all other existing work are preserved; comparison alone does not re-adjudicate their external FAC claims.

The existing corpus index was queried once for affine blowup. Its two hits were routed to original TeX. Bhatt--Scholze, Projectivity of the Witt vector affine Grassmannian, lines 251--284, supply contextual valuation-cover definitions and examples, with their cited dependencies not adopted here. Bhatt--Lurie, Absolute Prismatic Cohomology, lines 11899--11920, cover a regular-immersion blowup setup and only the start of the syntomic formula, not its proof. Bezrukavnikov--Finkelberg--Mirkovic, Equivariant (K-)homology of affine Grassmannian and Toda lattice, lines 1695--1814, cover the earlier symbolic-chart proposition through its proof, comments and the next proof opening. Its root-theoretic and flatness dependencies are not fully reviewed, and it is not adopted as a proof of the present algebraic claims. In particular symbolic powers in that setting are not substituted for the ordinary Rees powers here.

The proofs used for this batch are the exact supplied Stacks arguments and the explicit comparisons above. No exhaustive reading, novelty claim, global synthesis, rendering or publication is claimed. Broader synthesis follows the core report review.

\hypertarget{algebra_depth_tor_projective_notes_20260927}{}
\subsection{Editorial evidence for Algebra 18022--18862}\label{algebra-proof-008-editorial-evidence-for-algebra-1802218862}

Primary source: Stacks Project authors, algebra.tex, commit a04446e57ec1fbc252a871afcec7752fb2807b14, SHA-256 FA8BB92E58A4F78A2BD01B3B6A4A87DE0A0D279F5DD90641B574DD5FBFFFA4F3. The entire interval and seventeen due reports were read. The next direct-sum lemma, 18863--18870, is only partial lookahead. Source corrections below are proposals; the full arguments and strengthenings stay in source-linked editorial material.

\subsubsection{Finite-map depth and the infinite boundary}\label{algebra-proof-008-finite-map-depth-and-the-infinite-boundary}

Source: lemma-depth-goes-down-finite, 18047--18082; OCC-00433.

Use \(\min\varnothing=\infty\). If \(S=0\), then \(N=0\), the set of maximal ideals is empty, and both sides are infinity. If \(N=0\) over nonzero \(S\), each displayed depth is infinity and equality again holds directly. These cases do not enter an induction on a finite integer.

For \(N\ne0\), finite \(S\) over Noetherian local \(R\) is Noetherian and semilocal. Some maximal localization of \(N\) is nonzero: choose a prime in its nonempty support and then a maximal ideal containing it; localization at the latter cannot vanish because the support of a finite module is closed. Every nonzero such localization has finite depth by the preceding local dimension bound. Thus the minimum \(k\) is a finite nonnegative integer.

At \(k=0\), the source associated-prime and contraction argument gives \(\mathfrak m\in\operatorname{Ass}_R(N)\), hence depth zero. For \(k>0\), no maximal ideal of \(S\) is associated to \(N\); the established associated-prime comparison gives \(\mathfrak m\notin\operatorname{Ass}_R(N)\). Choose the source nonzerodivisor \(f\in\mathfrak m\). At each nonzero maximal localization, its image is in the local maximal ideal, multiplication is injective, and the quotient is nonzero by Nakayama. Its depth drops by one. Zero localizations remain zero and keep infinite depth. The same finite depth drop holds for the nonzero finite \(R\)-module \(N\). Therefore the finite minimum drops from \(k\) to \(k-1\), and induction proves the equality. No arithmetic of the form an unqualified finite induction at infinity is used.

The preceding localization inequality retains its source proof: if \(\mathfrak p\subset\mathfrak q\in\operatorname{Ass}(M)\), then the quotient map \(R/\mathfrak p\to R/\mathfrak q\) gives \(\dim(R/\mathfrak p)\geq\dim(R/\mathfrak q)\geq\operatorname{depth}(M)\). Under the opposite strict inequality prime avoidance supplies the required nonzerodivisor in \(\mathfrak p\); both nonzero depths drop by one. Vanishing local modules retain depth infinity.

\subsubsection{Ext base change and the original splitting obstruction}\label{algebra-proof-008-ext-base-change-and-the-original-splitting-obstruction}

Sources: 18087--18203; OCC-00434, OCC-00435 and OCC-01392.

For the original flat map \(R\to R'\) and original free resolution \(F_\bullet\), the degreewise adjunction sends \[
u:F_i\otimes_RR'\to N'
\quad\longmapsto\quad (x\mapsto u(x\otimes1)).
\] Its inverse sends \(v:F_i\to N'\) to \(x\otimes r'\mapsto r'v(x)\). Balancing and \(R'\)-linearity prove both maps well-defined, and evaluating on \(x\otimes r'\) proves they are inverse. They commute with precomposition by each original differential, so the resulting map is on cohomology. The two Ext operator changes are typographical consistency, not changes of functor. The missing article in ``is an injection'' is grammatical.

For the splitting lemma keep \(\varphi:N\to M\), the original finite modules, and \(I\). Injectivity of the reductions puts \(\ker\varphi\) in every sufficiently small \(I\)-power submodule: cofinally large \(n\) implies containment in every \(I^mN\). When \(I\) lies in the Jacobson radical, the original Krull-intersection result makes this kernel zero.

After this injectivity, let \(Q=M/N\) and retain the displayed initial segment of a resolution by finite free modules. Choose the stated \(\alpha:F_0\to M\). Its composite with \(d_1\) lands in \(N\) and defines \(\beta:F_1\to N\); \(\beta d_2=0\). If \(\beta=\gamma d_1\), then \(\alpha-\varphi\gamma\) vanishes on \(\operatorname{im}d_1\), hence descends to \(s:Q\to M\). The quotient map sends \(s\) to the identity. Conversely a section of \(M\to Q\) makes \(\alpha-s(F_0\to Q)\) factor uniquely through \(N\), giving such a \(\gamma\). This proves the exact obstruction assertion.

For a split reduction choose its retraction \(r_n:M/I^nM\to N/I^nN\). The composite with the reduction of \(\alpha\) is \(\overline\gamma_n\), and \(\overline\gamma_n d_1=r_n\overline\varphi\,\overline\beta=\overline\beta\). Finite freeness of \(F_0\) lifts that map to \(\gamma_n:F_0\to N\). Finite freeness of \(F_1\) gives the exact equality \[
\operatorname{Hom}_R(F_1,I^nN)
=I^n\operatorname{Hom}_R(F_1,N):
\] in a fixed finite free basis both sides consist of tuples with every entry in \(I^nN\). Thus the class of \(\beta\) in the original finite cokernel \[
C=\operatorname{coker}\bigl(\operatorname{Hom}_R(F_0,N)
\longrightarrow\operatorname{Hom}_R(F_1,N)\bigr)
\] belongs to every \(I^nC\). The same intersection theorem makes it zero. This supplies the retraction, without assuming a finite-length free resolution or a compatible inverse system of retractions.

\subsubsection{Splitting on one principal neighbourhood}\label{algebra-proof-008-splitting-on-one-principal-neighbourhood}

There is a stronger geometric conclusion for an arbitrary ideal \(I\) of a Noetherian ring: under the same cofinally many split reductions, there exists \(s\in1+I\) such that the original localized map \(N_s\to M_s\) is a split injection. The original global theorem keeps its Jacobson-radical hypothesis.

First let \(D=\bigcap_n I^nN\). The complete determinant/Artin--Rees argument in ALGEBRA-RECON-344 gives \(s_1\in1+I\) with \(s_1D=0\). Since \(\ker\varphi\subset D\), the map \(\varphi_{s_1}\) is injective. Put \(A=R_{s_1}\), \(J=IA\), and retain \(N_A,M_A,Q_A=M_A/N_A\). Localization of the original split reductions gives split maps modulo the same cofinal powers of \(J\).

Repeat the finite-free construction of the preceding section over \(A\). It gives the actual class \(c\) of its \(\beta_A\) in the finite cokernel \(C_A\), with \(c\in\bigcap_nJ^nC_A\). The same uniform-annihilator result gives \(s_2\in1+J\) annihilating this intersection. Therefore \(c\) becomes zero over \(A_{s_2}\), and the explicit section construction splits \(\varphi\) there.

Write the actual \(s_2\) as \(1+u/s_1^k\) with \(u\in I\), \(k\geq0\); this is possible since \(J=IR_{s_1}\). Set \[
s=s_1(s_1^k+u).
\] Both factors are congruent to one modulo \(I\), so \(s\in1+I\). Inverting their product in a commutative ring inverts each factor; conversely both factors are units in \(A_{s_2}\), since \(s_1^k+u=s_1^ks_2\). The two universal localization maps therefore give inverse isomorphisms \(R_s\cong(R_{s_1})_{s_2}\), preserving the original module fractions and map. The split injection is consequently on the one original principal open \(D(s)\), which contains \(V(I)\).

If \(I=R\), take \(s=0\); the neighbourhood is empty and the localized modules are zero. If \(R=0\), this handles it as well. If \(I=0\), the hypothesis is already that \(\varphi\) splits, and \(s=1\) works. If an intermediate localization is zero, the same explicit localization conclusion is valid. For \(I\) contained in the Jacobson radical, every member of \(1+I\) is a unit, recovering the original global assertion.

The unrestricted global assertion would be false. For \(R=k\times k\), \(I=k\times0\), \(N=k\times0\), \(M=0\), and \(\varphi=0\), every positive-power reduction is \(0\to0\) and splits, while \(\varphi\) is not injective. The actual \(s=(0,1)\in1+I\) kills \(N\) after localization and gives exactly the proved neighbourhood conclusion. This strengthening receives the uniform annihilator and neighbourhood maps from groups 344--345 and remains separate editorial material.

\subsubsection{The zig-zag quotient with its original signs}\label{algebra-proof-008-the-zig-zag-quotient-with-its-original-signs}

Source: lemma-no-spectral-sequence, 18397--18522; OCC-00436, OCC-12070, OCC-00437, OCC-00438, OCC-01393 and OCC-00439.

Keep the commuting squares, \(d^2=\delta^2=0\), all bidegrees, and the original positive cycle equations. For a fixed \(n\geq0\), let \(Z_n\) be exactly the tuples \((a_{n,0},\ldots,a_{0,n})\) with \[
d a_{p+1,q}=\delta a_{p,q+1}\qquad(p+q=n-1).
\] The original denominator is the image \(B_n\) of \[
(b_{n+1,0},\ldots,b_{0,n+1})
\longmapsto
(d b_{n+1,0}+\delta b_{n,1},\ldots,
 d b_{1,n}+\delta b_{0,n+1}).
\] Every image is a cycle: \(d(d b+\delta b)=\delta d b\), whereas the adjacent \(\delta(d b+\delta b)=\delta d b\). Thus \(B_n\subset Z_n\) over every ring, with no characteristic assumption.

This image is generated by both endpoint tuples, whose only components are \(d b_{n+1,0}\) in position \((n,0)\) and \(\delta b_{0,n+1}\) in position \((0,n)\), and the adjacent pairs \[
(\ldots,\delta b_{t+1,n-t},d b_{t+1,n-t},\ldots),
\qquad0\leq t<n.
\] Their positions are respectively \((t+1,n-t-1)\) and \((t,n-t)\). The input has total degree \(n+1\), both outputs have total degree \(n\), and both signs are positive. For \(n=0\) only the two endpoints occur, in the same sole component \(A_{0,0}\), and \(B_0=dA_{1,0}+\delta A_{0,1}\).

The source paragraph has three linked defects: the reversed second index, the minus sign on one member of the pair, and missing endpoints. The sign is not a harmless alternative convention while the cycle equations remain as printed. Take \(A_{1,1}=A_{1,0}=A_{0,1}=A_{0,0}=\mathbf Z\), all four square maps the identity, and other terms zero. Every row and column is a resolution of its cokernel. For \(b_{1,1}=1\) the source pair \((-1,1)\) fails \(d a_{1,0}=\delta a_{0,1}\), while \((1,1)\) satisfies it. For the endpoint issue take only \(A_{1,0}=A_{0,0}=\mathbf Z\) with \(d=\mathrm{id}\). Again the hypotheses hold, and \(B_0=A_{0,0}\); a definition allowing only interior generators would incorrectly leave it unquotiented. The proposed paragraph corrects all three defects while retaining the displayed quotient and original cycle convention.

The right endpoint map is \[
\rho:Z_n/B_n\to H_n(R(A)_\bullet),\qquad
[a]\mapsto[\text{image of }a_{0,n}].
\] The cycle equation makes the image a cycle; right endpoint generators become \(\delta\)-boundaries, interior generators reaching the right end become \(d\)-images, and all others map to zero. Therefore this is well-defined.

For surjectivity lift a cycle \(r\in R(A)_n\) to \(a_{0,n}\). If \(n>0\), its \(\delta\)-image is a horizontal boundary, so lift it to \(a_{1,n-1}\). At each next step, commutation and \(\delta^2=0\) show \(\delta a_{p,n-p}\) is in the kernel of \(d\). Exactness of the row at positive horizontal degree lifts it to \(a_{p+1,n-p-1}\). This terminates at \(p=n\). For \(n=0\) the initial lift already is a tuple, without a negative degree.

For injectivity suppose the class of \(a_{0,n}\) is a boundary. With the correct source indices this means \[
a_{0,n}=\delta b_{0,n+1}+d b_{1,n}.
\] Subtract the right endpoint generator and the generator contributed by \(b_{1,n}\) (a left endpoint if \(n=0\)); the last component becomes zero. If \(n>0\), the cycle equation now says \(d a_{1,n-1}=0\), and exactness gives \(a_{1,n-1}=d b_{2,n-1}\). Subtract its adjacent-pair generator, or the left endpoint when \(n=1\). Continue: when all components with horizontal index less than \(p\) are zero, \(d a_{p,n-p}=0\); choose \(b_{p+1,n-p}\) with \(d b=a_{p,n-p}\), and subtract its generator. The other affected component has higher horizontal index, so no cleared component returns. The final \(p=n\) step uses the left endpoint. The entire tuple is a boundary. This proves injectivity.

The top endpoint map to \(H_n(U(A)_\bullet)\) is proved by the same construction with the two indices and \(d,\delta\) interchanged; the equations and positive boundary formula are invariant under that interchange, and columns replace rows. A morphism of double complexes acts componentwise, preserves \(Z_n,B_n\), and commutes with both endpoint maps, proving their functoriality. No total-complex sign convention has replaced the original data.

The unmatched parenthesis is already fixed by MC-STK-ERR-0413. The remaining minimal changes call a homological representative a cycle, close the parenthetical calculation, and restore the three \(R(A)\) cokernels to their actual degrees. In particular the preceding boundary group is \(\operatorname{coker}(A_{1,n+1}\to A_{0,n+1})\), not a transposed \(U(A)\) term.

\subsubsection{Tensor comparison, swap and finiteness}\label{algebra-proof-008-tensor-comparison-swap-and-finiteness}

Sources: 18533--18765; OCC-00440, OCC-01394, OCC-00441, OCC-12071 and OCC-01395.

For \(A_{p,q}=F_p\otimes_RG_q\), the horizontal and vertical maps land in \(F_{p-1}\otimes_RG_q\) and \(F_p\otimes_RG_{q-1}\). Writing the two omitted subscripts makes the already specified base ring explicit; it does not change the tensor product. The corrected zig-zag quotient proves the source natural comparison between \(F_\bullet\otimes_RN\) and \(M\otimes_RG_\bullet\). Tensoring chain homotopies keeps their same covariance and degrees; it must not copy the contravariant arrow order of Hom. The original tensor swap is \(x\otimes y\mapsto y\otimes x\). No assertion is made here that the author's unsigned convention is the usual signed symmetry of every derived construction.

For a basis \(e_1,\ldots,e_n\) of \(V\), preserve the decomposition \[
V\otimes_kV
=\bigoplus_{i=1}^n k(e_i\otimes e_i)
\ \oplus\!
\bigoplus_{1\leq i<j\leq n}
\bigl(k(e_i\otimes e_j)\oplus k(e_j\otimes e_i)\bigr).
\] The swap fixes the first \(n\) lines and exchanges the two vectors in every remaining pair. When \(\operatorname{char}k\ne2\), the sum and difference vectors form a basis of each pair, with eigenvalues \(+1,-1\); their coordinate change has determinant \(-2\), a unit. Thus the multiplicities are \(n(n+1)/2\) and \(n(n-1)/2\).

In characteristic two, each pair has \(\tau-1\) of rank one and square zero: with \(u=e_i\otimes e_j\), \(v=e_j\otimes e_i\), \((\tau-1)u=u+v=(\tau-1)v\), and \((\tau-1)(u+v)=0\). The basis \((u+v,u)\) gives one size-two Jordan block with eigenvalue one. The diagonal lines give \(n\) size-one blocks. Hence the exact form is \(n\) blocks of size one and \(n(n-1)/2\) blocks of size two; it is not diagonalizable precisely for \(n\geq2\). At \(n=0,1\) it is the identity. The source statement that this swap is always nonidentity also needs ``need not be''; the higher Tor example does not remove these degree-zero exceptions.

The Tor-finiteness homology wording is already corrected as MC-STK-ERR-0414. The filtered-colimit sentence only needs its missing verb. Tensor products commute with the given filtered colimit and exactness of that colimit commutes with kernels and images, so it also commutes with the original homology quotient.

\subsubsection{Flat resolutions compute the same Tor}\label{algebra-proof-008-flat-resolutions-compute-the-same-tor}

The corrected double-complex argument gives further generality without rewriting the source definition: any resolution \(P_\bullet\to M\) by flat \(R\)-modules computes the original Tor through \(H_n(P_\bullet\otimes_RN)\).

Choose an original free resolution \(G_\bullet\to N\). In \(A_{p,q}=P_p\otimes_RG_q\), each row resolves \(M\otimes_RG_q\), since \(G_q\) is flat, and each column resolves \(P_p\otimes_RN\), since \(P_p\) is flat. Both cokernel identifications send the original simple tensors through the corresponding augmentation maps. The proved endpoint isomorphisms give \[
H_n(P_\bullet\otimes_RN)
\cong H_n(M\otimes_RG_\bullet)
\cong\operatorname{Tor}_n^R(M,N),
\] where the last map uses exactly the preceding original Tor comparison. Thus no projective lifting property for the flat terms is assumed. A map between flat resolutions and a lifted map between the free \(G\)-resolutions give a map of double complexes; functoriality of the endpoints proves naturality for such maps. Changing the free resolution of \(N\) gives inverse homotopy comparisons; tensoring those homotopies preserves them, so the isomorphism is independent of that choice. The existing projective comparison in group 460 is a receiving special case, but supplies no lifting assumption for arbitrary flat modules.

\subsubsection{Projectors and the existing projective definition}\label{algebra-proof-008-projectors-and-the-existing-projective-definition}

Sources: 18780--18862; OCC-00442 and OCC-01396.

For \(F=P\oplus Q\), retain the actual two idempotents \(a(p,q)=(p,0)\) and \(b(p,q)=(0,q)\). The augmentation is \(\pi(p,q)=p\). Then \(\ker\pi=Q=\operatorname{im}b\), \(\ker b=P=\operatorname{im}a\), and \(\ker a=Q=\operatorname{im}b\). These equalities prove the alternating displayed free resolution in every degree. The leading ellipsis needs its arrow, and the prose needs the plural ``projectors \ldots{} respectively''; neither repair changes the maps.

After applying Hom the decomposition into maps from \(P\) and from \(Q\) gives exactness in each positive degree. In the converse, the source map \(s:F_0\to\ker(F_0\to P)\) restricts to the identity on that kernel because its composite with the surjection \(F_1\to\ker(F_0\to P)\) is that surjection. Therefore \[
F_0\longrightarrow\ker s\oplus\ker(F_0\to P),\quad
f\longmapsto(f-s(f),s(f))
\] and addition are inverse. The augmentation restricts to an isomorphism \(\ker s\to P\), proving the source characterization with its actual maps. This also verifies that the lifting convention used in group 460 agrees exactly with the now-read Hom-exactness definition. The next finite-projective criterion uses the same single finite kernel, as printed; its author's announced extensions are not new claims of this review.

\subsubsection{Reading and preservation limits}\label{algebra-proof-008-reading-and-preservation-limits}

The current interval differs from the authority by two already recorded canonical operations and one FAC reference/history pair at the projective definition. Those additions are retained, without a fresh certification of the external FAC citation.

Two bounded corpus queries, first literal terms and then the exact phrase ``double complex'', were used. PDF and unrelated local hits remain unread routing records. Antti J. Harju's original TeX, Twisted K-theory constructions in the case of a decomposable Dixmier-Douady class II: Topological and Equivariant Models, lines 295--312, was read for the groupoid double-complex conventions and section context. Its Morita-invariance and de Rham assertions cite dependencies not read here and are not adopted. Its cohomological groupoid complex is not substituted for the source's commuting homological square.

The proofs above use the exact source definitions and the named previously proved intersection result. There is no claim of novelty, exhaustive literature coverage, rewritten translation, formal certification, rendering or publication. Broader combined-results synthesis remains after core completion.

\hypertarget{algebra_projective_lifting_notes_20260927}{}
\subsection{Editorial evidence for Algebra 18863--19374}\label{algebra-proof-009-editorial-evidence-for-algebra-1886319374}

Primary source: Stacks Project authors, algebra.tex, commit a04446e57ec1fbc252a871afcec7752fb2807b14, SHA-256 FA8BB92E58A4F78A2BD01B3B6A4A87DE0A0D279F5DD90641B574DD5FBFFFA4F3. The complete interval and ten due reports were read. The proofs and consequences below are editorial evidence, distinct from the proposed minimal source corrections. They are not replacements for the translated exposition.

\subsubsection{Original projectors and nilpotent lifting}\label{algebra-proof-009-original-projectors-and-nilpotent-lifting}

Sources: 18863--18958; earlier idempotent lemmas 6318--6328 and 6385--6399 were consulted at their actual source loci.

If each original projective module \(P_\alpha\) is a summand of \(F_\alpha\), the direct sum of the inclusions and projections has composite identity on \(\bigoplus P_\alpha\). Its ambient \(\bigoplus F_\alpha\) is free on the disjoint union of the original bases. This proves the direct-sum assertion, including the empty sum.

For the arbitrary free module \(F=\bigoplus_{\alpha\in A}R\), lift each column of the original projector \(\overline p\) with its finite support. The resulting endomorphism \(p\) is well-defined on the direct sum. Put \(x=p^2-p\). If \(I^n=0\), then \(x(F)\subset IF\), and induction gives \(x^j(F)\subset I^jF\); hence \(x^n=0\). This argument uses the original uniform nilpotence and does not replace an infinite matrix by a finite one.

The cited noncommutative lemma applies to this actual \(x\). Its proof uses the homomorphism \(\mathbf Z[e]/((e^2-e)^n)\to\operatorname{End}_R(F)\) sending \(e\) to \(p\); it is a ring homomorphism because all polynomials in one endomorphism commute. In that commutative source ring, the ideal \((e^2-e)\) is nilpotent. The earlier idempotent-lifting lemma supplies an idempotent congruent to \(e\), giving exactly an idempotent \(p'\) congruent to \(p\) modulo endomorphisms with image in \(IF\). Thus \[
F=\operatorname{im}p'\oplus\operatorname{im}(1-p'),
\qquad
(\operatorname{im}p')/I(\operatorname{im}p')\cong\operatorname{im}\overline p.
\] For the last equality, reduce the displayed direct-sum decomposition modulo \(I\); the induced idempotent is precisely \(\overline p\). This also proves that the quotient map is the original reduction, rather than an unspecified isomorphism between modules of the same rank.

For a finite projective quotient, the image of finitely many generators in a free direct sum uses a finite subset of its basis. The original retraction restricts to that finite free submodule because its image is the same projective module contained there; consequently it remains a direct summand. Retain the source matrix \(p\in\operatorname{Mat}(n\times n,R)\). There are finitely many entries \(c_{ij}\) of \(p^2-p\), so a common positive \(t\) satisfies \(c_{ij}^t=0\). Every entry of \((p^2-p)^{tn^2}\) is a finite sum of products of \(tn^2\) such entries. In each product one of the \(n^2\) entries occurs at least \(t\) times. Since \(R\) is commutative, that product is zero. The printed bound therefore works without assuming the entire ideal is nilpotent. When \(n=0\), both the module and its lift are zero. Applying the same polynomial correction preserves the original reduction and gives a summand of the same finite free module.

In the flat lifting lemma, for each original exponent \(a\), flatness of a module \(T\) identifies \[
(I^a/I^{a+1})\otimes_RT\longrightarrow I^aT/I^{a+1}T,
\quad [u]\otimes t\longmapsto[ut]
\] as an isomorphism: tensor the two inclusions of ideals into \(R\), and then the exact sequence \(0\to I^{a+1}\to I^a\to I^a/I^{a+1}\to0\). Because \(I\) kills the first factor, this tensor product is also canonically \((I^a/I^{a+1})\otimes_{R/I}(T/IT)\), with the displayed formula unchanged. Apply this to the actual lift \(P\to M\). Its degree-zero reduction is an isomorphism, so every map on successive quotients is an isomorphism. Starting at \(I^nP=I^nM=0\), injectivity and surjectivity in each short exact sequence of successive filtration steps prove, by downward induction, that \(P\to M\) is an isomorphism. No unbounded filtration argument is being used.

\subsubsection{Patching over the two original ideals}\label{algebra-proof-009-patching-over-the-two-original-ideals}

Source: 18960--18989; OCC-01397 and OCC-01398.

Keep the source surjection \(p:F\to P\), ideals \(I,J\) with \(I\cap J=0\), and original maps \(f,q,f',g\). To prove the omitted surjectivity of \(q\), let \((\overline v,\overline w)\) be in its target fibre product. Choose representatives \(v\in F\), \(w\in P\). Compatibility says \(w-p(v)\in(I+J)P\). Write it as \(x+y\) with \(x\in IP\), \(y\in JP\). The equality \(p(IF)=IP\) supplies \(u\in IF\) with \(p(u)=x\). The class of \(v+u\) in \(F/JF\) maps to the specified pair: its first component is \(\overline v\), and its image under \(p\) is \(w-y\), which has class \(\overline w\) modulo \(JP\). This proves surjectivity using exactly the original modules.

The map induced by \(f\) on \(P/(I+J)P\) is well-defined because the \(R\)-linear \(f\) sends the image of \(JP\) in \(P/IP\) into the image of \(JF\) in \(F/IF\). This defines \(f'\). Projectivity over \(R/J\) then gives the stated \(g\), and \(qg=f'\) is the exact compatibility of their reductions.

The natural map \[
F\longrightarrow(F/IF)\times_{F/(I+J)F}(F/JF)
\] is injective since, coordinate by coordinate on the original free basis, \(IF\cap JF=(I\cap J)F=0\). It is surjective: if representatives \(u,v\in F\) satisfy \(u-v=i+j\) with \(i\in IF\), \(j\in JF\), then \(u-i=v+j\) represents both required classes. The construction gives the original \(h:P\to F\), whose reductions are \(f,g\).

Set \(a=ph\) and \(t=a-\operatorname{id}_P\). The reductions of \(a\) are identities, so \(t(P)\subset IP\cap JP\). For each \(x_\alpha\in P\), the difference \[
j_\alpha=a(x_\alpha)-x_\alpha
\] is an element of \(JP\), not an element of the ideal \(J\). Write \(j_\alpha=\sum_\beta j_{\alpha\beta}p_{\alpha\beta}\) with \(j_{\alpha\beta}\in J\). For \(i_\alpha\in I\), every \(i_\alpha j_{\alpha\beta}\) is zero because \(IJ\subset I\cap J=0\). Therefore \(i_\alpha j_\alpha=0\), and \(t(IP)=0\). Since \(t(P)\subset IP\), it follows that \(t^2=0\).

In particular the omitted inverse has the exact formula \[
a^{-1}=\operatorname{id}_P-t=2\operatorname{id}_P-a,
\qquad
\sigma=h(2\operatorname{id}_P-a),\qquad p\sigma=\operatorname{id}_P.
\] Both products \((1+t)(1-t)\) and \((1-t)(1+t)\) are \(1-t^2=1\); this uses no division by two and works in every characteristic. The explicit section proves projectivity. The proposed source edits insert the missing ``is'' and correct membership to \(JP\); they do not replace the author's omitted details with this full proof.

\subsubsection{Finite-projective criteria and the localization element}\label{algebra-proof-009-finite-projective-criteria-and-the-localization-element}

Sources: 19000--19208; OCC-01310, OCC-00443, OCC-01399 and OCC-01400.

The tab on line 19024 is a whitespace byte in source, not an extra mathematical symbol or a rendered mathematical error. Removing it is source hygiene only. Uniqueness of rank still requires a nonzero ring, as the author states.

For OCC-00443 retain the original prime \(\mathfrak p\), element \(g\notin\mathfrak p\), and surjection \(\varphi:R_g^r\to M_g\). Its finite kernel has zero fibre at \(\mathfrak pR_g\), so localized Nakayama supplies \(g'\in R_g\setminus\mathfrak pR_g\) with \((\ker\varphi)_{g'}=0\). Write the actual \(g'\) as \(h/g^a\), where \(a\geq0\), \(h\in R\). The avoidance condition gives \(h\notin\mathfrak p\). Inverting \(gh\) inverts \(g\) and \(h\), hence \(g'=h/g^a\). Conversely in \((R_g)_{g'}\), the equality \(h=g^ag'\) makes \(gh\) invertible. The universal localization maps are inverse and preserve the original fractions. Tensoring them with the original module gives \[
(M_g)_{g'}\cong M_{gh}.
\] Thus \(D(gh)\) is an original-ring neighbourhood of \(\mathfrak p\), and \(M_{gh}\) is free of the same rank \(r\). Writing \(gg'\) as if \(g'\in R\) did not specify this original-ring localization. The proposed correction retains the actual fraction and its numerator, rather than silently changing the chosen ring.

The rest of the equivalence proof retains its hypotheses. A finite free summand is finite and projective; its complement is finite, giving a finite presentation of the original quotient. The original finite-presentation hypothesis permits the Hom-localization comparison. For a short exact sequence, the three localized Hom terms agree with Hom out of the locally free module; their exactness detects the original exactness on the open cover. The two grammar repairs in this paragraph affect neither that map nor its cover.

For (8), the source map \(\Psi:R_{fg}^r\to M_{fg}\) is surjective and, at every prime in \(D(fg)\), both source and target are free of rank \(r\). After choosing bases locally its matrix has a right inverse. Taking determinants makes its determinant a unit; the adjugate formula supplies the inverse. The zero-rank case is the unique map between zero modules. This validates the map referred to by the missing subject in OCC-01400.

In the reduced-ring criterion keep the same \(\Psi\). At each minimal prime the localized ring is the original residue field and the two finite dimensions coincide, so the kernel localizes to zero. A vector in the kernel has each coordinate zero in every such localization. The canonical injection of a reduced ring into the product of these fields kills no nonzero coordinate; hence the vector is zero. Surjectivity was already established. This proves the source's conclusion without deleting the reducedness assumption.

\subsubsection{The smooth-function counterexample}\label{algebra-proof-009-the-smooth-function-counterexample}

Source: 19210--19220; OCC-01401.

Keep \(R=\mathcal C^\infty(\mathbf R)\), \(\mathfrak m=\{f:f(0)=0\}\), and the printed ideal \(I\) of functions vanishing on some neighbourhood of zero. The kernel of \(R\to R_{\mathfrak m}\) is exactly \(I\). If \(gf=0\) with \(g(0)\ne0\), continuity makes \(g\) nonzero near zero, so \(f\) vanishes there. Conversely, for \(f\in I\), choose a smooth cutoff \(\chi\), equal to one near zero and supported inside a neighbourhood where \(f=0\). Then \(\chi f=0\) and \(\chi\notin\mathfrak m\).

This localization map is surjective. Given \(f/g\) with \(g(0)\ne0\), choose such a cutoff supported in the open set where \(g\ne0\), and define \(h=\chi f/g\) on that open set and zero outside. Its support is contained there, so the extension is smooth. Near zero, \(gh=f\), and therefore \(f-gh\in I\); its localization vanishes. Thus \(f/g=h/1\). These constructions prove the author's exact isomorphism \(R/I\cong R_{\mathfrak m}\), with the actual quotient map.

Consequently \(M\) is cyclic and flat. If it were projective, the quotient \(R\to R/I\) would split and its kernel \(I\) would be an idempotent-generated summand of \(R\). Every smooth idempotent takes only values zero and one, and continuity on connected \(\mathbf R\) forces it to be constant. But \(I\) is nonzero (take a nonzero bump function supported away from zero) and proper (it does not contain one). This contradiction proves the source counterexample. The accepted edit only joins the noun ``counterexample''; it is not a mathematical correction to this example.

\subsubsection{Local flatness, descent and semilocal bases}\label{algebra-proof-009-local-flatness-descent-and-semilocal-bases}

Sources: 19222--19347; OCC-00444 and OCC-01402.

For the local flatness proof, the equational criterion converts each finite relation \(\sum_i f_i x_i=0\) into the source equations \(x_i=\sum_j a_{ij}y_j\) and \(\sum_i f_i a_{ij}=0\). If the residues of \(x_1,\ldots,x_n\) are independent, then some \(a_{nj}\) is a unit. Substitution gives exactly \[
f_n=\sum_{i<n}(-a_{ij}/a_{nj})f_i,
\qquad
\sum_{i<n}f_i\bigl(x_i+(-a_{ij}/a_{nj})x_n\bigr)=0.
\] If a residue-field linear combination of the displayed \(n-1\) modified vectors vanishes, independence of the original \(n\) residues makes every coefficient zero. Induction therefore gives every \(f_i=0\), and the displayed first equality then gives \(f_n=0\). In degree one, the same criterion gives a unit coefficient annihilated by \(f\), hence \(f=0\). In degree zero the zero set is independent. Nakayama gives the original finite generating family from a residue basis; its independence proves that a finite flat module over a local ring is free.

For the local descent lemma, the residue field is \(\kappa(\mathfrak m_R)\), correcting the otherwise undeclared ideal. For the original local map \((R,\mathfrak m_R)\to(S,\mathfrak m_S)\), its residue-field map is injective. The residue of the free \(S\)-module \(M\otimes_RS\) is \[
(M/\mathfrak m_RM)\otimes_{\kappa(\mathfrak m_R)}\kappa(\mathfrak m_S).
\] The selected basis tensors to a residue basis. Nakayama makes it a generating set, and its size equals the rank of the free \(S\)-module. The determinant argument just given makes the corresponding square matrix invertible. Faithful flatness of the original local map then kills both kernel and cokernel of \(R^r\to M\), proving the asserted descent. The forward implication follows by tensoring a direct-summand decomposition.

For constant rank \(d\) over a nonzero semilocal ring, the finitely many maximal ideals \(\mathfrak m_i\) are pairwise comaximal. The Chinese remainder map \(M\to\prod_iM/\mathfrak m_iM\) is surjective: choose ring elements that are one at one maximal ideal and zero at all others, and form their finite linear combination of representatives. Lift the \(j\)-th vector of a basis at every maximal ideal simultaneously to \(x_j\in M\), for \(1\leq j\leq d\). The map \(R^d\to M\) is an isomorphism at every maximal ideal by Nakayama and the square-matrix determinant argument. Its kernel and cokernel are therefore zero, since any nonzero module has a nonzero localization at some maximal ideal. This proves the hinted argument. Over the zero ring the only module is zero, so the assertion is immediate. Connected nonempty spectrum makes the locally constant rank constant, as stated by the author. Only the singular ``ideal'' in the hint is corrected in the source proposal.

For the basis-in-submodule lemma retain the source quotients by \(I=\operatorname{Jac}(S)\). Because \(\mathfrak mN\subset IM\), the image of \(N\) is indeed a vector space over the original infinite field \(k=R/\mathfrak m\). For \(S\ne0\), write the source quotient as \(\prod_{i=1}^r k_i\). When \(n>0\), choose finitely many \(x_j\in N\) with a nonzero image in every factor among them; such a family exists because \(N\) generates \(M\). For each \(i\), the \(k\)-linear map \((c_j)\mapsto\sum_j c_j(x_j)_i\) has a proper kernel. A finite union of proper subspaces cannot fill the finite-dimensional vector space over infinite \(k\): choose a nonzero linear functional vanishing on each subspace, take the product of those linear forms, and use induction on the number of variables to see that a nonzero polynomial cannot vanish on all of \(k^m\). Thus some original coefficient vector gives \(y\in N\) nonzero in every factor. In each factor its vector spans a direct summand; together these give \(Sy\cong S\) with free quotient of rank \(n-1\). Apply induction to the image of \(N\) and lift the resulting basis. The base \(n=0\), and the case \(S=0\), are immediate. These details verify the source use of infinitude; they introduce no replacement hypothesis.

\subsubsection{Constant fibre rank and the exact nilradical kernel}\label{algebra-proof-009-constant-fibre-rank-and-the-exact-nilradical-kernel}

This is a separate consequence of the original finite-projective lifting lemma, reduced-ring criterion, and local flatness result. No source statement is replaced.

Let \(R\) be any commutative ring, \(\mathcal N=\sqrt{(0)}\) its actual nilradical, and \(M\) a finite \(R\)-module. The following are equivalent:

\begin{enumerate}
\def\labelenumi{\arabic{enumi}.}
\tightlist
\item
  The original function \(\rho_M(\mathfrak p)=\dim_{\kappa(\mathfrak p)}(M\otimes_R\kappa(\mathfrak p))\) is locally constant.
\item
  There exists a finite projective \(R\)-module \(P\) and a surjection \(\pi:P\to M\) whose actual kernel \(K\) is contained in \(\mathcal NP\).
\end{enumerate}

For (1), set \(\overline R=R/\mathcal N\) and \(\overline M=M/\mathcal NM\). Each prime of \(R\) contains \(\mathcal N\), so the quotient gives a homeomorphism of spectra with the corresponding residue fields unchanged. The explicit fibre map \[
(M/\mathcal NM)\otimes_{R/\mathcal N}\kappa(\mathfrak p/\mathcal N)
\longrightarrow M\otimes_R\kappa(\mathfrak p),
\quad (m+\mathcal NM)\otimes c\longmapsto m\otimes c
\] is well-defined and has the inverse induced by \(m\mapsto m+\mathcal NM\). Hence the same rank function is locally constant for \(\overline M\). The proved reduced-ring criterion makes \(\overline M\) finite projective over \(\overline R\). Every finite subset of \(\mathcal N\) generates a nilpotent ideal: if its generators have powers \(a_i^{t_i}=0\), every product of more than \(\sum_i(t_i-1)\) generators vanishes. Thus the finite-projective lifting lemma applies to the actual nilradical, even if that ideal is not nilpotent. It gives finite projective \(P\) and an isomorphism \(\alpha:P/\mathcal NP\to\overline M\).

Projectivity lifts the composite \(P\to P/\mathcal NP\xrightarrow\alpha\overline M\) through the actual quotient \(M\to\overline M\), producing \(\pi:P\to M\). Its cokernel \(C\) is finite and satisfies \(C/\mathcal NC=0\). Since \(\mathcal N\) lies in every maximal ideal, Nakayama gives \(C=0\). If \(p\in\ker\pi\), then \(\alpha(p+\mathcal NP)=0\); injectivity of \(\alpha\) gives \(p\in\mathcal NP\). This proves the asserted surjection and exact kernel containment.

Conversely, such a surjection gives \[
M/\mathcal NM=P/(K+\mathcal NP)=P/\mathcal NP.
\] The residue-field maps are the same maps induced by \(\pi\). Therefore \(\rho_M=\rho_P\), a locally constant function because \(P\) is finite projective. This proves (2) implies (1).

For every such \(\pi\), the module \(M\) is flat if and only if \(K=0\). One implication is immediate from \(M\cong P\). For the other, if \(M\) is flat, then at every prime both \(P_{\mathfrak p}\) and \(M_{\mathfrak p}\) are finite free by the original local theorem. Their ranks agree by the proved fibre comparison. The surjection \(\pi_{\mathfrak p}\) is consequently a square invertible matrix (including rank zero), so \(K_{\mathfrak p}=0\) for every prime. A nonzero element of \(K\) would have a proper annihilator contained in a maximal ideal and would survive there, a contradiction. Hence \(K=0\). This proves an exact obstruction statement about the original kernel, without discarding it or asserting that some common nilpotence exponent exists.

The obstruction can be nonzero. For \(R=k[\epsilon]/(\epsilon^2)\), \(M=R/(\epsilon)\), the unique prime gives constant fibre rank one. The original quotient \(R\to M\) has kernel \((\epsilon)\subset\mathcal NR\), which is nonzero. Thus \(M\) is not flat or projective. Directly, tensoring the inclusion \((\epsilon)\hookrightarrow R\) with \(M\) gives the zero map from a nonzero copy of \(k\). This verifies that the reducedness hypothesis of the original projectivity criterion cannot simply be erased.

A receiving consequence is the finite version of the earlier lifting theorem for any nil ideal \(I\): if \(M\) is finite flat and \(M/IM\) is projective over \(R/I\), then \(M\) is projective over \(R\). Indeed \(M/IM\) is finite projective, its rank is locally constant, and the quotient by the nil ideal identifies spectra and residue fields exactly as above. Thus \(\rho_M\) is locally constant, and the just-proved kernel criterion forces \(M\) to be finite projective. Conversely base change of a finite-projective summand is finite projective. This does not claim the same relaxation for arbitrary nonfinite modules, and the source's unrestricted nilpotent theorem remains intact.

\subsubsection{Evaluation and the converse dual-basis criterion}\label{algebra-proof-009-evaluation-and-the-converse-dual-basis-criterion}

Source: 19349--19369; OCC-01403. Inserting the missing article is the only proposed source edit here.

Retain the canonical map \[
\Phi:\operatorname{Hom}_R(M,N)\otimes_RL\longrightarrow
\operatorname{Hom}_R(M,N\otimes_RL),
\quad f\otimes\ell\longmapsto(m\mapsto f(m)\otimes\ell).
\] For finite projective \(M\), choose a splitting of a finite free module and let \(m_i\in M\), \(\lambda_i\in M^*=\operatorname{Hom}_R(M,R)\) be the images of its basis vectors and restrictions of its coordinate functionals. The composite of projection and inclusion is the identity, so \[
m=\sum_{i=1}^n\lambda_i(m)m_i\quad\text{for every }m\in M.
\] For \(T:M\to N\otimes_RL\), write \(T(m_i)=\sum_j n_{ij}\otimes\ell_{ij}\). Define \[
\Theta(T)=\sum_{i,j}(m\mapsto\lambda_i(m)n_{ij})\otimes\ell_{ij}.
\] This is independent of each tensor presentation: for fixed \(i\), the map \(N\times L\to\operatorname{Hom}(M,N)\otimes L\), \((n,\ell)\mapsto(m\mapsto\lambda_i(m)n)\otimes\ell\), is balanced and therefore factors through \(N\otimes_RL\). Evaluating on \(m\) gives \(\Phi\Theta(T)(m)=\sum_i\lambda_i(m)T(m_i)=T(m)\). Conversely \(\Theta\Phi(f\otimes\ell)=\sum_i(m\mapsto\lambda_i(m)f(m_i))\otimes\ell=f\otimes\ell\). Thus these are actual inverse maps for the source assertion.

There is also an exact converse for an arbitrary module \(M\). It suffices that the identity belongs to the image of the single canonical map \[
\eta:M^*\otimes_RM\longrightarrow\operatorname{End}_R(M),
\quad \lambda\otimes m\longmapsto(x\mapsto\lambda(x)m).
\] If \(\operatorname{id}_M=\eta(\sum_{i=1}^n\lambda_i\otimes m_i)\), set \[
s:M\to R^n,\quad x\mapsto(\lambda_1(x),\ldots,\lambda_n(x)),
\qquad
t:R^n\to M,\quad(a_i)\mapsto\sum_i a_i m_i.
\] Then \(ts=\operatorname{id}_M\), proving that the original \(M\) is a finite-projective summand. Its complement is \(\ker t\): the explicit mutually inverse decomposition maps are \((x,z)\mapsto s(x)+z\) and \(v\mapsto(t(v),v-st(v))\). Conversely a finite-projective splitting gives the identity tensor just constructed. This includes \(M=0\) using the empty sum.

Taking \(N=R\), \(L=M\) in \(\Phi\), and using the canonical map \(R\otimes_RM\to M\), identifies that instance with \(\eta\). Hence the source isomorphism for all \(N,L\) is equivalent to finite projectivity, and the weaker requirement that this one map hit the identity already suffices. The original forward lemma is retained; this proved converse is recorded separately as an underclaim consequence.

\subsubsection{Reading and preservation limits}\label{algebra-proof-009-reading-and-preservation-limits}

The complete current interval agrees with the authority after removing exactly one retained FAC reference/history pair at the finite-projective criterion. There are no earlier canonical correction operations in this interval. The FAC passage is retained without a new audit of its external historical assertions.

The next source section at 19375--19414 was read only as lookahead: it contains the complete first surjectivity-locus lemma and the beginning of the next statement. It has not been adjudicated in this batch.

Three bounded canonical-corpus queries were executed: literal terms ``finite projective'', that exact phrase, and the exact phrase ``projective module'' together with ``nilpotent''. The first two routed to unrelated geometry or local work; the last supplied only PDF-primary records. Those hits were not read or adopted, and no PDF fallback was used. The exact supplied Stacks TeX and the explicitly located prior source lemmas provide the proof dependencies here. These standard consequences are not asserted to be novel, and the bounded queries are not an exhaustive literature survey. The broader synthesis remains after the core work.

\hypertarget{algebra_open_loci_pure_descent_notes_20260927}{}
\subsection{Editorial evidence for Algebra 19375--20373}\label{algebra-proof-010-editorial-evidence-for-algebra-1937520373}

Primary source: Stacks Project authors, algebra.tex, commit a04446e57ec1fbc252a871afcec7752fb2807b14, SHA-256 FA8BB92E58A4F78A2BD01B3B6A4A87DE0A0D279F5DD90641B574DD5FBFFFA4F3. The complete interval and seven received reports were read. Substantive corrections and the consequences below remain source-linked editorial material. The translated source is not rewritten around them.

\subsubsection{The original isomorphism neighbourhood}\label{algebra-proof-010-the-original-isomorphism-neighbourhood}

Sources: 19408--19444; OCC-00445 and OCC-00446. Keep the presentation \(N=R^n/\sum_jRk_j\), its generators \(e_i\), and the actual map \(\varphi:M\to N\).

Surjectivity of \(\varphi_{\mathfrak p}\) supplies \(n_i/s_i\in M_{\mathfrak p}\), with \(n_i\in M\), \(s_i\notin\mathfrak p\), whose image is \(e_i/1\). Equality of fractions means there is \(t_i\notin\mathfrak p\) with \[
t_i\bigl(\varphi(n_i)-s_ie_i\bigr)=0\quad\text{in }N.
\] Put \(m_i=t_in_i\in M\) and \(f_i=t_is_i\notin\mathfrak p\). Then \(\varphi(m_i)=f_ie_i\) holds in the original module \(N\), not merely in its localization. For the original product \(f=\prod_i f_i\), every \(m_i/f_i\) now belongs to \(M_f\). The resulting \(\psi:R_f^n\to M_f\) has \(\varphi_f\psi\) equal to the original presentation map.

The element \(\psi(k_j)\) is in \(M_f\). Its image in \(M_{\mathfrak p}\) is zero because \(\varphi_{\mathfrak p}\) is injective. The map \(M_f\to M_{\mathfrak p}\) is localization at the image of \(R\setminus\mathfrak p\); therefore some original \(g_j\notin\mathfrak p\) annihilates \(\psi(k_j)\) already in \(M_f\). With \(g=\prod_jg_j\), the map factors to \(\chi:N_{fg}\to M_{fg}\) and \(\varphi_{fg}\chi=\operatorname{id}\).

The finite target \(M_{fg}\) permits the source surjectivity-locus lemma to be applied to \(\chi\). It gives a principal neighbourhood of the selected prime where \(\chi\) is onto. A fraction in \(R_{fg}\) determining that neighbourhood has an original numerator \(h\notin\mathfrak p\); inverting it in \(R_{fg}\) is canonically localization at \(fgh\). A surjective right inverse is also a left inverse: for \(x=\chi(y)\), one has \(\chi\varphi(x)=\chi\varphi\chi(y)=\chi(y)=x\). Thus the original \(\varphi_{fgh}\) is an isomorphism. Empty products when \(n=0\) or the list of relations is empty are one; the same argument includes the zero module. These are denominator and codomain repairs, not changed finiteness hypotheses.

\subsubsection{Split kernels and the original minors}\label{algebra-proof-010-split-kernels-and-the-original-minors}

Source: 19471--19532; OCC-00447 and OCC-12087. The source correction is already composed as MC-STK-ERR-0430, changing only the ambient summand from \(P_{2,f}\) to \(P_{1,f}\).

For the original split surjection \(u:P_{1,f}\to P_{2,f}\), choose its section \(s\). The mutually inverse maps are \[
\ker u\oplus P_{2,f}\longrightarrow P_{1,f},\quad(k,y)\mapsto k+s(y),
\qquad
x\longmapsto(x-su(x),u(x)).
\] The first coordinate of the inverse is in \(\ker u\), and both composites are identities because \(us=1\). Thus the kernel is a finite-projective summand of the domain. It need not be a summand of the codomain: the surjection \(R\to0\) over nonzero \(R\) already disproves that claim.

The injectivity-locus argument preserves the original matrix. On a chosen chart, after listing the selected \(n_1\) rows first, write it as \(\binom A B\), where \(A\) is the actual invertible \(n_1\times n_1\) minor and \(B\) contains all remaining rows. The map \((v,w)\mapsto w-BA^{-1}v\) has kernel exactly the original image of \(x\mapsto(Ax,Bx)\). It induces an isomorphism of the cokernel with \(R^{n_2-n_1}\), whose inverse sends \(w\) to the class of \((0,w)\). The left inverse on the image is \(A^{-1}\) applied to the selected coordinates. This proves both injectivity and the source cokernel assertion without replacing the matrix. When \(n_1=0\), the empty minor is one and the source module is zero; when \(n_1>n_2\), the fibre-injectivity locus is empty. On a principal open contained in that locus the source local assertions patch exactly as printed.

\subsubsection{Finite relations and universal exactness}\label{algebra-proof-010-finite-relations-and-universal-exactness}

Sources: 19583--19993. The supplied proofs of Lazard's theorem and the universally exact criteria were read through their ends; no source correction is required in those arguments.

The equational criterion writes \(x_i=\sum_j a_{ij}y_j\) with \(\sum_i f_i a_{ij}=0\). In the source map \(h:R^n\to R^m\), this means exactly \(h(e_i)=\sum_j a_{ij}e'_j\): the matrix entry in row \(j\), column \(i\), is \(a_{ij}\). Its composite with \(g(e'_j)=y_j\) recovers the original map, and it kills the original relation vector \((f_i)\). Successive compositions that kill one more relation retain every relation already killed. Thus a finite relation submodule can be killed after finitely many steps. For a finite presentation \(P=R^n/\operatorname{im}A\), applying this to the finitely many columns of \(A\) gives exactly a factorization of any map \(P\to M\) through a finite free module when \(M\) is flat. Conversely the presentation \(R^n/Rx\) for a single relation recovers the equational criterion from this factorization property.

In the source directed system \(E\), a finite vector in the free module has finite support. Two pairs \((J,N)\) have a common upper bound using the union of their finite supports and sum of their finite relation modules. An element mapping to zero in \(M\) is killed at the larger pair obtained by adjoining its one kernel vector to the relation module. This proves the displayed colimit map is both onto and injective. In the flat case, the new basis labels chosen from \(M\times\mathbf Z\) lie outside the original finite \(J\). The map onto the actual finite free \(F\) splits on these new basis vectors, so its kernel is the image of a projector on a finite free module and is finite. The original relations lie in that kernel, and the diagram to \(M\) commutes. Therefore the constructed free stages are cofinal. The empty basis and zero module cases obey the same construction.

For the universal-exactness theorem, the original cokernel \(Q\) of the matrix \((a_{ij})\) turns its finite system of equations into the class of \((x_i)\) in \(M_1\otimes Q\). Injectivity into \(M_2\otimes Q\) makes that class zero, supplying the \(z_j\) in the same original module \(M_1\). For a finite presentation \(R^n\xrightarrow{g_1}R^m\xrightarrow{g_2}P\), the printed subtraction \(h'_2=h_2-f_1k_1\) kills the relations since \(k_1g_1=h_1\); then \(f_2h'_2=\varphi g_2\) proves its descended map is a lift. The subsequent fibre product \(M_{2,i}=M_2\times_{M_3}M_{3,i}\) has the exact kernel identification \(x\mapsto(f_1(x),0)\). A lift \(M_{3,i}\to M_2\) supplies a section into this fibre product. These sections need not be compatible with the transition maps: the exact sequences themselves form the required system, and each is split. Exact filtered colimits identify its first and third terms with the original ones and then its middle term. No compatibility of arbitrary chosen splittings is assumed. These observations verify the full source implication chain and its later splitting and flatness applications.

\subsubsection{The nonzero boundary in the examples}\label{algebra-proof-010-the-nonzero-boundary-in-the-examples}

Source: 19995--20038. The statement ``any torsion module'' includes the zero module, which is flat. To obtain the promised nonflat example over \(\mathbf Z\), it must read ``any nonzero torsion module''. This finding was detected in the source review independently of the received reports.

The first example is correct. The groups \(A=\bigoplus_{n\geq1}\mathbf Z\) and \(B=\prod_{n\geq1}\mathbf Z\) are torsion-free. Their quotient \(C\) is torsion-free too: if \(d b\) has finite support for nonzero integer \(d\), then \(b\) has finite support, since each integer coordinate has no nonzero torsion. For completeness, a torsion-free abelian group \(T\) is flat by the original equational criterion. For a relation \(a(x_1,\ldots,x_n)^t=0\), apply successive two-coordinate integral changes of basis to the row \(a\). For a pair \((u,v)\) of integers with positive gcd \(d\), Bezout coefficients \(r,s\) give \[
(u,v)\begin{pmatrix}r&-v/d\\s&u/d\end{pmatrix}=(d,0),
\qquad \det\begin{pmatrix}r&-v/d\\s&u/d\end{pmatrix}=1.
\] Iterating gives \(aU=(d,0,\ldots,0)\) with integral invertible \(U\). Write \(x=Uy\). Torsion-freeness gives \(y_1=0\), so the remaining columns of \(U\) express the original \(x_i\) as combinations with every original coefficient relation annihilated. If \(a=0\), use the identity matrix instead. This is exactly the equational criterion, including an empty relation. Thus the three source groups are flat.

The class \(x\in C\) of the original sequence \((2,2^2,2^3,\ldots)\) is nonzero because this sequence has infinite support. For each \(r\geq1\), truncate its first finitely many coordinates and divide the remaining coordinates by \(2^r\); this gives an element of \(B\) whose class multiplies to \(x\). Any homomorphism \(C\to B\) consequently sends \(x\) to an integer sequence each of whose coordinates is divisible by every \(2^r\), hence to zero. A section cannot do this to a nonzero quotient class. The original extension is therefore nonsplit.

For a nonzero torsion module \(T\), choose a nonzero element killed by a nonzero integer \(d\). Tensoring \(\mathbf Z\xrightarrow{d}\mathbf Z\) with \(T\) fails to preserve injectivity, proving nonflatness. The middle term \(T\oplus T\) is likewise nonflat, as witnessed in either coordinate. The split sequence is universally exact. Taking its direct sum with the first example preserves universal exactness and gives three nonflat terms. A section of this new quotient would, after inclusion of the original quotient and projection to its original middle term, give a section of the nonsplit original sequence. Thus the combined example is nonsplit as claimed. The extra word ``nonzero'' is the only required correction.

\subsubsection{Localization and the rejected proof-gap report}\label{algebra-proof-010-localization-and-the-rejected-proof-gap-report}

Source: 20160--20194; OCC-00448.

Retain the original map \(A\to B\), multiplicative sets \(S\subset A\), \(S'\subset B\), and condition \(\varphi(S)\subset S'\). For every \(A\)-module \(Q\), the map \[
(S')^{-1}(Q\otimes_AM)\longrightarrow Q\otimes_A(S')^{-1}M,
\quad (q\otimes m)/u\longmapsto q\otimes(m/u)
\] is an isomorphism: the inverse is induced by the balanced map \((q,m/u)\mapsto(q\otimes m)/u\). Localization of the original injection after tensoring is therefore the first conclusion of (1), as the source proves.

Part (2) is precisely the remaining required comparison. If \(M,M'\) are \((S')^{-1}B\)-modules, they are \(S^{-1}A\)-modules. For an \(S^{-1}A\)-module \(Q\), balancing by inverses of elements of \(S\) gives \(Q\otimes_AM\cong Q\otimes_{S^{-1}A}M\). For an arbitrary \(A\)-module \(Q\), the mutually inverse fraction maps give \[
Q\otimes_AM\cong(S^{-1}Q)\otimes_{S^{-1}A}M,
\quad q\otimes m\mapsto(q/1)\otimes m,
\quad(q/s)\otimes m\mapsto q\otimes s^{-1}m.
\] These identities, with their counterparts for \(M'\), prove both implications of (2). Apply (2) to the already localized pair \((S')^{-1}M,(S')^{-1}M'\) from (1). Its established universal injectivity over \(A\) is equivalent to universal injectivity over \(S^{-1}A\). Thus the second conclusion of (1) already follows from the next part of the printed proof. There is no missing hypothesis or mathematical gap. An explicit cross-reference would be optional exposition; the received claim of a required correction is rejected, with zero source operations.

The same fraction identities verify the preceding stalk comparison. Testing the tensor kernel at all maximal ideals of the original \(S\) detects its vanishing, including when the test module is initially only over \(R\). Tensor associativity proves permanence under tensoring; composition and the first-factor implication follow by composing the actual injective maps after each tensor test. These source omissions are not turned into rewrite assignments.

\subsubsection{The finite descent statement and coefficients}\label{algebra-proof-010-the-finite-descent-statement-and-coefficients}

Source: 20306--20343; OCC-00449 and OCC-00450.

The fourth enumerated item, ``add more here as needed'', is an authoring instruction with no theorem or proof content. The proposal removes that item and makes the third item end with a period. In a receiving translation its original occurrence must remain identifiable in editorial documentation; this is not a new mathematical descent property attributed to the author.

For the original finitely many \(S\)-generators \(y_j\) of \(M\otimes_RS\), write each as a finite sum of simple tensors. Take the union of all the original module entries of those sums, label it \(x_1,\ldots,x_n\), and sum coefficients for any repeated entry. Setting missing coefficients to zero gives exactly \[
y_j=\sum_{i=1}^n x_i\otimes f_{ij},\qquad f_{ij}\in S.
\] The second index is needed to describe this whole family; a single fixed list \(f_i\) need not produce every \(y_j\). The original map \(R^n\to M\) now becomes visibly surjective after base change. The source faithfully flat proof remains valid. Its kernel-base-change step uses flatness of \(S\); the broader argument below proves the stronger case separately instead of silently deleting that hypothesis from this step.

\subsubsection{Surjectivity needs only a finite cokernel}\label{algebra-proof-010-surjectivity-needs-only-a-finite-cokernel}

This is a separate consequence of the first source lemma, 19383--19406. For any original map \(u:M\to N\), suppose only that its actual cokernel \(C\) is finite. Then \[
U=\{\mathfrak p:u_{\mathfrak p}\text{ is onto}\}
=\{\mathfrak p:u\otimes_R\kappa(\mathfrak p)\text{ is onto}\}
=\operatorname{Spec}R\setminus V(\operatorname{Ann}_RC).
\] Indeed right exactness gives the original cokernels \(C_{\mathfrak p}\) and \(C\otimes_R\kappa(\mathfrak p)\); finite Nakayama makes one zero exactly when the other is zero. For finitely many generators \(c_i\) of \(C\), vanishing at \(\mathfrak p\) gives annihilators \(s_i\notin\mathfrak p\), whose product annihilates all \(C\). Thus \(C_{\mathfrak p}=0\) exactly when \(\operatorname{Ann}_RC\) contains an element outside \(\mathfrak p\), giving the stated open set. If \(D(f)\subset U\), the localized cokernel has zero localizations everywhere and hence is zero, proving the original localized surjection.

This open locus commutes with every ring base change \(R\to T\). For a prime \(\mathfrak q\subset T\) over \(\mathfrak p\subset R\), the actual cokernel fibre is \[
(C\otimes_RT)\otimes_T\kappa(\mathfrak q)
\cong (C\otimes_R\kappa(\mathfrak p))\otimes_{\kappa(\mathfrak p)}\kappa(\mathfrak q).
\] The isomorphism sends each original simple tensor to the same element after the residue-field maps. A field extension kills no nonzero vector space, so the two fibres vanish simultaneously. This proves the exact inverse-image assertion for \(U\). It does not assert equality of the two annihilator ideals. If \(C=0\), its annihilator is \(R\) and the open is all of the spectrum; zero rings and empty spectra are included. The source hypothesis that \(N\) is finite is sufficient but stronger than needed, and remains in its original statement.

\subsubsection{Descent along the original universally injective ring map}\label{algebra-proof-010-descent-along-the-original-universally-injective-ring-map}

The source definition calls \(R\to S\) universally injective when every canonical map \[
\eta_X:X\to X\otimes_RS,\qquad x\mapsto x\otimes1
\] is injective. Such a ring map is also called pure. Under this hypothesis, finite type, finite presentation, flatness, and finite projectivity of \(M\otimes_RS\) over \(S\) imply the corresponding property of the original \(M\) over \(R\). The converses hold by base change. Here is a full proof, which never assumes \(S\) is flat over \(R\).

First, the source reflection lemma has a direct proof from the actual unit maps. If \(u\otimes_RS\) is injective and \(u(x)=0\), then \(u_S(x\otimes1)=0\), hence \(x\otimes1=0\), hence \(x=0\). If \(u_S\) is onto, right exactness gives \((\operatorname{coker}u)\otimes_RS=0\), and injectivity of its unit map gives \(\operatorname{coker}u=0\). This proves reflection of injections, surjections and isomorphisms.

For finite type, use the original coefficient construction in the preceding section. The resulting \(u:R^n\to M\) is onto after tensoring with \(S\), hence onto by reflection. No flatness is required.

For finite presentation, first descend finite type and choose the resulting surjection \(p:F=R^n\to M\), with actual kernel \(K\). Right exactness alone identifies \[
L:=\ker(F\otimes_RS\to M\otimes_RS)
=\operatorname{im}(K\otimes_RS\to F\otimes_RS).
\] The hypothesis that \(M\otimes_RS\) is finitely presented makes \(L\) finite. This uses the elementary finite-presentation fact for a finite free surjection: comparing it with a fixed finite presentation and lifting the two finite generating families shows that its kernel is generated by the images of the finitely many original relations together with the finitely many differences of the composite lifts from the identity. More explicitly, if \(S^a\xrightarrow{A}S^b\xrightarrow{q}M_S\to0\) is a finite presentation, choose \(v:S^b\to F_S\), \(w:F_S\to S^b\) with \(p_Sv=q\), \(qw=p_S\). Then \(vA(S^a)\) and \((1-vw)(F_S)\) lie in \(L\), and for \(z\in L\), write \(wz=A t\) to obtain \(z=vAt+(1-vw)z\). Their finite sets of columns therefore generate \(L\).

Lift those finitely many generators of \(L\) to finite sums in \(K\otimes_RS\). Let \(K_0\subset K\) be generated by all original \(K\)-entries appearing in these sums. Then \(\operatorname{im}(K_0\otimes_RS\to F_S)=L\). Set \(C=F/K_0\). The induced surjection \(C\to M\) becomes an isomorphism after base change, since both resulting quotients of \(F_S\) are \(F_S/L\). Reflection makes \(C\to M\) an isomorphism. Hence \(K=K_0\) is finite and the original \(M\) is finitely presented. At no point was \(K\otimes_RS\to F_S\) assumed injective.

For flatness let \(M_S=M\otimes_RS\) be flat over \(S\). We prove the source finite-presentation factorization criterion for \(M\). Fix an arbitrary finite presentation \[
R^t\xrightarrow{A}R^n\xrightarrow{p}P\to0
\] and an arbitrary map \(f:P\to M\). Over \(S\), the criterion factors \(f_S\) through \(S^r\): write \(f_S=g_Sh_S\). Express the finitely many images of the basis of \(S^r\) under \(g_S\) as finite sums of original elements of \(M\) tensored with coefficients in \(S\). Collect these module elements as \(m_1,\ldots,m_q\), and define the actual \(R\)-linear map \[
G:R^q\to M,\qquad e_i\mapsto m_i.
\] There is a coefficient matrix \(B:S^r\to S^q\) with \(G_SB=g_S\). The map \(Bh_S:P_S\to S^q\), composed with \(p_S\), is a matrix \(H_S\in\operatorname{Mat}(q\times n,S)\) satisfying \[
H_SA=0,\qquad G_SH_S=(fp)_S.
\]

Retain all original coefficients of \(A,G,fp\). Form the actual \(R\)-linear map \[
\beta:\operatorname{Mat}(q\times n,R)\longrightarrow
M^n\oplus\operatorname{Mat}(q\times t,R),
\quad H\longmapsto(GH,HA),
\] and let \(W=\operatorname{coker}\beta\). The element \(c=(fp,0)\) of the target has image zero in \(W\otimes_RS\), because \(H_S\) is its preimage under \(\beta_S\); cokernels commute with tensoring by right exactness. Purity says precisely that \(W\to W\otimes_RS\) is injective. Thus the class of \(c\) was already zero in \(W\), so there is an actual \(H\in\operatorname{Mat}(q\times n,R)\) with \[
HA=0,\qquad GH=fp.
\] The first equation makes \(H\) descend to \(h:P\to R^q\); the second gives \(Gh=f\), because \(p\) is onto. This is the required factorization through a finite free module. The source criterion therefore proves \(M\) flat. Empty lists and zero matrix dimensions yield the same maps between zero modules; no nonzero-rank assumption is used.

If \(M_S\) is finite projective, its finite presentation and flatness descend by the two constructions just proved. Apply the proved factorization criterion to \(P=M\) and \(f=\operatorname{id}_M\). It gives actual maps \(h:M\to R^q\), \(G:R^q\to M\) with \(Gh=\operatorname{id}_M\). Thus the original \(M\) is a finite-projective summand. Explicitly, let \(\lambda_i\) be the coordinates of \(h\) and \(m_i=G(e_i)\). Then \(m=\sum_i\lambda_i(m)m_i\), giving the identity tensor of the earlier converse evaluation criterion, ALGEBRA-RECON-492. This propagates the earlier characterization through the descended maps. The result generalizes the source faithfully flat descent proposition. Faithfully flat maps are pure by the preceding source lemma: after tensoring, the map \(x\otimes s\mapsto x\otimes1\otimes s\) has the multiplication retraction, and faithful flatness reflects its injection. Thus all original faithfully flat statements are receiving cases of the proved result; their author hypotheses are preserved in the translation.

The new hypothesis really allows ring maps excluded by flatness. Take the original diagonal map \[
\mathbf Z\longrightarrow\mathbf Z\times\mathbf Z/2\mathbf Z,
\qquad n\longmapsto(n,\overline n).
\] Projection to the first coordinate is a ring retraction. For every \(X\), the map \(X\otimes_RS\to X\), \(x\otimes s\mapsto r(s)x\), retracts \(\eta_X\), proving purity. But multiplication by two on \(S\) kills its nonzero element \((0,\overline1)\), so this \(S\) is not flat over \(\mathbf Z\). This exact example verifies the additional generality. It does not claim effective descent of arbitrary module objects or arbitrary projectivity; those larger assertions have not been proved here.

\subsubsection{Source reading and preservation limits}\label{algebra-proof-010-source-reading-and-preservation-limits}

The current interval differs from the authority by exactly MC-STK-ERR-0430, with no other correction or editorial environment in this interval. The earlier correction is retained under that identity. The source's historical discussion of Raynaud--Gruson, Gruson and Kaplansky is retained; reading this discussion does not constitute reading or certifying those external works or the later infinite-projectivity proofs.

Two bounded corpus queries were executed, the exact phrase ``pure morphisms'' and the literal terms ``pure descent''. PDF-primary and unrelated local hits remain unread routing records. Neil Epstein and Hop D. Nguyen's original TeX, \emph{Algebra retracts and Stanley-Reisner rings}, canonical source PUBUNIT-E140D1A58DCB9288442C1AD5, was read at lines 246--266 and 368--383, with the relevant section structure checked. Its retraction and pure-subring definitions agree exactly with the displayed original unit map; the retraction tensor formula above proves this comparison and supplies the nonflat example. The local retraction comparison is given by the actual maps \(R_{\mathfrak p}\to S_{r^{-1}\mathfrak p}\to R_{\mathfrak p}\), whose composition is identity by the universal localization maps. No later theorem about normality, regularity or complete intersections is adopted.

The author source declares Latin-1. The first attempted whole-file UTF-8 read failed without modifying the file. The actual bounded passages were then read as their exact ASCII byte lines; the original source and archive were left intact. File SHA-256: 83131E01670F9EF5AB2E6480F668267364B9B1729CB612C9EAE469FFBF12FED1.

Lookahead 20374--20390 is only the introduction and first part of the transfinite-devissage definition. The definition is incomplete and unadjudicated. No chapter mutation, rendering, publication, exhaustive literature survey or novelty claim is made. Combined-results synthesis remains after the core work.

\hypertarget{algebra_transfinite_devissage_notes_20260927}{}
\subsection{Transfinite dévissage: source indices and the original stopping bound}\label{algebra-proof-011-transfinite-duxe9vissage-source-indices-and-the-original-stopping-bound}

This is separate editorial evidence for the Stacks Project authors' \texttt{algebra.tex}, authority commit \texttt{a04446e57ec1fbc252a871afcec7752fb2807b14}, lines 20374--20588. The exact authority SHA-256 is \texttt{FA8BB92E58A4F78A2BD01B3B6A4A87DE0A0D279F5DD90641B574DD5FBFFFA4F3}. The source follows Kaplansky; that attribution is retained. The cited Kaplansky publication has not been independently read in this batch. The actual source and all due received reports were read. Later passages through 21480 were read as lookahead and are not adjudicated by this note.

The complete current interval equals this authority interval. Two minimal index corrections are proposed; neither the current chapter nor a translation has been changed. In a receiving translation, preserve the original reading and visibly identify the source corrections. The supplementary proofs and stopping bound below belong in linked conceptual errata or a supplement. There is no claim of mathematical novelty or of completion of the later combined-results synthesis.

\subsubsection{Sections exist only at successor stages in the index}\label{algebra-proof-011-sections-exist-only-at-successor-stages-in-the-index}

Source: \texttt{definition-devissage}, lines 20384--20404, and \texttt{lemma-direct-sum-devissage}, lines 20409--20453. Its index is an ordinal \(S\), with \(M_0=0\), increasing submodules, union \(M\), continuity at limits, and a splitting of \(M_\alpha\subset M_{\alpha+1}\) when \(\alpha+1\in S\). The definition's occurrence of \(M_0\) requires that its indexing family include zero.

Line 20418 says that the section of \(M_{\alpha+1}/M_\alpha\) is chosen for each \(\alpha\in S\). If \(S\) has a last member, the next module is not part of the family. The precise correction is to quantify over \(\alpha+1\in S\), as the definition and displayed direct sum already do. No change to the lemma's statement is needed.

For each such successor choose a retraction \(r_\alpha:M_{\alpha+1}\to M_\alpha\) of the inclusion and define \[
s_\alpha:M_{\alpha+1}/M_\alpha\longrightarrow M_{\alpha+1},
\qquad [x]\longmapsto x-r_\alpha(x).
\] This is independent of the representative because \(r_\alpha(a)=a\) for \(a\in M_\alpha\). The quotient map composed with \(s_\alpha\) is the identity, and every \(x\in M_{\alpha+1}\) has the unique expression \(r_\alpha(x)+s_\alpha([x])\). Compose each section with the original inclusion in \(M\). Their sum is exactly the source map \[
f:\bigoplus_{\alpha+1\in S}M_{\alpha+1}/M_\alpha\longrightarrow M.
\] For \(\beta\in S\), its restriction \(f_\beta\) to terms with \(\alpha+1\leq\beta\) is an isomorphism onto the original submodule \(M_\beta\). At zero this is the unique map between zero modules. At a successor it follows from the unique decomposition just proved and the induction hypothesis for \(f_\alpha\). At a limit, any element of the domain has finite support, hence lies in the domain of some earlier \(f_\gamma\); the same assertion proves injectivity. Any element of \(M_\beta\) lies in an earlier \(M_\gamma\) by continuity and therefore has a preimage. Finally, the union property gives surjectivity of \(f\), and finite support gives injectivity by choosing a stage containing every nonzero component. This includes successor and limit values of \(S\).

The source's later remark, lines 20485--20486, that (3) and (5) are equivalent under (0)--(2) is also justified by these actual maps. Under (3), split the displayed direct sum into terms \(\alpha+1\leq\beta\) and \(\beta<\alpha+1\), and transport its projection through \(f\). This is a retraction \(M\to M_\beta\). Conversely, a retraction \(M\to M_\alpha\) restricts to a retraction \(M_{\alpha+1}\to M_\alpha\). This proves the equivalence without requiring the quotients to be countably generated.

\subsubsection{Include the terminal partial sum}\label{algebra-proof-011-include-the-terminal-partial-sum}

Source: \texttt{lemma-Kaplansky-devissage}, lines 20455--20467; received report \texttt{OCC-00451}. Keep the original direct sum \(M=\bigoplus_{i\in I}N_i\), with each \(N_i\) countably generated, and the chosen well-order on \(I\). Identify that ordered set with its ordinal, as the source does. Set \[
S=I+1,\qquad M_\alpha=\bigoplus_{j<\alpha}N_j
\quad(\alpha\in S,\text{ equivalently }\alpha\leq I).
\] All these sums are actual coordinate submodules of the original \(M\). We have \(M_0=0\) and \(M_I=M\), so the union is \(M\). For a limit \(\lambda\leq I\), an element of \(M_\lambda\) has finite coordinate support below \(\lambda\). There is \(\beta<\lambda\) larger than all those indices, so that element lies in \(M_\beta\). The reverse inclusion follows from increasing coordinate sets. At a successor \(\alpha+1\leq I\), \[
M_{\alpha+1}=M_\alpha\oplus N_\alpha.
\] The coordinate projection is the required retraction; \(N_\alpha\to M_{\alpha+1}/M_\alpha\), sending an element to its coset, is an isomorphism with inverse induced by the coordinate projection to \(N_\alpha\). Thus the quotient is countably generated and all five defining properties hold.

This verifies every ordinal boundary. If \(I=0\), then \(M=0\), \(S=1\), and the family consists of \(M_0=0\); there are no successor quotients. If \(I=\delta+1\), the original index omitted \(M_I\) and its union was \(\bigoplus_{j<\delta}N_j\); it fails when \(N_\delta\ne0\). For example, \(I=1\) and \(N_0=\mathbb Z\) give only the zero stage in the source's uncorrected family. For a nonzero limit \(I\), the original union already exhausts the direct sum, but the uniform terminal-stage construction is still valid.

The reverse implication of the source lemma follows from the preceding exact isomorphism \(f\): its summands are the original countably generated successor quotients. The correction does not strengthen that statement or replace its proof by an unrelated construction.

\subsubsection{The original projector closure and a bounded stopping stage}\label{algebra-proof-011-the-original-projector-closure-and-a-bounded-stopping-stage}

Source: \texttt{theorem-kaplansky-direct-sum}, lines 20469--20566, followed by \texttt{theorem-projective-direct-sum}, lines 20572--20585. Preserve \(M=P\oplus Q=\bigoplus_{i\in I}N_i\). Write \(p:M\to P\) and \(q:M\to Q\) for the original projections, viewed also as endomorphisms of \(M\). They satisfy \(p+q=1\), \(p^2=p\), \(q^2=q\), and \(pq=qp=0\).

The source constructs at each nonterminal successor a countable matrix \((x_{mn})\), first listing generators of the least missing \(N_j\), then adding generators of each coordinate summand occurring in the finite supports of \(p(x_{mn})\) and \(q(x_{mn})\). Its specified order is valid: the entry \((m,n)\) has position \[
t(m,n)=\frac{(m+n-2)(m+n-1)}2+m
\] in the successive antidiagonals. The first row is present initially, and processing position \(t\) creates row \(t+1\). Since \(t(m,n)\geq m\), row \(m\) has already been created when that entry is processed. Every pair occurs at a finite position. Each new row can be chosen countable because the two supports are finite and each involved \(N_i\) has a countable generating set. Pad finite or empty lists with zero entries. This is precisely the source's order \((1,1),(1,2),(2,1),(1,3),\ldots\).

Let \(B\) be the submodule generated by all entries, and let \(D\) consist of \(j\) and every coordinate appearing in the supports of the two projections of any entry. The construction includes a generating set for each \(N_i\), \(i\in D\), hence \(\bigoplus_{i\in D}N_i\subset B\). Conversely, \(x=p(x)+q(x)\) for each entry, so each entry has support in \(D\). Thus \[
B=\bigoplus_{i\in D}N_i,
\qquad p(B)\subset B,\qquad q(B)\subset B.
\] The inclusions follow first on the generating entries and then by linearity. Consequently, if \(M_\alpha\) is a coordinate sum stable under both projections, so is \(M_{\alpha+1}=M_\alpha+B\). Its quotient by \(M_\alpha\) is countably generated by the images of the entries. Coordinate projections give a retraction from \(M\) onto every such coordinate sum. At limits, a union of increasing coordinate sums is the coordinate sum of the union of their coordinate sets: every vector has finite support. Stability under \(p,q\) is preserved by that union. Zero coordinate modules cause no exception, since they are already contained in every stage.

There is a useful precise bound on the source's phrase ``some large enough ordinal.'' Write the chosen well-order as the ordinal \(I\). The same recursion need only be performed for \(\alpha\leq I\), and it satisfies \[
\bigoplus_{i<\alpha}N_i\subset M_\alpha
\qquad(\alpha\leq I).
\] At zero the assertion holds. At a successor, all earlier \(N_i\) are already contained by induction. If \(N_\alpha\) is not yet contained, it is therefore the least missing summand, and the actual next step includes it. If it is contained, the increasing family retains it. At a limit the claim follows from the union rule, since every \(i<\alpha\) occurs by stage \(i+1<\alpha\). This proves the formula by transfinite induction and in particular \(M_I=M\). Thus \textbf{the original construction works with \(S=I+1\)}, including \(I=0\). No enlargement of its original modules or replacement of its projector maps is involved. This bound is separate editorial information; the source's unquantified existence assertion is correct and receives no source edit.

For completeness of the propagation check, put \(P_\alpha=P\cap M_\alpha\) and \(Q_\alpha=Q\cap M_\alpha\). Stability gives \(M_\alpha=P_\alpha\oplus Q_\alpha\), with actual projections the restrictions of \(p,q\). The map \[
M_{\alpha+1}/M_\alpha\longrightarrow
(P_{\alpha+1}/P_\alpha)\oplus(Q_{\alpha+1}/Q_\alpha),
\quad [x]\longmapsto([p(x)],[q(x)])
\] is well-defined. Its inverse sends \(([a],[b])\) to \([a+b]\); changing representatives alters that sum by an element of \(M_\alpha\). Both composites are the identity. Projection onto the first factor shows the original \(P\)-quotient is countably generated. If \(\rho_\alpha:M\to M_\alpha\) is a coordinate retraction, then \(p\rho_\alpha:M\to P_\alpha\) fixes \(P_\alpha\); its restriction to \(P_{\alpha+1}\) proves the successor inclusion splits. Intersections with \(P\) commute with these increasing unions, so the zero, union, and limit properties also hold. The corrected direct-sum lemma therefore proves the original theorem for \(P\).

If \(P\) is projective, take its original realization as a summand of a free module. The resulting countably generated pieces are direct summands of \(P\), hence projective. Explicitly, their inclusion and projection compose to the identity; for any surjection and any map from one piece, compose with that projection, lift from \(P\) by projectivity, then restrict along the inclusion. This proves the lifting property for that piece. The result and its projective adjective are unchanged.

\subsubsection{Propagation and coverage}\label{algebra-proof-011-propagation-and-coverage}

Groups 502 and 503 record the two source-index repairs. Group 504 has zero source operations and records the exact stopping bound for the original projector construction. Its proof receives both repaired indexing conventions. The direct-sum equivalence, summand theorem, and projective corollary in the current interval remain valid, with their original objects and conclusions. Later local-ring and Mittag-Leffler statements remain pending adjudication even where their source text has been read. Do not infer full downstream review from the present propagation check.

The bounded disk-literature query \texttt{Kaplansky\ Mittag\ Leffler} returned two PDF-routing records, for \emph{Geometric Methods in Physics XL} and \emph{Constructing Projective Modules}. Neither was read or used as proof evidence; no PDF fallback was performed. The supplied original Stacks TeX is the primary treatment used here. This is not an exhaustive literature survey, and the explicit stopping bound is not described as a new result.

\hypertarget{algebra_ml_review_notes_20260927}{}
\subsection{Local projectives, Mittag-Leffler indices and tensor products}\label{algebra-proof-012-local-projectives-mittag-leffler-indices-and-tensor-products}

Separate editorial evidence for the Stacks Project authors' original \texttt{algebra.tex}, authority commit \texttt{a04446e57ec1fbc252a871afcec7752fb2807b14}, lines 20589--21441. Authority SHA-256: \texttt{FA8BB92E58A4F78A2BD01B3B6A4A87DE0A0D279F5DD90641B574DD5FBFFFA4F3}. All units in this interval and its five received reports were read. The current interval differs by exactly the two already composed operations of \texttt{MC-STK-ERR-0440}; no other change or FAC environment occurs here. The following proposition, beginning at 21442, remains outside this adjudication.

Four minimal operations are proposed: two stage-choice repairs, one adverb, and one coefficient projection. The existing tensor correction retains its identity. Complete supporting arguments and the two further consequences below are editorial material, with the source's statements, hypotheses and exposition kept identifiable. No chapter or translation was mutated; a replay is not an instruction to replace a translation by these notes. No novelty or exhaustive literature claim is made.

\subsubsection{The original local-projective proof}\label{algebra-proof-012-the-original-local-projective-proof}

Source 20600--20693 receives the already reviewed decomposition into countably generated projective summands. Its proof is correct. Here are the exact matrix and summand checks; they generate no source operation.

Retain the original free module \(F=P\oplus Q\), basis, and shortest expression \(x=\sum_{i=1}^n a_i e_i\). Such a shortest expression exists because its possible finite lengths are natural numbers. If \(a_j=\sum_{i\ne j}a_i c_i\), the basis change \(e_i\mapsto e_i+c_i e_j\) for \(i\ne j\), fixing \(e_j\) and the remaining basis, has inverse obtained by subtracting these same multiples. In that basis the coefficient at \(e_j\) is zero. Thus no \(a_j\) belongs to the ideal generated by the other coefficients.

Write the original projections as \(e_i=y_i+z_i\), and \(y_i=\sum_j b_{ij}e_j+t_i\). The equality \(x=\sum_i a_i y_i\) gives, for each \(j\), \[
0=\sum_{i=1}^n a_i(b_{ij}-\delta_{ij}).
\] If any coefficient \(b_{ij}-\delta_{ij}\) in this relation were a unit, it would express the corresponding \(a_i\) in terms of the others. Every such coefficient lies in the maximal ideal \(\mathfrak m\). Consequently \(B=(b_{ij})\) reduces to the identity modulo \(\mathfrak m\), and \(\det B\) is a unit. The adjugate formula gives its actual inverse over \(R\).

In the original decomposition \(F=R^n\oplus F_{\rm rest}\), let \(T:R^n\to F_{\rm rest}\) have columns \(t_i\). The endomorphism sending \(e_i\) to \(y_i\) has the formula \[
(v,w)\longmapsto(B^{\mathsf t}v,Tv+w),
\qquad
(v,w)\longmapsto((B^{\mathsf t})^{-1}v,w-T(B^{\mathsf t})^{-1}v)
\] for its inverse. The transpose is required by the source's indexing of \(b_{ij}\). It follows that \(N=\sum_i Ry_i\) is a finite free summand of \(F\). Its retraction restricts to a retraction \(P\to N\), and \(x=\sum_i a_i y_i\in N\). If \(x=0\), the zero free summand suffices.

For a countably generated projective \(P\), every complement of a finite free summand is projective: restrict the projective lifting property using that complement's inclusion and retraction. At the next generator, the preceding construction gives a free summand of the complement containing its projected value. That value has finite support in a basis of this free summand; the span of those finitely many basis elements is again a direct summand, so it can be used as the next finite free piece. The partial sums are split and exhaust the chosen generators. The map from their external direct sum to \(P\) is injective because every finite partial sum is direct, and surjective because every generator occurs. Finally, the earlier countable-projective decomposition and the union of the resulting bases prove the original unrestricted local-ring theorem.

\subsubsection{Stable images and countable cofinality}\label{algebra-proof-012-stable-images-and-countable-cofinality}

Source 20697--20877. The following is a separate weakening of the index hypothesis: \textbf{a countable cofinal subset of the original directed set suffices} for both nonemptiness of the inverse limit of nonempty ML sets and exactness of the inverse limit of the stated short exact sequences. The source's countable-index statements are preserved.

For the original maps \(\varphi_{ji}:A_j\to A_i\), write \[
A'_i=\bigcap_{j\geq i}\varphi_{ji}(A_j).
\] If the images stabilize beyond \(j_0\geq i\), their stable value equals this intersection: for any other \(j\geq i\), choose \(k\geq j,j_0\). The stable value is the image from \(k\), hence is contained in the image from \(j\). This also handles incomparable indices in the definition of the intersection.

If \(j\geq i\) and \(a\in A'_j\), then for every \(k\geq i\) choose \(\ell\geq k,j\) and a preimage of \(a\) in \(A_\ell\). Its image in \(A_i\) factors through \(A_k\); thus \(\varphi_{ji}(a)\in A'_i\). These are the original restricted maps. They are surjective under the ML condition: choose \(k\geq j\) beyond stabilization thresholds for both \(i\) and \(j\). Any \(a_i\in A'_i\) has a preimage \(a_k\in A_k\); then \(\varphi_{kj}(a_k)\in A'_j\) maps to \(a_i\). If every \(A_i\) is nonempty, each stable image is nonempty. A compatible original family belongs coordinatewise to every image defining \(A'_i\), so inclusion induces an actual bijection \(\lim A'_i\to\lim A_i\), whose inverse retains the same coordinates.

Let \(C\subset I\) be a countable cofinal subset of a nonempty directed set. Enumerate it as \(c_1,c_2,\ldots\), repeating elements when finite, and choose \(j_1=c_1\), \(j_n\geq j_{n-1},c_n\). The sequence is cofinal in the original \(I\). Restriction of a compatible family gives the map to the limit over this sequence. Conversely, for a compatible sequence \((a_n\in A_{j_n})\), define \[
a_i=\varphi_{j_n i}(a_n)\quad\text{for any }n\text{ with }j_n\geq i.
\] Any two choices give the same value after passing to their larger sequence index. The resulting family is compatible: for \(i\leq k\), use a single \(j_n\geq k\) and the original composition rule. Both composites with restriction are identities. The same formulas are group or module homomorphisms in those categories and commute with morphisms of systems.

Apply these maps to the stable-image system. Starting with any element of \(A'_{j_1}\), successively choose a lift along each surjective map \(A'_{j_{n+1}}\to A'_{j_n}\), using the same axiom of choice as the source. The resulting compatible sequence gives a family over \(I\). This proves nonemptiness. If an empty index is allowed under a different convention, its set limit is a singleton and its group limit is zero, so the corresponding assertions hold separately without an enumeration.

For exactness retain the original short exact sequences and all their maps: \[
0\longrightarrow A_i\xrightarrow{f_i}B_i\xrightarrow{g_i}C_i\longrightarrow0.
\] Write \(\varphi^A,\varphi^B,\varphi^C\) for the three original transition families. Coordinatewise injectivity of \(f_i\) gives injectivity on limits. If a compatible \((b_i)\) maps to zero, its unique preimages \(a_i\) under \(f_i\) are compatible by commutation and injectivity, proving exactness in the middle.

Given \((c_i)\in\lim C_i\), put \(E_i=g_i^{-1}(c_i)\). Each is nonempty and the original \(B\)-maps restrict to them. Fix \(i\), and let \(j\) be a stabilization index for the images from \(A\). For \(k\geq j\), choose \(e_j\in E_j\), \(e'_k\in E_k\), and set \(e'_j=\varphi^B_{kj}(e'_k)\). There is a unique \(a_j\in A_j\) with \(f_j(a_j)=e_j-e'_j\). Stabilization gives \(a_k\in A_k\) with \(\varphi^A_{ki}(a_k)=\varphi^A_{ji}(a_j)\). Then \[
e_k=e'_k+f_k(a_k)\in E_k,
\qquad
\varphi^B_{ki}(e_k)=\varphi^B_{ji}(e_j).
\] The last equality follows by inserting the defining difference and the commutative square for \(f\). The reverse image inclusion follows by factorization through \(j\), so \(E\) is ML. The just-proved nonemptiness result gives \(\lim E_i\ne\varnothing\), proving surjectivity onto \((c_i)\). This proves exactness with a countable cofinal subset. For \(I=\mathbb N\), the criterion at 20865--20867 is exactly stabilization of these same images, so the source sequential lemma is an unchanged receiving case. No claim removes countable cofinality altogether.

\subsubsection{Domination and the original pushout}\label{algebra-proof-012-domination-and-the-original-pushout}

Source 20905--20996. The finite-test lemma is correct: for an arbitrary tensor-test module, express it as a filtered colimit of finitely presented modules. Any kernel element is represented at a stage and its image vanishes at a later stage. Apply the assumed kernel inclusion there and pass to the colimit. Conversely, the all-module condition includes those tests. This element argument preserves the original tensor maps and does not require a tensor functor to preserve arbitrary injections.

The original pushout is \[
N'=(M'\oplus N)/\{(g(x),-f(x)):x\in M\},
\quad f'(m')=[m',0],\quad g'(n)=[0,n].
\] The map \(\ker f\to\ker f'\) is \(x\mapsto g(x)\). It is onto because \([m',0]=0\) means \((m',0)=(g(x),-f(x))\) for some \(x\in\ker f\); its kernel is \(\ker f\cap\ker g\). After tensoring with any test module, the resulting pushout has the same quotient presentation for the tensor maps, by right exactness. Applying this exact kernel argument to that new pushout gives the source sequence. It is not obtained by declaring the old kernel sequence preserved under tensoring. This proves the domination/universal-injectivity equivalence.

The cokernel map \([m',n]\mapsto[n]\) identifies \(\operatorname{coker}f'\) with the original \(\operatorname{coker}f\). Under finite presentation, a universally injective \(f'\) has a retraction by the preceding source split criterion. If that retraction is \(h'\), then \(h=h'g'\) satisfies \(hf=g\). Conversely, if \(hf=g\), the formula \[
h'([m',n])=m'+h(n)
\] is well-defined since it kills \((g(x),-f(x))\), and is a retraction of \(f'\). This supplies the actual two directions needed below.

\subsubsection{Equalize the lifted maps at a finite stage}\label{algebra-proof-012-equalize-the-lifted-maps-at-a-finite-stage}

Source \texttt{proposition-ML-characterization}, 20999--21103; the missing step is at 21031. Keep its directed system \((M_i,f_{ij})\), colimit \(M\), and canonical maps \(f_i\). Condition (1) supplies \(g:M_i\to Q\) with exactly the same tensor kernels as \(f_i\). The cokernel of \(g\) is finitely presented: choose a finite presentation of \(Q\) and add lifts of the images of finitely many generators of \(M_i\) as finitely many extra relations. The domination lemma therefore gives \(h:Q\to M\) with \(hg=f_i\). Since \(Q\) is finitely presented, choose \(j\geq i\) and \(h':Q\to M_j\) with \(f_jh'=h\).

At this point only \[
f_j f_{ij}=f_i=f_jh'g
\] has been proved. Equality after the colimit map does not assert \(f_{ij}=h'g\) in \(M_j\). Choose generators \(a_1,\ldots,a_m\) of the original \(M_i\). Each difference \(f_{ij}(a_r)-h'g(a_r)\) maps to zero in the colimit, so vanishes after some transition from \(M_j\). Choose a common upper stage \(j'\) for these finitely many stages. Then \[
f_{ij'}=f_{jj'}h'g.
\] When the generating list is empty, both maps from \(M_i=0\) agree already. Replacing the stage by \(j'\) and the lift by \(f_{jj'}h'\) makes precisely the printed diagram commute. Its factorization shows that the new \(f_{ij'}\) dominates \(g\), which dominates \(f_i\), proving (2). The proposed insertion states this stage enlargement and the resulting equality, with the proof above retained separately.

For the converse, factor \(f:P\to M\), with \(P\) finitely presented, as \(f_i g'\). Choose \(j\) from (2). For every test module \(T\), \(f_i\) factors through \(f_{ij}\), so their tensor kernels on \(M_i\otimes T\) coincide by (2). Taking their inverse images under \(g'\otimes1\) gives exactly the equality of kernels for \(f\) and \(f_{ij}g'\). The latter map has finitely presented target \(M_j\), proving (1).

Conditions (2) and (3) are equivalent with the original quantifiers. If (2) holds, \(f_i\) dominates each \(f_{ik}\); transitivity and finite presentation of \(\operatorname{coker}f_{ik}\) give the required factorization of \(f_{ij}\) through \(f_{ik}\). Conversely, every element of \(\ker(f_i\otimes1_T)\) vanishes at some stage \(k\geq i\), by the original tensor-colimit map. A factorization through that stage makes it vanish under \(f_{ij}\otimes1_T\). This proves (2) and retains every tensor kernel rather than replacing them by a different presentation.

\subsubsection{Use a common upper bound for incomparable indices}\label{algebra-proof-012-use-a-common-upper-bound-for-incomparable-indices}

The second missing choice is at 21101--21102 in (5) implies (3). For (3) implies (4), fix \(i\) and the one index \(j\) from (3). If \(k\geq j\), a map \(M_j\to N\) composed with the factor \(M_k\to M_j\) gives the inclusion of its image in the image from \(M_k\). The reverse inclusion follows from \(f_{ik}=f_{jk}f_{ij}\). Thus the same index stabilizes the original Hom images for every \(N\). Condition (4) includes (5).

For (5), keep \(N=\prod_{s\in I}M_s\) and its stabilization index \(j\). There is a map \(e_j:M_j\to N\) equal to the identity in coordinate \(j\) and zero in all other coordinates. For \(k\geq j\), equality of the two Hom images gives \(H_k:M_k\to N\) such that \(H_k f_{ik}=e_j f_{ij}\). Projection to coordinate \(j\) gives \(h_k:M_k\to M_j\) with \[
f_{ij}=h_k f_{ik}.
\] This proves exactly the factor needed from the source's product-coordinate step. It does not require any assertion about exchanging arbitrary images with a product.

The printed final sentence treats \(k\geq j\) and \(k\leq j\); these need not exhaust a directed set. For any original \(k\geq i\), choose \(\ell\geq k,j\). The factor just proved for \(\ell\) gives \(h_\ell:M_\ell\to M_j\). Set \[
h=h_\ell f_{k\ell}:M_k\to M_j.
\] Then \(h f_{ik}=h_\ell f_{i\ell}=f_{ij}\). This proves (3) for incomparable indices as well. The proposition's five statements and definition of an ML module are unchanged.

\subsubsection{Retain the two tensor corrections and the finite-free test}\label{algebra-proof-012-retain-the-two-tensor-corrections-and-the-finite-free-test}

Source 21116--21274. Reports \texttt{OCC-12097}, \texttt{OCC-00452} and \texttt{OCC-00453} describe the two operations already composed as \texttt{MC-STK-ERR-0440}. The actual second system is \((N_j,g_{jj'})\). Thus the colimit terms are \(M_i\otimes_R N_j\), and the second factor map is \(b:N_{j''}\to N_{j'}\). With these types, \[
(a\otimes b)(f_{ii''}\otimes g_{jj''})
=f_{ii'}\otimes g_{jj'}.
\] The original tensor product is the colimit over the directed product set: its comparison sends the class of \(m_i\otimes n_j\) to \(f_i(m_i)\otimes g_j(n_j)\). For the inverse, choose representatives for both factors, take their tensor class, and use a common later pair to prove independence; additivity and balancing hold at a common pair. The two maps are inverse on simple tensors. The terms are finitely presented as in the source. Condition (3), now checked for every later pair including incomparable coordinates, proves the original tensor-closure result. No new operation duplicates either retained edit.

For the flat-module remark, write \(D_i=\operatorname{Hom}_R(M_i,R)\) for the duals of the original finite free stages, and let \(L_i\subset D_i\) be the stable image. Each sufficiently late map \(D_j\to D_i\) is the composite of a surjection onto \(L_i\) with its inclusion. After tensoring with an arbitrary \(T\), right exactness keeps that first map surjective; the resulting image is exactly \(\operatorname{im}(L_i\otimes T\to D_i\otimes T)\). This expression is independent of the later stage. The original finite-free Hom/tensor isomorphism carries it to the Hom images, proving the remark. Injectivity of \(L_i\otimes T\to D_i\otimes T\) is not needed.

Report \texttt{OCC-00454} only requires ``is clearly weaker'' at 21173. The converse argument is valid. If \(F\twoheadrightarrow P\) is the original finite free surjection and \(F\to Q\) has the same tensor kernels as \(F\to M\), its kernel contains that of \(F\to P\), so it induces \(g:P\to Q\). After tensoring, \(F\otimes T\to P\otimes T\) is onto. The inverse images under this surjection of the two kernels from \(P\otimes T\) are equal by hypothesis, hence the kernels themselves are equal. This proves the finite-free test without any additional source correction.

Restriction of scalars at 21200--21219 uses finitely presented \(S\)-stages that remain finitely presented over \(R\), under the stated finite and finitely presented ring-map hypothesis. The same actual factor maps remain \(R\)-linear. Although the paragraph only records factors for \(k\geq j\), the common-upper-bound construction above extends them to all \(k\geq i\), so no additional edit is necessary. In the quotient lemma 21221--21245, the same finite-presentation fact and the equality of the original \(R\)-linear and \(S\)-linear Hom sets justify the proof. A wider one-way implication is proved separately below.

The dual-number example keeps its original action \(\epsilon(a,b)=(0,u(a))\) on \(F_1\oplus F_0\); its square is zero. The quotient map \((a,b)\mapsto(a,[b])\) has kernel \(0\oplus\operatorname{im}u=\epsilon M\), so it induces exactly \(M/\epsilon M\cong F_1\oplus M_0\). Direct summands of ML modules are ML: given \(f:P\to L\), compose with a split inclusion \(L\to M\). Tensoring that inclusion remains injective because its retraction also tensors; hence the two maps from \(P\) have identical tensor kernels. Apply condition (1) for \(M\). Thus a non-ML summand \(M_0\) forces this quotient to be non-ML. Free modules are ML because any map from a finitely presented module has image in a finite coordinate summand, whose inclusion remains split after tensoring. If the original \(M\) were ML over \(S\), tensor closure with the finitely presented \(S/(\epsilon)\) would make its quotient ML over \(S\), and the quotient lemma would make it ML over \(R\), a contradiction. These arguments verify the original example without changing its action or presentation.

\subsubsection{The coefficient projection and the original product criteria}\label{algebra-proof-012-the-coefficient-projection-and-the-original-product-criteria}

Source 21277--21441; report \texttt{OCC-00455}. The proof at 21349 defines \(p_x:M^M\to M\), whereas \(a_i\) belongs to \(R^M\). Keep this \(p_x\) and replace only its out-of-domain application \(p_x(a_i)\) at 21353 by the coefficient \((a_i)_x\in R\). The canonical tensor-product map then gives exactly \[
x=p_x(d)=\sum_i p_x(f(x_i\otimes a_i))=\sum_i(a_i)_x x_i.
\] The finitely many original \(x_i\) generate \(M\). The zero module permits an empty sum. Conversely, a finite free surjection onto \(M\) tensors onto each receiving module and onto the product. Finite free tensoring commutes with the product by its coordinate maps, and the product of the surjections is onto by choice. The source square therefore proves the surjectivity criterion.

For finite presentation, keep a finite presentation \(R^m\xrightarrow{A}R^n\to M\to0\). The finite-free product comparison isomorphisms carry the image of the original matrix \(A\) in the first row of the source diagram to its image in the second row; those rows are right exact. Consequently they induce inverse isomorphisms on the two cokernels, precisely the canonical map for \(M\). In the converse, retain \(K=\ker(F\to M)\) for the chosen finite free surjection. For any set \(A\) and \(k\in K^A\), include it in \(F^A\), and use the isomorphism \(F\otimes R^A\to F^A\) to lift it to \(z\). The image of \(z\) in \(M\otimes R^A\) maps to zero in \(M^A\); injectivity of that comparison makes this image zero. Right exactness lifts \(z\) from \(K\otimes R^A\). By the commutative square and injectivity \(K^A\to F^A\), that lift maps to the original \(k\). Thus the comparison for \(K\) is onto for every \(A\), and the preceding criterion makes \(K\) finite. No injection at the left of the tensor row was assumed.

The original rational examples are correct. The residue map \(\mathbb Z\to\prod_{n\geq1}\mathbb Z/n\) is injective since an integer divisible by every positive integer is zero, including the zero factor at \(n=1\). Tensoring by the original localization \(\mathbb Q\) preserves that injection; each \(\mathbb Q\otimes\mathbb Z/n\) is zero. For the second example, every finite sum in \(\mathbb Q\otimes\prod\mathbb Z\) admits one common nonzero integer denominator, and conversely \((a_n/m)\) is the image of \((1/m)\otimes(a_n)\). The sequence \((1/n!)_{n\geq1}\) has no such denominator: it would require \(m/n!\) to be an integer for all \(n\), impossible for a fixed nonzero \(m\).

The example's non-ML premise can also be realized explicitly without using the as-yet-unreviewed tensor characterization. Write \(\mathbb Q\) as the colimit of \(M_n=\mathbb Z\), \(n\geq0\), with \(f_{n,n+1}\) multiplication by \(n+1\), and stage map \(m\mapsto m/n!\). The maps are compatible; every rational denominator divides some factorial, giving surjectivity to \(\mathbb Q\), and these injective stage maps give injectivity of the induced map from the colimit. For the test target \(\mathbb Z\), the image of \(\operatorname{Hom}(M_k,\mathbb Z)\) in \(\operatorname{Hom}(M_j,\mathbb Z)\) is \((k!/j!)\mathbb Z\). These subgroups do not stabilize: for each sufficiently large \(k\), the next is the proper multiple \((k+1)(k!/j!)\mathbb Z\). Condition (4) therefore proves \(\mathbb Q\) is non-ML. A presentation matching the original example is \(F_1\xrightarrow{u}F_0\to\mathbb Q\to0\), with bases \(c_n,b_n\), \[
u(c_n)=b_n-(n+1)b_{n+1},\qquad b_n\longmapsto1/n!.
\] The composites are zero. Any finite sum in \(F_0\) reduces modulo these relations to an integer multiple of the last basis element in its finite support, by successive substitutions. If its rational image is zero, that remaining integer is zero. This proves exactness and retains every coefficient of the presentation. It is an editorial realization of the source example, not a replacement of its general statement.

Finally, \texttt{lemma-kernel-tensored-fp}, 21414--21440, is correct. Factor the original map from the finitely presented \(P\) through some \(M_j\); the given single element \(x\) has image \(y\) in \(M_j\otimes Q\). Since its colimit image is zero, the defining filtered-colimit equivalence gives a later stage where that same \(y\) is zero. The original composite \(P\to M_j\to M_{j'}\) is the required map into the finitely presented \(P'=M_{j'}\). This uses a witness for the one element, not a claim that an entire arbitrary kernel vanishes at a uniform stage.

\subsubsection{Mittag-Leffler base change and descent along pure ring maps}\label{algebra-proof-012-mittag-leffler-base-change-and-descent-along-pure-ring-maps}

This is a separate consequence of the corrected characterization, receiving the finite-equation method of \texttt{ALGEBRA-RECON-501}. For every ring map \(R\to S\), an ML \(R\)-module \(M\) gives an ML \(S\)-module \(M\otimes_R S\). If the ring map is pure, meaning that every original unit map \[
X\longrightarrow X\otimes_R S,\quad x\longmapsto x\otimes1
\] is injective, the converse also holds. No flatness of \(S\) is assumed. Here is the complete factorization calculation.

First consider actual \(R\)-linear maps \(u:A\to B\), \(v:A\to C\), where \(A\) is finitely generated and \(B\) is finitely presented; \(C\) is arbitrary. Suppose there exists \(h_S:B\otimes S\to C\otimes S\) with \(h_Su_S=v_S\). Choose generators \(a_1,\ldots,a_m\) of \(A\) and a finite presentation \[
R^t\xrightarrow{D}R^n\xrightarrow{q}B\to0.
\] For each \(\ell\) choose an actual lift \(\sum_{r=1}^n U_{r\ell}e_r\) of \(u(a_\ell)\), retaining the entire matrices \(D=(D_{rs})\) and \(U=(U_{r\ell})\). Set \(c_\ell=v(a_\ell)\). Define \[
\begin{split}
\beta:C^n&\longrightarrow C^t\oplus C^m,\\
(z_r)_{r=1}^n&\longmapsto
\left(\left(\sum_{r=1}^n D_{rs}z_r\right)_{s=1}^t,
\left(\sum_{r=1}^n U_{r\ell}z_r\right)_{\ell=1}^m\right),\\
c&=(0,(c_\ell)_{\ell=1}^m).
\end{split}
\] The values \(z'_r=h_S(q(e_r)\otimes1)\) solve \(\beta_S(z')=c\otimes1\): the first equations are the original relations in \(B\), and the second are the factorization on the chosen generators of \(A\). Tensoring commutes with these finite direct sums and with the cokernel \(W=\operatorname{coker}\beta\). Hence the class of \(c\) maps to zero in \(W\otimes S\). Purity of its unit map makes that class zero already in \(W\). Choose \(z\in C^n\) with \(\beta(z)=c\). The first equations make \(e_r\mapsto z_r\) descend to an actual map \(h:B\to C\), and the second give \(hu=v\) on every generator of \(A\), hence on all of \(A\). Empty presentations and zero modules use the same maps with their zero dimensions. No identification of a Hom module after base change is required.

Now choose any original directed presentation \(M=\operatorname{colim}_i M_i\) by finitely presented \(R\)-modules. Right exactness gives the base-changed finite presentations and tensor-colimit comparison, so \(M\otimes S=\operatorname{colim}_i(M_i\otimes S)\) with its actual transition maps. If \(M\otimes S\) is ML, then for each \(i\) one index \(j\geq i\) satisfies condition (3) over \(S\): for every \(k\geq i\), \((f_{ij})_S\) factors through \((f_{ik})_S\). Apply the preceding calculation to \[
A=M_i,\quad B=M_k,\quad C=M_j,\quad u=f_{ik},\quad v=f_{ij}.
\] It supplies an \(R\)-linear factor for every such \(k\), keeping the same \(i,j\). Thus the original system satisfies (3) over \(R\), proving descent. Conversely, for an ML \(M\), tensor each actual factor in (3) with \(S\); the resulting factors prove the criterion for the base-changed system. This converse uses no purity assumption and proves arbitrary base-change ascent.

This calculation propagates to the source quotient lemma without changing its printed assumptions. For any ideal \(I\) and any \(S=R/I\)-module \(M\), the maps \[
M\otimes_R S\longrightarrow M,\quad m\otimes\overline r\longmapsto rm,
\qquad M\longrightarrow M\otimes_R S,\quad m\longmapsto m\otimes1
\] are well-defined \(S\)-linear inverses because \(IM=0\) and tensor balancing gives \(rm\otimes1=m\otimes\overline r\). Consequently the implication from ML over \(R\) to ML over \(R/I\) does not need finite generation of \(I\). The other direction in the source retains its finite-generation hypothesis and its already checked restriction proof.

There is no conflict with the dual-number counterexample: it concerns the original \(S\)-module with its specified square-zero action, whereas the base-change theorem concerns \(M_R\otimes_R S\). For the example with underlying \(M_R=F_1\oplus F_0\), the latter is free over \(S\) with the base-changed basis; the former has the explicit non-ML quotient just computed. The actual \(S\)-linear comparison is \[
\mu:M_R\otimes_R S\longrightarrow M,\quad m\otimes s\longmapsto s\cdot m.
\] It is onto, has an \(R\)-linear section \(m\mapsto m\otimes1\), and satisfies \(\mu(m\otimes\epsilon)=(0,u(a))\) for \(m=(a,b)\). Thus the precise relationship is a quotient of the base-changed free module; ML is not asserted to survive arbitrary quotients. This both supplies the connecting map and checks the source example against the broader result.

Pure maps include genuinely nonflat examples. For \(\mathbb Z\to\mathbb Z\times\mathbb Z/2\), \(n\mapsto(n,\overline n)\), first-coordinate projection \(r\) is a ring retraction. The map \(X\otimes S\to X\), \(x\otimes s\mapsto r(s)x\), retracts every unit map, proving purity. Multiplication by two kills the nonzero element \((0,\overline1)\) of \(S\), so tensoring the injection \(\mathbb Z\xrightarrow{2}\mathbb Z\) with \(S\) fails to preserve injectivity. This proves nonflatness. The ML descent statement therefore uses the same additional class of ring maps as group 501, with its full finite-equation comparison proved above. It asserts neither arbitrary projectivity nor effective descent of arbitrary module objects.

\subsubsection{Propagation and reading limits}\label{algebra-proof-012-propagation-and-reading-limits}

Groups 505 and 506 repair the two stage choices in the original characterization. Group 507 retains \texttt{MC-STK-ERR-0440} for all three reports and propagates the corrected all-indices quantifier into its tensor factors. Groups 508 and 509 handle the adverb and scalar coefficient. Group 510 records countable-cofinal nonemptiness and exactness with zero source operations. Group 511 records arbitrary base-change ascent and pure descent, receives group 501 and the corrected characterization, and checks both the one-way quotient consequence and the actual comparison to the source's dual-number module. Later uses remain to be reviewed in source order; the broader combined-results paper remains after core completion.

The new bounded disk query \texttt{Mittag-Leffler\ pure} returns a TeX-routing record for \emph{Higher-Rank Mathieu Opers, Toda Chain, and Analytic Langlands Correspondence}, a research PDF-routing record for \emph{Operators, Inequalities and Approximation}, and the local unpublished PDF-routing record \texttt{collatz\_path\_algebra}. All three remain unread and unused as proof evidence; no PDF fallback was performed. The initial display omitted the local layer by using a mismatched layer name; inspection of the saved raw result identified that third hit before adjudication. The topic snapshot preserves all three actual results together with earlier query history. The original Stacks TeX and the fully written comparisons above supply this batch's evidence; the lookup is not a novelty survey. Partial lookahead 21442--21480 remains unadjudicated.

\hypertarget{algebra_ml_coherent_notes_20260927}{}
\subsection{Tensor criteria, coherent rings and countable projectivity}\label{algebra-proof-013-tensor-criteria-coherent-rings-and-countable-projectivity}

Separate editorial evidence for the Stacks Project authors' \texttt{algebra.tex}, authority commit \texttt{a04446e57ec1fbc252a871afcec7752fb2807b14}, lines 21442--22460. Authority SHA-256: \texttt{FA8BB92E58A4F78A2BD01B3B6A4A87DE0A0D279F5DD90641B574DD5FBFFFA4F3}. All nineteen due reports and all complete source units in this interval were read. The next theorem's statement at 22461--22465 was read only as lookahead; its proof remains unread.

The source's Chase citation and its attribution of the faithfully flat ML argument to Juan Pablo Acosta Lopez are retained. Those citations are not claims of independent reading of the cited works or correspondence. Four already composed corrections, \texttt{MC-STK-ERR-0451} through \texttt{0454}, account for eight reports and six operations. Ten further minimal operations are proposed, and the false article report is rejected. Full supporting proofs and the two further consequences remain separate editorial material; the source statements and original hypotheses remain identifiable. No chapter or translation mutation, novelty claim, or later synthesis completion is asserted.

\subsubsection{The tensor criterion and its receiving diagrams}\label{algebra-proof-013-the-tensor-criterion-and-its-receiving-diagrams}

Source 21442--21687. The tensor criterion is correct with the ML characterization previously repaired in groups 505--506. For its forward direction, represent the original kernel element by \(x_i\in M_i\otimes\prod Q_\alpha\), and choose one \(j\geq i\) for which \(f_{ij}\) and \(f_i\) have equal tensor kernels for every test module. The coordinates of the comparison image of \(x_i\) vanish in \(M\otimes Q_\alpha\), hence in \(M_j\otimes Q_\alpha\). The finite-presentation product comparison for \(M_j\) is injective, so the image of \(x_i\) at that stage is zero. Thus the original colimit element is zero.

For the converse, finite presentations are specified by finite matrices with entries in the set \(R\), so their isomorphism classes have a set of representatives. The source's set of pairs \((Q,x)\), with \(x\in\ker(P\otimes Q\to M\otimes Q)\), is therefore a set. Retain its entire product, including all such kernel elements. The finite-presentation comparison for \(P\) lifts the tuple of these elements to one element \(x\in P\otimes\prod Q_\alpha\). The assumed injection for \(M\) makes its image zero. The already reviewed kernel-witness lemma factors \(P\to M\) through a finitely presented \(P'\) killing that one product element. Projecting to every coordinate kills every original finite-test kernel. The reverse inclusion follows from the actual factorization through \(P'\), and the finite-test domination lemma extends equality to all test modules. This proves the criterion without interchanging an unrestricted colimit and product.

For the universally exact sequence at 21565, tensoring with the product gives an exact row, and the product of the individually exact rows is exact. If the middle comparison is injective, the first is injective by its inclusion into the middle row. If the first and third comparisons are injective and a middle element maps to zero, its image at the third top term is zero; lift it from the first top term, where injectivity of the first comparison makes the lift zero. These are precisely the two asserted ML permanence statements.

For the quotient result at 21591, lift an element of the third top term with zero comparison image to the second top term. Its comparison image comes from the first bottom term by exactness. Finite generation makes the first comparison surjective, so lift that preimage to the first top term and subtract its image from the chosen second-top lift. The difference has zero middle comparison, hence is zero by the ML hypothesis on \(M_2\). Its image at the third top term was the original element, so that element is zero. This uses the source's finite generation assumption exactly.

In a filtered system with universally injective transitions, each map from a stage to the colimit stays injective after every tensor test: an element mapping to zero vanishes at a later stage, and that transition is injective. Thus every product of the individual stage-to-colimit tensor maps is injective. Representing one tensor-product kernel element at a stage now proves the colimit lemma. Finite direct sums are handled by the universally exact split sequences; an arbitrary direct sum is the filtered union of its finite subsums with split transitions. Conversely, every summand inclusion is split after tensoring, so the pure-submodule result applies. These maps also justify the earlier free and projective ML examples.

For restriction along \(R\to S\) in 21663--21687, keep the source's actual composite \[
M\otimes_R\prod Q_\alpha
\cong M\otimes_S(S\otimes_R\prod Q_\alpha)
\longrightarrow M\otimes_S\prod(S\otimes_RQ_\alpha)
\longrightarrow\prod(M\otimes_RQ_\alpha).
\] The associativity maps send \(m\otimes q\) to \(m\otimes(1\otimes q)\), with inverse \(m\otimes(s\otimes q)\mapsto sm\otimes q\). The first displayed arrow preserves the ML injection for \(S\) because \(M\) is flat over \(S\); the second is the ML comparison for \(M\) over \(S\). Their composite on simple tensors is the original \(R\)-comparison, proving the stated restriction result.

\subsubsection{Coherent-module extensions and the retained image correction}\label{algebra-proof-013-coherent-module-extensions-and-the-retained-image-correction}

Source 21700--21844; \texttt{OCC-12108} and \texttt{OCC-00456} retain \texttt{MC-STK-ERR-0451}. The second report's named proposition is inaccurate, but its line range identifies the correct step in \texttt{lemma-coherent}. The left term of \[
0\longrightarrow\operatorname{im}\varphi\longrightarrow E'\longrightarrow E\longrightarrow0
\] is the finite image, not the kernel from a different sequence. The existing operation supplies exactly that image.

Here are the finiteness arguments with the original modules. The image of the finite \(N\) is finite, hence finitely presented because it lies in coherent \(M\). Take a finite free surjection onto \(N\). The kernel of its composite onto \(\operatorname{im}\varphi\) is finite, since that target is finitely presented; its image surjects onto \(\ker\varphi\), proving the latter finite. For a finite submodule \(E\) of the cokernel, choose lifts to \(E'\) of finitely many generators of \(E\) and adjoin generators of \(\operatorname{im}\varphi\). Their span is all \(E'\). It is therefore finitely presented as a finite submodule of \(M\). Quotienting a finite presentation by a finite generated submodule adds finitely many lifted generator relations, so \(E\) is finitely presented. This proves the corrected cokernel argument and coherence of the image and cokernel.

The coherent kernel is a finite submodule of coherent \(N\). In an extension with coherent end terms, the middle term is finite by lifting generators. For its finite submodule \(N_2\), the image \(N_3\) is finitely presented in the third coherent module; the preceding kernel argument makes \(N_1\) finite, hence finitely presented in the first coherent module. The original extension of these two finite presentations presents \(N_2\) finitely. Explicitly, lift finitely many generators of \(N_3\) to \(N_2\), adjoin generators of \(N_1\), lift the finitely many relations of \(N_3\) into those of \(N_1\), and adjoin the relations of \(N_1\). Any relation reduces first to a relation in the quotient and then to one in \(N_1\), proving these finitely many relations suffice. This completes all the stated coherent-module implications.

The coherent-ring criterion follows with the author's original product maps. A finite ideal \(I\) is finitely presented over a coherent ring, so \(I\otimes\prod Q_\alpha\to\prod(I\otimes Q_\alpha)\) is the original comparison isomorphism. Its composite into \(\prod Q_\alpha\) is injective if each \(Q_\alpha\) is flat, giving flatness of the product by the ideal criterion. Conversely, if every \(R^A\) is flat, the map \(I\otimes R^A\to R^A\) is injective. Its image is \(I^A\), because finite generation makes the product comparison onto. Thus the comparison \(I\otimes R^A\to I^A\) is an isomorphism for every \(A\); the preceding finite-presentation criterion makes \(I\) finitely presented. The zero ideal and empty products cause no exception. The valuation-ring and Noetherian examples retain their original hypotheses; their finite ideals are respectively principal free ideals (or zero), and finite presented ideals by Noetherianity.

\subsubsection{The finite-free content criterion and products over the source ring}\label{algebra-proof-013-the-finite-free-content-criterion-and-products-over-the-source-ring}

Source 21859--21993. \texttt{OCC-00457} inserts the second ``if''; \texttt{OCC-00458} inserts ``of''. A finite module is ML exactly when it is finitely presented, since the finite-generation comparison is onto, the ML comparison is injective, and the combined isomorphism is the finite-presentation criterion.

Reports \texttt{OCC-12109} and \texttt{OCC-00459} retain \texttt{MC-STK-ERR-0452}: there are \textbf{three} undefined occurrences of \(J\), at 21937 and twice at 21944, all already changed to \(I\) by its two operations. The smaller report's suggested count of two occurrences does not reduce this complete correction.

In the original flat ML criterion, let \(D=\operatorname{Hom}_R(M_i,R)\), and retain the finite free stages. Choose a basis \(e_1,\ldots,e_r\) of \(M_j\), and write the actual transition tensor as \[
f_{ij}=\sum_{a=1}^r \lambda_a\otimes e_a\in D\otimes M_j.
\] The coefficients \(\lambda_a\) are the images of the dual basis under \(\operatorname{Hom}(f_{ij},R)\). Thus \(Q_j=\sum_aR\lambda_a\). Because \(M_j\) is free, this tensor belongs to \(L\otimes M_j\) for a submodule \(L\subset D\) precisely when each \(\lambda_a\in L\). This proves the smallest-submodule assertion used in the source.

Let \(Q\) be the smallest submodule containing the canonical tensor \(f_i\) after tensoring by \(M\). Since \(M\) is flat, these tensor inclusions are actual injections. Each original \(Q_j\) contains \(Q\). A tensor representation of \(f_i\) in \(Q\otimes M\) is represented in \(Q\otimes M_k\) at some stage. Choose a common stage above \(i,k\); the image of this representative and the transition tensor from \(M_i\) have the same colimit image in \(D\otimes M\). At a further common stage they are equal. Hence the source's \(f_{ij}\) belongs to \(Q\otimes M_j\) there, so \(Q_j\subset Q\). This proves stabilization with the original maps and, by the earlier dual criterion, ML. The other direction follows from the minimal-submodule lemma.

For \(M=R^A\) over Noetherian \(R\), flatness follows from the coherent-ring criterion just checked. Every submodule \(F'\) of a finite free \(F\) is finitely presented. The product comparisons therefore identify the original inclusion \(F'\otimes R^A\to F\otimes R^A\) with \((F')^A\to F^A\). For the given tuple \((x_\alpha)\), its coordinate span \(F'=\sum_\alpha Rx_\alpha\subset F\) is the smallest such submodule and is finite by Noetherianity. \texttt{OCC-00460} repairs only the malformed apposition ``submodule of F' of F''. For the power-series statement, the coefficient map from the original \(R[[t_1,\ldots,t_n]]\) to \(R^{\mathbb N^n}\) and the inverse formal coefficient sum are \(R\)-linear inverse maps. For finite positive \(n\), the set of exponent tuples is countable; every original monomial and its coefficient is retained.

\subsubsection{The non-ML examples with their original elements}\label{algebra-proof-013-the-non-ml-examples-with-their-original-elements}

Source 21995--22040. Reports \texttt{OCC-00461} and \texttt{OCC-00462} require only ``generated by'' and the relation symbol \texttt{\textbackslash{}gg}. The mathematical assertions are correct, as follows.

For \(R=k[[x]]\), keep \(M=\prod_{n\geq1}R/(x^n)\) and the displayed element whose coordinate at \(n=2^m\), \(m\geq1\), is \(x^{2^{m-1}}\), and whose other coordinates are zero. This is exactly the displayed sequence starting \((0,x,0,x^2,\ldots)\). If the natural-number indexing also includes zero, that extra module \(R/(x^0)\) is zero and has no effect. Coordinatewise division of elements in the image of multiplication by \(x^l\) gives \[
M/x^lM\cong\prod_{n\geq1}R/(x^{\min(n,l)}).
\] Thus the annihilator of the original \(\xi\) is the intersection of the coordinate annihilators, namely \[
\operatorname{Ann}_R(\xi\bmod x^lM)=(x^{b(l)}),
\quad b(l)=\max_{m\geq1}\max\{\min(2^m,l)-2^{m-1},0\}.
\] Only finitely many terms are positive. For \(l=2^r\), \(r\geq1\), the maximum is \(2^{r-1}\): the term \(m=r\) attains it, earlier terms are smaller, and later terms vanish. This proves the exact annihilator claimed in the source.

For any finite \(R\)-module \(Q\) and its actual element \(\xi'\), let \(A=R\xi'\subset Q\). The earlier original Artin--Rees lemma, \texttt{lemma-Artin-Rees}, gives \(c\) such that \[
A\cap x^lQ=x^{l-c}(A\cap x^cQ)\qquad(l\geq c).
\] If \(\xi'\) is torsion with annihilator \((x^a)\), the right side is zero for \(l\geq c+a\), so its annihilator modulo \(x^lQ\) is exactly \((x^a)\). This includes \(\xi'=0\), with \(a=0\). Otherwise the map \(R\to A\), \(r\mapsto r\xi'\), is injective. Its inverse image of \(A\cap x^cQ\) is a principal ideal \((x^b)\), with \(0\leq b\leq c\), because it contains \(x^cR\). The displayed Artin--Rees equality then gives the annihilator \((x^{l-c+b})=(x^{l-a})\), where \(a=c-b\geq0\), for all \(l\geq c\). This retains the original submodule and every exponent; no replacement of \(Q\) by a presumed decomposition is used. Neither eventual form agrees with \((x^{2^{r-1}})\) for all large \(l=2^r\). Applying ML condition (1) to \(R\to M\), \(1\mapsto\xi\), would give a finitely presented, hence finite, \(Q\) with equality of all these tensor-test kernels for \(R/(x^l)\). The contradiction proves the example.

For the completion of \(D=\bigoplus_{n\geq1}R/(x^n)\), retain the exact completion \(\widehat D=\lim_l D/x^lD\). The coordinate maps identify it with the submodule of the above product consisting of sequences \((a_n)\) for which, for each \(l\), all but finitely many coordinates belong to \(x^lR/(x^n)\). Indeed a compatible finite-support residue family determines each coordinate once \(l\geq n\), and conversely such a sequence supplies finite-support residues at each \(l\); these maps are inverse. Moreover \[
\widehat D\cap x^lM=x^l\widehat D.
\] For the nontrivial inclusion, divide the original coordinate polynomials by \(x^l\), choosing zero for zero coordinates. If the original sequence is eventually divisible by \(x^{l+s}\), the resulting sequence is eventually divisible by \(x^s\), so it still lies in \(\widehat D\). The displayed equality makes \(\widehat D/x^l\widehat D\to M/x^lM\) injective. The original \(\xi\) lies in \(\widehat D\) because its nonzero valuations tend to infinity. Its annihilators modulo every \(x^l\) are therefore the same ones just computed, and the same finite-module contradiction applies.

For the last example, retain \(R=k[a,b]/(a^2,ab,b^2)\) and \(S=R[t]/(at-b)\). There is an exact comparison with pairs of polynomials: \[
f(t)+a g(t)+b h(t)\longmapsto(f(t),g(t)+t h(t)),
\quad (f,g)\longmapsto f+a g\text{ in }S,
\] where pair multiplication is \((f,g)(f',g')=(ff',fg'+gf')\). The first map kills \(at-b\), and its kernel before quotienting consists exactly of \((b-at)h\); thus these are inverse ring maps. Under them the original actions are \(a(f,g)=(0,f)\), \(b(f,g)=(0,tf)\).

Every \(R\)-linear endomorphism of the original \(S\), written using this comparison, has form \[
(f,g)\longmapsto(hf,Hf+hg)
\] for some polynomial \(h\in k[t]\) and some \(k\)-linear map \(H:k[t]\to k[t]\). To prove this, commuting with \(a\) forces the image of \((0,g)\) to be \((0,Fg)\), where \(F\) is the first component on \((g,0)\); commuting with \(b\) forces \(F(tf)=tF(f)\). Hence \(F\) is multiplication by \(h=F(1)\), while the second component \(H\) on \((f,0)\) is unrestricted. If the endomorphism is idempotent, then \(h^2=h\), so \(h=0\) or \(1\) in the domain \(k[t]\). The second component of the idempotence equation gives \(H=0\) in either case, in every characteristic. Only the zero and identity idempotents exist, proving indecomposability. The original elements \(t^n\) generate countably, while the quotient by \((a,b)\) is \(k[t]\), so the module is not finite.

The non-ML conclusion can be checked without relying on the later henselian splitting lemma. Let \(S_n=R\langle1,t,\ldots,t^n\rangle\subset S\). These are finite presented \(R\)-modules because \(R\) is finite dimensional over \(k\), hence Noetherian, and their union is \(S\). Their first polynomial coordinates have degree at most \(n\); their second coordinates have degree at most \(n+1\). For a proposed stabilization stage \(j\) at \(i=0\), set \(k=j+1\). A factor \(u:S_k\to S_j\) of the inclusion from \(S_0\) would satisfy \(u(1)=1\). Write \(F_r\) for the first coordinate of \(u(t^r)\). The original relations \(at^{r+1}=bt^r\) give \(F_{r+1}=tF_r\), so \(F_r=t^r\). At \(r=k\), this contradicts the degree bound \(j\) in \(S_j\). Thus condition (3) fails for every \(j\), proving non-ML directly while preserving the original example and its source attribution.

\subsubsection{Countable subsystem, parity and cofinal sequences}\label{algebra-proof-013-countable-subsystem-parity-and-cofinal-sequences}

Source 22052--22179. The two article reports \texttt{OCC-12110}, \texttt{OCC-00463} retain \texttt{MC-STK-ERR-0453}, which already says ``a countable set''. The two parity reports \texttt{OCC-12111}, \texttt{OCC-00464} retain \texttt{MC-STK-ERR-0454}. Its witness step acts on even stages and its directed-enlargement step on odd stages, exactly as in the explicit first two cycles.

A countable subset of a directed set can be enlarged to a countable directed subset by adjoining an upper bound for every pair, repeating countably many times, and taking the union. There are countably many pairs at each step. The alternating original construction therefore yields a countable directed \(I'\) containing the chosen generator stages and, for each of its stages \(i\), an index \(j\in I'\) whose map factors through every \(k\geq i\) in the original \(I\). The canonical colimit map is onto because its image contains every generator. If a represented element becomes zero in the full colimit, it vanishes at some such \(k\). The factor at \(j\) kills it already in the restricted system, proving injectivity. No cofinality of \(I'\) in the entire original \(I\) is asserted or needed.

Report \texttt{OCC-00465} correctly detects the reversed order, but changing the relation symbol alone leaves a boundary error: a finite directed set cannot contain a subset isomorphic to \(\mathbb N\). The proposed minimal replacement says to obtain \textbf{an increasing cofinal sequence indexed by \((\mathbb N,\leq)\)}. Repetitions are allowed, including a constant sequence at a greatest element. Given an enumeration of the countable directed set, repeat entries if needed and choose successive upper bounds. For the colimit comparison, send a representative at an original index \(i\) to its image at any sequence stage above \(i\). Different choices agree at a common later sequence stage. Any original equality is witnessed at an original upper stage and then at a sequence stage above it, so this map is well-defined and inverse to the canonical map from the sequential colimit. This proves the intended assertion for finite and infinite index sets.

The polynomial example at 22109--22114 is correct. Its ideal \((x_i)\) is not finite generated: a proposed finite list of polynomial generators uses only finitely many variables, and setting those variables to zero leaves another \(x_j\) nonzero. Hence the cyclic quotient is not finitely presented. Its sequential presentation by \(R/(x_1,\ldots,x_n)\) is still the exact colimit, because the union of the defining ideals is \((x_i)\). Sequential finite presentation alone therefore does not give the ML property.

In \texttt{lemma-ML-countable}, a prescribed \(y\in E_j\) lifts to the inverse limit of the stable-image system, not merely to an unspecified nonempty inverse limit. To verify the source's use of the earlier lemma, choose an increasing cofinal sequence starting above \(j\). Surjectivity of the stable-image transition gives a first lift of \(y\), and then successive lifts. Extend this sequence to the full compatible family using the exact receiving maps proved in group 510. It projects to \(y\). Its \(k\)-component is the original \(z_k:M\to M_k\) with \(z_k f_j=f_{jk}\). Hence \(\alpha=f_kz_k\) fixes \(\operatorname{im}f_j\), and therefore fixes \(f(P)\). Finite generation of \(P\) is used only to put its image in one stage; it need not supply a map \(P\to M_j\). This validates the source proof without an additional rewrite.

\subsubsection{Projectivity and the original countability boundary}\label{algebra-proof-013-projectivity-and-the-original-countability-boundary}

Source 22192--22299. Restrict the Lazard finite free presentation of a flat countably generated ML module to the countable subsystem just proved. For any short exact sequence \(0\to N_1\to N_2\to N_3\to0\), finite freeness gives the source exact Hom sequence at every stage. The ML condition stabilizes the inverse system with target \(N_1\). The exact inverse-limit theorem then gives surjectivity onto \(\lim\operatorname{Hom}(M_i,N_3)\). The maps \[
\operatorname{Hom}(M,N)\longrightarrow\lim_i\operatorname{Hom}(M_i,N),
\quad h\longmapsto(hf_i)_i
\] have inverse defined on a colimit class by \([m_i]\mapsto h_i(m_i)\); compatibility proves independence and linearity. Thus \(\operatorname{Hom}(M,-)\) preserves this arbitrary short exact sequence, proving projectivity.

The example \(\mathbb Z[[x]]\) retains its essential countability boundary. For the original submodule \(N\) of coefficient sequences tending to zero \(p\)-adically, the map \((a_i)\mapsto\sum_i a_i p^i x^i\) from \(\mathbb Z^{\mathbb N}\) is injective, so \(N\) is uncountable. Reduction of coefficients maps \(N\) onto \(\bigoplus_{i\geq0}\mathbb F_p x^i\), since every reduced sequence has finite support and every finite-support sequence lifts to a polynomial. Its kernel is exactly \(pN\): division of a sequence divisible coordinatewise by \(p\) preserves the defining decay, using the original condition at exponent \(m+1\). Thus \(N/pN\) has countably infinite dimension. A free abelian group of countable rank is countable, while a free group of uncountable rank has uncountable-dimensional reduction modulo \(p\). Therefore \(N\) is not free.

The assertion that subgroups of free abelian groups are free has a proof compatible with the previously reviewed original filtration. Well-order the original basis of a free group \(F\), and intersect a subgroup \(H\) with the initial coordinate sums. Each successor quotient of the intersections injects into the corresponding \(\mathbb Z\) coordinate, so is an ideal \(d\mathbb Z\) or zero; it is free of rank at most one. Lift a generator to split the successor extension. At a limit the intersections commute with the coordinate union by finite support. The original direct-sum dévissage lemma now makes \(H\) a direct sum of those free quotients. A projective abelian group is a subgroup of a free group and is therefore free; its subgroups would also be free. The explicit \(N\) contradicts this for \(\mathbb Z[[x]]\).

The source projectivity characterization consequently keeps all three conditions. In a direct sum, flatness and ML pass to every summand; countably generated summands satisfying both are projective by the preceding proof. A direct sum of projectives is projective because, for any surjection, lift each summand's map separately and combine on finite-support elements. The converse is the already reviewed Kaplansky decomposition and permanence of flatness and ML under summands. The pure-submodule and power-series corollaries follow by the actual pure injection, the source flatness descent for its domain, and the checked tensor criterion.

\texttt{OCC-00466} inserts the article in ``be a universally injective map''. \texttt{OCC-00467} is rejected: authority line 22289 already reads \texttt{let\ \$M\$\ be\ an\ \$R\$-module}, and the current interval has that same reading. No operation is warranted by its inaccurate quotation.

\subsubsection{Source ascent, descent and the countable image construction}\label{algebra-proof-013-source-ascent-descent-and-the-countable-image-construction}

Source 22303--22460. Base change preserves flatness by the actual tensor associativity test on injections of \(S\)-modules, and it preserves the original splitting that witnesses projectivity. It commutes with direct sums and sends a countable list of generators to a countable list of tensor generators. The ML ascent proof agrees exactly with group 511's tensoring of the original factor maps. None of these arguments changes the source module into an unrelated \(S\)-module.

Reports \texttt{OCC-00468} and \texttt{OCC-00469} supply ``to be'' and lowercase ``is''. In the faithfully flat ML proof, retain its pushout \[
N=(M_j\oplus M)/\{(f_{ij}(x),-f_i(x)):x\in M_i\}.
\] Tensoring this cokernel constructs the exact displayed pushout over \(S\). ML of \(M\otimes S\) gives \(j\) for which \(M_j\otimes S\to N\otimes S\) is universally injective. For any original test module \(T\), use the \(S\)-test \(T\otimes S\) and the canonical associativity maps. The resulting injection is the base change of \(M_j\otimes T\to N\otimes T\). Faithful flatness reflects this injection, proving universal injectivity over \(R\) and hence domination of \(f_i\) by \(f_{ij}\).

In fact the reflection step also holds for the pure maps of group 511: for any map \(u:X\to Y\), if \(u_S\) is injective and \(u(x)=0\), then \(x\otimes1=0\), and purity of \(X\to X\otimes S\) gives \(x=0\). Apply this to each of the original tensor-test maps just specified. This supplies a second exact proof of the already recorded pure ML descent, propagated through the source's actual pushout. The source faithfully flat hypothesis and its human attribution remain in the translated statement.

For countable generation, the source collects all original entries \(x_{ij}\) in finite tensor expressions of a countable generating family \(y_i\). Their span \(M'\) is countably generated and the image of \(M'\otimes S\to M\otimes S\) contains every \(y_i\), hence is onto. Right exactness gives \((M/M')\otimes S=0\), so faithful flatness makes \(M'=M\). The same finite-entry construction for a countably generated submodule \(Q\subset M\otimes S\) only asserts containment in the image of \(P\otimes S\); it never assumes that map injective. \texttt{OCC-00470} inserts ``to'' in this latter construction.

For the complete adapted-submodule lemma, let \(M\otimes S=\bigoplus_i Q_i\) be its original decomposition. The image of a countably generated \(N'_\ell\) is generated over \(S\) by a countable list, each with finite coordinate support. Thus the set \(I'_\ell\) of its nonzero coordinate projections is countable. The direct sum of those countably generated \(Q_i\) is countably generated, so the previous finite-entry construction gives \(P\) whose tensor image contains it. Set \(N'_{\ell+1}=N'_\ell+P\). Images into \(M\otimes S\) commute with this increasing union, by the finite expression of each tensor. Every such image is contained in the sum of its coordinate supports, and every \(Q_i\) in a stage support is contained in the next image. Hence the union image is exactly \(\bigoplus_{i\in\bigcup I'_\ell}Q_i\), as claimed. This proves the lemma for arbitrary ring maps while retaining its use of images. The next projectivity-descent theorem has not been adjudicated.

\subsubsection{Arbitrary intersections and base change of the smallest submodule}\label{algebra-proof-013-arbitrary-intersections-and-base-change-of-the-smallest-submodule}

This is a separate consequence of \texttt{lemma-minimal-contains} and \texttt{lemma-flat-ML-criterion}, with complete original-map proofs. For a flat module \(M\), the ML property is equivalent to preservation of arbitrary intersections of submodules of every \(F\): \[
\left(\bigcap_\alpha F_\alpha\right)\otimes_R M
=\bigcap_\alpha(F_\alpha\otimes_R M)
\quad\text{inside }F\otimes_RM.
\] For the forward implication, the kernel of \(F\to\prod_\alpha F/F_\alpha\) is the displayed intersection. Flatness identifies the kernel after tensoring with that intersection tensored by \(M\). The ML product-comparison injection shows this is also the kernel of the map into \(\prod_\alpha((F/F_\alpha)\otimes M)\), whose kernel is exactly the right-hand intersection by flatness and right exactness at each coordinate. Empty families give \(F\) on both sides. Conversely, the asserted intersection identity supplies the smallest submodule for any tensor in a finite free \(F\), by intersecting all submodules that contain it after tensoring. This family is nonempty because it contains \(F\). The criterion proved above makes \(M\) ML. In either direction the smallest submodule is finite: write its tensor element as a finite sum using its own entries, take their span, and apply minimality.

There is also an exact arbitrary-base-change formula for this smallest submodule. Write \(C_R(x)\subset F\) for it, for the original \(x\in F\otimes_RM\). For any ring map \(R\to S\), let \(x_S\) be its image in \(F_S\otimes_S M_S\). Then \[
C_S(x_S)=\operatorname{im}(C_R(x)\otimes_RS\longrightarrow F\otimes_RS).
\] Here the image is essential: no injection of \(C_R(x)\otimes S\) is assumed.

To prove this with original coefficients, use a finite free Lazard presentation \(M=\operatorname{colim}M_i\). Represent \(x\) by \(x_i\in F\otimes M_i\). Choose \(j\geq i\) for which \(f_{ij}\) dominates \(f_i\). In an original basis \(e_1,\ldots,e_r\) of \(M_j\), write its image exactly as \[
x_j=\sum_{a=1}^r y_a\otimes e_a,
\qquad C=\sum_{a=1}^r Ry_a\subset F.
\] Its image gives \(x\in C\otimes M\). If \(x\in F'\otimes M\), its image in \((F/F')\otimes M\) is zero. Domination with this test module makes the image of \(x_j\) zero in \((F/F')\otimes M_j\). Freeness forces each original \(y_a\in F'\), so \(C\subset F'\). Thus \(C=C_R(x)\).

After arbitrary base change, both stages remain finite free. For every \(S\)-module \(T\), the tensor-test maps for these stages are exactly the original maps with the underlying \(R\)-module \(T\), by associativity. Hence the same domination holds over \(S\). Also \(M_S\) is flat and ML, by the established ascent statements. Repeat the preceding coefficient argument in \(F_S\), using the original coefficients \(y_a\otimes1\) and the base-changed basis. It gives \(C_S(x_S)=\sum_a S(y_a\otimes1)\), which is precisely the stated image. This includes zero ranks, zero tensors and nonflat ring maps. It strengthens the original smallest-submodule statement without silently replacing it in a translation.

\subsubsection{Pure descent of countable generation and the resulting projective cases}\label{algebra-proof-013-pure-descent-of-countable-generation-and-the-resulting-projective-cases}

This is a separate propagation of groups 501 and 511 through the source results just checked. If \(R\to S\) is pure and \(M_S\) is countably generated, the original finite-entry construction gives countably generated \(M'\subset M\) with \((M/M')\otimes S=0\). Purity of the unit map of the \textbf{actual quotient} \(M/M'\) gives \(M/M'=0\). Thus countable generation descends along pure maps. More generally, the same proof descends generation by at most an infinite cardinal \(\kappa\): a union of \(\kappa\) finite sets has size at most \(\kappa\). This is not a claim of preservation of an exact finite number of generators.

If \(M_S\) is countably generated and projective, it is flat and ML over \(S\). Group 501's complete matrix-cokernel argument descends flatness to the original \(M\), and group 511, or the pushout reflection proof above, descends ML. Countable generation descends as just proved. The original countable-projectivity proof therefore makes \(M\) countably generated and projective. Conversely, countable generation and an actual projective splitting both ascend along every ring map, by their original tensor maps. Thus countably generated projectivity is equivalent before and after a pure ring map.

There is a second, explicitly bounded receiving case. Suppose the original \(R\)-module \(M\) is a direct sum of countably generated \(R\)-modules, and \(M_S\) is projective. The same two descents give flatness and ML over \(R\). With the given original decomposition, all three hypotheses of the source projectivity characterization hold, so \(M\) is projective. The reverse implication again ascends along any ring map. This argument retains that decomposition hypothesis; it does not claim descent of an arbitrary decomposition from \(S\) or unrestricted projectivity descent along every pure ring map. In particular, it does not try to apply the source pure-submodule theorem to \(M_S\) as an \(R\)-module, where flatness of \(S\) was not assumed.

\subsubsection{Propagation and reading limits}\label{algebra-proof-013-propagation-and-reading-limits}

Groups 512--526 account for all nineteen reports, including four retained canonical units and one rejected quotation. The cofinal-sequence correction includes the finite-index boundary missed by the report's proposed symbol-only edit. Group 527 records intersection preservation and the exact image formula for base change of the smallest submodule, with zero source operations. Group 528 records pure descent of countable generation and the two bounded projectivity consequences, receiving groups 501 and 511 and the original projectivity criterion. The earlier pure ML result is also propagated through the source's actual pushout argument, giving an independent proof with the same conclusion. The original Noetherian product hypotheses and the uncountable-projectivity counterexample remain intact.

The new disk query \texttt{coherent\ Mittag-Leffler} returned two research PDF-routing hits: \emph{Éléments de géométrie algébrique III, première partie} and \emph{Basic Oka Theory in Several Complex Variables}. Both remain unread and unused here; no PDF fallback or new EGA coverage is claimed. The exact supplied original Stacks TeX provides the current primary treatment. Earlier topic snapshots and external-reading records remain historical evidence. The next theorem's statement is only lookahead, and broader combined-results synthesis remains after core completion.

\hypertarget{algebra_projectivity_completion_notes_20260927}{}
\subsection{Projectivity descent and completion}\label{algebra-proof-014-projectivity-descent-and-completion}

Separate editorial evidence for the Stacks Project authors' algebra.tex, authority commit a04446e57ec1fbc252a871afcec7752fb2807b14, lines 22461--23009. Authority SHA-256: FA8BB92E58A4F78A2BD01B3B6A4A87DE0A0D279F5DD90641B574DD5FBFFFA4F3. The complete source units in this interval and all eleven due reports were read. The theorem beginning at 23010 remains unadjudicated.

The source's Matlis citation, Poonen email attribution, and Lenstra--de Smit attribution are retained without claiming independent reading of those works or correspondence. The five operations already composed as MC-STK-ERR-0455 and MC-STK-ERR-0456 resolve four reports. The proposed rewriting of induction on the well-ordered set M is rejected. Seven further minimal copyedits or missing declarations are recorded. Full arguments and the consequences below remain separate editorial material; no cumulative chapter or translation source is changed. These are source-review consequences, without a novelty claim or completion of the later synthesis.

\subsubsection{Faithfully flat projectivity descent and the retained corrections}\label{algebra-proof-014-faithfully-flat-projectivity-descent-and-the-retained-corrections}

Source 22461--22515. Reports OCC-12112 and OCC-00471 retain MC-STK-ERR-0455; OCC-12113 and OCC-00473 retain the four ordinal-symbol operations of MC-STK-ERR-0456. Report OCC-00472 is rejected as a mathematical objection.

Keep the original faithfully flat map \(R\to S\), the module \(M\), and the decomposition \[
M\otimes_R S=\bigoplus_{i\in I}Q_i
\] into countably generated projective summands. Tensoring an inclusion \(M_\alpha\subset M\) with the flat module \(S\) identifies its image as an actual submodule. Write that image as a coordinate sum \(A_\alpha=\bigoplus_{i\in J_\alpha}Q_i\). Zero summands may be included consistently in every \(J_\alpha\); they affect none of the images or arguments. At a limit stage, the original union commutes with tensor product, and every element has finite support. The resulting submodule is precisely the sum over the union of the preceding coordinate subsets.

At a successor stage with \(M_\alpha\ne M\), retain the least missing element \(x\) in the fixed well-order. Flatness and the quotient map identify \[
(M/M_\alpha)\otimes_R S
 \cong (M\otimes_R S)/A_\alpha
 \cong\bigoplus_{i\notin J_\alpha}Q_i .
\] The adapted-submodule lemma gives a countably generated \(P\subset M/M_\alpha\) containing the image of \(x\), whose tensor image is a coordinate sum \(B\) in this displayed quotient. Flatness makes \(P\otimes_R S\to(M/M_\alpha)\otimes_R S\) injective. This is the map repaired by MC-STK-ERR-0455. There is no previously chosen section to \(M\), so the old codomain \(M\otimes_R S\) was not justified.

The module \(P\otimes_R S\) is projective because it is a coordinate summand. The preceding countable projectivity descent lemma gives projectivity of \(P\). Let \(M_{\alpha+1}\) be its original preimage in \(M\). Tensoring its kernel description with the flat module \(S\) shows that its image in \(M\otimes_R S\) is exactly the preimage of \(B\), namely \(A_\alpha\oplus B\). Thus the coordinate-sum property is maintained, not merely an abstract isomorphism to a sum. The quotient \(M_{\alpha+1}/M_\alpha=P\) is projective, so its exact sequence splits. If \(M_\alpha=M\), use the stationary step and zero quotient.

The source explicitly chooses a well-order on the underlying set \(M\). Therefore induction on that well-ordered set is legitimate. To verify exhaustion and its precise indexing, let \(\lambda\) be its order type and write its elements \(x_\beta\), \(\beta<\lambda\). We claim that \(x_\beta\in M_\alpha\) for every \(\beta<\alpha\leq\lambda\). The assertion at zero is empty. At a successor, all earlier elements are already present. If \(x_\alpha\) is missing it is the least missing element and is inserted; if it is present it remains present. At a limit, each earlier element belongs to an earlier stage and hence to the union. This proves the claim and \(M_\lambda=M\). Choosing the fresh ordinal \(\gamma=\lambda+1\) retains the terminal stage and supplies the index required in the source. The existing correction renames all four ordinal occurrences and leaves the ring \(S\) intact.

For completeness, choose each actual quotient section \(s_\alpha:M_{\alpha+1}/M_\alpha\to M_{\alpha+1}\), for \(\alpha<\lambda\). The resulting map from the direct sum of these quotients to \(M\) sends a finite tuple to the sum of its images under the sections. A zero image has zero component in its largest successor quotient, then in the next largest, and hence every component is zero. Surjectivity follows by induction: successor stages add that quotient, and limit stages are unions. Thus the original \(M\), with its actual submodules, is isomorphic to the displayed sum of projective quotients. This also checks the terminal-stage correction from group 503 and the least-missing-element argument from group 504. No unrestricted descent of an arbitrary \(S\)-decomposition along a merely pure map is asserted; the preimage calculation here uses flatness.

\subsubsection{Completion exactness and the original finite-generation hypothesis}\label{algebra-proof-014-completion-exactness-and-the-original-finite-generation-hypothesis}

Source 22528--22699. The completion of \(M\) is the original inverse limit \(\widehat M=\lim_{n\geq1}M/I^nM\), with its canonical unit and tensor comparison. All inverse systems in this section are indexed by the positive integers, so the countable lifting arguments from group 510 apply.

Suppose \(\varphi:M\to N\) is onto modulo \(I\). Its cokernel modulo \(I^n\) equals its own multiple by \(I/I^n\); iterating \(n\) times gives zero. This proves surjectivity \(M/I^nM\to N/I^nN\), without assuming either module is finite. Let \(K_n=\{x:\varphi(x)\in I^nN\}\). For \(x\in K_n\), write the finite expression \(\varphi(x)=\sum_jz_jn_j\), with \(z_j\in I^n\), and use the modulo-\(I\) surjectivity to write \(n_j=\varphi(m_j)+\sum_kz_{jk}n_{jk}\), with \(z_{jk}\in I\). Then \[
x'=x-\sum_jz_jm_j\in K_{n+1},
\qquad x'-x\in I^nM .
\] The kernels \(K_n/I^nM\) therefore have surjective transition maps. Compatible lifts are constructed recursively by subtracting lifts from these kernels. This proves surjectivity on completions. If \(0\to K\to M\to N\to0\) is exact and \(N\) is flat, tensoring with \(R/I^n\) is exact also on the left. The kernels are then \(K/I^nK\), whose transitions are onto, proving the source's completed short exact sequence. If \(M\) is finite, completing its finite free surjection gives the asserted surjection \(M\otimes_R\widehat R\to\widehat M\).

For a finitely generated ideal \(I\), fix \(n\geq1\) and actual generators \(f_1,\ldots,f_r\) of \(I^n\). The completion of the surjection \(M^{\oplus r}\to I^nM\) is onto. The internal filtration on \(I^nM\) is \(I^k(I^nM)=I^{n+k}M\). Consequently its completion is the same inverse limit over the tail \(m\geq n\) of \(I^nM/I^mM\). Each \(M/I^mM\to M/I^nM\) is onto, with that kernel. Thus this tail limit is precisely the kernel of \(\widehat M\to M/I^nM\), embedded through the original inclusion. The image of the completed generator map is exactly \(I^n\widehat M\). This proves \[
I^n\widehat M=(I^nM)^\wedge
 =\ker(\widehat M\to M/I^nM).
\] The induced isomorphisms \(\widehat M/I^n\widehat M\to M/I^nM\) carry the canonical unit to the identity compatible family; passing to their limits proves completeness of \(\widehat M\) through that unit. Empty generating lists for \(I^n=0\) give the zero kernel and the same conclusion.

For the torsion-quotient lemma, keep \(I^cQ=0\) in \(0\to M\to N\to Q\to0\). Then \(I^cN\subset M\), and for \(n\geq c\) \[
I^nM\subset M\cap I^nN\subset I^{n-c}M.
\] The natural map from the \(I\)-adic completion of \(M\) to \(\lim M/(M\cap I^nN)\) reduces each coordinate using the first inclusion. The inverse obtains its coordinate modulo \(I^aM\) from coordinate \(a+c\) using the second inclusion. Compatibility makes both composites the identity. The systems of kernels \(M/(M\cap I^nN)\) have onto transitions. Applying the same recursive lifting to the quotient system, which is constantly \(Q\) for \(n\geq c\), gives the asserted exact sequence. If \(c=0\), \(Q=0\) and the same formulas remain valid.

\subsubsection{The exact split defect of a second completion}\label{algebra-proof-014-the-exact-split-defect-of-a-second-completion}

Source 22701--22731; propagation from group 510. Retain the original ideals and modules, with no finite-generation assumption on \(I\). Set \[
K_n=\ker(\widehat M\to M/I^nM),\qquad
D_n=K_n/I^n\widehat M,\qquad D=\lim_nD_n .
\] These are the source's defect quotients, not replacements for its completion. The inclusion \(I^n\widehat M\subset K_n\) follows coordinatewise. The projection \(\widehat M\to M/I^nM\) is onto: choose a representative in \(M\) and its canonical compatible family. This gives the source's levelwise short exact sequences.

For \(k\in K_n\), its next coordinate belongs to \(I^nM/I^{n+1}M\). Write a representative as \(\sum_j a_jm_j\) with \(a_j\in I^n\), and set \(u=\sum_j a_j\eta_M(m_j)\in I^n\widehat M\). Then \(k-u\in K_{n+1}\). This proves \(K_{n+1}+I^n\widehat M=K_n\) and surjectivity of \(D_{n+1}\to D_n\). Taking limits gives exact maps \[
0\longrightarrow D\xrightarrow{\iota}
  \widehat{\widehat M}\xrightarrow{\pi}\widehat M
  \longrightarrow0.
\] Here \(\iota\) is the limit of the actual inclusions, and \(\pi\) sends a family of classes in \(\widehat M/I^n\widehat M\) to its images in \(M/I^nM\). The canonical completion map \(\eta_{\widehat M}:\widehat M\to\widehat{\widehat M}\) satisfies \(\pi\eta_{\widehat M}=\mathrm{id}_{\widehat M}\) on every original coordinate. The exact sequence therefore splits canonically as \(R\)-modules, with maps \[
\widehat M\oplus D\longrightarrow\widehat{\widehat M},
\quad (m,d)\longmapsto\eta_{\widehat M}(m)+\iota(d),
\] \[
t\longmapsto
 \bigl(\pi(t),\ \iota^{-1}(t-\eta_{\widehat M}\pi(t))\bigr).
\] The second component is defined because its argument has zero \(\pi\)-image. The two displayed composites are identities by exactness and the section identity. Every construction commutes with a map of the original modules. This records an explicit structural consequence of the source proof, without a novelty claim.

Because the transitions of \(D_n\) are onto, any specified class at any level has a compatible family of lifts: lift to higher levels recursively and reduce to lower levels. Thus \(D=0\) if and only if every \(D_n=0\). It follows that the canonical map \(\eta_{\widehat M}\) is an isomorphism if and only if every \(K_n=I^n\widehat M\), exactly the source criterion. Also \(\bigcap_n I^n\widehat M=0\), since it is contained in \(\bigcap_nK_n=0\). The defect measures failure of completeness, not failure of this separatedness. The finitely generated case in the preceding section annihilates each \(D_n\), so all the maps agree with that case. No infinite-generation case is silently identified with it.

\subsubsection{Units and the principal-generator completion argument}\label{algebra-proof-014-units-and-the-principal-generator-completion-argument}

Source 22733--22822. An element \(x=(x_n)\in\widehat R\) invertible modulo \(I\) is invertible at every level: \(I/I^n\) is nilpotent, and the finite geometric sum in a nilpotent error gives an inverse lift. Uniqueness of the inverse ensures that the \(x_n^{-1}\) form a compatible family. This proves the first unit assertion and the two assertions about \(1+I\) and \(1+I\widehat R\). If \(a\) belongs to the kernel of \(\widehat R\to R/I\), then \(1-ba\) maps to \(1\) for every \(b\in\widehat R\), so it is a unit. This is the Jacobson radical criterion. It applies also to \(I\widehat R\) but does not identify that ideal with the kernel when \(I\) is arbitrary.

Reports OCC-00474 and OCC-00475 correctly request the missing declarations that \(M\) is an \(A\)-module. These are the only new source operations in the two principal-generator lemmas.

There is also an unreported boundary in the first proof: for \(r=0\), the ideal \(I\) and each empty-generated \(J_n\) are zero, while the displayed \(I^{rn}=I^0=A\) need not be contained in \(J_n\). The theorem remains true, since \(M\to\lim M/(0)M\) is the identity compatible-family map. This separate editorial verification handles the boundary without imposing \(I\ne0\), changing the statement, or substituting a new proof for the translation.

For \(r>0\), set \(J_n=(f_1^n,\ldots,f_r^n)\). Retain both inclusions \(I^n\supset J_n\supset I^{rn}\). The second follows because a monomial of degree \(rn\) in \(r\) generators has an exponent at least \(n\). Thus the systems are cofinal with the displayed exponents. Given the source element \(\xi\in\lim M/J_nM\), choose representatives \(y_n\in M\) of its coordinates for \(n\geq1\), and put \(y_0=0\). Compatibility gives \[
y_{n+1}-y_n=\sum_{i=1}^r f_i^n x_{n,i}\quad(n\geq0).
\] For \(n=0\), this uses \(J_0=A\), since \(r>0\), so it also retains the initial term of the source series. Each individual series \(\sum_{n\geq0}f_i^nx_{n,i}\) defines an element of the \(f_i\)-adic completion. By hypothesis choose \(x_i\in M\) mapping to it. Modulo \(J_NM\), the difference between \(x_i\) and its first \(N\) terms lies in \(f_i^NM\). Therefore \(\sum_i x_i\) agrees with \(y_N\) modulo \(J_NM\), for every \(N\). Its image is \(\xi\). This proves the asserted surjectivity with the original coefficients and initial term.

For the subideal lemma, \(I^nM\subset J^nM\) gives separatedness. The preceding surjectivity argument reduces to \(I=(f)\), including \(f=0\). Keep the original sequence \(x_n\) and its \(J\)-adic limit \(x\), and put \(u_n=x_n-x\). Then \(u_n\in J^nM\), and write \(u_n-u_{n+1}=f^nz_n\). The partial sums \[
v_{n,c}=\sum_{k=0}^{c-1}f^kz_{n+k}
\] form a \(J\)-adic Cauchy sequence, hence have a limit \(v_n\) in \(M\). Multiplication by \(f^n\) preserves the \(J\)-adic congruences, while \[
f^nv_{n,c}=u_n-u_{n+c}.
\] The second term tends to zero. Separatedness therefore gives \(u_n=f^nv_n\). This proves the required \(f\)-adic congruences for the original limit \(x\), without discarding it during the temporary subtraction.

For the change-of-ideal lemma retain \(I^c\subset J\), \(J^d\subset I\), \(c,d>0\). The exact inclusions \(I^n\subset J^{\lfloor n/c\rfloor}\) and \(J^m\subset I^{\lfloor m/d\rfloor}\) give the printed maps. For example, a coordinate modulo \(J^aM\) of the first limit map can be taken from coordinate \(ca\) of the \(I\)-system. Its inverse obtains coordinate \(b\) from coordinate \(db\). Composing uses coordinate \(cdb\), which reduces to the original coordinate \(b\) by compatibility. The other composite is identical in the opposite order. Floor value zero gives the zero quotient by \(I^0M=M\) or \(J^0M=M\), so it causes no boundary defect.

\subsubsection{Quotients and finite modules over a complete ring}\label{algebra-proof-014-quotients-and-finite-modules-over-a-complete-ring}

Source 22824--22883. For complete \(M\) and \(K\subset M\), the composite \(M\to M/K\to\widehat{M/K}\) is onto by completion surjectivity. Its kernel is \[
\overline K=\bigcap_{n\geq1}(K+I^nM).
\] Consequently the actual map identifies \(\widehat{M/K}\) with \(M/\overline K\), and the unit on \(M/K\) is an isomorphism precisely when \(\overline K=K\), the source claim. The submodule \(\overline K\) is closed: for every \(n\), \(\overline K+I^nM\subset K+I^nM\), whence taking intersections gives closure of \(\overline K\) contained in \(\overline K\); the reverse containment is immediate. Applying the proved equivalence shows \(M/\overline K\) is complete even without finite generation of \(I\). No claim about all arbitrary module completions follows from this quotient-of-a-complete-module case.

When \(R\) is complete and \(M\) finite, the canonical map \(M\otimes_R\widehat R\to\widehat M\) is onto and its source is canonically \(M\). Its kernel is exactly \(\bigcap_n I^nM\). Thus the source's additional separatedness hypothesis makes it an isomorphism. For the next lemma choose the original lifts \(x_1,\ldots,x_n\) and their span \(M'\subset M\). Since \(I^aM'\subset I^aM\), the submodule \(M'\) is separated. It is finite and hence complete by the just-proved lemma. Completion of \(M'\to M\) is onto because it is onto modulo \(I\). Hence every element of \(M\) has the same image in \(\widehat M\) as some element of \(M'\); injectivity of \(M\to\widehat M\) proves it is that same element. Therefore \(M=M'\), and both are complete. For \(n=0\), this argument says \(M'=0\) and proves \(M=0\), so the zero-generator case is retained. Reports OCC-00476 and OCC-00477 concern only the article and the missing ``by''.

\subsubsection{Noetherian exactness and the tensor comparison}\label{algebra-proof-014-noetherian-exactness-and-the-tensor-comparison}

Source 22900--23009. Reports OCC-00478 and OCC-00479 request the missing ``by'' in the completion notation. The identical naming omission at 22978 is independently propagated as one further copyedit.

For an exact sequence \(0\to K\to N\to M\to0\) of finite modules over the original Noetherian ring \(R\), the quotient sequence has left term \(K/(I^nN\cap K)\). Its transitions are onto, so the original inverse limit sequence is exact. Artin--Rees supplies the actual bound \[
I^nK\subset I^nN\cap K\subset I^{n-c}K\quad(n\geq c).
\] The natural map from \(\widehat K\) to the limit of the induced quotients and its inverse using coordinate \(a+c\) for coordinate \(a\) are the same explicit maps verified above. Thus completion is exact on these finite modules.

For the tensor comparison choose the source's presentation \(0\to K\to R^{\oplus t}\to M\to0\). Noetherianity makes \(K\) finite. The tensor row is right exact; no injection of \(K\otimes_R\widehat R\) is assumed. The completed row is exact, and the comparison for \(K\) is onto while the comparison for \(R^{\oplus t}\) is an isomorphism. To prove injectivity of the comparison for \(M\), lift an element in its kernel to \(R^{\oplus t}\otimes_R\widehat R\). Its completed image lifts from \(\widehat K\), which in turn lifts from \(K\otimes_R\widehat R\). Subtracting this latter image gives zero under the middle isomorphism and hence zero. The original element is therefore zero. Surjectivity was already proved. This is the source's diagram argument without presupposing flatness of \(\widehat R\).

Every ideal \(J\subset R\) is finite, so \(J\otimes_R\widehat R\to\widehat R\) identifies with the now-proved injection \(\widehat J\to\widehat R\). The ideal criterion for flatness gives flatness of \(R\to\widehat R\). The source's exactness assertion is exactly the preceding finite-module result. Finally, completeness of \(\widehat M\) for arbitrary \(M\) follows from the already-proved finitely generated ideal lemma, because \(I\) is finite in a Noetherian ring. This last statement is explicitly identified by the source as a special case; it is not a newly discovered extension to all rings.

\subsubsection{Faithful flatness criterion and projectivity after completion}\label{algebra-proof-014-faithful-flatness-criterion-and-projectivity-after-completion}

Source 22733--22753 and 22955--23008, receiving the exact projectivity descent proof at 22461--22515. For Noetherian \(R\), \[
R\longrightarrow\widehat R\text{ is faithfully flat}
\quad\Longleftrightarrow\quad I\subset\operatorname{Jac}(R).
\] The forward implication is an additional consequence; the reverse is the source lemma.

First prove necessity, which does not require Noetherianity once faithful flatness is assumed. A faithfully flat map reflects units: if the image of \(a\) is invertible, then \((R/aR)\otimes_R\widehat R=0\), so faithfulness gives \(R/aR=0\), and \(a\) is a unit. For \(i\in I\) and \(r\in R\), the image of \(1-ri\) is a unit by the original completion unit lemma. Thus \(1-ri\) is a unit in \(R\) for every \(r\), proving \(i\in\operatorname{Jac}(R)\).

Conversely, Noetherianity supplies flatness as proved above. If \(I\subset\operatorname{Jac}(R)\), every maximal ideal \(\mathfrak m\) contains \(I\). The composite \(\widehat R\to R/I\to R/\mathfrak m\) is onto and its kernel contracts to \(\mathfrak m\). To check faithfulness directly, take a nonzero element \(x\) of any nonzero \(R\)-module \(X\), set \(\mathfrak a=\operatorname{Ann}_R(x)\), and choose a maximal ideal \(\mathfrak m\) containing the proper ideal \(\mathfrak a\). Flatness preserves the injection \(R/\mathfrak a\to X\) after tensoring. Its source after tensoring is \(\widehat R/\mathfrak a\widehat R\), which is nonzero because it maps onto \(R/\mathfrak m\). Hence \(X\otimes_R\widehat R\ne0\). For the zero ring there is no nonzero \(X\), and the identity map satisfies the assertion.

Under these equivalent conditions, any \(R\)-module \(M\) satisfies \[
M\text{ is projective over }R
\quad\Longleftrightarrow\quad
M\otimes_R\widehat R\text{ is projective over }\widehat R .
\] A splitting of a free module remains a splitting after tensoring, proving ascent; the exact faithfully flat theorem above proves descent. If \(M\) is finite, the proved tensor comparison identifies the right-hand module with its actual completion. Thus a finite \(R\)-module is finite projective if and only if its completion is projective over \(\widehat R\); the completion is finite because it is the base change of the finite module. For arbitrary \(M\), the tensor comparison need not be an isomorphism, and no replacement of \(M\otimes_R\widehat R\) by \(\widehat M\) is made.

These consequences connect the existing projectivity review to completion while retaining every finite-generation and faithful-flatness boundary. They remain editorial consequences and will be compared with the rest of the work as source-ordered review continues. They do not replace the source's statement or initiate the later combined-results paper.

\hypertarget{algebra_graded_completion_notes_20260927}{}
\subsection{Noetherian completion and graded inverse systems}\label{algebra-proof-015-noetherian-completion-and-graded-inverse-systems}

Separate editorial evidence for the Stacks Project authors' algebra.tex, authority commit a04446e57ec1fbc252a871afcec7752fb2807b14, lines 23010--23414. Authority SHA-256: FA8BB92E58A4F78A2BD01B3B6A4A87DE0A0D279F5DD90641B574DD5FBFFFA4F3. All complete source units in this interval and all five received reports were read. The following flatness section was read only through line 23452; its first proof is incomplete and unadjudicated.

The current interval has no previously composed canonical correction. Five received reports and one independently found reference correction yield nine minimal operations in six groups. The graded stabilization repair also handles zero generators and avoids nonpositive indices. The complete proofs and three further consequences below remain separate editorial material. The original hypotheses and exposition stay identifiable; there is no chapter or translation mutation and no novelty claim. The bounded corpus lookup records all four routing hits, including the local-work hits, without claiming content reading or using PDFs.

\subsubsection{Associated-graded generation and the completion map}\label{algebra-proof-015-associated-graded-generation-and-the-completion-map}

Source 23010--23087; OCC-00480. Retain the original ideal \(I\subset R\) with \(R/I\) Noetherian and \(I\) finitely generated. Write \(B=\widehat R\) and \(L=IB\). The already proved canonical comparison gives \(B/L=R/I\), \(L^n=I^nB\), and completeness of \(B\). Thus the source's temporary replacement has these exact maps; it changes none of the original conclusion about \(\widehat R\).

For its argument in a complete ring, keep the source notation \(R,I,J,f_i,g_j,d_j\). The polynomial map \(R/I[T_1,\ldots,T_t]\to\bigoplus_{n\geq0}I^n/I^{n+1}\) sends \(T_i\) to the actual class of \(f_i\). Every degree-\(n\) ideal product is a finite sum of degree-\(n\) monomials in these generators with coefficients in \(R\); reducing the coefficients modulo \(I\) proves surjectivity. The associated graded ring is therefore Noetherian. The submodule \[
\bigoplus_{n\geq0}(J\cap I^n)/(J\cap I^{n+1})
\] is a homogeneous ideal of that ring: multiplication by a class in \(I^a/I^{a+1}\) sends its degree-\(n\) part into degree \(n+a\). Choose its finitely many homogeneous generators \(\overline g_j\) and original lifts \(g_j\in J\cap I^{d_j}\). In degree \(n\) its exact expression is \[
(J\cap I^n)/(J\cap I^{n+1})
 =\sum_{j:\,d_j\leq n}
    \overline g_j\,(I^{n-d_j}/I^{n-d_j+1}).
\] The restriction is necessary; no negative ideal power is defined. This is the report's minimal correction. For \(d_j>n\) choose the corresponding coefficient zero, which satisfies the source's stated bound \(I^{\max(0,n-d_j)}\).

Starting with \(x\in J\), choose coefficients \(a_{j,n}\) inductively to remove the degree-\(n\) residual. The residual remains in \(J\cap I^{n+1}\), and \(a_{j,n}\in I^{\max(0,n-d_j)}\). For each fixed \(j\), the partial sums form a Cauchy sequence: terms with \(n\geq b+d_j\) lie in \(I^b\). Completeness gives \(A_j=\sum_{n\geq0}a_{j,n}\). Multiplication by the fixed \(g_j\) preserves all \(I\)-adic congruences. Taking limits in the source's finite sum over \(j\), with residual in \(I^{N+1}\), gives \(x=\sum_jA_jg_j\). Separatedness gives actual equality. If the graded ideal has no generators, each \(x\in J\) lies in every \(I^n\), hence \(x=0\); the empty case is included. This proves finite generation of every \(J\) and hence the original Noetherian-completion assertion.

The alternative formal-power-series map in 23079 also has its stated meaning. At level \(R/I^n\), evaluate the terms of total degree \(<n\) in \(R[[x_1,\ldots,x_t]]\) at the original \(f_i\); every higher-degree monomial vanishes. There are finitely many monomials below each total degree. These compatible ring maps give the printed map into \(\widehat R\). For \(\xi\in\widehat R\), choose \(y_n\in R\) representing its \(n\)-th coordinate. Set \(y_0=0\). The initial \(y_1\) is a constant coefficient; for \(n\geq1\), \(y_{n+1}-y_n\in I^n\) is a finite linear combination of the degree-\(n\) monomials \(f^\alpha\). Use the same coefficients on \(x^\alpha\). The resulting power series evaluates to \(\xi\) at every level. With no generators \(I=0\), the map is the identity \(R\to R\). This verifies the omitted detail as editorial evidence, without making a rewritten source proof a requirement.

\subsubsection{Local completion and the redundant maximal-ideal assumption}\label{algebra-proof-015-local-completion-and-the-redundant-maximal-ideal-assumption}

Source 23089--23132; OCC-00481 requests only ``Hence it has dimension 0''. Keep the local homomorphism \(R\to S\), maximal ideals \(\mathfrak m,\mathfrak n\), and \(k=R/\mathfrak m\). Let \[
C=S/\mathfrak mS,\qquad
\mathfrak r=\mathfrak n/\mathfrak mS,\qquad
d=\dim_k C<\infty .
\] This finite-dimensional local algebra is Artinian. Its maximal ideal is nilpotent: its descending powers stabilize because they are subspaces of a finite-dimensional space; at a stable power \(V=\mathfrak rV\), the ordinary Nakayama lemma over the local Artinian ring \(C\), applied to the finite module \(V\), gives \(V=0\). Until zero each step is a strict inclusion. Therefore \(\mathfrak r^d=0\), and \[
\mathfrak n^d\subset\mathfrak mS\subset\mathfrak n.
\] Here \(d\geq1\) for the nonzero local rings in the source convention. Consequently the two adic filtrations on \(S\) are cofinal with the actual bound \(\mathfrak n^{dn}\subset\mathfrak m^nS\subset\mathfrak n^n\). Coordinate reduction from levels \(dn\) and \(n\) gives mutually inverse canonical maps between the two completions, as in the already proved change-of-ideal lemma.

Only finite generation of \(\mathfrak m\) is needed as a separate hypothesis. To prove this, let \(e=\dim_k(S/\mathfrak n)\). The vector space \(\mathfrak r\) has dimension \(d-e\). Lift a basis \(z_1,\ldots,z_{d-e}\) to \(\mathfrak n\), and let \(a_1,\ldots,a_h\) generate \(\mathfrak m\). Every element of \(\mathfrak n\) differs from an \(R\)-linear combination of the \(z_j\) by an element of \(\mathfrak mS\): lift its basis coefficients from \(k\) to \(R\). Therefore \[
\mathfrak n=(a_1,\ldots,a_h,z_1,\ldots,z_{d-e})S .
\] The source's separate finite-generation assumption on \(\mathfrak n\) follows, with this explicit bound. No Noetherianity of \(R\) or \(S\) is used.

The generator hypothesis on \(\mathfrak m\) gives completeness of \(\widehat R\) and of \(\widehat S=\lim S/\mathfrak m^nS\) as modules for the original \(\mathfrak m\)-action, together with \[
\widehat R/\mathfrak m\widehat R=k,\qquad
\widehat S/\mathfrak m\widehat S=C.
\] The canonical action of \(\widehat R\) on \(\widehat S\) is defined at each quotient level. Powers of \(\mathfrak m\widehat R\) acting on \(\widehat S\) are precisely \(\mathfrak m^n\widehat S\), so the separatedness and completeness required for the finite-module lemma hold for that ring and ideal. Choose \(d\) actual lifts of a \(k\)-basis of \(C\). The source's finite-over-complete-ring argument proves that these elements generate \(\widehat S\) over \(\widehat R\). Their number is minimal: any generating family maps to a spanning family of the \(d\)-dimensional quotient \(C\). Also \(\widehat R\) is local, because an element outside the kernel of its map to \(k\) is a unit by the completion unit lemma. The preceding Noetherian-completion theorem applies to \((R,\mathfrak m)\), so \(\widehat R\) is Noetherian even when \(R\) is not. Its finite algebra \(\widehat S\) is then Noetherian. These are proved editorial consequences, not replacement hypotheses silently inserted in the translated source.

\subsubsection{The finite-extension product with its localization maps}\label{algebra-proof-015-the-finite-extension-product-with-its-localization-maps}

Source 23134--23174. Retain \(R\to S\) finite, \(R\) Noetherian, and the original prime \(\mathfrak p\). Put \(T=R_{\mathfrak p}\), \(U=S\otimes_R T\), and \(\mathfrak a=\mathfrak pT\). The canonical tensor isomorphism \[
\widehat T\otimes_R S
 \longrightarrow\widehat T\otimes_T U
\] sends \(b\otimes s\) to \(b\otimes(s\otimes1)\); its inverse sends \(b\otimes(s\otimes t)\) to \(tb\otimes s\). Both maps respect the balancing relations and their composites fix every simple tensor. The finite-module completion comparison now identifies the second tensor with the \(\mathfrak a\)-adic completion of \(U\).

The primes \(\mathfrak q_i\) of \(S\) over \(\mathfrak p\) correspond to the maximal ideals of \(U\) over \(\mathfrak a\). The actual localizations identify \(U_{\mathfrak q_iU}\) with \(S_{\mathfrak q_i}\): writing a fraction in \(U\) as \(s/u\), the map sends \((s/u)/(t/v)\) to \(sv/(ut)\); all denominators lie outside \(\mathfrak q_i\), since the original \(u,v\) lie in \(R\setminus\mathfrak p\). Its inverse sends \(s/t\) to \((s/1)/(t/1)\). The fraction equivalence relations give inverse ring maps.

For each \(n\geq1\), \(U/\mathfrak a^nU\) is a finite module over the Artinian local ring \(T/\mathfrak a^n\), hence is Artinian. Its maximal ideals are exactly the images of the \(\mathfrak q_i\). The already proved Artinian-ring product theorem identifies this ring, by its canonical localization map, with \[
\prod_i U_{\mathfrak q_iU}/\mathfrak a^nU_{\mathfrak q_iU}.
\] The maps commute with reduction from \(n+1\) to \(n\). Taking the limit and commuting a finite product with the compatible-family construction gives the product of the \(\mathfrak a\)-adic completions. Each local factor has finite-dimensional residue quotient over \(T/\mathfrak a\): it is the localization of the finite-dimensional algebra \(U/\mathfrak aU\). The preceding local lemma identifies this completion with its maximal-ideal completion. These exact maps produce the source's whole displayed comparison.

If no prime lies over \(\mathfrak p\), the finite-dimensional fibre \(U/\mathfrak aU\) has no prime and is the zero ring. Nakayama applied to the finite \(T\)-module \(U\) gives \(U=0\). The tensor, completion and empty product are then all the zero ring. No nonempty-fibre assumption has been inserted.

\subsubsection{Compatible splittings and successive completion}\label{algebra-proof-015-compatible-splittings-and-successive-completion}

Source 23176--23237. In the split-completed-sequence lemma retain \(0\to K\to P\to M\to0\), flatness of \(M\), and projectivity of \(M/IM\). Flatness gives the original levelwise short exact sequences. For every \(n\geq1\), \(M/I^nM\) is flat over \(R/I^n\), its reduction modulo the nilpotent ideal \(I/I^n\) is the given projective module, and the source's nilpotent projectivity-lifting lemma at 18923--18958 proves its projectivity.

Choose the original \(s_1\). Given \(s_n\), use exactly the source's submodule \[
P_{n+1}=\{p\in P:p\bmod I^nP\in\operatorname{im}s_n\}.
\] It maps onto \(M\). Indeed, lift \(m\in M\) to \(p\in P\); the difference between \(p\bmod I^nP\) and \(s_n(m\bmod I^nM)\) is in the kernel \(K/I^nK\). Subtract a lift \(k\in K\). Then \(p-k\in P_{n+1}\) and still maps to \(m\). Hence the printed quotient map \(P_{n+1}/I^{n+1}P_{n+1}\to M/I^{n+1}M\) is onto. Projectivity gives its section \(t\). After mapping to \(P/I^{n+1}P\), the image reduces modulo \(I^n\) into \(\operatorname{im}s_n\); its image in \(M/I^nM\) is the original argument, so that reduction is exactly \(s_n\). Thus the \(s_n\) are compatible. Their limit is a section of the completed quotient map, linear even over \(\widehat R\) by checking its \(R/I^n\)-linear coordinates. The source's exactness lemma supplies the left term \(\widehat K\). No choice of a noncompatible family of splittings has been used.

For the following lemma, retain the original Noetherian \(A\), ideals \(I,J\), and \(B=\varprojlim A/(I+J)^n\). The canonical finite-module tensor-completion comparison, applied to \(A/I\), gives \[
B/IB
 \cong (A/I)\otimes_A B
 \cong \lim_n A/(I+(I+J)^n)
 =\lim_n A/(I+J^n).
\] The last equality follows because every mixed monomial in \(I+J\) lies in \(I\), leaving \(J^n\) modulo \(I\). This is precisely the \(J\)-adic completion of \(A/I\). The source cites lemma-completion-flat at 23224; the direct supporting statement is lemma-completion-tensor. The reference correction is recorded independently of the received grammatical report OCC-00482.

By hypothesis the last limit is canonically \(A/I\). The module \(B\) is \((I+J)\)-adically complete and hence \(I\)-adically complete by the finite-subideal lemma. As \(A\) is \(I\)-adically complete, the finite-over-complete-ring lemma makes \(B\) finite over \(A\). The quotient \(B/\operatorname{im}(A)\) is finite and is zero modulo \(I\). The completion unit lemma places \(I\) in \(\operatorname{Jac}(A)\), so Nakayama gives surjectivity \(A\to B\). Completion is flat here. Every maximal ideal of \(A\) contains \(I\), and the isomorphism \(B/IB=A/I\) supplies a prime above it. The faithful-flatness criterion therefore gives faithful flatness, and the canonical unit map is injective. Thus \(A\to B\) is an isomorphism. Its unit is the completion map, so \(A\) is \((I+J)\)-adically complete, and the finite-subideal lemma gives \(J\)-adic completeness. The title's reference to nilpotence is not a new hypothesis.

\subsubsection{Limits and their canonical completeness maps}\label{algebra-proof-015-limits-and-their-canonical-completeness-maps}

Source 23253--23297; receiving group 539 and its exact section. Let \(M=\lim M_n\), with \(I^nM_n=0\) and \(I\) finitely generated. The coordinate maps factor through \(M/I^nM\). Their inverse limit gives \(\pi:\widehat M\to M\), and \(\pi\eta_M=\mathrm{id}_M\) on every coordinate. Hence \(M\) is a direct summand of \(\widehat M\), through these specific maps. The latter is complete by finite generation of \(I\). A direct summand of a complete module is complete: for a decomposition \(X=Y\oplus Z\), one has \(I^nX=I^nY\oplus I^nZ\), and both the quotients and their limits split into these two factors; the canonical unit on \(X\) is the direct sum of the two units. If it is an isomorphism, both component units are isomorphisms. This proves the source lemma without identifying abstractly isomorphic modules in place of the canonical units.

In the next lemma the transition maps identify \(M_n=M_{n+1}/I^nM_{n+1}\). They are onto, and \(I^nM_n=0\). Recursive lifting gives surjectivity \(M\to M_n\). Set \(N_n=\ker(M\to M_n)\), so \(I^nM\subset N_n\). If \(u\in N_n\), its coordinate in \(M_{n+1}\) lies in \(I^nM_{n+1}\). Write it as a finite sum of products \(a_jv_j\), \(a_j\in I^n\), and lift each \(v_j\) to \(M\). Subtracting the resulting element of \(I^nM\) places \(u\) in \(N_{n+1}\). This proves \(N_n=N_{n+1}+I^nM\) and surjective transitions on \(N_n/I^nM\). The source's expressions with intersections agree because \(I^nM\subset N_n\).

Their limit injects into the kernel of the canonical \(\pi:\widehat M\to M\). By the first lemma \(\eta_M\) is an isomorphism, and the section identity makes \(\pi\) its inverse. Thus that kernel is zero. Prescribed-coordinate lifting in the onto defect system proves each \(N_n/I^nM=0\), and so the actual projection induces \(M/I^nM\cong M_n\). This is an application of the previously proved defect mechanism, not a replacement of the original inverse system.

\subsubsection{Correct stabilization of graded modules}\label{algebra-proof-015-correct-stabilization-of-graded-modules}

Source 23299--23333; OCC-00483. Graded rings have nonnegative degrees, while the original modules may have nonzero components in any integer degree. Keep the original homogeneous ideal \(I\subset A_+\), transition maps, lifted generators \(x_{n,i}\), and their degrees \(d_i\). For \(k\geq n\), composing the transitions gives \[
N_n\cong N_k/I^nN_k .
\] This follows inductively: passing successively through quotient by \(I^{k-1}\) and then by \(I^n\), with \(n\leq k-1\), has kernel exactly \(I^n\). In particular \(I^nN_n=0\) and \(N_n/IN_n=N_1\).

If \(N_1=0\), then \(N_n=IN_n\). Iterating \(n\) times gives \(N_n=I^nN_n=0\); take \(N=0\). Otherwise choose \(r>0\) homogeneous generators of \(N_1\) and compatible homogeneous lifts \(x_{n,i}\). Such lifts exist by the surjective graded transitions. The quotient of \(N_n\) by their span is equal to its own multiple by \(I\), and is killed by \(I^n\); it is zero. Thus these original \(r\) elements generate every \(N_n\), without an additional finiteness argument.

Put \(\delta=\min_i d_i\). A homogeneous product in \(I^mN_k\) has degree at least \(m+\delta\), since every homogeneous element of \(I\) has degree at least one. Therefore \[
(N_{n+1})_d\longrightarrow(N_n)_d
\quad\hbox{is an isomorphism when }n>d-\delta.
\] The strict inequality is essential. Set the positive index \[
b(d)=\max\{1,1+d-\delta\}.
\] Then the stable tail is \[
\cdots=(N_{b(d)+2})_d=(N_{b(d)+1})_d=(N_{b(d)})_d.
\] For \(d<\delta\) all the degree-\(d\) parts are zero. At \(d=\delta\), \(b(d)=1\), so no \(N_0\) is invoked. The source's equality at level \(d-\delta\) is unjustified; its later generation check needs the same corrected index.

Define \(N_d=\lim_n(N_n)_d\) and \(N=\bigoplus_{d\in\mathbb Z}N_d\). Multiplication by \(a\in A_e\) acts coordinatewise and sends \(N_d\) into \(N_{d+e}\); distributivity, unit and associativity follow at every original coordinate. This gives an actual graded \(A\)-module. Each \(x_i=(x_{n,i})_n\) lies in \(N_{d_i}\). For \(x\in N_d\), its image in \((N_{b(d)})_d\) is a homogeneous expression \(\sum_i a_ix_{b(d),i}\), with \(a_i\in A_{d-d_i}\), using zero when \(d-d_i<0\). The same coefficients define \(\sum_i a_ix_i\in N_d\). It has the same stable coordinate as \(x\), so it equals \(x\). Thus the \(x_i\) generate \(N\).

The projection \(N\to N_n\) is onto because it contains the original generating family. Its kernel is exactly \(I^nN\). One inclusion follows from \(I^nN_n=0\). For the reverse, take a homogeneous \(x\in N_d\) mapping to zero and choose \(k\geq n,b(d)\). Its image in \(N_k\) belongs to \((I^nN_k)_d\). As \(N_k\) is generated by the \(x_{k,i}\), it can be written \(\sum_i a_ix_{k,i}\) with \(a_i\in (I^n)_{d-d_i}\): expand the finite products and then retain their actual degree-\(d\) terms. The same coefficients give an element of \((I^nN)_d\), with the same stable coordinate as \(x\). They are equal. For a general kernel element, apply this to its finitely many homogeneous components, since the projection is graded. Hence the canonical maps induce \(N/I^nN\cong N_n\).

The premature equality has a direct source-compatible counterexample. Take \(A=k[t]\) with \(t\) in degree one, \(I=(t)\), and \(N_n=A/(t^n)\), generated in degree zero. In degree one, the transition from \(N_2\) to \(N_1\) maps the nonzero class of \(t\) to zero. Thus equality at level \(d-\delta=1\) is false, while the repaired stable level is \(b(1)=2\). In degree zero, the repaired level is one and never invokes the undefined \(N_0\).

\subsubsection{Graded algebraization without Noetherianity}\label{algebra-proof-015-graded-algebraization-without-noetherianity}

The proof in the preceding section never uses Noetherianity of \(A\) or finite generation of \(I\). It only uses the nonnegative grading, the homogeneous containment \(I\subset A_+\), the specified quotient transitions, and finite generation of \(N_1\). Thus the following strengthening is established as editorial material:

For every nonnegatively graded ring \(A\), every homogeneous ideal \(I\subset A_+\), and every system of graded modules with transition maps \(N_{n+1}\to N_n\) inducing \(N_n\cong N_{n+1}/I^nN_{n+1}\), if \(N_1\) is finitely generated, there is a finitely generated graded \(A\)-module \(N\) with compatible canonical isomorphisms \(N/I^nN\cong N_n\). All the modules \(N_n\) are already finitely generated by the same number of compatible lifted generators. Its construction and all comparison maps are proved above, including \(N_1=0\).

This gives an equivalence with the category of such systems, with morphisms compatible graded maps, from the category of finite graded \(A\)-modules. Here is the remaining proof. For a finite graded \(U\), choose homogeneous generators of least degree \(\delta_U\). Its degree-\(d\) part maps isomorphically to \((U/I^nU)_d\) for \(n>d-\delta_U\), so \[
U_d\cong\lim_n(U/I^nU)_d .
\] For \(U=0\) this holds directly. Given compatible graded maps \(f_n:U/I^nU\to V/I^nV\), apply them to these compatible families in each degree. The resulting maps \(U_d\to V_d\) define \(f:U\to V\) on their direct sums. Additivity and compatibility with each homogeneous element of \(A\) follow coordinatewise; hence \(f\) is a graded \(A\)-module map. Its reduction is \(f_n\) for every \(n\), because it agrees on each original coordinate and on all homogeneous elements. Conversely any graded \(f\) gives these reductions. If all reductions vanish, each homogeneous image lies in every \(I^nV\) and is zero by the degree bound. Thus this correspondence on morphisms is injective as well as surjective, and respects composition and identities. The construction above proves essential surjectivity, completing the equivalence.

The direct sum of degreewise limits used here is not identified with the unrestricted module limit. For \(A=k[t]\), \(I=(t)\), and \(N_n=k[t]/(t^n)\), the reconstructed graded module is \(k[t]\), while the unrestricted limit is \(k[[t]]\). The compatible series with every coefficient one has infinitely many nonzero degrees and therefore is not in that direct sum. This exact example retains the distinction necessary for the source's next lemma. No finite-presentation claim is made after removing Noetherianity.

\subsubsection{The shifted degree bound and separate injectivity and surjectivity}\label{algebra-proof-015-the-shifted-degree-bound-and-separate-injectivity-and-surjectivity}

Source 23335--23401; OCC-00484. Retain \(A'=\lim A/I^n\), the graded modules \(G_n\) with \(I^nG_n=0\), the original \(M\), and compatible graded maps \(\varphi_n\). Coordinate actions give the canonical map \[
\varphi:M\otimes_A A'\longrightarrow\lim_nG_n.
\]

First verify the canonical unit \(M\to M\otimes_AA'\) is injective for every graded \(M\), with no finite-generation assumption. Suppose a nonzero \(x\in M\) maps to zero. Express \(M\) as the filtered union of its finitely generated homogeneous submodules. Tensor products commute with this union's colimit, so one such submodule \(M'\) contains \(x\) and witnesses its zero tensor image; a finite homogeneous generating family exists by taking components of the finitely many generators and enlarging to another stage if necessary. Let \(\delta\) be the least generator degree and \(D\) the greatest degree in the finite homogeneous support of \(x\). Then \(D\geq\delta\), and choose \(n=D-\delta+1>0\). Every component of \(I^nM'\) has degree at least \(\delta+n=D+1\). The quotient \(M'\to M'/I^nM'\) is therefore injective on the direct sum of degrees at most \(D\). But this quotient factors as \[
M'\longrightarrow M'\otimes_AA'
 \longrightarrow M'\otimes_A(A/I^n)
 \cong M'/I^nM',
\] so the assumed zero tensor image forces \(x=0\), a contradiction.

Applied to a homogeneous element of degree \(d\), this proof uses \(n=d-\delta+1\) and the module-degree bound \(e<\delta+n\). The source's coefficient bound \(e\leq n-1\) does not account for generator shifts. Replacing its module bound by \(e<\min(d_i)+n\), explicitly including degree \(d\), is the minimal report correction. If \(x=0\) the claim is immediate, so the nonempty generating family and positive \(n\) can be chosen as just described.

It follows that injectivity of \(\varphi\) alone implies injectivity of \(M_d\to\lim_nG_{n,d}\): a homogeneous element with zero image has zero tensor image by that injectivity and hence is zero by the unit argument. The source's full isomorphism hypothesis is stronger than needed for this direction.

For the surjectivity direction, the homogeneous coefficient construction is also exact without Noetherianity. Since \(I\subset A_+\), the degree-\(j\) part of \(A/I^n\) agrees with \(A_j\) for \(n>j\). An element \(f'\in A'\) therefore has a well-defined coefficient in every \(A_j\). These coefficients identify \(A'\) as a submodule of \(\prod_{j\geq0}A_j\): if all coefficients vanish, its image in each graded quotient \(A/I^n\) has every homogeneous component zero and is zero. Conversely no arbitrary infinite product is assumed to define an element of \(A'\). Every finite sum of actual homogeneous coefficients is in \(A\subset A'\); the latter inclusion is injective because a nonzero finite homogeneous sum cannot lie in all \(I^n\).

Suppose only that \(\varphi\) is surjective and take \(y\in\lim_nG_{n,d}\). Lift it to a finite tensor sum \(\sum_i x_i\otimes f'_i\), splitting each \(x_i\) into homogeneous parts, with degrees \(d_i\). Put \(c_i=d-d_i\). If \(c_i\geq0\), decompose using the original coefficients \[
f'_i=\sum_{0\leq j<c_i}f_{i,j}+f_i+g'_i,
\quad f_{i,j}\in A_j,\quad f_i\in A_{c_i},
\] where every coefficient of \(g'_i\) in degrees at most \(c_i\) is zero. If \(c_i<0\), take the finite lower sum and \(f_i\) to be zero and put \(g'_i=f'_i\). This handles the negative-degree boundary of the printed expression without inventing negative coefficients in \(A\).

At every quotient level, a lower term \(f_{i,j}\varphi_n(x_i)\) has degree \(j+d_i<d\). The image of \(g'_i\) in \(A/I^n\) has only finitely many homogeneous components, all of degree greater than \(c_i\), or of nonnegative degree if \(c_i<0\). Multiplying by \(\varphi_n(x_i)\) therefore gives only degrees greater than \(d\). Consequently the degree-\(d\) component of the lifted tensor image is exactly \(\varphi_n(\sum_i f_ix_i)\) at every \(n\). Since \(y\) has degree \(d\) at every coordinate, \(\sum_i f_ix_i\in M_d\) maps to \(y\). Surjectivity of \(\varphi\) alone thus gives surjectivity on every degree. Together the two directions prove the source lemma and the separate stronger implications.

The unit result concerns graded modules and is not purity on every ungraded module. For \(A=k[t]\), \(I=(t)\), the ungraded \(U=A/(t-1)\) is nonzero, but \(U\otimes_A k[[t]]=k[[t]]/(t-1)=0\), since \(t-1\) has inverse \(-\sum_{n\geq0}t^n\). This example prevents an unsupported propagation to the pure-descent results elsewhere. Likewise the finite tensor lift in the surjectivity proof is essential; no identification of an arbitrary tensor product with an unrestricted product of homogeneous pieces is used.

The strengthened algebraization and degreewise implications are linked to their exact source passages in the claim ledger. They preserve all original generator degrees and completion maps. Later source review must check whether these consequences already appear elsewhere before any broader classification or novelty assessment. The combined-results synthesis remains after the core correction work.

\hypertarget{algebra_flatness_criteria_notes_20260928}{}
\subsection{Local flatness criteria and their exact Tor maps}\label{algebra-proof-016-local-flatness-criteria-and-their-exact-tor-maps}

Separate editorial evidence for the Stacks Project authors' algebra.tex, authority commit a04446e57ec1fbc252a871afcec7752fb2807b14, lines 23415--24176. Authority SHA-256: FA8BB92E58A4F78A2BD01B3B6A4A87DE0A0D279F5DD90641B574DD5FBFFFA4F3. All complete source units and all ten received reports in this interval were read. The next section heading at 24177--24180 is lookahead only.

The two retained canonical units MC-STK-ERR-0485 and MC-STK-ERR-0486 resolve four reports. Seventeen minimal operations are proposed for the other six reports and the two independently found much-greater-than symbols. Full proofs and three further consequences are separate editorial material. The source statements retain their original hypotheses and remain identifiable; no chapter or translation mutation, formal certification or novelty claim is asserted. The bounded library query records all four hits as unread, including the local-work hits.

\subsubsection{Injection modulo the maximal ideal and its original quotients}\label{algebra-proof-016-injection-modulo-the-maximal-ideal-and-its-original-quotients}

Source 23423--23477. The naming correction at 23427 is part of OCC-00485. Keep the local map \(R\to S\), maximal ideal \(\mathfrak m\subset R\), finite \(S\)-module \(N\), flat \(R\)-module \(M\), and original \(R\)-linear map \(u:N\to M\). No \(S\)-module structure on \(M\), or \(S\)-linearity of \(u\), is assumed.

Assume \(u_1:N/\mathfrak mN\to M/\mathfrak mM\) is injective. For each \(n\), the map on the left tensor terms in the source diagram is \(u_1\otimes_{R/\mathfrak m}\mathrm{id}_{\mathfrak m^n/\mathfrak m^{n+1}}\). It is injective because tensoring over a field is exact. The lower left term injects into \(M/\mathfrak m^{n+1}M\) by flatness of \(M\). To prove the inductive step, take an element in the kernel of \(u_{n+1}\). Its image modulo \(\mathfrak m^n\) is zero by injectivity of \(u_n\), so it lifts from the upper left tensor term. The image of that lift in the lower middle term is zero, and lower left injectivity then kills its image in the lower left term. Injectivity of the left vertical map kills the lift itself. Thus \(u_{n+1}\) is injective.

For \(x\in\ker u\), all these injections give \(x\in\bigcap_n\mathfrak m^nN\). The ideal \(\mathfrak mS\) lies in the maximal ideal of the Noetherian local ring \(S\); \(N\) is finite over \(S\). The source's Krull-intersection result gives this intersection zero, hence \(u\) is injective.

For a proper ideal \(J\subset R\), \(J\subset\mathfrak m\), and \(JS\) lies in the maximal ideal of \(S\). Consequently \(R/J\to S/JS\) is a local map of nonzero local rings, \(S/JS\) is Noetherian, \(N/JN\) is finite over it, and \(M/JM\) is flat over \(R/J\). Reduction of the quotient map modulo \(\mathfrak m/J\) is precisely \(u_1\). Applying the just-proved injection assertion shows \(N/JN\to M/JM\) injective. For \(J=R\) both quotient modules are zero and the injection is immediate; the zero ring is not incorrectly called a local ring. This supplies the boundary implicit in the source's ideal argument without changing its theorem.

The long exact sequence from \(0\to N\to M\to Q\to0\), with \(Q=M/u(N)\), now identifies \[
\operatorname{Tor}_1^R(Q,R/J)
 =\ker(N/JN\to M/JM)=0
\] for every ideal \(J\). Thus \(Q\) is flat. More generally tensoring that exact sequence with any module \(X\) remains exact on the left because \(\operatorname{Tor}_1^R(Q,X)=0\). The injection \(u\) is therefore universally injective as a map of \(R\)-modules. Also \(N\) is flat: for an injection \(X\to Y\), its tensor with \(N\) sits as a map of submodules of \(X\otimes_RM\to Y\otimes_RM\), with both horizontal inclusions just proved; flatness of \(M\) makes the latter injective. These consequences use no Noetherianity of \(R\).

\subsubsection{The early corollaries and the exact-complex induction}\label{algebra-proof-016-the-early-corollaries-and-the-exact-complex-induction}

Source 23479--23552. The remaining two naming corrections of OCC-00485 occur at 23482 and 23494. Multiplication by the original \(f\) on \(S\) is an \(R\)-linear map from a finite \(S\)-module to an \(R\)-flat module. Its reduction is injective by the given fibre nonzerodivisor hypothesis. The preceding lemma gives both injectivity of multiplication by \(f\) and flatness of its exact cokernel \(S/fS\).

For a regular sequence, put \(S_i=S/(f_1,\ldots,f_i)\). At stage \(i\), multiplication by \(f_i\) on \(S_{i-1}\) has exactly the specified injective fibre map. Apply the same lemma with the original Noetherian \(S\) as the auxiliary ring and the finite \(S\)-module \(S_{i-1}\) as both domain and target. It gives the next injection and \(R\)-flatness of \(S_i\). The source definition of regularity also requires \(S_c\ne0\); this follows because its reduction modulo \(\mathfrak m\) is the given nonzero final fibre quotient. Thus no unit or zero-quotient case has been lost.

For the free-fibre lemma, \(M\ne0\) finite over local \(S\) implies \(M/\mathfrak mM\ne0\) by Nakayama, since \(\mathfrak mS\) is inside the maximal ideal. The fibre basis has finite positive size \(n\). The source map \(S^{\oplus n}\to M\) is injective by the injection lemma, and its finite cokernel vanishes modulo \(\mathfrak m\), hence vanishes by Nakayama. This gives the actual chosen-basis isomorphism. Since \(n>0\), \(S\) is a direct summand of \(M\) as an \(R\)-module, so it is \(R\)-flat.

For the finite complex with \(e\geq1\), apply the injection lemma to \(F_e\to F_{e-1}\). Its cokernel \(C\) is finite over \(S\) and \(R\)-flat. Right exactness identifies \(C/\mathfrak mC\) with the fibre cokernel. The assumed fibre exactness therefore makes \(0\to C/\mathfrak mC\to F_{e-2}/\mathfrak mF_{e-2}\to\cdots\) exact. This is the required induction hypothesis on the shorter original quotient complex. It proves exactness at every term at which the displayed complex asserts exactness and flatness of \(\operatorname{coker}(F_1\to F_0)\). If a complex of length zero is included by the convention \(F_1=0\), its final flatness assertion is just the given flatness of \(F_0\). No missing end map is treated as a surjectivity assertion.

\subsubsection{Finite length, Artin--Rees and the local criterion}\label{algebra-proof-016-finite-length-artinrees-and-the-local-criterion}

Source 23559--23641; OCC-00486 changes the incorrectly described descending induction to ordinary induction, and OCC-00487 introduces the quantified ideal \(J\). Both printed occurrences of \(n>>0\) receive the intended TeX relation \(\gg\).

For a finite-length \(N\), length zero means \(N=0\), whose Tor groups vanish. Length one identifies \(N\) with the residue field \(\kappa\). For length greater than one, use the source exact sequence with both nonzero end modules of smaller length. The exact segment \[
\operatorname{Tor}_1^R(N',M)\longrightarrow
\operatorname{Tor}_1^R(N,M)\longrightarrow
\operatorname{Tor}_1^R(N'',M)
\] and the two induction hypotheses force the middle term to zero. This proves the preparation lemma over every local \(R\).

In the local flatness criterion retain both Noetherian hypotheses. For every ideal \(J\) of finite colength, the preparation lemma and the exact sequence \(0\to J\to R\to R/J\to0\) show \(J\otimes_RM\to M\) injective. Fix the source ideal \(I\subset R\) and \(n\geq1\). Its displayed short exact sequence can be equipped consistently with the actual maps \[
x\longmapsto(x,-x),\qquad (a,b)\longmapsto a+b.
\] Use the same maps in the ambient split sequence. Both \(I+\mathfrak m^n\) and \(\mathfrak m^n\) have finite colength, since \(R\) is Noetherian local. If \(z\in\ker(I\otimes_RM\to M)\), then \((z,0)\) maps to zero in \((I+\mathfrak m^n)\otimes_RM\), because the latter maps injectively into \(M\). Tensor right exactness supplies a lift \(w\in(I\cap\mathfrak m^n)\otimes_RM\). Its component in \(\mathfrak m^n\otimes_RM\) is zero, so its multiplication image in \(M\) is zero. Thus this lift belongs to the precise kernel asserted in the source and maps to \(z\).

Artin--Rees for \(I\subset R\) gives \(c\) with \(I\cap\mathfrak m^n\subset\mathfrak m^{n-c}I\) for \(n\geq c\). The image of its tensor with \(M\) is contained in \(\mathfrak m^{n-c}(I\otimes_RM)\), without assuming that tensoring the inclusion is injective. Hence the original kernel lies in every sufficiently high power of \(\mathfrak m\) on \(I\otimes_RM\). This module is finite over \(S\): choose the finite \(R\)-generators of \(I\) and the finite \(S\)-generators of \(M\), and use their simple tensors. Krull intersection over \(S\) therefore makes the kernel zero. Equivalently, as the source argues, the finite-module comparison identifies its completion with tensoring by the faithfully flat \(\mathfrak mS\)-adic completion of \(S\), whose unit is injective. Both routes retain the original ideal and uniform Artin--Rees shift. Flatness follows from the ideal criterion.

\subsubsection{The correct free index and the ideal-adic Tor criterion}\label{algebra-proof-016-the-correct-free-index-and-the-ideal-adic-tor-criterion}

Source 23650--23748. OCC-00488 concerns a substantive rank problem, not merely a glyph collision. The original ideal \(I\) cannot be the prescribed index set for a free presentation of an arbitrary \(R/I\)-module. For instance, for a nonzero field \(R\), \(I=0\) has one underlying element and its indexed free module cannot surject onto \(R^2\). Choose an actual sufficiently large set \(\Lambda\), for example the underlying set of \(N\), and the evaluation surjection \[
(R/I)^{(\Lambda)}\longrightarrow N,\qquad e_\lambda\longmapsto \lambda .
\] With its actual kernel \(K\), this supplies the source exact sequence. All six free-sum occurrences must use the same new index, not just the two mentioned by the report. The separate malformed base \(R/\) is already corrected to \(R/I\) by MC-STK-ERR-0485; both OCC-12142 and OCC-00489 retain it.

Direct sums commute with the Tor construction, by a free resolution and componentwise homology. Thus \(\operatorname{Tor}_1^R((R/I)^{(\Lambda)},M)=0\). For every module \(X\) annihilated by \(I\), the canonical map \[
X\otimes_R M\longrightarrow X\otimes_{R/I}(M/IM)
\] sends \(x\otimes m\) to \(x\otimes\overline m\). Its inverse sends \(x\otimes\overline m\) to \(x\otimes m\); changing \(m\) by a finite sum of terms \(am'\), \(a\in I\), changes this by terms \(ax\otimes m'=0\). These inverse balanced maps preserve the exact presentation and its kernel tensor. Flatness of \(M/IM\) therefore shows \(K\otimes_RM\to(R/I)^{(\Lambda)}\otimes_RM\) injective, and the Tor exact sequence gives \(\operatorname{Tor}_1^R(N,M)=0\).

If \(I^mN=0\), the case \(m=0\) gives \(N=0\). The case \(m=1\) is proved. For \(m>1\), retain \(N'=IN\) and \(N''=N/IN\); both are annihilated by \(I^{m-1}\). The same Tor exact segment proves the assertion inductively. Applied to every exact sequence of \(R/I^n\)-modules, this vanishing shows exactness of tensoring with \(M/I^nM\), through the canonical comparison just given with \(I^n\) in place of \(I\). Thus \(M/I^nM\) is \(R/I^n\)-flat.

For the following graded-piece criterion, work at stage \(R/I^{n+1}\) with its ideal \(I^n/I^{n+1}\). Its square is zero for \(n\geq1\), and the ideal tensor exact sequence identifies \[
\operatorname{Tor}_1^{R/I^{n+1}}(M/I^{n+1}M,R/I^n)
 =\ker\bigl(M\otimes_R I^n/I^{n+1}\to I^nM/I^{n+1}M\bigr).
\] The target on the right is the actual multiplication image; replacing it by its containing module does not change the kernel. Assuming the previous stage flat, the nilpotent case of the just-proved criterion makes the next stage flat. This proves both statements with the original stage indices.

\subsubsection{The two local variants and the support boundary}\label{algebra-proof-016-the-two-local-variants-and-the-support-boundary}

Source 23750--23843. In the first proof of the ideal variant, the proper ideal \(I\) is inside \(\mathfrak m\), so \(\kappa\) is annihilated by \(I\). The preceding criterion gives its Tor vanishing, and the proved local criterion applies.

For the second proof retain every original \(f_i,x_i,a_{ij},y_j\). Put \(v_i=x_i-\sum_j a_{ij}y_j\in IM\) and \(b_j=\sum_i f_i a_{ij}\in I\). For each \(i\) choose a finite expression \(v_i=\sum_\ell c_{i\ell}z_{i\ell}\), with \(c_{i\ell}\in I\). Then \[
w=\sum_{i,\ell}f_ic_{i\ell}\otimes z_{i\ell}
       +\sum_j b_j\otimes y_j\in I\otimes_RM
\] maps to the source's tensor \(\sum_i f_i\otimes x_i\) in \(\mathfrak m\otimes_RM\). The equality uses the two balancing identities, with the original plus signs and the subtracted \(a_{ij}y_j\) restored in \(v_i\). Its product in \(M\) is zero. The assumed Tor vanishing makes \(I\otimes_RM\to M\) injective, so \(w=0\) and the original tensor is zero. For \(I=(x)\) with \(x\) a nonzerodivisor of \(R\), the original length-one free resolution with differential multiplication by \(x\) identifies the stated Tor group with the actual annihilator \(M[x]\).

In the powers lemma, retain \(K=\ker(I\otimes_RM\to M)\). The level-\(n\) flatness assumption kills its image in \[
(I/I^n)\otimes_{R/I^n}(M/I^nM)
 \cong (I\otimes_RM)/I^{n-1}(I\otimes_RM).
\] The final quotient follows from tensor right exactness and \(I^n=I^{n-1}I\); the image of \(I^n\otimes_RM\) is precisely the displayed power submodule. At \(n=1\) both descriptions are zero, as required. Hence \(K\) lies in their intersection for all \(n\). The module \(I\otimes_RM\) is finite over Noetherian \(S\), and the earlier uniform Krull-intersection annihilator from group 344 gives one \(s\in1+IS\) annihilating that intersection. For every \(\mathfrak q\supset IS\), \(s\notin\mathfrak q\), so \(K_{\mathfrak q}=0\). This proves the localization assertion without an unsupported interchange of intersection and localization.

For completeness, put \(\mathfrak p=\mathfrak q\cap R\). The localized hypotheses give the original local map \(R_{\mathfrak p}\to S_{\mathfrak q}\), proper ideal \(IR_{\mathfrak p}\), and finite module \(M_{\mathfrak q}\). Localization and the tensor comparisons identify its Tor kernel with \(K_{\mathfrak q}\), and its quotient is flat over \(R_{\mathfrak p}/IR_{\mathfrak p}\). The local variant proves flatness over \(R_{\mathfrak p}\), hence over \(R\).

The locus restriction is necessary. For \(R=S=k[t]\), \(I=(t)\) and \(M=R/(t-1)\), every \(M/I^nM\) is zero and therefore flat, while \(M\) is not \(R\)-flat: multiplication by the nonzero \(t-1\) on \(R\) becomes the zero map on the nonzero \(M\). At the only prime containing \(t\), the localization of \(M\) is zero, in agreement with the proved result. No global flatness is inferred solely from these formal quotients.

\subsubsection{Tor comparisons and the exact change-of-rings sequence}\label{algebra-proof-016-tor-comparisons-and-the-exact-change-of-rings-sequence}

Source 23845--23918. MC-STK-ERR-0486 already changes ``cohomology'' to ``homology''; both OCC-00490 and OCC-12143 retain it. The original complexes compute first homology.

The two lemmas have the following common extension for every ring map \(R\to R'\), \(R\)-module \(M\), and \(R'\)-module \(X\). Put \(M'=M\otimes_RR'\). There is a natural right-exact sequence \[
\operatorname{Tor}_1^R(M,R')\otimes_{R'}X
 \xrightarrow{a}\operatorname{Tor}_1^R(M,X)
 \xrightarrow{b}\operatorname{Tor}_1^{R'}(M',X)
 \longrightarrow0.
\] No injectivity of \(a\) is asserted.

Here is a construction and proof with all original maps retained. Take a free resolution \(F_\bullet\to M\) over \(R\), and put \(C_\bullet=F_\bullet\otimes_RR'\). Let \(Z=\ker(C_1\to C_0)\) and \(B=\operatorname{im}(C_2\to C_1)\). Then \(Z/B=\operatorname{Tor}_1^R(M,R')\), and \(C_1\to C_0\to M'\to0\) is exact. Choose a free \(R'\)-module \(G\) surjecting onto \(Z\). Replace \(C_2\) by \(C_2\oplus G\), sending \((c,g)\) to \(d_2(c)+g\) through this chosen surjection. Keep degrees one and zero exactly \(C_1,C_0\). This is exact in degrees one and zero; choose further free modules onto its successive kernels to obtain a resolution \(P_\bullet\to M'\). The identity in degrees zero and one and the inclusion \(C_2\to C_2\oplus G\) extend to a chain map \(C_\bullet\to P_\bullet\): at every higher step, the previous compatibility puts the required image in the next kernel, and freeness lifts it along the chosen surjection.

After tensoring with \(X\), both complexes have the same degree-one cycles \[
T=\ker(C_1\otimes_{R'}X\to C_0\otimes_{R'}X).
\] The first homology of \(C_\bullet\otimes X=F_\bullet\otimes_RX\) is \(T/\operatorname{im}(C_2\otimes X)\). That of \(P_\bullet\otimes X\) is \(T/\operatorname{im}(Z\otimes X)\), because \(G\to Z\) is onto and tensoring is right exact. The quotient map gives the surjective \(b\), and its kernel is the image of \(Z\otimes X\) modulo the old boundaries. The map from \(Z\otimes X\) to the first homology kills \(\operatorname{im}(B\otimes X)\); tensor right exactness therefore factors it through \((Z/B)\otimes X\), giving \(a\) with exactly that image. Explicitly \(a([z]\otimes x)\) is the class of the cycle \(z\otimes x\). This proves exactness at every asserted term.

Changing the free resolutions uses their comparison maps lifting the identity. Such maps exist by degreewise free lifting, and two choices are chain homotopic by the same kernel-lifting induction. Tensoring preserves the equation \(f-g=dh+hd\), so the induced homology maps agree. The same construction for module maps commutes on homology. Thus this is the natural sequence, not a choice-dependent replacement. The earlier complete comparison-homotopy argument in group 460 applies to these exact maps.

If \(M'\) is flat, the last Tor group is zero and \(a\) is onto for every \(R'\)-module \(X\). Taking \(X=R''\) recovers the source's first lemma; it actually proves the stronger arbitrary-module version. The source proof's tensor exactness is also justified directly: the image of \(C_1\to C_0\) is flat, as the kernel of a surjection from a free module onto flat \(M'\), and tensoring the two short exact sequences preserves the stated exactness. Taking arbitrary \(M'\) and \(X=R'/IR'\) recovers the second lemma; its surjection \(b\) in fact holds for every \(X\).

\subsubsection{The localized comparison criterion and the fibre criterion}\label{algebra-proof-016-the-localized-comparison-criterion-and-the-fibre-criterion}

Source 23920--24023. OCC-00491 only inserts ``by'' in the notation declaration. Let \(T=S\otimes_RR'\) and write \(S'=W^{-1}T\), retaining the source localization and all original maps. Then \[
M'=(M\otimes_RR')\otimes_T S',
\] by the map on simple tensors \(m\otimes(s\otimes r')\mapsto sm\otimes r'\) and its balanced inverse. Its reduction modulo \(I'=IR'\) is the corresponding localization of \((M/IM)\otimes_{R/I}(R'/I')\), which is flat over \(R'/I'\). Localization of a flat module with an auxiliary algebra action remains flat over the base: each localization is a filtered colimit of copies of the base-flat module with its multiplication transitions, and filtered colimits preserve tensor exactness.

Apply \(a\) in the preceding sequence first to \(R\to R/I\) and \(X=R'/I'\). Since \(M/IM\) is flat, it is surjective. Then apply \(b\) to \(R\to R'\) and this same \(X\). Their composite is exactly \[
\operatorname{Tor}_1^R(M,R/I)\otimes_R R'
 \longrightarrow \operatorname{Tor}_1^{R'}(M\otimes_RR',R'/I'),
\] using the canonical equality of tensoring an \(I\)-annihilated module by \(R'\) and by \(R'/I'\). It is onto. Localizing the target in \(T\) identifies it with the Tor group for \(M'\); this follows either from a free resolution of \(R'/I'\) tensored with \(M\otimes_RR'\), or from exactness of localization on its resulting homology. The source assumption that the original Tor map is zero makes every image generator, and hence its localization, zero. The local criterion over the original Noetherian \(R',S'\) then gives the asserted flatness.

For the fibre criterion retain \(I=\mathfrak mS\). Multiplication gives a surjection \(\mathfrak m\otimes_RS\to I\), and tensoring with \(M\) gives \[
\mathfrak m\otimes_RM\longrightarrow I\otimes_SM\longrightarrow M.
\] The first map is onto and the composite is injective by \(R\)-flatness of \(M\). Thus the first map is also injective, and the second is injective: any element in its kernel lifts along the first map and the composite kills that lift. It follows that \(\operatorname{Tor}_1^S(S/I,M)=0\). The ideal variant for the original Noetherian local map \(S\to S'\) gives \(S\)-flatness of \(M\). Since \(M\ne0\) is finite over local \(S'\), its residue quotient at the maximal ideal of \(S'\) is nonzero by Nakayama. Its quotient at the smaller maximal ideal from \(S\) maps onto that quotient and is nonzero. Flatness over the local ring \(S\) therefore makes \(M\) faithfully flat.

The original source then tensors the exact Tor-kernel sequence with this flat \(M\) and uses faithfulness to kill \(\operatorname{Tor}_1^R(\kappa,S)\). This is valid under its Noetherian \(R\) hypothesis. The next section gives the stronger conclusion by testing arbitrary ideals, thereby avoiding that particular extra hypothesis.

\subsubsection{Removing the unused base Noetherian hypothesis}\label{algebra-proof-016-removing-the-unused-base-noetherian-hypothesis}

This is editorial propagation, not a change of the source theorem statements. The injection lemma assumes only \(S\) Noetherian, and its proof above uses no Noetherianity of \(R\). It follows that the three early corollaries and the finite-complex assertion remain valid for arbitrary local \(R\), retaining local Noetherian \(S\). The individual arguments in the second section explicitly check every application and cokernel. This does not remove Noetherianity from the separate local Tor criterion, whose Artin--Rees argument uses it.

There is a module version of the regular-sequence consequence. Let \(R\to S\) be local with \(S\) Noetherian, let \(Q\) be finite over \(S\) and flat over \(R\), and suppose the original elements \(f_1,\ldots,f_c\in S\) are regular on \(Q/\mathfrak mQ\), including the nonzero final quotient required by the source definition. Set \(Q_i=Q/(f_1,\ldots,f_i)Q\). Multiplication by \(f_i\) on \(Q_{i-1}\), with the unchanged auxiliary ring \(S\), satisfies the injection lemma whenever the preceding quotient is \(R\)-flat. Induction gives each injection and flatness of every \(Q_i\). The final fibre quotient is nonzero; each preceding quotient maps onto it, so every \(Q_i\) is nonzero. Each is finite over local \(S\). Nakayama makes \(Q_i/\mathfrak mQ_i\ne0\); the flat local criterion for faithful modules then makes every \(Q_i\) faithfully flat over \(R\). Thus this strengthens the ring case while retaining the original elements, quotients and nonzero condition.

The fibre criterion at 23977 also works for arbitrary local \(R\), with \(S,S'\) still Noetherian local. The preceding section already proves \(M\) faithfully flat over \(S\) without Noetherianity of \(R\): its local criterion is over \(S\). For any ideal \(J\subset R\), put \[
L_J=\ker(J\otimes_RS\to S).
\] Flatness of \(M\) over \(S\) identifies \(L_J\otimes_SM\) with the kernel of \(J\otimes_RM\to M\), which is zero by \(R\)-flatness of \(M\). Faithfulness over \(S\) gives \(L_J=0\). The ideal criterion therefore makes \(S\) flat over arbitrary \(R\). This is the exact map replacing the Noetherian-only last step, and proves the extended fibre criterion. Moreover the local map \(R\to S\) is faithfully flat because its residue fibre is nonzero. No finiteness of the ideal \(J\) was assumed in this argument.

\subsubsection{Elementwise ideal-power torsion and the associated graded map}\label{algebra-proof-016-elementwise-ideal-power-torsion-and-the-associated-graded-map}

Under exactly the source hypotheses \(M/IM\) flat over \(R/I\) and \(\operatorname{Tor}_1^R(R/I,M)=0\), the preceding bounded-power vanishing extends to every module \(N\) whose every element is annihilated by some power of \(I\), with no single power required for all elements. Define \[
N[n]=\{x\in N:I^nx=0\}.
\] These are increasing submodules and their union is exactly \(N\). Each has the proved Tor vanishing. A fixed free resolution of \(M\), tensored with this filtered system, computes the Tor groups. Tensor products commute with filtered colimits, and filtered colimits of modules are exact, so kernels modulo images, and hence homology, commute with this system. It follows that \[
\operatorname{Tor}_1^R(N,M)=
 \varinjlim_n\operatorname{Tor}_1^R(N[n],M)=0.
\] Thus tensoring with \(M\) preserves any short exact sequence of such modules. No finite generation of \(I\) is assumed, and no closure of this class under arbitrary extensions is asserted when \(I\) is infinitely generated.

The same hypotheses give canonical isomorphisms in every associated-graded degree: \[
(I^n/I^{n+1})\otimes_{R/I}(M/IM)
 \xrightarrow{\;\mu_n\;} I^nM/I^{n+1}M,\qquad
(a\bmod I^{n+1})\otimes(m\bmod IM)\longmapsto am\bmod I^{n+1}M.
\] For \(n=0\) this is the usual identity comparison. For \(n\geq1\), the bounded-power vanishing applied to \(R/I^n\) identifies \(I^n\otimes_RM\) injectively with its actual multiplication image \(I^nM\). The same holds for \(n+1\). Tensor right exactness on \(I^{n+1}\to I^n\to I^n/I^{n+1}\to0\) then identifies the quotient with \(I^nM/I^{n+1}M\). The source-balanced comparison for the \(I\)-annihilated module \(I^n/I^{n+1}\) gives the displayed map. Its inverse is induced by these quotient isomorphisms, so it is independent of every representative. Multiplication by an original class \(b\in I^a/I^{a+1}\) commutes with the maps, since both composites send \(b\otimes(a\otimes m)\) to \(bam\). Hence they assemble into an isomorphism of graded modules over the original associated graded ring.

Conversely, flatness of \(M/IM\) and injectivity of all these original graded-piece multiplication maps imply flatness of every \(M/I^nM\), by the source's proved square-zero induction. This converse is only a statement about all finite quotients. It does not by itself prove \(\operatorname{Tor}_1^R(R/I,M)=0\) or flatness of \(M\); the preceding support example shows why global flatness cannot be inferred from the formal quotients alone.

\subsubsection{The localization gluing lemmas and faithful flatness}\label{algebra-proof-016-the-localization-gluing-lemmas-and-faithful-flatness}

Source 24025--24168; OCC-00492 repairs only the two sentence fragments explaining Tor vanishing. For the single element \(f\), the original exact free resolution \(0\to A\xrightarrow{f}A\to A/fA\to0\) and injectivity of \(f\) on \(M\) give \(\operatorname{Tor}_1^A(M,A/fA)=0\), while length one gives the vanishing in degree two. The flat quotient \(M/fM\) and the ideal-adic criterion give Tor degree one zero for every \(f\)-annihilated module. Presenting such a module by a sum of copies of \(A/fA\) and using the exact Tor segment in the source gives degree two zero as well.

For any original \(K\), use the actual maps \(K\to K'=K/K[f]\) and \(K'\to K\), the latter sending the class of \(k\) to \(fk\). This is well defined, injective, and has image \(fK\); their composite is multiplication by \(f\). The source's two exact Tor segments show that both induced maps on degree one are injective, because their preceding terms vanish. Thus multiplication by \(f\) on \(\operatorname{Tor}_1^A(M,K)\) is injective. The localization map from this Tor module is injective: any element in its kernel is killed by a finite power of \(f\), and repeated injectivity forces it to zero. Exact localization of the free resolution identifies that localized Tor module with \(\operatorname{Tor}_1^{A_f}(M_f,K_f)=0\). Therefore every original Tor group vanishes and \(M\) is flat.

For the multi-generator lemma retain \(I=(f_1,\ldots,f_r)\), \(r\geq1\), and the full assumed range \(1,\ldots,r+1\). First, presenting any \(I\)-annihilated \(K\) by a free \(A/I\)-module, whose kernel is again \(I\)-annihilated, gives the stated vanishing successively in each degree through \(r+1\). Define the induction assertion for \(0\leq j\leq r\) to be vanishing in degrees \(1,\ldots,j+1\) on modules annihilated by \(f_1,\ldots,f_j\). The assertion at \(r\) is just proved. For \(j\geq1\), take \(K\) annihilated by the first \(j-1\) elements. Its submodule \(K[f_j]\) and quotient \(K/f_jK\) are both annihilated by all first \(j\) elements. For each \(1\leq i\leq j\), the two source sequences and the inductive vanishings in degrees \(i\) and \(i+1\) make multiplication by \(f_j\) injective on \(\operatorname{Tor}_i^A(M,K)\). Its localization is zero by the given \(A_{f_j}\)-flatness of \(M_{f_j}\). Thus it is zero. This proves the assertion at \(j-1\), retaining both endpoints of the degree range. At \(j=0\) it is precisely Tor degree one vanishing for every \(K\), which proves flatness.

Here is the faithful-flatness argument omitted by both source proofs. Once flatness is established, take any maximal ideal \(\mathfrak q\) of \(A\). If some \(f_i\notin\mathfrak q\), the assumed faithful flatness of \(M_{f_i}\) gives \(M\otimes_A\kappa(\mathfrak q)
 \cong M_{f_i}\otimes_{A_{f_i}}\kappa(\mathfrak q)\ne0\). If all \(f_i\in\mathfrak q\), use instead the faithfully flat \(M/IM\) over \(A/I\) to obtain the same nonzero fibre. The single-element case is this argument with \(I=(f)\). The nonzero-fibre criterion for a flat module now gives faithful flatness over \(A\). Conversely faithful flatness over \(A\) passes to each displayed localization and quotient by its base-change definition. Unit ideals, empty spectra and zero rings are covered by these same base-change maps; the source explicitly keeps \(r\geq1\). This completed argument stays in editorial evidence and does not replace the source's omitted-proof sentence.

All receiving consequences above are linked in the claim ledger. The current scope is source-ordered correction accounting and the exact implications it reveals. Broader synthesis, claims of novelty and a combined-results paper remain after the core work.

\hypertarget{algebra_nilpotent_flatness_notes_20260928}{}
\subsection{Base change and nilpotent flatness descent}\label{algebra-proof-017-base-change-and-nilpotent-flatness-descent}

Separate editorial evidence for the Stacks Project authors' algebra.tex, authority commit a04446e57ec1fbc252a871afcec7752fb2807b14, lines 24177--24575. Authority SHA-256: FA8BB92E58A4F78A2BD01B3B6A4A87DE0A0D279F5DD90641B574DD5FBFFFA4F3. Every complete source unit and all eight received reports in this interval were read. The following complex section was read through line 24680; its unit-entry splitting lemma is complete in that reading, but no report or result beyond 24575 is adjudicated in this batch.

All eight reports retain the four existing canonical units MC-STK-ERR-0487 through MC-STK-ERR-0490. One separate grammatical operation is proposed. The complete original arguments and two further consequences below are editorial material, with no chapter or translation mutation and no novelty claim. The new corpus query records four unread routing hits; its historical Bhatt--Scholze source coverage is not enlarged by rediscovering its index entry.

\subsubsection{The original base-change and local-ring diagrams}\label{algebra-proof-017-the-original-base-change-and-local-ring-diagrams}

Source 24184--24278. OCC-12144 and OCC-00493 retain MC-STK-ERR-0487: a \emph{nonzero} kernel stays nonzero under faithful flatness, which is the assertion required by the argument.

Keep the original square \(R\to R'\), \(S\to S'\), the original multiplicative subset \(W\subset S\otimes_RR'\), and \(S'=W^{-1}(S\otimes_RR')\). The associativity map identifies \[
S'\otimes_SM \cong W^{-1}(M\otimes_RR').
\] Before localization its inverse balanced maps are \((s\otimes r')\otimes m\mapsto sm\otimes r'\) and \(m\otimes r'\mapsto(1\otimes r')\otimes m\); the same formulas on fractions give the localized comparison. If \(M\) is \(R\)-flat, base change makes \(M\otimes_RR'\) flat over \(R'\). Localizing this module in the auxiliary algebra preserves \(R'\)-flatness, since it is a filtered colimit of copies of the \(R'\)-flat module with its multiplication maps.

If in addition \(R\to R'\) is flat, base changing it and then localizing proves flatness of \(S\to S'\). This is also a local map, so it is faithfully flat. For an ideal \(J\subset R\), flatness of \(R'\) identifies \(J\otimes_RR'\) with its actual image \(JR'\). The source's associativity isomorphism \[
(J\otimes_RM)\otimes_SS'\cong JR'\otimes_{R'}M'
\] sends \((j\otimes m)\otimes s'\) to \(\varphi(j)\otimes s'm\), using the canonical \(S'\)-module image of \(m\). Equivalently, before replacing \(J\otimes_RR'\) by its image it is the tensor associativity map, whose inverse moves the \(R'\)-coefficient into \(M'\). It intertwines multiplication into \(M'\). If the original map has kernel \(K\), \(S\)-flatness identifies the new kernel with \(K\otimes_SS'\). Faithfulness makes this nonzero when \(K\ne0\). But \(R'\)-flatness of \(M'\) makes the new multiplication map injective. Thus \(K=0\), proving \(R\)-flatness by the ideal criterion. Taking \(M=S\) gives the two ring statements.

For the next square no localization assumption is made. Instead retain the source's two horizontal flat maps and \(\mathfrak m_RR'=\mathfrak m_{R'}\). The dependency lemma at 8929--8955, reread here, shows \(M'=M\otimes_SS'\) is \(R\)-flat: for every injection of \(R\)-modules the tensor kernel is the original kernel tensored over the flat map \(S\to S'\). Also \(M'\) is finite over \(S'\). The actual injection \(\mathfrak m_R\otimes_RM'\to M'\) identifies, by associativity and flatness of \(R'\), with \(\mathfrak m_{R'}\otimes_{R'}M'\to M'\). Consequently the residue-field Tor group over \(R'\) vanishes. The original local criterion applies because \(R'\) and \(S'\) are Noetherian. Neither Noetherianity of \(R,S\) nor a tensor-product-localization presentation is silently added to this lemma.

\subsubsection{Nilpotent basis lifting with arbitrary index sets}\label{algebra-proof-017-nilpotent-basis-lifting-with-arbitrary-index-sets}

Source 24294--24381. In every use of Nakayama here nilpotence supplies the result even for infinitely generated modules. If \(I^c=0\) and \(X=IX\), iteration gives \(X=I^cX=0\).

For the first lemma, use the actual map \(F=R^{(A)}\to M\) with \(e_\alpha\mapsto x_\alpha\). Its cokernel is equal to its product by the nilpotent maximal ideal, so it is zero. Let \(K\) be the kernel. Flatness of \(M\) makes \[
0\to K/\mathfrak mK\to F/\mathfrak mF\to M/\mathfrak mM\to0
\] exact, and the basis hypothesis makes the middle map an isomorphism. Thus \(K=\mathfrak mK\), and nilpotence gives \(K=0\). This proves the equivalence for the originally chosen elements, including an empty basis when \(M=0\). Taking any vector-space basis of the residue quotient proves the source equivalence of flat, free and projective for a local ring with nilpotent maximal ideal. No Artinian or finite-generation hypothesis is needed, and the source already states this generality.

For lemma-lift-basis, retain arbitrary \(R\), nilpotent \(I\), the original chosen elements, and \(\operatorname{Tor}_1^R(R/I,M)=0\). The same cokernel iteration proves \(F\to M\) is onto. The long exact sequence with \(F\) free and the assumed Tor vanishing gives \(0\to K/IK\to F/IF\to M/IM\to0\), hence \(K=IK\) and \(K=0\). This is a direct check of the source conclusion.

The source's coefficient iteration is also valid, with the following exact comparison filling its implicit step. For every \(a\geq1\), \[
\ker(F/I^aF\to M/I^aM)=(K+I^aF)/I^aF,
\] because the inverse image of \(I^aM\) under the surjection \(F\to M\) is \(K+I^aF\). Thus the assertion that every relation in \(M\) has all coefficients in \(I^a\) is exactly what is needed for the images of the chosen generators to be a basis over \(R/I^a\).

For a finite relation \(\sum_\alpha f_\alpha x_\alpha=0\), its coefficients first lie in \(I\). The assumed injectivity of \(I\otimes_RM\to M\) makes \(\sum_\alpha f_\alpha\otimes x_\alpha=0\). If the basis has already been proved modulo \(I^a\), map this tensor to \[
(I/I^{a+1})\otimes_{R/I^a}(M/I^aM).
\] Its first factor is indeed an \(R/I^a\)-module, since \(I^aI\subset I^{a+1}\). The original basis identifies this tensor module with a direct sum of copies of \(I/I^{a+1}\), so every \(f_\alpha\) lies in \(I^{a+1}\). The kernel comparison proves the next basis assertion. Nilpotence terminates the process. Only finitely supported relations are used; infinite index sets cause no change. This supplies the source's actual missing comparison as separate proof evidence, not an automatic replacement of its translated proof.

\subsubsection{The contracted square and all quotient comparison maps}\label{algebra-proof-017-the-contracted-square-and-all-quotient-comparison-maps}

Source 24383--24442. Keep \(\varphi:R\to R'\), \(I\), \(M\), \(M/IM\) flat over \(R/I\), and \(R'\otimes_RM\) flat over \(R'\). Define the original ideal \[
I_2=\varphi^{-1}\bigl(\varphi(I^2)R'\bigr).
\] Write \(J=\varphi(I)R'\); then \(J^2=\varphi(I^2)R'\) by finite sums of products. Put \[
\overline R=R/I_2,\quad \overline M=M/I_2M,\quad
\overline R'=R'/J^2,\quad \overline I=(I+I_2)/I_2 .
\] The induced \(\overline R\to\overline R'\) is injective by the definition of \(I_2\), and \(\overline I^2=0\), since \(I^2\subset I_2\). The residue module is \[
\overline M/\overline I\overline M
 =M/(I+I_2)M
 \cong (M/IM)\otimes_{R/I}R/(I+I_2),
\] which is flat over \(R/(I+I_2)\) by base change. Also the canonical map \[
\overline M\otimes_{\overline R}\overline R'
 \cong (M\otimes_RR')\otimes_{R'}(R'/J^2)
\] is an isomorphism: on simple tensors both sides send the original \(m\) and \(r'\) to \(m\otimes(r'\bmod J^2)\); balancing by \(I_2\) is permitted because its image is in \(J^2\). The right side is flat over \(\overline R'\). Thus every original hypothesis is transported by an explicit quotient map.

In the reduced situation \(I^2=0\) and \(\varphi\) injective, the ideal \(I'=IR'\) is an \(R/I\)-module because \(I\cdot I'=I^2R'=0\), and \(I\to I'\) is injective. Flatness of \(M/IM\) gives injection \[
I\otimes_{R/I}M/IM\longrightarrow I'\otimes_{R/I}M/IM.
\] The canonical balanced isomorphisms identify these modules with \(I\otimes_RM\) and \(I'\otimes_{R'}(M\otimes_RR')\), respectively. The latter multiplication map into \(M\otimes_RR'\) is injective by its flatness over \(R'\). The original kernel \(\ker(I\otimes_RM\to M)\) maps to zero under this composite and is therefore zero. Applying the source's nilpotent Tor criterion proves flatness of \(M\) in the reduced situation, hence flatness of \(M/I_2M\) over \(R/I_2\) in the original one.

For lemma-lift-flatness, retain exactly \[
I_1=I,\qquad
I_{n+1}=\varphi^{-1}\bigl(\varphi(I_n)^2R'\bigr).
\] The preceding lemma applies successively because its second flatness hypothesis always concerns the same original \(M\otimes_RR'\), while its first is the previously obtained flat quotient. Thus every \(M/I_nM\) is flat over \(R/I_n\). If \(\varphi\) is injective and \(I\) is nilpotent, the source's powers force some \(I_n=0\), giving the desired flatness. The sharper exact recurrence and kernel quotient, without silently changing \(I_1\), are proved in the later section.

\subsubsection{The Artinian statement and the retained module corrections}\label{algebra-proof-017-the-artinian-statement-and-the-retained-module-corrections}

Source 24447--24517. The Artinian-local ideal criterion follows from the source's previously proved nilpotent-ideal criterion, since every proper ideal lies in its nilpotent maximal ideal. The printed proof also works: the quotient has nilpotent maximal ideal, its flat module is free by the preceding source theorem, and its basis lifts by the just-proved lemma. The general nilpotent criterion is already source content, not a newly discovered generalization here.

For arbitrary Artinian \(R\), its Jacobson radical \(I\) is nilpotent and \(R/I\) is a finite product of fields. A module over that finite product is the finite direct sum of its idempotent components, each a vector space and hence flat. Therefore \(M/IM\) is flat, and the injective-map nilpotent lifting lemma proves the first source proof. The introductory phrase receives the independent grammatical correction ``injective homomorphisms''.

In the second proof, OCC-12145 and OCC-00496 retain MC-STK-ERR-0488, replacing the repeated \(R\) by the module \(M\). Let \(1=\sum_{i=1}^t e_i\) be the original orthogonal central idempotents of \(R=\prod_i R_i\). The actual maps \[
S\longrightarrow\prod_i e_iS,\quad s\longmapsto(e_is)_i,
\qquad
M\longrightarrow\prod_i e_iM,\quad m\longmapsto(e_im)_i
\] have inverses given by the finite sums of their components. Orthogonality proves both composites are identities. Each \(R_i=e_iR\to e_iS\) is injective. Moreover \[
e_i(M\otimes_RS)\cong(e_iM)\otimes_{R_i}(e_iS)
\] via \((e_im)\otimes(e_is)\mapsto e_im\otimes s\) and the inverse obtained by applying \(e_i\) to both entries. Flatness of \(M\otimes_RS\) over \(S\) passes to these factors, so the local statement applies. Flatness of all \(e_iM\) over \(R_i\) gives flatness of \(M\) by these same finite decompositions and tensor exactness componentwise.

OCC-12146 and OCC-00494 retain MC-STK-ERR-0489: the chosen \(x_\alpha\) belong to \(M\), indexed by \(A\). When \(R\) is local Artinian, nilpotent Nakayama makes these lifts of a residue basis generate \(M\). For \(N=S\otimes_RM\) and \(I=\mathfrak mS\), the canonical comparison \[
N/IN\cong(S/\mathfrak mS)\otimes_{R/\mathfrak m}(M/\mathfrak mM)
\] carries \(1\otimes x_\alpha\) to the base-changed basis. The ideal \(I\) is nilpotent, and flatness of \(N\) gives its required Tor vanishing. The lift-basis lemma says the same chosen elements are an \(S\)-basis of \(N\). Any finite relation \(\sum r_\alpha x_\alpha=0\) in \(M\) becomes a relation with coefficients \(\varphi(r_\alpha)\) in this \(S\)-basis; all vanish, and injectivity of \(R\to S\) makes all \(r_\alpha=0\). Thus the chosen lifts are the required \(R\)-basis. This verifies the original second proof with its actual elements and coefficients.

\subsubsection{The nilpotent fibre criterion and its faithful support}\label{algebra-proof-017-the-nilpotent-fibre-criterion-and-its-faithful-support}

Source 24528--24567. OCC-12147 and OCC-00495 retain MC-STK-ERR-0490: the exact quotient argument is \(S/IS\) in the first Tor slot.

Flatness over \(R\) injects the original \(I\otimes_RM\) into \(M\). Multiplication first factors through the surjection \(I\otimes_RM\to IS\otimes_SM\), induced by \(I\otimes_RS\to IS\). Hence both arrows in the source factorization are injective, the first being an isomorphism. The exact sequence \(0\to IS\to S\to S/IS\to0\) identifies \[
\ker(IS\otimes_SM\to M)=\operatorname{Tor}_1^S(S/IS,M)=0.
\] Together with the given \(S/IS\)-flatness of \(M/IM\) and nilpotence of \(IS\), this proves \(S\)-flatness of \(M\) by the original criterion. Although eventual \(S\)-flatness also makes the misprinted Tor group zero, that fact is not available as the stated justification; the retained quotient correction repairs the actual inference.

Let \(\mathfrak q\subset S\) have nonzero original fibre \(M\otimes_S\kappa(\mathfrak q)\). Then \(M_{\mathfrak q}\) is flat over \(S_{\mathfrak q}\) and has nonzero residue quotient, hence is faithfully flat over that local ring. It is also \(R\)-flat: localization in \(S\) is a filtered colimit of copies of the \(R\)-flat \(M\). For each ideal \(J\subset R\), let \[
K_J=\ker(J\otimes_RS_{\mathfrak q}\to S_{\mathfrak q}).
\] Tensoring with the \(S_{\mathfrak q}\)-flat \(M_{\mathfrak q}\) identifies \(K_J\otimes_{S_{\mathfrak q}}M_{\mathfrak q}\) with the kernel of \(J\otimes_RM_{\mathfrak q}\to M_{\mathfrak q}\), which is zero. Faithfulness gives \(K_J=0\). The ideal criterion proves \(S_{\mathfrak q}\) flat over \(R\), using the exact maps from the preceding batch's group 559. No nonzero-fibre conclusion is imposed at the other primes. The unused additional action through \(S'\) is retained as part of the source diagram.

\subsubsection{Exact contraction powers and the kernel quotient}\label{algebra-proof-017-exact-contraction-powers-and-the-kernel-quotient}

The square-contraction argument gives a stronger, quantitative statement even when \(\varphi\) is not injective. Continue to assume \(I^c=0\) for some integer \(c\geq1\), \(M/IM\) flat over \(R/I\), and \(M\otimes_RR'\) flat over \(R'\). Put \(J=\varphi(I)R'\) and \(H=\ker\varphi\). Retain the original \(I_n\).

For any ideal \(L\subset R\), its contraction \(C=\varphi^{-1}((\varphi(L)R')^2)\) satisfies \[
\varphi(C)R'=(\varphi(L)R')^2.
\] One inclusion is the definition of contraction. For the reverse, \(L^2\subset C\), and \(\varphi(L^2)R'=(\varphi(L)R')^2\), so extending \(C\) contains the entire square. Induction therefore gives exact equalities \[
\varphi(I_n)R'=J^{2^{\,n-1}}\quad(n\geq1),
\qquad
I_n=\varphi^{-1}(J^{2^{\,n-1}})\quad(n\geq2).
\] The second formula need not hold for \(n=1\), because \(I\) need not be contracted. In particular \(I_2\subset I_1\) is not assumed; only the contracted stages from \(n=2\) onward are visibly decreasing.

Choose \(n\geq2\) with \(2^{n-1}\geq c\), for example \(n=\max\{2,1+\lceil\log_2c\rceil\}\). Then \(J^{2^{n-1}}=0\), so the exact formula gives \(I_n=H\). Iteration of the already proved square-contraction lemma has shown \(M/I_nM\) flat over \(R/I_n\) at every stage. Hence the actual quotient \[
M/HM\quad\hbox{is flat over }R/H.
\] Its unit map and scalar action are the original quotient maps; no flatness over \(R\) is inferred unless further hypotheses justify it. If \(\varphi\) is injective, this proves the source conclusion with an explicit termination bound. If \(I=0\), the initial hypothesis already makes \(M\) flat over \(R\); the stated later-stage formula still holds.

The quotient cannot be discarded. For \(R=k[\epsilon]/(\epsilon^2)\), \(R'=k\), \(I=(\epsilon)\), and \(M=k\), both source flatness hypotheses hold and \(H=I\). The quotient \(M/HM=k\) is flat over \(R/H=k\). But \(M\) is not \(R\)-flat: \(I\cong k\) as an \(R\)-module, so \(I\otimes_RM\cong k\ne0\), while its multiplication map to \(M\) is zero. This retains the exact obstruction supplied by the kernel.

\subsubsection{Injective descent beyond Artinian source rings}\label{algebra-proof-017-injective-descent-beyond-artinian-source-rings}

A separate consequence of the original nilpotent lifting lemma is the following. Suppose \(R\) has a nilpotent ideal \(I\) such that every \(R/I\)-module is flat. Then for every injective map \(R\to S\) and every \(R\)-module \(M\), flatness of \(M\otimes_RS\) over \(S\) implies flatness of \(M\) over \(R\). Indeed the residue hypothesis makes the original \(M/IM\) flat, and the proved nilpotent lifting argument applies without any Artinian assumption. The reverse implication is ordinary flat base change. Thus, under this stated property of the source ring, flatness is equivalent before and after every injective extension.

This applies when \(R/I\) is an arbitrary product of fields, not only a finite product. Here is a proof that does not incorrectly decompose every module over an infinite product into its coordinate modules. Let \(B=\prod_{t\in T}k_t\). For \(a=(a_t)\in B\), set \(b_t=a_t^{-1}\) when \(a_t\ne0\), and \(b_t=0\) otherwise. Then \(e=ab\) is idempotent, \(a=ea\), and \((a)=eB\). For finitely many generators \(a_1,\ldots,a_m\), form their idempotents \(e_i\) and set \[
e=1-\prod_{i=1}^m(1-e_i).
\] This is idempotent, satisfies \(ee_i=e_i\), and belongs to the ideal generated by the \(e_i\), as seen by expanding the finite product. Consequently the original finitely generated ideal is exactly \(eB\). The zero ideal is the empty case \(e=0\).

For any \(B\)-module \(X\), multiplication identifies \(eB\otimes_BX\) with \(eX\). Its inverse sends \(ex\) to \(e\otimes x\); if \(e(x-y)=0\), then \(e\otimes(x-y)=e\otimes e(x-y)=0\), so the inverse is well defined. Thus the map of this tensor into \(X\) is injective for every finitely generated ideal of \(B\). The ideal criterion makes \(X\) flat, as required.

For an explicit source ring beyond the Artinian case, take \(R=\prod_{n\geq1}k[\epsilon]/(\epsilon^2)\) with its original ideal \(I=\prod_{n\geq1}(\epsilon)\). Then \(I^2=0\) and \(R/I=\prod_{n\geq1}k\), so the result applies. The ring is not Artinian: the ideals whose first \(n\) coordinates vanish form a strictly descending chain, with the coordinate idempotent at position \(n+1\) witnessing strictness. This example and the ideal-tensor maps prove a substantive extension of the reviewed Artinian statement, without a claim of novelty or an assertion that later source sections have been surveyed.

Both new consequences are retained in the claim ledger with their exact maps and boundaries. They remain separate editorial results; broader cross-work synthesis and its short paper follow the core correction accounting.

\hypertarget{algebra_complex_exactness_notes_20260928}{}
\subsection{Complex exactness: source review and conceptual errata}\label{algebra-proof-018-complex-exactness-source-review-and-conceptual-errata}

Primary source: the Stacks Project, algebra.tex at authority commit a04446e57ec1fbc252a871afcec7752fb2807b14, SHA256 FA8BB92E58A4F78A2BD01B3B6A4A87DE0A0D279F5DD90641B574DD5FBFFFA4F3. Read and adjudicated: 24576--24951. The depth lemma at 17904--17937 is reread as a dependency. The chapter explicitly permits the zero ring. Lookahead 24952--25120 is read but unadjudicated.

The source cites Buchsbaum--Eisenbud, WhatExact, and Peskine--Szpiro, Lemma 1.8. Those original references are retained. Their papers have not been independently read in this batch. The actual corpus query is saved separately; its sole hit supplies PDF paths, no accessible original TeX in that routing record, and is not read. This is a bounded source review, not a novelty survey.

These complete arguments are editorial evidence. They do not replace translated source passages. Minimal proposed corrections preserve identifiable original readings. No chapter or translation has been mutated. The wider synthesis remains after core completion.

\subsubsection{The original pivot and its exact splitting maps}\label{algebra-proof-018-the-original-pivot-and-its-exact-splitting-maps}

At 24607--24654 write the original map as \[
 \varphi_i(e_j)=\sum_k a_{jk}f_k.
\] Choose the actual invertible coefficient \(u=a_{pq}\); no coefficient or pivot factor is suppressed. Put \[
 E_p=e_p,\quad E_j=e_j-a_{jq}u^{-1}e_p\ (j\ne p),
 \qquad F_q=\sum_k a_{pk}f_k,\quad F_k=f_k\ (k\ne q).
\] These are bases, with inverse comparisons \[
 e_p=E_p,\quad e_j=E_j+a_{jq}u^{-1}E_p,\qquad
 f_q=u^{-1}\left(F_q-\sum_{k\ne q}a_{pk}F_k\right),\quad f_k=F_k.
\] Direct substitution, retaining every term, gives \[
 \varphi_i(E_p)=F_q,\qquad
 \varphi_i(E_j)=\sum_{k\ne q}
 (a_{jk}-a_{jq}u^{-1}a_{pk})F_k\quad(j\ne p).
\] For an incoming vector with original coordinates \(c_j\), the complex identity says \(\sum_j a_{jq}c_j=0\). Its \(E_p\)-coordinate is therefore \(c_p+\sum_{j\ne p}a_{jq}u^{-1}c_j=0\). The outgoing map kills \(F_q=\varphi_i(e_p)\). Consequently the spans of \(E_p,F_q\) form precisely the identity disk and their complementary spans, with all other original terms, form a subcomplex. The inclusion and projection maps in these bases are inverse direct-sum comparisons.

For completeness the disk projections in original coordinates are \[
 Q_i\!\left(\sum_j c_je_j\right)
   =u^{-1}\!\sum_j a_{jq}c_j\,e_p,\qquad
 Q_{i-1}\!\left(\sum_k b_kf_k\right)
   =u^{-1}b_q\sum_k a_{pk}f_k .
\] The degree-one map \(h_{i-1}(\sum b_kf_k)=u^{-1}b_qe_p\), with all other components zero, satisfies \(dh+hd=Q\); the two displayed formulas give the only nonzero components. Thus the complementary projection is also an explicit chain deformation retraction. The incoming and outgoing complex identities prove compatibility at the neighboring degrees.

In the source's already chosen coordinates \(p=q=1,u=1,a_{1k}=0\) for \(k>1\), these are exactly its \(e'_j\) and \(f_k\). The span at 24649 is an \(R\)-submodule, and the literal original unprimed \(e_j\) must be replaced there by \(e'_j\). For example over \(\mathbb Z\), take the map \((c_1,c_2)\mapsto c_1+c_2\) and incoming generator \(e_2-e_1\). The image lies in the span of \(e'_2=e_2-e_1\), not the span of the original \(e_2\). Reading the source as silently relabelling the changed basis gives the intended argument; the proposed notation repair makes that intent explicit.

Over a nonzero local ring a coefficient outside the maximal ideal is a unit. Each splitting reduces the sum of free ranks by two, so finitely many splittings leave a complex whose differentials have entries in that ideal. Identity disks have zero homology. Degree-zero singleton summands retain exactly the original unrestricted degree-zero homology.

\subsubsection{Depth-zero splitting and the notation corrections}\label{algebra-proof-018-depth-zero-splitting-and-the-notation-corrections}

At 24677--24713 the hypothesis supplies \(0\ne x\in R\) with \(\mathfrak m x=0\). If there are no positive-degree terms, the complex is a sum of degree-zero singletons; this includes the empty, zero complex. Otherwise choose the largest \(i>0\) with \(n_i>0\). Exactness makes \(\varphi_i\) injective. If all its entries belonged to \(\mathfrak m\), then \(xe_1\ne0\) would lie in its kernel. Hence a unit entry exists. Split its identity disk by the preceding exact maps and repeat. The sum of positive-degree free ranks strictly decreases. This proves the original statement, including the case in which the displayed highest degrees have rank zero.

For an Artinian nonzero local ring, its maximal ideal is nilpotent. If \(\mathfrak m=0\), take \(x=1\). Otherwise take a nonzero element of the last nonzero power of \(\mathfrak m\). It is annihilated by \(\mathfrak m\). The source's Artinian corollary therefore follows. The article at 24704 is corrected to ``an''; both notation declarations at 24693 and 24777 receive ``by''. These are grammatical changes, not altered assumptions.

\subsubsection{Determinantal ranks and the zero-ring boundary}\label{algebra-proof-018-determinantal-ranks-and-the-zero-ring-boundary}

For a nonzero ring the zeroth exterior power of every free-module map is the identity of \(R\), so the set in Definition 24721--24732 is nonempty. For positive degrees the exterior matrices are the actual minors. The rank is therefore the largest size of a nonzero minor, with rank zero precisely for the zero map and the zero-size minor equal to \(1\).

The zero ring is explicitly permitted by the chapter. Over it all exterior powers, including the zeroth, are zero modules, and every exterior map is zero. The displayed maximum is then a maximum of the empty set and is undefined. Merely declaring rank zero there would not repair the numerical statement of 24738--24763: the displayed identity disk with \(n_1=n_0=1\) would have rank zero, whereas its required alternating rank is one. Moreover all free modules of all displayed ranks are the same zero module. We therefore propose the explicit nonzero-ring hypothesis in the definition and in that numerical lemma. Over the zero ring all complexes in the situation are zero and exact; no numerical rank identity based on arbitrary chosen free presentations is asserted. All subsequent local-ring uses are already over nonzero rings and retain their hypotheses.

Here is the numerical proof over a nonzero ring. Let \(t_i\) be the number of identity disks on degrees \(i,i-1\), put \(t_{e+1}=0\), and let \(s\) be the number of degree-zero singleton summands. Then \[
 n_i=t_i+t_{i+1}\ (1\le i\le e),\qquad n_0=s+t_1.
\] The matrix of \(\varphi_i\) contains an identity block of size \(t_i\) and is zero elsewhere. Its exterior rank is exactly \(t_i\), since the corresponding minor is \(1\ne0\) and all larger minors vanish. Induction gives \[
 t_i=n_i-n_{i+1}+\cdots+(-1)^{e-i}n_e .
\] The maximal-minor ideal is \(R\), also when \(t_i=0\). The stray hyphen in ``\(r_{i+1}\)-basis elements'' is a grammatical error.

For later use let \(A\) have rank \(r\). Adding a zero row or column leaves its nonzero minors and maximal-minor ideal unchanged. Adding an identity disk changes the relevant matrix to \({\rm diag}(A,1)\). Its minors of size \(r+1\) consist of minors of \(A\) of size \(r\), with the pivot row and column, and minors of \(A\) of size \(r+1\), which are zero. Thus its rank is exactly \(r+1\) and its maximal-minor ideal is the original \(I(A)\). Minors of larger size vanish for the same reason. Changes of bases preserve these ideals: the minors of a product are sums of products of minors by the exterior-power multiplication formula, giving one inclusion; applying the inverse change of basis gives the reverse inclusion. Permutation signs and unit factors remain in the actual matrix comparisons; only equality of the generated ideals is used.

\subsubsection{Reduction by a nonzerodivisor and the acyclicity proof}\label{algebra-proof-018-reduction-by-a-nonzerodivisor-and-the-acyclicity-proof}

For 24765--24790, multiplication by the original \(x\) is injective on each finite free \(F_j\). Hence the displayed short exact sequence of complexes is exact degree by degree. The connecting map is explicit: lift a cycle \(\bar z\in F_i/xF_i\) to \(z\), write \(dz=xw\), and send it to \([w]\). Injectivity of \(x\) on \(F_{i-2}\) gives \(dw=0\). A different lift changes \(w\) by a boundary, and a lifted boundary changes it by a boundary as well. This gives the long exact homology sequence with the source's signs. For \(i\ge2\) both neighboring positive homology modules of \(F\) vanish, so the asserted quotient homology vanishes. In degree one the more precise statement is \[
 H_1(F/xF)\ \simeq\ \ker(x:H_0(F)\longrightarrow H_0(F)).
\] This identifies why exactness in degree one is not part of the source's conclusion. These maps also show that neither locality nor \(x\in\mathfrak m\) is needed for this auxiliary lemma, although the following proposition uses exactly those original hypotheses.

For the acyclicity lemma put \(d_j:M_j\to M_{j-1}\), \(B_j=\operatorname{im}d_{j+1}\), \(Z_j=\ker d_j\), and \(B_e=0\). Let \(i>0\) be the largest degree of nonzero homology. For \(j>i\), exactness gives \(Z_j=B_j\), and there are actual short exact sequences \[
 0\longrightarrow B_j\longrightarrow M_j\longrightarrow B_{j-1}
 \longrightarrow0 .
\] Starting at \(j=e\), where \(B_{e-1}\simeq M_e\), the source depth lemma proves \(\operatorname{depth}B_i\ge i+1\). If \(i=e\), this bound means \(B_i=0\); if \(i=e-1\), it is the initial isomorphism. There is no need for nonexistent intermediate terms in either boundary case.

To handle zero modules without applying a nonzero-module statement to them, interpret these lower bounds as the vanishing of the corresponding \(\operatorname{Ext}^q_R(R/\mathfrak m,-)\). Zero modules have all these groups zero. When a short exact sequence contains a zero term, its other two terms are isomorphic and the required bounds follow directly. For nonzero terms use the three inequalities in 17904--17937.

The source's sequence \(0\to Z_i\to M_i\to M_{i-1}\) is exact where a condition is displayed but need not be surjective at its last term. Its actual image factor gives \[
 0\longrightarrow Z_i\longrightarrow M_i\longrightarrow B_{i-1}
 \longrightarrow0,\qquad B_{i-1}\subset M_{i-1}.
\] Since \(\operatorname{depth}M_i\ge i\ge1\), an \(M_i\)-regular element acts injectively on its submodule \(Z_i\). Thus \(Z_i\), if nonzero, has depth at least one. Finally \[
 0\longrightarrow B_i\longrightarrow Z_i\longrightarrow H_i
 \longrightarrow0
\] and the depth inequality give depth \(H_i\ge1\). Changing ``short exact sequences'' to ``exact sequences'' repairs the source's description of its original displayed list. The explicit image factor needed for the depth argument is supplied here without replacing that list with a reconstructed translated proof. ``Cohomology'' is corrected to ``homology'' under the source's lowering-degree convention.

\subsubsection{The exactness criterion with all rank-zero cases retained}\label{algebra-proof-018-the-exactness-criterion-with-all-rank-zero-cases-retained}

At 24833--24924 let \(F_j=R^{n_j}\). Identity-disk removal preserves exactness. The determinantal calculation above shows that it also preserves condition (2), with precisely the corresponding increase or decrease of \(r_j\) and no change of the maximal-minor ideal. The other alternating ranks are unchanged. Repeated removal reaches the source's minimal matrices.

Discard only zero terms at the left. If no positive-degree term remains, every remaining positive map has rank zero, \(I=R\), and both conditions hold. Otherwise keep the actual top nonzero degree \(e\ge1\), with \(n_e>0\). This is the missing boundary convention in the proof; it changes no module or nonzero differential.

Assume first exactness. At an associated prime \(\mathfrak q\), the localized ring has a nonzero socle, so the depth-zero lemma gives a sum of trivial complexes. Thus every localized rank is the required alternating \(r_i\) and every localized maximal-minor ideal is the entire localized ring. The standard injection \(R\to\bigoplus_{\mathfrak q\in\operatorname{Ass}R}R_{\mathfrak q}\) then makes every minor of size greater than \(r_i\) zero. At least one minor of size \(r_i\) survives some localization; for \(r_i=0\) it is the zeroth exterior identity. Hence the global rank is exactly \(r_i\). This verifies the use of the original determinantal ideals under localization.

For clarity the injection can be proved directly. If \(a\ne0\), the finite nonzero module \(Ra\) has an associated prime \(\mathfrak q\), so it contains \(0\ne ba\) with annihilator \(\mathfrak q\). This is an associated prime of \(R\) as well. If \(a_{\mathfrak q}=0\), some \(s\notin\mathfrak q\) kills \(a\), and hence kills \(ba\), a contradiction. The union of these finitely many associated primes is the set of zero divisors: an annihilator of a nonzero element lies in an associated prime of its cyclic submodule, while an associated-prime annihilator exhibits a killed nonzero element. Thus the source's prime-avoidance choice of \(x\) is a nonzerodivisor.

In fact every positive differential in this exact minimal complex has positive rank. If \(\varphi_j=0\), exactness gives \(F_j=\operatorname{im}\varphi_{j+1}\subset\mathfrak mF_j\), so Nakayama gives \(F_j=0\). Repeating this argument upward would give \(F_e=0\), contrary to the retained top nonzero term. Thus every \(I(\varphi_i)\) is proper. This fact is special to the exact minimal complex; it does not justify the source's converse assertion before exactness has been proved.

Each \(I(\varphi_i)\) avoids every associated prime, so their finite product does also; prime avoidance produces the original \[
 x\in I(\varphi_e)\cdots I(\varphi_1)
       \setminus\bigcup_{\mathfrak q\in\operatorname{Ass}R}\mathfrak q .
\] The minimal top matrix has positive rank \(r_e=n_e\), so its ideal lies in \(\mathfrak m\); therefore \(x\in\mathfrak m\). Every factor contains \(x\). For \(i\ge2\), reduce modulo this exact \(x\) and delete only the degree-zero term. The preceding quotient-complex lemma gives the required positive exactness of this shorter complex. The inductive rank formula, with all degrees lowered by one, is still \(n_i-n_{i+1}+\cdots+(-1)^{e-i}n_e=r_i\). Consequently its maximal-minor ideal is exactly \(I(\varphi_i)/xR\), not a differently sized determinantal ideal.

Induction supplies a regular sequence of length \(i-1\) in every proper one of those quotient ideals. Lifting it inside \(I(\varphi_i)\) and prepending the original \(x\) gives a regular sequence of length \(i\). For \(i=1\), \(x\) itself suffices whenever the ideal is proper. All these ideals are proper by the preceding exactness argument, so the source's phrase ``all the ideals'' at 24890 is justified once the zero top terms have been removed. No additional correction to that sentence is needed. Before cancellation, any unit ideal already satisfies (2)(b); disk removal retains that alternative through its unchanged ideal.

Conversely assume (2) for the minimal complex, with \(e,n_e\) as above. The positive rank \(r_e=n_e>0\) implies \(I(\varphi_e)\subset\mathfrak m\). Condition (2)(b) therefore gives \(\operatorname{depth}R\ge e\). It is not valid to infer this containment for every \(i\): a rank-zero map has \(I(\varphi_i)=R\), even when all its entries lie in \(\mathfrak m\).

Proceed by induction on the dimension of the nonzero Noetherian local ring. It is finite because its finitely generated maximal ideal bounds dimension by the height theorem. If the dimension is zero, no proper ideal contains a regular element; the conditions force every \(r_i=0\), so backward recursion forces every positive \(n_i=0\). This is the already treated zero-tail case. At any prime \(\mathfrak p\), each \(I(\varphi_i)\) contains a nonzerodivisor (take \(1\) if it is the unit ideal); that element cannot vanish upon localization. Larger minors were already zero and at least one maximal minor survives, so ranks are unchanged. A regular sequence inside a proper localized ideal remains regular: injectivity localizes, and the final quotient is nonzero because all its entries belong to the local maximal ideal. If the localized ideal is the unit ideal, its other alternative holds. Thus (2) localizes exactly.

At a nonmaximal prime, the dimension of \(R_{\mathfrak p}\) is smaller: each chain there extends strictly by \(\mathfrak m\). The inductive conclusion gives positive exactness. Every positive homology module of the original finite complex is therefore finite with support contained in \(\{\mathfrak m\}\). A nonzero such module has finite length and a nonzero socle, hence depth zero. On the other hand every nonzero \(F_j\) has depth at least \(e\ge j\), by the regular sequence just obtained, and the acyclicity lemma says that the highest nonzero positive homology would have depth at least one. This contradiction proves exactness.

For Remark 24926--24940 the maximal-minor ideals are invariant under all the preceding disk cancellations. In the minimal complex, \(I(\varphi_j)=R\) if and only if \(\varphi_j=0\). If \(\varphi_j=0\), exactness at \(F_j\) says that \(\varphi_{j+1}:F_{j+1}\to F_j\) is onto. Its entries lie in \(\mathfrak m\), so \(F_j=\mathfrak m F_j\), and Nakayama gives \(F_j=0\). The same argument then gives \(F_{j+1}=0,F_{j+2}=0,\ldots\). Thus the indices with unit ideal form a terminal interval, including the possibility of the empty interval with threshold \(e+1\). This is a separate direct justification of the original remark, not a replacement of its argument by longer regular sequences.

\subsubsection{Splitting characterized beyond Noetherian depth zero}\label{algebra-proof-018-splitting-characterized-beyond-noetherian-depth-zero}

Let \(R\ne0\) be local with maximal ideal \(\mathfrak m\). The following are equivalent.

\begin{enumerate}
\def\labelenumi{(\alph{enumi})}
\item
  Every finitely generated proper ideal of \(R\) has a nonzero annihilator.
\item
  Every finite complex of finite free \(R\)-modules which is exact in positive degrees is a direct sum of identity disks and degree-zero singletons.
\end{enumerate}

Assume (a). At the highest nonzero positive degree, the differential is injective. If all its coefficients lie in \(\mathfrak m\), let \(J\) be the ideal generated by the coefficients in its first column. This ideal is finitely generated and proper. A nonzero \(x\in\operatorname{ann}J\) makes \(xe_1\ne0\) lie in the kernel, a contradiction. There is therefore a unit coefficient. Apply the explicit splitting maps and induct on the sum of positive free ranks. This proves (b), without Noetherianity and without assuming a single element annihilates the entire maximal ideal.

Conversely let \(J=(a_1,\ldots,a_t)\) be proper with zero annihilator. Then \(t>0\), and the actual map \[
 R\longrightarrow R^t,\qquad r\longmapsto(a_1r,\ldots,a_tr)
\] is injective. By (b), its positive-degree term lies entirely in identity disks, so this injection has a left inverse. Any \(R\)-linear map \(R^t\to R\) is given by coefficients \(b_1,\ldots,b_t\); being a left inverse requires \(\sum b_ja_j=1\), contradicting properness. Hence (a). The zero ideal itself has annihilator \(R\ne0\), covering the empty generating family. In particular (a) implies that every module admitting a finite resolution by finite free modules is free: it is the degree-zero homology in (b). The converse follows by splitting such a cokernel and the same left-inverse argument.

A nonzero socle implies (a), so this propagates the source's depth-zero lemma beyond the Noetherian hypothesis. The condition can hold even with zero socle. Here is a fully specified example. For a field \(k\), let \[
 R=k[\mathbb Q_{\ge0}]/(t^1).
\] Its elements have unique finite sums of monomials \(t^q\) with \(0\le q<1\); products with exponent at least one are zero. The positive-exponent ideal \(\mathfrak m\) is the unique maximal ideal: every element of it is nilpotent, while an element \(c+v\), \(c\ne0\), has inverse \(c^{-1}\sum_{j=0}^{N-1}(-c^{-1}v)^j\) when \(v^N=0\). A nonzero finitely generated proper ideal has generators with only positive exponents. Let \(\delta>0\) be the least exponent occurring among them. Then \(0<\delta<1\), and the nonzero element \(t^{1-\delta}\) annihilates every generator. Thus (a) holds. For any nonzero \(w\in R\), choose rational \(0<q<1-\max(\operatorname{supp}w)\). The product \(t^qw\ne0\), since every exponent remains below one and the distinct shifted exponents remain distinct. Thus \(\operatorname{ann}\mathfrak m=0\). Finally \(\mathfrak m\) is not finitely generated: if its generators had least exponent \(\delta>0\), their ideal could not contain \(t^{\delta/2}\). This is a non-Noetherian example, with all monomial support data retained.

This consequence receives the source unit-entry lemma and propagates it to the source depth-zero and Artinian splitting statements. It is separate editorial mathematics; no novelty is claimed and no hypothesis is silently removed from those source statements.

\subsubsection{A stronger depth conclusion for the highest nonzero homology}\label{algebra-proof-018-a-stronger-depth-conclusion-for-the-highest-nonzero-homology}

Under exactly the source acyclicity hypotheses, let \(i>0\) be the highest index with \(H_i\ne0\). Then the argument actually gives \[
 \operatorname{depth}H_i\ \ge\ \min\{i,2\}.
\] The preceding tail argument already proves \(\operatorname{depth}B_i\ge i+1\). If \(i=1\), the source proof supplies the desired bound one. If \(i\ge2\), the hypothesis on \(M_{i-1}\) gives it depth at least one. Its submodule \(B_{i-1}\), when nonzero, also has depth at least one, because any regular element on the ambient module remains injective on the submodule. The actual image-factor short exact sequence gives \[
 \operatorname{depth}Z_i\ge
 \min\{\operatorname{depth}M_i,\operatorname{depth}B_{i-1}+1\}\ge2.
\] If \(B_{i-1}=0\), then \(Z_i=M_i\), giving the same inequality directly. Applying the quotient depth inequality to \(0\to B_i\to Z_i\to H_i\to0\) now gives depth \(H_i\ge2\); if \(B_i=0\), it follows from \(H_i=Z_i\). These cases cover every zero term without inventing a nonzero module for the depth lemma.

Consequently, if all positive homology modules have depth zero when nonzero, the complex is exact in positive degrees, as used in the source criterion. More generally, if every \(H_j\), \(j\ge2\), has depth at most one when nonzero, all those \(H_j\) vanish: a highest nonzero one would contradict the stronger bound. If also \(H_1\) has depth zero when nonzero, then it vanishes by the original degree-one bound. Finiteness of the complex supplies the highest index in each assertion.

The corrected image-factor argument is the precise additional input to this strengthening. The source exactness criterion only needs the original bound one and remains unchanged by this editorial consequence. Check later occurrences before treating the strengthening as absent from the whole corpus; it is not a novelty claim.

\hypertarget{algebra_cm_modules_notes_20260928}{}
\subsection{Cohen--Macaulay modules and rings: source review and conceptual errata}\label{algebra-proof-019-cohenmacaulay-modules-and-rings-source-review-and-conceptual-errata}

Primary source: the Stacks Project, algebra.tex at authority commit a04446e57ec1fbc252a871afcec7752fb2807b14, SHA256 FA8BB92E58A4F78A2BD01B3B6A4A87DE0A0D279F5DD90641B574DD5FBFFFA4F3. Complete new review: 24952--25455, both Cohen--Macaulay sections. Lookahead 25456--25480 is unadjudicated and ends before the last proof terminates.

Dependencies reread: 9487--9508, 14694--14710, 15297--15326, 16639--16687, 17765--17787, and 17961--18045. The original topology.tex at the same commit, SHA256 C6BAC8DCF8AD96DC47416BF34CB45BA4A10B894E40D67D3E1FA68D8EF0D9F872, is checked at 1397--1413 for the empty-space convention. The earlier complete proof in ALGEBRA\_BLOWUP\_EXT\_DEPTH\_NOTES\_20260927.md, section ``Depth quotients and propagation to large powers'', is reread and received explicitly below (group 463).

The original EGA Chapter 0, Proposition 16.5.4 reference is retained; this batch does not claim a fresh reading of EGA. The saved corpus query returns three routing hits. Two have PDF paths only; the third has original TeX but is not read in this batch. The supplied original Stacks TeX is the primary treatment used. No novelty or exhaustive literature coverage is asserted.

All complete arguments and further consequences here are source-linked editorial material. Proposed source corrections remain identifiable separately; the cumulative chapter and all translated branches are unchanged. The broader synthesis follows core completion.

\subsubsection{The zero module and the good-element boundary}\label{algebra-proof-019-the-zero-module-and-the-good-element-boundary}

The source sets \(\operatorname{depth}(0)=+\infty\) at 17772 and 17779--17782, while its topology definition gives dimension of the empty support as \(-\infty\). Therefore its literal equality at 24964 does not call the zero module Cohen--Macaulay. Later it asserts localization is Cohen--Macaulay at every prime and that checking maximal ideals suffices. These formulations are incompatible with that omission. For an explicit example let \(R=k[t]_{(t)}\), \(M=R/(t)\), and \(\mathfrak p=(0)\). The nonzero module \(M\) has dimension and depth zero, but \(M_{\mathfrak p}=0\). The later localization assertion thus fails under the literal initial definition.

The proposed conceptual correction states the zero-module exception explicitly: \(M\) is Cohen--Macaulay if \(M=0\), or if its support dimension equals its depth. It does not identify \(+\infty\) with \(-\infty\) or change either convention. Where an actual regular sequence is asserted, the nonzero hypothesis must still be retained: the source regular-sequence definition requires a nonzero final quotient, even for the empty sequence. Thus the proposition at 25054 is made explicitly about a nonzero Cohen--Macaulay module. The numerical depth definition of maximal Cohen--Macaulay at 25138 still excludes zero, so its explanatory ``largest possible depth'' comparison is restricted to nonzero finite modules. In the polynomial proof a zero localization is dealt with before choosing a finite length parameter sequence.

Independently, the given regular sequence in the good-element setup already forces \(M\ne0\), but it allows \(d=0\). Then the conditions indexed by \(0\le i\le d-1\) are empty, so every \(g\in\mathfrak m\) is good. Taking \(g=0\) on \(M=k\) over \(R=k\) contradicts the claimed injectivity. The minimal repair is \(d>0\) in the lemma, and induction beginning at \(d=1\). The definition of goodness itself is retained.

For \(d=1\) let \(K=\ker(g:M\to M)\). Localization away from \(V(g)\) makes \(g\) invertible, and \(K\subset M\); hence \[
 \operatorname{Supp}K\subset
 \operatorname{Supp}M\cap V(g)=\{\mathfrak m\}.
\] The containment is essential: the source writes equality even though its goal is \(K=0\), whose support is empty. This gets a separate minimal wording repair within the same lemma. A finite module supported at the closed point has finite length: a power of \(\mathfrak m\) annihilates it, and its finite filtration by powers of \(\mathfrak m\) has finite-dimensional residue-field factors. Multiplication by \(f_1\) is injective on \(K\), so equality of lengths makes it surjective. Nakayama gives \(K=0\). Also \(M/gM\ne0\) by Nakayama, has zero-dimensional support, and has depth zero.

For \(d>1\), the actual quotient \(M/f_1M\) has depth and support dimension \(d-1\) by the earlier depth-drop and dimension lemmas. Goodness descends through the support identity \[
 \operatorname{Supp}(M/(f_1,\ldots,f_i)M)
 =\operatorname{Supp}M\cap V(f_1,\ldots,f_i).
\] Induction gives \(g\) regular on \(M/f_1M\) and regularity of the stated tail on \(M/(f_1,g)M\). The original permutation lemma then gives the regular sequence \(g,f_1,\ldots,f_{d-1}\) on \(M\). The support of \(M/gM\) has dimension \(d-1\) by the \(i=0\) goodness condition. The tail has that length, so its depth reaches the dimension bound; the quotient is Cohen--Macaulay. This verifies the full original induction without replacing it.

\subsubsection{One-element cuts and complete parameter sequences}\label{algebra-proof-019-one-element-cuts-and-complete-parameter-sequences}

In the one-element lemma, first take nonzero \(M\) and let \(d=\dim\operatorname{Supp}M\). If its displayed dimension-drop hypothesis holds, then \(d\ge1\): for \(d=0\), Nakayama keeps \(M/gM\ne0\), so its support cannot have dimension \(-1\). The \(d=1\) case is the corrected good-element lemma.

For \(d>1\), every original quotient \(M/(f_1,\ldots,f_i)M\) is nonzero Cohen--Macaulay of dimension \(d-i\). The finitely many closed sets listed in 25034--25037 all have positive dimension. None has the maximal ideal as a minimal prime: each contains a nonmaximal prime below it, and the maximal ideal contains every prime of the local ring. Prime avoidance therefore chooses the source's \(h\in\mathfrak m\) avoiding every listed minimal prime. The one-equation dimension lemma, applied to their actual annihilator quotients, decreases each dimension exactly by one. Thus \(h\) is good for the original data and \(\dim(\operatorname{Supp}M\cap V(h,g))=d-2\).

The corrected good-element lemma makes \(M/hM\) Cohen--Macaulay. Induction makes \(g\) regular there and its quotient Cohen--Macaulay. Permuting the regular pair \(h,g\) gives \(g,h\), so \(g\) is regular on \(M\) and \(h\) on \(M/gM\). The latter quotient has support dimension \(d-1\), while its quotient by \(h\) has depth \(d-2\). The earlier depth-drop formula gives depth \(d-1\) for \(M/gM\), proving the assertion. With the explicit zero-module exception, the zero case is separate and tautological: all modules in that case are zero, and no finite-dimensional induction or regular sequence is asserted. Its infinite depth remains infinite.

For the parameter proposition retain \(M\ne0\), \(g_j\in\mathfrak m\), and the original equality of dimensions. Every quotient is nonzero by Nakayama. Put \[
 M_i=M/(g_1,\ldots,g_i)M,\qquad d_i=\dim\operatorname{Supp}M_i.
\] The maps are the actual quotient maps with kernel \((g_1,\ldots,g_i)M/(g_1,\ldots,g_{i-1})M\). The one-equation inequality gives \(d_i\le d_{i-1}\le d_i+1\). Since \(d_0=d\) and \(d_c=d-c\), every difference is exactly one. Repeatedly apply the one-element lemma to these exact modules. This proves regularity, Cohen--Macaulayness of each quotient, and depth \(d-i\). The omitted terminal \(M\) at 25068 is therefore a genuine type correction.

If \(c<d\), avoid the finitely many minimal primes of \(\operatorname{Supp}M_c\) to choose the next element in \(\mathfrak m\). The support dimension drops by one, and the same argument extends the sequence. It terminates at length \(d\); the depth bound forbids extension beyond this length. This includes \(c=0\) for nonzero \(M\).

The two article reports correct the original notation at 25027 and 25249. They do not change the chosen elements or their localizations.

\subsubsection{Quotients, associated primes, chains and localization}\label{algebra-proof-019-quotients-associated-primes-chains-and-localization}

If \(x\in\mathfrak m\) is regular on a nonzero finite \(M\), both its depth and its support dimension drop exactly by one. Hence the Cohen--Macaulay equivalence at 25081--25094 follows. If \(M=0\), both sides are zero and the corrected convention supplies the equivalence directly, without subtracting finite numbers from an empty support dimension in a proof.

For the omitted proof at 25096--25105 let \(R\to S=R/J\) be the given surjection. For a finite \(S\)-module \(N\), \(J\) annihilates \(N\). The exact correspondence \(\mathfrak q\mapsto\varphi^{-1}(\mathfrak q)\) is an inclusion-preserving bijection from \(\operatorname{Supp}_S N\) onto \(\operatorname{Supp}_R N\). Its inverse is \(\mathfrak p\mapsto\mathfrak p/J\). Localizations identify by \[
 N_{\mathfrak q}\longrightarrow N_{\mathfrak p},\qquad n/\bar s\longmapsto n/s,
\] where \(s\) is any lift outside \(\mathfrak p\); another lift differs by an element of \(J\), which acts as zero. The reverse map is given by reduction of denominators. Thus chains and their lengths are preserved exactly. A sequence in \(\mathfrak m_R\) and its image in \(\mathfrak m_S\) have the same successive multiplication maps and quotient modules. Conversely every finite sequence in \(\mathfrak m_S\) lifts into \(\mathfrak m_R\). Regularity, including the nonzero final quotient, is identical. Depths are therefore equal, and the support dimensions are equal. The zero case follows from the explicit exception. This is an editorial proof of the original equivalence, not replacement text.

For nonzero Cohen--Macaulay \(M\), each \(\mathfrak p\in\operatorname{Ass}M\) satisfies \[
 d=\operatorname{depth}M\le\dim R/\mathfrak p\le d.
\] Thus equality holds. If a support prime strictly contained \(\mathfrak p\), a chain of maximal length above \(\mathfrak p\) could be extended downward, contradicting dimension \(d\) of the support. Every associated prime is therefore minimal in the support, and the earlier theorem that minimal support primes are associated gives equality of the two sets. This proves the cited EGA consequence from the actual preceding Stacks results.

Now assume the original full-support hypothesis of 25150--25185. Let \(\mathfrak p_0\subsetneq\cdots\subsetneq\mathfrak p_n=\mathfrak m\) be a maximal chain. The first prime is minimal and every adjacent inclusion is saturated. If \(d=0\), \(n=0\). If \(d>0\), prime avoidance supplies \(x\in\mathfrak p_1\) outside all minimal primes of \(R\). The support and dimension statements imply \(x\) is \(M\)-regular and the quotient is Cohen--Macaulay, with support \(V(x)\) and dimension \(d-1\).

The prime \(\mathfrak p_1\) is minimal over \(\mathfrak p_0+(x)\), since \(x\notin\mathfrak p_0\) and no prime lies strictly between the two chain terms. Use the original cyclic submodule \(N_0\cong R/\mathfrak p_0\subset M\) and the Artin--Rees constant from group 463. For any integer \(a\) at least its bound, \(\mathfrak p_1\in\operatorname{Ass}(M/x^aM)\). The actual power \(x^a\) is still regular on \(M\). The module \(M/x^aM\) is Cohen--Macaulay over the distinct ring \(R/(x^a)\); its support is all that spectrum and its dimension is \(d-1\). We do not identify that ring or module with the quotients by \(x\); only their support sets coincide under the original prime correspondence.

The just-proved associated-prime result makes \(\mathfrak p_1/(x^a)\) a minimal prime of \(R/(x^a)\). The original remaining chain gives a maximal chain there of length \(n-1\). Induction gives \(n-1=d-1\). Thus every original maximal chain has length \(d\). For any prime \(\mathfrak p\), concatenate maximal chains below and above it. Their lengths are respectively \(\dim R_{\mathfrak p}\) and \(\dim R/\mathfrak p\), since every maximal concatenation has the same length. This proves the exact dimension formula at 25187--25199.

For localization at 25201--25225, first treat \(M_{\mathfrak p}=0\) by the explicit zero convention. This is the actual trivial case that must precede dimension or regular-sequence arguments. Otherwise \(M\ne0\) and \(\mathfrak p\in\operatorname{Supp}M\). The case \(\mathfrak p=\mathfrak m\) is identical to the original module. For the other primes, choose a saturated finite chain from \(\mathfrak p\) to \(\mathfrak m\) and induct along it. For an immediate predecessor, \(\dim R/\mathfrak p=1\). A chain in \(\operatorname{Supp}M_{\mathfrak p}\) extends by \(\mathfrak m\), giving \[
 \dim\operatorname{Supp}M_{\mathfrak p}\le d-1
 =\operatorname{depth}M-1
 \le\operatorname{depth}M_{\mathfrak p}.
\] The last inequality is the earlier depth-localization theorem, and the opposite dimension/depth inequality holds for the nonzero finite localized module. Equality proves the immediate step. Applying it successively proves the source localization theorem. The proof's opening clause is corrected to \(M_{\mathfrak p}\) being nonzero, also repairing the received missing verb.

For a general Noetherian ring, every prime lies in a maximal ideal and iterated localization identifies \((M_{\mathfrak m})_{\mathfrak pR_{\mathfrak m}}\) with \(M_{\mathfrak p}\) via the original fractions. Thus checking the maximal localizations suffices, now consistently including zero localizations.

\subsubsection{The complete polynomial-module argument}\label{algebra-proof-019-the-complete-polynomial-module-argument}

It suffices to add one variable at a time. Put \(T=R[x]\), take a maximal ideal \(\mathfrak n\subset T\), and put \(\mathfrak p=\mathfrak n\cap R\). If \(M_{\mathfrak p}=0\), then \((M\otimes_R T)_{\mathfrak n}=0\), and there is no parameter sequence to choose. Assume \(M_{\mathfrak p}\ne0\). Let \[
 A=R_{\mathfrak p},\quad B=T_{\mathfrak n},\quad
 L=M_{\mathfrak p}/(f_1,\ldots,f_d)M_{\mathfrak p}.
\] The original flat local map \(A\to B\) is faithfully flat. Indeed localization and polynomial extension are flat, the map is local, and its closed residue fibre is nonzero. Equivalently, every nonzero cyclic module has nonzero residue after this local flat base change, which is the usual faithful-flatness criterion. The chosen \(f_1,\ldots,f_d\) remain regular on \(M_{\mathfrak p}\otimes_A B\) by tensoring their injective maps; the nonzero successive final quotient is preserved by faithfulness.

The actual quotient \(Q=L\otimes_A B\) is nonzero. Its support is precisely \(V(\mathfrak pB)\): \(L\) is nonzero finite length over \(A\), so its annihilator has radical \(\mathfrak pA\). For a flat extension, the annihilator of a finite module extends exactly. To see this without assuming a tensor injection falsely, choose finitely many generators \(l_1,\ldots,l_s\). The annihilator is the kernel of the actual map \(A\to L^s\), \(a\mapsto(al_1,\ldots,al_s)\). Flatness identifies its tensor kernel with the kernel of \(B\to(L\otimes_A B)^s\), whose entries generate the extended module. This proves the assertion and hence the support formula.

The residue map gives a surjection \(\kappa(\mathfrak p)[x]\to T/\mathfrak n=\kappa(\mathfrak n)\). Its kernel \(\mathfrak q\) is a nonzero maximal ideal. Surjectivity holds because the image of \(R/\mathfrak p\) in the field contains its fractions after extending scalars, and the original field quotient is generated by the image of \(x\). A zero kernel would identify a polynomial ring in one indeterminate with a field, which is impossible. This nonzero maximal ideal is generated by an irreducible polynomial \(P\); division using its actual invertible leading coefficient gives the basis \(1,x,\ldots,x^{\deg P-1}\) for the quotient field over \(\kappa(\mathfrak p)\). This also verifies the source's finite residue-field extension. Thus \[
 B/\mathfrak pB=\kappa(\mathfrak p)[x]_{\mathfrak q}
\] has dimension one. In particular the original \(Q\) has support dimension one. Clearing the finitely many denominators in a nonzero polynomial of \(\mathfrak q\) gives \(f\in\mathfrak n\subset R[x]\) whose leading coefficient does not belong to \(\mathfrak p\). Such an \(f\) has positive degree.

Multiplication by \(f\) on \[
 L\otimes_R R[x]=\bigoplus_{j\ge0}Lx^j
\] is injective: for a nonzero polynomial vector, its highest nonzero coefficient times the unit leading coefficient of \(f\) remains nonzero and is the unique coefficient at the highest resulting degree. Localization preserves this injection on \(Q\). Also \(Q/fQ\ne0\) by Nakayama because \(f\in\mathfrak nB\) and \(Q\) is finite and nonzero. Hence depth \(Q\ge1\), and the dimension bound makes it exactly one. The original module localized at \(\mathfrak n\) has depth \(d+1\) by reversing the \(d\) depth drops, and support dimension \(d+1\) by reversing the corresponding dimension drops. It is Cohen--Macaulay. Induction proves the stated finite-variable module theorem.

The ring corollary says ``Any polynomial algebra'', although the definition requires Noetherian rings and the cited theorem has finitely many variables. The chapter elsewhere permits infinitely many variables. For a field \(k\), the chain \((x_1)\subsetneq(x_1,x_2)\subsetneq\cdots\) in \(k[x_1,x_2,\ldots]\) is strictly increasing: the next variable survives after setting the preceding variables to zero. Thus the ring is not Noetherian and is not Cohen--Macaulay under this chapter's definition. The correction explicitly says finitely many variables; it does not replace that definition by a different non-Noetherian notion.

\subsubsection{The ring specializations and the exact syzygy depths}\label{algebra-proof-019-the-ring-specializations-and-the-exact-syzygy-depths}

For the local ring results use the original module \(M=R\). Its support is the entire spectrum and it is nonzero. Regularity gives each successive quotient depth and dimension one smaller; conversely the required dimension drop gives regularity by the module proposition. All intermediate quotient rings are Cohen--Macaulay and the sequence extends to length \(d\). The article before \(R\)-regular is corrected.

The maximal-chain statement requires prime ideals. For a concrete adverse example, \(k[\epsilon]/(\epsilon^2)\) is a zero-dimensional Cohen--Macaulay local ring and its maximal chain of arbitrary ideals is \(0\subsetneq(\epsilon)\subsetneq R\), of length two, not zero. It has exactly those ideals because every element with a nonzero constant coefficient is a unit and every nonzero element of \((\epsilon)\) spans it. Thus ``prime'' is necessary, not merely stylistic. The two reports on the opening sentence share the same article-and-noun correction.

In the depth-shift lemma, if \(M=0\), its first alternative holds. If \(M\ne0\) and \(K=0\), depth \(K=+\infty>\operatorname{depth}M\), so the second alternative holds. Otherwise all required modules in the short exact sequence are nonzero, and write \(e=\operatorname{depth}M\le d\). The depth lemma gives \[
 \operatorname{depth}K\ge\min\{d,e+1\}.
\] In fact equality holds. If \(e=d\), the dimension bound on nonzero \(K\) gives equality \(d\). If \(e<d\) and depth \(K>e+1\), the other depth inequality would give \[
 e=\operatorname{depth}M\ge
 \min\{d,\operatorname{depth}K-1\}\ge e+1,
\] a contradiction. Thus for every nonzero kernel in this exact sequence \[
 \operatorname{depth}K=\min\{d,e+1\}.
\]

For the resolution lemma require \(M\ne0\), so \(0\le e\le d\) is a finite integer. If \(e=d\), take the actual identity complex \(0\to K=M\xrightarrow{1}M\to0\); there is no free segment and no \(F_{-1}\). If \(e<d\), put \(L_0=M\). For each \(j=0,\ldots,d-e-1\), choose a finite set of generators of \(L_j\), let \(G_j\to L_j\) be its finite free cover, and use \[
 F_j=G_j\oplus R\longrightarrow L_j,\qquad (v,r)\longmapsto v\text{ mapped to }L_j.
\] Its kernel \(L_{j+1}\) contains the actual extra summand \(R\), so it is nonzero. It is finite because \(R\) is Noetherian. This harmless freedom in the choice of a free cover is needed because the source definition of maximal Cohen--Macaulay excludes the zero module. The preceding exact formula gives \(\operatorname{depth}L_j=\min\{d,e+j\}\). Splicing the kernel inclusions and covers yields exactly the source complex, with \(K=L_{d-e}\) nonzero and depth \(d\). This supplies the complete existence argument without assuming a nonzero syzygy from an arbitrary minimal resolution.

\subsubsection{The original image-parameter construction and its module extension}\label{algebra-proof-019-the-original-image-parameter-construction-and-its-module-extension}

At 25425--25447 the radical hypothesis makes the original map local, and prime avoidance may choose \(f\) inside \(\mathfrak m_A\), not merely in \(A\). For positive dimension the minimal primes of the Cohen--Macaulay ring \(B\) are all strictly below \(\mathfrak m_B\). None contains \(\varphi(\mathfrak m_A)B\), since its radical is \(\mathfrak m_B\). The finitely many contracted prime ideals therefore do not contain \(\mathfrak m_A\), and prime avoidance chooses the required \(f\). Its image avoids all associated primes of \(B\), which are minimal, so it is regular and reduces the dimension by one. The same original map followed by \(B\to B/\varphi(f)B\) retains the radical hypothesis. Induction constructs the complete sequence. Dimension zero uses the empty sequence, whose final module \(B\ne0\) satisfies the source regularity convention.

More generally let \(A\to B\) be a local ring map, with \(B\) Noetherian local, and let \(N\ne0\) be a finite Cohen--Macaulay \(B\)-module of support dimension \(d\). It suffices to assume the weaker, support-relative condition \[
 \operatorname{Supp}\bigl(N/\varphi(\mathfrak m_A)N\bigr)
       =\{\mathfrak m_B\}.
\] Equivalently \(\sqrt{\varphi(\mathfrak m_A)B+\operatorname{ann}_B N}=\mathfrak m_B\); the equivalence is the actual support-of-quotient formula.

For \(d>0\), none of the finitely many minimal support primes is maximal. If one contained every image of \(\mathfrak m_A\), the radical equality would force it to contain \(\mathfrak m_B\), a contradiction. Prime avoidance in the original \(A\) chooses \(f_1\in\mathfrak m_A\) whose image avoids all those primes. The one-element module lemma makes \(\varphi(f_1)\) regular on \(N\) and gives a Cohen--Macaulay quotient of dimension \(d-1\). The original image ideal still cuts its support down to the closed point, because this new quotient's support is a subset of the old one, contains the maximal ideal, and is intersected with the same closed image-ideal set. Repeat on the actual successive quotients. This constructs \(f_1,\ldots,f_d\in\mathfrak m_A\) whose images are \(N\)-regular and whose final quotient has dimension zero. For \(d=0\), the empty sequence suffices. No Noetherianity of \(A\) or Cohen--Macaulayness of \(B\) itself is assumed. Taking \(N=B\) and the source radical hypothesis recovers its construction.

\subsubsection{Support-relative dimension and chain formulas}\label{algebra-proof-019-support-relative-dimension-and-chain-formulas}

Let \(R\) be Noetherian local and \(M\ne0\) finite Cohen--Macaulay of dimension \(d\). Retain the original module and put \(J=\operatorname{ann}_R M\), \(A=R/J\) solely to prove the comparison. The quotient-map argument above shows that the same \(M\) is Cohen--Macaulay over \(A\), with full support \(\operatorname{Spec}A\). For every \(\mathfrak p\in\operatorname{Supp}_R M\), write \(\bar{\mathfrak p}=\mathfrak p/J\). The actual prime and localization maps give \[
 A_{\bar{\mathfrak p}}\simeq R_{\mathfrak p}/JR_{\mathfrak p},\qquad
 A/\bar{\mathfrak p}\simeq R/\mathfrak p,\qquad
 \operatorname{Spec}A_{\bar{\mathfrak p}}\simeq
 \operatorname{Supp}_{R_{\mathfrak p}}M_{\mathfrak p}.
\] For the final identity, annihilators commute with this flat localization: apply the finite-generator kernel argument used above. Each comparison is given by the original quotient classes and denominators.

Applying the source's full-support chain theorem and dimension formula to these actual comparison objects proves \[
 d=\dim\operatorname{Supp}_{R_{\mathfrak p}}M_{\mathfrak p}
       +\dim R/\mathfrak p
   =\operatorname{depth}_{R_{\mathfrak p}}M_{\mathfrak p}
       +\dim R/\mathfrak p .
\] The second equality uses the already proved localization theorem at this support prime. No assertion that \(R\) itself is Cohen--Macaulay is made. At a prime outside the support the localized module is zero; its two infinite conventions cannot be substituted into this finite formula.

Every maximal chain inside \(\operatorname{Supp}M\) has length \(d\), through the same prime correspondence. For \(\mathfrak p\subset\mathfrak q\) in that support, fix maximal chains below \(\mathfrak p\) and above \(\mathfrak q\). Inserting any maximal chain between them gives a maximal chain in the support. Its total length is \(d\), so all interval chains have the same finite length. Similarly every maximal chain from \(\mathfrak p\) to \(\mathfrak m\) has length \(\dim R/\mathfrak p\). Subtracting the lengths of the actual concatenated chains gives the interval length \[
 \dim R/\mathfrak p-\dim R/\mathfrak q.
\] Thus this closed support is catenary with its exact dimension function. This is a proved consequence of the current module results, not a claim that the entire ambient spectrum is catenary. The following source catenary section remains unadjudicated.

\subsubsection{Every positive parameter power and its associated primes}\label{algebra-proof-019-every-positive-parameter-power-and-its-associated-primes}

Let \(M\ne0\) be Cohen--Macaulay over the same Noetherian local \(R\), and let \(g_1,\ldots,g_c\in\mathfrak m\) be \(M\)-regular, with \(0\le c\le d\). Each successive quotient is Cohen--Macaulay of dimension \(d-i\), by the original regular-element equivalence. Fix any positive integers \(a_1,\ldots,a_c\) and any permutation \(\sigma\). Keep the original elements and their powers: \[
 J_{\mathbf a}=(g_1^{a_1},\ldots,g_c^{a_c}),\qquad
 h_i=g_{\sigma(i)}^{a_{\sigma(i)}} .
\] These ideals have the same radical as \(J_{\mathbf 1}\), since a prime contains \(g_j^{a_j}\) exactly when it contains \(g_j\). Therefore the support of the actual quotient \(M/J_{\mathbf a}M\) equals \(\operatorname{Supp}M\cap V(g_1,\ldots,g_c)\), of dimension \(d-c\). The corrected nonzero-module parameter proposition applies to the ordered tuple \(h_1,\ldots,h_c\) and proves it regular. It also proves every intermediate quotient Cohen--Macaulay of its required dimension. All final quotients are nonzero by Nakayama. In particular \[
 \operatorname{Ass}_R(M/J_{\mathbf a}M)
 =\operatorname{Min}\bigl(\operatorname{Supp}M\cap V(g_1,\ldots,g_c)\bigr)
\] for every positive exponent vector, and every prime in this set has \(\dim R/\mathfrak p=d-c\). The empty tuple gives the original \(\operatorname{Ass}M=\operatorname{Min}\operatorname{Supp}M\).

At any prime of this final support, the original quotient and localization maps identify \((M/J_{\mathbf a}M)_{\mathfrak p}\) with \(M_{\mathfrak p}/J_{\mathbf a}M_{\mathfrak p}\). The powers remain regular there: localization preserves injectivity, and the last quotient is nonzero since the prime belongs to its support. Thus \[
 \dim\operatorname{Supp}(M/J_{\mathbf a}M)_{\mathfrak p}
 =\dim\operatorname{Supp}M_{\mathfrak p}-c .
\] These equalities retain the full exponent vector; no quotient modules for different vectors are identified with one another.

This receives group 463 concretely. For a single \(M\)-regular \(x\in\mathfrak m\), that earlier proof puts every prime minimal over \(\mathfrak p+(x)\), \(\mathfrak p\in\operatorname{Ass}M\), in \(\operatorname{Ass}(M/x^nM)\) for all sufficiently large \(n\). The equality of associated-prime sets just proved holds for every \(n\ge1\). Hence every such prime is already associated to \(M/xM\), and remains associated for every positive power. The large-power bound is removable in this Cohen--Macaulay regular-element setting. It is not removed from the earlier theorem for arbitrary finite modules or arbitrary \(x\).

\hypertarget{algebra_catenary_regular_notes_20260928}{}
\subsection{Catenary and regular local rings: exact source arguments and further consequences}\label{algebra-proof-020-catenary-and-regular-local-rings-exact-source-arguments-and-further-consequences}

Primary source: the Stacks Project, algebra.tex at authority commit a04446e57ec1fbc252a871afcec7752fb2807b14, SHA256 FA8BB92E58A4F78A2BD01B3B6A4A87DE0A0D279F5DD90641B574DD5FBFFFA4F3. Complete new review: 25456--25861. Lookahead 25862--25930 is read but unadjudicated; it ends at the start of the local epimorphism proof.

Dependencies read again: algebra.tex 13964--13982, the statement 14521--14536, and 14614--14626. The earlier exact Hilbert argument and its zero-variable cases in ALGEBRA\_GRADED\_PROJ\_HILBERT\_NOTES\_20260927.md, ``Hilbert bound notation'', are reread. The previous CM notes supply their fully proved parameter, quotient and depth results, with their actual zero-module conventions retained.

Original topology.tex at the same authority commit has SHA256 C6BAC8DCF8AD96DC47416BF34CB45BA4A10B894E40D67D3E1FA68D8EF0D9F872. Reading: 1534--1544 and 3839--3984. Its old proof reverses the specialization and dimension-function differences at 3880 and 3887--3888; the already composed MC-STK-ERR-0041 supplies the correct direction and both signs. That dependency is retained, not allocated again. The algebra dimension-function statement itself has the correct sign.

The two new corpus hits are unread routing records. No external theorem or novelty claim is taken from them. Complete arguments below are separate editorial evidence, with only minimal proposed operations in the review. Neither the cumulative chapter nor any translation is rewritten.

The current chapter already contains a FAC reference and history environment at lemma-regular-quotient-regular, referring to Chapter III, §5, no. 75, proof of Theorem 3, pages 269--270. The complete current block has been read. It records the existing connection between the kernel generators, codimension and successive colon equalities. Both environments are retained exactly. They are the only difference in this interval from the authority, and are not new correction operations or a fresh reading of FAC.

\subsubsection{Prime intervals, localization and the omitted universal argument}\label{algebra-proof-020-prime-intervals-localization-and-the-omitted-universal-argument}

For a ring \(R\), the map \(\mathfrak p\mapsto V(\mathfrak p)\) is a bijection from prime ideals to irreducible closed subsets, reversing inclusion. Its inverse is the ideal of functions vanishing on the closed subset, which is prime for a nonempty irreducible closed set. Reversing a strict prime chain preserves its number of inclusions. This proves precisely the equality of the ring and topological catenarity definitions, including the empty spectrum of the zero ring.

For a multiplicative subset \(S\), extension and contraction identify primes of \(S^{-1}R\) with primes of \(R\) disjoint from \(S\). If \(\mathfrak p\subset\mathfrak q\) are both disjoint from \(S\), every intermediate prime is disjoint too. Thus the entire closed order interval is identified, not merely its endpoints. All strict inclusions, saturation conditions, chain lengths and their bounds are preserved. Similarly primes of \(R/I\) correspond to primes of \(R\) containing \(I\), and if the lower endpoint contains \(I\), so does every intermediate prime. This proves both ordinary localization and quotient permanence.

For the universal localization statement let \(C\) be a finite type \(S^{-1}R\)-algebra and choose its actual finite algebra generators \(c_j\). Let \(B\) be the image of \(R[T_1,\ldots,T_n]\to C\), \(T_j\mapsto c_j\). It is a finite type \(R\)-algebra. The map \[
 S^{-1}B\longrightarrow C,\qquad b/s\longmapsto b\,s^{-1}
\] is onto by the generating assumption and is injective: \(B\subset C\) and every \(s\in S\) is a unit of \(C\). Hence it is an isomorphism. If \(R\) is Noetherian universally catenary, \(B\) is catenary, and the interval comparison gives catenarity of \(C\). The localization is Noetherian as required by the definition. Compositions of finite type maps are finite type, and the same argument after localizing proves the source essentially-finite-type result. A finite type algebra is a quotient of a finite-variable polynomial algebra, which justifies the polynomial test without assuming infinitely many variables.

For local checking of ordinary catenarity choose a maximal ideal containing the upper endpoint of any original interval and use its exact localization bijection above. For universal checking take a finite type \(R\)-algebra \(B\), a prime \(\mathfrak q\) of \(B\), its contraction \(\mathfrak p\), and a maximal ideal \(\mathfrak m\supset\mathfrak p\). The actual ring \(R_{\mathfrak p}\) is a localization of \(R_{\mathfrak m}\), while \(B_{\mathfrak q}\) is a localization of the finite type \(R_{\mathfrak p}\)-algebra \((R\setminus\mathfrak p)^{-1}B\). Universal catenarity at \(\mathfrak m\) therefore makes \(B_{\mathfrak q}\) catenary, and ordinary local checking gives catenarity of \(B\). The received correction at 25559 is the single word ``primes'' to ``prime''.

For the omitted universal assertion at 25589--25604, let \(B\) be any finite type \(R\)-algebra and let \(\mathfrak a\subset\mathfrak b\) be primes of \(B\). Choose a minimal prime \(\mathfrak p\) of \(R\) contained in \(\mathfrak a\cap R\). The entire interval \([\mathfrak a,\mathfrak b]\) contains \(\mathfrak pB\), so quotienting identifies it with its interval in \(B/\mathfrak pB\). This is a finite type algebra over \(R/\mathfrak p\), hence catenary by the assumed universal property. The interval bounds and lengths follow. This proves the omitted direction; the converse follows by quotient permanence. It does not assume a minimal prime of \(B\) contracts to a minimal prime of \(R\), which is false in general.

\subsubsection{Arbitrary closed covers and nil-ideal invariance}\label{algebra-proof-020-arbitrary-closed-covers-and-nil-ideal-invariance}

Let \(\{I_\lambda\}\) be any family of ideals of a ring \(R\) such that \(\operatorname{Spec}R=\bigcup_\lambda V(I_\lambda)\). Then \[
 R\text{ is catenary}\quad\Longleftrightarrow\quad
 R/I_\lambda\text{ is catenary for every }\lambda.
\] Necessity is quotient permanence. For sufficiency, in any interval \(\mathfrak p\subset\mathfrak q\), choose \(\lambda\) with \(I_\lambda\subset
\mathfrak p\). Every intermediate prime contains this same ideal. Extension of quotient classes and contraction thus identify the whole interval with one in \(R/I_\lambda\), preserving all lengths and saturation. No finite cover and no Noetherian hypothesis are used.

If \(R\) is Noetherian, the same equivalence holds for universal catenarity. For a finite type algebra \(B\) and an interval with lower endpoint \(\mathfrak a\), apply the cover to \(\mathfrak a\cap R\). Then \(I_\lambda B\subset\mathfrak a\), and the interval is one in the finite type \(R/I_\lambda\)-algebra \(B/I_\lambda B\). Its catenarity proves the claim. Necessity again follows by quotient permanence.

Every prime of any ring contains a minimal prime. Indeed the intersection of a chain of primes contained in a given prime is prime: if \(ab\) lies in all the chain members but neither \(a\) nor \(b\) lies in their intersection, choose members missing each one; the smaller of those two comparable primes misses both, a contradiction. The reversed-inclusion form of Zorn's lemma supplies a minimal member. Thus the source ordinary minimal-prime criterion holds for arbitrary rings, without Noetherianity. The universal criterion retains Noetherianity because it is part of the source definition.

If \(I\) is a nil ideal, every prime contains \(I\), so the one-member cover proves catenarity invariance under \(R\to R/I\). For Noetherian \(R\) it also proves universal catenarity invariance. Alternatively, in a finite type algebra \(B\), every element of \(IB\) is a finite sum of multiples of nilpotent elements. If their exponents are \(e_1,\ldots,e_t\), every term of a power of total degree \(1+\sum_j(e_j-1)\) contains a vanishing factor. Hence \(IB\) is nil and the spectra of \(B\) and \(B/IB\) have the same prime intervals. No ring or nilpotent contribution is identified with zero in the actual algebra; this is a proved comparison of their spectra.

\subsubsection{The local dimension function without Noetherianity}\label{algebra-proof-020-the-local-dimension-function-without-noetherianity}

Let \(A\ne0\) be any local ring with maximal ideal \(\mathfrak m\). Then \(A\) is catenary if and only if \[
 \delta(\mathfrak p)=\dim(A/\mathfrak p)
\] is a finite integer for every prime and is a dimension function in the source sense: it strictly decreases under proper specialization and decreases by exactly one under immediate specialization.

Suppose \(A\) is catenary. For a fixed \(\mathfrak p\), its interval to \(\mathfrak m\) has bounded finite chain length. Any finite chain in \(\operatorname{Spec}(A/\mathfrak p)\) extends at both ends to that interval. Thus its dimension is finite and is the common length of maximal \(\mathfrak p\)-to-\(\mathfrak m\) chains. Such chains exist: successive strict refinements must terminate because of the integer bound. For \(\mathfrak p\subset\mathfrak q\), concatenate a maximal interval chain from \(\mathfrak p\) to \(\mathfrak q\) with one from \(\mathfrak q\) to \(\mathfrak m\). The concatenation is saturated and has the common total length \(\delta(\mathfrak p)\). Hence \[
 \ell[\mathfrak p,\mathfrak q]
 =\delta(\mathfrak p)-\delta(\mathfrak q).
\] This gives both dimension-function conditions.

Conversely, along any strict interval chain each integer decrease is at least one, so the length is bounded by the displayed difference. A maximal chain is finite and saturated; each step is immediate and decreases \(\delta\) by one. Its length is exactly that difference. This proves catenarity and the precise interval formula without Noetherianity. The values may be unbounded over different primes; no finite global dimension assumption has been inserted.

This direct argument proves the original Noetherian local result with its correct signs. In the source's topological proof, an open neighborhood of the closed point in a local spectrum is the whole spectrum: any basic open containing the maximal ideal is defined by a unit. Thus the cited local dimension-function construction really supplies a global function there. The corrected MC-STK-ERR-0041 sign is \(\delta(\text{generic})-\delta(\text{special})\), consistent with the algebra formula above. The original algebra statement needs no operation.

\subsubsection{Universal catenarity of module supports and covering families}\label{algebra-proof-020-universal-catenarity-of-module-supports-and-covering-families}

The source full-support CM theorem is proved by its original two comparisons. For a finite-variable polynomial extension, the previous CM-module theorem preserves Cohen--Macaulayness. Its support is the entire polynomial spectrum: the polynomial map is flat and the finite-generator annihilator kernel comparison in the previous notes gives the extended annihilator; an original full-support module has nil annihilator, so every polynomial prime contains that extension. Localizing at any prime gives a nonzero CM module with full local support. The original maximal-chain theorem then fixes the lengths of prime intervals. All polynomial algebras are catenary; finite type quotients give universal catenarity. Taking \(M=R\) recovers the ring case.

More generally, let \(R\) be Noetherian and \(M\) a finite CM module. Retain the original \(R\) and \(M\), and set \(A=R/\operatorname{ann}_R M\). At corresponding primes the exact quotient and localization maps of group 590 identify the same nonzero module and all regular-sequence multiplication maps. It is CM over \(A\), with full support. Hence \(A\) is universally catenary by the current source theorem. If \(M=0\), then \(A=0\), and the claim holds directly: its spectrum and those of its unital algebras are empty. For nonzero local \(M\), this strengthens group 590's support-catenarity conclusion to universal catenarity of its actual support ring; its original dimension and interval formulas stay unchanged.

If a family of finite CM \(R\)-modules has supports covering \(\operatorname{Spec}R\), each quotient by its annihilator is universally catenary. The proved arbitrary closed-cover criterion therefore makes \(R\) universally catenary. No direct sum is assumed to be CM: the argument uses the actual individual quotient rings and their interval maps. The source's single full-support module is the one-member case.

\subsubsection{Regular graded rings, domains and parameter quotients}\label{algebra-proof-020-regular-graded-rings-domains-and-parameter-quotients}

At 25690--25712, if \(d=0\), the regular maximal ideal is generated by the empty family and is zero. Thus \(R=\kappa\), its associated graded ring is \(\kappa\) in degree zero, and the claimed comparison is the identity. No degree-\(-1\) polynomial or negative binomial lower index is used. This boundary is editorial evidence, as in the earlier group 384.

For \(d>0\), retain the original \(x_1,\ldots,x_d\) and the graded map \[
 \theta:\kappa[X_1,\ldots,X_d]\longrightarrow
 \bigoplus_{n\ge0}\mathfrak m^n/\mathfrak m^{n+1},\quad
 X_j\longmapsto x_j+\mathfrak m^2.
\] It is surjective: products of the actual generators span each ideal power, and coefficients reduce to their residue classes in that graded piece. The dimension theorem and the earlier Hilbert calculation give its target the original eventual degree \(d-1\). A nonzero homogeneous kernel element of degree \(e\) would force the upper bound \[
 \dim_\kappa(\operatorname{gr}_{\mathfrak m}R)_n
 \le {n-1+d\choose d-1}-{n-e-1+d\choose d-1}
\] for all sufficiently large \(n\). For \(d\ge2\) the highest coefficients cancel, giving degree at most \(d-2\); for \(d=1\) the expression is zero. Both contradict the original degree. Thus the actual \(\theta\) is injective.

For the domain proof take \(f,g\ne0\). Krull intersection makes their actual orders finite: each lies in exactly the initial segment of the powers through some maximal \(a,b\). Their nonzero initial forms have nonzero product in the polynomial graded domain. Hence \[
 fg\in\mathfrak m^{a+b}\setminus\mathfrak m^{a+b+1},
\] contradicting \(fg=0\). If either factor is zero the domain implication is already satisfied; the received correction explicitly excludes that case before choosing orders.

For 25731--25747 the \(d=0\) field case has the empty regular sequence and needs no \(x_1\). For \(d>0\), \(x_1\notin\mathfrak m^2\) and in particular \(x_1\ne0\); the domain result makes it regular. The quotient has dimension at least \(d-1\) by the one-equation lemma and its maximal ideal has at most \(d-1\) generators. The earlier dimension bound gives equality and regularity. Induction with the actual images of \(x_2,\ldots,x_d\) proves every asserted quotient dimension and the full regular sequence. This supplies omitted boundary steps as editorial evidence.

For a regular quotient \(R/I\), \(I\subset\mathfrak m\) and both rings share \(\kappa\). There is the actual exact sequence of \(\kappa\)-vector spaces \[
 0\longrightarrow (I+\mathfrak m^2)/\mathfrak m^2
 \longrightarrow \mathfrak m/\mathfrak m^2
 \longrightarrow (\mathfrak m/I)/(\mathfrak m/I)^2\longrightarrow0 .
\] Thus the source's third dimension must also have the \(\kappa\) subscript. Choose the first \(c\) basis lifts in the original \(I\) and complete them to the original minimal generator set of \(\mathfrak m\). The map \(R/(x_1,\ldots,x_c)\to R/I\) is onto between regular local rings of the same dimension. If its kernel contained \(0\ne y\), the source domain result makes \(y\) regular, and the proper quotient by \(y\) has dimension one smaller. Factoring the map through that quotient contradicts the equal dimension of its target. Therefore its kernel is zero and the ideal equality holds, including \(c=0\) and \(c=d\). The notation declaration at 25760 receives the missing ``by''.

\subsubsection{Free lifting, maximal CM modules and regular lifting}\label{algebra-proof-020-free-lifting-maximal-cm-modules-and-regular-lifting}

In 25779--25797 choose the original finite basis of \(M/xM\) and lifts \(m_i\). The quotient of \(M\) by their span is finite and equal to its multiple by \(x\in\mathfrak m\), so Nakayama makes the lift map \(F=R^r\to M\) onto. Let \(K\) be its actual kernel. For \((a_i)\in K\), reduction modulo \(x\) gives \(a_i=xb_i\). The relation becomes \(x\sum b_im_i=0\). Regularity of \(x\) on \(M\) gives \((b_i)\in K\). Hence \(K=xK\). Over the original Noetherian ring \(K\) is finite, and Nakayama gives \(K=0\). If \(r=0\), the same Nakayama argument gives \(M=0\), a free module of rank zero. The article before \(R/xR\)-basis is corrected.

For the maximal-CM theorem over regular local \(R\), first handle \(d=0\): \(R\) is a field and every finite module is free. Otherwise the original regular local ring is a domain. A nonzero maximal-CM module has dimension \(d\), and its minimal support primes all have quotient dimension \(d\). In the local domain only the zero prime can have that dimension: a strictly larger prime shortens chains by at least one. Hence the module has full support. The chosen parameter \(x\) consequently cuts that support by one dimension, and the previous nonzero-module CM proposition makes \(x\) regular on \(M\). Its quotient is a nonzero maximal-CM module over the regular local ring \(R/xR\). Induction makes the quotient free, and the proved original lifting lemma makes \(M\) free. This verifies the source's implicit full-support and field-base cases without rewriting its proof.

For 25815--25831, let \(x\in\mathfrak m\) be regular on \(R\), and suppose \(R/xR\) regular of dimension \(d-1\). The one-equation equality gives \(\dim R=d\). Lift its \(d-1\) minimal maximal-ideal generators and adjoin the actual \(x\). Their images generate \(\mathfrak m/xR\), so they and \(x\) generate \(\mathfrak m\). The dimension lower bound prevents fewer than \(d\) generators. This is exactly regularity of \(R\); induction through the original successive quotients proves the sequence version.

\subsubsection{Directed colimits with noninjective transition maps}\label{algebra-proof-020-directed-colimits-with-noninjective-transition-maps}

Let the source maps be \(\varphi_{ii'}:R_i\to R_{i'}\), all local. The direct limit \(R\) is nonzero: equality \(1=0\) would hold at a finite stage, contradicting a nonzero local ring. The subset of elements represented by maximal-ideal elements is an ideal \(\mathfrak m\). An element represented outside \(\mathfrak m_i\) is a unit already at that stage. An element represented inside it cannot become a unit: an inverse relation would occur at a later stage, where the original element still lies in the maximal ideal. Thus \(R\) is local and its maximal ideal is the colimit of the original maximal ideals. The plural wording is corrected.

The induction parameter is the finite vector-space dimension \[
 d=\dim_{R/\mathfrak m}(\mathfrak m/\mathfrak m^2),
\] not an ungrouped quotient of a dimension symbol. Noetherianity of the given limit makes \(\mathfrak m\) finite. If \(d=0\), Nakayama gives \(\mathfrak m=0\), so \(R\) is a field.

If \(d>0\), choose the actual \(x\in\mathfrak m\setminus\mathfrak m^2\) and a representing \(x_i\in\mathfrak m_i\). For every \(i'\ge i\), its image \(x_{i'}=\varphi_{ii'}(x_i)\) is outside \(\mathfrak m_{i'}^2\): membership there would imply \(x\in\mathfrak m^2\). It is therefore part of a minimal generator set by the residue tangent-space basis and Nakayama. The source parameter theorem makes \(R_{i'}/x_{i'}R_{i'}\) regular local. The induced transition maps remain local. Their limit is the actual quotient \(R/xR\), with map \[
 a_{i'}+x_{i'}R_{i'}\longmapsto a+xR.
\] Surjectivity follows from representatives. If a representative maps to zero, write \(a=xb\) in the limit, choose a common later stage for \(a,x,b\), and then a later stage where this equality holds. This proves injectivity.

Its tangent space is \[
 (\mathfrak m/xR)/(\mathfrak m/xR)^2
 \simeq\mathfrak m/(\mathfrak m^2+xR),
\] whose dimension is exactly \(d-1\), since the class of \(x\) is nonzero in \(\mathfrak m/\mathfrak m^2\). Thus induction applies to this Noetherian quotient limit.

A directed colimit of the nonzero domains \(R_i\) here is a domain even when the transition maps are not injective. If a product is zero in the limit, represent its factors at a common stage and then pass to a stage where their product is zero. One factor is zero at that domain stage, hence zero in the limit. Together with \(1\ne0\) this proves the claim. The direct source citation is lemma-regular-domain, as the reports state. Now \(x\ne0\), so it is regular on \(R\), and the previous lifting lemma concludes the source colimit theorem.

\subsubsection{Two weakenings of Noetherianity in free lifting}\label{algebra-proof-020-two-weakenings-of-noetherianity-in-free-lifting}

The free-lifting result remains true for a local ring \(R\), \(x\in\mathfrak m\), and \(M\) finite, with \(x\) injective on \(M\) and \(M/xM\) free, if either:

\begin{enumerate}
\def\labelenumi{(\alph{enumi})}
\item
  \(M\) is finitely presented, or
\item
  the original \(x\)-adic filtration of \(R\) is separated: \(\bigcap_{n\ge0}x^nR=0\).
\end{enumerate}

The same chosen quotient basis is finite: tensor it with the nonzero residue field and use finite generation of \(M/xM\) to bound its cardinality. The preceding proof gives the surjection \(F=R^r\to M\) and \(K=xK\) verbatim.

For (a), the kernel of a finite free surjection onto a finitely presented module is finite. Explicitly, take a finite presentation \(G\to H\to M\) with \(H\) finite free and finitely generated kernel \(L\). Lift the maps from \(F\) to \(M\) through \(H\), and from \(H\) to \(M\) through \(F\), obtaining \(u:F\to H\), \(v:H\to F\) compatible with the two surjections. For \(k\in K\), \[
 k=(1-vu)k+v(uk).
\] Here \(uk\in L\), and \((1-vu)F\subset K\), \(vL\subset K\). Consequently \(K=(1-vu)F+vL\), a sum of two finite submodules. Nakayama now gives \(K=0\).

For (b), \(K=xK\) implies \(K\subset x^nF\) for every \(n\). In the original finite coordinates \(\bigcap_n x^nF=(\bigcap_n x^nR)^r=0\), so again \(K=0\). No finite generation of \(K\) is assumed in this second proof. Neither alternative requires \(x\) to be a nonzerodivisor on \(R\); the needed injectivity is on the original \(M\). Both include \(M=0\).

\subsubsection{The valuation determined by the original maximal-ideal powers}\label{algebra-proof-020-the-valuation-determined-by-the-original-maximal-ideal-powers}

Let \(R\) be regular local of dimension \(d>0\), with the original parameters \(x_1,\ldots,x_d\), residue field \(\kappa\), and fraction field \(K\). For \(0\ne f\in R\) put \[
 v(f)=\max\{n:f\in\mathfrak m^n\}.
\] The graded-domain argument proves the exact equality \(v(fg)=v(f)+v(g)\). Also \(v(f+g)\ge\min(v(f),v(g))\) when the sum is nonzero, with equality if the two orders differ; the lowest nonzero graded term then cannot cancel. Extend by \[
 v(a/b)=v(a)-v(b)\quad(a,b\ne0),\qquad v(0)=+\infty.
\] If \(a/b=c/e\), the exact equality \(ae=bc\) and additivity on \(R\) prove that the definition is independent of representatives. Multiplication and the sum inequality follow after a common denominator, retaining every factor in that comparison. Since \(v(x_d)=1\), its image is all \(\mathbb Z\).

The ring \(V=\{z\in K:v(z)\ge0\}\) is a valuation ring: for any nonzero \(z\), either it or its inverse has nonnegative value. Its units have value zero, and its maximal ideal consists of positive-valued elements, exactly \(x_dV\). Every nonzero ideal of \(V\) has a least value \(n\ge0\); if \(z\) in that ideal has value \(n\), every other element divided by \(z\) lies in \(V\), so the ideal is \(zV=x_d^nV\). Thus \(V\) is a discrete valuation ring with the actual uniformizer \(x_d\). Its maximal ideal contracts to \(\mathfrak m\), and \(R\subset V\).

The residue field is explicitly \[
 V/x_dV\simeq
 \kappa(X_1/X_d,\ldots,X_{d-1}/X_d).
\] For \(a/b\) of value zero, send its residue to \(\theta^{-1}(\operatorname{in}a)/\theta^{-1}(\operatorname{in}b)\); the two homogeneous polynomials have the same degree. Positive-valued elements map to zero. Equality of fractions is respected by the exact initial-form product identities. For addition, use \(ae+bc\) as numerator: if its lowest terms cancel, both the proposed residue sum and the new residue are zero; otherwise the initial numerator is their actual sum. This proves a ring map on \(V\), with kernel exactly the positive-valued ideal.

Every ratio of homogeneous polynomials of the same degree is a rational function in \(X_j/X_d\), by dividing each full polynomial by the same power of \(X_d\). Conversely any rational function in those ratios is a ratio of same-degree homogeneous polynomials after multiplying numerator and denominator by the required powers of \(X_d\). The graded isomorphism \(\theta\) supplies actual representatives of every homogeneous polynomial in the corresponding \(\mathfrak m^n/\mathfrak m^{n+1}\). Their fraction has value zero. This proves surjectivity and the residue-field isomorphism, with all original initial forms and parameter comparisons retained.

For \(d=1\), \(V=R\): every nonzero \(f\) is \(x_1^{v(f)}u\) with \(u\in R^\times\), because \(\mathfrak m=(x_1)\); any fraction of nonnegative value lies in \(R\). For \(d>1\), \(x_1/x_d\in V\setminus R\). An equality \(x_1=x_dr\) in \(R\) would put \(x_1\) in \(\mathfrak m^2\) if \(r\in\mathfrak m\), or would make the two independent tangent-space classes proportional if \(r\) were a unit. Both are impossible. For \(d=0\), the original ring is a field and its order valuation is trivial; no discrete rank-one valuation is asserted. This is a proved standard consequence of the source graded comparison, not a novelty claim or a replacement of the original local ring.

\hypertarget{algebra_epimorphism_notes_20260928}{}
\subsection{Ring epimorphisms: source review and separate editorial arguments}\label{algebra-proof-021-ring-epimorphisms-source-review-and-separate-editorial-arguments}

Authority: Stacks Project authors, git a04446e57ec1fbc252a871afcec7752fb2807b14, algebra.tex:25862--26191. SHA-256 FA8BB92E58A4F78A2BD01B3B6A4A87DE0A0D279F5DD90641B574DD5FBFFFA4F3. Read the entire section and all twelve received reports. Source dependencies 20236--20269 were reread. The current complete epimorphism section equals the authority; no canonical correction or later editorial environment occurs within it.

This is separate, source-linked editorial material. The original author's statements, proof omissions, order of exposition and mathematical objects remain identifiable. The nine proposed text operations correct wording or bindings; the full arguments, typed matrix correction and further consequences below do not replace the translated section. Both current and diplomatic translations remain translations. Broader synthesis follows the core work. No novelty claim is made.

\subsubsection{Tensor characterization and permanence}\label{algebra-proof-021-tensor-characterization-and-permanence}

At 25869--25882 write the original map as \(\varphi:R\to S\), and retain \[
j_1(s)=s\otimes1,\quad j_2(s)=1\otimes s,\quad
\mu(s\otimes t)=st .
\] These are ring maps and \(\mu j_1=\mu j_2=1_S\). If \(\varphi\) is an epimorphism, the two maps \(j_i\) coincide because their composites with \(\varphi\) coincide. Conversely, if \(j_1=j_2\), any ring maps \(a,b:S\to A\) with \(a\varphi=b\varphi\) define \(F:S\otimes_RS\to A\) by \(F(s\otimes t)=a(s)b(t)\); the balancing relation is exactly equality of those composites. Hence \(a=Fj_1=Fj_2=b\). If \(j_1=j_2\), then \(j_1\mu(s\otimes t)=(s\otimes1)(t\otimes1)=(s\otimes1)(1\otimes t)=s\otimes t\). Thus \(j_1\) and \(\mu\) are inverse. If either \(j_i\) is an isomorphism, its left inverse \(\mu\) is its inverse; \(\mu j_{3-i}=1\) forces equality of the two \(j_i\). If \(\mu\) is an isomorphism, both \(j_i\) equal its inverse. This proves all four source conditions, including zero rings.

For composition (25884--25891), cancel the first epimorphism and then the second in the two composite equalities. For 25903--25911, two maps out of \(C\) agreeing on \(B\) agree on \(A\), hence coincide if \(A\to C\) is epic. In particular the map from the actual image \(\varphi(R)\) to \(S\) is epic whenever \(\varphi\) is; this does not identify \(R\) itself with its image.

For base change (25893--25901), set \(S'=R'\otimes_RS\). The actual isomorphism \[
S'\otimes_{R'}S'\longrightarrow R'\otimes_R(S\otimes_RS),
\quad (r'\otimes s)\otimes(t'\otimes u)\longmapsto r't'\otimes(s\otimes u)
\] has inverse \(r'\otimes(s\otimes u)\mapsto(r'\otimes s)\otimes(1\otimes u)\). Both formulas respect the balancing relations, and their composites are the identity on elementary tensors. Under this map multiplication becomes \(1_{R'}\otimes\mu\), so remains an isomorphism. A localization \(R\to U^{-1}R\) is epic directly: two maps agreeing on \(R\) must send \(r/u\) to the same \(a(r)a(u)^{-1}\). This also covers the zero localization, for which any existing maps out of it have zero target.

Report OCC-00517 removes the misplaced comma at 25914; this has no mathematical effect.

\subsubsection{Local comparison and finite or field targets}\label{algebra-proof-021-local-comparison-and-finite-or-field-targets}

At 25920--25941 the forward local implication is the preceding base change. Conversely \(a,b:S\to A\) with equal composites on \(R\) give \(A\) its common \(R\)-algebra structure. They induce maps \(S_{\mathfrak p}\to A_{\mathfrak p}\). Equality of these maps for every prime forces \(a(s)-b(s)=0\): if an \(R\)-module element \(x\ne0\), its annihilator is a proper ideal and is contained in a maximal ideal \(\mathfrak m\); the relation \(x/1=0\) at \(\mathfrak m\) would give \(u x=0\) for some \(u\notin\mathfrak m\), contradicting that containment. Thus the localization comparison is injective. For \(R=0\), all unital \(R\)-algebras are zero, and the assertion holds as well. OCC-00518 makes the equality of composites explicit instead of saying that two maps equalize a single map.

For the finite epimorphism result (25943--25974), use the actual \(R\)-module \(C=\operatorname{coker}(\varphi)\), which is what \(S/R\) denotes after passing to the image. Tensoring \(R\to S\to C\to0\) with \(S\) gives \(C\otimes_RS=0\) because its first map is \(j_2\), an isomorphism. The surjection \(S\to C\) then gives \(C\otimes_RC=0\). If the finite module \(C\) were nonzero it would have a maximal proper submodule: a chain of proper submodules cannot cover a finite generating set unless one member is already the whole module. Zorn's lemma applies. Its simple quotient is \(R/\mathfrak m\) for a maximal ideal \(\mathfrak m\) (take a nonzero element of the simple quotient and its annihilator). Tensoring the two quotient surjections gives a surjection \(C\otimes_RC\to(R/\mathfrak m)\otimes_R(R/\mathfrak m)=R/\mathfrak m\ne0\), a contradiction. Conversely a surjective map is epic and its target is cyclic as an \(R\)-module.

At 25976--25985 faithful flatness reflects the isomorphism obtained by tensoring \(\varphi\) with \(S\). More explicitly, flatness identifies the base changes of its kernel and cokernel with the corresponding kernel and cokernel after tensoring; both vanish, and faithfulness detects their vanishing.

At 25987--25998 a nonzero \(k\)-algebra \(S\) is faithfully flat over the field \(k\): it is a vector space with a nonempty basis, and tensoring any vector space with it is a direct sum of copies of that vector space. The preceding result gives the actual structural isomorphism \(k\to S\). For the alternate argument, if \(s\notin k1\), choose a \(k\)-linear functional \(\lambda:S\to k\) with \(\lambda(1)=0,\lambda(s)=1\), by extending \(\{1,s\}\) to a basis. The map \(1\otimes\lambda\) kills \(j_1(S)\) but takes \(1\otimes s\) to \(1\ne0\), so \(j_1\) is not surjective. This proves the stated dimension bound, with \(S=0\) treated separately. OCC-12192 says ``three'' occurrences of \(R\), but the clause actually contains four: three at 25996 and one at 25997. Both reports receive the correction of all four bindings to \(S\). Multiplication itself is always surjective; the relevant map here is \(j_1\) or \(j_2\), not multiplication.

At 26000--26020 retain each original fibre \[
T=S\otimes_R\kappa(\mathfrak p)
=(R\setminus\mathfrak p)^{-1}S/
 \mathfrak p(R\setminus\mathfrak p)^{-1}S .
\] The base changed field map is epic, so \(T=0\) or its structural map from \(\kappa(\mathfrak p)\) is an isomorphism. Primes \(\mathfrak q\) of \(S\) contracting to \(\mathfrak p\) correspond exactly to primes of \(T\), by first localizing and then quotienting. The zero case has none; the field case has one. For that prime the localization of \(T\) at zero is \(T\), and the usual localization and quotient maps identify it with \(S_{\mathfrak q}/\mathfrak qS_{\mathfrak q}=\kappa(\mathfrak q)\). This identifies the actual induced field map, not just the abstract isomorphism types.

\subsubsection{Finite relations and the original matrix construction}\label{algebra-proof-021-finite-relations-and-the-original-matrix-construction}

At 26022--26058 let \(K=\ker(\bigoplus_JR\to N)\). Right exactness gives the displayed \(K\otimes_RM\to\bigoplus_JM\to N\otimes_RM\to0\). A preimage of the tuple \((m_j)\) is a finite sum \(\sum_{\ell=1}^t k_\ell\otimes u_\ell\). Express each \(u_\ell\) as a finite linear combination of the original \(x_i\). Collecting its coefficients gives \(\sum_i\widetilde k_i\otimes x_i\) with only finitely many nonzero \(\widetilde k_i\in K\). Each such vector has finite support in \(J\). Writing \(\widetilde k_i=(a_{ij})_j\) gives a matrix with finite total support, the required column equalities \(m_j=\sum_i a_{ij}x_i\), and row relations \(\sum_j a_{ij}y_j=0\). Conversely these equations give \(\sum_jm_j\otimes y_j=\sum_i x_i\otimes\sum_j a_{ij}y_j=0\). The three reports at 26057 delete the superfluous conjunction, without replacing this source argument.

For 26060--26097 retain every coefficient through \(\varphi\). Suppose \(g=\sum_{ij}\varphi(x_{ij})y_i z_j\) and choose \(u_j,v_i\in R\) such that \(\sum_i\varphi(x_{ij})y_i=\varphi(u_j)\) and \(\sum_j\varphi(x_{ij})z_j=\varphi(v_i)\). Then \[
g\otimes1
=\sum_j\varphi(u_j)z_j\otimes1
=\sum_j z_j\otimes\varphi(u_j)
=\sum_{ij}z_j\otimes\varphi(x_{ij})y_i
=\sum_i\varphi(v_i)\otimes y_i
=1\otimes g .
\] Balancing and commutativity justify each displayed equality.

Conversely apply the proved relation lemma to the source generators \(s_0=1,s_1=g\); the indices 0 and 1 may designate equal elements. Choose a finite set \(K_0\) of indices containing 0, 1 and every row or column index of the finite support of \(a_{ij}\). Put \(H=K_0\setminus\{0\}\), retaining all original generators and coefficients. The original column and row equations give \[
h:=\sum_{i,j\in H}\varphi(a_{ij})s_i s_j
=-g+\varphi(a_{00}),\qquad
\sum_{j\in H}\varphi(a_{ij})s_j=-\varphi(a_{i0})\ (i\in H).
\] Indeed the sum over all \(i,j\in K_0\) is zero by the row relations; the row with \(i=0\) is zero as well, whereas the remaining \(j=0\) contribution is \(g-\varphi(a_{00})\). For the column sums over \(H\), the \(j=1\) column is \(-1-\varphi(a_{01})\); every other \(j\in H\) gives \(-\varphi(a_{0j})\). All these sums lie in \(\varphi(R)\), precisely as the source statement requires. This proves the received correction at 26092; no injection of \(R\) is assumed.

For an entirely explicit representation of \(g\), use \[
X=\operatorname{diag}(a_{00},-(a_{ij})_{i,j\in H}),\quad
Y=(1,(s_i)_{i\in H}),\quad Z=(1,(s_j)_{j\in H})^{\mathsf T}.
\] Then \(Y\varphi(X)Z=\varphi(a_{00})-h=g\); every entry of \(Y\varphi(X)\) and \(\varphi(X)Z\) belongs to \(\varphi(R)\) by the displayed row and column equations. This completes the omitted step without changing any source hypothesis.

The source subalgebra hint also holds with exact matrices. Scalars use \(n=1\), \(X=(r)\), \(Y=Z=(1)\). The empty \(n=0\) representation gives zero. Sums use block diagonal \(X\)'s and concatenated \(Y,Z\); additive inverses negate \(X\). Products of representations use row index \((i,k)\), column index \((j,l)\), coefficients \(x_{ij}x'_{kl}\), row entries \(y_i y'_k\), and column entries \(z_j z'_l\). Their sum is the product of the two original sums; each row or column contraction is the product of the corresponding two original contractions in \(\varphi(R)\). This proves closure. OCC-01405 changes the singular subject's verb to ``forms''.

\subsubsection{Typed triples and the cardinality argument}\label{algebra-proof-021-typed-triples-and-the-cardinality-argument}

The two reports at 26112--26113 correctly change the vector names \(y,z\) to the declared matrices \(Y,Z\). There is a further source type issue in 26116--26122: for a general, possibly noninjective \(R\to S\), the products are matrices over \(\varphi(R)\) inside \(S\), not literally over \(R\). The original printed formulas are retained and linked to this conceptual erratum.

Here is their precise interpretation without replacing \(R\) by a quotient. An associated triple consists of \[
P\in\operatorname{Mat}(n\times n,R),\
U\in\operatorname{Mat}(1\times n,R),\
V\in\operatorname{Mat}(n\times1,R)
\] and witnesses \(Y,Z\) over \(S\) such that \[
g=Y\varphi(P)Z,\qquad
Y\varphi(P)=\varphi(U),\qquad
\varphi(P)Z=\varphi(V).
\] The preceding construction supplies \(P\); choose entrywise lifts \(U,V\) of its contractions. These lifts need not be unique, and no uniqueness is used. For two elements with the same triple and witnesses \(Y,Z\) and \(Y',Z'\), the source's comparison becomes \[
g=Y\varphi(P)Z=\varphi(U)Z
=Y'\varphi(P)Z=Y'\varphi(V)
=Y'\varphi(P)Z'=g'.
\] Thus every triple has at most one associated element. When \(R\) is infinite, the set of all finite triples has cardinality at most \(|R|\): for fixed \(n\) it has \(|R|^{n^2+2n}\), and a countable union has cardinality \(|R|\), using the same choice-based infinite cardinal arithmetic cited by the source to Kunen. Well-order the triples and select the least associated triple for each \(g\); the displayed uniqueness gives an injection. For an epimorphism every \(g\in S\) has a representation, proving the infinite case. This also bounds the epicenter for any map from an infinite \(R\); no finite-epicenter bound is inferred from this counting step.

For finite \(R\ne0\), let \(J\) be its Jacobson radical. There are finitely many maximal ideals \(\mathfrak m_1,\ldots,\mathfrak m_t\), pairwise comaximal, and the Chinese remainder map gives \(R/J\cong\prod_iR/\mathfrak m_i\). Finiteness makes \(J^n\) stabilize, say \(J^n=J J^n\). This finite module is zero by the elementary determinant form of Nakayama: for finite generators \(v\), write \(v=A v\) with entries of \(A\) in \(J\); multiplying by the adjugate gives \(\det(1-A)v=0\), and \(\det(1-A)\in1+J\) is a unit, since it belongs to no maximal ideal. Hence \(J^n=0\). The nilpotent result proved below makes every epimorphism \(R\to S\) surjective, whence \(|S|\leq|R|\). If \(R=0\), the unital target \(S\) is also zero. This supplies the omitted finite case without inserting a new proof into the translation.

\subsubsection{Modules and the restriction functor}\label{algebra-proof-021-modules-and-the-restriction-functor}

At 26158--26188 an \(S\)-linear map is automatically \(R\)-linear, so full faithfulness is exactly equality of the Hom sets. Suppose \(\varphi\) is epic and let \(f:N_1\to N_2\) be \(R\)-linear. For any \(x\in N_1\), the map \[
F_x:S\otimes_RS\to N_2,\quad s\otimes t\mapsto s f(tx)
\] is well-defined: for \(r\in R\), \(s\varphi(r)f(tx)=s f(\varphi(r)tx)\) by \(R\)-linearity. The equality \(s\otimes t=st\otimes1=1\otimes st\) then gives \(s f(tx)=st f(x)=f(stx)\), in particular \(f(sx)=s f(x)\). Conversely take the left and right \(S\)-actions on the original \(T=S\otimes_RS\). Their restrictions to \(R\) coincide, so the identity on \(T\) is \(R\)-linear between these modules. If every such map is \(S\)-linear, applying its linearity to \(1\otimes1\) gives \(s\otimes1=1\otimes s\), and hence the epimorphism criterion. All formulas include zero modules.

\subsubsection{Maximal local tests and pure descent}\label{algebra-proof-021-maximal-local-tests-and-pure-descent}

These are separate corollaries of the original tensor argument, not replacements for the source's primewise criterion or faithfully flat theorem.

For every ring map \(R\to S\), set \(D=\operatorname{coker}(j_1:S\to S\otimes_RS)\). The map \(j_1\) is split injective with retraction \(\mu\). Therefore the original map is epic exactly when \(D=0\). Right exactness and the explicit tensor comparison above identify \(D\otimes_RB\) with the analogous cokernel for every base change \(R\to B\). Consequently it is enough to test epimorphisms at all maximal ideals of \(R\), by the annihilator localization argument already proved; the converse follows from base change. No Noetherianity or finite generation is used.

Suppose now \(R\to B\) is pure, meaning the original unit \(M\to M\otimes_RB\), \(m\mapsto m\otimes1\), is injective for every \(R\)-module \(M\). If \(B\to B\otimes_RS\) is epic, \(D\otimes_RB=0\); purity on this exact \(D\) yields \(D=0\). Thus epimorphisms descend along pure base changes; ascent holds for every base change.

A pure epimorphism \(\varphi:R\to S\) is itself an isomorphism. Purity on \(R\) makes \(\varphi\) injective. Put \(C=\operatorname{coker}\varphi\). As before, epicity gives \(C\otimes_RS=0\); purity on \(C\) gives \(C=0\). Hence \(\varphi\) is surjective as well. This strengthens the source faithfully flat result because 20254--20269 proves every faithfully flat map pure. It receives the original unit-map reflection argument of ALGEBRA-RECON-501, in ALGEBRA\_OPEN\_LOCI\_PURE\_DESCENT\_NOTES\_20260927.md, section ``Descent along the original universally injective ring map'', SHA-256 07F5EE25587956218DB7F89CE8F82C9B054F75E0C79C1E50ADA24FAC25C8C7A5. That earlier argument and the exact source dependency were reread; they do not require flatness to reflect zero modules. No effectiveness assertion about descent of arbitrary objects is made.

\subsubsection{Surjectivity over nilpotent extensions of finite products of fields}\label{algebra-proof-021-surjectivity-over-nilpotent-extensions-of-finite-products-of-fields}

Let \(I\subset R\) be an ideal with \(I^N=0\) for some positive integer \(N\), and suppose \(R/I\) is a finite product of fields (the empty product, the zero ring, is allowed). Then every ring epimorphism \(\varphi:R\to S\) is surjective, without Noetherianity or finiteness assumptions on either module \(I\) or \(S\).

First prove the assertion for \(A=\prod_{i=1}^t k_i\). Let \(e_i\) be its original orthogonal idempotents. Any \(A\)-algebra \(B\) decomposes by the actual maps \[
B\to\prod_i e_iB,\quad b\mapsto(e_i b)_i,\qquad
(b_i)_i\mapsto\sum_i b_i .
\] Their composites are identities by \(e_i e_j=0\) for \(i\ne j\) and \(\sum_i e_i=1\). The ring \(e_iB\) is the base change \(k_i\otimes_AB\); if \(A\to B\) is epic, so is \(k_i\to e_iB\). The field result makes each component zero or makes its structural map an isomorphism. Each component map is therefore surjective, as is the finite product map \(A\to B\). For \(t=0\), the unital target is zero.

For the given \(R,I,S\), the base change \(R/I\to S/IS\) is epic, so it is surjective by the preceding paragraph. Thus \(S=\varphi(R)+IS\) as an \(R\)-module. With \(C=S/\varphi(R)\), the image of \(IS\) in \(C\) is exactly \(IC\), hence \(C=IC\). Iterating this equality \(N\) times gives \(C=I^NC=0\), without applying Nakayama to an infinite module. This proves surjectivity and completes the source's finite-ring case above.

The hypothesis extends beyond Artinian rings. Fix a field \(k\) and the vector space \(V=\bigoplus_{n\ge1} k e_n\). On the original abelian group \(R=k\oplus V\), define \[
(a,v)(b,w)=(ab,aw+bv),\qquad 1_R=(1,0).
\] Associativity follows by expanding either triple product to \((abc,ab u+ac w+bc v)\); commutativity and the unit follow from the displayed formula. The ideal \(I=0\oplus V\) has \(I^2=0\) and \(R/I=k\). Every epimorphism out of this \(R\) is consequently surjective. Yet its ideals \(0\oplus\operatorname{span}\{e_n:n\ge m\}\) form a strictly descending infinite chain, so \(R\) is not Artinian. The original source Artinian observation is retained; this example verifies that the separate nilpotent criterion genuinely has additional scope.

\subsubsection{Reading limits and provenance}\label{algebra-proof-021-reading-limits-and-provenance}

The inherited immutable literature routes preserve earlier reading. A bounded query for ``ring epimorphisms epicenter Mazet'' gave no hits in either corpus layer. A broader ``epimorphisms'' query gave two original-TeX routing hits, Mesablishvili's Morita theory for quantales and Schäppi's Graded-Tannakian categories of motives. Neither was content-read; neither supports a claim here. These limited results do not establish an exhaustive search or novelty. The mathematical arguments used are the original Stacks TeX, the reread earlier pure-descent proof and the explicit derivations above. The source citations Autour, Mazet and Kunen remain intact; reading their citations is not claimed as reading those works.

The next unadjudicated section starts at 26192, Pure ideals. Earlier lookahead through 26260 stops inside its proof and is not adjudicated. The source chapter, translations, retained additions and unrelated files are unchanged. No TeX compilation or publication was performed.

\hypertarget{algebra_pure_global_notes_20260928}{}
\subsection{Pure ideals and finite global dimension: editorial review}\label{algebra-proof-022-pure-ideals-and-finite-global-dimension-editorial-review}

Primary source: Stacks Project authors, algebra.tex at git a04446e57ec1fbc252a871afcec7752fb2807b14, complete lines 26192--26913, SHA-256 FA8BB92E58A4F78A2BD01B3B6A4A87DE0A0D279F5DD90641B574DD5FBFFFA4F3. All twelve reports in that interval were read. The complete current interval contains exactly the earlier wording correction MC-STK-ERR-1599 and 112 lines of separately admitted FAC material, captured verbatim in ALGEBRA\_PURE\_GLOBAL\_RETAINED\_EXTRAS\_20260928.json. Those additions were read and are retained exactly.

All complete arguments, boundary clarifications and consequences in this file are editorial material. They do not replace the source's translation, including the current branch. Minimal wording, notation and reference proposals remain separately enumerated. An omitted source case is not classified as a translation error. No novelty or complete literature-survey claim is made. Later synthesis remains after the core intake.

\subsubsection{The eleven pure-ideal conditions}\label{algebra-proof-022-the-eleven-pure-ideal-conditions}

For 26199--26286 retain the original ring \(R\), ideal \(I\), and every auxiliary ideal \(J\). There are natural maps \[
J\otimes_RR/I\longrightarrow J/JI,\quad j\otimes\bar r\mapsto rj+JI,
\qquad j+JI\longmapsto j\otimes1
\] which are well-defined inverses. The map to \(R/I\) is \(j+JI\mapsto j+I\), with kernel \((J\cap I)/JI\) and cokernel \(R/(I+J)\). Thus the flatness criterion in lemma-flat proves (1), (2), (3) equivalent; the criterion tests injectivity for all finitely generated ideals. The two received reports at 26232 restore the actual quotient parentheses.

Conditions (2) and (3) imply (4) by taking principal ideals. If \(x\in I\), then \(x\in(x)\cap I=xI\), giving (5). For the finite list in (6), choose \(x_i=x_i y_i\), \(y_i\in I\), and retain exactly \[
y=1-\prod_{i=1}^n(1-y_i)\in I .
\] Each \(x_i(1-y)\) contains the vanishing factor \(x_i(1-y_i)\). For the empty list the product is 1 and \(y=0\). Thus (5) and (6) are equivalent.

The set \(1+I\) is multiplicative: \((1+a)(1+b)=1+(a+b+ab)\). In (5), \(x(1-y)=0\) shows \(x\) vanishes after localization. Conversely if \(x\) vanishes, some \(1+a\), \(a\in I\), kills it, and \(x=(-a)x\). For any ideal \(I\) the localization kernel lies in \(I\), because every \(1+a\) maps to 1 in \(R/I\) and the quotient map factors through the localization. This proves the equivalence with (9), with the zero localization included.

For a prime \(\mathfrak p\notin V(I)\), any element of \(I\setminus\mathfrak p\) becomes a unit, so \(IR_{\mathfrak p}=R_{\mathfrak p}\). For \(\mathfrak p\in V(I)\), the equation \(x(1-y)=0\) has \(1-y\notin\mathfrak p\); hence every \(x/1\) vanishes and \(IR_{\mathfrak p}=0\). The received repeated-verb correction at 26259 expresses this consequence without changing the equation. Condition (7) makes every localized quotient either \(R_{\mathfrak p}\) or zero, both flat, so the local flatness criterion proves (1). For (8), outside \(V(I)\) the ideal localization is the nonzero local ring, and on \(V(I)\) it is zero; conversely (8) supplies the zero case and the unit argument supplies the other case. If \(R=0\) there are no primes and every module under discussion is zero.

Under these conditions the explicit isomorphism in (11) is \[
R/I\longrightarrow(1+I)^{-1}R,\quad r+I\mapsto r/1,
\qquad r/(1+a)\longmapsto r+I .
\] The second map is well-defined by the localization universal property. In the localization every \(a\in I\) is zero, so \(1+a\) equals 1; both composites are identities. This gives the source surjectivity assertion at 26277--26279 directly, without assuming a finite ideal. Condition (11) implies (10); every localization is flat, proving (10) implies (1).

\subsubsection{Closed sets, exact kernels and idempotents}\label{algebra-proof-022-closed-sets-exact-kernels-and-idempotents}

For 26288--26364 a pure ideal is recovered by the source's actual kernel \[
I=\ker\left(R\longrightarrow\prod_{\mathfrak p\in V(I)}R_{\mathfrak p}\right).
\] The inclusion from left to right follows from the preceding local calculation. If \(x\notin I\), its nonzero class in \(R/I\) has a proper annihilator. A maximal ideal containing that annihilator contains \(I\), and the class remains nonzero there: a denominator killing it would lie in the annihilator and outside that maximal ideal. Since \(I_{\mathfrak m}=0\), \(x\) itself is nonzero in \(R_{\mathfrak m}\). This proves the reverse inclusion. For \(I=R\) the empty product is the zero ring and its kernel is \(R\), as required.

The vanishing locus of a pure ideal is closed under generalizations directly. If \(\mathfrak q\subset\mathfrak p\), \(I\subset\mathfrak p\), and \(x\in I\), choose \(x=xy\), \(y\in I\). Then \(1-y\notin\mathfrak p\), hence not in \(\mathfrak q\); primality in \(x(1-y)=0\) forces \(x\in\mathfrak q\).

Conversely retain the source closed generalization-stable set \(Z=V(J)\), with \(J\) radical, and \(I=\{x\in R:\exists f\in J,\ x=xf\}\). If \(x=xf\) and \(y=yg\), then \((x+y)(f+g-fg)=x+y\): for \(x\), the terms \(xg-xfg\) cancel because \(xf=x\); for \(y\), \(yf-yfg\) cancel because \(yg=y\). The multiplier belongs to \(J\). Scalar multiples and additive inverses retain the same multiplier, so \(I\) is an ideal. Also \(I\subset J\).

For a prime \(\mathfrak p\) containing \(I\), consider exactly \(S=(R\setminus\mathfrak p)(1+J)\). If \(t(1+j)=0\) with \(t\notin\mathfrak p\), then \(t=t(-j)\), hence \(t\in I\), a contradiction. Thus \(0\notin S\). A maximal ideal of the nonzero localization \(S^{-1}R\) contracts to a prime \(\mathfrak q\) disjoint from \(S\); consequently \(\mathfrak q\subset\mathfrak p\) and \(\mathfrak q\cap(1+J)=\varnothing\). The latter says \(\mathfrak q+J\ne R\), since \(1=q+j\) would put \(1-j=q\) in that intersection. Choose a maximal ideal \(\mathfrak m\supset\mathfrak q+J\). Generalization stability gives \(\mathfrak q\in Z\) from \(\mathfrak m\in Z\), and closedness gives \(\mathfrak p\in Z\). Therefore \(V(I)=V(J)\) and \(\sqrt I=J\). For \(x=xy\) with \(y\in J\), choose \(n>0\) with \(y^n\in I\); iteration gives \(x=xy^n\), proving purity by condition (5). This proves both existence and uniqueness, not only uniqueness.

For 26366--26399, if \(I\) is pure and generated by \(x_1,\ldots,x_n\), condition (6) gives \(e\in I\) acting as the identity on all generators. It acts as the identity on \(I\), so \(e^2=e\) and \(I=eR\). Conversely \(R=eR\oplus(1-e)R\), and \(R/I\cong(1-e)R\) is projective; \(V(e)=D(1-e)\). If \(R/I\) is projective, choose a splitting of \(R\to R/I\). The complementary projection onto \(I\) is \(R\)-linear, hence multiplication by \(e\), its value on 1. Being a projection gives \(e^2=e\) and image \(I=eR\). If \(I\) is pure and \(V(I)\) is open, the source clopen-idempotent correspondence supplies \(e\) with \(V(I)=V(e)\); uniqueness of pure ideals then gives \(I=eR\). The zero ideal and unit ideal correspond to \(e=0,1\), also in the zero ring.

\subsubsection{Finite flat modules and their rank strata}\label{algebra-proof-022-finite-flat-modules-and-their-rank-strata}

For 26405--26469, the implication from finite locally free modules to open closed-generalization-stable subsets requires the existence result lemma-pure-open-closed-specializations. The cited uniqueness lemma alone cannot produce an ideal for an arbitrary such set. The two reports at 26419 receive the reference correction.

Here are the exact rank and relation arguments for the converse, which will also be used below. A finite flat \(M\) localizes to a finite free module at every prime by lemma-finite-flat-local. Its support is \(V(\operatorname{Ann}M)\), and if \(\mathfrak p\subset\mathfrak p'\), a nonzero free \(M_{\mathfrak p'}\) remains a nonzero free module after localization to the nonzero ring \(R_{\mathfrak p}\). Thus support is generalization-stable. The punctuation report at 26429 terminates exactly that localization equality.

For a generating list of length \(r\), each \(\bigwedge^i M\) is finite, and its localization is the exterior power of the finite free localization of \(M\). It is therefore flat by the local criterion. Under condition (1) each support \(W_i=\operatorname{Supp}(\bigwedge^iM)\) is clopen. Put \(W_0=\operatorname{Spec}R\), \(W_{r+1}=\varnothing\), and \(U_i=W_i\setminus W_{i+1}\). These are the original rank-\(i\) strata: at a prime where \(M_{\mathfrak p}\) has rank \(q\), \(\bigwedge^i M_{\mathfrak p}\) is nonzero exactly when \(i\le q\). Thus the \(U_i\) partition the spectrum into clopen sets.

At \(\mathfrak p\in U_i\), choose \(i\) of the original generators whose residues form a basis. Nakayama makes their map \(R_{\mathfrak p}^i\to M_{\mathfrak p}\) onto. Its finite cokernel vanishes after localizing at one \(f\notin\mathfrak p\): kill a finite generating set by a product of denominators. Shrink \(D(f)\) inside the open \(U_i\). For \(i=0\) this already gives \(M_f=0\), which is the required rank-zero isomorphism. For \(i>0\), the surjection \(R_f^i\to M_f\) gives a cyclic flat \(\bigwedge^iM_f\cong R_f/J\), with support all of \(\operatorname{Spec}R_f\). The preceding kernel criterion gives \(J=0\).

Explicitly, if \(\sum_{j=1}^i a_jm_j=0\), wedge this equality with the ordered list of all \(m_k\) except \(m_j\). Every term except the \(j\)th vanishes; that term is \((-1)^{j-1}a_j(m_1\wedge\cdots\wedge m_i)\), up to the fixed choice of putting the omitted-list wedge after the relation. Its coefficient is zero because this top wedge is a basis. Thus every \(a_j=0\), proving the original map injective. (With the relation first and the remaining indices increasing, the sign is exactly \((-1)^{j-1}\).)

This proves finite local freeness. The argument also proves the precise reusable statement: a finite flat module is finite locally free if its local rank function is locally constant. Indeed for finite \(r\) its rank fibres \(U_i\) are then clopen, and the same maps apply. Conversely local freeness makes the ranks locally constant. Finite locally free is equivalent to finite projective by lemma-finite-projective.

\subsubsection{Schanuel maps with their original signs}\label{algebra-proof-022-schanuel-maps-with-their-original-signs}

The two wording reports at 26495 are already corrected by MC-STK-ERR-1599. No new operation is allocated.

In 26484--26566, retain \(c_1,c_2,p_1,p_2,a,b\) and \(N=\ker(p_1,p_2)\), where the map is the sum, not a difference. Projection \(N\to P_1\) is onto by choosing a lift of \(-p_1(u)\) in \(P_2\); its kernel is \(c_2(L)\). Projection to \(P_2\) is onto by the corresponding lift of \(-p_2(v)\); its kernel is \(c_1(K)\). Their projective targets split these sequences.

The source matrices and inverses are \[
S=\begin{pmatrix}1&0\\a&1\end{pmatrix},\
S^{-1}=\begin{pmatrix}1&0\\-a&1\end{pmatrix},\qquad
T=\begin{pmatrix}1&b\\0&1\end{pmatrix},\
T^{-1}=\begin{pmatrix}1&-b\\0&1\end{pmatrix}.
\] Direct multiplication gives identity. Moreover \((0,p_2)S=(p_1,p_2)\) and \((p_1,0)T=(p_1,p_2)\). Thus the map from the top exact sequence to the bottom one is \[
ST^{-1}=\begin{pmatrix}1&-b\\a&1-ab\end{pmatrix},
\qquad
TS^{-1}=\begin{pmatrix}1-ba&b\\-a&1\end{pmatrix}
\] for its inverse; the orders \(ab\) and \(ba\) are retained.

On \(K\oplus P_2\) the first map sends \((k,v)\), through its specified inclusion, to \((c_1k-bv,\ ac_1k+(1-ab)v)\). The second component belongs to \(\ker p_2=c_2(L)\), since applying \(p_2\) gives \(0+p_2v-p_1bv=0\). Conversely for \((u,\ell)\in P_1\oplus L\), the inverse pair in \(P_1\oplus P_2\) is \(((1-ba)u+bc_2\ell,-au+c_2\ell)\), whose first component lies in \(\ker p_1=c_1(K)\). The injectivity of \(c_1,c_2\) identifies the actual kernel coordinates uniquely. This proves all vertical isomorphisms and commutativity in the source diagram.

\subsubsection{Syzygies, augmented indices and the reported boundary cases}\label{algebra-proof-022-syzygies-augmented-indices-and-the-reported-boundary-cases}

For 26584--26627 use induction on an upper bound \(d\) for projective dimension, so a resolution of length less than \(d\) is also covered. If \(d=0\), \(F_0\to M\) splits and its kernel is projective; repeat this along the remaining finite complex to reach the stated final kernel. For \(d>0\), let \(K_F=\ker(F_0\to M)\) and \(K_P=\ker(P_0\to M)\), using a projective resolution of length at most \(d\). Schanuel gives \(P_0\oplus K_F\cong F_0\oplus K_P\). The right module has a projective resolution of length at most \(d-1\), obtained by adjoining \(F_0\) to the last surjection onto \(K_P\). For \(e=0\) the hypotheses give \(d\le1\); the displayed isomorphism immediately makes \(K_F\) projective. For \(e>0\), the source exact sequence ending in \(P_0\oplus K_F\) has projective terms and length \(e-1\). Its final kernel is the original kernel for \(e\ge2\). For \(e=1\) its final map is the identity on \(P_0\) plus \(F_1\to K_F\), so its kernel is exactly \(\ker(F_1\to F_0)\). Induction proves the claim. This verifies the bound and all cases used in 26629--26675.

The two reports at 26690--26698 identify omitted cases in an otherwise valid theorem, rather than a false theorem or translation error. They are resolved here as a source-linked proof clarification, with zero source operations. For \(d=0\), \(M\) is already finite projective and the resolution is \(0\to M\xrightarrow{1}M\to0\). For \(d=1\), take \(P_1=\ker(F_0\to M)\); it is projective by the earlier criterion and finite by Noetherianity, so \(0\to P_1\to F_0\to M\to0\) is the desired resolution. For \(d\ge2\) the source's original \(P_d=\ker(F_{d-1}\to F_{d-2})\) and displayed sequence apply unchanged. Over a local ring each finite projective term is finite free, proving the stated local variant. None of this supplies undefined negative-index free modules.

For the Ext argument at 26724--26759, the required notation clarification is \[
d_0:F_0\to M,\quad d_e:F_e\to F_{e-1}\ (e\ge1),\quad
P_{-1}=M,\quad P_e=\ker d_e\ (e\ge0).
\] Then \(0\to P_e\to F_e\to P_{e-1}\to0\) is exact for every \(e\ge0\). If \(\operatorname{pd}M\le n\), each \(P_e\) for \(e\ge n-1\) is projective, with \(P_{-1}=M\) in the \(n=0\) case. Whenever \(e\ge n\), projectivity of \(P_{e-1}\) gives the actual split sequence \(F_e\cong P_e\oplus P_{e-1}\). Its differential projects to \(P_{e-1}\) and includes that summand in the next term. Applying \(\operatorname{Hom}_R(-,N)\) makes the cochain tail split exact above degree \(n\), so the required Ext groups vanish. The two small source operations specify these indices; they do not replace the proof.

Conversely, if \(\operatorname{Ext}^{n+1}_R(M,N)=0\) for every \(N\), the case \(n=0\) makes \(M\) projective: a surjection from a free module has extension class zero and splits. If \(n>0\), a free cover with kernel \(K\) gives \(\operatorname{Ext}^n_R(K,N)\cong\operatorname{Ext}^{n+1}_R(M,N)\) by the long exact Ext sequence and vanishing of positive Ext from the free module. Induction gives \(\operatorname{pd}K\le n-1\), and adjoining the free cover gives the required resolution of \(M\).

For the three short-exact-sequence bounds at 26761--26782, apply that criterion to the original contravariant long exact sequence. The first uses \(\operatorname{Ext}^i(M,N)\to\operatorname{Ext}^i(M',N)\to\operatorname{Ext}^{i+1}(M'',N)\) for \(i\ge n+1\); the second uses \(\operatorname{Ext}^i(M'',N)\to\operatorname{Ext}^i(M,N)\to\operatorname{Ext}^i(M',N)\); the third uses \(\operatorname{Ext}^{i-1}(M',N)\to\operatorname{Ext}^i(M'',N)\to\operatorname{Ext}^i(M,N)\) for \(i\ge n+2\). The outside terms vanish with exactly the source bounds.

\subsubsection{Transfinite filtrations and the cyclic-module test}\label{algebra-proof-022-transfinite-filtrations-and-the-cyclic-module-test}

For 26799--26853 the \(n=0\) section maps give the source \(\psi:\bigoplus_eP_e\to M\). A nonzero finite-support element has a largest index; its image modulo the union of earlier submodules is its nonzero last coordinate, proving injection. For surjection, transfinite induction subtracts the chosen lift of the current quotient class, leaving an element at an earlier stage. No infinite sum is taken. If the index set is empty, \(M=0\) and the claim holds.

For \(n>0\), retain \(F_e\) as the free module on all elements of \(M_e\), including its zero element. Its transition maps include basis sets, and its kernel \(K_e\) maps injectively into later kernels. The unions give an exact sequence \(0\to K\to F\to M\to0\): every element has finite support and therefore occurs in one stage. At a limit, the same finite-support argument makes the earlier union sequence exact.

The quotient \(F_e/F_{<e}\) is free on the set of elements of \(M_e\) not belonging to any \(M_{e'}\) with \(e'<e\). At the least index there are no earlier sets, so its basis is all of \(M_e\), still including zero. The zero module used as the empty union of submodules must not be mistaken for an earlier underlying basis set. With this precise convention the source free-basis argument is valid; no defect or new source correction is inferred from it.

The injective maps of the earlier and current exact sequences give, by the kernel/cokernel sequence, exactly \[
0\to K_e/K_{<e}\to F_e/F_{<e}\to M_e/M_{<e}\to0.
\] The middle term is free and the last has projective dimension at most \(n\). The first short-exact-sequence bound gives projective dimension at most \(n-1\) for the kernel quotient. Induction gives the same bound for \(K\), and adjoining \(F\) gives the bound \(n\) for \(M\).

For 26855--26878 choose the original well-ordered generating set and its generated submodules. Each successive quotient is cyclic (zero is allowed as \(R/R\)), so the proved filtration result transfers the uniform cyclic bound to every module. The other implications are direct. For 26880--26903 localization preserves exactness and projective summands, hence the original localized resolution. For any \(S^{-1}R\)-module \(N\), the maps \(S^{-1}N\to N\), \(n/s\mapsto s^{-1}n\), and \(n\mapsto n/1\) are inverse, giving the global bound.

\subsubsection{Finite resolution data in place of Noetherianity}\label{algebra-proof-022-finite-resolution-data-in-place-of-noetherianity}

Here is a separate weakening of the hypotheses of 26677--26722. Let \(R\) be any ring, \(d\ge0\), and suppose \(M\) has projective dimension at most \(d\). For \(d=0\), suppose \(M\) is finite. For \(d\ge1\), suppose the original module has an exact partial resolution \[
F_d\longrightarrow F_{d-1}\longrightarrow\cdots
\longrightarrow F_0\longrightarrow M\longrightarrow0
\] with all \(F_i\) finite free. Then \(M\) has a resolution of length at most \(d\) by finite projective modules, with no Noetherianity assumption.

For \(d=0\) this is the finite projective module itself. For \(d\ge1\), put \(K_d=\ker(F_{d-1}\to F_{d-2})\), using \(M\) as the target when \(d=1\). The already proved projective-dimension criterion makes \(K_d\) projective. Exactness makes it the image of the finite module \(F_d\), hence finite. Truncate the same original complex at \(K_d\); all lower terms remain the original \(F_i\). Over a local ring this is a finite free resolution, since finite projectives are free. The partial-resolution hypothesis is often called finite presentation through degree \(d\); its full displayed meaning, rather than that terminology, is what is assumed here.

There is a corresponding extension of the retained local FAC addition. If \(R\) is any nonzero local ring and \(M\ne0\) has finite projective dimension \(d\) and the finite resolution data just specified, then \(\operatorname{Ext}^d_R(M,R)\ne0\). The unit removal below uses the complete original-coordinate proof of ALGEBRA-RECON-569, ALGEBRA\_COMPLEX\_EXACTNESS\_NOTES\_20260928.md, section ``The original pivot and its exact splitting maps'', SHA-256 02D6389E6F820C53832BD9691965D4D8479783E8B80D9BDD761B077E29B7ECC3, reread for this application. For the actual coefficient \(u=a_{pq}\), its bases are \(E_p=e_p\), \(E_j=e_j-a_{jq}u^{-1}e_p\), \(F_q=\sum_k a_{pk}f_k\), \(F_k=f_k\) for the other indices. The complementary coefficients are exactly \(a_{jk}-a_{jq}u^{-1}a_{pk}\); the incoming complex equation kills the \(E_p\)-coordinate and the outgoing map kills \(F_q\). Its inverse basis formulas and chain homotopy retain every coefficient in the cited proof. Thus its identity-disk removal preserves the original module and all homology, not only a matrix rank. To prove it, take the resulting finite free resolution. Any unit matrix entry can be made a one-term identity block by invertible row and column operations; the adjacent differential equations force the incoming and outgoing entries for that block to vanish. Remove this contractible pair. The total finite rank strictly decreases, so the process ends with all differential entries in the maximal ideal. Minimality of \(d\) forces the last term \(F_d\ne0\). For \(d\ge1\), the image of \(\operatorname{Hom}(F_{d-1},R)\to\operatorname{Hom}(F_d,R)\) lies in \(\mathfrak m\operatorname{Hom}(F_d,R)\); its cokernel surjects onto the nonzero residue vector space of the last dual free module. For \(d=0\), \(M\) is nonzero finite free and \(\operatorname{Hom}(M,R)\ne0\) directly. Positive Ext above degree \(d\) vanishes by the resolution, so testing Ext against \(R\) alone detects every nonnegative bound, under this finite-data hypothesis. This proof retains the actual complex and does not assume that arbitrary finite modules over non-Noetherian rings have such a resolution.

\subsubsection{Pure ideals under base change and pure descent}\label{algebra-proof-022-pure-ideals-under-base-change-and-pure-descent}

For any \(\varphi:R\to S\), if \(I\) is pure then \(IS\) is pure: the canonical quotient isomorphism \[
S\otimes_R(R/I)\longrightarrow S/IS,\quad
s\otimes(r+I)\longmapsto s\varphi(r)+IS
\] has inverse \(s+IS\mapsto s\otimes1\), and base change preserves flatness. The local-unit proof is equally explicit. Express finitely many generators of any chosen element of \(IS\) using finitely many coefficients from \(I\); choose one \(y\in I\) fixing all those coefficients by condition (6). Then \(\varphi(y)\in IS\) fixes the original element.

Conversely if \(R\to S\) is pure and \(IS\) is pure, \(S/IS\) is flat over \(S\); the flatness-descent argument of ALGEBRA-RECON-501 applied to the original module \(R/I\) proves it flat over \(R\). Hence \(I\) is pure. Purity also makes the actual map \(R/I\to (R/I)\otimes_RS=S/IS\) injective, which says exactly \(\varphi^{-1}(IS)=I\). No injectivity of a different kernel tensor map and no flatness of \(S\) is assumed. The earlier full argument is bound at ALGEBRA\_OPEN\_LOCI\_PURE\_DESCENT\_NOTES\_20260927.md, ``Descent along the original universally injective ring map'', SHA-256 07F5EE25587956218DB7F89CE8F82C9B054F75E0C79C1E50ADA24FAC25C8C7A5.

\subsubsection{Nil-ideal invariance with the actual quotient maps}\label{algebra-proof-022-nil-ideal-invariance-with-the-actual-quotient-maps}

Let \(N\subset R\) be a nil ideal, without a uniform nilpotence exponent. Every prime contains every element of \(N\), so the inverse-image bijection \(\operatorname{Spec}(R/N)\to\operatorname{Spec}R\) preserves inclusions and identifies the vanishing set of any ideal image with its original vanishing set. It is a homeomorphism by those same equations for closed sets.

Consequently \(I\mapsto(I+N)/N\) gives a bijection between pure ideals of \(R\) and of \(R/N\). Base change proves the image pure. For surjection, let \(\bar J\) be pure in \(R/N\), put \(L=\pi^{-1}(\bar J)\), and use the explicit inverse \[
I=\{x\in R:\exists y\in L,\ x=xy\}.
\] Indeed \(V(L)\) is closed and generalization-stable. The earlier constructive proof with \(J=\sqrt L\) gives a pure ideal by that formula with \(y\in\sqrt L\). If \(x=xy\) and \(y^m\in L\), iteration gives \(x=xy^m\); therefore the same ideal is given by the displayed formula with \(L\) itself. Its vanishing set is \(V(L)\), so its image, a pure ideal with the same vanishing set as \(\bar J\), equals \(\bar J\). Uniqueness of the pure ideal from its vanishing set proves injectivity as well.

This does not identify \(I\) with \(\bar J\): its actual kernel is \(I\cap N=IN\), by condition (2) for the pure ideal \(I\), so the exact comparison is \(I/IN\cong\bar J\). Finite generation of corresponding pure ideals is equivalent, because the source criterion identifies it with openness of their common vanishing set. Thus idempotent-generated pure ideals correspond as well.

The property that every finite flat module is finite projective is also invariant under \(R\to R/N\), by the original closed-generalization-stable-set criterion and this homeomorphism. More precisely, for an individual finite flat \(R\)-module \(M\), \[
M\text{ finite projective over }R
\quad\Longleftrightarrow\quad
M/NM\text{ finite projective over }R/N .
\] The forward implication is base change. For the reverse, the residue fields at corresponding primes are the same via the quotient map, and \((M/NM)\otimes_{R/N}\kappa(\mathfrak p/N)
\cong M\otimes_R\kappa(\mathfrak p)\). These isomorphisms send \((m+NM)\otimes\bar r\) to \(m\otimes r\). The rank of the finite free stalk of \(M\) equals this fibre dimension; hence its rank function is locally constant when the quotient is projective. The exact rank-stratum argument above proves \(M\) finite locally free, hence finite projective. The original flatness hypothesis on \(M\) is essential to this argument and is retained; no assertion of projectivity descent for arbitrary modules along a nil quotient is made.

\subsubsection{Global dimension along flat epimorphisms}\label{algebra-proof-022-global-dimension-along-flat-epimorphisms}

The source localization bound extends to every flat ring epimorphism \(\varphi:R\to S\). For any \(S\)-module \(Q\), the multiplication map \[
S\otimes_RQ\longrightarrow Q,\quad s\otimes q\mapsto sq
\] is an isomorphism with inverse \(q\mapsto1\otimes q\). To verify its \(S\)-linearity and inverse identities, tensor the already proved isomorphism \(S\otimes_RS\to S\) over its second \(S\)-action with \(Q\). The explicit comparison is \((s\otimes t)\otimes q\mapsto s\otimes tq\), with inverse \(s\otimes q\mapsto(s\otimes1)\otimes q\). The equality \(s\otimes1=1\otimes s\) yields \(s\otimes q=1\otimes sq\), exactly the needed identity.

A projective \(R\)-resolution of the underlying \(Q\) stays exact after tensoring with the flat \(S\); its terms stay projective because tensoring preserves a given direct-summand decomposition of a free module. The last term is the original \(Q\) by the preceding map. Thus \[
\operatorname{pd}_S(Q)\le\operatorname{pd}_R(Q)
\] whenever the right side is finite. In particular a bound \(n\) for the global dimension of \(R\) is also a bound for that of \(S\). The zero target is included. This receives the exact tensor-factor characterization of ALGEBRA-RECON-606, rather than merely an abstract resemblance to localization.

For the present section's pure ideal \(I\), the map \(R\to R/I\) is flat by definition and epic by surjectivity. Hence the same inequality applies to every \(R/I\)-module with its actual restricted \(R\)-action. The source identifies this quotient with \((1+I)^{-1}R\); the two proofs agree through that explicit isomorphism, and no new multiplication or quotient presentation is substituted.

\subsubsection{Retained FAC additions and their degree-zero clarification}\label{algebra-proof-022-retained-fac-additions-and-their-degree-zero-clarification}

The current source adds a reference/history block at definition-finite-proj-dim, citing FAC Chapter III, §5, no. 74, pp.~268--269, and two Ext lemmas after lemma-projective-dimension-ext, citing FAC no. 74, Lemma 1, and no. 77, Lemma 3. The entire 16-line block and 96-line pair of lemmas and proofs are captured and preserved byte for byte. Reading these local additions is not a fresh reading of FAC.

The local added lemma's theorem is valid. Its proof at current algebra.tex:27340 and following, however, writes the top Ext cokernel using \(F_{d-1}\) after handling only \(M=0\). For a nonzero projective \(M\), \(d=0\) and \(F_{-1}\) is undefined in that dual complex. This is recorded as a conceptual boundary clarification of the retained addition, with no source operation against the upstream authority. Handle \(d=0\) by the nonzero finite free module \(\operatorname{Hom}(M,R)\); use the stated top-cokernel proof only for \(d\ge1\), exactly as proved above. The augmented convention \(F_{-1}=M\) from the resolution argument is not a valid substitute inside this Hom cochain complex. This distinction prevents a false degree-zero cokernel formula.

The global added lemma at current label line 27390 also holds. For \(d=0\), choose a maximal ideal in the support of the nonzero finite projective module. For \(d>0\), if every maximal localization had projective dimension at most \(d-1\), the finite syzygy testing that bound would be projective at every maximal localization and hence finite projective over \(R\), by lemma-finite-projective and Noetherian finite presentation. When \(d=1\), the module testing the zero bound is \(M\) itself. This contradicts the definition of \(d\). Thus some maximal localization attains \(d\). A finite projective resolution computes Ext, its Hom terms commute with localization, and exact localization commutes with its cohomology; therefore \(\operatorname{Ext}^d_R(M,R)_{\mathfrak m}
\cong\operatorname{Ext}^d_{R_{\mathfrak m}}(M_{\mathfrak m},R_{\mathfrak m})\ne0\). The corrected local argument supplies the last nonvanishing.

The finite-projective-dimension hypothesis in both additions must be retained. For a field \(k\), let \(A=k[\epsilon]/(\epsilon^2)\) and \(M=k=A/(\epsilon)\). The exact infinite free resolution has every differential multiplication by \(\epsilon\), since its kernel and image are both \(\epsilon A\). Tensoring with \(k\) makes all these differentials zero, so \(\operatorname{Tor}^A_i(k,k)=k\) for every \(i\ge0\); consequently \(k\) has infinite projective dimension. But applying \(\operatorname{Hom}_A(-,A)\) gives the same multiplication-by-\(\epsilon\) maps, exact in every positive degree, so \(\operatorname{Ext}_A^i(k,A)=0\) for all \(i>0\). This explicitly checks the adverse case rather than silently removing the finite-dimension hypothesis.

\subsubsection{Source and reading limits}\label{algebra-proof-022-source-and-reading-limits}

The two new bounded corpus queries yielded eight routing records across both layers, including two PDF-primary records and duplicate-topic or unrelated material. None was content-read or used for a theorem. They remain routing records, not proof or evidence of novelty. Earlier immutable routes and their actual reading coverage are retained.

Authority 26914--27280 was read as lookahead, including the regular/global-dimension section and the beginning of Auslander--Buchsbaum; it is not adjudicated. The last proof is incomplete at the lookahead boundary. The next review starts at 26914. The original human references Lazard, Bkouche, DeMarco and Auslander, the retained FAC materials, and unrelated dirty files remain unchanged. No translation/source mutation, rendering or publication occurred.

\hypertarget{algebra_regular_ab_notes_20260928}{}
\subsection{Regular rings and Auslander--Buchsbaum: source review}\label{algebra-proof-023-regular-rings-and-auslanderbuchsbaum-source-review}

Primary source: Stacks Project authors, algebra.tex at git a04446e57ec1fbc252a871afcec7752fb2807b14, complete lines 26914--27304, SHA-256 FA8BB92E58A4F78A2BD01B3B6A4A87DE0A0D279F5DD90641B574DD5FBFFFA4F3. All twelve due reports were read. The current complete interval is exactly the authority with MC-STK-ERR-1600 retained and the captured FAC addition at the end. Its 49-line lemma/proof and surrounding whitespace are preserved by ALGEBRA\_REGULAR\_AB\_RETAINED\_EXTRA\_20260928.json.

These full arguments, boundary clarifications and consequences are separate editorial material. The source statements and exposition remain identifiable, including in the current translation. Four proposed text operations are separately enumerated. An omitted low-degree proof case is not treated as a translation error. No novelty claim or later combined-results synthesis is asserted.

\subsubsection{Finite depth and the maximal-localization criterion}\label{algebra-proof-023-finite-depth-and-the-maximal-localization-criterion}

At 26921--26936 the indexed resolution requires \(M\ne0\). The original depth definition, reread at 17765--17787, assigns depth \(+\infty\) to zero. Thus \(F_{d-e}\) is undefined for \(M=0\). The minimal proposal adds ``nonzero''; the zero module itself has the zero free resolution and creates no exception to the global-dimension bound.

For nonzero finite \(M\), \(e\) is an integer with \(0\le e\le d\). The full original syzygy construction in ALGEBRA-RECON-589 gives a maximal Cohen--Macaulay terminal module at step \(d-e\), including the actual identity complex for \(e=d\). The source regular-ring theorem makes that finite module free. The resulting original resolution gives projective dimension at most \(d-e\). Every finite module therefore has projective dimension at most \(d\), including zero separately; the cyclic-module criterion then gives the stated bound for every module.

At 26938--26963 localization proves the forward implication. The two wording reports receive ``We saw in Lemma''. For the converse, let \(M\) be finite. If \(n=0\), put \(K_0=M\); every maximal localization is projective, and the finite-presentation local criterion makes \(M\) finite projective. This supplies the omitted case in a linked proof clarification, without imposing a negative-index convention on the displayed source complex. If \(n\ge1\), take the original finite free terms up to \(F_{n-1}\). Their terminal kernel \(K_n\) is finite by Noetherianity. Its maximal localizations are projective by the resolution-independence criterion; hence \(K_n\) is finite projective over \(R\). Splicing it to the same free terms proves \(\operatorname{pd}M\le n\). Apply the cyclic criterion again. Over the zero ring all modules are zero, so this bound and its vacuous maximal-localization condition both hold.

\subsubsection{The Koszul map, its signs and a characteristic-free left inverse}\label{algebra-proof-023-the-koszul-map-its-signs-and-a-characteristic-free-left-inverse}

For 26965--27069 keep the source elements \(x_1,\ldots,x_n\), their original basis classes in \(\mathfrak m/\mathfrak m^2\), and \[
d(e_{j_1}\wedge\cdots\wedge e_{j_i})
=\sum_{a=1}^i(-1)^{a+1}x_{j_a}
 e_{j_1}\wedge\cdots\widehat e_{j_a}\cdots\wedge e_{j_i}.
\] For \(a<b\), removing \(a\) and then \(b\) has sign \((-1)^{a+1}(-1)^b\); removing \(b\) and then \(a\) has sign \((-1)^{b+1}(-1)^{a+1}\). Their coefficients \(x_{j_a}x_{j_b}\) coincide and their signs are opposite. Thus \(d^2=0\), also in characteristic two. Nakayama makes the \(x_j\) generate \(\mathfrak m\), so the degree-zero cokernel is the original residue field \(\kappa\).

Choose a minimal free resolution \(F_\bullet\to\kappa\). In the finite-projective-dimension case it may be taken finite; otherwise a minimal resolution exists degree by degree, since each kernel over the Noetherian local ring is finite. At each step choose a minimal generating set of the kernel. Its relations have all coefficients in \(\mathfrak m\), because a relation with a unit coefficient would express one chosen generator in terms of the others. Consequently all differentials have entries in \(\mathfrak m\), and \(F_0\) has rank one.

A comparison \(\alpha:K_\bullet\to F_\bullet\) lifting \(1_\kappa\) can be constructed directly, without assuming the Koszul complex is a resolution. Lift the augmentation from \(K_0=R\). At each next degree the image of \(\alpha_{i-1}d_K\) is a cycle in the exact target complex; the surjection \(F_i\to Z_{i-1}(F)\) and freeness of \(K_i\) lift it to \(\alpha_i\). In degree one use the augmentation kernel. These maps commute with exactly the original differentials.

Let \(V=\kappa^n\), with basis \(v_j\), and \(W=\mathfrak m/\mathfrak m^2\), with basis \(w_j=x_j+\mathfrak m^2\). The first-order Koszul map is \[
\Delta_i(v_{j_1}\wedge\cdots\wedge v_{j_i})
=\sum_{a=1}^i(-1)^{a+1}
(v_{j_1}\wedge\cdots\widehat v_{j_a}\cdots\wedge v_{j_i})\otimes w_{j_a}.
\] For \(i\ge1\), define a \(\kappa\)-linear map \(L_i:\bigwedge^{i-1}V\otimes W\to\bigwedge^iV\) on the ordered basis by \[
L_i(v_{l_1}\wedge\cdots\wedge v_{l_{i-1}}\otimes w_j)
=\begin{cases}
v_j\wedge v_{l_1}\wedge\cdots\wedge v_{l_{i-1}},
 &j<l_1,\\
0,&\text{otherwise}.
\end{cases}
\] For \(i=1\) the list is empty and the first case applies to every \(j\). For \(i>n\) the source of \(\Delta_i\) is zero and the target exterior power of \(L_i\) is zero. In \(\Delta_i(v_{j_1}\wedge\cdots\wedge v_{j_i})\) with increasing indices, \(L_i\) keeps exactly the summand \(a=1\); its sign is \(+1\). Thus \(L_i\Delta_i=1\). No division by \(i\), factorial or characteristic-dependent coefficient is used.

The source's square modulo \(\mathfrak m^2\) now proves by induction that every \(\bar\alpha_i\) is injective. Its right vertical map is \(\bar\alpha_{i-1}\otimes1_W\), injective over the field. The upper horizontal map is \(\Delta_i\), just proved injective. If \(\bar\alpha_i(v)=0\), commutativity forces \(\Delta_i(v)=0\), so \(v=0\). At degree zero the nonzero augmentation composite proves injection. When the resolution is finite, \(F_n\ne0\) follows and gives \(\operatorname{pd}\kappa\ge n\), exactly as needed by the source.

The same argument gives a separate consequence without a finite-projective-dimension hypothesis. For every Noetherian local \(R\) of embedding dimension \(n\), \[
\beta_i^R(\kappa)=\dim_\kappa\operatorname{Tor}_i^R(\kappa,\kappa)
=\operatorname{rank}_RF_i\ \ge\ \binom ni\qquad(0\le i\le n).
\] The Tor equality follows because tensoring this minimal resolution with \(\kappa\) makes every differential zero. This proves the precise further bound carried by the source argument; it is not attributed as a new theorem of the source or claimed novel.

\subsubsection{The dimension bound, exact global dimension and zero rings}\label{algebra-proof-023-the-dimension-bound-exact-global-dimension-and-zero-rings}

MC-STK-ERR-1600 already changes ``have to property'' to ``have the property'' at 27084; both reports retain that identity.

For 27071--27093, when \(n=0\), the inequality \(\dim R\ge0\) holds because the local ring has its maximal prime. No positive differential \(\varphi_0\) or determinantal ideal is needed. This is a separate boundary clarification. When \(n>0\), remove identity disks using the exact original pivot and chain maps of ALGEBRA-RECON-569. The last nonzero degree is exactly \(n\), since tensoring the resulting minimal resolution with \(\kappa\) gives nonzero \(\operatorname{Tor}_n(\kappa,\kappa)\). The top free module has positive rank. The exactness criterion and the fully checked top-rank argument of ALGEBRA-RECON-576 give \[
\operatorname{depth}_{I(\varphi_n)}R\ge n,\qquad
I(\varphi_n)\subset\mathfrak m .
\] The containment is valid for this positive-rank top map; it is not inferred for a rank-zero map, whose zeroth-minor ideal is the unit ideal. Every regular sequence in that proper ideal is also in \(\mathfrak m\), so \(\dim R\ge\operatorname{depth}R\ge n\). The earlier zero-tail correction therefore propagates to this receiving proof.

Together with the Koszul bound, finite projective dimension of the residue field gives \(\dim R\ge\operatorname{pd}\kappa\ge\dim_\kappa(\mathfrak m/\mathfrak m^2)\). The reverse dimension inequality is the original height bound for a minimal generating set of \(\mathfrak m\). Thus equality holds and \(R\) is regular. Conversely the source regular-ring resolution bound gives finite global dimension. It is at most \(\dim R\) and at least \(\operatorname{pd}\kappa\), so all three dimensions equal. This proves the regular/local/global statements and their localization consequences at 27095--27161.

For the exact global characterization at 27163--27186 the ring must be nonzero. Under the existing conventions the zero ring has only the zero module, hence global dimension 0, but its spectrum is empty and its Krull dimension is \(-\infty\). The latter convention was reread in the same authority topology.tex:1397--1413, particularly 1410. The ``at least one prime/maximal ideal'' clauses are false for the zero ring. Adding ``nonzero'' is the minimal proposed source guard, while the zero-ring case remains explicitly documented here.

For a nonzero ring and a fixed nonnegative integer \(n\), the precise equivalence says that global dimension is exactly \(n\), regular Krull dimension is exactly \(n\), and all prime (or maximal) localizations are regular with dimensions at most \(n\), with at least one attaining \(n\). Every finite chain of primes is contained in some maximal ideal, and localization there preserves its strict inclusions. Conversely every chain in a localization is such a chain in \(R\). Therefore \(\dim R=\sup_{\mathfrak m}\dim R_{\mathfrak m}\); a finite nonnegative integer supremum of this nonempty set is attained. Use the proved uniform local global-dimension criterion and the local equalities to obtain each implication. The bound is uniform; pointwise finite local dimensions with no uniform bound are not claimed to give finite global dimension.

\subsubsection{Flat descent with the actual low-degree kernels}\label{algebra-proof-023-flat-descent-with-the-actual-low-degree-kernels}

For 27188--27213 a flat local map is faithfully flat. Indeed its closed fibre is nonzero because \(\mathfrak m_RS\subset\mathfrak m_S\); a flat module with nonzero closed fibre over a local ring is faithful, by the source ideal/fibre criterion.

Retain \(F_\bullet\to\kappa\), and use \(K_1=\ker(F_0\to\kappa)\), \(K_d=\ker(F_{d-1}\to F_{d-2})\) for \(d\ge2\). Flatness gives the actual kernel after tensoring. Since \(\operatorname{gldim}S=d\), these \(d\)th syzygies over \(S\) are projective and finite, hence free; finite projectivity descends under faithful flatness, so the original \(K_d\) is finite free over \(R\). The completed resolution of \(\kappa\) proves \(R\) regular.

For \(d=0\), \(S\) is a field. Locality sends \(\mathfrak m_R\) to zero, and faithful flatness makes \(R\to S\) injective. Thus \(\mathfrak m_R=0\) and \(R\) is a field. Equivalently the zeroth module is \(K_0=\kappa\), but no negative-index free term is placed in the displayed complex. The two low-degree reports receive this separate proof clarification, which handles both \(d=0\) and \(d=1\).

The equality at 27205--27206 is correct. The report is resolved by citing proposition-finite-gl-dim-regular, which proves that equality, rather than weakening the claim to the bound proved by the currently cited proposition. This is one reference operation.

There is a further reduction of hypotheses, proved here separately. If \(R\to S\) is any flat local map and \(S\) is regular local, the assumption that \(R\) is Noetherian is unnecessary. Faithful flatness and Noetherianity of \(S\) imply it: for any ideal \(I\subset R\), the ideal \(IS\) is generated by finitely many elements. Express those generators as finite sums of elements of \(I\) times coefficients of \(S\), and let \(I_0\subset I\) be generated by the finitely many original \(I\)-elements appearing. Then \(I_0S=IS\). Faithful flatness gives injective quotient unit maps \(R/I_0\to S/I_0S\), so contraction gives \(I_0=I\). Thus every ideal of \(R\) is finite. The already proved local descent theorem then applies.

More generally, if \(R\to S\) is faithfully flat and \(S\) is a Noetherian regular ring, then \(R\) is Noetherian regular. The same ideal proof gives Noetherianity. For every \(\mathfrak p\), the nonzero fibre \(S\otimes_R\kappa(\mathfrak p)\) has a prime, yielding \(\mathfrak q\) contracting to \(\mathfrak p\). The actual localized map \(R_{\mathfrak p}\to S_{\mathfrak q}\) is flat and local, with regular target, so the original \(R_{\mathfrak p}\) is regular. This asserts no uniform dimension bound when none is supplied for \(S\).

\subsubsection{Auslander--Buchsbaum with every short case}\label{algebra-proof-023-auslanderbuchsbaum-with-every-short-case}

For 27228--27295 retain nonzero finite \(M\), its original finite projective dimension \(e\), and a minimal finite free resolution. Write \(t=\operatorname{depth}R\). The top-minor argument above gives \(t\ge e\) if \(e>0\), and \(t\ge0\) if \(e=0\).

In the depth-zero base case, if \(e=0\), \(M\) is a nonzero finite free module and its regular sequences are exactly those of \(R\), with the same nonzero final quotients; hence \(\operatorname{depth}M=t=0\). If \(e=1\), the exact sequence \(0\to F_1\to F_0\to M\to0\) and the depth lemma give \(0=\operatorname{depth}M\ge t-1\), while the top-minor bound gives \(t\ge1\); therefore \(t=1\). For \(e\ge2\), the original sequence decomposition and iterative depth inequality give \(\operatorname{depth}M\ge t-e\), so \(t\le e\); together with the other bound this is equality. At the short end \(e=2\), the decomposition consists of \(0\to F_2\to F_1\to K_0\to0\) and \(0\to K_0\to F_0\to M\to0\); no intermediate \(K_{-1}\) is introduced by the ellipsis. These cases resolve the two reports as a separate proof clarification.

If \(\operatorname{depth}M>0\), the same top-minor bound gives \(\operatorname{depth}R>0\) when \(e>0\), and the nonzero free case gives it when \(e=0\). Thus neither \(M\) nor \(R\) has \(\mathfrak m\) as an associated prime. Finite prime avoidance for the union of their associated primes gives the original \(x\in\mathfrak m\) regular on both \(R\) and \(M\).

The source announces this choice, justifies its existence, and then restates the conclusion. OCC-00535 treats this as a second unnecessary choice, but there is no distinct element or conflicting hypothesis. The exposition is mathematically sound; deleting this explanation is optional editing, not a required correction. The source wording is retained.

Multiplication by the actual \(x\) is injective on every free term. The short exact sequence of complexes \(0\to F_\bullet\xrightarrow{x}F_\bullet\to F_\bullet/xF_\bullet\to0\) has homology sequence with the only possibly positive homology of the quotient equal to \(\ker(x:M\to M)=0\). Its degree-zero homology is \(M/xM\). Thus it is exactly the displayed free resolution over \(R/xR\); tensoring was not assumed flat. This quotient is nonzero by Nakayama. All differential coefficients remain in the quotient maximal ideal, so the resolution is minimal; its last free term remains nonzero since \(R/xR\ne0\). Tensoring with the common residue field therefore detects projective dimension exactly \(e\).

The depth-drop lemma lowers both depths by one. For a module annihilated by \(x\), regular sequences in the original maximal ideal correspond exactly, by quotient and lifting, to regular sequences in the quotient maximal ideal: the actions and successive quotient modules coincide. This proves the equality of the two depth conventions used in the source, without changing rings silently. Induction over \(R/xR\) now gives \[
\operatorname{depth}R-1=e+\operatorname{depth}M-1,
\] which is the desired formula.

\subsubsection{The retained FAC quotient result and all Betti numbers}\label{algebra-proof-023-the-retained-fac-quotient-result-and-all-betti-numbers}

The captured current lemma lemma-projective-dimension-modulo-nonzerodivisor cites FAC Chapter III, §5, no. 76, Lemma 2, pp.~270--271. Its hypothesis is nonzero finite \(M\) of finite projective dimension and \(x\in\mathfrak m\) regular on \(M\), without assuming that \(x\) is regular on \(R\). Its exact sequence, depth drop and the just-checked Auslander--Buchsbaum formula prove \(\operatorname{pd}_R(M/xM)=\operatorname{pd}_R(M)+1\). All 49 lines remain unchanged. Reading this retained source/history block is not a fresh reading of FAC.

The exact same module quotient permits a stronger separate calculation, retaining \(R\) throughout. Let \(M\ne0\) be any finite module over a Noetherian local ring, without finite projective dimension, and let \(x\in\mathfrak m\) be regular on \(M\). Choose its original minimal free resolution \((F_i,d_i)\), finite or infinite. Write \(p:F_0\to M\) for its augmentation and put \(F_{-1}=0\), \(d_0:F_0\to0\) for the following unaugmented complex and cone. Define \[
C_i=F_i\oplus F_{i-1},\qquad
D_i(a,b)=(d_i a+x b,\ -d_{i-1}b)
\] for \(i\ge1\), with augmentation \(C_0=F_0\to M/xM\). Terms involving \(F_{-1}\) are zero. Direct computation gives \(D_{i-1}D_i=0\): the first component is \(d^2a+x\,db-x\,db=0\), and the second is \(d^2b=0\).

Here is exactness with the original signs. For a cycle \((a,b)\) in positive degree \(i\), \(db=0\). If \(i\ge2\), exactness of \(F\) gives \(b=dc\). If \(i=1\), applying the augmentation to \(da+xb=0\) gives \(x\bar b=0\) in \(M\); regularity gives \(\bar b=0\), so again \(b=dc\). Then \(d(a+xc)=0\), hence \(a+xc=du\), and \[
D_{i+1}(u,-c)=(du-xc,dc)=(a,b).
\] At degree zero, an element mapping to zero in \(M/xM\) has augmentation \(x\bar b\); lift \(\bar b\) to \(b\in F_0\), write \(a-xb=dc\), and obtain \(a=D_1(c,b)\). Thus this is a free resolution of the actual \(M/xM\).

All its matrix entries lie in \(\mathfrak m\), including the retained diagonal \(x\). It is minimal, so tensoring with \(\kappa\) gives zero differentials and yields the exact formula \[
\beta_i^R(M/xM)=\beta_i^R(M)+\beta_{i-1}^R(M)
\quad(i\ge0),\qquad \beta_{-1}^R(M)=0 .
\] If \(M\) has finite projective dimension \(e\), \(C_{e+1}=F_e\ne0\) and there are no higher terms; this proves the retained FAC formula directly. If its projective dimension is infinite, its minimal resolution has nonzero terms in arbitrarily large degrees and so does \(C\); the quotient still has infinite projective dimension. For every \(a\ge1\), \(x^a\) is regular on \(M\) by repeated injectivity, so all these formulas also hold for the original quotient \(M/x^aM\).

For an \(M\)-regular sequence \(x_1,\ldots,x_c\) in \(\mathfrak m\), apply exactly this construction successively to the nonzero modules \(M/(x_1,\ldots,x_j)M\). Each is nonzero by Nakayama. Induction on the displayed two-term rank formula gives \[
\beta_i^R\bigl(M/(x_1,\ldots,x_c)M\bigr)
=\sum_{j=0}^c\binom cj\beta_{i-j}^R(M),
\] where terms with negative index are zero. Equivalently the formal power series is multiplied by \((1+t)^c\), with no convergence assertion. Finite projective dimension increases by exactly \(c\); infinite projective dimension remains infinite. No regularity of this sequence on \(R\) is used. This extends the retained quotient result through explicit maps while keeping the translated source and the current added lemma intact.

\subsubsection{Bounded literature reading and a separate source overclaim}\label{algebra-proof-023-bounded-literature-reading-and-a-separate-source-overclaim}

The existing corpus lookup for Auslander--Buchsbaum returned two original-TeX records. The DG-ring item, PUBUNIT-41E1EB5390E0880F61FC0B92, remains unread. Neil Epstein and Hop D. Nguyen's Algebra retracts and Stanley-Reisner rings, arXiv:1301.3967, canonical PUBUNIT-E140D1A58DCB9288442C1AD5, was previously read only at 246--266 and 368--383. This turn read the complete additional native TeX range 390--456, with section routing checked. The exact author file is F:/user/Documents/arxiv\_latex/\_expanded\_by\_topic/motives\_motivic\_cohomology/6e7bbb412e444a77/1301.3967/1301.3967.tex, SHA-256 83131E01670F9EF5AB2E6480F668267364B9B1729CB612C9EAE469FFBF12FED1. It declares Latin-1 and was read accordingly; its source archive and files are unchanged.

The introduction and background conventions at 203--246 were also read, and a bounded file search checked for any blanket exclusion of zero modules; none was found. The convention at 242 requires Noetherian commutative rings with unit, which the example below satisfies. The displayed theorem at lines 427--430, labelled regularity\_ci\_descent and attributed there to Apassov, asserts descent from a regular target if there exists a finite target module of finite flat dimension over the source. As printed it does not require that module to be nonzero. The literal statement fails. Let \[
R=k[\epsilon]/(\epsilon^2),\quad S=k,\quad
\varphi(a+b\epsilon)=a,\quad M=0 .
\] This is a local map of Noetherian local rings; \(S\) is regular, and zero is a finite \(S\)-module of finite flat dimension over \(R\). But \(R\) has its unique prime/maximal ideal \((\epsilon)\), hence dimension zero, while \(\dim_k(\mathfrak m/\mathfrak m^2)=1\); it is not regular. This is an exact counterexample to the statement as printed. It shows a nonzero-module restriction is necessary; it does not certify that adding that restriction alone gives the exact theorem proved by Apassov.

No result depending on the unread Apassov proof, the unread Herzog Poincaré-series identity, or the later complete-intersection arguments is adopted. For the flat local maps treated above, the receiving module is explicitly \(M=S\), finite free of rank one over \(S\), nonzero, and flat over \(R\). The exact preceding proof establishes this special case independently. The printed external overclaim is recorded separately in the claim ledger, not converted into an upstream Stacks edit. This preserves authorship and the distinction between the consulted author's quotation of another theorem and a reading of its original proof.

\subsubsection{Propagation and remaining work}\label{algebra-proof-023-propagation-and-remaining-work}

The nonzero-depth guard receives ALGEBRA-RECON-589. The positive top-rank/depth inference receives ALGEBRA-RECON-576, with the original pivot maps of group 569 retained. The low-degree syzygy clarifications receive group 621. The flat regularity argument uses the exact faithful unit and finite-projectivity descent of group 501. The cone computation checks and extends the current FAC quotient lemma without changing it. Each receiving link is recorded in the ledger.

Lookahead through 27420 includes Homomorphisms and dimension and stops at the beginning of a new lemma; it is unadjudicated. The next review starts at 27305. No source or translation mutation, rendering or publication occurred. Wider combined-results synthesis remains after core completion.

\hypertarget{algebra_dimension_maps_notes_20260928}{}
\subsection{Homomorphisms, fibre dimensions and the dimension formula}\label{algebra-proof-024-homomorphisms-fibre-dimensions-and-the-dimension-formula}

Primary authority: \texttt{algebra.tex} 27305--27647 at commit \texttt{a04446e57ec1fbc252a871afcec7752fb2807b14}, SHA-256 \texttt{FA8BB92E58A4F78A2BD01B3B6A4A87DE0A0D279F5DD90641B574DD5FBFFFA4F3}. The entire current interval is byte-identical to this authority. No earlier canonical operation or added environment occurs in it.

These are source-linked editorial arguments. Both translation branches retain the source's statements and exposition; none of the expanded proofs below is a replacement translation. Proposed minimal corrections retain their old readings in the review ledger. Further consequences are not novelty claims.

\subsubsection{Chains, integral maps and the original fibre rings}\label{algebra-proof-024-chains-integral-maps-and-the-original-fibre-rings}

Source 27312--27388. Write the ring map as \(\varphi:R\to S\). For a finite strict chain \(\mathfrak p_0\subsetneq\cdots\subsetneq
\mathfrak p_e\), surjectivity on spectra supplies a prime over the first endpoint for going up, or the last endpoint for going down. Apply the respective lifting property successively. Every lifted inclusion is strict, since its contractions are distinct. Taking the supremum of these finite lengths proves the dimension inequality, also when the supremum is infinite. If \(R=0\), a unital map forces \(S=0\); both spectra are empty. If \(R\ne0\), surjectivity rules out an empty target spectrum.

If \(\mathfrak q\) is maximal and contracts to \(\mathfrak p\), every prime \(\mathfrak p'\supset\mathfrak p\) lifts by going up to \(\mathfrak q'\supset\mathfrak q\). Maximality gives \(\mathfrak q'=\mathfrak q\), hence \(\mathfrak p'=\mathfrak p\). This proves the assertion labelled ``Trivial'' in the source. OCC-12242 and OCC-00536 concern only its broken if/then punctuation.

For an integral map, incomparability makes contractions of any strict chain in \(S\) strict in \(R\), giving \(\dim S\leq\dim R\); going up gives the closed-point conclusion. More exactly, let \(K=\ker\varphi\). The original map factors through the injection \(R/K\to S\). Integrality and lying over give surjectivity \(\operatorname{Spec}S\to\operatorname{Spec}(R/K)\), and the two chain arguments prove \(\dim S=\dim(R/K)\). The prime correspondence is inverse image under \(R\to R/K\); it preserves every inclusion. The source inclusion \(R\subset S\) is the case \(K=0\). If \(S=0\), then \(K=R\) and both spectra in the equality are empty. No nonzero-ring convention has silently replaced this case.

Fix \(\mathfrak q\in\operatorname{Spec}S\) and \(\mathfrak p=\varphi^{-1}\mathfrak q\). Put \(A=R_{\mathfrak p}\), \(B=S_{\mathfrak q}\), \(\mathfrak m=\mathfrak pA\), and \(\mathfrak n=\mathfrak qB\). The map is explicitly \(r/u\mapsto\varphi(r)/\varphi(u)\); its denominator is a unit because \(u\notin\mathfrak p\). Put \(C=S\otimes_R\kappa(\mathfrak p)\). First \(C=(R\setminus\mathfrak p)^{-1}(S/\mathfrak pS)\) by \(s\otimes(\bar r/\bar u)\mapsto\overline{s\varphi(r)}/\bar u\). Its inverse takes \(\bar s/\bar u\) to \(s\otimes1/\bar u\). The prime induced by \(\mathfrak q\) is \(\tilde{\mathfrak q}=(R\setminus\mathfrak p)^{-1}
(\mathfrak q/\mathfrak pS)\). There are inverse maps \[
 B/\mathfrak pB\longleftrightarrow C_{\tilde{\mathfrak q}},
 \qquad
 [s/t]\longmapsto(s\otimes1)/(t\otimes1),\quad t\notin\mathfrak q.
\] To describe the reverse map, \(s\otimes(\bar r/\bar u)\) maps to \([s\varphi(r)/\varphi(u)]\). An element \(\bar s/\bar u\) outside \(\tilde{\mathfrak q}\) has \(s\notin\mathfrak q\), so its image is invertible. The localization universal property therefore supplies the reverse map; both composites fix the original numerators and every inverted denominator. The comparison with \((S/\mathfrak pS)_{\mathfrak q/\mathfrak pS}\) uses the same \([s/t]\). This proves all displayed fibre identifications without discarding their prime labels.

\subsubsection{The parameter bound and the zero-length going-down case}\label{algebra-proof-024-the-parameter-bound-and-the-zero-length-going-down-case}

Source 27390--27456. For the original Noetherian local \(A,B\), put \(d=\dim A\), \(e=\dim(B/\mathfrak mB)\). Choose parameters \(x_1,\ldots,x_d\) in \(A\) and lifted fibre parameters \(y_1,\ldots,y_e\) in \(B\). The source chooses representatives in \(R,S\); this is possible by clearing each localization denominator, which is a unit in the respective local ring. Multiplying a parameter by that unit does not change its generated ideal.

Here is the claimed nilpotence explicitly. There are positive integers \(a,b\) such that \[
 \mathfrak m^a\subset(x_1,\ldots,x_d)A,\qquad
 \mathfrak n^b\subset\mathfrak mB+(y_1,\ldots,y_e)B.
\] The first follows from the ideal of definition, the second from the fibre ideal of definition. With \(J=(\varphi(x_1),\ldots,\varphi(x_d),
y_1,\ldots,y_e)B\), the image of \(\mathfrak mB\) in \(B/J\) has \(a\)-th power zero. The second inclusion therefore gives \(\mathfrak n^{ab}\subset J\). Thus \(J\) is an ideal of definition generated by \(d+e\) elements, proving \(\dim B\leq d+e\). Empty parameter sequences, including \(d=0\) or \(e=0\), obey exactly the same ideal calculation.

For the reverse inequality assume going down. In the source's indexing choose \[
 \mathfrak pS\subset\mathfrak q_0\subsetneq\cdots
 \subsetneq\mathfrak q_d=\mathfrak q,\qquad
 \mathfrak p_0\subsetneq\cdots\subsetneq\mathfrak p_e=\mathfrak p,
\] where now \(d=\dim(B/\mathfrak pB)\) and \(e=\dim A\), exactly as at 27442--27447. Every \(\mathfrak q_i\) contracts to \(\mathfrak p\): it contains \(\mathfrak pS\) and is contained in \(\mathfrak q\). If \(e=0\), the existing fibre chain already has length \(d+e\). If \(e\geq1\), use going down \(e\) times to obtain \(\mathfrak q_{-j}\) over \(\mathfrak p_{e-j}\), \(1\leq j\leq e\). Each new inclusion is strict by contraction. The concatenation has length \(d+e\). This proves equality and resolves OCC-12243 in a separate proof clarification, without rewriting the source's exposition.

There are two precise extensions of these calculations. For arbitrary rings and a fixed prime \(\mathfrak q\), going down alone gives the lower bound \[
 \dim S_{\mathfrak q}\ \geq\
 \dim R_{\mathfrak p}+\dim(S_{\mathfrak q}/\mathfrak pS_{\mathfrak q})
\] in the extended nonnegative integers: concatenate arbitrary finite chains as above; if either dimension is infinite, their finite lengths are unbounded. All these local spectra are nonempty, so no undefined sum with \(-\infty\) occurs. No Noetherian hypothesis enters this bound.

For the upper bound, \(A\) need not be Noetherian. It suffices that \(B\) is Noetherian local and that there are \(c\) actual elements \(x_1,\ldots,x_c\in\mathfrak m\) with \(\sqrt{(x_1,\ldots,x_c)A}=\mathfrak m\). Then \(\mathfrak mB\subset\sqrt{(x_1,\ldots,x_c)B}\): for every \(z\in
\mathfrak m\), an actual positive power of \(z\) lies in the ideal in \(A\). Choose \(e=\dim B/\mathfrak mB\) fibre parameters. Every prime of \(B\) containing their lifts and the \(x_i\) contains \(\mathfrak mB\), and hence is \(\mathfrak n\). Thus the lifted \(c+e\) elements generate an ideal of definition and \(\dim B\leq c+e\). This argument does not assume a uniform power of \(\mathfrak m\) in non-Noetherian \(A\). If \(c=\dim A<\infty\) and going down holds, the upper and lower bounds give equality.

\subsubsection{Regular ascent under going down and the resulting flatness}\label{algebra-proof-024-regular-ascent-under-going-down-and-the-resulting-flatness}

Source 27458--27480. Suppose \(R\to S\) is a local map of Noetherian local rings, \(R\) is regular, and \(S/\mathfrak m_RS\) is regular. Put \(d=\dim R\), \(e=\dim(S/\mathfrak m_RS)\). Choose the original \(d\) generators \(x_i\) of \(\mathfrak m_R\) and lift the original \(e\) generators \(y_j\) of the fibre maximal ideal. Then \[
 \mathfrak m_S=(\varphi(x_1),\ldots,\varphi(x_d),y_1,\ldots,y_e)S.
\] Indeed, reducing an element of \(\mathfrak m_S\) modulo \(\mathfrak m_RS\) expresses it in the \(\bar y_j\); subtracting the lifts leaves an element of \(\mathfrak m_RS\), which is generated by the displayed \(\varphi(x_i)\).

The following three conditions are equivalent under these hypotheses: the map is flat; it has going down; and \(\dim S=d+e\). Flatness gives going down and the preceding chain argument gives the dimension equality. Conversely, if the equality holds, the displayed generating set shows that the embedding dimension of \(S\) is at most \(\dim S\). The reverse inequality follows from the ideal-of-definition bound. Thus \(S\) is regular and this generating set is a regular system of parameters. Regular local rings are Cohen--Macaulay, so the ordered sequence \(\varphi(x_1),\ldots,\varphi(x_d),y_1,\ldots,y_e\) is regular. In particular the initial \(d\) elements are regular on \(S\).

For completeness the flatness implication uses the original rings and their successive quotients. Put \(R_i=R/(x_1,\ldots,x_i)\) and \(S_i=S/(\varphi(x_1),\ldots,\varphi(x_i))S\). Then \(R_d=\kappa(R)\), so \(S_d\) is \(R_d\)-flat. Working backwards from \(i=d\) to \(i=1\), multiplication by \(x_i\) is injective on \(R_{i-1}\), and multiplication by \(\varphi(x_i)\) is injective on \(S_{i-1}\). The exact two-term free resolution of \(R_i\) over \(R_{i-1}\) therefore gives \(\operatorname{Tor}_1^{R_{i-1}}(S_{i-1},R_i)=0\). The local criterion at authority 23750--23805, with \(M=S_{i-1}\) finite over itself, lifts flatness of \(S_i\) to \(S_{i-1}\). This reaches \(R\to S\). Every ideal used is proper because the maps are local. If \(d=0\), \(R\) is already a field and the backwards induction is empty. This also proves the result for \(e=0\).

Thus the source's flatness hypothesis can be weakened to going down, and flatness is recovered as a conclusion. This agrees with the later regular-parameter flatness proof at 32693--32715, read in full as a bounded receiving comparison. Reading that later passage does not adjudicate its received reports or advance the source-order cursor past it.

\subsubsection{Cohen--Macaulay ascent with an arbitrary Cohen--Macaulay fibre}\label{algebra-proof-024-cohenmacaulay-ascent-with-an-arbitrary-cohenmacaulay-fibre}

Source 27482--27503. The received copula correction OCC-12244 is accepted. The original finite-flat case is justified as follows. A flat local map is faithfully flat, hence injective. Finiteness implies integrality, so the integral dimension equality gives \(\dim S=\dim R=d\). A length-\(d\) regular sequence in \(R\) stays regular on \(S\): flatness preserves each multiplication injection after quotienting, and the final quotient is nonzero by faithful flatness (or by locality and Nakayama). In the other source case going down gives \(\dim S\geq d\), while the stated hypothesis gives the reverse inequality. In both cases \(d\leq\operatorname{depth}S\leq\dim S=d\), proving the assertion.

More generally, let \(R\to S\) be flat local between Noetherian local rings. If \(R\) and \(F=S/\mathfrak m_RS\) are Cohen--Macaulay, then \(S\) is Cohen--Macaulay and \(\dim S=\dim R+\dim F\), without assuming \(\dim F=0\). Here is the complete regular-sequence argument. Put \(d=\dim R\), \(e=\dim F\). Choose a system of parameters \(x_1,\ldots,x_d\) of \(R\), which is regular since \(R\) is Cohen--Macaulay. Flatness makes its image an \(S\)-regular sequence. Put \(A=R/(x_1,\ldots,x_d)\) and \(Q=S/(x_1,\ldots,x_d)S\). Base change gives \(Q\) flat over \(A\), and its residue fibre is the original \(F\). Choose fibre parameters \(\bar y_1,\ldots,\bar y_e\), which form an \(F\)-regular sequence, and take lifts \(y_j\) in \(Q\).

For \(Q_j=Q/(y_1,\ldots,y_j)Q\), begin with the \(A\)-flat \(Q_0=Q\). Multiplication by \(y_j\) on \(Q_{j-1}\) is injective after tensoring with \(\kappa(A)\), by the original regular sequence on \(F\). The proved injection-modulo-the-maximal-ideal lemma applies with the Noetherian local auxiliary ring \(Q\) and finite modules \(Q_{j-1}\). It proves that multiplication is injective and \(Q_j\) is \(A\)-flat. Induction proves the entire lifted sequence regular. Its final quotient surjects onto the nonzero \(F/(\bar y_1,\ldots,\bar y_e)\), so the nonzero condition is retained at every step. Thus \(x_1,\ldots,x_d,y_1,\ldots,y_e\) is \(S\)-regular. The dimension equality already proved gives \(\dim S=d+e\); the depth inequality gives the opposite bound to the constructed sequence, proving that \(S\) is Cohen--Macaulay.

This receives group 559's exact module and quotient maps, in ALGEBRA\_FLATNESS\_CRITERIA\_NOTES\_20260928.md, section ``Removing the unused base Noetherian hypothesis''. The auxiliary ring remains Noetherian and the module finite as required. The construction works with either empty sequence. The source lemma is the case of zero-dimensional fibre, which is automatically Cohen--Macaulay as a nonzero zero-dimensional Noetherian local ring. No converse about arbitrary fibres has been assumed.

\subsubsection{The dimension formula with the original tower and height-one prime}\label{algebra-proof-024-the-dimension-formula-with-the-original-tower-and-height-one-prime}

Source 27519--27596. Let \(R\subset S\) be the stated finite-type extension of domains with \(R\) Noetherian, and retain \(\mathfrak p,\mathfrak q\). All heights here are finite, since the corresponding local rings are Noetherian. For \(R\subset S'\subset S\), where \(S'\) is generated by an initial part of the fixed generator list, put \(\mathfrak q'=\mathfrak q\cap S'\). Write \[
 g_1=\operatorname{trdeg}_{\operatorname{Frac}R}\operatorname{Frac}S',
 \quad g_2=\operatorname{trdeg}_{\operatorname{Frac}S'}\operatorname{Frac}S,
 \quad t_1=\operatorname{trdeg}_{\kappa(\mathfrak p)}\kappa(\mathfrak q'),
 \quad t_2=\operatorname{trdeg}_{\kappa(\mathfrak q')}\kappa(\mathfrak q).
\] Adding the two inequalities gives \[
 \operatorname{ht}\mathfrak q\leq
 \operatorname{ht}\mathfrak q'+g_2-t_2\leq
 \operatorname{ht}\mathfrak p+g_1+g_2-t_1-t_2.
\] Field-tower additivity identifies both sums with those in the source formula. If \(R\) is universally catenary, each finite-type \(S'\) is universally catenary: every finite-type \(S'\)-algebra is finite type over \(R\). Consequently both inductive equalities are justified.

For the one-generator polynomial case, put \(T=R[x]\) and \(\mathfrak r=\mathfrak pT\). The original fibre is \(\kappa(\mathfrak p)[x]\). A prime over \(\mathfrak p\) corresponds either to its zero prime, exactly when it equals \(\mathfrak r\), or to an ideal generated by a nonconstant irreducible polynomial. In the first case its residue field is \(\kappa(\mathfrak p)(x)\), of transcendence degree one, and the local fibre has dimension zero. In the second case its residue field is a finite algebraic extension and the local fibre is a one-dimensional discrete valuation ring: ideals in the PID \(\kappa(\mathfrak p)[x]\) are principal, and the primes below that nonzero prime are just zero and itself. Flatness and the preceding dimension equality prove the source formula with equality in both cases.

For the algebraic case \(S=R[x]/\mathfrak n\), \(\mathfrak n\) is a nonzero prime with \(\mathfrak n\cap R=0\). Let \(F=\operatorname{Frac}R\). Localization at \(R\setminus\{0\}\) gives a nonzero prime \(\mathfrak nF[x]\): a nonzero polynomial of \(\mathfrak n\) stays nonzero because \(R[x]\) is a domain. The polynomial relation also shows that the image of \(x\) is algebraic over \(F\), so the generic transcendence degree is zero. Every prime below \(\mathfrak n\) is disjoint from \(R\setminus\{0\}\). The localization bijection therefore preserves the entire interval \([0,\mathfrak n]\), whose length is one since \(F[x]\) is a PID. In particular \(\operatorname{ht}\mathfrak n=1\), without assuming catenarity.

Let \(\mathfrak q'\) be the preimage of \(\mathfrak q\) and put \(C=R[x]_{\mathfrak q'}\). Quotient and localization give the exact isomorphism \[
 C/\mathfrak nC\longrightarrow S_{\mathfrak q},\qquad
 [(f/g)]\longmapsto \bar f/\bar g
\] with inverse given by any polynomial representatives; two choices differ by the localized kernel. The residue fields at \(\mathfrak q'\) and \(\mathfrak q\) are thereby the same. Every chain in \(C/\mathfrak nC\) extends by \(0\subsetneq\mathfrak nC\), so \(\dim(C/\mathfrak nC)\leq\dim C-1\). The already proved polynomial formula supplies \(\dim C=\operatorname{ht}\mathfrak p+1-
\operatorname{trdeg}_{\kappa(\mathfrak p)}\kappa(\mathfrak q)\). This proves the inequality. If \(C\) is catenary, every saturated chain from \(\mathfrak nC\) to its maximal ideal, extended by this same height-one interval, has length \(\dim C\). Hence the quotient dimension is exactly \(\dim C-1\). Universal catenarity of \(R\) implies the required catenarity of \(C\), proving the source equality.

OCC-00537 is the missing ``by'' in the prime declaration and is accepted. OCC-00538 overstates a stylistic preference: ``in case R is \ldots{}'' is a valid conditional construction. OCC-00539 likewise proposes a smoother order of modifiers but the existing comparison is intelligible and mathematically correct: append \(0\subsetneq\mathfrak n\) to the quotient chain. Both optional rewrites remain unapplied; they are not counted as defects.

\subsubsection{Essential finite type, local catenarity and the full codimension-one fibre}\label{algebra-proof-024-essential-finite-type-local-catenarity-and-the-full-codimension-one-fibre}

The dimension formula above extends as follows. Suppose \(R\) is a Noetherian domain and \(S=T^{-1}B\), where \(R\subset B\) is a finite-type domain and \(S\ne0\). For \(\mathfrak q\in\operatorname{Spec}S\), let \(\mathfrak q_0\) be its contraction to \(B\). Every prime below \(\mathfrak q_0\) is disjoint from \(T\). Thus localization identifies the entire interval below \(\mathfrak q_0\) with the interval below \(\mathfrak q\), preserving height. It also gives the explicit isomorphism \(B_{\mathfrak q_0}\simeq S_{\mathfrak q}\) by the unchanged fractions. Fraction fields and residue fields are unchanged. Applying the preceding formula to \(R\subset B\) proves it for \(R\subset S\).

For equality it suffices that \(R_{\mathfrak p}\) is universally catenary. Indeed apply the finite-type formula after localization at \(R\setminus\mathfrak p\). Every prime below \(\mathfrak p\) and below \(\mathfrak q_0\) survives this localization, so both heights, both fraction fields and both residue fields in the formula retain their original values. The localized finite-type algebra is over the stipulated universally catenary ring \(R_{\mathfrak p}\). This proves equality without requiring universal catenarity away from \(\mathfrak p\). This is a sufficient local hypothesis; no necessity claim is made.

Now retain the source's generically finite extension \(A\subset B\), and allow \(B\) essentially of finite type. Let \(\operatorname{ht}\mathfrak p=1\). For each prime \(\mathfrak q\) above it, the dimension formula and generic transcendence degree zero give \[
 1\leq\operatorname{ht}\mathfrak q
 \leq1-\operatorname{trdeg}_{\kappa(\mathfrak p)}\kappa(\mathfrak q).
\] The first inequality uses a nonzero \(a\in\mathfrak p\), whose image is nonzero in the domain \(B\) and belongs to \(\mathfrak q\); thus \(0\subsetneq\mathfrak q\). Hence \(\operatorname{ht}\mathfrak q=1\) and the residue transcendence degree is zero. Put \(D=B\otimes_A\kappa(\mathfrak p)\). It is an essentially finite-type Noetherian algebra over the field \(\kappa(\mathfrak p)\). Every residue field at a prime of \(D\) is a finite extension of that field: represent \(B\) by a finite-type algebra before localization, retain the same residue field, and use that an algebraic finitely generated field extension is finite. Each corresponding prime of that finite-type field algebra is maximal by the field-algebra lemma, so each prime of \(D\) is maximal. There are finitely many of them by the corrected spectrum argument in the next section.

In fact the entire fibre algebra \(D\), including its nilpotents, is finite dimensional over \(\kappa(\mathfrak p)\). If \(D=0\), this is immediate and there is no fibre point. Otherwise its nilradical \(J\) is finitely generated, say by nilpotent elements \(z_1,\ldots,z_t\), with \(z_i^{a_i}=0\). Every monomial of total degree \(N=1+\sum_i(a_i-1)\) contains such a zero factor, so \(J^N=0\). For \(J=0\) use \(N=1\). The finitely many distinct maximal ideals are pairwise comaximal, with intersection \(J\); the Chinese remainder maps give \(D/J\simeq\prod_{\mathfrak q}\kappa(\mathfrak q)\), a finite-dimensional algebra over \(\kappa(\mathfrak p)\). Each \(J^i/J^{i+1}\) is a finite \(D/J\)-module (its generators are the degree-\(i\) monomials in the original \(z_j\)), so it too is finite dimensional. The exact sequences of the finite filtration by the \(J^i\) prove the assertion for \(D\). No reduced-fibre assumption is used. This strengthens the source's finite-point conclusion while retaining its original finite-type case and its possibly empty fibre.

\subsubsection{The missing topological hypothesis and its exact repair}\label{algebra-proof-024-the-missing-topological-hypothesis-and-its-exact-repair}

Source 27624--27637. The statement ``A Noetherian topological space consisting of closed points is finite'' is false for general spaces. Let \(X\) be an infinite set with the cofinite topology. Closed subsets are \(X\) and finite subsets. Every descending chain of closed subsets stabilizes: once it leaves \(X\), it is a chain inside one finite set. Thus \(X\) is Noetherian, and all its points are closed, but \(X\) is infinite. Moreover any two nonempty opens meet, so \(X\) is irreducible. It has no generic point, since each point has singleton closure. This identifies the exact missing structure.

The corrected assertion is true for a sober Noetherian space with all points closed. A Noetherian space has finitely many irreducible components: otherwise choose a minimal closed counterexample by the descending-chain condition; it is reducible, and its two proper closed pieces each have finite decompositions, a contradiction. In a sober space each component is the closure of its generic point. If that point is closed, the component is a singleton. Hence the space is finite, including the empty case.

For \(X=\operatorname{Spec}D\), sobriety follows directly. An irreducible closed subset \(V(I)\) has prime radical \(P=\sqrt I\): if \(ab\in P\), then \(V(I)\subset V(a)\cup V(b)\), and irreducibility puts it in one of these sets, so \(a\in P\) or \(b\in P\). Its generic point is precisely \(P\), since \(\overline{\{P\}}=V(P)=V(I)\); uniqueness follows because distinct primes have distinct closures. This proves the hypothesis for the original fibre, so the source theorem remains valid. The minimal proposal adds ``sober'' to its topological assertion and keeps the cited Topology lemma, which states finiteness of irreducible components at authority 1252--1301, not finiteness of arbitrary Noetherian \(T_1\) spaces.

\subsubsection{The original literature counterexample and propagation boundaries}\label{algebra-proof-024-the-original-literature-counterexample-and-propagation-boundaries}

The indexed original author source is Katharina Heinrich, ``Some remarks on biequidimensionality of topological spaces and Noetherian schemes'', arXiv:1403.5814, canonical ID \texttt{PUBUNIT-EB76FB876CD8CEEBBA2086CE}. Exact file: \texttt{F:/user/Documents/arxiv\_latex/\_expanded\_by\_topic/krull\_dimension\_theory/6637bf65ee5385b5/1403.5814/1403.5814.tex}, SHA-256 \texttt{CB4305F104ABA334E06B0556137AA89484EE797E31D3558B5D6AEB244841A911}. Actual reading: lines 1--343, 502--667 and 701--741. The construction section 344--501 and the later example 668--700 remain unread. No claim is made to have independently checked every result in the read ranges or any EGA passage quoted by this paper.

Here is an independent verification of the exact example \texttt{ex:glue2}, lines 531--550, used to test propagation. Keep the author's ring and both prime ideals: \[
 B=k[u,v,w,x,y,z]/(uy,uz,vy,vz,wy,wz),\qquad
 M_1=(u,v,w,y,z),\quad M_2=(u,v,w,x,y-1,z-1),
\] and put \(T=B\setminus(M_1\cup M_2)\), \(A=T^{-1}B\). The complement is multiplicatively closed because \(M_1,M_2\) are prime. A prime disjoint from \(T\) is contained in \(M_1\cup M_2\), hence in one of them: if an ideal contained in their union were contained in neither, choose one element outside each and use their sum for a contradiction. The two primes are incomparable. Thus these are exactly the two maximal ideals after localization.

In the original polynomial ring, \((u,v,w)\cap(y,z)=(u,v,w)(y,z)\): a monomial belongs to the intersection exactly when it has a factor from each of the two disjoint variable sets. Their product is exactly the six displayed generators. Consequently \(B\) is reduced with precisely two minimal primes \(P=(u,v,w)\), \(Q=(y,z)\). Every prime contains one of them: in a prime containing their product, failure to contain either would give a product of two elements outside the prime. The component rings before localization are \(B/P=k[x,y,z]\) and \(B/Q=k[u,v,w,x]\).

The \(P\)-component has maximal primes \((y,z)\) and \((x,y-1,z-1)\), of heights two and three. The \(Q\)-component has only the maximal prime \((u,v,w)\), of height three: \(M_2\) does not contain \(Q\). These height counts can be checked with the exact coordinate ideals. At \((y,z)\) the coefficient \(x\) passes to \(k(x)\), leaving the two-variable coordinate maximal ideal; at \((u,v,w)\) the coefficient \(x\) likewise passes to \(k(x)\), leaving three variables. At \((x,y-1,z-1)\) the three displayed generators give the upper bound three, and their successive prime ideals give the lower bound three. Thus both irreducible components of \(\operatorname{Spec}A\) have dimension three, and both maximal ideals have height three; at \(M_1\) the two component heights are two and three, while at \(M_2\) only the height-three \(P\)-component remains. The equality of dimension with the maximum of component dimensions follows by putting each prime chain inside a component containing its first prime.

These component polynomial rings are catenary. One can check this without using a general global dimension formula: their local rings are Cohen--Macaulay, as successive polynomial extension and localization of a field preserve the already proved Cohen--Macaulay property. The earlier CM chain theorem proves each such local ring catenary. Every fixed interval of primes lies in the localization at its upper prime and retains exactly its intermediate primes, so the whole component is catenary. The same holds after localization. The closed-cover criterion proved in group 594 therefore makes \(A\) catenary.

Now keep the source prime \(\mathfrak p=(u,v,w,y)A\). Before localization it is prime with quotient \(k[x,z]\), it is disjoint from \(T\), and it lies under \(M_1\) alone. It contains \(P\) and does not contain \(Q\), since \(z\notin\mathfrak p\). Thus every prime below it contains \(P\); on that component it corresponds to the height-one prime \((y)\) of \(k[x,y,z]\). The entire lower interval survives localization, so \(\operatorname{ht}_A\mathfrak p=1\). The primes of \(A/\mathfrak p\) correspond to the primes of \(k[x,z]\) contained in \((z)\), giving dimension one, with its strict chain \(0\subsetneq(z)\). Therefore, retaining all original objects, \[
 \operatorname{ht}_A\mathfrak p+\dim(A/\mathfrak p)=1+1=2<3=\dim A.
\] This verifies the cited counterexample, rather than treating the abstract as proof.

It checks the limitations of earlier groups 590 and 595: they retained a local ring, and group 590 also retained a nonzero Cohen--Macaulay module and its actual support. Equidimensionality, equal heights of maximal ideals and catenarity of a nonlocal ring do not replace those hypotheses. No contradiction with the source's relative dimension formula for an extension of domains follows from this example. Both original statements remain valid in their checked scopes; no earlier formula is silently broadened. This is use of an existing human result, not a new result or a claim that the EGA counterexample construction has been exhaustively read.

The other new routing hit, Epstein--Shapiro arXiv:1501.03411, remains unread. The topic route preserves its source identity without using its theorems.

\hypertarget{algebra_affine_dimension_notes_20260928}{}
\subsection{Affine dimension, triangular parameters and equidimensional pieces}\label{algebra-proof-025-affine-dimension-triangular-parameters-and-equidimensional-pieces}

Authority: \texttt{algebra.tex} 27648--27919, commit \texttt{a04446e57ec1fbc252a871afcec7752fb2807b14}, SHA-256 \texttt{FA8BB92E58A4F78A2BD01B3B6A4A87DE0A0D279F5DD90641B574DD5FBFFFA4F3}. The entire current interval equals the authority after MC-STK-ERR-1601, together with the captured eighteen-line FAC reference and history at \texttt{proposition-finite-gl-dim-polynomial-ring}. Those additions are retained.

Complete arguments and consequences below are editorial evidence linked to the source, not replacement translations. Minimal wording changes retain the original readings in the ledger. No novelty claim is made.

\subsubsection{The original triangular polynomials and the corrected induction}\label{algebra-proof-025-the-original-triangular-polynomials-and-the-corrected-induction}

Source 27659--27720. Retain \(A=k[x_1,\ldots,x_n]\), its maximal ideal \(\mathfrak m\), the finite extension \(\kappa=A/\mathfrak m\), the original \(\alpha_i\), and the fields \(\kappa_i=k(\alpha_1,\ldots,\alpha_i)\) with \(\kappa_0=k\). Each \(\alpha_i\) is algebraic over \(k\). The algebra generated by finitely many such elements is a finite-dimensional domain over a field and is a field: multiplication by a nonzero element is injective on its finite vector space, hence surjective, giving an inverse. Therefore \(\kappa_i=k[\alpha_1,\ldots,\alpha_i]\), as the source requires.

Let \(P_i(T)\in\kappa_{i-1}[T]\) be the original monic minimal polynomial. Lift each of its coefficients through \(k[x_1,\ldots,x_{i-1}]\to\kappa_{i-1}\), retaining the leading coefficient one and the degree \(d_i=\deg P_i\). This gives the specified \(Q_i(T)\); put \(f_i=Q_i(x_i)\). The evaluation map kills each \(f_i\) and is onto \(\kappa_i\). Suppose the induced \(\psi_i\) is an isomorphism. It identifies \[
 k[x_1,\ldots,x_{i+1}]/(f_1,\ldots,f_{i+1})
 \ \simeq\ \kappa_i[x_{i+1}]/(P_{i+1}(x_{i+1}))
 \ \simeq\ \kappa_{i+1}.
\] The first map sends each original variable to its quotient class and uses \(\psi_i\) on its coefficients; its inverse lifts those coefficient classes. The second map evaluates at \(\alpha_{i+1}\), with kernel exactly the minimal polynomial ideal by Euclidean division. It follows that the map proved injective at the end of the step is \(\psi_{i+1}\). MC-STK-ERR-1601 already makes that correction. Both OCC-12261 and OCC-00542 retain it, without a new operation. The two earlier wording reports concern only the declarations of \(\alpha_i\) and \(\kappa_i\).

The same coefficient maps, leaving the remaining variables fixed, give \[
 A/(f_1,\ldots,f_i)\simeq
 \kappa_i[x_{i+1},\ldots,x_n].
\] Thus each prefix ideal is prime. Each next inclusion is strict, since the monic positive-degree \(P_{i+1}\) is nonzero in that polynomial domain and is not a unit. For \(i=n\) the quotient is exactly \(\kappa\), so the prefix ideal is \(\mathfrak m\). The resulting chain from zero to \(\mathfrak m\) has length \(n\). Localizing keeps all its primes and strict inclusions. The maximal ideal of \(A_{\mathfrak m}\) has \(n\) generators, giving the upper dimension bound \(n\); the chain gives the lower bound. Hence \(A_{\mathfrak m}\) is regular of dimension \(n\). For \(n=0\), \(A=k\), \(\mathfrak m=0\), \(\kappa=k\), and the empty chain has length zero. Every construction above then has its empty interpretation.

\subsubsection{Explicit freeness over the original parameter algebra}\label{algebra-proof-025-explicit-freeness-over-the-original-parameter-algebra}

The triangular construction gives more than the source's dimension statement. Put \(D=\prod_{i=1}^n d_i=[\kappa:k]\), with empty product one. Use new independent variables \(t_i\) solely to prove the actual comparison: \[
 C=k[t_1,\ldots,t_n,X_1,\ldots,X_n]/
       (f_i(X_1,\ldots,X_i)-t_i:1\leq i\leq n).
\] Eliminating the \(t_i\) defines an isomorphism \(C\to A\) by \(X_i\mapsto x_i,\ t_i\mapsto f_i(x_1,\ldots,x_i)\). The inverse sends \(x_i\) to the class of \(X_i\); each composite fixes every named generator using the defining equations.

Regard the same quotient successively as an algebra over \(k[t_1,\ldots,t_n]\). Its \(i\)-th defining equation is monic of degree \(d_i\) in \(X_i\), with coefficients in the preceding algebra. For any ring \(E\), division by a monic degree-\(d\) polynomial in \(E[X]\) gives a unique remainder of degree less than \(d\): subtract its leading monomial successively for existence; a nonzero multiple has the same nonzero leading coefficient as that of its multiplier, so cannot have degree less than \(d\), proving uniqueness. This proof remains valid when \(E\) has zero divisors. Induction consequently makes \(C\) a free \(k[t_1,\ldots,t_n]\)-module with the exact basis \[
 X_1^{b_1}\cdots X_n^{b_n},\qquad 0\leq b_i<d_i.
\] The basis contains \(1\), and therefore the coefficient map is injective: a coefficient annihilating \(1\) has all its basis coordinates zero. Under the isomorphism with \(A\), this proves that the original \(f_i\) are algebraically independent and that \(k[f_1,\ldots,f_n]\subset A\) is finite free of rank \(D\), with basis \(x_1^{b_1}\cdots x_n^{b_n}\). In particular this is an explicit finite faithfully flat map, not a coordinate replacement for \(A\).

For positive integers \(a_1,\ldots,a_n\), base change along \(k[t]\to k[t]/(t_1^{a_1},\ldots,t_n^{a_n})\) gives the actual quotient \(A/(f_1^{a_1},\ldots,f_n^{a_n})\), with \(k\)-basis \[
 f_1^{c_1}\cdots f_n^{c_n}x_1^{b_1}\cdots x_n^{b_n},
 \qquad 0\leq c_i<a_i,\quad 0\leq b_i<d_i.
\] Its \(k\)-dimension is \(D\prod_i a_i\). The radical of the quotient ideal is \(\mathfrak m=(f_1,\ldots,f_n)\): a prime containing all the powers contains all the \(f_i\). Thus the quotient is a finite-dimensional local \(k\)-algebra with residue field \(\kappa\). Its composition factors are all \(\kappa\): each simple module over a commutative local ring is its residue field, and finite \(k\)-dimension ensures a composition series. Additivity of \(k\)-dimension then gives its module length \(\prod_i a_i\).

For each prefix of the powered sequence, the same free base-change description is free over \(k[t]/(t_1^{a_1},\ldots,t_{i-1}^{a_{i-1}})\). Multiplication by \(t_i^{a_i}\) on that coefficient ring is injective, as is seen from its monomial basis in the still unrestricted variable \(t_i\); it is therefore injective on the free module. The final quotient has the nonempty displayed basis. This proves directly that every positive-power sequence \(f_1^{a_1},\ldots,f_n^{a_n}\) is \(A\)-regular and remains so at \(\mathfrak m\). No power or coefficient has been absorbed into a new parameter.

There is also the exact graded comparison \[
 \kappa[T_1,\ldots,T_n]\longrightarrow
 \bigoplus_{r\geq0}\mathfrak m^r/\mathfrak m^{r+1},
 \qquad T_i\longmapsto f_i\bmod\mathfrak m^2.
\] Here a coefficient \(\bar a\in\kappa=A/\mathfrak m\) multiplies a degree-\(r\) class by any lift \(a\in A\); a different lift differs by \(\mathfrak m\), so the class is well defined. Writing \(A=\bigoplus_b k[t]x^b\) by the proved freeness yields \(\mathfrak m^r=\bigoplus_b(t)^r x^b\). For each \(r\), both sides of the comparison have \(k\)-basis indexed by \(b\) and the exponent vectors \(c\) with \(\sum c_i=r\). The displayed map takes one such basis to the other, proving bijectivity in every degree, and it respects multiplication by its definition. No ring section \(\kappa\to A\) is assumed, which matters for inseparable residue extensions.

If \(n\geq1\), the number of degree-\(r\) exponent vectors is \(\binom{n+r-1}{n-1}\). Consequently, for every \(r\geq1\), \[
 \dim_k(A/\mathfrak m^r)=D\binom{n+r-1}{n},
 \qquad
 \operatorname{length}_A(A/\mathfrak m^r)=\binom{n+r-1}{n}.
\] The sum of the degree counts for \(0\leq j<r\) gives the binomial coefficient, by partitioning degree-at-most-\(r-1\) exponent vectors according to their last coordinate (or by the elementary Pascal recurrence). For \(n=0\), \(\mathfrak m=0\), the graded algebra is \(k\) in degree zero and all positive powers give \(A/\mathfrak m^r=k\), of dimension and length one. This is treated directly, without factorials with negative indices.

\subsubsection{Global dimension and the exact prime-interval lengths}\label{algebra-proof-025-global-dimension-and-the-exact-prime-interval-lengths}

Source 27722--27779. Since \(A\) is Noetherian and all its maximal localizations are regular of dimension \(n\), the previously checked finite-global-dimension criterion gives \(\operatorname{gldim}A=n\). Every prime lies below a maximal ideal; localization of the corresponding regular local ring is regular. This gives the claimed regularity at all primes. The field case \(n=0\) has global dimension zero. The eighteen retained FAC lines cite the graded Hilbert-syzygy use and the localized projective-space bound. They remain editorial attribution; the source comparison does not silently alter those passages.

For primes \(\mathfrak p\subset\mathfrak q\) of \(A\), localize at \(\mathfrak q\). This is a regular local domain, hence Cohen--Macaulay. Every saturated chain between zero and its maximal ideal has length \(\dim A_{\mathfrak q}=\operatorname{ht}_A\mathfrak q\). Every saturated chain from zero to \(\mathfrak p\) has length \(\operatorname{ht}_A\mathfrak p\), by localizing at \(\mathfrak p\) and the same theorem. These intervals preserve all primes under localization. Concatenating such a zero-to-\(\mathfrak p\) chain with an arbitrary saturated \(\mathfrak p\)-to-\(\mathfrak q\) chain is saturated. Subtracting the two exact finite lengths gives \(\operatorname{ht}\mathfrak q-\operatorname{ht}\mathfrak p\). For equal endpoints the interval has length zero.

Now let the original finite-type domain be \(S=A/P\). Its primes are the primes of \(A\) containing \(P\), with all intervals preserved. A maximal chain of \(S\) begins at its zero prime and ends at a maximal ideal; the preimage of that maximal ideal is maximal in \(A\), of height \(n\). The interval formula gives every such chain length \(n-\operatorname{ht}_A P\). It follows that \[
 \dim S=\dim S_{\mathfrak m}=n-\operatorname{ht}_A P
\] for every maximal ideal \(\mathfrak m\) of \(S\).

For later use, if \(\mathfrak p\) is any prime of this same domain, with preimage \(P'\subset A\), the identical interval calculation gives \[
 \dim S_{\mathfrak p}=\operatorname{ht}_A P'-\operatorname{ht}_A P,
 \qquad
 \dim(S/\mathfrak p)=n-\operatorname{ht}_A P'.
\] Adding retains the exact identity \(\dim S=\dim S_{\mathfrak p}+\dim(S/\mathfrak p)\). This was proved by the original polynomial presentation; it does not assume that an arbitrary essentially finite-type ring has maximal ideals of the same height.

\subsubsection{Dimension at a point and the closed-specialization minimum}\label{algebra-proof-025-dimension-at-a-point-and-the-closed-specialization-minimum}

Source 27786--27867. Let \(S\) be a finite-type \(k\)-algebra, and let \(Z_i=V(\mathfrak q_i)\) be its finitely many irreducible components. Fix \(x\leftrightarrow\mathfrak p\), put \(I'=\{i:\mathfrak q_i\subset
\mathfrak p\}\), and retain \(T=\bigcup_{i\notin I'}Z_i\), \(U=X\setminus T\). For any open \(W\) containing \(x\) and contained in \(U\), its components are exactly the nonempty \(W_i=Z_i\cap W\), \(i\in I'\). A containment between two of them would, after closure in \(X\), give a containment of the two original components; these are distinct and maximal irreducible subsets.

Each \(W_i\) contains a point closed in \(X\), since the component is the spectrum of a finite-type domain and its closed points are dense. Take such a point \(z\). All generalizations of \(z\) lie in \(W\): if a closed complement contained a generalization, it would also contain its specialization \(z\). Therefore every prime chain in the local ring of \(Z_i\) at \(z\) lies inside \(W_i\), retaining its length \(\dim Z_i\) by the preceding domain result. The reverse inequality follows by taking closures of a strict chain of irreducible closed subsets of \(W_i\) in \(Z_i\); intersecting back with the open set recovers the chain, so strictness is retained. Hence \(\dim W_i=\dim Z_i\), and \[
 \dim_x X=\max_{i\in I'}\dim Z_i.
\] The equality of the dimension of a finite union of components with the maximum follows because the largest member of an irreducible chain is contained in one component; its whole chain lies there.

For a maximal \(\mathfrak m\supset\mathfrak p\), every component through \(x\) also passes through the closed point \(x_0\leftrightarrow\mathfrak m\). The minimal primes of \(S_{\mathfrak m}\) are exactly \(\mathfrak q_i S_{\mathfrak m}\) for the components through \(x_0\). Their quotient local rings are \((S/\mathfrak q_i)_{\mathfrak m/\mathfrak q_i}\), by the exact map \([s/t]\mapsto\bar s/\bar t\). Each has dimension \(\dim Z_i\), so \(\dim S_{\mathfrak m}\) is their maximum and is at least \(\dim_xX\). The nonempty open subset \(V(\mathfrak p)\setminus T\) of \(V(\mathfrak p)\) contains a closed point of \(X\), again by the finite-type Jacobson property. For its maximal ideal \(\mathfrak m\), no extra component occurs. This realizes the minimum and proves all three source equalities. For a closed \(x\), the only maximal ideal above \(\mathfrak p\) is \(\mathfrak p\) itself, proving the subsequent closed-point lemma.

These formulas also give the exact upper-semicontinuity statement: for any integer \(d\), \[
 \{x:\dim_xX\geq d\}=\bigcup_{\dim Z_i\geq d}Z_i,
\] a closed set, and the equality-\(d\) locus is its difference from the corresponding equality with threshold \(d+1\). These statements refer to the source's topological dimension at a point, not the height of a nonmaximal prime.

\subsubsection{The decomposition with its weaker exact hypothesis}\label{algebra-proof-025-the-decomposition-with-its-weaker-exact-hypothesis}

Source 27869--27903. The original proof applies the CM maximal-chain theorem to each \(S_{\mathfrak m}\). Therefore all its minimal-prime components have dimension \(\dim S_{\mathfrak m}\). By the preceding domain result, these are exactly the dimensions of the global components \(Z_i\) containing \(\mathfrak m\). Thus two components meeting at a closed point have equal dimension. If two components meet at all, their nonempty closed intersection contains a closed point (indeed a maximal ideal containing their sum). They therefore have equal dimension. Group the finitely many components by their dimension: \[
 T_d=\bigcup_{\dim Z_i=d}Z_i.
\] Each \(T_d\) is closed, and its complement is a finite union of the other groups, also closed. Distinct groups are disjoint. Every nonempty \(T_d\) has all its components of dimension \(d\).

The exact criterion requires less than Cohen--Macaulayness. For a finite-type \(k\)-algebra \(S\ne0\), the following conditions are equivalent: (1) all local rings at maximal ideals are equidimensional; (2) all local rings at all primes are equidimensional; (3) any two intersecting global irreducible components have equal dimension; (4) the function \(x\mapsto\dim_x\operatorname{Spec}S\) is locally constant; (5) there is the displayed finite decomposition into open-and-closed equidimensional sets \(T_d\).

Here are all implications. Condition (2) implies (1). The argument just given for the source proof uses only local equidimensionality at maximal ideals, so (1) implies (3). Grouping components proves (3) implies (5); the point-dimension formula makes the function equal to \(d\) on \(T_d\), giving (4). If (4) holds and \(z\) belongs to components with generic points \(\eta_i\), then \(\dim_{\eta_i}X=\dim Z_i\). The inverse image of any value of the locally constant function is open and closed. Its closedness forces it to contain every specialization of each point it contains. Since \(z\) specializes every such \(\eta_i\), all these component dimensions must agree. Thus (4) implies (3).

Finally assume (3), and fix a prime \(\mathfrak p\). For each minimal prime \(\mathfrak q_i\subset\mathfrak p\), the domain \(S/\mathfrak q_i\) and the formula in the preceding section give \[
 \dim (S/\mathfrak q_i)_{\mathfrak p/\mathfrak q_i}
   =\dim(S/\mathfrak q_i)-\dim(S/\mathfrak p).
\] All the first terms on the right are equal by (3), since these components meet at \(\mathfrak p\). The second term is independent of \(i\). These quotient local rings are precisely the components of \(S_{\mathfrak p}\), so that local ring is equidimensional. This proves (2), completing the equivalence with all original quotient maps retained.

The source CM hypothesis implies (1), but is stronger. For example retain \(S=k[x,y]/(x^2,xy)\). Its nilradical is \((x)\) and its reduced quotient is the domain \(k[y]\), so its spectrum has just one one-dimensional irreducible component and satisfies (3). At \(\mathfrak m=(x,y)\), \(x\) is nonzero: a polynomial multiplier outside \((x,y)\) has a nonzero constant term, so its product with \(x\) cannot belong to \((x^2,xy)\). Every element of the maximal ideal annihilates \(x\), since \(x^2=xy=0\); hence the local depth is zero, whereas its dimension is one. This local ring is not CM. Thus the weaker criterion has actual examples beyond the source hypothesis.

Finite type must be retained in these equivalences. In Heinrich's original semilocal example verified in group 652, both global components have dimension three, yet the local ring at the original \(M_1\) has components of dimensions two and three. That essentially finite-type example satisfies (3) but fails (1). The earlier evidence retains the full original six-variable presentation; it supplies a concrete boundary, not an authorization to weaken the present finite-type hypothesis.

\subsubsection{Product maps, gaps in the dimensions and empty spectra}\label{algebra-proof-025-product-maps-gaps-in-the-dimensions-and-empty-spectra}

The source invokes the disjoint-decomposition/product comparison. Here it is with the actual nilpotents retained. A clopen set \(T\subset\operatorname{Spec}S\) and its complement are quasi-compact, so write them \(D(I)\) and \(D(J)\) for finitely generated ideals \(I,J\). Disjointness gives \(IJ\subset\sqrt0\), and the covering gives \(I+J=S\). Since \(IJ\) is finitely generated by nilpotents, some power \((IJ)^N\) is zero: choose nilpotence exponents for its finite generators and use the same pigeonhole bound as in the fibre proof. The ideals \(I^N,J^N\) are comaximal, since any prime containing their sum would contain both \(I,J\). Therefore choose actual \(a\in I^N,b\in J^N\) with \(a+b=1\). Their product is zero, so \(a^2=a\), and \(D(a)=T\): a prime outside \(D(I)\) contains \(a\), while on \(D(I)\) it contains \(J\), hence \(b\), so it cannot contain \(a=1-b\). Thus \(e=a\) is the desired idempotent.

For the finitely many nonempty \(T_d\), take these idempotents \(e_d\). For distinct \(d,d'\), their product is idempotent and has empty standard open, so it is nilpotent, hence zero. Their sum is idempotent with standard open all of \(X\); thus its complement is a nilpotent idempotent and is zero. Hence \(\sum e_d=1\). The inverse product maps are \[
 S\longrightarrow\prod_d e_dS,\quad s\longmapsto(e_ds)_d,
 \qquad (s_d)_d\longmapsto\sum_d s_d.
\] Orthogonality and the sum-one identity prove both composites are identities; each factor ring has identity \(e_d\). These are exact maps of the original rings, not just maps of their reductions.

When a dimension between zero and \(\dim S\) does not occur, one may include a zero factor \(S_d=0\). Its spectrum is empty and every assertion about heights of its maximal ideals is vacuous; its dimension remains \(-\infty\), not \(d\). The union statement is interpreted using nonempty \(T_d\), and the source's full product index permits these zero factors. If \(S=0\), there are no components or maximal ideals, the union is empty and the product over the empty set is the zero ring in the category of unital commutative rings. No integer range ending at \(-\infty\) is needed. This is a separate proof clarification, not a proposed alteration of the source's dimension convention.

\subsubsection{Source use and continuation}\label{algebra-proof-025-source-use-and-continuation}

All six reports in 27648--27919 are adjudicated. MC-STK-ERR-1601 is retained. The FAC reference and history are read and preserved exactly; their presence is explicitly captured rather than stripped during comparison. The bounded corpus query has two routing records for the Bhatt--Hochster--Ma title, one TeX and one PDF-primary. Both remain unread for this dependency; there is no PDF fallback and no claim they are distinct mathematical works. Heinrich's already read original TeX supplies the precise finite-type boundary through group 652; no additional reading of that paper is claimed. The triangular free-basis proof, graded maps, component criterion and exact product maps above are established directly from the current source. Continue at authority 27920. Lookahead through 28030 is read but unadjudicated.

\hypertarget{algebra_noether_dimension_notes_20260928}{}
\subsection{Noether normalization, field extensions and the graded vertex}\label{algebra-proof-026-noether-normalization-field-extensions-and-the-graded-vertex}

Authority: \texttt{algebra.tex} 27920--28447 at commit \texttt{a04446e57ec1fbc252a871afcec7752fb2807b14}, SHA-256 \texttt{FA8BB92E58A4F78A2BD01B3B6A4A87DE0A0D279F5DD90641B574DD5FBFFFA4F3}. The complete current interval equals this authority byte for byte. There are no prior canonical operations or added environments in the interval.

These complete arguments and further consequences are separate editorial evidence. They preserve the original polynomials, automorphisms, units, prime labels and torsion. Neither translation branch is replaced by these expanded proofs. Proved consequences carry no novelty claim.

\subsubsection{Finite supports and explicit positive weights}\label{algebra-proof-026-finite-supports-and-explicit-positive-weights}

Source 27927--27960. Let \(N\subset\mathbb Z_{\geq0}^n\) be the original finite nonempty support, \(n\geq1\), and retain \(A_i=\max_{\nu\in N}\nu_i-\min_{\nu\in N}\nu_i\). The required inequalities are indexed by \(i=2,\ldots,n\); there is no \(e_0\). This resolves OCC-12271 and OCC-00545.

An explicit choice is \[
 e_n=1,\qquad
 e_i=1+\sum_{j=i+1}^n A_je_j
     =\prod_{j=i+1}^n(A_j+1)\quad(1\leq i<n).
\] The second equality follows downwards from \(\sum_{j=i+1}^n A_j\prod_{\ell=j+1}^n(A_\ell+1)
=\prod_{\ell=i+1}^n(A_\ell+1)-1\), a telescoping sum. Every weight is a positive integer. If \(\nu,\nu'\) first differ in coordinate \(i\), then \[
 |e_i(\nu_i-\nu_i')|\geq e_i
 >\sum_{j>i}A_je_j
 \geq\left|\sum_{j>i}e_j(\nu_j-\nu_j')\right|.
\] The first differing coordinate therefore determines the sign of the weighted difference, which cannot vanish. This proves injectivity on the entire original support, retaining its actual coordinate ranges. For \(n=1\), \(e_1=1\) gives the assertion directly, with no tail inequality. The optional phrase ``with \(a_\nu\in R\) not zero'' is mathematically and grammatically intelligible; OCC-00546 does not warrant a required source edit.

For any finite family of nonconstant polynomials in these same variables, choose the weights using the union of their finite supports. This one choice separates the support of every member simultaneously.

\subsubsection{The original substitutions, full coefficients and one relation}\label{algebra-proof-026-the-original-substitutions-full-coefficients-and-one-relation}

Source 27962--28022. Write the original polynomial \(g=\sum_{\nu\in N}a_\nu x^\nu\), with all listed \(a_\nu\ne0\). Keep the automorphism \[
 \sigma(x_i)=x_i+x_n^{e_i}\ (i<n),\qquad \sigma(x_n)=x_n.
\] Its inverse sends \(x_i\) to \(x_i-x_n^{e_i}\) for \(i<n\) and fixes \(x_n\); the two composites fix each generator. The full expansion of its \(\nu\)-summand is \[
 a_\nu\sum_{\substack{0\leq b_i\leq\nu_i\\i<n}}
 \left(\prod_{i<n}\binom{\nu_i}{b_i}x_i^{b_i}\right)
 x_n^{\,\nu_n+\sum_{i<n}e_i(\nu_i-b_i)}.
\] The term with every \(b_i=0\) is \(a_\nu x_n^{e\cdot\nu}\). Every other term has strictly smaller \(x_n\)-degree, since all weights are positive. The unique support vector \(\nu^*\) maximizing \(e\cdot\nu\) therefore gives \[
 \sigma(g)=a_{\nu^*}x_n^{d}
              +\sum_{j=0}^{d-1}c_j(x_1,\ldots,x_{n-1})x_n^j,
 \qquad d=e\cdot\nu^*>0 .
\] This is valid over arbitrary coefficient rings, including rings with nilpotents, because the leading coefficient is the single original nonzero coefficient \(a_{\nu^*}\); no cancellation or division has been used. All other binomial summands remain in the displayed expansion.

For a ring map \(R\to R'\), the same integer substitution and its inverse commute with coefficient base change. Some original coefficients may vanish. If the image polynomial is nonconstant, its surviving support is a subset of \(N\), so the same weights still separate it and its new leading coefficient is the image of one specified surviving original coefficient. If all nonconstant terms vanish, the image is constant or zero, and no positive-degree conclusion is asserted. This also proves the simultaneous assertion for the finite family above under any base change, with its precise surviving-support boundary.

Over the original field \(k\), choose nonzero \(f\in I\) with \(I\) proper. Then \(f\) is nonconstant, since a nonzero constant is a unit. Put \(y_i=x_i-x_n^{e_i}\) for \(i<n\). These actual polynomials have integer coefficients, and \(k[y_1,\ldots,y_{n-1},x_n]=k[x_1,\ldots,x_n]\), with inverse expressions \(x_i=y_i+x_n^{e_i}\). In the quotient by \(f\), retain the complete equation \[
 a x_n^d+\sum_{j<d}c_j(y_1,\ldots,y_{n-1})x_n^j=0,\qquad a\in k^\times.
\] It gives the explicit recurrence \(x_n^d=-a^{-1}\sum_{j<d}c_j(y)x_n^j\). Induction on the exponent reduces every power to the span of \(1,x_n,\ldots,x_n^{d-1}\). Hence the quotient by \((f)\), and then its quotient \(S\), is finite over the actual image of \(k[y_1,\ldots,y_{n-1}]\). The inverse unit \(a^{-1}\) has not been dropped from the calculation. For \(n=1\) there are no \(y_i\), and this is the same univariate proof. For \(n=0\) a field has no nonzero proper ideal, so the premise does not occur. OCC-00547 is the subject--verb correction ``there exist integers''.

\subsubsection{The full induction and a neighborhood of the original point}\label{algebra-proof-026-the-full-induction-and-a-neighborhood-of-the-original-point}

Source 28024--28071. For \(n=0\), the proper-ideal quotient is \(k\); take \(r=0\). For \(I=0\), take the original variables and \(r=n\). Otherwise use the preceding \(n-1\) integer polynomials and let \(S'\) be their actual image subalgebra of \(S\). This is nonzero and is generated by \(n-1\) elements. Apply the induction hypothesis to a polynomial presentation of that image. Its generators \(z_j\) are integer polynomials in the \(y_i\); substituting their actual integer-polynomial expressions in the original \(x_i\) retains integer coefficients. The composition \(k[z_1,\ldots,z_r]\to S'\to S\) is injective and finite. To check the latter directly, products of a finite generating family for \(S'\) over \(k[z]\) and a finite generating family for \(S\) over \(S'\) generate \(S\) over \(k[z]\).

The dimension statement follows from injective integral dimension equality and \(\dim k[z_1,\ldots,z_r]=r\). The integer-polynomial choice follows from the explicit construction, not from that dimension lemma. Thus OCC-12272 and OCC-00548 replace the misleading ``last assertion'' by the assertion that \(r\) is the dimension of \(S\).

For a point \(x\leftrightarrow\mathfrak q\) of the original affine finite-type algebra, the integer \(\dim_xX\) is the minimum dimension of its open neighborhoods. Basic opens form a basis: if a neighborhood attains the minimum, some \(D(g)\) with \(g\notin\mathfrak q\) lies inside it, and its dimension is squeezed between that minimum and the chosen neighborhood dimension. Thus \(\dim S_g=\dim_xX\). The localization remains finite type, with presentation \(S_g=S[t]/(gt-1)\), retaining \(g\) and its inverse. The preceding induction applies to this actual ring. The zero ring has no such point, so no point case is missing.

\subsubsection{The refined map, every unit and a finite free strengthening}\label{algebra-proof-026-the-refined-map-every-unit-and-a-finite-free-strengthening}

Source 28073--28123. Retain the original prime \(\mathfrak q\subset A=k[x_1,\ldots,x_n]\) and \(r=\operatorname{trdeg}_k\kappa(\mathfrak q)\). If \(\mathfrak q=0\), then \(r=n\); take the identity map. Conversely a nonzero polynomial relation among the images of the \(n\) generators makes their field transcendence degree at most \(n-1\): a maximal algebraically independent subfamily has fewer than \(n\) elements, and every omitted generator is algebraic over its field. Thus \(r=n\) forces \(\mathfrak q=0\). For a nonzero prime choose the original \(0\ne f\in\mathfrak q\).

Apply the actual \(\sigma\) above and put \(h=\sigma(f)\), \(\mathfrak q^\sigma=\sigma(\mathfrak q)\). Retain \(h=a x_n^d+\sum_{j<d}c_j(x_1,\ldots,x_{n-1})x_n^j\), \(a\in k^\times\). The coefficient ring is \(k[x_1,\ldots,x_{n-1}]\), as corrected in OCC-00550 and OCC-12273. The source's monic wording requires multiplication by \(a^{-1}\), in addition to the automorphism. Rather than silently replacing \(f\), track both choices: its monic version is \(h_0=a^{-1}h\in\mathfrak q^\sigma\), while the following proof works with the unscaled \(h\) and keeps \(a\) in every relation. Changing the last source generator by \(Y_n\mapsto aY_n\) identifies the two finite maps, and \((Y_{r+1},\ldots,Y_n)\) is unchanged by this unit.

Define \(B=k[Y_1,\ldots,Y_n]\to A\) by \(Y_i\mapsto x_i\) for \(i<n\) and \(Y_n\mapsto h\). There is the exact presentation \[
 A\simeq B[T]\big/
 \left(T^d+\sum_{j<d}a^{-1}c_j(Y_1,\ldots,Y_{n-1})T^j-a^{-1}Y_n\right),
\] sending \(T\mapsto x_n\). Its inverse sends \(x_i\) to \(Y_i\), \(i<n\), and \(x_n\) to \(T\); the defining equation identifies the image of \(h\) with \(Y_n\), so both composites fix all generators. Monic division proves that \(A\) is free over \(B\) on \(1,T,\ldots,T^{d-1}\). The coefficient map is injective because a free module basis containing \(1\) makes its coefficient zero whenever that coefficient times \(1\) is zero.

The contraction \(Q=(B\to A)^{-1}(\mathfrak q^\sigma)\) contains \(Y_n\), so \(Q=\mathfrak pB+(Y_n)\) with \(\mathfrak p\subset k[Y_1,\ldots,Y_{n-1}]\) prime. The finite injection \(B/Q\to A/\mathfrak q^\sigma\) gives a finite extension of fraction fields, hence \(\operatorname{trdeg}_k\kappa(\mathfrak p)=r\). Inductively choose the refined finite map on \(n-1\) variables carrying \(\mathfrak p\) back to \((Z_{r+1},\ldots,Z_{n-1})\). Extend it polynomially by \(Z_n\mapsto Y_n\). The preimage of \(Q\) is now \((Z_{r+1},\ldots,Z_n)\): modulo \(Z_n,Y_n\) this is exactly the induction statement, and the last generator lies in the contracted ideal. Composing with \(B\to A\) and then \(\sigma^{-1}:A\to A\) carries this ideal to the original \(\mathfrak q\), not to a substituted prime.

This proves a stronger assertion: the refined finite map may be chosen injective and finite free of positive rank. The inductive map has a finite free basis by the stronger induction; polynomial extension preserves that basis. If it is \(b_1,\ldots,b_D\), the composite before \(\sigma^{-1}\) has basis \(b_jT^\ell\), \(1\leq j\leq D,\ 0\leq\ell<d\). Transport that exact basis by \(\sigma^{-1}\) for the original target. This proves finite freeness and injectivity at every step, including rank one at \(n=0\) or \(\mathfrak q=0\).

Write the final original images as \(a_i=\varphi(Z_i)\). Then \(a_{r+1},\ldots,a_n\) is a regular sequence in \(A\): the coordinate sequence in the polynomial source is regular, and freeness preserves each injection and each quotient. Its final quotient is nonzero since the original \(\mathfrak q\) contains the sequence. Moreover \(A/(a_{r+1},\ldots,a_n)\) is finite free of the same rank over \(k[Z_1,\ldots,Z_r]\). The prime \(\mathfrak q/(a_{r+1},\ldots,a_n)\) is minimal in this quotient: its contraction is zero, any smaller prime also contracts to zero, and incomparability for the integral extension prohibits strict inclusion. Thus the original prime is a minimal prime of this explicitly constructed regular-sequence quotient. No assertion that \(\mathfrak q\) itself is generated by that sequence has been substituted.

\subsubsection{Generic comparison over a domain without losing torsion}\label{algebra-proof-026-generic-comparison-over-a-domain-without-losing-torsion}

Source 28125--28161. Let the original injection \(R\to S\) be finite type with \(R\) a domain, \(K=\operatorname{Frac}R\), and fixed generators \(x_i\). The algebra \(S_K\) is nonzero, since \(R\to S\) is injective. The field construction supplies \(y_1,\ldots,y_d\) which are integer polynomials in those \(x_i\), so they are actual elements of \(S\). Any polynomial relation over \(R\) among these elements maps to a relation over \(K\); injectivity of the polynomial map over \(K\) and \(R\to K\) show all its original coefficients vanish. Thus \(R[y]\to S\) is injective. The source's change of dummy index from \(i\) to \(j\) refers to the same already specified finite family; OCC-00551 is an optional stylistic change.

Choose the source monic polynomials \(P_i(T)=T^{m_i}+\sum_{j<m_i}c_{ij}(y)T^j\in K[y][T]\) with \(P_i(x_i)=0\) in \(S_K\), necessarily \(m_i\geq1\). Choose nonzero \(g\in R\) clearing all coefficients \(c_{ij}\). The element \(gP_i(x_i)\) is then defined in \(S\) and maps to zero in \(S_K\); localization gives nonzero \(t_i\in R\) with \(t_i gP_i(x_i)=0\) in \(S\). Put \(f=g\prod_i t_i\), nonzero because \(R\) is a domain. Then every \(fP_i\) has coefficients in \(R[y]\), and \(fP_i(x_i)=0\) in \(S\).

Retain \(x_i'=fx_i\) and the exact polynomial \[
 Q_i(T)=f^{m_i}P_i(T/f)
       =T^{m_i}+\sum_{j<m_i} f^{m_i-j}c_{ij}(y)T^j .
\] Every coefficient is in \(R[y]\), since \(m_i-j\geq1\) and \(fc_{ij}\) already is. Evaluation in \(S\) gives \(Q_i(x_i')=f^{m_i-1}(fP_i(x_i))=0\). These formulas retain every factor, including those killing localization torsion. For \(S'=R[y,x_1',\ldots,x_n']\subset S\), repeated use of the monic equations shows that the products \(\prod_i(x_i')^{b_i}\), \(0\leq b_i<m_i\), generate \(S'\) as an \(R[y]\)-module. They are not claimed linearly independent.

The localized inclusion \(S'_f\to S_f\) is injective: if \(s'/f^a\) maps to zero, some \(f^b s'=0\) in \(S\), and the same equation holds in the actual subring \(S'\). It is onto because each original generator \(x_i=x_i'/f\) lies in its image. This gives both directions of the claimed identification, with every original denominator retained.

Injectivity \(R\to S\) does not make \(S\to S_K\) injective. For example \(R=k[t]\), \(S=R[z]/(tz)\) has injective coefficient map (split by \(z\mapsto0\)), but \(z\ne0\) in \(S\) and \(z=0\) in \(S_K\). Here the generic relation \(P(T)=T\) must be multiplied by \(t\); with \(f=t\), \(x'=tz=0\), \(S'=R\), and \(S'_t=S_t\). This verifies exactly why the source's annihilator step cannot be omitted.

\subsubsection{Dimension, transcendence degree and the zero-length chains}\label{algebra-proof-026-dimension-transcendence-degree-and-the-zero-length-chains}

Source 28177--28287. For a finite-type domain \(S\), take the source finite injection \(k[y_1,\ldots,y_d]\to S\), \(d=\dim S\). Tensor with the fraction field \(k(y)\). The resulting finite-dimensional domain over \(k(y)\) is a field, since multiplication by a nonzero element is an injective endomorphism of a finite-dimensional vector space. It contains \(S\) and the inverses of every nonzero element of \(S\), so it is its fraction field \(K\). This proves the required field extension is finite and gives \(\operatorname{trdeg}_kK=d\). The earlier exact prime-interval calculation gives the same dimension at every maximal ideal. The k-algebra wording at 28180 is corrected minimally.

If \(\mathfrak q\subsetneq\mathfrak q'\), work in the finite-type domain \(S/\mathfrak q\). A maximal chain in its proper prime quotient \(S/\mathfrak q'\) extends at its bottom by \(0\subsetneq\mathfrak q'/\mathfrak q\). Thus the dimension strictly decreases, and the just-proved equality with residue-field transcendence degree gives the source specialization inequality. Finiteness of these dimensions is essential to the strict integer comparison.

For a general finite-type \(S\), retain \(\mathfrak p\), the minimal primes \(\mathfrak q_j\subset\mathfrak p\), and \(r=\operatorname{trdeg}_k\kappa(\mathfrak p)=\dim S/\mathfrak p\). For every \(j\), the earlier polynomial-presentation chain formula gives \[
 \dim(S/\mathfrak q_j)
   =\dim (S/\mathfrak q_j)_{\mathfrak p/\mathfrak q_j}+r.
\] Taking the maximum over these actual components yields \[
 \dim_x\operatorname{Spec}S
       =\dim S_{\mathfrak p}+\operatorname{trdeg}_k\kappa(\mathfrak p).
\] This is precisely the source's concatenated-chain argument. If \(r=0\), the upper chain consists of \(\mathfrak p\) alone. If its lower interval has length \(s(j)=0\), then \(\mathfrak q_j=\mathfrak p\), and no \(\mathfrak p'_0\) is introduced. These are separate proof clarifications of the displayed chains.

For a surjection of finite-type \(k\)-algebras \(S'\to S\), the corresponding residue fields \(\kappa(\mathfrak p')\to\kappa(\mathfrak p)\) are isomorphic: localization and quotient give the same fraction field of \(S'/\mathfrak p'=S/\mathfrak p\). Subtract the two preceding formulas to get the exact stated difference of point dimensions and prime heights. The two other occurrences of k algebra(s), at 28224 and 28272, receive only the missing noun hyphens.

\subsubsection{Field extension and the common original fibre ring}\label{algebra-proof-026-field-extension-and-the-common-original-fibre-ring}

Source 28289--28387. If \(S=0\), then \(K\otimes_k S=0\) and both dimensions are \(-\infty\); this case needs no polynomial ring on \(d=-\infty\) variables. For \(S\ne0\), tensor the finite injection furnished above with \(K\). Flatness of the field extension preserves injectivity and finite generation, giving the original finite injection \(K[y_1,\ldots,y_d]\to K\otimes_kS\), and integral dimension equality proves the asserted global equality.

For the pointwise assertion, keep the source presentation \(A=k[x_1,\ldots,x_n]\to S=A/I\), \(B=K[x_1,\ldots,x_n]\), and \(S_K=B/IB\). Retain the original primes \(\mathfrak q'\subset A\), \(\mathfrak q_K'\subset B\) and their quotient primes \(\mathfrak q,\mathfrak q_K\). Both vertical localized maps are flat. Their common local fibre is exactly \[
 B_{\mathfrak q_K'}/\mathfrak q'B_{\mathfrak q_K'}
 \ \simeq\
 (S_K)_{\mathfrak q_K}/\mathfrak q(S_K)_{\mathfrak q_K}.
\] The map sends \([b/u]\) to \([\bar b/\bar u]\); its inverse lifts quotient representatives. It is well defined and bijective because \(I\subset\mathfrak q'\), so quotienting by \(IB\) before \(\mathfrak q'B\) has the same final kernel. All denominators avoid the specified primes.

Writing the fibre dimension as \(h\), the local base-plus-fibre equality gives \[
 \operatorname{ht}(\mathfrak q_K')
     =\operatorname{ht}(\mathfrak q')+h,\qquad
 \operatorname{ht}(\mathfrak q_K)
     =\operatorname{ht}(\mathfrak q)+h.
\] The source's quotient codimension formula subtracts these two equalities. The point dimension of either ambient affine \(n\)-space is \(n\), including at nonclosed points, by its single component and the earlier point-dimension result. Hence \(\dim_xX=\dim_{x_K}X_K\). This proves the source cancellation through the exact same fibre ring, not merely through equally named dimensions. The duplicated ``Denote'' reports receive ``Denote by'' at 28343.

Combining the local dimension equality and the pointwise formula gives \[
 h=\dim(S_K)_{\mathfrak q_K}-\dim S_{\mathfrak q}
   =\operatorname{trdeg}_k\kappa(\mathfrak q)
      -\operatorname{trdeg}_K\kappa(\mathfrak q_K).
\] The missing finite verb at 28377 is corrected as reported twice. A prime minimal over \(\mathfrak qS_K\) exists because \(S_K\) is Noetherian and that ideal is proper. Properness follows from faithful flatness of the field base change. Such a prime contracts to \(\mathfrak q\): going down supplies a smaller prime over \(\mathfrak q\), and it still contains \(\mathfrak qS_K\), so minimality forces equality. The localized fibre then has only its minimal maximal prime and dimension zero. This proves the final source existence assertion with the required contraction checked.

\subsubsection{Exact maximum fibre dimension and both height-jump extremes}\label{algebra-proof-026-exact-maximum-fibre-dimension-and-both-height-jump-extremes}

Let \(F/k\) be the original finitely generated residue field \(F=\kappa(\mathfrak q)\), put \(r=\operatorname{trdeg}_kF<\infty\), and let \(K/k\) be arbitrary, of transcendence degree \(s\), possibly infinite. Then the actual tensor ring, including its possible nilpotents, satisfies \[
 \dim(K\otimes_kF)=\min(r,s).
\] Here is a direct proof. Choose an actual transcendence basis \(t_1,\ldots,t_r\) in \(F\). The extension \(F/k(t_1,\ldots,t_r)\) is finite: \(F\) is finitely generated as a field, and the remaining generators are algebraic. If \(D\) is its degree, tensoring its vector space basis gives a finite free, hence integral and injective, comparison \[
 A=(k[t_1,\ldots,t_r]\setminus\{0\})^{-1}K[t_1,\ldots,t_r]
      \ \longrightarrow\ K\otimes_kF .
\] It is obtained from \((K\otimes_k k(t))\otimes_{k(t)}F\simeq K\otimes_kF\); \((a\otimes b)\otimes c\mapsto a\otimes bc\) and \(a\otimes c\mapsto(a\otimes1)\otimes c\) are inverse. No reducedness or separability assumption is introduced. Integral dimension equality reduces only this dimension calculation to \(A\); the original tensor ring remains present.

Every prime \(P\subset K[t_1,\ldots,t_r]\) surviving in \(A\) contracts to zero in \(k[t_1,\ldots,t_r]\), so its residue field \(L\) contains \(k(t_1,\ldots,t_r)\). The earlier affine-domain dimension formula gives \(\operatorname{ht}P=r-\operatorname{trdeg}_K L\). This is at most \(r\). If \(s<\infty\), the actual field tower gives \(\operatorname{trdeg}_k L=s+\operatorname{trdeg}_K L\geq r\), so \(\operatorname{ht}P\leq s\). The entire interval below \(P\) survives the localization, giving the upper bound \(\dim A\leq\min(r,s)\).

For the lower bound put \(m=\min(r,s)\) and choose algebraically independent \(u_1,\ldots,u_m\in K\) over \(k\). In \(K[t_1,\ldots,t_r]\), take \(P=(t_1-u_1,\ldots,t_m-u_m)\). Its quotient is \(K[t_{m+1},\ldots,t_r]\). Evaluation of \(k[t_1,\ldots,t_r]\) into this quotient is injective: expand in the last \(r-m\) variables and use algebraic independence of the \(u_i\) on each coefficient. Thus \(P\) is disjoint from the localization set. Its height is \(m\), either by the exact dimension formula for this polynomial quotient or by its successive coordinate prime chain and its \(m\)-generator local maximal ideal bound. All those lower primes survive. This proves the lower bound, and hence the displayed equality. For \(r=0\) or \(m=0\), the argument uses empty transcendence bases and the zero prime directly.

Consequently every source height jump \(h\) obeys \(0\leq h\leq\min(r,s)\), and both extremes occur among primes over the fixed \(\mathfrak q\). For zero, use a minimal prime of the nonzero Noetherian tensor ring. For the maximum, a finite-dimensional spectrum attains its integer dimension: a finite strict chain of maximal possible length has a last prime whose local dimension equals that length. The prime correspondence between \(\operatorname{Spec}(K\otimes_kF)\) and primes of \(S_K\) contracting to \(\mathfrak q\) preserves their local fibre rings. Thus both choices really lie over the original \(\mathfrak q\). For an algebraic extension \(K/k\), \(s=0\), so every height jump is zero.

As an exact example retaining distinct base and point dimensions, take \(S=k[x]\), \(\mathfrak q=0\), \(K=k(u)\). Both primes \(0\) and \((x-u)\) of \(K[x]\) contract to zero in \(k[x]\); their local heights are respectively zero and one, while the topological dimension at both points remains one. This realizes both extremes without suggesting that point dimension and local height are identical.

\subsubsection{The graded vertex, its Hilbert function and its full completion}\label{algebra-proof-026-the-graded-vertex-its-hilbert-function-and-its-full-completion}

Source 28402--28434. Keep the original finitely generated standard graded \(k\)-algebra \(S\), \(S_0=k\), and \(\mathfrak m=S_+\). Its quotient is \(k\), so \(\mathfrak m\) is maximal. Every minimal prime is homogeneous by the source homogeneous-prime lemma. A proper homogeneous ideal contains no nonzero degree-zero scalar, so every such minimal prime lies in \(\mathfrak m\). All components therefore pass through the vertex. The earlier point dimension formula gives \(\dim S=\dim_{\mathfrak m}\operatorname{Spec}S=\dim S_{\mathfrak m}\).

Standard generation gives the exact ideals \[
 \mathfrak m^d=\bigoplus_{j\geq d}S_j.
\] Every product of \(d\) positive-degree elements has degree at least \(d\). Conversely each degree-\(j\) monomial in the degree-one generators, \(j\geq d\), factors as a product of \(d\) positive-degree monomials; linear combinations give the reverse inclusion. Thus the degree-\(d\) projection identifies \(\mathfrak m^d/\mathfrak m^{d+1}\) with \(S_d\), including \(d=0\).

Put \(R=S_{\mathfrak m}\), \(\mathfrak n=\mathfrak mR\). For \(u\notin\mathfrak m\), write its original homogeneous decomposition \(u=c+v\), \(c\in k^\times\), \(v\in\mathfrak m\). On \(\mathfrak m^d/\mathfrak m^{d+1}\), multiplication by \(u\) is multiplication by \(c\), with inverse \(c^{-1}\). This proves that localization gives an isomorphism onto \(\mathfrak n^d/\mathfrak n^{d+1}\), with every denominator accounted for. The source local Hilbert function is exactly \(\operatorname{length}_R(\mathfrak n^d/\mathfrak n^{d+1})\), as read at 13994--14015; it is not the cumulative length. Since these quotients are vector spaces over \(k\), that length is \(\dim_k S_d\). This proves the claimed equality of functions.

If the eventual polynomial \(P\) is nonzero, its leading coefficient is positive (its values are eventual nonnegative dimensions), and summing from zero through \(n\) increases its degree by one. The source local dimension theorem applied to that cumulative polynomial therefore gives \(\dim R=\deg P+1\). If \(P=0\), all sufficiently high \(S_d\) vanish, so \(S\) is finite dimensional, \(\mathfrak m\) is nilpotent and \(R\) is nonzero Artinian of dimension zero. This gives \(\deg(0)+1=-1+1=0\) exactly. No subtraction of infinities or negative binomial indices is used. The definitions and the prior group 384's zero-degree check agree.

In fact the preceding degree maps are multiplicative and prove the complete graded algebra identification \[
 S\longrightarrow\operatorname{gr}_{\mathfrak n}R,\qquad
 a\in S_d\longmapsto a/1\bmod\mathfrak n^{d+1}.
\] Every degree map is bijective by the explicit localization inverse above, and multiplication preserves degrees, so this is an isomorphism of graded \(k\)-algebras.

For every \(r\geq1\), the full localized quotient map \(S/\mathfrak m^r\to R/\mathfrak n^r\) is also an isomorphism. Indeed the inverse of \(u=c+v\) modulo \(\mathfrak m^r\) is the exact finite expression \[
 c^{-1}\sum_{j=0}^{r-1}(-c^{-1}v)^j.
\] Multiplication by \(c+v=c(1+c^{-1}v)\) gives \(1-(-c^{-1}v)^r\), which equals one in that quotient. This supplies an inverse for every localized denominator without assuming one in the original ring. The quotient is the direct sum of its original degrees \(0,\ldots,r-1\). Passing to compatible quotient systems therefore gives inverse maps \[
 \widehat{R}\ \simeq\ \prod_{d\geq0}S_d .
\] On the right multiplication is the actual convolution \((a b)_d=\sum_{i=0}^d a_i b_{d-i}\); this sum is finite in each degree. A sequence of coefficients maps to its finite truncations, and a compatible system maps to its uniquely stabilized coefficients in each degree. These operations are inverse, preserve addition and convolution, and identify the respective adic topologies. All homogeneous nilpotents and their products are retained. This is a proved comparison with the original completion, not a replacement of the graded ring or its Hilbert function.

\subsubsection{Corpus and continuation boundaries}\label{algebra-proof-026-corpus-and-continuation-boundaries}

The two bounded corpus queries return eight routing matches across both layers and five distinct canonical IDs. They include physics records, local programme notes and a PDF-primary algebra volume. None is read or adopted as a theorem for this batch; no PDF fallback is used. The existing primary Stacks source and earlier source-bound proofs supply the exact calculations above. No new claim of human-literature novelty or exhaustive topic coverage is made. The previous original Heinrich and other bounded reading records are retained unchanged in the topic route.

All seventeen reports through 28447 receive dispositions. Required edits are minimal; the two stylistic reports remain unapplied. Continue at 28448, Generic flatness. Lookahead through 28460 has been read but remains unadjudicated.

\hypertarget{algebra_generic_freeness_notes_20260928}{}
\subsection{Generic freeness, exact presentations and the original good locus}\label{algebra-proof-027-generic-freeness-exact-presentations-and-the-original-good-locus}

Authority: Stacks Project authors, commit a04446e57ec1fbc252a871afcec7752fb2807b14, algebra.tex:28448--28876; SHA-256 FA8BB92E58A4F78A2BD01B3B6A4A87DE0A0D279F5DD90641B574DD5FBFFFA4F3. The entire current interval equals this authority. Seven received reports are read in full. The next section starts at 28877; lookahead ends at 28895 and is unadjudicated.

This is separate editorial evidence. Both translation branches retain the original structure and arguments. The six proposed minimal operations concern five reported wording/type loci and the separately detected empty-spectrum dimension convention. The expanded arguments and two consequences below are not replacement translation text. No novelty claim is made.

\subsubsection{The Noetherian argument and the empty generic fibre}\label{algebra-proof-027-the-noetherian-argument-and-the-empty-generic-fibre}

Source: 28457--28524. Write \(K=\operatorname{Frac}(R)\), \(S_K=K\otimes_R S\). If \(S_K=0\), the equality \(1/1=0\) in the localization supplies \(0\ne f\in R\) with \(f1_S=0\). Then \(S_f=0\) and \(M_f=0\), which is free of rank zero. This also treats \(S=0\), taking \(f=1\). The dimension of this spectrum is \(-\infty\), not \(-1\): Topology, definition-Krull, authority topology.tex:1391--1419, explicitly says this at 1410. Algebra repeats the convention at 43890--43895. Hence the proposed source correction at 28494 is precisely \(d=-1\) to \(d=-\infty\); the start of the induction remains the same.

For nonempty generic fibre put \(d=\dim S_K\), a nonnegative integer. Since \(R\) is Noetherian and \(S\) is of finite type, \(S\) is Noetherian and the finite \(S\)-module \(M\) has the source's finite prime filtration. Its factors are \(S/\mathfrak q\). A finite product of the nonzero localizing elements for these factors remains nonzero in \(R\). Localization preserves the filtration. For an exact sequence with free quotient, choose a lift of each quotient basis vector; this defines a module section and splits the sequence. Thus freeness of all factors yields freeness of \(M\), with no finite-rank assumption.

Each domain factor with nonzero kernel from \(R\) is killed by a nonzero element of that kernel. For a remaining factor, retain its name \(S\) and the injection \(R\subset S\). The original normalization-over-a-domain argument (28125--28161, verified in group 676) gives, after one localization in \(R\), the original finite extension \[
 A=R[y_1,\ldots,y_d]\ \subset\ S.
\] Here the \(y_i\) are algebraically independent. Localizing \(S\) by \(A\setminus\{0\}\) gives a finite-dimensional domain over \(\operatorname{Frac}(A)\), hence a field: multiplication by a nonzero element is an injective endomorphism of a finite-dimensional vector space, therefore surjective. This field is \(\operatorname{Frac}(S)\). Choose its basis \(z_1,\ldots,z_r\) from the spanning images of elements of \(S\). The original map \[
 A^{\oplus r}\longrightarrow S,\qquad (b_j)\longmapsto\sum_j b_jz_j
\] is injective by that basis independence. Its cokernel \(N\) is finite over \(A\) and becomes zero over \(\operatorname{Frac}(A)\). For a finite generating set of \(N\), choose a nonzero annihilator for each generator after localization; their product is a nonzero \(g\in A\) annihilating \(N\). No common annihilator for an arbitrary infinite family is presumed.

The image of \(g\) in \(K[y_1,\ldots,y_d]\) is nonzero. If it is a unit, the quotient has empty spectrum. Otherwise every prime containing it is nonzero, so every chain above such a prime extends strictly downward by \((0)\); the polynomial ring has dimension \(d\), and its quotient by \(g\) has dimension at most \(d-1\). For \(d=0\) the quotient is necessarily zero. The lower-dimensional induction applies to \(N\) as a module over \(A/(g)\). After the resulting localization \(N_f\) is free over \(R_f\). The original \(A_f^{\oplus r}\) is also free over \(R_f\), on the monomials in the specified \(y_i\) and the \(r\) basis positions. The displayed exact sequence therefore splits over \(R_f\) and \(S_f\) is free. If successive localizations use a fraction \(b/a^v\), their common base is \(R_{ab}\), with \(ab\ne0\); this retains the actual numerator and denominator.

Report OCC-00557 concerns the two fused clauses at 28522. Replacing ``also we conclude'' by ``also, we conclude'' gives the required clause boundary with a minimal edit. It does not change this argument.

\subsubsection{Finite coefficient models and the original relation kernel}\label{algebra-proof-027-finite-coefficient-models-and-the-original-relation-kernel}

Source: 28526--28625. Retain the presentation \[
 S=R[x_1,\ldots,x_n]/(g_1,\ldots,g_m),\qquad
 M=S^{\oplus t}/\sum_i S n_i,\qquad
 n_i=(\overline g_{i1},\ldots,\overline g_{it}).
\] Let \(R_0\) be the subring generated over the prime subring by every original coefficient in these finitely many polynomials. It is the image of a polynomial ring over \(\mathbf Z\) with finitely many variables, hence Noetherian, and is a domain because \(R_0\subset R\). Define \(S_0,M_0\) by the same polynomials and the same matrices, now over \(R_0\).

The coefficient map gives \(R\otimes_{R_0}S_0\to S\), \(r\otimes \overline p\mapsto r\overline p\). Its inverse sends each coefficient \(r\) to \(r\otimes1\) and each original variable to \(1\otimes\overline{x_i}\); every \(g_i\) vanishes in both directions. These maps are inverse on the algebra generators. On modules, tensor the actual presentation of \(M_0\). Right exactness gives the presentation with the same relation rows over \(S\), so the map \(r\otimes\overline v\mapsto r\overline v\) identifies it with \(M\). No flatness of \(R\) over \(R_0\) is used. A nonzero \(f\in R_0\) from the preceding argument stays nonzero in \(R\); a basis of \((M_0)_f\) tensors to a basis of \(M_f\).

For the arbitrary finite-type case first take \(P=R[x_1,\ldots,x_n]\) and a finite \(P\)-module \(M\), with its actual surjection \(P^t\to M\) and kernel \(N\). Localization at the nonzero elements of \(R\) is exact: for example, a localized vector with zero image has a nonzero multiplier sending its numerator into the original kernel. Thus \(N_K\) is the kernel of \(P_K^t\to M_K\). The ring \(P_K\) is Noetherian, so choose finitely many generators of \(N_K\), each a fraction with numerator in \(N\). The finitely many original numerators \(n_1,\ldots,n_r\) generate \(N_K\) after localization. Put \[
 M'=P^t/\sum_{i=1}^r Pn_i.
\] The map \(M'\to M\) is surjective and becomes an isomorphism over \(K\). The finite coefficient argument gives \(M'_f\) free over \(R_f\). A free module over the domain \(R_f\) injects into its tensor product with \(K\): every vector has finite support and each coefficient injects. Consequently every element of the kernel of \(M'_f\to M_f\) is zero, since it is already zero over \(K\). This proves the isomorphism without assuming that \(N\) was finitely generated before localization. It proves both freeness and finite presentation after localization.

For \(P\to S\) surjective, \(S\) and \(M\) are finite \(P\)-modules. Apply this result to both and take the product of the two nonzero base elements. Write \(J=\ker(P_f\to S_f)\). Finite presentation of the cyclic \(P_f\)-module \(S_f\) implies that this particular kernel \(J\) is finitely generated (the finite-presentation kernel property proved in algebra.tex:390--475). Thus \(S_f\) is a finitely presented \(R_f\)-algebra. A finite presentation of \(M_f\) over \(P_f\), tensored along \(P_f\to S_f\), remains a presentation of \(M_f\) over \(S_f\): the multiplication map \(S_f\otimes_{P_f}M_f\to M_f\) has inverse \(m\mapsto1\otimes m\), since \(J\) acts by zero. This proves the last source inference.

OCC-00558 inserts ``by'' at 28545. OCC-00559 changes the same original line 28593 to ``Denote by \(K\) the fraction field of \(R\). Set'', followed by the existing displayed definitions. No relation or map changes.

\subsubsection{The good locus and its localization maps}\label{algebra-proof-027-the-good-locus-and-its-localization-maps}

Source: 28627--28723. Keep exactly the source definition \(U(R\to S,M)\): it is the union of \(D(f)\) on which \(S_f\) is a finitely presented \(R_f\)-algebra, \(M_f\) is a finitely presented \(S_f\)-module, and both are free over \(R_f\). Freeness may have infinite rank. This definition is stronger than the statement that each stalk is free.

If \(0\to M_1\to M_2\to M_3\to0\) is exact and a point has good neighborhoods \(D(f_1)\) and \(D(f_3)\) for the outer terms, localize at the actual product \(f_1f_3\). The two outer modules remain free, and the sequence splits as a sequence of base modules by lifting a basis of \(M_3\). Finite presentation of the middle \(S\)-module follows from algebra.tex:390--475; explicitly lift generators of \(M_3\), add generators of \(M_1\), and lift their finitely many relations. The relations for \(M_1\), together with the lifted relations for \(M_3\), generate the middle relation module: subtract the latter to make the projection zero, then use the former. This proves the source inclusion of good loci.

For the localization equality, retain \(g=a/f^n\in R_f\) and \(h=af\). The comparison \(R_h\to(R_f)_g\) fixes \(R\); \(a=g f^n\) and \(f\) are invertible on the right. Its inverse sends \(1/f\) to \(a/h\) and \(1/g\) to \(f^{n+1}/h\) in \(R_h\). Indeed \(1/a=f/h\), so the latter is precisely \(f^n/a\), also for \(n=0\). These formulas prove both composites on all original generators and inverses. The same formulas give the ring maps for \(S\) and the module maps on fractions of \(M\). Hence \[
 U(R_f\to S_f,M_f)=D(f)\cap U(R\to S,M).
\] Conversely a witness \(g\in R\) is restricted by replacing it with \(fg\); freeness and both presentations persist under this localization. Zero localizations yield empty opens and zero modules, so cause no exception.

For the dense-open reduction, let \(W\) be a nonempty open of \(\operatorname{Spec}R\). It meets the dense covered open \(U\). Some \(W\cap D(f_i)\) is nonempty, and density of the good locus within \(D(f_i)\) makes it meet that locus. The preceding equality then makes \(W\) meet the original good locus. This proves density directly for arbitrary spectra and arbitrary covering index sets.

\subsubsection{Reduced bases, coefficients and finite presentation across the substitution}\label{algebra-proof-027-reduced-bases-coefficients-and-finite-presentation-across-the-substitution}

Source: 28725--28870. A finite \(P=R[x_1,\ldots,x_n]\)-module with generators \(m_1,\ldots,m_t\) has the filtration \(M_i=\sum_{j\le i}Pm_j\). Its \(i\)-th factor is \(P/J_i\), where \(J_i\) is the annihilator of the image of \(m_i\) in that quotient. In the source notation \(P=S\); thus the three reports OCC-12286, OCC-00560 and OCC-01408 identify one type error, corrected once: \(R/J_i\) at 28771 becomes \(S/J_i\). No contraction of \(J_i\) to \(R\) gives the same module in general.

A finite intersection of dense opens is dense: intersect a nonempty open successively with each dense open; at every finite step the intersection remains nonempty and open. The extension result reduces the proof to \(P/J\).

Let \(I\subset R\) be the ideal generated by all coefficients of elements of \(J\). The set of coefficients itself is already an ideal. Scalar multiples are coefficients after multiplying the polynomial by that scalar. For addition, if \(a,b\) are the coefficients at multi-indices \(A,B\) of \(p,q\in J\), choose a coordinatewise upper bound \(L\) of \(A,B\). The coefficient at \(L\) of \[
 x^{L-A}p+x^{L-B}q\in J
\] is exactly \(a+b\); no other monomial of either polynomial contributes at \(L\). Negatives and zero are included. When \(n=0\), this is simply addition in the ideal \(J\subset R\). Thus \(1\in I\) is an actual coefficient of an actual element of \(J\).

Put \(U_1=D(I)\) and \(U_2=\operatorname{Spec}R\setminus\overline{U_1}\). Their union is dense: every nonempty open either meets \(U_1\), or is disjoint from its closure and lies in \(U_2\). On a basic open within \(U_1\), the localized coefficient ideal is the unit ideal because it lies in no prime. On a basic open within \(U_2\), the localized coefficient ideal lies in every prime, and hence is zero because \(R_f\) remains reduced. For completeness, if \((a/f^r)^v=0\), some \(f^s a^v=0\); then \((f^s a)^v=0\), so \(f^s a=0\) and \(a/f^r=0\). In the latter case \(J_f=0\) and \(M_f=P_f\).

In the unit-ideal case keep an actual polynomial \(g=\sum_{K\in E}a_Kx^K\in J\) with a unit coefficient. Take \(E\) to be its finite nonzero support, and let \(m\) count the nonunit coefficients. If \(m>0\), select a nonunit coefficient \(a_K\) and perform the same dense-open division using \(D(a_K)\) and the complement of its closure. On each basic open of the first part \(a_K\) becomes a unit; on each basic open of the second it becomes zero, by reducedness. Original unit coefficients stay units unless the localization is the zero ring, whose spectrum is empty. Other nonunit coefficients can become units or zero, but no original unit becomes a nonunit in a nonzero localization. Delete only coefficients proved zero. Therefore \(m\) strictly decreases on every nonempty branch.

For \(m=0\), a constant support makes \(g\) a unit and \(M=0\). This treats \(n=0\). Otherwise use exactly the source substitution \[
 x_i=y_i+y_n^{e_i}\quad(i<n),\qquad x_n=y_n,
\] with inverse \(y_i=x_i-x_n^{e_i}\), \(y_n=x_n\). Group 672 supplies explicit weights from the original support ranges, \(e_n=1\), \(e_i=\prod_{j>i}(A_j+1)\), where \(A_j=\max_E k_j-\min_E k_j\). Their strict separation proof applies to this same support. The entire polynomial is \[
 G=\sum_{K\in E}a_K
 \sum_{0\le b_i\le k_i,\ i<n}
 \left(\prod_{i<n}\binom{k_i}{b_i}y_i^{b_i}\right)
 y_n^{\,k_n+\sum_{i<n}e_i(k_i-b_i)}.
\] The unique largest exponent has its original coefficient \(a\), which is a unit. Write \(G=aT^d+\sum_{j<d}c_j(y)T^j\), \(T=y_n\), \(d\ge1\), and \(A=R[y_1,\ldots,y_{n-1}]\). The original quotient is a module over \[
 B=A[T]/(a^{-1}G).
\] Monic division gives the precise free \(A\)-basis \(1,T,\ldots,T^{d-1}\): subtract the leading term repeatedly, retaining all coefficients \(a^{-1}c_j\); a nonzero multiple of the monic polynomial cannot have degree below \(d\), so remainders are unique. Since \(G\in J\), the same quotient \(M\) is finite over \(A\).

We now justify the source equality of good loci at 28833, rather than presuming finite presentation transfers automatically. For any finite \(B\)-module \(L\), \[
 L\text{ is finitely presented over }A[T]
 \quad\Longleftrightarrow\quad
 L\text{ is finitely presented over }B
 \quad\Longleftrightarrow\quad
 L\text{ is finitely presented over }A.
\] The first forward implication tensors an \(A[T]\)-presentation with \(B\), using that \(G\) kills \(L\). The reverse lifts the finitely many \(B\)-relation vectors and adds the relation \(a^{-1}G\) in every generator position. For the second forward implication, replace the finite free \(B\)-modules in a presentation by their displayed finite free \(A\)-bases, giving an \(A\)-presentation. Conversely choose a finite \(A\)-presentation with generators \(v_j\) and relation vectors \(r_\ell\). Write \(Tv_j=\sum_i b_{ij}v_i\), \(b_{ij}\in A\). Over \(B\), the finitely many relations \(r_\ell\) and \(T e_j-\sum_i b_{ij}e_i\) present \(L\): repeated use of the action relations reduces every coefficient polynomial to an \(A\)-linear vector, which vanishes exactly when it is a combination of the \(r_\ell\). This proves injectivity of the presented map as well as surjectivity.

These arguments persist after every localization in \(R\). Both polynomial algebras over \(R_f\) are free and finitely presented. The condition that \(M_f\) be free over \(R_f\) is literally identical on the two sides. Thus their good loci are equal, and induction on \(n\) proves density.

Finally for \(P\to S=P/J\), at each base localization finite presentation of \(S\) as an \(R\)-algebra is equivalent to finite generation of this presentation ideal \(J\), hence to finite presentation of \(S\) as a \(P\)-module. One direction of this standard presentation independence follows by lifting a second finite set of algebra generators to polynomials in the original generators, and the original generators to polynomials in the second set. Adjoin both sets and the finitely many equations giving those polynomial expressions. Eliminating either set by substitution identifies the resulting quotient with the original algebra. The finitely many defining relations in the second presentation, together with these expression relations, therefore generate the kernel in the first. The reverse direction uses the presentation \(P/J\) itself.

If \(J\) is finitely generated, a finite \(S\)-presentation for \(M\) becomes a finite \(P\)-presentation by lifting its relations and adding \(J\) times each module generator; the converse tensors with \(S\). The freeness conditions in the source then show exactly \[
 U(R\to S,M)=U(R\to P,S)\cap U(R\to P,M).
\] This proves the final displayed equality and its density consequence. OCC-00561 repairs the clause order at 28868--28869; it does not add a different topological hypothesis.

\subsubsection{Finite complexes and arbitrary base change after one localization}\label{algebra-proof-027-finite-complexes-and-arbitrary-base-change-after-one-localization}

Consequence of the source generic-freeness argument, with its full finite coefficient construction retained. Let \(R\) be a domain, \(S\) a finite-type \(R\)-algebra and \(C_\bullet\) a complex with only finitely many nonzero terms, each a finite \(S\)-module. There exists \(0\ne f\in R\) such that \(S_f\) is finitely presented over \(R_f\), and every term, cycle, boundary and homology module of \(C_{\bullet,f}\) is free over \(R_f\) and finitely presented over \(S_f\). For every \(R_f\)-algebra \(T\), the natural maps identify cycles and boundaries after tensoring and give \[
 T\otimes_{R_f}H_i(C_{\bullet,f})
 \ \xrightarrow{\ \cong\ }\
 H_i(T\otimes_{R_f}C_{\bullet,f})
\] in every degree. No flatness of \(T\) is required and no finite rank over \(R_f\) is asserted.

Here is a proof without assuming that the original kernels over \(S\) are finite. First apply the source theorem to \(S\) and each of the finitely many terms, and take the product \(a\ne0\) of the resulting elements. Work temporarily over \(R_a\); \(S_a\) is finitely presented and each \(C_{i,a}\) has a finite \(S_a\)-presentation. Choose matrices \(A_i:S_a^{p_i}\to S_a^{q_i}\) with cokernel \(C_{i,a}\), and matrices \(D_i:S_a^{q_i}\to S_a^{q_{i-1}}\) lifting the actual differentials. Because these are well-defined differentials, there are finite matrices \(E_i,F_i\) with \[
 D_i A_i=A_{i-1}E_i,\qquad
 D_{i-1}D_i=A_{i-2}F_i.
\] Choose polynomial representatives for every matrix entry and for a finite algebra presentation of \(S_a\). Each of the finitely many entrywise equalities above has a witness as a finite sum of polynomial multiples of the defining algebra relations. Include every coefficient of these witness polynomials as well as every original presentation and matrix coefficient in a finitely generated subring \(R_0\subset R_a\).

The identical algebra presentation over \(R_0\) defines a Noetherian algebra \(S_0\). The same matrices and the retained witness equations define an actual complex \(C_{\bullet,0}\) over \(S_0\), not merely a sequence of unrelated maps. Its base change to \(R_a\) is the original complex \(C_{\bullet,a}\), with the specified maps. Set \[
 Z_{i,0}=\ker d_{i,0},\quad B_{i,0}=\operatorname{im}d_{i+1,0},
 \quad H_{i,0}=Z_{i,0}/B_{i,0}.
\] These modules are finite, and hence finitely presented, over the Noetherian ring \(S_0\). Apply Noetherian generic freeness to \(S_0\), all \(C_{i,0}\), all \(Z_{i,0},B_{i,0},H_{i,0}\), and take the product \(b\ne0\) of their finitely many witnesses in \(R_0\). All those localized modules are free over \((R_0)_b\). In particular both exact sequences \[
 0\to Z_{i,0,b}\to C_{i,0,b}\to B_{i-1,0,b}\to0,\qquad
 0\to B_{i,0,b}\to Z_{i,0,b}\to H_{i,0,b}\to0
\] split as sequences of \((R_0)_b\)-modules, by lifting bases of their free quotients.

Write the actual image of \(b\) in \(R_a\) as \(c/a^v\), \(c\in R\), \(c\ne0\), and take \(f=ac\). The original inclusion \(R_0\subset R_a\) sends \(b\) to a unit in \(R_f\). Tensor the two split sequences with \(R_f\). They remain exact even though this coefficient extension need not be flat. The differential factors through the displayed quotient onto \(B_{i-1}\) and its inclusion into \(C_{i-1}\), so the resulting cycle, boundary and homology modules are precisely those of the original \(C_{\bullet,f}\). They are free over \(R_f\), because a free basis tensors to a free basis, and finitely presented over \(S_f\), because their finite \(S_{0,b}\)-presentations tensor to presentations over \(S_f\).

The same two sequences now split over \(R_f\). Tensoring them with an arbitrary \(T\) proves the claimed identifications of kernels and images. The displayed homology map sends \(t\otimes[z]\) to \([t\otimes z]\); exactness proves it is well-defined and bijective. These identifications are the natural ones and do not depend on which sections were used to prove exactness. Zero terms and zero cycles are free on the empty basis. A finite list of maps can be treated by applying this proof to their two-term complexes and taking a common product of the witnesses; thus kernels, images and cokernels for every map in that list have the same base-change property.

\subsubsection{A dense good locus with no dense principal open}\label{algebra-proof-027-a-dense-good-locus-with-no-dense-principal-open}

This exact example records the boundary of the source's reduced-base conclusion and of replacing neighborhood freeness by stalkwise freeness. Fix a field \(k\), let \[
 R=\prod_{n\ge1}k,\qquad I=\bigoplus_{n\ge1} k e_n\subset R,\qquad
 S=R,\quad M=R/I,
\] where \(e_n\) is the original coordinate idempotent. The ideal \(I\) consists of the sequences with finite support. The ring \(R\) is reduced, \(S\) is a finite-type \(R\)-algebra and \(M\) is cyclic.

For \(f=(f_n)\in R\), put \(A=\{n:f_n\ne0\}\) and define \(g_n=f_n^{-1}\) on \(A\), \(g_n=0\) otherwise. Then \(fg=e_A\), the coordinate indicator of \(A\). Projection gives an isomorphism \[
 R_f\longrightarrow R_A:=\prod_{n\in A}k,\qquad
 r/f^j\longmapsto (r_n f_n^{-j})_{n\in A}.
\] Its inverse extends a sequence by zero outside \(A\), and sends that extension to its class in \(R_f\). The two composites agree because \(e_A=1\) in \(R_f\) and \(f\) is invertible on \(A\). This includes \(A=\varnothing\), with the zero ring as the empty product. The image of \(I_f\) is exactly the finite-support ideal \(I_A\): localization cannot add coordinates to a support, and every finite-support vector is the projection of one in \(I\). Thus \(M_f=R_A/I_A\).

If \(A\) is finite, \(I_A=R_A\), so \(M_f=0\), free and finitely presented. If \(A\) is infinite, \(I_A\) is a nonzero proper ideal: it contains each \(e_n\) for \(n\in A\), while the unit has infinite support. It is not finitely generated, since a finite collection of its generators has support in one finite subset, and multiplication by arbitrary sequences cannot enlarge that union. Hence its cyclic quotient is not finitely presented. It is not free either: its annihilator is the nonzero ideal \(I_A\), whereas a nonzero free module over a nonzero ring is faithful. We have proved the exact equality \[
 U(R\to R,R/I)=\bigcup_{n\ge1}D(e_n).
\]

Each \(D(e_n)\) is a singleton, identified through \(R_{e_n}=k\) with the coordinate prime. The union is dense: any nonempty basic open \(D(f)\) has at least one nonzero coordinate \(f_n\) and contains that coordinate prime. But no dense principal open is contained in it. Indeed, if \(f\) vanishes at any coordinate \(n\), the nonempty open \(D(e_n)\) is disjoint from \(D(f)\), so \(D(f)\) is not dense. If \(f\) has no zero coordinate, its coordinatewise inverse belongs to \(R\); it is a unit and \(D(f)=\operatorname{Spec}R\). This spectrum is not the displayed union: the nonzero ring \(R/I\) has a maximal ideal, whose preimage contains every \(e_n\). Thus its prime lies outside every coordinate open. Equivalently a principal witness for the good locus has finite support, which always misses another coordinate.

Nevertheless every stalk of \(M\) is free. To see this without a classification by ultrafilters, first note that \(R_{\mathfrak p}\) is a field for any prime \(\mathfrak p\). For \(r\in\mathfrak p\), the coordinate indicator \(e\) of its support satisfies \(e=rg\) for the coordinatewise partial inverse \(g\). Thus \(e\in\mathfrak p\), and \(1-e\notin\mathfrak p\) annihilates \(r\); hence \(r/1=0\) in the localization. The local ring therefore has zero maximal ideal. If some \(e_n\notin\mathfrak p\), it becomes a unit and \(I_{\mathfrak p}=R_{\mathfrak p}\), so \(M_{\mathfrak p}=0\). If all \(e_n\in\mathfrak p\), each finite-support element of \(I\) becomes zero and \(I_{\mathfrak p}=0\), so \(M_{\mathfrak p}=R_{\mathfrak p}\). This proves stalkwise freeness in all cases and the strict difference from the good-neighborhood conclusion.

Reducedness cannot be omitted altogether: for \(R=k[\epsilon]/(\epsilon^2)\), \(S=R\), \(M=R/(\epsilon)\), the spectrum has one point. Every nonempty basic open inverts a unit, leaving \(M\) unchanged. Its nonzero annihilator \((\epsilon)\) precludes freeness. The good locus is empty, hence is not dense in this nonempty spectrum.

\subsubsection{Reading and propagation record}\label{algebra-proof-027-reading-and-propagation-record}

The two new indexed queries return eight routing records across both corpus layers. All remain unread and unadopted. Their contents include PDF-primary entries and unrelated programme records; no PDF fallback is used. The one query carrying quotation marks still reports literal\_terms\_and behavior and is not described as an exact-phrase search.

Group 672's explicit weights are applied to the original reduced-base substitution, with every binomial coefficient and the leading unit retained. Group 676's generic-domain proof supplies the precise normalization input and preserves its annihilator factors. The original presentation and filtration dependencies at algebra.tex:390--498 were reread. Topology:1391--1419 and Algebra:43890--43895 fix the empty-spectrum convention; no conclusion about the later Algebra section is made beyond these lines.

The finite-complex consequence and the exact product-ring good locus belong in source-linked editorial supplements and the claim ledger. They do not recast the source theorem or its translation. Further combined-results synthesis follows completion of the core correction work.

\hypertarget{algebra_krull_akizuki_notes_20260928}{}
\subsection{Krull--Akizuki: original maps, local modifications and length bounds}\label{algebra-proof-028-krullakizuki-original-maps-local-modifications-and-length-bounds}

Source: Stacks Project authors, authority commit a04446e57ec1fbc252a871afcec7752fb2807b14, algebra.tex:28877--29346, SHA-256 FA8BB92E58A4F78A2BD01B3B6A4A87DE0A0D279F5DD90641B574DD5FBFFFA4F3. All 23 reports in this interval are read. The current interval is exactly the authority with nine previously composed correction operations, MC-STK-ERR-0621 through 0629, and no added environment. Those operations are retained, not proposed again.

These are separate source-linked editorial arguments. Three new minimal operations supply ``by,'' give the local ring to which the local modification lemma applies, and parenthesize the residue field in the associated graded polynomial ring. Neither translation is replaced by these expanded proofs. Four consequence groups are recorded below without claiming novelty.

\subsubsection{The original domination chart}\label{algebra-proof-028-the-original-domination-chart}

At 28903--28919 choose the source valuation ring \(A\subset K\) dominating the original local Noetherian domain \(R\), and its valuation \(v\) in its actual ordered value group. No discreteness or rank-one assumption is made about \(A\). Since \(R\) is not a field, its maximal ideal \(\mathfrak m\) has a nonempty finite minimal generating set \(x_1,\ldots,x_r\), all nonzero. Reorder this same set so \(v(x_r)\) is minimal. Each quotient \(x_i/x_r\) belongs to \(A\), and \(v(x_r)>0\). Thus the specified ring \[
 S=R[x_1/x_r,\ldots,x_{r-1}/x_r]\subset K
\] satisfies \(\mathfrak mS=x_rS\ne S\). It is a Noetherian domain of finite type over \(R\). A prime \(\mathfrak q\) minimal over \(x_rS\) has height at most one by the principal ideal theorem, and at least one because it contains the nonzero element \(x_r\) in a domain. Its contraction contains \(\mathfrak m\), so equals \(\mathfrak m\). Consequently \(R'=S_{\mathfrak q}\subset K\) has dimension one, is essentially of finite type, dominates \(R\), and has the same fraction field \(K\). The last equality follows from \(R\subset R'\subset K=\operatorname{Frac}(R)\).

The existing MC-STK-ERR-0621 supplies the singular verb at 28905. OCC-00562 additionally supplies ``by'' in the original valuation declaration at 28906.

\subsubsection{The four local cases and an explicit annihilator exponent}\label{algebra-proof-028-the-four-local-cases-and-an-explicit-annihilator-exponent}

At 28922--29029 keep Kollár's four alternatives exactly as stated in the source. The source attributes the lemma to János Kollár's ``Variants of normality for Noetherian schemes''; that attribution is retained. No separate reading of that manuscript is claimed here.

Assume the local Noetherian ring \((R,\mathfrak m)\) is not Artinian. The union \(J=\bigcup_{n\ge0}(0:_R\mathfrak m^n)\) stabilizes because \(R\) is Noetherian; hence some power of \(\mathfrak m\) kills \(J\). It is proper, since otherwise a power of \(\mathfrak m\) would be zero, making \(R\) Artinian. If \(J\ne0\), \(R\to R/J\) is a nonisomorphism with the required finite kernel. The quotient has no nonzero element killed by any power of \(\mathfrak m\): if \(\mathfrak m^a r\subset J\) and \(\mathfrak m^b J=0\), then \(\mathfrak m^{a+b}r=0\), so \(r\in J\). In particular \(\mathfrak m\) is not an associated prime of \(R/J\), and this quotient is nonzero.

Now let \(J=0\). By the associated-prime criterion there is a nonzerodivisor \(x\in\mathfrak m\). If \(R/xR\) also has a nonzerodivisor in \(\mathfrak m\), this gives depth at least two. Otherwise choose the original \(y\notin xR\) with \(y\mathfrak m\subset xR\).

Suppose first \(y\mathfrak m\subset x\mathfrak m\). Put \(z=y/x\in R_x\) and retain the map \[
 \varphi:\mathfrak m\longrightarrow\mathfrak m,\qquad t\longmapsto yt/x.
\] Let \(r\) be the number of elements of a chosen minimal generating set of \(\mathfrak m\). Write the image of each generator under \(\varphi\) as an \(R\)-linear combination of these same generators, producing an \(r\)-by-\(r\) matrix \(C\). The adjugate identity for \(TI-C\) gives the monic polynomial \(P(T)=\det(TI-C)\), of degree \(r\), with \(P(\varphi)\mathfrak m=0\). In \(R_x\), \(\varphi\) is multiplication by the actual element \(z\). Since \(x\in\mathfrak m\) is invertible in \(R_x\), \(xP(z)=0\) implies \(P(z)=0\). Thus \[
 R'=R[z]\subset R_x
\] is finite, generated over \(R\) by \(1,z,\ldots,z^{r-1}\). The injection \(R\to R_x\) and \(y\notin xR\) show \(R'\ne R\). Also \(x\) remains a nonzerodivisor on \(R'\), so \(\mathfrak m\notin\operatorname{Ass}_R(R')\) and \(R'\ne0\).

There is an explicit bound stronger than ``a power'': since \(\mathfrak m z\subset R\), \[
 \mathfrak m^i z^i\subset R,\qquad
 \mathfrak m^{r-1}(R'/R)=0.
\] For the first inclusion write a generator of \(\mathfrak m^i\) as a product of \(i\) elements of \(\mathfrak m\), and group each factor with one \(z\). For \(i<r-1\), the remaining factors lie in \(R\). Applying this to the displayed generators proves the second inclusion. Necessarily \(r\ge2\), since a monic degree-one equation would give \(z\in R\). The cokernel is generated by at most \(r-1\) elements over \(R/\mathfrak m^{r-1}\), so also \[
 \operatorname{length}_R(R'/R)
 \le (r-1)\operatorname{length}_R(R/\mathfrak m^{r-1}).
\] This is the first separate consequence group.

If instead \(y\mathfrak m\not\subset x\mathfrak m\), the original argument gives \(t_0\in\mathfrak m\) and a unit \(u\in R\) with \(yt_0=ux\). Keep the unit explicit and set \(t=u^{-1}t_0\), so \(yt=x\). If \(ya=0\), then \(xa=tya=0\), hence \(a=0\). Thus \(y\) is a nonzerodivisor. For each \(t'\in\mathfrak m\), write \(yt'=xs\); then \(y(t'-st)=0\), so \(t'=st\). Hence \(\mathfrak m=(t)\). The ring is not Artinian and the principal ideal theorem gives dimension at most one, so it has dimension one and is regular.

The alternatives are disjoint. A nonzero Artinian local ring has depth zero; every nonzero finite module over it has a simple submodule, hence associated prime \(\mathfrak m\), excluding the fourth alternative. A regular one-dimensional local ring has depth one, excluding depth at least two. In that regular case, a finite target as in alternative four has depth at least one, so the earlier source theorem lemma-regular-mcm-free makes it free. Its generic fibre is one-dimensional, so its rank is one. Multiplication gives an injective map \(R'\to\operatorname{End}_R(R')\cong R\); it is injective because an element annihilating \(1\) is zero, and its image contains every scalar from \(R\), hence is all of \(R\). This map is inverse to \(R\to R'\). For depth at least two, a kernel killed by \(\mathfrak m^N\) is zero, since an \(R\)-regular element in \(\mathfrak m\) acts injectively on it. The depth inequality for the resulting exact sequence gives depth at least one to its cokernel; a nonzero finite module killed by \(\mathfrak m^N\) has a simple submodule and depth zero. The cokernel is therefore zero as well.

MC-STK-ERR-0622 retains exactly the map hypotheses while fixing their relative clauses. MC-STK-ERR-0623 fixes the preposition at the generic point. Neither changes this classification.

\subsubsection{Globalizing the specified local modification}\label{algebra-proof-028-globalizing-the-specified-local-modification}

The report OCC-12323 correctly identifies that the lemma invoked at 29075 is local, while the remark starts with a semilocal domain. The minimal text correction applies it to \((R_{\mathfrak m},\mathfrak mR_{\mathfrak m})\). Here is the actual finite global extension; no local hypothesis is imposed on the original \(R\).

More generally let \(R\) be any Noetherian domain of dimension one, \(K=\operatorname{Frac}(R)\), and \(\mathfrak m\) a maximal ideal. Let \[
 R_{\mathfrak m}\subset B\subset K
\] be a finite overring. Choose integral generators \(b_j\) of \(B\) over \(R_{\mathfrak m}\). In their finitely many monic equations clear coefficient denominators using \(s\in R\setminus\mathfrak m\). Choose \(N\) large enough that all coefficients \(s^{N(d-i)}c_i\) of the equations for \(s^Nb_j\) lie in \(R\). Explicitly an equation \(b^d+\sum_{i<d}c_i b^i=0\) becomes \[
 (s^Nb)^d+\sum_{i<d}s^{N(d-i)}c_i(s^Nb)^i=0.
\] Put \(B_0=R[s^Nb_1,\ldots,s^Nb_t]\subset K\). It is finite over \(R\), and \((B_0)_{\mathfrak m}=B\), since \(s\) is a unit there.

The finite module \(Q=B_0/R\) is torsion, so a nonzero \(d\in R\) kills it: take a common denominator for its finite generating set. Its support lies in the finite set of maximal ideals containing \(d\), because \(R/dR\) is a zero-dimensional Noetherian ring. There is an integer \(a\ge1\) such that \(\mathfrak n^a Q_{\mathfrak n}=0\) at each ideal in that finite support. The Chinese remainder theorem provides \(c\in R\) with \(c\equiv1\pmod{\mathfrak m^a}\) and \(c\equiv0\pmod{\mathfrak n^a}\) for every other ideal in the support. If the support does not contain \(\mathfrak m\), adjoin that ideal to the finite CRT list. Define \[
 B_1=R[c s^Nb_1,\ldots,c s^Nb_t]\subset B_0.
\] Each generator is integral over \(R\), so \(B_1\) is finite. At \(\mathfrak m\), \(c,s\) are units and \((B_1)_{\mathfrak m}=B\). At any other maximal ideal, \(cQ_{\mathfrak n}=0\) if \(Q_{\mathfrak n}\ne0\), so all these generators belong to \(R_{\mathfrak n}\); if \(Q_{\mathfrak n}=0\), the same assertion already holds. Thus \((B_1)_{\mathfrak n}=R_{\mathfrak n}\) at every other maximal ideal. In particular \(B_1/R\) is killed by a power of \(\mathfrak m\), and \(B_1\ne R\) whenever \(B\ne R_{\mathfrak m}\).

For the remark, the local construction in the preceding section over the one-dimensional domain \(R_{\mathfrak m}\) has \(J=0\), depth one, and produces exactly such a strict integral overring \(B\subset K\). Thus the required global extension exists with unchanged fraction field. If \(R\) is semilocal, a finite algebra over it is semilocal: its finitely many fibres over maximal ideals are finite-dimensional algebras over fields, hence have finitely many maximal ideals, and every maximal ideal lies over one of the original maximal ideals by integrality. Dimension remains one by the integral-extension prime-chain theorem.

The remark's optional blowup construction is also finite. The blowup is \[
 X'=\operatorname{Proj}\bigoplus_{n\ge0}\mathfrak m^n,
\] with original affine charts \(R[\mathfrak m/a]\subset K\). This is a projective morphism, is an isomorphism away from \(\mathfrak m\), and makes \(\mathfrak m\mathcal O_{X'}\) invertible. These precise assertions and chart maps are in the same authority's divisors.tex, lemma-blowing-up-affine, lemma-blowing-up-gives-effective-Cartier-divisor, and lemma-blowing-up-projective (bounded reading recorded below). Its fibre over \(\mathfrak m\) is \(\operatorname{Proj}\operatorname{gr}_{\mathfrak mR_{\mathfrak m}}R_{\mathfrak m}\). For each original nonzero chart generator \(a\in\mathfrak m\), the chart \(S_a=R[\mathfrak m/a]\subset K\) has \(\mathfrak mS_a=aS_a\). Localizing this chart over \(R_{\mathfrak m}\) keeps it a subring of \(K\). The finite-length argument proved below, applied to this original submodule of \(K\), says that its quotient by \(\mathfrak m\) is finite dimensional over \(\kappa(\mathfrak m)\). This is exactly the chart of the displayed fibre; localizing the coefficients outside \(\mathfrak m\) does not change that quotient. There are finitely many such charts from a finite generating set of \(\mathfrak m\), so the fibre has finitely many points. All other fibres are singletons because the blowup is an isomorphism away from \(\mathfrak m\). More on Morphisms, lemma-characterize-finite, authority more-morphisms.tex:12543--12589, proves that a proper morphism with finite fibres is finite, so \(X'\) is affine and gives a finite birational overring of \(R\). It is strict at a nonregular \(R_{\mathfrak m}\): if the blowup were an isomorphism there, the original maximal ideal would be invertible, hence principal in that local ring, making it regular of dimension one. This verifies the optional construction without identifying it with every extension produced by the earlier lemma.

\subsubsection{Termination, a numerical bound and the completion obstruction}\label{algebra-proof-028-termination-a-numerical-bound-and-the-completion-obstruction}

The current MC-STK-ERR-0624 removes the unresolved placeholder and specifies completion at each maximal ideal. The underlying implication has a precise source already present later in the authority: algebra.tex, lemma-analytically-unramified-easy, 45807--45887, especially assertion (4). Here is its application and the numerical consequence for the original procedure.

For a Noetherian local domain \(T\) with reduced completion \(\widehat T\), that lemma proves the integral closure \(\overline T\) in \(\operatorname{Frac}(T)\) finite. The actual argument uses faithful flatness \(T\to\widehat T\) and the finite normalization of the reduced complete Noetherian ring \(\widehat T\). The latter is a finite product of the normalizations of \(\widehat T/\mathfrak p_i\) at its finitely many minimal primes; finiteness follows from lemma-Noetherian-complete-local-Nagata, 45754--45778. Since every nonzerodivisor of \(T\) remains one on \(\widehat T\), the map \[
 Q(T)\otimes_T\widehat T\longrightarrow Q(\widehat T)
\] is the injective localization map retaining those same denominators. Flatness embeds \(\overline T\otimes_T\widehat T\) in it. This tensor ring is integral over \(\widehat T\), hence lies in its finite normalization, and is a finite module since \(\widehat T\) is Noetherian. Choose finitely many original elements of \(\overline T\) whose tensors generate it. Their quotient module tensors to zero, so faithful flatness shows it was zero. This proves finiteness over \(T\), with the complete-local Nagata theorem as the explicitly cited source dependency, not a newly proved theorem.

Now retain the original semilocal domain \(R\) and assume each \(\widehat{R_{\mathfrak m}}\) reduced. Integral closure in \(K\) commutes with localization (the earlier source result), so \(\overline R_{\mathfrak m}\) is finite over \(R_{\mathfrak m}\) at each of the finitely many maximal ideals. Choose finitely many elements of \(\overline R\) whose localizations generate at each ideal, and let \(N\) be their combined \(R\)-span. Then \((\overline R/N)_{\mathfrak m}=0\) at every maximal ideal. Any nonzero element of this quotient has a proper annihilator contained in a maximal ideal and survives there, a contradiction. Thus \(N=\overline R\) and the global normalization is finite.

The finite module \(\overline R/R\) is torsion and has finite length, as it is killed by a nonzero denominator and \(R/dR\) is Artinian. Put \[
 \delta=\operatorname{length}_R(\overline R/R).
\] Every finite birational integral modification \(R_i\subset R_{i+1}\) in the original fraction field lies inside \(\overline R\), and \[
 0\longrightarrow R_{i+1}/R_i
 \longrightarrow \overline R/R_i
 \longrightarrow \overline R/R_{i+1}
 \longrightarrow0
\] is the actual sequence of \(R\)-modules. A strict step decreases the last length by at least one. Hence any such procedure has at most \(\delta\) strict steps. If its current semilocal ring is nonregular somewhere, the globalization above supplies a further strict step. It must therefore stop at a regular semilocal ring. That terminal ring is normal by the DVR characterization below, so, being intermediate and integral, equals \(\overline R\). This is a quantitative consequence of the original claim, independent of which admissible strict modifications are chosen.

Conversely, if some completion \(\widehat{R_{\mathfrak m}}\) is nonreduced, no finite birational extension \(B\) of \(R\) can be a regular semilocal ring. If it were, exactness of completion on finite modules would embed \(\widehat{R_{\mathfrak m}}\) into the \(\mathfrak m\)-adic completion of \(B\otimes_R R_{\mathfrak m}\). The latter is the product of the completions of its finitely many DVR localizations. To justify the product, the radical of \(\mathfrak m B\) is the intersection of those finitely many maximal ideals, a power of that radical lies in \(\mathfrak mB\), and the two filtrations are cofinal; CRT then gives the product in the inverse limit. A DVR completion is a domain: each nonzero series in its inverse limit has a least uniformizer order, and the leading graded coefficients multiply nontrivially in the residue field, so orders add. The product is reduced, and so is every subring, contradicting the assumed completion. Thus the source's nontermination implication is valid for every chain of finite birational extensions that would otherwise terminate in a regular ring.

\subsubsection{The two original examples and their exact completions}\label{algebra-proof-028-the-two-original-examples-and-their-exact-completions}

First, at 29031--29047 retain \(R=k[[x,y]]/(y^2)\), written uniquely as \(k[[x]]\oplus k[[x]]y\). In \(R_x\) put \(z_n=y/x^n\) and \[
 R_n=k[[x]]\oplus k[[x]]z_n,\qquad z_n^2=0.
\] The original injection sends \(a(x)+b(x)y\) to \(a(x)+x^nb(x)z_n\). The transition \(R_n\to R_{n+1}\) sends \(z_n\) to \(xz_{n+1}\). Every \(R_n\) is finite over \(R\), and \[
 R_n/R\cong k[[x]]/(x^n)\cdot z_n,\qquad
 \operatorname{Ann}_R(R_n/R)=(x^n,y),\qquad
 \operatorname{length}_R(R_n/R)=n.
\] These formulas follow by comparing the two displayed coordinates; \(y\) kills the quotient and \(a(x)\) kills it precisely when divisible by \(x^n\). Each inclusion is strict. Multiplication by \(x\) is injective, while the nilradical is the square-zero \(k[[x]]y\), so the original ring has dimension one and depth one, as stated.

Next retain exactly the source's field \(k\) of characteristic \(p>0\) with \([k:k^p]\) infinite, and \[
 A=\left\{\sum_{i\ge0}a_i x^i:
       [\,k^p(a_0,a_1,\ldots):k^p\,]<\infty\right\}.
\] It is the union of \(F[[x]]\) over intermediate fields \(k^p\subset F\subset k\) finite over \(k^p\). The compositum of two such fields is still finite, so this is a ring. For each nonzero member, remove its exact initial power \(x^n\); the remaining unit and its inverse both have coefficients in the same \(F\). Thus \(A\) has uniformizer \(x\), residue field \(k\), and every nonzero ideal is generated by a power of \(x\), by choosing its least order. It is a DVR. Truncation gives \(A/x^nA=k[x]/(x^n)\): all finitely many coefficients belong to a finite extension of \(k^p\), and divisibility by \(x^n\) does not change the field of coefficients. Hence \(\widehat A=k[[x]]\). Its fraction field is \[
 K=\bigcup_F F((x)),
\] so \(K\cap k[[x]]=A\). Every Laurent series has its \(p\)-th power in \(k^p((x^p))\subset K\); therefore \(k((x))/K\) is purely inseparable of exponent one.

For completeness, this extension has infinite degree, as asserted in the source. Choose an infinite \(p\)-basis of \(k/k^p\), and for each prescribed finite \(r\) choose disjoint sequences \(t_{j,n}\) from it, \(1\le j\le r,n\ge1\). Set \(f_j=\sum_{n\ge1}t_{j,n}x^n\). Then \(f_j^p\in K\). If a nonzero polynomial relation of degrees less than \(p\) in each \(f_j\) existed over \(K\), choose one of least total degree. Its finitely many coefficients belong to \(F((x))\) for a field \(F\) contained in \(k^p\) with finitely many \(p\)-basis elements adjoined. For a variable occurring with positive exponent, choose a corresponding \(t_{j,n}\) outside that finite basis set. The coefficientwise \(p\)-basis derivation \(\partial/\partial t_{j,n}\) on Laurent series kills all relation coefficients and all other \(f_i\), and sends \(f_j\) to \(x^n\). Dividing the differentiated relation by the nonzero \(x^n\) gives its nonzero formal partial derivative as a relation of smaller total degree. It is nonzero because all positive exponents are less than \(p\). This is impossible; a constant relation is also impossible. The derivation used here exists by defining it on the \(p\)-basis monomials of exponents \(0,\ldots,p-1\), extending over \(k^p\), and differentiating quotients; on series the coefficient of each product involves a finite sum, so the Leibniz rule remains valid. Therefore \([K(f_1,\ldots,f_r):K]=p^r\) for every \(r\), proving infinite degree.

Now take the source's arbitrary \(f\in k[[x]]\setminus A\), not merely one of these auxiliary choices. Write \(a=f^p\in k^p[[x^p]]\subset A\), and \(a_0=f(0)\in k\). Since \(f\notin K\) but \(f^p\in K\), its minimal polynomial over \(K\) is \(T^p-a\). Thus the original domain \(R=A[f]\) has the exact presentation \[
 R=A[T]/(T^p-a)
\] and free \(A\)-basis \(1,T,\ldots,T^{p-1}\). It is Noetherian and integral of dimension one over \(A\). Modulo \(x\) the equation is \((T-a_0)^p\), so it has exactly one maximal ideal, \(\mathfrak m=(x,T-a_0)\). All maximal ideals lie over \((x)\) by integrality. Since \(\mathfrak m^p\subset xR\), the maximal-adic and \(x\)-adic filtrations are cofinal. Taking inverse limits in the displayed finite free \(A\)-basis gives \[
 \widehat R=k[[x]][T]/(T^p-f^p)
       \ \cong\ k[[x]][\epsilon]/(\epsilon^p),\qquad T\longmapsto f+\epsilon.
\] The inverse sends \(\epsilon\) to \(T-f\). Thus the original generator \(f\in R\) maps to \(f+\epsilon\) in this presentation of its completion; the subsequent map to \(k[[x]]\) kills \(\epsilon\). The element \(\epsilon^{p-1}\) is nonzero in the displayed free basis, while \(\epsilon^p=0\). This proves the source's nonreduced completion, including its precise nilpotence.

There is a further exact consequence. With \(L=K(f)\), the integral closure of \(R\) is \[
 V=L\cap k[[x]].
\] Indeed \(z\in V\) has \(z^p\in K\cap k[[x]]=A\), so it is integral over \(A\), hence over \(R\). Conversely an element integral over \(A\) cannot have negative \(x\)-order: in a monic equation its highest power would have strictly smaller valuation than every other term. Integrality over \(R\) implies integrality over \(A\) by transitivity. These prove both inclusions. The ring \(V\) is a DVR with the same \(x\), residue field \(k\), fraction field \(L\), and completion \(k[[x]]\); each assertion follows from its actual order and the inclusion \(A\subset V\). It is not finite over \(R\). Otherwise it would be finite torsion-free over the DVR \(A\), hence free of rank \([L:K]=p\). Completion of this finite free module would have rank \(p\) over \(k[[x]]\), whereas the actual completion of \(V\) is \(k[[x]]\), of rank one. This contradiction proves nonfiniteness without confusing the Noetherian conclusion of Krull--Akizuki with finiteness of normalization.

\subsubsection{Uniformizers and the original discrete valuation}\label{algebra-proof-028-uniformizers-and-the-original-discrete-valuation}

At 29103--29158 the DVR characterization retains all five conditions. MC-STK-ERR-0625 makes the exponent range explicit; the proposed new operation at 29135 writes \((A/\mathfrak m)[T]\). The direct-sum degree is already specified by ``for all \(n\ge0\)'' and does not need a second correction.

For the substantive regular-local implication, keep \(\mathfrak m=(\pi)\), with \(\pi\ne0\), in the original Noetherian local domain. A nonzero \(a\in A\) cannot lie in every \(\pi^nA\): otherwise write \(a=\pi^n a_n\), and use cancellation to obtain \(a_n=\pi a_{n+1}\). The ascending chain \((a_n)\) stabilizes, giving \(a_{n+1}=b a_n\); then \((1-\pi b)a_n=0\), and \(1-\pi b\) is a unit, a contradiction. There is therefore a unique largest \(n\ge0\) with \(a\in\pi^nA\). Write \(a=u\pi^n\); \(u\notin\mathfrak m\) is a unit. If \(u\pi^n=v\pi^m\), unequal exponents would make a unit a multiple of \(\pi\), so \(n=m\) and cancellation gives \(u=v\).

The map \((A/\mathfrak m)[T]\to\bigoplus_{n\ge0}\mathfrak m^n/\mathfrak m^{n+1}\) sends \(\overline cT^n\) to \(c\pi^n\) modulo \(\pi^{n+1}\). It is well-defined, multiplicative and onto; each homogeneous kernel is zero by cancellation and the fact that \(\pi\) is not a unit. Hence it is an isomorphism. The actual order satisfies \(v(ab)=v(a)+v(b)\); the formula \(v(a/b)=v(a)-v(b)\) is well-defined by cross-multiplication. The original ring is exactly the nonnegative part, giving the valuation ring and PID characterizations. Conversely the earlier source result for Noetherian valuation rings gives discreteness unless the ring is a field. The finite one-dimensional modification argument shows a normal Noetherian local domain of dimension one must be regular: its finite target has no torsion, since all nonzero torsion would be supported at \(\mathfrak m\) and have associated prime \(\mathfrak m\); its generic fibre is the original fraction field, so it embeds there, and normality forces equality.

\subsubsection{Length inequalities with the original ambient vector space retained}\label{algebra-proof-028-length-inequalities-with-the-original-ambient-vector-space-retained}

At 29170--29290 keep \(M\subset K^{\oplus r}\) and put \(\rho=\dim_K(KM)\le r\). This defines the actual span inside the given ambient coordinates; it does not replace that ambient space. If \(x\) is a nonzero unit, all asserted lengths are zero. Otherwise \(R/xR\) is an Artinian local ring because its only prime is the maximal ideal.

If \(M\) is finite, choose original elements \(m_1,\ldots,m_\rho\in M\) which form a \(K\)-basis of \(KM\), and let \(L=\bigoplus Rm_i\subset M\). For each remaining member of a finite generating set of \(M\), clear the coefficients in its expression in that basis; one nonzero product of the denominators kills \(C=M/L\). Thus \(C\) is a finite torsion module supported at the maximal ideal and has finite length. Multiplication by \(x\) is injective on both \(L\) and \(M\), since they remain actual submodules of \(K^r\). The original inclusion and quotient maps give, by their multiplication-by-\(x\) diagram, the exact sequence \[
 0\to C[x]\to L/xL\to M/xM\to C/xC\to0.
\] The first map sends a class represented by \(m\in M\) with \(xm\in L\) to \(xm\bmod xL\); changing \(m\) by an element of \(L\) changes this by \(xL\), and injectivity follows from cancellation in \(M\). The remaining maps are induced by inclusion and quotient. Since the kernel and cokernel of multiplication by \(x\) on the finite-length module \(C\) have equal length, the sequence proves exactly \[
 \operatorname{length}_R(M/xM)=\rho\operatorname{length}_R(R/xR).
\]

This also verifies the source's original denominator-clearing and sandwich argument without silently changing coordinates. A scalar \(d\ne0\) clearing the coordinates of a finite generating set gives an isomorphism \(M\to dM\), \(m\mapsto dm\), with inverse division by the same \(d\), including on quotients modulo \(x\). When the span is all of \(K^r\), the quotient \(C'=R^r/dM\) has finite length, hence \(x^cR^r\subset dM\) for some \(c\). For \(n\ge c\), the actual inclusions \[
 x^{n+c}R^r\subset x^n dM\subset x^nR^r,\qquad
 x^cR^r\subset dM\subset R^r
\] bound the length of \(dM/x^ndM\) between \(r(n-c)\ell\) and \(r(n+c)\ell\), where \(\ell=\operatorname{length}_R(R/xR)\). For example the upper bound follows by embedding in \(R^r/x^{n+c}R^r\); for the lower bound, the image of \(x^cR^r\) in \(dM/x^ndM\) has kernel contained in \(x^nR^r\). Successive multiplication by \(x\) identifies each \(x^iM/x^{i+1}M\) with \(M/xM\), so dividing the inequalities by \(n\) and letting \(n\) grow proves equality, exactly as in the source.

For arbitrary \(M\), retain any finite chain of \(k\) strict inclusions in \(M/xM\) and lift the source witnesses \(m_i\) of its successive differences to \(M\). The submodule \(M'=\sum_iRm_i\subset M\) is finite. The inverse images of that chain under \(M'/xM'\to M/xM\) remain strict because of these witnesses. No injectivity of this quotient map is assumed. The finite case gives \(k\le\operatorname{length}_R(M'/xM')\le\rho\ell\). If \(M/xM\) had length greater than \(\rho\ell\), such a chain with \(k>\rho\ell\) would exist. Hence its length is at most \(\rho\ell\). This treats \(\rho=0\), when \(M=0\). MC-STK-ERR-0626 fixes ``vectors space''; 0627 restores the missing inclusion symbol while allowing \(N_0=0\).

For a nonlocal one-dimensional Noetherian domain, the quotient \(R/xR\) is Artinian with finitely many maximal ideals \(\mathfrak m_i\). Its orthogonal idempotents split every \(R/xR\)-module, including \(M/xM\), into the product of its actual localizations at those ideals. Finite-length modules supported at one \(\mathfrak m_i\) have identical lengths over \(R\) and \(R_{\mathfrak m_i}\), since the same simple factors are the residue field. Summing the local formulas gives the stronger global bound \[
 \operatorname{length}_R(M/xM)
 \le \rho\operatorname{length}_R(R/xR),
\] with equality if \(M\) is finite. This is stronger than the global finiteness statement in the source, with the original module and all local contributions retained.

\subsubsection{Residue fields, ramification and Krull--Akizuki}\label{algebra-proof-028-residue-fields-ramification-and-krullakizuki}

At 29240--29264 let \(R\to S\) and \(K\subset L\) be the source domains and finite fraction-field extension, \(n=[L:K]\). A \(K\)-basis of the original \(L\) supplies an explicit linear isomorphism with \(K^n\), so the preceding inequality applies to the original \(R\)-submodule \(S\subset L\). The ring \(B=S/\mathfrak mS\), a quotient of \(S/xS\), has \[
 \dim_{\kappa(\mathfrak m)}B
 \le n\,\operatorname{length}_R(R/xR).
\] It may be zero. If nonzero it is Artinian and splits into its local factors \(B_i\). Their composition factors are their residue fields, so the exact formula is \[
 \dim_{\kappa(\mathfrak m)}B
 =\sum_i \operatorname{length}_{B_i}(B_i)
          [\kappa(\mathfrak n_i):\kappa(\mathfrak m)].
\] Every nilpotent multiplicity is retained. This proves the source's three assertions and their quantitative bound. MC-STK-ERR-0628 corrects only the plural verb.

If \(R\) is a DVR and \(S\) its integral closure in \(L\), Krull--Akizuki below makes \(S\) Noetherian, and the residue-field assertion makes it semilocal. Each \(S_{\mathfrak n_i}\) is a DVR. With the original uniformizer \(\pi\) of \(R\), put \(e_i=v_i(\pi)\), \(f_i=[\kappa(\mathfrak n_i):\kappa(\mathfrak m)]\). The filtration by the first \(e_i\) powers of a local uniformizer shows \[
 \sum_i e_i f_i
 =\operatorname{length}_R(S/\pi S)\le [L:K].
\] Equality holds if \(S\) is finite over \(R\), by the finite-module length equality. It need not hold merely because \(S\) is Noetherian: the source's characteristic-\(p\) example above has \(A\subset V\), \([K(f):K]=p\), and a single prime with \(e=f=1\), giving \(1<p\). Thus the completed example sharpens the precise limit of the length consequence.

For the Krull--Akizuki assertion at 29292--29309, let \(R\subset A\subset L\) and \(0\ne I\subset A\). Choose \(0\ne a\in I\). Its algebraicity over \(K\) gives an equation \(c_0+c_1a+\cdots+c_da^d=0\) with \(c_i\in R\), \(c_0\ne0\), by starting with its minimal polynomial over \(K\) and clearing every denominator. Then \(c_0=-a(c_1+\cdots+c_da^{d-1})\in I\cap R\). Put \(x=c_0\). The global length result makes \(A/xA\) a finite-length \(R\)-module, so \(I/xA\) has finitely many \(R\)-generators. Lift these to \(I\), and add \(x\). They generate \(I\) as an \(A\)-ideal: subtract the lifted \(R\)-linear combination to leave an element of \(xA\). The zero ideal is generated by the empty set. Thus every ideal of \(A\) is finitely generated. This proof asserts Noetherianity, not module finiteness over \(R\).

\subsubsection{A DVR in the original finitely generated extension}\label{algebra-proof-028-a-dvr-in-the-original-finitely-generated-extension}

For 29311--29338 retain the original local Noetherian domain \(R\), \(K\), and finitely generated extension \(L/K\). If necessary choose actual elements \(x_1,\ldots,x_r\in L\) forming a transcendence basis and let \[
 T=R[x_1,\ldots,x_r]_{(\mathfrak m_R,x_1,\ldots,x_r)}.
\] The polynomial ring is the specified subring of \(L\), its displayed ideal is maximal with residue field \(R/\mathfrak m_R\), and its contraction is \(\mathfrak m_R\). Hence \(T\) dominates \(R\), is Noetherian, and has fraction field \(K(x_1,\ldots,x_r)\); \(L\) is finite over this field. The chart construction in the first section gives a one-dimensional local domain \(T'\) inside that same fraction field dominating \(T\).

Let \(C\) be the integral closure of \(T'\) in the original \(L\). Its fraction field is \(L\): for any \(z\in L\), clear a monic equation's coefficient denominators by multiplying \(z\) by a nonzero \(d\in T'\); the equation for \(dz\) has coefficients \(d^{m-i}c_i\), all in \(T'\) for a sufficiently divisible \(d\). Thus \(dz\in C\) and \(z=(dz)/d\). Krull--Akizuki makes \(C\) Noetherian, and lying over gives a prime \(\mathfrak q\) over the maximal ideal of \(T'\). Integral prime-chain comparison gives dimension one, and the localization \(C_{\mathfrak q}\) is a normal one-dimensional Noetherian local domain, hence a DVR by the characterization. Its fraction field is the original \(L\). Contractions through the three local maps show it dominates the original \(R\). MC-STK-ERR-0629 restores the article without changing this construction.

\subsubsection{Reading and consequence record}\label{algebra-proof-028-reading-and-consequence-record}

The indexed query ``Krull Akizuki'' returned two research-layer routing records, one PDF-primary and one original-TeX Bhatt survey; the local-unpublished layer returned no matches. Both hits remain unread and unadopted. The arguments above use the read original Stacks TeX and the identified existing source results; no PDF fallback or novelty claim is made.

Bounded later dependencies: algebra.tex:45732--45887, including the complete-local Nagata proof and the analytically unramified normalization lemma; more-morphisms.tex:12543--12589 for the proper finite-fibre criterion; divisors.tex:8321--8341,8382--8435,8622--8648 for the stated blowup charts, exceptional ideal and projectivity. The foundational theorems cited within those source proofs are identified there; this is not a claim to have reread their entire dependency closure. Later report hits are not adjudicated by this bounded dependency reading.

The four editorial consequence groups record the explicit local annihilator exponent, globalization with the termination bound, global length and ramification bounds, and the precise completion/normalization consequences of the two source examples. These implications are propagated together: the positive-characteristic example proves the strict case of the ramification inequality and the limit of module-finiteness claims. Wider synthesis follows completion of the core task. Next source: 29347, Factorization; bounded lookahead through 29370 is read but unadjudicated.

\hypertarget{algebra_factorization_notes_20260928}{}
\subsection{Factorization with the original units, primes and ideal maps}\label{algebra-proof-029-factorization-with-the-original-units-primes-and-ideal-maps}

Authority: Stacks Project authors, commit a04446e57ec1fbc252a871afcec7752fb2807b14, algebra.tex:29347--29778, SHA-256 FA8BB92E58A4F78A2BD01B3B6A4A87DE0A0D279F5DD90641B574DD5FBFFFA4F3. All 30 reports are read. The current interval equals the authority plus MC-STK-ERR-0630 through 0642, comprising fourteen operations. There is no added environment.

These full arguments and three consequence groups are editorial material, not replacement translations. The two new unit-factor operations below complete a residual case in the already composed MC-STK-ERR-0633 while preserving its valid zero-case correction. Reported omissions of routine steps and the source's expressly omitted uniqueness proof are completed here without automatically replacing the source proofs. No novelty claim is made.

\subsubsection{Prime elements, associates and the rejected factorization reports}\label{algebra-proof-029-prime-elements-associates-and-the-rejected-factorization-reports}

At 29363 the existing MC-STK-ERR-0630 excludes zero from prime elements. This is mathematically necessary: \((0)\) is prime in every domain, but including zero in the generating primes would make the multiplicative set in Nagata's criterion contain zero and would invalidate the later fraction-field identification. A prime element is also a nonunit, since its principal ideal is proper. A nonzero prime element \(p\) is irreducible: if \(p=ab\), primality implies, say, \(a=pc\); cancellation gives \(1=cb\), so \(b\) is a unit.

For equality \((x)=(y)\), retain \(x=fy,y=gx\) with \(f,g\in R\). If \(x=0\), also \(y=0\), and the unit \(1\) witnesses association. Otherwise \(x=fgx\) gives \(fg=1\). These are precisely the valid current corrections MC-STK-ERR-0631 and 0632. The converse retains both maps \(x=uy\) and \(y=u^{-1}x\).

The proposed deletions in OCC-00571 and OCC-12334 are rejected as nondefects. They use only reducibility of \(x\), but the source at 29396 assumes the stronger fact that \(x\) has no finite factorization into irreducibles. Starting from this same \(x\), split a factor which has no irreducible factorization into two nonzero nonunits. At least one new factor again has no irreducible factorization, since otherwise concatenating the two finite factorizations would factor the selected factor. Repeat until the actual equality \[
 x=a b c d
\] has four nonzero nonunit factors. Put \(y=ab,z=cd\). A product of two nonunits in a commutative ring cannot be a unit: an inverse to their product gives an inverse to each factor. Neither \(y\) nor \(z\) is irreducible, since each has the displayed factorization into two nonunits. Thus the source's stronger assertion is true.

At least one of these \(y,z\) still lacks an irreducible factorization. Choose that factor as \(x_2\), and retain the other as \(q_1\), so \(x_1=q_1x_2\), with \(q_1\) a nonunit. Repeating gives \(x_i=q_i x_{i+1}\) with nonzero \(x_i\) and nonunit \(q_i\). If \((x_i)=(x_{i+1})\), the associates result would force \(q_i\) to be a unit, by cancellation. Hence the principal ideals form the strict ascending chain in the source, contradicting ACC. This verifies the original argument without weakening its assertion.

\subsubsection{The UFD criterion with every unit factor retained}\label{algebra-proof-029-the-ufd-criterion-with-every-unit-factor-retained}

At 29435 retain the valid current disposal of \(a=0\) or \(b=0\) in \(ab=cx\), where \(x\) is irreducible. Then \(a,b,c\ne0\). An arbitrary unit is not the empty product \(1\). The current phrase ``allowing the empty product for a unit,'' followed by unit-free equations for \(a,b,c\), therefore needs the explicit scalar factors \[
 a=u_a a_1\cdots a_n,\quad
 b=u_b b_1\cdots b_m,\quad
 c=u_c c_1\cdots c_r,\qquad u_a,u_b,u_c\in R^\times,
\] and the actual comparison \[
 u_a u_b a_1\cdots a_n b_1\cdots b_m
       =u_c c_1\cdots c_r x.
\] The two minimal new operations modify precisely these two displayed source lines; neither overlaps MC-STK-ERR-0633. Empty products now apply only to the lists of irreducibles, while the original units remain in the equations.

Uniqueness for these equations follows from the source UFD definition. A product of nonunits cannot be a unit, so if either irreducible list is empty, both must be. Otherwise multiply the first irreducible on each side by its displayed unit. This maps that factor to an associate and has the explicit inverse multiplication by the inverse unit; the other factors and both original unit multipliers stay recorded. The resulting lists are factorizations of the same original nonunit and the UFD axiom matches their associates. Applied to the displayed equation, \(x\) is associate to one of \(a_1,\ldots,a_n,b_1,\ldots,b_m\). Hence it divides \(a\) or \(b\). The restored full range is already MC-STK-ERR-0642.

Conversely assume every irreducible is prime. Prove uniqueness also with an arbitrary recorded unit multiplier \(u\): \[
 a_1\cdots a_m=u b_1\cdots b_n.
\] If one list is empty, the other must be empty. If both are nonempty, primality of \(a_1\) and invertibility of \(u\) imply \(a_1\mid b_j\) for some \(j\). Permute that original list to put this factor first and write \(b_1=a_1v\). Irreducibility of \(b_1\) makes \(v\) a unit, and cancellation gives \[
 a_2\cdots a_m=(uv)b_2\cdots b_n.
\] The induction retains the full unit product \(uv\) and the original endpoints. If either remaining list is empty, the other is empty. Otherwise apply induction on \(m+n\). The final association for the first factor retains \(v\), and the permutation extends that of the remaining indices. This verifies MC-STK-ERR-0634, including its endpoint repair and one-factor cases, without discarding its unit.

For the height-one criterion at 29458--29482, a height-one prime in a UFD contains a nonzero element \(x\), hence an irreducible factor \(p\). The prime \((p)\) is nonzero and lies in that height-one prime, so they coincide. This implication does not require Noetherianity. Conversely the stated Noetherian hypothesis gives ACC and finite irreducible factorizations. A prime minimal over the nonzero principal ideal of an irreducible has height one by the principal ideal theorem and the domain hypothesis. If it is generated by \(y\), write \(x=yz\); irreducibility makes \(z\) a unit, so the original \((x)\) itself is prime.

\subsubsection{Nagata cancellation with repeated factors}\label{algebra-proof-029-nagata-cancellation-with-repeated-factors}

At 29484--29540 retain the source's citation to Nagata-UFD, Lemma 2, and its actual multiplicative set \(S\), generated by nonzero prime elements. Hence \(0\notin S\), and the localization embeds in the original fraction field.

For \(\alpha=a/s,\beta=b/s'\) with \(x=\alpha\beta\), injectivity of that embedding gives the exact equality \[
 p_1\cdots p_r x=ab,\qquad ss'=p_1\cdots p_r.
\] Cancel these particular prime factors successively, including repetitions. At step \(i\), the surviving equation is \(p_i\cdots p_r x=a_{i-1}b_{i-1}\). Primality of \(p_i\) makes it divide one of the two current factors; divide that factor by this \(p_i\), leave the other unchanged, and cancel the same nonzero \(p_i\) from both sides. After \(r\) steps, \[
 x=a_r b_r,\quad a=s''a_r,\quad b=s'''b_r,\quad s''s'''=ss',
\] where \(s''\) and \(s'''\) are the products of the selected occurrences, not products of distinct primes. Thus \(\alpha=(s''/s)a_r\) and \(\beta=(s'''/s')b_r\); both displayed scalar fractions are units of \(S^{-1}A\). Irreducibility of the original \(x\) makes one of \(a_r,b_r\) a unit, proving the first assertion when the image of \(x\) is not already a unit.

If \(x\) is prime, the quotient \((S^{-1}A)/(x)\cong S^{-1}(A/(x))\) is a domain or zero, proving that its image is prime or a unit. The quotient map sends \(a/s\) to \(\overline a/\overline s\); its kernel consists precisely of fractions with an \(S\)-multiple of the numerator in \((x)\), which is the localized ideal.

For the reverse unit case \(yx=p_1\cdots p_r\), if any current \(p_i\) divides \(x\), irreducibility gives an association, so \(x\) is prime. Otherwise that occurrence divides the current \(y\); replace \(y\) by the quotient and cancel \(p_i\). If all occurrences are cancelled this way, the surviving equality is \(y_r x=1\), contradicting that \(x\) is a nonunit. This includes \(r=0\) directly. In the reverse prime-image case, \(ab\in(x)\) implies \(sa=xy\) or \(sb=xy\) for original \(s\in S,y\in A\). In the first case run the same process on \(p_1\cdots p_r a=xy\). A prime dividing \(x\) again proves \(x\) prime. Otherwise divide the successive current \(y\)'s and cancel every occurrence; the last equality is \(a=xy_r\), so \(a\in(x)\). The second case gives \(b\in(x)\). If \(a\) or \(b\) is zero the conclusion is immediate.

This supplies exactly the omitted multiplicity tracking in OCC-12348. The theorem is true; the full successive equalities are separate clarification rather than a wholesale replacement of its proof. The final equivalence follows because a UFD localizes to a UFD by its finite prime factorizations, with inverted primes becoming units and all others remaining prime; in the other direction atomicity of \(A\) and the two proved assertions make every original irreducible prime.

\subsubsection{ACC, polynomial rings and the full normality calculation}\label{algebra-proof-029-acc-polynomial-rings-and-the-full-normality-calculation}

For UFD ACC, MC-STK-ERR-0635 handles an all-zero chain or its initial zero segment. If the first nonzero generator is a unit, all following ideals are already \(R\). Otherwise retain \[
 a_1=u\prod_{j=1}^r p_j^{e_j},\qquad u\in R^\times.
\] Every later generator divides \(a_1\), so is a unit times \(\prod_jp_j^{c_j}\), with \(0\le c_j\le e_j\). Equality of exponent vectors is equality of the principal ideals, and there are only finitely many vectors. This proves stabilization while preserving the possible original unit even when it is absent from the source's displayed notation.

For polynomial ACC, MC-STK-ERR-0636 similarly disposes of zero polynomials before degrees are taken. In the nonzero tail write the actual equality \(f_i=h_i f_{i+1}\). Degrees are nonincreasing and eventually constant, so then \(h_i\in R\). The leading coefficients satisfy \(a_i=h_i a_{i+1}\). ACC in \(R\) makes \((a_i)=(a_{i+1})\) eventually; cancellation in the domain makes \(h_i\) a unit. Hence the polynomial ideals stabilize. Repetition proves the result for finitely many variables. For arbitrarily many variables, every divisor of the fixed first nonzero polynomial uses only its finite set of variables: the degree in any other variable is zero, and degrees in that variable add in a domain. Thus the entire tail and the quotient factors lie in a fixed finite-variable ring, to which the proved argument applies.

MC-STK-ERR-0637 correctly restricts both factorization assertions to nonzero nonunits. For a UFD \(R\), \(S\) generated by its nonzero primes gives \(S^{-1}R=K=\operatorname{Frac}(R)\): if \(b=u\prod p_i^{e_i}\ne0\), then \(1/b=u^{-1}/\prod p_i^{e_i}\) is the required fraction with its original unit. Every such prime remains prime in \(R[X]\), since the quotient is \((R/(p))[X]\), a domain. More generally, a prime of a finite-variable polynomial ring remains prime after adjoining any additional variables: its quotient is the polynomial ring in those additional variables over the original quotient domain. Thus factorizations obtained in a finite-variable subring remain prime factorizations in the full polynomial ring. The coefficientwise isomorphism \(S^{-1}R[X]\to K[X]\) sends \(g/s\) to its actual coefficient fractions. To invert it, take the product of finitely many coefficient denominators; equality is verified by multiplying both numerators by their common denominator. Over a field, division by the nonzero leading coefficient yields the Euclidean division algorithm, hence the polynomial PID and UFD property. Nagata's criterion and polynomial ACC prove the source result, including arbitrary variable sets by the finite-support argument.

For normality, MC-STK-ERR-0638 separates \(x=0\). Retain the original nonzero fraction and equation \[
 x=u p_1^{e_1}\cdots p_r^{e_r},\quad
 x^d-\sum_{i=1}^d a_i x^{d-i}=0,
\] where every original exponent is an integer, \(u\) is a unit and the \(p_j\)'s are pairwise nonassociate irreducibles. If \(r=0\), \(x=u\in R\). If \(e_1<0\), choose an integer \(N\ge\max(\{0\}\cup\{-d e_j:2\le j\le r\})\). Multiplying the equation exactly as in the source gives \[
 u^d\prod_{j=2}^r p_j^{de_j+N}
 =\sum_{i=1}^d a_i u^{d-i}p_1^{-i e_1}
                  \prod_{j=2}^r p_j^{(d-i)e_j+N}.
\] All exponents on the right are nonnegative, and the exponent of \(p_1\) is positive for every summand. The left has no \(p_1\) factor and its coefficient \(u^d\) is a unit. Primality of \(p_1\) gives a contradiction. Thus every \(e_j\ge0\), proving \(x\in R\) with all constants and summands retained.

\subsubsection{Explicit dual maps for a product of ideals}\label{algebra-proof-029-explicit-dual-maps-for-a-product-of-ideals}

At 29660--29683 let \(I,J\subset A\) and \(IJ=fA\) with the original nonzerodivisor \(f\). Choose the original finite expression \(f=\sum_{i=1}^n x_i y_i\), \(x_i\in I,y_i\in J\). For each \(z\in I\), there is a unique \(\lambda_i(z)\in A\) with \[
 f\lambda_i(z)=y_i z.
\] Uniqueness uses precisely that \(f\) is a nonzerodivisor; it also proves \(A\)-linearity. Cancelling \(f\) from \(f\sum_i x_i\lambda_i(z)=\sum_i x_i y_i z=fz\) gives \(\sum_i x_i\lambda_i(z)=z\). This is an explicit finite dual frame, proving finite generation and projectivity of \(I\); exchanging \(I,J\) proves the same for \(J\).

To prove rank one and retain the source local choices, write \(x_i y_i=a_i f\); then \(\sum a_i=1\). On \(A_{a_i}\), keep the original \(x_i\) and set \(y'_i=y_i/a_i\). Now \(x_i y'_i=f\); both factors are nonzerodivisors there. For \(z\in I_{a_i}\), write \(z y'_i=b f=b x_i y'_i\); cancellation gives \(z=b x_i\). Hence \(I_{a_i}=x_i A_{a_i}\), freely of rank one, with inverse map \(z\mapsto b\). The symmetric map identifies \(J_{a_i}=y'_i A_{a_i}\). The units \(a_i\) have not been suppressed.

There are stronger canonical identifications. Inside \(A_f\) define \(J/f=\{y/f:y\in J\}\). Then \[
 J/f\xrightarrow{\;\cong\;}\operatorname{Hom}_A(I,A),
 \qquad y/f\longmapsto(z\mapsto yz/f).
\] For \(\phi:I\to A\) the inverse is \(\sum_i\phi(x_i)y_i/f\). Indeed its product with \(z\) is \(\sum_i\phi(x_i)\lambda_i(z)=\phi(z)\). If \(tI=0\) with \(t\in J/f\), then \(tf=0\) since \(f=\sum x_i y_i\); invertibility of \(f\) in \(A_f\) gives \(t=0\), proving injectivity. The actual multiplication map \[
 I\otimes_A J\xrightarrow{\;\cong\;}fA
\] has inverse \(af\mapsto a\sum_i x_i\otimes y_i\). Its composite on \(z\otimes w\) is \(\sum_i x_i\otimes y_i zw/f=\sum_i x_i\lambda_i(z)\otimes w=z\otimes w\), and its composite on \(af\) is \(af\).

This extends to an arbitrary invertible ideal \(L=IJ\), without a chosen global generator: on a finite basic-open cover trivializing \(L\), its ideal generator is a nonzerodivisor, because the map from the local ring to this free rank-one ideal is injective. The preceding construction makes \(I,J\) finite locally free of rank one on that cover. Finiteness over the whole ring follows by taking numerators for the finitely many local generators on the finite cover; the quotient by their span vanishes on the cover and hence is zero. The multiplication isomorphism is local and therefore global, and it induces the canonical identification \[
 I^\vee\cong J\otimes_A L^\vee.
\] All maps are contractions of the same multiplication pairing, so agree on overlaps. Conversely if \(I,J\) are invertible ideals, locally they have nonzerodivisor generators \(x,y\), their product has generator \(xy\), and the same multiplication maps show \(IJ\) invertible. This proves the slogan with its actual module morphisms over arbitrary commutative rings, including rings with zero divisors.

\subsubsection{Dedekind domains and the full uniqueness proof}\label{algebra-proof-029-dedekind-domains-and-the-full-uniqueness-proof}

At 29637--29658, the unit ideal has the empty prime-ideal product \(R\). This conventional empty-product case is made explicit here for OCC-12349. In a PID, a nonzero ideal \(aR\) has \(a=u p_1\cdots p_r\) and \[
 aR=(p_1)\cdots(p_r);
\] the unit acts by the bijection \(R\to R\), \(b\mapsto ub\), whose inverse is multiplication by \(u^{-1}\). Each \((p_i)\) is a prime ideal. For a unit \(a\) the ideal is \(R\) and the prime list is empty. This states the precise prime-ideal map requested by OCC-12341; the source's contextual word ``primes'' does not assert a different object. OCC-00574 supplies the missing colon introducing the equivalences.

Assume the original prime-ideal factorization definition, without presupposing Noetherianity. MC-STK-ERR-0639 treats \((0)\) separately as already finitely generated. For a nonzero prime \(\mathfrak p\), uniqueness implies \(\mathfrak p\ne\mathfrak p^2\), since otherwise a product of one nonzero prime equals a product of two. Choose the original \(x\in\mathfrak p\setminus\mathfrak p^2\). For any \(y\in\mathfrak p\), factor the nonzero proper ideal \((x,y)\). At least one factor lies in \(\mathfrak p\), because the product does; two such factors would force \(x\in\mathfrak p^2\). Therefore exactly one factor \(\mathfrak p_1\) lies in \(\mathfrak p\), and \((x,y)R_{\mathfrak p}=\mathfrak p_1R_{\mathfrak p}\) is prime, all other factors becoming the unit ideal.

Apply this with \(y^2\). Since the resulting localized ideal is prime and contains \(y^2\), it contains \(y\), giving the source equality \(y=ax+by^2\). The coefficient \(by\) lies in \(\mathfrak pR_{\mathfrak p}\), so \(1-by\) is a unit; its inverse gives \(y=(1-by)^{-1}ax\). This is the omitted unit step in OCC-12343, completed without changing the theorem. Hence \(\mathfrak pR_{\mathfrak p}=(x)\).

Factoring \((x)\) again yields a factor \(\mathfrak p_1\subset\mathfrak p\) with equal localizations. Contracting these prime ideals from \(R_{\mathfrak p}\) gives \(\mathfrak p_1=\mathfrak p\), because the inverted set is disjoint from both primes. The product-of-ideals lemma makes \(\mathfrak p\) finitely generated. Thus every prime ideal, including zero, is finitely generated, and the source Cohen criterion makes \(R\) Noetherian. Each \(R_{\mathfrak p}\) for nonzero \(\mathfrak p\) now has a principal nonzero maximal ideal, so is a DVR by the already verified characterization. In particular MC-STK-ERR-0640 correctly states the required conclusion at every nonzero maximal ideal.

The equivalence with a Noetherian normal domain of dimension at most one preserves the field case: if \(R\) is a field there are no nonzero maximal ideals and the only nonzero ideal is \(R\). Otherwise the stated DVR localizations imply normality by maximal-local normality, and preclude a chain of length greater than one by localizing such a chain at a containing maximal ideal. Conversely a nonzero maximal localization of a normal Noetherian domain of dimension at most one has dimension exactly one and is a DVR.

For existence of factorization under these conditions, MC-STK-ERR-0641 starts with a nonzero proper ideal \(I\); \(R\) itself is the empty product. Choose an original maximal ideal \(\mathfrak p\supset I\), and retain \[
 J=(I:\mathfrak p)=\{z\in R:z\mathfrak p\subset I\}.
\] Formation of this colon commutes with localization because \(\mathfrak p\) is finitely generated: if a localized element sends each of its finitely many generators into the localized \(I\), take the product of the finitely many membership denominators to obtain a single witness in the original colon. At a different maximal ideal, \(\mathfrak p\) becomes the unit ideal and \(J_{\mathfrak m}=I_{\mathfrak m}\). At \(\mathfrak p\), write \(I_{\mathfrak p}=(\pi^n)\), \(n\ge1\), in the original DVR. Then \[
 J_{\mathfrak p}=(\pi^{n-1}),\qquad
 (J\mathfrak p)_{\mathfrak p}=I_{\mathfrak p}.
\] Thus \(J\mathfrak p=I\) after testing at every maximal ideal, and \(J\supsetneq I\) because the inclusion is strict at \(\mathfrak p\). The latter is exactly the induction step requested by OCC-12346. Noetherian induction on larger ideals now terminates and gives a product of nonzero prime ideals.

For the uniqueness expressly omitted in the source, every nonzero prime is maximal: its localization is a DVR, so a strict chain of two nonzero primes would contradict the dimension bound. Localize two products for the same \(I\) at each original maximal ideal \(\mathfrak m\). All factors except \(\mathfrak m\) become units, and the two products become \((\pi^{a_{\mathfrak m}})\) and \((\pi^{b_{\mathfrak m}})\). Equality forces \(a_{\mathfrak m}=b_{\mathfrak m}\) by cancellation and the fact that \(\pi\) is not a unit. Thus the original multisets of prime factors are equal. This supplies OCC-12347 as a separate editorial proof, not an automatic replacement of the author's explicit omission.

Finally at 29753--29771 the original integral closure \(B\) in finite \(L/K\) has fraction field \(L\). For each \(z\in L\), retain a monic equation \(z^n+\sum_{i=1}^n c_i z^{n-i}=0\), with \(c_i\in K\). Choose \(0\ne d\in A\) with every \(dc_i\in A\). The exact equation \[
 (dz)^n+\sum_{i=1}^n(d^i c_i)(dz)^{n-i}=0
\] has every coefficient in \(A\), so \(dz\in B\) and \(z=(dz)/d\in\operatorname{Frac}(B)\). If an element of \(L\) is integral over \(B\), transitivity of integrality makes it integral over \(A\), hence already in \(B\). Thus \(B\) is normal. Together with Krull--Akizuki and integral dimension equality this proves the source Dedekind conclusion. Lying over and the previous finite residue-fibre theorem give the remaining assertions. This supplies the implicit facts in OCC-12350 without assuming that \(B\) is finite over \(A\); the strict example in the preceding review shows why that stronger assertion would be false.

\subsubsection{Exact localization exponents and the atomicity boundary}\label{algebra-proof-029-exact-localization-exponents-and-the-atomicity-boundary}

Here is a separate consequence for the original UFD \(R\), its fraction field \(K\), and any multiplicative set \(S\subset R\) with \(0\notin S\). Choose one representative from each association class of nonzero prime elements and denote the set by \(\mathcal P\). For \(z\in K^\times\), retain its exact factorization \[
 z=u\prod_{p\in\mathcal P}p^{v_p(z)},\qquad u\in R^\times,
\] with integer exponents of finite support. Let \(\mathcal P_0=\{p:(p)\cap S=\varnothing\}\). Then, as actual subsets of the same \(K\), \[
 S^{-1}R=\{0\}\cup
 \{z\in K^\times:v_p(z)\ge0\text{ for every }p\in\mathcal P_0\}.
\] One inclusion follows by applying each prime exponent to the original numerator and denominator. For the other, for every negative exponent at an inverted prime \(p\), choose an actual \(s_p\in S\cap(p)\) and use \(s=\prod_{v_p(z)<0}s_p^{-v_p(z)}\in S\). All exponents of \(sz\) are now nonnegative, so \(sz\in R\), and the fraction \(z=(sz)/s\) is the required inverse construction. This retains every original unit and each chosen denominator witness.

An element is a unit in \(S^{-1}R\) precisely when all surviving exponents are zero, by applying the same criterion to it and its inverse. The nonunit irreducible classes there correspond exactly to \(\mathcal P_0\): the other primes become units, while a surviving prime remains prime because the localized quotient domain is nonzero. Consequently restriction of exponent coordinates realizes the exact sequence \[
 0\longrightarrow\bigoplus_{p\notin\mathcal P_0}\mathbf Z
 \longrightarrow K^\times/R^\times
 \longrightarrow K^\times/(S^{-1}R)^\times
 \longrightarrow0,
\] where the middle and last groups identify with the direct sums over all primes and over surviving primes respectively. The first map is the actual product of the indicated prime powers modulo \(R^\times\); surjectivity and its kernel follow from the explicit exponent criterion.

Atomicity in the converse of Nagata's criterion cannot be dropped. Fix a field \(k\) and the actual subring \[
 A=k[[t]]+u\,k((t))[[u]]
 \subset K=k((t))((u)).
\] For a nonzero Laurent series, let its value be the ordered pair consisting of its \(u\)-order and the \(t\)-order of its leading coefficient, ordered lexicographically. The value of a product is the sum, since the first nonzero coefficients multiply. A nonzero sum has value at least the minimum of the two values: different \(u\)-orders retain the lower leading term; at equal \(u\)-order the sum of the leading coefficients has \(t\)-order at least their minimum, or vanishes and raises the \(u\)-order. The displayed \(A\) is exactly the nonnegative part: positive \(u\)-order allows any integer \(t\)-order, and zero \(u\)-order requires nonnegative \(t\)-order of its constant coefficient. It is a valuation domain with fraction field \(K\), since one of a nonzero series and its inverse has nonnegative value.

Its maximal ideal is exactly \(tA\): division by the original \(t\) takes every positive value to a nonnegative value, and multiplication by \(t\) makes it positive. Thus \(t\) is a nonzero prime element. The actual localization is \[
 A[1/t]=k((t))[[u]],
\] a DVR: a series on the right can have its constant coefficient placed in \(k[[t]]\) by multiplication by one power of \(t\); all higher coefficients already satisfy the displayed definition of \(A\). Yet \(u\) is divisible by \(t^n\) in \(A\) for every \(n\). Every irreducible of \(A\) has value \((0,1)\): a positive value larger than this factors as \(t\) times another nonunit, and an element of value \((0,1)\) is associate to \(t\). A finite product of these has \(u\)-order zero and cannot equal the original \(u\), whose \(u\)-order is one. Thus \(A\) is not atomic, while its localization at the multiplicative set generated by the prime \(t\) is a UFD. This verifies the exact boundary of the source's hypothesis.

\subsubsection{Fractional ideal coordinates and the semilocal consequence}\label{algebra-proof-029-fractional-ideal-coordinates-and-the-semilocal-consequence}

For the original Dedekind domain \(R\) and its fraction field \(K\), a nonzero fractional ideal means a finite \(R\)-submodule \(I\subset K\) with \(cI\subset R\) for some nonzero \(c\in R\). At every nonzero maximal ideal \(\mathfrak p\), its actual localization is \(I_{\mathfrak p}=\pi_{\mathfrak p}^{\,v_{\mathfrak p}(I)}R_{\mathfrak p}\). Only finitely many exponents are nonzero: factor the integral ideal \(cI\) and the principal ideal \(cR\); the difference of their finite exponent lists gives the specified integers.

Every nonzero prime ideal has the actual fractional inverse \(\mathfrak p^{-1}=(R:_K\mathfrak p)\). Indeed choose \(0\ne f\in\mathfrak p\) and factor \((f)=\mathfrak p J\). The dual-map calculation gives \(\mathfrak p^{-1}=J/f\): one inclusion follows from multiplying by \(\mathfrak p\), and conversely for \(z\mathfrak p\subset R\), writing \(f=\sum x_i y_i\) gives \(z=\sum(zx_i)y_i/f\in J/f\). In particular \(\mathfrak p\mathfrak p^{-1}=R\).

Thus the map from nonzero fractional ideals to \(\bigoplus_{\mathfrak p\ne0}\mathbf Z\) sending \(I\) to its original local exponent list is bijective, with inverse \[
 (n_{\mathfrak p})\longmapsto\prod_{\mathfrak p}\mathfrak p^{\,n_{\mathfrak p}}.
\] For equality, clear denominators in two candidate fractional ideals and test the resulting finite modules at every maximal ideal. At the zero prime both become \(K\); if there are no nonzero maximal ideals, \(R\) is a field and the only such fractional ideal is \(R=K\). The operations are exactly \[
 v_{\mathfrak p}(IJ)=v_{\mathfrak p}(I)+v_{\mathfrak p}(J),\quad
 v_{\mathfrak p}(I+J)=\min(v_{\mathfrak p}(I),v_{\mathfrak p}(J)),\quad
 v_{\mathfrak p}(I\cap J)=\max(v_{\mathfrak p}(I),v_{\mathfrak p}(J)).
\] They follow by localizing the original module operations and comparing the actual powers of the uniformizer. The full fractional colon has exponent difference; for the integral colon \((I:_R J)\) of nonzero integral ideals the exponent is \(\max(v_{\mathfrak p}(I)-v_{\mathfrak p}(J),0)\). Finite generation of \(J\) supplies the common denominator for colon-localization, as proved above. These formulas preserve the original intersection, sum and product maps.

If \(R\) is semilocal, every ideal is principal. For a nonzero ideal \(I\) and each of the finitely many maximal ideals \(\mathfrak m_i\), choose \(z_i\in I\) generating \(I_{\mathfrak m_i}\); a numerator of a localized generator supplies such an element. CRT gives \(a_i\in R\) with \(a_i\equiv1\pmod{\mathfrak m_i}\) and \(a_i\equiv0\pmod{\mathfrak m_j}\) for \(j\ne i\). Set the actual element \(z=\sum_i a_i z_i\). At \(\mathfrak m_i\), its first term has precisely the exponent of \(I\), and all other terms have strictly larger exponent. Therefore \(zR_{\mathfrak m_i}=I_{\mathfrak m_i}\). Equality at all maximal ideals gives \(zR=I\). The zero ideal is generated by zero, and in the field case every nonzero ideal is generated by \(1\). Hence the original semilocal Dedekind domain is a PID, with a constructive generator for each ideal.

\subsubsection{Reading and propagation}\label{algebra-proof-029-reading-and-propagation}

The indexed ``Nagata factoriality'' query produced two research routing records, an original-TeX factorization-monoid paper and a PDF-primary algebraic-geometry volume; both remain unread and unadopted. The local-unpublished layer returned no hits. There is no PDF fallback or novelty claim.

The retained prime-element correction is propagated through Nagata's multiplicative set and the fraction-field construction. The residual unit issue in MC-STK-ERR-0633 is repaired by two disjoint additional source operations; its existing zero-case contribution is preserved. The two proposed weakenings of the ACC proof are rejected by the exact four-factor derivation. The omitted Dedekind arguments remain editorial proofs. The preceding Krull--Akizuki review supplies the already verified normality/fraction-field and nonfinite-normalization comparisons; no stronger finiteness is inferred.

Three consequence groups concern arbitrary UFD localizations and the atomicity counterexample, canonical ideal duality and multiplication over rings with zero divisors, and exact fractional ideal coordinates with the semilocal PID construction. Wider synthesis remains after core completion. Next source: Orders of vanishing, 29779. Source lookahead through 30080 was read but is unadjudicated.

\hypertarget{algebra_orders_notes_20260928}{}
\subsection{Orders of vanishing with the original lattices, determinants and norms}\label{algebra-proof-030-orders-of-vanishing-with-the-original-lattices-determinants-and-norms}

Authority: Stacks Project authors, commit a04446e57ec1fbc252a871afcec7752fb2807b14, algebra.tex:29779--30092, SHA-256 FA8BB92E58A4F78A2BD01B3B6A4A87DE0A0D279F5DD90641B574DD5FBFFFA4F3. All eleven reports are read. The whole current interval equals this authority plus MC-STK-ERR-0643, 0644 and 0645, comprising eight operations, with no added environment.

The complete arguments and consequences below are editorial material, not replacement translations. Three proposed operations are punctuation or grammar. The additional source-citation mismatch is recorded visibly below rather than silently reconstructing the translated proof. The source's signed distance and all original rings, modules, maps and multipliers remain identifiable. No novelty claim is made.

\subsubsection{Nonzerodivisor lengths and the citation hypothesis}\label{algebra-proof-030-nonzerodivisor-lengths-and-the-citation-hypothesis}

At 29794--29799, the map in the original exact sequence is \[
 R/(a)\longrightarrow R/(ab),\qquad r+(a)\longmapsto br+(ab),
\] followed by the quotient map to \(R/(b)\). It is well-defined since \(b(a)\subset(ab)\). If \(br=abq\), the original nonzerodivisor \(b\) cancels and gives \(r=aq\); hence this map is injective. The kernel of the quotient map consists precisely of the classes \(br+(ab)\). The quotient map is surjective. This proves the entire displayed short exact sequence.

There is an unreported mismatch in the finiteness citation at 29795. The cited lemma-finite-length-global, authority 29266--29290, assumes that its base ring is a domain. The present lemma allows a semilocal Noetherian ring of dimension one with zero divisors. Its conclusion is still correct; the cited domain result by itself does not prove that generality.

Here is the missing argument for the original ring, and in fact for any Noetherian ring \(R\) of dimension at most one, without semilocality. A nonzerodivisor \(a\) belongs to no minimal prime \(\mathfrak p\). Indeed \(R_{\mathfrak p}\) is a nonzero zero-dimensional Noetherian local ring and hence Artinian; its maximal ideal is nilpotent. If \(a\in\mathfrak p\), the localized \(a\) is nilpotent and still a nonzerodivisor. Iterated cancellation in an equality \(a^n=0\) would give \(1=0\), a contradiction. Every prime containing \(a\) is therefore maximal: a longer chain above a minimal prime would contradict the original dimension bound. Thus \(R/(a)\) is zero-dimensional Noetherian, or the zero ring if \(a\) is a unit, and has finite length. The same argument applies to the original \(b\) and \(ab\), since a product of nonzerodivisors is a nonzerodivisor. This uses the source's zero-dimensional Noetherian/Artinian and finite-length results, with the zero quotient kept separate.

Additivity of length in the proved sequence now gives exactly \[
 \ell_R(R/(ab))=\ell_R(R/(a))+\ell_R(R/(b)).
\] The grammar operation for OCC-00576 inserts ``that'' after ``Use''; it does not replace the source proof with this expanded editorial argument.

There is also a precise total-fraction-ring version. Let \(S\) be the set of nonzerodivisors of \(R\) and \(Q(R)=S^{-1}R\). Every unit of \(Q(R)\) can be written \(a/s\) with \(a,s\in S\). To check the numerator assertion, if \((a/s)(b/t)=1\), some \(v\in S\) satisfies \(v(ab-st)=0\); cancellation gives \(ab=st\), so \(a\) is a nonzerodivisor. Define \[
 o_R(a/s)=\ell_R(R/(a))-\ell_R(R/(s)).
\] If \(a/s=a'/s'\), cancellation of a common nonzerodivisor gives \(as'=a's\), and the proved length identity gives equal differences. Applying that same identity to both numerator and denominator proves that \(o_R:Q(R)^\times\to\mathbf Z\) is a group homomorphism. In the original domain case \(Q(R)=K\), this is precisely the source's \(\operatorname{ord}_R\). Units of \(R\) have order zero. No valuation inequality is asserted.

\subsubsection{The original lattices and their signed distance}\label{algebra-proof-030-the-original-lattices-and-their-signed-distance}

Keep the source's original domain \(R\), fraction field \(K\), vector space \(V\), and finite submodules of that same \(V\). The proofs in this section also work when \(R\) is any one-dimensional Noetherian domain, without locality. A finite torsion \(R\)-module has finite length: the product of annihilators of a finite generating list supplies a single nonzero annihilator \(f\), and \(R/(f)\) has finite length by the preceding argument.

A lattice contains a \(K\)-basis of \(V\), since a finite spanning list contains a basis. Conversely a finite submodule containing a basis localizes to all of \(V\). The map \(K\otimes_R M\to V\) sends \(a/b\otimes m\) to \(am/b\); it is injective because localization of the inclusion \(M\subset V\) remains injective, and surjective by spanning.

If \(M\subset M'\subset V\) and \(M'\) is finite, express each original generator \(y_i\) as \(x_i/f_i\), with \(x_i\in M\) and \(0\ne f_i\in R\). Their actual product \(f=\prod_i f_i\) satisfies \(fM'\subset M\). Therefore \(M'/M\) is finite torsion and has finite length. Conversely, a finite-length quotient is finitely generated, and lifting its finitely many generators and adjoining the original generators of \(M\) generates \(M'\). Since \(M'\) contains \(M\), it spans \(V\). This proves all three equivalences in part (1).

If \(M'\subset M\), it is finite by Noetherianity. When \(M'\) is a lattice, the preceding argument gives finite length of \(M/M'\). Conversely a finite-length module over a one-dimensional domain is torsion: each of its finitely many simple factors is \(R/\mathfrak m_i\), with \(\mathfrak m_i\ne0\); choose an actual \(0\ne f_i\in\mathfrak m_i\). The product of those elements annihilates the module, by its composition filtration. Thus \(fM\subset M'\) for a nonzero \(f\), and multiplying the original basis in \(M\) by \(f\) gives a basis in \(M'\). This proves part (2), retaining the original module inclusions.

The intersection \(M\cap M'\) is finite as a submodule of \(M\). For an original basis \(x_i\in M\), choose \(0\ne f_i\in R\) with \(f_ix_i\in M'\); the vectors \(f_ix_i\) lie in the intersection and are still a \(K\)-basis. The sum is generated by the union of the two finite generating lists and contains \(M\). This proves part (3), including \(V=0\), when every lattice is zero.

For \(M\subset M'\subset M''\), the actual inclusion and quotient maps give \[
 0\longrightarrow M'/M\longrightarrow M''/M
   \longrightarrow M''/M'\longrightarrow0.
\] All terms have finite length, proving part (4). OCC-00578 adds the missing period to this display; OCC-00577 and OCC-01416 add the earlier sentence-ending period.

Put \(P=M\cap M'\). If \(N\subset P\), the two exact sequences with common submodule \(P/N\) give \[
 \ell_R(M/P)-\ell_R(M'/P)
 =\ell_R(M/N)-\ell_R(M'/N).
\] The maps \(m+P\mapsto m+M'\) and \(m'+P\mapsto m'+M\) give the explicit isomorphisms \[
 M/P\cong(M+M')/M',\qquad M'/P\cong(M+M')/M.
\] Their kernels are precisely \(P\), and the sums show surjectivity. If \(M+M'\subset N'\), quotienting \(N'/M'\) by its submodule \((M+M')/M'\), and doing the same with \(M\), gives a common remaining quotient \(N'/(M+M')\). Cancelling its length proves the last equality in part (5). These are the six grouping corrections already supplied by MC-STK-ERR-0643 for OCC-12352.

Consequently the original signed distance \[
 d_R(M,M')=\ell_R(M/(M\cap M'))-\ell_R(M'/(M\cap M'))
\] can be computed using any common lower lattice \(N\). For three lattices take \(N=M\cap M'\cap M''\); the two intermediate terms \(\ell_R(M'/N)\) cancel, giving \[
 d_R(M,M'')=d_R(M,M')+d_R(M',M'').
\] It follows that \(d_R(M,M)=0\) and \(d_R(M,M')=-d_R(M',M)\). This supplies the proof explicitly omitted at 29961--29963 as separate editorial material. The source's term ``distance'' denotes this signed integer and does not impose the axioms of a metric.

\subsubsection{The original determinant identity and every generator}\label{algebra-proof-030-the-original-determinant-identity-and-every-generator}

A \(K\)-linear isomorphism \(\varphi:V\to V\) carries intersections to intersections and gives \(R\)-linear isomorphisms on the two quotient modules. Hence \[
 d_R(\varphi M,\varphi M')=d_R(M,M').
\] Using the just-proved cocycle identity in the exact order written in the source yields \[
\begin{aligned}
 d_R(M,\varphi M)
 &=d_R(M,M')+d_R(M',\varphi M')+
        d_R(\varphi M',\varphi M)\\
 &=d_R(M',\varphi M').
\end{aligned}
\] Thus \(D_R(\varphi)=d_R(M,\varphi M)\) is independent of the original lattice used to compute it. Also \[
 D_R(\varphi\psi)
 =d_R(M,\psi M)+d_R(\psi M,\varphi\psi M)
 =D_R(\psi)+D_R(\varphi).
\] MC-STK-ERR-0644 correctly closes the source's extra grouping around \(M\). Its resulting expression is balanced and has the same composite. The alternative proposal deleting that redundant grouping is not another required operation.

The other side, \(\operatorname{ord}_R(\det_K\varphi)\), is additive under composition by determinant multiplicativity and the proved order homomorphism. Retain a chosen original basis \(e_1,\ldots,e_n\) of \(V\). For \(n=0\), the lattice is zero, the determinant is the empty determinant \(1\), and both sides are zero.

For \(n\ge1\), the group is generated by the transvections \(E_{ij}(\lambda)\), \(i\ne j\), and diagonal scalings \(E_i(a)\), \(a\ne0\). Gaussian elimination proves this without assuming a nonzero first pivot: exchange the first row with a row having a nonzero first-column entry, retain that entry as the diagonal scaling factor, clear the other entries by the corresponding actual transvections, and apply induction to the remaining invertible block. Clear the remaining first-row entries with further transvections. Row exchange is itself a product of the stated generators; on the affected two coordinates, \[
 \begin{pmatrix}0&1\\1&0\end{pmatrix}
 =E_{12}(1)E_{21}(-1)E_{12}(1)
       \begin{pmatrix}-1&0\\0&1\end{pmatrix}.
\] Direct multiplication verifies this identity, including characteristic two, and retains its determinant factor \(-1\).

If \(\lambda=0\), the transvection is the identity. If \(\lambda\ne0\), retain the actual new lattice with basis \(f_i=\lambda e_i\) and \(f_k=e_k\) for \(k\ne i\). The original transformation sends \(f_j\) to \(f_j+f_i\), fixes the other basis vectors, and its inverse subtracts \(f_i\). Hence it carries this lattice to itself. This is the source basis comparison with its multiplier \(\lambda\) and inverse \(\lambda^{-1}\) kept explicit. Its distance is zero and its determinant is \(1\).

For a diagonal factor \(a=f/g\), with the original nonzero \(f,g\in R\), compare \[
 M=Re_1\oplus\cdots\oplus Re_n,\qquad
 E_1(a)M=aRe_1\oplus Re_2\oplus\cdots\oplus Re_n
\] using the common lower lattice \(N=fRe_1\oplus Re_2\oplus\cdots\oplus Re_n\). The quotients are \(R/fR\) and \(aR/fR\), and \[
 R/gR\xrightarrow{\cong}aR/fR,\qquad r+gR\longmapsto ar+fR.
\] It is surjective, and \(ar=fq=agq\) gives \(r=gq\) by cancellation in \(K\), proving injectivity. Thus \[
 d_R(M,E_1(a)M)=\ell_R(R/fR)-\ell_R(R/gR)
              =\operatorname{ord}_R(a).
\] This verifies the original equality with its actual fraction, without an assumption that \(a\in R\). MC-STK-ERR-0645 corrects only the undefined rank \(b\) to the original \(n\). Additivity on all the displayed generators proves \[
 d_R(M,\varphi M)=\operatorname{ord}_R(\det_K\varphi).
\] The proof has established this for every one-dimensional Noetherian domain, with no local or semilocal assumption.

\subsubsection{Finite modules, torsion terms and exact-sequence indices}\label{algebra-proof-030-finite-modules-torsion-terms-and-exact-sequence-indices}

Let \(R\) be a one-dimensional Noetherian domain, \(M\) any finite \(R\)-module, and \(\varphi:M\to M\) an endomorphism for which \(\varphi_K\) on the original \(V=K\otimes_R M\) is invertible. The actual kernel and cokernel are finite torsion modules, hence have finite length. Let \(T\subset M\) be its torsion submodule, preserved by \(\varphi\), and let \(\overline M=M/T\), embedded canonically in \(V\). It is a lattice. The induced \(\overline\varphi\) is injective since it is injective after passage to \(V\). The exact kernel/cokernel sequence for \(0\to T\to M\to\overline M\to0\) is therefore \[
 0\to\ker\varphi_T\to\ker\varphi\to0
 \to\operatorname{coker}\varphi_T
 \to\operatorname{coker}\varphi
 \to\operatorname{coker}\overline\varphi\to0.
\] Both equalities \(\ell_R(T)=\ell_R(\ker\varphi_T)+\ell_R(\operatorname{im}\varphi_T)\) and \(\ell_R(T)=\ell_R(\operatorname{im}\varphi_T)+\ell_R(\operatorname{coker}\varphi_T)\) hold, so the torsion kernel and cokernel have equal lengths. Subtracting in the exact sequence gives the full formula \[
 \ell_R(\operatorname{coker}\varphi)-\ell_R(\ker\varphi)
 =\ell_R(\overline M/\overline\varphi\overline M)
 =\operatorname{ord}_R(\det_K\varphi_K).
\] The original \(M\) has not been replaced by its torsion-free quotient: both original torsion contributions are present and their cancellation is proved. For multiplication by \(0\ne x\in R\), with \(r=\dim_K V\), this gives \[
 \ell_R(M/xM)-\ell_R(M[x])=r\,\ell_R(R/xR).
\] It recovers the previously verified torsion-free length equality in ALGEBRA-RECON-700 and specifies exactly the extra term when torsion is allowed. If \(r=0\), the determinant is \(1\) and the formula is equality of the two lengths.

There is also exact compatibility with subspaces. Let \(W\subset V\) be a \(K\)-subspace and \(\pi:V\to V/W\) the actual quotient map. For lattices \(M,M'\), the submodules \(M\cap W,M'\cap W\) are lattices in \(W\): they are finite by Noetherianity, and denominators can be cleared in each vector of a basis of \(W\) to place it in the corresponding original lattice. The images \(\pi M,\pi M'\) are finite and span the quotient. Choose a common lower lattice \(N\subset M\cap M'\). The actual maps give \[
 0\longrightarrow (M\cap W)/(N\cap W)
 \longrightarrow M/N\longrightarrow \pi M/\pi N\longrightarrow0.
\] If \(\pi m=\pi n\) with \(n\in N\), then \(m-n\in M\cap W\); this proves exactness at the middle term and identifies its kernel without assuming a splitting. The analogous sequence for \(M'\), followed by subtraction of lengths, proves \[
 d_V(M,M')=
 d_W(M\cap W,M'\cap W)+
 d_{V/W}(\pi M,\pi M').
\] If \(\varphi W=W\), then \(\varphi M\cap W=\varphi(M\cap W)\), while \(\pi(\varphi M)=\overline\varphi(\pi M)\). Thus the determinant-distance identity is additive in this exact sequence, with determinant product \(\det(\varphi|_W)\det(\overline\varphi)\). An adapted original basis represents \(\varphi\) by an upper block triangular matrix; expanding its determinant retains exactly this product, with no extra sign or scalar.

\subsubsection{The original finite-extension norm and its global version}\label{algebra-proof-030-the-original-finite-extension-norm-and-its-global-version}

At 30037--30078 retain the original finite injection \(A\subset B\) of domains and the corresponding fraction fields \(K\subset L\). The ring \(K\otimes_A B\) embeds in \(L\), is a finite-dimensional domain over \(K\), and hence is a field: multiplication by any nonzero element is an injective endomorphism of a finite-dimensional vector space, so is surjective and gives an inverse. It contains \(B\) and \(K\), so equals the original fraction field \(L\). Consequently \(B\) is an \(A\)-lattice in that same \(L\).

For local \(A\) of dimension one, finiteness makes \(B\) integral and Noetherian of dimension one. Every maximal ideal of \(B\) contracts to \(\mathfrak m_A\). The algebra \(B/\mathfrak m_AB\) is finite-dimensional over \(\kappa(\mathfrak m_A)\), and hence has finitely many maximal ideals and finite residue fields. These are precisely the maximal ideals of \(B\); none is zero, since a field integral over \(A\) would force \(A\) to be a field. Thus each \(B_{\mathfrak m_i}\) has dimension one, as required for its original order function.

Write the original \(y\in L^\times\) as \(b/b'\), with \(b,b'\in B\) nonzero. Both the order functions and the norm are multiplicative in the exact sense needed to subtract the two cases, so first take \(y=b\in B\setminus\{0\}\). The finite-length \(B\)-module \(B/yB\) has a composition series. Each simple factor is one of the residue fields \(\kappa(\mathfrak m_i)\), and as an \(A\)-module its length is the actual field degree \([\kappa(\mathfrak m_i):\kappa(\mathfrak m_A)]\). Localizing the composition series at \(\mathfrak m_i\) counts precisely the factors of that type. This proves the restriction-of-scalars formula, as in the source lemma-pushdown-module, authority 12602--12639: \[
 \ell_A(B/yB)
 =\sum_i[\kappa(\mathfrak m_i):\kappa(\mathfrak m_A)]
          \ell_{B_{\mathfrak m_i}}((B/yB)_{\mathfrak m_i}).
\] On the other hand the proved determinant identity for multiplication by the same \(y\) on the same \(L\) gives \[
 \ell_A(B/yB)=d_A(B,yB)
 =\operatorname{ord}_A(\det_K(L\xrightarrow{\,y\,}L))
 =\operatorname{ord}_A(\operatorname{Nm}_{L/K}(y)).
\] Subtracting the expressions for \(b\) and \(b'\) proves the source formula for all \(y\), preserving every residue degree. Separability and normality of \(B\) are not needed.

For a finite injection of arbitrary one-dimensional Noetherian domains \(A\subset B\), the identical argument gives the global expression \[
 o_A(\operatorname{Nm}_{L/K}(y))
 =\sum_{\mathfrak q\in\operatorname{Max}(B)}
 [\kappa(\mathfrak q):\kappa(\mathfrak q\cap A)]
       \operatorname{ord}_{B_{\mathfrak q}}(y).
\] For \(b/b'\), only the finitely many maximal ideals in the supports of \(B/bB\) and \(B/b'B\) contribute. Each contraction is maximal by integrality, and the restriction of the corresponding simple factor has the displayed length over \(A\). The global determinant identity supplies the left side. In particular, for an original \(a\in K^\times\), multiplication on \(L\) is the scalar \(a\) on \([L:K]\) dimensions, so \[
 [L:K]\,o_A(a)=
 \sum_{\mathfrak q}
 [\kappa(\mathfrak q):\kappa(\mathfrak q\cap A)]
       \operatorname{ord}_{B_{\mathfrak q}}(a).
\] All weights and the actual extension degree remain in this equation.

\subsubsection{Two exact boundaries: valuations and finite extensions}\label{algebra-proof-030-two-exact-boundaries-valuations-and-finite-extensions}

First, the original order homomorphism need not be a valuation, even on elements of its local domain. Work with the following explicit original rings and elements over \(\mathbf Q\): \[
 x=t^2-1,\quad y=t(t^2-1),\quad
 R_0=\mathbf Q[x,y]\subset\mathbf Q[t],\qquad
 R=(R_0)_{(x,y)}.
\] The exact presentation is \[
 R_0\cong\mathbf Q[X,Y]/(Y^2-X^2(1+X)),
 \qquad X\mapsto x,\quad Y\mapsto y.
\] To prove its kernel, divide by the monic polynomial in \(Y\). A remaining expression \(F(X)+YG(X)\) maps to \(F(t^2-1)+t(t^2-1)G(t^2-1)\). Its even and odd parts both vanish, by applying \(t\mapsto-t\) and adding or subtracting. Since \(\mathbf Q[X]\to\mathbf Q[t]\) is injective and \(t(t^2-1)\ne0\), both \(F\) and \(G\) vanish. Thus the presentation is exact. It makes \(R_0\) finite over \(\mathbf Q[X]\), and hence a one-dimensional Noetherian domain; \(R\) is its one-dimensional local domain at the stated maximal ideal. Its fraction field is the original \(\mathbf Q(t)\), since \(t=y/x\).

The ring \(\mathbf Q[t]=R_0+R_0t\) is finite over \(R_0\), because \(t^2=1+x\). Localize it by \(R_0\setminus(x,y)\) to obtain \(B\). The fibre is \[
 \mathbf Q[t]/(t^2-1)\cong\mathbf Q\times\mathbf Q,
\] with points \(t=1,-1\). Hence \(B\) has exactly two maximal ideals and its localizations are the two original DVRs \(\mathbf Q[t]_{(t-1)}\) and \(\mathbf Q[t]_{(t+1)}\), both with residue field \(\mathbf Q\). The fraction-field extension has degree one, so its norm is the identity. The proved formula therefore gives, for the same \(z\in\mathbf Q(t)^\times\), \[
 \operatorname{ord}_R(z)=v_{t-1}(z)+v_{t+1}(z).
\]

This failure of a single valuation can be seen entirely in \(R\), with every length explicit. Take \[
 f=y-x,\qquad g=-y-x,\qquad f+g=-2x.
\] Substitution in the exact presentation gives \[
 R/(y-x)\cong\mathbf Q[X]_{(X)}/(X^3),\qquad
 R/(y+x)\cong\mathbf Q[X]_{(X)}/(X^3),\qquad
 R/(x)\cong\mathbf Q[Y]_{(Y)}/(Y^2).
\] The filtrations by \(X\) and \(Y\) have respectively three, three and two simple quotients, all \(\mathbf Q\). Since \(-1,-2\) are units, the original orders are \[
 \operatorname{ord}_R(f)=3,\qquad
 \operatorname{ord}_R(g)=3,\qquad
 \operatorname{ord}_R(f+g)=2<3.
\] The valuation inequality fails. Equivalently \(t-1=f/x\) and \(-t-1=g/x\) both have order one, while their sum \(-2\) has order zero. These actual maps and numbers prevent extending the earlier DVR valuation arguments to every order homomorphism.

Second, module finiteness in the source norm formula cannot be replaced by integrality plus Noetherianity. Reuse the exact characteristic-\(p\) example already proved in ALGEBRA\_KRULL\_AKIZUKI\_NOTES\_20260928.md, section ``The two original examples and their exact completions,'' and group ALGEBRA-RECON-701: \[
 A=\bigcup_{[F:k^p]<\infty}F[[x]],\qquad
 K=\operatorname{Frac}(A),\quad
 f\in k[[x]]\setminus A,\quad L=K(f),\quad
 V=L\cap k[[x]],
\] where \(k^p\subset F\subset k\), \([k:k^p]=\infty\), and \(\operatorname{char}(k)=p>0\). That proof establishes that \(A,V\) are DVRs, \(V\) is integral but not finite over \(A\), both original uniformizers are \(x\), both residue fields are \(k\), and \([L:K]=p\). On the original \(K\)-space \(L\), multiplication by \(x\) is the scalar \(x\) on \(p\) dimensions, so \[
 \operatorname{ord}_A(\operatorname{Nm}_{L/K}(x))
 =\operatorname{ord}_A(x^p)=p,
 \qquad
 [k:k]\operatorname{ord}_V(x)=1.
\] These unequal integers violate the proposed formula without finiteness. This propagates the earlier strict ramification example to the exact norm identity without changing the original field extension or its uniformizer.

\subsubsection{Reading and propagation}\label{algebra-proof-030-reading-and-propagation}

The bounded corpus queries ``lattice determinant length'' and ``Noetherian lattice determinant'' were inspected in both research and local-work layers. Each gives two research and two local routing records. The broad query is dominated by Euclidean-lattice material; the refined query still does not supply an adopted primary result. All eight new routing records remain unread and unadopted, including the two local TeX hits. There is no PDF fallback or claim of exhaustive literature coverage or novelty.

The complete Orders of vanishing section and all eleven reports are reviewed. The domain-only finiteness citation is a separate conceptual source correction; its full proof and the source's omitted cocycle argument are editorial. The three retained canonical units are checked against their actual current text, including the valid redundant-parentheses repair. Three consequence groups record the general order/determinant and torsion formulas, exact-sequence and global norm compatibility, and the two explicit boundaries. The earlier rank-sensitive length equality and characteristic-\(p\) normalization example receive the precise implications proved above. Broader synthesis remains after core completion.

Next source is Quasi-finite maps at 30093. Lookahead through 30140 is read but unadjudicated.

\hypertarget{algebra_quasifinite_notes_20260928}{}
\subsection{Quasi-finite fibres with their original idempotents, maps and denominators}\label{algebra-proof-031-quasi-finite-fibres-with-their-original-idempotents-maps-and-denominators}

Authority: Stacks Project authors, commit a04446e57ec1fbc252a871afcec7752fb2807b14, algebra.tex:30093--30453, SHA-256 FA8BB92E58A4F78A2BD01B3B6A4A87DE0A0D279F5DD90641B574DD5FBFFFA4F3. All thirteen reports are read. The whole current interval equals the authority plus MC-STK-ERR-0646, 0647 and 0648, comprising four operations, with no added environment.

The complete arguments and consequences below are separate editorial material. The proposed source operations are four small language or punctuation changes. The fronted-object sentence in OCC-00581 is intelligible and grammatically possible; its proposed reordering is optional and is not imposed. Current and diplomatic translations retain the original statements, arguments and sequence. No novelty claim is made.

\subsubsection{Isolated points and the exact fibre maps}\label{algebra-proof-031-isolated-points-and-the-exact-fibre-maps}

At 30105--30193, let \(S\) be the original finite-type \(k\)-algebra and \(\mathfrak q\) its original prime. Because \(S\) is Jacobson, its closed points meet every nonempty open subset. If \(\{\mathfrak q\}\) is open, it therefore contains a closed point, which must be \(\mathfrak q\). The singleton is consequently both open and closed. The source product decomposition gives an idempotent \(e\in S\) and the actual map \[
 S\xrightarrow{\cong}eS\times(1-e)S,\qquad s\longmapsto(es,(1-e)s),
\] whose inverse is addition. Here \(D(e)=\{\mathfrak q\}\). The ring \(eS\) has one prime, is Noetherian of dimension zero and hence Artinian local, and its localization at that prime is itself. Thus the canonical localization identifies \(eS\) with the original \(S_{\mathfrak q}\). This finite-type zero-dimensional \(k\)-algebra is finite over \(k\), by the stated Noether normalization result.

Conversely, if \(S_{\mathfrak q}\) is finite over \(k\), it is Artinian local. In the original sequence \[
 0\to K_0\to S\to S_{\mathfrak q}\to Q_0\to0
\] the kernel is finite because \(S\) is Noetherian, and the cokernel is finite over \(S\) because it is a quotient of a finite-dimensional \(k\)-space. Both localizations at \(\mathfrak q\) vanish. For each member of fixed finite generating lists choose a multiplier outside \(\mathfrak q\) killing it; the product \(g\) of these actual multipliers kills both modules. The localized sequence identifies \(S_g\) with \(S_{\mathfrak q}\), since \(g\) is already a unit in the latter. Hence \(D(g)=\{\mathfrak q\}\). Conversely any such principal open isolates \(\mathfrak q\).

If the dimension at the point is zero, choose a principal open neighbourhood of dimension zero. Its Noetherian ring is Artinian, so its spectrum is finite and discrete, and the point is isolated. Isolation itself gives a zero-dimensional open neighbourhood. To compare conditions (5) and (6), finite residue-field extension makes \(\mathfrak q\) a closed point, while every closed point of a finite-type algebra over a field has finite residue extension. At a closed point, the source's previously proved equality of point dimension and local-ring dimension gives the remaining equivalence. All six conditions and both final assertions follow. The dimension at a point must not be replaced by the dimension of its local ring alone: in \(k[T]\), the prime \((0)\) has local ring \(k(T)\) of dimension zero but is not isolated, and its residue extension over \(k\) is not finite.

For the original finite-type map \(R\to S\), retain \[
 A=S\otimes_R\kappa(\mathfrak p)
   =(R\setminus\mathfrak p)^{-1}S/
            \mathfrak p(R\setminus\mathfrak p)^{-1}S.
\] The primes of this ring correspond exactly to primes of \(S\) contracting to \(\mathfrak p\). At the original corresponding prime \(\overline{\mathfrak q}\), the canonical isomorphism is \[
 A_{\overline{\mathfrak q}}\xrightarrow{\cong}
 S_{\mathfrak q}/\mathfrak pS_{\mathfrak q}.
\] It sends the image of \(s/r\), \(r\in R\setminus\mathfrak p\), to the same fraction, and then inverts the images of the original elements of \(S\setminus\mathfrak q\). Its inverse sends the class of \(s/u\) to \((s\otimes1)/(u\otimes1)\). The two maps agree on the original generators and all their specified inverses, so are mutually inverse.

Every element of \(A\) has the form \(\overline s/r\) with \(s\in S\) and \(r\in R\setminus\mathfrak p\). Since \(r\) is a unit in this fibre, its principal open is the principal open of the original numerator \(s\). Thus the isolating element in the fibre lifts to precisely the element \(g\in S\) used in condition (3); no new localization of the base is hidden. This proves all six fibre conditions. The two proposed edits at 30205 and 30206 supply ``by'' and a colon, respectively.

For all points of the fibre at once, a discrete spectrum is quasi-compact and hence finite. Conversely, a finite-type \(k\)-algebra with finitely many primes is Artinian: its Jacobson property writes each prime as an intersection of finitely many maximal ideals, and the product of those ideals lies in that prime. Primality forces one maximal ideal to be contained in it, so the prime is maximal. The ring is zero-dimensional Noetherian, hence Artinian, and finite over \(k\) by Noether normalization. A finite \(k\)-algebra has finite discrete spectrum. The zero algebra gives the empty fibre and satisfies all three conditions. This proves the source's fibrewise and global quasi-finiteness equivalences.

Localizing at the original \(f\notin\mathfrak p\) and \(g\notin\mathfrak q\) restricts the original fibre to its principal open \(D(\overline g)\). A singleton is open in this open neighbourhood exactly when it is open in the whole fibre. This proves the stated localization criterion with the actual point and fibre maps.

\subsubsection{The original four-ring, composition and base-change proofs}\label{algebra-proof-031-the-original-four-ring-composition-and-base-change-proofs}

Retain the source's commutative square and all four prime contractions. The maps induce field embeddings \(\kappa(\mathfrak p)\to\kappa(\mathfrak p')\) and \(\kappa(\mathfrak q)\to\kappa(\mathfrak q')\): an element outside the contracted prime stays outside the receiving prime, so the residue maps extend to the indicated fields.

Put \[
 A=S\otimes_R\kappa(\mathfrak p),\quad
 k=\kappa(\mathfrak p),\quad k'=\kappa(\mathfrak p'),\quad
 B=S'\otimes_{R'}k'.
\] The exact tensor map \[
 A\otimes_k k'=S\otimes_R k'\longrightarrow B
\] is obtained from the given comparison \(S\otimes_RR'\to S'\) by tensoring over \(R'\) with \(k'\). In the original lemma it is surjective, because tensoring preserves surjections. Quasi-finiteness at \(\mathfrak q\) gives the original idempotent \(e\) with \(eA=A_{\overline{\mathfrak q}}\) finite over \(k\). Its image \(e'\in B\) yields the exact product \(B=e'B\times(1-e')B\), and the restricted tensor map \[
 (eA)\otimes_k k'\longrightarrow e'B
\] is still surjective: multiply a preimage by \(e\). Thus \(e'B\) is finite over \(k'\). The image of \(e\) in \(\kappa(\mathfrak q)\) is \(1\), hence its image in \(\kappa(\mathfrak q')\) is also \(1\). The receiving point lies in the first factor, whose finite discrete spectrum isolates it. This proves the original conclusion. The proposed operation at 30304 supplies ``is''; MC-STK-ERR-0646 already supplies the required article at 30334.

For composition \(A\to B\to C\), retain the original witnesses \(b\notin\mathfrak q\) and \(c\notin\mathfrak r\). If \(c'\in C\) is the actual image of \(b\), then \(c'\notin\mathfrak r\), so \(cc'\notin\mathfrak r\). Any prime in \(D(cc')\) over \(\mathfrak p\) contracts to a prime in \(D(b)\) over \(\mathfrak p\), which must be \(\mathfrak q\); it then lies in \(D(c)\) over \(\mathfrak q\), so is \(\mathfrak r\). This proves isolation for the composite. Finite-type generators for both maps give finite-type generators for the composite. The deletion of ``then'' at 30353 changes only the causal sentence.

For base change, the original fibre at \(\mathfrak p'\) is exactly \[
 (S\otimes_R\kappa(\mathfrak p))
       \otimes_{\kappa(\mathfrak p)}\kappa(\mathfrak p').
\] The source lemma-dimension-at-a-point-preserved-field-extension, reread at authority 28306--28356, proves equality of dimension at the original point and its lifted point. It retains the polynomial presentation and compares the two flat local maps with the same fibre ring; their local-dimension increments cancel in the two height differences. Combined with the proved isolated-point criterion, this gives the stated equality of quasi-finite loci. This is a comparison of point dimensions, not an assertion that all local-ring dimensions or nilpotent multiplicities are unchanged.

For the surjectivity assertion, choose \(\mathfrak p'\) over the contraction of a given \(\mathfrak q\). The ring \[
 \kappa(\mathfrak q)\otimes_{\kappa(\mathfrak p)}
       \kappa(\mathfrak p')
\] is nonzero: tensoring a nonzero vector space with a field extension is faithful, as is seen by extending a nonzero vector to a basis. Choose a prime of this ring and contract along the actual map \(S'=R'\otimes_RS\) into it. The resulting \(\mathfrak q'\) contracts to \(\mathfrak q\) and \(\mathfrak p'\), since each of the two field maps into the prime quotient is injective. Thus all needed points lift, without a flatness assumption on \(R\to R'\). This proves descent under surjectivity on spectra. Finite-type stability follows by tensoring an original finite generating list. MC-STK-ERR-0647 already makes the noun phrase grammatical; the longer proposed rewording is not a further necessary operation.

For permanence under an intermediate ring \(A\to B\to C\), any original \(A\)-algebra generating list for \(C\) also generates it as a \(B\)-algebra. The set of primes over \(\mathfrak q\) is a subset of the set over \(\mathfrak p\). The same original isolating \(c\notin\mathfrak r\) therefore works. This proves the source statement without requiring \(A\to B\) to be of finite type.

\subsubsection{Generic finiteness with actual coefficient lifts and nilpotence powers}\label{algebra-proof-031-generic-finiteness-with-actual-coefficient-lifts-and-nilpotence-powers}

Retain the arbitrary original ring \(R\), the finite-type \(R\)-algebra \(S\), its generators \(x_1,\ldots,x_n\), and its minimal prime \(\mathfrak p\). No Noetherian assumption on \(R\) is introduced. The assumed finite fibre has finite dimension over \(\kappa(\mathfrak p)\), by the preceding result.

The ring \(R_{\mathfrak p}\) has exactly one prime, so its maximal ideal \(\mathfrak pR_{\mathfrak p}\) is its nilradical. Every element of this ideal is nilpotent, without a claim of a common exponent for the entire ideal. The same is true of its extension to \(S_{\mathfrak p}\). Explicitly, for a fixed element \[
 z=\sum_{j=1}^m r_js_j,\qquad
 r_j\in\mathfrak pR_{\mathfrak p},\quad s_j\in S_{\mathfrak p},
\] choose actual exponents \(n_j\ge1\) with \(r_j^{n_j}=0\). In the full multinomial expansion of \(z^N\), with \(N=1+\sum_j(n_j-1)\), every term contains some \(r_j^{n_j}\) and is zero. This proves local nilpotence with every original summand retained.

If the fibre is zero, then \(1\in\mathfrak pS_{\mathfrak p}\). The preceding argument makes \(1\) nilpotent, so \(S_{\mathfrak p}=0\). Some original \(g\notin\mathfrak p\) kills \(1\) in \(S\); therefore \(S_g=0\), which is finite over \(R_g\). Suppose now the fibre is nonzero. For each generator choose a monic polynomial over \(\kappa(\mathfrak p)\), of positive degree \(d_i\), annihilating its fibre image. Lift each coefficient through \(R_{\mathfrak p}\to\kappa(\mathfrak p)\), retaining a monic \(P_i\in R_{\mathfrak p}[T]\) of that degree. Its actual value \(P_i(x_i)\) lies in \(\mathfrak pS_{\mathfrak p}\). The preceding finite-sum argument gives an exponent \(e_i\ge1\) for each such value, so \(P_i(x_i)^{e_i}=0\). MC-STK-ERR-0648 already restores the index \(i\) in both source occurrences; neither the family nor its exponents can be replaced by an unindexed polynomial.

Let \(g_1\notin\mathfrak p\) be the product of denominators in these finitely many original coefficients. Choose actual monic lifts \(\widetilde P_i\in R_{g_1}[T]\) mapping to \(P_i\). The map \(R_{g_1}\to R_{\mathfrak p}\) need not be injective; these are specified lifts, not an unproved identification of the two rings. The elements \(\widetilde P_i(x_i)^{e_i}\in S_{g_1}\) vanish after the further localization at \(R\setminus\mathfrak p\). For each of these finitely many elements, choose a multiplier \(h_i\in R\setminus\mathfrak p\) killing it already in \(S_{g_1}\). Put \[
 g_2=\prod_i h_i,\qquad g=g_1g_2.
\] Now the actual monic equations \[
 Q_i(x_i)=0\text{ in }S_g,\qquad
 Q_i(T)=\widetilde P_i(T)^{e_i}\in R_g[T],\qquad
 \deg Q_i=e_id_i
\] hold. Repeated substitution of these individual monic equations expresses every original monomial in the \(x_i\)'s as an \(R_g\)-linear combination of \[
 \prod_{i=1}^n x_i^{a_i},
 \qquad 0\le a_i<e_id_i.
\] Each substitution replaces the highest power by its actual lower-degree terms and coefficients; the total exponent decreases until it is in the indicated ranges. Thus the map from the free module on these monomials onto \(S_g\) is surjective, and \(S_g\) is generated by at most \(\prod_i(e_id_i)\) elements over \(R_g\). No linear independence or flatness of these generators is asserted. If \(n=0\), the original algebra is already a quotient of \(R\), and the empty monomial \(1\) supplies its one generator.

Every fibre over \(D(g)\) is consequently a finite-dimensional algebra of dimension at most this same bound. Its number of points and its local multiplicities are accounted for by its Artinian decomposition; the dimension bound is not merely a bound on its set of points.

\subsubsection{The canonical finite-dimensional factor and field extension}\label{algebra-proof-031-the-canonical-finite-dimensional-factor-and-field-extension}

Here is a separate consequence for any finite-type algebra \(A\) over a field \(k\). Let \(I\) be the set of its isolated points. Every such point is both an open singleton and a closed irreducible component. There are only finitely many, since \(A\) is Noetherian. Their clopen singleton idempotents \(e_{\mathfrak q}\) are pairwise orthogonal: the idempotent product has empty principal open, hence is nilpotent and therefore zero. Thus \[
 e=\sum_{\mathfrak q\in I}e_{\mathfrak q}
\] is an actual idempotent, and localization gives the canonical isomorphism \[
 eA\xrightarrow{\cong}\prod_{\mathfrak q\in I}A_{\mathfrak q}.
\] For each component the inverse is the inclusion of its idempotent factor, and their sum gives the inverse for the product. The complementary factor \((1-e)A\) has no isolated points: its spectrum is clopen in \(\operatorname{Spec}(A)\), so an isolated point there would already belong to \(I\).

This idempotent is uniquely determined by its open support. If idempotents \(e,f\) have the same principal open, then \(e(1-f)\) and \(f(1-e)\) are idempotents with empty support. Each is nilpotent and hence zero; therefore \(e=ef=f\). Any finite-dimensional direct factor of \(A\) has finite discrete clopen spectrum, whose points lie in \(I\). Its idempotent \(h\) consequently satisfies \(h=he\), and its factor is a direct factor of \(eA\). Thus \(eA\) is the unique largest finite-dimensional direct factor in this precise sense. For no isolated points it is the zero factor; for a finite algebra it is all of \(A\).

For any field extension \(K/k\), the source point-dimension comparison proves that a point of \(\operatorname{Spec}(K\otimes_k A)\) is isolated exactly when its contracted point of \(\operatorname{Spec}(A)\) is isolated. Thus its isolated locus is \(D(1\otimes e)\), and uniqueness proves that its canonical finite-dimensional factor is exactly \[
 K\otimes_k eA.
\] In particular its vector-space dimension remains \(d=\dim_k(eA)\). If its Artinian local factors are \(B_j\), a composition series for each local factor gives \[
 d=\sum_j\ell_{B_j}(B_j)[\kappa(B_j):K].
\] Every simple factor is its residue field, so this identity retains all nilpotent multiplicities and all residue degrees. It also gives a bound of \(d\) on the number of isolated points after any field extension, when \(d>0\); for \(d=0\) the number is zero.

\subsubsection{A finite four-ring comparison and its full multiplicity bound}\label{algebra-proof-031-a-finite-four-ring-comparison-and-its-full-multiplicity-bound}

In the original four-ring lemma, the surjectivity assumption on \[
 \pi:S\otimes_RR'\longrightarrow S'
\] can be replaced by finiteness of this ring map. Keep every ring, prime and contraction as in the source, and suppose \(S'\) is generated by \(m\) specified elements as a module over the original tensor ring. It is of finite type over \(R'\): a finite generating list for \(S\) over \(R\), together with these module generators, generates \(S'\) as an \(R'\)-algebra.

With the original notation \(A,k,k',B,e,e'\) from the source proof above, tensoring the module surjection gives \[
 (A\otimes_k k')^{\oplus m}\twoheadrightarrow B.
\] Multiplication by the same idempotent restricts it to \[
 ((eA)\otimes_k k')^{\oplus m}\twoheadrightarrow e'B.
\] Since \(d=\dim_k(eA)\) is finite, this proves \(\dim_{k'}(e'B)\le md\). The receiving point still lies in this factor because the residue-field map still sends \(e'\) to \(1\); therefore it is isolated. This proves the strengthened four-ring conclusion on the original objects without surjectivity or flatness. If \(\pi\) is surjective, one module generator, the unit, gives \(m=1\).

Writing the full Artinian decomposition \(e'B=\prod_j B_j\) gives the exact bound \[
 \sum_j\ell_{B_j}(B_j)[\kappa(B_j):k']
       =\dim_{k'}(e'B)\le m\,\dim_k(eA).
\] For a pure base change, \(e'B=(eA)\otimes_k k'\), so the dimension is exactly \(d\); a surjective quotient can decrease it. A quotient can also create new isolated points outside this retained component: the map \(k[T]\to k\), \(T\mapsto0\), has an isolated point in its target lying over the nonisolated point \((T)\). Thus the preservation argument does not assert reflection for arbitrary quotients.

The nilpotent multiplicities in the formula are essential even for field extension. Let \[
 k=\mathbf F_p(s),\qquad K=\mathbf F_p(t),
 \qquad s\longmapsto t^p.
\] The elements \(1,t,\ldots,t^{p-1}\) form a \(k\)-basis of \(K\). For independence, clear the coefficient denominators in a proposed relation: the exponents in the resulting polynomial in \(t\) lie in distinct residue classes modulo \(p\), so all coefficients vanish. For spanning, any polynomial is a sum in those residue classes, and the inverse of a nonzero polynomial \(F(t)\) equals \(F(t)^{p-1}/F(t)^p\), with the denominator in \(k^\times\). Therefore the original finite \(k\)-algebra \[
 A=k[u]/(u^p-s)
\] is isomorphic to \(K\) by \(u\mapsto t\), and has dimension \(p\). After the stated field extension, \[
 K\otimes_k A
 =K[u]/(u^p-t^p)
 \cong K[\epsilon]/(\epsilon^p),\qquad u\longmapsto t+\epsilon.
\] The inverse sends \(\epsilon\) to \(u-t\); the equality \((u-t)^p=u^p-t^p\) holds because every intermediate binomial coefficient is zero in characteristic \(p\). The original local length and residue degree are \(1,p\), while after extension they are \(p,1\). Their product remains \(p\), even though the base-changed ring has nonzero nilpotents. No reduced quotient or separable substitute has been used.

\subsubsection{The finite locus and the boundary of the minimal-prime assumption}\label{algebra-proof-031-the-finite-locus-and-the-boundary-of-the-minimal-prime-assumption}

For a finite-type map \(R\to S\), let \[
 U=\{\mathfrak p\in\operatorname{Spec}(R):
                     S_{\mathfrak p}\text{ is finite over }R_{\mathfrak p}\}.
\] This set is open without a Noetherian assumption. At a point of \(U\), each of a fixed finite list of original algebra generators is integral over \(R_{\mathfrak p}\). Choose the corresponding monic equations, lift their finitely many coefficients to one \(R_{g_1}\), and kill their finitely many relation defects by a product \(g_2\notin\mathfrak p\), exactly as in the preceding original construction. The resulting monic relations make \(S_{g_1g_2}\) finite, so \(D(g_1g_2)\subset U\). Conversely, such a finite localization remains finite on further localization, proving precisely the asserted neighbourhood property.

If every minimal-prime fibre has finitely many points, the source theorem, with its original nilpotence and denominator arguments, places every minimal prime in \(U\). Every prime contains a minimal prime, so the union of their closures is the whole spectrum; hence \(U\) is dense. This is an open finite locus; no single common denominator across infinitely many minimal primes is assumed.

The minimal-prime qualification cannot be dropped merely because a fibre is finite and nonempty. Take the actual finitely presented map \[
 R=\mathbf Q[t]\longrightarrow
 S=R\times R[1/t],\qquad r\longmapsto(r,r).
\] An explicit presentation is \[
 S\cong R[e,z]/(e^2-e,\ ez-z,\ tz-e),
 \qquad e\longmapsto(0,1),\quad z\longmapsto(0,t^{-1}).
\] The idempotent decomposition of the presented ring has its first factor \(e=0,z=0\), namely \(R\), and its second factor \(e=1,tz=1\), namely \(R[1/t]\); these give the inverse isomorphism. If \(t\notin\mathfrak p\), the fibre is \(\kappa(\mathfrak p)\times\kappa(\mathfrak p)\). At \(\mathfrak p=(t)\), it is \(\mathbf Q\times0\cong\mathbf Q\). Thus this map is quasi-finite everywhere, and the latter fibre is finite and nonempty.

For every \(g\notin(t)\), the ring \(R_g/(t)\) is still \(\mathbf Q\), so \(t\) is not a unit in \(R_g\). If \(S_g\) were finite over \(R_g\), its quotient \(R_g[1/t]\) would be finite, and \(1/t\) would satisfy a monic equation \[
 t^{-n}+a_{n-1}t^{-(n-1)}+\cdots+a_0=0,\qquad a_i\in R_g.
\] Multiplication by the original \(t^n\) gives \(1+a_{n-1}t+\cdots+a_0t^n=0\), contradicting the nonzero quotient \(R_g/(t)\). Hence no base neighbourhood of \((t)\) makes this map finite. In contrast, at the minimal prime \((0)\), the actual choice \(g=t\) gives \(S_t=R_t\times R_t\), finite free of rank two. This identifies exactly the boundary of the source's generic-finiteness conclusion while leaving its hypotheses intact.

\subsubsection{Reading and propagation}\label{algebra-proof-031-reading-and-propagation}

The indexed query ``quasi-finite isolated'' produced two PDF-primary research routing records and no local-work hits. Both records remain unread and unadopted; there is no PDF fallback or novelty claim. The original Stacks section is fully read. The already reviewed point-dimension comparison was reread at authority 28306--28356; its proof is used with its exact point-dimension convention and hypotheses.

The three retained canonical units are semantically verified against their actual current text. One proposed prose reordering remains optional. Four proposed operations supply ``by,'' a colon, ``is,'' and remove the redundant ``then.'' Full source arguments and three consequence groups remain editorial. The canonical finite factor, its field-extension identity and all local multiplicities propagate into the finite four-ring bound; the original nilpotence powers and coefficient lifts give the generator bound and finite-locus result. The nonempty-fibre example prevents extending generic finiteness beyond its proved hypothesis. Broader synthesis remains after core completion.

Next source is Zariski's Main Theorem at 30454. Lookahead through 30550 is read but unadjudicated.

\hypertarget{algebra_zariski_main_notes_20260928}{}
\subsection{Zariski's Main Theorem: original maps and separate editorial consequences}\label{algebra-proof-032-zariskis-main-theorem-original-maps-and-separate-editorial-consequences}

Authority: Stacks Project authors, commit a04446e57ec1fbc252a871afcec7752fb2807b14, algebra.tex:30454--30953, SHA-256 FA8BB92E58A4F78A2BD01B3B6A4A87DE0A0D279F5DD90641B574DD5FBFFFA4F3. All 23 reports are read. The complete current interval is compared with the authority and the eleven operations in MC-STK-ERR-0649 through 0658. In particular, 0654 adds two references to integrality transitivity. The original human citation keys Henselian and Peskine and the acknowledgement of Thierry Coquand remain intact. Those human sources have not separately been read for this batch.

The source statements, hypotheses and exposition remain identifiable. Every complete argument and strengthening below is separate editorial evidence, not replacement translation text. Three proposed operations add a period and close two occurrences of ``sub algebra''. One contextual sentence-fragment proposal is optional. An omitted argument is not classified as a translation mistake. No novelty claim is made.

\subsubsection{Integral equations with the original localization kernels}\label{algebra-proof-032-integral-equations-with-the-original-localization-kernels}

At 30481--30496 let \(\varphi:R\to S\) and retain the original equation \[
 \sum_{i=0}^{n}\varphi(a_i)t^i=0.
\] If \(n=0\), then \(\varphi(a_0)=0\), so \(\varphi(a_0)t=0\) is integral over \(R\). This is the existing correction 0649; it avoids an unauthorized negative exponent. For \(n\geq1\), put \(z=\varphi(a_n)t\). Multiplying the original equation by \(\varphi(a_n)^{n-1}\) gives exactly \[
 z^n+\sum_{i=0}^{n-1}\varphi(a_i a_n^{\,n-1-i})z^i=0.
\] All exponents are nonnegative, and the equation is monic in \(z\), even when a displayed leading coefficient vanishes or is a zero divisor. Correction 0650 supplies ``by''.

For the trick at 30503--30538 retain the original element \(t_0\), the monic polynomial \(p\), and \(t_0\varphi(p)=\varphi(r)\). Divide \(r=qp+r'\) with \(\deg r'<\deg p\), and define the actual element \[
 \tau=t_0-\varphi(q),\qquad \tau\varphi(p)=\varphi(r').
\] Substitution in the original monic equation for \(t_0\) proves that \(\tau\) is integral over \(R[x]\). If \(p=1\), then \(r'=0\) and \(\tau=0\). Otherwise, in \(S_\tau\) the polynomial \(p(X)-\tau^{-1}r'(X)\) is monic and annihilates the image of \(\varphi(x)\). Thus that image is integral over the subring generated by the image of \(\varphi(R)\) and \(\tau^{-1}\). The image of \(\tau\) is integral over this subring with \(\varphi(x)\) adjoined, by its original integral equation. Transitivity therefore gives, in \(S_\tau\), \[
 \tau^d+\sum_{i=0}^{d-1}\sum_{j=0}^{n_i}
               \varphi(r_{ij})\tau^{i-j}=0,\qquad r_{ij}\in R,\ d\geq1.
\] Choose \(N_0\geq\max(0,j-i)\) over this finite list. The resulting element \[
 E=\tau^{d+N_0}+
       \sum_{i=0}^{d-1}\sum_{j=0}^{n_i}
             \varphi(r_{ij})\tau^{i+N_0-j}\in S
\] maps to zero in \(S_\tau\). By the definition of localization there is \(M\geq0\) with \(\tau^M E=0\) in \(S\). For \(N=N_0+M\) this is the monic equation \[
 \tau^{d+N}+\sum_{i=0}^{d-1}\sum_{j=0}^{n_i}
                \varphi(r_{ij})\tau^{i+N-j}=0
\] in the original ring. Every other exponent is smaller than \(d+N\). No injectivity of \(S\to S_\tau\), reducedness or cancellation was assumed.

For 30541--30568 retain \(a=\varphi(a_k)\). If \(a\) is nilpotent in \(S\), choose \(n\) with \(a^n=0\) and \(q=0\); the required element is zero. Otherwise localize \(R\) at \(a_k\) and \(S\) at \(a\), retaining the actual maps \(\varphi'\) and \(t'\). The polynomial \(p/a_k\) is monic over \(R_{a_k}\), and the preceding argument gives \(q'\) such that \(t'-\varphi'(q')\) is integral. Choose \(n,q\) with \(a_k^n q'\) equal to the image of \(q\). Then \[
 b=a^n t-\varphi(q)\in S
\] has integral image in \(S_a\). Here is the denominator and kernel calculation behind the cited localization lemma. Write its integral equation as \[
 b^d+\sum_{i=0}^{d-1}\frac{\varphi(r_i)}{a^{e_i}}b^i=0
 \quad\text{in }S_a.
\] Choose \(N\) such that \(N(d-i)\geq e_i\) for all \(i\). For \(y=a^N b\) the element \[
 E=y^d+\sum_{i=0}^{d-1}\varphi(r_i)a^{N(d-i)-e_i}y^i
\] maps to zero. Choose \(M\geq0\) with \(a^M E=0\). Multiplication by \(a^{Md}\), which also kills \(E\), gives for \(z=a^M y\) \[
 z^d+\sum_{i=0}^{d-1}\varphi(r_i)a^{(N+M)(d-i)-e_i}z^i=0.
\] Consequently the actual element \[
 z=a^{N+M+n}t-\varphi(a_k^{N+M}q)
\] is integral over \(R\). This retains both the polynomial denominator and the possible localization kernel. The original localization result at 7596--7627 was reread; its original multipliers remain valid.

\subsubsection{The conductor and the coefficient argument}\label{algebra-proof-032-the-conductor-and-the-coefficient-argument}

In the original situation at 30571 set \(A=\operatorname{Im}(\varphi:R[x]\to S)\). The conductor is \[
 J=\{g\in S:gS\subseteq A\}.
\] Because \(1\in S\), every member lies in \(A\). Addition and multiplication by \(S\) preserve its defining condition, so it is an ideal of \(S\) and of \(A\), and \[
 J=\operatorname{Ann}_A(S/A).
\] The module \(S/A\) is finite over \(A\). For any specified \(g\in A\), \[
 A_g=S_g\quad\Longleftrightarrow\quad
                  g^N\in J\text{ for some }N\geq0.
\] Indeed localization of \(0\to A\to S\to S/A\to0\) is exact. If its last term vanishes, each member of a finite generating list is killed by a power of \(g\); their maximum kills the whole module. The converse follows by inverting \(g\). This is a statement about the whole principal localization, not about an individual pair of local rings at a chosen prime.

The ``In Situation \ldots'' phrase at 30584 is a conventional contextual fragment; changing its punctuation is optional. Correction 0651 already states correctly that \(\varphi(R)\) is integrally closed in \(S\). For the leading-coefficient lemma, choose a finite \(A\)-module generating list \(t_1,\ldots,t_\ell\) for \(S\), adjoining \(1\) if necessary so that the list is nonempty even for the zero ring. Every \(ut_i\) is integral over \(R[x]\). If \(u\varphi(P)\in J\), then \[
 \varphi(P)ut_i\in A.
\] The preceding lemma gives \(n_i,q_i\) such that \[
 \varphi(a_k^{n_i})ut_i-\varphi(q_i)
\] is integral over \(R\), hence belongs to \(\varphi(R)\). Thus \(\varphi(a_k^{n_i})ut_i\in A\), and \(m=\max_i n_i\) gives \(u\varphi(a_k)^m\in J\). Multiplication by the remaining powers takes place in \(A\); no division occurs.

For 30607--30623 suppose \(u\varphi(P)\in\sqrt J\). Choose the actual \(n\geq1\) with \(u^n\varphi(P)^n\in J\). The leading coefficient of \(P^n\) is \(a_k^n\), so the previous argument gives \[
 u^n\varphi(a_k)^{nm}\in J
\] for some \(m\geq1\), increasing \(m\) if needed. Multiplying by \(u^{n(m-1)}\) gives \((u\varphi(a_k))^{nm}\in J\). Subtract \(u\varphi(a_k)x^k\) from \(u\varphi(P)\) and induct on \(k\); for \(k=0\) the assertion is the hypothesis. Each original coefficient therefore satisfies \(u\varphi(a_i)\in\sqrt J\). This uses powers of specified elements, not a claim that the entire radical ideal has one nilpotence exponent.

\subsubsection{Strong transcendence and the reducedness boundary}\label{algebra-proof-032-strong-transcendence-and-the-reducedness-boundary}

The source definition at 30629 is retained: for an inclusion \(R\subseteq S\), \(x\in S\) is strongly transcendental if \[
 uP(x)=0,\quad P(T)=\sum_{i=0}^k a_iT^i\in R[T]
 \quad\Longrightarrow\quad ua_i=0\text{ for every }i.
\] Equivalently, for every such original polynomial, \[
 \operatorname{Ann}_S(P(x))
       =\bigcap_{i=0}^k\operatorname{Ann}_S(a_i).
 \tag{ST}
\] The right side is always contained in the left. The reverse is exactly the definition. These are the kernels of the two actual \(S\)-linear maps \[
 S\longrightarrow S,\ u\longmapsto uP(x),
 \qquad
 S\longrightarrow S^{k+1},\ u\longmapsto(ua_0,\ldots,ua_k).
\] Tensoring their kernel sequences by a flat \(S\)-algebra \(T\) shows that (ST) is preserved for the image of \(x\) over the image of the same original \(R\). Every coefficient in that image has an \(R\)-lift, so this verifies all its polynomials. If \(T\) is faithfully flat, equality of these kernels after tensoring forces equality before tensoring: the quotient of the first kernel by the second has zero faithful tensor. Thus strong transcendence is reflected as well. These assertions do not claim invariance under enlarging the coefficient ring by an arbitrary \(R\)-base change.

In particular localization at any multiplicative subset \(W\subset S\) preserves strong transcendence over the image of \(R\). Directly, if \((u/w)P(x)=0\), some \(w'\in W\) kills \(uP(x)\) in \(S\); hence \(w'ua_i=0\) for every \(i\), so every localized coefficient product vanishes.

\textbf{Editorial strengthening of 30644--30663, ALGEBRA-RECON-762.} Reducedness is unnecessary for passage to a minimal prime. Let \(\mathfrak q\) be any minimal prime of \(S\), and put \(\mathfrak p=R\cap\mathfrak q\). Suppose \[
 uP(x)\in\mathfrak q.
\] If \(u\in\mathfrak q\), every \(ua_i\) already belongs to \(\mathfrak q\). Otherwise \(u\) is a unit in \(S_{\mathfrak q}\), so \(P(x)\) belongs to \(\mathfrak qS_{\mathfrak q}\). This local ring has exactly one prime, because primes of the localization correspond to primes contained in the minimal prime \(\mathfrak q\). Its nilradical is therefore \(\mathfrak qS_{\mathfrak q}\). Choose \(N\geq1\) with \(P(x)^N=0\) there.

Let \(R_0\) be the image of \(R\) in \(S_{\mathfrak q}\), retaining its possible nilpotents. Localized strong transcendence applied with multiplier \(1\) to the polynomial \(P(T)^N\) says that all its coefficients are zero in \(R_0\). Hence \(P(T)\) is nilpotent in \(R_0[T]\). For every prime \(\mathfrak r\) of \(R_0\), its image in the polynomial ring over the domain \(R_0/\mathfrak r\) must be zero; a polynomial ring over a domain has no nonzero nilpotents, since leading coefficients of products do not vanish. Thus every coefficient of \(P\) lies in every \(\mathfrak r\), and is nilpotent in \(R_0\). Its image in \(S_{\mathfrak q}\) is nilpotent and belongs to \(\mathfrak qS_{\mathfrak q}\). Contraction gives \(a_i\in\mathfrak q\), and hence \(ua_i\in\mathfrak q\). This proves strong transcendence in \(S/\mathfrak q\) over \(R/\mathfrak p\), without discarding any nilpotent term in the argument.

The reduced source proof remains valid and shorter: \(S_{\mathfrak q}\) is then a field, and the original multiplier \(u'\notin\mathfrak q\) kills \(uP(x)\); strong transcendence yields \(uu'a_i=0\), and primality yields \(ua_i\in\mathfrak q\). The general argument is not substituted into that translation.

When \(S\) is reduced, this also gives the exact componentwise characterization: \[
 x\text{ strongly transcendental over }R
 \quad\Longleftrightarrow\quad
 x\bmod\mathfrak q\text{ transcendental over }
 R/(R\cap\mathfrak q)
 \text{ for every minimal }\mathfrak q.
\] Here transcendence over the domain is equivalent to transcendence over its fraction field, after clearing denominators. Forward implication follows from the theorem. Conversely \(uP(x)=0\) implies \(ua_i\in\mathfrak q\) for every minimal prime, by the domain criterion, and their intersection is the nilradical, zero by reducedness.

Reducedness cannot be dropped from this converse. Take \[
 R=k,\qquad S=k[X,\epsilon]/(\epsilon^2,X\epsilon),\qquad x=X.
\] As a \(k\)-module this is \(k[X]\oplus k\epsilon\), with multiplication \[
 (f,a\epsilon)(g,b\epsilon)
  =(fg,(f(0)b+g(0)a)\epsilon).
\] The direct expression proves \(\epsilon\ne0\), \(\epsilon^2=0\) and \(X\epsilon=0\). Its unique minimal prime is \((\epsilon)\), because the quotient is \(k[X]\), and every prime contains the nilpotent \(\epsilon\). The image of \(X\) in that quotient is transcendental. Nevertheless \(\epsilon X=0\) violates strong transcendence for \(P(T)=T\), whose coefficient \(1\) is not annihilated by \(\epsilon\). Arbitrary prime quotients also need not preserve strong transcendence: \(X\) is strongly transcendental in the domain \(k[X]\) over \(k\), but its image in \(k[X]/(X)\) is zero.

\subsubsection{The original nowhere quasi-finite argument}\label{algebra-proof-032-the-original-nowhere-quasi-finite-argument}

At 30666--30711 suppose \(R\subset S\) are domains, \(x\) is transcendental over \(R\), and \(S\) is finite over the actual embedded polynomial ring \(R[x]\). First assume \(R\) normal. The original polynomial-normality lemma at 8246--8257 makes \(R[x]\) normal. If \(R\to S\) were quasi-finite at \(\mathfrak q\), its residue extension over \(\mathfrak p=R\cap\mathfrak q\) would be finite. For \(\mathfrak r=R[x]\cap\mathfrak q\), the intermediate residue extension is also finite, so \[
 \mathfrak pR[x]\subsetneq\mathfrak r.
\] Equality would give the transcendental residue field \(\kappa(\mathfrak p)(x)\). Going down for the integral inclusion \(R[x]\subset S\) produces \(\mathfrak q_0\subsetneq\mathfrak q\) over \(\mathfrak pR[x]\). Both primes lie in the same \(R\)-fibre. Every open containing \(\mathfrak q\) contains its generization \(\mathfrak q_0\); thus \(\mathfrak q\) is not isolated in that fibre, contradicting quasi-finiteness.

The going-down statement at 8517--8572 was reread. Its previously recorded proof corrections in ALGEBRA-RECON-216 through 218 retain the minus sign \(g^n=-\sum_{j<n}a_jg^j\) and handle \(z=0\) before scaling the minimal polynomial. Those proof corrections are dependencies, not newly allocated operations or instructions to expand a translation.

For arbitrary \(R\), retain \(K=\operatorname{Frac}(R)\), \(L=\operatorname{Frac}(S)\), the integral closure \(\overline R\subset K\), and \(\overline S=S[\overline R]\subset L\). The actual multiplication map \[
 S\otimes_R\overline R\longrightarrow\overline S,\qquad
 s\otimes r\longmapsto sr
\] is surjective. Hence \(\overline S\) is finite over \(\overline R[x]\), by transporting the original finite module generators. Transcendence is preserved because \(\operatorname{Frac}(\overline R)=K\). The inclusion \(S\subset\overline S\) is integral: any specified element lies in an algebra generated by finitely many elements integral over \(R\), hence integral over \(S\). Lying over provides a prime above every \(\mathfrak q\). The original four-ring lemma would transfer quasi-finiteness to that prime, contradicting the normal case. No finiteness of the whole normalization is assumed.

For the source reduced-ring lemma at 30714--30732, choose a minimal prime \(\mathfrak q_0\subseteq\mathfrak q\), pass to the domain inclusion \[
 R/(R\cap\mathfrak q_0)\ \subseteq\ S/\mathfrak q_0,
\] and use the preceding domain theorem and the four-ring lemma. Correction 0652 fixes only ``be'' to ``is''.

The general minimal-prime result in ALGEBRA-RECON-762 proves the same nowhere quasi-finite conclusion for any inclusion \(R\subseteq S\), with no reducedness assumption, provided \(x\) is strongly transcendental and \(S\) is finite over \(R[x]\). The finite module generators descend to the displayed domain quotient, and the quotient comparison map in the four-ring lemma is surjective. Thus hypothetical quasi-finiteness at \(\mathfrak q\) would give the forbidden quasi-finiteness at its quotient prime. This is the precise propagation of the strengthening to the next source lemma; the original reduced hypotheses remain in the translated statements.

\subsubsection{Monogenic localization and the explicit Horner witness}\label{algebra-proof-032-monogenic-localization-and-the-explicit-horner-witness}

At 30735--30776 let \(S=R[x]/I\), let \(B\) denote the source integral closure \(S'\) of \(R\) in \(S\), and suppose \(R\to S\) is quasi-finite at \(\mathfrak q\). There must be an original \(f=\sum a_iX^i\in I\) with some \(a_i\notin\mathfrak p=R\cap\mathfrak q\): otherwise the fibre ideal is zero and its ring is \(\kappa(\mathfrak p)[X]\), which the domain theorem shows has no quasi-finite point over \(\kappa(\mathfrak p)\). Its coefficients give a relation \[
 \sum_{i=0}^m b_i x^i=0,\qquad b_i\in B,
 \qquad b_j\notin\mathfrak q\text{ for some }j.
\] When \(m=0\) this is impossible; correction 0653 removes the duplicated copula, without altering that argument.

For \(m\geq1\) the first integral-equation lemma shows that \(b_mx\) is integral over \(B\). Since \(B\) is integral over \(R\), transitivity makes \(b_mx\) integral over \(R\), and its membership in \(S\) then puts it in \(B\). These are precisely the two transitivity insertions already composed as 0654. If \(b_m\notin\mathfrak q\), the inclusion \(B_{b_m}\subset S_{b_m}\) is surjective because \(x=(b_mx)/b_m\); it is injective by exactness of localization of the original inclusion. Otherwise set \(b'_{m-1}=b_mx+b_{m-1}\in B\), and retain the relation \[
 b'_{m-1}x^{m-1}+b_{m-2}x^{m-2}+\cdots+b_0=0.
\] The replacement leading coefficient has the same residue as \(b_{m-1}\) at \(\mathfrak q\). At least one coefficient is still outside that prime, and induction applies. The proposed period at 30770 is punctuation only.

\textbf{Editorial consequence ALGEBRA-RECON-763.} The source induction produces a specified witness from the original relation. Put \[
 j=\max\{i:b_i\notin\mathfrak q\}.
\] One has \(j\geq1\), since otherwise reduction of the original relation modulo \(\mathfrak q\) would assert \(b_0=0\) there. Define \(c_m=b_m\) and, successively for \(r=m,m-1,\ldots,j+1\), \[
 c_{r-1}=c_rx+b_{r-1}.
\] At each stage the retained relation is \[
 c_rx^r+\sum_{i=0}^{r-1}b_ix^i=0.
\] The integral-equation lemma and transitivity prove \(c_rx\in B\), and hence \(c_{r-1}\in B\). For \(r>j\), induction also gives \(c_r\in\mathfrak q\). Therefore the final elements are \[
 g=c_j=\sum_{i=j}^m b_i x^{i-j}\in B\setminus\mathfrak q,
 \qquad h=gx\in B.
\] The assertion about \(h\) uses the last relation of degree \(j\geq1\). Every summand and the original coordinate \(x\) are retained.

The actual finite \(R\)-subalgebra \(A=\operatorname{Im}(R)[g,h]\subset S\) satisfies \(A_g=S_g\): its localization contains \(x=h/g\), and it already contains the image of \(R\). The canonical localization inclusion is injective, giving the inverse to this generator description. If specified monic equations for \(g,h\) have degrees \(r,s\), successive substitution of their leading powers expresses every element of \(A\) as an \(R\)-linear combination of \[
 g^a h^b,\qquad 0\leq a<r,\quad 0\leq b<s.
\] Thus \(rs\) is a bound on the length of this specified generating list; linear independence is not asserted.

\subsubsection{The main induction with both final denominators}\label{algebra-proof-032-the-main-induction-with-both-final-denominators}

The theorem at 30779 retains \(S\) as the finite-type \(R\)-algebra, as already corrected by 0655. Corrections 0656--0658 retain ``theorem'', the example comma and the localization explanation. Adding a verb to the elliptical example is optional; it is not a remaining mathematical defect. The two proposed spelling changes occur at 30791 and 30857; the differently arranged phrase at 30894 is not automatically changed.

Write \(B\subset S\) for the original integral closure of \(R\), and choose the original \(x_1,\ldots,x_n\) with \(S\) finite over \(R[x_1,\ldots,x_n]\). If \(n=0\), then \(B=S\), and \(g=1\) works.

For \(n=1\), the permanence lemma allows the actual intermediate base inclusion \(B\subset S\), and the same module generators show finiteness over \(B[x_1]\). Prove the assertion with this base, denoting it by \(R\) in the source comparison: \(R\subset S\) is now integrally closed. This is an explicit intermediate-ring map, not an identification of the old base ring with its integral closure. Let \(A=\operatorname{Im}(R[x]\to S)\) and retain its conductor \(J\). The coefficient lemma proves that the image of \(x\) in \[
 S/\sqrt J
\] is strongly transcendental over the embedded ring \(R/(R\cap\sqrt J)\). These two rings are reduced. The original reduced nowhere quasi-finite lemma therefore suffices; the stronger editorial version is not needed to replace this proof. The four-ring lemma implies \(\mathfrak q\notin V(J)\).

Choose the actual \(s\in J\setminus\mathfrak q\), \(s=\varphi(f)\in A\). The inclusion \(A_s\to S_s\) is an isomorphism: any \(y/s^n\) is the image of \((sy)/s^{n+1}\), whose numerator is in \(A\), and localization preserves the inclusion. This explicit formula is independent of the representation because it is an inverse to the canonical inclusion. The local fibre criterion transfers quasi-finiteness to the corresponding prime of \(A\). The monogenic lemma, using that \(R\) is integrally closed in \(A\), gives \(h\in R\setminus\mathfrak q\) with \(R_h=A_h\).

Write the image of the original \(s\) in \(A_h=R_h\) as \(a/h^N\), \(a\in R\). Since \(s\) and \(h\) avoid \(\mathfrak q\), so does \(a\). Set \(g=ha\). In any localization where \(g\) is invertible, both \(h,a\) are invertible and the equality \(s=a/h^N\) makes \(s\) invertible. Conversely, after inverting \(h,s\), the element \(a=s h^N\) and hence \(g\) are invertible. The universal properties of these localizations give the canonical equalities \[
 S_g=S_{hs}=A_{hs}=R_{ha}=R_g.
\] All equalities mean the maps induced by the original inclusions and these stated inverse fractions. No unsupported cancellation in the original rings occurs.

For \(n>1\), retain the source \(R'\), the integral closure of \(R[x_1,\ldots,x_{n-1}]\) in \(S\). Permanence transfers quasi-finiteness, and the same generators make \(S\) finite over \(R'[x_n]\). The case \(n=1\) gives \(g'\in R'\setminus\mathfrak q\) with \(R'_{g'}=S_{g'}\). This integral closure need not be of finite type over \(R\).

Choose a finite \(R\)-algebra generating list for \(S_{g'}\) in the form \(y_i/(g')^{n_i}\), and adjoin \(g'\) to it if needed. Define the actual \(R\)-subalgebra \[
 R''=\operatorname{Im}(R)[x_1,\ldots,x_{n-1},y_1,\ldots,y_N,g']
 \subset R'.
\] Localization of its inclusion is injective; the specified generators prove \(R''_{g'}=S_{g'}\). Each adjoined generator is integral over \(R[x_1,\ldots,x_{n-1}]\), so its monic equation supplies a finite bounded-power spanning list, and their products prove that \(R''\) is finite over that original polynomial-image ring. For \(\mathfrak q''=R''\cap\mathfrak q\), one has \(R''_{\mathfrak q''}=S_{\mathfrak q}\). Since \(R''\) is of finite type over \(R\), the local fibre criterion transfers quasi-finiteness again.

Induction gives \(C=R'''\subset R''\), integral over \(R\), and \(g''\in C\setminus\mathfrak q\) with \(C_{g''}=R''_{g''}\). Retain the source numerator \(g'''\in C\) and exponent \(N\) in \[
 g'=\frac{g'''}{(g'')^N}\quad\text{in }R''_{g''}.
\] The numerator avoids \(\mathfrak q\), by reduction in the localization at that prime. Set \(g=g''g'''\). After inverting \(g\), both \(g''\) and \(g'\) become units. Localizing the two previously proved canonical equalities gives \[
 C_g=R''_g=S_g.
\] Every element of \(C\) is integral over \(R\), so \(C\subseteq B\subseteq S\); localization preserves these inclusions and gives \(B_g=S_g\). This proves the source induction with both final denominators retained.

\subsubsection{Finite envelopes and the exact open immersion}\label{algebra-proof-032-finite-envelopes-and-the-exact-open-immersion}

For the openness lemma at 30880 and the global lemma at 30906, retain the original \(B=S'\). Given a local witness \(B_g=S_g\), choose the original finite \(R\)-algebra generators \(z_1,\ldots,z_r\) of \(S\). Express their images in \(B_g\) as \(a_j/g^{n_j}\), \(a_j\in B\), and set \[
 A=\operatorname{Im}(R)[g,a_1,\ldots,a_r]\subseteq B.
\] Every generator is integral over \(R\). Given monic equations of degrees \(d_0,d_1,\ldots,d_r\), the monomials \[
 g^{e_0}a_1^{e_1}\cdots a_r^{e_r},
 \qquad 0\leq e_i<d_i,
\] span \(A\) over \(R\), by replacing one excessive exponent at a time with its stated monic equation. This is a finite list of size \(\prod_i d_i\), with no freeness claim. The actual inclusion identifies \(A_g=S_g\), since all \(z_j\) and \(g^{-1}\) lie in \(A_g\). Finite \(A/R\) has finite fibres, and localization at \(g\) preserves quasi-finiteness at each of its primes. Thus every point of \(D_S(g)\) is quasi-finite over \(R\), proving openness.

Now assume \(S/R\) is quasi-finite at every prime. Choose finitely many original witnesses \(g_i\in B\) with \[
 B_{g_i}=S_{g_i},\qquad (g_1,\ldots,g_n)S=S.
\] For \(S=0\), also \(B=0\); take \(A=0\), whose spectrum is empty and whose map to \(S\) is the identity. Its finite presentation over itself, flatness and epimorphism assertions below are immediate. For \(S\ne0\), choose \(n\geq1\) as above.

For every \(i\), clear denominators in the same finite list of \(R\)-algebra generators of \(S\), obtaining elements \(a_{ij}\in B\) with \(z_j=a_{ij}/g_i^{n_{ij}}\) in \(S_{g_i}\). Put \[
 A=\operatorname{Im}(R)[g_1,\ldots,g_n,(a_{ij})_{i,j}]\subseteq B.
\] The preceding bounded-power argument proves \(A\) finite over \(R\), and \(A_{g_i}=S_{g_i}\) for every \(i\). No assertion that the \(g_i\) generate the unit ideal in \(A\) is made.

\textbf{Editorial consequence ALGEBRA-RECON-764.} The actual map \[
 \operatorname{Spec}(S)\longrightarrow\operatorname{Spec}(A)
\] is an open immersion onto \(U=\bigcup_iD_A(g_i)\). Every source prime avoids some \(g_i\), since their ideal in \(S\) is the unit ideal. On that chart the canonical ring isomorphism \(A_{g_i}=S_{g_i}\) gives a homeomorphism and isomorphisms of all local rings. Conversely every prime in \(D_A(g_i)\) has its unique preimage on this chart. If two source primes had the same image, choosing such a \(g_i\) for the image puts both in this one chart and identifies them. These chart inverses agree on intersections, because their ring maps on \(A_{g_ig_j}=S_{g_ig_j}\) are all induced by the same inclusion. They therefore give the inverse homeomorphism and inverse maps on the structure sheaves. This proves the open immersion, not merely its underlying topological assertion. Replacing \(A\) by \(B\) proves the analogous statement for \(\operatorname{Spec}(B)\).

Here is the full module argument for source assertion (2). If \(D_A(g)\subseteq U\), then the \(g_i\) generate the unit ideal in \(A_g\); otherwise a maximal ideal containing them would be a point of \(D_A(g)\) outside \(U\). Localizing \(A_g\to S_g\) further at each \(g_i\) gives an isomorphism. Its kernel and cokernel vanish after each localization. For any module \(M\) over a ring \(D\) and finite elements \(v_i\) generating the unit ideal, \(M_{v_i}=0\) for all \(i\) implies \(M=0\): for a specified \(m\in M\), choose \(N_i\) with \(v_i^{N_i}m=0\). The powers \(v_i^{N_i}\) still generate the unit ideal, since a maximal ideal containing them would contain all \(v_i\). Taking a linear combination equal to \(1\) proves \(m=0\). Apply this to the kernel and cokernel. Thus \(A_g=S_g\), and the same proof works for \(B\). This completes the source's omitted details, without adding a Noetherian hypothesis. The original cover lemma at 4133--4217 was reread.

Both \(A\to S\) and \(B\to S\) are flat epimorphisms of finite presentation. We prove each property with the original \(g_i\). Write \(D=A\) or \(B\). To prove flatness, tensor any injection of \(D\)-modules by \(S\). Its kernel, viewed as an \(S\)-module, vanishes after localizing at each \(g_i\), because \(S_{g_i}=D_{g_i}\) is a flat localization of \(D\). Since the \(g_i\) generate the unit ideal in \(S\), the preceding module argument makes the kernel zero.

For the epimorphism, consider \[
 \mu:S\otimes_D S\longrightarrow S,\qquad s\otimes t\longmapsto st.
\] The elements \(g_i\otimes1\) generate the unit ideal in its source: transport an actual equality \(1=\sum c_i g_i\) in \(S\) through \(s\mapsto s\otimes1\). After localizing at \(g_i\otimes1\), the source is \[
 S_{g_i}\otimes_D S=D_{g_i}\otimes_D S=S_{g_i},
\] and \(\mu\) is the identity under these canonical maps. Its kernel and cokernel therefore vanish by the same finite-cover argument. The inverse is \(s\mapsto s\otimes1\), and this also proves \(s\otimes1=1\otimes s\). If two ring maps from \(S\) agree on \(D\), applying their induced map from this tensor product to that equality shows they agree on every \(s\); this is the categorical epimorphism assertion.

For finite presentation, choose actual \(D\)-algebra generators \(z_1,\ldots,z_r\) of \(S\); the original \(R\)-generators suffice. In the unit equality \(1=\sum_i c_i g_i\), choose polynomials \(C_i\in D[X_1,\ldots,X_r]\) with \(C_i(z)=c_i\). Put \[
 F(X)=\sum_i g_i C_i(X)-1,\qquad
 Q=D[X_1,\ldots,X_r]/(F).
\] The evaluation map \(Q\to S\) is surjective, and the \(g_i\) generate the unit ideal of \(Q\). For each \(i,j\), write \(z_j=a_{ij}/g_i^{n_{ij}}\) in \(S_{g_i}=D_{g_i}\), retaining the corresponding original numerator and denominator. Choose \(m_{ij}\geq0\) killing the image of \(g_i^{n_{ij}}X_j-a_{ij}\) under evaluation; this is possible because that image vanishes after inverting \(g_i\). Thus \[
 K_{ij}=g_i^{m_{ij}}(g_i^{n_{ij}}X_j-a_{ij})
\] belongs to the actual evaluation kernel \(K\subset Q\). In \(Q_{g_i}\), the finitely many \(K_{ij}\) for this fixed \(i\) generate \(K_{g_i}\): the original polynomial evaluation into \(D_{g_i}\) has kernel generated by \(X_j-a_{ij}/g_i^{n_{ij}}\), as follows by substituting the variables one at a time in each polynomial. The relation \(F\) already evaluates to zero. Let \(K_0\) be the ideal of \(Q\) generated by all \(K_{ij}\). Then \((K/K_0)_{g_i}=0\) for every \(i\), and the finite unit cover in \(Q\) gives \(K=K_0\). We have the explicit finite presentation \[
 S\cong
 D[X_1,\ldots,X_r]\big/
 \left(F,\,
       \bigl(g_i^{m_{ij}}
             (g_i^{n_{ij}}X_j-a_{ij})\bigr)_{i,j}\right).
\] Every localization-kernel power remains present. This proves finite presentation over \(D=A\) or \(B\); it does not assert finite presentation over the original \(R\), since the finite algebra \(A/R\) need not be finitely presented.

Finally the same charts give the exact ring equalizer \[
 S\longrightarrow\prod_iD_{g_i}
       \rightrightarrows\prod_{i,j}D_{g_ig_j},
 \tag{C}
\] where the two arrows restrict the \(i\)- and \(j\)-components, and the first map uses the original inclusions. Injectivity follows from the unit-cover module argument in \(S\). For existence, identify \(D_{g_i}=S_{g_i}\) and let \((s_i)\) be a compatible tuple. Clear the finitely many denominators to choose \(N_0\) and \(b_i^0\in S\) representing \(g_i^{N_0}s_i\). In \(S_{g_j}\), the difference \(b_i^0-g_i^{N_0}s_j\) vanishes after inverting \(g_i\), by compatibility. Choose one \(M\) at least as large as all the finitely many required annihilating exponents, and set \[
 N=N_0+M,\qquad b_i=g_i^M b_i^0.
\] Then \(b_i=g_i^N s_j\) in \(S_{g_j}\) for every \(i,j\). The elements \(g_i^N\) generate the unit ideal in \(S\); choose actual \(d_i\in S\) with \(\sum_i d_i g_i^N=1\). The element \(s=\sum_i d_i b_i\) restricts to \(s_j\) for every \(j\), proving (C) with the original gluing maps.

\subsubsection{Reading, propagation and remaining scope}\label{algebra-proof-032-reading-propagation-and-remaining-scope}

The new corpus query ``Zariski Main conductor'' was inspected in both layers. Its two research records route to a TeX and a PDF record of the Bianchi modular-forms item by Vlad Serban; they are not counted as two independently read works. The local-path layer routes to an original Connes--Consani TeX copy and an Arithmetic Time PDF. All four hits remain unread and unadopted. No PDF fallback or novelty inference occurs. The route preserves the previous seven separately recorded human-source readings and updates only the bounded Stacks coverage.

ALGEBRA-RECON-762 propagates its minimal-prime strengthening to the source's next nowhere quasi-finite lemma and records the precise reducedness boundary. ALGEBRA-RECON-763 combines the retained integral-equation correction with the original induction, producing the actual Horner numerator and a two-integral-generator finite algebra. ALGEBRA-RECON-764 uses the original main theorem and finite-cover arguments to prove the finite envelope, actual open immersion, flat epimorphism, finite presentation and equalizer. Their complete proofs remain here, with zero translation replacement operations.

The reviewed four-ring results ALGEBRA-RECON-741, 742 and 747 supply the original comparison and its already recorded finite-map strengthening; only the source surjective version is needed in the main argument. The finite-locus boundary in ALGEBRA-RECON-748 remains intact: these open immersions are into the integral closure and finite intermediate algebras, not an assertion that every quasi-finite map is finite near every point of the original base. Nor does this batch claim a global finite envelope for a possibly non-quasi-compact quasi-finite locus of an arbitrary finite-type map. The constructed global envelope is proved under the source's stated hypothesis that all of \(S/R\) is quasi-finite.

Next is Applications of Zariski's Main Theorem, authority 30954. Lookahead 30954--30975 is not adjudicated. Broader synthesis follows the core task. No source or translation mutation, TeX build or publication is performed by this review.

\hypertarget{algebra_zariski_applications_notes_20260928}{}
\subsection{Applications of Zariski Main: missing points and completed finite factors}\label{algebra-proof-033-applications-of-zariski-main-missing-points-and-completed-finite-factors}

Authority: Stacks Project authors, commit a04446e57ec1fbc252a871afcec7752fb2807b14, algebra.tex:30954--31139, SHA-256 FA8BB92E58A4F78A2BD01B3B6A4A87DE0A0D279F5DD90641B574DD5FBFFFA4F3. All three reports in this interval are read. The complete current interval is the authority plus MC-STK-ERR-0659, which adds that \(B\) is not a field. Two new proposed language/punctuation operations are separate from all editorial proofs below. The original statements and source exposition remain identifiable in both translated branches.

\subsubsection{The field exception and the original order formula}\label{algebra-proof-033-the-field-exception-and-the-original-order-formula}

At 30964 let \(A\subset B\) be the original domains, with \(A\) Noetherian local of dimension one, \(B\) of finite type over \(A\), and finite fraction-field extension \(L/K\). The generic fibre \(B\otimes_A K\) is \(L\): it is a finitely generated \(K\)-algebra contained in the finite-dimensional field extension \(L\), so is a finite-dimensional domain, hence a field, and has fraction field \(L\). Thus its only prime is zero. The earlier lemma at 27602--27638 says there are only finitely many primes above \(\mathfrak m_A\), each of height one and with algebraic residue extension. Their residue extensions are finite by finite type and the field case of Zariski's lemma, as used in the source. These primes are maximal because every nonzero prime has height one and no longer chain can occur; the only other prime of \(B\) is zero.

If \(B\) is not a field, zero is not maximal, so all maximal ideals of \(B\) lie above \(\mathfrak m_A\), have height one, and have finite residue extensions. The ring \(B\) is then semi-local of dimension one and \(\mathfrak m_A B\ne B\). Conversely, if \(B\) is a field, it is \(L\), and each nonzero element of \(\mathfrak m_A\) becomes a unit; hence \(\mathfrak m_A B=B\). This proves the equivalence of the field exclusion with the stated closed-fibre condition in this particular setting.

The source's original hypotheses do allow the field case. For \[
 A=k[t]_{(t)},\qquad B=k(t)=A[1/t],
\] the extension is of finite type and \(L=K\), yet the unique maximal ideal of \(B\) is zero. Its local ring has dimension zero, so the cited one-dimensional definition of \(\operatorname{ord}_{B_{(0)}}\) does not apply. Nor is there a residue-field map from \(A/\mathfrak m_A\) to \(B\) induced by the given inclusion: \(t\) is zero in the former and invertible in the latter. Even an independently chosen inclusion of constants does not give the required finite residue extension. Its contraction in a finite integral intermediate algebra is the generic prime, not a maximal ideal. MC-STK-ERR-0659 excludes exactly this case. It is retained as one already composed correction, not allocated again.

The above fibre description also proves \(A\to B\) is globally quasi-finite, including when \(B=L\): the generic fibre is the finite field \(L/K\), and the closed fibre is a finite-type zero-dimensional algebra over \(\kappa(\mathfrak m_A)\), hence Artinian finite-dimensional (or zero). Thus the source's use of the global finite-envelope lemma at 30990 has its needed hypothesis.

Under the added nonfield assumption, let \(C\subset B\) be the finite \(A\)-subalgebra from that lemma. The actual map \(\operatorname{Spec}(B)\to\operatorname{Spec}(C)\) is an open immersion by ALGEBRA-RECON-764. For each maximal ideal \(\mathfrak m_i\) of \(B\), its contraction \(\mathfrak n_i\) in \(C\) lies above \(\mathfrak m_A\), and is therefore maximal by integrality of \(C/A\). These contractions are distinct and \[
 C_{\mathfrak n_i}=B_{\mathfrak m_i}
\] under the original inclusions. Their residue fields and local orders consequently agree. Also \(\operatorname{Frac}(C)=L\): a nonzero principal chart on which \(C_g=B_g\) identifies their fraction fields. All remaining maximal ideals of \(C\) are the source's missing \(\mathfrak p_j\).

For \(0\ne x\in\mathfrak m_A\), multiplication by \(x\) on the \(K\)-space \(L\) is scalar multiplication. In any basis its matrix is \(x\) times the identity, so its determinant, with no factor omitted, is \(x^{[L:K]}\). Apply the original finite norm/order formula at 30037--30078 to \(A\subset C\): \[
 [L:K]\operatorname{ord}_A(x)
 =
 \sum_i[\kappa(\mathfrak n_i):\kappa(\mathfrak m_A)]
          \operatorname{ord}_{C_{\mathfrak n_i}}(x)
 +
 \sum_j[\kappa(\mathfrak p_j):\kappa(\mathfrak m_A)]
          \operatorname{ord}_{C_{\mathfrak p_j}}(x).
 \tag{1}
\] Each local ring here is a one-dimensional Noetherian domain. The quotient by \(x\) is nonzero (since \(x\) belongs to its maximal ideal) and Artinian (all primes containing nonzero \(x\) are maximal). Therefore its length, which is the displayed order, is a positive finite integer. Each residue degree is also a positive finite integer. Dropping the second sum yields the original inequality, and equality forces no missing maximal ideals.

With no missing ideals, fix any original \(b\in B\) and set \[
 I_b=\{c\in C:cb\in C\}.
\] For every maximal ideal \(\mathfrak n_i\) of \(C\), the local equality writes \(b=d_i/s_i\) with \(d_i,s_i\in C\) and \(s_i\notin\mathfrak n_i\). The equality holds in \(L\), so \(s_i b=d_i\) holds in the original \(B\); hence \(I_b\) avoids each maximal ideal. It follows that \(I_b=C\), and \(b\in C\). Thus \(B=C\) is finite over \(A\). Conversely, if \(B/A\) is finite, the same finite norm formula directly gives equality. Equality for one chosen \(0\ne x\in\mathfrak m_A\) already suffices.

\subsubsection{The exact deficit, including the field case}\label{algebra-proof-033-the-exact-deficit-including-the-field-case}

\textbf{ALGEBRA-RECON-769, editorial extension.} The field case can be included without assigning a one-dimensional order to a field. Retain all three original hypotheses, and index the first sum only by the primes of \(B\) lying above \(\mathfrak m_A\). Define \[
 \Delta_B(x)
 =
 [L:K]\operatorname{ord}_A(x)
 -
 \sum_{\mathfrak q\cap A=\mathfrak m_A}
       [\kappa(\mathfrak q):\kappa(\mathfrak m_A)]
       \operatorname{ord}_{B_{\mathfrak q}}(x),
 \qquad 0\ne x\in\mathfrak m_A.
 \tag{2}
\] All these local rings have dimension one by the argument above. If \(B=L\), the sum is empty, rather than a term at its zero maximal ideal. Take any finite envelope \(C\subset B\) with its actual open immersion. Let \(M_C\) be the maximal ideals of \(C\) absent from its image. Equation (1), with these indexing conventions, proves the exact identity \[
 \Delta_B(x)=
 \sum_{\mathfrak p\in M_C}
     [\kappa(\mathfrak p):\kappa(\mathfrak m_A)]
     \operatorname{length}_{C_{\mathfrak p}}
                          (C_{\mathfrak p}/xC_{\mathfrak p}).
 \tag{3}
\] This includes every missing residue degree and local multiplicity. In particular \[
 \Delta_B(x)\geq
 \sum_{\mathfrak p\in M_C}
       [\kappa(\mathfrak p):\kappa(\mathfrak m_A)]
 \geq |M_C|.
\] There is no cancellation because every length is positive. The left side shows that the full weighted sum, although its presentation uses an envelope, is independent of the choice of \(C\).

For all these \(B\), including fields, \(\Delta_B(x)=0\) for one such \(x\) if and only if \(B/A\) is finite, and then it vanishes for every such \(x\). Indeed zero deficit means no missing maximal ideals. The denominator-ideal argument then gives \(B=C\), exactly as above. Conversely finiteness gives (1) with no missing points. A field \(B=L\) cannot be finite integral over the nonfield \(A\): if a field is integral over a domain, then for \(0\ne a\in A\), a monic equation for \(a^{-1}\), multiplied by \(a^{d-1}\), puts \(a^{-1}\) in \(A\), making \(A\) a field. For this field case (2) is explicitly \([L:K]\operatorname{ord}_A(x)>0\).

For the original example \(A=k[t]_{(t)}, B=k(t)\), choose \(C=A\). For \(x=t^n u\), \(n\geq1\) and \(u\in A^*\), the sole missing prime is \((t)\). Its residue degree is one and the quotient \(A/(t^n u)\) has the filtration with successive quotients \(k\) given by \(1,t,\ldots,t^{n-1}\), so \(\Delta_B(x)=n\). The original unit \(u\) is retained; multiplication by it induces the actual equality of ideals \((t^n u)=(t^n)\).

This exact deficit strengthens the previously reviewed rank bound in ALGEBRA-RECON-700 within the present finite-type hypotheses. It does not remove the module-finiteness requirement from the norm formula for an arbitrary integral overring: the nonfinite normalization counterexample in ALGEBRA-RECON-701 remains outside that conclusion. The source's corrected nonfield theorem stays unchanged in the translation; (2)--(3) are editorial statements with a different indexing convention.

\subsubsection{The local finite-algebra realization}\label{algebra-proof-033-the-local-finite-algebra-realization}

At 31023 retain the local map \((R,\mathfrak m_R)\to(S,\mathfrak m_S)\), the essentially finite-type hypothesis, the finite residue extension and the zero-dimensional closed fibre. Choose the actual finite-type presentation \(S=S'_{\mathfrak q'}\) used in the source. Its local fibre is \[
 S'_{\mathfrak q'}/\mathfrak m_R S'_{\mathfrak q'}
 =S/\mathfrak m_R S.
\] It has dimension zero and the given finite residue extension, so \(S'/R\) is quasi-finite at \(\mathfrak q'\). Openness provides \(g'\notin\mathfrak q'\) whose principal open is wholly quasi-finite. Applying the finite-envelope construction to this localized finite-type algebra gives a finite \(R\)-algebra \(S''\) and an open immersion of its spectrum into \(\operatorname{Spec}(S'')\). Choose one principal chart through the point. The chart identifies a localization of \(S'_{g'}\) with a localization of \(S''\), and localizing at the corresponding prime \(\mathfrak q''\) gives the actual \(R\)-algebra isomorphism \[
 S=S'_{\mathfrak q'}\cong S''_{\mathfrak q''}.
\] Each arrow comes from the chosen inclusion and the explicit numerator/denominator maps in ALGEBRA-RECON-764. This proves the source conclusion; no whole-proof replacement is proposed.

The finite residue extension is essential even with zero-dimensional local fibre. For the local map \(k\to k(t)\), the target is a localization of \(k[t]\), so is essentially of finite type, and its fibre is the zero-dimensional field \(k(t)\). It cannot be a localization of a finite \(k\)-algebra: such an algebra is finite-dimensional, has only finite residue extensions, and every localization is a finite product of its finite-dimensional Artinian local factors. Thus the conclusion would make \(k(t)/k\) finite, a contradiction. This counterexample verifies why the source retains both conditions rather than just zero fibre dimension.

\subsubsection{The completion factor and the citation scope}\label{algebra-proof-033-the-completion-factor-and-the-citation-scope}

At 31052 write \[
 D=\widehat{R_{\mathfrak p}},\qquad T=D\otimes_R S.
\] The hypotheses make \(R\) Noetherian, \(S/R\) of finite type, and the selected \(\mathfrak q\) quasi-finite above \(\mathfrak p\). They do not say that \(S/R\) is quasi-finite everywhere.

The first paragraph, 31065--31069, cites the global quasi-finite open-integral-closure lemma for the existence of a finite \(R\)-subalgebra \(C\subset S\) and a witness \(g\in C\setminus\mathfrak q\) with \(C_g=S_g\). That particular global hypothesis has not been supplied. \textbf{ALGEBRA-RECON-768 records this citation scope problem as a separate conceptual correction.} The claimed finite algebra does exist, directly from the main theorem and the finite-subalgebra construction in the proof of the openness lemma, as written fully in ALGEBRA-RECON-764: if \(B_g=S_g\) for the integral closure \(B\), lift the finite list of original algebra generators as \(a_j/g^{n_j}\), and take \(C=\operatorname{Im}(R)[g,a_1,\ldots,a_r]\). It is finite by the original monic equations and \(C_g=S_g\). Thus no stronger hypothesis or normalized source statement is needed.

Put \(\mathfrak q'=C\cap\mathfrak q\). Then \(C_{\mathfrak q'}=S_{\mathfrak q}\). The finite-completion lemma, reread at 23134--23174, gives \[
 D\otimes_R C=E\times F,\qquad
 E=\widehat{C_{\mathfrak q'}}.
\] Its proof uses the actual Artinian product decomposition modulo each \(\mathfrak p^n\), followed by inverse limit, and the equivalence of the maximal-ideal and \(\mathfrak p\)-adic topologies on these finite local factors. Let \(e=(1,0)\) be this specified idempotent, and let \(\bar e\) be its image in \(T\). The actual ring maps \[
 T\longrightarrow \bar eT\times(1-\bar e)T,\qquad
 t\longmapsto(\bar et,(1-\bar e)t)
\] have inverse addition. Write \(A=\bar eT\) and \(B=(1-\bar e)T\), as in the source. The element \(g\) has image \((u,v)\) in \(E\times F\), with \(u\) a unit in \(E\); its inverse is the image in \(E\) of \(g^{-1}\in C_{\mathfrak q'}\). Therefore the image of \(g\) in \(A\) is already a unit, with inverse transported from \(E\).

Tensoring \(C_g=S_g\) with \(D\) and localizing gives \[
 (D\otimes_R C)_g=T_g.
\] Restrict to the specified idempotent component. Localization leaves \(E\) and \(A\) unchanged because their \(g\)-images are units, so this canonical map gives \(E\cong A\). Finally \(C_{\mathfrak q'}=S_{\mathfrak q}\) induces the canonical isomorphism of their completions. Hence \[
 T=\widehat{S_{\mathfrak q}}\times B
\] with the claimed completed factor.

To identify the projection, denote it by \(\phi:T\to\widehat{S_{\mathfrak q}}\). On \(d\otimes1\), it is the natural map from \(D\), by the finite-completion construction. For any original \(y\in S\), the equality \(C_g=S_g\) writes \(y=c/g^n\) in \(S_g\), \(c\in C\). Retain the possible localization kernel: some \(g^m\) kills \(g^ny-c\) in \(S\), so \[
 g^{m+n}y=g^m c
\] in the original \(S\). The projection agrees with the natural completion map on \(g^mc\), because this is an element of \(C\). Since the image of \(g\) is a unit in \(\widehat{S_{\mathfrak q}}\), multiplying by its inverse power gives the correct natural value of \(\phi(1\otimes y)\). Pure tensors generate the tensor product additively, so this proves the naturality on all of \(T\). No injectivity of \(S\to S_g\) was assumed.

The period proposed at 31080 terminates the displayed product decomposition. The preposition proposed at 31097 changes ``by the composition'' to ``under the composition''. Neither changes a ring map or replaces the source argument.

\subsubsection{All isolated completed factors and their maximality}\label{algebra-proof-033-all-isolated-completed-factors-and-their-maximality}

\textbf{ALGEBRA-RECON-770, editorial consequence.} Let \(R\) be Noetherian, \(S/R\) of finite type, and fix \(\mathfrak p\subset R\). Let \(Q\) be the set of primes of \(S\) over \(\mathfrak p\) at which \(S/R\) is quasi-finite. These correspond to the isolated points of the finite-type fibre \[
 H=S\otimes_R\kappa(\mathfrak p).
\] They form a finite set: by ALGEBRA-RECON-746, the isolated points of a finite-type field algebra are finitely many clopen points, with actual Artinian local factors. Each \(\mathfrak q\in Q\) gives from the preceding proof an idempotent \(e_{\mathfrak q}\in T=D\otimes_R S\) and the natural factor \[
 e_{\mathfrak q}T\cong\widehat{S_{\mathfrak q}}.
\] This factor is finite over \(D\), because it is a direct factor of the finite algebra \(D\otimes_R C\), and it is a nonzero local ring.

Its reduction modulo the maximal ideal \(\mathfrak m_D\) is the original fibre local ring \[
 (e_{\mathfrak q}T)/\mathfrak m_D(e_{\mathfrak q}T)
 \cong S_{\mathfrak q}/\mathfrak pS_{\mathfrak q}
 =H_{\bar{\mathfrak q}}.
 \tag{4}
\] Indeed the right side is Artinian local. Its maximal ideal is nilpotent, so \((\mathfrak qS_{\mathfrak q})^N\subseteq\mathfrak pS_{\mathfrak q}\) for some \(N\), while \(\mathfrak pS_{\mathfrak q}\subseteq\mathfrak qS_{\mathfrak q}\). Thus completing at \(\mathfrak q\) and then reducing modulo \(\mathfrak p\) yields that same Artinian quotient. Exactness of completion for finite modules over the Noetherian local ring \(S_{\mathfrak q}\) justifies the quotient map; the local topologies and every original ideal are specified here.

If \(\mathfrak q\ne\mathfrak r\), then \(e_{\mathfrak q}e_{\mathfrak r}\) is an idempotent in the local ring \(e_{\mathfrak q}T\), hence is either zero or its identity \(e_{\mathfrak q}\). In a local ring an idempotent \(a\) satisfies \(a(1-a)=0\), and either \(a\) or \(1-a\) is a unit, proving this dichotomy. By (4), the two idempotents reduce to different isolated-point idempotents in \(H\), whose product is zero, so the product cannot be the identity. They are pairwise orthogonal. Put \[
 e=\sum_{\mathfrak q\in Q}e_{\mathfrak q}.
\] Then \(e^2=e\), and the actual mutually inverse maps are \[
 eT\longrightarrow\prod_{\mathfrak q\in Q}e_{\mathfrak q}T,\quad
 t\longmapsto(e_{\mathfrak q}t)_{\mathfrak q},
 \qquad
 (t_{\mathfrak q})_{\mathfrak q}\longmapsto\sum_{\mathfrak q}t_{\mathfrak q}.
\] Consequently \[
 T\cong
 \left(\prod_{\mathfrak q\in Q}\widehat{S_{\mathfrak q}}\right)
 \times(1-e)T.
 \tag{5}
\] For \(Q=\varnothing\), take \(e=0\), the empty product is the zero ring, and (5) is the identity \(T=0\times T\).

This is the unique largest direct factor of \(T\) finite over \(D\). To prove maximality let \(f\in T\) be any idempotent such that \(fT\) is a finite \(D\)-module. Its special fibre \(fH\) is finite-dimensional over \(\kappa(\mathfrak p)\); as a direct factor of \(H\), every one of its points is isolated and therefore lies in \(Q\). Thus \[
 f(1-e)H=0.
\] The direct factor \(f(1-e)T\) is finite over the local ring \(D\), and its reduction modulo \(\mathfrak m_D\) is zero. Nakayama's lemma gives \(f(1-e)T=0\), so \(f=fe\) and \(fT\) is a direct factor of \(eT\). Since \(eT\) itself is finite, this proves maximality and uniqueness: two largest such factors have identities \(e,e'\) satisfying \(e=ee'\) and \(e'=ee'\), hence \(e=e'\). In particular \((1-e)T\) has no nonzero direct factor finite over \(D\). This is a statement about finite direct factors, not a claim that its special fibre has no points.

Each individual factor is also independent of the chosen finite envelope. If two choices give \(e_{\mathfrak q},e'_{\mathfrak q}\), their product has nonzero special fibre \(H_{\bar{\mathfrak q}}\). Since both factors are local, the product equals each identity, so the idempotents agree. The projection in (5) is therefore exactly the tuple of the natural completion maps proved above.

If in addition the original \(R\)-module \(S\) is flat, then \(T\) is flat over \(D\), and so is its finite direct factor \(eT\). A finite flat module over the Noetherian local ring \(D\) is finite free. Its rank is its special-fibre dimension, namely the full sum \[
 \operatorname{rank}_D(eT)
 =
 \sum_{\mathfrak q\in Q}
   \operatorname{length}_{H_{\bar{\mathfrak q}}}
                         (H_{\bar{\mathfrak q}})
   [\kappa(\mathfrak q):\kappa(\mathfrak p)].
 \tag{6}
\] For each Artinian local factor, take a composition series; every simple quotient is its residue field, so vector-space dimension is the displayed length times residue degree. This proves (6) including nonreduced fibres and inseparable residue extensions. Without flatness no rank claim is made.

\subsubsection{Coverage and propagation}\label{algebra-proof-033-coverage-and-propagation}

The source dependencies reread in this batch are algebra.tex:27602--27638, 23134--23174 and 30037--30078. The exact missing-point formula retains the weighted norm/order calculation from ALGEBRA-RECON-736 and compares with the boundary in 701. The completed decomposition receives the finite-envelope argument in 764 and the canonical isolated fibre factor in 746. All supplementary statements have complete arguments above and stay separate from the translated source. No novelty is asserted.

The corpus query ``quasi-finite completion'' returned two research-layer PDF routing records (SGA I and Mazur's Modular curves and the Eisenstein ideal), and no local-layer hits. Both records remain unread and unadopted. No PDF fallback was used. These routing results do not establish literature coverage or novelty.

Next source section is Dimension of fibres at 31140. The opening 31140--31170 and reports whose first loci lie in 31171--31280 were seen as lookahead only; no disposition for those reports is claimed. No source or translation mutation, TeX build or publication occurs in this review.

\hypertarget{algebra_fibre_dimension_notes_20260928}{}
\subsection{Dimension of fibres: original charts, universal loci and valuation coefficients}\label{algebra-proof-034-dimension-of-fibres-original-charts-universal-loci-and-valuation-coefficients}

Authority: Stacks Project authors, commit a04446e57ec1fbc252a871afcec7752fb2807b14, algebra.tex:31140--31433, SHA-256 FA8BB92E58A4F78A2BD01B3B6A4A87DE0A0D279F5DD90641B574DD5FBFFFA4F3. All twenty reports in this section are read. The complete current interval equals the authority plus MC-STK-ERR-0660 through 0666, seven units with ten operations. Three proposed operations supply two terminal periods and ``by'' after ``denote''. Complete proofs and all extended statements below are separate editorial material; neither current nor diplomatic translations are rewritten around them.

\subsubsection{The original polynomial chart and every localization}\label{algebra-proof-034-the-original-polynomial-chart-and-every-localization}

The relative dimension at \(\mathfrak q\) is the dimension of the fibre as a topological space at its point, not just the dimension of the local fibre ring. For example, at the generic prime of \(k[X]\) the local ring \(k(X)\) has dimension zero, but every open neighbourhood in \(\operatorname{Spec}(k[X])\) has dimension one. This distinction is retained in every use below.

At 31170 the already composed 0660 defines \(\mathfrak p\) as the inverse image of \(\mathfrak q\). Set \[
 H=S\otimes_R\kappa(\mathfrak p),\qquad
 n=\dim_{\bar{\mathfrak q}}\operatorname{Spec}(H).
\] This is a finite-type field algebra, and \(\bar{\mathfrak q}\) exists, so the fibre is nonzero and \(n\) is a finite nonnegative integer. Choose an open neighbourhood \(U\) of dimension \(n\). The fibre topology is induced by the quotient and localization maps from \(S\): each fibre element is a fraction of the image of an element of \(S\) by an element from \(R\setminus\mathfrak p\), and that denominator is a unit on the entire fibre. Thus the principal fibre opens have the form \(D(\bar a)\), \(a\in S\). Choose \(a\notin\mathfrak q\) with \(D(\bar a)\subset U\). Its dimension is at most \(n\) and at least the minimum neighbourhood dimension \(n\); hence it is exactly \(n\). Keep \(S^{(1)}=S_a\), with this explicit localization map.

Choose original generators \(z_1,\ldots,z_N\) of \(S^{(1)}\) over \(R\), including \(a^{-1}\) if needed. The Noether normalization lemma at 28024--28054 gives polynomials with integer coefficients in these original generators, whose images \(y_1,\ldots,y_n\) make \[
 \kappa(\mathfrak p)[t_1,\ldots,t_n]\longrightarrow
                  H_{\bar a},\qquad t_i\longmapsto\bar y_i
\] finite. Their integer coefficients lift without choosing an embedding of \(\kappa(\mathfrak p)\) into \(R\). Hence they define the actual map \[
 P=R[t_1,\ldots,t_n]\longrightarrow S^{(1)}.
\] Let \(\mathfrak r\) be the inverse image of \(\mathfrak qS_a\). It contracts to \(\mathfrak p\), and its residue field is an algebra over \(\kappa(\mathfrak p)\). The fibre of \(S^{(1)}\) over \(\mathfrak r\) is the base change of the displayed finite map to \(\kappa(\mathfrak r)\), so it is finite. Thus \(P\to S^{(1)}\) is quasi-finite at \(\mathfrak qS_a\).

Openness provides \(v\in S_a\) outside that prime such that \(S_a[1/v]\) is quasi-finite over \(P\). Write the actual \(v=b/a^M\), \(b\in S\), with \(b\notin\mathfrak q\). Then \[
 S_a[1/v]\cong S_{ab};
\] one direction sends \(v^{-1}\) to \(a^M/b\), and the other uses the fact that \(a\) and \(v\) are units to invert \(b=a^M v\). Set \(g=ab\). This proves the source conclusion with a witness in the original \(S\), including \(n=0\). Corrections 0661 and 0662 retain the actual barred prime and the polynomial-variable comma; no map or coefficient is changed.

\subsubsection{The refined chart and its contracted prime}\label{algebra-proof-034-the-refined-chart-and-its-contracted-prime}

At 31211 retain \(n=\dim_{\mathfrak q}(S/R)\) and \[
 r=\operatorname{trdeg}_{\kappa(\mathfrak p)}\kappa(\mathfrak q).
\] Use the preceding principal chart \(P=R[t_1,\ldots,t_n]\to S_a\), and denote the corresponding prime of \(P\) by \(\mathfrak q'\). Its residue extension into \(\kappa(\mathfrak q)\) is finite because the map is quasi-finite. Consequently the corresponding prime \(\overline{\mathfrak q}'\) of \(\kappa(\mathfrak p)[t_1,\ldots,t_n]\) has residue transcendence degree \(r\), and \(0\leq r\leq n\).

The refined normalization theorem at 28073--28123 gives a finite map \[
 \kappa(\mathfrak p)[x_1,\ldots,x_n]\longrightarrow
 \kappa(\mathfrak p)[t_1,\ldots,t_n],\qquad x_i\longmapsto\bar h_i,
\] whose contracted prime is \((x_{r+1},\ldots,x_n)\). The earlier proof audit ALGEBRA-RECON-667 and 673 retains the original leading coefficient, its unit inverse, and the forward and inverse polynomial-coordinate maps. Those exact maps supply this dependency; no source change of coordinate is used to erase its original data.

There are finitely many coefficients in the \(\bar h_i\). Write them as original fractions in \(\operatorname{Frac}(R/\mathfrak p)\), choose lifts of each numerator and denominator to \(R\), and let \(f\notin\mathfrak p\) be the product of the denominators. Then each \(\bar h_i\) lifts to a specified \(h_i\in R_f[t_1,\ldots,t_n]\). Define \[
 \psi:R_f[x_1,\ldots,x_n]\longrightarrow R_f[t_1,\ldots,t_n],
 \qquad x_i\longmapsto h_i.
\] Its fibre over \(\mathfrak p\) is the displayed finite map. The same fibre calculation as above proves quasi-finiteness at \(\mathfrak q'R_f[t_1,\ldots,t_n]\). Choose the original principal witness \(v\) outside this prime for its quasi-finite neighbourhood. The second map to \(S_{af}\), being a base change and restriction of the first quasi-finite chart, is quasi-finite, and remains so after inverting the image of \(v\). Their composite is quasi-finite.

To express the conclusion in the original \(S\), write that image as \[
 \varphi(v)=b/(a^M\varphi(f)^N)\in S_{a\varphi(f)},
 \qquad b\notin\mathfrak q.
\] Set \(g=a\varphi(f)b\in S\setminus\mathfrak q\). The stated inversions give exactly \(S_g\); none of the three required factors is suppressed. Since \(f\) has invertible image there, this is an \(R_f\)-algebra.

The contracted prime of the composite is exactly \[
 \mathfrak pR_f[x_1,\ldots,x_n]+
                      (x_{r+1},\ldots,x_n).
 \tag{1}
\] Indeed reduction to \(\kappa(\mathfrak p)[x_1,\ldots,x_n]\) gives the coordinate-tail ideal by construction. Its contraction to \(R_f[x_1,\ldots,x_n]\) is (1): after deleting monomials involving a tail variable, membership says that each remaining original coefficient maps to zero in \(\kappa(\mathfrak p)\), which is equivalent to its membership in \(\mathfrak pR_f\). The quotient \((R_f/\mathfrak pR_f)[x_1,\ldots,x_r]\) is a domain, and its inclusion into \(\kappa(\mathfrak p)[x_1,\ldots,x_r]\) is injective, confirming the contraction without a hidden denominator. The proposed period at 31226 only terminates this statement. OCC-01428 has an inaccurate quoted phrase, but its specified terminal locus does have this punctuation omission; the review does not adopt that quote as source text.

\subsubsection{Dimension bounds and exact base change}\label{algebra-proof-034-dimension-bounds-and-exact-base-change}

For 31267, localize the original map at \(\mathfrak p\) and put \(T=(R\setminus\mathfrak p)^{-1}S\). Its integral closure \(B\subset T\) over the image of \(R_{\mathfrak p}\) is an integral \(R_{\mathfrak p}\)-algebra. The main theorem at the original prime gives \[
 S_{\mathfrak q}\cong B_{\mathfrak r},
 \qquad \mathfrak r=B\cap\mathfrak qT.
\] This is the prime corrected by 0663. In particular one must not assert \(B\subset S\) before that base localization, nor reverse the direction of the prime map. An integral ring map, even if not injective, contracts every strictly increasing prime chain strictly, by integral incomparability after quotienting its kernel. Thus the original integral dimension result at 27347--27358 gives \[
 \dim S_{\mathfrak q}
 =\dim B_{\mathfrak r}
 \leq\dim B
 \leq\dim R_{\mathfrak p}.
 \tag{2}
\] No Noetherian hypothesis is needed. For a quasi-finite map \(k[t_1,\ldots,t_n]\to A\) with \(A\) of finite type over a field \(k\), apply (2) to every prime. Any finite prime chain in \(A\) lies in the localization at its last prime, so \(\dim A\) is the supremum of those local dimensions, at most \(n\). Injectivity of the polynomial map is unnecessary. Correction 0664 changes the inclusion sign to a map; this is needed by the later fibre applications, whose maps may have kernels.

The polynomial chart through \(\mathfrak q\) now proves upper semicontinuity of pointwise relative dimension. At any base prime \(\mathfrak p'\), the chart fibre is quasi-finite over \(\kappa(\mathfrak p')[t_1,\ldots,t_n]\). The bound just proved shows its entire dimension at most \(n\), so all its point dimensions are at most \(n\). The principal chart is therefore an open neighbourhood on which the source's relative dimension is at most its original value at \(\mathfrak q\).

For a base change \(R\to R'\), let \(S'=R'\otimes_R S\), let \(\mathfrak q'\subset S'\) contract to \(\mathfrak q\subset S\), and let \(\mathfrak p'\subset R'\) and \(\mathfrak p\subset R\) be their base primes. Keep the induced field map \(\kappa(\mathfrak p)\to\kappa(\mathfrak p')\). Associativity and the original quotient/localization maps give \[
 S'\otimes_{R'}\kappa(\mathfrak p')
 \cong S\otimes_R\kappa(\mathfrak p')
 \cong
 (S\otimes_R\kappa(\mathfrak p))
                   \otimes_{\kappa(\mathfrak p)}\kappa(\mathfrak p').
 \tag{3}
\] On a pure tensor the first isomorphism sends \((a'\otimes s)\otimes c\) to \(s\otimes a'c\); its inverse sends \(s\otimes c\) to \((1\otimes s)\otimes c\). The second uses the same scalar balancing, and its inverse inserts \(1\). These maps match the corresponding primes. The original point-dimension field-extension result at 28306--28356 therefore gives equality of the two relative dimensions, with no assertion that local-ring heights or nilpotent multiplicities stay unchanged. This proves the stated inverse-image identity. The proposed ``denote by'' at 31337 corrects the actual source sentence, not the inaccurate ``composition'' quoted in OCC-01430.

For finite presentation, retain \[
 S=R[z_1,\ldots,z_N]/(F_1,\ldots,F_m).
\] The bound \(n\) is independent of the number \(N\) of variables, as corrected in both occurrences by 0665. Let \(R_0\) be the actual integer subalgebra generated by all original coefficients in the \(F_i\), and define \(S_0\) by those same polynomials. The natural map \(R\otimes_{R_0}S_0\to S\), sending \(r\otimes z_i\) to \(rz_i\), is an isomorphism by the polynomial presentation. The ring \(R_0\) is Noetherian as a quotient of a finitely generated integer polynomial ring, and so is \(S_0\). Hence its open relative-dimension locus \(V_{0,n}\) is quasi-compact. Express it as a finite union of principal opens. By (3), its inverse image is exactly \(V_n\), and the inverse image of each principal open is the corresponding principal open in \(\operatorname{Spec}(S)\), which is quasi-compact. This proves the assertion directly without an extra finiteness condition on \(R_0\to R\).

\subsubsection{Minimal chart dimension and finite charts}\label{algebra-proof-034-minimal-chart-dimension-and-finite-charts}

\textbf{ALGEBRA-RECON-781, editorial consequence.} For any finite-type \(R\to S\) and \(\mathfrak q\subset S\), \(\dim_{\mathfrak q}(S/R)\) equals the least nonnegative integer \(d\) for which some original \(g\in S\setminus\mathfrak q\) admits a quasi-finite \(R\)-algebra map \[
 R[t_1,\ldots,t_d]\longrightarrow S_g.
 \tag{4}
\] The source chart proves existence for \(d=n=\dim_{\mathfrak q}(S/R)\). Conversely base changing (4) to \(\kappa(\mathfrak p)\) and using (2) gives dimension at most \(d\) on that fibre chart. The point dimension at \(\mathfrak q\) is unchanged by restriction to an open neighbourhood, by its definition as the minimum over such neighbourhoods. Hence \(n\leq d\). This proves the minimum statement, including \(n=0\), without changing the original coordinate system of \(S\). By (3), this minimum is unchanged at every point above \(\mathfrak q\) under arbitrary base change.

Suppose in addition \(S/R\) is finitely presented, and fix \(n\geq0\). The quasi-compact open \(V_n\) has a finite cover by such principal charts with integers \(d_i\leq n\). If desired to use the same number of variables in each chart, define the explicit quotient map \[
 R[t_1,\ldots,t_n]\longrightarrow R[t_1,\ldots,t_{d_i}],
 \quad t_j\longmapsto
 \begin{cases}t_j&j\leq d_i,\\0&j>d_i.\end{cases}
\] It is finite as a module map because it is surjective, and composition preserves quasi-finiteness. The zero images are retained, not treated as an injective polynomial inclusion. For each chart, ALGEBRA-RECON-764 then constructs a finite algebra over its polynomial base into whose spectrum the chart is an open immersion. This is a finite list of actual finite-envelope factorizations covering \(V_n\), not a claim of one globally chosen coordinate map.

\subsubsection{Universal dimension ideals and the finite-presentation boundary}\label{algebra-proof-034-universal-dimension-ideals-and-the-finite-presentation-boundary}

\textbf{ALGEBRA-RECON-782, editorial consequence.} For each finite-type map and \(n\geq0\), define the actual ideal \[
 I_n=\{a\in S:D_S(a)\subseteq V_n\}.
 \tag{5}
\] This set is an ideal: \(D(ab)\subseteq D(a)\), and \(D(a+b)\subseteq D(a)\cup D(b)\). It is radical since \(D(a^r)=D(a)\) for every positive integer \(r\). Every point of the open \(V_n\) has a principal neighbourhood contained in it, so \[
 D(I_n)=V_n,\qquad V(I_n)=\operatorname{Spec}(S)\setminus V_n.
\] Thus (5) is the unique radical ideal of that closed complement. The ideals increase with \(n\), because the open sets do.

For any base change \(R\to R'\), (3) identifies the open \(V_n'\) with the inverse image of \(V_n\). Consequently their closed complements have radical ideals related by the exact equality \[
 I_n(S'/R')=\sqrt{I_n(S/R)\,S'}.
 \tag{6}
\] This does not assert equality of the extended ideal before taking its radical. If \(S/R\) is finitely presented, quasi-compactness supplies finitely many original \(a_1,\ldots,a_\ell\in I_n\) whose principal opens cover \(V_n\), and \[
 I_n=\sqrt{(a_1,\ldots,a_\ell)S}.
\] For any ring map \(S\to T\), the induced spectrum map has image in \(V_n\) exactly when \(I_nT=T\): this is equivalent to saying no prime of \(T\) contains its extension, and any proper ideal is contained in a maximal ideal. This gives a universal algebraic test for that open, without imposing finite type on \(T\).

If \(S\) is generated by \(N\) elements over \(R\), every fibre is a quotient of a polynomial ring in \(N\) variables over a field and has dimension at most \(N\). Thus \(V_N=\operatorname{Spec}(S)\). In the finitely presented case the exact relative-dimension sets are \[
 V_0,\quad V_1\setminus V_0,\quad\ldots,\quad V_N\setminus V_{N-1}.
\] Each is a finite Boolean combination of principal opens, using the finite lists just constructed; hence these give a finite constructible stratification. Formula (3) shows that every one of these strata pulls back exactly under arbitrary base change.

Finite type alone does not imply quasi-compactness of \(V_0\). Retain the product ring and coordinate idempotents from the earlier boundary example in ALGEBRA-RECON-685: \[
 R=\prod_{i\geq1}k,\qquad
 I=\bigoplus_{i\geq1}k e_i,\qquad
 S=R[X]/(e_iX:i\geq1).
 \tag{7}
\] This is a finite-type \(R\)-algebra, with its full infinite relation list retained. If a prime \(\mathfrak p\) of \(R\) avoids some \(e_i\), its localization has \(e_i=1\), hence is \(k\) via the \(i\)-th coordinate; such a prime is uniquely the coordinate kernel. The corresponding fibre of (7) is \(k[X]/(X)=k\), of dimension zero. If \(\mathfrak p\) contains all \(e_i\), then it contains \(I\), and the fibre is \(\kappa(\mathfrak p)[X]\), with point dimension one at every prime. Such primes exist because \(I\) is proper: a finite-support element cannot equal \(1\), so \(I\) lies in a maximal ideal.

It follows that \[
 V_0=\bigcup_{i\geq1}D_S(e_i).
\] Each displayed open is the nonempty one-point spectrum of \(k\), and they are pairwise disjoint because \(e_ie_j=0\) for \(i\ne j\). This cover has no finite subcover; hence \(V_0\) is not quasi-compact. In fact its radical ideal is \(IS\), with no suppressed relation: \(R/I\) is reduced because a power of a sequence has finite support only when the sequence has finite support, and \[
 S/IS=(R/I)[X]
\] is reduced. Therefore \(IS\) is radical, its complement open is exactly the displayed union, and (5) gives \(I_0=IS\). This counterexample proves why the finite-presentation step needs finite coefficient data; it does not weaken the source openness statement, which holds for all finite-type maps.

\subsubsection{The valuation-ring proof with the original coefficient factor}\label{algebra-proof-034-the-valuation-ring-proof-with-the-original-coefficient-factor}

At 31380 retain the valuation domain \(R\), its fraction field \(K\), its residue field \(k\), and the original finite-type domain \(S\) containing \(R\). Put \[
 d=\dim(S\otimes_R K),\qquad
 n=\dim_{\mathfrak q}(S/R)
\] for any \(\mathfrak q\supseteq\mathfrak m_R S\). The three substitutions in 0666 are necessary: \(S\) is not generally a \(k\)-algebra. For \(R=\mathbf Z_{(p)}\), \(S=R\), a unital map \(\mathbf F_p\to S\) cannot exist, since it would make \(p1_S=0\). Only \(S\otimes_R k\) has the stated residue-field algebra structure.

Upper semicontinuity supplies a principal open through \(\mathfrak q\) with relative dimension at most \(n\). Since \(S\) is a domain, this open contains its generic prime zero. At zero the fibre is the generic domain \(S\otimes_R K\); its point dimension is \(d\), since a finite-type field domain has the same dimension on every nonempty open, equal to its fraction-field transcendence degree. Thus \(d\leq n\).

Use the original chart \(R[t_1,\ldots,t_n]\to S_g\), \(g\notin\mathfrak q\). Its generic fibre remains a domain with dimension \(d\). If \(d<n\), the generic polynomial map has a nonzero kernel: injectivity would make the images of \(t_1,\ldots,t_n\) algebraically independent in \(\operatorname{Frac}(S)\), forcing transcendence degree at least \(n\). Choose the original nonzero kernel polynomial \[
 f=\sum_{\alpha}a_\alpha t^\alpha\in K[t_1,\ldots,t_n].
\] The finitely many nonzero coefficients have a member \(a_\beta\) of least valuation, because the divisibility order of a valuation ring is total. Retain both this coefficient and the comparison polynomial \[
 F=\sum_\alpha (a_\alpha/a_\beta)t^\alpha\in R[t_1,\ldots,t_n],
 \qquad f=a_\beta F,\qquad
 a_\beta\ne0,\quad [t^\beta]F=1.
 \tag{8}
\] Equation (8) is an explicit factorization of the original polynomial, not permission to replace it without its factor. In \(S_g\otimes_R K\), the image of \(a_\beta\) is a unit, so the original equation \(f=0\) gives \(F=0\). The map \(S_g\to S_g\otimes_R K\) is injective since \(S_g\) is a domain containing \(R\); therefore the same \(F=0\) holds in \(S_g\) itself. Its residue polynomial \(\bar F\in k[t_1,\ldots,t_n]\) is nonzero, because its coefficient at \(t^\beta\) is \(1\).

The special-fibre chart consequently factors through \[
 k[t_1,\ldots,t_n]/(\bar F)\longrightarrow S_g\otimes_R k.
\] This map is quasi-finite: its fibres are the same as those of the original chart over primes containing \(\bar F\), and it is of finite type. A nonzero polynomial quotient has dimension at most \(n-1\): for any prime containing \(\bar F\), its image domain has fraction field generated by \(n\) elements satisfying the nonzero relation \(\bar F\), so its transcendence degree is at most \(n-1\); the dimension theorem for finite-type field domains then bounds every irreducible component. Applying (2) to this quasi-finite map bounds the entire special chart by \(n-1\). Its point at the original \(\mathfrak q\) still has point dimension \(n\), since the chart is an open neighbourhood. This is the contradiction. The case \(n=0\) needs no polynomial argument, since \(d\geq0\) already prevents \(d<n\).

Therefore every point of the nonempty special fibre has point dimension \(d\). This proves equidimensionality as well: for each irreducible component of the finite-type field fibre, its generic point lies on no other component, and an open neighbourhood obtained by removing the other finitely many components has the dimension of that component. Its point dimension is therefore that same dimension, which was proved to equal \(d\). The special fibre is nonempty, so its dimension is \(d\). The proposed period at 31386 only terminates this source equality.

\subsubsection{Every nonempty valuation fibre and the generically finite consequence}\label{algebra-proof-034-every-nonempty-valuation-fibre-and-the-generically-finite-consequence}

\textbf{ALGEBRA-RECON-783, editorial extension.} Under the same hypotheses, every nonempty fibre over every prime \(\mathfrak p\subset R\) is equidimensional of dimension \(d=\dim(S\otimes_R K)\). Indeed \(R_{\mathfrak p}\) is again a valuation ring, with fraction field \(K\) and residue field \(\kappa(\mathfrak p)\). The actual localization \[
 S_{\mathfrak p}=(R\setminus\mathfrak p)^{-1}S
\] is a domain of finite type containing \(R_{\mathfrak p}\), and \[
 S_{\mathfrak p}\otimes_{R_{\mathfrak p}}\kappa(\mathfrak p)
       =S\otimes_R\kappa(\mathfrak p),\qquad
 S_{\mathfrak p}\otimes_{R_{\mathfrak p}}K=S\otimes_R K
\] by the original scalar maps. Apply the proved special-fibre theorem to this pair. At every existing point \(\mathfrak q\) of \(S\), its fibre is nonempty, so its relative point dimension is exactly \(d\). Thus the source function is constant on \(\operatorname{Spec}(S)\), even when some base fibres are empty.

If \(\operatorname{Frac}(S)/K\) is algebraic, then the finite-type generic field algebra has dimension zero, so \(d=0\). Conversely \(d=0\) gives a finite algebraic fraction-field extension. The constant-dimension conclusion makes \(R\to S\) globally quasi-finite. ALGEBRA-RECON-764 then supplies a finite \(R\)-subalgebra \(A\subset S\) such that \(\operatorname{Spec}(S)\to\operatorname{Spec}(A)\) is an open immersion. This includes the actual finite-envelope maps; it is not a claim that \(S\) itself must be finite over \(R\).

Two boundary examples retain the source hypotheses that matter. For \(R=k[t]_{(t)}\) and \(S=R[1/t]=k(t)\), the closed fibre is zero while the generic fibre is nonzero; one must restrict the equality to nonempty fibres. If ``domain'' is weakened only to a reduced algebra with injective base map, take \[
 S=R[X]\times R/(t),\qquad
 r\longmapsto(r,\bar r).
\] This map is injective because of the first factor. The algebra is reduced, and its actual finite presentation is \[
 R[e,z]/(e^2-e,ez-z,t(1-e))\longrightarrow R[X]\times R/(t),
 \quad e\longmapsto(1,0),\quad z\longmapsto(X,0).
\] The quotient splits by its complementary idempotents \(e,1-e\). Its \(e\)-factor is \(R[z]\), and its \(1-e\)-factor has \(z=0\) and \(t=0\), hence is \(R/(t)\). These two descriptions give the inverse and prove the asserted presentation. Its generic fibre is \(K[X]\), of dimension one, while the special fibre is \(k[X]\times k\), with components of dimensions one and zero. Thus equidimensionality and constant point dimension fail. The source domain hypothesis excludes that vertical component; the proof did not silently discard it.

\subsubsection{Reading and propagation}\label{algebra-proof-034-reading-and-propagation}

The exact bounded dependencies reread are algebra.tex:28024--28054, 28073--28123, 28306--28356 and 27347--27358. Earlier groups 665, 667 and 673 supply the checked normalization arguments with all original coefficient factors, and 674 supplies the point-dimension comparison. Group 781 combines the source chart with the corrected noninjective polynomial bound. Group 782 carries the exact base-change equality into radical ideals and finite constructible strata, and reuses the original product-ring data from group 685 for the finite-presentation boundary. Group 783 applies the source valuation argument at every original base prime and carries its zero-dimensional case into the finite-envelope theorem 764.

The corpus query ``valuation fibre dimension'' has two research and two local-layer routing hits. The Bhatt--Scholze item is already present in the reading ledger with only lines 251--284 read earlier; this repeated hit adds no reading coverage or adopted result. The Bouis source and the two Split Support Geometry records remain unread and unadopted. No PDF fallback or novelty claim occurs. The existing source identities and exact historical coverage are retained.

Next is Algebras and modules of finite presentation at 31434. Lookahead 31434--31460 and its three displayed report hits remain unadjudicated. No source or translation mutation, TeX build or publication occurs in this review.

\hypertarget{algebra_finite_presentation_notes_20260928}{}
\subsection{Finite presentation: source review and separate editorial arguments}\label{algebra-proof-035-finite-presentation-source-review-and-separate-editorial-arguments}

Authority: Stacks Project authors, algebra.tex at commit a04446e57ec1fbc252a871afcec7752fb2807b14, lines 31434--31758; whole-file SHA256 FA8BB92E58A4F78A2BD01B3B6A4A87DE0A0D279F5DD90641B574DD5FBFFFA4F3. Bounded primary dependencies 539--563 and 26366--26399 were reread. All twenty-one received reports whose first locus is in this section were read. The previous reading summary missed OCC-00612; the source-ordered inventory includes it here.

This document is editorial evidence. The translated source retains the author's statements, hypotheses, argument structure and omitted-proof wording. The existing six canonical units 0667--0672 remain identifiable; their ten operations correct two index switches, the selected product factor, three variable bounds, two verbs and two coordinated replacement clauses. One terminal period is proposed. One stylistic rearrangement is declined as optional. All supplementary proofs and stronger statements below remain separate from both translation branches. No novelty claim is made. Broader synthesis follows completion of the core intake.

\subsubsection{The original descent families}\label{algebra-proof-035-the-original-descent-families}

Source 31442--31487; groups 784--785. In the first proof retain \[
y_j=\sum_i a_{ij}\otimes x_{ji},\qquad
A=R[x_{ji}:j,i]\subset S.
\] The finite families need not have interchangeable index ranges. MC-STK-ERR-0667 changes only the subsequent reference to this original family. Flatness injects \(R'\otimes_R A\) into \(R'\otimes_R S\), where its image contains every \(y_j\); hence the cokernel \((S/A)\otimes_R R'\) vanishes. Faithful flatness gives \(S/A=0\). The converse follows by tensoring an actual finite generating family.

In the second proof retain the original polynomial algebra \(P=R[x_1,\ldots,x_n]\), its kernel \(I\), and \[
g_j=\sum_i a_{ij}\otimes f_{ji},\qquad
J=(f_{ji}:j,i)\subset I.
\] MC-STK-ERR-0668 restores the reference to this family. Flatness identifies both tensor ideals inside \(R'[x_1,\ldots,x_n]\); since the first contains the generators \(g_j\) of the second, they coincide. Faithfulness applied to \(I/J\) yields \(I=J\). All coefficients and all components of the two tensor expressions are retained.

\subsubsection{The original module lifts and stalk specialization}\label{algebra-proof-035-the-original-module-lifts-and-stalk-specialization}

Source 31489--31578; group 786. For \(R'=S^{-1}(R/I)\) and \(M'=(R')^n/K'\), let \(K\) be the inverse image of \(K'\) in \(R^n\). The image of \(I R^n\) is zero, hence \(IR^n\subset K\). Each element of \(K'\) is a vector over \(R'\); multiplying its finitely many entries by one original denominator \(s\in S\) lifts it to \(v\in R^n\). Since \(K'\) is an \(R'\)-submodule, \(v\in K\), and the original element is the image of \(v/s\). Thus \(S^{-1}K\) maps onto \(K'\), and the quotient \(M=R^n/K\) has \(S^{-1}(M/IM)=M'\).

For the finitely presented case retain the actual matrix relation \[
A_{R'}=sA',\qquad s\in S,
\] and the actual generators \(x_i\mapsto x'_i\). Multiplication by \(s\) is an automorphism of \((R')^m\), so \(\operatorname{im}(A_{R'})=\operatorname{im}(A')\). Their cokernels are identified by the identity of the original \((R')^n\), which proves the asserted map. The factor \(s\) is neither discarded nor hidden in a replacement definition. Zero source or target ranks give the same argument with zero matrices.

For the approximation of an arbitrary \(M\), numerators \(x_i\in M\) of a finite localized generating family generate \(S^{-1}M\). The map \(R^n\to M\) need not be onto before localization. After localization it is onto, with kernel \(S^{-1}K\) by exactness of localization. Finite presentation of \(S^{-1}M\) makes that kernel finite. Choose numerators of its generators in the actual \(K\) and let \(K'\) be their finite span. Then \(M'=R^n/K'\to M\) induces the required isomorphism after localization. Taking \(S=R\setminus\mathfrak p\) proves the stalk statements. The only proposed text operation is the missing period after the lemma reference at 31577.

\subsubsection{Local isomorphisms with all original denominators}\label{algebra-proof-035-local-isomorphisms-with-all-original-denominators}

Source 31580--31660; groups 787--788. Write the original presentation \[
S=R[x_1,\ldots,x_n]/(g_1,\ldots,g_m)
\] and the images of the original \(x_i\) under \(S_{\mathfrak q}\cong R_{\mathfrak p}\) as \(a_i=b_i/f_0\), with \(f_0\notin\mathfrak p\). In \(S_{f_0}\) define \(\xi_i=x_i-a_i\). This is an explicit coordinate comparison; the original presentation remains recorded. The new relation is exactly \[
G_j(\xi_1,\ldots,\xi_n)=g_j(\xi_1+a_1,\ldots,\xi_n+a_n).
\] Retain every term of this polynomial and its constant \(c_j=g_j(a_1,\ldots,a_n)\). Since \(c_j\) vanishes in \(R_{\mathfrak p}\), choose an annihilator \(h_j\in R\setminus\mathfrak p\) of \(c_j\) in \(R_{f_0}\), and set \(f_1=f_0\prod_jh_j\). Evaluation at \(\xi_i=0\), equivalently at the original \(x_i=a_i\), defines \[
v:S_{f_1}\to R_{f_1},\qquad J=\ker v=(\xi_1,\ldots,\xi_n).
\] The composite into \(R_{\mathfrak p}\) agrees with the original isomorphism on \(S\), so the inverse image of \(\mathfrak pR_{f_1}\) is the original \(\mathfrak q\) localized. It extends to \(S_{\mathfrak q}\), since every element outside \(\mathfrak q\) maps to a unit of \(R_{\mathfrak p}\). Consequently \(S_{\mathfrak p}/J_{\mathfrak p}\cong R_{\mathfrak p}\cong S_{\mathfrak q}\) as \(S_{\mathfrak p}\)-algebras.

The last quotient is a localization of \(S_{\mathfrak p}\), hence flat. The ideal \(J_{\mathfrak p}\) is finitely generated and pure. The source lemma 26366--26399 gives an idempotent \(e\) with \(J_{\mathfrak p}=eS_{\mathfrak p}\). To record its actual spreading, write \(e=z/t\), \(t\in R\setminus\mathfrak p\). Its equation \(e^2-e=0\) means an additional \(u\in R\setminus\mathfrak p\) kills \(z^2-tz\). After inverting \(f_1tu\), \(e'=z/t\) is idempotent. Each of the finitely many \(\xi_i\) lies in \(e'S_{\mathfrak p}\), and \(e'\) is a finite combination of them there. Take the product of all coefficient denominators and all annihilators of the resulting differences, each outside \(\mathfrak p\). With this further factor included in \(f\), the actual ideals satisfy \(J_f=e'S_f\).

Thus the exact maps are \[
S_f\longrightarrow R_f\times e'S_f,\quad z\longmapsto(v(z),e'z),
\qquad
(r,c)\longmapsto(1-e')\varphi(r)+c.
\] They are inverse: \(v(e')=0\), the kernel of \(v\) is \(e'S_f\), and \((1-e')\varphi(v(z))=(1-e')z\). The ring \(e'S_f\) has unit \(e'\) and its \(R_f\)-structure is \(r\mapsto e'\varphi(r)\).

For the next lemma, retain the original \(S,S'\) and \(\psi:S_{\mathfrak q}\to S'_{\mathfrak q'}\). If \(\psi(x_i)=h_i/g_i\), put \(d=\prod_i g_i\) in \(S'\); it is outside \(\mathfrak q'\). The polynomial presentation of \(S'_d\) adjoins a variable \(Y\) with relation \(dY-1\), and the image of \(x_i\) is \(h_i(\prod_{k\ne i}g_k)Y\). Each original \(f_j\) evaluated at these images becomes zero at \(\mathfrak q'\). Kill its numerator by an element outside that prime and multiply those finitely many annihilators to get a further denominator. These are exactly the two coordinated replacement operations corrected by MC-STK-ERR-0672; they do not alter the original isomorphism.

This constructs \(u:S\to S'_D\) for some original \(D\notin\mathfrak q'\). The finite-presentation composition lemma 539--563 applies. The preceding local result supplies \(a\in S\setminus\mathfrak q\) and \[
(S'_D)_{u(a)}\cong S_a\times C.
\] Write \(u(a)=b/D^r\), so the left ring is \(S'_{Db}\) and \(b\notin\mathfrak q'\). Its relevant prime is on the first factor: it is the inverse image of the original prime under the evaluation map used above. Therefore \(e=(1,0)\) lies outside that prime. MC-STK-ERR-0669 selects this \(S_a\) factor; \((0,1)\) selects \(C\), the wrong factor. Write \(e=c/(Db)^s\) with \(c\in S'\). Then \(c\notin\mathfrak q'\), and inverting \(e\) is precisely passage from \(S'_{Db}\) to \(S'_{Dbc}\). The result uses the original principal neighborhoods \(S_a\cong S'_{Dbc}\), not unnamed replacement rings.

\subsubsection{Nilpotent lifting and the source automorphism}\label{algebra-proof-035-nilpotent-lifting-and-the-source-automorphism}

Source 31662--31695; group 789. For nilpotent \(I\), lift the original finite generating family of \(S/IS\) and let \(A\subset S\) be its \(R\)-subalgebra. The \(R\)-module \(M=S/A\) satisfies \(M=IM\). Iterating gives \(M=I^NM=0\) for the original nilpotence exponent \(N\); finite generation of \(M\) is not needed.

In the locally nilpotent argument keep the source presentation \(S'=R[x_1,\ldots,x_m]/K\) and all polynomials \[
g_j=x_j+\sum_J\delta_{j,J}x^J,\qquad \delta_{j,J}\in I.
\] MC-STK-ERR-0670 changes the three unbound \(n\)'s to the original \(m\). Only finitely many displayed coefficients are nonzero. Enumerate them as \(\delta_1,\ldots,\delta_t\), choose \(n_\ell\ge1\) with \(\delta_\ell^{n_\ell}=0\), and set \[
D=(\delta_1,\ldots,\delta_t),\qquad N=1+\sum_{\ell=1}^t(n_\ell-1).
\] Every degree-\(N\) product in these ideal generators contains some \(\delta_\ell\) at least \(n_\ell\) times, so \(D^N=0\). For an empty coefficient list take \(D=0,N=1\).

Let \(F:R[x_1,\ldots,x_m]\to R[x_1,\ldots,x_m]\) send \(x_j\) to the full \(g_j\), and let \(\Delta=F-\operatorname{id}\) as an \(R\)-linear map. For every monomial, replacing each \(x_j\) by \(x_j+\sum_J\delta_{j,J}x^J\) shows that the difference lies in \(D R[x]\). Since \(F\) fixes \(R\), \[
\Delta(D^kR[x])\subset D^{k+1}R[x],\qquad \Delta^N=0.
\] The finite sum \[
H=\sum_{k=0}^{N-1}(-1)^k\Delta^k
\] satisfies \(HF=FH=\operatorname{id}\) by multiplication of these exact additive operators; the terminal term is \((-1)^{N-1}\Delta^N=0\). Thus \(F\) is bijective. Its inverse is a ring homomorphism: applying injective \(F\) to \(H(ab)\) and to \(H(a)H(b)\) gives \(ab\) in both cases, and the same argument proves addition and the unit. Its coordinates are the explicit polynomials \(H(x_j)\). Hence the \(g_j\) generate the original polynomial ring, so their classes in the image of \(S\to S'\) generate \(S'\). This supplies the omitted source argument as an editorial proof. It does not replace the author's Noetherian-subring argument in the translation.

\subsubsection{The original isomorphism criteria}\label{algebra-proof-035-the-original-isomorphism-criteria}

Source 31697--31753; groups 790--791. The two missing verbs are retained as MC-STK-ERR-0671. In the local lemma let \(J=\ker(S\to S')\). The source composition lemma makes the surjection finitely presented, so \(J\) is a finitely generated ideal. Localization gives the exact sequence \[
0\to J_{\mathfrak q}\to S_{\mathfrak q}\to S'_{\mathfrak q}\to0.
\] Tensoring over the original \(R\) with \(R/I\) remains left exact at \(J_{\mathfrak q}/IJ_{\mathfrak q}\) because the stated \(R\)-flatness of \(S'_{\mathfrak q}\) makes the preceding Tor group vanish. The assumed quotient isomorphism therefore gives \(J_{\mathfrak q}/IJ_{\mathfrak q}=0\). Since \(IS_{\mathfrak q}\) is contained in its maximal ideal and \(J_{\mathfrak q}\) is finite, Nakayama gives \(J_{\mathfrak q}=0\). For actual generators \(j_1,\ldots,j_t\) of \(J\), choose \(u_i\in S\setminus\mathfrak q\) with \(u_i j_i=0\). Their product \(g=\prod_i u_i\) kills \(J\), proving \(S_g\cong S'_g\). If the generator list is empty use \(g=1\).

In the globally locally nilpotent case the preceding surjectivity result applies since \(S'\) is finitely presented, in particular finite type. Every prime contains \(IS\), so the local result applies everywhere. If \(j\in J\) were nonzero, its proper annihilator would be contained in a maximal ideal \(\mathfrak q\), and \(j/1\ne0\) there: an element outside \(\mathfrak q\) cannot annihilate it. This contradicts the local vanishing. Thus the global map is an isomorphism. The source's phrase ``for example by'' has an unambiguous inference and citation; OCC-00612 proposes an optional rearrangement, not a mathematical or grammatical defect. No operation is adopted for it.

\subsubsection{Algebra descent along the original pure map}\label{algebra-proof-035-algebra-descent-along-the-original-pure-map}

Source 31442--31487; group 792, receiving the exact unit-map mechanism of ALGEBRA-RECON-501. Suppose \(R\to R'\) is pure, meaning every \(R\)-module \(M\) has an injective unit map \(M\to R'\otimes_R M\). For an arbitrary \(R\)-algebra \(S\), finite type and finite presentation each hold for \(S/R\) if and only if they hold for \(S'=R'\otimes_R S\) over \(R'\).

For finite type use precisely the source's \(y_j=\sum_i a_{ij}\otimes x_{ji}\) and \(A=R[x_{ji}]\). The image of \(R'\otimes_R A\to S'\) contains every \(y_j\), so it is all of \(S'\). Right exactness implies \(R'\otimes_R(S/A)=0\). Purity of the unit on the original module \(S/A\) gives \(S=A\). No injection of \(R'\otimes_R A\) was presumed.

For finite presentation use the resulting \(P=R[x_1,\ldots,x_n]\twoheadrightarrow S\), kernel \(I\). Its base change has kernel \[
L=\operatorname{im}(R'\otimes_R I\to R'\otimes_R P),
\] by right exactness. This is a finitely generated ideal because the target \(S'\) is finitely presented. To justify this for the chosen original generators, take any finite presentation \(S'=R'[y_1,\ldots,y_b]/(h_1,\ldots,h_c)\). Express the images of \(x_i\) as \(u_i(y)\), and the images of \(y_j\) as \(v_j(x)\). The kernel of \(R'[x]\to S'\) is generated by the finite list \[
h_\ell(v_1(x),\ldots,v_b(x)),\qquad x_i-u_i(v_1(x),\ldots,v_b(x)).
\] Indeed these vanish in \(S'\); in their quotient, \(y_j\mapsto v_j(x)\) defines an inverse to the map onto \(S'\), with the composites fixed on every original generator. This proves the finite-kernel claim without flatness.

Choose generators \(g_j\) of \(L\), express each as the image of \(\sum_i a_{ij}\otimes f_{ji}\), \(f_{ji}\in I\), and let \(J=(f_{ji})\subset I\). The image of \(R'\otimes_R J\) is exactly \(L\). Consequently the actual surjection \(P/J\to S\) becomes an isomorphism after tensoring. An element in its kernel has zero image in \(R'\otimes_R(P/J)\), because that last map is an isomorphism. Purity on \(P/J\) makes the element zero. Hence \(I=J\), and \(S\) is finitely presented. The converse in both cases is direct base change of the finite generators and relations.

This strictly weakens faithful flatness: \(R=\mathbf Z\to R'=\mathbf Z\times\mathbf Z/2\mathbf Z\), \(n\mapsto(n,\bar n)\), has a ring retraction to the first coordinate, so every module unit is split injective. It is not flat, since multiplying by two kills the nonzero \((0,\bar1)\). All original source hypotheses remain in the translation; this is a separate generalization, with the tensor-kernel injection deliberately replaced by its proved image formula.

\subsubsection{Finite inverse for endomorphisms over a nil ideal}\label{algebra-proof-035-finite-inverse-for-endomorphisms-over-a-nil-ideal}

Source 31677--31695; group 793. Let \(A\) be a finite-type \(R\)-algebra, \(I\subset R\) a nil ideal, and \(\alpha:A\to A\) an \(R\)-algebra endomorphism inducing the identity on \(A/IA\). Then \(\alpha\) is an automorphism.

Fix actual generators \(a_1,\ldots,a_m\) and write \[
\alpha(a_j)-a_j=\sum_{\ell=1}^{t_j}d_{j\ell}b_{j\ell},\qquad
d_{j\ell}\in I,\ b_{j\ell}\in A.
\] Let \(D\) be the ideal generated by all these actual \(d_{j\ell}\). The same sum of individual nilpotence exponents as above gives \(D^N=0\). The maps \(\alpha\) and the identity agree on all generators modulo \(DA\), hence on all of \(A/DA\). With \(\Delta=\alpha-\operatorname{id}\), \(R\)-linearity gives \[
\Delta(D^kA)\subset D^{k+1}A,\quad
\alpha^{-1}=\sum_{k=0}^{N-1}(-1)^k\Delta^k.
\] The exact operator multiplication and multiplicativity argument in the preceding source-proof section apply to this original \(A\), including all its relations. This proves the statement without finite presentation or a globally nilpotent \(I\). If \(I^N=0\) itself, finite type is unnecessary: the same filtration works for any \(R\)-algebra \(A\).

For the source polynomial map let \(d=\max(1,\deg g_1,\ldots,\deg g_m)\). Expanding substitution gives \(\deg\Delta(p)\le d\deg p\) for positive-degree \(p\), and \(\Delta\) kills constants. Induction gives \(\deg\Delta^k(x_j)\le d^k\). Thus the displayed exact inverse has degree at most \(d^{N-1}\) on each generator. This bound supplements the full finite inverse; it does not replace its coefficients or terms. This result asserts no lifting of arbitrary automorphisms of \(A/IA\): an actual endomorphism \(\alpha\) with the stated reduction is part of the hypothesis.

\subsubsection{The exact Tor obstruction and weaker sufficient hypotheses}\label{algebra-proof-035-the-exact-tor-obstruction-and-weaker-sufficient-hypotheses}

Source 31697--31753; group 794. Retain an arbitrary \(R\)-algebra surjection \(S\to S'\) with finitely generated kernel \(J\), a prime \(\mathfrak q\supset IS\), and the source isomorphism modulo \(I\) at \(\mathfrak q\). The long exact tensor sequence contains the actual map \[
\delta:\operatorname{Tor}_1^R(S'_{\mathfrak q},R/I)
\longrightarrow J_{\mathfrak q}/IJ_{\mathfrak q}
\longrightarrow S_{\mathfrak q}/IS_{\mathfrak q}
\longrightarrow S'_{\mathfrak q}/IS'_{\mathfrak q}\longrightarrow0.
\] The last arrow is an isomorphism, so \(\delta\) is onto. Therefore \(\delta=0\) if and only if \(J_{\mathfrak q}/IJ_{\mathfrak q}=0\). By the finite-generation and local Jacobson hypotheses, Nakayama makes this equivalent to \(J_{\mathfrak q}=0\). The original finite list of kernel generators and its annihilator product then show that it is equivalent to the existence of \(g\notin\mathfrak q\) with \(S_g\cong S'_g\). Conversely such an isomorphism makes \(J_{\mathfrak q}=0\), hence \(\delta=0\). This proves the exact obstruction, not only a one-way sufficient condition.

In particular \(\operatorname{Tor}_1^R(S'_{\mathfrak q},R/I)=0\) suffices in place of \(R\)-flatness, and finite generation of \(J\) suffices in place of the pair of finite-type/finite-presentation assumptions used to obtain it. Vanishing of the whole Tor group is stronger than necessary: for \(R=k[\epsilon]/(\epsilon^2)\), \(I=(\epsilon)\), and the identity \(S=S'=k\), the kernel and \(\delta\) vanish, although \(\operatorname{Tor}_1^R(k,k)=k\).

The flatness assumption cannot simply be removed with no substitute. For that same original \(R,I\), take instead \(S=R\to S'=k\), \(\mathfrak q=(\epsilon)\). Both finite-presentation assumptions hold and reduction modulo \(I\) is an isomorphism. The resolution \[
\cdots\xrightarrow{\epsilon}R\xrightarrow{\epsilon}R
\xrightarrow{\epsilon}R\longrightarrow k\longrightarrow0
\] is exact because both the kernel and image of multiplication by \(\epsilon\) equal \((\epsilon)\). Tensoring with \(k\) gives zero differentials, so the first Tor group is \(k\). For the usual boundary convention, lifting its generator to \(1\in R\) and applying the differential gives \(\delta(1)=\epsilon\) in \(J/IJ=(\epsilon)\). The map is nonzero; every \(g\notin\mathfrak q\) is a unit, so no such localization kills this kernel.

The global version is also exact. For a surjection with finite kernel, a global isomorphism modulo \(I\), and \(IS\subset\operatorname{Jac}(S)\), the global boundary \[
\operatorname{Tor}_1^R(S',R/I)\longrightarrow J/IJ
\] is onto, and it vanishes if and only if \(J=0\), by the same Nakayama argument. Consequently the source locally nilpotent theorem still holds if its flatness hypothesis is replaced by \(\operatorname{Tor}_1^R(S',R/I)=0\): the preceding source lemma first proves surjectivity, finite presentation gives finite \(J\), and the nil ideal \(IS\) lies in every maximal ideal. This propagates the local weakening through the actual global receiving result.

\subsubsection{Finite type does not replace finite presentation in the local splitting}\label{algebra-proof-035-finite-type-does-not-replace-finite-presentation-in-the-local-splitting}

Source 31580--31660; group 795, using and extending the exact product-ring example in ALGEBRA-RECON-685. Fix a field \(k\) and retain \[
R=\prod_{i\ge1}k,\qquad I=\bigoplus_{i\ge1}ke_i,\qquad S=R/I.
\] The quotient \(R/I\ne0\) has a maximal ideal; let \(\mathfrak p\subset R\) be its preimage and \(\mathfrak q=\mathfrak p/I\). Every finite-support element \(z\in I\) is killed by \(1-e_F\), where \(F\) contains its support and \(e_F=\sum_{i\in F}e_i\in I\subset\mathfrak p\). Thus \(1-e_F\notin\mathfrak p\), and \(I_{\mathfrak p}=0\). The original quotient map induces \[
R_{\mathfrak p}\xrightarrow{\sim}S_{\mathfrak q}.
\] The \(R\)-algebra \(S\) is finite type, even cyclic as an \(R\)-module. It is flat: it is the filtered colimit of \(R/(e_F)\) over finite \(F\), each an idempotent direct summand of \(R\). Tensoring an injection with each such summand preserves it, and filtered colimits preserve injections of modules, proving the flatness assertion. It is not finitely presented, since a finite family in \(I\) has support in one finite \(F\) and cannot generate any \(e_i\) outside it.

For any \(f=(f_i)\notin\mathfrak p\), its support \(A=\{i:f_i\ne0\}\) is infinite, since every finite-support element is in \(\mathfrak p\). The exact localization maps from group 685 give \[
R_f\cong\prod_{i\in A}k,\qquad I_f\cong\bigoplus_{i\in A}ke_i.
\] Here \(I_f\ne0\). If an \(R_f\)-algebra isomorphism \(S_f\cong R_f\times C\) existed, the surjection \(R_f\to S_f\) would make the structural map \(r\mapsto(r,\eta(r))\) onto the product. A preimage of \((1,0)\) must be \(r=1\), giving \(1_C=0\). Hence \(C\) would be the zero ring, and that same structural map would force \(I_f=0\), a contradiction. Thus even a flat finite-type map with the given stalk isomorphism need not split on any such base principal neighborhood. The finite-presentation hypothesis is preserved; the failure of this stronger assertion is an editorial boundary, not an error in the source theorem.

\subsubsection{Corpus use and continuation}\label{algebra-proof-035-corpus-use-and-continuation}

The existing canonical disk index was queried in both layers for nilpotent polynomial automorphism. The Bell--Hamidizadeh--Huang--Venegas TeX routing record, the PDF-only Nonstandard analysis record and the two local Split Support records are routing hits only; their contents were not read or adopted. The seven earlier externally read source records retain exactly their prior bounded coverage. No PDF fallback or literature-novelty claim is made.

The current section and twenty-one reports are adjudicated by groups 784--795. Source lookahead 31759--31990 was read earlier but remains unadjudicated, with no report dispositions inferred from it. Continue at Colimits and maps of finite presentation, authority 31759. No source or translation is changed by this review, and no TeX build or publication is claimed.

\hypertarget{algebra_colimit_categories_notes_20260928}{}
\subsection{Colimits and categories of finite presentations}\label{algebra-proof-036-colimits-and-categories-of-finite-presentations}

Authority: Stacks Project authors, algebra.tex at a04446e57ec1fbc252a871afcec7752fb2807b14, lines 31759--32217, SHA256 FA8BB92E58A4F78A2BD01B3B6A4A87DE0A0D279F5DD90641B574DD5FBFFFA4F3. Categories dependencies 2130--2147 and 2229--2243 at that same commit were read. The latter gives a reference to Kashiwara--Schapira, Proposition 3.2.4; that cited book passage was not read here. The full argument needed below is proved directly.

All fifty received reports in this interval were read. The source continues with local-ring approximations at 32218; the already read source interval 32218--32380 remains unadjudicated. This batch is the first part of the longer section, not its completion.

All extended arguments below are editorial. The translated author text, its sequence and its hypotheses remain identifiable. Twelve canonical correction units 0399--0410 are retained, comprising sixteen operations. Eight new small source operations are proposed. Two proposed changes of the original B-colimit are rejected. Two requests to spell out an implicit composite map are optional; the exact finite-stage witness construction is recorded separately. No novelty is claimed.

\subsubsection{Filtered presentations with their actual structural maps}\label{algebra-proof-036-filtered-presentations-with-their-actual-structural-maps}

Source 31759--31877 and 31946--31960; groups 796--803. The opening count noun becomes ``results''. The first object of the category of finitely presented \(R\)-algebras mapping to \(A\) is the actual structural map \(R\to A\), not \(R\to R\). This fixes an unreported wrong target at 31782. The source set-size convention can be implemented by choosing presentations with finitely many named variables within the countable list; tensor products and quotients are replaced by their specified isomorphic presentations when needed.

Given two objects \(A'\to A,A''\to A\), the object \(A'\otimes_R A''\to A\) has structure map \(a'\otimes a''\mapsto f(a')g(a'')\). Tensoring the finite presentations and retaining both lists of generators and relations proves finite presentation. The two factor maps give a common successor. For parallel maps \(\sigma,\tau:A'\to A''\), keep the source generators \(x_1,\ldots,x_r\) and quotient by the actual differences \(\sigma(x_i)-\tau(x_i)\). This quotient has a finite presentation over \(R\), maps to \(A\), and equalizes the maps on every polynomial in the original generators. The indexing category is therefore filtered in exactly the variance of Categories 2130--2147. MC-STK-ERR-0400 corrects ``cofiltered''; 0399 restores the index set \(\Lambda\).

For completeness, the set colimit of a filtered ring diagram has a ring structure: represent finitely many elements at a common successor, add or multiply there, and pass to their equivalence classes. Two choices agree at a further successor by filteredness, including equalization of parallel arrows. The ring laws hold at such common successors. This construction has the ring-colimit universal property because a compatible family of ring maps preserves the operations on all representatives. Equality with zero in this colimit means vanishing after an arrow to a successor.

Every \(a\in A\) is the image of \(X\) under \(R[X]\to A\). If an element represented by \(a'\in A'\) maps to zero, pass to the actual finitely presented quotient \(A'/(a')\to A\), which kills it. Thus the colimit map is bijective as a set map and an isomorphism of rings.

In the explicit partially ordered construction retain \(A_{S,E}=R[X_s:s\in S]/(f_e:e\in E)\). Unions of two finite generator sets and the images of their two relation lists give a common upper bound. Every element of \(A\) has a representative \(X_a\). A polynomial in a finite generator set that maps to zero can be added as one more original relation, proving injectivity of the colimit map. For finite-type \(A=R[x_1,\ldots,x_n]/I\), fix the actual finite generating set and increase only the finite relation list; every element of \(I\) is eventually included, so the resulting colimit is \(A\) and all transition maps are surjective. This restricted indexing set need not be cofinal in the set of all arbitrary finite subsets of \(A\); the direct generator-and-relation argument is what proves its colimit. The remaining new operations here are ``be'', ``equal to'', ``a given type'' and the source's standard \(R\)-algebra spelling.

\subsubsection{The Hom comparison, retractions and cofinality}\label{algebra-proof-036-the-hom-comparison-retractions-and-cofinality}

Source 31879--31980. For a finitely presented algebra \(S=R[x_1,\ldots,x_n]/(f_1,\ldots,f_m)\), two maps into possibly different stages with equal colimit maps agree at some common stage: first put the stages above a common upper bound, then kill the finitely many differences of the images of the original \(x_i\). This proves injectivity of the Hom comparison. A map into the colimit is represented on all \(x_i\) at one stage; each of the finitely many actual \(f_j\) then has zero colimit value, so a later common stage kills them all. This proves surjectivity. Empty generator and relation lists cause no problem because the directed indexing set is nonempty.

Conversely, if the comparison is surjective for every directed system, express the original \(S\) as the colimit of finitely presented algebras just constructed. Its identity factors as \[
S\xrightarrow{i}B\xrightarrow{p}S,\qquad pi=\operatorname{id}_S,
\] with \(B\) finitely presented over \(R\). To prove the retract statement explicitly, write \(B=R[b_1,\ldots,b_n]/(F_1,\ldots,F_m)\) and choose actual polynomials \(H_j(b)\) representing \(i(p(b_j))\). The kernel of \(p\) is generated by \[
b_j-i(p(b_j))=b_j-H_j(b),\qquad 1\le j\le n.
\] Indeed for any polynomial \(Q\), telescoping replacement of its variables shows \(Q(b)-Q(i p(b))\) belongs to this ideal, term by term. For \(z\in\ker p\), this difference is \(z-i p(z)=z\). The reverse inclusion follows from \(pi=\operatorname{id}\). Hence the original \(S\) has the finite presentation with all \(F_\ell\) and all \(b_j-H_j(b)\). This proves the source conclusion and the finite-presentation comparison without concealing a kernel assumption.

For the class \(\mathcal E\), the factorization property implies that its full subcategory \(\mathcal J\) inside the filtered presentation category \(\mathcal I\) is nonempty. Common successors for two objects of \(\mathcal J\) are obtained first in \(\mathcal I\) and then mapped to an object of \(\mathcal J\); the same procedure after an equalizing arrow proves filteredness. For every \(i\in\mathcal I\), the category of arrows \(i\to j\), \(j\in\mathcal J\), is nonempty and connected: given two such arrows, take a common successor, then equalize their two composites out of \(i\), and map the resulting object to \(\mathcal J\). Thus compatible cocones on \(\mathcal J\) extend uniquely to \(\mathcal I\): define the value at \(i\) via any such arrow; connectedness makes it independent, and the same argument gives naturality. This is the required cofinality and proves the source's colimit criterion.

\subsubsection{Module maps and the exact finite-stage witnesses}\label{algebra-proof-036-module-maps-and-the-exact-finite-stage-witnesses}

Source 31982--32100; groups 804--808. The variance in the introductory categorical sentence is colimit; the proposed operation changes ``limit of the category'' to ``colimit of the categories''. The precise categorical construction is given below rather than silently replacing the source exposition.

For fixed \(A\)-modules \(M,N\) and directed \(A\)-algebras \(R_i\) with colimit \(R\), tensoring identifies \(N\otimes_A R\) with \(\operatorname{colim}_i(N\otimes_A R_i)\), and likewise for \(M\). One can verify this from the tensor universal property, or from the presentation of a tensor product by generators and its finite bilinearity relations. If \(M\) has generators \(x_1,\ldots,x_m\), equality of two maps after tensoring with \(R\) kills the finite list \((u-u')(x_j)\otimes1\) at a common stage. These elements lie in \(N\otimes_A R_i\), as retained correction 0401 states.

For surjectivity, keep the original finite generators \(y_j\) of \(N\). Surjectivity at \(R\) gives actual \(w_j\in M\otimes_A R\) with \(u_R(w_j)=y_j\otimes1\). The colimit for \(M\) represents all \(w_j\) by \(z_j\in M\otimes_A R_i\) at one stage. Then form the actual errors \[
d_j=u_i(z_j)-y_j\otimes1\in N\otimes_A R_i.
\] The colimit for \(N\) kills these errors at a later common stage \(k\ge i\). Thus \(u_k(z_j\otimes1)=y_j\otimes1\) in \(N\otimes_A R_k\), and \(u_k\) is surjective. Images in the source prose already mean the evident composite \(M_i\to N_i\to N_R\); spelling out that composite is optional. However, lifting preimages and annihilating their finitely many errors are distinct steps. The supplementary proof records both.

The first representing stage need not itself work. Take \(A=k[t]\), \(M=N=A\), \(u\) multiplication by \(1+t\), and the two-stage system \(R_0=A\to R_1=A/(t)=R\). The element \(z=1\in M\otimes_A R_0\) maps under \(u\) to \(1\) in \(N\otimes_A R\), but \(u_0\) is not onto since \(1+t\) is not a unit of \(k[t]\). The later stage kills the error \(t\). This does not contradict the source existence theorem; it identifies the required finite refinement.

For map lifting, use exactly the original presentation \[
A^t\xrightarrow{(k_{sj})}A^m\longrightarrow N\longrightarrow0,
\qquad (a_1,\ldots,a_m)\longmapsto\sum_j a_jy_j.
\] Correction 0402 restores \(y_j\). Lift \(v(y_j\otimes1)\) to \(z_j\) at one stage and kill all the actual relation errors \(\xi_s=\sum_jk_{sj}z_j\) at a later stage. Right exactness of tensoring, not injectivity of the tensor of the kernel, then gives the required map. For the isomorphism assertion, first descend the inverse. Over a later stage apply the equality assertion with that stage as the new coefficient ring to each of the two composites and its identity. The finite module \(M\) and finite module \(N\) provide the two finite generator lists. Correction 0403 keeps both tensor bases equal to \(A\).

Finally, a presentation matrix of a finitely presented \(R\)-module has finitely many original entries. Lift them together to \(R_i\), with no discarded entry or factor; its cokernel is the descended module. The preceding map and equality arguments establish all three source categorical assertions.

\subsubsection{Algebra maps, rejected target changes and inverse identities}\label{algebra-proof-036-algebra-maps-rejected-target-changes-and-inverse-identities}

Source 32102--32217; groups 809--816. For equality of algebra maps \(u,u':B\to C\), the finite list of generator images lies in \(C\), so correction 0404 properly changes both the colimit and equality target to \(C\).

For surjectivity, the original \(B\)-colimit at 32143 is correct and needed to lift actual preimages of the target generators. Pick \(y_j\) generating \(C\); surjectivity over \(R\) gives \(w_j\in B\otimes_A R\) mapping to them. The \(B\)-colimit puts these preimages at one stage. The \(C\)-colimit subsequently kills the finitely many errors \(u_i(z_j)-y_j\otimes1\). Thus the reports OCC-01451 and OCC-12061 incorrectly treat the \(B\)-colimit as a wrong target. Replacing it by only the \(C\)-colimit removes the stated reason that the preimages themselves descend. Both colimits participate, with these different exact maps. The requested explicit composite in OCC-00626 is optional and is covered in the separate proof.

The same early-stage pitfall occurs with \(A=k[t]\), \(B=A[x]\), \(C=A[y]\), \(u(x)=(1+t)y\), and \(R_0=A\to R_1=A/(t)\). The element \(x\) maps to \(y\) after base change to \(R_1\), but \(u_0\) is not onto: applying \(A\to k\), \(t\mapsto-1\), would make an alleged preimage of \(y\) express \(y\) as a constant in \(k[y]\). The refinement kills the actual error \(ty\).

For map lifting retain the presentation \[
A[x_1,\ldots,x_m]\longrightarrow C,\quad x_j\longmapsto c_j,
\qquad K=(f_1,\ldots,f_t).
\] Corrections 0405 and 0406 restore the codomain \(C\) and the individual generator images. In particular, sending every variable to the sum is not an equivalent presentation. Lift the images of \(c_j\) under \(v\) to \(z_j\in B\otimes_A R_i\), then annihilate the original finite list \(f_s(z_1,\ldots,z_m)\). Correction 0407 aligns the subscript \(s\). Right exactness identifies the quotient by the image of \(K\otimes_A R_i\); no left-exactness assumption is used.

The inverse to \(u_R:B_R\to C_R\) is \(v:C_R\to B_R\), as correction 0408 states. Descend this map and apply generator-wise eventual equality over the later stage to both \(v_i u_i\) and \(u_i v_i\). Their identities have the original base \(A\), as correction 0409 records. The supplementary minimal clarification at 32124 explicitly declares the \(A\)-algebra map \(u:B\to C\), which item (4) uses but did not quantify in that item. It supplies no new hypothesis beyond the kind of map the displayed tensor formulas require.

For the final categorical lemma lift every coefficient of every original relation polynomial \(f_j\) to one stage, preserving its monomials, signs and coefficients. The resulting finite presentation tensors to the original one. Arrow descent and eventual equality follow from the proved map statements. Correction 0410 retains the originally quantified \(\varphi_\lambda,\psi_\lambda\) in all four occurrences.

\subsubsection{Finite type and the boundary between unions and arbitrary colimits}\label{algebra-proof-036-finite-type-and-the-boundary-between-unions-and-arbitrary-colimits}

Source 31879--31944; group 817. An \(R\)-algebra \(S\) is of finite type if and only if, for every directed system of \(R\)-algebras with injective transition maps, the canonical map \[
\operatorname{colim}_i\operatorname{Hom}_R(S,A_i)
\longrightarrow\operatorname{Hom}_R(S,\operatorname{colim}_i A_i)
\] is bijective. Surjectivity for these systems alone already suffices.

If \(S\) has finite generators, equality of maps descends by the finite list argument. For a map to the union, put all generator images at one stage. Every original polynomial relation vanishes in the union, hence at that same stage because its map into the union is injective. This kills even an infinite relation ideal without asserting finite presentation; the original quotient map therefore factors through this stage. To prove the converse, write \(S\) as the directed union of its finitely generated \(R\)-subalgebras. The identity factors through one of them by the assumed surjectivity. Since the structural map of that stage is the actual inclusion into \(S\), its image is all of \(S\), proving finite type.

The distinction from arbitrary transition maps is strict. Use the earlier original ring \(R=\prod_{n\ge1}k\), its ideal \(I=\bigoplus_{n\ge1}ke_n\), and \(S=R/I\). For finite \(F\subset\mathbf N\), let \(J_F=(e_n:n\in F)\) and \(A_F=R/J_F\). Quotient transitions have colimit \(R/I=S\), because every finite-support element is eventually killed. An \(R\)-algebra map \(S\to A_F\) would imply \(I\subset J_F\), as it must take the image of each \(r\in R\) to its original image modulo \(J_F\). For \(n\notin F\), \(e_n\in I\setminus J_F\), a contradiction. Thus every \(\operatorname{Hom}_R(S,A_F)\) is empty, while the target Hom set contains \(\operatorname{id}_S\). The finite-type algebra \(S\) preserves the directed-union comparison but fails the arbitrary-colimit one. This proves the exact role of finite presentation without weakening the source theorem.

\subsubsection{Descent of finite diagrams and finite equations}\label{algebra-proof-036-descent-of-finite-diagrams-and-finite-equations}

Source 31988--32216; group 818. Let \(R=\operatorname{colim}_iR_i\). Define a category \(\mathcal C\) whose objects are pairs \((i,M_i)\), with \(M_i\) a finitely presented \(R_i\)-module. For two objects its morphism set is \[
\operatorname{Hom}_{\mathcal C}((i,M_i),(j,N_j))
 =
\operatorname{colim}_{k\ge i,j}
\operatorname{Hom}_{R_k}(M_i\otimes_{R_i}R_k,N_j\otimes_{R_j}R_k).
\] Composition is represented at a common successor and is independent of representatives, since a later stage preserves each equality. Identities are the original identities. The functor to finitely presented \(R\)-modules sends an object to its actual base change and a morphism to the base-changed map. It is essentially surjective by matrix descent, full by map descent over a common starting stage, and faithful by eventual equality. This is the precise equivalence intended by the source categorical colimit. It keeps stage objects and comparison isomorphisms explicit. The identical construction with finitely presented algebras and their algebra maps gives the algebra statement.

More generally, fix a category \(D\) presented by a finite directed graph and finitely many equalities between paths with the same endpoints. A diagram of finitely presented \(R\)-modules or \(R\)-algebras indexed by \(D\) descends to some \(R_i\). Indeed descend the finitely many objects to a common stage; descend all generating arrows at a further common stage; then kill the differences in the finite list of path equations at a final common stage, using the finite source generators of each path. The path equalities then hold exactly, so the arrows define a functor from the original presented \(D\). No enumeration of all composite paths is needed: equality respects composition, so every consequence of the defining relations is retained.

Natural transformations descend as well. Descend their finitely many components, then enforce naturality for each generating arrow; this implies naturality for every path by composition. Two natural transformations that agree at \(R\) agree at a common later stage on the finite object list. An isomorphism of diagrams descends with its inverse and the two finite lists of identity equations. Thus the corresponding functor from the category of stage diagrams and eventual morphisms to diagrams over \(R\) is an equivalence.

For module diagrams the same proof handles finitely many equations that are finite \(R\)-linear combinations of parallel paths. Descend all coefficients in these equations together with the objects and arrows. Their differences are actual maps between the descended modules. Eventual equality on finite source generators makes them zero at one further stage. This retains every coefficient and every path; it is not a replacement by a different presentation. The source results are therefore applicable to specified finite systems of operators, commutation equations and additive identities.

\subsubsection{Split structures descend but unrestricted exactness need not}\label{algebra-proof-036-split-structures-descend-but-unrestricted-exactness-need-not}

Source 31988--32216; group 819. The preceding proof has concrete receiving cases. If \(e:M\to M\) is an idempotent on a finitely presented \(R\)-module, descend \(M,e\) and the equation \(e^2=e\). The stage decomposition is the actual isomorphism \[
M_i\longrightarrow e_iM_i\oplus(1-e_i)M_i,\quad
m\longmapsto(e_i m,(1-e_i)m),
\] with inverse \((x,y)\mapsto x+y\). Each summand is finitely presented: for example \(e_iM_i=M_i/(1-e_i)M_i\), and the denominator submodule is generated by the images of the finite generator list of \(M_i\). Appending lifts of those generators to a finite presentation gives a finite presentation of the quotient. Tensoring the direct-sum decomposition recovers the original one. In particular a finite projective \(R\)-module, presented as the image of an idempotent on \(R^n\), descends to the actual image of a descended idempotent matrix on \(R_i^n\).

A specified split exact sequence descends with all its witnesses. Write its maps as \(a:L\to M\), \(b:M\to N\), \(r:M\to L\), \(s:N\to M\) satisfying \[
ra=1_L,\quad bs=1_N,\quad ba=0,\quad rs=0,\quad ar+sb=1_M.
\] Descend this finite diagram and all five equations. They imply that \((a,s):L\oplus N\to M\) and \((r,b):M\to L\oplus N\) are inverse, so the sequence is split exact at that stage. The exact equations, not tensor exactness for a possibly nonflat transition, prove the assertion.

Likewise, for a bounded cochain complex of finitely presented \(R\)-modules with specified contracting homotopy, descend every original \(d^n:M^n\to M^{n+1}\), \(h^n:M^n\to M^{n-1}\) and the finite lists \[
d^{n+1}d^n=0,\qquad
d^{n-1}h^n+h^{n+1}d^n=1_{M^n}.
\] Only finitely many objects and equations occur. These same equations hold at a sufficiently late stage, proving contractibility there and preserving the original complex after base change.

Unrestricted injectivity or exactness has a different boundary. Let \[
A=k[t,x_1,x_2,\ldots]/(t x_j:j\ge1),\qquad
R_n=A/(x_1,\ldots,x_n),\qquad
R=\operatorname{colim}_{n\ge0}R_n=k[t].
\] Keep the full infinite relation list. For each \(n\), consider the complex \[
0\longrightarrow R_n\xrightarrow{\,t\,}R_n
\longrightarrow R_n/(t)\longrightarrow0.
\] Every term is finitely presented over \(R_n\), and the maps commute with the transition maps. At the limit the complex is the exact sequence \(0\to k[t]\xrightarrow{t}k[t]\to k\to0\). At every finite stage, \(x_{n+1}\ne0\) and \(t x_{n+1}=0\); nonvanishing is witnessed by the homomorphism \(R_n\to k[X]\) taking \(t\) and every other \(x_j\) to zero and \(x_{n+1}\) to \(X\). Thus the first map is never injective at a finite stage. This proves that the source isomorphism-descent assertions and the finite-diagram result must not be strengthened to arbitrary exactness descent. The split and contractible cases work because their exact additional witnesses descend.

\subsubsection{Corpus, propagation and continuation}\label{algebra-proof-036-corpus-propagation-and-continuation}

The existing disk index was queried in both layers for filtered colimit finite presentation. The Ben-Zvi--Nadler--Preygel TeX record and two local Split Support records remain unread routing records. The repeated Bhatt--Scholze record retains only its previously recorded bounded definition/example reading; this query adds no content coverage or adopted result. All seven earlier external source records remain unchanged. There is no PDF fallback or novelty claim.

Groups 796--819 account for the fifty reports and these three proved consequence groups. The retract proof, both finite-stage error lists and the exact diagram comparison provide receiving arguments for the following local approximation section without changing its unread or unadjudicated status. Next is lemma-limit-no-condition-local at 32218. No source or translation mutation, build or publication is claimed.

\hypertarget{algebra_noetherian_approximation_notes_20260928}{}
\subsection{Noetherian approximation: exact systems and their boundaries}\label{algebra-proof-037-noetherian-approximation-exact-systems-and-their-boundaries}

Authority: Stacks Project authors, algebra.tex at a04446e57ec1fbc252a871afcec7752fb2807b14, lines 32218--32642; SHA256 FA8BB92E58A4F78A2BD01B3B6A4A87DE0A0D279F5DD90641B574DD5FBFFFA4F3. Bounded dependencies 5990--6025 and 7502--7515 were reread. All twelve received reports in this interval were read. This finishes the remaining part of the Colimits section; next is More flatness criteria at 32643.

The original author exposition, including arguments left to the reader, remains identifiable in both translations. The following complete arguments, system-choice clarification and stronger consequences are separate editorial material. One canonical correction 0411 is retained. Eight small source operations are proposed in six groups; the source's example and all its original relations are retained. No novelty is claimed.

\subsubsection{Subrings, localizations and the actual finite-type maps}\label{algebra-proof-037-subrings-localizations-and-the-actual-finite-type-maps}

Source 32218--32343 and 32483--32503, 32539--32564; group 820. Let the original local map be \(\varphi:R\to S\), with maximal ideals \(\mathfrak m,\mathfrak n\), so \(\varphi^{-1}(\mathfrak n)=\mathfrak m\). For each finite pair \((A,B)\) with \(\varphi(A)\subset B\), retain \[
R'_\lambda=\mathbf Z[A]\subset R,\quad S'_\lambda=\mathbf Z[B]\subset S,
\quad
\mathfrak m_\lambda=R'_\lambda\cap\mathfrak m,\quad
\mathfrak n_\lambda=S'_\lambda\cap\mathfrak n.
\] The empty pair is allowed and unions give common successors. The map \(R'_\lambda\to S'_\lambda\) is the restriction of the original \(\varphi\). If \(r\notin\mathfrak m_\lambda\), then \(\varphi(r)\notin\mathfrak n_\lambda\), so it extends to a local map \[
R_\lambda=(R'_\lambda)_{\mathfrak m_\lambda}
 \longrightarrow
S_\lambda=(S'_\lambda)_{\mathfrak n_\lambda}.
\] The maps to \(R,S\) are injections on the separate source and target rings: denominators map to units, and a zero fraction has zero numerator in the corresponding original subring. No injectivity of \(\varphi\) is assumed. A larger pair induces the actual inclusion maps on subrings and their fraction maps; contraction of the two fixed maximal ideals makes these transition maps local.

The images of the \(R_\lambda\) cover \(R\), since any \(r\) can be included in \(A\). The images of the \(S_\lambda\) cover \(S\), since any \(s\) can be included in \(B\). Thus the colimit is the original map \(R\to S\), not merely a pair of abstractly isomorphic rings. Each \(R_\lambda\) is essentially of finite type over \(\mathbf Z\). The finite list \(B\) generates \(S'_\lambda\) over the image of \(R'_\lambda\), so \(S_\lambda\) is a localization of a finite-type \(R_\lambda\)-algebra. Removing the prime localizations gives the complete proof of the nonlocal assertion.

For the essentially finite-type case, retain the original \(x_1,\ldots,x_n\in S\) and the image subalgebra \(T=\varphi(R)[x_1,\ldots,x_n]\subset S\). If \(S=U^{-1}T\) is local, then with \(\mathfrak p=T\cap\mathfrak n\) the natural maps \(T_{\mathfrak p}\to S\) and \(S\to T_{\mathfrak p}\) are inverse: \(U\) avoids \(\mathfrak p\), and every element outside \(\mathfrak p\) is already a unit in \(S\). Thus every original fraction of \(S\) can be represented with a numerator and denominator in \(T\), the latter outside \(\mathfrak n\).

For a finite \(A\subset R\), set \(R'_\lambda=\mathbf Z[A]\) and \(S'_\lambda=\varphi(R'_\lambda)[x_1,\ldots,x_n]\), and use the same contracted-prime localizations. Any numerator and denominator in \(T\) involves finitely many original coefficients in \(R\), so appears at some stage, with its original denominator still outside the contracted prime. This proves the colimit assertion for \(S\); it does not presume \(S\) is generated as an algebra by the \(x_i\) without localization.

For \(\lambda\le\mu\), the exact map \[
b\otimes r\longmapsto b\varphi(r):
S'_\lambda\otimes_{R'_\lambda}R'_\mu\longrightarrow S'_\mu
\] is onto because its image contains both \(\varphi(R'_\mu)\) and every original \(x_i\). Let \(Q\) be its actual kernel. The ring \(S_\lambda\otimes_{R_\lambda}R_\mu\) is obtained from the source tensor ring by inverting the images of \(S'_\lambda\setminus\mathfrak n_\lambda\) and \(R'_\mu\setminus\mathfrak m_\mu\). These images avoid the inverse image of \(\mathfrak n_\mu\). Quotient by the extension of \(Q\), then invert the remaining elements outside \(\mathfrak n_\mu\); the result is exactly \(S_\mu\). This proves ``localization of a quotient'' with the original kernel and all denominator sets. Without the prime localizations, the same original map is a surjection and proves the nonlocal finite-type assertion.

Both occurrences of ``Denote'' receive the missing ``by''; none of these full arguments is inserted in place of the author's abbreviated proof.

\subsubsection{Finite presentations require a compatible local base system}\label{algebra-proof-037-finite-presentations-require-a-compatible-local-base-system}

Source 32345--32392; groups 821--824. The statement's plural becomes ``local homomorphisms''. Keep the original chosen isomorphism and full equations \[
\Phi:(R[X_1,\ldots,X_n]/(f_1,\ldots,f_m))_{\mathfrak q}\xrightarrow{\sim}S.
\] Here each \(f_j\) is the original polynomial, with every coefficient and sign retained. The inverse image of \(\mathfrak q\) in \(R\) is \(\mathfrak m\); this is a contraction under the structural map, without any assumption that \(R\) embeds into the quotient.

Choose the base system \(R_\lambda\) from the preceding subring construction, or any directed system as in the source with local transition maps. The explicit locality condition matters. A directed colimit of local rings and local transition maps has local structure maps: an element of the maximal ideal of a stage remains a nonunit at every later stage. If its image in the colimit were a unit, its inverse and the equality of their product with one would be represented at a later stage, making it a unit there, a contradiction. An element outside the maximal ideal is already a unit and stays so. The images of the maximal ideals form an ideal after passage to common stages, and every element outside it is a unit. This is the maximal ideal of the colimit. In particular every structural \(R_\lambda\to R\) is local.

Requiring only that each ring be local is insufficient for the ``any way'' wording in this proof. In the two-stage diagram \(R_0=\mathbf Z_{(2)}\to R_1=\mathbf Q=R\), both rings are local and essentially of finite type over \(\mathbf Z\), and the colimit is \(\mathbf Q\). Take \(S=\mathbf Q\) with its identity local map from \(R\). Using the empty polynomial presentation, the source construction at stage zero localizes \(R_0\) at the inverse image \((0)\) of the prime of \(\mathbf Q\), giving \(S_0=\mathbf Q\). The resulting \(R_0\to S_0\) is not local: the inverse image of the zero maximal ideal is \((0)\), not \((2)\). The existence theorem remains true; a compatible local base system must be selected. The proposed short clarification adds ``with local transition maps'' to the base-system choice, and the original subring construction supplies it.

All coefficients of the finite list \(f_j\) descend to one stage \(\lambda_0\). Choose their actual representatives \(f_{j,\lambda_0}\) and transport them unchanged to every later stage. Put \[
B_\lambda=R_\lambda[X_1,\ldots,X_n]/(f_{1,\lambda},\ldots,f_{m,\lambda}),
\quad
\mathfrak q_\lambda=(B_\lambda\to B)^{-1}(\mathfrak q),\quad
S_\lambda=(B_\lambda)_{\mathfrak q_\lambda},
\] where the middle expression denotes inverse image under the actual map \(B_\lambda\to B=R[X]/(f_1,\ldots,f_m)\), not a tensor product of ideals. Locality of \(R_\lambda\to R\to S\) makes the contraction of \(\mathfrak q_\lambda\) the maximal ideal of \(R_\lambda\), so \(R_\lambda\to S_\lambda\) is local.

The presentation gives the exact base-change isomorphism \(B_\lambda\otimes_{R_\lambda}R_\mu\cong B_\mu\). Let \(T_\lambda=B_\lambda\setminus\mathfrak q_\lambda\), transported into \(B_\mu\). Its elements avoid \(\mathfrak q_\mu\), and \[
S_\lambda\otimes_{R_\lambda}R_\mu
\cong T_\lambda^{-1}B_\mu,\qquad
S_\mu\cong(T_\lambda^{-1}B_\mu)_{\mathfrak q_\mu T_\lambda^{-1}B_\mu}.
\] These are the original maps induced by \(b/t\otimes r\mapsto br/t\) and subsequent localization. They prove the exact prime-localization assertion.

Finite polynomial expressions give \(\operatorname{colim}B_\lambda=B\): elements are represented at a stage, and a relation in the ideal \((f_1,\ldots,f_m)\) is a finite sum of polynomial multiples, whose coefficients also appear at a later stage. For fractions, any \(b/t\in B_{\mathfrak q}\) lifts together with \(t\notin\mathfrak q\). If a represented fraction becomes zero, there is \(h\notin\mathfrak q\) with \(hb=0\) in \(B\); lift \(h\) and the finite equality to a later stage, where \(h\) is still outside the contracted prime and kills the numerator. This proves \(\operatorname{colim}S_\lambda=S\), without claiming that the individual maps \(S_\lambda\to S\) are injective.

The two remaining wording operations add ``by'' to the transported polynomial's declaration and ``to be'' to the inverse-image prime's definition. All displayed coordinate and localization constructions remain source-linked editorial evidence.

\subsubsection{The original wrong system and every later transition}\label{algebra-proof-037-the-original-wrong-system-and-every-later-transition}

Source 32394--32424; group 825 and consequence 828. Keep \(k=\mathbf F_2\) and exactly \[
R=\bigl(k[z,y_1,y_2,\ldots]/(y_i^2-z y_{i+1}:i\ge1)\bigr)_{(z,y_1,y_2,\ldots)},
\quad S=R/zR,
\] \[
R_n=\bigl(k[z,y_1,\ldots,y_n]/(y_i^2-z y_{i+1}:1\le i<n)\bigr)_{(z,y_1,\ldots,y_n)},
\quad S_n=R_n/(z,y_n^2),\quad n\ge1.
\] The original polynomial rings, the full relation lists and the maximal-ideal localizations remain explicit. Their directed colimit is \(R\): any original polynomial fraction uses finitely many variables, and a zero relation uses only finitely many of the listed generators of the ideal. The colimit of the \(S_n\) is \(S\), because after setting \(z=0\) every original \(y_i^2\) is zero.

For \(m>n\), exact right base change gives \[
S_n\otimes_{R_n}R_m=R_m/(z,y_n^2)=R_m/zR_m=:T_m.
\] The second equality is justified by the original relation \(y_n^2=z y_{n+1}\) present in \(R_m\). The comparison with \(S_m\) is the actual quotient map with kernel \[
\ker(T_m\to S_m)=(y_m^2).
\] Retaining the original presentation identifies \[
T_m=
\bigl(k[y_1,\ldots,y_m]/(y_i^2:1\le i<m)\bigr)_{(y_1,\ldots,y_m)}.
\] The element \(y_m^2\) is nonzero: the local homomorphism to \(k[Y]_{(Y)}\) which sends \(y_m\) to \(Y\) and every earlier \(y_i\) to zero sends it to \(Y^2\ne0\). It is a non-zero-divisor as well. Before localization, this ring is the polynomial algebra in \(y_m\) over \(k[y_1,\ldots,y_{m-1}]/(y_i^2)\); multiplication by \(y_m^2\) shifts its coefficient list and is injective. Localization preserves that injection.

Both \(T_m\) and \(S_m\) are local and the quotient map is local. If this map were a localization, every inverted element would have unit image, hence lie outside the maximal ideal of \(T_m\), hence already be a unit. Such a localization is an isomorphism, contrary to the computed nonzero kernel. This proves failure for every \(m>n\), not just adjacent indices. Consequently no cofinal subsequence of these stages makes the transition condition of the essentially finite-presentation lemma true. The obstruction at each later stage is the exact principal ideal generated by the original \(y_m^2\), despite its disappearance in the full colimit.

The source's corrected system \(S_n'=R_n/zR_n\) instead satisfies \[
S_n'\otimes_{R_n}R_m=R_m/zR_m=S_m'
\] for all \(m\ge n\), with the canonical maps. The finite-presentation existence theorem guarantees an appropriate choice; it does not force every finite-type approximation to have this property. The wording ``None of the maps \ldots{} is a localization'' makes the source's intended universal negative explicit. The expanded calculation is editorial and keeps all nilpotents.

\subsubsection{Modules and the nonlocal finite-presentation constructions}\label{algebra-proof-037-modules-and-the-nonlocal-finite-presentation-constructions}

Source 32426--32481 and 32566--32628; consequence 827. For the original finitely presented \(S\)-module use the source presentation \[
S^{\oplus s}\xrightarrow{A}S^{\oplus t}\longrightarrow M\longrightarrow0.
\] Each of the finitely many entries of the actual \(t\times s\) matrix \(A\) is a fraction at a local stage. Lift all its numerators and denominators, with their original positions and values, to one common \(\lambda_2\ge\lambda_1\). Transport this same matrix to every later stage and set \(M_\lambda=\operatorname{coker}A_\lambda\). Right exactness gives the canonical isomorphisms \[
M_\lambda\otimes_{S_\lambda}S_\mu\cong M_\mu,\qquad
M_\lambda\otimes_{S_\lambda}S\cong M.
\] No flatness of a stage transition is required. Tensoring commutes with the directed colimit by the universal property, so the underlying directed colimit of these modules is also \(M\). Each \(M_\lambda\) is finitely presented, stronger than the displayed assertion that it is finite. If \(s=0\) or \(t=0\), the same actual zero-dimensional matrices and their cokernels give the corresponding assertions.

For the nonlocal finite-presentation statement choose finite-type subrings \(R_\lambda\subset R\) containing all coefficients of the original finite list \(f_j\), and put \[
S_\lambda=R_\lambda[X_1,\ldots,X_n]/(f_{1,\lambda},\ldots,f_{m,\lambda}).
\] Every transition gives the exact isomorphism \(S_\lambda\otimes_{R_\lambda}R_\mu\cong S_\mu\), and the same finite-relation argument gives colimit \(S\). Lift the finitely many entries of the module presentation matrix to one of these algebras and proceed as above. Now, for every stage in the restricted cofinal tail, \[
S_\lambda\otimes_{R_\lambda}R\cong S,\qquad
M_\lambda\otimes_{R_\lambda}R
\cong M_\lambda\otimes_{S_\lambda}(S_\lambda\otimes_{R_\lambda}R)
\cong M.
\] The canonical map on elementary tensors is \(m\otimes r\mapsto m\otimes(1\otimes r)\), with inverse \(m\otimes(s\otimes r)\mapsto sm\otimes r\). These exact maps justify the final source equality. In the local version one must retain the further prime localization; this unrestricted base-change isomorphism is not asserted there.

\subsubsection{Integral approximation with the original monic equations}\label{algebra-proof-037-integral-approximation-with-the-original-monic-equations}

Source 32505--32537; group 826. For a pair \((E,F)\), retain finite \(E\subset R,F\subset S\) and an actual monic equation for every \(f\in F\), \[
P_f(T)=T^{d_f}+\sum_{j=0}^{d_f-1}a_{f,j}T^j,\qquad
a_{f,j}\in E,\qquad P_f(f)=0.
\] Unions give upper bounds because the original witness equations remain present. The empty pair is allowed. The rings \(R_\lambda=\mathbf Z[E]\subset R\) and \(S_\lambda=\varphi(R_\lambda)[F]\subset S\) have the original structure maps and inclusion transitions. Every \(r\in R\) occurs by enlarging \(E\), and every \(s\in S\) occurs by taking its actual integral equation over \(R\) and including its coefficients in \(E\). The latter is the reason that canonical correction 0411 changes ``over S'' to ``over R''.

Here is the finite-module calculation without discarding the equations. For each \(f\), the exact relation \[
f^{d_f}=-\sum_{j=0}^{d_f-1}\varphi(a_{f,j})f^j
\] reduces an exponent at least \(d_f\). Repeatedly substitute this whole sum, preserving its signs and coefficients, and induct on the excess exponents. Every product in the \(f\)'s is an \(R_\lambda\)-linear combination of \[
\prod_{f\in F}f^{e_f},\qquad 0\le e_f<d_f.
\] Thus \(S_\lambda\) is generated by at most \(\prod_{f\in F}d_f\) elements as a module. For \(F=\varnothing\), the product is one and its image is the structural unit. If the target is the zero ring, the same generating assertion remains valid; one can choose degree-one witness polynomials. This is a spanning family, not a claim of linear independence, freeness or constant rank. It proves the source finite-map statement and preserves the exact original monic polynomials.

\subsubsection{Simultaneous Noetherian models for finite diagrams}\label{algebra-proof-037-simultaneous-noetherian-models-for-finite-diagrams}

Source 32345--32392, 32426--32481 and 32566--32628; group 827, receiving groups 818--819. All \(R_\lambda\) in the constructions are Noetherian: \(\mathbf Z\) is Noetherian, finite-type extensions are Noetherian and localization preserves this property, by the reread source 5990--6025. Every \(S_\lambda\), being essentially of finite type over \(R_\lambda\), is Noetherian as well. Therefore the local stage maps are essentially of finite presentation: present their finite-type precursor as a quotient of a polynomial algebra over the Noetherian \(R_\lambda\); its actual kernel is finitely generated. Globally the stage maps are of finite presentation. Every finite stage module is finitely presented, because the kernel of a finite free surjection is a submodule of a finite free module over a Noetherian ring. For the chosen \(M_\lambda\) the original matrix already provides its explicit presentation.

More generally take a diagram of finitely presented \(S\)-modules given by finitely many objects and generating arrows, with finitely many path or finite linear path equations. First construct the original ring system \(R_\lambda\to S_\lambda\) just proved. Apply the explicit diagram descent of group 818 to \(S=\operatorname{colim}S_\lambda\). It produces a common stage \(\lambda_0\) with the finitely many actual modules, arrows and equations, including all coefficient representatives. Restrict the base system to \(\lambda\ge\lambda_0\) and define each diagram at \(\lambda\) by tensoring that entire chosen diagram with \(S_\lambda\). Composition, addition and all equations are preserved; right exactness gives its finite presentations. Its transitions have the exact isomorphisms \(M_{\lambda,v}\otimes_{S_\lambda}S_\mu\cong M_{\mu,v}\) at every vertex \(v\), and all arrows commute with these isomorphisms. The colimit is the original entire diagram and its equations.

Any specified finite set of natural transformations, inverse maps, idempotents, splitting witnesses or bounded contracting homotopies can be included among these finitely many data, as proved in group 819. The ring system continues to satisfy the source's exact base-change/localization conditions. This strengthens the single-module statement to a simultaneous finite diagram without assuming flat transitions or arbitrary exactness descent. The counterexample in group 819 remains a boundary.

\subsubsection{Persistent defects and uniformly finite approximations}\label{algebra-proof-037-persistent-defects-and-uniformly-finite-approximations}

Group 828 is the all-\(m>n\) result proved in the original wrong-system section: the comparison kernel is exactly the nonzero regular principal ideal \((y_m^2)\), and no cofinal subsequence repairs it. Its receiving comparison with group 819 distinguishes disappearance of each individual defect in a colimit from a finite stage satisfying every desired transition property.

For group 829, the integral spanning bound above upgrades each finite stage to a finitely presented \(R_\lambda\)-module and finitely presented \(R_\lambda\)-algebra, because \(R_\lambda\) is Noetherian. There is an additional uniform construction when the original map \(R\to S\) is finite. Fix an actual finite \(R\)-module generating family \(f_1,\ldots,f_r\) containing \(1_S\); for the zero ring include \(f_1=0=1_S\). Choose the actual full multiplication table \[
f_if_j=\sum_{k=1}^r\varphi(a_{ijk})f_k,\qquad a_{ijk}\in R.
\] Take the directed finite-type subrings \(R_\lambda\subset R\) containing every original \(a_{ijk}\), and put \[
S_\lambda=\sum_{k=1}^r\varphi(R_\lambda)f_k\subset S.
\] This is an actual subring: it contains the unit, is closed under addition, and the displayed table proves closure under multiplication term by term. The structural map has image in it because it contains \(1_S\). Each \(S_\lambda\) is generated as an \(R_\lambda\)-module by these same \(r\) elements. For \(\lambda\le\mu\), the natural map \[
S_\lambda\otimes_{R_\lambda}R_\mu\to S_\mu,\qquad s\otimes a\mapsto s\varphi(a)
\] is onto: every element of \(S_\mu\) is a linear combination of the same \(f_k\). The subrings cover \(S\), since the finitely many coefficients of any such original linear combination eventually belong to \(R_\lambda\). Their colimit and that of the bases are the original map. All stages are finite and finitely presented as modules and algebras over their Noetherian bases, and the generator bound \(r\) is uniform. This makes no freeness, flatness or injective tensor-map claim. The exact multiplication coefficients are retained, not replaced by a different algebra model.

\subsubsection{Corpus and continuation}\label{algebra-proof-037-corpus-and-continuation}

Both canonical corpus layers were queried for Noetherian approximation local rings. The Aitken, Murayama, Connes--Consani and Arithmetic Time records remain unread routing records and supply no theorem here. The seven earlier external source records retain exactly their prior bounded coverage; no PDF fallback or novelty claim.

Groups 820--829 dispose the twelve reports and record three separate consequence groups. Next is More flatness criteria, authority 32643. No later source or report is counted as read by this batch. No source or translation mutation, TeX build or publication is claimed.

\hypertarget{algebra_more_flatness_notes_20260928}{}
\subsection{More flatness criteria: source review and separate editorial proofs}\label{algebra-proof-038-more-flatness-criteria-source-review-and-separate-editorial-proofs}

Source: Stacks Project authors, algebra.tex, authority commit a04446e57ec1fbc252a871afcec7752fb2807b14, complete interval 32643--33226. The original TeX file has SHA-256 FA8BB92E58A4F78A2BD01B3B6A4A87DE0A0D279F5DD90641B574DD5FBFFFA4F3. This note supplies editorial arguments and consequences; it does not replace the original exposition in a translation. Groups 830--853 bind the individual findings to their source passages. No novelty is claimed.

All 21 received reports OCC-00638--OCC-00658 were read. Dependencies reread: algebra.tex:23423--23477 (Noetherian injectivity criterion), 23750--23810 (local flatness criterion and its remark), 23845--24023 (the two Tor comparisons, change-of-base criterion and Noetherian fibre criterion), and 9195--9262 (flatness and localization). The preceding approximation arguments are recorded separately in ALGEBRA\_NOETHERIAN\_APPROXIMATION\_NOTES\_20260928.md, especially groups 827--829. These are exact local proof locators, not claims of new public publication.

\subsubsection{Source dispositions and translation treatment}\label{algebra-proof-038-source-dispositions-and-translation-treatment}

Restore ``is'' at 32662 and the comma ending the display at 32674. Supply ``by'' in the original declarations at 32772, 32794, 32871, 32888, 32925 and 33011; the occurrence at 32794 is an additional matching defect outside the report's three listed declarations. Change the index bound at 32813 to weak inequality: a directed preorder can have a largest element or a single element, so strict enlargement is not available.

At 32841--32846 choose an actual fraction representative in the earlier module, retaining its tensor base. At 32927 use the quotient base ring and at 32933 state that base for flatness. At 32975 acknowledge the localization after base change. At 33040 retain both original relation families; at 33049 place the matrix over the primed ring. Both endpoint tensor equalities at 33108--33110 need localization. At 33136 use the finite-type hypothesis, not the finite-presentation conclusion being proved. At 33138 use ``Let''; at 33158 use the two distinct ideals.

The final proof must choose a prime in the support of the module, and must then consider all prime localizations, including those where the module is zero. The complete argument appears below. The proposed short source repair makes these two points visible without substituting the generalized theorem.

The reordered modifier in OCC-00641, abbreviated proofs in OCC-00644 and OCC-00645, ``whence'' in OCC-00649, and double relative clause in OCC-00658 are optional prose changes. In OCC-00656 the plural ``ideals'' correctly quantifies over the family of finitely generated ideals having the specified property. None is a mathematical defect. Their source text is retained.

All expanded arguments in this note, including proofs of source omissions, belong in source-linked editorial material. All consequences in groups 851--853 have zero translation replacement operations. Minimal proposed source changes are review records and are not yet source or translation mutations.

\subsubsection{Regular bases and eventual flatness}\label{algebra-proof-038-regular-bases-and-eventual-flatness}

For miracle flatness, retain the original local Noetherian map \(R\to S\), with \(R\) regular, \(S\) Cohen--Macaulay, and \[
\dim S=\dim R+\dim(S/\mathfrak m_RS).
\] If \(\dim R=0\), \(R\) is a field. Otherwise every minimal prime \(\mathfrak q_i\) of the Cohen--Macaulay local ring \(S\) satisfies \(\dim(S/\mathfrak q_i)=\dim S\), by the source's maximal-chain lemma. If \(\mathfrak m_RS\subset\mathfrak q_i\), passage to that quotient would give \(\dim(S/\mathfrak q_i)\leq\dim(S/\mathfrak m_RS)\), a contradiction. Thus the contractions \(\mathfrak p_i\) are strictly smaller than \(\mathfrak m_R\).

Prime avoidance, in precisely the form of the cited lemma-silly, gives \(x\in\mathfrak m_R\setminus\mathfrak m_R^2\) outside every \(\mathfrak p_i\). It avoids all associated primes of \(S\), since a Cohen--Macaulay local ring has no embedded associated primes. Hence multiplication by \(x\) on \(S\) is injective. The source regular-ring and Cohen--Macaulay results give \[
\dim(R/xR)=\dim R-1,\qquad
\dim(S/xS)=\dim S-1,
\] with \(R/xR\) regular and \(S/xS\) Cohen--Macaulay. The closed fibre is unchanged. Induction gives flatness of \(S/xS\) over \(R/xR\). The exact resolution \(0\to R\xrightarrow{x}R\to R/xR\to0\), together with injectivity of multiplication on \(S\), gives \(\operatorname{Tor}_1^R(S,R/xR)=0\). The source local flatness criterion now proves \(S\) flat over \(R\).

For lemma-flat-over-regular, a regular sequence of parameter images generates a proper ideal in the local ring \(S\). Thus all parameter images lie in its maximal ideal. Since those parameters generate \(\mathfrak m_R\), the ring map is local. Starting with the field quotient by all parameters, apply the same exact multiplication resolution and local criterion successively to the quotients by \(x_1,\ldots,x_{i-1}\). This proves the stated flatness without adding a hidden locality assumption.

For eventual flatness retain the systems and all six comparison properties of lemma-limit-module-essentially-finite-presentation. Fix \(\lambda\). Write \(\mathfrak m_\lambda\) for the maximal ideal of \(R_\lambda\). The module \[
K_\lambda=\operatorname{Tor}_1^{R_\lambda}
 (M_\lambda,R_\lambda/\mathfrak m_\lambda)
 =\ker(\mathfrak m_\lambda\otimes_{R_\lambda}M_\lambda\longrightarrow M_\lambda)
\] is finite over the Noetherian ring \(S_\lambda\): both the ideal and the module have finite generating sets. Choose generators \(\xi_1,\ldots,\xi_n\). For \(\mu\geq\lambda\) let \(J_\mu=\mathfrak m_\lambda R_\mu\). The filtered colimit of \(J_\mu\otimes_{R_\mu}M_\mu\) is \(\mathfrak m_\lambda R\otimes_R M\). Indeed every tensor is a finite sum with finitely many coefficients, and both its representatives and the defining additive and balanced relations appear at a common stage. The images of all \(\xi_i\) in this colimit vanish because \(M\) is \(R\)-flat. Each image therefore vanishes at a sufficiently large common stage \(\mu\).

The original quotient \(M_\lambda/\mathfrak m_\lambda M_\lambda\) is flat over the field \(R_\lambda/\mathfrak m_\lambda\). Apply the source change-of-base local criterion at 23920--23968 to the original local maps \(R_\lambda\to R_\mu\), \(S_\lambda\to S_\mu\), ideal \(\mathfrak m_\lambda\), and module \(M_\lambda\). Its base-changed module is localized to \(M_\mu\), as required by that criterion. The induced first Tor map is zero by the preceding generator calculation. It follows that \(M_\mu\) is \(R_\mu\)-flat. Every later module is a localization of its base change, so all later modules are flat over their respective bases.

For later uses, filtered colimits of these varying-base flat modules are flat over the colimit ring. To verify this directly, take a finite relation \(\sum_i a_i m_i=0\) in the limit module. Lift its finitely many coefficients and elements to one stage and enlarge the stage until the relation is zero there. Stage flatness gives finitely many elements \(n_j\) and coefficients \(b_{ij}\) with \(m_i=\sum_jb_{ij}n_j\) and \(\sum_i a_i b_{ij}=0\) for every \(j\). Their images give the flatness relation criterion in the limit. No injectivity of transition maps is required.

\subsubsection{Fibre injectivity with actual fraction representatives}\label{algebra-proof-038-fibre-injectivity-with-actual-fraction-representatives}

For the source mod-injective-general lemma retain \(u:M\to N\), with both modules finitely presented over the same essentially finitely presented local \(R\)-algebra \(S\), \(N\) flat over \(R\), and \(\overline u\) injective on the closed fibre. Group 827 supplies compatible models of this entire map. Eventual flatness permits restricting to a cofinal tail on which all \(N_\lambda\) are \(R_\lambda\)-flat.

Put \(\kappa_\lambda=R_\lambda/\mathfrak m_\lambda\), \(\kappa=R/\mathfrak m\), \(A_\lambda=S_\lambda/\mathfrak m_\lambda S_\lambda\) and \(V_\lambda=M_\lambda/\mathfrak m_\lambda M_\lambda\). The ring \(A_\lambda\otimes_{\kappa_\lambda}\kappa\) is a localization of a finite-type algebra over the field \(\kappa\), so is Noetherian. Its finite module \(V_\lambda\otimes_{\kappa_\lambda}\kappa\) has a finite kernel \[
K_\lambda=\ker\bigl(\Psi_\lambda:
 V_\lambda\otimes_{\kappa_\lambda}\kappa\longrightarrow M/\mathfrak mM\bigr).
\] Lift a finite generating list of this kernel to \(V_\lambda\otimes_{\kappa_\lambda}\kappa_\mu\). Each has zero image in the colimit of the \(V_\nu\); enlarging \(\mu\) kills all these images in \(V_\mu\). Consequently \(K_\lambda\) is generated, after extension to \(\kappa\), by the image of \[
\ker\bigl(\Psi_{\lambda,\mu}:
 V_\lambda\otimes_{\kappa_\lambda}\kappa_\mu\longrightarrow V_\mu\bigr).
\] The reverse containment holds because this comparison factors the map to the final module.

The original ring construction gives an actual multiplicative set \(W\subset S_\lambda\otimes_{R_\lambda}R_\mu\) and an isomorphism \[
V_\mu\otimes_{\kappa_\mu}\kappa
 \cong W^{-1}(V_\lambda\otimes_{\kappa_\lambda}\kappa).
\] Here \(W\) in the right side denotes its image after the displayed scalar extensions; every such element maps to a unit in \(S\) and in the closed-fibre target. If \(x\in\ker\Psi_\mu\), choose an actual representative \(x=y/w\), with \(y\in V_\lambda\otimes_{\kappa_\lambda}\kappa\) and \(w\in W\). Since \(w\) acts invertibly on the target, \(\Psi_\lambda(y)=0\). The generators just constructed for \(K_\lambda\) all vanish in \(V_\mu\), hence their scalar-extended images vanish in \(V_\mu\otimes_{\kappa_\mu}\kappa\). Therefore \(y/1=0\) and \(x=0\). This does not identify an earlier module with its localization.

For any later \(\nu\), write \(V_\nu\otimes_{\kappa_\nu}\kappa\) as a localization of \(V_\mu\otimes_{\kappa_\mu}\kappa\). The same fraction argument, using injectivity already established at \(\mu\), proves injectivity at \(\nu\). Because the field extension \(\kappa_\nu\to\kappa\) is faithfully flat, \(V_\nu\to V_\nu\otimes_{\kappa_\nu}\kappa\) is injective. Hence \(V_\nu\to M/\mathfrak mM\) is injective.

If an element of \(V_\nu\) maps to zero under \(\overline{u_\nu}\), its image in the final closed fibre maps to zero under \(\overline u\), and thus is zero; the preceding injection makes the element itself zero. The Noetherian injectivity criterion at 23423--23477 gives injectivity of \(u_\nu\) and flatness of \(Q_\nu=N_\nu/u_\nu(M_\nu)\) over \(R_\nu\). Exactness of filtered colimits gives injectivity of \(u\) and identifies its cokernel with \(\operatorname{colim}Q_\nu\). The finite-relation calculation in the preceding section gives its \(R\)-flatness.

\subsubsection{The local Tor criterion with its quotient bases}\label{algebra-proof-038-the-local-tor-criterion-with-its-quotient-bases}

Retain all four hypotheses of lemma-variant-local-criterion-flatness-general, including the proper ideal \(I\subset R\), finite presentation of \(S/R\) in the stated essential sense, finite presentation of \(M/S\), vanishing \(\operatorname{Tor}_1^R(M,R/I)\), and flatness of \(M/IM\) over \(R/I\). Let \(I_\lambda\) be the inverse image of \(I\) in \(R_\lambda\). The correct quotient system is \[
R_\lambda/I_\lambda\longrightarrow S_\lambda/I_\lambda S_\lambda,
\qquad M_\lambda/I_\lambda M_\lambda.
\] Its colimit is the original quotient data. Its local transitions and module tensor isomorphisms follow by taking these exact quotients of the original system. Eventual flatness gives \(M_\lambda/I_\lambda M_\lambda\) flat over \(R_\lambda/I_\lambda\) on a cofinal tail.

Fix \(\lambda\) there. The module \[
K_\lambda=\operatorname{Tor}_1^{R_\lambda}(M_\lambda,R_\lambda/I_\lambda)
 =\ker(I_\lambda\otimes_{R_\lambda}M_\lambda\to M_\lambda)
\] is finite over \(S_\lambda\); choose its full generating list \(\xi_1,\ldots,\xi_n\). The ideal systems have colimit \(I\). Exactness of filtered colimits and the finite tensor-relation argument identify the colimit of these kernels with \(\operatorname{Tor}_1^R(M,R/I)=0\). Enlarge to one stage \(\mu\geq\lambda\) at which every \(\xi_i\) has zero image.

The source's three comparison maps, with every original base retained, are \[
\begin{aligned}
R_\mu\otimes_{R_\lambda}\operatorname{Tor}_1^{R_\lambda}
 (M_\lambda,R_\lambda/I_\lambda)
&\longrightarrow \operatorname{Tor}_1^{R_\lambda}
 (M_\lambda,R_\mu/I_\lambda R_\mu)\\
&\longrightarrow \operatorname{Tor}_1^{R_\mu}
 (M_\mu,R_\mu/I_\lambda R_\mu)\\
&\longrightarrow \operatorname{Tor}_1^{R_\mu}
 (M_\mu,R_\mu/I_\mu).
\end{aligned}
\] The first is onto by the source Tor comparison at 23845--23878 and quotient flatness at \(\lambda\). For the second, use the base-change comparison proved at 23880--23918, then the actual localization from \(S_\lambda\otimes_{R_\lambda}R_\mu\) to \(S_\mu\). Localization commutes with the relevant Tor: take a free \(R_\mu\)-resolution of its second argument, tensor with the module, and use exactness of localization on that complex and on its homology.

For the third arrow the needed quotient module is \(M_\mu/I_\lambda M_\mu\). It is a localization of \[
(M_\lambda/I_\lambda M_\lambda)
 \otimes_{R_\lambda/I_\lambda}(R_\mu/I_\lambda R_\mu),
\] and hence is flat over \(R_\mu/I_\lambda R_\mu\). Apply 23845--23878 again. Thus the composite is onto after tensoring its source with the actual \(S_\mu\) localization. The images of the full original generating list \(\xi_i\) are zero there, so the last Tor module is zero. The Noetherian local flatness criterion proves \(M_\mu\) flat over \(R_\mu\); passage to the final localized base change proves \(M\) flat over \(R\).

\subsubsection{The fibre criterion with all relations and localizations}\label{algebra-proof-038-the-fibre-criterion-with-all-relations-and-localizations}

Retain the local maps \(R\to S\to S'\), the nonzero finitely presented \(S'\)-module \(M\), both essential finite-presentation hypotheses over \(R\), \(R\)-flatness of \(M\), and \(S/\mathfrak mS\)-flatness of \(M/\mathfrak mM\). The relative map \(S\to S'\) is essentially finitely presented as well. Here is the finite presentation argument needed for the source construction.

Write \(S=(R[x_1,\ldots,x_n]/(f_1,\ldots,f_u))_{\mathfrak q}\) and write \(S'\) as a prime localization of \(R[z_1,\ldots,z_l]/(h_1,\ldots,h_v)\). In the second presentation express the image of each \(x_i\) as an actual fraction \(a_i(z)/b_i(z)\), with \(b_i\) outside its specified prime. Over \(S\), impose the finite equations \(h_j(z)=0\) and \(b_i(z)x_i-a_i(z)=0\), then localize at the prime induced by \(S'\). Each \(b_i\) is invertible there. These equations give precisely the same local ring as \(S'\): the mutually inverse maps send every \(z_j\) to its original element and every \(x_i\) to its stated fraction, and elements outside the contracted prime become units. This keeps the full denominators and source relations; it does not erase the \(f_i\).

Choose the compatible Noetherian systems of the preceding source construction, retaining all coefficients and contracted primes. The stage presentation over \(R_\lambda\) must be \[
S'_\lambda=
\bigl(R_\lambda[x_1,\ldots,x_n,y_1,\ldots,y_m]/
 (f_{1,\lambda},\ldots,f_{u,\lambda},
  g_{1,\lambda},\ldots,g_{v,\lambda})\bigr)_{\mathfrak q'_\lambda}.
\] The \(f\)-relations already defining \(S_\lambda\) cannot be omitted on changing its presentation to one over \(R_\lambda\). All finite coefficient equalities hold after enlarging the stage. Descend the entire original matrix \(A\in\operatorname{Mat}(t\times s,S')\) to a matrix over \(S'_{\lambda_3}\), retaining its entries and both dimensions, and set \[
M_\lambda=\operatorname{coker}\bigl((S'_\lambda)^{\oplus s}
 \xrightarrow{A_\lambda}(S'_\lambda)^{\oplus t}\bigr).
\] The module comparison \(M_\lambda\otimes_{S'_\lambda}S'\cong M\) is an isomorphism. In particular every stage module on this tail is nonzero, since \(M\ne0\).

The original transition maps for \(S_\lambda\) and \(S'_\lambda\) are localizations of their stated base changes; the module transition maps over \(S'_\lambda\) are isomorphisms. Eventual flatness first makes \(M_\lambda\) flat over \(R_\lambda\). Apply it again to the exact closed-fibre systems \[
S_\lambda/\mathfrak p_\lambda S_\lambda
\longrightarrow S'_\lambda/\mathfrak p_\lambda S'_\lambda,\qquad
M_\lambda/\mathfrak p_\lambda M_\lambda.
\] These systems satisfy the same localization and module conditions by quotienting the original maps. Their limit module is flat over \(S/\mathfrak mS\) by hypothesis. Thus on a common cofinal tail the stage closed-fibre module is flat over \(S_\lambda/\mathfrak p_\lambda S_\lambda\). All rings there are Noetherian and all stage modules finite and nonzero. The Noetherian fibre criterion at 23977--24023 gives \(S_\lambda\) flat over \(R_\lambda\) and \(M_\lambda\) flat over \(S_\lambda\).

At the endpoint, \(S_\lambda\otimes_{R_\lambda}R\to S\) is a localization. Therefore \(S\) is \(R\)-flat. Moreover \[
M_\lambda\otimes_{S_\lambda}S
 \cong M_\lambda\otimes_{S'_\lambda}(S'_\lambda\otimes_{S_\lambda}S)
\] is \(S\)-flat, and \(S'_\lambda\otimes_{S_\lambda}S\to S'\) is a localization. Localizing this module gives \(M\), proving its \(S\)-flatness. These are localization maps, not generally the two tensor equalities printed in the source.

For an explicit boundary take \(R_0=\mathbf Q\), \(R_1=\mathbf Q(t)\), \(S_0=\mathbf Q[x]_{(x)}\), and \(S_1=\mathbf Q(t)[x]_{(x)}\). The element \(1+tx\) is a unit in \(S_1\). It is not a unit in \(S_0\otimes_{\mathbf Q}\mathbf Q(t)\): that ring only inverts the nonzero rational polynomials with nonzero constant term, and \(1+tx\) divides none of them, since a nonzero rational polynomial cannot vanish at the transcendental element \(-1/t\). This disproves the first unlocalized equality. To test the module equality independently, use \(S_i=R_i\), take \(S'_i\) to be those same localized polynomial rings and \(M_i=S'_i\). The asserted module tensor equality would give the same false ring equality. Both corrections are necessary.

\subsubsection{Finite presentation from the finite-type hypothesis}\label{algebra-proof-038-finite-presentation-from-the-finite-type-hypothesis}

In lemma-criterion-flatness-fibre-fp-over-ft start only with the stated essential finite type of \(S/R\). Write \(S=B/J\), with \(B=R[x_1,\ldots,x_n]_{\mathfrak q}\), retaining the original inverse-image prime and ideal. The closed fibre \(B/\mathfrak mB\) is Noetherian. Choose \(f_1,\ldots,f_k\in J\) whose images generate the image of \(J\) in that fibre, and put \(J_0=(f_1,\ldots,f_k)\).

For each finitely generated \(J'\) with \(J_0\subset J'\subset J\), the local ring \(B/J'\) is essentially finitely presented over \(R\), and its closed fibre agrees with that of \(S\). The original nonzero module \(M\), regarded through the maps \(R\to B/J'\to S'\), has exactly the other fibre-criterion hypotheses. Hence \(B/J'\) is \(R\)-flat and \(M\) is \(B/J'\)-flat.

If \(J'\subset J''\subset J\) are two such ideals, apply mod-injective-general to the common algebra \(B\) and the map of \(B\)-modules \(B/J'\to B/J''\). Both modules are finitely presented over \(B\), the target is \(R\)-flat, and the map on closed fibres is an isomorphism. This verifies every hypothesis of that lemma; essential finite presentation of two unrelated algebras alone would not state its module hypotheses. The map is injective, and being the quotient map it is also surjective. Its kernel \(J''/J'\) is zero.

For any \(f\in J\), take \(J'=J_0\), \(J''=J_0+(f)\). Both satisfy the required conditions, so \(f\in J_0\). Thus \(J=J_0\) is finitely generated. It follows that \(S\) is essentially finitely presented over \(R\), and the preceding fibre criterion applied to the original \(R\to S\to S'\) gives the stated flatness conclusions. No finite-presentation conclusion was assumed at the start.

\subsubsection{The support prime required by the nil-ideal proof}\label{algebra-proof-038-the-support-prime-required-by-the-nil-ideal-proof}

Retain the original arbitrary ring diagram \(R\to S\to S'\), finite type of \(S/R\), finite presentation of \(S'/R\) and \(M/S'\), \(R\)-flatness of \(M\), and \(S/IS\)-flatness of \(M/IM\). First consider a prime \(\mathfrak q'\in\operatorname{Supp}_{S'}M\) whose contraction \(\mathfrak p\) to \(R\) contains \(I\), and write \(\mathfrak q\) for its contraction to \(S\).

The maps \(R_{\mathfrak p}\to S_{\mathfrak q}\to S'_{\mathfrak q'}\) are local. The first is essentially of finite type and the composite essentially of finite presentation. The module \(M_{\mathfrak q'}\) is finitely presented and nonzero over the local ring \(S'_{\mathfrak q'}\), and is \(R_{\mathfrak p}\)-flat. Since \(I\subset\mathfrak p\), \[
M/\mathfrak pM\cong (M/IM)\otimes_{S/IS}(S/\mathfrak pS)
\] is \(S/\mathfrak pS\)-flat. Further localization at the elements defining \(R_{\mathfrak p},S_{\mathfrak q},S'_{\mathfrak q'}\) preserves this flatness. The closed-fibre hypothesis of fp-over-ft is therefore satisfied. That lemma gives \(M_{\mathfrak q'}\) flat over \(S_{\mathfrak q}\), and \(S_{\mathfrak q}\) flat and essentially finitely presented over \(R_{\mathfrak p}\), hence over \(R\). As an \(S\)-module \(M_{\mathfrak q'}\) is flat as well.

For the source nil-ideal hypothesis, every prime of \(R\) contains \(I\). Thus the preceding argument applies to every nonzero stalk \(M_{\mathfrak q'}\). A zero stalk is flat automatically. The flat-localization criterion of 9195--9262 now proves \(M\) flat over \(S\). Equivalently, tensor an arbitrary injection of \(S\)-modules with \(M\); its kernel is an \(S'\)-module whose localization at each prime is zero by the stalk flatness just proved. A module with all prime localizations zero is zero.

If \(M\otimes_S\kappa(\mathfrak q)\ne0\), this is a nonzero finite module over \(S'\otimes_S\kappa(\mathfrak q)\). Choose a prime in its support, not merely any prime of that fibre ring. Its inverse image \(\mathfrak q'\) in \(S'\) lies over \(\mathfrak q\) and has \(M_{\mathfrak q'}\ne0\). The preceding local argument then proves the asserted properties of \(S_{\mathfrak q}\).

An arbitrary prime over \(\mathfrak q\) need not have this nonvanishing property: for example take \(S=k\), \(S'=k\times k\), \(M=k\times0\). The module fibre is nonzero, but its localization at the prime supported on the second factor is zero. This is why the source's choice must be made in the support. The source's final phrase ``arbitrary prime'' also needs the separate zero-stalk case, since its earlier choice only supplied one support prime.

\subsubsection{Universal exactness and regular sequences}\label{algebra-proof-038-universal-exactness-and-regular-sequences}

Group 851. Under precisely the source mod-injective-general hypotheses, write \(Q=N/u(M)\). The preceding proof gives a short exact sequence \[
0\longrightarrow M\xrightarrow{u}N\longrightarrow Q\longrightarrow0
\] with \(N,Q\) flat over \(R\). For any \(R\)-module \(T\), the tensor long exact sequence has \(\operatorname{Tor}_1^R(Q,T)=0\); thus this short exact sequence remains exact after tensoring with \(T\).

The module \(M\) is flat over \(R\). Indeed for any injection \(T_1\to T_2\), the two just-proved short exact tensor sequences form a commutative diagram. An element of \(M\otimes_R T_1\) killed in \(M\otimes_R T_2\) maps to an element of \(N\otimes_R T_1\) killed in \(N\otimes_R T_2\); \(N\)-flatness and the left injection in the first row force it to be zero. Also \(Q\) is finitely presented over \(S\): retain a finite presentation of \(N\), lift images of a finite generating list of \(M\) to its finite free module, and append those finitely many lifted vectors to its existing relation list. Its cokernel is exactly \(Q\).

Under the source regular-sequence-general hypotheses let \(S_i=S/(f_1,\ldots,f_i)\), including \(S_0=S\). Induction by grothendieck-general makes every \(S_i\) flat over \(R\), and every multiplication sequence \[
0\longrightarrow S_{i-1}\xrightarrow{f_i}S_{i-1}
\longrightarrow S_i\longrightarrow0
\] exact. These sequences remain exact after tensoring with every \(R\)-module because their right terms are \(R\)-flat. The original closed-fibre regular sequence generates a proper ideal, so each \(S_i\) is a nonzero local ring and \(R\to S_i\) is local and flat. Hence it is faithfully flat: any proper ideal of the local ring \(R\) lies in \(\mathfrak m\) and its extension remains in the maximal ideal of \(S_i\); the flatness criterion for faithful flatness applies.

For every nonzero \(R\)-algebra \(R'\), all \(S_i\otimes_R R'\) are nonzero by that faithfulness. Their multiplication sequences stay exact, and their final quotient is nonzero. Thus the images of the same \(f_1,\ldots,f_c\) form a regular sequence in \(S\otimes_R R'\). Every quotient is faithfully flat over \(R'\), by base change. If \(R'=0\), the exactness statements still hold; no proper regular-sequence assertion is made in the zero ring.

\subsubsection{Faithfulness and a pure receiving map}\label{algebra-proof-038-faithfulness-and-a-pure-receiving-map}

Group 852. Under either local fibre criterion in this section, \(M\) is faithfully flat over \(S\) and \(S\) is faithfully flat over \(R\). The second assertion follows because \(R\to S\) is local and flat. For the first, \(M\) is flat over \(S\), finite nonzero over local \(S'\), and the extended ideal \(\mathfrak m_SS'\) lies in the maximal ideal of \(S'\). Nakayama's lemma gives \(M/\mathfrak m_SM\ne0\). For a proper ideal \(H\subset S\), the quotient \(M/HM\) maps onto this nonzero module. Flatness and nonvanishing of these quotients give faithful flatness over \(S\).

The original map \(S\to S'\) is pure, even though no flatness of \(S'\) over \(S\) is asserted. For any \(S\)-module \(T\), tensor its unit map \[
\eta_T:T\longrightarrow S'\otimes_S T,\qquad t\longmapsto1\otimes t
\] with the faithfully flat \(S\)-module \(M\). The resulting map has the explicit retraction \[
S'\otimes_S T\otimes_S M\longrightarrow T\otimes_S M,\qquad
s'\otimes t\otimes m\longmapsto t\otimes(s'm).
\] This is well-defined by the \(S'\)-module structure of \(M\), and its composite with \(\eta_T\otimes_S M\) is the identity. Thus that tensor map is injective. Flatness identifies its kernel with \((\ker\eta_T)\otimes_S M\); faithfulness forces \(\ker\eta_T=0\). This proves purity with the actual map and action.

In particular \(S/H\to S'/HS'\) is injective for every ideal \(H\subset S\), so \(HS'\cap S=H\), including injectivity of \(S\to S'\) at \(H=0\). Also \(\kappa(\mathfrak q)\to S'\otimes_S\kappa(\mathfrak q)\) is injective for every prime \(\mathfrak q\subset S\). Each fibre ring is nonzero and has a prime, so \(\operatorname{Spec}(S')\to\operatorname{Spec}(S)\) is surjective. These are consequences of the original module criterion; replacing them by a claim that \(S'/S\) is flat would exceed the proof.

\subsubsection{Nilness on module support suffices}\label{algebra-proof-038-nilness-on-module-support-suffices}

Group 853. In the global diagram of the locally nilpotent fibre criterion, replace the requirement that \(I\subset R\) be a nil ideal by \[
IS'\subset\sqrt{\operatorname{Ann}_{S'}M}.
\] Keep every other original hypothesis: \(S/R\) finite type, \(S'/R\) finite presentation, \(M/S'\) finite presentation, \(M\) flat over \(R\), and \(M/IM\) flat over \(S/IS\). Then exactly the source conclusions hold: \(M\) is \(S\)-flat, and \(S_{\mathfrak q}\) is flat and essentially finitely presented over \(R\) whenever \(M\otimes_S\kappa(\mathfrak q)\ne0\).

To prove this, finiteness of \(M/S'\) gives \(\operatorname{Supp}_{S'}M=V(\operatorname{Ann}_{S'}M)\). Explicitly, if \(M_{\mathfrak q'}=0\), a product of finitely many elements outside \(\mathfrak q'\), one killing each chosen generator, annihilates \(M\); conversely an annihilator outside \(\mathfrak q'\) kills the localization. The displayed hypothesis is therefore equivalent to every prime in this support containing \(IS'\). Its contraction \(\mathfrak p\subset R\) contains \(I\). The complete local proof in the preceding support-prime section now applies to every nonzero stalk, while zero stalks are flat. It proves global \(S\)-flatness and the asserted properties at each receiving prime.

For completeness the exact receiving locus is \[
U=\{\mathfrak q\in\operatorname{Spec}S:
M\otimes_S\kappa(\mathfrak q)\ne0\}
=\operatorname{image}(\operatorname{Supp}_{S'}M\longrightarrow\operatorname{Spec}S).
\] The forward inclusion follows by choosing a support prime of the nonzero finite fibre module. Conversely, let \(\mathfrak q'\in\operatorname{Supp}_{S'}M\) contract to \(\mathfrak q\). The finite nonzero local module \(M_{\mathfrak q'}\) has nonzero quotient by \(\mathfrak qM_{\mathfrak q'}\), by Nakayama, since \(\mathfrak qS'_{\mathfrak q'}\) lies in its maximal ideal. This quotient is the corresponding localization of \(M\otimes_S\kappa(\mathfrak q)\), so that fibre is nonzero. This proves equality of the loci; it does not assert openness of \(U\).

The new condition is strictly weaker than nilness of \(I\) in \(R\). Let \(R=k\times k\), \(I=k\times0\), \(S=S'=k\) via the second projection, and \(M=k\). The ideal \(I\) contains a nonzero idempotent and is not nil, but annihilates \(M\). Both algebra presentations are finite, \(M\) is a projective \(R\)-module and a free \(S\)-module, and all the hypotheses and conclusions hold.

Some support condition is needed for this extension. Let \(R=k[t]\), \(I=(t)\), \(S=R[x]\), \(S'=R[1/t]\) with \(x\mapsto0\), and \(M=S'\). The algebra \(S'/R\) has the finite presentation \(R[u]/(tu-1)\); \(S/R\) is polynomial, and \(M/S'\) is free of rank one. The module \(M\) is \(R\)-flat, and \(M/IM=0\) is \(S/IS\)-flat. But \(M\) is not \(S\)-flat: the injective multiplication map \(S\xrightarrow{x}S\) becomes the zero map on the nonzero module \(M\). Here \(IS'=S'\), so the support condition fails. This retains both the original base and every localization denominator.

\subsubsection{Propagation and limits of the review}\label{algebra-proof-038-propagation-and-limits-of-the-review}

Groups 834--835 repair the exact injection argument used by the local flatness and finite-presentation results; groups 838--839 and 841 repair its quotient-base successor; groups 842--844 repair the original fibre criterion; groups 845 and 848 keep the finite-type-to-finite-presentation proof noncircular; group 850 repairs the support and all-stalk reasoning of its global application. Their dependencies are traced above. Group 851 records universally exact module sequences and the full regular-sequence base-change conclusion. Group 852 records faithful flatness and purity of the receiving algebra map. Group 853 records the strictly weaker support condition and its exact receiving locus and counterexample.

No translated source theorem is strengthened by substitution. The generalizations and full proofs stay in this editorial note and the durable claim ledger. Later source review must check dependent uses of the source criteria, and the broader cross-finding synthesis remains follow-on work after the core task. The corpus query ``fibre flatness criterion'' returned four unread routing hits in its two layers. No additional external text, PDF or claim of literature exhaustiveness or novelty is part of this batch.

\hypertarget{algebra_open_flat_notes_20260928}{}
\subsection{Openness of the flat locus: corrections and separate consequences}\label{algebra-proof-039-openness-of-the-flat-locus-corrections-and-separate-consequences}

Source: Stacks Project authors, algebra.tex:33227--33514, authority a04446e57ec1fbc252a871afcec7752fb2807b14, file SHA-256 FA8BB92E58A4F78A2BD01B3B6A4A87DE0A0D279F5DD90641B574DD5FBFFFA4F3. All six reports OCC-00659--OCC-00664 are reviewed. Source lookahead 33515--33550 has been read but is not adjudicated here. No earlier canonical correction overlaps this section.

Dependencies reread in the original TeX: 23495--23645, 24825--24970, and 28208--28296. The rank-zero correction and complete criterion proof in ALGEBRA\_COMPLEX\_EXACTNESS\_NOTES\_20260928.md, section ``The exactness criterion with all rank-zero cases retained'', were reread. Use that proved criterion, including groups 572, 576 and 577, rather than repeating its earlier overclaim that every determinantal ideal is proper. The cited Buchsbaum--Eisenbud bibliography entry is human-source attribution, not a new claim to have read that paper. The preceding Noetherian approximation and flatness proofs remain the receiving dependencies.

All full arguments here are editorial supplements, not replacement translations. Source statements, formulas, matrix ranks, ideals, prime contractions, residue fields and localization denominators remain identifiable. No novelty is claimed.

\subsubsection{The nonempty zero locus and its exact equivalence}\label{algebra-proof-039-the-nonempty-zero-locus-and-its-exact-equivalence}

In lemma-CM-dim-finite-type the assertion of equality needs \(V(f_1,\ldots,f_i)\ne\varnothing\). For \(S=k[x]\), \(d=i=1\), \(f_1=1\), the zero locus is empty and has dimension strictly less than zero under the source convention. It satisfies the displayed upper bound but not the asserted equality. The regular-sequence conclusion over primes of that empty set is vacuous. A nonempty-zero-locus hypothesis repairs the dimension statement; it must be visible as a source correction.

Here is the complete proof and its converse. Let \(S\) be a finite-type, Cohen--Macaulay, equidimensional \(k\)-algebra of dimension \(d\), and let \(J=(f_1,\ldots,f_i)\) be proper. Then the following are equivalent:

\begin{enumerate}
\def\labelenumi{\arabic{enumi}.}
\tightlist
\item
  \(\dim(S/J)\leq d-i\).
\item
  The original ordered sequence is regular in \(S_{\mathfrak q}\) for every \(\mathfrak q\in V(J)\).
\item
  \(S/J\) is Cohen--Macaulay and equidimensional of dimension \(d-i\).
\end{enumerate}

For (1), every maximal ideal \(\mathfrak m\) containing \(J\) has \(\dim S_{\mathfrak m}=d\), by the source disjoint-decomposition result for Cohen--Macaulay finite-type algebras. The height theorem gives \[
d-i\leq\dim(S_{\mathfrak m}/JS_{\mathfrak m})\leq\dim(S/J)\leq d-i.
\] Thus equality holds throughout. The source characterization of parameter subsequences in a nonzero Cohen--Macaulay local module gives regularity of the specified \(f_1,\ldots,f_i\) and the Cohen--Macaulay quotient of dimension \(d-i\) at each such maximal ideal. If \(\mathfrak q\in V(J)\), choose a maximal ideal \(\mathfrak m\supset\mathfrak q\). Localizing the multiplication injections from \(S_{\mathfrak m}\) to \(S_{\mathfrak q}\) preserves them; the final quotient is nonzero because \(J\subset\mathfrak q\). This proves (2).

Conversely, (2) at those maximal ideals gives Cohen--Macaulay local quotients of dimension \(d-i\) by successive regular-element quotients. All local rings of \(S/J\) are localizations of these quotients, so the quotient is Cohen--Macaulay. Its equidimensional decomposition over the field, by the same source decomposition result, has components of the dimensions of its closed local rings, all equal to \(d-i\). It is nonzero since \(J\) is proper. This proves (3), which gives (1).

The source's ``forms'' is changed to ``form''. The phrase ``prime of \(V(J)\)'' is understandable as a prime of the quotient defining that closed subset, so the suggested replacement of ``of'' is optional prose. It does not identify a mathematical defect.

For any field extension \(K/k\), the original sequence localized at a prime over \(V(J)\) stays regular: the relevant map from \(S_{\mathfrak q}\) to the localized base change is flat and local, and the final quotient remains nonzero by faithful flatness. This preserves the exact ordered sequence. No assertion is made about regularity at primes where its ideal is the unit ideal.

\subsubsection{Openness of the fibre regular-sequence condition}\label{algebra-proof-039-openness-of-the-fibre-regular-sequence-condition}

Retain the source finite-type map \(R\to S\), equidimensional Cohen--Macaulay fibres of dimension \(d\), and the elements \(f_1,\ldots,f_i\). At a prime \(\mathfrak q\in V(J)\) with contraction \(\mathfrak p\), let \[
A=S_{\mathfrak q}/\mathfrak pS_{\mathfrak q},\qquad
r=\operatorname{trdeg}_{\kappa(\mathfrak p)}\kappa(\mathfrak q).
\] If the sequence is regular in \(A\), successive regular quotients give \(\dim(A/JA)=\dim A-i\). The original point-dimension formula at 28222--28262 therefore gives \[
\dim_{\mathfrak q}(S/R)=\dim A+r=d,\qquad
\dim_{\mathfrak q}((S/J)/R)=\dim(A/JA)+r=d-i.
\] Upper semicontinuity of the source relative-dimension function on \(\operatorname{Spec}(S/J)\) provides a principal neighborhood \(D(g)\) of \(\mathfrak q\) on which its values are at most \(d-i\).

Fix any prime \(\mathfrak q'\in V(J)\cap D(g)\), with contraction \(\mathfrak p'\). The fibre \((S_g)\otimes_R\kappa(\mathfrak p')\) is nonempty because it contains that point. It remains Cohen--Macaulay and equidimensional of dimension \(d\): each surviving component is a nonempty open subset of an original \(d\)-dimensional finite-type component over the field. The closed subscheme defined by the original \(f_j\) is nonempty and has dimension at most \(d-i\), since every point in this quotient fibre satisfies the preceding bound. The corrected nonempty-locus lemma applies. It gives regularity at \(\mathfrak q'\), proving the required openness in \(V(J)\). Empty fibres elsewhere in this localization require no dimension or regularity assertion.

\subsubsection{Exact fibre complexes under the equidimensional hypothesis}\label{algebra-proof-039-exact-fibre-complexes-under-the-equidimensional-hypothesis}

The printed lemma-exact-on-fibres-open assumes Cohen--Macaulay fibres of a fixed dimension, without saying equidimensional, but its proof invokes the preceding lemma which does require equidimensionality. These hypotheses differ: over a field, \(k[x]\times k\) is Cohen--Macaulay of dimension one and has a zero-dimensional component. Merely inserting ``equidimensional'' would narrow the original claim. Instead this note first proves the restricted case needed for the following theorem, then proves the full original claim with weaker hypotheses after that theorem. This supplies a noncircular editorial repair and preserves the original statement.

For now assume \(R\) Noetherian, \(S\) finite type and flat over \(R\), and all nonempty fibres Cohen--Macaulay and equidimensional of dimension \(d\). Retain the actual finite free complex \[
0\longrightarrow S^{n_e}\xrightarrow{\varphi_e}\cdots
\xrightarrow{\varphi_2}S^{n_1}\xrightarrow{\varphi_1}S^{n_0}.
\] Exactness here and below means exactness in the positive homological degrees, as in the cited source criterion; no surjectivity at degree zero is required.

If the exact-fibre locus is empty, openness is immediate. At a point \(\mathfrak q\) in that locus, with contraction \(\mathfrak p\), the source local complex criterion at 23528--23551 applies to \(R_{\mathfrak p}\to S_{\mathfrak q}\): all terms are \(R_{\mathfrak p}\)-flat and the closed-fibre complex is exact. Thus the original complex is exact at \(\mathfrak q\). Its positive homology modules are finite over Noetherian \(S\), so finitely many annihilators of their generating sets give \(g\notin\mathfrak q\) killing them all. Localize at their product, retaining the original maps.

For \(1\leq j\leq e\) set, with all terms and signs retained, \[
r_j=n_j-n_{j+1}+\cdots+(-1)^{e-j}n_e,\qquad
I_j=I_{r_j}(\varphi_j).
\] Positive exactness at every local ring and the corrected source rank criterion show that these \(r_j\) are nonnegative and all minors of size \(r_j+1\) vanish. They vanish in the ring itself because an element zero in all prime localizations is zero. Use the convention \(I_0(\varphi_j)=S\); rank-zero maps do not have proper determinantal ideals.

At \(\mathfrak q\), for each \(j\), either \((I_j)_{\mathfrak q}=S_{\mathfrak q}\), an open condition, or its image in \(A=S_{\mathfrak q}/\mathfrak pS_{\mathfrak q}\) contains a regular sequence of length \(j\). In the second case choose actual lifts \(a_l/s_l\in(I_j)_{\mathfrak q}\), \(s_l\notin\mathfrak q\), of that ordered sequence. Localize at \(\prod_l s_l\). Each original lift is now an element of \(I_j\) in this ring, with its exact fraction retained. This is the missing localization before treating these elements as global equations.

By the preceding regular-sequence openness lemma, the sequence is fibrewise regular on a neighborhood of \(\mathfrak q\) within its zero locus. The complement of that neighborhood is closed in the zero locus, hence closed in the ambient spectrum; shrink once more to avoid it. Outside the zero locus at least one chosen element is invertible, so \(I_j\) is the unit ideal. Inside it, \(I_j\) contains the required fibre regular sequence. For each fibre local ring these alternatives, together with the already-vanishing larger minors, give exactly the original ranks and grade conditions: a regular element in a positive-size minor ideal is nonzero, so rank cannot drop; a unit minor also cannot vanish. Apply the corrected exactness criterion. Taking the finite product of all these localization elements proves openness near \(\mathfrak q\). The complex with no positive-degree terms has the whole spectrum as its locus.

\subsubsection{Openness of flatness with the zero and one variable cases}\label{algebra-proof-039-openness-of-flatness-with-the-zero-and-one-variable-cases}

Retain the original finitely presented \(R\)-algebra \(S\), finitely presented \(S\)-module \(M\), and a point \(\mathfrak q\) where \(M_{\mathfrak q}\) is \(R\)-flat. Let \(\mathfrak p\) be its contraction. The preceding global Noetherian approximation of the entire presentation gives \(R_\lambda,S_\lambda,M_\lambda\), with exact global tensor isomorphisms for \(S\) and \(M\). Localize those stage rings at the contractions of \(\mathfrak p,\mathfrak q\). Eventual flatness from the previous section supplies a stage at which \((M_\lambda)_{\mathfrak q_\lambda}\) is \((R_\lambda)_{\mathfrak p_\lambda}\)-flat, equivalently \(R_\lambda\)-flat. A principal flat neighborhood at that stage pulls back, by the actual module tensor isomorphism and base change, to the requested neighborhood. Thus it is enough to prove the assertion when \(R,S\) are finite type over \(\mathbf Z\), hence Noetherian.

Choose the original surjection \(P=R[x_1,\ldots,x_n]\to S\), and let \(\widetilde{\mathfrak q}\) be the inverse image of \(\mathfrak q\). Regard the same module \(M\) through that surjection. It is finite, hence finitely presented, over Noetherian \(P\). Its localization at \(\widetilde{\mathfrak q}\) is the original \(M_{\mathfrak q}\). A principal element of \(P\) producing flatness maps to an element of \(S\) outside \(\mathfrak q\) and gives exactly the same module localization. This proves the legitimacy of the source polynomial presentation reduction without replacing the original quotient object in the result. In the following polynomial calculation write \(S=P\), as in the source.

If \(n=0\), \(S=R\) and \(M_{\mathfrak q}\) is a finite flat module over the Noetherian local ring \(R_{\mathfrak q}\), hence free, including rank zero. A finite basis and its inverse matrix lift after a localization, and the finitely many resulting errors are killed by another element outside \(\mathfrak q\). More explicitly, lift basis vectors to \(M_g\), let \(R_g^r\to M_g\) be the induced map, and kill the finite kernel and cokernel, which vanish at \(\mathfrak q\), by a product of annihilators outside \(\mathfrak q\). This gives a finite free principal neighborhood, proving this case.

Assume \(n\geq1\) and choose the original finite-free resolution of \(M\). Define \[
K_1=\ker(F_0\to M),\qquad
K_n=\ker(F_{n-1}\to F_{n-2})\quad(n\geq2).
\] The source's displayed formula alone has an undefined \(F_{-1}\) when \(n=1\); the preceding \(n=0\) case also cannot be put into that formula. The two boundary repairs preserve the original formula for every \(n\geq2\).

At \(\mathfrak q\), all these syzygies are \(R\)-flat by induction from \(M_{\mathfrak q}\) and the short exact sequences with the \(R\)-flat finite free terms. Thus tensoring the original truncated resolution with \(\kappa(\mathfrak p)\) remains exact there. The fibre local ring \[
A=S_{\mathfrak q}/\mathfrak pS_{\mathfrak q}
\] is a localization of a polynomial ring in \(n\) variables over \(\kappa(\mathfrak p)\); its global dimension is at most \(n\), not necessarily equal to \(n\). For instance the generic point of a polynomial line over a field gives a field of global dimension zero. Dimension shifting along the exact fibre resolution gives \(\operatorname{Ext}_A^1((K_n)_{\mathfrak q}/\mathfrak p(K_n)_{\mathfrak q},T)
\cong\operatorname{Ext}_A^{n+1}(M_{\mathfrak q}/\mathfrak pM_{\mathfrak q},T)=0\) for every \(A\)-module \(T\). Hence that finite fibre syzygy is projective, and therefore free over local \(A\).

If \((K_n)_{\mathfrak q}=0\), it is free of rank zero. Otherwise the source free-fibre-flat-free lemma at 23507--23524 applies and makes it free over \(S_{\mathfrak q}\). Since \(K_n\) is finite over Noetherian \(S\), the same finite-presentation spreading argument gives \(g\notin\mathfrak q\) for which \((K_n)_g\) is finite free. Retain its basis, the original differential into \(F_{n-1}\), and all other original matrices.

The fibres of this localized polynomial algebra are Cohen--Macaulay and equidimensional of dimension \(n\) whenever nonempty. Apply the restricted exact-fibre openness lemma just proved to \[
0\longrightarrow (K_n)_g\longrightarrow (F_{n-1})_g
\longrightarrow\cdots\longrightarrow(F_0)_g.
\] It has an exact fibre at \(\mathfrak q\), so a further principal localization makes its fibres exact everywhere. At every remaining prime the local complex criterion then gives \(R\)-flatness of its degree-zero cokernel, which is the original \(M\). For \(n=1\) this cokernel is that of \(K_1\to F_0\), equal to that of the original \(F_1\to F_0\) because both have the same image. Since flatness can be checked at all prime localizations of the module's acting algebra, \(M_g\) is \(R\)-flat. This proves the theorem without using the unproved, non-equidimensional instance of the preceding lemma.

\subsubsection{Exact-fibre openness without Cohen--Macaulay assumptions}\label{algebra-proof-039-exact-fibre-openness-without-cohenmacaulay-assumptions}

Keep \(R\) Noetherian and \(S\) a finite-type \(R\)-algebra. More generally than the printed lemma, let \[
0\longrightarrow F_e\longrightarrow\cdots\longrightarrow F_0
\] be a finite complex of finite \(S\)-modules, each flat over \(R\); neither freeness over \(S\), flatness of \(S/R\), nor any Cohen--Macaulay or fibre-dimension condition is needed. Put \(C=\operatorname{coker}(F_1\to F_0)\), taking \(C=F_0\) when \(e=0\), and let \(H_j\) be its positive homology modules.

The exact-fibre locus is precisely \[
E=
\left(\operatorname{Spec}S\setminus\bigcup_{j=1}^e\operatorname{Supp}_S H_j\right)
\cap
\{\mathfrak q:C_{\mathfrak q}\text{ is flat over }R\}.
\] If the fibre is exact at \(\mathfrak q\), the source local complex criterion gives both positive exactness of the original localized complex and \(R_{\mathfrak p}\)-flatness of \(C_{\mathfrak q}\); this is equivalent to \(R\)-flatness. Conversely assume those two properties. Beginning with \(0\to K_1\to(F_0)_{\mathfrak q}\to C_{\mathfrak q}\to0\), flatness of the two right modules proves flatness of \(K_1\) and exactness after every \(R\)-module tensor. Continue upward using positive exactness. Every syzygy is \(R\)-flat and every short exact sequence stays exact under tensor. Tensoring with \(\kappa(\mathfrak p)\) proves exactness of the fibre complex. This also proves the assertion for all zero terms and for \(e=0\).

All \(H_j\) are finite over Noetherian \(S\), so their supports are closed. The module \(C\) is finitely presented, and the map \(R\to S\) is finitely presented because it is finite type over Noetherian \(R\). The now-proved flat-locus theorem makes its flat locus open. This proves openness of \(E\), establishes the original printed exact-fibre lemma with its original hypotheses, and removes all its Cohen--Macaulay and dimension restrictions. The proof order is: equidimensional restricted lemma, flat-locus theorem, then this general lemma; there is no circular dependence.

On any open part of \(E\), the augmented complex ending in \(C\) is exact and universally exact over \(R\), by the same syzygy argument. In particular it stays exact after every \(R\)-module tensor, not only on residue-field fibres.

There is also an exact base-change identity for the fibre locus itself, without restrictions on the new base. For any map \(R\to R'\), put \(S'=S\otimes_R R'\) and use the original base-changed complex. At a prime \(\mathfrak q'\subset S'\), let \(\mathfrak p'\subset R'\), \(\mathfrak q\subset S\), \(\mathfrak p\subset R\) be its contractions. The map \[
A=S_{\mathfrak q}\otimes_R\kappa(\mathfrak p)
\longrightarrow
B=S'_{\mathfrak q'}\otimes_{R'}\kappa(\mathfrak p')
\] is a local flat map: it is the localization induced by the field extension \(\kappa(\mathfrak p)\to\kappa(\mathfrak p')\) on the original fibre. The contraction of its maximal ideal is the maximal ideal of \(A\), so it is faithfully flat. The new fibre complex is the old localized fibre complex tensor \(B\) over \(A\). Exactness therefore ascends and descends, giving \(E'=(\operatorname{Spec}S'\to\operatorname{Spec}S)^{-1}(E)\). This equality does not require Noetherianity of \(R'\) or flatness of \(R'/R\).

\subsubsection{The flat open and its radical ideal under base change}\label{algebra-proof-039-the-flat-open-and-its-radical-ideal-under-base-change}

For the original arbitrary \(R\), finitely presented \(R\)-algebra \(S\) and finitely presented \(S\)-module \(M\), let \(U\) be the flat locus of the theorem. Define the actual ideal \[
J_M=\{s\in S:M_s\text{ is flat over }R\}
    =\{s\in S:D(s)\subset U\}.
\] The equality holds by the flat-localization criterion. This is a radical ideal: \(0\) belongs to it; \(D(as)\subset D(s)\); \(D(s+t)\subset D(s)\cup D(t)\); and \(D(s^m)=D(s)\) for every positive integer \(m\). Consequently \[
U=D(J_M),\qquad
J_M=\bigcap_{\mathfrak q\notin U}\mathfrak q,
\] with empty intersection equal to \(S\). No finite-generation or quasi-compactness assertion is needed.

For a flat base change \(R\to R'\), put \(S'=S\otimes_R R'\), \(M'=M\otimes_R R'\), and let \(U'\) be its \(R'\)-flat locus. Then \(U'\) is exactly the inverse image of \(U\), and \[
J_{M'}=\sqrt{J_M S'}.
\] Forward inclusion of the loci holds by flat base change and localization, even when \(R'\) is not flat over \(R\). For the converse under flatness, take a prime \(\mathfrak q'\subset S'\) over \(\mathfrak q\) and write \(A=S_{\mathfrak q}\), \(B=S'_{\mathfrak q'}\). The map \(A\to B\) is local, flat and faithfully flat. If \(M'_{\mathfrak q'}\) is \(R'\)-flat, it is \(R\)-flat because \(R'\) is \(R\)-flat. For any injection \(T_1\to T_2\) of \(R\)-modules, the kernel of \(M_{\mathfrak q}\otimes_R T_1\to M_{\mathfrak q}\otimes_R T_2\) is an \(A\)-module. Tensoring it with flat \(B\) identifies it with the zero kernel for the \(R\)-flat module \(M_{\mathfrak q}\otimes_A B=M'_{\mathfrak q'}\). Faithfulness makes the original kernel zero. Thus \(M_{\mathfrak q}\) is \(R\)-flat. The radical-ideal identity follows from equality of the corresponding basic-open unions.

For arbitrary, possibly nonflat base change only the inclusion \(\sqrt{J_M S'}\subset J_{M'}\) follows in general. It can be strict: take \(R=S=k[t]\), \(M=R/(t)\), and \(R'=k=R/(t)\). The original flat locus is \(D(t)\), since the module is zero there and is not flat at a prime containing \(t\): multiplication by the nonzerodivisor \(t\) becomes zero on a nonzero module. After base change \(M'=k\) is free everywhere. Thus \(J_M=(t)\), \(J_MS'=0\), and \(J_{M'}=k\). This also distinguishes ordinary flat-locus base change from the unrestricted exact-fibre-locus identity above.

\subsubsection{Propagation and reading boundary}\label{algebra-proof-039-propagation-and-reading-boundary}

The nonempty-locus repair is received by the relative regular-sequence proof only at fibres with an actual point of the zero locus. The earlier rank-zero correction is received by the determinantal argument with \(I_0=S\) intact. The equidimensionality gap is repaired editorially through the noncircular proof order above, preserving the original full statement. The syzygy boundary repairs retain the original formula for \(n\geq2\), and the rank-zero syzygy needs no nonzero-module citation.

The three consequence groups give the full dimension/regularity/quotient equivalence, exact-fibre openness with much weaker hypotheses and unrestricted fibre-locus base change, and the actual flat-open radical ideal with exact flat-base-change law and a counterexample outside that hypothesis. These receive the preceding flatness results rather than replacing their source formulations.

The query ``openness flat locus'' has two research routing hits and no local-work hits. Neither was read as mathematical evidence; the first is indexed only by PDF content and no PDF fallback was used. Seven earlier external source records retain their prior bounded coverage. Broader synthesis and novelty review remain after the core task.

\hypertarget{algebra_cm_open_notes_20260928}{}
\subsection{Cohen--Macaulay loci: original charts, base changes and component maps}\label{algebra-proof-040-cohenmacaulay-loci-original-charts-base-changes-and-component-maps}

Source: Stacks Project authors, algebra.tex:33515--33787, authority a04446e57ec1fbc252a871afcec7752fb2807b14, SHA-256 FA8BB92E58A4F78A2BD01B3B6A4A87DE0A0D279F5DD90641B574DD5FBFFFA4F3. All twelve reports OCC-00665--OCC-00676 were read. There are no earlier canonical operations in this interval. Lookahead 33788--33890 is read but unadjudicated.

Original TeX dependencies reread: 27475--27533, 28048--28105, 24180--24254, 28300--28365, 4376--4438, and 31160--31226. The preceding source-bound notes for fibre dimension, local flatness and openness of flatness supply the already proved receiving arguments. All complete proofs and stronger consequences below are separate editorial material. No translation is silently recast; no novelty is claimed.

\subsubsection{Report dispositions and the empty-fibre boundary}\label{algebra-proof-040-report-dispositions-and-the-empty-fibre-boundary}

The omitted complementizer in ``prove the set \ldots{} is open'' and the abbreviated ``Trivial from \ldots'' proof are optional prose, not false mathematics. Retain those passages. Supply the two short predicate verbs in the displayed sets, use ``Set'' for the contraction at 33540, specify \(\operatorname{Spec}(S)\) at 33575, repair the \(k\)-algebra spelling, retain the fixed relative dimension in the proof's local goal, add its intervening comma, supply ``by'' in the two final declarations, and make the three map predicates grammatically parallel.

At 33762 explicitly say ``all nonempty fibres''. A factor in the product can have empty fibres, even when the original algebra has a point over each base prime. For example let \(R=k\times k\) and \(S=k[x]\times k\), with the componentwise structural map. It is flat and finitely presented over \(R\), with Cohen--Macaulay fibres. Its dimension-one factor \(k[x]\) has empty fibre over the second component of the base; its dimension-zero factor \(k\) has empty fibre over the first. The intended component-dimension assertion applies to each nonempty fibre and keeps the zero factors and empty spectra harmless. This clarification does not replace any nonempty fibre by another object.

\subsubsection{The original quasi-finite flat chart criterion}\label{algebra-proof-040-the-original-quasi-finite-flat-chart-criterion}

Let \(k\) be a field, \(S\) a finite-type \(k\)-algebra and \(P=k[y_1,\ldots,y_d]\to S\) the original quasi-finite map. For \(\mathfrak q\subset S\), retain the actual contraction \(\mathfrak p\subset P\), and write \(A=P_{\mathfrak p}\), \(B=S_{\mathfrak q}\). These are nonzero Noetherian local rings and their map is local. Quasi-finiteness gives \(\dim(B/\mathfrak pB)=0\), the dimension inequality \(\dim B\leq\dim A\), and a finite residue-field extension \(\kappa(\mathfrak q)/\kappa(\mathfrak p)\).

If \(B\) is \(P\)-flat, it is \(A\)-flat. The source lemma at 27489--27503 has a second alternative: flatness together with \(\dim B\leq\dim A\) suffices. This alternative applies since \(A\) is regular, hence Cohen--Macaulay. One must not pretend that the localized quasi-finite map \(A\to B\) is necessarily finite. It follows that \(B\) is Cohen--Macaulay and \(\dim B=\dim A\). Retaining the source point-dimension formula and the exact residue fields gives \[
\dim_{\mathfrak q}(S/k)
=\dim B+\operatorname{trdeg}_k\kappa(\mathfrak q)
=\dim A+\operatorname{trdeg}_k\kappa(\mathfrak p)=d.
\]

Conversely, if \(B\) is Cohen--Macaulay and the left side is \(d\), equality of the two transcendence degrees implies \(\dim B=\dim A\). Together with \(\dim(B/\mathfrak pB)=0\), this is the dimension identity required by miracle flatness. The original \(A\) is regular. The source local criterion therefore proves \(B\) flat over \(A\), and then over \(P\), since \(A\) is a localization of \(P\). This proves both directions of the displayed source locus identity with the original objects and maps.

For the field case of openness, at a Cohen--Macaulay prime use the original Noether normalization at that point: some \(S_g\) has a finite injective map from \(k[y_1,\ldots,y_d]\), where \(d=\dim_{\mathfrak q}(S/k)\). Apply the just-proved criterion at that prime, then the proved openness-of-flatness theorem to obtain a further principal neighborhood flat over the same polynomial algebra. The map remains quasi-finite after this localization. The criterion at each remaining prime makes every local ring Cohen--Macaulay.

Every minimal prime of a finite-type algebra over a field belongs to this open: its local ring is zero-dimensional Noetherian, hence Artinian and Cohen--Macaulay. The spectrum of such an algebra is Noetherian, and every nonempty open meets a minimal prime of a component it meets. Thus the open is dense, including the vacuous assertion when the spectrum is empty.

\subsubsection{The original relative proof with fixed dimension retained}\label{algebra-proof-040-the-original-relative-proof-with-fixed-dimension-retained}

Let \(R\to S\) be flat of finite presentation and let \(\mathfrak q\) have Cohen--Macaulay fibre local ring and relative dimension \(d\). Write \(\mathfrak p=R\cap\mathfrak q\). The source quasi-finite chart lemma at 31170--31208 gives a principal neighborhood of \(\mathfrak q\) and an actual quasi-finite map \[
P=R[t_1,\ldots,t_d]\longrightarrow S.
\] Retain these original parameter images and set \(\mathfrak q'=P\cap\mathfrak q\). On the residue-field fibre the preceding field criterion gives flatness of \[
P_{\mathfrak q'}/\mathfrak pP_{\mathfrak q'}
\longrightarrow S_{\mathfrak q}/\mathfrak pS_{\mathfrak q}.
\] Apply the local fibre criterion to \(R_{\mathfrak p}\to P_{\mathfrak q'}\to S_{\mathfrak q}\) and \(M=S_{\mathfrak q}\). The last module is nonzero and free of rank one over its own local ring, is flat over \(R_{\mathfrak p}\), and has the closed-fibre flatness just proved. Both displayed algebras are essentially finitely presented over \(R_{\mathfrak p}\). Thus \(S_{\mathfrak q}\) is \(P_{\mathfrak q'}\)-flat.

The map \(P\to S\) is finitely presented, a hypothesis worth verifying before applying openness of flatness. Write the actual finite presentation \(S=R[x_1,\ldots,x_m]/(h_1,\ldots,h_l)\), including any variables and inverse equations introduced by the principal localization. Choose polynomial representatives \(F_j(x)\) for the actual images of \(t_j\). Then \[
S\cong
P[x_1,\ldots,x_m]/
(h_1,\ldots,h_l,\ t_1-F_1(x),\ldots,t_d-F_d(x)).
\] The mutually inverse maps send each original \(x_a\) and \(t_j\) to its specified element; every original relation is retained. This proves finite presentation over \(P\).

Openness of flatness now supplies a further \(g\notin\mathfrak q\) for which \(S_g\) is flat over \(P\). At every \(\mathfrak r\subset S_g\), contract to \(\mathfrak r'\subset P\) and \(\mathfrak p'\subset R\). Base change of the local flat map to the fibre over \(\mathfrak p'\), followed by the field criterion, proves both that its fibre local ring is Cohen--Macaulay and that \(\dim_{\mathfrak r}(S_g/R)=d\). The proof's intermediate goal must retain this fixed-dimension conclusion; Cohen--Macaulayness alone would not identify the displayed locus.

The union over \(d\) of these open strata is the fibre Cohen--Macaulay locus \(W\). On each fibre it is exactly the open Cohen--Macaulay locus of the corresponding finite-type algebra over its residue field, hence dense there. It is also dense in the total spectrum: any nonempty open contains a point, and its intersection with that point's fibre is nonempty open in the fibre, so it meets \(W\).

\subsubsection{Field extension and the actual local tensor comparison}\label{algebra-proof-040-field-extension-and-the-actual-local-tensor-comparison}

Retain \(S_K=K\otimes_k S\) and a prime \(\mathfrak q_K\) over \(\mathfrak q\). The original pointwise Noether normalization provides a principal localization with finite injective \(P=k[x_1,\ldots,x_d]\to S\), where \(\dim_{\mathfrak q}(S/k)=d\). Flat extension of the field preserves this finite injection. The source point-dimension result at 28307--28355 proves \(\dim_{\mathfrak q_K}(S_K/K)=d\).

Write \(\mathfrak p=P\cap\mathfrak q\), \(P_K=K[x_1,\ldots,x_d]\), and \(\mathfrak p_K=P_K\cap\mathfrak q_K\). In the original square of local rings, the bottom map \(P_{\mathfrak p}\to(P_K)_{\mathfrak p_K}\) is flat. The top-right ring \((S_K)_{\mathfrak q_K}\) is the localization of \[
S_{\mathfrak q}\otimes_{P_{\mathfrak p}}(P_K)_{\mathfrak p_K}
\] at the prime induced by \(\mathfrak q_K\), not an unqualified tensor equality. Every denominator from \(S\setminus\mathfrak q\) and \(P_K\setminus\mathfrak p_K\) is a unit there, so the universal property gives the original ring with precisely that remaining prime localization.

The source flat-up-and-down result at 24185--24236 therefore proves that the two vertical maps are flat simultaneously. Its downward implication is faithful descent through the top local flat map and retains the actual localization. Apply the field chart criterion on each side to conclude \[
S_{\mathfrak q}\text{ is Cohen–Macaulay}
\quad\Longleftrightarrow\quad
(S_K)_{\mathfrak q_K}\text{ is Cohen–Macaulay}.
\] This proof does not assume that \(K/k\) is algebraic, finite or separable.

For any finite-type map \(R\to S\), any \(R\to R'\), and \(S'=R'\otimes_R S\), a prime \(\mathfrak q'\subset S'\) gives contractions \(\mathfrak q\subset S\), \(\mathfrak p'\subset R'\), \(\mathfrak p\subset R\). The new fibre local ring is a prime localization of the old fibre algebra \(S\otimes_R\kappa(\mathfrak p)\) after the exact field extension \(\kappa(\mathfrak p)\to\kappa(\mathfrak p')\). The preceding equivalence proves the source equality \(W'=f^{-1}(W)\). In addition, the point-dimension theorem proves equality of the relative dimensions at \(\mathfrak q'\) and \(\mathfrak q\). Neither equality needs flatness of the base change.

\subsubsection{A maximal open without assuming global flatness}\label{algebra-proof-040-a-maximal-open-without-assuming-global-flatness}

Group 878. Suppose only that \(R\to S\) is finitely presented. Define \[
C_d=\{\mathfrak q:S_{\mathfrak q}\text{ is }R\text{-flat},
\ S_{\mathfrak q}\otimes_R\kappa(\mathfrak p)\text{ is Cohen–Macaulay},
\ \dim_{\mathfrak q}(S/R)=d\},
\qquad \mathfrak p=R\cap\mathfrak q.
\] Then each \(C_d\) is open. At a point of \(C_d\), choose the original quasi-finite \(d\)-parameter chart \(P=R[t_1,\ldots,t_d]\to S_g\). Within that chart the exact equality is \[
C_d\cap D(g)
=\{\mathfrak r\in D(g):(S_g)_{\mathfrak r}
       \text{ is flat over }P\}.
\] Indeed \(P\)-flatness implies \(R\)-flatness because \(P\) is \(R\)-flat, and the field chart criterion on every fibre proves Cohen--Macaulayness and relative dimension \(d\). In the other direction, apply the local fibre criterion exactly as in the preceding relative proof, using the assumed flatness of the one local ring \((S_g)_{\mathfrak r}\) over \(R\), not global flatness of \(S_g/R\). The same finite presentation over \(P\) proved there allows openness of its flat locus. This proves the equality and openness near every point of \(C_d\).

Let \(C=\bigcup_d C_d\). It is the largest open part on which the original morphism is flat with Cohen--Macaulay fibres. Any such open has every local ring flat over \(R\) and every fibre local ring Cohen--Macaulay, so is contained in \(C\). Conversely those two properties hold at every point of \(C\), which is precisely their local meaning for the restricted morphism. The \(C_d\) are disjoint and open; each is closed in \(C\), since its complement there is the union of the others. Finite type bounds the fibre dimension by the number of generators in any fixed algebra presentation, so only finitely many \(d\) occur.

For a flat base change \(R\to R'\), the \(C_d\) pull back exactly. The ordinary flat-locus equality from group 865 and the just-proved arbitrary-base-change invariance of fibre Cohen--Macaulayness and relative dimension prove all three conditions. For arbitrary base change the inverse image is contained in \(C'_d\), but equality can fail: take \(R=k[t]\), \(S=k=R/(t)\). The source algebra is finitely presented and every nonempty fibre is Cohen--Macaulay of relative dimension zero, but its sole local ring is not \(R\)-flat because multiplication by \(t\) kills a nonzero module. Thus \(C=\varnothing\). After base change to \(R'=k\), the map is the identity of \(k\), and \(C'_0=\operatorname{Spec}k\).

\subsubsection{Universal dimension strata and their radical ideals}\label{algebra-proof-040-universal-dimension-strata-and-their-radical-ideals}

Group 879. For any finite-type \(R\to S\), let \(W_d\) consist of the primes whose fibre local ring is Cohen--Macaulay and whose relative dimension is \(d\). The field-extension proof gives, for every ring map \(R\to R'\), \[
W'_d=f^{-1}(W_d)
\] with no flatness or finite-presentation hypothesis for this set identity. Flatness and finite presentation of \(S/R\) make these sets open by the source theorem; the base-changed map remains flat and finitely presented, so the new sets are open too.

Under these latter hypotheses define the actual radical ideal \[
J_d=\{s\in S:D(s)\subset W_d\}
    =\bigcap_{\mathfrak q\notin W_d}\mathfrak q .
\] The principal-open union and prime-intersection identities prove \(W_d=D(J_d)\); closure under addition and multiplication and radicality follow from \(D(a+b)\subset D(a)\cup D(b)\), \(D(ab)\subset D(a)\), and \(D(a^n)=D(a)\). The exact set identity under arbitrary base change therefore gives \[
J'_d=\sqrt{J_d S'}
\] for every \(R\to R'\), including nonflat maps. The same proof applies to the union \(W\). This differs from the ordinary flat-open ideal law, which needs a flat base change for equality. No finite generation of these ideals is asserted.

On \(W=\coprod_d W_d\) relative dimension is a locally constant function with finite range, bounded by a fixed number of algebra generators. Each \(W_d\) is closed in \(W\), not necessarily closed in \(\operatorname{Spec}S\). For a flat finitely presented map, \(W\) is dense in every fibre and in the total spectrum as proved above. These assertions retain the original relative point-dimension function, not the varying Krull dimensions of the fibre local rings.

\subsubsection{Canonical dimension idempotents with all nilpotents retained}\label{algebra-proof-040-canonical-dimension-idempotents-with-all-nilpotents-retained}

Group 880. Assume \(R\to S\) is flat, finitely presented, and has Cohen--Macaulay fibres. Choose any finite algebra generating list of size \(n\). Every fibre is a quotient of an \(n\)-variable polynomial ring over a field, so at every point \(0\leq\dim_{\mathfrak q}(S/R)\leq n\). Thus \[
\operatorname{Spec}S=\coprod_{d=0}^n W_d .
\] Each \(W_d\) is open and its complement is the finite union of the other opens, hence it is closed. The original disjoint-decomposition lemma gives an idempotent \(e_d\in S\) with \(D(e_d)=W_d\).

These idempotents are unique, including empty strata. If \(e,f\) are idempotents defining the same open, then \(e(1-f)\) and \(f(1-e)\) are idempotents with empty basic opens. Such an idempotent lies in every prime and is nilpotent, hence is zero. Therefore \(e=ef=f\). The same argument proves \(e_de_c=0\) for \(d\ne c\), and \(1-\sum_d e_d=0\). This reasoning does not replace \(S\) by its reduced ring.

Define the original component algebra \[
S_d=e_dS=S/(1-e_d)S\cong S_{e_d},
\] with unit \(e_d\) and structural map \(r\mapsto e_d\varphi(r)\). The exact mutually inverse product maps are \[
S\longrightarrow\prod_{d=0}^nS_d,\quad s\longmapsto(e_ds)_d,
\qquad
(s_d)_d\longmapsto\sum_{d=0}^n s_d.
\] All original nilpotents survive in their corresponding factors. Since \(S_d\) is an \(R\)-module direct summand of the flat module \(S\), it is \(R\)-flat. Its displayed quotient adds the one relation \(1-e_d\) to a finite algebra presentation, so it is finitely presented over \(R\). Every nonempty fibre is Cohen--Macaulay and equidimensional of dimension \(d\): all of its points have the original point-dimension value \(d\), in particular its generic points; their component closures therefore have dimension \(d\). Empty fibres and zero factors require no finite numerical dimension assertion.

For every \(R\to R'\), the image \(e'_d=e_d\otimes1\) has basic open \(f^{-1}(W_d)=W'_d\), by the preceding arbitrary-base-change result. Uniqueness makes it the dimension-\(d\) idempotent of \(S'=S\otimes_R R'\). The component comparison \[
S_d\otimes_R R'\longrightarrow e'_dS',
\qquad (e_ds)\otimes r'\longmapsto e'_d(s\otimes r')
\] is an isomorphism: tensor the direct decomposition \(S=e_dS\oplus(1-e_d)S\) and retain its two projections. Thus the full finite product decomposition commutes with every base change. A component may become zero; it is not reassigned a different dimension.

The radical ideal \(J_d\) of the preceding section is generally not \(e_dS\). Its exact formula is \[
J_d=\sqrt{e_dS}=e_dS+\sqrt{(0)_S}.
\] If \(x^m\in e_dS\), then \(((1-e_d)x)^m=0\), and \(x=e_dx+(1-e_d)x\) lies in the last ideal. Conversely its image modulo \(e_dS\) is nilpotent, giving the reverse inclusion. For \(S=k[\epsilon]/(\epsilon^2)\), the empty dimension-one stratum has \(e_1=0\) but \(J_1=(\epsilon)\ne0\). This illustrates why replacing the radical ideal by the idempotent ideal would erase relevant nilpotent data.

Every \(S\)-module has the exact component decomposition \[
M\longrightarrow\prod_{d=0}^n e_dM,\quad m\longmapsto(e_dm)_d,
\qquad (m_d)_d\longmapsto\sum_d m_d.
\] The composites are identities by orthogonality and the sum-to-one equation. Every \(S\)-linear map preserves these components, and these decompositions commute with tensor base change by the same split projections. Finally, if \(\operatorname{Spec}S\) is nonempty and connected, exactly one \(e_d\) is nonzero, and all its nonempty fibres are equidimensional of that one dimension. Connectedness of the base is insufficient: over the field \(k\), the flat finitely presented algebra \(k[x]\times k\) has dimensions one and zero on its two components.

\subsubsection{Propagation and reading boundary}\label{algebra-proof-040-propagation-and-reading-boundary}

The chart criterion uses the actual second alternative of the earlier Cohen--Macaulay flatness lemma, rather than a false finite-local-map assertion. The local fibre criterion and the corrected openness-of-flatness proof are used with their full hypotheses. Group 878 propagates these to the maximal flat Cohen--Macaulay open without a global flatness assumption. Group 879 adds the dimension-preserving universal set and radical-ideal identities. Group 880 gives the canonical component maps, their universal comparisons and the exact nilpotent correction to the radical ideals.

All three consequence groups have zero translation replacement operations. The source product proof was checked directly: its finite range follows either from the generator bound just proved or compactness of the affine spectrum. No new hypothesis is needed to make the product finite.

The query ``Cohen Macaulay locus base change'' returns two research routing hits and one local-work hit. None is read as mathematical evidence; the local hit is a PDF and was not used. Earlier external-source coverage is unchanged. No claim of novelty, literature exhaustiveness or source mutation accompanies this batch. Broader synthesis remains follow-on work after core completion.

\hypertarget{algebra_differentials_notes_20260928}{}
\subsection{Differentials: source review and separate editorial arguments}\label{algebra-proof-041-differentials-source-review-and-separate-editorial-arguments}

Authority: Stacks Project authors, \texttt{algebra.tex} at commit \texttt{a04446e57ec1fbc252a871afcec7752fb2807b14}, SHA256 \texttt{FA8BB92E58A4F78A2BD01B3B6A4A87DE0A0D279F5DD90641B574DD5FBFFFA4F3}, lines 33788--34386. All twenty received reports were read. This note supplies editorial proofs and consequences; it does not replace the translated exposition. The current source interval has no earlier canonical operation. Subsequent de Rham material at 34387--34520 is read lookahead, not adjudicated here.

All rings are commutative with identity, and all module actions are through the specified ring maps. Zero rings, zero modules, positive characteristics, and zero divisors are allowed throughout unless an example explicitly specifies a nonzero field. The three consequence sections below make no novelty claim. The indexed query \texttt{Kahler\ differentials\ conormal\ splitting} returned two research records and one local record; none was read or used as evidence. Earlier external reading records retain their exact coverage.

\subsubsection{Report dispositions and preservation of the source}\label{algebra-proof-041-report-dispositions-and-preservation-of-the-source}

OCC-00677 changes the display-final period to a comma at 33967. OCC-00678 is optional: the demonstrative in ``functoriality of this'' can refer to the preceding presentation. OCC-00679 corrects the lifting corner at 34001 from lower right to lower left; both relation modules occur on the left of the displayed square. OCC-00680 is optional: 34008 already explicitly cites the direct sum decomposition; the detailed reduction is supplied below. OCC-00681 removes the ill-typed application of \(\varphi:S\to S'\) to \(s'_l\in S'\) at 34014. OCC-00682 supplies ``the'' at 34025. OCC-00683 replaces ring elements by their differentials in the kernel description at 34037. These are local source corrections, not rewritten proofs.

OCC-00684 supplies ``an'' at 34096; OCC-00696 adds a comma at 34093. The derivation used to construct the localization inverse must extend the original universal derivation. The source's next displayed formula does precisely this, so the argument does not require a replacement. OCC-00691 adds the period at 34115, OCC-00685 supplies ``by'' at 34123, and OCC-00686 corrects ``show'' to ``shows'' at 34130. OCC-00695 requests an optional complementizer. OCC-00692 and OCC-00693 add periods at 34184 and 34270. OCC-00687 restores the established differential notation at 34278. OCC-00688 requests expansion of understandable abbreviated explanations; OCC-00689 requests an optional complementizer. Both are retained. OCC-00694 adds the period at 34360; OCC-00690 supplies the article at 34371. The independently noticed ``there exist an'' at 34046 becomes ``there exists an''.

Thus the twenty reports produce fifteen proposed local operations, with five optional proposals retained unchanged; the unreported agreement correction adds one operation. Full exercises, reconstructed arguments, and the three stronger consequences stay in this note. The fixed-base convention is examined explicitly below without silently treating it as a different theorem.

\subsubsection{Universal presentation, functoriality and directed colimits}\label{algebra-proof-041-universal-presentation-functoriality-and-directed-colimits}

At 33794--33880 let \(F=\bigoplus_{a\in S}S[a]\). Let \(L\) be generated by every \([a+b]-[a]-[b]\), \([fg]-f[g]-g[f]\), and \([\alpha(r)]\). The map \(d(a)=[a]+L\) is additive, kills \(\alpha(R)\), and satisfies the Leibniz rule. Conversely an \(R\)-derivation \(D:S\to M\) gives the unique \(S\)-linear map \(F\to M\), \(\sum s_a[a]\mapsto\sum s_aD(a)\), which kills exactly the required relations and therefore factors uniquely through \(F/L\). Composition with \(d\) recovers \(D\); a linear map is recovered because the \(da\) generate. These two maps commute with every module homomorphism. Also \(D(1)=D(1)+D(1)\), hence \(D(1)=0\); an \(R\)-linear Leibniz map therefore kills \(\alpha(R)\). This proves the equivalent definition, including zero rings. If \(\alpha\) is onto, every generator \(da\) is zero.

For the square at 33894--33955, \(\varphi\alpha=\beta\psi\). The formula \(s\,da\mapsto\varphi(s)d\varphi(a)\) carries additive and product relations to the corresponding relations, and \(d\varphi\alpha(r)=d\beta\psi(r)=0\). It is \(\varphi\)-semilinear and compositions agree on every generator.

For 33958--33975, put \(R=\operatorname{colim}R_i\), \(S=\operatorname{colim}S_i\), and \(M=\operatorname{colim}\Omega_{S_i/R_i}\) as abelian groups with the induced \(S\)-action. The action is defined by representing a scalar and a differential at a common stage; independence follows by moving two representatives to a further common stage where they agree. Define \(d:S\to M\) by a representative \(s_i\mapsto[d_i s_i]\). An equality in a directed colimit holds at some further stage, so this is independent of the representative. The derivation laws and vanishing on \(R\) follow at a common stage. This gives \(u:\Omega_{S/R}\to M\). The compatible functorial maps give \(v:M\to\Omega_{S/R}\). Both composites are identities on every represented differential, hence on all elements. This proves the asserted colimit with its actual varying-ring actions.

\subsubsection{The surjective square and both kernel arguments}\label{algebra-proof-041-the-surjective-square-and-both-kernel-arguments}

At 33977--34054, let \(I=\ker\varphi\), assume \(\varphi:S\to S'\) onto, and let \(K\subset\Omega_{S/R}\) be the \(S\)-submodule generated by \(da\) with \(\varphi(a)\in\beta(R')\). Its image in \(\Omega_{S'/R'}\) is zero. Surjectivity follows by lifting both the coefficient and the argument of each \(s'da'\).

For the first proof, lift each additive or product relation, including its coefficient, from the upper-left free relation module to the lower-left one. Subtract its image from the free representative of the original differential. The remaining equality has only base relations on the right. Decompose the equality into the finitely many basis indices \(a'\in S'\) that occur. In each component, choose a lift \(a\). If \(a_i\) has this image, then \(a_i=a+x_i\), \(x_i\in I\), and \(da_i=da+dx_i\equiv da\pmod K\). That component is therefore \(s\,da\) modulo \(K\). Its upper coefficient is \(\varphi(s)=\sum s'_l\), where the sum runs over base relations with \(\beta(r'_l)=a'\); it is zero if that set is empty. If such a base relation occurs, then \(da\in K\). Otherwise \(s\in I\), and \[
s\,da=d(sa)-a\,ds\in K,
\] because both \(s\) and \(sa\) belong to \(I\). Thus every component, and hence the entire kernel, belongs to \(K\). This also proves \(I\Omega_{S/R}\subset K\); it is not an implicit cancellation of the coefficients.

For the second proof let \(R''=S\times_{S'}R'\), with its given map from \(R\) and projection to \(S\). Quotienting \(\Omega_{S/R}\) by the differentials of the image of \(R''\) has precisely the universal property of \(\Omega_{S/R''}\). The image consists of those \(a\) with \(\varphi(a)\in\beta(R')\), so this quotient is \(\Omega_{S/R}/K\). Every \(i\in I\) is the projection of \((i,0)\in R''\); consequently \(di=0\) in \(\Omega_{S/R''}\). For every \(s\), \(0=d(is)=i\,ds\), since \(is\in I\). Therefore the whole universal module is annihilated by \(I\). Its derivation descends to \(S'\). If \(r'\in R'\), choose \(s\) with \(\varphi(s)=\beta(r')\). Then \((s,r')\in R''\), so the descended derivation kills \(\beta(r')\). It induces \(\Omega_{S'/R'}\to\Omega_{S/R''}\), inverse to the canonical map because both composites fix every universal differential. This completes the omitted argument without changing the source presentation.

For 34056--34075, apply this kernel formula to \(S=S'=C\), \(R=A\), \(R'=B\). The kernel is generated by differentials of elements in the image of \(B\). It is exactly the image of \(C\otimes_B\Omega_{B/A}\to\Omega_{C/A}\), \(c\otimes db\mapsto c\,d(b)\). This proves transitivity and does not assert injectivity of its first arrow.

\subsubsection{Localization with the complete denominators}\label{algebra-proof-041-localization-with-the-complete-denominators}

At 34077--34100, if \(t\in A\) maps to a unit, \(0=D(1)=\varphi(t)D(\varphi(t)^{-1})\), so \(D(\varphi(t)^{-1})=0\). The derivations over \(A\) and over its indicated localization are identical, proving (1).

For (2), write \(T\subset B\) for the source's multiplicative set to avoid confusing it with a ring. Work in \(N=T^{-1}\Omega_{B/A}\), retaining the formula \[
d'(b/s)=s^{-1}db-bs^{-2}ds.
\] If \(b/s=c/t\), there is \(u\in T\) with \(u(bt-cs)=0\). In \(T^{-1}B\), \(bt-cs=0\); differentiating the original equality and then passing to \(N\) yields \(t\,db+b\,dt-s\,dc-c\,ds=0\), because \(u\) is a unit. Multiplying the difference of the two displayed formulas by the unit \(s^2t^2\) gives \[
st^2db-bt^2ds-s^2t\,dc+cs^2dt
=t(sc-bt)ds+s(cs-bt)dt=0.
\] Here the equality uses \(st(t\,db-s\,dc)=st(c\,ds-b\,dt)\). Thus the formula is independent of representatives, without a non-zero-divisor assumption.

For addition substitute \((bt+cs)/(st)\), use \(d(bt+cs)=t\,db+b\,dt+s\,dc+c\,ds\) and \(d(st)=t\,ds+s\,dt\). The terms containing \(cs\,ds\) and \(bt\,dt\) cancel with the matching numerator terms, leaving \(s^{-1}db-bs^{-2}ds+t^{-1}dc-ct^{-2}dt\). For products substitute \(bc/(st)\); the result is \[
\frac{c}{st}db+\frac{b}{st}dc-\frac{bc}{s^2t}ds-\frac{bc}{st^2}dt
=\frac ct d'(b/s)+\frac bs d'(c/t).
\] The formula kills the image of \(A\) and restricts to \(d\) on \(B\). It therefore induces an inverse to the canonical localized differential map: one composite fixes \(db/s\), and the other fixes \(d(b/s)\) by differentiating \(s^{-1}\) with the Leibniz rule. If \(0\in T\), both the localized ring and module are zero and the same conclusions hold.

\subsubsection{Fixed base, conormal maps and powers of the ideal}\label{algebra-proof-041-fixed-base-conormal-maps-and-powers-of-the-ideal}

The standard fixed-base interpretation of 34102--34155 is \(R'=R\) with \(\psi=\mathrm{id}_R\); in particular \(\varphi\alpha=\beta\). Equality of the two ring objects alone would not suffice. For example take \(R=\mathbf F_p[t]\), \(S=S'=\mathbf F_p[x]\), \(\alpha(t)=x^p\), \(\beta(t)=x\), \(\psi(t)=t^p\), and \(\varphi=\mathrm{id}\). The square commutes, but \(I=0\), \(\Omega_{S/R}=S\,dx\), and \(\Omega_{S'/R}=0\), so the asserted exact sequence with zero left term would fail. The first equality follows because derivatives of all polynomials in \(x^p\) vanish and \(d/dx\) remains a nonzero derivation; the second follows because \(\beta\) is onto. This records the precise map convention, rather than counting an ordinary fixed-base convention as a proved source error. In the split statement an actual \(R\)-algebra section already forces \(\varphi\alpha=\beta\).

Henceforth set \(B=S/I\) with its quotient \(R\)-structure, \(N=I/I^2\), \(E=\Omega_{S/R}\otimes_S B\), and \(Q=\Omega_{B/R}\). Define \(\delta:N\to E\) by \([f]\mapsto df\otimes1\), and \(\pi:E\to Q\) by \(ds\otimes b\mapsto b\,d\bar s\). If \(f,g\in I\), \(d(fg)=f\,dg+g\,df\) maps to zero; thus \(\delta\) is well defined. Also \(d(sf)=s\,df+f\,ds\) maps to \(\bar s\delta[f]\), proving \(B\)-linearity. The surjective-square result says that the kernel of \(\Omega_{S/R}\to Q\) is generated by \(da\) whose image comes from \(R\). Writing \(a=\alpha(r)+(a-\alpha(r))\) gives \(da=d(a-\alpha(r))\), with the parenthesized element in \(I\). Thus that kernel is generated by \(dI\), and passage modulo \(I\Omega_{S/R}\) gives \(\operatorname{im}\delta=\ker\pi\). In particular \(N\to E\to Q\to0\) is exact, with no left injectivity assertion.

For an \(R\)-algebra section \(\beta:B\to S\), set \(D(s)=[s-\beta(\bar s)]\in N\). If \(u=s-\beta(\bar s)\) and \(v=t-\beta(\bar t)\), then \[
st-\beta(\bar s\bar t)=\beta(\bar s)v+\beta(\bar t)u+uv.
\] The last term lies in \(I^2\); hence \(D(st)=\bar sD(t)+\bar tD(s)\). Additivity and vanishing on \(R\) follow directly. It gives \(\tau:E\to N\) with \(\tau(ds\otimes b)=bD(s)\), and \(\tau\delta[f]=[f]\). The section also gives \(\sigma:Q\to E\), \(db\mapsto d\beta(b)\otimes1\). On generators \(\pi\sigma=1\), \(\tau\sigma=0\), and \(\delta\tau+\sigma\pi=1_E\). Thus \((n,q)\mapsto\delta n+\sigma q\) and \(e\mapsto(\tau e,\pi e)\) are inverse \(B\)-linear maps. This proves all omitted verifications in the source split lemma.

For 34193--34210, every element of \(I^{n+1}\) is a finite sum of products of \(n+1\) elements of \(I\) (with a coefficient absorbed into one factor). Leibniz gives all \(n+1\) terms, each in \(I^n\Omega_{S/R}\). The conormal sequence for \(S\to S/I^{n+1}\), after tensoring with \(S/I^n\), is right exact and its left arrow is zero. Its middle-to-right arrow is therefore an isomorphism. No division by \(n+1\) or characteristic restriction is used.

\subsubsection{Base change, the diagonal and finite presentations}\label{algebra-proof-041-base-change-the-diagonal-and-finite-presentations}

For 34212--34233 put \(S'=S\otimes_R R'\). The map \(D'(\sum a_i\otimes x_i)=\sum da_i\otimes x_i\) is defined by an \(R\)-balanced bilinear rule. On elementary tensors its product rule is \[
D'((a\otimes x)(b\otimes y))=(a\,db+b\,da)\otimes xy
=(a\otimes x)D'(b\otimes y)+(b\otimes y)D'(a\otimes x).
\] It kills \(1\otimes R'\). An \(R'\)-derivation restricted to \(S\otimes1\) factors uniquely through \(\Omega_{S/R}\); extending that factor by \(\omega\otimes x\mapsto x\,f(\omega)\) gives the unique \(S'\)-linear factor. Thus the asserted base-change isomorphism holds for arbitrary \(R'\), without flatness.

For 34235--34304 write \(A=S\otimes_R S\), \(\mu(a\otimes b)=ab\), \(J=\ker\mu\), and \(\Delta(b)=1\otimes b-b\otimes1\). The two actions of \(S\) on \(J/J^2\) agree, since their difference on \(j\in J\) is \((s\otimes1-1\otimes s)j\in J^2\). The map \(b\mapsto[\Delta(b)]\) is additive, kills \(R\), and satisfies the derivation law modulo \(J^2\): the product defect is \(\Delta(a)\Delta(b)\). It induces \[
\theta:\Omega_{S/R}\longrightarrow J/J^2,\qquad a\,db\longmapsto[a\otimes b-ab\otimes1].
\] All three sorts of defining relations, including additivity, have now been checked.

Define \(F:A\to\Omega_{S/R}\) by \(F(\sum a_i\otimes b_i)=\sum a_i\,db_i\). It is well defined by the tensor balancing rule and satisfies \(F(uv)=\mu(u)F(v)+\mu(v)F(u)\), as is checked on elementary tensors and then expanded additively. Thus it kills \(J^2\) and induces \(\bar F:J/J^2\to\Omega_{S/R}\). One composite fixes every \(a\,db\). If \(j=\sum a_i\otimes b_i\in J\), then \(\sum a_ib_i=0\), so \[
\theta\bar F[j]=\left[\sum_i(a_i\otimes b_i-a_ib_i\otimes1)\right]=[j].
\] These are exact inverses. For the source's first proof, relative differentials over the first factor kill \(a\otimes1\); the base-change tensor with \(\Omega_{S/R}\) uses the second factor. Differentiating \(a\otimes b-ab\otimes1\) relative to the first factor and then applying \(\mu\) yields \(a\,db\), the same inverse just proved. No swap of the two factor actions is suppressed.

For 34309--34339 and \(P=R[x_1,\ldots,x_n]\), write each polynomial uniquely as \(f=\sum_\nu a_\nu x^\nu\) with finite support. Define \[
\partial_i f=\sum_{\nu_i>0}\nu_i a_\nu x^{\nu-e_i}.
\] The integer coefficients retain their actual images in \(R\). On two monomials, \(\nu_i+\lambda_i\) is the sum of their two integer exponents; this proves Leibniz, and additivity extends it to all polynomials. Repeated Leibniz gives \(df=\sum_i\partial_i f\,dx_i\). The induced coordinate functionals take \(dx_j\) to \(\delta_{ij}\); together they are inverse to the free-module map with basis \(dx_i\). Empty variable lists and the zero polynomial cause no exception.

For \(S=P/(f_1,\ldots,f_m)\), every \(f\in I\) is \(\sum h_jf_j\), and modulo \(I\Omega_{P/R}\) its differential is \(\sum\bar h_j\,df_j\). Consequently the exact finite presentation is \[
S^m\xrightarrow{(\overline{\partial_i f_j})_{i,j}}S^n\longrightarrow\Omega_{S/R}\longrightarrow0.
\] No injectivity of the matrix is inferred. For finite type without finite presentation, the same conormal sequence with an arbitrary ideal makes \(\Omega_{S/R}\) a quotient of \(S^n\), which proves finite generation. For example \(S=\mathbf F_p[x]/(x^p)\) retains the coefficient \(p=0\); the conormal generator maps to zero while \(I/I^2\ne0\). Indeed multiplication by \(x^p\) identifies \(\mathbf F_p[x]/(x^p)\) with \((x^p)/(x^{2p})\). This demonstrates why a left zero cannot be added to the general conormal sequence.

\subsubsection{First-order sections are exactly conormal retractions}\label{algebra-proof-041-first-order-sections-are-exactly-conormal-retractions}

Editorial consequence, receiving the fixed-base conormal sequence, its split proof, and the ideal-power comparison at 34102--34210. Retain \(B=S/I\), \(N=I/I^2\), \(E=\Omega_{S/R}\otimes_S B\). The conormal map \(\delta:N\to E\) has a \(B\)-linear left inverse if and only if the quotient \(q:A=S/I^2\to B\) has an \(R\)-algebra section. Thus an actual section into \(S\), assumed by the source, can be replaced in the editorial statement by a section into this first-order quotient. The equivalence also specifies all such retractions.

Let \(\tau\delta=1_N\). Define \(D:S\to N\), \(D(s)=\tau(ds\otimes1)\). This is an \(R\)-derivation with \(D(i)=[i]\) for \(i\in I\). It kills \(I^2\) by Leibniz, so induces \(D_A:A\to N\), restricting to the identity on \(N\subset A\). For \(b\in B\), choose any \(a\in A\) above it and set \[
\beta(b)=a-D_A(a)\in A.
\] If \(a\) changes by \(n\in N\), its derivative changes by that same \(n\), proving independence. Additivity, \(\beta(1)=1\), and agreement on \(R\) follow from the derivation laws. For \(a,a'\in A\), \(N^2=0\) gives \[
(a-D_A(a))(a'-D_A(a'))
=aa'-q(a)D_A(a')-q(a')D_A(a)
=aa'-D_A(aa').
\] Hence \(\beta\) is multiplicative and \(q\beta=1_B\).

Conversely, given such a section, define \(D_A(a)=a-\beta(q(a))\in N\). The same product expansion as in the original split proof makes it an \(R\)-derivation with identity restriction to \(N\). Composing with \(S\to A\) gives \(D:S\to N\), hence a unique \(\tau:E\to N\) with \(\tau\delta=1_N\). The two constructions are inverse by their displayed formulas. The power comparison identifies \(E\) with \(\Omega_{A/R}\otimes_A B\), so a section into \(A\) also supplies the full split exact sequence and its complementary map from \(\Omega_{B/R}\). These arguments do not assume finite generation, flatness, or a characteristic restriction.

The weakening is strict. Let \(k\) be a nonzero field, \(S=k[x]\), \(h=x^2-x\), \(I=(h)\), and identify \(B=S/I\) with \(k\times k\) by evaluation at \(0,1\). Surjectivity follows from \((a,b)\mapsto a(1-x)+bx\) modulo \(h\), and a polynomial vanishing at both points is divisible by \(x(x-1)\); thus this is an isomorphism. Put \[
e=x-(2x-1)h=3x^2-2x^3,
\qquad e^2-e=(4h-3)h^2.
\] Therefore \(e\) is idempotent in \(S/I^2\), with image \((0,1)\) in \(B\). The map \((a,b)\mapsto a(1-e)+be\) is a \(k\)-algebra section: the complementary idempotents multiply to zero and sum to one. All coefficients remain valid even in characteristics \(2\) and \(3\). There is no section into \(k[x]\), because a section would send \((0,1)\) to an idempotent with that image, whereas the domain \(k[x]\) has only idempotents \(0,1\). This proves an actual strengthening of the section hypothesis without altering the source lemma.

\subsubsection{All sections and the exact change of splitting}\label{algebra-proof-041-all-sections-and-the-exact-change-of-splitting}

Editorial consequence receiving the preceding equivalence. Suppose one \(R\)-algebra section \(\beta:B\to A=S/I^2\) exists. For every other section \(\beta'\), the difference \(h(b)=\beta'(b)-\beta(b)\) lies in \(N\). It is additive, kills \(R\), and \[
h(bb')=b\,h(b')+b'\,h(b),
\] because the product of the two differences lies in \(N^2=0\). Conversely every \(R\)-derivation \(h:B\to N\) makes \(\beta+h\) an \(R\)-algebra section by that product calculation. The action is free and transitive, so the sections form a torsor under \(\operatorname{Der}_R(B,N)=\operatorname{Hom}_B(\Omega_{B/R},N)\). This statement asserts a torsor only when a section exists.

Let \(\bar h:Q\to N\) correspond to \(h\). Let \((\tau,\sigma)\) and \((\tau',\sigma')\) be the conormal splitting maps associated with \(\beta\) and \(\beta+h\), using \(E=\Omega_{A/R}\otimes_A B\). Their exact formulas are \[
\tau'=\tau-\bar h\pi,\qquad
\sigma'=\sigma+\delta\bar h.
\] The first follows on \(da\) from \(a-\beta'(q(a))=a-\beta(q(a))-h(q(a))\); the second follows on \(db\) by differentiating \(\beta'(b)=\beta(b)+h(b)\). Thus the old coordinates \((n,q)\) of an element of \(E\) become \((n-\bar h(q),q)\) in the new splitting; conversely the new coordinates \((n',q)\) become \((n'+\bar h(q),q)\) in the old one. Every sign and map is fixed. If sections into \(S\) itself exist, their reductions give this formula; no claim is made that every first-order section lifts to \(S\), as the preceding example disproves.

\subsubsection{The diagonal quotient as an algebra, its flip and base change}\label{algebra-proof-041-the-diagonal-quotient-as-an-algebra-its-flip-and-base-change}

Editorial consequence receiving 34212--34304. With \(A=S\otimes_R S\), \(J=\ker\mu\), the original quotient algebra \(A/J^2\), not merely its kernel module, admits an explicit comparison with the square-zero algebra \(S\oplus\Omega_{S/R}\). The latter has multiplication \((s,\omega)(t,\eta)=(st,s\eta+t\omega)\), unit \((1,0)\), and the \(R\)-structure through the first summand. Define \[
\Phi\left(\left[\sum_i a_i\otimes b_i\right]\right)
=\left(\sum_i a_ib_i,\sum_i a_i\,db_i\right).
\] The previously proved rule for \(F\) proves multiplicativity. Its value on \(J^2\) is zero in both coordinates, so it is defined on the quotient. Define the reverse map \[
\Psi(s,\omega)=[s\otimes1]+\theta(\omega),
\] where \(\theta(\omega)\in J/J^2\) is included in \(A/J^2\). This is a ring map because the kernel is square zero and its \(S\)-action is the one already proved. The composite \(\Phi\Psi\) fixes both summands. For \([a\otimes b]\), the other composite gives \([ab\otimes1]+[a\otimes b-ab\otimes1]=[a\otimes b]\). Thus these are inverse algebra maps; the original quotient and all its tensor coordinates are retained.

The first structural map \(s\mapsto[s\otimes1]\) becomes \((s,0)\); the second \(s\mapsto[1\otimes s]\) becomes \((s,ds)\). The factor flip \(a\otimes b\mapsto b\otimes a\) therefore becomes \[
T(s,\omega)=(s,ds-\omega).
\] Indeed \(b\,da=d(ab)-a\,db\) on elementary tensors. This formula is additive and multiplicative by Leibniz, and \(T^2=1\). On the kernel it acts by \(-1\), including the case where that equals \(+1\) in characteristic \(2\). It interchanges the two \(S\)-structures; it is not claimed to be linear for the first structure alone.

Finally let \(R\to R'\) be arbitrary and \(S'=S\otimes_R R'\). The natural tensor rearrangement sends \((a\otimes b)\otimes r'\) to \((a\otimes1)\otimes_{R'}(b\otimes r')\), identifying \(A\otimes_R R'\) with \(A'=S'\otimes_{R'}S'\). The sequence \(0\to J\to A\xrightarrow\mu S\to0\) splits as \(R\)-modules via \(s\mapsto s\otimes1\); consequently after any base change its kernel is exactly \(J\otimes_R R'\), identified with \(J'\). The image of \(J^2\otimes_R R'\) in \(A'\) is \((J')^2\): products of finite tensor sums expand into images of products from \(J\), and each image of a product belongs to that square. Right exactness therefore gives \[
(A/J^2)\otimes_R R'\ \cong\ A'/(J')^2.
\] No injectivity of \(J^2\otimes_R R'\to A'\) is asserted or needed. Under the already proved differential base-change map, \((a\,db)\otimes r'\mapsto(a\otimes r')d(b\otimes1)\), the maps \(\Phi,\Psi\), both structural sections, and \(T\) commute with this comparison on every elementary tensor. This proves the universal compatibility without a flatness assumption.

\subsubsection{Propagation and remaining work}\label{algebra-proof-041-propagation-and-remaining-work}

The wrong lifting corner, ill-typed coefficient, and omitted differentials are corrected locally in the review proposals. The fixed-base convention is recorded with its exact failure example. The surjection proof supplies transitivity and conormal exactness; the conormal and ideal-power arguments supply the first-order section criterion; that criterion supplies the torsor and both signed changes of splitting. The diagonal comparison retains both factor actions and supplies the quotient algebra, flip, and arbitrary-base-change maps. The polynomial proof retains all exponent coefficients and supplies the finite Jacobian presentation without adding unjustified injectivity.

The three consequence groups contain no source-text operations. Later uses, beginning with the de Rham complex, remain to be reviewed against these statements. Broader cross-work synthesis follows the core review; this note makes neither a comprehensive literature claim nor a publication claim. No translated source has been rewritten, and no build has been performed.

\hypertarget{algebra_de_rham_notes_20260928}{}
\subsection{The de Rham complex: source review and editorial consequences}\label{algebra-proof-042-the-de-rham-complex-source-review-and-editorial-consequences}

Authority: Stacks Project authors, \texttt{algebra.tex} at \texttt{a04446e57ec1fbc252a871afcec7752fb2807b14}, SHA256 \texttt{FA8BB92E58A4F78A2BD01B3B6A4A87DE0A0D279F5DD90641B574DD5FBFFFA4F3}, lines 34387--34650. All twelve reports are adjudicated here. The current interval equals the authority, with no prior canonical correction. Lines 34651--34780 are read lookahead in the following differential-operator section and remain unadjudicated. The exterior presentation dependency 2240--2404 was reread. The preceding differential arguments are in \texttt{ALGEBRA\_DIFFERENTIALS\_NOTES\_20260928.md}, SHA256 \texttt{0EC7E677030EF57E9AAA46CCBC90104CB59F57E386BF08F1ADE48A13937A62F9}.

This is separate editorial material. Source statements and exposition remain identifiable; none of the expanded proofs or stronger consequences is a silent translation replacement. The corpus query \texttt{de\ Rham\ differential\ ideal\ quotient} returned two research and two local records. The Bhatt--Lurie record retains only its previously recorded bounded reading at 11899--11920; no new content from any query hit was read or used. The other three hits remain unread, including the PDF-only record. There is no novelty or exhaustive-reading claim.

\subsubsection{Dispositions and the additive convention}\label{algebra-proof-042-dispositions-and-the-additive-convention}

OCC-00697 corrects the opening declaration: \(d:B\to\Omega_{B/A}\) denotes the universal derivation into the module, rather than denoting the module itself. OCC-00698 supplies a period at 34414. OCC-00699 requests explicit bounds for the sum at 34465. The source has exactly \(p\) displayed tensor positions and uses \(i\) for the differentiated position; the bounds are determined by that notation. This optional clarification is retained unchanged, while the proof below states the bounds explicitly. OCC-00700 supplies a comma at 34498, OCC-00701 a comma at 34535, OCC-00702 a period after the diagram at 34545, and OCC-00703 a period at 34549. OCC-00704 supplies ``by'' at 34577. OCC-00705 supplies a period at 34591. OCC-00706 supplies a period after the diagram at 34607 and capitalizes the next sentence at 34609. OCC-00707 changes ``left most'' to ``leftmost''; OCC-00708 supplies the comma at 34633. These are twelve operations in eleven groups, with one optional report left unchanged.

The absolute-module discussion at 34525--34537 retains additivity from the definition of derivation and the universal presentation. With that convention, Leibniz gives \(D(1)=2D(1)\), so \(D(1)=0\); additivity then gives \(D(n)=0\) for every integer, including negative integers. Thus additive Leibniz maps are precisely \(\mathbf Z\)-derivations. Leibniz alone without additivity is a different condition: on \(\mathbf R\), set \(D(0)=0\), \(D(x)=x\log|x|\) for \(x\ne0\). For nonzero \(x,y\), the logarithm product identity gives \(D(xy)=xD(y)+yD(x)\); if either is zero both sides are zero. But \(D(2)=2\log2\ne D(1)+D(1)=0\). This verifies the precise retained convention, without treating the source's reference to its existing additive derivation as an error.

\subsubsection{Construction in degree one}\label{algebra-proof-042-construction-in-degree-one}

Write \(M=\Omega_{B/A}\), with universal derivation \(d_0\). On the free \(B\)-module with generators \([b]\), define an additive \(A\)-linear rule \[
\sum_i b'_i[b_i]\longmapsto\sum_i d_0b'_i\wedge d_0b_i.
\] It is \(A\)-linear because \(d_0(ab')=a\,d_0b'\) for \(a\) in the image of \(A\). It sends \(b[a]\), with \(a\) in that image, to zero. Additive relations map to \(d_0b\wedge(d_0(b'+b'')-d_0b'-d_0b'')=0\). A multiplied Leibniz relation maps to \[
d_0b\wedge(b'\,d_0b''+b''\,d_0b')
-(b\,d_0b'+b'\,d_0b)\wedge d_0b''
-(b\,d_0b''+b''\,d_0b)\wedge d_0b'=0.
\] The two remaining terms are \(-b(d_0b'\wedge d_0b''+d_0b''\wedge d_0b')=0\). Anticommutativity here follows by expanding \((u+v)\wedge(u+v)=0\), without division by two. Thus the rule descends to \(d_1:M\to\bigwedge_B^2M\) over every base ring.

For \(\omega=\sum b_i\,d_0c_i\), expansion gives the useful identity \[
d_1(f\omega)=d_0f\wedge\omega+f\,d_1\omega.
\] Indeed both sides equal \(\sum (b_i\,d_0f+f\,d_0b_i)\wedge d_0c_i\). Also \(d_1d_0b=d_01\wedge d_0b=0\). These are identities of the actual universal modules, with no freeness assumption.

\subsubsection{Higher degrees, alternating signs and the graded product rule}\label{algebra-proof-042-higher-degrees-alternating-signs-and-the-graded-product-rule}

For \(p\ge2\) define the \(A\)-multilinear map \[
\gamma(\omega_1,\ldots,\omega_p)
=\sum_{i=1}^p(-1)^{i+1}\omega_1\wedge\cdots\wedge d_1\omega_i\wedge\cdots\wedge\omega_p.
\] All wedge products are over \(B\), even though the initial tensor product is over \(A\). First, swapping two adjacent positions reverses the sign. Terms differentiated outside these positions reverse sign by swapping two degree-one factors. If the positions are \(j,j+1\), their two terms are \[
(-1)^{j+1}P\wedge d_1\omega_j\wedge\omega_{j+1}\wedge Q
+(-1)^{j+2}P\wedge\omega_j\wedge d_1\omega_{j+1}\wedge Q.
\] After the swap, commuting a degree-two factor past a degree-one factor has sign \((-1)^2=1\); the resulting coefficients are the negatives of the respective original coefficients. If the adjacent factors are equal, the outside terms contain their repeated wedge and vanish; the two displayed terms cancel because the degree-two factor commutes with that degree-one factor. An arbitrary repetition can be brought adjacent by the already proved swap identity, so \(\gamma\) is alternating. The cancellation is valid in characteristic two as well.

For a scalar \(f\in B\) inserted in position \(i\), the product identity in degree one gives \[
\gamma(\omega_1,\ldots,f\omega_i,\ldots,\omega_p)
=f\gamma(\omega_1,\ldots,\omega_p)
+d_0f\wedge\omega_1\wedge\cdots\wedge\omega_p.
\] In the extra term, moving \(d_0f\) past \(i-1\) preceding degree-one factors changes its sign by \((-1)^{i-1}\), which cancels \((-1)^{i+1}\). The expression is independent of \(i\). Thus \(\gamma\) kills every scalar-balancing relation in 2353--2367 as well as every alternating relation. Those relations generate the kernel: first quotient the \(A\)-tensor power by scalar transfers to obtain the \(B\)-tensor power, then quotient by repeated tensors to obtain its exterior power. This proves that \(\gamma\) descends to \(d_p\).

For the source's original generators, all \(d_1d_0b_i\) terms vanish and \[
d_p(b_0d_0b_1\wedge\cdots\wedge d_0b_p)
=d_0b_0\wedge d_0b_1\wedge\cdots\wedge d_0b_p.
\] Every exterior power is additively generated by these expressions, so this proves uniqueness. Applying the next differential writes the result with coefficient \(1\), and hence gives zero because \(d_01=0\). Degree zero was handled above; therefore \(d_{p+1}d_p=0\) for every \(p\ge0\).

Let \(\alpha\) have degree \(r\) and \(\beta\) degree \(s\). If \(r,s\ge1\), expand them into wedges of one-forms and split the sum defining \(\gamma\) between the first \(r\) positions and the last \(s\). The last positions contribute a factor \((-1)^r\), giving \[
d(\alpha\wedge\beta)=d\alpha\wedge\beta+(-1)^r\alpha\wedge d\beta.
\] If one degree is zero, this is the scalar identity just proved, with the degree-one differential moved across \(r\) factors when necessary. If both are zero it is the universal Leibniz rule. This proves the full graded product rule used at 34620 and 34631, including all degree and sign conventions. The proof is editorial; no expanded argument is substituted for the source construction.

\subsubsection{Functoriality, base change and the quotient construction}\label{algebra-proof-042-functoriality-base-change-and-the-quotient-construction}

For a commutative square with maps \(f:B\to B'\) and \(A\to A'\), the degree-\(p\) map sends \(b_0db_1\wedge\cdots\wedge db_p\) to \(f(b_0)df(b_1)\wedge\cdots\wedge df(b_p)\). It is well defined because the degree-one map respects every universal relation and because its induced tensor map respects the alternating relations and scalar transfers. The displayed differential formula proves commutation with \(d\) on generators; compositions and identities do so as well.

If \(B'=B\otimes_A A'\), the preceding differential note proves \(M\otimes_A A'\cong\Omega_{B'/A'}\). Exterior powers commute with this base change directly from their presentation: tensor powers commute by the maps on elementary tensors, and the image of the submodule of alternating relations is exactly the submodule generated by the base-changed alternating relations. For the latter assertion, squares of sums expand into repeated terms and pairs \(u\otimes v+v\otimes u\), each a consequence of squares. Equivalently the alternating multilinear universal property on base-changed generators gives inverse maps. Thus \[
(\bigwedge_B^pM)\otimes_A A'\cong\bigwedge_{B'}^p(M\otimes_A A')
\] without flatness. In degree zero this is \(B\otimes_A A'=B'\). In every degree, on a generating tensor the differential acts on the original \(b_i\) and kills coefficients from \(A'\), so these isomorphisms assemble into an isomorphism of complexes and of graded algebras.

For the source quotient lemma let \(\pi:M\twoheadrightarrow Q\), \(K=\ker\pi\), and \(q_p:\bigwedge^pM\twoheadrightarrow\bigwedge^pQ\). Its kernel is the degree-\(p\) part of the ideal generated by \(K\), namely the image of \(K\otimes_B\bigwedge^{p-1}M\) under the wedge map for \(p\ge1\). To see this, quotienting the tensor algebra by \(K\) gives the tensor algebra of \(Q\), and then impose the same alternating relations. This proves the exterior-presentation assertion at 2240--2267 and its use at 34626--34628 without an injection claim for \(K\otimes_B\bigwedge^{p-1}M\).

If \(K\) is generated by \(\kappa_i\) with \(q_2(d\kappa_i)=0\), then \[
q_2d\left(\sum b_i\kappa_i\right)
=\sum q_2(db_i\wedge\kappa_i)+\sum b_iq_2(d\kappa_i)=0.
\] For \(\kappa\in K\) and \(\eta\in\bigwedge^{p-1}M\), the graded product rule gives \(d(\kappa\wedge\eta)=d\kappa\wedge\eta-\kappa\wedge d\eta\), whose projection vanishes. Hence every square in the source diagram exists. Since every \(q_p\) is onto, the projected differential is unique; its square is zero by lifting an element to the original complex. The displayed rule follows by applying \(q\) to each original generator.

\subsubsection{The exact obstruction to an exterior quotient complex}\label{algebra-proof-042-the-exact-obstruction-to-an-exterior-quotient-complex}

Editorial consequence receiving 34574--34638. The map \[
c:K\longrightarrow\bigwedge_B^2Q,\qquad c(\kappa)=q_2(d\kappa)
\] is \(B\)-linear, since \(q_2d(b\kappa)=q_2(db\wedge\kappa)+bq_2d\kappa=bc(\kappa)\). The source's condition on some generating set is equivalent to \(c=0\), and this condition is both necessary and sufficient for the exterior algebra of \(Q\) to carry the differential specified in the source.

Sufficiency is the preceding proof. For necessity, any such differential commutes with projection in every degree because the generator rule characterizes it. In particular \(q_2d\kappa=d(q_1\kappa)=d0=0\). This also proves uniqueness and identifies the exact module map whose nonvanishing prevents the construction, not merely a difference of presentations.

The obstruction can be nonzero even when the kernel is generated by a single one-form. Take a nonzero field \(k\), \(B=k[x,y]\), \(A=k\), and \(K=B(x\,dy)\subset M=B\,dx\oplus B\,dy\). Then \[
Q=B\,dx\oplus(B/(x))\,dy,
\qquad\bigwedge^2_BQ=(B/(x))\,(dx\wedge dy).
\] The second assertion follows by expanding alternating tensors: the two cyclic summands have zero second exterior power, and their cross term is \(B/(x)\); the coefficient map taking \(dx\wedge dy\) to \(1\bmod x\) supplies the inverse. Now \(c(x\,dy)=dx\wedge dy\ne0\). A hypothetical exterior quotient differential would therefore send the zero relation \(x\,dy\) to this nonzero class. This counterexample retains the original quotient module and proves precisely where the required differential fails.

If \(A\to A'\) is any base change, set \(B'=B\otimes_A A'\), \(Q'=Q\otimes_A A'\), and let \(K'\) be the actual kernel of the resulting surjection \(M\otimes_A A'\to Q'\). Right exactness says \(K'\) is the image of \(K\otimes_A A'\), not necessarily an isomorphic copy. For \(\kappa\otimes a'\), the base-changed obstruction is \(c(\kappa)\otimes a'\). Consequently vanishing ascends under every base change. If \(A\to A'\) is faithfully flat, it also descends: the map \(\bigwedge^2Q\to(\bigwedge^2Q)\otimes_A A'\) is injective, since faithful flatness detects the nonzero cyclic submodule generated by any proposed kernel element. Thus its vanishing forces \(c(\kappa)=0\) for every \(\kappa\).

\subsubsection{The smallest differential quotient when the obstruction is nonzero}\label{algebra-proof-042-the-smallest-differential-quotient-when-the-obstruction-is-nonzero}

Editorial consequence retaining the same objects and the exact obstruction. In \(\Omega^\bullet=\bigwedge_B^\bullet M\), form the graded ideal \(\mathcal J\) generated by \(K\) in degree one and \(dK\) in degree two. Explicitly, \[
\mathcal J^p=K\wedge\Omega^{p-1}+dK\wedge\Omega^{p-2},
\] where a term with negative degree is zero and the notation means the \(B\)-span of the indicated wedges. For each generator \(\kappa\wedge\alpha\), its differential is \(d\kappa\wedge\alpha-\kappa\wedge d\alpha\). For each \(d\kappa\wedge\beta\), it is \(d\kappa\wedge d\beta\), because \(d^2\kappa=0\). Coefficients are already allowed inside \(\alpha,\beta\); their derivatives give the same included terms. Hence \(d\mathcal J\subset\mathcal J\).

Every differential ideal containing \(K\) must also contain \(dK\), so \(\mathcal J\) is the smallest such ideal. The quotient \(\Omega^\bullet/\mathcal J\) is therefore a well-defined differential graded algebra. Its degree zero is still \(B\) and its degree one is still \(Q\). As a graded algebra it is exactly \[
\left(\bigwedge_B^\bullet Q\right)\big/\left(c(K)\right),
\] with the ideal on the right generated by the actual degree-two image of \(c\). This follows by first quotienting by the ideal generated by \(K\), and then by the images of \(dK\). Thus the original obstruction is retained as the precise additional relation module. When \(c=0\), this recovers the source exterior quotient complex; when \(c\ne0\), the exterior algebra itself cannot carry that differential, but this smaller graded algebra does.

The construction is universal: a morphism of differential graded \(A\)-algebras out of \(\Omega^\bullet\) that kills \(K\) kills \(dK\) by commutation with the differential, hence factors uniquely through the quotient. Conversely every map through the quotient kills \(K\). In the preceding example, \(c(K)\) generates all of \(\bigwedge^2Q\); higher exterior powers are zero because \(Q\) has two generators. The resulting complex is therefore exactly \(B\to Q\to0\), with the original projected derivative in degree zero. This supplies the actual construction determined by the obstruction.

Under any base change \(A\to A'\), the image of \(\mathcal J\otimes_A A'\) is the ideal generated by \(K'\) and \(dK'\). Indeed every element of \(K'\) is a finite sum of images of \(\kappa\otimes a'\), whose derivative is the image of \(d\kappa\otimes a'\), and all their products are images of the indicated generators. Conversely those images lie in that ideal. Right exactness and the already proved base-change isomorphism of de Rham algebras now give \[
(\Omega^\bullet/\mathcal J)\otimes_A A'
\cong \Omega_{B'/A'}^\bullet/\mathcal J'.
\] No injectivity assertion for a tensor product of ideals is used. This is an exact base-change comparison for the constructed quotient even when the obstruction survives.

\subsubsection{Complex base change and the exact polynomial cohomology boundary}\label{algebra-proof-042-complex-base-change-and-the-exact-polynomial-cohomology-boundary}

Editorial consequence receiving 34559--34572 and the original polynomial differential calculation at 34309--34339. For every ring \(A\), retain the original polynomial coordinate \(x\) and \(B=A[x]\). The complex has exactly two possibly nonzero degrees: \[
A[x]\xrightarrow d A[x]\,dx,
\qquad d\left(\sum_{n\ge0}a_nx^n\right)=\sum_{n\ge1}n a_n x^{n-1}dx,
\] with finite support and every integer coefficient retained in \(A\). Higher powers vanish because \(\Omega^1\) is free of rank one. Comparing coefficients independently gives the exact \(A\)-module descriptions \[
H^0=A\ \oplus\!\bigoplus_{n\ge1}\operatorname{Ann}_A(n)x^n,
\qquad
H^1=\bigoplus_{n\ge1}(A/nA)x^{n-1}dx,
\qquad H^p=0\ (p\ge2).
\] The kernel description follows because \(n a_n=0\) for each \(n\); the image in the \(x^{n-1}dx\) coefficient is exactly \(nA\), giving the cokernel description. Direct sums, not products, record the finite polynomial support. Terms such as \(A/A=0\) are retained in the indexing.

For a ring map \(A\to A'\), the comparison \(H^0\otimes_A A'\to H^0(\Omega_{A'[x]/A'}^\bullet)\) is the identity on the constant summand and, on each positive degree, the actual map \[
\operatorname{Ann}_A(n)\otimes_A A'\longrightarrow\operatorname{Ann}_{A'}(n),\qquad a\otimes a'\longmapsto\operatorname{image}(a)a'.
\] No general injectivity or surjectivity is asserted. In degree one the comparison is always an isomorphism, since \((A/nA)\otimes_A A'\cong A'/nA'\), with the displayed coordinates fixed. For flat base change, kernels commute with tensoring the exact sequence \(0\to\operatorname{Ann}_A(n)\to A\xrightarrow n A\), so the degree-zero comparison is also an isomorphism. More generally flatness gives cohomology base change for any complex: tensor the exact sequences from cycles to the terms to boundaries, and then from boundaries to cycles to cohomology. Exactness preserves both kernels and images in those sequences.

For the explicit nonflat base change \(\mathbf Z\to\mathbf F_p\), with \(p\) prime, the formulas give \[
H^0(\Omega_{\mathbf Z[x]/\mathbf Z}^\bullet)=\mathbf Z,
\qquad H^1=\bigoplus_{n\ge1}(\mathbf Z/n\mathbf Z)x^{n-1}dx,
\] whereas \[
H^0(\Omega_{\mathbf F_p[x]/\mathbf F_p}^\bullet)=\mathbf F_p[x^p],
\qquad H^1=\bigoplus_{m\ge1}\mathbf F_p\,x^{mp-1}dx.
\] The degree-zero comparison is the inclusion of constants, with cokernel \(\bigoplus_{m\ge1}\mathbf F_p x^{mp}\). It is not surjective because \(x^p\) is not a constant polynomial, even though its derivative is \(p x^{p-1}dx=0\). The degree-one comparison remains the coordinatewise isomorphism just proved. Thus the source's arbitrary-base-change statement for complexes is correct and must not be silently promoted to arbitrary cohomology base change. The failed comparison is measured here with every exponent and coefficient intact.

\subsubsection{Propagation and continuation}\label{algebra-proof-042-propagation-and-continuation}

The degree-one calculation proves the scalar product rule; the higher-degree construction proves the alternating descent and full graded product rule used by the quotient lemma. The exact kernel map \(c\) makes the source hypothesis necessary as well as sufficient. When it does not vanish, its image defines the additional higher relations in the universal differential quotient, with explicit maps back to the original objects. The polynomial cohomology calculation provides the exact boundary of the source base-change statement.

These three consequence groups have no translation operations. Their future use must retain the actual quotient module, the additional degree-two relations when required, and the distinction between complexes and cohomology. The next source section begins at 34651; its material, including the order-one differential-operator connection, remains to be adjudicated. A broader synthesis remains follow-on work after the core review. No translation rewrite, TeX build, publication, new task or delegation occurred.

\hypertarget{algebra_diff_operators_notes_20260928}{}
\subsection{Finite-order differential operators: source review and editorial arguments}\label{algebra-proof-043-finite-order-differential-operators-source-review-and-editorial-arguments}

Authority: Stacks Project authors, \texttt{algebra.tex} at \texttt{a04446e57ec1fbc252a871afcec7752fb2807b14}, SHA256 \texttt{FA8BB92E58A4F78A2BD01B3B6A4A87DE0A0D279F5DD90641B574DD5FBFFFA4F3}, lines 34651--35115. All twenty-four reports, including the duplicate diagonal-statement report OCC-12206, are adjudicated. The whole current interval equals the authority, with no earlier canonical operation. The naive cotangent complex at 35116--35330 is read lookahead, not adjudicated here.

Bounded dependencies were read at \texttt{categories.tex:536–621} (including the complete Yoneda lemma and representability discussion), authority SHA256 \texttt{62F7611AF4C3FEEBD041DB4728B42C7112004CFBB9FA5ECB643C6F5D90DB3F25}, and \texttt{algebra.tex:291–376} (including the snake lemma). The cited original Yoneda, Grothendieck and Cartan--Eilenberg works were not newly read. Arguments needed here are supplied below. The preceding differential and de Rham notes are retained with SHA256 \texttt{0EC7E677030EF57E9AAA46CCBC90104CB59F57E386BF08F1ADE48A13937A62F9} and \texttt{E7527FC4B77594FB1F50FB94471756AC92EF5B8D3BACD94B1A23C4D3707E5992} respectively.

The query \texttt{principal\ parts\ differential\ operators} returned two research PDF records and two local TeX records; none was read or adopted as evidence. There was no PDF fallback. Earlier external reading records retain their exact coverage. All proofs and stronger consequences here are separate editorial material, with no novelty claim or authorization to rewrite the source exposition.

\subsubsection{Dispositions and the order convention}\label{algebra-proof-043-dispositions-and-the-order-convention}

OCC-00709 adds a period at 34680. OCC-00710 corrects the reversed universal factorization at 34716 to \(D=\alpha\circ D_{univ}\), with the original domains and codomains retained. OCC-00711 identifies an unfinished citation placeholder. The proposed local correction removes that parenthetical; it does not invent a reference. The source's explicit generators-and-relations construction already proves existence and remains in place. OCC-00712 and OCC-00713 add periods at 34744 and 34770. OCC-00714 adds a period at 34814; OCC-00715 supplies ``to'' at 34819. OCC-00716 and OCC-00717 punctuate the two functoriality diagrams; OCC-00718 adds the comma before ``but''.

OCC-00719 and OCC-12206 describe the same missing-verb defect at 34938--34939 and receive one shared correction: the single universal operator corresponds to the specified class map. OCC-00720 and OCC-00721 are handled together at 34943, supplying the tensor base \(R\) and the sentence-final period. OCC-00722 adds a period at 34984. OCC-00723 makes \(k\ge1\) explicit in the generator criterion, since the source has not defined negative order. OCC-00724 punctuates the product commutator. OCC-00727 punctuates the inference at 35039; OCC-00728, OCC-00729, OCC-00730 and OCC-00731 punctuate the localization formulas. OCC-00725 replaces both undefined \(R\)-linear occurrences at 35057--35058 by \(A\)-linear. OCC-00726 supplies the actual order-zero tensor argument before the induction step. These are twenty-three local operations in twenty-two report groups; all twenty-four reports receive exactly one disposition.

The source's ``order \(k\)'' is the filtered convention expressed by its inclusions \(\operatorname{Diff}^0\subset\operatorname{Diff}^1\subset\cdots\); it does not assert minimal order exactly \(k\). All statements below use those original nonnegative indices. Whenever a commutator would have negative order in an induction, the corresponding order-zero commutator is zero and is treated directly.

\subsubsection{Commutators, composition and the original universal presentation}\label{algebra-proof-043-commutators-composition-and-the-original-universal-presentation}

For an \(R\)-linear map \(D:M\to N\), write \(\delta_gD=D\circ g-g\circ D\), where multiplication on each side acts on its specified module. The operators \(\delta_g\) commute: both \(\delta_g\delta_hD\) and \(\delta_h\delta_gD\) equal \[
D\circ gh-g\circ D\circ h-h\circ D\circ g+gh\circ D.
\] Induction on \(k\) gives the precise criterion: \(D\) has order \(k\) if and only if every \((k+1)\)-fold iterated commutator vanishes. For \(k=0\) this is exactly \(S\)-linearity; the induction step applies the criterion to each \(\delta_gD\). The full expansion, for any \(r\ge1\), is \[
(\delta_{g_1}\cdots\delta_{g_r}D)(m)
=\sum_{U\subset\{1,\ldots,r\}}(-1)^{r-|U|}
\left(\prod_{j\notin U}g_j\right)
D\left(\left(\prod_{j\in U}g_j\right)m\right).
\] Induction proves the expansion by splitting subsets according to membership of the last index. Empty products equal \(1\), and every summand is retained. The source's displayed relation with initial positive coefficient is \((-1)^r\) times this expression for \(r=k+1\), so its vanishing is exactly equivalent, including characteristic two.

The identity \[
\delta_g(D'\circ D)=D'\circ\delta_gD+(\delta_gD')\circ D
\] proves the composition lemma by induction on \(k+k'\). If both orders are zero, the composite is \(S\)-linear. If just one is zero, its displayed commutator term is zero; the other term has total order at most \(k+k'-1\). If both are positive, both terms have that bound by induction. Sums and multiplication by scalars preserve each filtered piece by the same commutator criterion.

At 34700--34751, quotient the original free module \(F=\bigoplus_{m\in M}S[m]\) by its additive relations, its \(R\)-linearity relations, and the full alternating \((k+1)\)-fold relations. A linear map \(F\to N\) annihilates the first two sets exactly when \(m\mapsto L([m])\) is \(R\)-linear, and the last set exactly when that map has order \(k\), by the expansion just proved. Therefore this quotient represents the stated operator functor. The universal map is \(j_k(m)=[m]\), and its unique factorization is \(D=L_D\circ j_k\). In particular \(L_D\) follows \(j_k\), not the reverse. Two representing pairs have inverse linear maps because both composites fix all \(j_k(m)\), which generate. This proves the universal assertion directly; no unspecified category-theory existence theorem is needed.

The identity on generators gives the quotient maps \(P^k(M)\twoheadrightarrow P^{k-1}(M)\), since vanishing of every \(k\)-fold commutator implies vanishing of every \((k+1)\)-fold one. Their composites are the original tower maps. All constructions are functorial in \(M\): a linear map sends the generator \([m]\) to the corresponding generator for its image and respects every displayed relation.

\subsubsection{The diagonal comparison with both actions retained}\label{algebra-proof-043-the-diagonal-comparison-with-both-actions-retained}

Put \(C=S\otimes_R S\), \(J=\ker(\mu:C\to S)\), and \(W=S\otimes_R M\). The first factor supplies the left \(S\)-action, while the second acts on \(M\). Set \(\Delta(g)=1\otimes g-g\otimes1\). The ideal \(J\) is generated by all \(\Delta(g)\): if \(\sum a_i\otimes b_i\) has product sum zero, then it equals \(\sum(a_i\otimes1)\Delta(b_i)\). This proof keeps the original multiplication map and tensor coordinates.

For any \(R\)-linear \(D:M\to N\), define \(L_D:W\to N\) by \(L_D(a\otimes m)=aD(m)\). It exists by \(R\)-balancing. Expanding the product of \(\Delta(g_i)\) gives \[
L_D\bigl(\Delta(g_1)\cdots\Delta(g_r)(1\otimes m)\bigr)
=(\delta_{g_1}\cdots\delta_{g_r}D)(m).
\] Thus the commutator criterion is equivalent to \(L_D\) killing \(J^{k+1}W\). The source uses the negative generators \(g\otimes1-1\otimes g=-\Delta(g)\); its full product has the factor \((-1)^{k+1}\), exactly matching its alternating relation. Every element of \(J^{k+1}W\) is a sum of scalar multiples of these products applied to \(1\otimes m\), so the criterion checks the entire submodule, not merely its named generators.

Consequently the maps \[
P^k_{S/R}(M)\longrightarrow W/J^{k+1}W,\quad j_k(m)\longmapsto[1\otimes m],
\] and \[
W/J^{k+1}W\longrightarrow P^k_{S/R}(M),\quad[a\otimes m]\longmapsto a\,j_k(m)
\] are well defined and inverse on their generating classes. This supplies the omitted inverse verification. Writing \(C_k=C/J^{k+1}\), there is also the exact comparison \[
P^k_{S/R}(M)\cong C_k\otimes_{S,\mathrm{right}}M.
\] Indeed \(C\otimes_{S,\mathrm{right}}M\to W\), \((a\otimes b)\otimes m\mapsto a\otimes bm\), has inverse \(a\otimes m\mapsto(a\otimes1)\otimes m\); passing to the quotient uses right exactness, not flatness.

\subsubsection{The first principal-parts sequence for every module}\label{algebra-proof-043-the-first-principal-parts-sequence-for-every-module}

For order one, an operator \(D:S\to N\) satisfies \(\delta_gD(h)=h\delta_gD(1)\). Thus \(\sigma_D(g)=D(g)-gD(1)\) obeys \(\sigma_D(gh)=g\sigma_D(h)+h\sigma_D(g)\), and is \(R\)-linear with value zero on \(R\). Conversely a derivation plus \(g\mapsto gn\) is an order-one operator. The two constructions give inverse maps \(\operatorname{Diff}^1(S,N)\cong\operatorname{Der}_R(S,N)\oplus N\), proving the source example.

For arbitrary \(M\), the injection in the claimed sequence is \[
i:\Omega_{S/R}\otimes_SM\longrightarrow P^1(M),\qquad
dg\otimes m\longmapsto j_1(gm)-g\,j_1(m).
\] The right-hand side defines a derivation in \(g\) by the order-one commutator law and is \(S\)-linear in \(m\), since \(\delta_gj_1\) is linear. It therefore gives the displayed tensor map. The augmentation \(\epsilon:P^1(M)\to M\) sends \(j_1(m)\) to \(m\). These formulas are precisely the maps constructed via Yoneda in the source.

To prove injectivity without assumptions on \(M\), use the previously proved identification \(J/J^2\cong\Omega_{S/R}\), \(dg\mapsto\Delta(g)\). The short exact sequence of right \(S\)-modules \[
0\longrightarrow J/J^2\longrightarrow C_1\xrightarrow\mu S\longrightarrow0
\] splits by the actual second-factor map \(s\mapsto[1\otimes s]\). Tensoring this split sequence over that right action with \(M\) is exact on the left as well as the right, because a split direct-sum decomposition remains a direct sum after tensoring. The two \(S\)-actions on \(J/J^2\) agree, and the left action on \(C_1\) supplies the indicated actions on the resulting sequence. Hence it becomes exactly \[
0\longrightarrow\Omega_{S/R}\otimes_SM\xrightarrow{i}P^1(M)\xrightarrow{\epsilon}M\longrightarrow0.
\] This proves the source assertion for arbitrary modules. The right-action splitting before tensoring is not a claim that the resulting sequence splits for the left \(S\)-action.

The source free-module argument is also valid without assuming the desired conclusion for \(M'\). A presentation \(0\to M'\to F\to M\to0\) gives right exact columns for \(\Omega\otimes-\) and \(P^1(-)\), the latter from the displayed right-tensor description. Suppose \(x\in\Omega\otimes M\) maps to zero. Lift it to \(y\in\Omega\otimes F\). Its image in \(P^1(F)\) comes from some \(z\in P^1(M')\). The image of \(z\) in \(M'\) maps to zero in \(F\), hence is zero because \(M'\to F\) is injective. Right exactness of the top row gives \(z=i(y')\) for some \(y'\in\Omega\otimes M'\). Injectivity for the free module \(F\), which follows by direct sums from the source example, implies that \(y\) is the image of \(y'\). Therefore \(x=0\). The leading zero drawn in the top row is not needed for this chase; no circular hypothesis is used.

\subsubsection{Functoriality, base change and order-one de Rham maps}\label{algebra-proof-043-functoriality-base-change-and-order-one-de-rham-maps}

For a commutative ring square and a compatible map of modules, the maps \(a\otimes m\mapsto f(a)\otimes h(m)\) carry each \(\Delta(g)\) to \(\Delta(f(g))\). They carry every indicated ideal power into its target ideal power and so induce the source tower maps. Composition is exact on tensors, and the formulas for \(i,\epsilon\) show compatibility with the first principal-parts sequence.

For \(B'=B\otimes_A A'\) and \(M'=M\otimes_A A'\), the base-changed universal map is \(j_k\otimes1\). Its \((k+1)\)-fold commutators vanish: it is enough to test the \(A'\)-algebra generators \(b\otimes1\), for which the commutator is the base change of the original one; the generator criterion is proved below. For \(k=0\) linearity is direct. Every target operator restricted to \(M\) factors uniquely through \(j_k\), and extending the factor by \(u\otimes a'\mapsto a'\gamma(u)\) gives its unique \(B'\)-linear factor. Equality and uniqueness can be checked on \(j_k(m)\otimes a'\), which generate. This proves the source arbitrary-base-change isomorphism without a flatness condition.

The earlier graded Leibniz proof gives the exact de Rham commutator \([d,b](\omega)=db\wedge\omega\). This is \(B\)-linear in \(\omega\), so \(d\) has order one in every degree, including degree zero. Expanding \(d(bb'b_0)-b\,d(b'b_0)-b'\,d(bb_0)+bb'\,d(b_0)\) gives zero: the coefficients of \(db\), \(db'\) and \(db_0\) each cancel. Wedge with the retained \(db_1,\ldots,db_i\) to obtain the source's entire displayed calculation.

For the algebra-generator criterion at 34998--35023 assume \(k\ge1\). The set of \(g\) with \(\delta_gD\) of order \(k-1\) contains the image of \(A\), since those commutators vanish. It is closed under sums and additive inverses. The identity \(\delta_{gg'}D=(\delta_gD)\circ g'+g\circ\delta_{g'}D\), together with composition with order-zero multiplication, proves closure under products. It is therefore an \(A\)-subalgebra containing the named generators, hence all of \(B\). At order zero one instead checks that all generator commutators vanish; the same identities prove linearity directly.

For the tensor lemma at 35088--35109, if \(k=0\), then \(D\) is \(A\)-linear, and for every elementary scalar \(a\otimes b\), \[
(D\otimes1)((a\otimes b)(m\otimes n))=D(am)\otimes bn
=aD(m)\otimes bn=(a\otimes b)(D\otimes1)(m\otimes n).
\] Sums of elementary tensors give full \(A\otimes_RB\)-linearity. For \(k\ge1\), \(D\otimes1_N\) is \(B\)-linear and its commutator with \(a\otimes1\) equals \((\delta_aD)\otimes1_N\); induction and the generator criterion give the stated order. No flatness of \(N\) enters.

\subsubsection{Localization and the complete fraction comparison}\label{algebra-proof-043-localization-and-the-complete-fraction-comparison}

Use the original notation \(A\to B\), multiplicative set \(T\subset B\), and modules \(M,N\), so that the source's unnamed \(R\) at 35057--35058 is correctly read as \(A\). Set \(B_T=T^{-1}B\). In the induction, \(L_b=\delta_bD\) has order \(k-1\), with unique extension \(E_b\). The original identities give \(E_{b'b}=E_{b'}\circ b+b'\circ E_b\) and \([E_{b'},b]=[E_b,b']\), by the induction's uniqueness. Define \[
E(m/g)=g^{-1}\bigl(D(m)-E_g(m/g)\bigr).
\] For a common rescaling \(m/g=g'm/(g'g)\), substitute \(D(g'm)=g'D(m)+L_{g'}(m)\) and \[
E_{g'g}(g'm/(gg'))=E_{g'}(m)+g'E_g(m/g).
\] The second identity follows by applying \(E_{g'g}=E_{g'}\circ g+g'\circ E_g\) to the actual element \(g'm/(gg')=m/g\). Since \(E_{g'}\) extends \(L_{g'}\), the two extra terms cancel and the original fraction formula is unchanged. If \(m/g=m'/g'\) only after a torsion factor, choose \(u\in T\) with \(u(g'm-gm')=0\). Rescaling both fractions to denominator \(ugg'\) gives equal numerators \(ug'm=ugm'\). The just-proved rescaling invariance therefore establishes full representative independence even with torsion and zero divisors.

The formula is additive by choosing a common denominator and using additivity of \(D,E_g\); multiplication by an element of \(A\) commutes by their \(A\)-linearity. Taking \(g=1\), with \(E_1=0\), shows that \(E\) extends \(D\). For \(b\in B\), use \([E_g,b]=[E_b,g]\) to obtain exactly \[
E(bm/g)-bE(m/g)
=g^{-1}\bigl(L_b(m)-E_b(m)+gE_b(m/g)\bigr)=E_b(m/g).
\] For \(g'\in T\), the product identity for \(E_{g'g}\) gives \[
E(m/(g'g))-(g')^{-1}E(m/g)=-(g')^{-1}E_{g'}(m/(g'g)).
\] These are compositions of an order-\((k-1)\) operator with order-zero maps. Since \(B_T\) is generated over \(A\) by the images of \(B\) and the inverses of \(T\), the generator criterion proves that \(E\) has order \(k\). If \(E^*\) were another extension of that order, its commutator with \(g\) would extend \(L_g\), hence equal \(E_g\) by induction. The equation \(E^*(m)=gE^*(m/g)+E_g(m/g)\) forces the same displayed formula, proving uniqueness. The order-zero case is ordinary localization of a linear map. If \(0\in T\), both localized modules are zero, so the unique zero extension satisfies the same conclusion.

\subsubsection{Symbols and the exact connection criterion}\label{algebra-proof-043-symbols-and-the-exact-connection-criterion}

Editorial consequence retaining the source principal parts and the specified left action. For \(k\ge1\), the tower has the exact sequence \[
0\longrightarrow J^kW/J^{k+1}W\longrightarrow P^k(M)\longrightarrow P^{k-1}(M)\longrightarrow0.
\] Products of the original \(\Delta(g)\) give a surjection \[
\operatorname{Sym}^k_S(\Omega_{S/R})\otimes_SM\twoheadrightarrow J^kW/J^{k+1}W,
\quad dg_1\cdots dg_k\otimes m\longmapsto[\Delta(g_1)\cdots\Delta(g_k)(1\otimes m)].
\] To prove that it is defined on differentials, additivity and vanishing on \(R\) follow for each \(\Delta\), and the exact product identity \[
\Delta(ab)=(a\otimes1)\Delta(b)+(b\otimes1)\Delta(a)+\Delta(a)\Delta(b)
\] shows Leibniz after multiplication by \(k-1\) other \(\Delta\)'s and reduction modulo \(J^{k+1}\). The factors commute, giving symmetry. Moving a scalar from \(M\) to the first-factor coefficient changes the expression by one more \(\Delta\), hence by zero in this quotient. Generation of \(J\) proves surjectivity. Injectivity is not asserted.

An operator \(D\) induces its symbol \[
dg_1\cdots dg_k\otimes m\longmapsto(\delta_{g_1}\cdots\delta_{g_k}D)(m).
\] It is the composite of the last surjection with \(L_D\) on the kernel of the tower map. It vanishes exactly when \(L_D\) factors through \(P^{k-1}(M)\), equivalently when \(D\) has order \(k-1\). Thus there is an injection \[
\operatorname{Diff}^k(M,N)/\operatorname{Diff}^{k-1}(M,N)
\hookrightarrow\operatorname{Hom}_S(\operatorname{Sym}^k\Omega_{S/R}\otimes_SM,N).
\] This construction retains every commutator and its positive \(\Delta\) sign.

At order one, an \(S\)-linear splitting \(s:M\to P^1(M)\) of \(\epsilon\) is equivalent to an \(R\)-linear connection \(\nabla:M\to\Omega_{S/R}\otimes_SM\) with \[
\nabla(am)=a\nabla(m)+da\otimes m.
\] Given the splitting, define \(\nabla=i^{-1}(j_1-s)\), using the proved injectivity and exactness. The identity \(j_1(am)-a j_1(m)=i(da\otimes m)\) proves the product rule. Conversely \(s=j_1-i\nabla\) is \(S\)-linear and \(\epsilon s=1\); these constructions are inverse. If one connection exists, all are obtained by adding an arbitrary \(S\)-linear map \(M\to\Omega\otimes M\), since subtraction cancels the displayed product-rule defect. This is a torsor statement only for the nonempty set of connections; no flatness or integrability of a connection is claimed.

For a concrete failure take a nonzero field \(k\), \(S=k[x]\), \(R=k\), and \(M=N=S/(x)\). Every \(k\)-linear map \(k\to k\) is \(S\)-linear because the \(S\)-action is evaluation at zero. Thus \(\operatorname{Diff}^r(M,M)=k\) for every integer \(r\ge0\), with every positive symbol zero, whereas \(\operatorname{Hom}_S(\operatorname{Sym}^r\Omega\otimes M,M)=k\) for every integer \(r\ge1\). The symbol injection is therefore not generally onto. Nor can a connection exist: applying its product rule to \(x\cdot1=0\) gives \(0=dx\otimes1\), a nonzero generator of \(\Omega\otimes M\). Explicitly \(P^1(M)=k[x]/(x^2)\): in the diagonal presentation the second coordinate is zero and \(\Delta(x)=-x\). The injection sends \(dx\otimes1\) to \(-x\), the universal class of \(1\) is \(1\), and the augmentation is reduction modulo \(x\). A section would lift \(1\) to \(1+ax\), whose product with \(x\) is the nonzero class \(x\). This proves the nonsplitting with the exact maps and signs.

\subsubsection{Localized principal parts and finite denominator formulas}\label{algebra-proof-043-localized-principal-parts-and-finite-denominator-formulas}

Editorial consequence receiving the diagonal comparison and source localization theorem. Return to \(R\to S\), \(T\subset S\), and \(C_k=(S\otimes_RS)/J^{k+1}\). After inverting the first-factor elements \(g_L=g\otimes1\) for \(g\in T\), the second-factor element is \(g_R=g_L+\Delta(g)\). Since \(\Delta(g)^{k+1}=0\) in \(C_k\), it has the exact inverse \[
g_R^{-1}=\sum_{i=0}^k(-1)^i g_L^{-i-1}\Delta(g)^i.
\] Multiplication by \(g_L+\Delta(g)\) cancels every intermediate term and leaves \(1+(-1)^kg_L^{-k-1}\Delta(g)^{k+1}=1\). The zero-ring localization causes no exception. Thus localizing the first factor already inverts the required second-factor elements.

Localizing both factors of \(S\otimes_RS\) gives \(S_T\otimes_RS_T\). Localization is exact, so the kernel of its multiplication map is the extended ideal \(J\): applying the two localizations to the exact multiplication sequence gives that kernel and target \(S_T\). The image of \(J^{k+1}\) is the corresponding power, by expansion of products of localized generators. Therefore the preceding finite inverse gives the canonical isomorphism \[
S_T\otimes_{S,\mathrm{left}}P^k_{S/R}(M)
\cong P^k_{S_T/R}(M_T),
\] with \([a\otimes m]\) carried to the same localized class. A reverse map sends \(1\otimes(m/g)\) to \(g_R^{-1}(1\otimes m)\); the displayed inverse makes it defined on localization relations and the diagonal ideal power. Composites fix these generators, proving the isomorphism for arbitrary \(M\).

Apply the localized \(L_D\) to that formula. The unique extension of an order-\(k\) operator is \[
E(m/g)=\sum_{i=0}^k(-1)^i g^{-i-1}(\delta_g^iD)(m).
\] Every term has the original denominator power and iterated commutator. This agrees with the source's recursive formula because the extension of \(\delta_gD\) has the same formula with order \(k-1\); substituting it yields exactly the terms indexed \(1\) through \(k\).

For an integer \(n\ge1\), multiplication of \(n\) copies of the finite inverse gives \[
E(m/g^n)=\sum_{i=0}^k(-1)^i\binom{n+i-1}{i}g^{-n-i}(\delta_g^iD)(m).
\] Indeed for a fixed \(i\le k\), the coefficient counts ordered \(n\)-tuples of nonnegative exponents with sum \(i\); inserting \(n-1\) separators among \(i\) entries gives \(\binom{n+i-1}{i}\). Terms of total exponent above \(k\) vanish by the ideal power. These are integer coefficients, not divisions by factorials in the ring. The formula thus holds in all characteristics with zero divisors and torsion retained. It supplies a finite explicit denominator expansion of the original extension, not a modified operator.

\subsubsection{Polynomial principal parts with all characteristic coefficients}\label{algebra-proof-043-polynomial-principal-parts-with-all-characteristic-coefficients}

Editorial consequence for \(S=R[x_1,\ldots,x_n]\). Let \(\xi_i=1\otimes x_i-x_i\otimes1\). The two inverse \(S\)-algebra maps identify the original diagonal tensor algebra with \(S[\xi_1,\ldots,\xi_n]\): send the second coordinate \(1\otimes x_i\) to \(x_i+\xi_i\), and send \(\xi_i\) back to the displayed tensor difference. Multiplication sets every \(\xi_i\) to zero, so \(J=(\xi_1,\ldots,\xi_n)\). Consequently \[
P^k_{S/R}(S)\cong S[\xi_1,\ldots,\xi_n]/(\xi_1,\ldots,\xi_n)^{k+1}
\] is free for the left \(S\)-action on the explicitly retained monomials \(\xi^\alpha\), \(|\alpha|\le k\). The quotient has precisely those coefficients because its ideal is spanned by monomials of higher total degree. Their number is \(\binom{n+k}{k}\), counted by introducing the extra nonnegative coordinate \(k-|\alpha|\). For \(n=0\) only the empty monomial remains.

For \(f=\sum_\nu a_\nu x^\nu\) with finite support, define \[
\partial^{[\alpha]}f
=\sum_{\nu\ge\alpha}a_\nu\left(\prod_{j=1}^n\binom{\nu_j}{\alpha_j}\right)x^{\nu-\alpha}.
\] Binomial expansion, with every integer coefficient mapped into \(R\), gives the exact universal operator \[
j_k(f)=[f(x+\xi)]=\sum_{|\alpha|\le k}(\partial^{[\alpha]}f)\xi^\alpha.
\] All omitted terms are zero in the specified quotient because their total \(\xi\)-degree exceeds \(k\); no coefficient of a retained term is suppressed. For every \(S\)-module \(N\), the universal property therefore gives the explicit bijection \[
\bigoplus_{|\alpha|\le k}N\longrightarrow\operatorname{Diff}^k_{S/R}(S,N),
\qquad(n_\alpha)_\alpha\longmapsto\left[f\mapsto\sum_{|\alpha|\le k}(\partial^{[\alpha]}f)n_\alpha\right].
\] The inverse takes the associated linear map \(P^k(S)\to N\) on each basis monomial \(\xi^\alpha\). Equivalently its coefficient is \((\delta_{x_1}^{\alpha_1}\cdots\delta_{x_n}^{\alpha_n}D)(1)\), by the original diagonal calculation. The order drops to \(k-1\), for \(k\ge1\), exactly when all coefficients of total degree \(k\) vanish, because those basis monomials span the kernel of the tower projection.

The comparison with the ordinary iterated partial derivative is \[
\partial_1^{\alpha_1}\cdots\partial_n^{\alpha_n}
=\left(\prod_j\alpha_j!\right)\partial^{[\alpha]}.
\] On each monomial it follows from the exact identity \(\nu_j(\nu_j-1)\cdots(\nu_j-\alpha_j+1)=\alpha_j!\binom{\nu_j}{\alpha_j}\), and then by additivity. The multiplier is retained; it is not inverted. In particular over \(\mathbf F_p[x]\), \(\partial^{[p]}(x^p)=1\), while \(\partial^p=p!\partial^{[p]}=0\). Thus ordinary derivatives alone do not recover these original finite-order operators in positive characteristic.

\subsubsection{Propagation and continuation}\label{algebra-proof-043-propagation-and-continuation}

The corrected universal factorization determines the diagonal inverse and every functorial map. The diagonal description proves the first principal-parts sequence without flatness, supplies its exact symbol and connection maps, and gives the localized comparison with finite denominator expansions. The full commutator expansion proves the source free presentation and every sign in the denominator formula. The polynomial comparison retains all binomial and factorial coefficients, including their vanishing in positive characteristic. The previous de Rham graded product rule now receives an exact principal symbol: \(db\otimes\omega\mapsto db\wedge\omega\).

The three consequence groups have no translation operations. Full proofs and extensions remain editorial. The failed symbol-surjectivity and connection examples are exact boundaries, not a claim that the structures are unrelated. Next is the naive cotangent complex at 35116; its read lookahead remains to be adjudicated. Broader synthesis follows core completion. No source or translation mutation, TeX build, publication, new task or delegation occurred.

\hypertarget{algebra_naive_cotangent_notes_20260928}{}
\subsection{Naive cotangent complexes: source review and separate editorial arguments}\label{algebra-proof-044-naive-cotangent-complexes-source-review-and-separate-editorial-arguments}

Authority: Stacks Project authors, commit \texttt{a04446e57ec1fbc252a871afcec7752fb2807b14}, \texttt{algebra.tex:35116–35804}. Original source: https://github.com/stacks/stacks-project/blob/a04446e57ec1fbc252a871afcec7752fb2807b14/algebra.tex\#L35116 . All statements below retain the original rings, presentation ideals, tensor actions and degree conventions. These are supporting editorial arguments, not replacement translation paragraphs or claims of novelty. Proposed local source fixes are recorded separately. A corrected translation must retain the source statement and its place in the original exposition.

Reading includes the entire section and reports OCC-00732 through OCC-00770. Bounded dependencies are \texttt{cotangent.tex:80–178,1830–2017}, \texttt{homology.tex:3605–3722}, \texttt{more-algebra.tex:8439–8555} and \texttt{simplicial.tex:6728–6775,6813–6849,6875–6991}. Additional read material in Cotangent 2019--2090 and More on Algebra 8556--8595 is lookahead, not an adjudicated chapter. A requested Simplicial read beyond its last line raised an index error after displaying the complete relevant examples; actual coverage, rather than the requested endpoint, is recorded. Two disk-index hits for ``cotangent complex Tor base change'' remain unread routing records. No PDF was used.

\subsubsection{Presentations and the exact homotopy maps}\label{algebra-proof-044-presentations-and-the-exact-homotopy-maps}

For a ring map \(R\to S\), put \(P=R[S]\), with variable \([s]\) for every element of \(S\), and augmentation \(\epsilon([s])=s\). Its monomial basis includes the empty monomial. The ideal generated by \[
[s+t]-[s]-[t],\qquad [s][t]-[st],\qquad [r]-r
\] is exactly \(I=\ker\epsilon\): modulo these relations the map \(s\mapsto[s]\) is a unital ring map inverse to the induced augmentation. In the last expression \([r]\) denotes the variable of the image of \(r\) in \(S\). This also handles a noninjective base map and the zero ring.

For any polynomial presentation \(\alpha:P\twoheadrightarrow S\), with kernel \(I\), write \[
N(\alpha)_1=I/I^2,\quad N(\alpha)_0=\Omega_{P/R}\otimes_P S,
\qquad d(\overline i)=di\otimes1.
\] The product rule makes this well defined on \(I/I^2\). A map from its cokernel to an \(S\)-module \(M\) is exactly an \(R\)-derivation \(P\to M\) vanishing on \(I\), hence exactly an \(R\)-derivation \(S\to M\). Thus its cokernel is the original \(\Omega_{S/R}\), with its actual universal map. The degree-zero term is the free \(S\)-module on the differentials of the polynomial variables.

For the source square \(R\to R'\), \(\phi:S\to S'\), and a presentation lift \(\varphi:P\to P'\), the two maps are \[
\varphi_1(\overline i)=\overline{\varphi(i)},\qquad
\varphi_0(dp\otimes s)=d\varphi(p)\otimes\phi(s).
\] They are \(S\)-linear using \(\phi\) on the target and commute with \(d\). Each polynomial variable may be lifted through the surjection \(P'\to S'\), proving existence of \(\varphi\). Composition of these two formulas is exactly the formula for the composite lift.

For two lifts \(\varphi,\varphi'\), their difference takes values in \(I'\). The induced map \(D:P\to I'/(I')^2\) is an \(R\)-derivation because \[
(\varphi-\varphi')(pq)=\varphi(p)(\varphi-\varphi')(q)
+(\varphi-\varphi')(p)\varphi'(q),
\] and the two actions coincide modulo \((I')^2\). That action factors through \(S\). Therefore \[
h(dp\otimes s)=\phi(s)\,\overline{\varphi(p)-\varphi'(p)}
\] is a well-defined \(S\)-linear map \(N(\alpha)_0\to N(\alpha')_1\). On \(\overline i\) and \(dp\otimes s\), respectively, the formulas give \[
hd=\varphi_1-\varphi'_1,\qquad dh=\varphi_0-\varphi'_0.
\] If both base and target maps are isomorphisms, lift the reverse square and apply these identities to both composites and the identity lifts. This proves a homotopy equivalence, without any projectivity assertion about the conormal module. The reference to part (3) at line 35384 also needs part (2), which supplies these homotopies.

Taking the identity presentation of a polynomial algebra gives \((0\to\Omega_{S/R})\). Taking the zero-variable presentation of a surjection gives \((I/I^2\to0)\). The comparisons just proved are homotopy equivalences, not literal identifications of the canonical presentation complexes. In particular, for a nonzero ring \(T\), the canonical complex for \(T\to T\) has degree-zero term \(\bigoplus_{t\in T}T\,d[t]\), which is nonzero, although the complex is contractible. This justifies clarifying ``zero'' to ``homotopy equivalent to zero'' at line 35586 while retaining the original argument.

\subsubsection{The original truncation really is isomorphic}\label{algebra-proof-044-the-original-truncation-really-is-isomorphic}

The standard resolution uses \(P_0=R[S]\), \(P_1=R[P_0]\), \(P_2=R[P_1]\), with the original, unnormalized associated chain complex \(L_n=\Omega_{P_n/R}\otimes_{P_n}S\). We do not replace this with a normalized complex. On a variable \([q]\in P_1\), the faces are \[
d_0([q])=q,\qquad d_1([q])=[\epsilon(q)].
\] On \([F]\in P_2\), they are \(F,[d_0F],[d_1F]\). These are the maps in Simplicial, \texttt{example-polynomial-algebra-maps}, lines 6875--6906. Let \(Q=\operatorname{coker}(L_2\to L_1)\), and denote the image in \(Q\) of \(d[q]\otimes1\) by \(e_q\). The exact relations defining \(Q\), one for each polynomial \(F\in R[P_0]\), are \[
\sum_q \epsilon_1\!\left(\frac{\partial F}{\partial[q]}\right)e_q
-e_{d_0F}+e_{d_1F}=0.\tag{1}
\] Every sum is finite and includes its original integer polynomial coefficients. Here \(\epsilon_1([q])=\epsilon(q)\).

Define \(u:L_1\to I/I^2\) by \[
u(d[q]\otimes1)=\overline{q-[\epsilon(q)]}.
\] More generally \(u(dF\otimes s)=s\overline{d_0F-d_1F}\). This is well defined because \(d_0-d_1\) is a derivation modulo \(I^2\), for the common action through \(S\), by the same product calculation as above. The simplicial identities give \((d_0-d_1)(d_0-d_1+d_2)=0\); explicitly the terms pair as \(d_0d_0=d_0d_1\), \(d_1d_0=d_0d_2\), and \(d_1d_1=d_1d_2\). Hence \(u\) kills all the relations (1) and descends to \(Q\). Its differential is exactly the degree-one differential into \(L_0\).

Here is an inverse, with no derived-category inference. First (1) with \(F=[[s]]\) gives \(e_{[s]}=0\). If \(i\in I\) and \(q\in P_0\), use \(F=[q+i]-[q]\). Its two faces are \(i\) and \(0\), so \[
e_{q+i}-e_q=e_i-e_0.\tag{2}
\] For \(i\in I,q\in P_0\), use \(F=[q][i]-[q][0]\). Its two faces are \(qi\) and \(0\), and the coefficients of \(e_q\) cancel. Thus \[
e_{qi}-e_0=\epsilon(q)(e_i-e_0).\tag{3}
\] Set \(v(i)=e_i-e_0\). Equation (2) makes \(v\) additive, and (3) makes it \(P_0\)-linear for the action through \(S\). It annihilates \(I^2\), and hence induces \(v:I/I^2\to Q\). Directly, \[
u(v(i))=\overline{i-[0]}-\overline{0-[0]}=\overline i.
\] Applying (2) with \(q=[\epsilon(q_0)]\), \(i=q_0-[\epsilon(q_0)]\), gives \[
v(u(e_{q_0}))=e_{q_0}-e_{[\epsilon(q_0)]}=e_{q_0}.
\] Thus \(u\) and \(v\) are inverse isomorphisms. Together with the identity on \(L_0\), they give the actual canonical isomorphism \[
\tau_{\leq1}L_{S/R}\ \xrightarrow{\ \cong\ }\ N(R[S]\to S).
\] This is the chain truncation of Homology 3694--3722: its degree-one term is precisely \(Q\). Cotangent 129--133 uses cochain degree \(-n\) for this same term, so its \(\tau_{\geq-1}\) is the same construction with that stated indexing. The weaker wording ``quasi-isomorphic'' in Cotangent 1937--1975 does not refute the stronger statement in Algebra 35170. No correction to that stronger statement is warranted. The explicit map agrees with Cotangent 1977--2017, and the proof here establishes its low-degree isomorphism directly, without asserting that all the sheaf-derived arguments in that chapter have been checked.

There is a different qualification needed at Algebra 35161. The standard augmentation is a homotopy equivalence of underlying simplicial sets, as Simplicial 6908--6991 explicitly states; it need not be a homotopy equivalence of simplicial \(R\)-algebras. For \(R=\mathbf Z\), \(S=\mathbf F_p\), a reverse simplicial algebra map would in degree zero give a unital map \(\mathbf F_p\to\mathbf Z[\mathbf F_p]\). This is impossible because \(p\cdot1\ne0\) in the polynomial ring. The actual set section sends \(s\) to the iterated bracket \([\ldots[s]\ldots]\). The extra degeneracy \(t_n(q)=[q]\) satisfies \(d_0t_n=\mathrm{id}\), \(d_{i+1}t_n=t_{n-1}d_i\), and \(s_{i+1}t_n=t_{n+1}s_i\), with the augmentation in degree \(-1\); these follow by deleting or inserting the indicated bracket. They contract the augmented underlying simplicial sets as in the cited construction. Adding ``on underlying simplicial sets'' retains the precise true claim.

\subsubsection{Jacobi-Zariski and the tensor comparison sequence}\label{algebra-proof-044-jacobi-zariski-and-the-tensor-comparison-sequence}

Keep the source presentations \(P=A[x_s]\twoheadrightarrow B\), \(Q=B[y_t]\twoheadrightarrow C\), and \(T=A[x_s,y_t]\twoheadrightarrow C\), with kernels \(I,J,K\). Write \(W=IT\). Then \(T/W=Q\), \(J=K/W\), and \[
J/J^2=K/(K^2+W),\qquad
(I/I^2)\otimes_B C\cong W/KW.
\] For the latter formula, use \(T=P[y_t]\), so \(W=I\otimes_P T\), and then tensor over \(T\) with \(C=T/K\). Since \(I^2T\subset KW\), this gives the stated quotient. Thus the bottom row in the source is right exact, with map \(W/KW\to K/K^2\) induced by inclusion. No injectivity of that map is asserted. The top row is the split sequence of free \(C\)-modules on \(dx_s\), on \(dx_s,dy_t\), and on \(dy_t\), with their indicated polynomial tensor bases. Differentiation commutes with these maps.

For a cycle \(\overline j\in J/J^2\), choose a lift \(\overline k\in K/K^2\). The \(dy_t\) part of \(dk\) is zero, so \(dk\) comes from a unique vector \(v\in\bigoplus_s Cdx_s\). Send the cycle to the class of \(v\) in \(\Omega_{B/A}\otimes_B C\). Changing the lift by the image of \(W/KW\) changes \(v\) by the differential of that image's preimage; hence this boundary is well defined. A zero boundary allows subtraction of such a preimage from \(\overline k\), making \(\overline k\) a cycle. Conversely a lifted cycle has zero boundary. A cycle in \(K/K^2\) mapping to zero in \(J/J^2\) lifts from \(W/KW\), and injectivity of the top-left arrow forces that lift's differential to be zero. A vector from \(\bigoplus_s Cdx_s\) becomes a differential in the middle precisely when it arises from this boundary construction. Finally the right-exact cokernel sequence follows by lifting a differential on the right to the middle. These statements prove every exactness position of the displayed six-term sequence. Comparing presentations with the already proved homotopies supplies the canonical homology groups in the source.

The tensor proviso has an exact refinement. For any two-term complex \(N_1\xrightarrow d F\), where \(F\) is a free \(B\)-module, set \(H=\ker d\), \(M=\operatorname{im}d\), and \(\Omega=\operatorname{coker}d\). Tensor the exact sequences \(0\to M\to F\to\Omega\to0\) and \(0\to H\to N_1\to M\to0\) with a \(B\)-module \(C\). The first yields \[
\operatorname{Tor}_1^B(M,C)\cong\operatorname{Tor}_2^B(\Omega,C),
\quad
0\to\operatorname{Tor}_1^B(\Omega,C)\to M\otimes_B C\to F\otimes_B C.
\] Let \(\delta\) be the tensor connecting map from \(\operatorname{Tor}_1(M,C)\) to \(H\otimes C\). The surjection \(N_1\otimes C\to M\otimes C\) maps \(\ker(d\otimes1)\) onto \(\ker(M\otimes C\to F\otimes C)\). Its kernel is the image of \(H\otimes C\); the kernel of the map out of \(H\otimes C\) is \(\operatorname{im}\delta\). Consequently the natural exact sequence is \[
\operatorname{Tor}_2^B(\Omega,C)\xrightarrow\delta H\otimes_B C
\longrightarrow H_1(N\otimes_B C)
\longrightarrow\operatorname{Tor}_1^B(\Omega,C)\longrightarrow0.\tag{4}
\] In particular the comparison is an isomorphism exactly when \(\operatorname{Tor}_1(\Omega,C)=0\) and the particular map \(\delta\) is zero. Vanishing of the whole \(\operatorname{Tor}_2\) module, as assumed in the source, is sufficient. Formula (4) supplies the exact obstruction rather than silently weakening the source theorem's hypotheses in a translation.

For the source nullhomotopy, let \(\phi:B\to C\), \(P=A[B]\), \(Q=B[C]\), \(J=\ker(Q\to C)\). Compare \(f:P\to Q\), \([b]\mapsto[\phi(b)]\), with \(e:P\to Q\), \([b]\mapsto b\) as a coefficient. Their difference modulo \(J^2\) is a derivation, giving \[
h(d[b]\otimes b')=\phi(b')\,\overline{[\phi(b)]-b}.
\] Here the scalar \(b'\) acts via \(B\to C\). On \(I=\ker(P\to B)\), \(e(I)=0\), so \(hd=f_1\). Relative to \(B\), \(de(p)=0\), so \(dh=f_0\). This proves the composite is nullhomotopic. For a surjection \(A\to C\), applying the six-term sequence and the zero-variable presentations gives exactly \(I/I^2\to J/J^2\to\Omega_{B/A}\otimes_B C\to0\), retaining the original ideals.

\subsubsection{Base change with its full obstruction module}\label{algebra-proof-044-base-change-with-its-full-obstruction-module}

Let \(R\to S\), \(P\twoheadrightarrow S\) with kernel \(I\), and an arbitrary ring map \(R\to R'\) be given. Put \(P'=P\otimes_R R'\), \(S'=S\otimes_R R'\), and \(I'=\ker(P'\to S')\). Since \(P\) is free over \(R\), tensoring \(0\to I\to P\to S\to0\) gives \[
0\to T:=\operatorname{Tor}_1^R(S,R')
\longrightarrow I\otimes_R R'\longrightarrow I'\longrightarrow0.
\] The image of \(I^2\otimes_R R'\) in \(I'\) is exactly \((I')^2\): products of sums of images of tensors are sums of such products, and coefficients in \(R'\) can be put in either tensor factor. Taking quotients therefore gives \[
T\longrightarrow (I/I^2)\otimes_R R'\longrightarrow I'/(I')^2\longrightarrow0.
\] Define \(O\) to be the image of the first map, with its induced \(S'\)-module structure. It is the exact kernel of the second map, not necessarily all of \(T\). The universal derivation identifies the degree-zero terms: \[
(\Omega_{P/R}\otimes_P S)\otimes_R R'
\cong\Omega_{P'/R'}\otimes_{P'}S'.
\] For polynomial \(P\), both sides are the free \(S'\)-module on the original differentials, and the map is their identity. An element of \(O\) has zero differential, since its image under this degree-zero isomorphism is the differential of zero. Thus the original comparison fits in a short exact sequence of complexes \[
0\to (O\to0)\to N(\alpha)\otimes_R R'\to N(\alpha')\to0.
\] Its full nontrivial homology assertion is \[
0\to O\to H_1(N(\alpha)\otimes_R R')
\to H_1(N(\alpha'))\to0,
\qquad H_0(N(\alpha)\otimes_R R')\cong H_0(N(\alpha')).\tag{5}
\] Hence the comparison is a quasi-isomorphism exactly when \(O=0\). Since its degree-zero map is an isomorphism and its degree-one map is surjective with kernel \(O\), this is also exactly when it is an isomorphism of these chosen presentation complexes. It suffices that \(\operatorname{Tor}_1^R(S,R')=0\); flatness of \(R'\) over \(R\) is unnecessary for that implication. Comparison with the canonical presentations gives the source's homotopy equivalence in this more general Tor-vanishing case. Tensoring an existing homotopy preserves its equations, without flatness.

For a strict extension of the flat-base hypothesis take \(R=\mathbf Z\), \(S=\mathbf Z[x]\), \(R'=\mathbf F_p\): \(S\) is free over \(R\), so the Tor group is zero, whereas \(R'\) is not flat (multiplication by \(p\) on \(\mathbf Z\) ceases to be injective after tensoring). For failure of unrestricted base change take \(R=\mathbf Z\), \(S=R'=\mathbf F_p\), and the zero-variable presentation \(P=\mathbf Z\), \(I=(p)\). The original complex after tensoring is \((\mathbf F_p\to0)\), while \(P'=S'=\mathbf F_p\) has zero kernel and zero complex. Here \(O=\mathbf F_p\), so the degree-one homology is genuinely lost. These examples verify both the stronger sufficient hypothesis and the necessary obstruction; they do not replace the flat-base theorem in the source.

For a directed system \(R_\lambda\to S_\lambda\), every polynomial and every equality between finite polynomial expressions occurs at some common stage. Hence \(R[S]=\operatorname{colim}R_\lambda[S_\lambda]\), and an element mapping to zero in \(S\) is already zero at a later stage, giving \(I=\operatorname{colim}I_\lambda\). Every element of \(I^2\) is a finite sum of finite products, so it and every relation in the quotient are similarly witnessed at a stage. The same argument applies to the free modules on the elements of \(S_\lambda\). It gives both terms, their scalar actions, and their differentials in the source's asserted colimit identity. Exactness of filtered colimits of modules follows by the same finite-element test, and will be used for localization below.

\subsubsection{Localization and why the reported splitting is valid}\label{algebra-proof-044-localization-and-why-the-reported-splitting-is-valid}

For \(U\subset A\) multiplicative, flat base change to \(U^{-1}A\) compares \(N_{U^{-1}A/A}\) with the canonical identity complex over \(U^{-1}A\); the latter is contractible by the zero-variable presentation. Every term of the former is already a \(U^{-1}A\)-module, so tensoring it with \(U^{-1}A\) is an isomorphism. This proves the source localization complex is contractible. If \(U^{-1}A\to B\), localizing a presentation \(P\to B\) gives the presentation \(U^{-1}P\to B\). Its conormal module is \(U^{-1}(I/I^2)=I/I^2\), as all these scalars already act invertibly through \(B\); the differential module agrees for the same reason, by the polynomial differential formula. This proves the base-localization comparison with the original actions.

For the principal localization retain exactly \(\alpha:P\to B\), \(I\), \(f\in P\) mapping to \(g\in B\), and \(\beta:P[x]\to B_g\), \(x\mapsto g^{-1}\). Its kernel is \(J=(I,fx-1)\), because the indicated quotient is \(B[x]/(gx-1)=B_g\). In \(B[x]\), the polynomial \(gx-1\) is a nonzerodivisor: if \((gx-1)\sum_{i=0}^n b_ix^i=0\), the constant coefficient gives \(b_0=0\), and successive coefficients give \(b_i=gb_{i-1}=0\). Thus its conormal module is free on its class over \(B_g\).

The map \(\pi:P[x]\to(P/I^2)_f\), \(x\mapsto1/f\), sends \(J\) into the square-zero ideal \((I/I^2)_g\), kills \(J^2\), and induces a \(B_g\)-linear retraction of \((I/I^2)_g\to J/J^2\). On a fraction \(\overline i/g^n\) its composite is exactly \(\overline i/g^n\); hence the retraction proves injectivity for all fractions, including torsion. The conormal and differential rows consequently split as modules, and \[
J/J^2=(I/I^2)_g\oplus B_g\overline{fx-1},\quad
\Omega_{P[x]/A}\otimes_{P[x]}B_g
=(\Omega_{P/A}\otimes_P B_g)\oplus B_gdx.
\] Keep the full derivative \[
d(fx-1)=g^{-1}df+g\,dx=g\,(dx+g^{-2}df).
\] The change from \(dx\) to \(e=dx+g^{-2}df\) has inverse \(dx=e-g^{-2}df\), so it gives an isomorphism of the original complex with \[
N(\alpha)\otimes_B B_g\ \oplus\ (B_g\xrightarrow{g}B_g).
\] No \(df\) or denominator is discarded. The last summand contracts by multiplication by \(g^{-1}\) from degree zero to degree one. All formulas remain meaningful if \(B_g\) is the zero ring.

Report OCC-00761 incorrectly interprets ``Such a short exact sequence'' as referring to every sequence of complexes. In context the quotient is the displayed disk \((B_g\xrightarrow g B_g)\), where \(g\) is a unit. Any short exact sequence with that quotient splits in complexes: choose a lift \(v\) of its degree-one generator (the degree-one quotient is free), and define the lift of the degree-zero generator to be \(g^{-1}dv\). Its projection is the required generator and its differential condition is \(dv=g(g^{-1}dv)\). Extending linearly gives a chain section. This proves the general claim for exactly the displayed type of quotient, independently of the explicit splitting. The report is rejected and the true source claim is retained.

For arbitrary multiplicative \(U\subset B\), the rings \(B_g\), \(g\in U\), form a filtered system under divisibility and have colimit \(U^{-1}B\). Divisibility is a directed preorder, sufficient for the filtered-colimit argument; no antisymmetry assumption is required. The canonical comparison maps commute at every stage. Applying the proven principal result and filtered exactness gives the source's quasi-isomorphism. A stronger homotopy assertion is proved in the next section and retained as editorial material.

\subsubsection{Two-term quasi-isomorphisms and stronger localization}\label{algebra-proof-044-two-term-quasi-isomorphisms-and-stronger-localization}

Let \(f:C\to D\) be a quasi-isomorphism of complexes in degrees one and zero, and suppose \(D_0\) is projective. The sequence \[
0\to C_1\xrightarrow{a=(f_1,-d_C)}D_1\oplus C_0
\xrightarrow{b=(d_D,f_0)}D_0\to0\tag{6}
\] is exact. Indeed \(a\) is injective by injectivity of \(H_1(f)\), and \(b\) is surjective by surjectivity of \(H_0(f)\). If \(d_Dv+f_0u=0\), injectivity of \(H_0(f)\) gives \(u=d_Cw\); then \(v+f_1w\) is a cycle. Surjectivity of \(H_1(f)\) supplies \(z\in\ker d_C\) with \(f_1z=v+f_1w\). The element \(z-w\) maps under \(a\) to \((v,u)\), proving exactness in the middle.

Choose a section \((s,t):D_0\to D_1\oplus C_0\) of \(b\), so \(d_Ds+f_0t=1\). Put \(g_0=t\). For \(v\in D_1\), the pair \((v-sd_Dv,-td_Dv)\) is in \(\ker b\); its unique lift under \(a\) defines \(g_1v\). Thus \[
f_1g_1=1-sd_D,\qquad d_Cg_1=td_D,
\] so \(g\) is a chain map and \(fg-1\) has homotopy \(-s\). For \(u\in C_0\), the pair \((-sf_0u,u-tf_0u)\) is in \(\ker b\); its unique lift defines \(hu\). Then \[
f_1h=-sf_0,\qquad d_Ch=tf_0-1.
\] Applying \(a\) to \(g_1f_1-1\) and \(hd_C\) gives the same two components, using the chain identities, so its injectivity gives \(g_1f_1-1=hd_C\). Thus \(gf-1=dh+hd\), and \(f\) is a homotopy equivalence. This proof imposes no projectivity on \(C_1,D_1,C_0\).

Apply this to the canonical arbitrary-localization map of the preceding section: the target degree-zero term is free over \(U^{-1}B\) on all its canonical variables. The established quasi-isomorphism is therefore a homotopy equivalence, even for an arbitrary multiplicative set. This strengthens the conclusion at Algebra 35694--35695. No compatible choices of homotopies at the filtered stages are assumed or needed. Keep this strengthened statement separate from the translation. It also shows how the source's later use of localization up to homotopy in More on Algebra 8585--8592 can be justified without silently treating all quasi-isomorphisms as homotopy equivalences.

\subsubsection{Stable conormal isomorphisms with both products}\label{algebra-proof-044-stable-conormal-isomorphisms-with-both-products}

Retain the source complexes \(A_1\to A_0\), \(B_1\to B_0\), maps \(\varphi:A\to B\), \(\psi:B\to A\), and homotopies \[
1-\psi_1\varphi_1=hd_A,\quad1-\psi_0\varphi_0=d_Ah,
\quad1-\varphi_1\psi_1=h'd_B,\quad1-\varphi_0\psi_0=d_Bh'.
\] The displayed matrices have the exact types \[
T=\begin{pmatrix}\psi_1&h\\-d_B&\varphi_0\end{pmatrix}:
B_1\oplus A_0\to A_1\oplus B_0,
\qquad
U=\begin{pmatrix}\varphi_1&-h'\\d_A&\psi_0\end{pmatrix}:
A_1\oplus B_0\to B_1\oplus A_0.
\] Multiplying all four entries gives \[
TU=\begin{pmatrix}1&E\\0&1\end{pmatrix},\quad E=-\psi_1h'+h\psi_0,
\qquad
UT=\begin{pmatrix}1&F\\0&1\end{pmatrix},\quad F=\varphi_1h-h'\varphi_0.
\tag{7}
\] For example the lower-left entries are \(-d_B\varphi_1+\varphi_0d_A=0\) and \(d_A\psi_1-\psi_0d_B=0\), by the chain identities. The diagonal entries are the four displayed homotopy identities. Both shears are invertible, with off-diagonal entries \(-E\) and \(-F\). Consequently \[
T^{-1}=U(TU)^{-1}=(UT)^{-1}U,\qquad
U^{-1}=(TU)^{-1}T=T(UT)^{-1}.
\] Indeed each first expression is a right inverse and each second expression a left inverse; associativity gives equality of the left and right inverses. This proves the original stable isomorphism for arbitrary modules, not only finite free modules.

At line 35738 the source's \(-h'\circ\psi_1\) is ill typed; the required entry is \(-\psi_1\circ h'\). Report OCC-00766 correctly identifies that the displayed product alone is insufficient. For a concrete boundary, on \(V=\bigoplus_{n\geq0}ke_n\), let \(Se_n=e_{n+1}\), \(Le_0=0\), \(Le_{n+1}=e_n\). Then \(LS=1\), but \(SL(e_0)=0\ne e_0\). Thus neither factor need be invertible merely because one product is. The reverse calculation in (7) repairs the source proof without changing its theorem. The proposed local text only states that the reverse product too has identity diagonal; the complete computations remain here.

For two finite polynomial presentations of \(S/R\), the presentation homotopy equivalence and (7) give exactly \[
I/I^2\oplus S^{\oplus m}\cong J/J^2\oplus S^{\oplus n}.
\] For the source localization \(S_g\), use the proved homotopy equivalence from \(N(\alpha)\otimes_SS_g\) to the adjoined-variable presentation, and then to \(N(\beta)\). Applying the same two-term result gives exactly \[
(I/I^2)_g\oplus S_g^{\oplus m}\cong J/J^2\oplus S_g^{\oplus n}.
\] The reference required at line 35796 is \texttt{lemma-sum-two-terms}; \texttt{lemma-conormal-module} is the earlier statement for two polynomial presentations of the same algebra, whereas the tensor-localized complex in this argument need not itself arise from a polynomial presentation over \(R\). In particular, cancelling a free summand from a different stable isomorphism would be unjustified for arbitrary modules. The direct matrix comparison supplies the asserted result, including zero ranks and the zero ring.

\subsubsection{Propagation and bounded follow-up}\label{algebra-proof-044-propagation-and-bounded-follow-up}

The presentation homotopy and the exact truncation comparison protect the source's genuine claims from proposed weakenings. The disk splitting rejects OCC-00761. The full tensor sequence (4), the obstruction (5), and the two-term inverse construction are separate consequences with zero translation operations. The original flat-base and quasi-isomorphism formulations remain identifiable. Previous conormal and base-change work in \texttt{ALGEBRA\_DIFFERENTIALS\_NOTES\_20260928.md} supplies the receiving comparison; the current arguments retain its actual ideals and module actions.

A source dependency also needs a separate follow-up: More on Algebra 8502 and 8504 write \((I/I^2)\otimes_A C\) where the exact quotient \(IT/IK\) requires \((I/I^2)\otimes_B C\). For \(A=k\), \(B=C=k[x]/(x^2)\), \(P=k[x]\), \(I=K=(x^2)\), and the zero-variable identity presentation \(B\to C\), the left expression with base \(k\) has dimension four over \(k\), whereas \(I/I^2\) and the required quotient have dimension two. The natural multiplication map is onto and has a two-dimensional kernel, so the printed identification is false. This explicit counterexample and the corrected map follow from the quotient calculation above. Record the two exact loci as a pending More on Algebra correction rather than editing an unadjudicated chapter or crediting its entire proof as checked. The cited Koszul-regular intersection theorem was not independently read in this batch.

The truncation isomorphism and the localization splitting need no corrective source operations. Copyediting preferences for hyphens, terse ``Follows from'' proofs, the causal ``being'', and two optional commas are not mathematical defects. The report about ``each \ldots{} are quasi-isomorphisms'' needs both verb and complement corrected to ``is a quasi-isomorphism''. Tensor-base clarifications retain the existing maps. Broader joint synthesis, literature comparison for any novelty claim, and a resulting short paper follow core completion. No chapter mutation, translation rewrite, TeX build, publication, or complete-project claim is made here.

\hypertarget{algebra_lci_notes_20260928}{}
\subsection{Local complete intersections: source review and editorial consequences}\label{algebra-proof-045-local-complete-intersections-source-review-and-editorial-consequences}

Source: Stacks Project authors, \texttt{algebra.tex:35805–36426} at commit \texttt{a04446e57ec1fbc252a871afcec7752fb2807b14}, https://github.com/stacks/stacks-project/blob/a04446e57ec1fbc252a871afcec7752fb2807b14/algebra.tex\#L35805 . This note retains the original presentations, ideals, local rings, dimensions and residue fields. It supports the twenty received reports and records further exact consequences. Complete supplementary arguments remain editorial material; they do not silently replace a translation. No novelty is claimed.

Primary dependencies read: Algebra 16880--17001, 23423--23504, 25051--25075, 25749--25777, 27420--27480, 27781--27852, 28269--28356 and 31623--31660. The last dependency includes the previously adjudicated wrong-factor sentence at 31659: use the already composed correction MC-STK-ERR-0669 selecting the S factor, and the full proof in ALGEBRA\_FINITE\_PRESENTATION\_NOTES\_20260928.md. The preceding cotangent note supplies the corrected exact stable-conormal matrices, not an unjustified cancellation. Four disk-index hits for ``complete intersection conormal'' remain unread routing records. Syntomic 36427--36470 is lookahead, with its first proof incomplete and unadjudicated.

\subsubsection{Original finite-type presentations and local criteria}\label{algebra-proof-045-original-finite-type-presentations-and-local-criteria}

Write the original global presentation as \(S=k[x_1,\ldots,x_n]/(f_1,\ldots,f_c)\), with \(\dim S=n-c\), and first suppose \(S\ne0\). Every maximal ideal \(\mathfrak m'\) of the polynomial ring containing the presentation ideal corresponds to a maximal ideal \(\mathfrak m\) of \(S\). The localized polynomial ring is regular local of dimension \(n\). Successive principal quotients have dimension at least \(n-c\), whereas \(\dim S_{\mathfrak m}\leq\dim S=n-c\); equality follows. In particular the parameter criterion for a Cohen--Macaulay module (Algebra 25051--25075) proves that the original ordered sequence is regular at each such maximal ideal, and that its quotient is Cohen--Macaulay. Localizing these regular sequences proves the same statements at every prime of \(S\). This also proves the Cohen--Macaulay claim for a locally complete-intersection cover.

For equidimensionality, each minimal prime over the ideal has height at most \(c\) by the height theorem for \(c\) generators. Its component dimension is \(n\) minus that height, hence at least \(n-c\); since no component exceeds \(\dim S=n-c\), all have dimension \(n-c\). A nonempty principal open meets at least one component. Every nonempty open in an irreducible finite-type scheme over a field has the full component dimension, by the closed-point and chain argument in Algebra 27781--27852. Thus \(\dim S_g=n-c\) whenever \(S_g\ne0\). For an original polynomial lift \(g'\) of \(g\), keep the actual presentation \[
S_g=k[x_1,\ldots,x_n,x_{n+1}]/(f_1,\ldots,f_c,x_{n+1}g'-1).
\] The number of variables and equations both increase by one. If \(S_g=0\), the source's explicit zero-ring convention applies instead. Localizing every member of a local complete-intersection cover proves its localization assertion.

For an arbitrary presentation \(P=k[x_1,\ldots,x_n]\twoheadrightarrow S\) with kernel \(I\) and primes \(\mathfrak q'\mapsto\mathfrak q\), put \[
R=P_{\mathfrak q'},\quad \overline R=S_{\mathfrak q},\quad
c=\dim R-\dim\overline R=n-\dim_{\mathfrak q}\operatorname{Spec}S.
\] The last equality is the source codimension formula, which retains topological dimension at the point rather than substituting local-ring dimension. Set \(J=IR\subset\mathfrak m_R\). The least number of generators of \(J\) equals that of \(J/J^2\) over \(R/J\): both residue vector spaces are \(J/\mathfrak m_RJ\), since \(J^2\subset\mathfrak m_RJ\). Nakayama applies because \(R\) is Noetherian and \(J\) is finite. If \(c\) elements generate either module, the Cohen--Macaulay parameter criterion shows they form an ordered regular sequence. A regular sequence of length \(c\) makes \(J/J^2\) free of rank \(c\), via the source regular-to-quasi-regular theorem. Conversely freeness of rank \(c\) gives \(c\) generators. Any regular sequence generating \(J\) has length \(\dim R-\dim R/J=c\), by successive dimension drops in the Cohen--Macaulay ring. This proves conditions (2)--(5) and the assertion about any residue-basis lifts in the source lemma.

To obtain the actual global neighbourhood, write each selected element of \(IR\) as \(a_i/d_i\), with \(a_i\in I\) and \(d_i\notin\mathfrak q'\). Multiplication by these units retains both generation and regularity locally. The finite module \(I/(a_1,\ldots,a_c)\) vanishes at \(\mathfrak q'\), so a product of finitely many annihilating denominators gives \(g'\notin\mathfrak q'\) with \(I_{g'}=(a_1,\ldots,a_c)_{g'}\). Choose \(g''\notin\mathfrak q'\) excluding every irreducible component not passing through \(\mathfrak q\). Then \(\dim S_{g''}=\dim_{\mathfrak q}\operatorname{Spec}S=n-c\). Each component through \(\mathfrak q\) survives further localization at \(g'\), so the same dimension holds for \(S_g\), where \(g\) is the image of \(g'g''\). The exact presentation is \[
S_g=P[x_{n+1}]/(a_1,\ldots,a_c,x_{n+1}g'g''-1),
\] with \(n+1\) variables and \(c+1\) equations. No denominator has been suppressed.

Conversely, take the source global presentation \(S_g=k[y_1,\ldots,y_m]/J\), with \(t=m-\dim S_g\) equations. At the prime in question its conormal is free of rank \(t\), by the local parameter argument. The corrected stable-conormal comparison from the preceding cotangent note gives \[
(J/J^2)_{\mathfrak q}\oplus S_{\mathfrak q}^{\oplus n}
\cong (I/I^2)_{\mathfrak q}\oplus S_{\mathfrak q}^{\oplus m}.
\] It follows that the finite module \((I/I^2)_{\mathfrak q}\) is projective. A finite projective module over a local ring is free: lift a basis of its residue vector space to a surjection from a finite free module, split the surjection, and use Nakayama on the finite complementary summand to show it is zero. Comparing residue ranks in the displayed isomorphism gives rank \(t+n-m=n-\dim S_g=c\). This proves condition (1) implies (4) without free-summand cancellation over an arbitrary ring.

\subsubsection{Changing regular presentations and the quantitative defect}\label{algebra-proof-045-changing-regular-presentations-and-the-quantitative-defect}

Keep the source surjections \(A\twoheadrightarrow B\twoheadrightarrow C\) of local rings with \(A,B\) regular. Write \(\mathfrak m=\mathfrak m_A\), \(J=\ker(A\to B)\), \(I=\ker(A\to C)\), and \(r=\dim A-\dim B\). The original regular-quotient lemma (25749--25777) gives parameters \(x_1,\ldots,x_d\) of \(A\) with \(J=(x_1,\ldots,x_r)\). In particular \(J/\mathfrak mJ\to\mathfrak m/\mathfrak m^2\) is injective. Since \(I\subset\mathfrak m\), this implies \(J\cap\mathfrak mI=\mathfrak mJ\): one inclusion is immediate, and the other follows from \(\mathfrak mI\subset\mathfrak m^2\) and that injectivity. Thus there is an exact sequence of residue-field vector spaces \[
0\to J/\mathfrak mJ\to I/\mathfrak mI
\to (I/J)/\mathfrak m_B(I/J)\to0.
\tag{1}
\] Consequently \[
\mu_A(I)=r+\mu_B(I/J),\quad
\mu_A(I)-(\dim A-\dim C)
=\mu_B(I/J)-(\dim B-\dim C).\tag{2}
\] This proves the equivalence of the two generator counts in the source. Each is equivalent to generation by a regular sequence by the Cohen--Macaulay parameter criterion. It also proves the source induction step exactly: a nonzero residue class from \(J\) extends to a minimal generating set of \(I\), and is part of the regular system of parameters of \(A\); quotienting by it decreases both relevant counts by one. The zero-dimensional difference gives \(J=0\), which is the induction base.

For a finite-type \(k\)-algebra define, at every prime \(\mathfrak q\), \[
\delta_S(\mathfrak q)=
\dim_{\kappa(\mathfrak q)}\big((I/I^2)\otimes_S\kappa(\mathfrak q)\big)
-n+\dim_{\mathfrak q}\operatorname{Spec}S.\tag{3}
\] This integer retains the full conormal fibre and the original number of variables. It is nonnegative because an ideal with \(\mu\) generators decreases local dimension by at most \(\mu\). The local criteria above show \(\delta_S(\mathfrak q)=0\) exactly at complete-intersection local rings.

It is independent of the finite polynomial presentation. Indeed the stable-conormal isomorphism for two presentations, after tensoring with the residue field, says \(\mu(I/I^2)+m=\mu(J/J^2)+n\); the last dimension in (3) is intrinsic. The localized stable comparison proves invariance under a principal localization around \(\mathfrak q\). Formula (2) proves its counterpart for arbitrary regular local presentations related by a surjection, and the source local-isomorphism neighbourhood construction permits comparison through polynomial presentations in the essentially finite-type setting.

For a field extension \(K/k\) and \(\mathfrak q_K\) over \(\mathfrak q\), flat tensoring gives the exact original ideal \(I_K=K\otimes_k I\) and conormal \(I_K/I_K^2=(I/I^2)\otimes_S S_K\). Their residue vector spaces differ only by scalar extension along \(\kappa(\mathfrak q)\to\kappa(\mathfrak q_K)\), so their dimensions are equal. The source equality of topological dimensions at the corresponding points follows from the two local polynomial/quotient squares: their flat vertical arrows have the same closed fibre, so the local dimension additions cancel, giving equal height differences and hence equal topological dimensions. Therefore \[
\delta_{S_K}(\mathfrak q_K)=\delta_S(\mathfrak q).\tag{4}
\] This proves equality of the entire nonnegative defect, extending the source's equivalence of its vanishing. It does not assert equality of the two local-ring dimensions individually.

\subsubsection{The closed loci of each defect value}\label{algebra-proof-045-the-closed-loci-of-each-defect-value}

Keep a finite presentation \(P\to S\), and put \(M=I/I^2\). For \(0\leq d\leq n\), let \(C_d\) be the union of irreducible components of \(\operatorname{Spec}S\) of dimension at least \(d\). It is closed because there are finitely many components. Algebra 27781--27852 identifies it with the set where \(\dim_{\mathfrak q}\operatorname{Spec}S\geq d\). For an integer \(s\geq1\), define \[
Z_s=\bigcup_{d=0}^{n}
\left(C_d\cap V(\operatorname{Fitt}_{n-d+s-1}(M))\right).\tag{5}
\] To specify all ideals, choose a finite module presentation \(S^a\xrightarrow D S^b\to M\to0\). The Fitting ideal \(\operatorname{Fitt}_i(M)\) consists of all \((b-i)\)-minors of this full matrix; it is \(S\) for \(i\geq b\), and zero for \(i<0\). No row, coefficient or minor is dropped. At a residue field the rank calculation gives \[
\mathfrak q\in V(\operatorname{Fitt}_i(M))
\quad\Longleftrightarrow\quad
\dim_{\kappa(\mathfrak q)}M\otimes\kappa(\mathfrak q)>i.
\] If \(\mathfrak q\) belongs to the term of (5) indexed by \(d\), its actual point dimension is at least \(d\), so (3) gives \(\delta_S(\mathfrak q)\geq s\). Conversely take \(d\) equal to the actual point dimension of a prime with defect at least \(s\); it belongs to that term. Thus \(Z_s=\{\delta_S\geq s\}\). These are closed, and the defect is upper semicontinuous. In particular the complete-intersection locus is the open complement of \(Z_1\); exact defect strata are \(Z_s\setminus Z_{s+1}\), with the zero stratum treated separately.

For an explicit reduced defining ideal let \(a_d\) be the intersection of the minimal primes defining the components in \(C_d\), with empty intersection \(S\). Then \[
J_s=\bigcap_{d=0}^{n}\sqrt{a_d+\operatorname{Fitt}_{n-d+s-1}(M)}
\] defines \(Z_s\). Equations (3)--(4) prove that these closed sets are independent of presentation and their inverse images under every field extension are exactly the corresponding sets for \(S_K\). Accordingly their reduced ideals satisfy \[
J_{s,S_K}=\sqrt{J_{s,S}\,S_K}.
\] The radical is retained; no assertion about an unreduced scheme structure after inseparable base change is made. For the zero algebra the spectrum is empty and all these reduced defining ideals are the unit ideal. These consequences belong in editorial material rather than being inserted in the original definitions.

\subsubsection{Local presentations, neighbourhoods and field change}\label{algebra-proof-045-local-presentations-neighbourhoods-and-field-change}

For the source essentially finite-type local \(k\)-algebra, write a regular local presentation as \(R=(k[x_1,\ldots,x_n]/J)_{\mathfrak q}\twoheadrightarrow S\), with polynomial kernel \(I\supset J\). The local ring \(k[x]_{\mathfrak q}\) and its quotient \(R\) are regular. Formula (1) transfers the generator count to \(I_{\mathfrak q}\), and the original local criterion constructs a global complete-intersection principal neighbourhood whose local ring is \(S\). This proves (3) implies (4) in Algebra 36079--36157 with the actual presentation and prime.

In the converse, the given isomorphism \(S\cong A_{\mathfrak a}\) with \(A\) a finite-type local complete intersection extends to isomorphic principal neighbourhoods of the two finite-presentation \(k\)-algebras. The proof at 31623--31660 lifts the finitely many generators, clears the finitely many relation denominators, and uses the local-isomorphism product decomposition. Its previously composed correction selects the idempotent \((1,0)\) for the desired algebra, not \((0,1)\) for the other factor. Adding one variable with equation \(zg'-1\) expresses the relevant principal localization as a polynomial quotient. The local criterion then proves its polynomial kernel regular, and the regular-presentation equivalence proves that the kernel of the original \(R\to S\) is regular. This supplies (5) implies (2), retaining the given isomorphism and its denominators. The other implications follow directly from the definitions and the local parameter criterion.

The pointwise criterion now gives the three global equivalences in the source: checking maximal ideals suffices because every prime is contained in a maximal ideal, and a principal open containing a maximal ideal also contains each of its generizations. The resulting opens cover the spectrum; quasi-compactness yields a finite subcover if desired. The zero spectrum is covered by the empty family. Equation (4) gives the field-change equivalence at every pair of corresponding primes. Every prime downstairs has a prime above it because field extension is faithfully flat, so it also gives the global equivalence.

The alternative proof in the source is valid: regular sequences ascend through the flat localized polynomial map; at a maximal ideal downstairs the finite residue extension \(\kappa/k\) gives a nonzero finite \(K\)-algebra \(K\otimes_k\kappa\). Its finitely many primes are maximal and form a nonempty set. A chosen prime above it has zero-dimensional closed fibre. The flat local dimension formula consequently gives equal local dimensions in this special maximal-ideal situation. Conormal residue dimensions agree by the exact field-base-change map, and the local generator criterion gives descent. This does not extrapolate that local-dimension equality to arbitrary primes.

\subsubsection{The permanence proof has a genuine hypothesis gap}\label{algebra-proof-045-the-permanence-proof-has-a-genuine-hypothesis-gap}

In Algebra 36360--36420 the square has local rings, with \(R\) and \(\overline S=S/\mathfrak m_RS\) regular, but does not assume \(S\) Noetherian. The cited regular-sequence lifting lemma at 23492 requires both local rings Noetherian. The injectivity lemma at 23424 requires the target Noetherian. The dimension formula at 27421 also requires Noetherian rings. Regularity of the base and closed fibre with flatness does not supply the missing hypothesis.

Here is an exact boundary example, including all maps. Let \(k\) be a field, \[
R=k[t]_{(t)},\quad F=k((t)),\quad
S=k[[t]]+uF[[u]]\subset F[[u]],\quad P=uF[[u]].
\] The map \(R\to S\) is inclusion. The ring \(S\) is local: a series with constant coefficient a unit in \(k[[t]]\) has its formal \(u\)-series inverse in \(S\); the remaining elements form \(\mathfrak m_S=tk[[t]]+P=tS\). Its residue field is \(k\), so \(S/\mathfrak m_RS=k\) is regular. The inclusion is flat because \(S\) is torsion free over the principal ideal domain \(R\): it is the directed union of its finitely generated torsion-free, hence free, \(R\)-submodules.

The ideal \(P\) is not finitely generated. For a finite list of putative generators, their coefficients of \(u\) have \(t\)-valuations bounded below. Multiplying by elements of \(S\) only multiplies those first coefficients by elements of \(k[[t]]\). A term \(u/t^N\) with sufficiently large \(N\) violates that bound. Also \(\bigcap_{n\geq1}t^nS=P\ne0\). There is a strict prime chain \(0\subset P\subset\mathfrak m_S\). It exhausts the primes: primes containing \(P\) correspond to the two primes of \(k[[t]]\); if a nonzero prime contained in \(P\) contains a nonzero element of \(u\)-order \(a\), then a sufficiently high power of every element of \(P\) is divisible by that element with quotient in \(P\subset S\), and hence the prime contains all of \(P\). A prime not contained in \(P\) contains an element with nonzero constant coefficient in \(tk[[t]]\), which is a unit times a power of \(t\) in \(S\), and is therefore maximal. Thus \(\dim S=2\), whereas \(\dim R+\dim(S/tS)=1\).

Take \(A=R\), \(B=S\), \(I=J=0\). Both ideals in condition (3) are generated by the empty regular sequence and both vertical maps are flat. All source assumptions hold, while the dimension equation invoked at 36415 is false. This refutes that step of the printed proof, not the lemma's conclusion (which is true in this example). The next two sections prove the original conclusion without adding Noetherianity of \(S\), and in fact remove two unused regularity hypotheses. The source proof needs a visible conceptual erratum, not an unproved assertion that its whole theorem is false or a silent restriction of its statement.

\subsubsection{A Tor-product criterion with a complete proof}\label{algebra-proof-045-a-tor-product-criterion-with-a-complete-proof}

Let \((R,\mathfrak m,k)\) be a Noetherian local ring, \(I\subset\mathfrak m\), and \(A=R/I\). If the graded algebra \(\operatorname{Tor}^R_*(A,k)\) is generated by its degree-one part, then \(I\) is generated by a regular sequence. Only the surjectivity of products onto degree two is needed.

Choose a minimal generating list \(f_1,\ldots,f_r\) of \(I\), and keep its Koszul differential \[
d(e_i)=f_i,\qquad d(e_i\wedge e_j)=f_i e_j-f_j e_i.
\] The finite module \(H_1(K_R(f))\) is annihilated by \(I\), because multiplication by each \(f_i\) is homotopic to zero through wedge multiplication by \(e_i\). Choose a minimal list of its generators, represented by cycles \(z_1,\ldots,z_a\in K_1\). Every coefficient of every \(z_j\) is in \(\mathfrak m\): a relation \(\sum b_i f_i=0\) reduces in \(I/\mathfrak mI\) to a relation on the basis given by the minimal generators, so each \(b_i\) is zero modulo \(\mathfrak m\).

Extend the augmented differential graded algebra \(K_R(f)\to A\) to an \(R\)-free differential graded algebra resolution by adjoining generators \(y_j\) of degree two with \(dy_j=z_j\), then adjoining generators in successive degrees whose differentials represent generators of the remaining positive homology. One may use free associative graded algebra adjunctions with \(R\) central; no division by factorials is required. The original exterior relations among the \(e_i\) are retained. In each step adjoining a generator in degree \(n+1\) kills its specified degree-\(n\) class and changes no lower degree. Taking the union after successively killing all positive degrees gives an augmented resolution. Its underlying modules are free: free products over \(R\) have the basis of alternating words in the chosen free augmentation-ideal bases. This construction provides the resolution rather than assuming a missing acyclicity assertion.

In degrees at most two the resulting resolution \(T\) has \[
T_0=R,\quad T_1=\bigoplus_i Re_i,\quad
T_2=\bigwedge^2R^r\oplus\bigoplus_{j=1}^a Ry_j.
\] Its first two differentials vanish after tensoring with \(k\). For every element \(w+\sum a_jy_j\in\ker(T_2\to T_1)\), the equation \(d_Kw+\sum a_jz_j=0\) gives \(\sum a_j[z_j]=0\) in \(H_1(K_R(f))\). Minimality of the chosen homology generators implies \(a_j\in\mathfrak m\) for every \(j\). Every image from degree three is such a cycle, including differentials of products and of the subsequently adjoined degree-three generators. Consequently projection onto the \(y_j\) coordinates induces a surjection \[
\operatorname{Tor}^R_2(A,k)=H_2(T\otimes_R k)\longrightarrow k^a.
\] All degree-one products come from \(e_i e_j\), so this projection kills every such product. The Tor product here is the usual one obtained by tensoring an \(R\)-flat differential graded algebra resolution with \(k\); comparison through free algebra resolutions preserves the multiplication. Hence generation by degree one forces \(a=0\), and Nakayama gives \(H_1(K_R(f))=0\).

Finally, vanishing of this Koszul homology implies regularity in the Noetherian local ring. Induct on \(r\). Write \(K'=K_R(f_1,\ldots,f_{r-1})\). The mapping-cone exact sequence contains \[
H_1(K')\xrightarrow{f_r}H_1(K')\to H_1(K_R(f))
\to H_0(K')\xrightarrow{f_r}H_0(K').
\] The zero middle homology makes the first map surjective. The finite module \(H_1(K')\) is therefore zero by Nakayama. Induction gives regularity of the first \(r-1\) elements, and the last arrow is injective, proving regularity of \(f_r\) on the preceding quotient. The final quotient is nonzero because all generators lie in \(\mathfrak m\). The empty list gives \(I=0\), with the same conclusion. This proves the criterion in every characteristic.

\subsubsection{Permanence without the unused regularity assumptions}\label{algebra-proof-045-permanence-without-the-unused-regularity-assumptions}

Consider the source square \(R\to S\), \(A=R/I\to B=S/J\), with local maps, \(R\) Noetherian local, and both \(R\to S\) and \(A\to B\) flat. Assume \(J\) is generated by a regular sequence of length \(c\), and its image \[
\overline J=J/(\mathfrak m_RS\cap J)\subset\overline S=S/\mathfrak m_RS
\] is generated by a regular sequence of length \(b\). Put \(k=R/\mathfrak m_R\) and \(\overline B=B/\mathfrak m_RB=\overline S/\overline J\). These are the original image ideal and fibre rings; the two already composed denominator corrections MC-STK-ERR-0443 and 0444 remain under their existing identities.

Let \(z_1,\ldots,z_c\) be the specified regular generators of \(J\). Their Koszul algebra is a finite resolution of \(B\) by free \(S\)-modules, each flat over \(R\). Its tensor product with \(k\) is exactly \(K_{\overline S}(\overline z_1,\ldots,\overline z_c)\). To compute its entire homology algebra, choose regular generators \(h_1,\ldots,h_b\) of \(\overline J\). A regular sequence has free conormal module on its classes, so these form a minimal generating list over the local ring \(\overline S\). Write the \(\overline z_i\) as linear combinations of the \(h_j\). Their coefficient matrix has rank \(b\) modulo the maximal ideal, since the \(\overline z_i\) generate \(\overline J\). An invertible \(b\)-minor, permutation, its matrix inverse, and elementary row subtractions give a matrix in \(\mathrm{GL}_c(\overline S)\) carrying this list to \((h_1,\ldots,h_b,0,\ldots,0)\). This also proves \(c\geq b\). The induced exterior-power maps give an isomorphism of Koszul algebras, keeping all matrix coefficients and differentials. Since \(h\) is regular, \[
\operatorname{Tor}^R_*(B,k)
\cong\bigwedge^*_{\overline B}(\overline B^{\oplus(c-b)}).
\tag{6}
\] This calculation uses no Noetherian property of \(S\) or \(\overline S\).

For the second computation, take an \(R\)-free differential graded algebra resolution \(G\to k\), constructed by successive cycle adjunctions as above. The complex \(G\otimes_R A\) consists of free \(A\)-modules and has homology \(\operatorname{Tor}^R_*(A,k)\). Flatness of \(B\) over \(A\) commutes with its kernels and images, giving the multiplicative isomorphism \[
\operatorname{Tor}^R_*(A,k)\otimes_k\overline B
\ \xrightarrow{\ \cong\ }\ \operatorname{Tor}^R_*(B,k).
\tag{7}
\] The scalar reduction is exact: the \(A\)-action on these Tor groups factors through \(k\), and \(B\otimes_A k=\overline B\). The map sends a Tor class tensored with \(b\) to \(b\) times its image under \(A\to B\), so it respects products. Locality makes \(\overline B\ne0\), and hence it is faithfully flat over the field \(k\).

In (6), the intrinsic map from the exterior algebra of degree-one homology to all homology is an isomorphism. Combining (6) and (7), and descending its kernel and cokernel along the faithful field extension to the nonzero \(k\)-algebra \(\overline B\), shows that \(\operatorname{Tor}^R_*(A,k)\) is generated by degree one. The preceding criterion now proves that \(I\) is generated by a regular sequence. Also \[
\operatorname{Tor}^R_1(A,k)=I/\mathfrak m_RI,
\qquad \mu_R(I)=c-b,
\] because (6) has degree-one rank \(c-b\); tensoring further to the residue field of \(\overline B\) compares these finite ranks. Any minimal generating list of \(I\) is regular by the same Koszul criterion. Thus the original conclusion holds for arbitrary \(S\), and the assumptions that \(R\) and \(\overline S\) are regular can be replaced by the sole requirement that \(R\) be Noetherian local, while retaining all the flatness and regular-sequence assumptions. In particular no missing hypothesis has been inserted to make the printed proof work.

The source's equation \(J=(f_1,\ldots,f_b)+IS\), for lifts of regular fibre generators, also admits a direct proof without the cited Noetherian-target injectivity lemma. Its kernel module \(L=J/((f_1,\ldots,f_b)+IS)\) is finite over \(S\), because \(J\) is finitely generated. Tensor the exact sequence \(0\to L\to S/((f)+IS)\to B\to0\) over \(A\) with \(k\). Flatness of \(B\) over \(A\) preserves injectivity on the left, and the two quotient rings have identical closed fibres. Hence \(L/\mathfrak m_RL=0\), so Nakayama gives \(L=0\). The later claims of flatness of \(S/(f)\) and the dimension calculations still require a different justification in the source's stated generality; (6)--(7) supply a complete proof of the theorem without relying on them.

\subsubsection{Report dispositions and propagation}\label{algebra-proof-045-report-dispositions-and-propagation}

Localization of an algebra at a prime of its base polynomial ring means inversion of the complement of that prime; it is well defined and equals localization at the corresponding quotient prime when the kernel is contained in it. Thus the notation in OCC-00771 is not an undefined localization. A maximal prime is a maximal ideal, so OCC-00773 is a terminology preference. Finite products and sums of modules coincide, making the ordinary powers in OCC-00774 legitimate finite free modules. The TeX forms q-subscript-K-prime and q-prime-subscript-K render identically, so OCC-00779 does not identify a mathematical or rendered inconsistency. A comma in ``Let {[}diagram{]} be \ldots'' would interrupt the verb construction; OCC-00785 is rejected. Source-only spacing in OCC-00788 and terse transitions remain optional. Missing designation prepositions, one finite verb, the singular condition wording, and terminal punctuation receive local proposals. The two repeated mathematical denominator findings retain their two existing canonical units and three operations, with no double application.

The precise stable-conormal result from groups 991--992 supports the presentation independence and local criterion here. The original conormal base-change result and group 995 support the exact residue comparison; no Tor defect is lost because field extension is flat. Groups 1017--1018 record the quantitative defect and its closed strata, and group 1019 records the generalized permanence theorem and exact length formula. Group 1016 records the printed proof's actual hypothesis gap together with its full resolution, rather than asserting a false counterexample to the theorem. All three strengthenings have zero source operations. The pending More on Algebra tensor-base correction from group 997 remains separate. Broader synthesis and a short paper follow completion of the core work.

\hypertarget{algebra_syntomic_notes_20260928}{}
\subsection{Syntomic morphisms: source review and separate editorial consequences}\label{algebra-proof-046-syntomic-morphisms-source-review-and-separate-editorial-consequences}

Authority: Stacks Project authors, commit \texttt{a04446e57ec1fbc252a871afcec7752fb2807b14}, \texttt{algebra.tex:36427–37054}, SHA256 \texttt{FA8BB92E58A4F78A2BD01B3B6A4A87DE0A0D279F5DD90641B574DD5FBFFFA4F3}. The additional bounded reading at 37760--37775 verifies the other two operations of the already composed localization unit MC-STK-ERR-0446; it does not adjudicate the intervening smooth section. The next source interval begins at 37055.

This is editorial evidence. It does not replace the source statements, original presentations, proofs or exposition in a translation. Minimal text operations are recorded separately in the review. The source's omitted details do not automatically count as errors. The consequences below are proved comparisons and deductions, without a novelty claim. The bounded corpus query \texttt{syntomic\ complete\ intersection} returned two research records, including one PDF-only record; both remain unread routing records. No PDF reading or exhaustive literature survey is claimed.

\subsubsection{Descent, base change and the original presentations}\label{algebra-proof-046-descent-base-change-and-the-original-presentations}

For a faithfully flat map \(R\to R'\), write \(S'=R'\otimes_R S\). Finite presentation and flatness descend and ascend as stated at 36462--36464. For every prime \(\mathfrak p'\) above \(\mathfrak p\), the original fibre comparison is \[
(R'\otimes_R S)\otimes_{R'}\kappa(\mathfrak p')
\longrightarrow
(S\otimes_R\kappa(\mathfrak p))\otimes_{\kappa(\mathfrak p)}\kappa(\mathfrak p'),
\qquad (r'\otimes s)\otimes a\longmapsto(s\otimes1)\otimes r'a.
\] Its inverse sends \((s\otimes b)\otimes a\) to \((1\otimes s)\otimes ba\). The balancing relations give well-defined mutually inverse algebra maps. Faithful flatness supplies a prime above every \(\mathfrak p\); the previously checked field-extension equivalence for local complete intersections then gives both directions. For arbitrary base change the same comparison gives ascent at every new prime. The finite principal-cover assertion follows because flatness, finite presentation and the complete-intersection property of the fibre local rings are local on the target. Empty spectra and the zero algebra satisfy the stated conventions.

Retain a relative global complete-intersection presentation exactly as given: \[
P=R[x_1,\ldots,x_n],\qquad I=(f_1,\ldots,f_c),\qquad S=P/I.
\] Every nonempty fibre has dimension \(n-c\). Localizing at the image \(g\) of \(h\in P\) uses the original presentation \[
S_g=R[x_1,\ldots,x_n,x_{n+1}]/(f_1,\ldots,f_c,hx_{n+1}-1).
\] The forward map sends \(x_{n+1}\) to \(g^{-1}\); the inverse extends \(P\to P[x_{n+1}]/(I,hx_{n+1}-1)\), where \(g\) is invertible. A nonempty open of a field global complete intersection has the same component dimension \(n-c\), by the original complete-intersection criterion. Thus this additional variable and additional equation retain the required difference. The presentation and these maps commute with arbitrary \(R\to R'\).

For Huber's lemma, finite presentation makes \(I\) finitely generated. If \(S\ne0\), a free finite \(S\)-module has finite rank: a finite generating family involves only finitely many basis coordinates, so any further basis vector could not lie in its span. Choose the specified basis lifts \(f_1,\ldots,f_c\). In the finite module \(M=I/(f_1,\ldots,f_c)\), one has \(IM=M\). The determinant proof of Nakayama, authority 3615--3668, supplies \(g\in1+I\) annihilating \(M\). Therefore \(I_g=(f_1,\ldots,f_c)_g\) and \(S_g=S\). For the kernel \(J=(f_1,\ldots,f_c,gz-1)\) of \(P[z]\to S\), the principal-localization conormal comparison at 35604--35661, with its already checked continuation, identifies \[
J/J^2\cong (I/I^2)_g\oplus S(gz-1).
\] The comparison sends each original \(f_i\) to its class and the last basis vector to \(gz-1\). Hence these actual equations form the claimed basis. If \(S=0\), the presentation \(R/(1)\) suffices; its conormal module is zero over the zero ring, so the basis assertion has the stipulated empty-spectrum interpretation.

\subsubsection{Neighbourhoods, approximation and local regularity}\label{algebra-proof-046-neighbourhoods-approximation-and-local-regularity}

Authority 36670--36720 uses the open locus where fibre dimension is at most \(n-c\), established at 31307--31327. In case (1), write that open as \(\operatorname{Spec}(S)\setminus V(J)\). Disjointness \(V(J)\cap V(IS)=\varnothing\) means \(J+IS=S\); an identity \(g+a=1\), \(g\in J\), \(a\in IS\), supplies the required \(g\). In case (2), \(J(S\otimes_R\kappa(\mathfrak p))\) is the unit ideal. A finite identity there can be written \[
1=\sum_{i=1}^r \overline{j_i}\,\overline{s_i}/\overline{t_i},
\qquad j_i\in J,\ s_i\in S,\ t_i\in R\setminus\mathfrak p.
\] Put \(t=\prod_i t_i\) and \(g=\sum_i j_i s_i\prod_{\ell\ne i}t_\ell\). Then \(g\in J\) and its image in the fibre is the nonzero scalar \(\overline t\), hence is a unit. The empty finite sum causes no issue in a zero fibre: take \(g=0\), whose image is the unit of the zero ring. In case (3), choose a principal neighbourhood of the given prime within the same dimension open. In every case the displayed localization has \(n+1\) variables and \(c+1\) equations. The height theorem in the polynomial ring over each residue field gives dimension at least \(n-c\) for each of its nonempty components. The constructed upper bound therefore gives equality. No equation or denominator has been dropped.

For 36726--36764, begin with the subring \(R_0\) generated over \(\mathbf Z\) by every coefficient of the original \(f_i\). The source field-extension argument and quasi-compactness give finitely many \(g_j\in S_0\) whose distinguished opens have the required dimension bound and cover the image of \(\operatorname{Spec}(S)\). Choose polynomial lifts \(G_j\) and representatives \(A_j\in R[x_1,\ldots,x_n]\) of a unit-ideal identity in \(S\). There are polynomials \(H_i\in R[x_1,\ldots,x_n]\) with \[
1-\sum_j A_jG_j=\sum_i H_i f_i.
\] Adjoin every coefficient of all \(A_j,H_i\) to \(R_0\). This finite enlargement makes the identity hold already over the enlarged ring, with the same original \(f_i\). Fibre dimensions on each localized chart remain bounded after that base change. The charts now cover its entire spectrum, and the height lower bound gives equality. If \(S=0\), use the certificate \(1=\sum_i H_i f_i\); the enlarged \(S_0\) is then zero. This proves the approximation without assuming a nonempty spectrum or inserting an extra relation.

In 36766--36840, over a Noetherian base the original fibre regular sequence lifts by the local Noetherian result at 23491--23504, since the localized polynomial ring is flat over the local base. It also gives flatness of each displayed partial quotient. For an arbitrary base take all finitely generated \(\mathbf Z\)-subalgebras containing the just-constructed \(R_0\). Each is Noetherian and each carries the same complete-intersection presentation. At the contractions of a fixed original prime, the local polynomial rings, their partial quotients and the multiplication maps by each \(f_{i+1}\) have the stated filtered colimits. Exactness of filtered colimits preserves those injections; the quotient by all \(f_i\) is nonzero at the original prime. The flat-colimit criterion gives flatness over \(R\) of each partial quotient. This use of the Noetherian lemma has its hypothesis at every finite stage and does not apply it directly to the possibly non-Noetherian final ring.

The map \(S^c\to I/I^2\), \((a_i)\mapsto\sum_i a_i f_i\bmod I^2\), is an isomorphism at every prime of \(S\), by regular implies quasi-regular, authority 16895--16978. Its kernel and cokernel consequently vanish. Flatness of \(S\) follows from the proved flatness of \(S_{\mathfrak q}\) for every \(\mathfrak q\): for any injection of \(R\)-modules, the kernel after tensoring with \(S\) has zero localization at every prime of \(S\), and is therefore zero. The source's phrase ``syntomic, i.e., flat'' is valid in this relative-global-complete-intersection context, where finite presentation and the fibre condition are already established.

\subsubsection{Ordered roots with all coefficients and signs}\label{algebra-proof-046-ordered-roots-with-all-coefficients-and-signs}

The basis asserted at 36604--36636 can be proved without replacing the original polynomial by a rescaled one. Start with \[
Q_0(x)=x^n+a_1x^{n-1}+\cdots+a_n\in R[x],\qquad
R=\mathbf Z[a_1,\ldots,a_n].
\] Set \(T_0=R\). Inductively put \[
T_i=T_{i-1}[\beta_i]/(Q_{i-1}(\beta_i)),\qquad
Q_{i-1}(x)=(x-\beta_i)Q_i(x)\quad\hbox{in }T_i[x].
\] The second equality exists uniquely by division by the monic linear polynomial. To retain its full coefficients, if \(Q_{i-1}(x)=x^d+u_1x^{d-1}+\cdots+u_d\), write \(Q_i(x)=x^{d-1}+v_1x^{d-2}+\cdots+v_{d-1}\), where \(v_0=1\) and \(v_j=u_j+\beta_i v_{j-1}\) for \(1\le j\le d-1\). The final constant equality is \(u_d=-\beta_i v_{d-1}\), precisely the relation \(Q_{i-1}(\beta_i)=0\). Monic polynomial division gives \(T_i\) the free \(T_{i-1}\)-basis \(1,\beta_i,\ldots,\beta_i^{n-i}\). Thus \(T_n\) has the original ordered monomial basis with \(0\le e_i\le n-i\), and rank \(n!\).

There is a map \(T_n\to\mathbf Z[\alpha_1,\ldots,\alpha_n]\) taking \(\beta_i\mapsto-\alpha_i\), \(a_j\mapsto e_j(\alpha_1,\ldots,\alpha_n)\). Each successive factor relation is satisfied. Conversely send \(\alpha_i\mapsto-\beta_i\). In \(T_n[x]\), the iterated divisions give \(Q_0(x)=\prod_i(x-\beta_i)\); comparison of every coefficient gives \(a_j=e_j(-\beta_1,\ldots,-\beta_n)\). The two maps are inverse on all generators, including the original \(a_j\). The signs change each basis monomial by the unit \((-1)^{\sum e_i}\), and the displayed \(\alpha\)-basis follows.

For the source polynomial \(P(x)=x^n+b_1x^{n-1}+\cdots+b_n\in A[x]\), base change along \(a_i\mapsto b_i\) gives precisely \(A'=A\otimes_R S\) and roots \(\beta_i=-1\otimes\alpha_i\). It is free of rank \(n!\). Its basis contains \(1\), so the map \(A\to A'\) is injective and tensoring with \(A'\) is a direct sum of \(n!\) copies on underlying modules; it reflects zero modules. Its syntomic property follows by the checked base-change result. For \(n=0\), the polynomial is \(1\), the extension is \(A'=A\), the basis is the single empty monomial, and the factorization is the empty product. No factorial has been inverted in any coefficient ring.

\subsubsection{Finite locally free factorization algebras and their exact rank}\label{algebra-proof-046-finite-locally-free-factorization-algebras-and-their-exact-rank}

This is a separate consequence of the source example 36550--36602 and the proved flatness theorem 36842--36854. Retain \(n,m\ge1\), \(N=n+m\), and the original rings \[
R=\mathbf Z[a_1,\ldots,a_N],\qquad
S=\mathbf Z[b_1,\ldots,b_n,c_1,\ldots,c_m].
\] The original coefficient map is exactly \(a_\ell\mapsto\sum_{i+j=\ell}b_i c_j\), with \(b_0=c_0=1\), \(b_i=0\) outside \(0\le i\le n\), and \(c_j=0\) outside \(0\le j\le m\). This convention interprets the displayed \(b_2,c_2\) when \(n=1\) or \(m=1\) and agrees with the source's defining polynomial identity.

The map is finite. Indeed \(S\) is a domain, so apply the monic-divisor result at 8432--8456 in its fraction field to each factor of the original monic degree-\(N\) polynomial. Each \(b_i,c_j\) is integral over the image of \(R\), and the finite list of integral generators makes \(S\) finite as an \(R\)-module. Every geometric fibre is nonempty: split the specialized polynomial over an algebraic closure and allocate any \(n\) of its \(N\) roots, with multiplicities, to the first factor. Finiteness gives dimension zero for each nonempty fibre. The source presentation has \(N\) variables and \(N\) equations, hence is a relative global complete intersection and is flat. Since \(R\) is Noetherian, its finite module \(S\) is finitely presented. It is therefore finite projective, equivalently finite locally free.

Its rank is exactly \(\binom{N}{n}\). The rank of a finite projective module is locally constant; \(\operatorname{Spec}(R)\) is connected since \(R\) is a domain. It suffices to compute after extending the fraction field of \(R\) to an algebraic closure \(E\). The generic polynomial has distinct roots: its discriminant is a nonzero polynomial, as witnessed by the characteristic-zero specialization \(\prod_{i=1}^N(x-i)\). There are precisely \(\binom Nn\) monic factorizations of degrees \(n,m\) over \(E\), one for each size-\(n\) subset of these distinct roots.

To count scheme lengths as well as points, fix such a factorization \(F=BC\). Its Zariski tangent vectors are coefficient lists of polynomials \(\dot B,\dot C\), with degrees less than \(n,m\), satisfying \[
C\dot B+B\dot C=0.
\] The factors \(B,C\) are relatively prime. Reduction modulo \(B\) gives \(\dot B=0\bmod B\), hence \(\dot B=0\) by its degree; then \(\dot C=0\). This calculation is the linearization of all \(N\) original coefficient equations, with no discarded constant term. The finite local \(E\)-algebra at the point thus has maximal ideal \(\mathfrak m\) with \(\mathfrak m/\mathfrak m^2=0\). Nakayama gives \(\mathfrak m=0\); its residue field is \(E\), so its length is one. Decomposition of the finite Artinian fibre into its local factors proves that its dimension over \(E\) is \(\binom Nn\), as claimed.

For every map \(R\to A\), the same algebra \(A\otimes_R S\), with the same coefficient equations, is finite locally free of constant rank \(\binom Nn>0\), and is faithfully flat and syntomic. The rank remains the same at repeated-root fibres, counting their full algebraic lengths; reducedness at those fibres is not asserted. Faithfulness follows because localization at each prime is a nonzero finite free module of that positive rank. This argument asserts finite local freeness, not an unproved global choice of basis for the factorization algebra.

\subsubsection{Universal Koszul resolution and every infinitesimal neighbourhood}\label{algebra-proof-046-universal-koszul-resolution-and-every-infinitesimal-neighbourhood}

Keep the original relative global complete-intersection presentation \(P=R[x_1,\ldots,x_n]\), \(I=(f_1,\ldots,f_c)\), \(S=P/I\). Its Koszul complex has \[
K_j=\bigwedge^j_P P^c,\qquad
d(e_{i_1}\wedge\cdots\wedge e_{i_j})
=\sum_{r=1}^j(-1)^{r-1}f_{i_r}
 e_{i_1}\wedge\cdots\widehat e_{i_r}\cdots\wedge e_{i_j},
\] for \(i_1<\cdots<i_j\). Every pair of terms in \(d^2\) has opposite signs and the same coefficient \(f_{i_r}f_{i_s}\), so \(d^2=0\).

This complex resolves \(S\). At a prime containing \(I\), the source has proved that the specified \(f_i\) form a regular sequence. The regular-sequence Koszul assertion follows by induction: adjoining \(f_j\) forms the mapping cone of multiplication by \(f_j\) on the preceding Koszul complex; its only possible positive homology is the kernel of multiplication by \(f_j\) on \(P/(f_1,\ldots,f_{j-1})\), and that kernel is zero. At a prime not containing \(I\), some \(f_i\) is invertible. The degree-one operator \(h(\omega)=f_i^{-1} e_i\wedge\omega\) obeys \(dh+hd=1\) by the displayed differential. Thus the localized complex is contractible there. Localizing its homology at every prime now proves exactness globally, with degree-zero cokernel \(S\).

Every \(K_j\) is \(R\)-flat, and \(S\) is \(R\)-flat by the source theorem. In the augmented exact complex, let \(Z_0=S\) and let \(Z_{j+1}\) be the kernel of the surjection \(K_j\to Z_j\). In \[
0\longrightarrow Z_{j+1}\longrightarrow K_j\longrightarrow Z_j\longrightarrow0,
\] flatness of \(K_j,Z_j\) implies flatness of \(Z_{j+1}\), by the long exact Tor sequence. Induction shows that all these sequences remain exact after tensoring with every \(R\)-module \(M\). Consequently \(K\otimes_R M\to S\otimes_R M\) is exact in positive degrees. In particular every algebra base change retains the resolution, with the original images of all \(f_i\).

Tensoring the \(P\)-free resolution with \(S\) makes every differential zero. Its multiplication is the exterior multiplication already specified, so there is an isomorphism of graded \(S\)-algebras \[
\operatorname{Tor}^{P}_*(S,S)\cong\bigwedge^*_S S^c,
\] with degree-one basis labelled by the original equations. This assertion includes its algebra structure and commutes with the algebra base changes just established.

The entire associated graded algebra is also explicit: \[
S[T_1,\ldots,T_c]\longrightarrow\operatorname{gr}_I(P)
=\bigoplus_{d\ge0}I^d/I^{d+1},\qquad T_i\longmapsto f_i\bmod I^2.
\] At primes of \(S\), it is the regular-to-quasi-regular isomorphism whose full coefficient induction is at 16907--16977. For completeness, its coefficient test says that if \(\sum_{|\nu|=d}a_\nu f^\nu\in I^{d+1}\), then every \(a_\nu\in I\). Subtract a representation of the element of \(I^{d+1}\) to obtain a zero relation. With \(J=(f_1,\ldots,f_{c-1})\), group it by the highest occurring power \(f_c^\ell\). Induction on \(c\) identifies each \(J^r/J^{r+1}\) as free over \(P/J\); multiplication by \(f_c\) is injective there. The highest group therefore has every coefficient in \(J\). Write those coefficients as sums of multiples of \(f_1,\ldots,f_{c-1}\), and move one factor \(f_c\) into the coefficient in the group with power \(\ell-1\). Induction on \(\ell\) then puts every coefficient in \(I\). For \(c=1\), cancel the nonzerodivisor power \(f_1^d\); for \(c=0\), the assertion is the identity in degree zero and zero in positive degrees. This proves the localized coefficient test. The kernel and cokernel of each graded piece are \(S\)-modules and vanish at all its primes, so the displayed map is globally an isomorphism.

In particular \(I^d/I^{d+1}\) is free over \(S\) on every monomial \(f_1^{\nu_1}\cdots f_c^{\nu_c}\) with \(\nu_i\in\mathbf Z_{\ge0}\) and \(\sum_i\nu_i=d\). Each such module is \(R\)-flat. The sequences \[
0\longrightarrow I^d/I^{d+1}\longrightarrow P/I^{d+1}\longrightarrow P/I^d\longrightarrow0
\] show inductively that every \(P/I^d\), \(d\ge1\), is \(R\)-flat. From \(0\to I^d\to P\to P/I^d\to0\) it follows that \(I^d\) is \(R\)-flat too. For an arbitrary algebra map \(R\to R'\), put \(P'=R'[x_1,\ldots,x_n]\) and \(I'=IP'\). Tensoring that last short exact sequence gives an injection \(I^d\otimes_R R'\to P'\) whose image is exactly \((I')^d\). Thus \[
I^d\otimes_R R'\cong(I')^d,\quad
(P/I^d)\otimes_R R'\cong P'/(I')^d,\quad
\operatorname{gr}_I(P)\otimes_R R'\cong\operatorname{gr}_{I'}(P').
\] The maps send each original power product to the identical power product of its images, and respect multiplication in every degree. No flatness of \(R'\) over \(R\) is needed. These strengthenings follow for this particular relative-complete-intersection presentation; they do not assert arbitrary nonflat base change for all quasi-regular ideals.

Global regularity of the ordered list in the whole polynomial ring is not inferred. For example, take \(R=k\times k\), \(P=R[x,y]\cong k[x,y]\times k[x,y]\), \(f_1=(x,0)\), \(f_2=(y,1)\). Then \(P/(f_1,f_2)\cong k\times0\); its nonempty fibre has dimension \(0=2-2\), so this is a relative global complete intersection. But \(f_1(0,1)=0\), and \((0,1)\ne0\). The original ordered list is not regular in \(P\). The local statements, full Koszul resolution and associated graded conclusion above all still hold.

\subsubsection{Local criteria, composition and lifting}\label{algebra-proof-046-local-criteria-composition-and-lifting}

For 36877--36943, after the permitted principal localization choose \(P\to S=P/I\) with \(I\) finite, and retain the prime \(\mathfrak q'\subset P\) above \(\mathfrak q\), with contraction \(\mathfrak p\subset R\). Choose the original lifts \(f_1,\ldots,f_c\in I\) of a basis of the fibre ideal modulo its local maximal ideal. They generate that localized fibre ideal by Nakayama. Put \(S'=P/(f_1,\ldots,f_c)\), \(J=I/(f_1,\ldots,f_c)\). The ideal \(J\) is finite because \(I\) is finite. Localization gives \[
0\to J_{\mathfrak q'}\to S'_{\mathfrak q'}\to S_{\mathfrak q}\to0.
\] The assumed flatness over \(R_{\mathfrak p}\), followed by flatness of \(R_{\mathfrak p}\) over \(R\), makes the last module \(R\)-flat. Tensoring with \(\kappa(\mathfrak p)\) therefore preserves the initial injection. The last map is an isomorphism on that fibre by choice of the \(f_i\), so \(J_{\mathfrak q'}/\mathfrak pJ_{\mathfrak q'}=0\). Nakayama in the local ring \(S'_{\mathfrak q'}\) gives \(J_{\mathfrak q'}=0\). For each of a finite set of generators of \(J\), choose a killing denominator outside \(\mathfrak q'\); their product gives the one \(g\) with \(J_g=0\). The relative dimension at the point is \(n-c\) by the fibre complete-intersection criterion. Applying the proved neighbourhood lemma gives the required further principal localization, still avoiding the given prime. This establishes the source implication without assuming flatness of \(S'\) in advance.

The conormal projectivity statement at 36945--36964 uses a finite cover by those relative complete intersections. On each chart, the corrected stable comparison at 35775--35797 says that the localized original conormal module plus a finite free module is isomorphic to the chart conormal module plus another finite free module. The chart conormal module is finite free. The original conormal module is consequently a direct summand of a finite free module and is finite projective on that chart. The finite-cover criterion then gives the result over \(S_g\). No cancellation claiming freeness of a stably free module is used.

For composition retain both presentations and all chosen lifts: \[
S=R[x_1,\ldots,x_n]/(f_1,\ldots,f_c),\quad
S'=S[y_1,\ldots,y_m]/(h_1,\ldots,h_d),
\] \[
S'=R[x_1,\ldots,x_n,y_1,\ldots,y_m]/(f_1,\ldots,f_c,h'_1,\ldots,h'_d).
\] On a fibre over a field, localize at a maximal ideal of \(S'\) and contract to \(S\). The source base/fibre dimension inequality bounds that local dimension by \((n-c)+(m-d)\). Taking the supremum over maximal ideals bounds the total fibre dimension by the same number. The polynomial height theorem for the displayed \(c+d\) equations gives the reverse bound on every nonempty component. This proves the exact relative global complete-intersection condition, including the permitted empty fibres. For syntomic maps use the corrected primes \(\mathfrak q'\subset S'\), \(\mathfrak q\subset S\), choose \(g'\in S'\setminus\mathfrak q'\) and \(g\in S\setminus\mathfrak q\), and retain the actual base change \(S_g\to S'_{gg'}\). Their composition is \(R\to S'_{gg'}\). These are exactly the three operations of MC-STK-ERR-0445. Quasi-compactness gives a finite principal cover of \(\operatorname{Spec}(S')\), completing the source argument.

In the lifting lemma, the chosen elements lie in \(\overline S\) and are \(\overline g_i\); the target is therefore \(\overline S_{\overline g_i}\). Choose a finite subcover, keep its given presentations, and lift all coefficients of every \(\overline f_i\) to \(R\). For each chart the resulting \(S\) has \(S/IS\cong\overline S\). The earlier neighbourhood construction supplies \(g\equiv1\bmod IS\) with \(S_g\) a relative global complete intersection. Its reduction is canonically the same \(\overline S\), since the image of \(g\) is \(1\). This proves the chartwise lifting asserted by the source; no global gluing of those lifts is claimed. The two occurrences here and the two in 37760--37775 are the four operations of MC-STK-ERR-0446. The later smooth proof as a whole remains unadjudicated.

\subsubsection{Report dispositions and propagation}\label{algebra-proof-046-report-dispositions-and-propagation}

Twenty-six received reports occupy twenty-two decision groups. Eleven groups propose fourteen minimal prose or punctuation operations. Seven source-style or optional wording preferences are retained unchanged: whitespace after a TeX comma, ``at most finitely'', the ordinal spelling, punctuation before an enumeration, the collective singular regular-sequence and basis constructions, and ``Proof of (2).'' as a proof transition. Two reports are rejected as nondefects: the context-specific equivalence with flatness, and the two TeX orders for the same prime/subscript notation. Four reports share the already composed composition unit; two share the already composed lifting unit. The latter is verified as a whole across its two source intervals.

Two further groups have no source operations. The factorization rank result receives the proved syntomic flatness and the finite-projective criterion, with every fibre multiplicity retained. The universal Koszul and infinitesimal-neighbourhood result receives the earlier regular/quasi-regular and Tor arguments and sharpens their base-change conclusion in the specified relative-complete-intersection setting. The stronger permanence theorem from the preceding LCI note is preserved but is not needed as a new hypothesis here. The More on Algebra tensor-base correction recorded earlier remains a separate pending chapter item. The source hypotheses and exposition remain identifiable throughout; combined-results synthesis and a short paper remain follow-on work after core completion.

\hypertarget{algebra_smooth_notes_20260928}{}
\subsection{Smooth ring maps: source decisions and editorial consequences}\label{algebra-proof-047-smooth-ring-maps-source-decisions-and-editorial-consequences}

Authority: Stacks Project authors, commit a04446e57ec1fbc252a871afcec7752fb2807b14, algebra.tex:37055--37794, SHA256 FA8BB92E58A4F78A2BD01B3B6A4A87DE0A0D279F5DD90641B574DD5FBFFFA4F3. The chapter's original objects, equations and exposition remain identifiable. These supporting proofs and further consequences are editorial material, not replacement translations. The next section starts at 37795; only its opening through 37805 is read. The bounded corpus query ``smooth Jacobian cotangent'' returned two TeX-routing hits whose titles concern integrable systems and Hitchin quantization; neither has been read or used as evidence. No novelty or exhaustive-literature claim is made.

\subsubsection{The source definitions and the split conormal map}\label{algebra-proof-047-the-source-definitions-and-the-split-conormal-map}

The reports OCC-00822 and OCC-00831 incorrectly call \(H_1(L_{S/R})\) unexplained. The source defines that very notation at 35152--35153 as \(H_1(\mathrm{NL}_{S/R})\), then defines the full cotangent complex and its truncation at 35157--35176. The previously checked explicit truncation comparison remains in ALGEBRA\_NAIVE\_COTANGENT\_NOTES\_20260928.md, group 993. Thus the original notation at 37528--37566 and 37736--37737 is valid and is retained.

For a finite presentation \(P=R[x_1,\ldots,x_n]\to S=P/I\), let \(E=I/I^2\), \(F=\bigoplus_iS\,dx_i\), and \(d:E\to F\) be the original differential. Smoothness says \(d\) is injective and \(\Omega_{S/R}=\operatorname{coker}(d)\) is finite projective. Choose a section \(s:\Omega_{S/R}\to F\) of the quotient \(p:F\to\Omega_{S/R}\). The map \(1-sp\) has image \(d(E)\), so define \(h:F\to E\) as \(d^{-1}(1-sp)\). Then \[
dh=1-sp,\qquad hd=1_E,\qquad ps=1_{\Omega_{S/R}}.
\] These equations give a homotopy equivalence between the original two-term complex and its degree-zero cokernel, with all chosen maps explicit. They do not give literal equality of the complexes. OCC-00815 therefore receives the minimal ``quasi-isomorphic'' qualification; the stronger homotopy statement is editorial evidence already supported by the earlier general comparison.

For arbitrary \(R\to R'\), put \(P'=R'[x_1,\ldots,x_n]\), \(S'=S\otimes_RR'\) and \(I'=\ker(P'\to S')\). The natural map \[
E\otimes_RR'\longrightarrow I'/(I')^2
\] is onto: tensor the sequence \(I\to P\to S\to0\) and then pass to conormal quotients. The top differential \(E\otimes_RR'\to F\otimes_RR'\) is split injective because the equations \(hd=1\) persist under tensoring. Its factorization through \(I'/(I')^2\) proves that the preceding surjection is also injective. Hence the conormal map is an isomorphism, the bottom differential is split injective, and its cokernel is \(\Omega_{S/R}\otimes_SS'\), a finite projective module. This proves base change without assuming flatness of \(S/R\) before the source establishes it. Localization follows by the earlier comparison; finite presentation also persists in each case.

The opening characteristic-\(p\) example keeps \(f=x^p+y^p\), with both derivatives zero. Since this polynomial is monic in \(x\), it is a nonzerodivisor in \(R[x,y]\); its conormal module is \(S\), and the differential \(S\to S\,dx\oplus S\,dy\) is zero. For \(R\ne0\), \(S\ne0\), so \(H_1=S\ne0\) although \(\Omega_{S/R}\) is free of rank two. Every nonempty fibre has dimension one. This verifies the stated failure of the criterion using only freeness of differentials.

\subsubsection{Standard smooth charts with exact Jacobian matrices}\label{algebra-proof-047-standard-smooth-charts-with-exact-jacobian-matrices}

Retain the source equations \(f_1,\ldots,f_c\in R[x_1,\ldots,x_n]\), \(n\ge c\), and \(S=P/(f_1,\ldots,f_c)\). With variables indexing rows and equations indexing columns, write the full derivative matrix as \[
J=\begin{pmatrix}A\\B\end{pmatrix},\quad
A_{ij}=\frac{\partial f_j}{\partial x_i}\ (1\le i,j\le c),\quad
B_{rj}=\frac{\partial f_j}{\partial x_{c+r}},\quad
\Delta=\det A.
\] If \(\Delta\) is a unit in \(S\), the map \(S^c\to I/I^2\to S^n\) has composite \(J\), and its projection onto the first \(c\) coordinates is the invertible matrix \(A\). The first arrow is already onto, so it is an isomorphism. The original differential has left inverse \[
(u,v)\longmapsto A^{-1}u=\Delta^{-1}\operatorname{adj}(A)u.
\] The quotient map to the actual remaining coordinates is \[
(u,v)\longmapsto v-BA^{-1}u,
\] with section \(v\mapsto(0,v)\). Its kernel consists exactly of \(J(A^{-1}u)\). Thus \(\Omega_{S/R}\) has the claimed basis \(dx_{c+1},\ldots,dx_n\), and no factor \(\Delta^{-1}\) or adjugate entry is suppressed.

For a localization at \(g\in S\), choose the source lift \(h\in P\) and add the equation \(hz-1\). Order the selected variables as \(x_1,\ldots,x_c,z\). The selected Jacobian matrix, with the same row convention, is \[
\begin{pmatrix}
A & (z\,\partial h/\partial x_i)_{i=1}^c\\
0 & h
\end{pmatrix}.
\] Its determinant is \(\Delta h\), which maps to the unit \(\Delta g\) in \(S_g\). The resulting presentation is \(P[z]/(f_1,\ldots,f_c,hz-1)\), with \(z\mapsto g^{-1}\), exactly as in the source. In Example 37339--37370 take \(h=\Delta\); the determinant is \(\Delta^2\), not merely \(\Delta\). Under a coefficient map \(R\to R'\), each formal derivative and each determinant maps coefficient by coefficient to the corresponding derivative or determinant. Base localization at an already invertible element of \(S\) preserves the same presentation.

The argument over an algebraically closed field at 37179--37233 uses the original coordinates \(x_i-a_i\). For each polynomial \(f\), expansion of its monomials at \(a=(a_i)\) gives \[
f(x)-f(a)-\sum_{i=1}^n\frac{\partial f}{\partial x_i}(a)(x_i-a_i)
\in(x_1-a_1,\ldots,x_n-a_n)^2.
\] Indeed in the full product expansion the constant term is \(f(a)\), the terms with exactly one coordinate factor are the displayed derivative terms, and every other term has at least two coordinate factors. For \(f\in I\), \(f(a)=0\). This proves the source diagram commutes. The split conormal injection makes the chosen \(f_i\) independent in the maximal ideal modulo its square. Extend them to a basis there and apply the original regular-local-ring lemma, 25731--25747. They form a regular sequence and generate the localized ideal by Nakayama. This establishes the source complete-intersection criterion at every maximal ideal and hence its local conclusion; the initial field extension descends it by the previously verified complete-intersection comparison.

\subsubsection{Empty fibres and the composition lifts}\label{algebra-proof-047-empty-fibres-and-the-composition-lifts}

The proof at 37326--37336 must restrict its dimension equalities to nonzero fibres or nonzero field algebras, as the definition of a relative global complete intersection does. For example, \[
R=\mathbf Z,\qquad S=\mathbf Z[z]/(2z-1)\cong\mathbf Z[1/2].
\] This is a standard smooth presentation with \(n=c=1\): its selected derivative is \(2\), a unit in \(S\). Its fibre over \((2)\) is the zero ring, with empty spectrum, so it is not a nonempty dimension-zero fibre. The theorem that standard smooth algebras are relative global complete intersections is unaffected. Two minimal qualifications restore the actual source definition. For a nonzero standard smooth algebra over a field, the preceding complete-intersection argument and the conormal rank \(c\) give dimension \(n-c\) on every component. These are exactly the fibres needed for the theorem.

For composition retain \[
S=R[x_1,\ldots,x_n]/(f_1,\ldots,f_c),\qquad
S'=S[y_1,\ldots,y_m]/(g_1,\ldots,g_d),
\] and the chosen lifts \(g'_j\in R[x_1,\ldots,x_n,y_1,\ldots,y_m]\). In the combined presentation order the selected variables as \[
x_1,\ldots,x_c,y_1,\ldots,y_d
\] before the other variables, leaving the equations ordered \(f_1,\ldots,f_c,g'_1,\ldots,g'_d\). The selected matrix is \[
\begin{pmatrix}
A&C\\0&D
\end{pmatrix},
\qquad
C_{ij}=\frac{\partial g'_j}{\partial x_i},\qquad
D_{\ell j}=\frac{\partial g'_j}{\partial y_\ell}.
\] All entries are then reduced in \(S'\). The \(y\)-derivatives reduce to the well-defined derivatives of \(g_j\in S[y_1,\ldots,y_m]\). In contrast, \(\partial/\partial x_i\) need not descend to \(S\), so the source's four displayed \(x\)-derivatives of \(g_j\) must use \(g'_j\). For example \(S=R[x]/(x)\) represents the zero coefficient both by \(0\) and by \(x\), whose \(x\)-derivatives differ. Thus this is a genuine type correction, not the harmless TeX ordering of a prime and subscript reported in OCC-00820.

The determinant is exactly \(\det(A)\det(D)\). One proof expands over permutations: the lower-left zero block forces the last \(d\) rows to use the last \(d\) columns, and the remaining permutation terms give the product, with no additional sign for the stated ordering. Both factors are units by the original hypotheses. If another lift replaces \(g'_j\) by \(g'_j+\sum_i a_{ij}f_i\), its \(x\)-column changes in \(S'\) by \(\sum_i a_{ij}\) times the corresponding \(f_i\)-columns; the \(y\)-column is unchanged because every \(f_i\) is independent of \(y\) and maps to zero. Column operations therefore leave the determinant unchanged. This proves composition with the original choices and also verifies the comparison between different lifts. Empty \(c\)- or \(d\)-blocks have determinant one.

\subsubsection{The local criterion and the correct splitting}\label{algebra-proof-047-the-local-criterion-and-the-correct-splitting}

The standard smooth cover at 37437--37508 follows from the original finite projective conormal module. At any prime, choose a subset of its given generators giving a basis of the residue vector space. Because it is locally free, the corresponding determinant is a unit on a neighbourhood; this gives the first finite principal cover. The localization conormal comparison and Huber presentation preserve the original algebra on each chart. The split injection from that now-free conormal module to the free differential module has rank \(c\) over every residue field, so some \(c\)-row minor is a unit locally. The finitely many corresponding minor opens cover. The preceding localization determinant \(\Delta g\) proves each further principal localization is standard smooth with the indicated coordinate ordering. This verifies every passage used to conclude smooth implies syntomic.

For OCC-00823, after making \(\Omega_{S/R}\) free on \(D(g)\), put \(E=(I/I^2)_g\), \(F=S_g^n\), \(C=\Omega_{S/R,g}\), and \(V=\operatorname{im}(d)\). The split sequence \[
0\longrightarrow V\longrightarrow F\longrightarrow C\longrightarrow0
\] makes \(V\) finite projective. It is the surjection \(E\to V\) that therefore has a right inverse \(t:V\to E\). The map \[
E\longrightarrow\ker(d),\qquad e\longmapsto e-t(d(e))
\] is a surjection and restricts to the identity on \(\ker(d)\). Hence \(H_1=\ker(d)\) is finite, as a direct summand of the finite module \(E\). This finiteness does not follow for an arbitrary kernel over a non-Noetherian ring. If \(H_{1,\mathfrak q}=0\), finitely many killing denominators, one per generator, yield a product \(g'\notin\mathfrak q\) with \(H_{1,gg'}=0\), proving the source local criterion. The map \(E\to F\) need not have a right inverse: for the polynomial algebra \(S=R[x]\), \(E=0\), \(F=S\,dx\ne0\) when \(R\ne0\). The correction names its actual image target.

Finite principal covers detect vanishing of \(H_1\) and finite projectivity of \(\Omega\); finite presentation also descends from the cover. This proves the target-locality assertion used in composition and in the product lemma. For a product \(S'\times S''\), the original idempotents \((1,0)\) and \((0,1)\) give the two open-and-closed charts and identify their localizations with the two factors, proving both implications including a zero factor.

The flat-fibre criterion at 37672--37706 applies the syntomic local criterion at the given prime. Its chosen \(g\) must and can satisfy \(g\notin\mathfrak q\); the report's added condition records that actual choice. On the resulting original relative-complete-intersection presentation, formal derivatives commute with passage to the residue field, and the same chosen minor is a unit at the original prime exactly when its fibre image is a unit at the fibre prime. This proves the criterion.

For flat base change, \(S_{\mathfrak q}\to S'_{\mathfrak q'}\) is flat and local, hence faithfully flat. Exactness of flat tensor products and the verified conormal comparison identify the first homology after base change. Faithful flatness detects its vanishing. The finite module \(\Omega_{S/R,\mathfrak q}\) becomes finite projective and is therefore finite projective by 19268--19288. This proves reflection in the source's stated scope.

The smooth lifting lemma retains each original standard presentation modulo the given ideal. Lift all its coefficients and invert the full lifted determinant by adjoining \(z\) and imposing \(z\Delta-1\). The resulting selected determinant is \(\Delta^2\), a unit; reduction is the original quotient algebra because \(\overline\Delta\) was already a unit. The reductions send \(z\) to \(\overline\Delta^{-1}\), and the inverse comparison is the original coefficient quotient map with that forced value of \(z\). The barred localizers in both source occurrences remain under the previously checked MC-STK-ERR-0446; its earlier report OCC-12103 is not counted again. No global gluing of the chart lifts is asserted.

\subsubsection{Smooth loci under arbitrary base change of a flat map}\label{algebra-proof-047-smooth-loci-under-arbitrary-base-change-of-a-flat-map}

The source at 37712--37744 assumes the base-change map is flat. A separate consequence allows every base change when the original map \(R\to S\) is flat and finitely presented. Let \(R\to R'\) be arbitrary, \(S'=S\otimes_RR'\), and \(\mathfrak q'\subset S'\) contract to \(\mathfrak q\subset S\). Let \(\mathfrak p'\subset R'\) and \(\mathfrak p\subset R\) be their base contractions. Then \[
S'\otimes_{R'}\kappa(\mathfrak p')
\cong
(S\otimes_R\kappa(\mathfrak p))\otimes_{\kappa(\mathfrak p)}\kappa(\mathfrak p').
\] The isomorphism sends \((s\otimes r')\otimes a\) to \((s\otimes1)\otimes r'a\), with inverse \((s\otimes b)\otimes a\mapsto(s\otimes1)\otimes ba\). Localization gives the corresponding primes on these fibres. Field extension is flat, so the already proved field-change criterion 37746--37758 says smoothness of these two fibre maps at the indicated primes is equivalent. Both original and base-changed maps are flat and finitely presented. Smoothness implies fibre smoothness by base change; the flat-fibre criterion proves the converse at each prime. Hence the smooth locus in \(\operatorname{Spec}(S')\) is exactly the inverse image of the smooth locus in \(\operatorname{Spec}(S)\).

The original-map flatness assumption cannot simply be omitted. Take \(R=k[t]\), \(S=R/(t)=k\), \(R'=k\). The map \(R\to S\) is not flat: tensoring the injection \(R\xrightarrow{t}R\) with \(S\) produces the zero map \(S\to S\), which is not injective. Since smooth maps are syntomic and therefore flat, its sole target point is not smooth. The base-changed map \(k\to k\) is smooth. Thus arbitrary base change can improve the locus for a nonflat original map. This preserves the source theorem while identifying the precise further scope.

\subsubsection{Intrinsic Jacobian defect and its full Fitting ideals}\label{algebra-proof-047-intrinsic-jacobian-defect-and-its-full-fitting-ideals}

For the original relative global complete-intersection presentation, use the full \(n\)-by-\(c\) Jacobian matrix \(J\) above without selecting a preferred minor. Put \(d=n-c\). At a prime \(\mathfrak q\subset S\), define \[
\sigma(\mathfrak q)=c-\operatorname{rank}_{\kappa(\mathfrak q)}J
=\dim_{\kappa(\mathfrak q)}(\Omega_{S/R}\otimes_S\kappa(\mathfrak q))-d.
\] The equality follows by taking the cokernel of the actual matrix over the residue field. Thus \(\sigma\ge0\); it also equals the dimension of the degree-one homology of the original two-term complex after tensoring with that residue field. This is not a claim that tensoring the original \(H_1\) alone computes that homology.

For \(s\ge1\), let \[
\mathcal J_s=\operatorname{Fitt}_{d+s-1}(\Omega_{S/R}).
\] By the source definition of Fitting ideals, More on Algebra 1355--1447, this is the ideal generated by every \((c-s+1)\)-by-\((c-s+1)\) minor of the full original matrix. For \(c-s+1\le0\) this means the unit ideal, and sizes exceeding the matrix dimensions generate zero. Over a residue field all these minors vanish exactly when the matrix rank is at most \(c-s\), which is exactly \(\sigma\ge s\). Hence \[
\{\sigma\ge s\}=V(\mathcal J_s),\qquad
\{\sigma=s\}=V(\mathcal J_s)\setminus V(\mathcal J_{s+1})\quad(s\ge1),
\] and \(\{\sigma=0\}=\operatorname{Spec}(S)\setminus V(\mathcal J_1)\) is precisely the smooth locus by the original Jacobian criterion. This proves upper semicontinuity and all locally closed defect strata.

The construction is independent of the chosen relative complete-intersection presentation. The integer \(d\) is the dimension of the nonempty fibre on each such chart, so two presentations of the same nonzero algebra have the same \(d\). The module \(\Omega_{S/R}\) and its Fitting ideals are intrinsic. For arbitrary \(R\to R'\), conormal and differential base change on these presentations give exactly the image matrix, while \(\Omega_{S'/R'}\cong\Omega_{S/R}\otimes_SS'\). Right exactness preserves its finite presentation, and each determinant commutes with the coefficient map. Thus, as actual ideals, not only as reduced closed subsets, \[
\mathcal J_s(S'/R')=\mathcal J_s(S/R)S'.
\] Ranks of matrices are preserved under the residue-field extension at any prime above \(\mathfrak q\), so \(\sigma(\mathfrak q')=\sigma(\mathfrak q)\). Empty spectra impose no extra rank convention. No radical is inserted in this ideal equality.

For a general syntomic \(R\to S\), use its finite principal cover by relative global complete intersections. Each chart has constant nonempty fibre dimension \(d\), so the function \(d\) is locally constant on the target and takes finitely many values by quasi-compactness. It gives a finite partition into open-and-closed subsets with rings \(S_d=e_dS\), where \(1=\sum_d e_d\) and \(e_de_{d'}=0\) for \(d\ne d'\). Define \(\sigma\) by the same intrinsic formula on each part and define the ideal with component \[
\mathcal J_sS_d=\operatorname{Fitt}_{d+s-1}(\Omega_{S_d/R}).
\] The chart calculation proves these definitions agree on overlaps. Equivalently, \(\mathcal J_s\) is the sum in \(S\) of these component ideals embedded by their original idempotents \(e_d\). All displayed locus and base-change conclusions persist: each base-changed chart has the same \(n,c,d\); the idempotents map to their original images; and finite-presentation Fitting ideals commute with the same base change. This gives intrinsic, possibly nonreduced defect subschemes for the full syntomic map.

This use of Fitting independence also exposed a bounded source-index typo outside the current chapter. At More on Algebra 1385, the added kernel vectors are written \(z'_{n+j'}\); the construction has first chosen \(n-k\) vectors and then adds \(n'\), so the added vectors must be \(z'_{n-k+j'}\), \(1\le j'\le n'\). To verify the comparison, append the actual \(n'\) coordinate unit vectors to the \(n-k\) original kernel vectors. The corresponding \((n+n'-k)\)-row minors are exactly the original \((n-k)\)-row minors, up to the row-selection signs. Conversely, expand any enlarged minor along the added-coordinate rows that it selects. Each remaining original-coordinate minor has size at least \(n-k\); expand it further to size \(n-k\) to see that every term belongs to the original minor ideal. This argument is for \(0\le k<n\). For \(k\ge n\), both Fitting ideals are the unit ideal: the added coordinate vectors give a unit minor of the required size, with the size-zero convention when needed. Both ideal inclusions follow. For two surjections \(\pi:A^n\to M\) and \(\rho:A^m\to M\), choose a lift \(U:A^m\to A^n\) with \(\pi U=\rho\). Then \((\pi,\rho)=(\pi,0)\circ T\), where \(T(a,b)=(a+Ub,b)\) has inverse \(T^{-1}(a,b)=(a-Ub,b)\). Interchanging the two surjections gives the other comparison with the common surjection \((\pi,\rho)\). Multilinearity of each determinant under \(T\) gives one ideal inclusion, and applying \(T^{-1}\) gives the reverse inclusion. This proves the comparison with the correct indices. The typo is recorded as a pending More on Algebra item, with no change to that chapter and no claim that it has been fully reviewed.

\subsubsection{Report decisions and propagation}\label{algebra-proof-047-report-decisions-and-propagation}

The twenty-one reports receive twelve one-operation corrections, five optional dispositions, three nondefect rejections, and one retained canonical-unit disposition. The harmless prime/subscript ordering is rejected alongside the two cotangent-notation reports. Slogan punctuation, enumeration punctuation, the implicit finite summation range and the phrase ``free rank'' remain optional. Two additional source groups propose the two nonempty-fibre qualifications and four required primes on the chosen lifts in \(x\)-derivatives. The valid \(y\)-derivatives of the original \(g_j\) remain unchanged.

The two consequence groups have zero translation operations: arbitrary base change of the smooth locus for flat finitely presented maps, and intrinsic Jacobian defect ideals for syntomic maps. The first keeps a counterexample outside its flatness scope; the second preserves all minors and their actual nonreduced ideals. The further More on Algebra index finding is separately pending and propagates into the proof just given, not into an unreviewed chapter edit. Source proofs and statements remain identifiable, and combined-results synthesis remains after completion of the core correction work.

\hypertarget{algebra_formal_smooth_notes_20260928}{}
\subsection{Formal smoothness: original lifts, exact splittings and nilpotent criteria}\label{algebra-proof-048-formal-smoothness-original-lifts-exact-splittings-and-nilpotent-criteria}

Authority: Stacks Project authors, commit a04446e57ec1fbc252a871afcec7752fb2807b14, algebra.tex:37795--38339, SHA256 FA8BB92E58A4F78A2BD01B3B6A4A87DE0A0D279F5DD90641B574DD5FBFFFA4F3. All source statements and original presentations remain identifiable. Complete arguments and extensions in this file are separate editorial material. The next source starts at 38340; only its opening through 38350 has been read. The bounded corpus query ``formally smooth lifting'' returned four routing records, three with PDF-only content and one with a TeX closure. They remain unread and are not mathematical evidence. No PDF fallback, exhaustive survey or novelty claim is made.

\subsubsection{Square-zero lifts and the original conormal criterion}\label{algebra-proof-048-square-zero-lifts-and-the-original-conormal-criterion}

Let \(R\to S\) be the original map. Formal smoothness requires a compatible lift \(S\to A\) of each \(R\)-algebra map \(S\to A/K\) whenever \(K^2=0\). For base change \(S'=R'\otimes_RS\), first obtain the compatible \(R\)-algebra lift \(\psi:S\to A\). With the specified map \(\eta:R'\to A\), the induced map is \[
R'\otimes_RS\longrightarrow A,\qquad r'\otimes s\longmapsto\eta(r')\psi(s).
\] The two maps agree on \(R\), so this map is well-defined; its reduction is the given map on \(S'\). For composition \(R\to S\to T\), first lift the restricted map \(S\to A/K\) over \(R\), then regard \(A\) as an \(S\)-algebra by that lift and lift \(T\to A/K\) over \(S\). The composite lift respects the original \(R\)-structure. For a polynomial algebra, including an arbitrary set of variables, choose the actual image lift of each variable and use its universal property.

Let \(P\to S=P/J\) be the original polynomial presentation. Formal smoothness gives an \(R\)-algebra section \(\sigma:S\to P/J^2\) by applying the definition to the square-zero ideal \(J/J^2\). Conversely, for a given lifting problem choose \(\psi:P\to A\) lifting \(P\to S\to A/K\). Then \(\psi(J)\subset K\) and \(\psi(J^2)=0\). The required lift is exactly \[
S\xrightarrow{\sigma}P/J^2\xrightarrow{\overline\psi}A.
\] The same proof works for any formally smooth \(P/R\), since it was only the existence of \(\psi\) that used polynomiality.

Put \(E=J/J^2\) and \(F=\Omega_{P/R}\otimes_PS\). The section criterion is equivalent to split exactness of \[
0\longrightarrow E\xrightarrow{d}F\xrightarrow{p}\Omega_{S/R}\longrightarrow0.
\] In the forward direction, reducing \(P\) modulo \(J^2\) does not change its differential module after tensoring with \(S\): every differential of a product in \(J^2\) becomes zero. A section of \(P/J^2\to S\) splits the square-zero conormal sequence, by the exact construction already verified in groups 903--904. Concretely, if \(D=P/J^2\) and \(\sigma:S\to D\) is the section, then \(a\mapsto a-\sigma(\bar a)\) is an \(R\)-derivation \(D\to J/J^2\): expanding the product leaves \(\sigma(\bar a)(b-\sigma(\bar b))+\sigma(\bar b)(a-\sigma(\bar a))\), since the product of two kernel elements is zero. Its induced map on \(F\) is a left inverse to \(d\).

For the reverse direction choose the source splitting \(u:\Omega_{S/R}\to F\). For each \(\lambda\in S\), choose \(x_\lambda\in P\) and \(f_\lambda\in J\) with \[
df_\lambda=dx_\lambda-u(d\lambda)\quad\hbox{in }F.
\] Injectivity of \(d:E\to F\) means \(h\in J\) with \(dh=0\) in \(F\) belongs to \(J^2\). It follows that the class \[
s(\lambda)=x_\lambda-f_\lambda\bmod J^2
\] is independent of both choices: the difference from another choice belongs to \(J\) and has differential zero. Its additive defect \[
x_\lambda+x_\mu-x_{\lambda+\mu}-f_\lambda-f_\mu+f_{\lambda+\mu}
\] also belongs to \(J\) and has differential zero. For multiplication, use the exact Leibniz relation \[
\lambda\,d\mu+\mu\,d\lambda-d(\lambda\mu)=0.
\] The actual polynomial-ring element to test is \[
H=x_\lambda x_\mu-x_{\lambda\mu}
-x_\mu f_\lambda-x_\lambda f_\mu+f_{\lambda\mu}\in J.
\] Its differential in \(F\) is \[
\mu(dx_\lambda-df_\lambda)+\lambda(dx_\mu-df_\mu)
-dx_{\lambda\mu}+df_{\lambda\mu}=0.
\] Terms with coefficients \(f_\lambda,f_\mu\) vanish in \(F\). Thus \(H\in J^2\). Adding the retained term \(f_\lambda f_\mu\in J^2\) proves \(s(\lambda)s(\mu)=s(\lambda\mu)\). This corrects the original ideal-membership formula, which uses \(\lambda,\mu\in S\) as though they were elements of \(P\): use their actual chosen lifts \(x_\lambda,x_\mu\) in that formula. For \(r\in R\), let \(\lambda\) be its image in \(S\). Then \(r-x_\lambda+f_\lambda\in J\) and its differential is zero, so \(s(\lambda)=r\bmod J^2\). This also proves preservation of \(1\) and the specified \(R\)-structure, even when \(R\to S\) is not injective.

The three original tensor-base errors at 37945, 38020 and 38138 retain MC-STK-ERR-0449, with tensor product over \(P\). The Leibniz product differential retains MC-STK-ERR-0450. Tensoring over \(R\) generally gives a different object: for \(R=k\), \(P=k[x]\), \(S=k=P/(x)\), the module \(\Omega_{P/R}\otimes_RS\) is \(k[x]\,dx\), whereas \(\Omega_{P/R}\otimes_PS=k\,dx\). The unscripted tensors in the other conormal formulas have the evident base \(P\) from their context and are retained as optional clarification. The source explicitly defines \(H_1(L_{S/R})\) at 35152--35176, so the repeated complaint about that notation is a nondefect.

\subsubsection{Projective differentials and the direction of the splittings}\label{algebra-proof-048-projective-differentials-and-the-direction-of-the-splittings}

The section and split-conormal criteria imply that formal smoothness is equivalent to \(H_1(\mathrm{NL}_{S/R})=0\) and projectivity of \(\Omega_{S/R}\), with no finite generation assumed. Use the canonical polynomial presentation: its \(F\) is free on the original variables, and a split cokernel is projective. Conversely projectivity of the cokernel and vanishing of the kernel make its conormal sequence split, giving the section just constructed.

For a formally smooth \(P/R\) and a surjection \(P\to S\), a chosen splitting of the polynomial presentation complex of \(P\) gives a homotopy retraction onto \(\Omega_{P/R}\) in degree zero. Tensoring these actual homotopy maps over \(P\) with \(S\) retains their identities; no flatness of \(S/P\) is needed. Thus \(H_1(\mathrm{NL}_{P/R}\otimes_PS)=0\). The original Jacobi--Zariski sequence then identifies the obstruction to injectivity of \(J/J^2\to\Omega_{P/R}\otimes_PS\) with \(H_1(L_{S/R})\). Its middle term is projective because \(\Omega_{P/R}\) is projective over \(P\). A split conormal sequence therefore gives projectivity of \(\Omega_{S/R}\) and the required vanishing. This verifies all directions of the source proposition without discarding the original tensor base.

For \(A\to B\to C\) with \(B\to C\) formally smooth, the Jacobi--Zariski sequence and \(H_1(L_{C/B})=0\) give the initial injection in \[
0\longrightarrow\Omega_{B/A}\otimes_BC\longrightarrow
\Omega_{C/A}\longrightarrow\Omega_{C/B}\longrightarrow0.
\] Projectivity of its last term gives the splitting. If instead \(A\to C\) is formally smooth and \(B\to C\) is onto with kernel \(J\), lift the identity of \(C\) to an \(A\)-algebra section \(C\to B/J^2\). The square-zero conormal construction above gives the source's split sequence even though \(B/A\) is not assumed formally smooth.

Now take the source hypotheses \(A\to C\) surjective, \(A\to B\) formally smooth, \(I=\ker(A\to C)\), \(J=\ker(B\to C)\). The lifting diagram supplies \[
\sigma:B\longrightarrow A/I^2
\] both lifting the given map \(B\to C=A/I\) and commuting with the map from \(A\). The first condition is necessary to ensure \(\sigma(J)\subset I/I^2\); it is made explicit in the minimal source correction. Its restriction gives a map \(r:J/J^2\to I/I^2\), since \(\sigma(J^2)=0\). If \(\alpha:I/I^2\to J/J^2\) is induced by \(A\to B\), then \[
r\alpha=1_{I/I^2}.
\] This is a left inverse, not a two-sided inverse. Indeed \(A=C=k\), \(B=k[x]\), with \(x\mapsto0\), gives \(I/I^2=0\), \(J/J^2\cong k\), and nonzero cokernel \(\Omega_{B/A}\otimes_BC\cong k\). The original split short exact sequence is correct; the ``inverse'' wording in its proof must retain its direction.

The two later uses of ``section'' in the square-zero base-lifting proof at 38142 and 38163 have the same direction: the map from the differential module to the conormal module is a left inverse to the differential. With \(d:E\to F\) and a lifted map \(r_0:F\to E\), put \[
\Delta=1_E-r_0d.
\] In the source setting \(\Delta(E)\subset IE\) and \(I^2=0\), so \(\Delta^2=0\). The correction can be chosen explicitly as \(\delta=\Delta r_0\). Then \[
(r_0+\delta)d=(1_E+\Delta)(1_E-\Delta)=1_E.
\] This supplies the exact corrected map asserted by the source, and proves its direction without an unspecified calculation.

\subsubsection{Nilpotent base ideals and the exact conormal obstruction}\label{algebra-proof-048-nilpotent-base-ideals-and-the-exact-conormal-obstruction}

Here is a separate extension of the source square-zero theorem, using the original base-change obstruction of group 995. Let \(R\to S=P/J\), where \(P\) is any polynomial \(R\)-algebra, with any set of variables. Let \(I\subset R\) satisfy \(I^N=0\) for some integer \(N\ge1\). Write \[
\bar R=R/I,\quad \bar P=P/IP,\quad \bar S=S/IS,\quad
\bar J=(J+IP)/IP,\quad E=J/J^2,\quad F=\Omega_{P/R}\otimes_PS.
\] The original conormal comparison is the surjection \[
E/IE=J/(IJ+J^2)\longrightarrow
\bar J/\bar J^2=J/(J^2+(J\cap IP)).
\] Its exact kernel is \[
\mathcal O_I(P,S)=\frac{J^2+(J\cap IP)}{J^2+IJ}.
\] The long exact Tor sequence of \(0\to J\to P\to S\to0\), using \(R\)-flatness of the actual polynomial ring \(P\), identifies \[
\operatorname{Tor}_1^R(S,R/I)\cong (J\cap IP)/IJ.
\] The map from this module to \(E/IE\) sends the original class of \(j\) to its class modulo \(IJ+J^2\); its image is exactly \(\mathcal O_I(P,S)\). Thus the obstruction is an image of Tor, not automatically the whole Tor group.

\textbf{Claim.} The original map \(R\to S\) is formally smooth if and only if \(\bar R\to\bar S\) is formally smooth and \(\mathcal O_I(P,S)=0\).

If \(R\to S\) is formally smooth, base change gives formal smoothness of \(\bar S/\bar R\). The map \(d:E\to F\) has a left inverse, so \(E/IE\to F/IF\) is still split injective. It factors through the displayed surjection to \(\bar J/\bar J^2\), which is therefore also injective. Hence \(\mathcal O_I=0\).

Conversely, assume these two conditions. The conormal comparison is now an isomorphism. The split conormal criterion over \(\bar R\) gives a left inverse \[
\bar r:F/IF\longrightarrow E/IE
\] to the reduced differential. Since \(F\) is the free \(S\)-module on the original polynomial variables, lift the image of each basis vector to obtain an \(S\)-linear map \(r_0:F\to E\). Define \(\Delta=1_E-r_0d\). Reduction modulo \(I\) is zero, so \(\Delta(E)\subset IE\). \(S\)-linearity then gives \(\Delta^q(E)\subset I^qE\) for every \(q\), by induction. In particular \(\Delta^N=0\). The finite formula \[
r=(1_E+\Delta+\Delta^2+\cdots+\Delta^{N-1})r_0
\] satisfies \[
rd=(1_E+\Delta+\cdots+\Delta^{N-1})(1_E-\Delta)=1_E-\Delta^N=1_E.
\] Thus the original conormal sequence splits, and the proved criterion gives formal smoothness of \(R\to S\). This argument needs no Noetherianity, no finite presentation and no finite rank for \(F\). It retains every term of the finite inverse and the original ideal power \(I^N\). Because the conclusion is intrinsic, vanishing of the obstruction under the stated reduced formal-smoothness hypothesis does not depend on the chosen polynomial presentation.

In particular, \(\operatorname{Tor}_1^R(S,R/I)=0\) is a sufficient replacement for the source's global flatness assumption, for every nilpotent \(I\), not only \(I^2=0\). The exact necessary condition is the vanishing of its displayed image. The example below proves this distinction can be strict.

\subsubsection{All lifts, finite relation errors and the finite-presentation boundary}\label{algebra-proof-048-all-lifts-finite-relation-errors-and-the-finite-presentation-boundary}

Fix an \(R\)-algebra map \(\varphi:S\to A/K\) with \(K^2=0\), and suppose it has one lift \(\psi:S\to A\). The \(S\)-action on \(K\) is induced by \(\varphi\): a choice of representative in \(A\) does not matter because two representatives differ by an element annihilating \(K\). For two lifts \(\psi,\psi'\), set \(D=\psi'-\psi\). Additivity and vanishing on \(R\) hold, and \[
D(st)=\varphi(s)D(t)+\varphi(t)D(s)
\] because the omitted product \(D(s)D(t)\) is exactly zero in \(K^2\). Conversely, for each such derivation \(D:S\to K\), the map \(\psi+D\) is multiplicative by this same equation, preserves \(R\), and preserves \(1\) since \(D(1)=D(1)+D(1)\) implies \(D(1)=0\). Consequently the lift set is a torsor under \[
\operatorname{Der}_R(S,K)=\operatorname{Hom}_S(\Omega_{S/R},K):
\] the action is addition, it is free, and every other lift differs by exactly one such map.

If \(K^N=0\), every formally smooth \(S/R\) has a lift through \(A/K\to A\). Inductively lift from \(A/K^j\) to \(A/K^{j+1}\), for \(j=1,\ldots,N-1\); the kernel \(K^j/K^{j+1}\) is square zero since \(2j\ge j+1\). Each fibre of this successive lifting problem is the torsor just described, with coefficient module \(K^j/K^{j+1}\). Its \(S\)-action factors through the original \(\varphi\), since \(K\cdot K^j\subset K^{j+1}\). For \(N=1\) the original map is already the required lift. No finite presentation is needed for this nilpotent-ideal statement.

For the source's stronger locally nilpotent-ideal theorem, keep a finite presentation \[
S=R[x_1,\ldots,x_n]/(f_1,\ldots,f_m)
\] of the smooth algebra and choose any original coordinate lifts \(a_i\in A\) of the given \(\varphi(x_i)\). Let \[
e_j=f_j(a_1,\ldots,a_n)\in K,\qquad Q=(e_1,\ldots,e_m)\subset A.
\] If each \(e_j^{b_j}=0\), then \[
Q^{\,1+\sum_{j=1}^m(b_j-1)}=0.
\] Indeed a monomial of that total degree in the \(m\) generators has some exponent at least \(b_j\); otherwise its total degree is at most \(\sum_j(b_j-1)\). Thus every product generator of this ideal power vanishes. If \(m=0\), take \(Q=0\) and exponent \(1\).

Evaluation at the actual \(a_i\) gives an \(R\)-algebra map \(S\to A/Q\), since every original relation vanishes there. Smoothness implies formal smoothness, and the finite nilpotent lifting just proved gives a lift \(S\to A\) whose coordinate corrections lie in \(Q\). Because \(Q\subset K\), it lifts the original \(\varphi\). This proves the source theorem with an explicit finite ideal and exponent while retaining all original equations and coordinates. It also proves the sharper assertion for an arbitrary ideal \(K\) whenever these particular relation errors are nilpotent. It is a supplementary argument; the author's Noetherian-approximation proof remains identifiable.

The finite-presentation distinction is real. Let \(k\) be a field and, for \(n\ge1\), put \[
R_n=k[t_n]/(t_n^{2^n}),\qquad
R_n\longrightarrow R_{n+1},\quad t_n\longmapsto t_{n+1}^{\,2}.
\] These maps are injective: the basis \(1,t_n,\ldots,t_n^{2^n-1}\) maps to distinct monomials \(1,t_{n+1}^{2},\ldots,t_{n+1}^{2^{n+1}-2}\) in the next basis. Let \(R=\varinjlim R_n\) using these maps and \(J=(t_1,t_2,\ldots)\subset R\). Every element of \(J\) belongs to \((t_N)\subset R_N\) for some sufficiently large \(N\), and is therefore nilpotent. Every finite set of such elements generates a nilpotent ideal by the preceding finite-exponent argument. But \(J\) is not nilpotent: given \(q\ge1\), choose \(2^n>q\); the element \(t_n^q\ne0\) survives every injective transition and lies in \(J^q\). Also \(J=J^2\), because \(t_n=t_{n+1}^2\), and \(J\ne0\).

The ring \(R\) is local with maximal ideal \(J\). Its quotient is \(k\), and every element with nonzero constant term is \(a(1+u)\), where \(a\in k^\times\), \(u\in J\) is nilpotent; the finite geometric sum in \(-u\) gives its inverse. Let \(S=R/J=k\). This quotient is formally smooth, with unique lifts in all square-zero problems: for any compatible \(R\to A\) and \(S\to A/K\), the image of \(J\) lies in \(K\), but \(J=J^2\) forces that image into \(K^2=0\). The original map \(R\to A\) therefore factors uniquely through \(S\).

Nevertheless the diagram with target \(A=R\), ideal \(K=J\), bottom map \(R\to R\) the identity, and top map \(S\to R/J\) the identity has no compatible lift. Such a lift would make the identity map of \(R\) factor through its quotient by the nonzero \(J\). Here \(K\) is locally nilpotent, so formal smoothness alone does not imply the source's stronger lifting conclusion. Moreover \(S\) is a finite type \(R\)-algebra but is not finitely presented: if the quotient kernel \(J\) were finitely generated, \(J^2=J\) and Nakayama in the local ring \(R\) would give \(J=0\), a contradiction.

This same example proves the exact obstruction criterion is sharper than Tor vanishing and flatness. Set \(I=(t_1)\subset R\), so \(I^2=0\). In the polynomial presentation with zero variables \(P=R\to S\), one has \(E=J/J^2=0\), hence \(\mathcal O_I(P,S)=0\). But \[
\operatorname{Tor}_1^R(S,R/I)\cong I/IJ\ne0.
\] The class of \(t_1\) is nonzero: if \(t_1\in IJ\), it can be written \(t_1=t_1v\) for one \(v\in J\); then \(1-v\) is a unit and forces \(t_1=0\), contrary to injectivity of the direct system. In particular \(S\) is not \(R\)-flat, since flatness would force this Tor group to vanish. The reduced map \(R/I\to S\) remains formally smooth by the already proved base-change property. Thus the necessary image obstruction vanishes although both proposed stronger requirements fail.

\subsubsection{Finite witnesses, descent and report accounting}\label{algebra-proof-048-finite-witnesses-descent-and-report-accounting}

For the source Noetherian approximation, retain every \(f_j\), every chosen section lift \(h_i\), and the actual polynomial identities \[
x_i-h_i=\sum_j b_{ij}f_j,\qquad
f_j(h_1,\ldots,h_n)=\sum_{j_1,j_2}a^{(j)}_{j_1j_2}f_{j_1}f_{j_2}.
\] The witnesses on the right of the second identity are chosen for each original \(j\); the superscript here records that dependence in the editorial argument. Adjoin all coefficients of this finite list of polynomials to \(\mathbf Z\) inside \(R\). The first identities make the resulting \(\sigma_0\) a right inverse; the second identities make it an algebra map on the original quotient. Hence the resulting algebra is finitely presented and formally smooth, therefore smooth. For a filtered colimit \(A=\varinjlim A_i\), this finite type \(\mathbf Z\)-subring is finitely presented over the Noetherian ring \(\mathbf Z\); its map to \(A\) factors through one \(A_i\), and base change of its smooth algebra gives the required \(B_i\).

For faithfully flat \(R\to R'\), the original kernel \(J\), its square and its conormal sequence base change exactly, by flatness. Faithful flatness detects the initial injection. The differential module becomes projective after this base change, so the previously proved descent theorem for arbitrary projective modules applies, not merely its finite version. The descended conormal sequence splits and the formal-smoothness criterion applies without adding finite-presentation assumptions.

Sixteen reports receive twelve decision groups: five minimal grammar operations, four optional preferences, one nondefect rejection and two already-composed units. Four reports share the three tensor-base operations of MC-STK-ERR-0449; two share the single Leibniz operation of MC-STK-ERR-0450. Two additional groups propose two actual coefficient lifts in the polynomial-ring membership formula and four repairs specifying the commuting lift and left-inverse directions. Two consequence groups record the nilpotent conormal-obstruction criterion and the lifting torsors, finite-error bound and exact finite-presentation counterexample. Both have zero source operations. The prior More on Algebra findings at 1385 and 8502/8504 remain separately pending. Broader synthesis and any short paper remain after completion of the core correction work.

\hypertarget{algebra_smooth_fields_notes_20260928}{}
\subsection{Smooth differentials, section completions and field criteria}\label{algebra-proof-049-smooth-differentials-section-completions-and-field-criteria}

Source: Stacks Project authors, authority commit a04446e57ec1fbc252a871afcec7752fb2807b14, algebra.tex:38340--38932, SHA256 FA8BB92E58A4F78A2BD01B3B6A4A87DE0A0D279F5DD90641B574DD5FBFFFA4F3. The entire interval and all twenty received reports were read. All supplementary arguments and extensions below remain identifiable editorial material; they do not replace the translated source. The next section starts at 38933; only its opening through 38943 is read. The bounded corpus query ``smooth conormal completion'' returned four unread routing records. One PDF-primary record and one TeX record concern the same Bhatt--Hochster--Ma title; the other two are local TeX records. No source-content reading, independent-work count, exhaustive literature survey or novelty claim follows from these hits.

\subsubsection{Exact sequences and the completion along the original section}\label{algebra-proof-049-exact-sequences-and-the-completion-along-the-original-section}

For the source maps \(A\to B\to C\) with \(B\to C\) smooth, the exact sequence at 38350 is in fact split: \[
0\longrightarrow C\otimes_B\Omega_{B/A}\longrightarrow
\Omega_{C/A}\longrightarrow\Omega_{C/B}\longrightarrow0.
\] The earlier Jacobi--Zariski comparison gives the initial injection because \(H_1(L_{C/B})=0\); the final module is finite projective over \(C\), so a \(C\)-linear section of the last surjection exists. The tensor swap from the earlier statement is the exact map \(c\otimes\omega\mapsto\omega\otimes c\), with its inverse. Thus the weaker statement ``exact'' is valid and needs no replacement.

Likewise, when \(A\to C\) is onto, \(A\to B\) is smooth, \(I=\ker(A\to C)\) and \(J=\ker(B\to C)\), the original sequence \[
0\longrightarrow I/I^2\longrightarrow J/J^2
\longrightarrow\Omega_{B/A}\otimes_B C\longrightarrow0
\] is split. Formal smoothness gives a map \(B\to A/I^2\) reducing to the given \(B\to C\) and commuting with \(A\). Restrict it to \(J\); it kills \(J^2\) and gives the left inverse \(J/J^2\to I/I^2\) proved in group 1083. For the other source sequence with \(A\to C\) smooth and \(B\to C\) onto, lift \(1_C\) to an \(A\)-algebra section \(C\to B/J^2\). The difference \(b-\sigma(\bar b)\) is the original square-zero derivation and gives a left inverse on the conormal module. These are the exact maps in all three source statements.

Keep the source smooth map \(\varphi:R\to S\), its specified retraction \(\sigma:S\to R\), and \(I=\ker\sigma\). Put \(M=I/I^2\). The split conormal sequence gives the canonical isomorphism \[
M\xrightarrow{\ d\ }\Omega_{S/R}\otimes_{S,\sigma}R,
\] so \(M\) is finite projective over \(R\). The following extends the source's free-module case without choosing new source coordinates.

Choose an \(R\)-linear lift \(\ell:M\to I\) of the quotient map; it exists by projectivity of \(M\). Let \[
T=\operatorname{Sym}_R(M),\qquad J=\bigoplus_{q\ge1}\operatorname{Sym}_R^q(M).
\] The map \(\ell\) induces an augmented \(R\)-algebra map \(\Psi:T\to S\). For every integer \(n\ge1\), it induces \[
\Psi_n:T/J^n\longrightarrow S/I^n.
\] These maps are isomorphisms, as follows.

For \(n=1\) both rings are the specified \(R\). For \(n=2\) the retraction identifies \(S/I^2\) with the square-zero algebra \(R\oplus M\), and \(\Psi_2\) is the identity on both summands. For any \(q\ge1\), products of elements of \(I\) generate \(I^q/I^{q+1}\). Replacing each factor modulo \(I^2\) by \(\ell(m)\) changes its product by an element of \(I^{q+1}\). Coefficients may be reduced modulo \(I\), so the original product map \[
\operatorname{Sym}_R^q(M)\longrightarrow I^q/I^{q+1}
\] is surjective. Induction on \(n\) therefore makes each \(\Psi_n\) surjective.

Suppose \(n\ge3\) and \(\Psi_{n-1}\) is invertible. The ideal \(J^{n-1}/J^n\) is square zero, since \(2(n-1)\ge n\). Formal smoothness lifts \[
S\longrightarrow S/I^{n-1}\xrightarrow{\Psi_{n-1}^{-1}}T/J^{n-1}
\] to an \(R\)-algebra map \(\tau:S\to T/J^n\). Its augmentation is \(\sigma\), so \(\tau(I)\subset J/J^n\), whence \(\tau(I^n)=0\). Write \(\bar\tau:S/I^n\to T/J^n\).

The endomorphism \(U=\bar\tau\Psi_n\) of \(T/J^n\) is the identity modulo \(J^{n-1}\). Write \(U=1+D\). It fixes \(R\) and takes each \(m\in M\) to \(m+D(m)\), with \(D(m)\in J^{n-1}/J^n\). Expansion of products shows that \(D\) kills \(J^2/J^n\): every changed term has degree at least \(n\). Also \(D\) is an \(R\)-derivation, since its product defect \(D(a)D(b)\) lies in \(J^{2n-2}/J^n=0\). As \(n-1\ge2\), \(D^2=0\). Consequently \(1-D\) is a ring map and is both inverses to \(1+D\), by direct multiplication. Thus \(U\) is invertible; \(\Psi_n\) is injective as well as surjective and hence is an isomorphism. Its inverse is uniquely determined and is compatible with the inverse at every earlier quotient.

Taking inverse limits gives the filtered \(R\)-algebra isomorphism \[
\widehat\Psi:\prod_{q\ge0}\operatorname{Sym}_R^q(M)
=\varprojlim_n T/J^n
\xrightarrow{\ \sim\ }\varprojlim_n S/I^n.
\] The product algebra uses the full convolution product: in degree \(q\), multiply by summing all products of degrees \(a,b\) with \(a+b=q\). This isomorphism depends on the chosen lift \(\ell\). Its associated graded isomorphism \[
\operatorname{Sym}_R(M)\xrightarrow{\ \sim\ }\bigoplus_{q\ge0}I^q/I^{q+1}
\] is the canonical original product map and is independent of that choice, since changing a lift by an element of \(I^2\) changes a degree-\(q\) product only modulo \(I^{q+1}\). If \(M\) is free on the classes of the source's \(f_1,\ldots,f_d\), this is precisely its power-series map \(x_i\mapsto f_i\). It also supplies the source's omitted finite-order inverse computation, with the negative corrections retained.

Every \(\operatorname{Sym}_R^q(M)\) is finite projective: choose a finite projective complement \(N\) with \(M\oplus N\cong R^r\); the degree-\(q\) decomposition of \(\operatorname{Sym}(M\oplus N)\) exhibits \(\operatorname{Sym}^q(M)\) as a direct summand of the finite free \(\operatorname{Sym}^q(R^r)\). Thus each \(S/I^n\) is finite projective over \(R\).

For every map \(R\to R'\), let \(S'=S\otimes_RR'\), with its induced retraction and kernel \(I'\). The split sequence \(0\to I\to S\to R\to0\) identifies \(I'\) with \(I\otimes_RR'\). Since \(S/I^n\) is \(R\)-flat, tensoring \(0\to I^n\to S\to S/I^n\to0\) remains exact. The image of \(I^n\otimes_RR'\) is exactly \((I')^n\): both are generated by the same original \(n\)-fold products and all their coefficients. Hence \[
(S/I^n)\otimes_RR'\cong S'/(I')^n,\qquad
I^n\otimes_RR'\cong(I')^n.
\] Symmetric powers commute with this base change by their explicit tensor-quotient definition. Thus all finite-order comparisons and all graded pieces commute with arbitrary base change. The assertion about completions uses the inverse limit of these finite-order base changes; it does not identify ordinary tensor product of an infinite product with that inverse limit.

\subsubsection{The regular-local hypotheses and the field smoothness criterion}\label{algebra-proof-049-the-regular-local-hypotheses-and-the-field-smoothness-criterion}

The short exact sequence at 38546--38550 is valid in its actual setting, contrary to OCC-00855. Put \(P=k[x_1,\ldots,x_n]\), \(\mathfrak m''=(x_1,\ldots,x_n)\), \(A=P_{\mathfrak m''}\), and \(B=S_{\mathfrak m}\). Both \(A\) and \(B\) are regular local, with dimensions \(n\) and \(d\). The earlier lemma at 25749--25777 shows that the kernel \(I_A\) is generated by \(c=n-d\) members of a regular system of parameters of \(A\). Here is its exact argument in these objects.

The kernel of \(\mathfrak m_A/\mathfrak m_A^2\to\mathfrak m_B/\mathfrak m_B^2\) has dimension \(n-d\). Choose \(f_1,\ldots,f_c\in I_A\) mapping to a basis of that kernel and extend their classes to a basis of \(\mathfrak m_A/\mathfrak m_A^2\). The corresponding minimal generators form a regular sequence, and \(A/(f_1,\ldots,f_c)\) is regular local of dimension \(d\), by 25731--25747. Its surjection onto \(B\) has zero kernel. Indeed both are domains; a nonzero prime kernel would let any chain above it be extended by the zero prime, reducing quotient dimension strictly below \(d\). Consequently \(I_A=(f_1,\ldots,f_c)\). Their images in \(I_A/\mathfrak m_A I_A\) generate, and any linear relation there would yield a relation between their independent images in \(\mathfrak m_A/\mathfrak m_A^2\). They are a basis, and the conormal map is injective. Localization at \(\mathfrak m''\) does not change the original \(k\)-vector spaces \(I/\mathfrak m'' I\) and \(\mathfrak m''/(\mathfrak m'')^2\). This proves the displayed short exact sequence with precisely the source hypotheses. The general counterexample of a singular quotient is irrelevant here.

The original construction of \(S'\) and its surjection to \(S\) remains intact. The regular-local argument proves \(S'_{\mathfrak m'}\to S_{\mathfrak m}\) is an isomorphism. Its finite kernel is killed by one element outside \(\mathfrak m'\): take a product of annihilating denominators for a finite generating set. This gives the indicated principal localization. The already composed prime on \(S'\) and the prime on \(g'\) are retained. Writing \(S_{g'}\) denotes localization at the image of \(g'\) under the specified map \(S'\to S\), a valid convention which does not require a new correction.

For an algebraically closed \(k\), the residue field of a maximal ideal of a finite type \(k\)-algebra is \(k\). The specified quotient map therefore has a \(k\)-algebra section, giving \(\mathfrak m/\mathfrak m^2\cong\Omega_{S/k}\otimes_S k\). Thus regularity gives exactly the rank condition on the original conormal matrix, and a nonvanishing \(c\)-minor gives a standard smooth neighbourhood. Conversely a standard smooth neighbourhood has differential rank equal to its constant local dimension, so its maximal local ring is regular.

For any field \(k\) and original prime \(\mathfrak q\), put \(d=\dim_x\operatorname{Spec}S\), preserving point dimension rather than replacing it by \(\dim S_{\mathfrak q}\). If the differential fibre dimension is at most \(d\), finite generation and Nakayama give a principal neighbourhood on which \(\Omega_{S/k}\) has at most \(d\) generators. Choose an embedding \(\kappa(\mathfrak q)\to K\) into an algebraically closed field. The map \[
K\otimes_kS\longrightarrow K,\qquad a\otimes s\longmapsto a\,\psi(\bar s)
\] is onto and defines the specified maximal ideal \(\mathfrak m\). Its contraction is \(\mathfrak q\). Point dimension is preserved by this field extension, and at this closed point equals the local dimension. The base-changed differential module therefore satisfies the algebraically closed criterion. Smoothness descends along the flat base change at this prime by 37713--37744. Regularity descends along the resulting faithfully flat local map as in the source's earlier regularity theorem. Conversely a standard smooth neighbourhood has precisely \(d\) independent differentials. This proves all three original conditions and the regularity conclusion, keeping the residue field and point dimension distinct.

\subsubsection{Residue-field sections without the Noetherian restriction}\label{algebra-proof-049-residue-field-sections-without-the-noetherian-restriction}

The source proof at 38686--38737 does not use Noetherianity. More generally, let \(A\) be any local \(k\)-algebra, with maximal ideal \(\mathfrak m\) and residue field \(\kappa\). Suppose there is a possibly infinite transcendence basis \(T\) such that \(\kappa\) is separable algebraic over \(F=k(T)\). Then the original conormal sequence is split exact: \[
0\longrightarrow\mathfrak m/\mathfrak m^2
\longrightarrow\Omega_{A/k}\otimes_A\kappa
\longrightarrow\Omega_{\kappa/k}\longrightarrow0.
\] Neither Noetherianity of \(A\) nor finite generation of \(\kappa/k\) is required under this stated separating-basis hypothesis.

Set \(D=A/\mathfrak m^2\) and \(\mathfrak n=\mathfrak m/\mathfrak m^2\). Choose an actual lift \(x_t\in D\) of every \(t\in T\). Every nonzero polynomial in a finite subset of \(T\) has nonzero residue on these lifts, and hence is a unit in the local ring \(D\). Evaluation therefore defines a \(k\)-algebra map \(F\to D\) reducing to the given inclusion in \(\kappa\).

For a finite separable intermediate extension \(E/F\), choose a primitive element \(\bar y\) with monic separable minimal polynomial \(P\in F[Y]\), and any lift \(y\in D\). The exact correction is \[
y^*=y-\frac{P^\varphi(y)}{(P')^\varphi(y)}.
\] The denominator is a unit, and \(P^\varphi(y)\in\mathfrak n\). Since \(\mathfrak n^2=0\), the full polynomial expansion is \[
P^\varphi(y+\delta)=P^\varphi(y)+(P')^\varphi(y)\delta
\quad(\delta\in\mathfrak n),
\] so \(P^\varphi(y^*)=0\). If two roots have the same residue, their difference \(\delta\in\mathfrak n\) satisfies the same equation with a unit derivative, so \(\delta=0\). This gives the unique \(F\)-algebra map \(E\to D\) lifting the inclusion of \(E\) in \(\kappa\). Uniqueness makes these maps compatible on all finite separable subextensions. Their union defines a \(k\)-algebra section \(s:\kappa\to D\).

The exact inverse maps for the resulting square-zero decomposition are \[
D\longrightarrow\kappa\oplus\mathfrak n,\quad
a\longmapsto(\bar a,a-s(\bar a)),\qquad
(\lambda,u)\longmapsto s(\lambda)+u,
\] where \((\lambda,u)(\mu,v)=(\lambda\mu,\lambda v+\mu u)\). Differentiating products of elements of \(\mathfrak m\) shows directly that \(\Omega_{A/k}\otimes_A\kappa\cong\Omega_{D/k}\otimes_D\kappa\). Under this identification the derivation \[
D\longrightarrow\mathfrak n\oplus\Omega_{\kappa/k},\qquad
a\longmapsto(a-s(\bar a),d\bar a)
\] induces an isomorphism of \(\kappa\)-modules. Its inverse sends \(u\in\mathfrak n\) to \(du\), and sends \(d\lambda\) to \(d(s(\lambda))\); these maps respect the derivation relations because \(\mathfrak n^2=0\). The two composites are the identity on \(du\) and \(d(s(\lambda))\), which generate. This proves the split exact sequence.

Finally \(\Omega_{\kappa/k}\) is free on the original \(dt\), \(t\in T\). Generation follows by differentiating the finite separable polynomial equation of each algebraic element: its nonzero derivative solves for that element's differential. Independence follows by extending each coordinate derivation of \(F\) to every finite separable extension, with the exact formula \[
D(\bar y)=-\frac{(D_{\mathrm{coeff}}P)(\bar y)}{P'(\bar y)}.
\] Uniqueness again makes these extensions compatible. These derivations detect the coefficients of every finite linear relation between the \(dt\).

In the source's finite type situation, \(T\) is finite whenever the residue extension is separable. Taking dimensions gives exactly \[
\dim_\kappa(\Omega_{S/k}\otimes_S\kappa)
=\dim_\kappa(\mathfrak m/\mathfrak m^2)+\operatorname{trdeg}_k\kappa.
\] The original dimension identity at 28222--28263 is \(\dim_x\operatorname{Spec}S=\dim S_{\mathfrak q}+\operatorname{trdeg}_k\kappa\). Subtracting these two complete expressions proves that, for the stated separable residue fields, the excess of differential dimension over point dimension equals the embedding-dimension excess of the local ring. In particular it vanishes exactly when that ring is regular, proving the source criterion. No separability hypothesis is silently removed.

\subsubsection{The characteristic-zero derivation and its exact boundary}\label{algebra-proof-049-the-characteristic-zero-derivation-and-its-exact-boundary}

The proof at 38783--38815 requires that the specified cyclic summand \(S\,df\) be free with basis \(df\): the coefficient \(\theta(a)\) must be uniquely defined and \(\theta(f)=1\). Under a literal cyclic-submodule reading, the displayed decomposition alone does not imply that. For \(R=S=\mathbf Q\), \(f=0\) and \(C=0\), it holds as \(0=0\oplus0\), but the conclusion that \(f\) is not nilpotent fails. The zero ring gives another boundary to that conclusion. The proposed source clarification explicitly requires \(S\ne0\) and that \(S\to S\,df,\ a\mapsto a\,df\), be an isomorphism. The original decomposition, derivation and conclusion are otherwise retained. The intended free-summand reading already has this property; no failure of the subsequent regularity theorem follows.

Equivalently, work with the exact datum actually used in the argument: a nonzero \(\mathbf Q\)-algebra \(S\), an \(R\)-derivation \(\theta:S\to S\), and an element \(f\) with \(\theta(f)=1\). If \(f^n=0\) with \(n\ge1\) minimal, then \[
0=\theta(f^n)=n f^{n-1}.
\] The integer \(n\) is invertible, and \(f^0=1\ne0\) includes the \(n=1\) case. This is a contradiction. Conversely such a derivation induces \(\Omega_{S/R}\to S\) carrying \(df\) to \(1\), so it exhibits \(S\,df\) as the specified free direct summand.

For every \(a\) satisfying \(fa=0\), the source induction proves \[
\operatorname{Ann}_S(f)\subseteq\bigcap_{n\ge1}f^nS.
\] Indeed \(0=\theta(fa)=a+f\theta(a)\) gives the first inclusion. If \(a=f^n b\), then \(f^{n+1}b=0\), so \[
0=(n+1)f^n b+f^{n+1}\theta(b),\qquad
a=-\frac{f^{n+1}\theta(b)}{n+1}.
\] All coefficients and the sign are retained. Thus \(f\) is a nonzerodivisor whenever \(S\) is \(f\)-adically separated; Noetherianity or locality is not needed for this implication.

The source's local Noetherian conclusion also extends to every nonzero Noetherian \(\mathbf Q\)-algebra \(S\). At each maximal ideal, the derivation extends by \[
\theta(a/u)=\frac{u\theta(a)-a\theta(u)}{u^2},
\] and still sends \(f\) to \(1\). If \(f\) becomes a unit its annihilator is zero; otherwise local Krull intersection and the preceding inclusion give zero annihilator. Hence the localization of every element of \(\operatorname{Ann}_S(f)\) is zero at every maximal ideal. An element vanishing at all those localizations is zero: if its annihilator were proper, choose a maximal ideal containing it, where its image could not vanish. Thus \(\operatorname{Ann}_S(f)=0\).

One cannot simply drop both separation and Noetherianity. Let \[
B=\mathbf Q[x],\qquad
V=\mathbf Q[x,x^{-1}]/\mathbf Q[x],\qquad
S=B\oplus V
\] with square-zero multiplication \((a,u)(b,v)=(ab,av+bu)\). Ordinary differentiation preserves both Laurent polynomials and the submodule of polynomials, so \(\theta(a,u)=(a',u')\) is a well-defined \(\mathbf Q\)-derivation of \(S\). For \(f=(x,0)\), \(\theta(f)=(1,0)\). Nevertheless \[
f(0,[x^{-1}])=(0,[1])=0,\qquad (0,[x^{-1}])\ne0.
\] Multiplication by \(x\) is onto \(V\), while \(\bigcap_nx^nB=0\), so \[
\bigcap_{n\ge1}f^nS=0\oplus V,\qquad
\operatorname{Ann}_S(f)=0\oplus\mathbf Q\cdot[x^{-1}].
\] Explicitly, \(x\sum_{j=1}^N a_jx^{-j}\) is a polynomial exactly when all \(a_j\), \(j\ge2\), vanish. The ideal \(0\oplus V\) is not finitely generated: multiplying finitely many Laurent classes by polynomials cannot create a larger negative exponent than already occurs. Hence \(S\) is not Noetherian. This example has the actual free differential summand, so its failure is not the earlier notation issue.

For the source regularity theorem in characteristic zero, use induction on \(e=\dim_\kappa\mathfrak m/\mathfrak m^2\). If \(e=0\), Nakayama gives a field. Otherwise choose the original \(f\in\mathfrak m\setminus\mathfrak m^2\). The separable residue-field injection makes \(df\) nonzero modulo the maximal ideal in the assumed finite free differential module. A coordinate is a unit; replacing one basis vector by \(df\) is an invertible basis change. Thus \(df\) has the required free summand and \(f\) is a nonzerodivisor. The exact conormal map for quotienting by \(f\) sends the basis of \((f)/(f^2)\cong S_{\mathfrak q}/fS_{\mathfrak q}\) to \(df\), so its cokernel is finite free. The quotient is still essentially of finite type over \(k\), and its embedding dimension is \(e-1\), since the original class of \(f\) is nonzero. Induction makes it regular, and the original regular-mod-a-nonzerodivisor lemma 25816--25831 makes \(S_{\mathfrak q}\) regular. This checks the exact use of the clarified hypothesis.

\subsubsection{Characteristic-p examples with every coefficient retained}\label{algebra-proof-049-characteristic-p-examples-with-every-coefficient-retained}

The already composed correction MC-STK-ERR-0463 specifies the source's missing parameter as \(t\). Its example is the case \(m=2\) of the following separate family, where the letter \(m\) here is a positive integer: \[
k=\mathbf F_p(t),\qquad
S_m=k[x,y]/(x^p+y^m+t),\qquad m\ge1.
\] Keep the original \(x,y,t\). Before inverting nonzero polynomials in \(t\), elimination of the equation gives the mutually inverse maps \[
\mathbf F_p[t,x,y]/(x^p+y^m+t)\ \longleftrightarrow\ \mathbf F_p[x,y],
\quad t\longmapsto-x^p-y^m,\quad x\longmapsto x,\ y\longmapsto y.
\] The map \(\mathbf F_p[t]\to\mathbf F_p[x,y]\) is injective: in a nonzero polynomial \(h(-x^p-y^m)\), the largest power of \(x\) has nonzero coefficient and cannot cancel. Thus localization at its nonzero elements gives exactly \(S_m\), a regular domain as a localization of \(\mathbf F_p[x,y]\). As a \(k[y]\)-module it is free on the original basis \(1,x,\ldots,x^{p-1}\), by monic division. It is finite integral over \(k[y]\), with dimension \(1\).

In \(S_m\), the ideal \(\mathfrak q=(y,x^p+t)\) equals \((y)\). Its quotient is \(k[x]/(x^p+t)\), a field: \(t\) is not a \(p\)th power in \(k\), since its order at \(t=0\) is \(1\), whereas orders of \(p\)th powers are divisible by \(p\); the same holds for \(-t\). A purely inseparable root of exponent one has degree either \(1\) or \(p\), so \(x^p+t\) is irreducible. Thus \(\mathfrak q\) is maximal and its local ring has dimension \(1\), consistent with its regularity and the generator \(y\).

The exact conormal differential over the original \(k\) is \[
d(x^p+y^m+t)=m y^{m-1}\,dy,
\] with \(p=0\) in \(k\) and \(dt=0\). Consequently \[
\Omega_{S_m/k}=S_m\,dx\ \oplus\
\bigl(S_m/(m y^{m-1})\bigr)\,dy.
\] The exact smooth locus is \(D(m y^{m-1})\). At such a point the displayed derivative is a unit and supplies the standard smooth criterion. At every other point the differential fibre has dimension \(2\), whereas the original point dimension is \(1\), so the field criterion excludes smoothness.

If \(p\nmid m\) and \(m>1\), this locus is \(D(y)\); if \(m=1\), it is the whole spectrum. If \(p\mid m\), it is empty and the differential module is nevertheless free of rank \(2\). In particular the source's \(m=2,p>2\) ring is regular and is nonsmooth exactly along \(V(y)\), while its omitted \(p=2\) case is also regular, has free differentials, and is nowhere smooth over \(k\). These are statements about the original full rings, including their inseparable residue fields; regularity is not being silently substituted for smoothness.

This family also tests the preceding residue-conormal extension. At \(\mathfrak q=(y)\), the residue field is always the same purely inseparable \(\kappa=k[x]/(x^p+t)\). If \(m=1\), then \(y=-x^p-t\), so \(dy=0\) over \(k\). The nonzero generator of \(\mathfrak q/\mathfrak q^2\) therefore maps to zero in the differential fibre, even though the local ring is regular and the algebra smooth. Thus separability cannot simply be omitted from the general residue-injection statement. If \(m\ge2\), that same conormal generator maps to the independent \(dy\), and the sequence is injective on the left despite the inseparable residue extension. In fact \(x^p+t=-y^m=0\) modulo \(y^2\), so \(x\mapsto x\bmod y^2\) defines an explicit section \(\kappa\to(S_m)_{\mathfrak q}/(y^2)\). Hence separability is a sufficient condition in the proved general result, not a necessary one for an individual conormal splitting. This records both boundaries using the original coordinates and exact first-order quotients.

The other source example \(S=\mathbf F_p[x]/(x^p)\) has \(\Omega_{S/\mathbf F_p}=S\,dx\), because the relation differential is zero. Its unique local ring has dimension zero and nonzero maximal ideal modulo its square, so it is not regular or smooth. This verifies both source counterexamples and their different failures.

\subsubsection{Generic smoothness and report dispositions}\label{algebra-proof-049-generic-smoothness-and-report-dispositions}

For the original injective finite type map of domains \(R\to S\), preserve the fraction fields \(K\subset L\). A smooth neighbourhood of the generic point base changes to a smooth \(K\)-algebra, and after every extension \(K'/K\) its local rings are regular by the preceding field criterion. Thus it is geometrically reduced. Localizing to \(L\) preserves geometric reducedness, and the source field-extension criterion identifies this with separability of \(L/K\).

Conversely, the original generic finite-presentation lemma at 5729--5767 supplies nonzero \(f\in R\), \(g\in S\) with \(R_f\to S_{fg}\) finitely presented. Its one-generator proof inverts the actual nonzero leading coefficient and uses monic division, and its induction preserves the successive nonzero denominators. These localizations retain the generic point and the same fraction fields. For this finitely presented map, flat base change to \(K\) detects smoothness at the corresponding generic point by 37713--37744. Its local ring is \(L\), regular, and its residue extension is the specified separable \(L/K\); the proved separable criterion applies. Composition with the localization \(R\to R_f\) gives the required original smooth neighbourhood. The already composed grammatical correction at 38905 changes none of these maps.

Twenty reports give sixteen decision groups: eight new copyediting groups with twelve operations, three nondefect rejections, one optional punctuation preference, and four retained canonical units with five operations. OCC-00849 repeats the false undefined-H\_1(L) complaint; OCC-00852 asks for a period already present at 38411; OCC-00855 overlooks the original regular-local hypotheses. OCC-00858's genuine missing prime is already corrected, while localization at the image of \(g'\) is valid shorthand. The original homology notation, existing period and short exact sequence are preserved.

One additional source group explicitly supplies the intended free-summand and nonzero-ring conditions and includes the \(n=1\) nilpotence case. Four separate consequence groups cover projective section completions and finite base change; residue-field sections beyond the Noetherian and finite-generation assumptions; derivation regularity and its exact non-Noetherian boundary; and the complete characteristic-\(p\) family. None of these four has a source replacement operation. They are not a claim of new mathematics. The earlier More on Algebra findings remain separately pending, and broader synthesis follows core completion.

\hypertarget{algebra_smooth_artin_notes_20260928}{}
\subsection{Small extensions, residue conormal tests and the smooth overview}\label{algebra-proof-050-small-extensions-residue-conormal-tests-and-the-smooth-overview}

Authority: Stacks Project authors, commit a04446e57ec1fbc252a871afcec7752fb2807b14, algebra.tex:38933--39235, SHA256 FA8BB92E58A4F78A2BD01B3B6A4A87DE0A0D279F5DD90641B574DD5FBFFFA4F3. The whole interval and thirteen received reports were read. Original statements and proof order remain identifiable; the strengthened tests below are separate editorial material. Only the next opening 39236--39247 has been read, without adjudicating the étale section. The query ``smooth small extensions Artin Rees'' returned two PDF-primary routing records and no local-work results. Both records remain unread and unused; no PDF fallback or novelty claim is made.

\subsubsection{The original small-extension argument}\label{algebra-proof-050-the-original-small-extension-argument}

Keep the source \(R\to S\), its prime \(\mathfrak q\), and its finite presentation \[
P=R[x_1,\ldots,x_n]\longrightarrow S=P/I,\qquad
I=(f_1,\ldots,f_m).
\] Write \(\mathfrak q'\subset P\) for the inverse image of \(\mathfrak q\), and put \[
A=S_{\mathfrak q},\quad \mathfrak a=\mathfrak q A,\quad
D=P_{\mathfrak q'}/I^2P_{\mathfrak q'},\quad
E=I_{\mathfrak q'}/I_{\mathfrak q'}^2,\quad
\mathfrak n=\mathfrak q'D,\quad \kappa=A/\mathfrak a.
\] Thus \(D\to A\) is the original square-zero extension with kernel \(E\), and \[
F=A\otimes_D\Omega_{D/R}\cong
\bigoplus_{j=1}^n A\,dx_j.
\] Localization of the polynomial differential module and the equation \(d(I^2)\subset I\Omega_{P/R}\) give this exact identification. The original conormal differential \(d:E\to F\) is the target map at 39031--39034. The already composed MC-STK-ERR-0466 retains its \(n\) basis elements, rather than replacing that number by the \(m\) equations.

In a small extension \(B'\to B\), its length-one kernel \(K\) is a simple module over the local Artinian ring \(B'\), hence a one-dimensional vector space over its residue field, annihilated by the maximal ideal. Since \(B\) is nonzero local, \(K\) lies in that maximal ideal. It follows that \(K^2=0\) and every nonzero element of \(K\) generates it.

For each positive integer \(N\), set \[
B'_N=D/\mathfrak n^N,\qquad
B_N=A/\mathfrak a^N,\qquad
E_N=\ker(B'_N\to B_N)=E/(E\cap\mathfrak n^N).
\] Under the source hypotheses these are local Artinian rings with the same residue field \(\kappa\), and \(E_N^2=0\). A composition series of the finite \(B'_N\)-module \(E_N\) is also a chain of ideals, since its submodules are stable under multiplication by \(B'_N\). If \[
0=J_0\subset J_1\subset\cdots\subset J_l=E_N
\] is such a series, then each map \(B'_N/J_{i-1}\to B'_N/J_i\) is a small extension, with exactly the same residue field \(\kappa\).

Starting with the specified map \(S\to B_N\), apply the source lifting hypothesis in reverse order along this chain. Every intermediate lift has maximal-ideal inverse image \(\mathfrak q\): its residue map is the original \(S\to\kappa\), since each quotient kernel is nilpotent. Thus it factors through \(A=S_{\mathfrak q}\), because every element of \(S\setminus\mathfrak q\) has nonzero residue and maps to a unit. Its local map \(A\to B'_N\) sends \(\mathfrak a^N\) to zero, so the final lift factors through \(B_N\). The resulting map \(B_N\to B'_N\) is an \(R\)-algebra section of the original quotient.

The square-zero difference derivation for that section gives the split injection \[
E_N\longrightarrow B_N\otimes_{B'_N}\Omega_{B'_N/R}.
\] Its full construction is the one in groups 903 and 1083: if \(s:B_N\to B'_N\) is the section, \(b'\mapsto b'-s(\bar b')\) is an \(R\)-derivation into the square-zero kernel and induces a left inverse to the displayed conormal map.

Choose an Artin--Rees constant \(c\ge1\) for the submodule \(E\subset D\) and the ideal \(\mathfrak n\). It gives \[
E\cap\mathfrak n^N
=\mathfrak n^{N-c}(E\cap\mathfrak n^c)
\subseteq\mathfrak a^{N-c}E\qquad(N\ge c).
\] The equality uses the original module \(E\subset D\); the action of \(\mathfrak n\) on \(E\) factors through \(\mathfrak a\) because \(E^2=0\). Consequently, for \(r=N-c\ge1\), \[
(A/\mathfrak a^r)\otimes_A E_N
\cong E/\mathfrak a^rE.
\] The quotient differential sequence for \(D\to D/\mathfrak n^N\) has relations generated by \(d(\mathfrak n^N)\). Every derivative of a product of \(N\) factors lies in \(\mathfrak n^{N-1}\Omega_{D/R}\), by the full Leibniz sum. Since \(c\ge1\), tensoring by \(D/\mathfrak n^{N-c}\) kills those relations. Further tensoring by \(A/\mathfrak a^r\) therefore identifies the original split injection with \[
(A/\mathfrak a^r)\otimes_A d:
E/\mathfrak a^rE\longrightarrow F/\mathfrak a^rF.
\] These are precisely the source's two module identifications, with both exponents \(N\) and \(N-c\) retained.

The source then applies the earlier splitting lemma. Its proof does not require compatible choices of splittings at different orders. For completeness, over the Noetherian local \(A\), the kernel of \(d\) lies in all \(\mathfrak a^rE\), and hence is zero by Krull intersection. Put \(Q=F/E\). Choose the original beginning \(F_2\to F_1\to F_0\to Q\to0\) of a resolution by finite free modules. A lift \(F_0\to F\) gives \(\beta:F_1\to E\); the extension splits exactly when \(\beta\) lies in the image of \[
\operatorname{Hom}_A(F_0,E)\longrightarrow\operatorname{Hom}_A(F_1,E).
\] Each finite-order splitting supplies a map \(F_0\to E\) whose image under this homomorphism differs from \(\beta\) by a map into \(\mathfrak a^rE\), since \(F_0\) is free and its reduced map lifts. The image of \(\beta\) in the finite cokernel of the Hom map therefore belongs to every \(\mathfrak a^r\) multiple. Krull intersection makes that image zero, giving the actual splitting. This verifies the original proof route without replacing it by the sharper argument below.

Conversely, a smooth neighbourhood \(S_g\), \(g\notin\mathfrak q\), is formally smooth. In any diagram from the source with maximal-ideal inverse image \(\mathfrak q\), the map \(S\to B\) factors through \(S_{\mathfrak q}\) and \(S_g\). Formal smoothness provides the lift and hence all three stated testing conditions. The notation \(S\cap\mathfrak m\) is the established contraction notation for this inverse image and does not assert that \(S\to B\) is injective.

\subsubsection{A residue conormal criterion and its full base-change defect}\label{algebra-proof-050-a-residue-conormal-criterion-and-its-full-base-change-defect}

First consider any local ring \((A,\mathfrak a,\kappa)\), a finite \(A\)-module \(E\), and a map \(d:E\to F\) with \(F\) finite free. Then \(d\) is a split injection if and only if \[
d_\kappa:E/\mathfrak aE\longrightarrow F/\mathfrak aF
\] is injective. One direction follows by tensoring an actual left inverse. For the converse, let \(e_1,\ldots,e_t\) lift a basis of \(E/\mathfrak aE\). Nakayama makes \(A^t\to E\) onto. The images \(d(e_i)\) are independent in \(F/\mathfrak aF\); extend them to a residue-field basis and lift the remaining vectors to \(F\). Relative to an original basis of \(F\), the resulting square matrix has determinant outside \(\mathfrak a\), hence a unit. It is an invertible matrix over \(A\). Thus the composite \(A^t\to E\to F\) is the injection of the first \(t\) basis vectors. In particular \(A^t\to E\) is also injective and is an isomorphism. Projection onto those first coordinates gives the required exact left inverse to \(d\). This includes \(t=0\), when Nakayama gives \(E=0\). No Noetherian hypothesis has been used.

Apply this to the original finite presentation \(P\to S\), now allowing any base ring \(R\). The conormal module \(I/I^2\) is finite, generated by the classes of the original \(f_i\); \(F\) retains all \(n\) original differentials. The source local criterion at 37523--37569 therefore gives \[
R\to S\text{ is smooth at }\mathfrak q
\quad\Longleftrightarrow\quad
\ker\left((I/I^2)\otimes_S\kappa(\mathfrak q)
\xrightarrow{d}\bigoplus_{j=1}^n\kappa(\mathfrak q)\,dx_j\right)=0.
\] The forward direction tensors the split conormal sequence. The reverse direction uses the explicit local splitting just proved; it makes \(H_1(L)_{\mathfrak q}=0\) and the differential module a finite free quotient. The earlier local smoothness criterion then applies with its finite-presentation hypothesis intact. Its previously corrected right-inverse target in group 1054 is retained.

Define the exact nonnegative integer \[
\tau(\mathfrak q)=
\dim_{\kappa(\mathfrak q)}
\ker\left((I/I^2)\otimes_S\kappa(\mathfrak q)
\longrightarrow\bigoplus_{j=1}^n\kappa(\mathfrak q)\,dx_j\right).
\] This number is independent of the finite polynomial presentation. The source's actual homotopy equivalences between presentation complexes remain homotopy equivalences after tensoring with the residue field; their degree-one homology is exactly this kernel. Thus no replacement by a derived tensor product or omission of a Tor term is implicit here.

It is upper semicontinuous on \(\operatorname{Spec}S\). Indeed, let \(M=I/I^2\), let \(\mu(\mathfrak q)=\dim_{\kappa(\mathfrak q)}(M\otimes_S\kappa(\mathfrak q))\), and let \(\rho(\mathfrak q)\) be the rank of the actual \(n\)-by-\(m\) matrix whose columns are \(df_i\). Then \(\tau=\mu-\rho\), with \(0\le\mu\le m\). The locus \(\mu\le a\) is open: choose \(a\) generators at a given localization, clear their denominators, and annihilate the finite cokernel on one principal neighbourhood. The locus \(\rho\le b\) is closed, cut out by all \((b+1)\)-minors of that matrix. Therefore \[
\{\tau\ge s\}=
\bigcup_{j=0}^m
\bigl(\{\mu\ge j\}\cap\{\rho\le j-s\}\bigr)
\] is closed, with the rank condition empty for a negative upper bound. In particular the smooth locus is exactly \(\{\tau=0\}\).

Now retain any base change \(R\to R'\), with \(S'=S\otimes_RR'\), the actual base-changed polynomial algebra \(P'\), its kernel \(I'\), and a prime \(\mathfrak q'\subset S'\) over \(\mathfrak q\). Set \[
M' = I'/(I')^2,\qquad
O=\ker(M\otimes_SS'\longrightarrow M').
\] The original comparison of group 995 identifies \(O\) as the image of \(\operatorname{Tor}_1^R(S,R')\) in \(M\otimes_SS'\). Put \(\kappa'=\kappa(\mathfrak q')\). Right exactness of tensor product gives the exact sequence \[
O\otimes_{S'}\kappa'\longrightarrow
M\otimes_S\kappa'\longrightarrow M'\otimes_{S'}\kappa'\longrightarrow0.
\] Let \(W\) be the image of its first arrow; do not identify it with the whole tensor product unless injectivity is separately proved. The degree-zero differential modules on the two sides are the same \((\kappa')^n\), and the comparison respects their differentials. Thus \(W\) lies in the old differential kernel. Every vector in the new kernel lifts to a vector in the old kernel, because its differential is measured in that same \((\kappa')^n\). This proves \[
0\longrightarrow W\longrightarrow
H_1(\mathrm{NL}_{S/R}\otimes_S\kappa')
\longrightarrow
H_1(\mathrm{NL}_{S'/R'}\otimes_{S'}\kappa')
\longrightarrow0.
\] All spaces here are finite-dimensional, since \(M\) is finite. Field extension from \(\kappa(\mathfrak q)\) to \(\kappa'\) preserves the old kernel dimension, so the exact formula is \[
\tau_{S/R}(\mathfrak q)-\tau_{S'/R'}(\mathfrak q')
=\dim_{\kappa'}W\ge0.
\] For flat \(R'\) over \(R\), or flat \(S\) over \(R\), the relevant Tor group vanishes; then \(W=0\) and the defect is preserved. The second case gives preservation under every base change when the original map is flat and finitely presented, agreeing with group 1067. On a relative complete-intersection chart the conormal module is free on \(c\) original equations, so this defect is exactly \(c-\operatorname{rank}J\), the syntomic defect of group 1068. Its previous full Fitting-ideal statement remains valid there. No new canonical scheme structure for the thresholds of the general \(\tau\) is asserted.

\subsubsection{A single Artinian order and the weaker local Noetherian hypothesis}\label{algebra-proof-050-a-single-artinian-order-and-the-weaker-local-noetherian-hypothesis}

Assume now that \(R\to S\) is finitely presented and that \(A=S_{\mathfrak q}\) is Noetherian; \(R\) itself need not be Noetherian. With the same original presentation and definitions, \(D\) is Noetherian. To see this, \(E\) is finite over the Noetherian ring \(A\) and \(E^2=0\). For any ideal \(K\subset D\), its image in \(A\) is finitely generated and \(K\cap E\) is a finite \(A\)-submodule of \(E\). Lift finite generators of the image and adjoin finite generators of the intersection. Their \(D\)-span is exactly \(K\). Thus every ideal of \(D\) is finitely generated. This proves the needed local Noetherianity directly; it is not assumed for the ambient polynomial ring.

Choose any Artin--Rees constant \(c\ge1\) for \(E\subset D\) with respect to \(\mathfrak n\), and set \(N=c+1\). Then the following are equivalent:

\begin{enumerate}
\def\labelenumi{\arabic{enumi}.}
\tightlist
\item
  The original map \(R\to S\) is smooth at \(\mathfrak q\).
\item
  The single canonical surjection \(B'_N=D/\mathfrak n^N\to B_N=A/\mathfrak a^N\) has an \(R\)-algebra section.
\item
  The original map \(S\to B_N\) has an \(R\)-algebra lift to \(B'_N\).
\end{enumerate}

The equivalence of 2 and 3 is exact. A section gives a lift by composition; a lift is local after factoring through \(A\), has \(\mathfrak a^N\) in its kernel, and hence gives the section. Condition 1 gives the lift by formal smoothness of a principal neighbourhood, as already verified.

For 2 implies 1, the section gives the split conormal injection for \(E_N\). Tensor it with \(\kappa\). Artin--Rees gives \(E\cap\mathfrak n^N\subset\mathfrak aE\), so \[
E_N\otimes_A\kappa=E/\mathfrak aE.
\] Since \(N=c+1\ge2\), the differential relation \(d(\mathfrak n^N)\subset\mathfrak n^{N-1}\Omega_{D/R}\) vanishes after tensoring with \(\kappa\), identifying its target with \(F/\mathfrak aF\). We obtain an injective \(d_\kappa\), and the preceding finite-module argument gives an actual splitting of \(d:E\to F\). Hence the original map is smooth at \(\mathfrak q\).

This proves all four source testing conditions under the weaker hypotheses ``finitely presented, with \(S_{\mathfrak q}\) Noetherian''. The universal small-extension condition with the fixed residue field supplies a lift through the composition series of this single \(E_N\), and the single-order criterion proves smoothness. The converse has no Noetherian requirement, since it comes from formal smoothness.

It also gives a bounded obstruction when smoothness fails. Use the same \(N=c+1\) and a composition series of \(E_N\) of length \(l=\operatorname{length}_{B'_N}E_N\). Beginning at \(S\to B_N\), make any available successive lifts. A full chain would give the forbidden section, so some step must fail after at most \(l\) small-extension tests. The source and target of every tested map are quotients of \(B'_N\); their maximal ideals have \(N\)th power zero and their residue fields are the specified \(\kappa\). Thus failure is witnessed by one small extension with this exact order bound and residue field, without constructing an infinite compatible tower. The bound depends on the actual Artin--Rees datum and the original presentation; no universal numerical bound independent of that datum is claimed.

\subsubsection{The overview's finite-presentation boundary}\label{algebra-proof-050-the-overviews-finite-presentation-boundary}

The overview at 39183--39185 omits ``finitely presented'', which is explicit in its cited lemma 37713--37744. Restore that hypothesis in the overview. Its opening instruction to consult precise statements explains the likely abbreviation, but the unqualified assertion is false, even for a flat finite type map. The following exact counterexample propagates the finite-support construction previously used in group 685.

Let \(k\) be any field and retain \[
R=\prod_{i\ge1}k,\qquad I=\bigoplus_{i\ge1}k,\qquad S=R/I.
\] Here \(I\) is the ideal of finitely supported sequences. It is proper, since \(1\) has infinite support, and \(S\ne0\). Every finitely generated ideal of \(R\) is generated by the idempotent of the union of the supports of its finitely many generators. For one element \(a=(a_i)\), define \(b_i=a_i^{-1}\) when \(a_i\ne0\), and \(b_i=0\) otherwise; then \(ab\) is that support idempotent and \(a=a^2b\). These exact coordinate formulas also show every localization \(R_{\mathfrak p}\) at a prime is a field: if \(a\in\mathfrak p\), then \(1-ab\notin\mathfrak p\) and \((1-ab)a=0\), so \(a/1=0\); elements outside \(\mathfrak p\) are invertible.

The quotient \(S\) is \(R\)-flat. Indeed \(I\) is the directed union of the idempotent ideals \(e_FR\) for finite sets \(F\subset\mathbf N\). Thus \(S\) is the filtered colimit of the projective modules \(R/e_FR\), with the actual quotient transition maps; filtered colimits of flat modules are flat. The ring map \(R\to S\) is finite type, generated as an algebra by the empty set.

For any nonzero \(\bar g\in S\), choose a representative \(g\in R\) and its support idempotent \(e\). That support is infinite, since \(\bar g\ne0\). The coordinate inverse above identifies \(S_{\bar g}=S_{\bar e}\), and the map from \(R\) is surjective, with kernel \[
K_e=I+(1-e)R.
\] For completeness, in the quotient by \(I+(1-e)R\), \(e=1\), so the universal localization map factors through it. Conversely localization kills \(1-e\), and if an element maps to zero then its \(e\)-part is finitely supported, placing it in \(K_e\). This proves the kernel equality.

The ideal \(K_e\) is not finitely generated. Each proposed generator restricts on the infinite support of \(e\) to finitely many nonzero coordinates. Finitely many generators cover only a finite subset \(F\) of that support. For any \(i\) in the support outside \(F\), the coordinate idempotent \(e_{\{i\}}\) lies in \(I\subset K_e\), but its \(i\)th coordinate cannot be obtained from those generators, whose \(i\)th coordinates are all zero. This is a contradiction. Hence \(S_{\bar g}\) is not finitely presented as an \(R\)-algebra: a surjective map \(R\to R/K_e\) is finitely presented exactly when \(K_e\) is finitely generated. It follows that the smooth locus of \(R\to S\) is empty, because every prospective principal neighbourhood at a prime uses a nonzero \(\bar g\).

Choose any maximal ideal of the nonzero ring \(S\), and let \(\mathfrak p\subset R\) be its inverse image. Each finite-support idempotent belongs to \(\mathfrak p\) and becomes zero in the field \(R_{\mathfrak p}\); thus \(I_{\mathfrak p}=0\). The flat base change \(R\to R'=R_{\mathfrak p}\) gives \[
S\otimes_RR'\cong R'/I_{\mathfrak p}=R'.
\] The base-changed map is the identity of the field \(R'\), smooth at its unique point. This nonempty smooth locus cannot be the inverse image of the original empty locus. The example retains even flatness and finite type of the original map; it isolates the missing finite-presentation condition.

There is no contradiction with the preceding defect calculation: it was proved for a finite presentation. Here \(I=I^2\) is infinitely generated, so its zero conormal module does not make the quotient a smooth algebra. This is another explicit boundary against using formal infinitesimal data as a substitute for the finite-presentation requirement.

In fact this same map is formally smooth, by the exact argument of group 1085: in any square-zero lifting diagram for \(R\to R/I\), the image of \(I\) in the target lies in the square-zero ideal \(K\), and \(I=I^2\) forces that image into \(K^2=0\). Thus the given map from \(R\) factors uniquely through \(R/I\), producing the lift. At every prime of \(S\), its local ring is the field \(R_{\mathfrak p}\) computed above. Consequently neither formal smoothness, flatness, finite type, nor Noetherianity of all local target rings removes the finite-presentation requirement from the strengthened Artinian criterion.

\subsubsection{Report accounting and propagation}\label{algebra-proof-050-report-accounting-and-propagation}

Thirteen reports give ten decision groups. Five copyediting units contain eight operations: the missing ``is'' in the finite type assumption; the designation of \(S'\); a period and capital letter between the differential calculation and its next operation; one maximal-ideal predicate; and three prime-ideal predicates in the overview. Colons before the enumerations remain optional. The request to insert ``that'' after ``we see'' is rejected as a nondefect: the complementizer can be omitted in this grammatical construction. The contraction notation receives an optional clarification disposition and is retained.

Three whole canonical units retain the original variable count, plural ``tensor products'', and the residue-field subscript. Their paired reports are accounted for once. A separate source correction restores finite presentation to the flat-base-change overview. Two zero-operation consequence groups contain the residue conormal defect and the single-order Artinian criterion with the weakened local Noetherian assumption. Their full proofs and boundaries remain editorial material. The other overview assertions retain their original field hypotheses, point dimensions, references and verified source criteria. The earlier More on Algebra corrections remain separately pending. Broader synthesis and any short paper follow completion of the core work.

\hypertarget{algebra_etale_notes_20260928}{}
\subsection{Étale maps, their infinitesimal kernels and polynomial factorizations}\label{algebra-proof-051-uxe9tale-maps-their-infinitesimal-kernels-and-polynomial-factorizations}

Private source-linked editorial evidence, 2026-09-28. Source: The Stacks Project authors, algebra.tex at authority a04446e57ec1fbc252a871afcec7752fb2807b14, lines 39236--39753; full source SHA-256 FA8BB92E58A4F78A2BD01B3B6A4A87DE0A0D279F5DD90641B574DD5FBFFFA4F3. The following arguments explain and extend the original passages; they do not replace the translated source or claim novelty.

\subsubsection{Original presentations and permanence arguments}\label{algebra-proof-051-original-presentations-and-permanence-arguments}

At 39264--39287 the source uses a finite polynomial presentation \(S=P/I\), \(P=R[x_1,\ldots,x_n]\). Étaleness says the original conormal differential \[
d:I/I^2\longrightarrow\bigoplus_{i=1}^n S\,dx_i
\] is an isomorphism. Choose \(f_i\in I\) lifting the inverse images of these particular \(n\) coordinate vectors. The proof of lemma-huber, 36520--36546, applies Nakayama to the finitely generated module \(I/(f_1,\ldots,f_n)\): there is \(v\in1+I\) with \(vI\subset(f_1,\ldots,f_n)\). Hence \[
S=P_v/(f_1,\ldots,f_n)
 =R[x_1,\ldots,x_n,z]/(f_1,\ldots,f_n,vz-1).
\] There are exactly \(n+1\) variables and \(n+1\) equations. The enlarged Jacobian has the original \(n\) by \(n\) derivative block, a zero column above its last entry, and last row \((z\partial v/\partial x_1,\ldots,z\partial v/\partial x_n,v)\). Its determinant is \(v\det(\partial f_i/\partial x_j)\), whose image in \(S\) is a unit. This proves the stated global square presentation while retaining the particular generators and the extra equation. It also avoids inferring equality of ranks over the zero ring: the number of chosen lifts was \(n\) from the outset. If \(S=0\), the alternative presentation \(R[x]/(1)\) has determinant \(0=1\) in \(S\), which is its unit. No extra nonzero hypothesis is needed.

For 39289--39387, composition follows from the differential sequence, base change from \(\Omega_{S'/R'}=S'\otimes_S\Omega_{S/R}\), and principal-cover locality from localization of that module and smoothness. In the flat-base-change statement finite presentation is explicitly present and retained. At a prime above a given prime, the induced map of local rings is faithfully flat; vanishing of the localized differential module descends. The smooth locus comparison is the earlier proved smoothness result. In 39368--39372 the descended smooth algebra has the clopen relative-dimension decomposition of 33757--33771. After the specified base change every component of positive dimension is empty, so retaining the dimension-zero factor gives exactly the original algebra. Finite presentation over the integers allows the coefficients and their finite equations to descend to a filtered stage. Principal localizations then give the multiplicative-set assertion. The differential module of a finite product is the product of the component differential modules.

At 39393--39432 the field case gives a finite product of finite separable field extensions; the empty product is allowed and is the zero ring. Nonempty factors are fields because they are zero-dimensional regular local rings. Finiteness follows from the finite type field theorem, and separability follows from smoothness at the generic point. Conversely the minimal polynomial of a primitive element is coprime to its derivative, so its conormal differential is an isomorphism. Base changing this description proves both conclusions at a prime in 39434--39453 and quasi-finiteness in 39455--39467.

For the converse at 39469--39502, isolate the specified fibre point by \(g\notin\mathfrak q\). The resulting fibre has one prime, is local, and equals \(S_{\mathfrak q}/\mathfrak pS_{\mathfrak q}=\kappa(\mathfrak q)\). The finite separable fibre and local flatness give a smooth neighbourhood. On its standard smooth presentation the original relative dimension is zero, so the numbers of equations and variables agree. In 39507--39530, an \(R\)-map between two étale algebras is finitely presented by 539--563. Apply 32991--33005 with \(M=S_{\mathfrak q}\): both local algebras are essentially finitely presented over \(R_{\mathfrak p}\), \(M\) is nonzero and finitely presented over itself, flat over the base, and its residue fibre is flat over the intermediate residue field. This proves the required local flatness. The residue extension is finite separable, so the preceding converse applies.

For a surjective flat finitely presented map \(R\to S\), its kernel \(K\) is finitely generated and tensoring \(0\to K\to R\to S\to0\) with the flat module \(S\) gives \(K/K^2=0\). The proof in 3922--3951 supplies an idempotent \(e\in K\) with \(K=(e)\). Thus \(S=R_{1-e}\), keeping the distinction between the generator of the kernel and the complementary idempotent inverted in the quotient.

In the arbitrary-ideal lifting argument at 39564--39587, retain every original lifted polynomial and write \[
S=R[x_1,\ldots,x_n,z]/(f_1,\ldots,f_n,z\Delta-1),
\qquad \Delta=\det(\partial f_i/\partial x_j).
\] With equation order \((f_1,\ldots,f_n,z\Delta-1)\) and variable order \((x_1,\ldots,x_n,z)\), the full derivative matrix is \[
\begin{pmatrix}
(\partial f_i/\partial x_j)_{i,j=1}^n&0\\
z\partial\Delta/\partial x_1\ \cdots\ z\partial\Delta/\partial x_n&\Delta
\end{pmatrix}.
\] Its determinant is \(\Delta^2\), a unit since \(z\Delta=1\). Modulo the original ideal the extra variable equals the already existing inverse of \(\overline\Delta\), so the reduction is precisely \(\overline S\). For \(n=0\), the empty determinant is \(1\), the extra equation is \(z-1\), and the conclusion remains valid. Existence here needs no nilpotence assumption.

\subsubsection{Canonical infinitesimal map and nil ideals}\label{algebra-proof-051-canonical-infinitesimal-map-and-nil-ideals}

Keep the exact diagram of 39589--39620: \(A'\to A\) and \(B'\to B\) are surjections with square-zero kernels \(I,J\); the vertical homomorphisms are \(A'\to B'\) and \(A\to B\). Multiplication defines the specific map \[
\mu:I\otimes_A B\longrightarrow J,\qquad
i\otimes b\longmapsto \iota(i)\widetilde b.
\] Here \(\iota:A'\to B'\) is the given map and \(\widetilde b\) is any lift of \(b\). A change of lift lies in \(J\), and \(\iota(I)J\subset J^2=0\), proving independence. Balancing over \(A\) follows by lifting an element of \(A\) to \(A'\); changes of that lift lie in \(I\), whose square is zero. Both sides are \(B\)-modules via the quotient. The source's equality \(J=I\otimes_A B\) must refer to this canonical identification.

Suppose \(A\to B\) is étale and \(\mu\) is an isomorphism. Lift \(B\) to an étale \(A'\)-algebra \(C\) by the preceding arbitrary-ideal construction, with a specified identification \(C/IC=B\). Formal smoothness lifts this specified quotient map to an \(A'\)-map \(\varphi:C\to B'\) reducing to the identity of \(B\). Flatness of \(C\) identifies \[
IC=I\otimes_{A'}C=I\otimes_A(C/IC)=I\otimes_A B.
\] Under this identification the restriction of \(\varphi\) to \(IC\) is exactly \(\mu\). It is therefore an isomorphism onto \(J\), and in particular \(J=IB'\). The quotient map is also an isomorphism. For injectivity of \(\varphi\), an element of its kernel first lies in \(IC\), where the restriction is injective. For surjectivity, lift the quotient class of any element of \(B'\), then lift its remaining difference in \(J\) through \(IC\). Thus \(\varphi\) is an isomorphism and \(A'\to B'\) is étale.

The converse also holds under the same square-zero diagram and the same assumption that \(A\to B\) is étale. If \(A'\to B'\) is étale, put \(D=B'/IB'\). Both \(D\) and \(B\) are étale over \(A\); hence the surjection \(D\to B\) is étale by 39507--39530. Its finitely generated kernel \(J/IB'\) is idempotent by flatness, but is square zero. Consequently it is zero. Thus \(J=IB'\), and flatness of \(B'\) over \(A'\) gives exactly the canonical isomorphism \(\mu\). This is a proved converse, recorded as an editorial consequence.

An abstract module isomorphism cannot replace \(\mu\). Let \(k\) be a field, \[
A'=k[\epsilon]/(\epsilon^2),\quad A=k,\quad
B'=k[\delta]/(\delta^2),\quad B=k,
\] with the map \(A'\to B'\) sending \(\epsilon\) to zero. Then \(I\otimes_A B\) and \(J\) are each one-dimensional over \(k\), but \(\mu=0\). Tensoring the inclusion \((\epsilon)\hookrightarrow A'\) with \(B'\) gives a zero map from the nonzero module \(B'\) to \(B'\). Thus \(B'\) is not flat over \(A'\) and the map is not étale. This clarifies the intended canonical reading; it does not assert that the source intended an arbitrary abstract isomorphism.

\textbf{Nil-ideal consequence.} For every nil ideal \(I\subset R\), reduction defines an equivalence between the categories of étale \(R\)-algebras and étale \(R/I\)-algebras. It restricts to finite étale algebras and preserves their ranks at corresponding primes. The whole ideal need not have a common nilpotence exponent.

Essential surjectivity is the source's lifting construction. For full faithfulness, take étale \(R\)-algebras \(S,T\). The ideal \(IT\) is nil: each of its elements is a finite sum of multiples of nilpotent elements of \(I\), and a sufficiently high power of such a finite sum is zero. Given a map \(S/IS\to T/IT\), choose a finite presentation \(S=R[x_1,\ldots,x_a]/(F_1,\ldots,F_b)\) and lifts \(t_i\in T\) of the images of the generators. Put \(e_j=F_j(t_1,\ldots,t_a)\in IT\) and choose positive integers \(d_j\) with \(e_j^{d_j}=0\). The finitely generated ideal \(Q=(e_1,\ldots,e_b)\) satisfies \[
Q^{\,1+\sum_{j=1}^b(d_j-1)}=0:
\] in each monomial of that total degree one exponent is at least its specified \(d_j\). If \(b=0\), then \(Q=0\) and exponent \(1\) suffices. The chosen coordinates define \(S\to T/Q\). Formal smoothness lifts it successively through the finitely many maps \(T/Q^{j+1}\to T/Q^j\), whose kernels are square zero for \(j\ge1\), until it gives \(S\to T\). Each coordinate correction lies in \(Q\subset IT\), so this is the required lift.

For uniqueness, let \(\phi,\psi:S\to T\) have the same reduction. For a finite algebra generating set \(s_1,\ldots,s_a\), set \(\eta_i=\phi(s_i)-\psi(s_i)\in IT\), and \(K=(\eta_1,\ldots,\eta_a)\). Choose \(e_i\ge1\) with \(\eta_i^{e_i}=0\); then \(K^{1+\sum(e_i-1)}=0\). Polynomial subtraction shows \(\phi(s)-\psi(s)\in K\) for every \(s\in S\). Modulo \(K^2\) this difference is an \(R\)-derivation into \(K/K^2\), with \(S\)-action induced by their common map to \(T/K\): the omitted product of two differences is in \(K^2\). Since \(\Omega_{S/R}=0\), the derivation is zero. Every \(\eta_i\) lies in \(K^2\), so \(K=K^2\); nilpotence gives \(K=0\). Thus \(\phi=\psi\). This proves full faithfulness. Lifts of inverse isomorphisms are inverse because of this uniqueness, completing the equivalence.

To check the finite subcategory, let the lifted \(S\) reduce to a finite \(R/I\)-module. Each of its finitely many algebra generators \(x_i\) satisfies a monic polynomial modulo \(IS\). One can obtain such a polynomial by writing multiplication by \(\bar x_i\) on finitely many module generators, taking the determinant of the matrix \(X\,1-M\), and using its adjugate; the resulting polynomial annihilates every generator and hence \(1\). Lift all its coefficients to \(R\), obtaining a monic \(p_i(X)\) with \(p_i(x_i)\in IS\). This element is nilpotent; choose \(N_i\ge1\) with \(p_i(x_i)^{N_i}=0\). The monic polynomial \(p_i(X)^{N_i}\) annihilates \(x_i\). Successive division by these monic polynomials shows that the monomials \[
\prod_i x_i^{a_i},\qquad 0\le a_i<N_i\deg p_i,
\] span \(S\) over \(R\). Retain all these actual degrees and exponents.

Moreover \(S\) is a finitely presented \(R\)-module. Indeed the intermediate algebra \[
U=R[X_1,\ldots,X_a]/(p_i(X_i)^{N_i}:1\le i\le a)
\] is finite free on the displayed bounded monomials. The kernel of \(U\to S\) is generated as a \(U\)-ideal by the finitely many relations of the original finite presentation; multiplying them by the finite \(R\)-basis of \(U\) gives finitely many \(R\)-module generators. Étale flatness now makes \(S\) finite projective over \(R\). Conversely reduction preserves finite étale algebras. Every prime of \(R\) contains \(I\), and its residue field agrees with that of the corresponding prime of \(R/I\). Tensoring the original module with these identical fields proves equality of ranks. No finite-degree or uniqueness assertion for arbitrary nonnil ideals follows from this argument.

\subsubsection{The full factorization matrix and its fibres}\label{algebra-proof-051-the-full-factorization-matrix-and-its-fibres}

Keep \(N=n+m\), \(R=\mathbf Z[a_1,\ldots,a_N]\), \(S=\mathbf Z[b_1,\ldots,b_n,c_1,\ldots,c_m]\), \(b_0=c_0=1\), and \(b_i=c_j=0\) outside their original index ranges. The map is \(a_k\mapsto\sum_{i+j=k}b_ic_j\). Put \[
g=x^n+b_1x^{n-1}+\cdots+b_n,\qquad
h=x^m+c_1x^{m-1}+\cdots+c_m.
\] The source's matrix \(J\) at 39640--39653 has rows ordered by \(b_1,\ldots,b_n,c_1,\ldots,c_m\) and columns by equations \(k=1,\ldots,N\). Its first \(n\) rows are the shifted coefficient rows of \(h\), and its last \(m\) rows are those of \(g\). Set \(\Delta=\det J\), exactly as in the source.

With descending monomial bases, \(J^t\) is the coefficient matrix of \[
L:S[x]_{<n}\oplus S[x]_{<m}\longrightarrow S[x]_{<N},
\qquad (u,v)\longmapsto hu+gv.
\] The map displayed in the source instead writes the summands in the other order: \[
T:S[x]_{<m}\oplus S[x]_{<n}\longrightarrow S[x]_{<N},
\qquad(a,b)\longmapsto ga+hb.
\] Let \(P(a,b)=(b,a)\). Then \(T=L\circ P\). Moving a block of \(m\) basis elements across a block of \(n\) elements gives exactly \(nm\) transpositions, so \[
\det T=(-1)^{nm}\det L=(-1)^{nm}\Delta.
\] Thus the source's transpose identification is exact after the specified permutation of the two domain blocks. Its displayed determinant and all its localization opens stay unchanged. For \(n=m=1\), \(J=\left(\begin{smallmatrix}1&c_1\\1&b_1\end{smallmatrix}\right)\), so \(\Delta=b_1-c_1\), whereas \(\det T=c_1-b_1\). This records the sign without replacing the source's resultant convention by an unstated convention.

At a prime of \(S\), work over its residue field \(K\). Let \(d\) be the monic greatest common divisor of the actual images of \(g,h\), of degree \(r\), and write \(g=dg_0,h=dh_0\). The kernel of the actual tangent coefficient map \(L_K\) is exactly \[
K[x]_{<r}\longrightarrow\ker L_K,\qquad
w\longmapsto(g_0w,-h_0w).
\] Indeed \(hu+gv=0\) implies \(h_0u=-g_0v\). Coprimality of \(g_0,h_0\), using their Bézout identity, gives \(g_0\mid u\) and \(h_0\mid v\); hence \(u=g_0w,v=-h_0w\). The two original degree bounds are each precisely \(\deg w<r\). Conversely those pairs satisfy the equation. Therefore the kernel dimension is \(r\), the rank is \(N-r\), and the cokernel dimension is \(r\). The zero polynomial \(w\) is included; when \(r=0\) the parameter space is zero.

The original universal factorization algebra is a relative global complete intersection of relative dimension zero, as proved with its full coefficients in group 1042. Consequently group 1068 applies with \(N\) variables and \(N\) equations: the Jacobian defect is exactly \(\deg\gcd(g,h)\). The locus where it is at least \(s\), for \(1\le s\le N\), has the previously defined scheme structure cut out by every \((N-s+1)\)-minor of the actual matrix \(J\), that is \(\operatorname{Fitt}_{s-1}(\Omega_{S/R})\). All these full ideals, not just their radicals, commute with every base change. The zero and unit boundary ideals for thresholds outside this range are retained. For \(s=1\), the ideal is \((\Delta)\), giving exactly the original étale open.

For completeness, the matrix equivalence also supplies the actual bounded Bézout coefficients. When \(\Delta\) is a unit, so is \(\det T=(-1)^{nm}\Delta\). Apply \((\det T)^{-1}\operatorname{adj}(T)\) to the coefficient vector of the constant polynomial \(1\) in the descending basis. The resulting vectors give \(a\in S[1/\Delta][x]_{<m}\), \(b\in S[1/\Delta][x]_{<n}\) with \(ag+bh=1\). Neither the determinant sign nor the polynomial degree bounds is discarded.

There are further precise consequences for the fibres. Over an algebraically closed field \(K\), write the base polynomial as \[
f=\prod_{i=1}^t(x-\alpha_i)^{e_i},
\quad \alpha_i\ne\alpha_j\ (i\ne j),\quad e_i\ge1.
\] A monic degree-\(n\) factor corresponds to integers \(r_i\) with \(0\le r_i\le e_i\), \(\sum_i r_i=n\). Its complementary factor has exponents \(e_i-r_i\). Their greatest common divisor has degree \[
\sum_{i=1}^t\min(r_i,e_i-r_i).
\] Thus they are coprime precisely when each entire multiplicity block goes into one factor. The number of étale points is \[
[z^n]\prod_{i=1}^t(1+z^{e_i}).
\] The fibre algebra is finite by group 1042. Its localization at \(\Delta\) remains finite-dimensional over \(K\): a finite-dimensional commutative algebra is Artinian and decomposes as finitely many local Artinian factors, and localizing keeps exactly the factors on which \(\Delta\) is a unit. Every retained factor is étale over \(K\), hence is \(K\) by the field classification above. Therefore the localized fibre is the product of exactly the displayed number of copies of \(K\). Zero copies means the zero ring, not a missing case. If all roots are distinct, this number is \(\binom{N}{n}\).

One can also retain the finite rank on the discriminant open of the base. Write \(D(f)\) for the monic discriminant, with root definition \(\prod_{i<j}(\rho_i-\rho_j)^2\), and similarly \(D(g),D(h)\). To check its full pullback, use the ordered-root extension from group 1042, putting \(g=\prod_{i=1}^n(x+u_i)\), \(h=\prod_{j=1}^m(x+v_j)\). Its map from \(S\) is injective: the two ordered-root towers are finite free with bases containing \(1\), so the coefficient map is an injection at every stage. Over the fraction field of the root polynomial ring, division by monic \(g\) decomposes every polynomial of degree \(<N\) uniquely as \(ga+r\), \(\deg a<m,\deg r<n\). In descending bases this division coordinate change is triangular with diagonal entries \(1\). The map \(T\) in these coordinates is block triangular with first diagonal block the identity and second block multiplication by \(h\) modulo \(g\). Evaluation at the distinct roots \(-u_i\) diagonalizes that second block, giving \[
\det T=\prod_{i=1}^n h(-u_i)
 =\prod_{i,j}(v_j-u_i)=(-1)^{nm}\prod_{i,j}(u_i-v_j).
\] The equality is polynomial and therefore holds in the entire root polynomial ring. It follows that \(\Delta=\prod_{i,j}(u_i-v_j)\). Separating the internal and cross pairs in the root product for \(f=gh\), and using injectivity, gives exactly \[
D(f)=D(g)D(h)\Delta^2.
\] No factorial or characteristic assumption is used. After inverting the actual base element \(D(f)\), \(\Delta\) is a unit because a factor of a unit in a commutative ring is a unit. Hence the entire base-changed factorization algebra is finite étale, of the same finite locally free rank \(\binom{N}{n}\) established in group 1042. This does not assert that \(S[1/\Delta]\) is finite over the original uninverted base; localization of a finite algebra need not have that property.

\subsubsection{Original factors and report accounting}\label{algebra-proof-051-original-factors-and-report-accounting}

In 39684--39748 the input factors need not be monic. Keep their original leading coefficients \(\overline b_0,\overline c_0\) with \(\overline b_0\overline c_0=1\). Choose \(a,s\in R\setminus\mathfrak p\) representing \(\overline b_0=\bar a/\bar s\) in \(\kappa(\mathfrak p)\). Work first over \(R[(as)^{-1}]\); both \(a,s\) are now units, and \(u=a/s\) is a unit lifting that specific leading coefficient. Apply the universal monic factorization to the auxiliary factors \(\overline g/\bar u\) and \(\bar u\overline h\), obtaining monic lifts \(G,H\) on its étale \(\Delta\)-open. Recover the required original factors by \[
g=uG,\qquad h=u^{-1}H.
\] They have exactly the original reductions and satisfy \(gh=GH=f\). If \(\alpha G+\beta H=1\), then \[
(\alpha/u)g+(\beta u)h=1.
\] The first \(n\) rows of the original coefficient matrix scale by \(u^{-1}\) and the last \(m\) rows by \(u\), so its determinant for \((g,h)\) is \(u^{m-n}\Delta(G,H)\); this explicit unit factor is retained. These auxiliary monic objects prove the statement about the original factors; they do not replace those factors.

The point in the universal algebra is the kernel of the coefficient-evaluation map into \(\kappa(\mathfrak p)\). Although that ring homomorphism need not be surjective, the residue field at its kernel is exactly \(\kappa(\mathfrak p)\): the image is a subring containing the image of \(R/\mathfrak p\), and hence its fraction field contains \(\operatorname{Frac}(R/\mathfrak p)=\kappa(\mathfrak p)\) while being contained in that same field. The nonzero determinant permits the same point after localization. For \(n=0\), take \(g=u,h=u^{-1}f\) already over \(R[(as)^{-1}]\); for \(m=0\), take \(h=u^{-1},g=uf\). The nonzero constant factor is a unit, so the Bézout property holds. If both degrees are zero then \(f=1\), and both formulas give the same unit factorization.

Exactly twelve received reports are accounted for in seven report groups:

\begin{itemize}
\tightlist
\item
  OCC-12127 retains the whole existing canonical correction MC-STK-ERR-0470 at 39440.
\item
  OCC-12128 and OCC-00874 retain MC-STK-ERR-0471 at 39464. The additional claim that ``fields finite separable over'' requires ``that are'' is not a defect; the original mathematical adjective phrase is valid.
\item
  OCC-00875 and OCC-12129 retain the whole multiline MC-STK-ERR-0472 at 39496.
\item
  OCC-00876 and OCC-12130 retain MC-STK-ERR-0473 at 39536.
\item
  OCC-00877 proposes only source whitespace in math mode at 39624; it is optional and is not integrated.
\item
  OCC-12131, OCC-00879 and explicitly duplicate OCC-01311 retain MC-STK-ERR-0474 at 39744, replacing the wrong pair \(\operatorname{Res}_x(f,g)\) by the actually inverted pair \(\operatorname{Res}_x(g,h)\). All three physical reports have a disposition, but there is one retained correction.
\item
  OCC-00878 supplies the single new grammar operation ``Denote'' to ``Denote by'' at 39735.
\end{itemize}

Groups 1127 and 1128 are separate zero-operation editorial consequence groups. The former receives group 1085's finite-error lifting argument and proves the canonical infinitesimal converse and nil-ideal equivalence, including finite étale ranks. The latter receives groups 1042 and 1068, identifies their actual Jacobian defect with the full gcd degree, proves the multiplicity-aware fibre count, and verifies the complete signed determinant and discriminant comparison. Source omissions of these details are not reclassified as translation errors. No silent replacement of the original presentation is proposed.

The complete current source interval is checked against the original after the five retained canonical operations. Dependency reading includes 539--563, 3922--3951, 32991--33005 (the flatness criterion's statement and its actual hypotheses), 33757--33771 and 36520--36546. The four new indexed query hits are unread routing records, not supporting literature; seven earlier external-source reading records retain their exact coverage. Relevant earlier proof notes are ALGEBRA\_FORMAL\_SMOOTH\_NOTES\_20260928.md, ALGEBRA\_SYNTOMIC\_NOTES\_20260928.md and ALGEBRA\_SMOOTH\_NOTES\_20260928.md, each hash-bound in the receipt. No external exhaustiveness or novelty claim is made.

Next authority line: 39754, Local structure of étale ring maps. Only the opening 39754--39765 has been read, without adjudication. The full claim ledger records the receiving links; broader synthesis and its paper remain follow-on work after the core review. No translation mutation, TeX run or publication is performed by this intake.

\hypertarget{algebra_etale_local_notes_20260928}{}
\subsection{Local étale presentations, retained scalar factors and finite covers}\label{algebra-proof-052-local-uxe9tale-presentations-retained-scalar-factors-and-finite-covers}

Private source-linked editorial evidence, 2026-09-28. Source: The Stacks Project authors, algebra.tex at a04446e57ec1fbc252a871afcec7752fb2807b14, lines 39754--40187, and the second locus 40463--40470 of report OCC-12163. Full authority SHA-256 FA8BB92E58A4F78A2BD01B3B6A4A87DE0A0D279F5DD90641B574DD5FBFFFA4F3. Original assertions, objects and proof order remain identifiable. The additional proofs below are editorial material; no novelty is claimed.

\subsubsection{Original local-structure proof and the missing scalar}\label{algebra-proof-052-original-local-structure-proof-and-the-missing-scalar}

For a standard étale presentation \(S=R[x]_g/(f)\), keep the given monic \(f\) and the specified invertible image of \(f'\). Writing \[
S=R[x,z]/(f,zg-1)
\] gives the derivative matrix with equation order \((f,zg-1)\), variable order \((x,z)\): \[
\begin{pmatrix} f'&0\\zg'&g\end{pmatrix},
\qquad \det=f'g.
\] Both factors are units in \(S\), proving étaleness. Under base change the same equations, their monicity and the inverse identities persist. If \(s=h/g^a\in S\), its principal localization is \(R[x]_{gh}/(f)\): inverting \(g\) makes \(s\) invertible exactly when \(h\) is invertible. This proves the first three properties at 39781--39792 without changing the definition.

For prescribed residue fields at 39809--39841, put \(K=\kappa(\mathfrak p)\), choose the original primitive element \(\alpha\) of the given finite separable \(L/K\), and let \[
F(X)=X^d+a_1X^{d-1}+\cdots+a_d.
\] Choose representations \(a_i=u_i/v_i\) in \(K\), with \(u_i,v_i\in R\), \(v_i\notin\mathfrak p\), and put \(c=\prod_i v_i\notin\mathfrak p\). If the degree is zero there is no field extension, so here \(d\ge1\). The auxiliary element \(\beta=\bar c\alpha\) has minimal polynomial \[
F_c(X)=\bar c^{\,d}F(X/\bar c)
=X^d+\bar c\,a_1X^{d-1}+\cdots+\bar c^{\,d}a_d.
\] Each coefficient is in the image of \(R\), because \(c^i a_i=u_i(c/v_i)c^{i-1}\). Retain these powers, and recover \(\alpha=\beta/\bar c\) in \(L\). The derivative is \(F_c'(X)=\bar c^{d-1}F'(X/\bar c)\), so \(F_c'(\beta)\ne0\). Lift \(F_c\) to a monic \(f\in R[X]\). The evaluation map \(R[X,1/f']/(f)\to L\) sends \(X\) to \(\beta\) and the inverse of \(f'\) to \(F_c'(\beta)^{-1}\). Its image contains the image of \(R/\mathfrak p\) and \(\beta\); its fraction field is therefore \(K(\beta)=L\). Thus the residue field at its kernel is the required original extension \(L/K\), without asserting that the evaluation map is surjective as a ring map.

In the original eleven-step proof at 39843--40062, Step 2 has a genuine target typo: the descended target is \(S=R\otimes_{R_0}S_0\), not \(R=R\otimes_{R_0}S_0\). This is recorded as a minimal correction, with the original approximation argument retained.

Step 3 uses 30906--30949 to choose a finite \(R\)-subalgebra and a principal open on which it agrees with the original algebra. Membership of its element in a prime of the target means membership of its image under the displayed map. This conventional notation is not an ill-typed assertion. Removing the repeated ``such that'' repairs the sentence while leaving that map explicit in the argument.

Steps 4--6 retain the fibre decomposition into local Artinian \(K\)-algebras, with distinguished factor \(A_1=\kappa(\mathfrak q)\). Choose a nonzero primitive element \(\alpha\) in that field. In the fibre choose \((\alpha,0,\ldots,0)\). Since the fibre is a localization of \(S/\mathfrak pS\), multiplying by the image of one element of \(R\setminus\mathfrak p\) makes this element lift to \(t\in S\); the resulting nonzero scalar multiple of \(\alpha\) remains primitive. Set \(T=R[t]\subset S\), using the source's \(S'\) for this ring when reading the passage. The other fibre primes contract to \(x=0\), whereas the distinguished one does not because \(\alpha\ne0\). Thus the distinguished prime is the only one over its contraction \(\mathfrak q'\subset T\).

Finiteness makes \(T\to S\) integral. The unique-prime localization result at 9797--9830 identifies \(S_{\mathfrak q'}=S_{\mathfrak q}\). The injective map \(T_{\mathfrak q'}\to S_{\mathfrak q}\) is finite, and \[
\mathfrak pS_{\mathfrak q}
=\mathfrak qS_{\mathfrak q}
=\mathfrak q'S_{\mathfrak q}.
\] The last equality follows because the middle ideal contains \(\mathfrak q'S_{\mathfrak q}\), which contains the first ideal. The image of \(\alpha\) shows \(\kappa(\mathfrak q')=\kappa(\mathfrak q)\). Nakayama applied to the finite cokernel now proves \(T_{\mathfrak q'}=S_{\mathfrak q}\). The source's Step 7 extends this to principal neighbourhoods. Its dependency 31623--31660 is read with the already retained corrections MC-STK-ERR-0669 (invert the idempotent selecting the \(S\) factor) and MC-STK-ERR-0672 (both missing polynomial-presentation replacement verbs). No duplicate correction is submitted here.

\textbf{The scalar error.} In Step 8, let \(H\in K[x]\) be the initial monic generator of the fibre ideal. The argument multiplies it by a nonzero scalar \(\lambda\in K^\times\) so that its replacement is the image of a particular \(h\in I\subset R[x]\). This scalar cannot subsequently disappear. With the original monic irreducible factors and their original exponents, the actual formula is \[
\overline h
=\lambda\overline h_1\overline h_2^{e_2}\cdots\overline h_n^{e_n}.
\] Set \(e_1=1\), making the following uniform formula for the rings \(A_i\) well-defined. The unit \(\lambda\) changes neither the ideal generated by \(\overline h\) nor its primary factors, but it does change its polynomial coefficients.

Keep the original monic \(m\in I\), its factorization \[
\overline m=\overline k\prod_{i=1}^n\overline h_i^{d_i},
\qquad d_1\ge1,\quad d_i\ge e_i\ (i\ge2),
\] and the original integer \(l\ge2\) with \(l\deg m>\deg h\). The actual \(f=m^l+h\) is monic, and its full reduction is \[
\begin{aligned}
\overline f
&=\lambda\overline h_1\prod_{i=2}^n\overline h_i^{e_i}
 +\overline k^{\,l}\prod_{i=1}^n\overline h_i^{ld_i}\\
&=\overline h_1
\left(
\lambda\prod_{i=2}^n\overline h_i^{e_i}
+\overline k^{\,l}\overline h_1^{ld_1-1}
 \prod_{i=2}^n\overline h_i^{ld_i}
\right)
=\overline h_1\overline w.
\end{aligned}
\] Modulo \(\overline h_1\), the second summand inside parentheses is zero because \(ld_1-1\ge1\), whereas the first is nonzero because \(\lambda\ne0\) and the original irreducibles are pairwise coprime. Therefore \(\overline w\) is coprime to \(\overline h_1\). The derivative retains both summands: \[
\overline{f'}=\overline h_1'\overline w+\overline h_1\overline w'.
\] At the specified residue point its value is \(\overline h_1'\overline w\ne0\), by separability. Hence the conclusion of Steps 9--10 remains valid with the scalar restored.

The missing factor is a literal polynomial error, not only a preferred notation. Take \(R=\mathbf Z\), \(\mathfrak p=(3)\), \(S=R[x]/(x-1)\), distinguished residue element \(1\), \(h=2(x-1)\), \(m=x-1\), \(l=2\). Then \(\lambda=2\) in \(\mathbf F_3\) and the actual \(f=m^2+h=x^2-1\). The uncorrected displayed expression instead gives \((x-1)+(x-1)^2=x^2-x\). Both satisfy a simple-root argument at the chosen point, but they are different polynomials. The repair therefore keeps the actual \(h,m,f\), inserts the missing \(\lambda\) at 39987, 40013 and 40018, and states \(\lambda\ne0\) at 39978.

In Step 11, retain \(g=f'\) and the separate element \(g'\) selecting an étale neighbourhood of \(S\). The resulting surjection \[
R[x]_{gg'}/(f)\longrightarrow S_{\varphi(gg')}
\] is between étale \(R\)-algebras, so is étale, flat and finitely presented. Its kernel is generated by an idempotent; the quotient is a principal localization at the complementary idempotent. Any element of the left side is represented by a polynomial divided by a power of \(gg'\), so further principal localization still has the original standard étale form. The local-structure theorem follows. The repaired formulas and the complete explanation remain visible editorial material; they are not grounds for replacing the original eleven-step translation by a new proof.

\subsubsection{Monogenic étale algebras over fields}\label{algebra-proof-052-monogenic-uxe9tale-algebras-over-fields}

For a field \(k\), a finite étale \(k\)-algebra \(A\) is standard étale if and only if it is generated by one element as a \(k\)-algebra. To prove the nontrivial forward assertion, write \(A=k[x]_g/(f)\), and let \(a\) be the image of \(x\). The subalgebra \(B=k[a]\subset A\) is finite-dimensional because \(f\) is monic. Multiplication by \(g(a)\) on \(B\) is injective, since \(g(a)\) is invertible in \(A\). It is consequently surjective on the finite-dimensional vector space \(B\), so \(g(a)b=1\) for some \(b\in B\). Thus \(g(a)^{-1}\in B\) and \(A=B\). Conversely, if \(A=k[a]\), write \(A=k[x]/(F)\) with \(F\) monic. The finite étale field classification implies that \(F\) is a product of distinct separable irreducible polynomials, so \(\gcd(F,F')=1\). Hence \(F'\) is invertible in \(A\), giving a standard étale presentation with \(g=1\). For \(A=0\), use \(F=1\), whose derivative is the unit \(0=1\) of the zero ring.

Let now \(k=\mathbf F_q\) and \[
A=\prod_{d\ge1}\bigl(\mathbf F_{q^d}\bigr)^{r_d},
\] with only finitely many nonzero \(r_d\). Let \(N_q(d)\) be the number of monic irreducible polynomials of degree \(d\) in \(\mathbf F_q[x]\). Then the exact criterion is \[
A\text{ is standard étale over }\mathbf F_q
\quad\Longleftrightarrow\quad
r_d\le N_q(d)\text{ for every }d.
\] Indeed an element generating \(A\) must generate each factor field, and its minimal polynomials in distinct factors must be pairwise different: the evaluation map onto the product is surjective exactly when its kernel ideals are pairwise comaximal. Distinct monic irreducible polynomials are comaximal, while equal ones are not. For each degree \(d\), therefore, one needs \(r_d\) distinct irreducibles of degree \(d\). Conversely choose precisely that many distinct irreducibles for every degree and use the Chinese remainder theorem on their product. Each quotient is the corresponding finite field, proving sufficiency.

The numerical bound is \[
N_q(d)=\frac1d\sum_{e\mid d}\mu(e)\,q^{d/e}.
\] Here \(\mu(e)\) is zero when a prime square divides \(e\), and otherwise equals \((-1)^s\) when \(e\) has \(s\) distinct prime factors. For completeness, the roots of \(x^{q^a}-x\) in an algebraic closure form a field with \(q^a\) elements: addition, multiplication and inverses are preserved by the \(q^a\)-power map, and the derivative \(-1\) gives exactly \(q^a\) distinct roots. Frobenius permutes the roots of a monic irreducible of degree \(d\) in one orbit of length \(d\). To justify the orbit length, its distinct Frobenius orbit product has coefficients fixed by \(z\mapsto z^q\), hence in \(\mathbf F_q\); minimality forces that product to be the given irreducible. Thus this irreducible divides \(x^{q^a}-x\) exactly when \(d\mid a\). Factoring that square-free polynomial gives \(q^a=\sum_{d\mid a}dN_q(d)\). Multiplying these equalities by \(\mu(e)\) over divisors \(e\mid d\) and using \(\sum_{e\mid m}\mu(e)=0\) for \(m>1\), \(=1\) for \(m=1\), yields the displayed formula. The divisor identity follows by expanding \(\prod_{\ell\mid m}(1-1)\) over the distinct prime divisors \(\ell\).

In the original example \(q=2,d=2\), \(N_2(2)=(4-2)/2=1\). Therefore \(\mathbf F_4^4\) is not standard étale over \(\mathbf F_2\), although \(\mathbf F_2\to\mathbf F_4\) has polynomial \(x^2+x+1\) and the diagonal \(\mathbf F_4\to\mathbf F_4^4\) has polynomial \(x^4-x\), with invertible derivative \(1\). The example is correct.

Two copies already suffice: the composition \(\mathbf F_2\to\mathbf F_4\to\mathbf F_4^2\) uses \(x^2+x+1\) for the first map and \(x(x-1)\) for the second. The derivative \(2x-1=1\) of the latter is invertible in characteristic two. But \(r_2=2>N_2(2)\), so the composite is not standard étale. This is a smaller valid counterexample, kept as a separate editorial consequence.

The obstruction disappears over infinite fields. Write a finite étale algebra as \(\prod_{i=1}^r L_i\) with each \(L_i/k\) finite separable, and choose a primitive element \(a_i\in L_i\). Choose constants \(c_i\in k\) successively so that the sets of conjugates of \(a_i+c_i\) are pairwise disjoint. At stage \(i\), the forbidden equalities are \[
\sigma(a_i)+c_i=\tau(a_j)+c_j,\qquad j<i,
\] over the finitely many embeddings into an algebraic closure. Each excludes at most one value of \(c_i\) in \(k\), so an infinite field supplies a permitted value. Within one factor the conjugates remain distinct. Therefore the shifted elements have pairwise different separable minimal polynomials, and the tuple \((a_i+c_i)_i\) generates the product by the Chinese remainder theorem. The original elements are recovered as \(a_i=(a_i+c_i)-c_i\). Hence every étale algebra over an infinite field is standard étale, using the earlier theorem that such algebras are finite. Compositions of standard étale maps over an infinite base field are accordingly standard étale. This statement concerns the base field and does not assert a general ring version.

\subsubsection{Explicit finite free covers and every chosen fibre point}\label{algebra-proof-052-explicit-finite-free-covers-and-every-chosen-fibre-point}

The end-of-section finite projective cover can be taken finite free, syntomic, and of a specified rank. For a given standard presentation \(S=R[x]_g/(f)\) of monic degree \(n\), construct the ordered-root algebra of the original polynomial, retaining all coefficients.

Put \(R_0=R\), \(Q_0(X)=f(X)\), and for \(j=1,\ldots,n\) define \[
R_j=R_{j-1}[\alpha_j]/(Q_{j-1}(\alpha_j)),
\qquad
Q_{j-1}(X)=(X-\alpha_j)Q_j(X)\text{ in }R_j[X].
\] The polynomial \(Q_j\) exists by monic division. If \(Q_{j-1}=X^d+u_1X^{d-1}+\cdots+u_d\), its quotient coefficients are \(v_0=1\), \[
v_k=u_k+\alpha_jv_{k-1}\ (1\le k<d),
\qquad u_d=-\alpha_jv_{d-1}.
\] Thus no coefficient or constant relation has been suppressed. Monic division gives the free \(R\)-basis \[
\alpha_1^{a_1}\cdots\alpha_n^{a_n},
\qquad 0\le a_j<n-j+1,
\] of \(R_n\), of rank \(n!\). The case \(n=0\) uses the empty tower \(R_n=R\), empty basis product \(1\), and rank \(0!=1\). Each step has a monic relation, hence a non-zero-divisor in its polynomial ring when that ring is nonzero, and after every residue-field base change has a zero-dimensional hypersurface fibre. Therefore each step is a relative global complete intersection, flat and syntomic. Composition is syntomic. The positive free rank makes \(R\to R_n\) faithfully flat. These assertions also follow from the exact ordered-root calculation in group 1042 and lemma-adjoin-roots at 36857--36875; the recursive proof here specifies the rank and basis used in the present cover.

Take \(D=R_n\). For every specified \(\mathfrak q\in\operatorname{Spec}S\) over \(\mathfrak p\in\operatorname{Spec}R\), and every \(\mathfrak q'\in\operatorname{Spec}D\) over the same \(\mathfrak p\), one can choose a principal neighbourhood of \(\mathfrak q'\) and a map from \(S\) whose contraction is exactly the specified \(\mathfrak q\). In fact put \(K=\kappa(\mathfrak p)\) and \(K'=\kappa(\mathfrak q')\). The prime \(\mathfrak q\) corresponds to an irreducible factor \(f_1\) of \(\bar f\in K[x]\) with \(f_1\nmid\bar g\). The polynomial \(\bar f\) splits in \(K'\) with roots the images of the actual \(\alpha_j\). Since \(f_1\) divides it, some root \(\bar\alpha_j\) satisfies \(f_1(\bar\alpha_j)=0\). Its minimal polynomial over \(K\) is \(f_1\). A Bézout identity for \(f_1,\bar g\), evaluated at this root, proves \(\bar g(\bar\alpha_j)\ne0\).

The map \(R[x]/(f)\to D\) sending \(x\) to \(\alpha_j\) therefore extends to \[
S\longrightarrow D_{\,g(\alpha_j)}.
\] Its contraction of \(\mathfrak q'D_{g(\alpha_j)}\) is \(\mathfrak q\), because its residue evaluation has exactly the irreducible kernel \((f_1)\) over \(K\). This proves both conditions the source leaves implicit at 40125--40127.

For an arbitrary étale \(R\)-algebra \(S\), choose a finite principal cover \(S_{s_i}\) by standard presentations of degrees \(n_i\), and perform the preceding construction for each, giving \(D_i\). Put \[
D=\bigotimes_{i=1}^r D_i.
\] The tensor products of the displayed ordered-root monomial bases give an actual free \(R\)-basis of cardinality \[
\prod_{i=1}^r n_i!.
\] This finite free map is faithfully flat and syntomic, and every construction commutes with arbitrary base change because it uses the same polynomial relations and monic divisions.

More precisely, no surjectivity assumption on \(\operatorname{Spec}S\to\operatorname{Spec}R\) is needed for the following statement. For each pair of primes \(\mathfrak q\in\operatorname{Spec}S\), \(\mathfrak q'\in\operatorname{Spec}D\) lying over the same \(\mathfrak p\), there is \(h\in D\setminus\mathfrak q'\) and an \(R\)-map \(S\to D_h\) whose contraction of \(\mathfrak q'D_h\) is exactly \(\mathfrak q\). Choose \(i\) with \(s_i\notin\mathfrak q\), and contract \(\mathfrak q'\) to \(D_i\). The previous paragraph supplies \(h_i\) outside that contracted prime and \(S_{s_i}\to(D_i)_{h_i}\) with the prescribed contraction. Its base change to \(D_{h_i}\), preceded by \(S\to S_{s_i}\), gives the required map. The element \(h_i\) remains outside \(\mathfrak q'\) by the definition of contraction.

Consequently the set of points of \(\operatorname{Spec}D\) admitting such a local factorization is exactly the inverse image of the image of \(\operatorname{Spec}S\). One inclusion is the construction just proved. Conversely any such factorization supplies a prime of \(S\) by contraction, lying over the same base prime. The image of an étale map is open, since it is flat and finitely presented, so this is an open subset of \(\operatorname{Spec}D\). When the original map is surjective on spectra it is the whole space, giving the source's last lemma with the stronger finite free, syntomic conclusion and the explicit rank above. If \(S=0\), use the empty cover and \(D=R\); the relevant open subset is empty and all assertions about specified pairs of primes are vacuous. No claim is made that \(D/R\) is étale, nor that a general nonfinite étale \(S\) becomes a finite product of copies of the base on these neighbourhoods.

\subsubsection{Direct finite-cokernel comparison over an arbitrary base}\label{algebra-proof-052-direct-finite-cokernel-comparison-over-an-arbitrary-base}

There is also a direct variant of the local-structure proof which avoids Noetherian approximation. This does not strengthen the already arbitrary-base conclusion of the source theorem; it weakens hypotheses in an intermediate argument and gives a separate proof of that same conclusion.

First suppose \(A\hookrightarrow B\) is a finite injective ring map, \(\mathfrak p\in\operatorname{Spec}A\) has just one prime \(\mathfrak q\) of \(B\) above it, and \(A_{\mathfrak p}\to B_{\mathfrak q}\) is an isomorphism. Finiteness makes \(B_{\mathfrak p}\) integral over the local ring \(A_{\mathfrak p}\). Its maximal ideals all lie over the maximal ideal of \(A_{\mathfrak p}\), so the specified uniqueness makes \(B_{\mathfrak p}\) local with maximal ideal induced by \(\mathfrak q\). Hence \(B_{\mathfrak p}=B_{\mathfrak q}\).

The cokernel \(C=B/A\) is a finitely generated \(A\)-module: any finite set of \(A\)-generators of \(B\) maps to generators of \(C\). Its localization at \(\mathfrak p\) is zero by the assumed isomorphism. For actual generators \(c_1,\ldots,c_a\), choose \(u_i\in A\setminus\mathfrak p\) with \(u_ic_i=0\), and set \(u=\prod_i u_i\). Then \(u\notin\mathfrak p\) and \(C_u=0\). Localization preserves the original injection, so \[
A_u\xrightarrow{\;\sim\;}B_u.
\] No Noetherian assumption and no finite-presentation hypothesis on \(A\) or \(B\) was used. The uniqueness of the prime is essential to this stated argument; a local isomorphism at just one of several points above \(\mathfrak p\) does not imply \(C_{\mathfrak p}=0\).

Apply this to the source after Step 1. The finite-algebra neighbourhood in Step 3 exists over an arbitrary base by 30906--30949. Choose the fibre element and \(t\) exactly as in Steps 4--5, and put \(T=R[t]\subset S\) for the resulting finite \(R\)-algebra \(S\). The element \(t\) is integral over \(R\), since \(S/R\) is finite; the powers of \(t\) up to one less than the degree of an actual monic relation span \(T\). Thus \(T/R\) is finite without invoking the generally false assertion that every submodule of a finite module is finite over a non-Noetherian ring.

The argument in Step 6 uses only finiteness, injectivity, uniqueness of the prime over \(\mathfrak q'\), the original étale residue-field conclusions and Nakayama for a finite cokernel. All remain valid over the arbitrary base. The finite-cokernel comparison just proved then replaces Step 7, with an explicit \(u\in T\setminus\mathfrak q'\) giving \(T_u=S_u\). Consequently \(T\) is étale at \(\mathfrak q'\), is finite over \(R\), and is a quotient \(R[x]/I\).

Steps 8--10 require only a principal ideal after passing to the field \(\kappa(\mathfrak p)\), a monic integral relation and a lift of a scalar multiple of that field-ideal generator. These hold without Noetherianity. To justify the lift explicitly, write the field generator as a finite combination of images of elements of \(I\), with polynomial coefficients in \(\kappa(\mathfrak p)[x]\); multiply by a product of their finitely many nonzero denominators from \(R/\mathfrak p\). This gives precisely the nonzero scalar \(\lambda\) and \(h\in I\) used above, with all three scalar factors retained. Step 11 uses a standard étale algebra and a principal neighbourhood of \(T\) already known to be étale; these are finitely presented by definition. Its surjection therefore has finitely generated idempotent kernel, completing the same standard étale neighbourhood construction.

One final localization detail keeps the statement's original source ring. When replacing a ring by an isomorphic principal neighbourhood, any further principal localization is still a principal localization of that original ring: if the first denominator is \(a\) and the new element is \(b/a^N\), their combined localization is the localization at \(ab\). The relevant point excludes both factors. This proves that the successive finite-algebra and monogenic comparisons return an actual \(S_g\) in the original theorem, rather than only an unrelated locally isomorphic algebra.

\subsubsection{Report accounting and propagation}\label{algebra-proof-052-report-accounting-and-propagation}

Nineteen received reports have sixteen report groups. No existing canonical operation lies in 39754--40187 or in 40463--40470; both current intervals agree with authority.

\begin{itemize}
\tightlist
\item
  OCC-00880: optional mathematical source whitespace, no edit.
\item
  OCC-12163: reject the alleged mismatch between extension notations \(L/K\) and \(K\subset L\). Both name the same field-extension data, not a quotient ring; the two source loci have been read. The second locus does not adjudicate the rest of the future section.
\item
  OCC-12164 and OCC-00881: one target-algebra correction at 39864, \(R\) to \(S\).
\item
  OCC-00882: ``Denote by'' at 39864.
\item
  OCC-12165, OCC-12166 and OCC-00883: one correction to the repeated connective at 39876. Membership is interpreted through the already specified ring map; that conventional shorthand is not a second mathematical defect.
\item
  OCC-00884: punctuation at the end of the fibre-product display, 39891.
\item
  OCC-00885: ``Denote by'' at 39931.
\item
  OCC-00886: reject the assertion that ``we have going up'' is grammatically malformed. It names a property, as does ``we have flatness''; the actual finite integral map has that property.
\item
  OCC-00887: supply ``is'' in the finite-algebra condition at 39968.
\item
  OCC-00888: explicitly set \(e_1=1\) at 39991 for the following formula indexed by all \(i\).
\item
  OCC-00889: ``Denote by'' at 40000.
\item
  OCC-00890: optional enumeration colon, no edit.
\item
  OCC-00891 and OCC-00892: separate ``Denote by'' corrections at 40117 and 40167.
\item
  OCC-00893: supply ``to be'' in the contracted-prime definition at 40171.
\item
  OCC-00894: optional restatement of the inherited nonmembership of \(g_i'\) in \(\mathfrak q_i'\), no source edit. The preceding lemma supplies that condition, and contraction proves the required nonmembership in \(\mathfrak q'\). An omitted repeated condition is not a missing proof here.
\end{itemize}

Group 1145 adds the four-operation scalar repair and propagates it through every displayed factorization of the same original polynomial. It does not change the theorem. Groups 1146--1148 have zero operations: the field-generator criterion and exact finite-field bound; the finite free syntomic cover with its rank and prescribed fibre points; and the finite-cokernel comparison with the direct arbitrary-base proof. The last is an alternative argument for the source's existing general theorem, not a newly generalized statement of that theorem.

Receiving dependencies are the source's étale field classification, finite flat quotient and local lifting arguments, group 1042's full ordered-root basis, and group 1128's factor and multiplicity analysis. Earlier idempotent and presentation corrections remain retained. Complete mathematical arguments stay in this note and are not substituted for the translation. Broader synthesis and its paper follow the core review.

Dependencies read: 9785--9834, 30896--30966, 31615--31670, 36848--36918; the exact cited portions are 9797--9830, 30906--30949, 31623--31660 and 36857--36875. The new indexed query returned two research routing records, both unread and unused. Seven earlier external-source reading records retain their exact coverage. No new external reading, PDF fallback or novelty claim is asserted.

Next authority line 40188, Étale local structure of quasi-finite ring maps; only opening 40188--40200 and the explicitly scoped second report locus 40463--40470 have been read ahead. No translation mutation, TeX run or publication is performed by this intake.

\hypertarget{algebra_etale_quasifinite_notes_20260928}{}
\subsection{Finite components after étale base change, with full fibre structure}\label{algebra-proof-053-finite-components-after-uxe9tale-base-change-with-full-fibre-structure}

Private source-linked editorial evidence, 2026-09-28. Source: The Stacks Project authors, algebra.tex at a04446e57ec1fbc252a871afcec7752fb2807b14, lines 40188--40496; full authority SHA-256 FA8BB92E58A4F78A2BD01B3B6A4A87DE0A0D279F5DD90641B574DD5FBFFFA4F3. The original integral closures, factor polynomials, prime contractions and fibre rings are retained. Additional arguments are editorial supplements, not replacements for the translated exposition. No novelty is claimed.

\subsubsection{The original finite-component construction}\label{algebra-proof-053-the-original-finite-component-construction}

For 40200--40225 retain \(R\to S'\to S\), the integral first map, the finite type composite, the element \(g\in S'\), the isomorphism \(S'_g\to S_g\), and \(\mathfrak p\subset R\). The image of \(V(g)\subset\operatorname{Spec}S'\) is closed by integrality. More explicitly it is \[
V\bigl(R\cap gS'\bigr),
\] where the intersection means inverse image along \(R\to S'\). Indeed \(R/(R\cap gS')\to S'/gS'\) is injective and integral, so lying over gives exactly that image. The unit condition on \(g\) in \(S'\otimes_R\kappa(\mathfrak p)\) says this image excludes \(\mathfrak p\). Choose an actual \[
f\in (R\cap gS')\setminus\mathfrak p,\qquad f=gt\text{ in }S'
\] for some \(t\in S'\). Then the inverse of \(g\) in \(S'_f\) is \(t/f\). Its image also inverts \(g\) in \(S_f\), so the original localization isomorphism gives \(S'_f=S_f\). This algebra is integral and finite type over \(R_f\), hence finite. Concretely, for algebra generators \(s_1,\ldots,s_a\) and monic integral equations of degrees \(d_1,\ldots,d_a\), the products \(\prod_j s_j^{b_j}\), \(0\le b_j<d_j\), span the algebra as an \(R_f\)-module. This is a spanning bound \(\prod_j d_j\), not a claim of freeness or exact rank.

For 40227--40359 put \(K=\kappa(\mathfrak p)\), \(F=S\otimes_R K\), and let \(S'\subset S\) be the original integral closure, with \(F'=S'\otimes_R K\). The prime \(\mathfrak q\) is quasi-finite exactly when its fibre point \(\overline{\mathfrak q}\) is isolated and closed. The Zariski Main Theorem element \(g\) identifies \(S'_g\) with \(S_g\), hence \(F'_g\) with \(F_g\). The corresponding point \(\overline{\mathfrak q}'\) of \(F'\) is closed because every domain integral over a field is a field: each nonzero element lies in a finite-dimensional domain subalgebra, where injective multiplication is surjective and supplies its inverse. The same point is open by the indicated open-neighbourhood isomorphism. Thus it determines a clopen factor \[
F'=F'_1\times F'_2,\qquad
\operatorname{Spec}F'_1=\{\overline{\mathfrak q}'\}.
\] The bar in the closedness clause at 40272 is essential: the claim concerns a point of the fibre spectrum, not a closed point of \(\operatorname{Spec}S'\). For example \((0)\subset\mathbf Z\) is not closed, whereas its point in the fibre \(\mathbf Q\) is closed.

The point \(\overline{\mathfrak q}\) is the only point of \(F\) mapping to \(\overline{\mathfrak q}'\), since any such point excludes \(g\) and belongs to the open on which the map is an isomorphism. Lift a scalar multiple of the idempotent \((1,0)\) to \(s'\in S'\): an element of the tensor fibre is represented by an element of \(S'\) divided by the image of an element of \(R\setminus\mathfrak p\). Thus there are \(r\in R\setminus\mathfrak p\) and \(s'\) whose actual image is \((\bar r,0)\), with \(\bar r\in K^\times\). This retains the scalar \(\bar r\); it is not silently replaced by one.

Keep a monic \(f(x)\in R[x]\) annihilating \(s'\). Its reduction factors as \(x^e\bar h\) with \(\bar h(0)\ne0\); multiplying the actual \(f\) by \(x\) when needed makes \(e\ge1\) and still annihilates \(s'\). Because \(f(\bar r)=0\) in the nonzero first factor and \(\bar r\ne0\), the factor \(\bar h\) has positive degree. The source's coprime-factor lifting produces an étale \(R\to R'\), a prime \(\mathfrak p'\) with its specified residue-field identity \(K=\kappa(\mathfrak p')\), and actual polynomials \[
f=hi,\qquad \bar i=x^e,\qquad \bar h=h\bmod\mathfrak p',
\qquad ah+bi=1.
\] All symbols refer to these actual polynomials and their evaluations, including the source's order of the two factors.

In \(D'=R'\otimes_R S'\), set \[
u=h(s'),\quad v=i(s'),\quad
\varepsilon=b(s')v,\quad 1-\varepsilon=a(s')u.
\] The relations \(uv=0\) and \(a(s')u+b(s')v=1\) give \(\varepsilon(1-\varepsilon)=0\), so \(\varepsilon^2=\varepsilon\). The exact product isomorphism is \[
D'\longrightarrow \varepsilon D'\times(1-\varepsilon)D',
\qquad z\longmapsto(\varepsilon z,(1-\varepsilon)z),
\] with inverse \((z_1,z_2)\mapsto z_1+z_2\). The first factor is \(D'_v\): in \(\varepsilon D'\), \(v\) has inverse \(b(s')\varepsilon\), and in \(D'_v\), the relation \(uv=0\) gives \(u=0\), hence \(\varepsilon=1\). Similarly the second factor is \(D'_u\).

The image of \(v\) in \(F'\) is exactly \((\bar r^{\,e},0)\). Therefore \(\varepsilon\) reduces to \((1,0)\), and the two factors reduce to the original \(F'_1,F'_2\) with their entire ring structures. The first factor \(S'_1=\varepsilon D'\) is integral over \(R'\), because it is a quotient of the integral algebra \(D'\). This quotient argument matters: arbitrary localizations of integral algebras need not remain integral.

The same idempotent maps into \(D=R'\otimes_R S\), giving \(D=S_1\times S_2\). Its first fibre factor has the one specified point, and the second has no point mapping to the original \(\mathfrak q\). The original \(g\) is a unit in \(S'_1\otimes_{R'}K\): it excludes the only prime of that fibre. The first lemma then applies to \(R'\to S'_1\to S_1\). Here \(S_1\) is finite type, being an idempotent quotient of \(D\), and the required \(g\)-localization isomorphism is inherited from \(S'_g=S_g\). Localizing \(R'\) at the actual resulting element outside \(\mathfrak p'\) makes \(S_1\) finite without changing the pointed residue field or the fibre decomposition.

For 40361--40418, the finite type fibre \(F\) is Noetherian. Each isolated closed point is an irreducible component, so there are only finitely many. Extract one, then repeat on the complementary algebra after the pointed étale base change. Finiteness of the already extracted algebra persists under every subsequent base change. Because each selected residue field is identified with the original \(K\), the fibre ring and the number of its remaining isolated closed points stay exactly the same except for the one removed component. This proves termination after the original finite number, and gives the displayed decomposition with no quasi-finite point of the remainder over the distinguished base point. It also checks the omitted verification at 40411--40417.

\subsubsection{Full Artinian fibres and a selected set of components}\label{algebra-proof-053-full-artinian-fibres-and-a-selected-set-of-components}

Let \(\mathfrak m_1,\ldots,\mathfrak m_n\) be the isolated closed points of \(F=S\otimes_R K\). Each corresponds to a clopen component, and \[
C_i=F_{\mathfrak m_i}
\] is a finite-dimensional local Artinian \(K\)-algebra. One way to see finite dimension is the finite-type isolated-point criterion: its clopen factor is a finite type zero-dimensional \(K\)-algebra, hence Artinian with finite residue extension and a finite composition series. In the one-point construction above the first fibre factor is exactly this clopen factor, not only a field with the same point. Let \(\eta_i\in F\) be its original indicator idempotent, and let \(\pi_i:F\to A_i\otimes_{R'}K\) be the fibre of the resulting projection. The explicit idempotent calculation gives \(\ker\pi_i=(1-\eta_i)F\). Thus \(\pi_i\) restricted to \(\eta_iF\) is an isomorphism with inverse \(\pi_i(x)\mapsto\eta_i x\). The original localization identifies the one-point factor \(\eta_iF\) with \(F_{\mathfrak m_i}\). Their composition gives the specified \(K\)-algebra isomorphism \[
A_i\otimes_{R'}K\xrightarrow{\;\sim\;}C_i
\] which sends \(\pi_i(x)\) to the image of \(\eta_i x\) in \(F_{\mathfrak m_i}\). It retains the original residue map, every nilpotent element, the maximal ideal and all its powers.

More generally choose any subset \(T\subset\{1,\ldots,n\}\). Repeating just those extraction steps gives a pointed étale base change with residue field \(K\) and \[
R'\otimes_R S=
\left(\prod_{i\in T}A_i\right)\times B_T,
\quad A_i/R'\text{ finite},\quad A_i\otimes_{R'}K=C_i,
\] where \(B_T\otimes_{R'}K\) is the complementary original clopen factor of \(F\). Each step removes precisely the chosen point, because its idempotent reduces to the prescribed indicator idempotent. Thus the other chosen points remain in the complementary fibre with their original local algebras, so induction applies. For \(T=\varnothing\), use \(R'=R\), the empty product \(0\), and \(B_T=S\). For all points, the complementary fibre has no isolated closed point. The identities of the fibres are canonical relative to the specified residue-field identity and chosen labelled components; the lifted algebras away from that fibre are not claimed unique.

In particular, if \(E_i=C_i/\mathfrak m_iC_i\) and \(\ell_i=\operatorname{length}_{C_i}C_i\), the finite block has \[
\dim_K(A_i\otimes_{R'}K)=\ell_i[E_i:K],
\] and the complete filtration by powers of its fibre maximal ideal agrees with that of \(C_i\). This does not assert that \(A_i\) is flat or projective over \(R'\), or identify this fibre dimension with a constant rank.

The remainder cannot be discarded merely because its distinguished fibre is empty. For a field \(k\), take \[
R=k[t],\quad \mathfrak p=(t),\quad
S=R\times R[1/t].
\] The map is étale, quasi-finite, and surjective on spectra. The source fibre at \(\mathfrak p\) is \(k\), and the component \(R\) already gives the desired finite block; the remainder \(R[1/t]\) has zero fibre there.

Nevertheless no pointed étale base change \(R\to R'\), \(\mathfrak p'\) over \((t)\), makes the entire \(R'\otimes_R S\) finite on a base neighbourhood of \(\mathfrak p'\). If it did, its direct factor would give a finite module \[
M=R'_{\mathfrak p'}[1/t]
\] over the local ring \(R'_{\mathfrak p'}\). Since \(t\) lies in that local ring's maximal ideal and acts invertibly on \(M\), its residue quotient is zero. Nakayama would imply \(M=0\). But the local map \(R_{(t)}\to R'_{\mathfrak p'}\) is flat, and \(t\) is a non-zero-divisor on the nonzero target. No power of \(t\) is zero, so its localization \(M\) is nonzero, a contradiction. Further base localization outside \(\mathfrak p'\) does not change this local argument. The exact direct factor and its punctured-base support explain the obstruction; this is a boundary on a possible overstatement, not a claimed error in the source's deliberately qualified introduction.

\subsubsection{Separable splitting, inseparable degree and every fibre length}\label{algebra-proof-053-separable-splitting-inseparable-degree-and-every-fibre-length}

In 40420--40489 set \(K=\kappa(\mathfrak p)\) and retain the local Artinian factors \(C_i\), their residue fields \(E_i=\kappa(\mathfrak q_i)\), and their lengths \(\ell_i\). Let \(K_i\) be the maximal subfield of \(E_i\) separable over \(K\), and put \[
s_i=[K_i:K],\qquad d_i=[E_i:K_i].
\] The extension \(E_i/K_i\) is purely inseparable. In characteristic \(p>0\), for each element \(\alpha\in E_i\), its minimal polynomial is a separable irreducible polynomial in \(X^{p^a}\) for some \(a\), so \(\alpha^{p^a}\) is separable over \(K\) and lies in \(K_i\). Finitely many generators give a finite purely inseparable extension. In characteristic zero, \(K_i=E_i\) and \(d_i=1\).

Take the source's finite Galois extension \(L/K\) containing the conjugates of all \(K_i\). Primitive-element factorization and the Chinese remainder theorem give \[
K_i\otimes_K L\cong
\prod_{\sigma:K_i\hookrightarrow L}L,\qquad
E_i\otimes_K L\cong
\prod_{\sigma:K_i\hookrightarrow L}
\bigl(E_i\otimes_{K_i,\sigma}L\bigr).
\] There are exactly \(s_i\) embeddings \(\sigma\). Every factor \[
E_{i,\sigma}=E_i\otimes_{K_i,\sigma}L
\] is a field purely inseparable of degree exactly \(d_i\) over \(L\); in particular there is no discarded nilpotent part in this residue-field tensor product.

Here is a proof of that field and degree assertion which covers non-simple purely inseparable extensions. Express \(E_i/K_i\) as a finite tower adjoining one purely inseparable element at a time. At a step the original minimal polynomial is \(X^{p^a}-c\). Over a separable extension of the preceding field it remains irreducible: if the minimal polynomial of its unique root shortened to \(X^{p^b}-u\) with \(b<a\), then \(u\) would lie in that separable extension while being purely inseparable over the original field, hence lie in the original field. That would contradict the original minimal degree. An irreducible polynomial with only one distinct root necessarily has the displayed pure-power form, by repeatedly removing zero derivatives. Thus adjoining the root still gives a field of the same degree. The base-changed extension of the preceding field is separable at each step, since separability is preserved by base change and by taking a field factor. Induction proves the assertion, retaining the product of all original tower degrees.

Now examine the entire \(C_i\), not only \(E_i\). Its maximal ideal \(\mathfrak n_i\) is nilpotent. The algebra \(C_i\otimes_K L\) has the nilpotent ideal \(\mathfrak n_i\otimes_K L\), and its quotient is the product of the fields \(E_{i,\sigma}\). The indicator idempotents of this product have unique orthogonal lifts, giving \[
C_i\otimes_K L=
\prod_{\sigma:K_i\hookrightarrow L} C_{i,\sigma}.
\] Each \(C_{i,\sigma}\) is local Artinian with residue field \(E_{i,\sigma}\).

For clarity the needed idempotent lifting uses no inverse of \(2\). If \(J^N=0\) in a commutative ring and \(a\) lifts an idempotent modulo \(J\), put \(r=a^2-a\in J\). The element \(u=2a-1\) is invertible since \[
u^2=1+4r,\qquad
(1+4r)^{-1}=\sum_{j=0}^{N-1}(-4r)^j.
\] Set \(a_{\mathrm{new}}=a-r/u\). Direct expansion gives \[
a_{\mathrm{new}}^2-a_{\mathrm{new}}=r^2/u^2.
\] Iteration reaches an idempotent once \(2^b\ge N\). If two idempotents \(e,f\) have the same reduction, then \(e(1-f)\) and \(f(1-e)\) are idempotents in the nilpotent ideal, hence zero; so \(e=f\). Products of lifts of distinct indicators therefore vanish, and their sum is the unique lift of \(1\), namely \(1\). This proves the full product decomposition, including in characteristic two.

Choose a composition series of the original \(C_i\)-module \(C_i\); all \(\ell_i\) simple factors are \(E_i\). Tensoring with the field \(L\) preserves exactness. Multiplication by each lifted indicator idempotent is an exact projection onto a direct summand, so projecting that series gives \(\ell_i\) nonzero simple factors, each \(E_{i,\sigma}\), in \(C_{i,\sigma}\). Thus \[
\operatorname{length}_{C_{i,\sigma}}C_{i,\sigma}=\ell_i,
\qquad
\dim_L C_{i,\sigma}=\ell_i d_i.
\] Summing the \(s_i\) factors recovers \[
\dim_L(C_i\otimes_K L)
=s_i\ell_i d_i=\dim_K C_i.
\] This keeps the separable degree, purely inseparable degree and every original Artinian multiplicity separately.

The residue extension \(L/K\) is realized by the actual standard étale construction of 39809--39841, retaining its coefficient-clearing powers and specified evaluation map. The fibre after this base change is \(F\otimes_K L\). Quasi-finiteness is preserved and reflected at every point by 30362--30390, so its isolated closed points are precisely those just described. Applying the residue-preserving component extraction from the preceding section now gives one finite block for every pair \((i,\sigma)\), with distinguished fibre exactly \(C_{i,\sigma}\), residue degree \(d_i\), and length \(\ell_i\). The second étale extension keeps the residue field \(L\) fixed, so none of these data changes. This proves the source's variant and its full multiplicity refinement.

There is a sharp limitation: further étale base changes cannot remove \(d_i>1\). More generally for any finite separable \(L'/K\), the algebra \(K_i\otimes_K L'\) is a product of finite separable fields \(M_j/L'\), each also separable over the corresponding embedded \(K_i\). The same tower proof makes \[
E_i\otimes_K L'=
\prod_j(E_i\otimes_{K_i}M_j)
\] a product of fields purely inseparable of degree \(d_i\) over \(M_j\). Within each factor, \(M_j\) is exactly its maximal separable subfield over \(L'\): an element separable over \(L'\) is also separable over \(M_j\), whereas the whole factor is purely inseparable over \(M_j\). Consequently the inseparable degree at every lifted point remains \(d_i\). The identical composition-series argument preserves length \(\ell_i\) for every local factor under this arbitrary separable base change as well.

Since pointed étale residue extensions are finite separable, one can make all selected residue fields equal to the new base residue field if and only if every original \(E_i/K\) is separable. Necessity follows from the invariant \(d_i\); sufficiency follows by choosing the finite Galois extension containing all \(E_i\). The original Artinian lengths can still exceed one. Rational residue fields do not mean reduced fibres.

Separable base change is necessary for the length assertion. Take \(K=\mathbf F_p(t)\), \(E=K(a)\) with \(a^p=t\), and base change by the purely inseparable extension \(L=E\). In \(E\otimes_K L\), the element \[
u=a\otimes1-1\otimes a
\] satisfies \(u^p=0\), and the map \(L[U]/(U^p)\to E\otimes_K L\), \(U\mapsto u\), is an isomorphism. Indeed the original basis \(1,a,\ldots,a^{p-1}\) tensored with \(L\) is transformed to \(1,u,\ldots,u^{p-1}\) by the triangular formulas \[
(a\otimes1)^j=
\sum_{k=0}^{j}\binom{j}{k}(1\otimes a)^{j-k}u^k,
\] with every diagonal entry \(1\). All coefficients and exponents are retained. The source field of length one becomes a local algebra of length \(p\), with rational residue field; the total dimension remains \(p\).

\subsubsection{Report accounting and propagation}\label{algebra-proof-053-report-accounting-and-propagation}

Twenty-five new received reports have nineteen report groups. The earlier report OCC-12163 about extension notation is already disposed of in group 1130, including its locus 40463--40470; it is not counted again. No existing canonical correction lies in the present source interval, which agrees with authority.

\begin{itemize}
\tightlist
\item
  OCC-00895 supplies ``is'' at 40210; OCC-00896 removes the incorrect indefinite article before \(f\) at 40212.
\item
  OCC-00897 repairs the punctuation of the specific closed-image citation across 40219--40220; adding ``that'' is not required.
\item
  OCC-00898 and OCC-00909 share the missing finite-type predicate repair at 40232 and 40424. The identical unreported omission at 40365 is corrected in the same group.
\item
  OCC-00899, OCC-00904, OCC-00905, OCC-00910 and OCC-00911 are optional enumeration-colon preferences.
\item
  OCC-12172 and OCC-00901 share the missing bar on the fibre prime at 40272.
\item
  OCC-00900 and OCC-00902 supply separate ``Denote by'' corrections at 40259 and 40295.
\item
  OCC-00903 is optional doubled source whitespace.
\item
  OCC-12173, OCC-00906 and OCC-00912 share the two missing ``is'' predicates for the remainder at 40381 and 40442.
\item
  OCC-00907 supplies ``Denote by'' at 40386, and OCC-00908 punctuates the completed displayed decomposition at 40409.
\item
  OCC-12174 and OCC-00914 share ``may'' to ``many'' at 40455.
\item
  OCC-12175 and OCC-00915 share the corrected base prime at 40457, replacing \(R\) by \(\mathfrak p\).
\item
  OCC-00913 supplies ``Denote by'' at 40452.
\end{itemize}

These give sixteen operations in thirteen proposed units and six optional groups. Group 1168 is a zero-operation editorial consequence recording exact Artinian fibres, any selected set of components, and the obstruction to discarding an empty-fibre remainder. Group 1169 records the full separable splitting, invariant inseparable degree and Artinian length, the exact criterion for rational residue fields after étale extension, and the inseparable-base-change counterexample. Original finite blocks are not relabelled as flat or étale algebras.

The receiving notes are ALGEBRA\_ETALE\_NOTES\_20260928.md for the actual coprime-factor lifting and nilpotent structure, and ALGEBRA\_ETALE\_LOCAL\_NOTES\_20260928.md for the prescribed residue-field construction and its full denominator powers. Source dependencies read include 4376--4413, 7774--7804, 9644--9680, 30226--30272 and 30355--30415; exact cited intervals are bound in the evidence receipt. Four new literature-index hits remain unread routing records, while seven earlier external-source reading records retain their coverage. No exhaustive-literature or novelty claim.

Next authority line: 40497, Local homomorphisms. Only opening 40497--40509 has been read, without adjudication. Complete proofs remain here as editorial material; the translation has not been rewritten. Broader synthesis follows the core review. No TeX run or publication is performed by this intake.

\hypertarget{algebra_local_integral_notes_20260928}{}
\subsection{Local homomorphisms and integral closure: source review and editorial arguments}\label{algebra-proof-054-local-homomorphisms-and-integral-closure-source-review-and-editorial-arguments}

\subsubsection{Source identity and reading scope}\label{algebra-proof-054-source-identity-and-reading-scope}

The source is the Stacks Project authors' algebra.tex, commit a04446e57ec1fbc252a871afcec7752fb2807b14, SHA256 FA8BB92E58A4F78A2BD01B3B6A4A87DE0A0D279F5DD90641B574DD5FBFFFA4F3. The complete interval 40497--40768 and its fifteen new reports have been read. The next opening, 40769--40778, is read but unadjudicated. Dependencies read here are 7517--7554, 7596--7627, 9183--9190, 20254--20269 and 31676--31695. The source's Lindel and KC citations at 40509 are retained; those original publications have not been read here. The indexed query ``integral closure smooth'' produced four routing hits, none read or used. Earlier source-reading records retain their actual bounded coverage. No novelty claim is made.

The proofs below are editorial material linked to the source. They do not replace its translated arguments. A compressed reduction, omitted argument or narrower theorem is not automatically a translation error. Earlier editorial arguments are in ALGEBRA\_ETALE\_LOCAL\_NOTES\_20260928.md, especially group 1147, and ALGEBRA\_SYNTOMIC\_NOTES\_20260928.md, group 1042.

\subsubsection{Local homomorphisms and the original parameter}\label{algebra-proof-054-local-homomorphisms-and-the-original-parameter}

At 40506--40554 the source assumes a local map \((R,\mathfrak m_R)\to(S,\mathfrak m_S)\), with \(S=T_{\mathfrak q}\) for an étale \(R\)-algebra \(T\), and the specified residue-field map an isomorphism. Localizing the original standard étale chart at an element outside \(\mathfrak q\) does not change \(S\). Write that chart \(T=R[X]_g/(f)\), with \(f\) monic of degree \(d\), and choose the source's lift \(a\in R\) of the residue of \(X\). The coordinate comparison is exactly \[
Y=X-a,\qquad F(Y)=f(Y+a),\qquad G(Y)=g(Y+a),\qquad
t=F(0)=f(a).
\] Its inverse sends \(X\) to \(Y+a\). Both the quotient relation and the denominator are transported by this map. The resulting prime is \((Y,\mathfrak m_R)\), \(t\in\mathfrak m_R\), and \(G(0)\) is a unit of \(R\), because its residue is nonzero. The original derivative being a unit implies that \(F'(0)=f'(a)\) has nonzero residue and is a unit of \(R\).

For \(d\geq2\), retain all coefficients in \[
F=Y^d+\sum_{i=1}^{d-1}a_iY^i+t,\qquad
H=\sum_{i=1}^{d-1}a_iY^{i-1}+Y^{d-1},\qquad F=t+YH.
\] Here \(a_1=F'(0)\) is a unit, so \(H\notin\mathfrak q\). Replacing the chart by \(T_1=R[Y]_{HG}/(F)\) leaves its local ring \(S\) unchanged. Modulo the actual \(t\), the relation becomes \(YH=0\). Since \(H\) is inverted, \(Y=0\); and the remaining denominator \(H(0)G(0)\) is a unit. Consequently the actual structure map \(R/tR\to T_1/tT_1\) is an isomorphism. For every \(n\geq1\), let \(M_n\) be the cokernel, as an \(R\)-module, of \(R/t^nR\to T_1/t^nT_1\). The mod-\(t\) surjectivity gives \(M_n=tM_n\); as \(t^nM_n=0\), iteration gives \(M_n=0\). This verifies the source's nilpotent surjectivity step without an unstated finiteness requirement.

The map \(R\to S\) is flat and local, hence faithfully flat. Its base change \(R/t^nR\to S/t^nS\) is faithfully flat and therefore injective. It factors through the surjection onto \(T_1/t^nT_1\), so this surjection has zero kernel. Finally \(S/t^nS\) is the localization of \(T_1/t^nT_1=R/t^nR\) at the prime over \(\mathfrak m_R/t^nR\). This prime is the maximal ideal itself; that localization changes nothing. All isomorphisms commute with the maps for decreasing \(n\).

The degree-one boundary requires its own interpretation because the printed \(a_1\) list is empty: then \(F=Y+t\), \(F'=1\), and \(H=1\). The quotient chart is \(R_{G(-t)}=R\), since \(G(-t)\) has the same nonzero residue as \(G(0)\). Thus \(S=R\), and the same \(t=F(0)\) works. A degree-zero monic polynomial is \(1\), whose zero quotient has no such prime. This explanation remains editorial, not a replacement proof.

At 40556--40580 the second lemma does not assume residue-field equality. Choose the source's standard étale chart \(T\) and its prime \(\mathfrak q\) above \(\mathfrak m_R\). Group 1147 gives a finite free faithfully flat syntomic \(R\)-algebra \(D\), of rank \(d!\) for this chart's monic polynomial degree \(d\). This is the ordered-root construction, with its original coefficients and recursive monic divisions; its free basis has exponents \(0\leq e_i<d-i+1\), for \(1\leq i\leq d\). Every maximal ideal \(\mathfrak n\) of \(D\) contracts to \(\mathfrak m_R\): a finite algebra is integral, and contraction of a maximal ideal along an integral map is maximal. Apply the prescribed-pair property of 1147 to \((\mathfrak q,\mathfrak n)\). It gives \(T\to D_u\), \(u\notin\mathfrak n\), whose contraction of the selected prime is exactly \(\mathfrak q\). Every element of \(T\setminus\mathfrak q\) therefore maps to a unit in \(D_{\mathfrak n}\). The universal property gives the required map \(S=T_{\mathfrak q}\to D_{\mathfrak n}\), over the original \(R\). This proves the stronger finite free syntomic choice; it does not claim that all denominators of \(S\) already invert in a single \(D_u\).

\subsubsection{Torsion, completion and an exact gluing square}\label{algebra-proof-054-torsion-completion-and-an-exact-gluing-square}

Keep the exact \(R,S,t\) above. More generally the following proof needs only a faithfully flat \(R\to S\) and the canonical isomorphisms \(R/t^nR\to S/t^nS\) for all \(n\geq1\). For any \(R\)-module \(M\) killed by \(t^n\), \[
S\otimes_R M=(S/t^nS)\otimes_{R/t^nR}M\cong M,
\] with the inverse to the actual unit \(m\mapsto1\otimes m\). If each element of \(M\) is killed by some power of \(t\), then \(M=\mathop{\rm colim}_n M[t^n]\); tensoring commutes with this union, so the same unit is an isomorphism. For an \(S\)-module \(N\) with that property, the action map \(S\otimes_R N\to N\) is inverse to the unit and is \(S\)-linear. These maps prove an equivalence between the two categories of elementwise \(t\)-power torsion modules, by extension and restriction of scalars. Neither finite generation nor a uniform exponent is imposed.

If an ideal \(J\subset R\) contains \(t^a\), quotienting the canonical isomorphism at exponent \(a\) by \(J/t^aR\) gives \(R/J\cong S/JS\). If \(t\in\sqrt I\), choose the actual exponent \(a\) with \(t^a\in I\); then \(t^{an}\in I^n\) for every \(n\). Thus \[
R/I^n\ \xrightarrow{\sim}\ S/I^nS,\qquad
\varprojlim_n R/I^n\ \xrightarrow{\sim}\ \varprojlim_n S/I^nS.
\] These are the original quotient maps, so compatibility in the inverse systems follows. In the local situation \(\mathfrak m_RS=\mathfrak m_S\): the étale fibre localized at the chosen prime is its residue field, so \(S/\mathfrak m_RS\) is a field. In particular the maximal-ideal completions agree, without a Noetherian assumption.

There is also an exact ring comparison retaining the punctured part: \[
\begin{matrix}
R&\longrightarrow&S\\
\downarrow&&\downarrow\\
R[1/t]&\longrightarrow&S[1/t].
\end{matrix}
\tag{LI-1}
\] To prove it is cartesian even when \(t\) is a zero divisor, set \(C=S/R\). Faithful flatness makes \(N\to S\otimes_R N\) injective for every \(N\): after tensoring once more with \(S\), the map has retraction \(n\otimes s\otimes s'\mapsto n\otimes ss'\); faithful flatness reflects injectivity. The exact sequence \(0\to R\to S\to C\to0\) and the Tor exact sequence give \(\operatorname{Tor}_1^R(N,C)=0\) for every \(N\), since \(S\) is flat and that unit is injective. Thus \(C\) is flat. The mod-\(t\) isomorphism gives \(C/tC=0\), or \(C=tC\). Tensor the exact sequence \(0\to\operatorname{Ann}_R(t)\to R\xrightarrow{t}R\) with the flat module \(C\). Its exactness says that \(\ker(t:C\to C)\) is generated by elements \(ac\), with \(at=0\). Write \(c=td\); then \(ac=atd=0\). Therefore multiplication by \(t\) on \(C\) is bijective, and \(C\to C[1/t]\) is an isomorphism.

Localization of \(0\to R\to S\to C\to0\) is exact. Comparing it with the original sequence through \(C=C[1/t]\) gives \[
0\longrightarrow R
\xrightarrow{\,r\mapsto(r,r/1)\,} S\oplus R[1/t]
\xrightarrow{\,(s,b)\mapsto s/1-b\,}S[1/t]
\longrightarrow0.
\tag{LI-2}
\] For surjectivity, the class of an element of \(S[1/t]\) in \(C[1/t]=C\) lifts to an \(s\in S\); subtracting \(s/1\) leaves an element of \(R[1/t]\). For the kernel, an \(s\) whose localized class is zero has zero class already in \(C\), hence comes from \(r\in R\). The injectivity of \(R[1/t]\to S[1/t]\) forces \(b=r/1\). The first map is multiplicative and identifies \(R\) with \(S\times_{S[1/t]}R[1/t]\), proving (LI-1). As \(S[1/t]\) is flat over \(R\), (LI-2) remains exact after tensoring with every \(R\)-algebra \(R_1\). Its kernel is the corresponding ring fibre product, so (LI-1) remains cartesian after every base change. For \(t=0\), the conclusions reduce correctly to \(S=R\) and zero localized rings.

Diagram (LI-1), with its explicitly labelled maps in (LI-2), shows the mechanism: the original rings agree modulo every power of \(t\), while their remaining difference is a module on which \(t\) acts invertibly. This last assertion is not a claim that its stalks vanish on \(V(t)\). No such support assertion is used.

\subsubsection{Integral closure, the derivative and smooth base change}\label{algebra-proof-054-integral-closure-the-derivative-and-smooth-base-change}

At 40589--40628 take the original \(R\to B\), a monic \(f\in R[X]\) of degree \(d\), and \(h\in B[X]/(f)\) satisfying the original monic equation \[
h^e+r_1h^{e-1}+\cdots+r_e=0,\qquad r_i\in R.
\] For \(d>0\), the ordered-root construction gives a finite free faithfully flat extension \(B\hookrightarrow B'\), including \(1\) in its basis, with \(f=\prod_{i=1}^d(X-\alpha_i)\). Each \(\alpha_i\) is integral over \(R\), by this same monic \(f\). Let \(h=\sum_{j=0}^{d-1}h_jX^j\) be its unique monic remainder. The exact identity in \(B'[X]/(f)\) is \[
f'h=\sum_{i=1}^d h(\alpha_i)\prod_{j\ne i}(X-\alpha_j).
\tag{LI-3}
\] Here is the universal argument without an assumption on the actual roots. Over \(U=\mathbf Z[A_1,\ldots,A_d,H_0,\ldots,H_{d-1}]\), take the monic remainder of the difference in (LI-3). Its degree is less than \(d\). In \(\operatorname{Frac}(U)\), the \(A_i\) are distinct; evaluation at each \(A_i\) gives zero because the derivative of the product at \(A_i\) is exactly \(\prod_{j\ne i}(A_i-A_j)\). A polynomial of degree less than \(d\) with those \(d\) roots is zero. Injectivity of \(U\to\operatorname{Frac}(U)\) makes every remainder coefficient zero already in \(U\). Monic division commutes with specialization, proving (LI-3) in \(B'\), including repeated roots and nilpotents. No actual Vandermonde denominator is inverted.

Evaluating the original monic equation proves that every \(h(\alpha_i)\) is integral over \(R\). Finite sets of integral elements generate finite modules: successively reduce each power by its own monic equation, keeping the bounded monomials. Their sums and products are integral, by the determinant argument for multiplication on that finite module containing \(1\). Thus all coefficients on the right of (LI-3) are integral over \(R\). The monic bases \(1,X,\ldots,X^{d-1}\) over \(B\) and \(B'\), and the injection \(B\to B'\), identify those coefficients with the coefficients in \(B\) of \(f'h\). Their monic equations descend through the injection. If \(d=0\), \(f=1\) and the quotient is zero, so there are no coefficients to check.

This proves a precise defect assertion. Set \[
C_0=R[X]/(f),\quad C=B[X]/(f),\quad
A=\overline R^{\,B},\quad D=A[X]/(f)\subset C,\quad
T=\overline {C_0}^{\,C}.
\] The displayed injection follows from the monic bases. Because \(C_0\) is integral over \(R\), transitivity says that \(T\) is also exactly the elements of \(C\) integral over \(R\). Conversely an \(R\)-monic equation is a \(C_0\)-monic equation, so both inclusions hold. Transitivity follows by adjoining the finitely many integral coefficients of a monic equation, then its root, and composing the two finite spanning sets. Every element of \(D\) is integral over \(R\), so \(D\subset T\). Hence \(E=T/D\) is a \(C_0\)-module, and (LI-3) proves \[
f'E=0.
\tag{LI-4}
\] In particular \(D[1/f']=T[1/f']\). No reducedness hypothesis or discarded factor is hidden in this conclusion.

For the source's standard étale \(S=(C_0)_g\), \(f'\) is invertible. Integral closure commutes with localization by 7596--7627, so \(\overline S^{\,S\otimes_RB}=T_g=D_g=S\otimes_RA\). For general étale \(S\), flatness first embeds \(S\otimes_RA\) in \(S\otimes_RB\); its elements are integral over \(S\). The standard étale open cover proves the reverse inclusion on each chart. The quotient module is zero because its localizations on that cover are zero, proving the original canonical isomorphism.

For the Fourier example keep \(R=\mathbf Z[1/p]\) and \[
R'=R[Z]/(Z^{p-1}+\cdots+Z+1),\quad \zeta=[Z],\quad
\alpha_i=\zeta^i\ (0\leq i\leq p-1).
\] The monic basis embeds \(R'\) in \(\mathbf Q[Z]/(\Phi_p)\): a monic remainder vanishing after rationalization has all coefficients zero in \(R\). The polynomial \(\Phi_p(Z+1)=\sum_{j=1}^p\binom pj Z^{j-1}\) is Eisenstein at \(p\): its leading coefficient is \(1\), all others are divisible by \(p\), and its constant coefficient is \(p\), not divisible by \(p^2\). Thus the rational quotient is a field. There \(\zeta\) has order \(p\) and its \(p\) powers are distinct roots of \(T^p-1\); comparing monic degree-\(p\) polynomials and using the injection proves \[
T^p-1=\prod_{j=0}^{p-1}(T-\alpha_j),\qquad
p\alpha_i^{p-1}=\prod_{\substack{0\leq j<p\\j\ne i}}
(\alpha_i-\alpha_j).
\tag{LI-5}
\] Both \(p\) and \(\alpha_i\) are units. In a commutative ring, a factor of a unit product is a unit, so every difference in (LI-5) is a unit. The printed endpoints at 40684 and 40686 are both \(\alpha_1\); the correct endpoints are \(\alpha_0\) and \(\alpha_{p-1}\), with the same omitted \(i\)-factor. This is a source correction. Separately, the argument supports every \(d\leq p\), not only the stated \(d<p\). That stronger boundary is editorial and does not replace the source claim.

For the polynomial step of smooth base change, let \(F=\sum_{i=0}^{d-1}b_iX^i\in B[X]\) be integral over \(R[X]\). For \(F=0\) the conclusion is immediate; otherwise choose any such positive \(d\). Choose distinct integer primes \(p,q>d\). For each \(\ell\in\{p,q\}\) separately, put \(R_\ell=R[1/\ell]\), \(B_\ell=B[1/\ell]\), and \[
R'_\ell=R_\ell[Z]/\Phi_\ell(Z),\qquad
B'_\ell=R'_\ell\otimes_{R_\ell}B_\ell.
\] The first is finite free of positive rank \(\ell-1\), with \(1\) in its basis; hence \(B_\ell\hookrightarrow B'_\ell\). The chosen \(d\) powers of \(\zeta\) have unit pairwise differences by the base change of (LI-5). The full interpolation formula is \[
F(X)=\sum_{i=1}^d F(\alpha_i)
\frac{\prod_{j\ne i}(X-\alpha_j)}
{\prod_{j\ne i}(\alpha_i-\alpha_j)}
\quad\hbox{in }B'_\ell[X].
\tag{LI-6}
\] To prove it over this ring, evaluate the difference at all \(\alpha_i\). Its degree is less than \(d\), and the determinant of its coefficient evaluation matrix \((\alpha_i^{j-1})_{i,j}\) is \(\prod_{i<j}(\alpha_j-\alpha_i)\), a unit. The adjugate therefore forces every coefficient to be zero. Every denominator in (LI-6) is retained and is a unit of \(R'_\ell\).

Evaluation of the original monic equation for \(F\) makes each \(F(\alpha_i)\) integral over \(R'_\ell\). Formula (LI-6) puts every image of \(b_i\) in the ring of integral elements over \(R'_\ell\). Since \(R'_\ell/R_\ell\) is finite, transitivity makes it integral over \(R_\ell\). The injection \(B_\ell\to B'_\ell\) then descends that monic equation to \(B_\ell\).

For the exact gluing back to \(R\), if \(b\in B\) is integral in both localizations, choose degrees \(e_p,e_q\) for those equations, let \(N=\max(e_p,e_q)\), and set \(M=\sum_{j=0}^{N-1}Rb^j\subset R[b]\). The quotient \(Q=R[b]/M\) has \(Q_p=Q_q=0\). For any \(v\in Q\) there are exponents \(a,c\) with \(p^av=q^cv=0\). Integer Bézout coefficients for the coprime integers \(p^a,q^c\) give \(v=0\). Thus \(R[b]=M\) is a finite module. Write multiplication by \(b\) on its listed generators as a matrix \(A\). The adjugate identity for \(bI-A\) implies that the monic \(\det(TI-A)\) evaluated at \(b\) annihilates every generator, including \(1\). This is the required monic equation over \(R\). Zero localized rings cause no exception to the argument. The converse follows from closure of integral elements under sums and products. Induction handles several variables; the original standard smooth charts are étale over polynomial rings, so the proved étale comparison finishes the smooth theorem.

The reduction at 40723--40729 has already stated the two localizations. Its reuse of \(R\) for the chosen localization is compressed notation, not a false assertion once that reduction is followed. The preceding argument makes both localizations and coefficient descent visible; it remains editorial instead of an automatic rewrite.

Finally integral closure commutes with filtered diagrams of ring maps \(R_i\to C_i\), even when transition maps are not injective. Let \(A_i\subset C_i\) denote their integral elements. Filtered colimits preserve these injections: if a representative in \(A_i\) becomes zero in the colimit of \(C_i\), it is zero in some later \(C_j\), hence in \(A_j\). Every monic equation at a stage persists in the colimit. Conversely, an element and the finitely many coefficients of its monic equation in the colimit lift to one common stage; the asserted zero relation becomes zero at some later stage. The element there is integral. This proves the reverse inclusion. Tensor products commute with filtered colimits by their defining finite sums and bilinear relations. Apply this to \(R_i=S_i,\ C_i=S_i\otimes_R B\), for smooth \(R\)-algebras \(S_i\), to obtain exactly the canonical isomorphism at 40748--40762. Also \(S=\mathop{\rm colim}_iS_i\) is flat over \(R\): tensor an injection of \(R\)-modules with each flat \(S_i\), then take the exact filtered colimit. Thus the injectivity of the canonical tensor map and the later conductor calculation use a proved flatness assertion here.

Flatness alone cannot replace smoothness. Keep \[
R=\mathbf Q,\ B=\mathbf Q(z),\ A=\mathbf Q,\qquad
S=\mathbf Q[\epsilon]/(\epsilon^2).
\] The only elements of \(\mathbf Q(z)\) algebraic over \(\mathbf Q\) are constants: if a nonconstant \(u=P(z)/Q(z)\) were algebraic, the nonzero polynomial \(P(T)-uQ(T)\) would make \(z\) algebraic over \(\mathbf Q(u)\), hence over \(\mathbf Q\), a contradiction. Thus \(A\) is the claimed integral closure. The ring \(S\) is finite free but is not smooth: the map from \(S\) to \(\mathbf Q[v]/(v^2)\) taking \(\epsilon\) to \(v\) cannot lift through \(\mathbf Q[v]/(v^3)\); every lift \(v+av^2\) has square \(v^2\ne0\). This is a square-zero extension, so formal smoothness fails. In \(S\otimes_RB=B\oplus B\epsilon\), reducing any integral element modulo \(\epsilon\) forces its constant term into \(\mathbf Q\). Conversely \(a+b\epsilon\), \(a\in\mathbf Q,\ b\in B\), satisfies \((T-a)^2=0\). The integral closure is therefore \(\mathbf Q\oplus\mathbf Q(z)\epsilon\), whereas \(S\otimes_RA=\mathbf Q\oplus\mathbf Q\epsilon\). The defect is exactly \((\mathbf Q(z)/\mathbf Q)\epsilon\). For the original \(f=X^2\), it is annihilated by the full derivative \(f'=2\epsilon\), as (LI-4) predicts; the factor \(2\) is retained.

\subsubsection{Finite conductors and an infinite boundary example}\label{algebra-proof-054-finite-conductors-and-an-infinite-boundary-example}

Suppose now that the original \(R\to B\) is injective, write \(A=\overline R^{\,B}\), and assume \(M=A/R\) is finitely generated. The conductor is the actual ideal \(\mathfrak c=\{r\in R:rA\subset R\}=\operatorname{Ann}_R(M)\); it is also an ideal of \(A\), because if \(rA\subset R\), then for \(a\in A\) one has \(ar\in R\) and \((ar)A\subset rA\subset R\). Let \(m_1,\ldots,m_s\) generate \(M\). The kernel of \[
R\longrightarrow M^s,\qquad r\longmapsto(rm_1,\ldots,rm_s)
\] is exactly \(\mathfrak c\). For a flat \(R\)-algebra \(S\), tensoring preserves this kernel; the elements \(1\otimes m_i\) generate \(S\otimes_RM\). Hence \[
\operatorname{Ann}_S(S\otimes_RM)=\mathfrak cS.
\tag{LI-7}
\] When \(S\) is smooth, or a filtered colimit of smooth algebras, the source theorems identify its actual integral closure in \(S\otimes_R B\) with \(A_S=S\otimes_RA\), while flatness identifies \(A_S/S=S\otimes_RM\). Thus (LI-7) is the exact conductor comparison, and the conductor closed subscheme is the base change \(\operatorname{Spec}(S/\mathfrak cS)\). No Noetherian or finite-presentation assumption on \(M\) was used. Finite generation of \(M\) implies that \(A\) is finite over \(R\), by lifting the \(m_i\) and adjoining \(1\).

Finite generation cannot be removed even for étale localization. Take \[
K_n=\mathbf Q(2^{1/2^n})\ (n\geq0),\qquad
K_\infty=\bigcup_{n\geq0}K_n,\qquad
R=\bigoplus_{n\geq0}K_nx^n\subset A=K_\infty[x],\qquad
B=K_\infty(x).
\] Choose the compatible positive real roots, so \(K_n\subset K_{n+1}\). Eisenstein at \(2\) gives \([K_n:\mathbf Q]=2^n\), hence the inclusions are strict. Products stay in \(R\), since \(K_nK_m\subset K_{\max(n,m)}\subset K_{n+m}\). For \(\alpha\in K_n\), both \(\alpha x^n\) and \(x^n\) lie in \(R\), so \(\alpha=(\alpha x^n)/x^n\) is in its fraction field. Therefore \(\operatorname{Frac}(R)=K_\infty(x)\). We have \(A=R[K_\infty]\), and every adjoined constant is algebraic over \(\mathbf Q\subset R\), hence integral. Each particular polynomial uses only finitely many such constants, so \(A/R\) is integral. The polynomial ring over \(K_\infty\) is integrally closed: write an integral fraction \(u/v\) with coprime polynomials. Multiplying its monic equation by \(v^e\) shows \(v\mid u^e\), so \(v\) is a unit. It follows that \(A\) is exactly the integral closure of \(R\) in \(B\).

For nonzero \(r=\sum c_nx^n\in R\), choose an index with \(c_n\ne0\) and choose \(\alpha\in K_\infty\setminus K_n\). Then \(\alpha c_n\notin K_n\), since \(c_n\) is nonzero in that field. Thus \(\alpha r\notin R\), and \(\mathfrak c=0\). In particular \(A\) is not finite over \(R\): if it had finitely many module generators in \(\operatorname{Frac}(R)\), the product of their nonzero denominators would give a nonzero \(c\in R\) with \(cA\subset R\), contradicting the computed conductor. Consequently \(A/R\) is not finitely generated either, by the lifting argument above. On the other hand \[
R[1/x]=K_\infty[x,x^{-1}]=A[1/x]:
\] the fraction \((\alpha x^n)/x^n\) displays every needed constant. Set \(S=R[1/x]\). This is an étale localization of \(R\). It is integrally closed in \(B\), by localization of the already proved integral-closure comparison. Its conductor in its own integral closure is therefore all of \(S\), whereas \(\mathfrak cS=0\). The rings are nonzero, so these ideals differ. This establishes the finiteness boundary in (LI-7), without contradicting the source theorem: integral closure still commutes with this localization.

\subsubsection{Report accounting and propagation}\label{algebra-proof-054-report-accounting-and-propagation}

Nine report groups account for fifteen reports exactly once: the paired article reports 12180/916; paired spelling 12181/917; paired conjunction 12182/918; display period 12183; basis agreement 12187; paired Fourier endpoints 12188/919; implicit localization 920; paired ``OK'' wording 12189/921; and paired colimit agreement 12190/922. The basis list names one basis, so singular agreement is defensible; changing it to ``form a basis'' is optional. ``OK'' is informal but can be intentional authorial voice; no evidence makes it a mathematical defect or a certain unremoved draft marker. The implicit localization is valid as proved above. These three groups remain optional and unintegrated.

The six other groups propose seven minimal operations: the article, spelling, conjunction, misplaced period, both Fourier endpoints, and the plural verb ``commute''. The full interval has no earlier canonical correction. Its mathematical formulas and maps have been checked above.

Three zero-operation editorial groups record the torsion and completion comparisons, universally cartesian gluing and finite free syntomic domination; the derivative-annihilated defect, Fourier boundary and exact flat counterexample; and the conductor comparison with its infinite boundary example. Group 1147 supplies the prescribed-point root cover, and 1042 supplies the ordered-root algebra used in (LI-3). The smooth integral-closure theorem receives the derivative calculation after localization; the filtered-colimit theorem receives the exact finite-equation argument; the conductor result receives that canonical integral-closure comparison and retains its finite-generation hypothesis. Both corrections and strengthenings enter the ledger. Broader combined synthesis and its paper remain after the core task.

\hypertarget{algebra_conormal_thickening_notes_20260928}{}
\subsection{Formal unramifiedness, universal thickenings and principal parts}\label{algebra-proof-055-formal-unramifiedness-universal-thickenings-and-principal-parts}

\subsubsection{Source identity and scope}\label{algebra-proof-055-source-identity-and-scope}

Source: Stacks Project authors, algebra.tex at commit a04446e57ec1fbc252a871afcec7752fb2807b14, SHA256 FA8BB92E58A4F78A2BD01B3B6A4A87DE0A0D279F5DD90641B574DD5FBFFFA4F3. Read and adjudicated: the complete 40769--41376 interval and its twenty-six reports. The next section was read through 41490 but remains unadjudicated. Dependencies read: 34102--34137, 34241--34292, 34657--34751, 34929--34970, 35411--35460, 37862--37902 and 38168--38184. Reading of a displayed dependency does not claim reading of its entire surrounding chapter. The earlier differential-operator note and formal-smoothness note retain their complete proofs and hashed reading coverage. Four indexed hits for ``formally unramified'' were routing records only; none was read or used. No novelty claim or comprehensive literature claim is made.

All supplementary arguments below remain source-linked editorial material. They do not replace the author's translated proofs. The visible source corrections retain all original rings, kernels, coefficients and tensor actions. In particular, the source graph error is corrected as a source error, not silently attributed to a translator.

\subsubsection{Uniqueness, colimits and the nilpotence boundary}\label{algebra-proof-055-uniqueness-colimits-and-the-nilpotence-boundary}

For \(R\to S\), let \(J=\ker(S\otimes_R S\to S)\). With the exact source sign convention, \[
ds\longleftrightarrow 1\otimes s-s\otimes1\pmod {J^2}.
\] Thus the maps \(\sigma_1(s)=s\otimes1\) and \(\sigma_2(s)=1\otimes s\) into \((S\otimes_R S)/J^2\) satisfy \(\sigma_2-\sigma_1=d\). Uniqueness of square-zero lifting forces their equality, hence every \(ds=0\), hence \(\Omega_{S/R}=0\). Conversely, if two lifts \(\tau_1,\tau_2:S\to A\) agree modulo an ideal \(I\) with \(I^2=0\), then \(\delta(s)=\tau_2(s)-\tau_1(s)\in I\) satisfies \[
\delta(st)=\tau_1(s)\delta(t)+\tau_1(t)\delta(s),
\qquad \delta(r)=0\quad(r\in R).
\] The omitted product is \(\delta(s)\delta(t)=0\). Hence \(\delta\) factors through \(\Omega_{S/R}\), proving uniqueness. Equivalently the source map \(s_1\otimes s_2\mapsto\tau_1(s_1)\tau_2(s_2)\) sends \(J\) into \(I\) and \(J^2\) to zero, so its quotient map is well defined. This supplies the source's omitted verification.

Base change preserves uniqueness by the tensor universal property. Localizing either ring gives the original localization isomorphisms of differentials, so the source's local and localization criteria follow by testing the zero module at every prime. No finite-generation assumption on \(\Omega\) is needed for vanishing at all primes: if a nonzero element \(m\) existed, choose a prime containing its proper annihilator; its localization there would remain nonzero.

The two source colimit lemmas hold for every small diagram of \(R\)-algebras, not just a directed system. If every object is formally unramified, two lifts from its colimit restrict to the same lift on every object and hence are equal by the colimit universal property. If every object is formally étale, each restriction has one lift. For an arrow \(u:S_i\to S_j\), the chosen lift on \(S_j\), composed with \(u\), and the chosen lift on \(S_i\) have the same reduction. Uniqueness on \(S_i\) makes them equal. Thus all lifts form a compatible cocone and give one lift from the colimit. The empty diagram has colimit \(R\) and also satisfies the conclusion. No injectivity, finite presentation or filtering assumption is used in this proof.

Square-zero uniqueness, and square-zero existence when available, extend to every nilpotent ideal by the successive maps \(A/I^{j+1}\to A/I^j\). Their kernels \(I^j/I^{j+1}\) have square zero for \(j\geq1\). A finite nilpotence exponent terminates the induction. This argument does not justify arbitrary nil ideals.

The exact example already constructed in group 1085 establishes that boundary even for a formally étale map. Fix a field \(k\) and take the injective system \[
R_n=k[t_n]/(t_n^{2^n}),\quad
t_n\longmapsto t_{n+1}^2\quad(n\geq1),\qquad
R=\mathop{\rm colim}_nR_n,\quad J=(t_n:n\geq1).
\] The monomials \(1,t_n,\ldots,t_n^{2^n-1}\) map to distinct nonzero even monomials, proving injectivity. Thus \(J\ne0\), while each element of \(J\) belongs to a nilpotent ideal at a finite stage and is nilpotent. The relations give \(J=J^2\), and \(S=R/J=k\). For any \(R\to A\) and \(S\to A/I\) with \(I^2=0\), the image of \(J\) lies in \(I\), and the image of \(J=J^2\) is therefore zero. The given \(R\to A\) factors uniquely through \(S\). Hence \(R\to S\) is formally étale. But the identity \(S\to R/J\) has no \(R\)-algebra lift \(S\to R\), since the composite \(R\to S\to R\) would have to be the identity while killing the nonzero ideal \(J\). Also \(J\) is not finitely generated: it is contained in the Jacobson radical, and \(J=J^2\) would make Nakayama force a finite \(J\) to zero. The quotient map is finite type but not finitely presented. This reinforces why the source's formally-étale/étale equivalence at 41200 keeps finite presentation.

\subsubsection{The universal thickening and its exact obstruction}\label{algebra-proof-055-the-universal-thickening-and-its-exact-obstruction}

Let \(R\to S\) be formally unramified. Retain the original presentation \(P=R[x_i]\twoheadrightarrow S\), \(x_i\mapsto z_i\), with kernel \(J\). The source uses \(I\) both for the variable index set and later for a square-zero ideal. The indices here are exactly its original indices; \(I\) in the lifting argument below denotes only that target ideal. There is no change to the set of variables. Set \[
E=J/J^2,\quad F=\Omega_{P/R}\otimes_PS=\bigoplus_iS\,dx_i,\quad
d:E\twoheadrightarrow F.
\] The last map is onto because \(\Omega_{S/R}=0\). Choose the source splitting \(\sigma:F\to E\). Let \(K=\ker d\), \(L=\sigma(F)\), and let \(J'\) be the inverse image of \(L\) in \(J\). Thus \[
E=K\oplus L,\quad d|_L:L\xrightarrow{\sim}F,\quad
U=P/J',\quad C=\ker(U\to S)=J/J'\cong K,\quad C^2=0.
\tag{CT-1}
\] The isomorphism \(K\to C\) is induced by the original quotient map: it is injective because \(K\cap L=0\), and surjective because \(E=K+L\).

Given \(R\to A,\ a:S\to A/I\), \(I^2=0\), choose lifts of all the \(z_i\) to define \(\beta:P\to A\). The ideal \(I\) is an \(S\)-module via \(a\): two lifts of a scalar differ by an element of \(I\), which annihilates \(I\). We have \(\beta(J)\subset I\), \(\beta(J^2)=0\), and the induced map \(\bar\beta:E\to I\) is \(S\)-linear. For \(\delta_i\in I\), replacing \(\beta(x_i)\) by \(\beta(x_i)+\delta_i\) gives the exact finite Taylor identity \[
\beta'(g)=\beta(g)+\sum_i\delta_i\,
\beta\!\left(\frac{\partial g}{\partial x_i}\right).
\tag{CT-2}
\] Each polynomial has only finitely many terms and variables. Every term with two \(\delta\)'s vanishes; no convergence or bounded number of variables is assumed. Define \[
\delta_i=-\bar\beta\bigl(\sigma(dx_i)\bigr).
\] Then on \(L\), the changed map is \(\bar\beta-\bar\beta\sigma d=0\). Hence \(\beta'(J')=0\), giving \(q:U\to A\). Any other choice which kills \(J'\) must have these same \(\delta_i\), since \(d|_L\) is an isomorphism onto the free basis. This proves the claimed existence and uniqueness, including the full missing summation in the source formula for \(dg\) at 40990.

The restriction \[
\lambda=q|_C:C\longrightarrow I
\tag{CT-3}
\] is the exact obstruction to a lift \(S\to A\): such a lift exists if and only if \(\lambda=0\), since this is precisely the condition that \(q\) factor through \(U/C=S\). When it exists it is unique. Under \(K\cong C\), (CT-3) is \(\bar\beta|_K\), independent of the chosen polynomial lift: a change adds an \(S\)-linear map composed with \(d\), which vanishes on \(K\). Different splittings give uniquely isomorphic \(U\)'s by applying their universal properties to each other; the composites must be the identity by the same uniqueness. Thus the resulting conormal module and obstruction do not depend on that choice.

For the quotient \(S=R/H\), take \(P=R,J=H,F=0\). Then \(J'=H^2\), \(U=R/H^2\), and \(C=H/H^2\), proving the omitted quotient lemma. Both localization assertions also follow directly: a lift of a unit modulo a square-zero ideal is a unit. Explicitly, if \(uv=1-e\), \(e^2=0\), its inverse is \(v(1+e)\). Apply this to the images of the specified multiplicative subset in \(R\), or to its full inverse image in \(U\) when the subset lies in \(S\). The unique universal map therefore extends uniquely to the indicated localization. The two comparison maps in the source are inverse because their composites are the unique maps of the same universal lifting diagram. Localization of the exact kernel sequence gives precisely the stated localized conormal modules.

There is a stronger base-change statement with an essential kernel qualification. For every \(R\to R_1\), set \(S_1=R_1\otimes_RS\), \(U_1=R_1\otimes_RU\). The map \(U_1\to S_1\) is surjective, with square-zero kernel the image of \(R_1\otimes_RC\). For a lifting diagram over \(R_1\), the universal map \(U\to A\) combines with the given \(R_1\to A\) to give the unique \(R_1\)-algebra map \(U_1\to A\). Thus \(U_1\) is the actual universal thickening of \(S_1/R_1\), without a flatness assumption. Its conormal module \(C_1\) fits in the exact sequence \[
\operatorname{Tor}_1^R(R_1,U)\longrightarrow
\operatorname{Tor}_1^R(R_1,S)\longrightarrow
R_1\otimes_RC\longrightarrow C_1\longrightarrow0.
\tag{CT-4}
\] This is the tensor long exact sequence of \(0\to C\to U\to S\to0\); the last map is the actual inclusion image in \(U_1\). For flat \(R_1\), or whenever the displayed middle Tor group vanishes, the conormal map is an isomorphism. In general it need not be: for \(R=\mathbf Z,\ S=\mathbf Z/p,\ U=\mathbf Z/p^2,\ R_1=\mathbf Z/p\), we have \(C=p\mathbf Z/p^2\), \(R_1\otimes_RC=\mathbf F_p\), but \(U_1=S_1=\mathbf F_p\) and \(C_1=0\). The connecting map in (CT-4) is then surjective onto \(\mathbf F_p\). Universal thickenings commute with this base change; their kernels require the exact Tor correction.

For completeness this obstruction also classifies every square-zero \(R\)-algebra extension of \(S\) with a fixed identified kernel \(M\), an arbitrary \(S\)-module. Given \(\lambda:C\to M\), form the ring \[
E_\lambda=(U\oplus M)/
\{(c,-\lambda(c)):c\in C\},\qquad
(u,m)(v,n)=(uv,\bar u n+\bar v m).
\tag{CT-5}
\] The denominator is an ideal: multiplication by \((u,m)\) sends its element indexed by \(c\) to that indexed by \(uc=\bar u c\). It intersects \(0\oplus M\) trivially. Thus (CT-5) has identified square-zero kernel \(M\) and quotient \(S\). The map \(U\to E_\lambda\) restricts on \(C\) to \(\lambda\). Conversely an extension \(0\to M\to E\to S\to0\) receives the unique map \(q:U\to E\); let \(\lambda=q|_C\). The map \((u,m)\mapsto q(u)+m\) factors through (CT-5), is the identity on kernel and quotient, and is an isomorphism: surjectivity follows by lifting the quotient element to \(u\), and injectivity follows by subtracting \((c,-\lambda(c))\) from a kernel representative. An automorphism fixing \(S,M,R\) fixes \(q\) by universality and fixes \(M\), hence is the identity. Under the Baer sum, use \(E_\lambda\times_SE_\mu\) and quotient its kernel \(M\oplus M\) by \(\{(m,-m)\}\); the map from \(C\) becomes \(\lambda+\mu\). Hence the classes, with this operation, are exactly \(\operatorname{Hom}_S(C,M)\). The zero class is the split extension.

In particular, for formally unramified \(R\to S\), formal étaleness is equivalent to \(C=0\). One direction follows from \(U=S\) and its lifting property. Conversely formal étaleness gives a section \(s:S\to U\); both \(1_U\) and \(s\pi\) are universal maps from \(U\) in the same diagram, so uniqueness forces \(s\pi=1_U\), and \(C=0\). This criterion is valid without finite presentation; it is not a claim that formal étaleness alone implies étaleness.

\subsubsection{The differential graph and the original base contribution}\label{algebra-proof-055-the-differential-graph-and-the-original-base-contribution}

At 41090--41154 retain the source rings \(R\to A\to B\), its presentation \(P=A[x_i]\twoheadrightarrow B\) with kernel \(J\), and its chosen \(f_i\in J\) such that \(d_{P/A}f_i=dx_i\) after tensoring with \(B\). Set \(J'=(f_i)+J^2\), \(B'=P/J'\), and \(N=J/J'\). The nilpotent ideal in \(B'\) is \(N\), not \(J'/J^2\). Indeed \(N^2=0\) and \(B'/N=B\). The cokernel of the differentials \(df_i\) in \(\Omega_{P/A}\otimes_PB'\) vanishes modulo \(N\), so it equals its multiple by \(N\), hence its multiple by \(N^2=0\). This proves \(\Omega_{B'/A}=0\) even with infinitely many variables. No finite-module version of Nakayama is being used.

For the \(R\)-relative statement, retain both summands \[
\Omega_{P/R}\otimes_PB=E_0\oplus F_0,\qquad
E_0=\Omega_{A/R}\otimes_AB,\qquad F_0=\bigoplus_iB\,dx_i.
\] If \(f_i=\sum_\nu a_{i,\nu}x^\nu\) is its full coefficient expansion, then \[
df_i=(\eta_i,dx_i),\qquad
\eta_i=\sum_\nu(da_{i,\nu})\otimes z^\nu\in E_0.
\tag{CT-6}
\] All multi-indices, coefficients and monomials are retained; the second component is \(dx_i\) by the chosen \(A\)-relative condition. The exact differential sequence, tensored with \(B\), identifies \(\Omega_{B'/R}\otimes_{B'}B\) with the quotient by the image of \(J'/(J')^2\otimes_{B'}B=J'/JJ'\). The image of \(J^2\) is zero, by the Leibniz rule after reduction modulo \(J\). Every element of \(J'\) is a finite sum of multiples of \(f_i\) plus an element of \(J^2\); its differential has image a finite \(B\)-linear combination of the pairs (CT-6).

Consequently this image is the graph of \(\eta:F_0\to E_0,\ dx_i\mapsto\eta_i\). Its projection to \(F_0\) is an isomorphism, but it need not be the second summand itself. The exact quotient map and its inverse are \[
(E_0\oplus F_0)/\{(\eta(v),v):v\in F_0\}
\xrightarrow{\sim}E_0,\quad [(e,v)]\mapsto e-\eta(v),
\qquad e\mapsto[(e,0)].
\tag{CT-7}
\] These are inverse because their difference on \((e,v)\) is the graph element \((\eta(v),v)\). Thus the canonical map \(E_0\to\Omega_{B'/R}\otimes_{B'}B\) is indeed an isomorphism. The theorem survives, but the printed claim at 41149--41150 that the submodule maps onto the summand must be replaced by projection of the image onto the summand.

For a concrete counterexample to the printed intermediate claim, take \(R=k,\ A=k[t],\ P=A[x],\ J=(x-t)\), and \(B=A\). Then \(f=x-t\) has \(d_{P/A}f=dx\), \(J'=J\), and \(B'=A\); but \(d_{P/k}f=dx-dt\). Its image is the graph of \(dx\mapsto-dt\), not \(0\oplus A\,dx\). Formula (CT-7) sends \(e+b\,dx\) to \(e+b\,dt\), preserving the original base differential and its sign.

\subsubsection{Formal étale jets and differential operators}\label{algebra-proof-055-formal-uxe9tale-jets-and-differential-operators}

The source's finite-presentation equivalence is valid: formal étaleness means formal smoothness and zero differentials; finite presentation turns formal smoothness into smoothness, and zero differentials then gives étaleness. In the other direction étale maps are finitely presented, formally smooth and have zero differentials. Localization lifts uniquely because units lift across square-zero ideals, with the inverse formula already given above.

For 41256--41288 keep its exact \(R\to S\), \(J\subset S\), and \(I=\ker(R\to S/J)\), with \(R/I\cong S/J\). Successive unique lifts \(v_n:S\to R/I^n\) exist by formal étaleness; the kernels \(I^n/I^{n+1}\) are square zero. Compatibility is forced by uniqueness. Each \(v_n\) sends \(J\) into \(I/I^n\), because it reduces to the specified isomorphism modulo \(I\), and so kills \(J^n\). It induces \(\bar v_n:S/J^n\to R/I^n\). The composite \(R/I^n\to S/J^n\to R/I^n\) is the identity because both maps are \(R\)-algebra maps. For the opposite composite, precompose with \(S\to S/J^n\). It and the canonical map agree modulo \(J\); uniqueness of lifts across the nilpotent ideal \(J/J^n\) makes them equal. This proves the canonical isomorphisms for every \(n\geq1\). Their commuting transition maps identify each kernel \(I^n/I^{n+1}\) with \(J^n/J^{n+1}\), including multiplication, giving the source graded-ring isomorphism. Their inverse limits give the corresponding completion isomorphism as a separate direct consequence.

For \(R\to S\to S'\) with \(S\to S'\) formally étale, define the actual diagonal ideals \[
J=\ker(S\otimes_RS\to S),\qquad
J'=\ker(S'\otimes_RS'\to S').
\] The printed \(R'\) at 41294 is undefined and is corrected to \(R\). On every diagonal algebra distinguish the left action \(s\otimes1\) and the right action \(1\otimes s\). Set \(T=S'\otimes_RS\) and \(T'=S'\otimes_RS'\). The map \(T\to T'\) is the base change of \(S\to S'\) through the right \(S\)-action and is formally étale. The map \(T\to T'/J'=S'\), \(s'\otimes s\mapsto s's\), is onto and has kernel \(I=JT\). To verify the kernel statement without flatness, the map from \(S\otimes_RS\) onto \(S\), as a left \(S\)-module map, has section \(s\mapsto s\otimes1\). Its split kernel sequence stays exact after tensoring on the left with \(S'\); its kernel image is precisely the ideal generated by \(J\). Also \(I^n=J^nT\), since products of generators give both inclusions. The preceding infinitesimal-isomorphism lemma yields \[
S'\otimes_{S,\mathrm{left}}
\bigl((S\otimes_RS)/J^{k+1}\bigr)
\xrightarrow{\sim}(S'\otimes_RS')/(J')^{k+1}.
\tag{CT-8}
\] It is the original map \((s'\otimes(a\otimes b))\mapsto s'a\otimes b\). It respects the right \(S\)-action and the quotient maps in \(k\). For \(k=1\), the split multiplication sequences identify their kernels with \(S'\otimes_S J/J^2\) and \(J'/(J')^2\). This proves the source differential isomorphism with its original sign convention.

Write \(Q_k=(S\otimes_RS)/J^{k+1}\) and \(Q'_k=(S'\otimes_RS')/(J')^{k+1}\). For any \(S\)-module \(M\), the source principal-parts construction is \(P^k_{S/R}(M)=Q_k\otimes_{S,\mathrm{right}}M\), with its left \(S\)-module structure. Retain the corrected extension \(M'=S'\otimes_SM\), not \(S'\otimes_RM\). For example, with \(R=k,\ S=S'=k[x]\) and \(M=S\), the identity map \(S\to S'\) is formally étale. At order zero the left module in the claimed comparison is \(S\), whereas the incorrectly defined \(M'\) is \(S\otimes_kS\), free of infinite rank as a left \(S\)-module. Thus the erroneous tensor base already fails at order zero. Tensoring (CT-8) with \(M\) on the right gives \[
S'\otimes_S P^k_{S/R}(M)
\cong Q'_k\otimes_{S,\mathrm{right}}M
\cong Q'_k\otimes_{S',\mathrm{right}}(S'\otimes_SM)
=P^k_{S'/R}(M').
\tag{CT-9}
\] The middle isomorphism sends \(q\otimes(s'\otimes m)\) to \(q(1\otimes s')\otimes m\), with inverse \(q\otimes m\mapsto
q\otimes(1\otimes m)\). Thus the left and right actions are never interchanged. The full module quotients use numerator \(S'\otimes_RM'\); the source's abbreviated quotient notation is interpreted accordingly. That convention is not a new mathematical defect. The \(R'\) at 41349, however, is undefined and must again be \(R\).

For every \(S'\)-module \(N'\), (CT-9), the principal-parts universal property, and scalar-extension adjunction give the exact bijection \[
\operatorname{Diff}^{\,k}_{S'/R}(M',N')
\xrightarrow{\sim}\operatorname{Diff}^{\,k}_{S/R}(M,N'),
\qquad D'\longmapsto D'\circ(m\mapsto1\otimes m).
\tag{CT-10}
\] To check that the adjunction is this map, the universal operator sends \(m\) to the class of \(1\otimes m\); (CT-9) sends that class to the corresponding class of \(1\otimes(1\otimes m)\). Hence restriction is exactly the asserted operation, not just an abstract bijection. No flatness or finite presentation of \(M\) is assumed.

Given the source's \(D:M\to N\), compose it with \(N\to S'\otimes_SN\). If \(D=\gamma\circ D_{\rm univ}\), its extension corresponds under (CT-9) to \(1\otimes\gamma\), tensoring all three relevant modules over \(S\). This proves the corrections at 41356--41357. It also proves uniqueness by (CT-10). The correspondence commutes with order truncation because (CT-8) commutes with all diagonal quotient maps. It preserves addition and identity maps. If \(D_1,D_2\) have orders at most \(k_1,k_2\), their composite has order at most \(k_1+k_2\), by the exact commutator identity \[
[D_2D_1,g]=D_2[D_1,g]+[D_2,g]D_1.
\] The extended composite and the composite of extensions restrict to the same map on \(M\). Applying (CT-10) at \(k_1+k_2\) makes them equal. Thus extension defines a functor on the category of \(S\)-modules with finite-order \(R\)-linear differential operators as morphisms, preserving the order filtration. In particular it extends complexes whose differentials have finite order: the square remains zero by this composition identity, not by an assertion of exactness of tensoring. It does not assert that arbitrary cohomology commutes with this possibly nonflat scalar extension.

The resulting endomorphism ring \(\operatorname{Diff}_{S/R}(M,M)\) is a possibly noncommutative \(R\)-algebra, not in general an \(S\)-algebra with central structural scalars. This requires a correction to the final sentence at 41369. For \(R=\mathbf Q,\ S=\mathbf Q[x],\ M=S,\ S'=S\), let \(\partial=d/dx\) and \(\mu_x\) be multiplication by \(x\). Then for every polynomial \(f\), \[
(\partial\mu_x-\mu_x\partial)(f)
 =\partial(xf)-x\partial(f)=f.
\] So the natural \(S\)-scalars are not central and multiplication is not bilinear over those \(S\)-scalars. All operators are \(R\)-linear, hence the image of \(R\) is central; addition and composition make the asserted \(R\)-algebra, and its action on \(M'=S'\otimes_SM\) is exactly the extension just constructed. The module conclusion remains valid with this corrected algebra base.

\subsubsection{Report decisions and propagation}\label{algebra-proof-055-report-decisions-and-propagation}

Twenty report groups cover all twenty-six reports. Five groups remain optional: the conventional ``Follows from'' proof opening (923), reuse of the index symbol in a typed context (925), renaming the individually specified inverse image \(g\) to \(g_i\) (929), the conventional ``With notation and assumptions'' openings (930), and the already grammatical plural ``unique maps'' (935). Report 12203 is rejected as a nondefect of the source's quotient-notation convention, consistently with the earlier quotient-notation adjudication; the actual quotient and its two actions are specified in (CT-8)--(CT-9).

The other fourteen report groups give sixteen operations: terminal punctuation at 40898, 40950, 40987, 41142, 41229, 41282 and 41325; the sum at 40990; both wrong kernel occurrences in the one line 41120; the ``given above,'' connection at 41148; tensor-base corrections at 41294, 41323, 41349, 41356 and twice at 41357. Two additional source-correction groups repair the differential graph argument at 41149--41150 and the algebra base at 41369. Those mathematical problems were found during source review, beyond the received reports.

Three zero-operation consequence groups record arbitrary small colimits with the exact nilpotence boundary; the universal obstruction, extension classification and arbitrary-base-change Tor comparison; and the order-preserving restriction bijections and differential-operator functor. Group 1085 receives the formally étale conclusion for its existing nil-idempotent example. The earlier differential-operator proof supplies the principal-parts universal map, corrected as \(D=\gamma\circ D_{\rm univ}\), and its full commutator relations. The graph correction is propagated into the conormal differential isomorphism, retaining every coefficient differential. The principal-parts tensor corrections are propagated through the operator extension and its composition action. All these source corrections and proved strengthenings are recorded separately in the ledger. The broader combined-results synthesis remains after core work.

\hypertarget{algebra_unramified_notes_20260928}{}
\subsection{Unramified ring maps and the exact diagonal idempotent}\label{algebra-proof-056-unramified-ring-maps-and-the-exact-diagonal-idempotent}

\subsubsection{Source and scope}\label{algebra-proof-056-source-and-scope}

Source: Stacks Project authors, algebra.tex at commit a04446e57ec1fbc252a871afcec7752fb2807b14, SHA256 FA8BB92E58A4F78A2BD01B3B6A4A87DE0A0D279F5DD90641B574DD5FBFFFA4F3. The full 41377--41711 interval and twelve received reports are read here; 41712--41723 is read but unadjudicated. Dependency reading is bounded to 539--567, 2427--2458, 3921--3953, 38608--38652, 32653--32686, 27420--27450 and 39393--39419. These bounded excerpts are not claims that all the associated proofs have been newly read. The already reviewed formal, smooth, étale and local-integral notes retain their prior evidence. Four indexed hits for ``unramified diagonal'' remain unread and unused. The source's Henselian citation and EGA terminology are retained; no new reading of those publications or novelty claim is made.

The source calls finite type plus \(\Omega_{S/R}=0\) unramified, and finite presentation plus the same condition G-unramified. Keep these definitions exactly. Neither a translation nor this review silently changes them to a later preferred convention. The fuller arguments and diagonal consequences below are separate editorial material.

\subsubsection{Source statements and the two reference repairs}\label{algebra-proof-056-source-statements-and-the-two-reference-repairs}

Formal unramifiedness is equivalent to \(\Omega_{S/R}=0\), as proved in the preceding section, so 41404--41421 follows with exactly the two different finiteness hypotheses. For the thirteen assertions at 41423--41547, base change uses the differential base-change isomorphism and the stated base-change finiteness lemma. Composition instead uses \[
T\otimes_S\Omega_{S/R}\longrightarrow\Omega_{T/R}
\longrightarrow\Omega_{T/S}\longrightarrow0
\] and lemma-compose-finite-type, which proves both composition finiteness statements at 539--563. The source reference at 41477 to the base-change finiteness lemma is replaced by this actual composition lemma; the maps and claim remain the same.

For a principal localization, every relative derivation kills \(R\) and kills \(f^{-1}\), since \(d(f^{-1})=-f^{-2}df=0\); the presentation \(R[X]/(fX-1)\) is finite. For \(R/H\) every relative derivation is zero, because it kills the surjective image of \(R\). The quotient is finite type and is finitely presented when \(H\) is finitely generated. Étale maps supply finite presentation and zero differentials as already proved.

For finite-type \(R\to S\), write \(S\) using finitely many generators \(x_1,\ldots,x_n\). Their differentials generate \(\Omega_{S/R}\). If its localization at \(\mathfrak q\) vanishes, choose one denominator outside \(\mathfrak q\) annihilating each \(dx_i\), and multiply those finitely many denominators to get \(g\notin\mathfrak q\) annihilating the module after localization. Then \(\Omega_{S_g/R}=0\). Finite type, and finite presentation when assumed, survive this principal localization. The residue-field criterion is equivalent by Nakayama for this finite localized module. Differential base change identifies the fibre criterion with the same residue vector space, so all four local tests in items (6)--(9) agree. Their finite-type hypothesis is used to pass from residue vanishing to an actual open neighbourhood.

For a finite principal cover \(D(g_j)\), vanishing of the localized differentials implies global vanishing; local finite type and local finite presentation glue by lemma-cover. This proves item (10). In item (11), neither global finiteness hypothesis has yet been given. Its hypothesis instead gives, at every prime, a principal open where the corresponding finite-type or finite-presentation condition and vanishing differentials hold. Quasi-compactness extracts finitely many such opens. They cover the spectrum, so the \(g_j\) generate the unit ideal, and item (10) applies. The reference to (6) and (7) at 41519 must therefore be replaced by the definition and (10). This is the actual argument, not an added global finiteness assumption. For \(S=0\), the map is already finitely presented as \(R/(1)\) and has zero differentials, so the empty-spectrum case is also covered.

For item (12), retain the exact finite presentation \(S=R[x_1,\ldots,x_n]/(g_1,\ldots,g_m)\). The vanishing of the differential cokernel gives all the actual identities \[
dx_i=\sum_j h_{ij}\,dg_j+\sum_{j,k}a_{ijk}g_j\,dx_k
\] in the original free differential module over \(R[x_1,\ldots,x_n]\). Adjoin to \(\mathbf Z\) every coefficient of \(g_j,h_{ij},a_{ijk}\) to get the subring \(R_0\subset R\). These finitely many identities hold over \(R_0\) itself, because its inclusion in \(R\) is injective. Thus the displayed \(S_0\) is finitely presented with \(\Omega_{S_0/R_0}=0\), and its original base change is \(S\). For item (13), use the source's full identities \[
dx_i=\sum_j h_{ij}\,dg_{ij}+\sum_k g'_{ik}\,dx_k,\qquad
g_{ij},g'_{ik}\in I.
\] Only finitely many of these polynomials occur. Keeping every coefficient gives the source \(S_0=R_0[x_1,\ldots,x_n]/(g_{ij},g'_{ik})\). It has zero relative differentials and its base change surjects onto the original \(S=R[x_1,\ldots,x_n]/I\). In fact the constructed \(S_0/R_0\) is already G-unramified, since its displayed relation list is finite. This is an explicit consequence of the existing proof, not an error in the source's weaker conclusion.

For 41567--41641, replace \(S\) only by the stated neighbourhood and then its actual fibre \(F=S\otimes_R\kappa(\mathfrak p)\). It is finite type over the field and has zero differentials. The field smoothness criterion at 38609 applies because zero is at most its nonnegative local component dimension; hence it is smooth with zero differentials, and is étale. The already proved field characterization identifies \(F\) with a finite product of finite separable fields. Its local ring at the original fibre prime is \[
F_{\bar{\mathfrak q}}=S_{\mathfrak q}/\mathfrak pS_{\mathfrak q}
=\kappa(\mathfrak q).
\] This gives both equality of maximal ideals and the exact finite separable residue extension. The fibre point is isolated, giving quasi-finiteness. Conversely, the two residue conditions give that same fibre local ring as a finite separable field; its differential module is zero, and the finite-type fibre test proves unramifiedness. No claim of flatness follows from these residue conditions alone.

The equivalence étale \(\Leftrightarrow\) flat plus G-unramified retains finite presentation. The previous étale criterion applies to the flat local maps with the just proved zero-dimensional separable fibres. The reverse implication uses the already proved flatness of étale maps.

Finally let the source map \(\varphi:k[x_1,\ldots,x_n]\to A,\ x_i\mapsto a_i\), be finite type. This is finite presentation because its polynomial base is Noetherian. If \(da_1,\ldots,da_n\) generate \(\Omega_{A/k}\), the exact differential sequence gives \(\Omega_{A/k[x_1,\ldots,x_n]}=0\), so \(\varphi\) is unramified. At a maximal ideal \(\mathfrak m\) of \(A\), its residue field is finite algebraic over \(k\). The image of the polynomial base in it is generated by finitely many algebraic elements; it is a finite dimensional domain over \(k\), hence a field. Thus \(\varphi^{-1}(\mathfrak m)\) is maximal, and its regular local polynomial ring has dimension \(n\). Under condition (a), the target local dimension is also \(n\); condition (b) bounds the differential residue dimension by \(n\). The field smoothness criterion makes \(A_{\mathfrak m}\) regular, hence Cohen--Macaulay. Its fibre over the polynomial local base is the finite separable residue field just proved, of dimension zero. All hypotheses of miracle flatness at 32654 are now met, including the exact dimension equality \(n=n+0\). This proves flatness at every maximal ideal and hence globally. Conversely, étaleness gives condition (b), flatness and the same zero-dimensional fibres. The original local dimension formula gives condition (a). For \(A=0\), the criteria are vacuous and the empty étale algebra still satisfies the statement; no maximal-ideal argument is needed.

\subsubsection{The diagonal, its determinant and equality of maps}\label{algebra-proof-056-the-diagonal-its-determinant-and-equality-of-maps}

For any map \(R\to S\), put \[
D=S\otimes_RS,\qquad \mu:D\to S,\quad a\otimes b\mapsto ab,\qquad
J=\ker\mu.
\] The exact diagonal-differential comparison gives \(J/J^2=\Omega_{S/R}\). When \(R\to S\) is finite type, \(J\) is generated by the actual elements \(u_i=x_i\otimes1-1\otimes x_i\), where the \(x_i\) are the chosen algebra generators. Indeed the quotient by these elements identifies the two structural maps from \(S\); the map \(S\to D/(u_i)\), \(s\mapsto s\otimes1\), is then inverse to \(\mu\). Thus an unramified map has a finitely generated ideal \(J=J^2\).

More generally only finite generation of \(J\), not finite type of \(S\), is needed in the following exact construction. Choose actual generators \(u_1,\ldots,u_m\) of \(J=J^2\), and write \[
u_i=\sum_{j=1}^m a_{ij}u_j,\quad a_{ij}\in J,\qquad
e=\det(I_m-(a_{ij}))
=\sum_{\pi\in\mathfrak S_m}\operatorname{sgn}(\pi)
\prod_{i=1}^m(\delta_{i,\pi(i)}-a_{i,\pi(i)}).
\tag{UR-1}
\] The adjugate applied to \((I_m-(a_{ij}))(u_i)_i=0\) gives \(eJ=0\). Every entry of the coefficient matrix lies in \(J\), so \(e-1\in J\), and \(e(e-1)=0\). Hence \[
e^2=e,\quad \mu(e)=1,\quad J=(1-e)D.
\tag{UR-2}
\] For the last equality, \(1-e\in J\) and each \(j\in J\) is \((1-e)j\). For \(m=0\), the empty determinant is \(1\), \(J=0\), and the same formulas hold. These formulas also hold for the zero ring.

The maps \[
S\longrightarrow eD,\quad s\longmapsto(s\otimes1)e,\qquad
eD\longrightarrow S,\quad z\longmapsto\mu(z)
\tag{UR-3}
\] are inverse: for \(d\in D\), \(d-\mu(d)\otimes1\in J\), annihilated by \(e\). The first map is \(D\)-linear when \(S\) is a \(D\)-module through \(\mu\), since that same difference is killed by \(e\). Consequently \(S\) is a direct summand of the rank-one free \(D\)-module \(D=eD\oplus(1-e)D\), hence is a finite projective \(D\)-module, and \(D_e\to S\) is the exact original multiplication quotient. Its inverse sends \(s\) to \((s\otimes1)/1\). In \(D_e\) the idempotent \(e\) is \(1\), so this agrees with (UR-3). This is a statement about the diagonal algebra \(D\), not projectivity or finiteness of \(S\) as an \(R\)-module.

The following three conditions on an arbitrary \(R\to S\) are equivalent: formal unramifiedness with \(J\) finitely generated; existence of \(e\) as in (UR-2); and projectivity of \(S\) as a \(D\)-module. The first implication is (UR-1), and the second gives projectivity by (UR-3). For the converse, a \(D\)-linear section \(\sigma\) of \(\mu\) exists. Put \(e=\sigma(1)\). Then \(\mu(e)=1\) and \(Je=0\). The same argument gives \(e^2=e\) and \(J=(1-e)D\), hence \(J\) is finitely generated and \(J/J^2=0\). Thus \(\Omega_{S/R}=0\). The idempotent is unique: if \(e'\) has the same properties, then \(e-e'\in J\), so \(e(e-e')=e'(e-e')=0\), giving \(e=ee'=e'\). The determinant construction is therefore independent of the selected generators and expressions, even though it preserves them as its explicit formula.

Every base change \(R\to R_1\) carries this idempotent to the diagonal idempotent for \(S_1=R_1\otimes_RS\), under \(R_1\otimes_RD\cong S_1\otimes_{R_1}S_1\). The multiplication map has the \(R\)-linear section \(s\mapsto s\otimes1\), so its kernel sequence remains exact after this base change. Alternatively the decomposition in (UR-2) itself tensors to the asserted kernel decomposition. Neither argument assumes \(R_1\) flat.

For two \(R\)-algebra maps \(f,g:S\to B\), let \(h:D\to B,\ h(a\otimes b)=f(a)g(b)\), and put \(d=h(e)\). The ideal generated by all differences \(f(s)-g(s)\) is exactly \[
\mathfrak a=h(J)B=(1-d)B,\qquad d^2=d.
\tag{UR-4}
\] Here \(J\) is generated by the differences in the two tensor factors, and (UR-2) gives the equality. Thus \(B/\mathfrak a\cong B_d\), and on this factor the two original maps are equal. It is the largest factor on which they agree: for every \(B\)-algebra \(B'\), equality of their composites is equivalent to the image of \(1-d\) being zero, equivalently the image of \(d\) being \(1\), equivalently factorization through \(B_d\). The equality subscheme is exactly the open and closed subset \(D(d)\) with this ring structure, including all nilpotents. Formula (UR-4) commutes with arbitrary further base change of \(B\). If \(\operatorname{Spec}(B)\) is nonempty and connected, then \(d\in\{0,1\}\); therefore equality at even one residue-field point forces \(d=1\) and \(f=g\) on the whole ring. Connectedness here does not require reducedness.

The finite diagonal condition has a real boundary. Take the already used tower \(K_n=\mathbf Q(2^{1/2^n})\) with compatible positive roots and \(L=\bigcup_nK_n\). The Eisenstein degrees \(2^n\) make this an infinite algebraic extension, and characteristic zero makes it separable. Every \(\alpha\in L\) satisfies a separable minimal polynomial \(F\); differentiating gives \(F'(\alpha)d\alpha=0\). The nonzero field element \(F'(\alpha)\) is a unit, hence \(\Omega_{L/\mathbf Q}=0\). But its diagonal admits no idempotent as above. If \(e=\sum a_i\otimes b_i\) existed, choose \(n\) with all \(b_i\in K_n\) and choose \(\alpha\in L\setminus K_n\). The equality \((\alpha\otimes1)e=(1\otimes\alpha)e\) would put the first side in \(L\otimes_{\mathbf Q}K_n\) and the second in \(L\otimes_{\mathbf Q}\alpha K_n\). These subspaces have zero intersection, because \(K_n\cap\alpha K_n=0\), and tensoring vector spaces over \(\mathbf Q\) preserves that direct sum. Both sides would be zero. Multiplication by \(1\otimes\alpha\) is invertible, forcing \(e=0\), contrary to \(\mu(e)=1\). Hence this formally unramified map has a diagonal ideal which is not finitely generated and no projective diagonal module. It is not a counterexample to the source's finite-type theorem.

\subsubsection{Report decisions and propagation}\label{algebra-proof-056-report-decisions-and-propagation}

Seven groups cover the twelve reports exactly once and give seven operations: paired article reports 12212/942; paired finite-presentation wording 12213/943; composition reference 944; local-to-global reference 945; paired maximal-ideal sentence 946/12214; paired plural verb 12215/947; and paired modal spelling 948/12216. The two reference repairs are mathematical corrections of the stated proof dependencies. The other changes preserve the exact hypotheses and conclusions as copyedits. No prior canonical correction lies in this source interval.

One zero-operation editorial group proves the exact determinant idempotent, the weaker finite-diagonal hypothesis and its converse via projectivity, arbitrary base change, the full equality factor for two maps, and the infinite separable-extension boundary. It receives the formal-unramified criterion from the preceding review and the exact field tower from group 1181, retaining every diagonal map and determinant coefficient. The stronger G-unramified choice already supplied by the finite-coefficient construction is recorded as existing proof content. All fuller arguments and consequences stay editorial; they do not replace the original exposition. Broader synthesis and its paper follow core completion.

\hypertarget{algebra_local_unramified_notes_20260929}{}
\subsection{Local unramified charts, retained scalars and their conormal objects}\label{algebra-proof-057-local-unramified-charts-retained-scalars-and-their-conormal-objects}

\subsubsection{Source identity and reading scope}\label{algebra-proof-057-source-identity-and-reading-scope}

Source: the Stacks Project authors, algebra.tex at a04446e57ec1fbc252a871afcec7752fb2807b14, SHA-256 FA8BB92E58A4F78A2BD01B3B6A4A87DE0A0D279F5DD90641B574DD5FBFFFA4F3. The complete source interval 41712--42061 and its twelve received reports have been read. The next opening, 42062--42084, is read but unadjudicated. Bounded dependency reading covers 566--624, 9797--9830, 30906--30949, 31623--31660 and 40419--40463. These readings do not claim that every dependency proof was newly read in full. Earlier checked proofs are retained in ALGEBRA\_ETALE\_LOCAL\_NOTES\_20260928.md, ALGEBRA\_ETALE\_QUASIFINITE\_NOTES\_20260928.md and ALGEBRA\_CONORMAL\_THICKENING\_NOTES\_20260928.md, together with the diagonal calculation in ALGEBRA\_UNRAMIFIED\_NOTES\_20260928.md.

Four indexed hits for ``unramified etale thickening'' are routing records only; none has been read or used here. No novelty or exhaustive literature-reading claim is made. The source's unramified convention is finite type with zero differentials, and its G-unramified convention adds finite presentation. Keep both conventions and the original eleven-step proof identifiable. All expanded proofs and consequences below are source-linked editorial material.

\subsubsection{The source reductions and all scalar factors}\label{algebra-proof-057-the-source-reductions-and-all-scalar-factors}

We verify proposition-unramified-locally-standard at 41720--41938 with its original rings and localizations. Step 1 replaces \(S\) by \(S_a\), with \(a\notin\mathfrak q\), on which the map is unramified. Step 2 uses the finite-coefficient model established in the previous section: there is an unramified \(R_0\to S_0\), with \(R_0\) of finite type over \(\mathbf Z\), and a surjection \(R\otimes_{R_0}S_0\to S_a\). Contract the specified prime to \(S_0\). A standard étale algebra surjecting onto a principal neighbourhood of this contraction remains standard étale after tensoring with \(R\). The base-changed map is still surjective, and its composite surjects onto the corresponding principal localization of \(S_a\). Thus the Noetherian reduction in the source is legitimate even though \(S_a\) need only be a quotient of the base-changed model.

Here and at the later reductions, principal localizations return to the original ring exactly. If the second denominator in \(S_a\) is \(b/a^N\), the two localization rings have inverse maps \[
(S_a)_{b/a^N}\ \longleftrightarrow\ S_{ab}.
\] Both maps preserve every original element of \(S\). In the left ring, \(b^{-1}=a^{-N}(b/a^N)^{-1}\), and hence \(ab\) is invertible. In the right ring, \(a^{-1}=b(ab)^{-1}\) and \(b^{-1}=a(ab)^{-1}\), and \((b/a^N)^{-1}=a^Nb^{-1}\). These formulas prove the two maps are inverse. The specified prime avoids \(a\) and \(b\), so it avoids \(ab\).

Step 3 uses the already proved quasi-finite finite-neighbourhood lemma at 30906--30949. If its finite algebra is \(T\), its map is \(\varphi:T\to S\), and its denominator is \(b\in T\), the condition is \(\varphi(b)\notin\mathfrak q\), equivalently \(b\notin\varphi^{-1}(\mathfrak q)\). Membership through this stated map is conventional notation, not a false mathematical assertion. The repeated ``such that'' is a prose defect; the proposed correction also spells out the image. The isomorphism \(T_b\cong S_{\varphi(b)}\) permits transfer of any later chart by localizing its source at a lift of \(b\). A principal localization of a standard étale algebra \(R[x]_g/(f)\), at an element \(h/g^N\), is \(R[x]_{gh}/(f)\), by the same inverse formulas. Thus the transfer keeps the required standard étale form and an actual principal open of the original target.

For Steps 4--6, write \(K=\kappa(\mathfrak p)\) and retain the actual finite fibre decomposition \[
S\otimes_R K=A_1\times\cdots\times A_n,\qquad
A_1=S_{\mathfrak q}/\mathfrak pS_{\mathfrak q}
    =\kappa(\mathfrak q).
\] The last equality and finite separability follow from unramifiedness at \(\mathfrak q\); the other factors are local Artinian rings, not necessarily fields. Choose a nonzero primitive element \(\alpha\) in the first factor. Since the fibre is the localization of \(S/\mathfrak pS\) at \(R\setminus\mathfrak p\), some scalar \(\bar c\in K^\times\), \(c\in R\setminus\mathfrak p\), makes \((\bar c\alpha,0,\ldots,0)\) lift to \(t\in S\). The source then uses this nonzero scalar multiple as its primitive element; \(K(\bar c\alpha)=K(\alpha)\), with inverse \(\alpha=\bar c^{-1}(\bar c\alpha)\).

Put \(T=R[t]\subset S\), the ring called \(S'\) in Steps 5--7. Every other fibre prime sends \(t\) to zero, whereas the distinguished prime does not. Consequently \(\mathfrak q\) is the only prime of \(S\) over its contraction \(\mathfrak q'\subset T\). The map \(T\to S\) is finite: a finite set of \(R\)-module generators of \(S\) also generates it as a \(T\)-module. Thus it is integral, and the unique-prime localization result gives \(S_{\mathfrak q'}=S_{\mathfrak q}\). The map \(T_{\mathfrak q'}\hookrightarrow S_{\mathfrak q}\) is finite and injective. The actual ideals satisfy \[
\mathfrak pS_{\mathfrak q}
\subset\mathfrak q'S_{\mathfrak q}
\subset\mathfrak qS_{\mathfrak q}
=\mathfrak pS_{\mathfrak q}.
\] The first residue field contains the chosen primitive element, so the induced inclusion \(\kappa(\mathfrak q')\subset\kappa(\mathfrak q)\) is equality. The finite \(T_{\mathfrak q'}\)-module \(S_{\mathfrak q}/T_{\mathfrak q'}\) therefore has zero quotient by \(\mathfrak q'\), and Nakayama makes it zero. This proves the exact local isomorphism used in Step 7.

In the source's Noetherian reduction, \(T\) is a finite \(R\)-module as a submodule of the finite module \(S\). Both algebras are finitely presented, so 31623--31660 gives principal neighbourhoods with the same local isomorphism. The earlier corrections of that dependency, including inversion of the idempotent selecting the correct factor, remain recorded in the prior review; no duplicate fix is made here. The missing copula in condition (b) at 41853 is a copyedit.

Now use the original Step 8 presentation \(S=R[x]/I\). The fibre ideal \(\bar I=I K[x]\) is nonzero because \(I\) contains a monic integral relation for \(x\), and it is proper because the specified fibre point exists. Let \(H\in K[x]\) be its monic generator. Its membership is in \(K[x]\), not in the coefficient field \(K\). Write an actual finite expression \[
H=\sum_{j=1}^r P_j\,\overline{a_j},
\qquad a_j\in I,\quad P_j\in K[x].
\] For every coefficient of each \(P_j\), choose a fraction with numerator in \(R/\mathfrak p\) and denominator the image of an element of \(R\setminus\mathfrak p\). Let \(c\) be the product of all these finitely many denominators (the empty product is 1), and put \(\lambda=\bar c\in K^\times\). Every coefficient of \(\lambda P_j\) is now in the image of \(R\). Choose lifts \(Q_j\in R[x]\), and define \[
h=\sum_{j=1}^r Q_j a_j\in I.
\] Then \(\bar h=\lambda H\). This establishes the required stronger choice \(h\in I\), not merely \(h\in R[x]\), while retaining the actual nonzero scalar. It does not assert that \(h\) can be monic.

Retain the source's distinct monic irreducibles, their original multiplicities, and its distinguished separable factor: \[
H=\overline h_1\overline h_2^{e_2}\cdots\overline h_n^{e_n},
\qquad e_1=1,\quad e_i\ge1\ (i\ge2).
\] Thus the exact replacement at 41872 is \[
\bar h=\lambda\overline h_1
                 \overline h_2^{e_2}\cdots\overline h_n^{e_n}.
\] The factor algebras are still \(A_i=K[x]/(\overline h_i^{e_i})\), including \(i=1\); the unit \(\lambda\) does not change these ideals but cannot be deleted from polynomial equalities.

Keep the chosen monic \(m\in I\), its full factorization \[
\bar m=\bar k\,\overline h_1^{d_1}\cdots\overline h_n^{d_n},
\quad d_1\ge1,\quad d_i\ge e_i\ (i\ge2),
\] where \(\bar k\) is prime to every \(\overline h_i\). Choose the source's integer \(l\ge2\) with \(l\deg(m)>\deg(h)\). Then \(f=m^l+h\) is monic and belongs to \(I\). Its full reduction is \[
\begin{aligned}
\bar f
&=\lambda\overline h_1\prod_{i=2}^n\overline h_i^{e_i}
  +\bar k^{\,l}\prod_{i=1}^n\overline h_i^{ld_i}\\
&=\overline h_1\left(
 \lambda\prod_{i=2}^n\overline h_i^{e_i}
 +\bar k^{\,l}\overline h_1^{ld_1-1}
                 \prod_{i=2}^n\overline h_i^{ld_i}\right)
=\overline h_1\bar w .
\end{aligned}
\] The terms involving \(\lambda\) at 41898 and 41903 are therefore also required. Modulo \(\overline h_1\), the second term in parentheses is zero since \(ld_1-1\ge1\). The first term is nonzero, since \(\lambda\ne0\) and the other irreducibles are coprime to \(\overline h_1\). Consequently \(\bar w\) is prime to \(\overline h_1\). The original derivative remains \[
g=f'=l m^{l-1}m'+h',\qquad
\bar g=\overline h_1'\bar w+\overline h_1\bar w'.
\] No coefficient is removed in positive characteristic. At the specified residue point the second summand is zero and the first is nonzero by separability. Therefore \(g\notin\mathfrak q\), and the exact map \[
R[x]_g/(f)\longrightarrow (R[x]/I)_g
\] is surjective. Its source is standard étale, since \(f\) is monic and the image of \(f'=g\) is invertible. This proves Steps 9--11 with the original \(m,h,l,f,g\) and all their coefficients retained.

For a concrete adverse example, take \(R=\mathbf Z\), \(\mathfrak p=(3)\), \(I=(x-1)\), \(h=2(x-1)\), \(m=x-1\), \(l=2\). Here \(\lambda=2\) in \(\mathbf F_3\), and the actual \(f=(x-1)^2+2(x-1)=x^2-1\). Dropping the scalar instead gives \((x-1)^2+(x-1)=x^2-x\), a different polynomial in \(\mathbf F_3[x]\). This propagates the same defect already proved in group 1145 to its distinct unramified occurrence; the earlier operation is not duplicated.

\subsubsection{The direct comparison and the separated fibre factors}\label{algebra-proof-057-the-direct-comparison-and-the-separated-fibre-factors}

The finite-cokernel argument of group 1148 also applies with the unramified residue-field conclusions just proved. In the finite neighbourhood of Step 3, \(t\) is integral over \(R\), so an actual monic relation of degree \(d\) makes \(1,t,\ldots,t^{d-1}\) span \(T=R[t]\). Thus \(T/R\) is finite over an arbitrary base. After the local isomorphism \(T_{\mathfrak q'}=S_{\mathfrak q}\), the finite \(T\)-module \(C=S/T\) vanishes at \(\mathfrak q'\). For generators \(c_1,\ldots,c_b\), choose \(u_i\notin\mathfrak q'\) with \(u_i c_i=0\) and put \(u=\prod_i u_i\). Then \(C_u=0\), and the original injection gives \(T_u=S_u\). This replaces the Noetherian and finite-presentation uses in Step 7. Steps 8--11 and their denominator formulas above apply unchanged. This is an alternative proof of the source's already arbitrary-base theorem, not a new generalization of its conclusion.

For 41942--42051, retain the earlier étale quasi-finite decomposition of 40420--40496, with its proved residue-field and finite-component construction in group 1168. At a chosen prime the finite factor \(A_i\) has a unique prime \(\mathfrak r_i\) over \(\mathfrak p\) and a purely inseparable residue extension. If the map is unramified there, the extension is also separable, hence trivial: a minimal polynomial that is both separable and a factor of a polynomial \(X^{p^a}-c\) has degree one; in characteristic zero purely inseparable already means trivial. The local fibre is the residue field and \(\mathfrak p(A_i)_{\mathfrak r_i}\) is its maximal ideal. Since \(A_i/R\) is finite, unique-prime localization gives \((A_i)_{\mathfrak p}=(A_i)_{\mathfrak r_i}\). The finite cokernel of \(R_{\mathfrak p}\to(A_i)_{\mathfrak p}\) has zero residue quotient and is zero by Nakayama.

For each finite factor to which this applies, choose actual generators of \(A_i/\operatorname{im}(R)\); each is annihilated by an element of \(R\setminus\mathfrak p\). The product of these denominators makes \(R_f\to(A_i)_f\) surjective. For all finitely many factors take the product of all their denominators, so one \(f\notin\mathfrak p\) works simultaneously. If no factor occurs, the product is 1. Localizing the étale base preserves étaleness and its specified prime. The image \(\mathfrak pA_i\) is prime with contraction exactly \(\mathfrak p\), because \(A_i=R/I_i\) and the surviving fibre implies \(I_i\subset\mathfrak p\). This supplies the precise simultaneous step implicit at 42049--42050.

For the one-point lemma, collect the other finite factors and the old remainder into the complementary factor. For the everywhere-unramified lemma, quasi-finiteness holds at every prime. Hence the remainder \(B\), which is not quasi-finite at any prime of the selected fibre, has no such prime at all. Equivalently \(B\otimes_R\kappa(\mathfrak p)=0\); a nonzero commutative ring has a prime ideal, so a nonzero fibre would contradict the assertion.

Writing the final surjections as \(R'\to R'/I_i\) makes the resulting decomposition exact: \[
R'\otimes_RS=\prod_{i=1}^n R'/I_i\times B.
\] For any \(R'\)-algebra \(T\), tensoring gives \[
T\otimes_RS=\prod_{i=1}^n T/I_iT\times(T\otimes_{R'}B).
\] The map is the coordinate map on pure tensors; its inverse uses the finite orthogonal idempotents selecting the original factors, whose sum is 1. Thus it is a ring isomorphism for arbitrary base change. At a prime \(\mathfrak t\) contracting to the specified \(\mathfrak p'\), each \(I_iT\subset\mathfrak t\) and each displayed quotient has fibre \(\kappa(\mathfrak t)\). The remaining fibre is \[
(B\otimes_{R'}\kappa(\mathfrak p'))
   \otimes_{\kappa(\mathfrak p')}\kappa(\mathfrak t)=0.
\] These claims concern the specified fibre and its pullbacks.

The remainder cannot in general be discarded on a neighbourhood. Take \(R=k[t]\), \(\mathfrak p=(t)\), and the unramified (indeed étale) algebra \(S=R\times R_t\). The second factor has zero fibre at \(\mathfrak p\), but for every \(f\notin(t)\) its localization is \(R_{tf}\ne0\), since \(tf\) is a nonzero polynomial in a domain. Thus the exact decomposition retains a nonzero remainder on every principal neighbourhood of the specified point.

\subsubsection{The exact conormal object on an etale quotient chart}\label{algebra-proof-057-the-exact-conormal-object-on-an-etale-quotient-chart}

Let \(E\) be a formally étale \(R\)-algebra and let \(A=E/I\), for any ideal \(I\). No finite-generation hypothesis is imposed here. The local standard étale charts above provide this situation for the original unramified map. Define the actual rings and module \[
U=E/I^2,\qquad C=I/I^2,\qquad
0\longrightarrow C\longrightarrow U\longrightarrow A
\longrightarrow0.
\] The ideal \(C\) is square-zero.

For an \(R\)-algebra \(B\), a square-zero ideal \(J\subset B\), and an \(R\)-map \(a:A\to B/J\), formal étaleness gives a unique lift \(\beta:E\to B\) of \(E\to A\xrightarrow a B/J\). For \(i\in I\), \(\beta(i)\in J\), so \(\beta(I^2)\subset J^2=0\). Hence \(\beta\) factors through a unique \(q:U\to B\), compatible with \(a\) on the quotients. Conversely any such \(q\) precomposed with \(E\to U\) is this unique lift. Thus \(U\to A\) has exactly the universal square-zero thickening property constructed in group 1205.

The restriction \[
\lambda=q|_C:C\longrightarrow J
\] is \(A\)-linear using \(a\) for the action on \(J\): changing a lift of an element of \(A\) changes its product with \(q(C)\subset J\) by an element of \(J^2=0\). A lift \(A\to B\) of \(a\) exists exactly when \(\lambda=0\), since this is exactly the condition that \(q\) kill \(\ker(U\to A)\). It is unique: any two lifts give two lifts from \(E\), which agree, and \(E\to A\) is surjective. This also proves that \(A/R\) is formally unramified.

If another formally étale presentation \(A=E'/I'\) is given, apply the universal property to \(B=E'/I'^2\), \(J=I'/I'^2\), and the identity on \(A\). It gives a unique map \(E/I^2\to E'/I'^2\) over \(A\). The reverse construction gives its inverse, since both composites are the unique endomorphism over the identity of \(A\). Restricting these maps gives inverse \(A\)-linear maps \(I/I^2\leftrightarrow I'/I'^2\). This constructs the exact comparison with the earlier universal conormal object, not merely a claimed similarity of presentations.

It follows that \(A/R\) is formally étale if \(I=I^2\). Conversely, if \(A/R\) is formally étale, lift its identity through \(E/I^2\to A\) to a section \(s:A\to E/I^2\). The maps \(E\to E/I^2\) given by the quotient and by \(E\to A\xrightarrow s E/I^2\) lift the same map into \(A\). Formal étaleness of \(E/R\) makes them equal. Every \(i\in I\) therefore maps to zero in \(E/I^2\), so \(I\subset I^2\). We have proved the exact criterion \[
A/R\text{ formally étale}\quad\Longleftrightarrow\quad I=I^2.
\]

For a standard étale \(E=R[x]_g/(f)\), use its unchanged presentation \[
E=R[x,z]/(f,zg-1).
\] Then \(A/R\) is finitely presented exactly when \(I\) is finitely generated as an ideal of \(E\). In one direction, lifts of finitely many generators of \(I\), together with \(f,zg-1\), present \(A\). In the other, finite-presentation independence at 566--581 makes the kernel of the given polynomial surjection \(R[x,z]\to A\) finitely generated; its image in \(E\) generates \(I\). Thus the source's G-unramified finiteness condition on the chart is precisely finite generation of this particular ideal.

If \(I=I^2\) and \(I\) is generated by \(u_1,\ldots,u_m\), choose \(a_{ij}\in I\) with \(u_i=\sum_j a_{ij}u_j\), and set \[
\delta=\det(\mathrm{Id}-(a_{ij}))
 =\sum_{\pi\in\mathfrak S_m}\operatorname{sgn}(\pi)
       \prod_{i=1}^m(\delta_{i,\pi(i)}-a_{i,\pi(i)}).
\] The adjugate identity gives \(\delta I=0\), and \(\delta-1\in I\). Therefore \(\delta^2=\delta\), and \(e=1-\delta\in I\) satisfies \(e^2=e\) and \(I=eE\): for every \(u\in I\), \(u=(1-\delta)u=eu\). The empty generator list gives \(\delta=1\), \(e=0\). The quotient and localization have inverse ring maps \[
E/I=E/eE\ \longleftrightarrow\ E_{1-e},
\] since \(1-e\) maps to 1 in the quotient, while inverting \(1-e\) forces \(e=0\). Hence a finitely presented formally étale chart is an actual principal localization of the original étale algebra. All determinant factors and signs are retained.

Finally let \(R\to R_1\) be arbitrary, put \(E_1=R_1\otimes_RE\), \(A_1=R_1\otimes_RA\), and \(I_1=\ker(E_1\to A_1)\). Right exactness of tensor identifies \(I_1\) with the image of \(R_1\otimes_R I\). Products of such images show that \(I_1^2\) is the image of \(R_1\otimes_R I^2\). Hence the canonical maps give \[
R_1\otimes_R U=E_1/I_1^2.
\] Formal étaleness is preserved: an \(R_1\)-algebra lifting problem for \(E_1\) is the corresponding \(R\)-problem for \(E\), and the unique lift extends to \(E_1\) by its prescribed \(R_1\)-action. Thus the displayed ring is the universal thickening of \(A_1\). Its actual kernel \(C_1=I_1/I_1^2\) is the image of \(R_1\otimes_R C\), with exact sequence \[
\operatorname{Tor}_1^R(R_1,U)\longrightarrow
\operatorname{Tor}_1^R(R_1,A)\longrightarrow
R_1\otimes_R C\longrightarrow C_1\longrightarrow0.
\] This is the tensor long exact sequence of \(0\to C\to U\to A\to0\), ending at the actual kernel of \(R_1\otimes_R U\to R_1\otimes_R A\). It is not permissible to replace this image by the tensor product without checking the preceding map. For \(R=E=\mathbf Z\), \(I=(p)\), \(A=\mathbf F_p\), the original \(U=\mathbf Z/p^2\) and \(C=p\mathbf Z/p^2\). After \(R_1=\mathbf F_p\), \(I_1=0\), \(U_1=A_1=\mathbf F_p\), and \(C_1=0\), whereas \(R_1\otimes_R C=\mathbf F_p\). This retains and confirms the earlier exact base-change boundary.

\subsubsection{Report accounting and propagation}\label{algebra-proof-057-report-accounting-and-propagation}

Groups 1215--1221 account for all twelve reports, with nine operations in seven proposed units. They respectively repair the repeated clause with an explicit image, the copula, the polynomial-ring membership, the nonzero scalar, all three occurrences of that scalar, the lift inside \(I\), and the explicit exponent \(e_1=1\). The report's strict typing objection to conventional prime-membership notation is not accepted as an independent mathematical defect. No optional or rejected group and no duplicate canonical operation is introduced in this interval.

Group 1219 receives group 1145's scalar correction at this distinct source locus. Group 1222 has no source operation: it receives 1148's finite-cokernel comparison, 1168's finite-fibre decomposition, 1205's universal conormal thickening and 1214's idempotent calculation. It records the direct arbitrary-base proof, exact separated quotient factors, universal chart thickening, obstruction map, presentation comparison, formal-étale and finite-presentation criteria, and the actual kernel after arbitrary base change. These arguments remain separate editorial material. The original translation, mathematical objects, eleven-step exposition and finite-type versus finite-presentation conventions remain identifiable. Broader synthesis and any paper follow completion of the core work.

\hypertarget{algebra_henselian_notes_20260929}{}
\subsection{Henselian lifting, finite algebras and the corrected module endomorphism}\label{algebra-proof-058-henselian-lifting-finite-algebras-and-the-corrected-module-endomorphism}

\subsubsection{Source and scope}\label{algebra-proof-058-source-and-scope}

Source: the Stacks Project authors, algebra.tex at a04446e57ec1fbc252a871afcec7752fb2807b14, SHA-256 FA8BB92E58A4F78A2BD01B3B6A4A87DE0A0D279F5DD90641B574DD5FBFFFA4F3. The complete interval 42062--42760 and all fourteen received reports were read. The next opening 42761--42785 is read but unadjudicated. No prior canonical correction operation occurs in the present interval. Bounded fresh dependencies: 2833--2867, 19221--19247, 21637--21669, 21858--21890, 22115--22182 and 32867--32893. Earlier source-linked notes retain the already checked local étale, unramified, conormal and Mittag-Leffler arguments. The bounded dependency excerpts are not claims of complete new reading of each surrounding proof.

The two indexed hits for ``henselian Mittag Leffler'' remain unread and unused. This is source correction and editorial derivation, with no novelty claim. Preserve the author's definitions and sequence of exposition; do not replace a translated proof by the fuller arguments below.

\subsubsection{Taylor formulas and all thirteen characterizations}\label{algebra-proof-058-taylor-formulas-and-all-thirteen-characterizations}

For the source polynomial \(f(T)=\sum_{i=0}^d c_iT^i\), retain every coefficient and use the exact identity \[
f(X+Y)-f(X)=f'(X)Y+Y^2G(X,Y),
\qquad
G(X,Y)=\sum_{i=2}^d c_i\sum_{j=2}^i
             {i\choose j}X^{i-j}Y^{j-2}.
\] It holds over any commutative ring by the binomial theorem, including all positive characteristics. The derivative is \(f'(X)=\sum_{i=1}^d i c_iX^{i-1}\); no term is silently dropped. If \(f(a)=f(b)=0\), \(b-a\in\operatorname{Jac}(R)\), and \(f'(a)\) is a unit, the identity gives \[
(b-a)f'(a)\bigl(1+f'(a)^{-1}(b-a)G(a,b-a)\bigr)=0.
\] The parenthesized factor is a unit because its difference from 1 is in every maximal ideal. Therefore \(a=b\). This proves the source's uniqueness lemma and its stated local-ring use, while also identifying the weaker Jacobson-radical hypothesis.

There are two zero-factor cases in the thirteen-condition lemma. The received report notices \(g_0=1,h_0=0\), but the stated hypotheses also permit \(g_0=0,h_0=1\). A degree comparison with a zero polynomial is undefined unless a convention is specified; with the usual \(\deg(0)=-\infty\), the original condition (6) is false. For example, in the henselian ring \(R=k[\epsilon]/(\epsilon^2)\), the nonzero constant polynomial \(f=\epsilon\) has residue factorization \(\bar f=0\cdot1\). No factorization with \(\deg(g)=\deg(0)\) can have product \(\epsilon\). Independently, this ring is henselian: lifting a simple residue root to \(b\in R\), the correction \(b-f(b)/f'(b)\) is an exact root, since \(f(b)\in(\epsilon)\) and \((\epsilon)^2=0\).

The minimal corrected condition (6) retains the factorization for every coprime \(g_0,h_0\), and requires \(\deg(g)=\deg(g_0)\) when \(g_0\ne0\). If \(g_0=0\), coprimality forces \(h_0\) to be a nonzero constant. Lift it to a unit \(h\in R\) and set \(g=h^{-1}f\). If \(h_0=0\), lift the nonzero constant \(g_0\) to a unit \(g\), and set \(h=g^{-1}f\). This case has the required degree zero. These constructions retain the exact prescribed residue factors. They repair both the omitted proof cases and the unrestricted degree assertion. Conditions (3) and (4), whose \(f\) is monic, have no zero residue factor.

We now verify the implications, including the maps omitted from the source prose. The immediate implications 2→1, 5→3, 6→4, 4→3, 6→5, 8→7 and 10→9 follow by restricting the polynomial class or forgetting the extra requirement. In 6→4 the monic residue polynomial is nonzero, so the corrected qualification in (6) is automatic. For a finite-type algebra over a field, a point is quasi-finite exactly when it is an isolated zero-dimensional fibre component. A zero-dimensional irreducible component of a Noetherian fibre is a single closed point disjoint from all other components, hence is open as well as closed. Conversely an isolated point is such a component. This proves 11↔12. In the quasi-finite situation all fibre points are of this kind, giving 11→13.

For 1→8, choose a principal neighbourhood of the specified prime on which the original étale algebra is standard: \(S_g=R[t]_h/(f)\), with \(f\) monic and \(f'\) invertible there. The residue point specifies \(a_0\in\kappa\) with \(\bar f(a_0)=0\), \(\bar h(a_0)\ne0\), and \(\bar f'(a_0)\ne0\). Henselianity gives its unique lift \(a\). Evaluation \(t\mapsto a\), \(h^{-1}\mapsto h(a)^{-1}\) gives the retraction on the neighbourhood and then on \(S\). Any retraction with the prescribed contraction sends \(g\) to a unit and therefore factors through this same localization. Its image of \(t\) has residue \(a_0\), so the uniqueness calculation above makes the two retractions equal.

For 7→8, there are finitely many other fibre primes, because an étale algebra over \(\kappa\) is a finite product of finite separable fields. For each other prime \(\mathfrak q_i\), choose \(g_i\in\mathfrak q_i\setminus\mathfrak q\) and use the product \(g=\prod_i g_i\); an empty product is 1. Then \(S_g\) has just the specified prime above the maximal ideal. Any retraction of \(S_g\) must contract that maximal ideal to this prime. It exists by (7). Its uniqueness follows from the preceding standard-neighbourhood argument, which uses uniqueness of an existing root and does not assume henselian existence. This avoids a circular use of 1→8.

For 8→11, use the previously proved étale finite-component decomposition \(R'\otimes_RS=A'\times B'\) at a specified \(\mathfrak m'\) with residue field \(\kappa\). The retraction \(\tau:R'\to R\) has \[
\mathfrak m'=\tau^{-1}(\mathfrak m),
\] not the reversed expression at 42232. The actual inverse maps \[
(S\otimes_RR')\otimes_{R',\tau}R\ \longleftrightarrow\ S
\] send \((s\otimes r')\otimes r\) to \(s\tau(r')r\), and send \(s\) to \((s\otimes1)\otimes1\). They split the finite product into \(A'\otimes_{R',\tau}R\) and \(B'\otimes_{R',\tau}R\). The residue map \(R'/\mathfrak m'\to R/\mathfrak m\) is the given isomorphism, so the new closed fibre is the same original fibre. Its quasi-finite points and non-quasi-finite components are therefore preserved exactly. This proves (11).

For 8→10, apply the same maps separately to every \(A'_i\). The last displayed factor at 42255 is \(A'_n\), not a repeated \(A'_1\). The resulting \(B\) is finite over the local ring \(R\). If it were nonzero, it would have a maximal ideal lying over \(\mathfrak m\) by integrality, and a finite map is quasi-finite there, a contradiction. Thus \(B=0\). Each finite \(A_i\) has exactly one maximal ideal, since all its maximal ideals lie over \(\mathfrak m\). For 9→10, a finite \(R\)-algebra has finitely many maximal ideals: they are the primes of its finite-dimensional closed fibre. A product of infinitely many nonzero local rings would have at least one distinct maximal ideal for each coordinate projection. Hence only finitely many nonzero factors occur.

For 10→1, the simple root \(a_0\) of the monic residue polynomial selects the factor with closed fibre \(\kappa[T]/(T-a_0)=\kappa\) in the finite algebra \(R[T]/(f)\). That factor is a direct summand of a finite free \(R\)-module and hence is free over the local ring. Its rank is one. Its unit element has nonzero residue, so in any rank-one basis its coefficient is a unit. The structure map from \(R\) is therefore an isomorphism. The image of the original \(T\) is the required lift. For 13→1, first use \[
S=(R[T]/(f))_g,\qquad \bar g=\bar f/(T-a_0).
\] This localization is flat over \(R\), and its closed fibre is \(\kappa\). The finite summand selected by (13) is finite flat, hence free over the local ring. Its rank is one, and the same unit argument gives the lift. All localizations retain the chosen \(g\).

For 8→2, the algebra \(R[T]_{f'}/(f)\) is étale by its invertible Jacobian \(f'\), whether or not \(f\) is monic. The residue point is the simple root of \(\bar f\), made explicit at 42311. The retraction sends \(T\) to the required root. For 3→1, the lifted factorization of the monic \(f\) gives \(gS+hS=S\) in \(S=R[T]/(f)\) by Nakayama on the finite module \(S/(g,h)\). Since \(gh=0\) there, the comaximal ideals have zero intersection. The Chinese remainder map and its usual inverse give \(S=S/(g)\times S/(h)\). The first factor is rank-one free as above and supplies the root.

Finally, for 1→6 handle both zero-factor cases first. When both factors are nonzero, keep their original leading scalar explicitly: let \(c_0\) be the leading coefficient of \(g_0\), choose a lift \(c\in R^\times\), and set \[
G_0=c_0^{-1}g_0,\qquad H_0=c_0h_0.
\] These are auxiliary factors only; the final factors below recover the originally prescribed \(g_0,h_0\). Now \(\bar f\ne0\), so at least one coefficient of \(f\) is a unit. Every fibre polynomial is nonzero, and \(R[T]/(f)\) has finite fibres. For any prime of this quotient, localize \(R[T]\) at its preimage and \(R\) at the contraction. The polynomial remains a nonzero divisor on the localized polynomial fibre, a localization of a domain. The hypotheses of 32869--32878 hold at these actual local rings: the polynomial ring localization is flat and essentially finitely presented over the localized base. Thus the quotient is flat at every prime, hence flat over \(R\).

Conditions (13) and (10) give the original finite local factors and the remainder of empty closed fibre. Select precisely the factors whose residue quotient is \(\kappa[T]/(G_0)\). Their product \(A\) is finite flat, hence free of rank \(d=\deg G_0\). Let \(G(T)\) be the characteristic polynomial of multiplication by the original class of \(T\) on \(A\). Reduction of that actual operator on \(\kappa[T]/(G_0)\) has characteristic polynomial \(G_0\), so \(G\) is monic of degree \(d\) and \(\bar G=G_0\). Cayley--Hamilton gives the original quotient map \(R[T]/(G)\to A\). It is a surjection between free modules of rank \(d\), with invertible reduction. Its matrix determinant is a unit, and its inverse is the adjugate divided by that determinant. Thus its kernel is zero; the kernel of \(R[T]\to A\) is exactly \((G)\). As the map also kills \(f\), there is a polynomial \(H\) with \(f=GH\). Reducing and cancelling the nonzero \(G_0\) in \(\kappa[T]\) gives \(\bar H=H_0\). Return to the original factors: \[
g=cG,\qquad h=c^{-1}H,\qquad
f=gh,\quad\bar g=g_0,\quad\bar h=h_0,\quad\deg g=\deg g_0.
\] This supplies the inverse comparison behind the source's monic reduction, without losing the chosen leading scalar. The rank-zero case uses \(A=0\), characteristic polynomial 1 and \(G=1\).

\subsubsection{Finite algebras, residue functors and arbitrary local targets}\label{algebra-proof-058-finite-algebras-residue-functors-and-arbitrary-local-targets}

A finite algebra over a henselian local ring is a finite product of local rings by the characterization. Each local factor is itself henselian: any algebra finite over it is finite over \(R\), hence is a product of local rings. A finite map between local rings is local because maximal ideals contract to maximal ideals under integral maps. The finite-type decomposition in (11) then shows that a quasi-finite closed-fibre localization is exactly one of these finite henselian factors. This proves all assertions at 42384--42447. Over a separably algebraically closed residue field, every remaining finite residue extension is purely inseparable; a separable part would otherwise be a nontrivial finite separable extension.

We verify the finite-étale equivalence at 42449--42505. The source assertion that a local factor \(A_i\) is finite étale over \(S\) is true: projection onto a factor is localization at its defining idempotent, and the factor is a finite projective \(S\)-module. It is also finite étale over \(R\) by composition. The report claiming that ``over \(S\)'' is false is rejected.

For local finite étale \(S_1,S_2\), put \(S_1\otimes_RS_2=\prod_i A_i\) with local factors. An \(R\)-map \(\varphi:S_1\to S_2\) gives the actual map \[
a\otimes b\longmapsto\varphi(a)b
 =\mu_{S_2}\bigl((\varphi\otimes\mathrm{id})(a\otimes b)\bigr).
\] The printed name \(\varphi\times1\) can denote this induced map. Replacing it by a bare \(\varphi\otimes1\) would not by itself specify the multiplication into \(S_2\). The notation change is therefore optional, with this exact formula recorded editorially. The map is surjective because it contains \(1\otimes S_2\). Idempotents select one local factor, so it factors through some \(A_i\to S_2\). This is a surjection of étale \(R\)-algebras, hence is étale and flat with finitely generated idempotent kernel. Since \(A_i\) is local and \(S_2\ne0\), the kernel is zero.

Conversely a factor for which \(S_2\to A_i\) is an isomorphism gives such a map. This happens exactly when the residue map \(\kappa_2\to A_i/\mathfrak mA_i\) is an isomorphism: \(A_i\) is finite flat over \(S_2\), so is free over this local ring, and the residue condition gives rank one with the unit as a basis. This proves the matching of index sets asserted in the source and the bijection on Hom sets for local objects. For products, each target local factor selects exactly one source factor, and these local Hom-set bijections combine uniquely. Their compositions are the original compositions of algebra maps. Thus the functor is fully faithful on all finite étale algebras.

For a finite separable residue extension, the prescribed-residue-field étale construction gives an étale \(R\)-algebra and the specified closed-fibre point. The decomposition needed next is lemma-mop-up, not an undefined ``part (1)'' at 42498. Its corresponding finite local factor remains étale and has the prescribed residue extension. Taking finite products proves essential surjectivity. Equivalently, choose a primitive element of the field extension, lift all coefficients of its monic minimal polynomial to \(R\), and form its finite free quotient. Its closed fibre is the given field, so henselian decomposition has one local factor; the derivative is a unit at its only maximal ideal, making this quotient finite étale. Different choices are compared by the unique lifted residue isomorphism just proved.

For a strictly henselian \(R\), every finite local factor of an unramified algebra has residue field \(\kappa\), by separability, and maximal ideal \(\mathfrak mA_i\). Nakayama on the finite cokernel of \(R\to A_i\) makes the map surjective. The complementary factor has no closed-fibre prime by quasi-finiteness. This proves 42507--42531 without adding finite presentation.

For the map-into-henselian lemma, the expression \(R\cap\mathfrak m_S\) is conventional contraction along the given map; no injectivity is assumed. Write it as its inverse image when checking the following maps. For the given residue embedding \(\tau\), the actual map defining \(\mathfrak q'\) is \[
A\otimes_RS\longrightarrow\kappa(\mathfrak m_S),
\qquad a\otimes s\longmapsto\tau(\bar a)\bar s.
\] It is onto because of the \(S\)-factor, so its kernel has exactly the required residue field. The unique henselian retraction \(\sigma:A\otimes_RS\to S\) gives \(f(a)=\sigma(a\otimes1)\). Its residue is the specified formula, hence its contraction is \(\mathfrak q\), and its induced residue map is \(\tau\). Conversely any eligible \(f\) gives \(\sigma_f(a\otimes s)=f(a)s\), with that same residue kernel. Uniqueness of the retraction proves uniqueness of \(f\).

The solution-bijection lemma at 42634--42690 admits a stronger target hypothesis. Let \(R\) be strictly henselian, let \(R\to S\) be any local map to a local ring, and let \(A\) be any étale \(R\)-algebra. Then \[
\operatorname{Hom}_R(A,R)\longrightarrow
\operatorname{Hom}_R(A,S)
\] is a bijection; \(S\) need not be henselian. Indeed write \(A=A_1\times\cdots\times A_n\times B\) by the source's quasi-finite decomposition. The finite local étale factors have residue field \(\kappa\), hence rank one and structure maps \(R\cong A_i\). The remaining \(B\otimes_R\kappa\) is zero. A map from this finite product to a nonzero local ring selects exactly one factor, since its only idempotents are 0 and 1. It cannot select \(B\): reduction to \(\kappa(S)\) would give a unital map from \[
B\otimes_R\kappa(S)
=(B\otimes_R\kappa)\otimes_\kappa\kappa(S)=0
\] to the nonzero field \(\kappa(S)\). Thus every map selects exactly one of the copies of \(R\), and on that copy is the prescribed structure map. The same description holds with target \(R\), so the displayed map preserves the identical finite index set and is bijective, including the case \(n=0\). This applies to the source's polynomial solution sets via their exact representing algebra, without changing its equations.

The relaxation is substantive: for an algebraically closed field \(k\) of characteristic zero, \(R=k\) is strictly henselian, while \(S=k[t]_{(t)}\) is not henselian. The simple residue root 1 of \(X^2-(1+t)\) cannot lift in \(S\), since a square in \(k(t)\) has even order at the irreducible polynomial \(t+1\), whereas \(1+t\) has order one there. Nevertheless the Hom-set bijection just proved holds for every étale \(k\)-algebra.

\subsubsection{Complete and nil-ideal lifting with exact errors}\label{algebra-proof-058-complete-and-nil-ideal-lifting-with-exact-errors}

The Newton argument at 42533--42594 does not require monicity. More generally let \(R\cong\varprojlim R/I^n\), with its stated separated completeness, and let \(a_0\in R/I\) satisfy \(f(a_0)=0\), with \(f'(a_0)\) a unit. At stage \(n\), choose any lift \(b\in R/I^{n+2}\) of \(a_n\in R/I^{n+1}\) and put \[
e=f(b),\quad u=f'(b),\quad a_{n+1}=b-eu^{-1}.
\] The derivative is invertible: its reduction modulo \(I\) is a unit, and the kernel \(I/I^{n+2}\) is nilpotent, so the inverse lifts by a finite geometric series. Also \(e\in I^{n+1}/I^{n+2}\), so \(e^2=0\) in this ring. The full Taylor identity gives \[
f(a_{n+1})=e-u(eu^{-1})
              +(eu^{-1})^2G(b,-eu^{-1})=0.
\] The correction lies in \(I^{n+1}\), proving compatibility. Completeness gives the element \(a\), and separatedness gives \(f(a)=0\), since all its quotient images vanish. For each \(z\in I\) and \(r\in R\), the convergent geometric series \(\sum_{j\ge0}(rz)^j\) in the original \(I\)-adic ring is an inverse of \(1-rz\). Hence \(I\subset\operatorname{Jac}(R)\), and the Taylor uniqueness argument proves uniqueness of \(a\). This includes the original maximal-ideal-adic local case with no Noetherian hypothesis.

If instead \(I\) is any nil ideal, choose an initial lift \(b_0\). The derivative is a unit because units lift modulo a nil ideal. Define \(b_{j+1}=b_j-f(b_j)f'(b_j)^{-1}\) in the original ring. Every derivative remains a unit, since \(b_j-b_0\in I\). If \(e_0=f(b_0)\), the same complete Taylor expression proves \[
f(b_j)\in(e_0)^{2^j},\qquad b_j-b_0\in I .
\] The element \(e_0\) is nilpotent, so some finite \(j\) makes \(f(b_j)=0\). A nil ideal is contained in the Jacobson radical, and uniqueness follows as before. The ideal itself need not be nilpotent. A local ring of dimension zero has nil maximal ideal, so this also proves the source's zero-dimensional assertion directly, retaining its full non-Noetherian generality.

There is an exact related improvement to the formal-étale lifting boundary in group 1204. If \(A\) is finitely presented and formally étale over \(R\), then every \(R\)-map \(A\to B/I\), with \(I\) nil, lifts uniquely to \(B\). Choose the original finite presentation \(A=R[x_1,\ldots,x_n]/(F_1,\ldots,F_s)\) and lifts \(b_i\in B\). The finitely many errors \(F_j(b_1,\ldots,b_n)\) lie in \(I\). If their nilpotence exponents are \(N_j\), the ideal \(J\) they generate satisfies \[
J^{\,1+\sum_{j=1}^s(N_j-1)}=0,
\] because every monomial of that total degree contains one generator to its stated nilpotence exponent. When there are no errors use \(J=0\). The chosen lifts give an actual map \(A\to B/J\), which lifts uniquely across this nilpotent ideal by finite iteration of the square-zero lifting property. Its reduction modulo \(I\) is the prescribed map. For uniqueness, the differences of the finitely many generator images of two lifts generate another finitely generated nil ideal, hence a nilpotent ideal by the same exponent calculation. The two maps agree modulo it, so formal unramifiedness makes them equal by the finite nilpotent filtration. This explains exactly why the earlier non-finitely-presented formally étale counterexample over a nonzero nil idempotent ideal does not contradict this result. It does not turn a general nil ideal into a globally nilpotent one.

\subsubsection{The original endomorphism and the missing square}\label{algebra-proof-058-the-original-endomorphism-and-the-missing-square}

Keep the source module \(P\), the maps \[
M\xrightarrow{\pi}P\xrightarrow{i}M,\qquad
\alpha=i\pi,\qquad a=\pi\alpha i,
\] and the original monic polynomial called \(P(T)\) that annihilates \(a\). The module and polynomial are different typed uses of the same letter; a renaming would be stylistic and would not repair the mathematical error.

Put \(b=\pi i\) only as an auxiliary composition. Then \(a=b^2\), and for every \(j\ge0\), \[
i a^j\pi=i b^{2j}\pi=\alpha^{2j+1}.
\] If \(P(T)=\sum_{j=0}^d c_jT^j\), with \(c_d=1\), it follows that \[
0=iP(a)\pi
 =\sum_{j=0}^d c_j\alpha^{2j+1}
 =\alpha P(\alpha^2).
\] Multiplying by \(\alpha\) gives the relation used for the corrected source construction: \[
\alpha^2P(\alpha^2)=0.
\] There is no deduction of \(\alpha^2P(\alpha)=0\).

For a concrete counterexample satisfying the fixed-vector hypothesis, let \(R=\mathbf Q\), \(M=P=\mathbf Q^2\), \(\pi=\mathrm{id}\), and \(i=\alpha=\operatorname{diag}(1,-1)\). The vector \(x=(1,0)\) is fixed, \(a=\mathrm{id}\), and its characteristic polynomial \(P(T)=(T-1)^2\) annihilates \(a\). But \[
\alpha^2P(\alpha)=\operatorname{diag}(0,4)\ne0.
\] The corrected relation is zero because \(\alpha^2=\mathrm{id}\). This disproves the printed intermediate assertion without altering the source's \(a,\pi,i,\alpha\) or the chosen polynomial.

The required repairs are therefore \(P(\alpha^2)\) at 42730 and 42733, the actual algebra \(R[T]/(T^2P(T^2))\) at 42731, and the full factorization \[
T^2P(T^2)=(T^2Q_1(T))Q_2(T)
\] at 42739. To construct it, assume \(x\ne0\); if \(x=0\), simply take the desired finite summand to be zero. Since \(\alpha x=x\), the corrected relation gives \(P(1)x=0\). The annihilator of the nonzero vector is a proper ideal of the local ring, so \(P(1)\in\mathfrak m\). Let \(e\ge1\) be the exact multiplicity of \(T-1\) in \(\overline{P(T^2)}\), including any increase in characteristic two. Apply the monic coprime-factor lifting already proved to \(P(T^2)\). It gives monic \(Q_1,Q_2\) with \[
P(T^2)=Q_1Q_2,\quad
\bar Q_2=(T-1)^e,\quad\bar Q_1(1)\ne0.
\] This makes the source's previously unquantified \(e\) explicit.

Set \(A(T)=T^2Q_1(T)\), \(B(T)=Q_2(T)\). Their reductions are coprime, because \(B\) has only the residue root 1 and \(Q_1(1)\) is a unit. The quotient \(R[T]/(B)\) is finite free since \(B\) is monic. Multiplication by \(A\) on it has invertible reduction, hence unit determinant and inverse given by its adjugate divided by that determinant. Apply the inverse to 1 and choose its representative \(U(T)\) of degree less than \(\deg B\). Monic division then gives the actual polynomial \(V(T)\) satisfying \[
U(T)A(T)+V(T)B(T)=1 .
\] Here the numerator \(1-UA\) is divisible by the original \(B\); no coefficient, constant term or power of \(T\) has been removed.

Because \(A(\alpha)B(\alpha)=0\), the endomorphisms \[
E=U(\alpha)A(\alpha),\qquad
1-E=V(\alpha)B(\alpha)
\] are complementary idempotents. In detail, their sum is the identity, their product in either order is zero, and multiplying the sum by \(E\) gives \(E^2=E\). Further \(EA(\alpha)=A(\alpha)\) and \((1-E)B(\alpha)=B(\alpha)\). Therefore the original images are \[
N=\operatorname{Im}(A(\alpha))=\operatorname{Im}(E),\qquad
K=\operatorname{Im}(B(\alpha))=\operatorname{Im}(1-E),
\qquad M=N\oplus K .
\] We have \(B(\alpha)N=0\), \(A(\alpha)K=0\). The fixed vector lies in \(N\), since \(A(\alpha)x=Q_1(1)x\), with \(Q_1(1)\) a unit.

The source's assertion that \(N\) is a quotient of the module \(P\) also needs its actual map. Because \[
A(\alpha)=\alpha^2Q_1(\alpha)
=i\bigl(\pi\alpha Q_1(\alpha)\bigr),
\] we have \(N\subset\operatorname{Im}(i)\). The map \(Ei:P\to N\) is surjective: for any \(n=A(\alpha)m\), choose \(p=\pi\alpha Q_1(\alpha)m\), so \(i(p)=n\) and \(Ei(p)=En=n\). Consequently \(N\) is finitely generated. Both summands are Mittag-Leffler by the source's direct-sum criterion, and a finitely generated Mittag-Leffler module is finitely presented by the already proved criterion at 21860--21868. This completes the local extraction argument.

It gives more than one-vector extraction. For any prescribed finite family \(x_1,\ldots,x_s\), apply lemma-ML-countable to the map \(R^s\to M\) with those images. Its one endomorphism \(\alpha\) fixes all of them and factors through one finitely presented module. If some \(x_j\ne0\), the same polynomial and idempotent \(E\) work, and \(A(\alpha)x_j=Q_1(1)x_j\) puts every vector in the same finitely presented direct summand. If all are zero, use the zero summand. Thus the entire finitely generated submodule they span is contained in one such summand by one explicitly constructed projection.

For the final countable decomposition, choose a sequence of generators of \(M\). After splitting off \(N_1,\ldots,N_n\), project the next generator into the remaining summand \(K_n\). This summand is countably generated, as a quotient of \(M\), and is Mittag-Leffler. Apply the extraction just proved to that projected generator. This yields nested direct sums containing successively all original generators. Their union is \(M\). The induced map \(\bigoplus_iN_i\to M\) is surjective, and is injective because every finite partial sum is a direct summand. This proves the original theorem with its original maps and the corrected polynomial argument.

\paragraph{Exact-map diagram and the rational test}\label{algebra-proof-058-exact-map-diagram-and-the-rational-test}

\begin{figure}
\centering
\pandocbounded{\includegraphics[keepaspectratio,alt={The original factorization and the corrected projection}]{HENSELIAN_PROJECTION_20260929.png}}
\caption{The original factorization and the corrected projection}
\end{figure}

The left panel records the original maps in the preceding proof. The right panel keeps the same rational counterexample and displays its corrected projection. Here the complete factors and coefficients are \[
A(T)=T^2(T+1)^2,\qquad B(T)=(T-1)^2,\qquad
U(T)=\frac{4-3T}{4},\qquad
V(T)=\frac{3T^3+8T^2+8T+4}{4}.
\] Indeed \[
4U(T)A(T)=-3T^5-2T^4+5T^3+4T^2,
\quad
4V(T)B(T)=3T^5+2T^4-5T^3-4T^2+4,
\] so their sum is exactly 4. Thus the displayed Bezout identity holds with every coefficient retained. Evaluation at the original \(\alpha=\operatorname{diag}(1,-1)\) gives \(A(\alpha)=\operatorname{diag}(4,0)\), \(B(\alpha)=\operatorname{diag}(0,4)\), and \(E=\operatorname{diag}(1,0)\). In particular the chosen fixed vector lies in the first summand. This is an illustration of the proved correction, with no change of the source's endomorphisms. Source: the Stacks Project authors, authority algebra.tex 42718--42760; the complete corrected argument is in the preceding section. Reproducible figure source: draw\_henselian\_projection\_20260929.py; vector edition: HENSELIAN\_PROJECTION\_20260929.svg. The PNG was inspected for legibility, exact formulas and arrow directions.

\subsubsection{Report dispositions and propagation}\label{algebra-proof-058-report-dispositions-and-propagation}

The fourteen reports form eleven report groups. Seven source-correction units result after adding the independently found missing-square correction: eleven operations in total. The inverse-image orientation and product index are repaired; the residue polynomial and positive exponent are clarified; both zero-factor cases and the degree qualification are repaired; the undefined reference is replaced by lemma-mop-up. The module-polynomial letter reuse and the tensor-map name remain optional style changes. The quotient-polynomial scope convention, contraction notation and the true finite-étale-over-S assertion are not mathematical defects and are rejected as such.

The missing-square correction is propagated through all four uses of the annihilating relation, the lifted factors, both image summands, the actual surjection from the original finite module and the countable direct-sum proof. The finite-family projection is a proved consequence of that corrected construction. Separate editorial groups record complete and nil-ideal lifting, the finite-presentation boundary for general nil-ideal formal étale lifting, and the stronger Hom-set statement allowing every local target of a strictly henselian base. All remain visibly separate from the translated original. The broader synthesis and any mathematical paper remain after the core integration work.

\hypertarget{algebra_indetale_notes_20260929}{}
\subsection{Ind-étale maps, exact colimit relations and nil-ideal lifting}\label{algebra-proof-059-ind-uxe9tale-maps-exact-colimit-relations-and-nil-ideal-lifting}

\subsubsection{Source, comparison and coverage}\label{algebra-proof-059-source-comparison-and-coverage}

Source: the Stacks Project authors, algebra.tex at a04446e57ec1fbc252a871afcec7752fb2807b14, SHA-256 FA8BB92E58A4F78A2BD01B3B6A4A87DE0A0D279F5DD90641B574DD5FBFFFA4F3. The complete interval 42761--42958 and its seven received reports were read. The following opening 42959--43060 was also read, without adjudication. No prior canonical correction operation occurs in 42761--42958. Fresh dependency reading: 31885--31990, 39289--39326 and 39507--39538. The previously checked source descent result at 39309--39311 and its editorial treatment are reused, without claiming that its surrounding proof was newly read in this interval.

The four indexed hits from the saved query ``ind etale'' remain unread. The source-reading route retains the seven earlier external readings. Earlier complete proof dependencies are the sections ``Canonical infinitesimal map and nil ideals'' in ALGEBRA\_ETALE\_NOTES\_20260928.md, the diagonal construction in ALGEBRA\_UNRAMIFIED\_NOTES\_20260928.md, ``Uniqueness, colimits and the nilpotence boundary'' in ALGEBRA\_CONORMAL\_THICKENING\_NOTES\_20260928.md, and the complete/nil lifting and local-target sections of ALGEBRA\_HENSELIAN\_NOTES\_20260929.md. The new arguments below receive groups1127,1204,1214,1235,1236. There is no novelty claim. All expanded proofs and strengthened conclusions remain editorial material, separate from the translation.

An R-algebra is called ind-étale here precisely when it is a filtered colimit of étale R-algebras. Rings and maps are commutative and unital; the zero algebra is retained. All diagrams are small. Finite-presentation arguments below refer to actual chosen presentations and every relation.

\subsubsection{The source colimits and their exact comparison maps}\label{algebra-proof-059-the-source-colimits-and-their-exact-comparison-maps}

If E=R{[}x\_1,\ldots,x\_n{]}/(F\_1,\ldots,F\_m) and B is the filtered colimit of B\_i, a map E→B descends by choosing one stage for all n generator images and then a later stage where all m relation errors vanish. Two such maps which agree in B agree at a common later stage: each of the finitely many generator-image differences is eventually zero. For a filtered category rather than a directed set, use the finite-cocone and parallel-arrow equalization axioms at these same finite steps. This proves the exact bijection

\[
\mathop{\rm colim}_i\operatorname{Hom}_R(E,B_i)
\ \longrightarrow\ \operatorname{Hom}_R(E,\mathop{\rm colim}_iB_i).
\]

Keep A=colim\_i A\_i and an arbitrary original map R→R'. The base-change map sends r'⊗a\_i to r'⊗{[}a\_i{]}. Its inverse is the map R'⊗A→colim\_i (R'⊗A\_i) induced by r'↦{[}r'⊗1{]} and {[}a\_i{]}↦{[}1⊗a\_i{]}. Filtered-colimit equalities prove that the second prescription is independent of the representative. The balancing relation over R holds at every stage. The two maps are inverse on every elementary tensor. Each R'⊗A\_i is étale over R', proving 42770--42780 with all original objects retained. Delete only the duplicated words ``so is'' in its conclusion.

For composition R→A→B, write A=colim\_i A\_i with A\_i étale over R, and B=colim\_j B\_j with B\_j étale over A. Given a finitely presented R-algebra P and a map P→B, first descend it to B\_j. Étale descent 39309--39311 provides i and an étale A\_i-algebra B'\emph{j with the specified isomorphism A⊗}(A\_i)B'\emph{j→B\_j. The original map P→B\_j then descends to A}(i')⊗\_(A\_i)B'\emph{j for a later i'. This algebra is étale over R: base change makes it étale over A}(i'), and composition gives the R-map. The displayed factorization is through this exact algebra; it neither asserts B'\_j=B\_j nor changes the given map to B. The finite-presentation factorization criterion 31951--31980 proves B ind-étale over R. At 42807 replace the incomplete initial ``As'' by ``The map'' and its second causal ``as'' by ``since''.

For use in the omitted details, here is the factorization criterion with its maps. Suppose every map P→B from a finitely presented R-algebra factors through an étale R-algebra. Let C\_B have objects (E,u), where E is étale over R and u:E→B; maps are R-maps respecting u. It is essentially small, since finite presentations have a set of possible finite tuples of coefficients and equations. It is nonempty by the factorization of R→B. Two objects map to a third by applying the hypothesis to their tensor product. To equalize two maps E⇉F over B, choose finite algebra generators e\_1,\ldots,e\_n of E and put

\[
D=F/(f(e_1)-g(e_1),\ldots,f(e_n)-g(e_n)).
\]

This is finitely presented over R and maps to B; factor D→B through an étale algebra. This gives the required equalizing arrow from F. Thus C\_B is filtered. Its canonical colimit map to B is surjective: the map R{[}T{]}→B for each b factors through an object. It is injective: if e∈E maps to zero, the finitely presented E/(e) maps to B and factors through another object, making e zero there. The same argument applied to a difference compares any two representatives. This proves the criterion as an isomorphism, with the original receiving maps.

Now retain the source's iterated system A\_i=colim\_(j∈J\_i) A\_(i,j), and A=colim\_i A\_i. A finite presentation P→A first descends to one A\_i and then to an A\_(i,j). The criterion proves that A is ind-étale. It also verifies the source's explicitly given category J. For two objects (i,j),(i',j'), choose a common later i'\,' in I; finite presentation sends both algebras into a common stage of J\_(i'\,`). For parallel arrows from (i,j) to (i',j'), their composites to A\_(i') agree. Their finitely many generator-image differences become zero at a later stage of J\_(i'), which equalizes them. Thus J is filtered. The map colim\_J A\_(i,j)→A sends each element to its original image. Every element of A has such a representative. If two representatives become equal, first compare them at a common A\_k, then at a common stage of J\_k, enlarging that stage for the equality itself. They therefore agree in colim\_J. This proves the claimed isomorphism. The word sequence ``category with objects \ldots{} and morphisms triples'' has two parallel noun phrases; a missing copula is not a defect. The proposed ``and morphisms are triples'' would break that construction.

For a filtered system of actual ring arrows R\_i→A\_i, with limits R,A, the source comparison is

\[
\theta:\mathop{\rm colim}_i(A_i\otimes_{R_i}R)\longrightarrow A,
\qquad [a_i\otimes r]\longmapsto [a_i]r.
\]

Its inverse sends {[}a\_i{]} to {[}a\_i⊗1{]}. For any r∈R, choose its representative r\_j and a common stage k for i,j; the equality a\_i⊗r=(a\_i r\_k)⊗1 holds in A\_k⊗\_(R\_k)R. This proves that the inverse respects the R-action, and that the two compositions are the identity. Base change and the iterated-colimit result apply to the exact rings in θ, proving 42842--42856.

Finally keep the original R-map u:A→B between ind-étale R-algebras, with presentations A=colim\_i A\_i and B=colim\_j B\_j. Form the category T of triples (i,j,t), where t:A\_i→B\_j represents u on that stage. A morphism uses transitions i→i',j→j' and the equality t' a\_(ii')=b\_(jj')t. Each t is étale by 39507--39530. Thus

\[
D_{i,j,t}=A\otimes_{A_i,t}B_j
\]

is étale over A, and it has the specified map a⊗b↦a u\_j(b) to B, where the product uses u:A→B and u\_j:B\_j→B. The category T is filtered. For two triples, choose a common later i', factor A\_(i')→B through some B\_k, and then enlarge k to receive both original B\_j's. The two compatibility equations hold in B; finite generation of the two A\_i's makes them hold at a still later stage. For two parallel arrows, first equalize their underlying index arrows in the filtered indexing categories, and use the same finite equality argument for the resulting factorization maps. Equivalently, for the directed-set presentations used in the source there is at most one such transition between fixed triples.

The map colim\_T D\_(i,j,t)→B is surjective. Every b∈B comes from some B\_j; start with any triple and enlarge its B-index to receive that B\_j. For injectivity, an element of a D\_(i,j,t) is a finite sum Σ a\_l⊗b\_l. Choose a later A\_(i') representing every a\_l, and then a later triple (i',j',t') as above. In that stage this sum is the element 1⊗Σ t'(a\_l)b\_l. If its image in B is zero, enlarge j' once more to kill this exact sum. It is then zero in the displayed colimit. These are the full comparison maps behind 42858--42878; no claim of injectivity of any stage map is used.

\subsubsection{Arbitrary small colimits and the full idempotent relations}\label{algebra-proof-059-arbitrary-small-colimits-and-the-full-idempotent-relations}

The preceding source results admit the following stronger conclusion: the full subcategory of ind-étale R-algebras is closed under every small colimit calculated in R-algebras. Finite diagrams of étale R-algebras have étale colimits. We construct both the original colimit and the maps proving these statements.

First recall the exact diagonal calculation, including its determinant. For an étale R-algebra E with algebra generators z\_1,\ldots,z\_n, put D=E⊗\emph{R E, μ(a⊗b)=ab, and I=ker μ. Its actual generators are δ\_i=z\_i⊗1−1⊗z\_i. The differential identification I/I²=Ω}(E/R)=0 gives I=I². Choose c\_ij∈I with δ\_i=Σ\_j c\_ij δ\_j and retain

\[
e_E=\det(1-C)
 =\sum_{\sigma\in S_n}\operatorname{sgn}(\sigma)
   \prod_{i=1}^n(\delta_{i,\sigma(i)}-c_{i,\sigma(i)}).
\]

The adjugate identity gives e\_E δ\_i=0, hence e\_E I=0. Every nonconstant determinant summand contains an entry of I, so e\_E−1∈I. Thus e\_E²=e\_E and I=(1−e\_E)D. The empty determinant is 1. The idempotent is independent of choices: if e and e' have the same properties, e(1−e')=e'(1−e)=0, whence e=ee'=e'.

For any two R-maps f,g:E→Q, let

\[
\rho:E\otimes_RE\to Q,\quad a\otimes b\mapsto f(a)g(b),
\qquad d_{f,g}=\rho(e_E).
\]

It is an idempotent. The exact ideal generated by all f(x)−g(x) is (1−d\_(f,g))Q: the generators are precisely the images of the diagonal generators, and I=(1−e\_E)D gives both inclusions. Consequently Q/(1−d) and Q\_d represent equality of f and g after a further Q-map. Their inverse comparison maps are q mod(1−d)↦q/1 and q/d\^{}m↦q mod(1−d). The relation d=1 in each ring proves that these prescriptions respect denominators, including d=0 and d=1.

Let B\_• be the given small diagram of ind-étale R-algebras, indexed by a category K. Its actual coproduct in R-algebras is

\[
Q=\bigotimes_{k\in\operatorname{Ob}K,R} B_k.
\]

The empty tensor product is R. A finite tensor product of ind-étale algebras is ind-étale, by the already proved base-change and composition results. Q is the filtered colimit of those tensor products indexed by finite subsets of Ob K. The iterated-colimit result makes Q ind-étale. Write ι\_k:B\_k→Q for the original coproduct maps.

For each arrow v:k→l, choose the given filtered presentation B\_k=colim\_λ E\_(k,λ) by étale R-algebras. Compose each stage map with the two original maps ι\_k and ι\_l B(v), obtaining f\_(v,λ),g\_(v,λ):E\_(k,λ)→Q. Define d\_(v,λ) by the full determinant construction above. The ordinary colimit of B\_• is Q/J, where

\[
J=(\,\iota_k(x)-\iota_l(B(v)(x)):
                    v:k\to l,\ x\in B_k\,).
\]

Every x has a stage representative, and each stage image consists of such x's. The preceding equality of ideals therefore proves

\[
J=\sum_{(v,\lambda)}(1-d_{v,\lambda})Q.
\]

For a finite subset F of these pairs, retain the actual product d\_F=∏\emph{(h∈F)d\_h, with d}∅=1. The associated finite sum of ideals is exactly (1−d\_F)Q. In one direction use, for an ordering F=\{h\_1,\ldots,h\_s\}, the full identity

\[
1-\prod_{a=1}^s d_{h_a}
 =\sum_{a=1}^s\left(\prod_{b<a}d_{h_b}\right)(1-d_{h_a}).
\]

In the other direction, (1−d\_h)d\_F=0 gives 1−d\_h=(1−d\_h)(1−d\_F). Thus the exact quotient presentation is

\[
Q/J\ \cong\ \mathop{\rm colim}_{F\ {m finite}}
     Q/(1-d_F)Q
\ \cong\ \mathop{\rm colim}_{F\ {m finite}} Q_{d_F}.
\]

For F⊆G the transition sends q/d\_F\^{}m to q/1 in Q\_(d\_G), because d\_F d\_G=d\_G. The map to Q/J has the same formula with a quotient class. Conversely Q→colim\_F Q\_(d\_F) kills every generator of J: at the stage containing its pair, the associated d\_h is 1. Hence it factors through Q/J. The two induced maps are inverse on Q and on every displayed denominator. This proves the isomorphisms, not merely a correspondence of points or radicals.

Each Q→Q\_(d\_F) is étale, so composition makes Q\_(d\_F) ind-étale over R. Its filtered colimit is ind-étale. This proves the stated arbitrary-colimit closure. If K is finite and every B\_k is itself étale, take E\_(k,λ)=B\_k with its one stage. Q is then a finite tensor product of étale algebras; there are only finitely many arrows and associated d\_h. The quotient is the single localization Q\_d for their product d, and is étale over R. Empty diagrams give R. Changing the chosen stage presentations gives canonically inverse comparison maps, since both constructions have the same specified colimit maps from every original B\_k; their composites fix those maps and are therefore identities by the colimit universal property.

\subsubsection{Nil ideals and the exact reduction equivalence}\label{algebra-proof-059-nil-ideals-and-the-exact-reduction-equivalence}

Let A=colim\_i A\_i be ind-étale over R and let N be any nil ideal of an R-algebra T. Every map A→T/N lifts uniquely to T. This conclusion does not require finite presentation of A or nilpotence of N.

Here are the full stage and compatibility arguments. For the finite presentation A\_i=R{[}x\_1,\ldots,x\_n{]}/(F\_1,\ldots,F\_m), lift the given coordinate images to t\_1,\ldots,t\_n in T. The errors e\_j=F\_j(t) lie in N. Choose b\_j≥1 with e\_j\^{}(b\_j)=0, and put Q=(e\_1,\ldots,e\_m). Then

\[
Q^{1+\sum_j(b_j-1)}=0,
\]

since every monomial of that total degree contains e\_j to at least its specified exponent. If m=0 use Q=0. The coordinate images give A\_i→T/Q. Its unique lift through the successive square-zero maps T/Q\textsuperscript{(a+1)→T/Q}a exists by formal étaleness. It reduces to the prescribed map to T/N. For two possible lifts, the finitely many generator differences generate a nilpotent ideal by exactly the same exponent formula. They agree modulo that ideal; repeated square-zero uniqueness makes them equal. This proves stagewise existence and uniqueness.

For every transition A\_i→A\_j, composing the lift for j gives a lift of the map for i. Uniqueness identifies it with the lift already constructed for i. The lifts therefore form a compatible cocone and define an actual map A→T. A second such map agrees on every stage and is identical. This also proves naturality in T and in A by comparison on each stage.

This is a precise strengthening of groups1204 and1235 for ind-étale sources. The former's formally étale counterexample R→R/J, where J≠0 is a nil idempotent ideal, remains valid. If R/J were ind-étale over R, the identity into R/J would lift across R→R/J by the preceding result. An R-algebra lift s:R/J→R would satisfy s∘q=id\_R, contradicting q(J)=0 and J≠0. Thus that particular formally étale source is not ind-étale. This proves the exact exclusion; it does not infer that every other formally étale map fails to be ind-étale.

Now let N be a nil ideal of R. Reduction A↦A/NA defines an equivalence between ind-étale R-algebras and ind-étale R/N-algebras. It preserves their given colimit constructions. Base change proves that reduction lands in the latter class. For any ind-étale A and any R-algebra B, the ideal NB is nil: each element is a finite sum of multiples of nilpotent elements and the same multinomial exponent argument kills that finite sum. Thus the unique lifting result gives the exact bijection

\[
\operatorname{Hom}_R(A,B)
 \ \longrightarrow\
\operatorname{Hom}_{R/N}(A/NA,B/NB).
\]

For essential surjectivity, write a given ind-étale R/N-algebra C as colim\_i C\_i with C\_i étale. The equivalence for étale algebras in group1127 gives an étale R-algebra E\_i and a specified isomorphism E\_i/NE\_i→C\_i for every i. Full faithfulness of that equivalence lifts each transition uniquely. The lifted identities are identities and the lifted composites are composites by uniqueness, so these are an actual diagram, not separate unrelated lifts. Put E=colim\_i E\_i. Base change gives the specified isomorphism E/NE→C. The Hom bijection just proved makes any two such lifts uniquely isomorphic over their specified reductions.

Finally the ideal and module data are retained exactly. An ind-étale A is flat over R. Indeed, for an injection U→V of R-modules, an element in the kernel of U⊗A→V⊗A has a representative in U⊗A\_i. Its zero image becomes zero at a later V⊗A\_j, and flatness of A\_j forces that representative to be zero in U⊗A\_j. Therefore the original tensor map is injective. Tensoring 0→N→R→R/N→0 with A gives the full exact sequence

\[
0\longrightarrow N\otimes_R A
 \xrightarrow{\ n\otimes a\mapsto na\ } A
 \longrightarrow A/NA\longrightarrow0.
\]

The image is the actual ideal NA. This is the exact comparison between the retained ideal and its tensor product, with no suppression of a kernel or assumption that N has one exponent.

\subsubsection{Residue maps, local targets and filtered local rings}\label{algebra-proof-059-residue-maps-local-targets-and-filtered-local-rings}

In 42882--42899 contraction notation is read along the given maps, not as an injectivity assertion. Write u:R→S and v\_i:A\_i→A. The primes are p=u\^{}(-1)(m\_S) and q\_i=v\_i\^{}(-1)(q). These equalities are what the conventional intersections R∩m\_S and A\_i∩q mean. Both reports alleging an undefined subring intersection are rejected.

The residue maps in the proof are the actual embeddings κ(q\_i)→κ(q) followed by the given τ:κ(q)→κ(m\_S). For a fraction a/b with b outside q\_i, its image is v\_i(a)/v\_i(b); the denominator is nonzero because q\_i is the contraction. Apply the finite-stage henselian-target lemma to obtain f\_i:A\_i→S. For each transition, its composite with f\_j has the same contracted prime and the same residue-field map, so the finite-stage uniqueness proves compatibility. The induced f:A→S has f\^{}(-1)(m\_S)=q: on a representative a\_i membership on either side is membership in q\_i. Its residue-field map sends a/b to the quotient of the residues of f(a),f(b), and equals τ by the stage formula. A second eligible map agrees at every stage and is equal to f.~Replace the shorthand ``f mod q=τ'' by the explicitly typed condition that f induces τ on residue fields. The fraction extension is unique; the source theorem is not false because of that shorthand.

For the two henselian local algebras S,S' with the prescribed common residue field K, apply this lifting result to produce f:S→S' and g:S'→S inducing id\_K. Both composites and the respective identity maps induce id\_K; uniqueness makes gf=id\_S and fg=id\_(S'). This proves the exact isomorphism in 42901--42919.

The local-target consequence of group1236 extends to every ind-étale algebra. Let R be strictly henselian and R→S any local map to a local ring. Then for every ind-étale R-algebra A the map

\[
\operatorname{Hom}_R(A,R)\longrightarrow\operatorname{Hom}_R(A,S)
\]

is bijective. Write A=colim\_i A\_i with A\_i étale. For each i, the finite-stage argument decomposes A\_i into copies of R and a factor of empty closed fibre. A map to S selects exactly one copy; the empty-fibre factor cannot map to κ(S), because its base change to that nonzero field is zero. Thus its map to S comes from a unique map to R. For a given f:A→S, apply this to each f\_i, obtaining g\_i:A\_i→R. For a transition, g\_j composed with it and g\_i have the same composite to S; the finite-stage bijection makes them equal. Their colimit map g satisfies u∘g=f.~Uniqueness follows on the same stages. No nonemptiness of these Hom sets is assumed. R=k with k algebraically closed of characteristic zero and S=k{[}t{]}\_(t) again gives a nonhenselian target: a square in k(t) has even order at t+1, so the simple residue root of X²−(1+t) cannot lift in S.

For completeness, the source's last local-colimit assertion keeps its original local transitions. For a filtered diagram of nonzero local rings R\_i, the residue-field transitions are embeddings. Their filtered colimit κ is a field and every κ\_i→κ is injective: an element that becomes zero does so at a stage, where a field embedding is injective. The map R=colim R\_i→κ is surjective, and its kernel is the ideal represented by the m\_i. If a class x∈R has nonzero residue, its representative in some R\_i has nonzero residue and hence is a unit there; its inverse maps to an inverse in R. Elements of the kernel cannot be units since κ is a nonzero field. Thus R is local with exactly this maximal ideal and residue field, even when the ring transitions are not injective.

If every R\_i is henselian, take the original monic polynomial f over R and simple residue root a\_0∈κ. Choose one stage for all its finitely many coefficients, keeping the leading coefficient equal to 1, and for a\_0. The root and nonzero derivative conditions hold at that common stage because its residue field embeds in κ. Henselianity supplies the required root there, and its image is the required root of the original f.~If every κ\_i is separably closed, a monic separable polynomial over κ descends together with the finitely many coefficients of a Bezout identity for it and its derivative. At that common stage it is separable, hence splits into linear factors in κ\_i. Passing those exact factors to κ shows that κ is separably closed. This proves the strictly henselian assertion as well.

\subsubsection{Diagram of the proved maps}\label{algebra-proof-059-diagram-of-the-proved-maps}

\begin{figure}
\centering
\pandocbounded{\includegraphics[keepaspectratio,alt={Exact colimit quotients and the nil-ideal sequence}]{INDETALE_COLIMIT_MAPS_20260929.png}}
\caption{Exact colimit quotients and the nil-ideal sequence}
\end{figure}

The top row is the actual coproduct Q, the finite idempotent localizations and their colimit Q/J from the arbitrary-colimit proof. The isomorphism in the middle uses the explicit mutually inverse maps proved there. The bottom row is the exact flat sequence for an ind-étale A and a nil ideal N of R. Both rows retain their original quotient and ideal data. Full proof locators are the two preceding colimit and nil-ideal sections; their human-source comparisons are with the Stacks Project authors at the authority locators listed above. Reproducible source: draw\_indetale\_colimit\_20260929.py; vector edition: INDETALE\_COLIMIT\_MAPS\_20260929.svg. The PNG was inspected for legibility, all arrow directions and mathematical labels.

\subsubsection{Findings and propagation}\label{algebra-proof-059-findings-and-propagation}

Seven reports give two copyedits, one typed residue-map clarification, and two rejected groups: the parallel noun phrases and the two contraction notations are not source defects. The source proofs of filtered closure and henselian lifting remain identifiable.

Three separate consequence groups receive earlier proofs. The first uses group1214's exact diagonal idempotent to prove all-small- colimit closure with the original tensor product, full difference ideal and every finite idempotent product. The second receives groups1127,1204,1235 and proves arbitrary nil-ideal lifting and reduction equivalence for ind-étale algebras, with the exact flat ideal sequence and the non-ind-étale counterexample boundary. The third receives group1236 and proves the Hom-set bijection for arbitrary local targets and all ind-étale sources. These are not substitutions for the original translated arguments. Broader combined-result synthesis and any paper remain after core work.

\hypertarget{algebra_henselization_notes_20260929}{}
\subsection{Henselization: source review and exact comparison maps}\label{algebra-proof-060-henselization-source-review-and-exact-comparison-maps}

Primary source: Stacks Project authors, \texttt{algebra.tex}, original authority commit \texttt{a04446e57ec1fbc252a871afcec7752fb2807b14}, lines 42959--43572, file SHA256 \texttt{FA8BB92E58A4F78A2BD01B3B6A4A87DE0A0D279F5DD90641B574DD5FBFFFA4F3}. This is a separate editorial validation of the original statements and their omitted details. It does not replace the source exposition or claim a new theorem. The source reading includes the next opening, 43573--43593; that opening remains unadjudicated.

The user's latest direction keeps the Algebra review first. Optional automorphism, invariant-ring and more general residue-extension investigations are deferred, along with additional paper mining and the later synthesis. Their consideration creates no proved-consequence group here. The eight minimal edits below account for fifteen received reports in nine groups.

The previously proved inputs used below are the finite-presentation descent and étale permanence statements, the henselian finite-algebra decomposition and equivalence of finite étale categories with residue-field finite étale categories, and the exact ind-étale comparison and henselian-target lifting statements. Their full arguments are in \texttt{ALGEBRA\_ETALE\_NOTES\_20260928.md}, \texttt{ALGEBRA\_HENSELIAN\_NOTES\_20260929.md} and \texttt{ALGEBRA\_INDETALE\_NOTES\_20260929.md}; the evidence receipt binds their bytes. These are mathematical inputs, not assumptions added to the source's hypotheses.

\subsubsection{Pointed systems and the original coequalizer}\label{algebra-proof-060-pointed-systems-and-the-original-coequalizer}

Let \((R,\mathfrak m,\kappa)\) be the original local ring. An ordinary object is precisely the source pair \((S,\mathfrak q)\), with \(S/R\) étale and its structural residue map \(\kappa\to\kappa(\mathfrak q)\) an isomorphism. Every arrow is an \(R\)-map contracting the selected target prime to the selected source prime. In the strict system a fixed inclusion \(\kappa\subset\kappa^{sep}\) is retained, and the original triple carries an embedding \(\alpha:\kappa(\mathfrak q)\to\kappa^{sep}\). An arrow must preserve this embedding. This condition is essential data even when the underlying arrow is also an arrow of pairs.

The systems are nonempty, containing \(R\) with its selected prime and residue map. For two ordinary objects the source ring is exactly \(S''=S\otimes_R S'\), with \(\mathfrak q''=\mathfrak qS''+\mathfrak q'S''\). Since \(\mathfrak m\) is maximal and \(S\) is étale, its closed fibre is a finite product of finite separable fields. In the ordinary case the selected factor is \(\kappa\); thus \(S/\mathfrak q=\kappa=S'/\mathfrak q'\) and \(S''/\mathfrak q''=\kappa\). Both structure maps contract \(\mathfrak q''\) as required.

For two strict objects use the same original tensor product, and the residue map \[
 S\otimes_R S'\longrightarrow\kappa^{sep},\qquad
 x\otimes y\longmapsto\alpha(\overline x)\alpha'(\overline y).
\] Its kernel is the selected prime. Its residue field is the compositum of the two embedded finite separable fields: the image contains both fields and is their finite-dimensional tensor image, a domain finite-dimensional over \(\kappa\), hence a field. This proves the residue assertion and both required contractions. Étaleness is preserved by this tensor product.

For parallel arrows \(\varphi,\psi:S\to S'\), keep the source's full ring, rather than replacing it in the construction: \[
 U=(S'\otimes_{\varphi,S,\psi}S')\otimes_{S'\otimes_R S'}S'.
\] The map from \(S'\otimes_R S'\) to the first factor is the canonical quotient map, and the map to the last factor is multiplication. Let \(J\subset S'\) be the ideal generated by every \(\varphi(s)-\psi(s)\), for \(s\in S\). The exact comparison maps are \[
 U\longrightarrow S'/J,\quad (a\otimes b)\otimes c\longmapsto[abc],
 \qquad
 S'/J\longrightarrow U,\quad[c]\longmapsto(1\otimes1)\otimes c.
\] The first map respects the \(S\)-balancing because the difference of the two images is \([ab(\varphi(s)-\psi(s))c]=0\). It respects the \(S'\otimes_R S'\)-balancing because multiplication of the two factors has the same image \([abc]\). In \(U\), balancing gives \((a\otimes b)\otimes c=(1\otimes1)\otimes abc\), and the original relation \(\varphi(s)\otimes1=1\otimes\psi(s)\) therefore kills \(J\) under the second map. The two displayed maps compose to the identity on every tensor generator and every quotient class. In particular the original map \(S'\to U\) equalizes \(\varphi\) and \(\psi\).

All étaleness claims here follow on the original factors. The arrows \(S\to S'\) are maps between étale \(R\)-algebras, hence étale. The first tensor factor is étale over \(S'\) by base change and over \(R\) by composition. Multiplication \(S'\otimes_R S'\to S'\) is also a map between étale \(R\)-algebras, hence étale; its base change to the first factor makes \(U\) étale over that factor and over \(R\).

Write \(\rho':S'\to\kappa\), or \(S'\to\kappa^{sep}\), for the selected residue map. The pair/triple conditions give \(\rho'\varphi=\rho'\psi\). Thus the exact map \[
 U\longrightarrow\kappa\text{ or }\kappa^{sep},\qquad
 (a\otimes b)\otimes c\longmapsto\rho'(a)\rho'(b)\rho'(c)
\] is defined. Its kernel contracts to \(\mathfrak q'\) in the last factor. In the ordinary local construction its kernel is exactly the source ideal \[
 (\mathfrak q'\otimes S'+S'\otimes\mathfrak q')\otimes S'
 +(S'\otimes_{\varphi,S,\psi}S')\otimes\mathfrak q'.
\] Indeed the comparison with \(S'/J\) sends this ideal to \(\mathfrak q'/J\), since \(J\subset\mathfrak q'\), and the quotient is \(S'/\mathfrak q'=\kappa\). The strict construction has the same contraction and the inherited embedding of \(\kappa(\mathfrak q')\). This proves filteredness, including the parallel-arrow condition, for the original systems.

For a general base prime \(\mathfrak p\subset R\), use the source fibre rings \(F=S\otimes_R\kappa(\mathfrak p)\) and \(F'=S'\otimes_R\kappa(\mathfrak p)\). They are finite products of finite separable fields. Their selected factor maps define precisely the same product and coequalizer residue maps. The original source formula is \[
 F''=(F'\otimes_{\varphi,F,\psi}F')
       \otimes_{F'\otimes_{\kappa(\mathfrak p)}F'}F'.
\] The selected map \(F'\to\kappa(\mathfrak p)\) belongs to \(\mathfrak q'\), not to an undefined \(\mathfrak p'\). Its extension sends \((a\otimes b)\otimes c\) to the product of the three selected images. This proves the correction at 43284 on the exact original ring. In the strict case its target is \(\kappa^{sep}\), and the triple compatibility gives the same balancing identity. The selected prime of the original ring is the inverse image of the fibre prime. One must not assume that the unlocalized map \(S\to\kappa(\mathfrak q)\) is surjective for a general base prime; its image has fraction field \(\kappa(\mathfrak q)\).

\subsubsection{Local colimits and the source henselian property}\label{algebra-proof-060-local-colimits-and-the-source-henselian-property}

For any selected \(\mathfrak q\) and any \(f\in S\setminus\mathfrak q\), the original system includes \(S_f\) with its selected prime and, in the strict case, the unchanged residue embedding. Hence every such \(f\) becomes invertible in the colimit \(H\). The canonical maps give inverse comparisons \[
 H=\mathop{\rm colim}S\longrightarrow\mathop{\rm colim}S_{\mathfrak q},
 \qquad s/f\longmapsto [s][f]^{-1}\in H.
\] The latter formula respects every localization relation and transition map; the composites fix \([s]\) and every \(s/f\). This proves the two colimit formulas with the original objects retained. In the localized system every transition map is local. The colimit of their residue fields is \(\kappa\) in the ordinary case and \(\kappa^{sep}\) in the strict case. In the latter assertion each element of \(\kappa^{sep}\) lies in a finite separable extension of \(\kappa\), and that extension occurs as a selected étale residue field by the previously proved prescribed-residue construction. Common tensor objects put every finite list in a single stage.

At each localized stage, étaleness gives \(\mathfrak pR_{\mathfrak p}S_{\mathfrak q}=\mathfrak qS_{\mathfrak q}\). The kernel of the colimit residue map is exactly \(\mathfrak pH\): an element with zero residue has zero residue already at its representing stage, since the residue transition maps are field embeddings, and then belongs to the displayed stage ideal. Conversely that ideal maps to zero. An element outside this kernel is represented by a unit in a localized stage. Thus \(H\) is local, its maximal ideal is precisely \(\mathfrak pH\), and its specified residue field is the indicated field colimit. This also proves that the zero ring cannot occur as a selected stage or final local ring.

For the original local-base argument at 43040--43056, all primes of the closed fibre are maximal, not only the selected prime. Distinct selected and unselected primes are therefore incomparable. Choose \(g_i\in\mathfrak q_i\setminus\mathfrak q\) and keep the full product \(g=\prod_{i=1}^r g_i\); for an empty list take \(g=1\). The localization at \(g\) removes every other closed-fibre factor. That fibre is now its single selected field, so \(\mathfrak mS_g=\mathfrak qS_g\), not merely an equality of radicals. This justifies the source's stronger ideal equality. The source inverse of \([f]\) is exactly the class of \(1/f\) in \(S_f\).

For the henselian property, let the original monic polynomial \(P\in H[T]\) and its specified simple residue root \(a_0\) be given. Filteredness places all coefficients at one stage, with a monic lift \(Q\). In the ordinary local construction \(a_0\in\kappa=S/\mathfrak q\), so choose its original lift \(a'\in S\). In the strict construction first enlarge to a common pointed stage containing the finite separable residue extension generated by \(a_0\), together with the coefficient stage. Since this stage lies over the maximal ideal of the local base, its selected prime is maximal and its quotient is its residue field; therefore \(a_0\) also has a lift \(a'\) there. This is the additional residue step implicit in the source's reference to the same proof.

Keep exactly \(S'=S[T]/(Q)\) and \(\mathfrak q'=\mathfrak qS'+(T-a')S'\). The quotient is the selected residue field. The derivative \(Q'(T)\) is nonzero there, since \(a_0\) is a simple root. Localizing at its class gives an étale \(S\)-algebra and preserves \(\mathfrak q'\); equivalently use the source's chosen \(g\notin\mathfrak q'\) in the étale locus. The class of \(T\) in that object maps to an element \(a\in H\) satisfying the exact equations \(P(a)=0\) and \(\overline a=a_0\). This proves the henselian property by the simple-root criterion already established in the preceding section. In the strict case the residue field is separably closed, so \(H\) is strictly henselian.

For an arbitrary \(R,\mathfrak p\), the source's descent from \(R_{\mathfrak p}\) uses the proved finite-presentation descent for étale algebras: each étale algebra over \(R_{\mathfrak p}\) descends to one over \(R_f\) for some \(f\notin\mathfrak p\). The composite from \(R\) is étale; contraction of the selected prime gives an object over \((R,\mathfrak p)\) with unchanged residue field. Conversely the base localization of every such object is an object over the local base. Every finite set of algebra generators, arrow images, and relations of two arrows descends after a further denominator \(g\notin\mathfrak p\), by finite presentation. Thus not only the objects but the equalities used in the colimit descend. The inverse comparison sends a representative at a descended object to its original colimit class; independence of the descent follows at a common further localization. This proves both bijectivity and ring compatibility of the source's canonical ordinary and strict comparisons. No injectivity of transition maps is assumed in these comparisons.

\subsubsection{Residue choices and functoriality}\label{algebra-proof-060-residue-choices-and-functoriality}

For a henselian local target \(B\), an étale \(R\)-algebra \(A\), a selected prime over the contracted maximal ideal, and a specified compatible embedding of its residue field into \(\kappa(B)\), the previously proved henselian lifting lemma gives one and only one map \(A\to B\) inducing these data. On a filtered system the maps are compatible: for each transition arrow the two candidate composites have identical base and residue maps, so stagewise uniqueness identifies them. The resulting colimit map is unique and local, because its reduction is the specified field embedding and the source maximal ideal is the kernel of its specified reduction. This proves all the functoriality assertions in 43186--43224 and 43337--43397 on the original systems.

In particular, the map in 43353 has target \(S^{sh}\). Its correctly typed contraction is \[
 f^{-1}(\mathfrak m_{S^{sh}})=\mathfrak q.
\] The source's \(\mathfrak m_{S^h}\) at 43360 belongs to a different ring, so the proposed replacement is a mathematical target correction. The residue condition still uses precisely the original chosen map \(\phi:\kappa(\mathfrak q)\to\kappa^{sep}\); no map has been discarded.

The unique-isomorphism assertions at 43143--43148 and 43157--43162 are understood in their specified residue-field data. Given two candidate rings with those data, the lifting lemma constructs maps in both directions, and uniqueness makes both composites identities. For the strict construction the given identification with \(\kappa^{sep}\) is part of that comparison, just as the arrows in the source explicitly preserve residue maps. This is a contextual reading of the source, not a new assertion that arbitrary residue automorphisms must be trivial. Investigation of the larger automorphism and fixed-ring structures remains deferred.

For the ordinary base-change statement, put \(H=(R_{\mathfrak p})^h\), \(K=(S_{\mathfrak q})^h\) and retain \(T=H\otimes_R S\). Its original comparison map is \[
 T\longrightarrow K,\qquad h\otimes s\longmapsto f(h)s.
\] The prime \(\mathfrak t\) is the kernel of its composite to \(\kappa(\mathfrak q)\), given by \(h\otimes s\mapsto\overline h\,\overline s\) using the structural residue extension. It is the unique prime contracting to \(\mathfrak m_H\) and \(\mathfrak q\): primes with those contractions correspond to primes of \[
 \kappa(\mathfrak p)\otimes_{\kappa(\mathfrak p)}\kappa(\mathfrak q)
 =\kappa(\mathfrak q),
\] which has just its zero prime. This correspondence is obtained by quotienting the two maximal/prime ideals and inverting all nonzero elements in both residue domains. The image of the unlocalized ring \(T\) need not be all of this field, but its fraction field is \(\kappa(\mathfrak q)\), because the image already contains \(S/\mathfrak q\). Hence this is exactly \(\kappa(\mathfrak t)\).

Both \(T\) and \(K\) are ind-étale over \(S\), so the original comparison makes \(K\) ind-étale over \(T\), by the earlier maps-between-ind-étale-algebras result. Every element outside \(\mathfrak t\) maps to a unit of \(K\); base change along \(T\to T_{\mathfrak t}\) consequently retains \(K\) itself. Thus \(K/T_{\mathfrak t}\) is ind-étale and local with equal residue field, and \(K\) is henselian. The henselian uniqueness lemma identifies it with the henselization of \(T_{\mathfrak t}\). The two inverse maps are the unique local maps inducing the identity on the common residue field; their composites are identities by that same uniqueness. This proves the source's identification, rather than only an abstract existence assertion.

For strict base change retain \(H^{sh}=(R_{\mathfrak p})^{sh}\), \(K^{sh}=(S_{\mathfrak q})^{sh}\), \(T^{sh}=H^{sh}\otimes_R S\), and the given \(\phi:\kappa_1^{sep}\to\kappa_2^{sep}\). The selected prime is precisely \[
 \mathfrak t_\phi=\ker\bigl(T^{sh}\longrightarrow\kappa_2^{sep},
            \ h\otimes s\longmapsto\phi(\overline h)\overline s\bigr).
\] Its residue field is the compositum \(L=\kappa(\mathfrak q)\,\phi(\kappa_1^{sep})\) inside \(\kappa_2^{sep}\). Indeed the fraction field of the image contains both indicated fields and is generated by them. The extension \(L/\kappa(\mathfrak q)\) is algebraic separable: each finite list of elements from \(\phi(\kappa_1^{sep})\) lies in a finite separable extension of \(\kappa(\mathfrak p)\), and after base change the selected field factor is finite separable over \(\kappa(\mathfrak q)\). Every element of \(\kappa_2^{sep}\) is algebraic separable over \(L\), because its separable minimal polynomial over \(\kappa(\mathfrak q)\) remains square-free and its minimal polynomial over \(L\) divides it. Since \(\kappa_2^{sep}\) is separably closed, it is therefore a separable closure of \(L\). The map to \(K^{sh}\) is ind-étale, and remains so after localization at \(\mathfrak t_\phi\), by the same exact base-change argument. The target is strictly henselian with the specified residue closure. Uniqueness with that residue embedding gives the source's strict-henselization identification and the inverse map. The source correctly says a prime, not a unique prime independent of \(\phi\).

\subsubsection{The finite étale tower and the original tensor formula}\label{algebra-proof-060-the-finite-uxe9tale-tower-and-the-original-tensor-formula}

Let \(A=(R_{\mathfrak p})^h\). For each finite separable subextension \(\kappa'\subset\kappa^{sep}\), the finite étale equivalence for the henselian ring \(A\) gives the local finite étale algebra \(A(\kappa')\), together with its specified residue identification. Every inclusion \(\kappa'\subset\kappa''\) lifts to a unique map \(A(\kappa')\to A(\kappa'')\). Uniqueness gives the identity and composition laws exactly. This map is finite because the target is finite over \(A\), and étale because it is a map between étale \(A\)-algebras. It is local: the maximal ideals are \(\mathfrak m_A A(\kappa')\) and \(\mathfrak m_A A(\kappa'')\), and the residue map is an injection. A flat local map is faithfully flat, hence injective. Consequently the source's union denotes an actual system of injective maps, not an unjustified injectivity assumption about the earlier arbitrary stages.

For completeness the faithful-flat injectivity used here follows on the original module map. Flatness sends \(I=\ker(A_1\to A_2)\subset A_1\) to an injection \(I\otimes_{A_1}A_2\to A_2\), whose image is zero; faithful flatness gives \(I=0\). Flatness is faithful for a local map because every proper ideal \(J\subset A_1\) is contained in its maximal ideal, and \(JA_2\) is contained in the target maximal ideal, so \(A_2/JA_2\ne0\). The established flatness criterion by quotients then gives faithful flatness.

Each tower stage is henselian, by the finite-over-henselian theorem. Their filtered colimit with local maps is local with residue \(\kappa^{sep}\) and is henselian. To verify the last assertion directly, put the coefficients of a monic polynomial and its selected residue root in one stage. The residue embeddings ensure that the derivative is already nonzero at that stage when it is nonzero in the limit. The stage's henselian property gives a root there with that residue, and its image is the required root in the original colimit. The colimit is ind-étale over \(A\), and by composition over \(R\). The specified residue comparison and uniqueness identify the tower colimit with the original \(R^{sh}\).

This also supplies all the omitted details at 43488--43505 without assuming that inclusion of an indexing subsystem implies injectivity on its colimit. Let \(C\) be the original triple colimit in that passage. The preceding colimit proof makes \(C\) local, with maximal ideal \(\mathfrak pC\) and residue \(\kappa^{sep}\). The ordinary subsystem gives a map \(A\to C\) compatible with \(\kappa\subset\kappa^{sep}\), hence local. Both \(A\) and \(C\) are ind-étale over \(R\); therefore \(C\) is ind-étale over \(A\). This map is flat and local, hence faithfully flat and injective by the preceding proof. This, rather than the subsystem alone, proves the source's containment.

To recover the finite stages needed for the source's henselian conclusion, take any étale \(A\)-stage \(E\to C\). The proved henselian decomposition is \(E=E_1\times\cdots\times E_n\times B\), with the \(E_i\) finite local étale over \(A\) and \(\mathfrak m_A B=B\). A map to the nonzero local ring \(C\) sends exactly one of the finitely many product idempotents to 1 and the rest to 0. It cannot choose \(B\), because a finite expression \(1=\sum a_jb_j\), with \(a_j\in\mathfrak m_A\), would put 1 in \(\mathfrak m_C\). Thus it factors through a selected \(E_i\). The factor map to \(C\) is local: its contracted prime lies over \(\mathfrak m_A\), and a finite local algebra has only its maximal ideal over that ideal. For any two stages or any equality of elements, take a common stage in the original filtered presentation and then its selected finite factor. The transition maps factor through the previously selected factors because the product idempotents have the same 0 or 1 images in that local factor. This proves that the selected finite local stages are filtered and have colimit exactly \(C\), on both elements and equalities. The same henselian colimit proof now applies. The finite étale equivalence identifies these stages and their exact residue embeddings with the stated tower.

Finally keep the exact objects in 43540--43558. Write \(A=(R_{\mathfrak p})^h\), \(B=(S_{\mathfrak q})^h\), and assume the original residue map \(\kappa(\mathfrak p)\to\kappa(\mathfrak q)\) is an isomorphism. Use that given map to identify the residue fields, as the source does, and use its chosen \(\kappa^{sep}\). The source tensor product is \[
 T=B\otimes_A(R_{\mathfrak p})^{sh}
   =\mathop{\rm colim}_{\kappa'}\bigl(B\otimes_A A(\kappa')\bigr).
\] The forward comparison sends \(b\otimes[a]\) to \([b\otimes a]\); its inverse sends each stage tensor to the same tensor in \(T\). Balancing and the transition equations make these inverse maps on every generator.

Each original stage \(B\otimes_A A(\kappa')\) is finite étale over \(B\), with closed fibre exactly \(\kappa(\mathfrak q)\otimes_{\kappa(\mathfrak p)}\kappa'=\kappa'\) by the original residue isomorphism. The finite-algebra henselian decomposition over \(B\) therefore has exactly one nonzero local factor: every maximal ideal of a finite algebra over a local ring contracts to the maximal ideal, and the closed fibre here is a field. The stage is local and henselian. Its maximal ideal is \(\mathfrak m_B\) times the stage, and the transition maps are local, with the specified inclusions of residue fields. Hence \(T\) is a henselian local ind-étale \(B\)-algebra with residue \(\kappa^{sep}\); it is strictly henselian. It is also ind-étale over \(S_{\mathfrak q}\) by composition. The original tensor comparison \[
 T\longrightarrow(S_{\mathfrak q})^{sh},\qquad
 b\otimes a\longmapsto b\,f(a)
\] is local and induces the identity on the chosen residue closure. The unique local map in the reverse direction with the same residue identification exists by henselian lifting on its original ind-étale stages. Both composites are identities by uniqueness. This proves precisely the original formula, with its original residue-isomorphism hypothesis. No stronger residue-extension claim is recorded in this batch.

\subsubsection{Report dispositions and the deferred follow-on work}\label{algebra-proof-060-report-dispositions-and-the-deferred-follow-on-work}

The fifteen reports give nine groups and eight source operations:

\begin{enumerate}
\def\labelenumi{\arabic{enumi}.}
\tightlist
\item
  \texttt{OCC-00968}, \texttt{OCC-12274}, 43084--43089: reject the asserted defect. In ``Now that triple \ldots{} determines'', ``now'' is an adverb and ``that triple'' is its demonstrative subject. The text does not need a subordinate ``now that'' clause or a compulsory replacement of ``that'' by ``the''. The class of \(T\) proves the stated mathematical conclusion as above.
\item
  \texttt{OCC-00969}: add the missing terminal period at 43139.
\item
  \texttt{OCC-00970}, \texttt{OCC-12278}: at 43284 replace the undefined \(\mathfrak p'\) by the actual \(\mathfrak q'\). The original fibre and coequalizer maps above verify the referent.
\item
  \texttt{OCC-00971}, \texttt{OCC-12282}: at 43360 replace \(\mathfrak m_{S^h}\) by \(\mathfrak m_{S^{sh}}\), the maximal ideal of the actual codomain.
\item
  \texttt{OCC-00972}: join the given-diagram clause to its main clause at 43381--43383, inserting the comma and changing ``There'' to ``there''.
\item
  \texttt{OCC-00973}, \texttt{OCC-12283}: at 43422 refer to the category of triples. All embedding data remain as originally stated.
\item
  \texttt{OCC-00974}, \texttt{OCC-12284}: at 43441 remove the stray article and add the participial comma: ``Combining the maps above, we''.
\item
  \texttt{OCC-00975}, \texttt{OCC-12285}: at 43461 call the arrows morphisms of triples. They also are morphisms of the underlying pairs; this terminological edit does not assert that the original mathematical arrow was invalid.
\item
  \texttt{OCC-01463}: join the given-diagram clause at 43522--43524, inserting the comma and changing ``The'' to ``the''. The report's claim that there is a post-display period is inaccurate; the original has no such period. The accepted edit repairs the actual capitalization and clause joining.
\end{enumerate}

The earlier ind-étale lifting, maps-between-algebras and henselian results have been checked on these receiving source constructions. Their original source hypotheses suffice. Full reconstructions stay in this note; only the eight minimal operations enter the proposed correction layer. The optional wider consequences and further external-source reading are queued after completion of the Algebra review. No formal proof-assistant certification, chapter admission, TeX build, publication or novelty claim is made by this note.

\hypertarget{algebra_quasifinite_henselization_notes_20260929}{}
\subsection{Quasi-finite henselization: verification of the received corrections}\label{algebra-proof-061-quasi-finite-henselization-verification-of-the-received-corrections}

Primary source: Stacks Project authors, original commit \texttt{a04446e57ec1fbc252a871afcec7752fb2807b14}, \texttt{algebra.tex:43573–43864}, SHA256 \texttt{FA8BB92E58A4F78A2BD01B3B6A4A87DE0A0D279F5DD90641B574DD5FBFFFA4F3}. Additional original-source reading: \texttt{algebra.tex:30292–30339}, \texttt{42384–42447}, \texttt{42842–42878}, and \texttt{fields.tex:1703–1724} at that same commit. Their bytes are bound in the evidence receipt. The preceding henselization, henselian and ind-étale validation notes provide the already established inputs.

This batch verifies two existing canonical corrections, accounting for four reports. It proposes no duplicate source operation and introduces no optional extension. All calculations below verify statements actually made in the original source. Additional paper mining and wider consequences remain deferred under the user's Algebra-first direction.

\subsubsection{The original primes and finite local comparisons}\label{algebra-proof-061-the-original-primes-and-finite-local-comparisons}

Retain the map \(R\to S\), the original primes \(\mathfrak q/\mathfrak p\), and the hypothesis that the map is quasi-finite at \(\mathfrak q\). Put \(H=R_{\mathfrak p}^h\) and \[
 U=H\otimes_R S,\qquad
 T=H\otimes_{R_{\mathfrak p}}S_{\mathfrak q}.
\] The second ring is the localization of the first at all \(1\otimes s\), \(s\notin\mathfrak q\). The forward map sends \((h\otimes s)/(1\otimes v)\) to \(h\otimes(s/v)\); the reverse map has that same formula read backwards. The \(R_{\mathfrak p}\)-balancing is valid because the images in \(H\) of all \(r\notin\mathfrak p\) are already units. The maps preserve products and sums and fix every tensor and denominator generator, so they are inverse.

The ideal \(\mathfrak qT\) is prime and maximal, since \[
 T/\mathfrak qT
 =H\otimes_{R_{\mathfrak p}}\kappa(\mathfrak q)
 =(H/\mathfrak pH)\otimes_{\kappa(\mathfrak p)}\kappa(\mathfrak q)
 =\kappa(\mathfrak q).
\] Here \(\mathfrak pH=\mathfrak m_H\), with the original specified identification \(H/\mathfrak m_H=\kappa(\mathfrak p)\). Its inverse image \(\mathfrak Q\) in \(U\) is the prime selected by these residue maps. Quasi-finiteness persists at \(\mathfrak Q\) by the four-rings lemma, whose exact proof retains the source fibre's finite factor and its complementary factor after base change. The finite-over-henselian theorem then makes \(U_{\mathfrak Q}\) a finite henselian local \(H\)-algebra. The preceding explicit localization identifies it with \(T_{\mathfrak qT}\). Its maximal ideal contracts to \(\mathfrak q\), and its residue field is exactly \(\kappa(\mathfrak q)\).

The earlier ordinary henselization base-change theorem identifies \(S_{\mathfrak q}^h\) with the henselization of \(T_{\mathfrak qT}\). A henselian local ring is its own henselization: the unique local map from its henselization back to the ring has identity residue map; its two composites with the structure map are identities by uniqueness. Thus this is the original finite local algebra \(U_{\mathfrak Q}\), proving 43580--43614. For an ideal \(I\subset\mathfrak m_R\), the quotient map \(R\to R/I\) is finite, the displayed tensor is exactly \(R^h/IR^h\), and this quotient is already local. Its selected localization changes nothing. This proves the source's quotient compatibility, with no finite-generation restriction on \(I\).

For the strict statement preserve its entire field diagram. Write \(k=\kappa(\mathfrak p)\), \(K=\kappa(\mathfrak q)\), \(\Omega=\kappa_2^{sep}\), and let \(L=\kappa_1^{sep}\) be precisely the elements of \(\Omega\) separable algebraic over \(k\). Quasi-finiteness implies that \(K/k\) is finite. The source's separable-first lemma, \texttt{fields.tex:1703–1724}, shows that \(L\) is a field and \(\Omega/L\) is purely inseparable. Indeed in characteristic \(p>0\), the irreducible polynomial of each algebraic \(x\) has the exact form \(Q(X^{p^e})\), with \(Q\) irreducible and separable, so \(x^{p^e}\) is separable over \(k\). In characteristic zero every algebraic extension is separable.

The field \(L\) is separably closed. Let a separable polynomial over \(L\) be given. Its finitely many coefficients lie in a finite separable subextension \(E/k\) of \(L/k\). Square-freeness remains valid over \(E\): a common nonconstant factor with the derivative over \(E\) would remain such a factor over \(L\). Since \(\Omega\) is separably closed, all roots lie in \(\Omega\); each is separable over \(E\), hence over \(k\), and therefore lies in \(L\). Thus \(L/k\) is the claimed separable closure.

Set \(M=LK\subset\Omega\), keeping the two original embedded fields. It is finite purely inseparable over \(L\). In positive characteristic choose one exponent \(p^e\) working for a finite set of generators of \(M/L\); then \(M^{p^e}\subset L\). To prove that \(M\) is separably closed, let \(x\) be separable algebraic over \(M\). A separable polynomial \(\sum a_iX^i\) for \(x\), after the Frobenius field isomorphism \(M\to M^{p^e}\), becomes the separable polynomial \(\sum a_i^{p^e}X^i\) for \(x^{p^e}\). Its coefficients lie in \(L\), and it stays separable after that field extension. Consequently \(x^{p^e}\) is separable algebraic over \(L\) and lies in \(L\). Hence \(x\) is purely inseparable over \(M\), as well as separable, so \(x\in M\). In characteristic zero \(M=L\). Thus \(M\) is separably closed in every case.

Every element of \(\Omega\) is separable algebraic over \(K\), and hence over the intermediate field \(M\). The preceding conclusion forces \(\Omega=M\). This proves exactly the source's assertions that \(\Omega/L\) is finite purely inseparable and that multiplication \[
 L\otimes_k K\longrightarrow\Omega,\qquad a\otimes b\longmapsto ab
\] is surjective. For the latter, the image is a finite-dimensional domain over \(L\), hence a field, containing \(L\) and \(K\); it is therefore their compositum \(M=\Omega\).

Now retain \(H^{sh}=R_{\mathfrak p}^{sh}\), \(U^{sh}=H^{sh}\otimes_R S\) and \(T^{sh}=H^{sh}\otimes_{R_{\mathfrak p}}S_{\mathfrak q}\). The prime in the statement is \[
 \mathfrak q'=\ker\bigl(T^{sh}\longrightarrow L\otimes_k K
                                     \longrightarrow\Omega\bigr).
\] The map is surjective: both the residue quotient/localization maps to \(L\) and \(K\) are surjective, as is the final multiplication map just proved. Thus \(\kappa(\mathfrak q')=\Omega\). The prime to which the four-rings lemma applies in \(U^{sh}\) is its inverse image under the original canonical localization map \(U^{sh}\to T^{sh}\). Denote that inverse image by \(\mathfrak Q'\); the primes have different ambient rings and are related by contraction and extension through this localization.

The canonical correction \texttt{MC-STK-ERR-0673} does exactly two things: names \(\mathfrak q'\) in the statement at 43650, and at 43683--43684 replaces the ambiguous reuse by the inverse image under the displayed canonical map. It preserves all original tensor factors. As above \((U^{sh})_{\mathfrak Q'}=(T^{sh})_{\mathfrak q'}\), through the inverse denominator maps already given. Quasi-finiteness and the finite-over-henselian theorem make this a finite henselian local \(H^{sh}\)-algebra with residue \(\Omega\), so it is strictly henselian. The earlier strict base-change theorem and uniqueness with this specified residue closure identify it with \(S_{\mathfrak q}^{sh}\). The induced residue extension is the original finite purely inseparable \(\Omega/L\). Applying this to the finite quotient \(R\to R/I\) gives the strict quotient formula in 43702--43712, with the original chosen residue closure retained.

\subsubsection{The local tensor ring and both localization maps}\label{algebra-proof-061-the-local-tensor-ring-and-both-localization-maps}

Keep the source's local maps \(A\to B\), \(A\to C\), the integral map \(A\to C\), and \(D=B\otimes_A C\). Write \(k_A,k_B,k_C\) for the original residue fields. Every maximal ideal of \(D\) contracts to \(\mathfrak m_B\), because \(D/B\) is integral and the contraction of a maximal ideal under an integral map is maximal. This last fact follows directly by passing to the quotient: an integral subring of a field over which that field is integral is itself a field, since a monic relation for the inverse of a nonzero element expresses that inverse in the subring.

The ring \(C/\mathfrak m_A C\) is integral over \(k_A\); all its primes are maximal by the same fact. It is local, so its only prime is \(\mathfrak m_C/\mathfrak m_A C\), a nil ideal \(N\). Under the exact field-base-change map \[
 k_B\otimes_{k_A}(C/\mathfrak m_A C)\longrightarrow k_B\otimes_{k_A}k_C
\] the kernel is the ideal generated by \(k_B\otimes N\), because tensoring over the field \(k_A\) is exact. This ideal is nil: each of its elements is a finite sum of multiples of finitely many nilpotent elements, and if their exponents are \(n_1,\ldots,n_r\), the sum ideal has exponent at most \(1+\sum(n_i-1)\). Thus this quotient does not change the spectrum.

If either original residue extension is purely inseparable, the spectrum of \(k_B\otimes_{k_A}k_C\) is one point. To verify the exact statement, suppose first that \(k_C/k_A\) is purely inseparable and use its unique embedding into an algebraic closure of \(k_B\). Multiplication maps the tensor ring to the compositum field \(k_Bk_C\). The image is a field: each of its elements has a suitable \(p\)-power in \(k_B\), and the inverse of a nonzero element \(z\) is \(z^{p^e-1}/z^{p^e}\). For an element \(\sum b_i\otimes c_i\) of the kernel choose a common exponent with all \(c_i^{p^e}\in k_A\). Its \(p^e\)-power is the scalar \(\sum b_i^{p^e}c_i^{p^e}=0\), so the kernel is nil. In characteristic zero the purely inseparable extension is the identity. If \(k_B/k_A\) is purely inseparable, interchange the two factors in this argument; the flip sends \(b\otimes c\) to \(c\otimes b\) and its square is the identity. The quotient field is nonzero, so the tensor ring has precisely that one prime. It follows that \(D\) has exactly one maximal ideal. Its contractions to \(B\) and \(C\) are the original maximal ideals, by the residue maps to this compositum. This proves all assertions in 43714--43739.

In case (1) of 43741--43842 keep \(D=B\otimes_A C\), the given map \(D\to\kappa\), and the exact chosen strict henselizations with residue fields \(\kappa_A,\kappa_B,\kappa_C,\kappa_D\) inside \(\kappa\). Put \[
 W=B\otimes_A C^{sh},\quad V=B\otimes_A A^{sh},\quad
 X=B^{sh}\otimes_{A^{sh}}C^{sh},\quad Y=D^{sh}.
\] The previous strict finite-local result gives the original descriptions \(B^{sh}=V_{\mathfrak q}\), \(Y=W_{\mathfrak r}\). If \(M\subset W\) is the multiplicative set consisting of images of \(V\setminus\mathfrak q\), then the exact map \[
 V_{\mathfrak q}\otimes_{A^{sh}}C^{sh}\longrightarrow M^{-1}W,
 \quad (v/u)\otimes c\longmapsto (\overline v(1\otimes c))/\overline u
\] is an isomorphism. The inverse sends each original tensor \(b\otimes c\in W\) to \((b\otimes1)/1\otimes c\), and sends the inverse of \(\overline u\) to the corresponding inverse in \(V_{\mathfrak q}\). Both composites fix tensors and the complete multiplicative set of denominators. Residue compatibility gives \((V\to W)^{-1}(\mathfrak r)=\mathfrak q\), so \(M\) avoids \(\mathfrak r\) and the comparison \(c:X\to Y\) is exactly the induced localization map.

The original maps \(A^{sh}\to B^{sh}\) and \(A^{sh}\to C^{sh}\) are local; the first is finite with purely inseparable residue extension. The just-proved tensor lemma makes \(X\) local. Both \(X\) and \(Y\) are finite over \(C^{sh}\), by the quasi-finite strict-henselization theorem and its base change, so \(Y\) is also finite over \(X\): any finite list of \(C^{sh}\)-module generators still generates as an \(X\)-module. In particular \(Y\) is integral over \(X\), and the contraction of its maximal ideal is the unique maximal ideal of \(X\). For each \(w\notin\mathfrak r\), its image in \(Y\) is outside the maximal ideal, hence its image in \(X\) is outside the maximal ideal of \(X\) and is a unit. Thus the explicit inverse to \(c\) is \[
 W_{\mathfrak r}\longrightarrow X,\qquad w/v\longmapsto w_X(v_X)^{-1}.
\] Localization's defining relation proves well-definedness; the composites fix every original tensor and every denominator. This supplies the omitted detail: locality alone for two unrelated localizations would not prove the claim, but the original finite comparison and the contracted maximal ideal do.

\subsubsection{The original filtered systems and the remaining cases}\label{algebra-proof-061-the-original-filtered-systems-and-the-remaining-cases}

In case (2), preserve the system \(B=\mathop{\rm colim}_i B_i\), all transition maps, \(D_i=B_i\otimes_A C\), and the primes and residue embeddings induced by the fixed map to \(\kappa\). All induced maps on the selected residue fields are embeddings. If \(K_i=\kappa(\mathfrak p_{B_i})\) and \(K=\kappa(\mathfrak p_B)\), their images in \(\kappa\) satisfy \(K=\bigcup_i K_i\). Every fraction of an element of \(B/\mathfrak p_B\) has its numerator and nonzero denominator at one stage, and conversely every stage fraction defines that fraction in \(K\); this proves both inclusions.

Let \(L_i\subset\kappa\) be the separable algebraic closure of \(K_i\) in the given separably closed field \(\kappa\). It is separably closed by the finite-coefficient argument above. The inclusions \(K_i\subset K_j\) give \(L_i\subset L_j\): an element separable algebraic over \(K_i\) remains so over \(K_j\), since its minimal polynomial over the latter divides a separable polynomial over the former. Their union equals precisely the separable closure of \(K\) inside \(\kappa\). One inclusion follows by the same separable-polynomial argument. For the other, the finite coefficients of a separable polynomial for a given element over \(K\) lie in one \(K_i\); that polynomial is already square-free over \(K_i\), since a common factor with its derivative would survive over \(K\). The element therefore belongs to \(L_i\). In particular this union is separably closed, with exactly the residue identification required by the source.

The functorial transition maps \(B_i^{sh}\to B_j^{sh}\) are local and induce these exact embeddings. Their colimit is local: an element outside the colimit maximal ideal is already a unit at a stage, and an element of the colimit maximal ideal has zero residue at its representing stage. Its residue field is \(\bigcup L_i\). It is henselian by the finite-coefficient simple-root proof in the preceding note: coefficients, a residue root and its nonzero derivative are realized at one stage, whose root lifts to the original colimit. Hence it is strictly henselian. The previously verified varying-base ind-étale lemma shows that this colimit is ind-étale over the original \(B=\mathop{\rm colim}B_i\). Uniqueness with the just-proved residue identification gives inverse local comparison maps \(\mathop{\rm colim}B_i^{sh}\leftrightarrows B^{sh}\); their composites are identities by uniqueness. The same proof applies to the entire original system \(D_i\), giving \(\mathop{\rm colim}D_i^{sh}=D^{sh}\).

The isomorphisms of case (1) commute with every transition, because all are the original tensor maps induced by the uniquely determined residue-compatible ring maps. Thus their colimit is an isomorphism. The final tensor/colimit comparison sends \([b_i\otimes c]\) to \([b_i]\otimes c\), with inverse \([b_i]\otimes c\mapsto[b_i\otimes c]\). Balancing, addition and equality of representatives hold at common later stages, so both maps are defined and inverse. This verifies the entire displayed formula at 43829--43830. The already composed \texttt{MC-STK-ERR-0674} correctly changes ``filtered colimit commute'' to ``filtered colimits commute'' at 43832. All three received reports describe this single operation.

For case (3) retain the comparison to the original rings instead of suppressing it. The maps to the given field select \(\mathfrak p_A,\mathfrak p_B,\mathfrak p_C,\mathfrak p_D\). The tensor \(B_{\mathfrak p_B}\otimes_{A_{\mathfrak p_A}}C_{\mathfrak p_C}\) is the localization of \(D\) inverting the images of \(B\setminus\mathfrak p_B\) and \(C\setminus\mathfrak p_C\); all images of \(A\setminus\mathfrak p_A\) are thereby inverted. The exact comparison is \[
 (b/u)\otimes(c/v)\longmapsto (b\otimes c)/((u\otimes1)(1\otimes v)).
\] The inverse sends \(b\otimes c\) to \((b/1)\otimes(c/1)\) and each of the two types of inverted generators to its corresponding inverse. The defining tensor and localization relations prove the inverse identities. All these denominators avoid \(\mathfrak p_D\); localization further at its extension is therefore exactly \(D_{\mathfrak p_D}\). The strict henselizations on both sides use that identical local ring and the identical residue closure in \(\kappa\), so their comparison maps are the identity under the unique specified isomorphism. This reduces case (3) to case (2) while retaining the original object and all maps.

For case (4), the integral \(A\)-algebra \(B\) is the directed union of its subalgebras generated by finite lists of its actual elements. Each such list is integral; retaining a monic equation for each generator proves finiteness by the spanning monomials whose exponent on that generator is smaller than its monic degree. Thus each subalgebra is finite over \(A\), hence quasi-finite. Their inclusions and union are the original \(B\); case (2) applies directly.

\subsubsection{Dispositions and scope}\label{algebra-proof-061-dispositions-and-scope}

\texttt{OCC-12382} is covered exactly once by the two existing operations in \texttt{MC-STK-ERR-0673}. Its kernel prime belongs to the localized tensor product and its inverse image belongs to the unlocalized tensor product; the explicit maps and contractions above verify both existing replacements.

\texttt{OCC-00976}, \texttt{OCC-12288} and \texttt{OCC-12383} all refer to the single plural correction already present in \texttt{MC-STK-ERR-0674}. The full filtered-system comparison above verifies that this grammatical edit describes the source's actual calculation.

These two canonical units are retained unchanged and now have semantic review. No newly proposed source operation, extension group, broad synthesis, formal certification, TeX build or publication is part of this batch. The full arguments remain separate source-linked editorial validation.

\hypertarget{algebra_serre_notes_20260929}{}
\subsection{Serre's criterion: source corrections and complete receiving arguments}\label{algebra-proof-062-serres-criterion-source-corrections-and-complete-receiving-arguments}

Source: Stacks Project authors, original commit \texttt{a04446e57ec1fbc252a871afcec7752fb2807b14}, \texttt{algebra.tex:43865–44126}, SHA256 \texttt{FA8BB92E58A4F78A2BD01B3B6A4A87DE0A0D279F5DD90641B574DD5FBFFFA4F3}. The source cites EGA IV, Theorem 5.8.6, for Serre's criterion; that citation is retained, without claiming a fresh reading of that EGA passage. The four-case local theorem at \texttt{algebra.tex:28922–29006} is attributed there to János Kollár, ``Variants of normality for Noetherian schemes''; its complete prior validation is in \texttt{ALGEBRA\_KRULL\_AKIZUKI\_NOTES\_20260928.md}, section ``The four local cases and an explicit annihilator exponent''.

This batch checks seven reports and the original proofs, retaining three existing canonical units with four operations. Two additional minimal operations are needed: the missing period at 44065, and the inclusion qualification at 44037. The latter keeps the zero quotient in the existing corrected Serre argument. All expanded arguments below remain separate editorial validation. Optional extensions and additional paper mining remain deferred until Algebra is finished.

\subsubsection{The depth conditions and reducedness}\label{algebra-proof-062-the-depth-conditions-and-reducedness}

Keep exactly the original conditions: \((R_k)\) requires regularity of \(R_{\mathfrak p}\) for every prime of height at most \(k\); \((S_k)\) for the original finite module \(M\) requires \[
 \operatorname{depth}_{R_{\mathfrak p}}M_{\mathfrak p}
 \geq \min\{k,\dim\operatorname{Supp}(M_{\mathfrak p})\}
\] at every prime. The source's zero-module conventions, depth \(+\infty\) and support dimension \(-\infty\), are retained. Conditions for rings use \(M=R\). These conditions localize: primes of \(R_{\mathfrak q}\) correspond to primes contained in \(\mathfrak q\), their further localizations are the original \(R_{\mathfrak p}\), and all primes below such a \(\mathfrak p\) are still contained in \(\mathfrak q\), so the relevant height and support chains are unchanged.

The established associated-prime localization and depth-zero criteria imply \(\operatorname{depth}M_{\mathfrak p}=0\) exactly when \(\mathfrak p\in\operatorname{Ass}_R(M)\). An embedded associated prime contains a strictly smaller minimal support prime, which is associated for a finite module over a Noetherian ring. The corresponding two-point chain survives in the localized support, so its dimension is at least one and \((S_1)\) fails there. Conversely failure of \((S_1)\) requires depth zero and support dimension at least one. There is then a smaller support prime; choosing a minimal support prime below it gives a smaller associated prime. Hence the original associated prime is embedded. This proves 43896--43932, including the zero module, where both sides hold and there are no associated primes.

If \(R\) is reduced, localization at a minimal prime is a reduced zero-dimensional local ring, hence a field, giving \((R_0)\). For a localization \((A,\mathfrak m)\) of positive dimension, an associated maximal ideal would be the annihilator of a nonzero \(x\in A\). Such an \(x\) cannot be a unit: that would give \(\mathfrak m=0\) and dimension zero. Thus \(x\in\mathfrak m\), whence \(x^2=0\), contradicting reducedness. Finite associated-prime avoidance supplies a nonzerodivisor in \(\mathfrak m\), proving \((S_1)\).

Conversely assume \((R_0)\) and \((S_1)\). For each \(\mathfrak p\), its Noetherian local ring has finite dimension: its maximal ideal has finitely many generators and the height theorem bounds its dimension by their number. We may therefore induct on this finite height separately for each prime; no global finite-dimensionality assumption on \(R\) is needed. At height zero \((R_0)\) makes the local ring a field. At positive height, \((S_1)\) gives the original nonzerodivisor \(t\in\mathfrak pR_{\mathfrak p}\). The map \(R_{\mathfrak p}\to R_{\mathfrak p}[1/t]\), \(a\mapsto a/1\), is injective: \(a/1=0\) means \(t^na=0\), which implies \(a=0\). The latter ring injects into the product of its prime localizations by the established local detection of zero. Its primes correspond to primes \(\mathfrak q\subsetneq\mathfrak p\) avoiding \(t\); the identification of their local rings sends the original fraction \((a/t^n)/v\) to that same quotient in \(R_{\mathfrak q}\), with inverse induced by localization. These primes have strictly smaller height, since any chain below \(\mathfrak q\) can be extended by \(\mathfrak p\). All factors are reduced by induction, so their subring \(R_{\mathfrak p}\) is reduced. Applying local detection of zero to \(R\) proves its reducedness.

\subsubsection{The original Serre argument and its zero quotient}\label{algebra-proof-062-the-original-serre-argument-and-its-zero-quotient}

Suppose first that \(R\) is normal. Its localizations are normal domains. In dimensions zero and one they are respectively fields and discrete valuation rings by the source's established DVR characterization; in particular they have the required regularity and depth. Let \((R,\mathfrak m)\) now be the source's local Noetherian normal domain of dimension at least two. In the four-case local theorem it is neither Artinian nor regular of dimension one. In case (4) retain the finite map \(R\to R'\), its kernel and cokernel killed by \(\mathfrak m^N\), and the assumptions \(R'\ne0\) and \(\mathfrak m\notin\operatorname{Ass}_R(R')\).

The kernel is zero: choose a nonzero \(u\in\mathfrak m\); its \(N\)th power kills every kernel element and is nonzero in the domain. Since \(R'\) is a finite module and \(\mathfrak m\) is not associated, finite prime avoidance gives the original \(x\in\mathfrak m\) that is a nonzerodivisor on \(R'\). Both maps \(R\to R_x\) and \(R'\to R'_x\) are injective. Their localizations are equal because \(x^N\) kills the original kernel and cokernel. Thus the exact composite \[
 R'\longrightarrow R'_x=R_x
 \longrightarrow\operatorname{Frac}(R)
\] is the map \(r'\mapsto r'/1\) into \(R'_x=R_x\), followed by \(a/x^n\mapsto a/x^n\) in the fraction field. The multiplicative set is \(\{1,x,x^2,\ldots\}\); every denominator is retained. This embeds \(R'\) in the original fraction field. Its elements are integral over \(R\) by finiteness, so normality forces \(R'=R\), contradicting the case-(4) nonisomorphism. Hence the remaining case gives depth at least two. This proves \((R_1)\) and \((S_2)\) for every original localization.

For the converse, localize the original \((R_1),(S_2)\) ring at each prime. It is reduced by the preceding argument, because these conditions imply \((R_0),(S_1)\). Induct on the finite dimension of this local ring. The zero- and one-dimensional cases are a field and a DVR. In dimension \(d\geq2\), keep the source element \(x\in Q(R)\) integral over \(R\), and the original subring \(R'=R[x]\subset Q(R)\). It is finite over \(R\), since the powers below the degree of a monic relation for \(x\) span it. No new fraction field is substituted for the total quotient ring.

For every nonmaximal \(\mathfrak p\), induction gives that \(R_{\mathfrak p}\) is a normal domain. Every nonzerodivisor of the original reduced \(R\) remains a nonzerodivisor in \(R_{\mathfrak p}\), since localization preserves the injection given by multiplication. The localization \(Q(R)_{\mathfrak p}\) is therefore the localization of this domain at the images of those original nonzerodivisors. Its map to \(\operatorname{Frac}(R_{\mathfrak p})\) sends \((a/s)/v\) to \((a/v)/s\), for \(v\notin\mathfrak p\) and an original nonzerodivisor \(s\). It is injective: vanishing of that fraction means \(a/v=0\) in the domain \(R_{\mathfrak p}\), which makes the original localized fraction zero. Localization also preserves the injection \(R'\subset Q(R)\). Normality of \(R_{\mathfrak p}\) now puts the image of the original integral \(x\) in \(R_{\mathfrak p}\), proving \(R'_{\mathfrak p}=R_{\mathfrak p}\) under these actual maps. Hence the correct statement for the actual quotient is \[
 \operatorname{Supp}_R(R'/R)\subseteq\{\mathfrak m\}.
\] The inclusion may be strict: choosing \(x=0\in R\subset Q(R)\) gives the monic equation \(X=0\), the original ring \(R'=R[0]=R\), quotient zero and empty support. This is an instance allowed by the source's choice of an integral \(x\). More generally every choice \(x\in R\) gives this empty support. Accordingly the proposed minimal correction at 44037 is ``is contained in'', retaining the original module and singleton. The already composed case distinction at 44040 does not retrospectively justify equality in the preceding sentence.

Because \(R'/R\) is finite and its support is contained in that singleton, some power of \(\mathfrak m\) annihilates it. One exact exponent is obtained by taking generators \(u_1,\ldots,u_s\) of \(\mathfrak m\), choosing \(n_i\geq1\) with \(u_i^{n_i}\) in the annihilator, and setting \(N=1+\sum_i(n_i-1)\). Every degree-\(N\) monomial contains one of these powers. The zero module also satisfies this annihilation statement, with every chosen exponent.

The existing \texttt{MC-STK-ERR-0675} then correctly separates \(x\in R\), where the conclusion is immediate, from \(x\notin R\). In the latter case \(R\to R'\) is injective and not an isomorphism. Since \(\operatorname{depth}R\geq2\), there is a nonzerodivisor \(t\in\mathfrak m\) on \(R\). Its inverse is part of the definition of the original \(Q(R)\), so multiplication by \(t\) is injective on \(R'\subset Q(R)\). Thus \(\mathfrak m\notin\operatorname{Ass}_R(R')\), and \(R'\ne0\) because it contains the nonzero local ring \(R\). These are exactly all case-(4) hypotheses, now verified on the original map. Case (3), depth at least two, excludes case (4). Equivalently the depth lemma applied to \(0\to R\to R'\to R'/R\to0\) gives depth at least one for a nonzero quotient, whereas a nonzero finite module killed by a power of \(\mathfrak m\) has depth zero. Therefore no \(x\notin R\) is possible, and \(R\) is integrally closed in its original \(Q(R)\).

To conclude that this reduced local ring is a normal domain, retain the exact source comparison at 8274--8309. There are finitely many minimal primes \(\mathfrak p_i\), and \(Q(R)=\prod_iR_{\mathfrak p_i}\) is their product of fields. Each coordinate idempotent \(e_i\) satisfies the complete monic equation \(e_i^2-e_i=0\), so belongs to \(R\) by integral closedness. The component maps are \(r\mapsto(e_ir)_i\), with inverse \((r_i)_i\mapsto\sum_i r_i\). A nonzero local ring has no idempotents other than 0 and 1: one of \(e,1-e\) is a unit, and \(e(1-e)=0\) then gives the claim. Thus only one of these nonzero field coordinates can occur. The reduced local ring is a domain and is integrally closed in its fraction field, as required. This proves the original criterion, including the zero-quotient case that the new minimal edit restores.

The corollary for regular rings follows with the original convention of Noetherian regularity: all local rings are regular and therefore Cohen--Macaulay, so they have depth equal to dimension. Both \((R_1)\) and \((S_2)\) follow on those same local rings, and the just-proved criterion applies.

\subsubsection{Height-one intersections with all original factors retained}\label{algebra-proof-062-height-one-intersections-with-all-original-factors-retained}

Let \(R\) be the original Noetherian normal domain with fraction field \(K\). For \(0\ne a\in R\), first allow the unit case: \(R/aR=0\), its associated-prime set is empty, and all stated restrictions are vacuous. Otherwise localize at a prime in its support. Multiplication by the actual \(a\) is injective on \(R_{\mathfrak p}\), giving \[
 0\longrightarrow R_{\mathfrak p}\xrightarrow{\ a\ }R_{\mathfrak p}
 \longrightarrow (R/aR)_{\mathfrak p}\longrightarrow0.
\] If the localized quotient support has dimension at least one, it contains a chain of two different primes containing \(a\). The smaller prime has height at least one because \(a\ne0\), so \(\operatorname{height}(\mathfrak p)\geq2\). The \((S_2)\) condition gives depth of \(R_{\mathfrak p}\) at least two; the exact depth lemma gives depth of the quotient at least one. At zero-dimensional support the \((S_1)\) bound is automatic, and outside the support the module is zero. Hence \(R/aR\) has \((S_1)\), and no embedded primes. Its associated primes are precisely its minimal support primes. The principal ideal theorem gives their height at most one, and \(a\ne0\) excludes height zero. This proves item (1) without a catenarity or dimension-drop assumption.

For \(b\in R\), apply the injective original localization map for the module \(R/aR\), proved at 15495--15516: \[
 R/aR\longrightarrow\prod_{\mathfrak p\in\operatorname{Ass}(R/aR)}(R/aR)_{\mathfrak p},
 \qquad [b]\longmapsto([b]/1)_{\mathfrak p}.
\] It says exactly that \(b\in aR\) if and only if \(b\in aR_{\mathfrak p}\) at those height-one primes. Consequently the original fraction \(b/a\) belongs to \(R\) if it belongs to every height-one localization; the converse is the structure inclusion. All fractions in \(K\) have such a representation. If the indexing set is empty, the intersection is taken inside the original \(K\); the same argument says every fraction belongs to \(R\), so \(R=K\). This proves item (2) with the empty case included.

For the original nonzero \(x\in K\), keep the actual representation \(x=a/b\) with both \(a,b\ne0\), and let \(\mathfrak p_1,\ldots,\mathfrak p_r\) be all minimal primes over \((ab)\). Each has height one. Every height-one prime containing \(a\) or \(b\) is among them: it contains a minimal prime over \((ab)\), and a strict inclusion would increase the height beyond one. At every other height-one prime both \(a\) and \(b\) are units, so \(xR_{\mathfrak p}=R_{\mathfrak p}\). Applying item (2) to the original fraction \(y/x\), for each \(y\in R\), gives the exact ideal equality \[
 I=R\cap xR=\bigcap_{i=1}^r I_i,\qquad I_i=R\cap xR_{\mathfrak p_i}.
\] Thus the original map \(R/I\to\bigoplus_{i=1}^rR/I_i\), \([y]\mapsto([y])_i\), has kernel \((\bigcap I_i)/I=0\). When \(r=0\), the intersection is \(R\), both this source and target are zero, and the same statement holds.

Each \(R_{\mathfrak p_i}\) is a DVR. Let \(v_i\) be its valuation and retain \(e_i=v_i(a)-v_i(b)=v_i(x)\). If a uniformizer \(\pi_i\) is used to express the comparison, retain also the exact unit \(u_i=x\pi_i^{-e_i}\), so \(x=u_i\pi_i^{e_i}\). Multiplication by \(u_i\) and by its inverse give both inclusions in \(xR_{\mathfrak p_i}=\pi_i^{e_i}R_{\mathfrak p_i}\). No unit factor or original numerator/denominator is discarded. For \(e_i\leq0\) the fractional ideal contains \(R_{\mathfrak p_i}\), hence \(I_i=R\) and the corresponding original summand is exactly zero. For \(e_i>0\), \(I_i=\mathfrak p_i^{(e_i)}\) by the source definition as a contraction. Thus the full original direct sum has those zero summands together with the stated positive-exponent quotients; removing the zero summands is the isomorphism given by projection, with inverse inserting zeros.

For positive \(e_i\), every element outside \(\mathfrak p_i\) acts injectively on \(R/\mathfrak p_i^{(e_i)}\): from \(sy\in\mathfrak p_i^{(e_i)}\) and \(s\notin\mathfrak p_i\), invert \(s\) in the original \(R_{\mathfrak p_i}\) to conclude \(y\in\mathfrak p_i^{e_i}R_{\mathfrak p_i}\), hence in its contraction. The support has radical \(\mathfrak p_i\), and the nonzero quotient has an associated prime. These two facts force its only associated prime to be \(\mathfrak p_i\), agreeing with the complete previously recorded symbolic-power proof. Associated primes of a submodule lie among those of its ambient module, and a finite direct sum has their union. Therefore every associated prime of the original \(R/(R\cap xR)\) is one of the original height-one \(\mathfrak p_i\). Distinct height-one primes in a domain cannot contain one another, so there are no embedded associated primes. This proves item (3), retaining all zero summands and all original fractional factors.

\subsubsection{Report accounting and propagation}\label{algebra-proof-062-report-accounting-and-propagation}

\texttt{OCC-12384} is covered by the retained operation in \texttt{MC-STK-ERR-0675}, together with the newly necessary support inclusion at 44037. The retained operation verifies nonisomorphism, zero kernel, nonzero target and the nonassociated maximal ideal. The added operation covers the \(x\in R\) branch that was already allowed before that operation. Its adverse example is the exact choice \(x=0\), not a change of ring or hypothesis.

\texttt{OCC-00977} requires the missing terminal period at 44065. \texttt{OCC-00978}, \texttt{OCC-12385} and \texttt{OCC-00980} are all accounted for by the two already composed adjective replacements in \texttt{MC-STK-ERR-0676}, at 44071 and 44108. \texttt{OCC-00979} and \texttt{OCC-12386} share the already composed terminal period in \texttt{MC-STK-ERR-0677}, at 44078.

The source's criterion and its following corollaries have been checked with the repaired support inclusion. The zero quotient does not change the annihilation assertion or the nontrivial case-(4) contradiction. The original height-one intersection and symbolic-power conclusions remain valid, with their complete comparisons above. The ledger records this as a correction of an overclaim about the exact support, not a new extension or a novelty claim. Full editorial proofs remain separate from the minimal source corrections and from any translation.

\hypertarget{algebra_field_formal_smooth_notes_20260929}{}
\subsection{Formal smoothness of fields: original-source validation}\label{algebra-proof-063-formal-smoothness-of-fields-original-source-validation}

Source: Stacks Project authors, \texttt{algebra.tex:44127–44565}, original commit \texttt{a04446e57ec1fbc252a871afcec7752fb2807b14}, file SHA256 \texttt{FA8BB92E58A4F78A2BD01B3B6A4A87DE0A0D279F5DD90641B574DD5FBFFFA4F3}. This note preserves the actual extensions, polynomials, coefficients, differential modules and maps in that source. It validates the original assertions and received corrections; optional generalizations and wider paper mining remain deferred until completion of the Algebra review.

\subsubsection{Finite generation and the original differential kernel}\label{algebra-proof-063-finite-generation-and-the-original-differential-kernel}

For the six equivalent conditions in 44136--44171 keep the stated finite generation of the field extension \(K/k\). The established characterizations of formal unramifiedness, étale field algebras and formal étaleness give all implications except \(\Omega_{K/k}=0\Rightarrow K/k\) finite separable. Choose the original finite-type domain \(A\subset K\) with fraction field \(K\), and retain the multiplicative set \(S=A\setminus\{0\}\). The exact localization map is \[
 S^{-1}\Omega_{A/k}\longrightarrow\Omega_{K/k},\qquad
 \omega/s\longmapsto s^{-1}\omega,
\] with the inverse determined on every original fraction by \(d(a/s)\mapsto (s\,da-a\,ds)/s^2\). The quotient and Leibniz rules give both inverse identities. If the target is zero, choose a finite generating list \(\omega_1,\ldots,\omega_m\) of \(\Omega_{A/k}\). For each generator its zero localization gives an actual \(s_i\in S\) with \(s_i\omega_i=0\). The full product \(f=\prod_i s_i\) is nonzero in the original domain and kills all generators. For an empty list take \(f=1\). Thus \(\Omega_{A_f/k}=0\) under the same comparison maps. This finite-type algebra is unramified by the previously proved differential criterion. Its selected prime is zero and its residue field is the unchanged \(K\), so the unramified-at-a-prime theorem makes \(K/k\) finite separable. No assertion is made here for an arbitrary extension without finite generation.

For 44173--44226 retain the perfect base \(k\) of characteristic \(p>0\), the extension \(K/k\), and the given \(a\in K\). The easy direction keeps the full derivative identity \(d(b^p)=p b^{p-1}db=0\). Conversely if \(da=0\), the directed colimit of the differential modules of finite generated intermediate fields realizes this particular zero relation at one such field containing \(a\). Indeed a class zero in a filtered module colimit is zero at a later stage, and the finite set of field elements needed for that stage generates a finite generated field extension. This reduces the proof to the original finite generated situation without assuming injective transition maps on differential modules.

Choose the source separating transcendence basis \(x_1,\ldots,x_r\) and put \(F=k(x_1,\ldots,x_r)\), retaining \(K/F\) finite separable. Its existence follows from the original definition of a perfect field, \texttt{algebra.tex:10488–10500}, and separability of finite generated extensions. The canonical derivations \(\partial_j\) on \(F\) give \(\Omega_{F/k}=\bigoplus_j F\,dx_j\); the inverse universal map sends \(df\) to \(\sum_j\partial_j(f)dx_j\), with the complete quotient derivative \(\partial_j(P/Q)=(Q\partial_jP-P\partial_jQ)/Q^2\).

Here is the source's rational-function claim with every coefficient retained. Perfection gives \(F^p=k(x_1^p,\ldots,x_r^p)\). The monomials \(x^I\), with all \(0\leq i_j<p\), form an \(F^p\)-basis of \(F\). Polynomial spanning follows by writing each exponent as a multiple of \(p\) plus its remainder and taking the actual coefficient's \(p\)th root in \(k\). For an original nonzero denominator \(Q\), keep the identity \(Q^{-1}=Q^{p-1}/Q^p\); this proves rational spanning without cancelling its factors. For independence, multiply a relation over \(k(x_1^p,\ldots,x_r^p)\) by the full product of its coefficient denominators. Distinct residue vectors \(I\) give disjoint monomial supports in the polynomial ring, so every coefficient vanishes. Accordingly write an arbitrary original \(f\in F\) uniquely as \[
 f=\sum_{0\leq i_j<p} b_I^p x^I.
\] In its derivative with respect to \(x_j\), the terms with \(i_j=0\) have coefficient zero and the remaining terms are \(i_jb_I^p x^{I-e_j}\). Their exponents still lie in the displayed basis range and are distinct. If all partial derivatives vanish, each coefficient with \(i_j>0\) is zero, since that integer is nonzero in the prime field. Every \(I\ne0\) has such a coordinate, so only \(b_0^p\) remains. For \(r=0\) the same formula consists of its single constant term.

The comparison to the original \(K\) is also explicit. A primitive generator \(\theta\) for the finite separable auxiliary extension \(K/F\) has original monic minimal polynomial \(M(T)=T^n+\sum_{i=1}^n c_iT^{n-i}\), with \(M'(\theta)\ne0\). Its differential relation is \[
 M'(\theta)d\theta+\sum_{i=1}^n\theta^{n-i}dc_i=0.
\] Thus the canonical map \(K\otimes_F\Omega_{F/k}\to\Omega_{K/k}\) is an isomorphism. Its inverse sends each \(dx_j\) to itself and sends \[
 d\theta\longmapsto-\bigl(M'(\theta)\bigr)^{-1}
                         \sum_{i=1}^n\theta^{n-i}dc_i.
\] The full relation above makes that map well-defined on the polynomial presentation and then on all fractions by the quotient rule. The two composites fix the differentials of the original generators, so are identities. No derivative coefficient or inverse factor is omitted.

Now retain the source monic minimal polynomial of the original \(a\), \(P(T)=T^d+\sum_{i=1}^d a_iT^{d-i}\in F[T]\). Differentiation gives the complete identity \[
 0=P'(a)da+\sum_{i=1}^d a^{d-i}da_i.
\] Only after retaining this formula use \(da=0\). The just-proved basis \(dx_j\) of \(\Omega_{K/k}\) gives \(\sum_{i=1}^d a^{d-i}\partial_j(a_i)=0\) for every \(j\). The powers \(1,a,\ldots,a^{d-1}\) are independent over \(F\) by minimality of the original \(P\). Hence all those partial derivatives vanish, and \(a_i=b_i^p\) in the original \(F\).

Choose the unique \(p\)th root \(\beta\) of \(a\) in an algebraic closure of \(K\), and retain the coefficient-root polynomial \(Q(T)=T^d+\sum_{i=1}^d b_iT^{d-i}\). The exact equality \(Q(\beta)^p=P(a)=0\) gives \(Q(\beta)=0\), so \([F(\beta):F]\leq d\). But \(a=\beta^p\) yields the original inclusion \(F(a)\subseteq F(\beta)\), and \([F(a):F]=d\). The tower degree therefore forces \(F(a)=F(\beta)\), placing \(\beta\) inside the original \(K\). This supplies the required root, consistently with the complete original source lemma \texttt{fields.tex:3746–3774}. It does not replace \(a\) or its minimal polynomial by a rescaled object.

For 44228--44253 let the source differentials \(da_1,\ldots\,da_n\) be independent in \(\Omega_{k/\mathbf F_p}\). The base \(\mathbf F_p\) is perfect, so the preceding result makes \(a_1\) a non-\(p\)th-power. The polynomial \(T^p-a_1\) is irreducible and gives degree \(p\). Inductively the preceding root adjunctions have their full monomial basis \(\prod_{j<n}a_j^{i_j/p}\), \(0\leq i_j<p\). If \(a_n\) became a \(p\)th power in that field, the actual putative root would have a unique expression with coefficients \(\lambda_I\in k\), giving exactly \[
 a_n=\sum_I\lambda_I^p\prod_{j<n}a_j^{i_j}.
\] Differentiation retains every term: \(d(\lambda_I^p)=0\), and the coefficient of \(da_j\) is the full sum \(\sum_{I:i_j>0}\lambda_I^p i_j a_j^{i_j-1}\prod_{\ell\ne j}a_\ell^{i_\ell}\). The omitted indices with \(i_j=0\) contribute exactly zero. This expresses \(da_n\) in the span of the preceding differentials, a contradiction. The next irreducible \(p\)th-root polynomial has degree \(p\), giving the original \(p^n\) by the tower law. The empty tower, when used, has degree \(1=p^0\).

\subsubsection{The corrected coefficient field and both differential maps}\label{algebra-proof-063-the-corrected-coefficient-field-and-both-differential-maps}

In 44255--44379 the lemma needs to introduce the actual extension \(K/k\) in its statement. The correction adds that datum after the original characteristic hypothesis. All subsequent maps then have their stated domains.

For the reduction to finite generated subfields keep the original system \(K=\mathop{\rm colim}K_i\). If the final map \(K\otimes_k\Omega_{k/\mathbf F_p}\to\Omega_{K/\mathbf F_p}\) is injective, the corresponding map at \(K_i\) is injective: its kernel maps to zero in the final module, while \(K_i\otimes_k\Omega_{k/\mathbf F_p}\to K\otimes_k\Omega_{k/\mathbf F_p}\) is injective because it is scalar extension along fields. Conversely if each stage map is injective, its filtered colimit is injective, since filtered colimits of modules preserve exactness. The source and target colimits are respectively the actual scalar extension and the actual differential module by their generator/Leibniz presentations. This proves both directions, without assuming that the target differential transition maps are injective.

For a finite generated separable \(K/k\), retain the source separating generators \(x_1,\ldots,x_r,x_{r+1}\) and its irreducible polynomial \(G\in k[X_1,\ldots,X_{r+1}]\). The evaluation map has kernel \((G)\): over the rational function field in the first \(r\) variables this is the minimal-polynomial ideal, and its primitive irreducible polynomial descends the same principal ideal to \(k[X_1,\ldots,X_r][X_{r+1}]\) by Gauss's lemma. Its image is the original \(k[x_1,\ldots,x_{r+1}]\) and its fraction field is \(K\). Thus the source domain is \(S=k[X_1,\ldots,X_{r+1}]/(G)\), with the original lower-case coefficient field. \texttt{MC-STK-ERR-0678} restores exactly this field and joins its defining sentence.

Retain the full polynomial \(G=\sum_I a_I X^I\). The universal differential presentation of that actual ring is \[
 \Omega_{S/\mathbf F_p}=
 \frac{(S\otimes_k\Omega_{k/\mathbf F_p})\oplus\bigoplus_j S\,dX_j}
 {S\bigl(\sum_I X^I\otimes da_I,\ \sum_j(\partial_jG)dX_j\bigr)}.
\] The map sends \(s\otimes da\) to \(s\,da\) and \(s\,dX_j\) to \(s\,d\overline X_j\). Its inverse is induced by the derivation on the polynomial ring obtained by differentiating every coefficient and every monomial; its value on the full original relation \(G\) is precisely the displayed vector. These constructions factor through the quotient and are inverse on every generator. Localizing at every nonzero element of the original \(S\), with the full quotient derivative already given above, yields \[
 \Omega_{K/\mathbf F_p}=
 \frac{(K\otimes_k\Omega_{k/\mathbf F_p})\oplus\bigoplus_j K\,dx_j}
 {K\bigl(\sum_I x^I\otimes da_I,\ \sum_j(\partial_jG)(x)dx_j\bigr)}.
\] Both occurrences of the source shorthand \(\Omega_k\) refer to \(\Omega_{k/\mathbf F_p}\); \texttt{MC-STK-ERR-0679} makes the base explicit without changing the module.

The derivative \(\partial_{r+1}G\) has nonzero image in \(K\). It is a nonzero polynomial of smaller degree in \(X_{r+1}\), hence is not divisible by the irreducible \(G\), which is the evaluation kernel. If \((\omega,0)\) in the localized numerator is a multiple \(t\) of the relation vector, its final coordinate is \(0=t(\partial_{r+1}G)(x)\). The nonzero field element forces \(t=0\), then \(\omega=0\). This proves injectivity of the original canonical map to \(\Omega_{K/\mathbf F_p}\), the codomain restored by \texttt{MC-STK-ERR-0680}. A map from \(K\otimes_k\Omega_{k/\mathbf F_p}\) into the unlocalized \(\Omega_{S/\mathbf F_p}\) does not in general have the needed scalar action. For example, with \(k=\mathbf F_p(t)\), \(S=k[X,Y]/(Y)=k[X]\) and \(K=k(X)\), the localized element \(X^{-1}\otimes dt\) has image \(X^{-1}dt\), which is not in the original free \(S\)-module \(S\,dt\oplus S\,dX\). The localization map between the two differential modules retains their exact relation.

\subsubsection{The original inseparable polynomial and its coefficient factor}\label{algebra-proof-063-the-original-inseparable-polynomial-and-its-coefficient-factor}

For the converse injectivity criterion choose the source transcendence basis \(x_1,\ldots,x_r\) with minimal finite inseparable degree of \(K/F\), \(F=k(x_1,\ldots,x_r)\). Suppose this degree is greater than one and choose the original inseparable \(\alpha\in K\). Its minimal polynomial is a polynomial in \(T^p\). Clearing its original denominators and taking a primitive irreducible polynomial yields exactly a \(G=\sum_{I,i} a_{I,i}X^IT^i\in k[X_1,\ldots,X_r,T]\) whose evaluation kernel is \((G)\), with \(\partial_TG=0\). Keep that actual \(G\), including an arbitrary chosen nonzero coefficient \(c=a_{I_0,i_0}\); do not replace it by a polynomial with that coefficient suppressed.

Assume also that all \(\partial_{X_j}G\) vanish. Every supported exponent in \(X\) and \(T\) is then divisible by \(p\). Write \(m_{I,i}=x^I\alpha^i\). Differentiation of the original relation, with all coordinate derivatives zero, gives \(\sum_{I,i}m_{I,i}\,da_{I,i}=0\) in \(\Omega_{K/\mathbf F_p}\). Injectivity puts this relation in \(K\otimes_k\Omega_{k/\mathbf F_p}\) itself.

The supported monomials other than \(m_{I_0,i_0}\) are linearly independent over \(k\). A relation among them is a polynomial in the evaluation kernel \((G)\), of total degree at most that of \(G\); it must therefore be a scalar multiple of \(G\). Its coefficient at the omitted supported monomial is zero, while that of \(G\) is \(c\ne0\), so this scalar and the entire relation vanish. Substituting the actual original relation \(c\,m_{I_0,i_0}=-\sum_{(I,i)\ne(I_0,i_0)}a_{I,i}m_{I,i}\) into the differential relation therefore gives, for every other coefficient, \[
 da_{I,i}-\frac{a_{I,i}}c\,dc=0,
 \qquad
 d(a_{I,i}/c)=\frac{c\,da_{I,i}-a_{I,i}\,dc}{c^2}=0.
\] The full terms involving \(dc\) and \(c^2\) remain present. The preceding differential-kernel theorem over the perfect prime field gives actual \(b_{I,i}\in k\) with \(a_{I,i}/c=b_{I,i}^p\). For the chosen coefficient take \(b_{I_0,i_0}=1\), and for a zero coefficient take zero. Since every supported exponent is divisible by \(p\), the original polynomial satisfies the complete factorization \[
 G=c\left(\sum_{I,i} b_{I,i}X^{I/p}T^{i/p}\right)^p.
\] The polynomial in parentheses is nonconstant because \(G\) is nonconstant; writing the right-hand side as \((cH)H^{p-1}\) gives two nonunits and contradicts irreducibility. The source's special choice \(c=1\) is recovered by that substitution, while this validation retains the actual coefficient throughout rather than silently changing the working polynomial.

Thus some \(\partial_{X_j}G\) is nonzero; keep that coordinate and use the source's indexing name \(X_1\) for it. The evaluation of this derivative is nonzero because its degree in \(X_1\) is smaller than that of the irreducible \(G\). The elements \(x_2,\ldots,x_r,\alpha\) are algebraically independent: a polynomial relation among them would belong to \((G)\) but have degree zero in \(X_1\), impossible since \(\deg_{X_1}G>0\). Put \(F'=k(x_2,\ldots,x_r,\alpha)\) and retain the common field \(E=F(\alpha)=F'(x_1)\). The original polynomial proves that \(x_1\) is separable algebraic over \(F'\), so \(E/F'\) is finite separable, whereas \([E:F]_i>1\) by the chosen inseparability of \(\alpha\). The two actual towers give \[
 [K:F]_i=[K:E]_i[E:F]_i>[K:E]_i
        =[K:E]_i[E:F']_i=[K:F']_i.
\] The finite inseparable-degree multiplication used here is the complete finite case of \texttt{fields.tex:1778–1803}: embeddings of the intermediate field into an algebraic closure each have exactly the same number of extensions to the top field, giving multiplication of separable degrees; divide the ordinary degree multiplication by that identity. All degrees here are finite, so the source's omitted infinite-degree case is not used. This strict inequality contradicts the original minimality. If \(r=0\), the coefficient argument already supplies the contradiction without selecting a coordinate. The differential injectivity criterion is proved on the original fields.

\subsubsection{Formal lifting and the characteristic-zero assertions}\label{algebra-proof-063-formal-lifting-and-the-characteristic-zero-assertions}

For the implication in 44381--44394, first retain its two characteristic cases. In characteristic zero every finite generated intermediate extension has a transcendence basis with finite algebraic remainder, and that remainder is separable because irreducible polynomials over a field of characteristic zero have nonzero derivative. Thus the extension is separable by the source definition. In characteristic \(p>0\), every derivation of \(k\) or \(K\) over \(\mathbf Z\) kills \(\mathbf F_p\), since it kills \(1\) and every integer multiple of \(1\). The identity on universal differentials therefore gives inverse comparison maps \(\Omega_{k/\mathbf Z}\leftrightarrows\Omega_{k/\mathbf F_p}\) and \(\Omega_{K/\mathbf Z}\leftrightarrows\Omega_{K/\mathbf F_p}\), each sending \(da\) to \(da\). Under these actual maps the short exact differential sequence for a formally smooth extension gives the injection just proved equivalent to separability. This validates the characteristic split and explicit base in the retained MC-STK-ERR-0681, rather than applying a positive-characteristic lemma to characteristic zero.

For 44396--44406 the earlier characterization, algebra.tex:37985--38037, says that formal smoothness is equivalent to \(H_1(L_{K/k})=0\) together with projectivity of \(\Omega_{K/k}\). The latter is a vector space over the actual field \(K\): choose a basis to obtain the isomorphism from the direct sum of copies of \(K\) sending each unit basis vector to the corresponding differential-module basis element. It is free and hence projective. Therefore the vanishing condition is exactly sufficient as well as necessary. The retained correction MC-STK-ERR-0682 changes only ``a vector spaces'' to ``a vector space''.

Here is the full lifting argument for 44408--44449. For a purely transcendental extension keep the original \(K=k(x_j;j\in J)\), the original \(k\)-algebra \(A\), the ideal \(I\) with \(I^2=0\), and the map \(\varphi:K\to A/I\). Choose the actual lifts \(a_j\in A\) of \(\varphi(x_j)\). Evaluation gives \(k[x_j;j\in J]\to A\). Every nonzero original polynomial \(Q(x)\) is a unit in \(K\), so \(Q(a)\) is a unit modulo \(I\). If \(b\in A\) lifts its inverse and \(Q(a)b=1-e\) with \(e\in I\), the exact inverse is \(b(1+e)\), because \[
 Q(a)b(1+e)=(1-e)(1+e)=1-e^2=1.
\] Thus evaluation extends uniquely to the original rational function field, with every fraction mapped by \[
 P(x)/Q(x)\longmapsto P(a)\,Q(a)^{-1}.
\] Cross multiplication proves that this map is independent of the chosen fraction representation. It reduces to \(\varphi\) and is unique after the lifts \(a_j\) have been chosen, since the images of every denominator inverse are then forced. Every polynomial uses only finitely many variables; no finiteness assumption on \(J\) is needed. If \(A/I=0\), then \(I=A\) and \(I^2=0\) force \(A=0\); the unique zero-ring map handles that case as well.

If \(K/k\) is separable algebraic, its finite intermediate extensions are finite étale \(k\)-algebras. Each has a unique lift against the given square-zero quotient. For nested finite intermediate fields, uniqueness makes the restriction of the larger lift equal to the smaller one. Their filtered union therefore gives a unique lift on \(K\). This is formal étaleness, which includes the existence required for formal smoothness.

For an arbitrary separable \(K/k\), take the actual directed system of finite generated \(k\)-subextensions \(K_i\subset K\). Each has a separating transcendence basis and a finite separable algebraic remainder. Composition of the two lifts just established gives formal smoothness of each \(K_i/k\), so \(H_1(L_{K_i/k})=0\). The filtered-colimit assertion used next is the specific cotangent assertion of algebra.tex:35556--35572. Its canonical presentation uses the polynomial ring \(P_i=k[K_i]\) on symbols for all elements of \(K_i\), its evaluation kernel \(J_i\), and \[
 J_i/J_i^2 \longrightarrow \bigoplus_{s\in K_i}K_i\,d[s].
\] Under an inclusion \(K_i\to K_j\), the symbol \([s]\) maps to the symbol of the same element in \(K_j\), and each polynomial relation maps to that same relation. Their colimits are \(P=k[K]\), its actual evaluation kernel \(J\), and the displayed presentation with \(K\) in place of \(K_i\). A polynomial relation and a finite sum of products of relations occur at a finite stage, so the relation modules also have colimit \(J/J^2\). Filtered exactness of modules identifies the first homology with the colimit of the first homologies, hence with zero. Since \(\Omega_{K/k}\) is projective over the field \(K\), the preceding criterion gives formal smoothness of \(K/k\). This is the source's argument; it does not assert that every filtered colimit of formally smooth algebras is formally smooth.

The positive-characteristic proposition now combines the original separability/reducedness criterion at 10360--10407 with the differential injection and formal-smoothness criteria proved here. Its maps retain the original fields. In particular the earlier reducedness criterion uses \[
 m:K\otimes_k k^{1/p}\longrightarrow K,\qquad
 \lambda\otimes\mu\longmapsto\lambda^p\mu^p,
 \qquad x^p=m(x)\otimes1.
\] If the tensor ring is reduced, \(m(x)=0\) forces \(x=0\). A relation \(\sum_i c_i a_i^p=0\) then lifts to the exact element \(\sum_i a_i\otimes c_i^{1/p}\) in its kernel. Independence of the \(a_i\) over \(k\) remains independence after the field extension \(k^{1/p}/k\), and forces all the actual \(c_i\) to vanish. The reverse separability implication, with its minimal-inseparable-degree argument, is the earlier reviewed criterion, not a new claim of this batch. The existing correction MC-STK-ERR-0683 supplies the missing comma and the exact homology criterion reference in the proposition's proof.

In characteristic zero \textbf{all five assertions of the proposition are true}. Separability, formal smoothness and homology vanishing follow above. The original differential injection over \(\mathbf Z\) follows from the formally smooth short exact sequence with that base. Geometric reducedness can also be checked directly: for a finite generated intermediate \(K_i/k\), let \(F=k(x_1,\ldots,x_d)\) be a separating rational subfield, so \(K_i/F\) is finite separable and hence finite étale. For every field extension \(k'/k\), the exact tensor ring \(F\otimes_k k'\) is the localization of \(k'[x_1,\ldots,x_d]\) at the original nonzero polynomials over \(k\), so is a domain. The ring \(K_i\otimes_k k'\) is finite étale over that domain and is flat; it injects into its scalar extension to the fraction field, which is a product of finite separable field extensions and is reduced. Thus it is reduced itself. Finally \(K\otimes_k k'\) is the filtered union of these tensor rings, with injective transitions since tensoring vector spaces by the field \(k'\) preserves injections. Every nilpotent lies in a stage and is zero there. This proves geometric reducedness with the actual tensor rings and localization maps.

Consequently the two reports alleging a missing predicate after ``If the characteristic of \(k\) is zero then'' are rejected. The five enumerated clauses are precisely the asserted conjunction, with ``and'' before the last clause. Replacing that assertion by mere equivalence would remove information: an equivalence by itself does not assert that any clause holds. The stronger, correct original statement is retained.

\subsubsection{The original triangular complete-intersection algebras}\label{algebra-proof-063-the-original-triangular-complete-intersection-algebras}

For 44511--44524 retain the finite generation of \(K/k\) and the finite type domain \(R\subset K\) with fraction field \(K\). The previously proved generic-point criterion says that the original map \(k\to R\) is smooth at its zero prime exactly when \(K/k\) is separable. Smoothness at that point gives a nonzero \(f\in R\) with \(R_f\) smooth, and \(K\) remains its fraction field. Conversely, if \(K\) is a localization of a smooth \(k\)-algebra, the smooth algebra is formally smooth and localization is formally étale. For the latter assertion, lift the map on the algebra and note that every specified denominator has unit image modulo a square-zero ideal, hence unit image in the target by the explicit inverse formula above; the lift to its localization is unique. Composition gives formal smoothness of \(K/k\), and the field criterion gives separability. This proves both directions with the original localization maps.

For 44526--44559 fix the source finite subset \(E\subset K\), its generated subfield \(L=k(E)\), the transcendence basis \(x_1,\ldots,x_d\), and \[
 F=k(x_1,\ldots,x_d),\qquad
 L=F[y_1,\ldots,y_r].
\] Keep \(d\), the number of transcendence generators, distinct from \(r\), the number of algebraic generators. Every intermediate \(F[y_1,\ldots,y_i]\) is a field: each generator is algebraic, and the inverse of a nonzero element of a finite-dimensional domain over a field exists by injectivity and then surjectivity of multiplication by that element.

Write the original monic minimal polynomial of \(y_i\) over the preceding field as \[
 T^{n_i}+\sum_{j=0}^{n_i-1}c_{ij}T^j.
\] Represent each actual coefficient \(c_{ij}\) by a polynomial \(C_{ij}\) over \(F\) in the preceding \(y\)'s, and keep its chosen full representative. Then \[
 P_i(Y_1,\ldots,Y_i)
   =Y_i^{n_i}+\sum_{j=0}^{n_i-1}C_{ij}(Y_1,\ldots,Y_{i-1})Y_i^j
\] is the exact triangular polynomial in the source. Induction gives the isomorphism \[
 F[Y_1,\ldots,Y_i]/(P_1,\ldots,P_i)
       \longrightarrow F[y_1,\ldots,y_i],
 \qquad Y_j\longmapsto y_j.
\] At the next step the preceding quotient is already the preceding field and the image of \(P_i\) is exactly the minimal polynomial. Polynomial division by that monic polynomial gives the inverse by the unique remainder of degree less than \(n_i\). Equivalently, each prefix quotient has basis \(\prod_{j\leq i}Y_j^{e_j}\) with \(0\leq e_j<n_j\). The remaining \(Y\) variables are polynomial variables, so multiplication by \(P_i\), monic in its new variable, is injective on the preceding prefix quotient: the leading coefficient of a nonzero polynomial remains the same nonzero coefficient after multiplication by a monic polynomial. The final quotient is the nonzero field \(L\). This proves regularity of the full sequence and the claimed quotient with the actual coefficient field \(F\).

Choose a nonzero \(h\in k[x_1,\ldots,x_d]\) as the full product of all denominators of all coefficients in the selected \(P_i\). Include also all denominators of the coefficients in polynomial expressions over \(F\) for each element of the original finite set \(E\) in the \(y_i\). There are finitely many such denominators and each is nonzero; for an empty list choose \(h=1\). Retain every factor of this actual product. Set \[
 D=k[x_1,\ldots,x_d,1/h],\qquad
 A=D[Y_1,\ldots,Y_r]/(P_1,\ldots,P_r).
\] The same monic division argument, now over \(D\), gives an exact free \(D\)-basis of \(A\) consisting of all \(\prod_iY_i^{e_i}\), \(0\leq e_i<n_i\). Each \(P_i\) remains a non-zero-divisor at its step by the monic leading-coefficient argument. In particular \(A\) is nonzero and torsion-free over the domain \(D\), and the exact localization map \[
 A\longrightarrow A\otimes_D F
       =F[Y_1,\ldots,Y_r]/(P_1,\ldots,P_r)=L
\] is injective: a nonzero coefficient in the free basis remains nonzero in \(F\). Its evaluation sends \(Y_i\) to \(y_i\), \(x_j\) to \(x_j\), and \(1/h\) to the original inverse in \(L\). The chosen coefficient denominators place every element of \(E\) in its image.

For a global polynomial presentation, retain the extra relation instead of absorbing the inverse into the notation. In \(k[X_1,\ldots,X_d,Z,Y_1,\ldots,Y_r]\) first impose \(h(X)Z-1\). Its quotient maps isomorphically to \(D[Y_1,\ldots,Y_r]\) by \(X_j\mapsto x_j\), \(Z\mapsto h(x)^{-1}\) and \(Y_i\mapsto Y_i\). The inverse exists because the original \(x_j\) are algebraically independent and the equation makes the image of \(h\) invertible. For every coefficient fraction \(a(X)/q(X)\) in a \(P_i\), the chosen product \(h\) is divisible by that actual \(q\); represent the fraction by \(a(X)(h(X)/q(X))Z\). This gives a polynomial lift \(\widetilde P_i\) with precisely the prescribed image \(P_i\), retaining the full coefficient and denominator factors. Therefore \[
 A\cong
 \frac{k[X_1,\ldots,X_d,Z,Y_1,\ldots,Y_r]}
 {(h(X)Z-1,\widetilde P_1,\ldots,\widetilde P_r)}.
\] The first relation is a nonzero element of a polynomial-ring domain, and every subsequent relation is a non-zero-divisor in the previous quotient because its image is the corresponding monic \(P_i\). The quotient is nonzero. This is the required global complete-intersection presentation. If \(d=0\), the actual nonzero constant \(h\) still has its inverse represented by \(Z\). If \(r=0\), the remaining sequence is empty and the same presentation gives \(D\). Neither boundary case changes the maps.

The set of actual global complete-intersection \(k\)-subalgebras of \(K\) just constructed is filtered under inclusion. Given two, take the union of their finite lists of algebra generators as \(E\); the construction supplies a third containing both as subalgebras of the original \(K\). Every element of \(K\) occurs by applying the construction to its singleton. Thus their filtered colimit with the actual inclusion maps is \(K\). No closure of complete intersections under an arbitrary tensor product is assumed.

In the separable case, \(L=k(E)\) is a finite generated separable subextension. Choose the smooth localization \(R_f\subset L\) above. Express every element of \(E\) as its original fraction in \(\operatorname{Frac}(R_f)=L\), and let \(g\) be the full product of their nonzero denominators. Then \((R_f)_g\subset L\subset K\) is a smooth \(k\)-algebra containing the actual \(E\). The same union-of-finite-generators argument proves filteredness and the exact colimit \(K\) in this case as well.

The retained MC-STK-ERR-0684 corrects ``minimum polynomial'' to ``minimal polynomial''. In MC-STK-ERR-0685, both occurrences of \(k(x_1,\ldots,x_r)\) must be \(k(x_1,\ldots,x_d)\), as the exact maps above show: \(d\) and \(r\) are independent quantities. The unit's existing change from ``\(P_1,\ldots,P_r\) is a regular sequence'' to ``are a regular sequence'' is a stylistic choice only; the singular wording can refer to the tuple and is not itself a mathematical defect. Retaining that already composed unit does not turn the stylistic wording into an additional defect.

\subsubsection{Report accounting and propagation}\label{algebra-proof-063-report-accounting-and-propagation}

Twenty physical reports have eleven dispositions. Eight previously composed units, MC-STK-ERR-0678 through MC-STK-ERR-0685, contribute eleven retained operations. Two new operations introduce the field extension in the differential criterion and change ``finitely generated sub \(k\)-extensions'' to ``finitely generated \(k\)-subextensions''. The two characteristic-zero predicate reports are one rejected group for the reason proved above. No optional extension group is added.

The corrected coefficient field, explicit prime-field differential base and localized target are used together in the exact presentation of \(\Omega_{K/\mathbf F_p}\). Its injectivity feeds the formal smoothness criterion, which feeds both the characteristic split and the smooth-colimit argument. The triangular presentation retains every original coefficient, every inverse factor and both independent generator counts. These are checked receiving arguments within the source interval. Full proof elaborations remain this separate source-linked editorial note; they do not replace the translated source. Additional paper mining and broader consequence synthesis remain after the core review.

\hypertarget{algebra_construct_flat_notes_20260929}{}
\subsection{Constructing flat ring maps: source review with the original maps}\label{algebra-proof-064-constructing-flat-ring-maps-source-review-with-the-original-maps}

Source: Stacks Project authors, algebra.tex:44566--44778, authority a04446e57ec1fbc252a871afcec7752fb2807b14, file SHA256 FA8BB92E58A4F78A2BD01B3B6A4A87DE0A0D279F5DD90641B574DD5FBFFFA4F3. The twelve existing correction units MC-STK-ERR-0686 through MC-STK-ERR-0697 are checked against the original constructions here. Full proof elaborations remain separate editorial material. This batch adds no optional extension and does not replace source translations.

\subsubsection{The original monogenic local extensions}\label{algebra-proof-064-the-original-monogenic-local-extensions}

Retain the original local ring \((R,\mathfrak m,k)\) and a monogenic extension \(k(\alpha)/k\). If \(\alpha\) is transcendental, take the source ring \[
 R'=R[x]_{\mathfrak mR[x]}.
\] The ideal \(\mathfrak mR[x]\) is prime because its quotient is the domain \(k[x]\). Thus \(R'\) is local, with maximal ideal \(\mathfrak mR'\), and reduction gives the exact isomorphism \[
 R'/\mathfrak mR'\longrightarrow k(\alpha),\qquad
 \overline{f(x)/s(x)}\longmapsto \overline f(\alpha)/\overline s(\alpha).
\] Here \(s\notin\mathfrak mR[x]\), so \(\overline s\) is a nonzero polynomial, and transcendence makes its value nonzero. Conversely every rational expression in \(\alpha\) comes from lifts of its numerator and nonzero denominator. Cross multiplication, with the nonzero reduced denominators retained, proves injectivity and independence of the lifts. The ring map \(R\to R'\) is flat because it is a polynomial extension followed by localization. It is injective directly: if a constant \(r\in R\) becomes zero, then \(sr=0\) for some polynomial \(s\notin\mathfrak mR[x]\); a coefficient of \(s\) lies outside \(\mathfrak m\) and is a unit, so its product with \(r\) being zero forces \(r=0\).

If \(\alpha\) is algebraic, retain its original monic minimal polynomial \[
 \overline F(T)=T^d+\sum_{i=1}^{d}\overline\lambda_iT^{d-i},
 \qquad
 F(T)=T^d+\sum_{i=1}^{d}\lambda_iT^{d-i},
 \qquad R'=R[T]/(F).
\] Each \(\lambda_i\) is the chosen lift in \(R\) of that actual coefficient. Monic division gives the free \(R\)-basis \(1,t,\ldots,t^{d-1}\), where \(t\) is the class of \(T\), so the original map is finite free and injective. Every maximal ideal of the integral \(R\)-algebra \(R'\) contracts to the maximal ideal \(\mathfrak m\). But \[
 R'/\mathfrak mR'\cong k[T]/(\overline F)
       \xrightarrow[\ T\mapsto\alpha\ ]{\ \sim\ } k(\alpha)
\] is a field. Thus \(\mathfrak mR'\) is the unique maximal ideal and \(R'\) is local. The inverse to the displayed field map represents each element by its unique polynomial remainder of degree less than \(d\). All original coefficients and the selected root are retained.

If the algebraic extension is separable, \(\overline F'(\alpha)\ne0\); hence \(F'(t)\) lies outside the unique maximal ideal of \(R'\) and is a unit. This makes the original monic quotient finite étale. The lifting assertion can be seen directly: for a square-zero ideal \(J\subset D\), a root modulo \(J\), and a chosen lift \(u\), the derivative \(F'(u)\) is a unit and the exact corrected root is \[
 u-F'(u)^{-1}F(u).
\] Taylor expansion retains the original coefficients; all terms with at least two copies of \(F(u)\in J\) vanish because \(J^2=0\). The resulting value of \(F\) is exactly zero. If two roots have the same reduction and differ by \(v\in J\), their polynomial values differ by \(F'(u)v\); invertibility forces \(v=0\). Thus the lift is unique. Finite presentation together with this formal étaleness gives étaleness, consistently with the preceding chapter criterion.

\subsubsection{The fixed base and the directed local colimit}\label{algebra-proof-064-the-fixed-base-and-the-directed-local-colimit}

The source category has objects \((k_i,R\to R_i,\phi_i)\), with \(\phi_i:R_i\to k_i\) inducing the prescribed residue-field isomorphism and \(k_i\subset K\). Its morphisms must be the \(R\)-algebra maps restored by MC-STK-ERR-0687. The residue-field diagram alone does not enforce this. For example, let \(R=k[\epsilon]/(\epsilon^2)\), let both objects have \(R_i=R\) with the identity structure map, and let \(\psi:R\to R\) send \(\epsilon\) to zero and fix \(k\). Its residue-field diagram commutes, but it is not an \(R\)-algebra map. A tower of these ring maps has colimit \(k\). The induced quotient \(R\to k\) is not flat: the inclusion \((\epsilon)\hookrightarrow R\) becomes the zero map from the nonzero module \((\epsilon)\otimes_R k\cong k\) into \(R\otimes_R k\cong k\). Thus the omitted base compatibility has a concrete effect on the asserted flatness.

MC-STK-ERR-0686 puts the displayed tuples in the same field-first order. Its shorter notation \((k_i,R_i,\phi_i)\) is read in the category of the fixed \(R\)-algebras: \(R_i\) carries the specified structure map. Restoring that already specified map gives \((k_i,R\to R_i,\phi_i)\); forgetting only its repeated display gives the inverse notation. No additional change of the mathematical object is needed.

Now let \(I\) be the original nonempty directed set of such objects. The residue diagram makes each transition local: an element is in the inverse image of the target maximal ideal exactly when its residue in \(k_i\) becomes zero in \(k_j\), and the field inclusion \(k_i\subset k_j\) is injective. Put \[
 R'=\operatorname*{colim}_{i\in I}R_i,\qquad k'=\bigcup_{i\in I}k_i.
\] The compatible \(\phi_i\) induce \(\phi':R'\to k'\). This map is surjective because an element of the union belongs to some \(k_i\) and lifts to \(R_i\). Its kernel is exactly \(\mathfrak mR'\): if an element represented by \(a_i\in R_i\) has zero residue in the union, injectivity of the field inclusion says \(\phi_i(a_i)=0\), hence \(a_i\in\mathfrak mR_i\). Conversely every element of \(\mathfrak mR'\) maps to zero. An element outside this kernel has a representative outside the corresponding maximal ideal of \(R_i\), so is already a unit at that stage and remains a unit. This proves that \(R'\) is local with maximal ideal \(\mathfrak mR'\) and with the exact stated residue field.

For every injection of \(R\)-modules \(M\hookrightarrow N\), flatness at each stage gives \(M\otimes_R R_i\hookrightarrow N\otimes_R R_i\). The identifications \[
 \operatorname*{colim}_i(M\otimes_R R_i)\longrightarrow
 M\otimes_R R',\qquad [m\otimes a_i]\longmapsto m\otimes[a_i],
\] and the analogous map for \(N\), are isomorphisms by the universal property of the tensor product. Their inverses send a tensor with a represented second factor to that stage and extend by additivity; finite tensor relations hold at a common stage. Filtered exactness preserves the injections. Thus the actual map \(R\to R'\) is flat. This argument needs flatness of each stage over the fixed \(R\); it does not infer flatness of arbitrary category transitions merely from the objects being flat over \(R\).

\subsubsection{The original transfinite construction and its endpoints}\label{algebra-proof-064-the-original-transfinite-construction-and-its-endpoints}

The coefficient base in MC-STK-ERR-0688 is necessary: the field \(K(x)\) must be generated \textbf{over \(k\)} by elements at most \(x\). Otherwise, for example, a well-order beginning with \(0\) in an extension of \(k=\mathbf Q(t)\) would initially generate only the prime field \(\mathbf Q\), which does not contain the given \(k\).

Keep the existing MC-STK-ERR-0691 choice of a well-order on the actual set \(K\) with a greatest element. It exists: select one element, well-order its complement, and put the selected element last. At a successor \(x\) of \(x'\), the correct equality is \[
 K(x)=K(x')(x),
\] as restored by MC-STK-ERR-0689. Square brackets would give only the generated algebra, which need not be a field: if \(x\) is transcendental over the preceding field, \(x^{-1}\) does not belong to that polynomial algebra. The monogenic construction above provides the required flat local extension \(R(x')\to R(x)\), with the specified extension of residue fields and the original new element \(x\).

At the least element, retain MC-STK-ERR-0690's actual initial pair \(R'(x)=R\), \(K'(x)=k\). An empty union of fields would not supply \(k\). At every nonzero limit position take the directed colimit \(R'(x)=\operatorname*{colim}_{y<x}R(y)\); the preceding proof gives its maximal ideal \(\mathfrak mR'(x)\) and residue field \(K'(x)=\bigcup_{y<x}K(y)\). Then use the same original monogenic construction for \[
 K(x)=K'(x)(x).
\] All transition maps produced in this construction are flat local and faithfully flat. For completeness, a flat local map \(A\to B\) is faithfully flat as follows. If \(M\ne0\), choose a nonzero cyclic submodule \(A/I\hookrightarrow M\). The ideal \(I\) is proper, so \(IB\) lies in the maximal ideal of \(B\). Hence \((A/I)\otimes_A B=B/IB\ne0\), and flatness embeds this nonzero module into \(M\otimes_A B\). Faithfulness together with flatness implies injectivity of \(A\to B\): if its kernel is \(J\), then the injection \(J\hookrightarrow A\) tensored with \(B\) has image zero, forcing \(J\otimes_A B=0\), then \(J=0\). Thus the constructed rings have the claimed compatible embeddings.

At the greatest element \(x_{\max}\), the original generated field is exactly \(K(x_{\max})=K\), and \(R(x_{\max})\) is the required ring. This verifies the actual retained endpoint correction. No assertion is made that an arbitrary well-order must have a largest stage. For example, a well-order of type \(\omega\) on the countable field \(\mathbf Q(t_1,t_2,\ldots)\) has only finite initial subsets, each generating a field of finite transcendence degree, whereas the full field has infinite transcendence degree.

The next lemma uses a continuous ordinal indexing of the same construction. To make the comparison exact, index the chosen ordered elements by \(x_\gamma\), \(\gamma<\lambda\), and define \(T_0=R\), \(T_{\gamma+1}=R(x_\gamma)\), and \(T_\delta=\operatorname*{colim}_{\gamma<\delta}T_\gamma\) for nonzero limits. At a limit \(\delta\), this is precisely the intermediate \(R'(x_\delta)\), before adjoining the next element; at a successor it is the original after-adjunction ring. These identifications commute with every transition by construction. The final \(T_\lambda\) equals the selected final \(R(x_{\max})\) when \(\lambda\) is a successor, and in the continuous formulation has residue field \(K\). This supplies the exact relation between the source's element-indexed and ordinal-indexed presentations; it does not silently assume that an after-adjunction ring at a limit position was already the preceding colimit.

\subsubsection{Finite étale local stages and descent of locality}\label{algebra-proof-064-finite-uxe9tale-local-stages-and-descent-of-locality}

Assume now that the original \(K/k\) is separable algebraic. Every successor residue extension in the preceding construction is finite separable, so every \(T_\gamma\to T_{\gamma+1}\) is the finite étale monic quotient already checked.

We use the following property of the finite étale rings that occur here. An \(R\)-algebra map \(A\to B\) between finite étale \(R\)-algebras is finite étale. It is finite because finitely many generators of \(B\) as an \(R\)-module also generate it as an \(A\)-module. It is finitely presented: if \(A=R[X_1,\ldots,X_s]/(f_1,\ldots,f_u)\), \(B=R[Y_1,\ldots,Y_t]/(g_1,\ldots,g_v)\), and the image of \(X_j\) is represented by \(h_j(Y)\), then the exact \(A\)-presentation of \(B\) is \[
 A[Y_1,\ldots,Y_t]/(g_1,\ldots,g_v,X_1-h_1(Y),\ldots,X_s-h_s(Y)).
\] Both inverse maps follow by substitution on these original generators. To verify formal étaleness over \(A\), start with an \(A\)-algebra square-zero lifting problem for \(B\). Formal étaleness of \(B/R\) gives a unique \(R\)-algebra lift. Its restriction to \(A\) and the prescribed \(A\)-structure have the same reduction; formal unramifiedness of \(A/R\) makes them equal. The lift is therefore \(A\)-linear and unique, proving the assertion by the finite-presentation criterion. If \(A\) and \(B\) are local, integrality implies that the maximal ideal of \(B\) contracts to the maximal ideal of \(A\). The map is flat local, hence faithfully flat and injective by the preceding argument.

Maintain inductively that each \(T_\gamma\) is a directed union of actual finite étale local \(R\)-subalgebras. At a successor write \(T_\gamma=\operatorname*{colim}_i R_i\), with those injections. The source uses a finite étale algebra \(R_{i,1}\) with \[
 T_{\gamma+1}=T_\gamma\otimes_{R_i}R_{i,1}.
\] The descent in this particular construction can be given explicitly. The successor is the original monic quotient \(T_\gamma[U]/(F)\). Its finitely many coefficients all belong to one \(R_i\); let \(F_i\in R_i[U]\) be that same polynomial and let \(R_{i,1}=R_i[U]/(F_i)\). Base change sends the class of \(U\) to the original class of \(U\), and coefficientwise evaluation gives the exact displayed isomorphism.

The residue field \(k_i\) of \(R_i\) embeds in the residue field \(k_\gamma\) of \(T_\gamma\). The image of \(\overline F_i\) over \(k_\gamma\) is the original irreducible separable minimal polynomial. A factorization over \(k_i\) would remain a nontrivial factorization over \(k_\gamma\), so \(\overline F_i\) is irreducible. A repeated irreducible factor, equivalently a nonconstant common divisor with its derivative, would likewise remain one after the field extension, so it is separable. The monogenic proof above now makes \(R_{i,1}\) a finite étale local ring. The same argument applies to the polynomial over every later \(R_{i'}\).

The locality assertion in the existing MC-STK-ERR-0692 also holds for an arbitrary finite étale descended algebra giving that exact base change. Indeed \(R_i\to T_\gamma\) is a filtered colimit of the finite étale local transition maps checked above. It is flat, and the residue-field injection makes it local, hence faithfully flat. Consequently \[
 R_{i,1}\longrightarrow T_\gamma\otimes_{R_i}R_{i,1}=T_{\gamma+1}
\] is faithfully flat. If a ring \(C\) has a faithfully flat map to a local ring \(D\), then \(C\) is local. To prove this, for each maximal ideal \(\mathfrak n\subset C\), the ring \(D/\mathfrak nD\) is nonzero by faithfulness. Choose a prime in it; its inverse image \(\mathfrak q\subset D\) contracts to \(\mathfrak n\), because its quotient contains the field \(C/\mathfrak n\) injectively. The unique maximal ideal \(\mathfrak m_D\) contains \(\mathfrak q\), so \(\mathfrak n\subseteq \mathfrak m_D\cap C\). The latter is proper, forcing equality. Every maximal ideal of \(C\) is therefore the same ideal. This proves the descent claim used by the correction.

For each \(i'\geq i\), retain the actual ring \[
 D_{i'}=R_{i'}\otimes_{R_i}R_{i,1}.
\] Its base change along the faithfully flat local \(R_{i'}\to T_\gamma\) is again \(T_{\gamma+1}\), so the same proof makes \(D_{i'}\) local. It is finite étale over \(R\) by base change and composition. Its map into \(T_{\gamma+1}\) is faithfully flat, hence injective, so it is an actual subalgebra. The tensor-colimit isomorphism \[
 \operatorname*{colim}_{i'\geq i}D_{i'}
      \longrightarrow T_\gamma\otimes_{R_i}R_{i,1}
\] sends \([a_{i'}\otimes b]\) to \([a_{i'}]\otimes b\). Its inverse is determined on simple tensors by a stage representing \(a\); all finite sums and tensor relations occur at a common later stage. This proves the successor assertion without suppressing a tensor factor or locality requirement.

At a nonzero limit \(\delta\), every finite subset of \(T_\delta\) belongs to some \(T_\gamma\), because there are finitely many representing stages and one larger ordinal below \(\delta\) contains all of them. At that stage the induction hypothesis places it in a finite étale local \(R\)-subalgebra. To verify directedness under inclusion, take the union of finite \(R\)-algebra generator sets of any two such subalgebras, and apply this finite-subset assertion; the resulting subalgebra contains both. The inclusions between these finite étale local \(R\)-algebras are faithfully flat by the argument above. Every element lies in one of them, so their directed union is exactly \(T_\delta\), with its actual ring operations. The initial stage \(R\) uses the identity finite étale extension. This proves the full lemma with the stated local rings, and verifies the receiving role of the existing locality correction. MC-STK-ERR-0693 repairs only the ``Since \ldots{} and we see'' coordination in the limit-stage sentence.

\subsubsection{Finite free realization with the original minimal polynomial}\label{algebra-proof-064-finite-free-realization-with-the-original-minimal-polynomial}

For 44674--44701 retain the original ring \(R\), prime \(\mathfrak p\), residue field \(F=\kappa(\mathfrak p)=\operatorname{Frac}(R/\mathfrak p)\), and finite field extension \(L/F\). The induction is on \([L:F]\), as MC-STK-ERR-0694 states. If the degree is one, the identity ring map suffices. A proper nontrivial intermediate field \(F\subsetneq L'\subsetneq L\) has both degrees strictly smaller. If \(R\to S'\) realizes \(L'/F\) and \(S'\to S\) realizes \(L/L'\), then the composite is finite free, its prime is \[
 (\mathfrak pS')S=\mathfrak pS,
\] and its residue-field isomorphism is the actual composite with \(L'\subset L\). A product of the two chosen free bases is a free \(R\)-basis, of rank \([L':F][L:L']=[L:F]\). These strict inequalities are the mathematical content of MC-STK-ERR-0695; the gerund correction is separate. If no such intermediate field exists and the degree exceeds one, choose \(\alpha\in L\setminus F\). Then \(F(\alpha)\) cannot be a proper intermediate field, so \(L=F(\alpha)\).

Keep this actual \(\alpha\) and its monic minimal polynomial \[
 P(T)=T^d+\sum_{i=0}^{d-1}a_iT^i.
\] Choose a single actual nonzero denominator \(\overline g\in R/\mathfrak p\) for the coefficients, lifted by \(g\in R\setminus\mathfrak p\), with \(a_i=\overline f_i/\overline g\). Such a common denominator is the full product of the finitely many original denominators; the numerators are multiplied by the corresponding products of the other factors. Keep those factors and retain the auxiliary element \(\beta=\overline g\alpha\) alongside \(\alpha\). Define \[
 Q(X)=X^d+\sum_{i=0}^{d-1}g^{d-i-1}f_iX^i\in R[X].
\] No negative exponent occurs because \(i<d\). Its image over \(F\) satisfies the complete comparison \[
 \overline Q(X)=\overline g^{\,d}P(X/\overline g),
 \qquad
 \overline Q(\beta)=\overline g^{\,d}P(\alpha)=0.
\] The nonzero factor \(\overline g^{\,d}\) and the inverse scale \(X/\overline g\) remain explicit. The inverse substitution gives \(P(T)=\overline g^{-d}\overline Q(\overline gT)\), so any factorization of \(\overline Q\) would factor \(P\). Thus \(\overline Q\) is irreducible and \[
 F[X]/(\overline Q)\xrightarrow{\sim}L,\quad [X]\longmapsto\overline g\alpha,
 \qquad \alpha\longleftrightarrow\overline g^{-1}[X].
\] This proves the exact relationship behind the source's generator replacement while keeping its original field, original root and polynomial available.

Take \(S=R[X]/(Q)\), with free basis \(1,x,\ldots,x^{d-1}\). Modulo \(\mathfrak p\) this is free over the domain \(R/\mathfrak p\), so it injects into its localization \[
 S/\mathfrak pS\longrightarrow
 (S/\mathfrak pS)\otimes_{R/\mathfrak p}F
       \cong F[X]/(\overline Q)\cong L.
\] The injection follows by the free-basis coefficients, each of which embeds into \(F\). Thus \(S/\mathfrak pS\) is a domain and \(\mathfrak q=\mathfrak pS\) is prime. Its fraction field is exactly \(L\): the displayed localization is already the field \(L\), and it contains both the embedded domain and inverses of its nonzero coefficient elements; its element \(\overline g^{-1}x\) is the original \(\alpha\). The induced map on residue fields sends each actual quotient \(u/v\) to the quotient of its two displayed images, with nonzero denominator preserved, and has the inverse determined by the coefficients in \(F\) and \(\alpha=\overline g^{-1}x\). This proves the stated extension isomorphism, not just equality of degrees.

\subsubsection{The cardinal bound and the actual relation witnesses}\label{algebra-proof-064-the-cardinal-bound-and-the-actual-relation-witnesses}

For 44703--44769 retain \(A\), the flat \(A\)-algebra \(B\), and the exact infinite cardinal \(\kappa=\max(|A|,\aleph_0)\). Given a subalgebra \(E=E_0\subset B\) of size at most \(\kappa\), the set \(S_k\) of tuples \[
 (n,a_1,\ldots,a_n,e_1,\ldots,e_n),\qquad
 n\geq1,\quad \sum_i a_ie_i=0
\] has cardinality at most \(\kappa\). For each fixed \(n\), it is a subset of \(A^n\times E_k^n\), with cardinality at most \(\kappa\); the countable union over \(n\) still has cardinality at most \(\kappa\).

Flatness and the equational criterion at algebra.tex:8994--9034 give exactly the source witnesses \[
 m_s\geq0,\quad b_{s,j}\in B,\quad a_{s,ij}\in A,\qquad
 e_i=\sum_{j=1}^{m_s}a_{s,ij}b_{s,j},\qquad
 \sum_i a_i a_{s,ij}=0\quad(1\leq j\leq m_s).
\] All coefficients, indices and zero cases remain. If \(m_s=0\), the first formula states that every \(e_i=0\), and no witness elements need be adjoined. For all other \(s\), adjoining the finitely many \(b_{s,j}\) for each of at most \(\kappa\) relations adds at most \(\kappa\) elements. The \(A\)-algebra generated by these elements and \(E_k\) consists of evaluations of finite polynomials in at most \(\kappa\) generators with coefficients in \(A\). The set of finite monomials has cardinality at most \(\kappa\), and so does the set of finite sums of such terms. Hence the actual \(E_{k+1}\) has cardinality at most \(\kappa\). The countable union \(B'=\bigcup_{k\geq0}E_k\) has the same bound.

Every finite relation \(\sum_i a_i b_i'=0\) in \(B'\) has all its entries in one \(E_k\), where it is that identical relation, since the rings are actual subalgebras of \(B\). Its chosen witnesses belong to \(E_{k+1}\subset B'\), and their displayed equalities hold unchanged. Thus every relation in \(B'\) is trivial in the precise equational sense, and \(B'\) is flat over \(A\). To retain the direction actually used from the criterion: if \(I=(a_1,\ldots,a_n)\subset A\) and an element \(\sum_i a_i\otimes b_i'\) of \(I\otimes_A B'\) maps to zero in \(B'\), the witness equalities give \[
 \sum_i a_i\otimes b_i'
 =\sum_{i,j}a_i\otimes a_{s,ij}b_{s,j}
 =\sum_j\left(\sum_i a_i a_{s,ij}\right)\otimes b_{s,j}
 =0.
\] The finitely generated ideal criterion then proves flatness.

If \(B\) is faithfully flat over \(A\), each such \(B'\) is faithfully flat too. For a maximal ideal \(\mathfrak n\subset A\), faithfulness gives \(\mathfrak nB\ne B\). The inclusion of the actual subalgebra implies \(\mathfrak nB'\subseteq\mathfrak nB\); hence \(1\notin\mathfrak nB'\) and \(B'/\mathfrak nB'\ne0\). Together with flatness, this gives faithfulness: a nonzero cyclic submodule \(A/I\) of any nonzero module has \(I\subseteq\mathfrak n\), so \(B'/IB'\ne0\), and flatness preserves its injection into the tensor of the full module. This includes the parenthetical assertion of the original lemma.

Finally, every finite subset of \(B\) generates an \(A\)-subalgebra of cardinality at most \(\kappa\). The algebra generated by two subalgebras of that size has that size bound by the same finite-polynomial count. Applying the witness construction to it gives a flat subalgebra containing both. Thus the actual collection of small flat subalgebras is directed under inclusion, contains every element of \(B\) at some stage, and has colimit \(B\) with the original maps. If \(B\) itself already satisfies the bound, it is the terminal member and the assertion remains valid. MC-STK-ERR-0696 corrects only ``choicse''; MC-STK-ERR-0697 corrects the word order of the same cardinal bound. The fuller cardinal and relation arguments here remain editorial.

\subsubsection{Report accounting and propagation}\label{algebra-proof-064-report-accounting-and-propagation}

There are twenty-one physical reports in this interval and thirteen review groups. Twelve existing units, MC-STK-ERR-0686 through MC-STK-ERR-0697, supply fifteen retained operations. The two reports requiring ``were'' in place of ``was'' are grouped as an optional stylistic proposal and rejected as a mandatory defect; they do not change the construction or its set-size issue.

The base-compatible category maps, generated-over-\(k\) residue fields, field adjunctions, initial pair and final selected stage are checked together at the construction's receiving maps. That constructed ring feeds the next lemma through the explicit continuous ordinal reindexing. Its finite étale local stages and exact tensor base changes validate the existing locality correction. The finite-free field realization retains the original minimal polynomial and all scaling factors in the comparison; the small-flat-subalgebra lemma retains every original relation witness and the exact cardinal bound. These checks produce no new optional consequence group. Broader paper mining and synthesis remain after the core correction review.

\hypertarget{algebra_cohen_notes_20260929}{}
\subsection{Cohen structure: original-source corrections and complete receiving arguments}\label{algebra-proof-065-cohen-structure-original-source-corrections-and-complete-receiving-arguments}

Source: Stacks Project authors, algebra.tex:44779--45127, original authority a04446e57ec1fbc252a871afcec7752fb2807b14, file SHA256 FA8BB92E58A4F78A2BD01B3B6A4A87DE0A0D279F5DD90641B574DD5FBFFFA4F3. This note validates the original rings, inverse systems, ideals and maps. The eight existing units MC-STK-ERR-0698 through MC-STK-ERR-0705 are retained. Additional paper mining and optional extensions remain deferred; this fuller argument is editorial material, not a replacement of the translated source.

\subsubsection{Completeness, proper quotients and finite products}\label{algebra-proof-065-completeness-proper-quotients-and-finite-products}

The source definition uses the actual map \[
 R\longrightarrow\varprojlim_{n\geq1}R/\mathfrak m^n,\qquad
 r\longmapsto(r\bmod\mathfrak m^n)_n.
\] Being an isomorphism includes both surjectivity and the exact kernel assertion \(\bigcap_n\mathfrak m^n=0\). For an Artinian local ring, some power of the maximal ideal is zero, so the inverse system is eventually the unchanged ring; the displayed map is an isomorphism.

The finite-generation qualification in the footnote is essential to its cited completeness argument. For a finitely generated ideal \(I\) and an \(R\)-module \(M\), retain the exact completion comparison from algebra.tex:22644--22678: \[
 I^n\widehat M
   =\ker(\widehat M\longrightarrow M/I^nM),\qquad
 \widehat M/I^n\widehat M\cong M/I^nM.
\] To see the stated kernel map, choose the actual finite generators \(f_1,\ldots,f_r\) of \(I^n\). Completion of the surjection \(M^{\oplus r}\to I^nM\), with coordinates mapped by \((f_1,\ldots,f_r)\), is onto. Its image in \(\widehat M\) is exactly \(I^n\widehat M\), and its target identifies with \(\varprojlim_{m\geq n}I^nM/I^mM\), which is the kernel of the displayed reduction. Taking inverse limits of the quotient isomorphisms proves completeness of \(\widehat M\). All maps send represented elements to the same residue classes. This is the earlier source completion-surjectivity argument; the finite generating list is not discarded.

If \(R/I\) is Noetherian and \(I\) is finitely generated, the original completion \(\widehat R\) is Noetherian by algebra.tex:23010--23062. The proof retains the map \[
 (R/I)[T_1,\ldots,T_t]\twoheadrightarrow
 \bigoplus_{n\geq0}I^n/I^{n+1},\qquad T_i\longmapsto[f_i]
\] and the actual initial forms of generators of an arbitrary ideal. In a complete ring, homogeneous generators \(\overline g_j\) of degrees \(d_j\) for the initial-form ideal are lifted to \(g_j\in J\cap I^{d_j}\). Successive correction of an original \(x\in J\) produces coefficients \[
 a_{j,n}\in I^{\max(0,n-d_j)},\qquad
 x-\sum_{n=0}^{N}\sum_j a_{j,n}g_j\in I^{N+1}.
\] For each \(j\), the sum \(A_j=\sum_{n\geq0}a_{j,n}\) converges, and separatedness gives the full identity \(x=\sum_j A_jg_j\). Thus \(J\) is finitely generated. Applied to the completion, the preceding exact quotient comparison supplies the required complete ring with the same Noetherian residue quotient. Applied to a complete local \(R\) with finitely generated \(\mathfrak m\), it makes \(R\) itself Noetherian. MC-STK-ERR-0705 only clarifies the footnote's attachment of these two already valid assertions.

For a Noetherian complete local \(R\), a finite \(R\)-algebra \(S\) is Noetherian and the finite-module completion comparison identifies its actual map \[
 S\longrightarrow\varprojlim_n S/\mathfrak m^nS
\] as an isomorphism. Here this is the original finite-module theorem, not an assertion about arbitrary modules. There are finitely many maximal ideals \(\mathfrak n_1,\ldots,\mathfrak n_s\) of \(S\), all over \(\mathfrak m\), because \(S/\mathfrak mS\) is a finite-dimensional algebra over \(R/\mathfrak m\). Put \(J=\bigcap_i\mathfrak n_i\). The radical of the Artinian ring \(S/\mathfrak mS\) is nilpotent, so for some actual integer \(c\geq1\), \[
 J^c\subseteq\mathfrak mS\subseteq J,\qquad
 J^{cn}\subseteq\mathfrak m^nS\subseteq J^n.
\] The two filtrations are cofinal, and the identity on \(S\) induces inverse maps on their inverse limits. The maximal ideals are pairwise comaximal, so \(J^n=\bigcap_i\mathfrak n_i^n\), and the Chinese remainder maps give \[
 S\cong\varprojlim_n S/J^n
   \cong\prod_{i=1}^{s}\varprojlim_n S/\mathfrak n_i^n
   \cong\prod_{i=1}^{s}\widehat{S_{\mathfrak n_i}}.
\] Each map is induced by the original quotient or localization map. In the last comparison, every element outside \(\mathfrak n_i\) is already a unit modulo \(\mathfrak n_i^n\): that quotient has a unique maximal ideal with nilpotent image, so the inverse is lifted by the finite geometric-series identity. Thus localization does not change that quotient. Each completion is Noetherian and complete for its finitely generated maximal ideal by the preceding completion statements. This verifies the source finite-product conclusion and its original comparison at 23134--23174. If \(S=0\), the product is empty and equals the zero ring.

For \(S=R/I\) with \(I\) proper, the unique maximal ideal is \(\mathfrak m/I\), its powers are the images of \(\mathfrak m^n\), and the finite-module comparison proves completeness for exactly this ideal. Hence the quotient is a Noetherian complete local ring. For \(I=R\), the quotient is zero and has no maximal ideal, so it is not a local ring under the source convention. MC-STK-ERR-0698's proper-ideal restriction corrects that boundary overclaim; it is not imposed on the finite-product statement. MC-STK-ERR-0699 only removes the doubled conditional wording in the latter statement.

\subsubsection{The original Cohen ring and all its finite quotients}\label{algebra-proof-065-the-original-cohen-ring-and-all-its-finite-quotients}

For a coefficient subring \(C\subset R\), the source conditions say that \(C\) is complete local, that its maximal ideal is \(C\cap\mathfrak m\), that its residue field maps isomorphically to \(R/\mathfrak m\), and that this maximal ideal equals \(pC\), with \(p\) the characteristic of that residue field. In residue characteristic zero \(p=0\), so \(C\) is a field. In equal positive characteristic \(pR=0\), again \(pC=0\), so \(C\) is a field. A local ring with residue characteristic zero contains \(\mathbf Q\): each nonzero integer has nonzero residue and is therefore a unit, giving the original inverse-denominator map into \(R\). Conversely any field contained in \(R\) injects into its residue field, since its nonzero elements remain units, and its characteristic agrees with that of the residue field.

For a complete local ring \(C\) with maximal ideal \(pC\), separatedness shows that every nonzero \(a\in C\) has a largest integer \(v\geq0\) with \(a\in p^vC\); write \(a=p^vu\). The actual factor \(u\) is a unit, since otherwise \(a\in p^{v+1}C\). If no power of \(p\) is zero, then products \[
 (p^vu)(p^ww)=p^{v+w}uw
\] are nonzero. Thus \(C\) is a domain, every nonzero ideal is generated by the smallest power occurring among its nonzero elements, and the only prime ideals are zero and \((p)\). It is a discrete valuation ring with the actual uniformizer \(p\), and it is complete by assumption. If instead \(p^e=0\), \(p^{e-1}\ne0\), then the finite filtration by \(p^iC\) has successive quotients of dimension at most one over \(C/pC\), because each is generated by the original \(p^i\). It gives finite length, so \(C\) is Artinian. This verifies the three cases without discarding unit factors or the nilpotent cases. MC-STK-ERR-0700 clarifies the wording that the uniformizer of a Cohen ring is the image of an actual prime integer.

For existence, retain the exact flat local extension \[
 \mathbf Z_p\longrightarrow R_*,\qquad
 \mathfrak m_{R_*}=pR_*,\qquad R_*/pR_*=k
\] constructed in the preceding section; \(R_*\) here names that auxiliary source ring, while \(R\) below remains the given ring of the structure theorem. Let \(\Lambda=\widehat{R_*}=\varprojlim_n R_*/p^nR_*\). The completion argument above gives \[
 \Lambda/p^n\Lambda\cong R_*/p^nR_*
\] for every \(n\), completeness, and Noetherianity, since \(pR_*\) is principal and its quotient is a field. The isomorphisms are induced by the original projection from the inverse limit. The maximal ideal of \(\Lambda\) is \(p\Lambda\): an element with nonzero residue has an inverse modulo \(p\), and the error is in \(p\Lambda\), so the convergent geometric series lifts the inverse in the \(p\)-adically complete ring.

Flatness of \(R_*/\mathbf Z_p\) makes multiplication by \(p\) injective on \(R_*\). It also proves injectivity on the actual completion without a suppressed limit argument. If \(a=(a_n)\in\Lambda\) satisfies \(pa=0\), then \(pa_{n+1}=0\) in \(R_*/p^{n+1}R_*\). The original injectivity on \(R_*\) implies \(a_{n+1}\in p^nR_*/p^{n+1}R_*\), so its projection \(a_n\) is zero. Every coordinate vanishes, hence \(a=0\). As the residue field is nonzero, \(p\) is a nonzero nonunit and no power vanishes. The preceding exact valuation argument now proves that \(\Lambda\) is a complete discrete valuation ring with uniformizer \(p\). This supplies the ring required by the source existence lemma.

The cited flatness argument in 44895--44898 is valid as well: with source ring \(\mathbf Z_p\), target \(\Lambda\), ideal \((p)\), and module \(\Lambda\), both rings are Noetherian, the target module is finite over itself, and every reduction is flat because it is the base change of the original flat \(\mathbf Z_p\to R_*\). The ideal lies in the target maximal ideal, so the precise final clause of algebra.tex:23807--23843 applies.

For every \(n\geq1\), put \(A_n=\mathbf Z/p^n\mathbf Z\) and \(B_n=\Lambda/p^n\Lambda\). The original map \(A_n\to B_n\) is flat. Indeed every ideal of \(A_n\) is \((p^i)\), \(0\leq i\leq n\), and for \(i<n\) the tensor comparison is \[
 (p^i)\otimes_{A_n}B_n
   \cong B_n/p^{n-i}B_n
   \longrightarrow B_n,\qquad [b]\longmapsto p^ib.
\] It is injective: if a lift \(\widetilde b\in\Lambda\) satisfies \(p^i\widetilde b\in p^n\Lambda\), injectivity of multiplication by \(p^i\) in \(\Lambda\) gives \(\widetilde b\in p^{n-i}\Lambda\). The zero ideal case has zero source. The ideal criterion proves flatness.

At \(n=1\), \(B_1=k\) is formally smooth over \(\mathbf F_p\), by the separable-field criterion already verified in ALGEBRA\_FIELD\_FORMAL\_SMOOTH\_NOTES\_20260929.md; the prime field is perfect. In the induction from \(n\) to \(n+1\), keep the exact square-zero ideal \[
 I=(p^n)\subset A_{n+1},\qquad I^2=0
\] because \(2n\geq n+1\). Its quotient map is precisely \(A_n\to B_n\), and the unquotiented map is flat by the preceding calculation. Thus all three hypotheses of the square-zero formal-smoothness lemma at 38099--38166 are satisfied.

The receiving proof uses the already corrected tensor base in MC-STK-ERR-0449 and the splitting directions of ALGEBRA-RECON-1083. For a polynomial presentation \(P\to S\) with kernel \(J\), keep \[
 E=J/J^2,\qquad F=\Omega_{P/A}\otimes_P S,\qquad d:E\to F.
\] The quotient by \(I\) supplies a left inverse, and its lifted map \(r_0:F\to E\) satisfies \(\Delta=1_E-r_0d\), \(\Delta(E)\subset IE\). Since \(I^2=0\) and \(\Delta\) is \(S\)-linear, \(\Delta^2=0\). The exact correction \(r=r_0+\Delta r_0\) then satisfies \[
 rd=(1_E+\Delta)(1_E-\Delta)=1_E.
\] This is a left inverse to the conormal differential, with its original domain and codomain; it is the split-conormal criterion for formal smoothness. The complete earlier proof is in ALGEBRA\_FORMAL\_SMOOTH\_NOTES\_20260928.md, ``Projective differentials and the direction of the splittings''. It proves formal smoothness of every actual \(A_n\to B_n\) used in the next construction.

\subsubsection{The compatible coefficient-ring lifting maps}\label{algebra-proof-065-the-compatible-coefficient-ring-lifting-maps}

Retain the source complete local ring \((R,\mathfrak m)\) and residue field \(\kappa\), without imposing Noetherianity for the coefficient-ring assertion. If \(\operatorname{char}(\kappa)=0\), the canonical \(\mathbf Q\)-map into \(R\) was proved above. Begin with \(\varphi_1:\kappa\to R/\mathfrak m=\kappa\) equal to the identity. For \(n\geq2\), the kernel of \[
 R/\mathfrak m^n\longrightarrow R/\mathfrak m^{n-1}
\] is \(\mathfrak m^{n-1}/\mathfrak m^n\), whose square is zero since \(2n-2\geq n\). Formal smoothness of the original \(\kappa/\mathbf Q\) therefore lifts the already chosen \(\varphi_{n-1}\) to a \(\mathbf Q\)-map \(\varphi_n\). This choice explicitly gives \(\pi_n\varphi_n=\varphi_{n-1}\). Hence the map \[
 \kappa\longrightarrow\varprojlim_n R/\mathfrak m^n=R,\qquad
 a\longmapsto(\varphi_n(a))_n
\] is a ring map splitting the original residue map. It is injective because its residue composite is the identity. Its image is a field with zero intersection with \(\mathfrak m\), complete for its zero maximal ideal, and has exactly the required residue isomorphism.

If \(\operatorname{char}(\kappa)=p>0\), retain the source Cohen ring \(\Lambda\) with the fixed identification \(\Lambda/p\Lambda=\kappa\). Since \(p\in\mathfrak m\), \(p^nR\subset\mathfrak m^n\), so every \(R/\mathfrak m^n\) is an \(A_n\)-algebra. Begin with that fixed identification for \(\varphi_1\). At stage \(n\geq2\), the actual already constructed arrow used in the lifting square is \[
 \psi_{n-1}:
 \Lambda/p^n\Lambda
 \xrightarrow{q_n}\Lambda/p^{n-1}\Lambda
 \xrightarrow{\varphi_{n-1}}R/\mathfrak m^{n-1}.
\] MC-STK-ERR-0701 explicitly introduces this composite before reusing the shorter arrow label. The square commutes as an \(A_n\)-algebra square. Formal smoothness of \(A_n\to B_n\), proved with all its flatness hypotheses above, gives a map \(\varphi_n:B_n\to R/\mathfrak m^n\) satisfying the exact identity \[
 \pi_n\varphi_n=\varphi_{n-1}q_n.
\] Thus the original inverse-limit map is \[
 \varphi:\Lambda=\varprojlim_n\Lambda/p^n\Lambda
       \longrightarrow\varprojlim_n R/\mathfrak m^n=R,\qquad
 a\longmapsto\bigl(\varphi_n(a\bmod p^n)\bigr)_n.
\] It is a ring map and its residue map is the selected identity on \(\kappa\).

Its image \(C=\varphi(\Lambda)\) is the actual coefficient subring. To verify every defining condition, let \(J=\ker\varphi\). The nonzero residue map makes \(J\) proper, and the Cohen-ring ideal classification gives either \(J=0\) or \(J=(p^e)\), \(e\geq1\). In the first case \(C\cong\Lambda\) is complete local; in the second \(C\cong\Lambda/(p^e)\) is Artinian local and complete for its nilpotent maximal ideal. The kernel of the induced residue map on \(C\) is exactly \(pC\), because an element \(\varphi(a)\) has zero residue precisely when \(a\in p\Lambda\). Therefore \(C\cap\mathfrak m=pC\), and \(C/pC\to\kappa\) is the required isomorphism. When \(pR=0\), \(J=(p)\), so \(C\) is a field; when no power of \(p\) vanishes in \(R\), \(J=0\); when \(p\) has exact nilpotence exponent \(e\), \(J=(p^e)\). All original cases are covered without treating an image as automatically complete.

\subsubsection{Power-series surjectivity with the respective ideals}\label{algebra-proof-065-power-series-surjectivity-with-the-respective-ideals}

Assume now that \(\mathfrak m=(y_1,\ldots,y_n)\) is finitely generated. The source ring for the presentation is \(A=\Lambda[[x_1,\ldots,x_n]]\), where in the equal-characteristic case \(\Lambda\) is the coefficient field, and in positive residue characteristic one may use the original Cohen ring mapping to \(R\), even when its image is an Artinian coefficient ring. Retain every coefficient and monomial. Define \[
 \Phi:A\longrightarrow R,\qquad
 \sum_{I\in\mathbf N^n}a_Ix^I
 \longmapsto
 \lim_{N\to\infty}\sum_{|I|<N}\varphi(a_I)y^I.
\] The terms of total degree at least \(N\) lie in \(\mathfrak m^N\), so completeness and separatedness of \(R\) give the limit. Modulo \(\mathfrak m^N\) only finitely many degrees contribute; finite coefficient convolution proves additivity and multiplication, hence \(\Phi\) is a ring map. It sends every original variable \(x_i\) to the specified \(y_i\).

Write \(\mathfrak a=(x_1,\ldots,x_n)\subset A\). The exact ideal comparison is \(\mathfrak aR=(y_1,\ldots,y_n)=\mathfrak m\), and \[
 A/\mathfrak a=\Lambda\longrightarrow R/\mathfrak m
\] is onto. The source is \(\mathfrak a\)-adically complete and the target is \(\mathfrak m\)-adically complete. This is MC-STK-ERR-0702's explicit wording. The original wording also admits the usual reading of \(\mathfrak a\)-adic completeness on the \(A\)-module \(R\), where \(x_i\) acts by \(y_i\); the retained edit is a clarification of the respective ideals, not a newly false completeness theorem.

For exact surjectivity, choose for a given \(r\in R\) a coefficient lifting its residue. Suppose polynomials through degree \(N-1\) have been chosen with residual error in \(\mathfrak m^N\). This ideal is generated by the original degree-\(N\) monomials \(y^I\). Write the error as a finite sum \(\sum_{|I|=N}b_Iy^I\), and lift each residue \(b_I\bmod\mathfrak m\) by a coefficient \(a_I\in\Lambda\). The degree-\(N\) polynomial \(\sum a_Ix^I\) corrects the error modulo \(\mathfrak m^{N+1}\). Iterating preserves every selected term and gives a series \(F\) with \(r-\Phi(F)\in\bigcap_N\mathfrak m^N=0\). Thus \(\Phi\) is onto and the original \(R\) is exactly \(A/\ker\Phi\). If \(n=0\), the same residue construction has no higher-degree terms and \(\mathfrak m=0\); it says that the coefficient map itself is onto the field \(R\).

\subsubsection{Regular power-series rings and the two finite-extension cases}\label{algebra-proof-065-regular-power-series-rings-and-the-two-finite-extension-cases}

For the uses in 45022--45122, retain either a field \(k\) or the actual Cohen ring \(\Lambda\). The ring \(k[[X_1,\ldots,X_d]]\) is complete for \((X_1,\ldots,X_d)\). The ring \(\Lambda[[X_1,\ldots,X_d]]\) is complete for \(\mathfrak q=(p,X_1,\ldots,X_d)\): modulo \(\mathfrak q^N\), the coefficient of \(X^I\), \(|I|<N\), is specified modulo \(p^{N-|I|}\); coefficients in degrees at least \(N\) contribute zero. Compatible systems recover each full coefficient by the original \(p\)-adic completeness of \(\Lambda\), giving inverse maps between the full series ring and this inverse limit.

Both series rings are Noetherian by the completion theorem applied to the polynomial ring over their Noetherian coefficient ring and the finitely generated variable ideal. Their units are exactly the series with unit constant term. For \(F=\sum a_IX^I\) with \(a_0\) a unit, the inverse coefficients are determined without omitting terms by \[
 b_0=a_0^{-1},\qquad
 b_I=-a_0^{-1}\sum_{0<J\leq I}a_Jb_{I-J}\quad(I\ne0).
\] If the constant term is a nonunit, the series lies in the stated maximal ideal, respectively \((X)\) or \((p,X)\). They are domains: the first nonzero total-degree homogeneous parts of two nonzero series multiply to a nonzero polynomial over the original coefficient domain, so their product series is nonzero. Prime chains obtained by successively adjoining the actual parameters give dimension at least \(d\), respectively \(d+1\); the bound by the number of generators of the maximal ideal gives the reverse inequalities. The parameter classes are independent modulo the square of the maximal ideal by coefficient comparison, including the nonzero class of \(p\) modulo \(p^2\). Thus these are the original complete regular local rings of dimensions \(d\) and \(d+1\).

The remark on universal catenarity follows with its original hypotheses. A regular local ring is Cohen--Macaulay by algebra.tex:25731--25747, and a Noetherian Cohen--Macaulay ring is universally catenary by 25606--25638. The latter proof computes the length of a maximal chain between \(\mathfrak p\) and \(\mathfrak q\) as \(\dim A_{\mathfrak q}-\dim A_{\mathfrak p}\), using the Cohen--Macaulay full-chain length theorem, and applies the same argument to polynomial algebras. A quotient \(R=A/J\) inherits universal catenarity because every finite type \(R\)-algebra, with its original presentation, is also a finite type \(A\)-algebra. The Cohen presentation above therefore has exactly the consequence asserted in the source remark.

For a complete regular local \(R\) with coefficient field \(k\), choose the original \(f_1,\ldots,f_d\) lifting a basis of \(\mathfrak m/\mathfrak m^2\). Nakayama's lemma makes them generators of \(\mathfrak m\), and regularity says \(d=\dim R\). The just-proved power-series construction gives a surjective \(k\)-algebra map \[
 k[[x_1,\ldots,x_d]]\longrightarrow R,\qquad x_i\longmapsto f_i.
\] If the kernel contained a nonzero \(F\), then \(F\) is in the source maximal ideal and is not in its zero minimal prime. The exact one-equation theorem, algebra.tex:14693--14710, gives source quotient dimension \(d-1\), forcing \(\dim R\leq d-1\), a contradiction. Thus the map is an isomorphism. This verifies both clauses of the regular-complete lemma, with the given \(k\)-algebra structure retained. Its wording about variable-adic completeness has the same legitimate module-action reading as above and does not require a new source correction.

Finally let the original \(R\) be a Noetherian complete local domain. Its coefficient subring is a field or a Cohen ring: an Artinian coefficient ring with a nonzero nilpotent \(p\) cannot embed in a domain. In Case I retain the coefficient field \(k\), set \(d=\dim R\), and choose the original \(x_1,\ldots,x_d\in\mathfrak m\) generating an ideal of definition \(I\). In Case II retain the coefficient Cohen ring \(\Lambda\), set \(d+1=\dim R\), and use its actual prime \(p\). It is a nonzero nonunit in the domain \(R\), so 14693--14710 gives \(\dim R/pR=d\). Choose the source \(x_1,\ldots,x_d\) whose images generate an ideal of definition in \(R/pR\); their lifts lie in \(\mathfrak m\). Then \(I=(p,x_1,\ldots,x_d)\) is an ideal of definition in \(R\).

In either case \(I\subseteq\mathfrak m\) and some \(\mathfrak m^c\subseteq I\); hence \(\mathfrak m^{cn}\subseteq I^n\subseteq\mathfrak m^n\). The identity map gives both inverse completion comparisons, so \(R\) is complete and separated for the actual \(I\)-adic filtration. Define, respectively, \[
 A=k[[X_1,\ldots,X_d]],\quad I_0=(X_1,\ldots,X_d),
\] or \[
 A=\Lambda[[X_1,\ldots,X_d]],\quad I_0=(p,X_1,\ldots,X_d).
\] The continuous map \(A\to R\) sends \(X_i\) to the original \(x_i\), and the coefficient ring by its original embedding. Then \(I_0R=I\). The ring \(R/I\) is Noetherian local of dimension zero, hence Artinian, and it has a finite filtration with factors its residue field \(k=A/I_0\). Thus it is a finite \(A/I_0\)-module.

Here is the exact finite-module lifting underlying algebra.tex:22859--22883 in this application. Choose elements \(e_1,\ldots,e_t\in R\) whose images generate \(R/I\) over \(k\). Given \(r\in R\), choose a first coefficient vector over \(A\) realizing its class modulo \(I\). An error in \(I^N=I_0^NR\) is a sum of degree-\(N\) monomials in the original generators of \(I_0\) times elements of \(R\). Lift the latter elements modulo \(I\) using the same \(e_j\). This gives a correction vector \(c_N\in I_0^N A^t\), leaving error in \(I^{N+1}\). The sum of the vectors converges in the complete module \(A^t\); its image under \((e_1,\ldots,e_t)\) is \(r\) by separatedness of \(R\). Hence \(A^t\to R\) is onto and the original map \(A\to R\) is finite.

Its kernel is prime because \(R\) is a domain. If it were nonzero, choose a nonzero \(F\) in it. The source \(A\) is a Noetherian local domain of dimension \(N=\dim R\), and \(F\) is a nonzero nonunit. The one-equation theorem gives \(\dim A/(F)=N-1\), so the quotient by the kernel has dimension at most \(N-1\). But \(R\) is integral over that quotient, and going up together with incomparability gives equality of their dimensions: prime chains lift and a strict chain in the integral extension has strict contracted inclusions. This contradicts \(\dim R=N\). Thus the map is injective, and its image is the exact subring \(R_0\subset R\) specified by MC-STK-ERR-0703. Its residue map is the original coefficient residue isomorphism. Case I has now been proved only in Case I; MC-STK-ERR-0704 correctly leaves the proof open until Case II is also completed. The argument retains the independent dimensions \(d\) and \(d+1\), all original parameters and both cases, including \(d=0\).

\subsubsection{Report dispositions and propagation}\label{algebra-proof-065-report-dispositions-and-propagation}

Ten physical reports give eight retained-unit groups and eight retained operations. The proper-ideal qualification corrects the zero-ring boundary. The typed lifting composite and the completion-ideal clarification preserve the original maps; the latter is not evidence that the source's module-adic interpretation was false. The remaining edits clarify terminology, sentence structure and which proof case has finished. No new operation or optional extension group is added.

The original flat-local residue-field construction of ALGEBRA\_CONSTRUCT\_FLAT\_NOTES\_20260929.md feeds the Cohen-ring existence proof. The field formal-smoothness theorem and the previously corrected square-zero lifting proof feed the compatible coefficient maps; their exact quotient arrows then feed power-series surjectivity. The same coefficient ring and completeness maps are retained in both finite-extension cases. These receiving effects, the complete source dependencies and the earlier proof-note hashes are bound in the evidence receipt. Source attribution includes the original completion reference to Matlis, Theorem 15, and the source's attribution of that argument to Bjorn Poonen; no new reading of those external materials is claimed.

\hypertarget{algebra_japanese_notes_20260929}{}
\subsection{Japanese rings: original fields, integral closures and finiteness maps}\label{algebra-proof-066-japanese-rings-original-fields-integral-closures-and-finiteness-maps}

Source: Stacks Project authors, algebra.tex:45128--45614, original authority a04446e57ec1fbc252a871afcec7752fb2807b14, file SHA256 FA8BB92E58A4F78A2BD01B3B6A4A87DE0A0D279F5DD90641B574DD5FBFFFA4F3. The thirteen received reports retain MC-STK-ERR-0706 through MC-STK-ERR-0714. A separate source-bound correction qualifies ``any localization'' by ``nonzero''. Complete receiving arguments below are editorial material; they do not replace the translated source. Additional paper mining and optional extensions remain deferred.

\subsubsection{Integral closures, localization and the zero-ring boundary}\label{algebra-proof-066-integral-closures-localization-and-the-zero-ring-boundary}

The definitions at 45139--45145 apply to a domain \(R\) with fraction field \(K\). N-1 requires the integral closure of \(R\) in \(K\) to be a finite \(R\)-module. N-2 requires this for its integral closure in every finite field extension \(L/K\), including \(L=K\).

The source allows zero in a multiplicative subset: 1041--1047 requires only the identity and closure under multiplication. Its fraction equivalence at 1049--1065 makes \(W^{-1}R=0\) whenever \(0\in W\), since the witness \(u=0\) equates any two fractions. Zero is not a domain: the source convention at 36--37 says it has no prime ideal, and 154 identifies domains with quotients by prime ideals. In particular, its zero ideal is not prime. Taking a field \(R=k\) and \(W=\{0,1\}\) therefore shows that ``any localization'' in 45175 includes an object outside the stated definitions of N-1 and N-2. The minimal correction is ``any nonzero localization''. The familiar intended reading quantifies over localizations which remain domains; the qualification makes that reading explicit without changing either definition or declaring the usual nonzero-localization theorem false.

For \(0\notin W\), the exact map \(W^{-1}R\to K\), \(r/w\mapsto r/w\), is injective, its image is a domain and its fraction field is the original \(K\). For \(C\), the integral closure of \(R\) in the original finite extension \(L/K\), the map \[
 W^{-1}C\longrightarrow L,\qquad c/w\longmapsto c/w
\] identifies its image with the integral closure of \(W^{-1}R\) in \(L\). Here is the coefficient calculation from 7596--7627 with all denominator factors. If \[
 c^d+\sum_{i=1}^{d}a_i c^{d-i}=0,
\] then \(c/f\) satisfies \[
 (c/f)^d+\sum_{i=1}^{d}(a_i/f^i)(c/f)^{d-i}=0.
\] Conversely, if \(z/f\) satisfies a monic equation with coefficients \(a_i/f_i\), put \(F=\prod_{j=1}^{d}f_j\). After clearing the equality in a localization there is \(f'\in W\) such that \[
 (f'Fz)^d+
 \sum_{i=1}^{d}
 f^i(f')^i\left(f_i^{i-1}\prod_{j\ne i}f_j^i\right)a_i
 (f'Fz)^{d-i}=0.
\] Thus \(f'Fz\in C\), and the original fraction is \(z/f=(f'Fz)/(ff'F)\). In the present field inclusion the clearing witness can be \(f'=1\); the displayed general formula retains it. If \(C\) is finite over \(R\), its localized generators generate \(W^{-1}C\), proving both corrected localization assertions.

For the finite-cover lemma at 45184--45199, retain the original elements \(f_1,\ldots,f_n\) generating the unit ideal and the original closure \(C\). The hypothesis ``each domain \(R_{f_i}\)'' already specifies a domain chart. Alternatively, any zero chart can be removed from this cover because \(R\) is a domain and its defining element must be zero. Choose finitely many elements of \(C\) whose localizations generate each \(C_{f_i}\), clearing the denominators of the chosen module generators. Their joint \(R\)-span \(N\) is finite and \((C/N)_{f_i}=0\). For each \(v\in C/N\), choose integers \(e_i\geq1\) with \(f_i^{e_i}v=0\). The ideal generated by these powers has radical the unit ideal, so there are \(b_i\in R\) with \(1=\sum_i b_i f_i^{e_i}\). Consequently \(v=0\), and \(C=N\). The later uses of localization in this section invert nonzero elements or complements of prime ideals; none uses the excluded zero ring.

For use throughout the section, every algebraic \(\alpha\in L\) has a nonzero multiplier \(g\in R\) with \(g\alpha\) integral. Write its original monic polynomial as \[
 P(T)=T^d+\sum_{i=0}^{d-1}a_iT^i\in K[T]
\] and choose \(g\ne0\) with every \(ga_i\in R\), for example the product of the actual coefficient denominators. Keep \(\beta=g\alpha\) alongside \(\alpha\). Its equation is \[
 g^dP(T/g)=T^d+\sum_{i=0}^{d-1}g^{d-i}a_iT^i\in R[T],
 \qquad \alpha=g^{-1}\beta.
\] This also proves that the fraction field of the integral closure in \(L\) is exactly \(L\).

\subsubsection{Finite-variable descent and finite extensions of domains}\label{algebra-proof-066-finite-variable-descent-and-finite-extensions-of-domains}

We verify the infinite polynomial example at 45153--45170 without substituting a different extension. Let \(K=k(x_1,x_2,\ldots)\), and choose an actual \(K\)-basis \(e_1=1,e_2,\ldots,e_d\) of the original finite \(L/K\). Its finitely many multiplication-table coefficients lie in \(K_0=k(x_1,\ldots,x_n)\) for some \(n\). The \(K_0\)-span \(L_0\) of these same basis elements inside \(L\) is closed under multiplication and contains 1. It is a finite-dimensional domain over \(K_0\); multiplication by any nonzero element is injective and hence surjective, so \(L_0\) is a field. The map \[
 K\otimes_{K_0}L_0\longrightarrow L,\qquad a\otimes b\longmapsto ab
\] is an isomorphism because it sends the displayed basis to the original basis. A polynomial relation in the tail variables with coefficients in \(L_0\), expanded in the \(e_i\), would give a relation in that \(K\)-basis, followed by a relation in the algebraically independent original variables over \(K_0\). All its coefficients vanish. Thus these tail variables are algebraically independent over \(L_0\) and \(L=L_0(x_{n+1},x_{n+2},\ldots)\).

Put \(R_0=k[x_1,\ldots,x_n]\) and let \(C_0\) be its integral closure in \(L_0\). The polynomial N-2 proof below, iterated from the field \(k\), proves that \(C_0\) is finite over \(R_0\). This is an editorial route within the present section; the source's forward citation to the ubiquity-of-Nagata proposition remains intact. The ring \(C_0[x_{n+1},x_{n+2},\ldots]\) is integral over \(R\) and has fraction field \(L\). It is normal: an integral equation and a fraction involve only finitely many variables and therefore lie in one of the finite polynomial rings over the normal domain \(C_0\), which is normal by the earlier polynomial-normality lemma. Hence it is the exact integral closure of \(R\) in \(L\). The same finite list of \(R_0\)-module generators for \(C_0\) generates it over \(R\).

For \(R=\mathbf Z[x_1,x_2,\ldots]\), take \(K_0=\mathbf Q(x_1,\ldots,x_n)\) and \(R_0=\mathbf Z[x_1,\ldots,x_n]\). The char-zero result below starts from the Noetherian normal domain \(\mathbf Z\); the same polynomial argument applies. Both infinite polynomial rings are non-Noetherian: the ideal generated by all \(x_i\) cannot be finitely generated. Any proposed finite generators use only a finite set of variables and have zero constant term; setting those variables to zero annihilates those generators but does not annihilate a remaining variable.

For the quasi-finite ascent at 45201--45232, \(R\subset S\) induces a finite extension \(K\subset L=\operatorname{Frac}(S)\): the generic fibre is a finite-type zero-dimensional domain over \(K\), hence a finite field extension. Let \(A\) be the integral closure of \(R\) in \(S\). It is an \(R\)-submodule of the integral closure of \(R\) in \(L\), which is finite by N-2; since \(R\) is Noetherian, \(A\) is finite. The earlier quasi-finite integral-closure theorem supplies the actual elements \(g_i\in A\) with \(A_{g_i}=S_{g_i}\) covering \(\operatorname{Spec}(S)\). Nonzero charts suffice. In the finite case, for a finite \(E/\operatorname{Frac}(S)\), also \(E/K\) is finite, and \[
 \operatorname{cl}_R(E)=\operatorname{cl}_S(E).
\] Indeed integrality over \(R\) implies integrality over \(S\), while integrality over the integral \(R\)-algebra \(S\) implies integrality over \(R\) by transitivity. The finite \(R\)-generators for this closure also generate it over \(S\). Thus finite extensions preserve N-2; applying this to \(A\), and then applying localization and the finite-cover argument, proves the quasi-finite assertion.

For Laurent descent at 45234--45254, retain \(C=\operatorname{cl}_R(K)\) and \(D=\operatorname{cl}_{R[z,z^{-1}]}(K(z))\). Since \(K[z,z^{-1}]\) is normal, \(D\subset K[z,z^{-1}]\), and \(C\subset D\). Write the given finite module generators as \(f_i=\sum_j a_{ij}z^j\). For \(c\in C\), express \(c=\sum_i h_i f_i\) with \(h_i=\sum_k r_{ik}z^k\) and all sums finite. The full coefficient at exponent zero is \[
 c=\sum_{i,j}r_{i,-j}a_{ij}.
\] Thus \(C\) lies in the finite \(R\)-module \(\sum_{i,j}Ra_{ij}\); Noetherianity makes \(C\) finite. We retain that original hypothesis and do not insert an optional generalization into the review.

The omitted finite-descent proof at 45256--45266 has the following exact field construction. For N-1, \(\operatorname{cl}_R(K)\) is an \(R\)-submodule of \(\operatorname{cl}_S(\operatorname{Frac}(S))\), finite over \(S\) and then over \(R\). For N-2, retain an arbitrary original finite extension \(M/K\). Choose a \(K\)-embedding \(\iota:\operatorname{Frac}(S)\to\overline M\) and put \(E=M\iota(\operatorname{Frac}(S))\). Both fields being finite over \(K\), \(E/\iota(\operatorname{Frac}(S))\) is finite. The actual inclusion \(M\to E\) sends \(\operatorname{cl}_R(M)\) into \(\operatorname{cl}_{\iota(S)}(E)\). N-2 for \(S\) makes the latter finite over \(\iota(S)\), hence over \(R\). Noetherianity makes the original submodule finite. No linear-disjointness assertion is required.

\subsubsection{The trace lattice and the original involution example}\label{algebra-proof-066-the-trace-lattice-and-the-original-involution-example}

For 45268--45300, let \(R\) be Noetherian and normal, \(K=\operatorname{Frac}(R)\), and \(L/K\) the original finite separable extension. Starting from a \(K\)-basis \(u_i\), the coefficient-clearing calculation above gives nonzero \(g_i\in R\) and integral elements \(x_i=g_i u_i\), retaining the inverse equations \(u_i=g_i^{-1}x_i\). These still form a \(K\)-basis.

The pairing is exactly \(\langle v,w\rangle=\operatorname{Tr}_{L/K}(vw)\). The nondegeneracy theorem is fields.tex:2527--2583. Its separable conclusion can also be read on the exact scalar-extension map \[
 \overline K\otimes_K L\longrightarrow\prod_{\sigma:L\hookrightarrow\overline K}\overline K,
 \qquad a\otimes b\longmapsto(a\sigma(b))_\sigma.
\] It is an isomorphism: adjoin finitely many separable generators successively; at each step the defining polynomial factors into distinct linear factors over \(\overline K\), and the Chinese remainder evaluation maps are inverse isomorphisms. In these coordinates multiplication by \(b\) is diagonal with entries \(\sigma(b)\); its trace is their sum. For the matrix \(B=(\sigma(x_i))_{\sigma,i}\), the basis comparison makes \(B\) invertible, and the trace matrix is \(G=B^{\mathsf t}B\), so \(\det(G)=\det(B)^2\ne0\).

If \(z\) is integral over \(R\), the coefficients of its minimal polynomial \(P_z(T)=T^d+c_{d-1}T^{d-1}+\cdots+c_0\) belong to \(R\) by the normal-domain minimal-polynomial lemma. The exact trace formula from fields.tex:2425--2441 is \[
 \operatorname{Tr}_{L/K}(z)=-[L:K(z)]c_{d-1}.
\] The sign and integer factor remain, including when the integer becomes zero in the field. Thus the trace lies in \(R\). Put \[
 v_j=\sum_l(G^{-1})_{lj}x_l,\qquad
 M=\{y\in L:\operatorname{Tr}_{L/K}(x_i y)\in R\text{ for every }i\}.
\] The exact inverse maps are \[
 M\longrightarrow R^{\oplus n},\quad y\longmapsto(\operatorname{Tr}(x_i y))_i,
 \qquad (r_j)_j\longmapsto\sum_j r_jv_j.
\] Matrix multiplication gives \(\operatorname{Tr}(x_i v_j)=\delta_{ij}\), proving both composites are identities. For \(S=\operatorname{cl}_R(L)\), each \(x_i y\) is integral, so \(S\subset M\). As \(R\) is Noetherian, \(S\) is finite.

In 45302--45314 retain \(A=\mathbf C[x_1,x_2,\ldots]\), the involution \(\sigma(x_i)=-x_i\), \(L=\operatorname{Frac}(A)\), and \(R=A^{\{1,\sigma\}}\). A polynomial is invariant precisely when every monomial with nonzero coefficient has even total degree, since the monomials are a basis and \(2\ne0\). Every even-degree monomial is a product of pairs \(x_i x_j\), proving the stated generators.

The fraction field is exactly \(K=L^{\{1,\sigma\}}\). For an invariant fraction \(a=u/v\), the original nonzero denominator product \(v\sigma(v)\) belongs to \(R\), and \[
 a=\frac{u\sigma(v)}{v\sigma(v)},\qquad
 u\sigma(v)=a\,v\sigma(v)\in A\cap L^{\{1,\sigma\}}=R.
\] This proves the required surjection from \(\operatorname{Frac}(R)\) onto the invariant field; the reverse inclusion is immediate. Since \(A\) is normal by the finite-variable argument, an element of \(K\) integral over \(R\) lies in \(A\) and remains invariant, hence lies in \(R\). Thus \(R\) is normal. Each \(x_i\) satisfies \(T^2-x_i^2\in R[T]\), so \(A\) is integral over \(R\); its normality and fraction field \(L\) identify it with \(\operatorname{cl}_R(L)\).

Moreover \(L=K(x_1)\), since \(x_i=(x_1x_i)/x_1\), and \(x_1\notin K\) because \(\sigma(x_1)=-x_1\ne x_1\). Thus \([L:K]=2\), and it is separable. Let \(R_+\) be the ideal of positive-degree invariants. Its extension \(R_+A\) consists exactly of polynomials with monomials of total degree at least two. Hence \[
 A/R_+A=\mathbf C\oplus\bigoplus_{i\geq1}\mathbf C x_i
\] as a \(\mathbf C=R/R_+\)-vector space, which is infinite-dimensional. Therefore \(A\) cannot be a finite \(R\)-module. Also \(R\) is non-Noetherian: the degree-two part of \(R_+\) has the infinite independent family \(x_i x_j\) with \(i\leq j\); a finite list of ideal generators could span only finitely many degree-two components, since every multiplier of positive degree raises degree to at least four. MC-STK-ERR-0706 repairs only the apposition; the original counterexample is valid.

\subsubsection{The derivation and every original denominator factor}\label{algebra-proof-066-the-derivation-and-every-original-denominator-factor}

For 45321--45367 let \(D:R\to R\) be the original derivation. Its unique extension to \(K\) is \[
 D(r/s)=\frac{sD(r)-rD(s)}{s^2},\qquad s\ne0.
\] Uniqueness follows by applying the product rule to \(s(r/s)=r\). If \(r/s=r'/s'\), differentiating \(rs'=r's\) and multiplying by \(ss'\) gives equality of these displayed fractions; addition and multiplication follow from the same product-rule expansion. This proves the typed extension specified by MC-STK-ERR-0707.

Retain the original \(a\in K\), \(\delta_a=D(a)\ne0\), and \(L=K[x]/(x^p-a)\). Since a derivation kills \(p\)th powers, \(a\notin K^p\). The polynomial \(T^p-a\) is irreducible: the minimal polynomial of its unique root in an algebraic closure is a purely inseparable polynomial of degree dividing \(p\); if its degree were one the root would belong to \(K\). Thus \(L\) is a field, with basis \(1,x,\ldots,x^{p-1}\).

Choose an actual nonzero \(f\in R\) with \(b=f^pa\in R\); if \(a=c/d\), take \(f=d\) and \(b=d^{p-1}c\). Keep \(\beta=fx\), \(b\), \(a\), and \(x\) simultaneously. The comparison \[
 K[T]/(T^p-b)\longrightarrow L,\quad T\longmapsto fx,
 \qquad x\longmapsto f^{-1}T
\] is an isomorphism: the two relations are related by \((fx)^p-b=f^p(x^p-a)\), and the inverse uses \(f^{-p}(T^p-b)\). The derivative is \[
 \delta_b=D(b)=p f^{p-1}D(f)a+f^pD(a)=f^p\delta_a\in R\setminus\{0\}.
\] The first summand vanishes because the field has characteristic \(p\), as displayed; its contribution has not been absorbed into the notation.

We prove by induction on \(0\leq i<p\) that if \(y=\sum_{j=0}^{i}c_j\beta^j\), \(c_j\in K\), is integral over \(R\), then \(\delta_b^i c_j\in R\) for all \(j\). For \(i=0\), normality gives \(c_0\in R\). For \(i>0\), normality first gives \[
 y^p=\sum_{j=0}^{i}c_j^p b^j\in R,
 \qquad
 D(y^p)=\delta_b\sum_{j=1}^{i}j c_j^p b^{j-1}\in R.
\] Put \(z=\delta_b\sum_{j=1}^{i}j c_j\beta^{j-1}\). Its full \(p\)th power is \[
 z^p=\delta_b^p\sum_{j=1}^{i}j^p c_j^p b^{j-1}
     =\delta_b^{p-1}D(y^p)\in R,
\] because \(j^p=j\) in the prime field. Thus \(z\) is integral, and the induction hypothesis for degree \(i-1\) gives \(\delta_b^i j c_j\in R\). The integers \(1,\ldots,i\) are nonzero prime-field elements and hence units of \(R\), so \(\delta_b^i c_j\in R\) for \(j>0\). Finally \[
 \delta_b^i c_0=\delta_b^i y-\sum_{j=1}^{i}\delta_b^i c_j\beta^j
\] is integral over \(R\) and belongs to \(K\); normality gives membership in \(R\). This is exactly the missing invocation made explicit by MC-STK-ERR-0708; integrality alone would not imply membership.

The entire integral closure consequently lies in the original finite lattice \[
 \sum_{j=0}^{p-1}R\,\delta_b^{-(p-1)}\beta^j
 =\sum_{j=0}^{p-1}R\,\frac{(fx)^j}{(f^pD(a))^{p-1}}\subset L.
\] For the original coordinates \(y=\sum_j a_jx^j\), the exact comparison is \(c_j=f^{-j}a_j\). Thus no original denominator factor has been lost by the auxiliary integral presentation. The closure is an \(R\)-submodule of this finite module, so is finite by Noetherianity.

\subsubsection{Purely inseparable reduction and polynomial N-2}\label{algebra-proof-066-purely-inseparable-reduction-and-polynomial-n-2}

For 45369--45408, N-2 implies N-1 by taking \(L=K\). In characteristic zero, if N-1 holds, \(R'=\operatorname{cl}_R(K)\) is a finite Noetherian normal \(R\)-algebra. Every finite \(L/K\) is separable; the trace proof gives \(\operatorname{cl}_{R'}(L)\) finite over \(R'\), hence over \(R\), and transitivity identifies it with \(\operatorname{cl}_R(L)\).

In positive characteristic, suppose all finite purely inseparable extensions have finite closures. Given the original finite \(L/K\), choose a finite normal extension \(M/K\) containing it. The field result fields.tex:3665--3681 supplies \(F=M^{\operatorname{Aut}(M/K)}\) with \(F/K\) purely inseparable and \(M/F\) Galois. In the finite case this follows directly from the finite-automorphism fixed-field theorem: \(M/F\) has group \(\operatorname{Aut}(M/K)\); every conjugate over \(K\) of an element of \(F\) lies in \(M\) and extends to an automorphism of the normal extension, so all its distinct conjugates coincide. Its irreducible polynomial therefore has only one distinct root, which says precisely that a \(p\)-power of the element lies in \(K\). Thus \(F/K\) is purely inseparable. The closure \(A=\operatorname{cl}_R(F)\) is finite, normal, Noetherian, and has fraction field \(F\) by the earlier multiplier calculation. The trace proof makes \(\operatorname{cl}_A(M)=\operatorname{cl}_R(M)\) finite. The original \(\operatorname{cl}_R(L)=L\cap\operatorname{cl}_R(M)\) is finite as an \(R\)-submodule. The converse is contained in the definition of N-2.

For polynomial N-1 at 45410--45440, \(R'[x]\), with \(R'=\operatorname{cl}_R(K)\), is finite over \(R[x]\), normal, and has fraction field \(K(x)\). Its integral inclusion and normality identify it exactly with the integral closure of \(R[x]\). For N-2, \(R'\) remains N-2 by finite ascent proved above; finite descent reduces the assertion to \(R'\), keeping the original comparison \(R[x]\subset R'[x]\). We describe the remaining construction for the resulting normal N-2 domain, denoting it by \(A\), with field \(F\), so the original \(R\) and \(R'\) remain identifiable. Characteristic zero follows from N-1 and the preceding equivalence.

In characteristic \(p\), retain a finite purely inseparable \(L/F(x)\). Choose its finitely many field generators \(u_j\) and a common \(q=p^e\) with \[
 u_j^q=P_j(x)/Q_j(x),\quad P_j,Q_j\in F[x],\quad Q_j\ne0.
\] In an algebraic closure containing \(L\), let \(F'/F\) adjoin the unique \(q\)th roots of every coefficient of every \(P_j,Q_j\); this is a finite purely inseparable extension. Choose \(t\) with \(t^q=x\), transcendental over \(F'\). Let \(\widetilde P_j(t)\) and \(\widetilde Q_j(t)\) retain exactly the exponents of \(P_j,Q_j\) and replace each coefficient by its chosen \(q\)th root. Then \[
 \left(\frac{\widetilde P_j(t)}{\widetilde Q_j(t)}\right)^q
 =\frac{P_j(x)}{Q_j(x)}=u_j^q.
\] The denominators remain nonzero because \(t\) is transcendental. Injectivity of the \(q\)th-power map in the ambient field gives \(u_j=\widetilde P_j(t)/\widetilde Q_j(t)\), proving the actual inclusion \(L\subset F'(t)\).

Put \(C=\operatorname{cl}_A(F')\). It is finite over \(A\), normal, and has fraction field \(F'\). The ring \(C[t]\) is normal and integral over \(A[x]\) under the original map \(x\mapsto t^q\), with fraction field \(F'(t)\). These facts identify it with the integral closure of \(A[x]\) in \(F'(t)\). If \(b_1,\ldots,b_m\) generate \(C\) over \(A\), its finite \(A[x]\)-generating list is exactly \[
 b_i t^j,\qquad 1\leq i\leq m,\quad 0\leq j<q,
\] since \(t^{aq+j}=x^a t^j\). The closure in the original field \(L\) is its intersection with this ring and therefore a finite \(A[x]\)-module, since \(A[x]\) is Noetherian. MC-STK-ERR-0709 specifies the correct named extension; MC-STK-ERR-0710 identifies the middle ring. Reuse of the source symbol \(R'\) is not an additional mathematical error.

\subsubsection{Normal loci and the exact equality of local rings}\label{algebra-proof-066-normal-loci-and-the-exact-equality-of-local-rings}

For 45442--45479 let \(R\) be a Noetherian domain and \(f\ne0\) with \(R_f\) normal. Use the already reviewed Serre criterion with its original hypotheses. A domain satisfies S1; the generic local ring is a field. Therefore a nonnormal \(R_{\mathfrak p}\) has a prime \(\mathfrak q\subset\mathfrak p\) such that either \(\operatorname{depth}R_{\mathfrak q}=1\) and \(\dim R_{\mathfrak q}\geq2\), or \(\dim R_{\mathfrak q}=1\) and \(R_{\mathfrak q}\) is not regular. MC-STK-ERR-0711 changes the connective between these alternatives to ``or'', retaining the two internal conjunctions. Since normality localizes and \(R_f\) is normal, every such \(\mathfrak q\) contains \(f\).

In the first case \(f\) is a nonzerodivisor in the maximal ideal of \(R_{\mathfrak q}\), so the depth of its quotient is zero. The localized associated-prime criterion then gives \(\mathfrak q\in\operatorname{Ass}_R(R/fR)\). Its height at least two makes it embedded: minimal primes over the nonzero principal ideal have height one by the principal ideal theorem. Conversely such an embedded associated prime has depth one in the original local domain. In the second case the quotient at the height-one prime is a nonzero zero-dimensional Noetherian local ring, hence has depth zero, again giving an associated prime. The precise finite set is \[
 E=\{\mathfrak q\in\operatorname{Ass}_R(R/fR):
 \operatorname{ht}\mathfrak q\geq2\ \text{or}\
 (\operatorname{ht}\mathfrak q=1\text{ and }R_{\mathfrak q}\text{ is not regular})\}.
\] Regular height-one associated primes are not inserted into \(E\). Serre's criterion gives one containment, and localization of normality gives the other in \[
 \operatorname{Spec}(R)\setminus U=\bigcup_{\mathfrak q\in E}V(\mathfrak q).
\] It is closed. If \(f\) is a unit, the quotient is zero, \(E\) is empty and \(U=\operatorname{Spec}(R)\), as required.

For the N-1 criterion at 45481--45532, finite closure \(C\subset K\) has actual module generators \(c_i\). Choose \(f_i\ne0\) with \(f_ic_i\in R\), and retain their full product \(f=\prod_i f_i\). Then \(R_f=C_f\), a normal ring. Local N-1 follows from nonzero localization.

Conversely assume \(R_f\) normal for \(f\ne0\) and every maximal localization N-1. Any finite intermediate \(R\subset A\subset K\) has \(A_f=R_f\), because it is integral over the normal ring \(R_f\) inside \(K\). Its nonnormal locus is closed by the preceding argument; the finite map of spectra is closed, so its image \(Z_A\subset\operatorname{Spec}(R)\) is closed. MC-STK-ERR-0712 corrects the wording assigning this image its symbol.

For a maximal \(\mathfrak m\), the finite closure \(C_{\mathfrak m}=\operatorname{cl}_{R_{\mathfrak m}}(K)\) is the localization of \(C\). Express its finitely many generators as \(a_i/s_i\) with \(a_i\in C\), \(s_i\notin\mathfrak m\). Their \(R_{\mathfrak m}\)-span is also generated by the \(a_i\). The finite ring \(A=R[a_i]\) therefore has \((R\setminus\mathfrak m)^{-1}A=C_{\mathfrak m}\), which is normal. Every prime of \(A\) above \(\mathfrak m\) consequently has a normal local ring, so \(\mathfrak m\notin Z_A\). The opens \(\operatorname{Spec}(R)\setminus Z_A\) cover all primes, since every prime lies in a maximal ideal and opens are stable under passage to smaller primes. Quasi-compactness gives finitely many \(A_i\) with \(\bigcap_i Z_{A_i}=\varnothing\). The ring \(A\subset K\) generated by all these finite integral algebras is finite over \(R\).

Here is the exact receiving step retained in MC-STK-ERR-0713. Given \(\mathfrak p'\in\operatorname{Spec}(A)\), put \(\mathfrak p=\mathfrak p'\cap R\), choose \(i\) with \(\mathfrak p\notin Z_{A_i}\), and put \(\mathfrak q_i=\mathfrak p'\cap A_i\). The ring \((A_i)_{\mathfrak q_i}\) is normal. Set \(W=A_i\setminus\mathfrak q_i\). The integral extension \(A_i\subset A\) gives \[
 (A_i)_{\mathfrak q_i}\subset W^{-1}A\subset K,
\] with the first extension integral and \(K\) the fraction field of \((A_i)_{\mathfrak q_i}\). Normality therefore gives \(W^{-1}A=(A_i)_{\mathfrak q_i}\). Under this equality \(W^{-1}\mathfrak p'\) contracts to, and thus equals, the maximal ideal of that local ring. Localizing further at \(\mathfrak p'\) leaves it unchanged and proves \[
 A_{\mathfrak p'}=(A_i)_{\mathfrak q_i}.
\] This argument first localizes both rings at the same multiplicative subset. It does not rely on the generally invalid assertion that arbitrary upper-prime localization preserves integrality over the lower local ring. All local rings of \(A\) are normal, so \(A\) is normal: an element of \(K\) integral over \(A\) belongs to each \(A_{\mathfrak p'}\), and their intersection in \(K\) is \(A\). For the last intersection assertion, a fraction outside \(A\) has a proper denominator ideal, contained in some maximal ideal, so remains outside that localization. Finally any element integral over \(R\) is integral over \(A\), so \(A=C\). This proves the original N-1 assertion.

\subsubsection{Tate finiteness with the original extension retained}\label{algebra-proof-066-tate-finiteness-with-the-original-extension-retained}

Retain the source hypotheses at 45534--45589: \(R\) is a Noetherian normal domain, \(R/xR\) is a domain and N-2, and the actual map \(R\to\varprojlim_n R/x^nR\) is an isomorphism. If \(x=0\), the N-2 hypothesis already concerns \(R\). Otherwise \(x\ne0\) and is not a unit, since the quotient is a domain. In characteristic zero the trace argument proves the conclusion. In characteristic \(p\), retain an arbitrary finite purely inseparable original \(L/K\), its closure \(S=\operatorname{cl}_R(L)\), and a power \(q=p^e\) with \(L^q\subset K\).

Choose \(y^q=x\) in an algebraic closure and retain the enlargement with distinct symbols \[
 E=L(y),\qquad T=\operatorname{cl}_R(E),\qquad S=T\cap L.
\] It is finite and purely inseparable over \(K\). The \(q\)th power of any rational expression in elements of \(L\) and \(y\) is the same rational expression with all coefficients and exponents raised to the \(q\)th power; its denominator stays nonzero, and \(L^q\subset K\), \(y^q=x\in K\) show \(E^q\subset K\). If \(T\) is finite over \(R\), its \(R\)-submodule \(S\) is finite by Noetherianity. This is the return map in MC-STK-ERR-0714; the original \(L,S\) remain present throughout the argument.

For every \(t\in T\), \(t^q\in K\) is integral over \(R\), so lies in \(R\). Set \(\mathfrak p=xR\). There is exactly one prime above it, namely \[
 Q=\{t\in T:t^q\in\mathfrak p\}.
\] The field Frobenius identities show this is an ideal, and primeness of \(\mathfrak p\) shows it is prime. Its contraction is \(\mathfrak p\), since \(r^q\in\mathfrak p\) exactly when \(r\in\mathfrak p\). Any prime over \(\mathfrak p\) contains \(t\) exactly when it contains \(t^q\), proving uniqueness. Moreover \[
 Q=yT.
\] For \(t\in Q\), write the actual equation \(t^q=xr\) with \(r\in R\). Then \((t/y)^q=r\), so \(t/y\) is integral and belongs to \(T\). Conversely \((yt)^q=xt^q\in\mathfrak p\). Division by \(y\ne0\) takes place in the unchanged enlarged field \(E\).

The nonzero principal prime \(\mathfrak p\) has height one. The normal Noetherian local ring \(R_{\mathfrak p}\) is a discrete valuation ring with maximal ideal generated by the actual \(x\). By the original Krull--Akizuki theorem, algebra.tex:29292--29309, the intermediate ring \(R_{\mathfrak p}\subset T_Q\subset E\) is Noetherian. It is normal, local, nonfield, and its maximal ideal is generated by \(y\); hence it is a discrete valuation ring. The source's finite-residue-extension theorem at 29240--29264 applies.

For completeness, its needed finiteness also has the following direct dimension proof preserving the coefficient factors. Lift any finite linearly independent residue list over \(\kappa(\mathfrak p)\) to \(b_1,\ldots,b_s\in T_Q\). Suppose \(\sum_i c_i b_i=0\) with \(c_i\in K\), not all zero. Take \(m=\min_{c_i\ne0}v_x(c_i)\in\mathbf Z\), and keep \(c_i=x^m d_i\), with \(d_i\in R_{\mathfrak p}\) and at least one unit. Then the original relation is \[
 x^m\sum_i d_i b_i=0.
\] Multiplication by the explicit inverse \(x^{-m}\in E\) gives \(\sum_i d_i b_i=0\). Reducing this equality modulo \(Q T_Q\) gives a nontrivial relation in the chosen residues, a contradiction. Hence the \(b_i\) are \(K\)-linearly independent in \(E\), so \(s\leq[E:K]\). Thus \([\kappa(Q):\kappa(\mathfrak p)]\leq[E:K]\) is finite.

The domain \(T/yT=T/Q\) embeds into \(\kappa(Q)\), is integral over \(R/xR\), and hence lies in its integral closure there. That closure is finite by the N-2 hypothesis on \(R/xR\). The quotient \(R/xR\) is Noetherian, so \(T/yT\) is finite. For each \(i\geq0\), multiplication by the actual \(y^i\) gives an isomorphism \[
 T/yT\longrightarrow y^iT/y^{i+1}T.
\] It is surjective by definition. If \(y^i t=y^{i+1}u\), cancellation of the nonzero \(y^i\) in the domain \(T\) gives \(t=yu\), proving injectivity. The \(q\) successive subquotients show that \(T/y^qT=T/xT\) is finite over \(R\).

If \(z\in\bigcap_{n\geq1}x^nT\), write \(z=x^nt_n\) for every \(n\). The full exponent identity is \[
 z^q=x^{nq}t_n^q\in x^{nq}R.
\] Separatedness of the actual completion map of \(R\) gives \(\bigcap_n x^{nq}R=0\), whence \(z^q=0\), and \(z=0\) in the field \(E\).

Apply the precise finite-module theorem at 22859--22883: completeness of the ring, separatedness of the module, and finiteness modulo the ideal suffice. No completeness assumption on \(T\) is inserted. Explicitly, let \(m_1,\ldots,m_r\in T\) lift generators modulo \(xT\). Given \(t\in T\), recursively write \[
 t-\sum_{n=0}^{N}\sum_{j=1}^{r}x^n a_{j,n}m_j\in x^{N+1}T,
 \qquad a_{j,n}\in R.
\] This follows by expressing the current remainder as \(x^N u\) and choosing coefficients for the image of \(u\) modulo \(xT\). The sums \(A_j=\sum_{n\geq0}x^n a_{j,n}\) converge in the original complete ring \(R\), with their tails in \(x^{N+1}R\). The actual finite sum \(\sum_j A_jm_j\) already belongs to \(T\); its difference from \(t\) lies in every \(x^{N+1}T\), hence is zero. Therefore \(T\) is finite over \(R\), and so is the original \(S=T\cap L\). The purely inseparable criterion proves N-2 for \(R\).

\subsubsection{Power series and receiving propagation}\label{algebra-proof-066-power-series-and-receiving-propagation}

For 45592--45608 retain the original Noetherian N-2 domain \(R\). Its normal closure \(R'\) in its fraction field is finite, normal, Noetherian, and N-2 by finite ascent. The map \(R[[x]]\subset R'[[x]]\) is finite: if \(b_1,\ldots,b_m\) generate \(R'\) over \(R\), then for an original series \(\sum_n c_nx^n\) choose \(c_n=\sum_i r_{i,n}b_i\), retaining the exact identity \[
 \sum_{n\geq0}c_nx^n=\sum_{i=1}^{m}b_i\left(\sum_{n\geq0}r_{i,n}x^n\right).
\] Noetherianity of the original power-series ring is the earlier power-series theorem.

The normality argument in algebra.tex:8149--8174 retains \(w=f/g\in\operatorname{Frac}(R'[[x]])\subset\operatorname{Frac}(R')((x))\). If \(w\) is integral, the finite module generated by its powers admits a common nonzero denominator \(h\in R'[[x]]\) with \(w^e h\in R'[[x]]\) for every \(e\geq0\). Write the original leading terms as \(w=a_nx^n+\cdots\), \(h=b_mx^m+\cdots\). The leading term of their product is \(b_m a_n^e x^{m+ne}\). Since this exponent is nonnegative for every \(e\), \(n\geq0\). Each \(b_m a_n^e\in R'\). The ring \(R'[a_n]\) is therefore an \(R'\)-submodule of \(b_m^{-1}R'\), a finite module. Noetherianity makes \(R'[a_n]\) finite, so the determinant argument makes \(a_n\) integral over \(R'\), and normality gives \(a_n\in R'\). Subtract this original term \(a_nx^n\) from \(w\); the remainder is still integral, and repeat. Every coefficient of the original Laurent expansion then lies in \(R'\), with no negative exponent, so \(w\in R'[[x]]\). The zero series is already in that ring. This proves its normality with all original leading factors retained.

The actual variable \(x\) in \(R'[[x]]\) has quotient \(R'\); the ring is complete for its powers. Tate's theorem applies to that ring and variable, yielding N-2 for \(R'[[x]]\). The finite-descent proof above then gives N-2 for the original \(R[[x]]\).

Propagation within the reviewed interval is complete. The nonzero-localization qualification is used by the finite-cover, quasi-finite-ascent and N-1-localization arguments; their specified localizations remain nonzero. The prior Serre proof feeds the exact normal-locus set and the local-equality step. The prior Krull--Akizuki and finite-over-complete-ring arguments feed Tate with the original extension and all \(q\)-power factors retained. The derivation lattice, trace lattice and finite-extension descent feed the polynomial and power-series conclusions. No new optional extension is claimed. The next source interval begins at 45615, Nagata rings; only its opening through 45644 has been read here.

\hypertarget{algebra_nagata_notes_20260929}{}
\subsection{Nagata rings: completion, finite fields and the original monogenic rings}\label{algebra-proof-067-nagata-rings-completion-finite-fields-and-the-original-monogenic-rings}

Source: Stacks Project authors, algebra.tex:45615--46292, authority a04446e57ec1fbc252a871afcec7752fb2807b14, SHA256 FA8BB92E58A4F78A2BD01B3B6A4A87DE0A0D279F5DD90641B574DD5FBFFFA4F3. The twenty-two received reports are checked against the original source and the ten existing units MC-STK-ERR-0715--0724. The following complete receiving proofs are editorial material. The source's structure, rings, generators and arguments remain identifiable. Optional extensions and further paper mining remain after Algebra.

\subsubsection{Reduced algebras and nonzero localization charts}\label{algebra-proof-067-reduced-algebras-and-nonzero-localization-charts}

The source defines a universally Japanese ring by requiring every finite-type algebra which is a domain to be N-2. A Nagata ring is Noetherian and has N-2 quotients by all prime ideals. The zero ring satisfies both definitions: it is Noetherian, has no prime ideals, and admits no unital map to a nonzero domain. These observations also cover empty spectra and empty finite covers below.

For 45639--45667 let \(R\) be Nagata and \(S\) an essentially finite-type reduced \(R\)-algebra. It is Noetherian. If its minimal primes are \(\mathfrak q_1,\ldots,\mathfrak q_m\), the maps to its minimal-prime localizations give an injection \[
 S\longrightarrow\prod_{i=1}^m K_i,
 \qquad K_i=S_{\mathfrak q_i}=\operatorname{Frac}(S/\mathfrak q_i).
\] Its kernel is the intersection of the minimal primes, which is zero; each localization is a field because it is reduced, local and zero-dimensional. Put \(\mathfrak p_i=\mathfrak q_i\cap R\). Each \(K_i\) is a finitely generated field extension of \(\kappa(\mathfrak p_i)=\operatorname{Frac}(R/\mathfrak p_i)\). The exact algebraic subfield \(L_i\subset K_i\) is finite over that base field by fields.tex:3583--3603. To recall its degree argument, choose an actual transcendence basis \(t_1,\ldots,t_r\) with \([K_i:\kappa(\mathfrak p_i)(t_1,\ldots,t_r)]=n\). For every finite algebraic subextension \(k'/\kappa(\mathfrak p_i)\) inside \(K_i\), adjoining those same independent variables preserves the original finite degree, so \[
 [k':\kappa(\mathfrak p_i)]
 =[k'(t_1,\ldots,t_r):\kappa(\mathfrak p_i)(t_1,\ldots,t_r)]\leq n.
\] The degree preservation uses the same basis of \(k'\) after rational extension: clear rational denominators in a putative relation, and compare the coefficients of the independent monomials. Algebraic extension of the constant field preserves independence of the variables. Among finite subextensions choose one of maximal degree; adjoining any further algebraic element cannot increase that degree, so it already contains the whole algebraic subfield. This proves finiteness.

An element of \(K_i\) integral over \(R\) is integral over its image \(R/\mathfrak p_i\), and therefore algebraic over \(\kappa(\mathfrak p_i)\). Thus its closure is precisely \(C_i=\operatorname{cl}_{R/\mathfrak p_i}(L_i)\), finite over \(R/\mathfrak p_i\) by the Nagata definition. The closure \(C\) of \(R\) in the original \(S\) maps injectively into \(\prod_i C_i\). Since \(R\) is Noetherian, this submodule of a finite module is finite. If \(S=0\), the product is empty and the closure is zero, also finite.

For the N-1 test at 45669--45688, let \(S\) be an original finite-type domain over \(R\), and let \(L/\operatorname{Frac}(S)\) be finite. Choose finitely many field generators \(\alpha_j\) of \(L\). For their monic polynomials \(P_j(T)=T^{d_j}+\sum_{i=1}^{d_j}a_{j,i}T^{d_j-i}\), choose actual nonzero \(c_j\in S\) clearing all coefficient denominators. Retain \(\beta_j=c_j\alpha_j\) and \[
 c_j^{d_j}P_j(T/c_j)
 =T^{d_j}+\sum_{i=1}^{d_j}c_j^i a_{j,i}T^{d_j-i},
 \qquad \alpha_j=c_j^{-1}\beta_j.
\] Then \(S'=S[\beta_1,\ldots,\beta_r]\subset L\) is finite over \(S\), has fraction field exactly \(L\), and is finite type over \(R\). Under the assumed N-1 test its closure \(S''\) in \(L\) is finite over \(S'\), hence over \(S\). Integrality in both directions identifies \(S''=\operatorname{cl}_S(L)\). This proves the N-2 property for the original \(S\), without a Noetherian hypothesis on \(R\).

The essentially finite-type assertion at 45690--45700 has the following exact descent. Write the original algebra as \(A=W^{-1}B\) with \(B\) finite type over a universally Japanese \(R\). Given an arbitrary finite-type domain \(T\) over \(A\), let \(C\subset T\) be the ring generated by the image of \(B\) and finitely many chosen \(A\)-algebra generators of \(T\). Then \(C\) is a finite-type domain over \(R\), and the actual fraction map identifies \(T=W^{-1}C\). The images of \(W\) are nonzero because they are units in the domain \(T\). The corrected nonzero-localization theorem from the Japanese-rings note gives N-2 for \(T\). Thus \(A\) is universally Japanese.

For quasi-finite ascent at 45702--45718, \(S\) is Noetherian because it is finite type over the Noetherian \(R\). For \(\mathfrak q\in\operatorname{Spec}(S)\), the exact injection \(R/\mathfrak p\to S/\mathfrak q\), \(\mathfrak p=\mathfrak q\cap R\), is quasi-finite: it is the quotient of the base change to \(R/\mathfrak p\) and has finite fibres. The original domain ascent theorem from the Japanese-rings note applies, since \(R/\mathfrak p\) is N-2. This proves the Nagata assertion. Localization at \(W\) preserves Noetherianity. Every prime in \(W^{-1}R\) comes from a prime \(\mathfrak p\) disjoint from \(W\), and its domain quotient is the nonzero localization \(W^{-1}(R/\mathfrak p)\); it is N-2. If zero belongs to \(W\), the result is the zero Nagata ring, not a domain to which N-2 is assigned.

For 45729--45752 retain \(\sum_i r_i f_i=1\) in \(R\). Given \(\varphi:R\to S\) with \(S\) a finite-type domain, the elements with \(\varphi(f_i)=0\) define zero charts. Remove exactly those indices from the identity \[
 1=\sum_{\varphi(f_i)\ne0}\varphi(r_i)\varphi(f_i).
\] For each surviving index, \(S_{\varphi(f_i)}\) is a finite-type domain over the universally Japanese \(R_{f_i}\), hence N-2 and N-1. The finite-cover argument proves N-1 for \(S\), and the preceding N-1 test proves the first conclusion. Accordingly the new minimal qualification at 45742 is ``each nonzero'' localization. This propagates ALGEBRA-RECON-1292 to the first covering assertion.

For the second assertion, Noetherianity descends along this finite standard cover: for every ideal \(I\subset R\), choose a finite set of elements of \(I\) generating each localization \(I_{f_i}\), and let \(I_0\) be their joint span. Every element of \(I/I_0\) is killed by powers of all \(f_i\); those powers still generate the unit ideal, so \(I=I_0\). For a prime \(\mathfrak p\subset R\), keep only indices \(f_i\notin\mathfrak p\). Their images still generate 1 in \(R/\mathfrak p\), and \[
 R_{f_i}/\mathfrak pR_{f_i}=(R/\mathfrak p)_{f_i}
\] is a domain and N-2. The domain cover lemma applies. This is precisely the already composed MC-STK-ERR-0715. Both parts now use nonzero domain charts, while the source's Nagata assertion itself retains the zero ring.

\subsubsection{Complete rings and integral closure in the total quotient ring}\label{algebra-proof-067-complete-rings-and-integral-closure-in-the-total-quotient-ring}

For 45754--45778 let \(R\) be a complete Noetherian local ring. A prime ideal is proper, so its quotient is again a complete Noetherian local domain by the proper-quotient result in the Cohen note. The original finite regular-subring theorem embeds either \(k[[X_1,\ldots,X_d]]\) or \(\Lambda[[X_1,\ldots,X_d]]\) in that domain, with finite extension. Its two cases, coefficient maps, actual prime \(p\) and all parameters remain those proved in the Cohen note. A field is N-2. For a Cohen ring \(\Lambda\), the actual element \(p\) generates its maximal ideal, its residue quotient is a field, and it is a complete Noetherian normal domain. Tate's theorem from the Japanese-rings note therefore proves N-2 for \(\Lambda\). Iterating that note's power-series theorem preserves all variables and gives N-2 for either regular subring. Finite domain ascent then gives N-2 for the original quotient. Hence \(R\) is Nagata. The empty list of parameters is included.

Write \(B=\widehat R=\varprojlim_n R/\mathfrak m^n\). The source completion theorems give its actual faithfully flat local map, Noetherianity and completeness. A ring is analytically unramified exactly when this \(B\) is reduced. For 45807--45887, faithful flatness makes \(R\to B\) injective, so reducedness of \(B\) implies reducedness of \(R\).

If \(\mathfrak q\) is a minimal prime of such an \(R\), it is associated, so choose the actual \(f\in R\) with annihilator \(\mathfrak q\). Tensoring the kernel sequence of multiplication by \(f\) gives \[
 \operatorname{Ann}_B(f)=\mathfrak qB,
 \qquad \widehat{R/\mathfrak q}=B/\mathfrak qB.
\] If \(z^2\in\mathfrak qB\), then \(fz^2=0\) and \((fz)^2=f(fz^2)=0\). Reducedness of \(B\) implies \(fz=0\), hence \(z\in\mathfrak qB\). This square-root criterion makes the ideal radical: repeated squaring reaches a power at least any given nilpotence exponent and then descends. Thus \(B/\mathfrak qB\) is reduced. Conversely, the actual injection \(R\to\prod_i R/\mathfrak q_i\) for a reduced Noetherian ring remains injective under completion, giving \[
 B\longrightarrow\prod_i\widehat{R/\mathfrak q_i}.
\] If every factor is reduced, so is the subring \(B\).

Assume now that \(B\) is reduced. Its total quotient ring is the product of the fields \(B_{\mathfrak p_i}\) at its finitely many minimal primes. The integral closure \(D\) of \(B\) in that product is \[
 D=\prod_i D_i,\qquad
 D_i=\operatorname{cl}_{B/\mathfrak p_i}\bigl(\operatorname{Frac}(B/\mathfrak p_i)\bigr).
\] For this equality, an integral tuple has integral coordinates. Conversely, lift to \(B[T]\) a monic equation for each coordinate; the product of the finitely many lifted monic polynomials vanishes on the whole tuple, coordinate by coordinate. Since \(B\) is Nagata, each \(D_i\) is finite over \(B/\mathfrak p_i\), hence \(D\) is finite over \(B\).

Keep the original total quotient ring \(Q(R)=U^{-1}R\), with \(U\) the original set of nonzerodivisors, and let \(R'=\operatorname{cl}_R(Q(R))\). For each \(u\in U\), flatness of \(B\) preserves the injection given by multiplication by \(u\). Thus its image is a nonzerodivisor in \(B\). The exact maps are \[
 R'\otimes_R B\ \hookrightarrow\ Q(R)\otimes_R B
 \xrightarrow{\ \cong\ }U^{-1}B\ \hookrightarrow\ Q(B),
 \qquad (r/u)\otimes b\longmapsto rb/u.
\] The first injection is flatness; the last is localization at a subset of the nonzerodivisors. Their image is integral over \(B\), since integrality is preserved by this base change. Hence \(R'\otimes_R B\subset D\), a finite \(B\)-module; Noetherianity makes the submodule finite. Choose module generators and express each as a finite sum \(\sum_j r'_j\otimes b_j\). The finite list of all \(r'_j\) generates the tensor product over \(B\). If \(N\subset R'\) is their \(R\)-span, then \((R'/N)\otimes_R B=0\); faithful flatness implies \(R'=N\). This proves finiteness of the original integral closure, and its domain case is N-1. MC-STK-ERR-0716 corrects the article at 45877; the new copyedit at 45872 supplies ``Denote by''. Neither changes these maps.

\subsubsection{Codimension one and the completion criterion}\label{algebra-proof-067-codimension-one-and-the-completion-criterion}

For 45889--45914, \(R\) need not itself be a domain. Let \(R_{\mathfrak p}\) be the original DVR and suppose \(B/\mathfrak pB\) is reduced. Its associated primes are minimal primes. For one such \(\mathfrak q\), flatness of \(R/\mathfrak p\to B/\mathfrak pB\) and going down show \(\mathfrak q\cap R=\mathfrak p\): if the contraction were larger, a prime below it over zero would contradict minimality in the quotient. Therefore the map \[
 R_{\mathfrak p}\longrightarrow B_{\mathfrak q}
\] is defined and flat. The quotient \((B/\mathfrak pB)_{\mathfrak q}\) is a reduced zero-dimensional local ring, hence a field; thus the maximal ideal of \(B_{\mathfrak q}\) is \(\mathfrak pB_{\mathfrak q}\). Choose \(x\in R\) whose image generates \(\mathfrak pR_{\mathfrak p}\); if an original uniformizer is given as a fraction, its nonzero numerator differs from it by the displayed denominator unit. Flatness of this localized map makes multiplication by \(x\) injective on \(B_{\mathfrak q}\). Its maximal ideal is \(xB_{\mathfrak q}\). A nonzero Noetherian local ring with maximal ideal generated by a nonzerodivisor is a DVR: separatedness of the maximal-ideal powers writes every nonzero element as \(x^n u\) with \(u\) a unit, and the full product \(x^{n+m}uv\) cannot vanish. This gives a domain and the discrete valuation by the actual exponent. No assertion that \(x\) is globally a nonzerodivisor of \(R\) is required.

For the criterion at 45916--45962 let \(R\) be the original Noetherian local domain, \(0\ne x\in\mathfrak m\), and assume the source's three conditions. The associated primes \(\mathfrak p_i\) of \(R/xR\) are its minimal support primes, all of height one by the principal ideal theorem. The regular local rings \(R_{\mathfrak p_i}\) are DVRs. For every \(\mathfrak q_{ij}\in\operatorname{Ass}_B(B/\mathfrak p_iB)\), the preceding paragraph proves that \(B_{\mathfrak q_{ij}}\) is a DVR. Apply the exact flat associated-prime formula of algebra.tex:15773--15851 to \(M=R/xR\) and \(N=B\): \[
 \operatorname{Ass}_B(B/xB)
 =\bigcup_i\operatorname{Ass}_B(B/\mathfrak p_iB)
 =\{\mathfrak q_{ij}\}_{i,j}.
\] Here all hypotheses are present: \(R\) is Noetherian and \(B\) is flat over \(R\). The formula's reverse inclusion follows by tensoring each injection \(R/\mathfrak p_i\to M\). For the forward inclusion, localize an associated prime of the tensor product and contract it to \(R\). If that contraction is not associated to \(M\), a nonzerodivisor in the contracted maximal ideal remains one after the flat tensor map, contradicting association. A finite prime filtration of \(M\), tensored by the flat module, then identifies the corresponding quotient \(B/\mathfrak p_iB\). This is the original proof with these specific inputs.

If \(y\in B\) satisfies \(y^2=0\), it vanishes in every domain \(B_{\mathfrak q_{ij}}\). Consequently its class vanishes in \(B/xB\): a nonzero cyclic submodule of this finite module has an associated prime of the module and cannot vanish at that localization. Thus \(y=xy_1\). Since multiplication by \(x\) on \(B\) is injective by flatness, \(x^2y_1^2=0\) implies \(y_1^2=0\). Iterating gives the actual equalities \(y=x^n y_n\) for all \(n\). The original local Noetherian \(B\) has \(x\) in its maximal ideal, so \(\bigcap_n x^nB=0\) and \(y=0\). A nonzero nilpotent would have a nonzero square-zero power, so \(B\) is reduced.

\subsubsection{Local Nagata domains and the finite local test}\label{algebra-proof-067-local-nagata-domains-and-the-finite-local-test}

For 45964--46015 induct on the finite dimension \(d\) of the original Noetherian local domain \(R\). In dimension zero it is a field. Let \(S\) be its closure in its original fraction field. N-2 for \(R\) makes \(S\) finite; finite ascent makes it Nagata, and it is normal. Its finitely many maximal ideals \(\mathfrak m_i\) all lie over \(\mathfrak m\). Put \(J=\bigcap_i\mathfrak m_i\). Since \(S/\mathfrak mS\) is Artinian, choose the actual integer \(c\geq1\) with \[
 J^c\subseteq\mathfrak mS\subseteq J,
 \qquad J^{cn}\subseteq\mathfrak m^nS\subseteq J^n.
\] The identity induces mutually inverse comparisons of the two completion limits. Chinese remainder and localization of each quotient give \[
 \widehat R\hookrightarrow\widehat S^{\,\mathfrak mS}
 \cong\prod_i\varprojlim_n S/\mathfrak m_i^n
 \cong\prod_i\widehat{S_{\mathfrak m_i}}.
\] The injection is exact completion of the inclusion of finite \(R\)-modules. Every map in the product comparison comes from the same residue or fraction class. The local domains \(S_{\mathfrak m_i}\) are normal Nagata rings of dimension at most \(d\). It suffices to prove reducedness of their completions.

For any such normal local Nagata domain \(A\) of positive dimension at most \(d\), choose the original nonzero \(x\) in its maximal ideal. The height-one intersection result, with its complete proof in the Serre note, gives no embedded primes for \(A/xA\); all associated primes \(\mathfrak p\) have height one and \(A_{\mathfrak p}\) is a DVR. The quotient \(A/\mathfrak p\) is a local Nagata domain of dimension strictly less than \(d\): append \((0)\subsetneq\mathfrak p\) to a chain of primes above \(\mathfrak p\) to obtain the strict bound. No catenary dimension formula is used. Induction makes \(\mathfrak p\) analytically unramified. The preceding criterion now proves that \(A\) is analytically unramified. This treats all the local factors, and their product contains the original completion injectively, proving the original assertion for \(R\).

For 46017--46060, condition (1), that \(R\) is Nagata, implies condition (2) by finite Nagata ascent, localization and the local-domain result just proved. Condition (2) implies (3) by taking the already local finite domain. The missing colon at 46020 and redundant ``then'' at 46026 are copyedits.

Assume (3), and retain an arbitrary original prime \(\mathfrak p\subset R\), \(A=R/\mathfrak p\), and finite \(L/\kappa(\mathfrak p)\). If \(\mathfrak p=\mathfrak m\), \(A\) is a field and its closure in \(L\) is \(L\), a finite module. Otherwise choose the actual \(g\in\mathfrak m\setminus\mathfrak p\) required by MC-STK-ERR-0717. Retain each original field generator \(x_i\) and its minimal polynomial \[
 P_i(T)=T^{d_i}+\sum_{j=1}^{d_i}a_{i,j}T^{d_i-j}.
\] Choose \(f_i\in R\setminus\mathfrak p\) such that all \(\overline f_i a_{i,j}\in A\). With bars denoting actual quotient images, put \[
 y_i=\overline g\,\overline f_i x_i,\qquad
 Q_i(T)=(\overline g\,\overline f_i)^{d_i}
 P_i\bigl(T/(\overline g\,\overline f_i)\bigr)
 =T^{d_i}+\sum_{j=1}^{d_i}
   \overline g^{,j}\overline f_i^{,j}a_{i,j}T^{d_i-j}.
\] All nonleading coefficients belong to \(\mathfrak m/\mathfrak p\). Every multiplier is nonzero in \(A\), and the inverse equation is \(x_i=(\overline g\,\overline f_i)^{-1}y_i\). Thus \(S=A[y_1,\ldots,y_n]\subset L\) is finite over \(R\), a domain, and has the original fraction field \(L\). The special fibre is the image of \[
 (R/\mathfrak m)[T_1,\ldots,T_n]/(T_1^{d_1},\ldots,T_n^{d_n})
 \longrightarrow S/\mathfrak mS,
 \qquad T_i\longmapsto[y_i].
\] This map is onto. The fibre is nonzero by lying over for the integral inclusion \(A\subset S\); the displayed source is local with nilpotent maximal ideal, so its nonzero quotient has one maximal ideal. Every maximal ideal of the finite \(R\)-algebra \(S\) lies over \(\mathfrak m\). Hence \(S\) is local. Condition (3) makes it analytically unramified, so its closure \(S'\) in \(L\) is finite. Transitivity identifies \(S'=\operatorname{cl}_A(L)\), and it is finite over \(A\). This proves Nagata for \(R\).

The new parentheses at 46052 and 46054 identify the original coefficient quotients; they do not change their rings. The suggestion in OCC-01009 to repeat ``in \(L\)'' is an optional clarification, not a mathematical defect: the immediately preceding sentence fixes \(L/\kappa(\mathfrak p)\), the generator construction takes place in that sole extension, and the last sentence explicitly identifies both closures in \(L\). This editorial proof states every ambient field, while the source wording remains. The \(\mathfrak p=\mathfrak m\) case is different: without MC-STK-ERR-0717 its prescribed multiplication could annihilate all generators and lose the given extension.

\subsubsection{Finite-type stability with every original ring retained}\label{algebra-proof-067-finite-type-stability-with-every-original-ring-retained}

For the proposition at 46066--46222, condition (3) implies condition (2) by essentially finite-type stability of universal Japaneseness and Noetherianity of finite-type algebras. Condition (2) implies (1) by the identity algebra. It remains to prove that Nagata implies finite-type stability. Choose actual algebra generators of the original \(R\)-algebra \(S\). Adjoining them successively factors \(R\to S\) into monogenic maps. At each step, and for each prime quotient of the monogenic target, pass to the corresponding prime quotient of the source. Nagata is preserved under that quotient directly from the definition. Thus it suffices to prove N-2 for an original injection of domains \(R\subset S=R[x]\), with \(R\) Nagata.

Write \(K=\operatorname{Frac}(R)\). If \(x\) is transcendental over \(K\), the evaluation \(R[X]\to S\), \(X\mapsto x\), is injective and onto. The polynomial N-2 result applies. Otherwise \(K(x)/K\) is finite. The source's wording ``purely transcendental'' at 46101 needs the qualification ``of positive transcendence degree'' for this first inference: fields.tex:3371--3384 permits a polynomial ring with no variables, so \(K/K\) is itself purely transcendental of degree zero. For example, \(\mathbf Z\subset\mathbf Z[1/2]\) has identical fraction fields, but the evaluation \(\mathbf Z[X]\to\mathbf Z[1/2]\), \(X\mapsto1/2\), has the nonzero relation \(2X-1\). The finite case below handles this example. The qualification selects exactly the transcendental-generator case and leaves all finite cases to the next step.

For the finite case retain \(L=\operatorname{Frac}(S)\), an arbitrary finite extension \(M/L\), and the same original generator \(x\). Define \[
 A=\operatorname{cl}_R(M),\qquad B=A[x]\subset M,
 \qquad C=\operatorname{cl}_S(M).
\] Since \(M/K\) is finite and \(R\) is N-2, \(A\) is finite over \(R\); its fraction field is exactly \(M\) by the multiplier calculation in the Japanese note. It is normal and Nagata by finite ascent. The actual inclusion \(S\subset B\) is finite: a finite \(R\)-generating list of \(A\) also generates \(B\) over \(S\). It is therefore integral, and \[
 C=\operatorname{cl}_B(M).
\] Indeed integrality over \(S\) implies integrality over \(B\), while transitivity gives the reverse implication. Finiteness of \(C\) over \(B\) implies finiteness over \(S\) by the displayed finite inclusion; conversely any finite \(S\)-generating list also generates it over \(B\). Thus it is enough to prove the following unchanged original subproblem: for a normal Nagata domain \(A\), an element \(x\in\operatorname{Frac}(A)\), and \(B=A[x]\) as a subring of that actual fraction field, \(B\) is N-1. No normality of the original \(R\) was needed. This verifies precisely the deletion in MC-STK-ERR-0718, rather than silently restricting the theorem's input.

We now prove that subproblem, initially writing its original normal Nagata domain as \(R\), its actual fraction field as \(K\), and its original generator as \(x=b/a_0\), \(a_0\ne0\). Then \(S=R[x]\subset K\), and the exact fraction maps give \(S_{a_0}=R_{a_0}\), which is normal. The N-1 criterion reduces to each original maximal ideal \(\mathfrak m\subset S\). Let \(\mathfrak n=\mathfrak m\cap R\). The residue field \(\kappa(\mathfrak m)\) is generated as an algebra over \(\kappa(\mathfrak n)\) by the image of \(x\): it is already generated over the subring \(R/\mathfrak n\), and contains its fraction field. By the field version of the Nullstellensatz, this extension is finite. Let \[
 \overline f(X)=X^d+\sum_{i=1}^{d}\overline a_iX^{d-i}
 \in\kappa(\mathfrak n)[X]
\] be its monic minimal polynomial, with the actual coefficients retained. Choose representatives of each coefficient as fractions modulo \(\mathfrak n\), and let \(c\in R\setminus\mathfrak n\) be the product of their chosen nonzero denominator representatives. Each coefficient lifts to \(a_i\in R_c\). Put \[
 f(X)=X^d+\sum_{i=1}^{d}a_iX^{d-i}\in R_c[X].
\] Its evaluation at the unchanged \(x\) belongs to \(\mathfrak mS_c\), because its residue is \(\overline f(\overline x)=0\). Localization preserves the original local ring: \(S_{\mathfrak m}\to(S_c)_{\mathfrak mS_c}\) sends every fraction to the same fraction, with inverse using that \(c\notin\mathfrak m\). The source uses its symbol \(a\) again for this localization element; here \(a_0\) and \(c\) are kept distinct so its two actual denominators cannot be confused. MC-STK-ERR-0719 supplies the missing subject, and MC-STK-ERR-0720 keeps the polynomial indeterminate \(X\) distinct from the element \(x\).

Choose a finite splitting field \(K''/K\) of \(f\), retain all its roots with their multiplicities as \(\rho_1,\ldots,\rho_d\), and set \[
 A=\operatorname{cl}_{R_c}(K''),\qquad B=A[x].
\] Each \(\rho_i\) lies in \(A\) because \(f\) is monic. The ring \(A\) is finite, normal and Nagata over \(R_c\). The original \(S_c\subset B\) is finite. Set \(W=S_c\setminus\mathfrak mS_c\). The ring \(D=W^{-1}B\) is finite over the local domain \((S_c)_{\mathfrak mS_c}=S_{\mathfrak m}\), and its maximal ideals correspond exactly to maximal ideals \(\mathfrak m'\) of \(B\) lying over \(\mathfrak mS_c\). Their local rings are the original \(B_{\mathfrak m'}\). Moreover, \(B_{a_0}=A_{a_0}\), and \(D_{a_0}\) is a localization of this normal domain. The actual nonzero \(a_0\) therefore supplies the generic normal open required by the N-1 criterion for \(D\).

If each of those \(B_{\mathfrak m'}\) is N-1, that criterion makes \(D\) N-1. Finite domain descent then makes the unchanged \(S_{\mathfrak m}\) N-1. This proves why MC-STK-ERR-0721 needs only maximal ideals over the fixed \(\mathfrak m\). The source's stronger antecedent involving every maximal ideal would be sufficient, but its following proof establishes exactly the restricted set needed here. The accompanying copyedit changes ``the polynomial \ldots{} split'' to ``the polynomial \ldots{} splits''.

For one such \(\mathfrak m'\), the full product \[
 f(x)=\prod_{i=1}^{d}(x-\rho_i)\in\mathfrak m'
\] puts some factor \(u=x-\rho_i\) in \(\mathfrak m'\). Keep its exact expression \[
 u=(b-a_0\rho_i)/a_0,\qquad x=u+\rho_i,
 \qquad A[u]=A[x]=B.
\] The polynomial comparison sends \(X\) to \(U+\rho_i\) and has inverse \(U\mapsto X-\rho_i\). It preserves the original ring and local ideal; it does not replace the original fraction or lose the root factor. If \(u=0\), then \(x=\rho_i\in A\) and \(B=A\), so \(B_{\mathfrak m'}\) is normal and N-1. This is the zero case made explicit by MC-STK-ERR-0722 in the source's reused variable. For the remaining case \(u\ne0\), we prove analytic unramifiedness of \(B_{\mathfrak m'}\) by the earlier criterion applied to this specific \(u\).

The original comma-list subscripts \(\sum_{i=1,\ldots,d}\) and \(\prod_{i=1,\ldots,d}\) specify the same full index set \(\{1,\ldots,d\}\) as lower and upper limits. They are valid mathematical notation. OCC-01014 and OCC-01312 do not identify a defect, and no notation-normalizing source edit is made. In particular all repeated roots and every coefficient remain in the displayed sum and product.

\subsubsection{The linear kernel, fractional ideal and local DVR map}\label{algebra-proof-067-the-linear-kernel-fractional-ideal-and-local-dvr-map}

Retain the actual evaluation \(A[U]\to B\), \(U\mapsto u\), with kernel \(I\). For every \(a\in A\) with \(au=b'\in A\), the linear form \(aU-b'\) belongs to \(I\). These forms generate \(I\). To prove this, take its original nonzero polynomial \[
 g(U)=c_dU^d+\cdots+c_1U+c_0,\qquad c_d\ne0.
\] A constant polynomial in the kernel is zero since \(A\subset B\), so \(d\geq1\). Multiplying the actual equation \(g(u)=0\) by \(c_d^{d-1}\) yields \[
 (c_du)^d+\sum_{j=0}^{d-1}c_jc_d^{d-1-j}(c_du)^j=0.
\] This is the full monic equation in algebra.tex:30480--30496; every leading-coefficient power is retained. Thus \(c_du\) is integral over the normal \(A\) and lies in its fraction field, so \(b'=c_du\in A\). The exact subtraction is \[
 g(U)=(c_dU-b')U^{d-1}+g_1(U),
 \quad g_1\in I,\quad\deg(g_1)<d.
\] In particular the new coefficient of \(U^{d-1}\) is the old coefficient plus \(b'\); lower coefficients remain. Induction proves generation. MC-STK-ERR-0723 restores the missing plus sign in the source's displayed polynomial, as independently reported by OCC-01313 and OCC-12434.

Put \(J=A\cap uA\), an ideal of \(A\). The coefficient-at-zero map sends each generating linear form to \(-b'\), and all elements \(b'\in J\) occur because they have the form \(b'=au\) with \(a\in A\). Thus the exact mutually inverse quotient maps give \[
 B/uB=A[U]/(I,U)\cong A/J,
 \qquad [a]\longmapsto[a],\quad [h(U)]\longmapsto[h(0)].
\] The sign \(-b'\) generates the same ideal through multiplication by the explicit unit \(-1\). The Serre note, section ``Height-one intersections with all original factors retained'', applies to the same nonzero fraction \(u=(b-a_0\rho_i)/a_0\). It proves that every associated prime of \(A/J\) has height one and that there are no embedded primes, retaining all original numerator, denominator and valuation factors. Under the displayed quotient isomorphism, associated primes correspond exactly: the annihilator of an element is the inverse image of its annihilator over the quotient ring. Thus \(B/uB\), and then its localization at \(\mathfrak m'\), have no embedded primes. This supplies criterion (2).

Let \(\mathfrak q\subseteq\mathfrak m'\) be one of the associated primes of \(B/uB\), and put \(\mathfrak p=\mathfrak q\cap A\). The correspondence makes \(\mathfrak p\) an associated prime of \(A/J\), hence of height one. The ring \(A_{\mathfrak p}\) is a DVR. The localized fractional-ideal equality is \[
 J_{\mathfrak p}=A_{\mathfrak p}\cap uA_{\mathfrak p}.
\] It follows by localizing the kernel sequence defining the intersection inside the original fraction field. It is a proper ideal, since \(\mathfrak p\) is in the support of \(A/J\). If \(v_{\mathfrak p}(u)\leq0\), the second fractional ideal contains \(A_{\mathfrak p}\), contradicting properness. Hence \[
 e=v_{\mathfrak p}(u)>0,\qquad u=w\pi^e,
 \quad w\in A_{\mathfrak p}^{\times},
\] for an actual uniformizer \(\pi\), with the actual unit \(w=u\pi^{-e}\) retained. In particular \(u\in A_{\mathfrak p}\). Localizing the actual ring \(B=A[u]\) at \(A\setminus\mathfrak p\) therefore gives exactly \(A_{\mathfrak p}\). The prime \(\mathfrak q\) becomes its maximal ideal, so further localization gives \[
 B_{\mathfrak q}=A_{\mathfrak p}.
\] This proves criterion (3)(a), including the precise map justifying the source's equality. It does not infer equality merely from both rings having dimension one.

The same quotient isomorphism gives \(B/\mathfrak q\cong A/\mathfrak p\). After localizing at \(\mathfrak m'\), its exact denominator images give a local ring of this quotient of \(A\). Finite Nagata ascent preserves the quotient, and localization preserves Nagata; the result is a local Nagata domain. The local-domain theorem above makes it analytically unramified, establishing criterion (3)(b). The remaining criterion (1) is the already separated case \(u\ne0\). Hence \(B_{\mathfrak m'}\) is analytically unramified and N-1. The finite semilocal descent in the preceding section returns this result to the original \(S_{\mathfrak m}\) and then to \(S\). The original finite-type stability proposition follows.

For 46224--46245, fields and complete Noetherian local rings have already been proved Nagata. A characteristic-zero Dedekind domain is Noetherian and normal; its nonzero prime quotients are fields, and its zero-prime quotient is N-2 by the original characteristic-zero trace argument. This includes \(\mathbf Z\). Finite-type stability now gives all the displayed finite-type extensions. These implications complete the source's earlier forward reference in the infinite-polynomial example; they do not change that example's original fields or polynomial generators.

\subsubsection{The original non-Japanese DVR and the completion root test}\label{algebra-proof-067-the-original-non-japanese-dvr-and-the-completion-root-test}

At 46247--46262 retain exactly the earlier example \[
 A=\left\{\sum_{i\geq0}a_i x^i\in k[[x]]:
 [k^p(a_0,a_1,\ldots):k^p]<\infty\right\},
 \qquad [k:k^p]=\infty.
\] The complete proof that \(A\) is a DVR with uniformizer \(x\), residue field \(k\), completion \(k[[x]]\) and \(\operatorname{Frac}(A)\cap k[[x]]=A\) is retained in the Krull--Akizuki note, section ``The two original examples and their exact completions''. Keep the source's arbitrary \(f=\sum_{i\geq0}a_i x^i\notin A\) and \(R=A[f]\). Since \(f^p\in k^p[[x]]\subset A\), the original \(R\) is finite over \(A\). For every \(n\geq1\), the actual elements are \[
 g_n=\sum_{0\leq i<n}a_i x^i\in A,
 \qquad h_n=\frac{f-g_n}{x^n}=\sum_{j\geq0}a_{n+j}x^j,
 \qquad h_n^p=\sum_{j\geq0}a_{n+j}^{p}x^{pj}\in k^p[[x]]\subset A.
\] The finite set of coefficients of \(g_n\) generates a finite extension of \(k^p\), so its claimed membership follows from the definition. Each \(h_n\) is integral over \(A\), hence belongs to the closure \(R'\) in the original fraction field of \(R\). If \(R'\) were finite over \(R\), it would be finite over \(A\). The exact identities \[
 f=g_n+x^n h_n
\] put the class of \(f\) in every \(x^n(R'/A)\). The quotient \(R'/A\) is then a finite module over the local Noetherian \(A\), and the intersection of these powers is zero. Hence \(f\in A\), a contradiction. This proves the stated failure of N-1 for \(R\). If \(A\) were N-2, finite domain ascent would make \(R\) N-2 and thus N-1, also a contradiction. A DVR's only prime quotients are itself and its residue field, so the displayed equivalence of Nagata and N-2 is exact.

For 46264--46286 let \((A,\mathfrak m)\) be the original local Nagata domain of positive characteristic \(p\), \(K=\operatorname{Frac}(A)\), \(a\in A\), and \(\alpha\in\widehat A\) with \(\alpha^p=a\). The preceding local-domain theorem has already proved that \(\widehat A\) is reduced; we use it with the actual completion map. Suppose no root belongs to \(A\). If \(\beta=b/c\in K\), \(b,c\in A\), \(c\ne0\), satisfied \(\beta^p=a\), retain the full identity in \(\widehat A\) \[
 (c\alpha-b)^p=c^p\alpha^p-b^p=c^pa-b^p=0.
\] Reducedness gives \(c\alpha=b\). Faithful flatness contracts the original principal ideal exactly: \(c\widehat A\cap A=cA\), since \(A/cA\to\widehat A/c\widehat A\) is injective. Thus \(b=cd\) for some \(d\in A\), and \(\beta=d\) is the forbidden root. This avoids treating \(c\) as a unit of \(\widehat A\); equivalently the fraction comparison takes place in its total quotient ring, into which the original nonzero \(c\) remains a nonzerodivisor by flatness.

Therefore \(T^p-a\) is irreducible over \(K\), and \[
 B=A[T]/(T^p-a)\hookrightarrow K[T]/(T^p-a)
\] is a domain. The injection follows by the unique representatives of degree less than \(p\) in the monic quotient, and the target is a field. It is a finite free \(A\)-module on \(1,T,\ldots,T^{p-1}\) and is Nagata by finite ascent. Its special fibre is \[
 B/\mathfrak mB=\kappa(\mathfrak m)[T]/(T^p-\overline a).
\] The polynomial either is irreducible or is \((T-\gamma)^p\) for its unique root \(\gamma\) in the residue field, so this ring has exactly one prime. Every maximal ideal of the finite integral algebra \(B\) lies over \(\mathfrak m\), proving that \(B\) is local with maximal ideal \(\mathfrak n\). In both cases \[
 \mathfrak n^p\subseteq\mathfrak mB\subseteq\mathfrak n,
 \qquad\mathfrak n^{pk}\subseteq\mathfrak m^kB\subseteq\mathfrak n^k.
\] Thus the identity gives inverse comparisons of the maximal-adic and \(\mathfrak mB\)-adic completion systems. These are the locality and topology facts made explicit by MC-STK-ERR-0724.

Completion in the displayed finite free basis gives the actual isomorphism \[
 \widehat B=\widehat A\otimes_A B
 =\widehat A[T]/(T^p-a).
\] In it, retain \(\eta=[T]-\alpha\). The full Frobenius equation is \(\eta^p=[T]^p-\alpha^p=a-a=0\). Yet \(\eta\ne0\), because its coefficient of \([T]\) in the free basis of this monic degree-\(p\) quotient is 1, and \(p\geq2\). More explicitly, the inverse presentation maps \(T\mapsto\alpha+\eta\) and \(\eta\mapsto[T]-\alpha\) give \(\widehat B\cong\widehat A[\eta]/(\eta^p)\). The original generator \([T]\) remains \(\alpha+\eta\). Thus \(\widehat B\) is not reduced, contradicting the proved analytic unramifiedness of the local Nagata domain \(B\). The original \(a\) must have its \(p\)th root in \(A\).

\subsubsection{Report dispositions and completed receiving checks}\label{algebra-proof-067-report-dispositions-and-completed-receiving-checks}

All ten existing units MC-STK-ERR-0715--0724 are retained, with both operations of 0717. Their roles are the nonzero quotient charts, article, maximal-prime boundary before changing generators, legitimate normality reduction, missing subject, polynomial indeterminate, restricted maximal ideals, zero-generator case, missing plus sign and locality/completion comparison. The two reports about comma-list operator indices are rejected as nondefects. The ambient-field repetition in OCC-01009 is recorded as an optional clarification not integrated. The remaining received copyedits add ``by'', the enumeration colon, remove the doubled conditional and parenthesize the two coefficient rings.

The full-source reading also propagates the earlier nonzero-localization correction to the first covering assertion, restricts the transcendental case to positive transcendence degree, and repairs the singular verb ``splits''. These are source-bound corrections, with no invented received report. The generic normal opens, faithful-flat total-quotient comparison, shared localization sets, all field-generator multipliers and inverse comparisons, the exact fractional ideal and its valuation factors, every completion filtration, and the nilpotent \([T]-\alpha\) have been verified above. The Cohen, Japanese, Serre and Krull--Akizuki arguments are used only through their proved original hypotheses. No new optional extension or novelty claim is made. The next interval starts at 46293, Ascending properties; only its opening through 46322 has been read here.

\hypertarget{algebra_ascent_notes_20260929}{}
\subsection{Algebra: ascending properties, with the original local rings retained}\label{algebra-proof-068-algebra-ascending-properties-with-the-original-local-rings-retained}

This is an editorial verification of the Stacks Project authors' argument, not a replacement translation or a novelty claim. The authority is commit \texttt{a04446e57ec1fbc252a871afcec7752fb2807b14}, \texttt{algebra.tex} lines 46293--46632, SHA-256 \texttt{FA8BB92E58A4F78A2BD01B3B6A4A87DE0A0D279F5DD90641B574DD5FBFFFA4F3}. The \href{https://github.com/stacks/stacks-project/blob/a04446e57ec1fbc252a871afcec7752fb2807b14/algebra.tex\#L46293}{original source} attributes the module depth formula to Grothendieck and Dieudonné, EGA IV, Proposition 6.3.1. The preceding normality criterion retains its EGA IV, Theorem 5.8.6 attribution. Those original attributions are preserved; no new reading of EGA is claimed here.

The source statement, hypotheses, notation, finite-depth induction, localization maps, and finite-type models remain identifiable. The long arguments below belong in linked editorial material. Four existing correction units, MC-STK-ERR-0725--0728, are retained. Only two additional copyedits are proposed. One is an explicit refinement of the postimage of MC-STK-ERR-0725-OP1, whose original receipt is immutable.

\subsubsection{The module depth formula, including zero modules}\label{algebra-proof-068-the-module-depth-formula-including-zero-modules}

Let \((R,\mathfrak m_R)\to(S,\mathfrak m_S)\) be the original local homomorphism of Noetherian local rings. Let \(M\) be a finite \(R\)-module and \(N\) a finite \(S\)-module, flat over \(R\). Put \(C=S/\mathfrak m_RS\). All three depths in the original assertion are taken over their stated rings: \[
 \operatorname{depth}_S(M\otimes_RN)
 =\operatorname{depth}_R(M)+\operatorname{depth}_C(N/\mathfrak m_RN).
\] Depth of the zero module is \(+\infty\), by the source definition at 17765--17776, and addition takes place in the nonnegative extended integers. If \(M=0\), the left side and the first summand are infinite. If \(N=0\), the left side and the second summand are infinite. These cases prove the formula directly. They must precede induction on a finite integer; no subtraction from infinity is used. This is the substantive content of retained unit 0725.

Assume henceforth that both modules are nonzero. Nakayama over \(S\), with \(\mathfrak m_RS\subseteq\mathfrak m_S\), gives \(N/\mathfrak m_RN\ne0\). Nakayama over \(R\) gives a nonzero finite-dimensional vector space \(M/\mathfrak m_RM\cong (R/\mathfrak m_R)^r\), with \(r\ge1\). Tensoring its quotient map with \(N\) gives a surjection \[
 M\otimes_RN\longrightarrow
 (M/\mathfrak m_RM)\otimes_RN
 \cong (N/\mathfrak m_RN)^r\ne0.
\] Thus the left module is nonzero as well. All these modules are finite over their respective Noetherian local rings. Their depths are finite: a regular sequence on a nonzero finite module has length at most its support dimension, and a Noetherian local ring has finite dimension, bounded by a finite generating set for its maximal ideal. Denote by \(n\) the right-hand sum. The additional copyedit inserts the missing word ``by'' in this designation, without altering the preceding zero-module correction.

For \(n=0\), the depth-zero associated-prime criterion supplies the original \(z\in M\), with annihilator \(\mathfrak m_R\), and \(\bar y\in N/\mathfrak m_RN\), with annihilator \(\mathfrak m_S/\mathfrak m_RS\) in \(C\). The map \[
 R/\mathfrak m_R\longrightarrow M,\qquad \bar r\longmapsto rz
\] is injective. Flatness of \(N\) over \(R\) therefore makes \[
 N/\mathfrak m_RN\longrightarrow M\otimes_RN,
 \qquad \bar v\longmapsto z\otimes v
\] injective; this formula is independent of the chosen lift of \(\bar v\). It is \(S\)-linear. If \(y\) is a lift of \(\bar y\), the annihilator of \(z\otimes y\) is consequently exactly \(\mathfrak m_S\). The target has depth zero.

Suppose that \(n>0\) and the fibre depth is positive. Choose the original \(f\in\mathfrak m_S\) whose image is a nonzerodivisor on \(N/\mathfrak m_RN\). Apply the source's fibre-injectivity lemma, 23423--23477, to multiplication by \(f:N\to N\). It is injective, and \(N/fN\) is flat over \(R\). To recall why the latter conclusion is available, the lemma first lifts injectivity modulo \(\mathfrak m_R\) through all \(\mathfrak m_R^j\) using \[
 (N/\mathfrak m_RN)\otimes_{R/\mathfrak m_R}
 (\mathfrak m_R^j/\mathfrak m_R^{j+1})
\] and flatness of the target. The tensor maps are injective over the residue field. Krull intersection for the finite \(S\)-module kills the intersection of the kernels. The same argument after quotienting by any ideal \(I\subset R\) proves injectivity modulo \(I\). The exact sequence \[
 0\to\operatorname{Tor}_1^R(N/fN,R/I)
 \to N/IN\xrightarrow{f}N/IN
\] then gives flatness of the cokernel by the ideal criterion. This preserves the actual multiplication map and the original ideals.

Nakayama shows that \(N/fN\ne0\). The quotient of the original fibre is \[
 (N/fN)/\mathfrak m_R(N/fN)
 =N/(f,\mathfrak m_R)N.
\] The depth-drop lemma gives \[
 \operatorname{depth}_C(N/\mathfrak m_RN)
 =\operatorname{depth}_C(N/(f,\mathfrak m_R)N)+1.
\] Depth here can equally be computed over the further quotient by \(\bar f\): all maximal-ideal actions on the module are the same, and regular sequences lift and descend through that quotient. The new right-hand sum is exactly \(n-1\), so induction applies to \(M,N/fN\). Flatness of \(N/fN\) makes the tensor sequence \[
 0\to M\otimes_RN\xrightarrow{1_M\otimes f}M\otimes_RN
 \to M\otimes_R(N/fN)\to0
\] exact. Its first map is the original scalar multiplication by \(f\), its cokernel is the displayed quotient, and the depth-drop lemma gives \[
 \operatorname{depth}_S(M\otimes_RN)
 =\operatorname{depth}_S(M\otimes_R(N/fN))+1
 =\operatorname{depth}_R(M)
   +\operatorname{depth}_C(N/(f,\mathfrak m_R)N)+1.
\] This is the required formula with every summand retained.

If the fibre depth is zero and \(n>0\), choose the original \(f\in\mathfrak m_R\) regular on \(M\). Flatness of \(N\) preserves the exact sequence \[
 0\to M\otimes_RN\xrightarrow{f}M\otimes_RN
 \to (M/fM)\otimes_RN\to0.
\] Both \(M/fM\) and the tensor cokernel are nonzero by the preceding Nakayama and tensor-quotient arguments. Now \(\operatorname{depth}_R(M)=\operatorname{depth}_R(M/fM)+1\). Induction for \(M/fM,N\), followed by depth drop over \(S\), proves \[
 \operatorname{depth}_S(M\otimes_RN)
 =\operatorname{depth}_R(M/fM)
  +\operatorname{depth}_C(N/\mathfrak m_RN)+1
 =\operatorname{depth}_R(M)+\operatorname{depth}_C(N/\mathfrak m_RN).
\] The two positive-depth branches exhaust the finite integer \(n>0\). Taking \(M=R,N=S\) proves the ring depth formula at 46373--46384.

\subsubsection{Cohen--Macaulay ascent and the original local fibre maps}\label{algebra-proof-068-cohenmacaulay-ascent-and-the-original-local-fibre-maps}

For a flat local homomorphism of Noetherian local rings, the source dimension formula gives \[
 \dim S=\dim R+\dim(S/\mathfrak m_RS).
\] Its proof at 27420--27456 retains a chain of length \(d\) in the closed fibre and a chain of length \(e\) below \(\mathfrak m_R\). Going down lifts the latter below the first prime of the former. Strictness follows from distinct contractions. This gives a chain of length \(e+d\) in \(S\), and the previously proved dimension upper bound gives equality. Combined with depth addition, it gives the exact identity \[
 \dim S-\operatorname{depth}S
 =\bigl(\dim R-\operatorname{depth}R\bigr)
 +\bigl(\dim(S/\mathfrak m_RS)
       -\operatorname{depth}(S/\mathfrak m_RS)\bigr).
\] Both right-hand summands are nonnegative finite integers. Thus the left side vanishes exactly when both right-hand summands vanish. This proves the original Cohen--Macaulay equivalence. The proposed colon after ``equivalent'' is punctuation only.

For the subsequent global assertions, keep the original map \(\varphi:R\to S\), prime \(\mathfrak q\subset S\), and its inverse image \(\mathfrak p\subset R\). Set \[
 A=R_{\mathfrak p},\qquad B=S_{\mathfrak q},\qquad
 F=S\otimes_R\kappa(\mathfrak p),\qquad C=B/\mathfrak pB.
\] Write \(U=R\setminus\mathfrak p\). The map \(R\to B\) inverts \(U\), so it factors uniquely as \(A\to B\), sending \(r/u\) to \(\varphi(r)/\varphi(u)\); it is flat and local. There is an explicit isomorphism \[
 F\xrightarrow{\sim}U^{-1}S/\mathfrak pU^{-1}S,
 \qquad s\otimes(\bar r/\bar u)\longmapsto
 \overline{s\varphi(r)/\varphi(u)}.
\] Its inverse takes the class of \(s/\varphi(u)\) to \(s\otimes1/\bar u\). Balanced tensor relations and the quotient by \(\mathfrak p\) make both maps well-defined; on each displayed generator their composites are the identity. Let \(\bar{\mathfrak q}\) be the prime of this ring induced by \(U^{-1}\mathfrak q\). Then \[
 F_{\bar{\mathfrak q}}\xrightarrow{\sim}C,
 \qquad
 \frac{\overline{s/\varphi(u)}}{\overline{t/\varphi(v)}}
 \longmapsto\overline{\frac{s\varphi(v)}{t\varphi(u)}}.
\] Here \(t\notin\mathfrak q\), and \(u,v\notin\mathfrak p\), so every displayed denominator is invertible in the specified target. The inverse sends the class of \(s/t\in B\) to \((s\otimes1)/(t\otimes1)\). Clearing these same invertible denominators verifies well-definedness, multiplicativity, additivity, and both inverse identities. This identifies the actual fibre localization, not a different field or local ring.

\subsubsection{Serre conditions, reducedness and normality in the Noetherian case}\label{algebra-proof-068-serre-conditions-reducedness-and-normality-in-the-noetherian-case}

Assume the original Noetherian hypotheses and flatness. Write \(a=\dim A\), \(c=\dim C\), \(b=\dim B\). The local fibre comparison just proved lets us apply the assumed fibre condition to \(C\). If the base and fibres satisfy \((S_k)\), for the original nonnegative integer \(k\), depth addition gives the source's full calculation \[
 \begin{aligned}
 \operatorname{depth}B
 &=\operatorname{depth}C+\operatorname{depth}A\\
 &\ge\min(k,c)+\min(k,a)\\
 &=\min(2k,c+k,k+a,c+a)\\
 &\ge\min(k,b).
 \end{aligned}
\] The middle equality distributes the two minima over all four sums; none is discarded. For the last inequality, the first three entries are at least \(k\) because \(k,a,c\ge0\), and the fourth is at least \(b\) by the dimension inequality (in fact equality). Thus every entry is at least \(\min(k,b)\), including \(k=0\). This proves \((S_k)\) for \(S\).

If the base and fibres satisfy \((R_k)\), take \(\mathfrak q\) with \(b\le k\). The exact dimension equality \(b=a+c\) forces \(a\le k,c\le k\), so \(A,C\) are regular. For completeness, the source regular-ascent argument at 27458--27480 chooses generators \(x_1,\ldots,x_a\) of the maximal ideal of \(A\), and lifts \(y_1,\ldots,y_c\in\mathfrak m_B\) of generators of the maximal ideal of \(C\). Any element of \(\mathfrak m_B\) is congruent modulo \(\mathfrak m_AB\) to a linear combination of the \(y_j\). Its difference is a combination of the images of the \(x_i\). Thus these \(a+c=b\) elements generate \(\mathfrak m_B\). The general inequality \(\dim B\le\dim_{\kappa(\mathfrak m_B)}\mathfrak m_B/\mathfrak m_B^2\), together with this generating set of size \(b\), gives equality and hence regularity of \(B\). Zero-sized parameter lists are permitted. This proves \((R_k)\) for \(S\).

The preceding source criteria are already proved with their zero-quotient and local-height boundaries in \texttt{ALGEBRA\_SERRE\_NOTES\_20260929.md}, sections ``The depth conditions and reducedness'' and ``The original Serre argument and its zero quotient''. For Noetherian rings, reducedness is exactly \((R_0)+(S_1)\), and normality is exactly \((R_1)+(S_2)\). Applying the two ascent assertions to these pairs gives both original Noetherian ascent theorems. The fibre rings are Noetherian, being quotients of localizations of the original Noetherian \(S\). The zero ring, when present globally, has no primes, so these local conditions are vacuous and it is reduced and normal under the source conventions. No exclusion is added. Unit 0726 merely supplies the two missing list commas.

\subsubsection{Smooth reduced ascent with the finite-stage witnesses retained}\label{algebra-proof-068-smooth-reduced-ascent-with-the-finite-stage-witnesses-retained}

Let \(R\to S\) be smooth and \(R\) reduced; no Noetherian assumption is imposed on \(R\). Smoothness gives flatness and regular, hence reduced, fibres. If \(R\) is Noetherian then \(S\), being finitely presented over \(R\), is Noetherian, and the preceding theorem applies.

In general the finite-presentation descent of smooth maps, as cited in source smooth overview item (10), gives a finitely generated \(\mathbf Z\)-subalgebra \(R_0\subset R\), a smooth \(R_0\)-algebra \(S_0\), and the original isomorphism \(S\cong R\otimes_{R_0}S_0\). Let \(R_\lambda\) run over finite-type \(\mathbf Z\)-subalgebras of \(R\) containing \(R_0\). This is a directed system under inclusion: adjoining the finite generating sets of any two gives a third. Its union is \(R\), and \[
 \underset{\lambda}{\operatorname{colim}}
 (R_\lambda\otimes_{R_0}S_0)
 \xrightarrow{\sim}R\otimes_{R_0}S_0
\] sends each \(r_\lambda\otimes s_0\) to that same tensor. Surjectivity holds because a tensor is a finite sum; injectivity holds because any finite tensor relation uses finitely many coefficients, all contained in one later stage. These are ring maps, since products of elementary tensors use the same finite products of coefficients.

If the original \(x\in S\) satisfies \(x^2=0\), first choose a stage containing a representative \(x_\lambda\). The equality \(x_\lambda^2=0\) in the colimit has a witness in one later stage: equality in a filtered colimit of rings is eventual equality. Thus the original two enlargements of \(R_0\) are justified separately. Call the resulting ring and element again \(R_0,S_0,x_0\), preserving the source's rebasing convention. Smoothness persists by base change, \(R_0\) is reduced as a subring of the reduced \(R\), and it is Noetherian as a finite-type \(\mathbf Z\)-algebra. The Noetherian theorem gives reduced \(S_0\), whence \(x_0=0\) and \(x=0\). Finally, if a nonzero element were nilpotent with minimal exponent \(e\ge2\), its nonzero power of exponent \(\lceil e/2\rceil<e\) would be square-zero. Hence absence of nonzero square-zero elements proves reducedness. Retained unit 0727 specifies the subring as the subject of ``is reduced''; it changes no ring or hypothesis.

\subsubsection{Smooth normal ascent and the exact filtered-colimit proof}\label{algebra-proof-068-smooth-normal-ascent-and-the-exact-filtered-colimit-proof}

Let the original \(R\to S\) be smooth with \(R\) normal. Normality means that each prime localization is an integrally closed domain. In particular \(R\) is reduced: a nilpotent vanishes at each prime, and an element vanishing at all prime localizations is zero (otherwise its annihilator is contained in a maximal ideal at which it does not vanish). Thus the preceding theorem gives reduced \(S\).

Fix \(\mathfrak q\subset S\) and \(\mathfrak p=\varphi^{-1}(\mathfrak q)\). With \(U=R\setminus\mathfrak p\), the comparison \[
 R_{\mathfrak p}\otimes_RS\xrightarrow{\sim}U^{-1}S,
 \quad (r/u)\otimes s\longmapsto\varphi(r)s/\varphi(u)
\] has inverse \(s/\varphi(u)\mapsto (1/u)\otimes s\). Further localization at \(U^{-1}\mathfrak q\) is \(S_{\mathfrak q}\); the maps send each represented quotient to the same quotient and are inverse on all generators. Smoothness is retained by this base change. The original source reduction to \(R_{\mathfrak p}\), a normal domain, therefore suffices to prove the desired assertion about the unchanged \(S_{\mathfrak q}\).

Now assume \(R\) is a normal domain and retain a finite model \(R_0\subset R,S_0\) as above. In the original fraction field \(K_0=\operatorname{Frac}(R_0)\), let \(R_0'\) be the integral closure of \(R_0\). The immediately preceding Nagata ubiquity theorem makes every finite-type \(\mathbf Z\)-domain Nagata, so \(R_0'\) is finite over \(R_0\). Since \(K_0\subset\operatorname{Frac}(R)\), every element of \(R_0'\) is integral over \(R\); normality of the domain \(R\) puts it in \(R\). Here local normality implies integral closedness in the fraction field: an integral fraction lies in every \(R_{\mathfrak p}\); if it did not lie in \(R\), the proper denominator ideal \(I=\{r\in R:rx\in R\}\) would lie in a maximal ideal \(\mathfrak m\), contradicting membership in \(R_{\mathfrak m}\). Thus the original inclusion \(R_0'\subset R\) is proved with its field specified.

Set \(S_0'=R_0'\otimes_{R_0}S_0\). The exact base-change comparison is \[
 R\otimes_{R_0'}S_0'\xrightarrow{\sim}R\otimes_{R_0}S_0,
 \quad r\otimes(a'\otimes s_0)\longmapsto ra'\otimes s_0,
\] with inverse \(r\otimes s_0\mapsto r\otimes(1\otimes s_0)\). Tensor balancing verifies both inverse identities and the preserved structure maps. The new finite model is smooth; \(R_0'\) is a normal Noetherian domain, and its algebra \(S_0'\) is normal by Noetherian ascent. It need not be a domain, exactly as the source notes.

For every finite subset \(E\subset R\), take the finite-type domain \(D_E=R_0'[E]\) and its integral closure \(A_E\) in \(\operatorname{Frac}(D_E)\). Nagata finiteness makes \(A_E\) finite over \(D_E\), hence Noetherian and finite type over \(\mathbf Z\). The preceding denominator argument gives \(A_E\subset R\). If \(E\subset E'\), then \(\operatorname{Frac}(D_E)\subset\operatorname{Frac}(D_{E'})\), and every element of \(A_E\) is integral over \(D_{E'}\); consequently \(A_E\subset A_{E'}\). This is a directed family of normal Noetherian domains whose union is exactly \(R\): each \(r\in R\) lies in \(D_{\{r\}}\subset A_{\{r\}}\). Put \(B_E=A_E\otimes_{R_0'}S_0'\). Each \(B_E\) is normal by Noetherian smooth ascent, and the same finite-tensor comparison identifies \[
 \operatorname{colim}_E B_E\xrightarrow{\sim}S.
\]

Here is the full normal-colimit argument needed for the source's final sentence. Let \(T=\operatorname{colim}_i T_i\) be a filtered colimit of normal rings with unital transition maps, and let \(\mathfrak q\subset T\). Put \(\mathfrak q_i\) equal to its inverse image and \(D_i=(T_i)_{\mathfrak q_i}\). These are normal domains. The transition maps \(D_i\to D_j\) are local, because denominators outside \(\mathfrak q_i\) remain outside \(\mathfrak q_j\). They need not be injective. There is a canonical map \[
 \operatorname{colim}_iD_i\longrightarrow T_{\mathfrak q},
 \quad [a_i/s_i]\longmapsto a/s.
\] Every target fraction has its numerator and denominator represented at a common stage, with the denominator outside the corresponding inverse-image prime, proving surjectivity. If \(a_i/s_i\) maps to zero, there is \(t\in T\setminus\mathfrak q\) such that \(ta_i=0\) in \(T\). Represent \(t\) at a common later stage and move once more to a stage where that equality holds. Its representative remains outside the inverse-image prime, so \(a_i/s_i\) becomes zero in the localized ring at that stage. This proves injectivity. Subtracting two fractions gives the corresponding equality test. The displayed map is therefore an isomorphism of the original local rings.

The target is a domain: a product equal to zero is represented by a product equal to zero at some later \(D_j\), and since that ring is a domain, one factor is already zero there and hence in the colimit. The target is nonzero because it is localization at a prime. To prove integral closedness, take an actual fraction \(a/b\in\operatorname{Frac}(T_{\mathfrak q})\), \(b\ne0\), satisfying a monic equation with coefficients \(c_1,\ldots,c_n\in T_{\mathfrak q}\): \[
 (a/b)^n+c_1(a/b)^{n-1}+c_2(a/b)^{n-2}
 +\cdots+c_{n-1}(a/b)+c_n=0.
\] All finitely many numerators, denominators and coefficients occur in one \(D_i\). Multiplication by the original \(b^n\), followed by one eventual-equality witness, gives in a later \(D_j\) the complete relation \[
 a_j^n+c_{1,j}a_j^{n-1}b_j+c_{2,j}a_j^{n-2}b_j^2
 +\cdots+c_{n-1,j}a_jb_j^{n-1}+c_{n,j}b_j^n=0.
\] The element \(b_j\) is nonzero, since its image is the fixed nonzero \(b\). Thus \(a_j/b_j\) is integral over the normal domain \(D_j\), so it equals some \(d_j\in D_j\). The equation \(a_j=b_jd_j\) maps to \(a=bd\) in the target, giving \(a/b=d\in T_{\mathfrak q}\). This argument does not assert a map between entire fraction fields when a transition map has a kernel; it uses the specific denominator whose image is nonzero. Hence every prime localization of \(T\) is a normal domain, proving normality of \(T\). Applied to the original \(B_E\), this proves the smooth normal-ascent theorem and fills the stated omitted details in editorial material.

\subsubsection{Regular smooth ascent and report accounting}\label{algebra-proof-068-regular-smooth-ascent-and-report-accounting}

For the source's Noetherian regular-ring convention, smoothness makes \(S\) Noetherian with regular fibres. A regular ring and all these fibres satisfy \((R_k)\) for every integer \(k\ge0\). By the established ascent theorem, \(S\) satisfies every such condition. At each fixed prime \(\mathfrak q\), the Noetherian local ring \(S_{\mathfrak q}\) has a finite dimension \(d\); applying \((R_d)\) proves it regular. No bound on the global dimension of \(S\) is assumed or needed. Retained unit 0728 makes the quantified integer explicit.

Ten physical reports are accounted for in six groups. OCC-01015 is the missing ``by'', explicitly refining the unchanged designation inside prior operation MC-STK-ERR-0725-OP1. OCC-12436 retains that prior operation's zero-module branch. OCC-01016 supplies the colon at 46389. OCC-01017, OCC-01018 and OCC-12437 retain both commas in unit 0726. OCC-01019 and OCC-12438 retain unit 0727. OCC-01020 and OCC-12439 retain unit 0728. These are two additional operations and five retained operations in four canonical units. The additional copyedit must never erase the zero-module paragraph, and inverse replay must recover the entire current chapter including that paragraph and unrelated admitted material.

The consequences checked here feed the existing ascent chain: the zero-module depth convention is verified before the finite induction; the preceding Serre criteria give the reduced and normal Noetherian cases; the preceding Nagata finite-type theorem supplies the finite normal domains in smooth ascent. The filtered-colimit argument is a complete standard proof of the source's asserted step, not a new theorem claim or a requested translation rewrite. Remaining source review begins at 46633. The later Brauer and other research packets remain deferred until the core Algebra review is complete.

\hypertarget{algebra_descent_notes_20260929}{}
\subsection{Algebra: descending properties and the original catenarity example}\label{algebra-proof-069-algebra-descending-properties-and-the-original-catenarity-example}

This editorial review covers the Stacks Project authors' \texttt{algebra.tex} 46633--46877 at authority commit \texttt{a04446e57ec1fbc252a871afcec7752fb2807b14}, SHA-256 \texttt{FA8BB92E58A4F78A2BD01B3B6A4A87DE0A0D279F5DD90641B574DD5FBFFFA4F3}. The twelve received reports retain six existing units, MC-STK-ERR-0729--0734. No additional change to Algebra source text is proposed. Complete arguments and the linked Examples correction remain editorial.

Primary-source dependencies read for this review are \texttt{algebra.tex} 25455--25639, 27188--27213, 27518--27596, 42001--42051, and \texttt{examples.tex} 1410--1540 at the same commit. The Examples construction retains its attribution to Nagata, Appendix, Example 2, and its Eakin citation. The proof below supplies a direct Noetherianity argument for this particular construction. The \href{https://stacks.math.columbia.edu/tag/0355}{current Stacks remark, Tag 0355}, was checked only to confirm that its local-étale wording agrees with the frozen source. The mathematical proof below uses the original TeX and the exact displayed rings, with no new source-version admission or novelty claim.

\subsubsection{Faithful contraction, Noetherianity and the regularity receiver}\label{algebra-proof-069-faithful-contraction-noetherianity-and-the-regularity-receiver}

Let the original map be \(\varphi:R\to S\). For every ideal \(I\subset R\), faithful flatness gives an injective quotient unit \[
 R/I\longrightarrow S/IS,\qquad r+I\longmapsto\varphi(r)+IS.
\] Consequently \(\varphi^{-1}(IS)=I\), and the case \(I=0\) makes \(\varphi\) injective. This proves the exact interpretation of the source's \(R\cap IS\) notation. If \(S\) is Noetherian and \(I_0\subset I_1\subset\cdots\) is the original sequence of ideals, the extended sequence stabilizes at some \(n\). Contracting each equality gives \(I_n=I_{n+1}=\cdots\), so \(R\) is Noetherian. The retained singular ``assumption \ldots{} is'' in unit 0729 is grammatical and adequate; pluralizing both words would be an alternative wording, not another correction to compose.

If \(S\) is reduced, an element \(r\in R\) with \(r^e=0\) maps to a nilpotent of \(S\), hence to zero, and injectivity gives \(r=0\). This proves the original reduced descent. In particular only injectivity is needed for that individual implication; no stronger new source theorem is inserted.

If \(S\) is regular in the source's Noetherian convention, the preceding argument first proves that \(R\) is Noetherian. For \(\mathfrak p\subset R\), faithful flatness makes \(S\otimes_R\kappa(\mathfrak p)\) a nonzero ring; a prime of it gives \(\mathfrak q\subset S\) contracting to \(\mathfrak p\). The map \[
 R_{\mathfrak p}\longrightarrow S_{\mathfrak q},\qquad
 r/u\longmapsto\varphi(r)/\varphi(u)
\] is flat and local. Its target is regular, so the complete local regularity descent of group 634 applies. In that proof, for \(d=\dim S_{\mathfrak q}\ge2\) the original kernel is \(K_d=\ker(F_{d-1}\to F_{d-2})\), for \(d=1\) it is \(K_1=\ker(F_0\to\kappa(\mathfrak p))\), and for \(d=0\) the target is a field and locality plus injectivity makes the source a field. Tensoring the finite-free resolution preserves these kernels; global dimension \(d\) of the regular target makes the last finite syzygy projective. Faithful flatness descends its projectivity, and a finite projective module over the local source is free. The residue field therefore has finite projective dimension, giving regularity of the source. These details and the corrected exact-global-dimension reference are preserved in \texttt{ALGEBRA\_REGULAR\_AB\_NOTES\_20260928.md}, ``Flat descent with the actual low-degree kernels'', groups 634--635.

The faithfully flat global regularity assertion in earlier group 639 is now identified as already explicit source content at 46714--46732. Its earlier proof remains valid. Its status in a whole-work comparison is recovered source content, while the separately proved redundancy of the Noetherian hypothesis in the earlier local statement remains a strengthening of that stated hypothesis. No claim of a new global regularity theorem is carried forward.

\subsubsection{Normal descent with the actual principal-ideal map}\label{algebra-proof-069-normal-descent-with-the-actual-principal-ideal-map}

Assume the original \(S\) is normal and \(R\to S\) faithfully flat. Normal rings are reduced by local detection of nilpotents, so the preceding argument makes \(R\) reduced. Fix \(\mathfrak p\subset R\). Base change gives the faithfully flat map \[
 R_{\mathfrak p}\longrightarrow S_{\mathfrak p}
 =(R\setminus\mathfrak p)^{-1}S,
\] where the equality sends \((r/u)\otimes s\) to \(\varphi(r)s/\varphi(u)\), and the inverse sends \(s/\varphi(u)\) to \((1/u)\otimes s\). Unit 0730 restores the missing word ``flat''. Choosing a prime of the nonzero closed fibre and localizing gives a flat local map from \(R_{\mathfrak p}\) to a prime localization of \(S\), hence a faithfully flat map to a normal domain. These are the actual localizations of the original rings. The resulting injection makes \(R_{\mathfrak p}\) a domain.

Write these local rings as \(R,S\), as the source does after its reductions. Let \(a/b\in\operatorname{Frac}(R)\), with the required \(b\ne0\), satisfy its monic equation over \(R\). The injective domain map makes \(b\) nonzero in \(S\), so the fraction is defined in \(\operatorname{Frac}(S)\), satisfies the image equation with all its coefficients, and lies in \(S\) by normality. Thus \(a\in bS\). The original multiplication map \[
 R\longrightarrow R/bR,\qquad r\longmapsto ra+bR
\] becomes zero after tensoring with \(S\), because its image of 1 is zero in \(S/bS\). The quotient unit \(R/bR\to S/bS\) is injective, so \(a+bR=0\). Hence \(a=bd\) for some \(d\in R\), and the unchanged fraction \(a/b=d\) belongs to \(R\). This proves normality of every original \(R_{\mathfrak p}\) and therefore of the original ring.

\subsubsection{The minimal-fibre choice in Serre descent}\label{algebra-proof-069-the-minimal-fibre-choice-in-serre-descent}

Assume \(R\to S\) faithfully flat and \(S\) Noetherian. For an original prime \(\mathfrak p\subset R\), put \(F=S\otimes_R\kappa(\mathfrak p)\). It is nonzero by faithful flatness and Noetherian, being a quotient of a localization of \(S\). Choose a minimal prime \(\bar{\mathfrak q}\) of \(F\) and its corresponding \(\mathfrak q\subset S\). The exact local fibre maps from the preceding ascent review identify \[
 A=R_{\mathfrak p},\qquad B=S_{\mathfrak q},\qquad
 B/\mathfrak pB\cong F_{\bar{\mathfrak q}}.
\] Explicitly, if \(U=R\setminus\mathfrak p\), then \[
 \frac{\overline{s/\varphi(u)}}{\overline{t/\varphi(v)}}
 \longmapsto\overline{\frac{s\varphi(v)}{t\varphi(u)}}
\] maps the localized fibre to the quotient local ring; its inverse takes the class of \(s/t\) to \((s\otimes1)/(t\otimes1)\). The conditions \(u,v\notin\mathfrak p\) and \(t\notin\mathfrak q\) make every denominator invertible. Both composites are the identity by the original fraction relations.

The ring \(F_{\bar{\mathfrak q}}\) is a nonzero zero-dimensional Noetherian local ring. Its maximal ideal is nilpotent; if it is nonzero, a last nonzero power contains a nonzero element annihilated by the maximal ideal, while if it is zero the ring is a field. In either case its depth is zero. Its dimension is zero by the minimal-prime choice. Thus \(\mathfrak m_AB\) has radical \(\mathfrak m_B\) and is an ideal of definition. The original dimension and depth formulas give \[
 \dim B=\dim A+0,\qquad
 \operatorname{depth}B=\operatorname{depth}A+0.
\] If \(S\) satisfies \((S_k)\), then \[
 \operatorname{depth}A=\operatorname{depth}B
 \ge\min(k,\dim B)=\min(k,\dim A),
\] at each original \(\mathfrak p\). If \(S\) satisfies \((R_k)\) and \(\dim A\le k\), then \(\dim B=\dim A\le k\) makes \(B\) regular, and the complete local descent including its \(d=0,1\) cases makes \(A\) regular. Both conclusions retain source Noetherianity proved by faithful descent; no regularity or reducedness of the chosen zero-dimensional fibre is assumed.

\subsubsection{Nagata descent and every total-quotient denominator}\label{algebra-proof-069-nagata-descent-and-every-total-quotient-denominator}

Keep the original smooth map \(R\to S\) surjective on spectra and the Nagata ring \(S\). Smoothness gives flatness, and surjectivity on spectra makes it faithful; Noetherianity descends to \(R\). For the original finite-type map into a domain \(R\to A\), put \(B=A\otimes_RS\). The map \(A\to B\) is faithfully flat and smooth by base change; \(B\) is a finite-type algebra over Nagata \(S\), hence Nagata by the already proved finite-type theorem. It is reduced by the preceding smooth reduced-ascent argument applied to the unchanged domain \(A\).

Let \(\mathfrak q_1,\ldots,\mathfrak q_t\) be its finitely many minimal primes. The ring \(B\ne0\), because \(A\ne0\) and the map is faithfully flat. Each \(\mathfrak q_i\) contracts to \((0)\subset A\): if its contraction contained a nonzero prime chain above \((0)\), flat going down would give a strictly smaller prime below \(\mathfrak q_i\). Unit 0731 terminates this sentence with a period. For the original total quotient ring, \[
 Q(B)=\prod_{i=1}^t L_i,\qquad L_i=\kappa(\mathfrak q_i).
\] The reduced Noetherian total-quotient comparison is the one already proved in the Nagata note. The integral closure \(B'\) equals the product of the closures \(B_i'\) of \(B/\mathfrak q_i\) in \(L_i\): an integral tuple has integral coordinates, while conversely the product of the finitely many lifted monic coordinate equations is a monic polynomial vanishing on the tuple in every coordinate. Each \(B_i'\) is finite over \(B/\mathfrak q_i\) by the Nagata property, and the product is therefore finite over \(B\).

Unit 0732 defines the previously unnamed \(A'\) to be the integral closure of the original \(A\) in \(Q(A)=\operatorname{Frac}(A)\). Each \(d\in A\setminus\{0\}\) acts injectively on \(A\); flatness makes its image act injectively on \(B\). Hence the map \[
 Q(A)\otimes_AB
 =(A\setminus\{0\})^{-1}A\otimes_AB
 \xrightarrow{\sim}(A\setminus\{0\})^{-1}B
 \longrightarrow Q(B)
\] sends \((a/d)\otimes b\) to \(\varphi(a)b/\varphi(d)\). The middle inverse sends \(b/\varphi(d)\) to \((1/d)\otimes b\). The last map is injective: if such a fraction is zero after inverting all nonzerodivisors of \(B\), a nonzerodivisor \(u\in B\) annihilates its numerator, forcing that numerator to be zero already. No nonzero denominator of \(A\) was discarded or assumed invertible in \(B\) itself.

Flatness injects \(A'\otimes_AB\) into \(Q(A)\otimes_AB\). The image lies in \(B'\), since the image of each element of \(A'\) satisfies its original monic equation over the image of \(A\), and the integral closure is a subring containing \(B\). This proves the original injection \[
 A'\otimes_AB\hookrightarrow B'.
\] As \(B\) is Noetherian and \(B'\) finite, its submodule \(A'\otimes_AB\) is finite. Express a finite set of its generators as finite sums of actual tensors \(x_i\otimes b_i\); collect all finitely many \(x_i\in A'\). The elements \(x_i\otimes1\) generate, because each original tensor generator is their \(B\)-linear combination. Put \(M=A'/\sum_iAx_i\). Right exactness gives \(M\otimes_AB=0\), and faithful flatness gives \(M=0\). Thus the original \(x_i\) generate \(A'\), as unit 0733 states in the present tense. Every finite-type domain algebra over \(R\) is N-1, so the previously proved criterion makes \(R\) universally Japanese and, together with its Noetherianity, Nagata.

\subsubsection{The original series construction and its missing coefficient factor}\label{algebra-proof-069-the-original-series-construction-and-its-missing-coefficient-factor}

Retain the Examples data: a field \(k\), a series \[
 z=\sum_{i=1}^{\infty}a_ix^i\in k[[x]]
\] transcendental over \(k(x)\), and \[
 z_j=x^{-j}\left(z-\sum_{i=1}^{j-1}a_ix^i\right)
     =\sum_{i=j}^{\infty}a_ix^{i-j},\qquad j\ge1.
\] Such data exist, for example with \(k=\mathbf Q\): the elements algebraic over the countable field \(\mathbf Q(x)\) form a countable set, while \(x\mathbf Q[[x]]\) is uncountable. Thus the construction really supplies counterexamples. No specific choice of its infinitely many coefficients is suppressed in the calculations.

For \(R=k[x,z_1,z_2,\ldots]\subset k[[x]]\), the exact relations are \[
 z=xz_1,\qquad xz_{j+1}+a_j=z_j,
 \qquad z_{j+1}=x^{-1}(z_j-a_j).
\] The last equality applies where \(x\) is inverted. \textbf{The source at \texttt{examples.tex:1481} writes \(x^{-1}z_j-a_j\), omitting the factor \(x^{-1}\) on \(a_j\).} Its difference from the correct expression is \(a_j(x^{-1}-1)\). At least one \(a_j\ne0\), since the stipulated series is transcendental, and for that index the difference is nonzero in \(k((x))\). This is a source correction in the cited Examples dependency, not an Algebra text operation. The construction and the claimed local-ring isomorphism remain valid with the exact recurrence.

Each \(z_j\) is transcendental over \(k(x)\), since its displayed formula would otherwise make \(z\) algebraic. Hence \(R_j=k[x,z_j]\) is a polynomial ring in the two indicated elements, and \[
 R=\operatorname{colim}_{j\ge1}R_j,\qquad
 R_j\to R_{j+1}:x\mapsto x,\ z_j\mapsto xz_{j+1}+a_j.
\] The maps are the actual injections inside \(k[[x]]\). Modulo \(x-1\), the transition is the invertible translation \(z_j\mapsto z_{j+1}+a_j\). Quotients commute with this colimit by their element relations. Therefore the quotient is exactly \(k[z]\), with \[
 x\mapsto1,\qquad z_j\mapsto z-\sum_{i=1}^{j-1}a_i.
\] The inverse sends the indeterminate \(z\) to the class of the original \(z=xz_1\). The formula verifies both composites at every generator. This supplies a direct proof of the source's quotient assertion, preserving every coefficient.

For the original \(\mathfrak m=(x)\), the relations show that reduction modulo \(x\) sends \(z_j\) to \(a_j\), so \(R/\mathfrak m=k\). If a nonzero \(r\in R\) has \(x\)-adic order \(v\) in \(k[[x]]\), vanishing of its constant term means \(r\in xR\). Repeating this exact divisibility \(v\) times gives \(r=x^vu\), with \(u\in R\setminus\mathfrak m\). Thus every nonzero element of \(R_{\mathfrak m}\) has this unique order and a unit factor, and any nonzero ideal is generated by an element of minimum order. It is a DVR with parameter \(x\) and residue field \(k\).

The quotient calculation shows that the original \[
 \mathfrak n=(x-1,z_1,z_2+a_1,z_3+a_1+a_2,\ldots)
\] equals \((x-1,z)\), and has residue field \(k\). Inverting the original \(x\) gives \[
 R[1/x]=k[x,x^{-1},z],
\] with every \(z_j\) given by its full displayed expression. Since \(x\notin\mathfrak n\), localizing gives the exact isomorphism \[
 k[x,x^{-1},z]_{(x-1,z)}\xrightarrow{\sim}R_{\mathfrak n}.
\] It sends \(x,z\) to the original elements; its inverse sends each \(z_j\) to \(x^{-j}(z-\sum_{i<j}a_ix^i)\) and each inverted denominator to its inverse. These formulas are mutually inverse on generators and respect every localization denominator. The ring is regular local of dimension 2, with parameters \(x-1,z\), and residue field \(k\).

Let the source's multiplicative set be \(W=R\setminus(\mathfrak m\cup\mathfrak n)\), and keep its original ring \(B=W^{-1}R\). Its maximal ideals are precisely \(\mathfrak mB,\mathfrak nB\): a prime disjoint from \(W\) is contained in their union, hence in one of them by the two-ideal prime-avoidance argument. Both ideals survive and are maximal. Their local rings are the original \(R_{\mathfrak m},R_{\mathfrak n}\), with fractions mapped to the same original fractions and inverse localization maps. To prove \(B\) Noetherian, take any ideal \(I\). At each of its two maximal ideals choose finitely many elements of \(I\) whose images generate the localized ideal. Let \(I_0\) be the ideal generated by the union of these two finite sets. Then \((I/I_0)\) vanishes at both maximal ideals and hence at every maximal ideal; an element of a nonzero module would have annihilator contained in some maximal ideal and would survive there. Thus \(I=I_0\). Every prime localization of \(B\) is a localization of one of its two regular local rings, so \(B\) is regular, a domain, and has dimension 2.

\subsubsection{The conductor ring and a finite étale counterexample}\label{algebra-proof-069-the-conductor-ring-and-a-finite-uxe9tale-counterexample}

Keep \(J=\mathfrak mB\cap\mathfrak nB\) and the original \(A=k+J\subset B\). The two maximal ideals are comaximal, so \(B/J\cong k\times k\); the image of \(A\) is the diagonal copy of \(k\). Choose \(b\in B\) whose residues are \((1,0)\), in that order. For an arbitrary \(v\in B\) with residues \((v_1,v_2)\), \[
 v-\bigl(v_2+(v_1-v_2)b\bigr)\in J\subset A.
\] Thus \(B=A+Ab\). Also \(c=b^2-b\in J\subset A\). These are exact elements of the original rings.

The ring \(A\) is local with maximal ideal \(J\): an element with nonzero common residue is a unit of \(B\), whose inverse has the same inverse residue in both coordinates and therefore belongs to \(A\). It is a domain. The module \(B/A\) is isomorphic to \(k\) by the residue difference and is killed by \(J\). Here is a direct proof that \(A\) is Noetherian. For an ideal \(I\subset A\), Noetherianity of \(B\) lets us choose \(i_1,\ldots,i_r\in I\) generating \(IB\) over \(B\), by expanding any finite generating set into the original \(I\)-elements. Put \(I_0=\sum Ai_j\). The module \(IB/I_0\) is a quotient of \((B/A)^r\), since the surjection \(B^r\to IB/I_0\), \((v_j)\mapsto\sum v_ji_j+I_0\), kills \(A^r\). It is consequently a finite-dimensional \(k\)-vector space. Its submodule \(I/I_0\) is another finite-dimensional \(k\)-vector space, so finitely many lifts together with the \(i_j\) generate \(I\) over \(A\). This proves Noetherianity without leaving a finiteness assertion to be assumed.

The nonzero element \(x(x-1)\) lies in \(J\); for any \(v\in B\), its product with this element belongs to \(J\subset A\). Hence \(\operatorname{Frac}(A)=\operatorname{Frac}(B)=K\). The extension is finite, since \(B=A+Ab\), so integral dimension comparison gives \(\dim A=\dim B=2\). If \(A\) were universally catenary, the source dimension formula at 27519--27596 for the original finite birational extension and the prime \(\mathfrak mB\) would give \[
 \operatorname{height}(\mathfrak mB)
 =\operatorname{height}(J)
   +\operatorname{trdeg}_{\operatorname{Frac}(A)}\operatorname{Frac}(B)
   -\operatorname{trdeg}_{\kappa(J)}\kappa(\mathfrak mB)
 =2+0-0=2.
\] The actual DVR calculation gives height 1. Thus the unchanged \(A\) is not universally catenary. The dimensions in the Algebra remark list the two maximal ideals in the opposite order; the two original local rings and their residue fields are the same ones, with this permutation made explicit.

Now form an explicit algebra over this original \(A\): \[
 C=A[T]/(T^2-T-c),\qquad t=[T],\qquad c=b^2-b.
\] Monic division proves that \(1,t\) is an \(A\)-basis. Thus \(C\) is finite free of rank 2 and faithfully flat: for every \(A\)-module \(M\), the basis identifies \(M\otimes_AC\) with \(M\oplus M\). The equality \[
 (2t-1)^2=1+4c
\] retains both numerical coefficients in every characteristic. Since \(c\in J\), \(1+4c\) is a unit of the local ring \(A\), and the derivative \(2t-1\) has inverse \((2t-1)/(1+4c)\) in \(C\). Consequently \(C\) is étale. Explicitly, for a square-zero ideal \(I\) in an \(A\)-algebra \(D\), a root modulo \(I\) has a lift \(r\); its unique root correction is \[
 \delta=-\frac{r^2-r-c}{2r-1}\in I.
\] The derivative is a unit because it is a unit modulo the nilpotent ideal. Expanding retains the full identity \[
 (r+\delta)^2-(r+\delta)-c
 =(r^2-r-c)+(2r-1)\delta+\delta^2=0.
\] The same identity proves uniqueness. Together with finite presentation, this proves formal étaleness and étaleness of the actual quadratic algebra.

There is an exact base-change decomposition \[
 B\otimes_AC=B[T]/\bigl((T-b)(T-(1-b))\bigr)
 \xrightarrow{\sim}B\times B,
 \qquad f(t)\longmapsto(f(b),f(1-b)).
\] The root difference \(2b-1\) is a unit in \(B\): its residues are \((1,-1)\). The inverse sends \((v,w)\) to \[
 v\frac{t-(1-b)}{2b-1}
 +w\frac{b-t}{2b-1}.
\] Evaluation proves both composites, and the two displayed fractions are the orthogonal idempotents summing to 1. This retains the actual numerator, sign and denominator in the splitting.

The two evaluation maps \(C\to B\), \(t\mapsto b\) and \(t\mapsto1-b\), are surjective because \(B=A+Ab=A+A(1-b)\). The product map \[
 C\hookrightarrow B\times B,\qquad
 a+dt\longmapsto(a+db,a+d(1-b))
\] is injective: equality to zero in both coordinates gives \(d(2b-1)=0\), hence \(d=0\) and \(a=0\). Let \(I_+,I_-\) be the two kernels. They are prime and have intersection zero. They are distinct and incomparable: write the original \(b=a_0/d_0\) with \(a_0,d_0\in A\), \(d_0\ne0\), using the proved equality of fraction fields. Then \[
 d_0t-a_0\in I_+\setminus I_-,\qquad
 d_0t-(d_0-a_0)\in I_-\setminus I_+,
\] because their other evaluations are respectively \(d_0(1-2b)\) and \(d_0(2b-1)\), both nonzero. Every prime of \(C\) contains one of these ideals: otherwise a product of elements chosen outside that prime, one from each ideal, would be zero. Thus these are its two minimal primes and both component quotients are the original regular domain \(B\).

The ring \(B\) is universally catenary. Its regular local rings are Cohen--Macaulay, and the source's Cohen--Macaulay catenarity theorem at 25606--25638 proves universal catenarity: a saturated prime chain has length equal to the difference of the two local dimensions, and polynomial extensions remain Cohen--Macaulay. It follows directly that \(C\) is universally catenary. Indeed, take any prime interval in \(C[X_1,\ldots,X_n]\). Its lower prime contracts to a prime of \(C\) containing \(I_+\) or \(I_-\); the whole interval then contains the extended chosen ideal. Quotienting identifies all its intermediate primes, inclusions and saturated chains with the same interval in \(B[X_1,\ldots,X_n]\). Lengths are bounded and all saturated chains have the same length there. This proves catenarity of every polynomial algebra over \(C\); quotient prime correspondence then gives it for every finite-type \(C\)-algebra. No locality of \(C\) or unproved finite presentation of a localization was used.

We have therefore proved the stronger precise counterexample: \textbf{the original non-universally-catenary Noetherian local domain \(A\) has a finite, faithfully flat, étale algebra \(C\) of rank 2 which is universally catenary.} This is a complete construction from the source's original \(A,B\), not a novelty claim.

\subsubsection{The local comparison and report accounting}\label{algebra-proof-069-the-local-comparison-and-report-accounting}

The algebra \(C\) has exactly two maximal ideals over \(J\), since it is finite over local \(A\) and \[
 C/JC=k[T]/(T(T-1))\cong k\times k.
\] Let \(\mathfrak r_0=(J,t)\). Its localization \(C_{\mathfrak r_0}\) is a local ring essentially étale over \(A\), and is universally catenary by localization of the proved \(C\). The exact split algebra becomes \[
 B\otimes_AC_{\mathfrak r_0}
 \cong B_{\mathfrak nB}\times B_{\mathfrak mB}.
\] For the first evaluation, \(t\mapsto b\), the inverse image of \(\mathfrak nB\) is \(\mathfrak r_0\), because \(b\) has residue zero there; for the second evaluation the corresponding prime is \(\mathfrak mB\). Surjectivity of the evaluation maps shows that the image of each denominator outside \(\mathfrak r_0\) ranges over the elements outside the appropriate target maximal ideal. Thus the localization maps send an original fraction to its evaluated numerator divided by its evaluated denominator; inverses lift each target numerator and denominator through the surjection, with independence following from the quotient relation. The two quotient local rings have the originally proved dimensions 2 and 1 and residue field \(k\). The two distinct minimal-prime quotients of \(C_{\mathfrak r_0}\) are these actual local rings.

The source's cited decomposition lemma supplies an étale algebra with a specified prime, and localization at that prime supplies a local essentially étale algebra. Arbitrary prime localization is not itself a justification of finite presentation. The explicit \(C\) above proves the stated étale non-descent conclusion, in the stronger finite form, and its exact localization proves the local comparison. The original stronger wording ``local étale ring map'' is not certified by silently treating ``essentially étale'' as ``étale''; nor is a falsehood claim about all possible local étale constructions inferred from this particular construction. The source passage stays identifiable and this clarification is editorial.

The twelve reports are paired under six existing units: OCC-01021/12440 under 0729, 01022/12441 under 0730, 01023/12442 under 0731, 01024/12443 under 0732, 01025/12444 under 0733, and 01026/12445 under 0734. The last unit's explicit comparison of each residue field is retained. Separate records hold the finite quadratic counterexample and the Examples recurrence correction. The regularity, depth, smooth-reducedness, Nagata and unramified-chart dependencies are propagated from their earlier verified records. Later synthesis and optional paper mining remain after the core review; the next authority section starts at 46878.

\hypertarget{algebra_geonormal_notes_20260929}{}
\subsection{Algebra: geometric normality and the original tensor localizations}\label{algebra-proof-070-algebra-geometric-normality-and-the-original-tensor-localizations}

This editorial verification covers \texttt{algebra.tex} 46878--47091 at Stacks Project authority \texttt{a04446e57ec1fbc252a871afcec7752fb2807b14}, SHA-256 \texttt{FA8BB92E58A4F78A2BD01B3B6A4A87DE0A0D279F5DD90641B574DD5FBFFFA4F3}. The \href{https://github.com/stacks/stacks-project/blob/a04446e57ec1fbc252a871afcec7752fb2807b14/algebra.tex\#L46878}{original source} remains the authority. Eight reports retain units MC-STK-ERR-0735--0737 with six existing operations. A designation receives one additional ``by''; the elliptic coordinated assumption is recorded as optional. Complete mathematical arguments remain linked editorial material, without replacing the translated source.

The source dependencies read here are 8217--8235, 8274--8309, 10062--10081, 10151--10194, 10220--10255, 10512--10542 and 44511--44524. Earlier complete proofs in groups 198, 204, 206, 268, 273, 275 and 282, and the completed smooth-normality and faithful-normality ascent/descent notes, are receiving dependencies, not new discoveries. No content of the four new literature-index hits was read.

\subsubsection{The four field-extension tests and the original inseparable tensor}\label{algebra-proof-070-the-four-field-extension-tests-and-the-original-inseparable-tensor}

Let \(k\) be the original field and \(A\) the original \(k\)-algebra. Normality means that every prime localization is an integrally closed domain. The zero algebra has no primes and is normal; every scalar extension of it is again zero, so all the stated tests hold for it. This case introduces no exclusion into the source theorem.

For an arbitrary field extension \(E/k\), the map \(A\to E\otimes_kA\) is faithfully flat. For every \(A\)-module \(M\), the actual comparison \[
 (E\otimes_kA)\otimes_AM\xrightarrow{\sim}E\otimes_kM,
 \quad (e\otimes a)\otimes m\longmapsto e\otimes am
\] has inverse \(e\otimes m\mapsto(e\otimes1)\otimes m\). Tensoring vector spaces over \(k\) is exact and reflects zero when \(E\ne0\); a basis of \(E\) identifies this functor with a nonempty direct sum of copies of the underlying \(M\). This also proves faithful flatness after any of the specified algebra base changes below.

The implications from all field extensions to finitely generated ones, from those to finite purely inseparable ones, and from all extensions to \(k^{\mathrm{perf}}\), follow by inclusion of the quantified classes. If \(k'/k\) is finite purely inseparable, its unique embedding into a chosen perfect closure is a map of the original \(k\)-extensions. Thus \[
 k'\otimes_kA\longrightarrow k^{\mathrm{perf}}\otimes_kA
 \cong k^{\mathrm{perf}}\otimes_{k'}(k'\otimes_kA)
\] is faithfully flat. The displayed associativity map sends \(u\otimes(v\otimes a)\) to \(uv\otimes a\), with inverse \(u\otimes a\mapsto u\otimes(1\otimes a)\). Normality of the target descends by the completed original normal-descent proof. In characteristic zero, the finite purely inseparable extensions and the perfect closure are the identity extension, so this step contains no positive-characteristic exponent assumption.

Every extension \(E/k\) is the directed union of its finitely generated intermediate fields: take the field generated by each finite set of elements, and use the union of two finite sets for a common upper index. The comparison \[
 \operatorname{colim}_i(k_i\otimes_kA)\xrightarrow{\sim}E\otimes_kA
\] sends a represented tensor to that same tensor. Each tensor and each tensor relation involve finitely many elements and hence occur at one sufficiently large field. This proves both surjectivity and injectivity and respects addition and multiplication. Normality passes to this colimit by the complete original prime-localization proof in group 206: contract a target prime, identify its local ring with the colimit of the normal local domains, descend each finite monic equation to a later stage, and transport the resulting equality \(a=bd\) in the stage ring. No map of entire fraction fields is required. Thus the finitely generated test implies the test for all fields.

It remains to derive the finitely generated test from the finite purely inseparable one. For the original finitely generated \(K/k\), the source field-diagram lemma supplies \[
 \begin{matrix}K&\longrightarrow&K'\\ \uparrow&&\uparrow\\ k&\longrightarrow&k',\end{matrix}
\] where \(k'/k\) and \(K'/K\) are finite purely inseparable, and \(K'/k'\) is separable. The exact construction is already proved in group 282. It permits \(K'=k'K\), but does not assert that \(k'\otimes_kK\) itself is a field. Retain the multiplication map \[
 \mu:k'\otimes_kK\longrightarrow k'K,\qquad
 \sum_i\lambda_i\otimes u_i\longmapsto\sum_i\lambda_iu_i.
\] In characteristic \(p>0\), choose the original finite exponent \(p^e\) with \((k')^{p^e}\subset k\). Its full Frobenius sum is \[
 \left(\sum_i\lambda_i\otimes u_i\right)^{p^e}
 =1\otimes\sum_i\lambda_i^{p^e}u_i^{p^e}.
\] All mixed multinomial coefficients vanish in the stated characteristic. If the original sum lies in \(\ker\mu\), the right-hand scalar is zero in the embedded original \(K\). Thus its \(p^e\)-th power is zero. Conversely every nilpotent maps to zero in the field. The image is a finite integral domain algebra over \(K\) containing \(K\) and \(k'\), hence is the compositum field: for a nonzero element its monic minimal equation over \(K\) has nonzero constant term and expresses its inverse in the image. Therefore \(\ker\mu\) is exactly the nilradical and \((k'\otimes_kK)_{\mathrm{red}}\cong k'K\). This records the quotient and both field inclusions; it does not discard the nilpotents of the original tensor. In characteristic zero the diagram uses the identity purely inseparable maps.

The finitely generated separable extension \(K'/k'\) is the fraction field of a smooth \(k'\)-domain \(B\). This is the source's smooth-model result 44511--44524: choose a finite-type domain with that fraction field and restrict to a nonzero smooth neighbourhood of its generic point. The complete generic smoothness argument is retained in \texttt{ALGEBRA\_FIELD\_FORMAL\_SMOOTH\_NOTES\_20260929.md}.

The finite purely inseparable test makes \(k'\otimes_kA\) normal. Base change of the actual smooth map makes \[
 k'\otimes_kA\longrightarrow B\otimes_{k'}(k'\otimes_kA)
\] smooth, so the target is normal by the preceding smooth-ascent theorem. The map \(b\otimes(\lambda\otimes a)\mapsto b\lambda\otimes a\), with inverse \(b\otimes a\mapsto b\otimes(1\otimes a)\), identifies it with \(B\otimes_kA\). Localizing at the images of all nonzero original \(b\in B\) gives \[
 K'\otimes_{k'}(k'\otimes_kA)\cong K'\otimes_kA.
\] On fractions, \((b/s)\otimes a\) corresponds to \((b\otimes a)/(s\otimes1)\); the inverse preserves the same \(s\). Normality is preserved under this localization, including a zero result if it occurs. Finally \(K\otimes_kA\to K'\otimes_kA\) is the base change of the field extension \(K'/K\), so it is faithfully flat and normality descends. This proves the last nontrivial implication among the original four conditions.

\subsubsection{Localization and separable field extensions}\label{algebra-proof-070-localization-and-separable-field-extensions}

For any multiplicative subset \(U\subset A\), any field extension \(E/k\) has the exact comparison \[
 E\otimes_kU^{-1}A\xrightarrow{\sim}(1\otimes U)^{-1}(E\otimes_kA),
 \quad e\otimes(a/u)\longmapsto(e\otimes a)/(1\otimes u).
\] Its inverse sends \((\sum_i e_i\otimes a_i)/(1\otimes u)\) to \(\sum_i e_i\otimes(a_i/u)\). Tensor relations and the same multiplicative-set equality witnesses prove both inverse identities. This also handles \(0\in U\), when both sides are zero. Since a prime localization of a localized ring is the corresponding original prime localization, normality localizes. Applying the four-condition definition proves the original localization assertion for geometric normality.

Let \(K/k\) be the source's separable field extension, without a finite-generation assumption. Group 273 proves the original reducedness assertion, so \(T=k^{\mathrm{perf}}\otimes_kK\) is reduced. The purely inseparable extension \(k^{\mathrm{perf}}/k\) is integral and radicial, and its field base change has a unique prime. One can check this with the original elements: for any \(z=\sum_i\lambda_i\otimes u_i\in T\), some finite exponent \(p^e\) puts every \(\lambda_i^{p^e}\) in \(k\), giving the same full Frobenius sum as above in the original \(K\). If this scalar is zero, \(z\) is nilpotent; if it is nonzero, \(z\) is a unit with inverse its \((p^e-1)\)-st power divided by that scalar. Thus every nonunit is nilpotent, and the nilradical is the only prime. Reducedness then makes every nonzero element a unit. The tensor is nonzero by faithful scalar extension, hence is a field. In characteristic zero it is just \(K\). The perfect-closure test proves geometric normality of the original \(K\).

\subsubsection{The original local reduction and normal finite subalgebras}\label{algebra-proof-070-the-original-local-reduction-and-normal-finite-subalgebras}

Assume now that \(A\) is geometrically normal over \(k\), and \(B\) is normal. If either algebra is zero, their tensor is zero and normal. Otherwise take an original prime \(\mathfrak r\subset A\otimes_kB\) and its contractions \(\mathfrak p,\mathfrak q\). Both factor maps to the target localization invert the elements outside these contractions. The map \[
 A_{\mathfrak p}\otimes_kB_{\mathfrak q}
 \longrightarrow (A\otimes_kB)_{\mathfrak r},\qquad
 (a/u)\otimes(b/v)\longmapsto (a\otimes b)/(u\otimes v)
\] therefore exists. The middle ring is the localization of \(A\otimes_kB\) at the multiplicative set generated by \(u\otimes1\), \(u\notin\mathfrak p\), and \(1\otimes v\), \(v\notin\mathfrak q\). The inverse to this first localization comparison sends these two types of inverse generators to \((1/u)\otimes1\) and \(1\otimes(1/v)\), respectively. Further localization at the prime induced by \(\mathfrak r\) yields exactly \((A\otimes_kB)_{\mathfrak r}\), with both composites fixing every original fraction. This proves the source reduction with its actual contractions. The additional ``by'' at 47000 only repairs their designation.

The localized \(A_{\mathfrak p}\) is geometrically normal by the previous result and is a normal domain by the identity field test; \(B_{\mathfrak q}\) is a normal domain. Retain the source's names \(A,B\) for these reduced cases and its fraction fields \(K=\operatorname{Frac}(A)\), \(L=\operatorname{Frac}(B)\). Unit 0736 corrects only ``fractions fields'' at both occurrences.

For every finite set \(E\subset B\), let \(D_E=k[E]\) and take its integral closure \(B_E\) in its own fraction field \(\operatorname{Frac}(D_E)\subset L\). Since fields are Nagata and the earlier finite-type theorem applies, \(B_E\) is finite over \(D_E\), hence a finite-type normal \(k\)-domain. Each element of \(B_E\) is integral over \(B\) and belongs to \(L\), so normality of the unchanged domain \(B\) puts it in \(B\). If \(E\subset E'\), the fraction-field inclusion sends any element integral over \(D_E\) to an element integral over \(D_{E'}\), giving the actual inclusion \(B_E\subset B_{E'}\). This is directed and its union is \(B\), since it contains every finite set \(E\). Consequently \[
 A\otimes_kB\cong\operatorname{colim}_E(A\otimes_kB_E)
\] by the same finite-tensor and finite-relation argument. Group 206 then reduces the desired normality to the case of the original finite-type normal domain \(B\). No Noetherian assumption on \(A\) is introduced. Unit 0737 closes all three occurrences of ``sub algebra'', including the locus omitted from one report's proposed count.

\subsubsection{Both tensor localizations inside the total quotient ring}\label{algebra-proof-070-both-tensor-localizations-inside-the-total-quotient-ring}

Write \(C=A\otimes_kB\), with the original domains and \(B\) now finite type over \(k\). It is nonzero: tensoring two nonzero vector spaces over a field does not kill \(1\otimes1\). It is flat over either factor. In particular every nonzero \(a\in A\) gives a nonzerodivisor \(a\otimes1\) on \(C\), and every nonzero \(b\in B\) gives a nonzerodivisor \(1\otimes b\). Let \(U\) and \(V\) be these respective multiplicative subsets. Keep the exact identifications \[
 U^{-1}C\cong K\otimes_kB=:D,\qquad
 V^{-1}C\cong A\otimes_kL=:E,\qquad
 (UV)^{-1}C\cong K\otimes_kL=:F.
\] They take the fractions \((a\otimes b)/(u\otimes1)\), \((a\otimes b)/(1\otimes v)\), and \((a\otimes b)/(u\otimes v)\) to \((a/u)\otimes b\), \(a\otimes(b/v)\), and \((a/u)\otimes(b/v)\), respectively; their inverses keep the displayed denominators. Every denominator is a nonzerodivisor in \(C\), so all four rings inject into the actual \(Q(C)\). Thus the comparison direction is \[
 C\subset D,E\subset F\subset Q(C).
\] No map from all of \(Q(C)\) into \(F\) is assumed. Such a map fixing \(C\) need not exist: for \(A=k[x],B=k[y]\), \(x+y\) is invertible in \(k(x,y)=Q(C)\), but is not a unit in the factor-denominator localization \(k(x)\otimes_k k(y)\). Indeed a relation \((x+y)h=f(x)g(y)\) with nonzero one-variable denominators would, after setting \(y=-x\), give \(0=f(x)g(-x)\) in the domain \(k[x]\), a contradiction. This verifies the direction of the actual embeddings used in the proof.

Reducedness of \(C\), needed by the source's finite-minimal-prime criterion, is available before invoking that criterion: geometric normality of \(A\) makes \(E=A\otimes_kL\) normal and therefore reduced. The injection \(C\hookrightarrow E\) makes \(C\) reduced. The normal ring \(E\) here uses the lowercase ground field, exactly the correction retained as MC-STK-ERR-0735. In general \(A\) is not even a \(K\)-algebra, so the printed uppercase tensor base has no justification.

The ring \(D=K\otimes_kB\) is a finite-type algebra over the field \(K\), hence Noetherian and has finitely many minimal primes. Every minimal prime of \(C\) contracts to zero in the domain \(A\): flat going down would otherwise put a strictly smaller prime below it. Hence every such prime avoids \(U\). Localization gives a bijection between minimal primes of \(C\) and minimal primes of \(D\). For the converse direction, contract a minimal prime of \(D\), take a minimal prime of \(C\) below it, and localize; minimality forces equality, and contraction recovers the original prime. Thus \(C\) has finitely many minimal primes, with no assumed Noetherianity of \(C\).

Here is the total-quotient comparison needed to use normality of both localizations. For any multiplicative subset \(W\) of nonzerodivisors of a ring \(C\), multiplication by a nonzerodivisor of \(C\) remains injective after localization. Conversely, if \(b/w\) is a nonzerodivisor on \(W^{-1}C\), then \(b\) is a nonzerodivisor on \(C\): from \(bc=0\) one gets \((b/w)(c/1)=0\), then \(c/1=0\), and injectivity of \(C\to W^{-1}C\) gives \(c=0\). Therefore the mutually inverse total-quotient maps are \[
 Q(C)\longrightarrow Q(W^{-1}C),\qquad
 a/b\longmapsto(a/1)/(b/1),
\] \[
 Q(W^{-1}C)\longrightarrow Q(C),\qquad
 (a/u)/(b/v)\longmapsto av/(ub).
\] In the second formula, \(u,v\in W\) and the just-proved criterion ensures that the original \(b\) is a nonzerodivisor; thus \(ub\) is an allowed denominator. Localization equality witnesses verify well-definedness, and the fraction formulas verify both inverse identities. Apply this to \(W=U\) and \(W=V\). These are identifications inside the original \(Q(C)\), rather than replacements of its objects by an unproved ambient field.

\subsubsection{Normality of the first localization and the exact intersection}\label{algebra-proof-070-normality-of-the-first-localization-and-the-exact-intersection}

Geometric normality of \(A\) implies geometric normality of its localization \(K\), hence geometric reducedness of \(K\). Write \(K=\bigcup_iK_i\) over all finitely generated intermediate fields. Each \(K_i\) is geometrically reduced: for every field extension of \(k\), its scalar extension injects into that of \(K\), because tensoring over the field preserves the original inclusion. A subring of a reduced ring is reduced. Thus \(K_i/k\) is separable in the source's equivalent field criterion.

For each original \(K_i\), choose the smooth \(k\)-domain \(D_i\) with fraction field \(K_i\) supplied by 44511--44524. Then \(B\to D_i\otimes_kB\) is the actual base change of a smooth map and is smooth. Since the original \(B\) is normal, smooth ascent makes this tensor normal. Localizing at the images of \(D_i\setminus\{0\}\) gives the normal ring \(K_i\otimes_kB\), with the map \((d/s)\otimes b\leftrightarrow(d\otimes b)/(s\otimes1)\) and its exact inverse. The natural transition maps are induced by the original field inclusions; no chosen smooth models need have compatible transition maps. Their tensor colimit is the original \(D=K\otimes_kB\). Group 206 proves it normal. This argument uses the earlier smooth-normality proof, which is independent of the present tensor-normality theorem, so there is no circular application.

Both \(D\) and \(E\) are now normal. Prove their intersection in the actual common ring \(F\) by tensoring the original vector-space sequence \[
 0\longrightarrow A\longrightarrow K\longrightarrow K/A\longrightarrow0
\] first with \(B\), then with \(L\). The comparison induced by \(B\hookrightarrow L\) is injective on \((K/A)\otimes_kB\), as tensoring over \(k\) is exact. If \(z\in K\otimes_kB\) has image in \(A\otimes_kL\), its image in \((K/A)\otimes_kL\) is zero. Injectivity forces its image in \((K/A)\otimes_kB\) to be zero, so exactness in the first row puts \(z\) in \(A\otimes_kB\). Conversely this last tensor lies in both rings. Hence, with all inclusions already proved, \[
 (K\otimes_kB)\cap(A\otimes_kL)=A\otimes_kB
\] inside \(F\), and therefore also inside \(Q(C)\).

Let the original \(z\in Q(C)\) be integral over \(C\). Its unchanged monic equation has coefficients in \(D\) as well as in \(E\). The exact total-quotient comparisons make it an element of both \(Q(D)\) and \(Q(E)\). A normal ring is integrally closed in its total quotient ring, by the complete original ideal and two-localization-multiplier proof in group 198. Therefore \(z\in D\) and \(z\in E\), and the proved intersection gives \(z\in C\).

We have proved that the original \(C\) is reduced, has finitely many minimal primes, and is integrally closed in its total quotient ring. The source criterion 8274--8309, with its complete product comparisons in group 204, therefore makes it normal. Its mechanism retains the actual coordinate idempotents of \(Q(C)=\prod_i\operatorname{Frac}(C/\mathfrak p_i)\): each satisfies \(T^2-T=0\) and so belongs to \(C\). The maps \(c\mapsto(e_ic)_i\) and \((c_i)_i\mapsto\sum_i c_i\) are inverse, giving \(C=\prod_i e_iC\). Each factor embeds in its original field and is integrally closed there, as a monic equation in that factor can be placed in the corresponding coordinate of the product with all other lower coefficients zero. Thus the factors are normal domains. Returning through the proved colimit of finite normal \(B_E\), and then through the original prime localizations, proves the full stated theorem for arbitrary \(A,B\) with the original hypotheses.

\subsubsection{Both scalar actions in the separable algebraic comparison}\label{algebra-proof-070-both-scalar-actions-in-the-separable-algebraic-comparison}

Let \(k'/k\) be the original separable algebraic extension and \(A\) a \(k'\)-algebra. For a finite purely inseparable \(L/k\), the ring \(L'=k'\otimes_kL\) is a field: separability gives reducedness and the purely inseparable field extension gives a single prime, as proved above. The exact maps \[
 A\otimes_kL\xrightarrow{\sim}A\otimes_{k'}L',
 \quad a\otimes\ell\longmapsto a\otimes(1\otimes\ell),
\] \[
 a\otimes(r\otimes\ell)\longmapsto ar\otimes\ell
\] are inverse by tensor balancing. Thus geometric normality over \(k'\) gives normality of every such \(A\otimes_kL\), and the original four-condition criterion gives geometric normality over \(k\).

Conversely assume geometric normality over \(k\), and take the original field extension \(K/k'\). Set \(F_0=k'\otimes_kk'\), acting on \(K\otimes_kA\) by \[
 (r\otimes s)(u\otimes a)=ur\otimes sa.
\] Let \(F_0\to k'\) be the original multiplication map. Then \[
 (K\otimes_kA)\otimes_{F_0}k'
 \xrightarrow{\sim}K\otimes_{k'}A,
 \quad (u\otimes a)\otimes c\longmapsto uc\otimes a,
\] with inverse \(u\otimes_{k'}a\mapsto(u\otimes a)\otimes1\). For well-definedness, multiplying on the left by \(r\otimes s\) gives \(urc\otimes sa=urcs\otimes a\), equal to multiplication of \(c\) by \(rs\) on the right. Conversely \((r\otimes1)-(1\otimes r)\) has zero image under multiplication, so the inverse respects the \(k'\)-balanced relation. These calculations verify both composites and both original scalar actions.

The earlier diagonal-localization proof, group 275, identifies \(F_0\to k'\) as a localization. In the finite separable case, retain a primitive element \(\alpha\), its monic minimal polynomial \(P(T)=(T-\alpha)Q(T)\), and its nonzero value \(Q(\alpha)\). The corrected original polynomial map is \[
 k'[T]/(P)\longrightarrow k'\otimes_kk',\qquad
 \sum_i a_iT^i\longmapsto\sum_i a_i\otimes\alpha^i.
\] It is an isomorphism by the basis \(1,\alpha,\ldots,\alpha^{d-1}\), retaining the summation fixed earlier. Inverting \(Q\) forces \(T=\alpha\), since \((T-\alpha)Q=0\), and gives exactly evaluation into \(k'\); the inverse is the coefficient inclusion. The full denominator rule is \(f(T)/Q(T)^n\mapsto f(\alpha)/Q(\alpha)^n\). For infinite separable algebraic \(k'/k\), take its finite intermediate fields and the multiplicative closure of all these original finite-stage sets in \(F_0\). Each element maps to a nonzero scalar. Every kernel element lies in a finite tensor stage and is killed there by one of the finite-stage denominators. Thus the localization map is injective after evaluation and surjective from the coefficient inclusion, proving the same comparison without a finite-generation assumption on \(k'/k\).

Base-changing this localization through the displayed \(F_0\)-action proves that \(K\otimes_{k'}A\) is a localization of the normal ring \(K\otimes_kA\). Hence it is normal. The original arbitrary \(K/k'\) proves geometric normality over \(k'\).

The algebraicity hypothesis cannot simply be removed from this equivalence. In characteristic \(p>0\), take the perfect field \(k=\mathbf F_p\), the rational fields \(k'=k(t)\), \(A=k(s)\), and the actual inclusion \(t\mapsto s^p\). The extension \(k'/k\) is separable, and \(A/k\) is separable, so \(A\) is geometrically normal over \(k\). But \[
 A\otimes_{k'}A\cong A[T]/(T^p-s^p)
                    =A[T]/((T-s)^p).
\] The first isomorphism sends the second factor's original \(s\) to \(T\) and the first factor's \(s\) to the coefficient \(s\). The polynomial \(X^p-t\) is irreducible over \(k(t)\), for instance by the Eisenstein argument at \(t\) in \(k[t]\); hence its degree and the indicated tensor basis are \(p\). The nonzero element \([T]-s\), nonzero in that monic basis, has \(p\)-th power zero. All intervening binomial coefficients vanish in characteristic \(p\). Thus this scalar extension is not reduced and not normal, proving that \(A\) is not geometrically normal over \(k'\). This verifies the source hypothesis and the need for its actual diagonal-localization argument; it is not a proposed extension or a new source correction.

\subsubsection{Report accounting and propagation}\label{algebra-proof-070-report-accounting-and-propagation}

OCC-01027 receives the minimal designation copyedit at 47000. OCC-01028 and OCC-12447 retain both operations of 0736. OCC-01029 and OCC-12448 retain all three operations of 0737. OCC-01031 and OCC-12446 retain the mathematically necessary tensor-base correction 0735. OCC-01030 is optional: the coordinated sentence already supplies the finite-type assumption through an ordinary elliptical predicate; adding ``is'' is readable but does not repair a missing mathematical hypothesis or an unintelligible sentence. It is not integrated as another required fix.

The source's omitted receiving comparisons have been proved with their actual rings, scalar actions, nonzero conditions, denominators, nilpotents and prime maps. The finite purely inseparable diagram receives group 282; reducedness and separability receive groups 268 and 273; the infinite diagonal localization receives group 275; the tensor-normality proof receives groups 198, 204 and 206, finite-type Nagata stability, smooth normal ascent and faithful normal descent. All full proofs remain editorial. No novelty or whole-chapter completion claim is made. Review continues at 47092.

\hypertarget{algebra_georeg_cm_notes_20260929}{}
\subsection{Algebra: geometric regularity and Cohen--Macaulay field extensions}\label{algebra-proof-071-algebra-geometric-regularity-and-cohenmacaulay-field-extensions}

This editorial verification covers the original Stacks Project \texttt{algebra.tex} 47092--47318 at authority \texttt{a04446e57ec1fbc252a871afcec7752fb2807b14}, SHA-256 \texttt{FA8BB92E58A4F78A2BD01B3B6A4A87DE0A0D279F5DD90641B574DD5FBFFFA4F3}. The \href{https://github.com/stacks/stacks-project/blob/a04446e57ec1fbc252a871afcec7752fb2807b14/algebra.tex\#L47092}{original source} remains identifiable. The two reports retain MC-STK-ERR-0738, changing only the left-hand field in 47214. No additional source operation is proposed. All arguments here are editorial.

Original source dependencies read in this review include 6165--6181, 27488--27503, 39289--39335, 46386--46400, 46609--46626 and 46714--46732. The completed field-formal-smoothness, geometric-normality, ascent, descent and dimension-map notes supply their already proved comparisons. Three new indexed routing hits were returned; none of their contents was read. This is not a new literature-reading or novelty claim.

\subsubsection{Noetherianity and the original regularity tests}\label{algebra-proof-071-noetherianity-and-the-original-regularity-tests}

The source uses a Noetherian regular ring, with every prime localization a regular local ring. It tests geometric regularity against finitely generated field extensions, or equivalently finite purely inseparable ones. These quantifiers remain unchanged. The zero ring is Noetherian and has no prime localizations; all scalar extensions of it are zero and satisfy the same conventions.

Let \(R\) be the original Noetherian \(k\)-algebra and let \(K/k\) be finitely generated as a field. Choose field generators \(u_1,\ldots,u_n\), let \(D=k[u_1,\ldots,u_n]\subset K\), and put \(U=D\setminus\{0\}\). The map \[
 (U^{-1}D)\otimes_kR\longrightarrow(U\otimes1)^{-1}(D\otimes_kR),
 \quad(d/u)\otimes r\longmapsto(d\otimes r)/(u\otimes1)
\] has inverse \((\sum_jd_j\otimes r_j)/(u\otimes1)\mapsto\sum_j(d_j/u)\otimes r_j\). The fraction equality witnesses and tensor relations establish both composites. The ring \(D\otimes_kR\) is a quotient of the original finite polynomial ring \(R[X_1,\ldots,X_n]\), hence Noetherian. A localized ideal is generated by the images of a finite generating set of its contraction, so its localization is Noetherian. Since \(U^{-1}D=K\), this proves every Noetherianity assertion needed for the finitely generated scalar extensions.

For the two regularity tests, one implication is inclusion of the quantified classes. Conversely let the finite purely inseparable test hold, and fix the source's finitely generated \(K/k\). The already proved make-separable diagram has \(k'/k\) and \(K'/K\) finite purely inseparable, with \(K'/k'\) separable. The geometric-normality note retains its actual multiplication map \(k'\otimes_kK\to k'K\), nilpotent kernel and original field embeddings. We do not replace this possibly nonreduced tensor by a field without that quotient map.

Let \(B\) be the original smooth \(k'\)-domain with fraction field \(K'\), supplied by the completed generic smooth-model proof. The ring \(k'\otimes_kA\) is regular by the assumed finite test. Smooth base change gives \[
 k'\otimes_kA\longrightarrow B\otimes_{k'}(k'\otimes_kA),
\] so regular smooth ascent, proved in the ascent note, makes the target regular. The exact comparison is \[
 b\otimes(\lambda\otimes a)\longmapsto b\lambda\otimes a,
 \qquad b\otimes a\longmapsto b\otimes(1\otimes a).
\] It identifies the target with \(B\otimes_kA\). Localizing at every original nonzero \(b\in B\) gives \(K'\otimes_kA\); the inverse fraction maps are those displayed above. Each prime localization is an original regular local ring, so this localization remains regular and Noetherian. Finally \[
 K\otimes_kA\longrightarrow K'\otimes_kA
\] is faithfully flat: for an original \(K\otimes_kA\)-module \(M\), its base change identifies with \(K'\otimes_KM\) by \((u\otimes a)\otimes m\mapsto u\otimes am\). The inverse sends \(u\otimes m\) to \((u\otimes1)\otimes m\). A nonempty field basis proves exactness and reflection of zero. The complete regular-descent argument in the descent note now proves the original \(K\otimes_kA\) regular. This proves the source equivalence without an unrestricted colimit assertion about regular rings.

If \(A\) is geometrically regular and \(K/k\) is finitely generated, then \(K\otimes_kA\) is geometrically regular over \(K\): a finitely generated \(L/K\) is finitely generated over \(k\), and \[
 L\otimes_K(K\otimes_kA)\cong L\otimes_kA,
 \quad \ell\otimes(u\otimes a)\longleftrightarrow\ell u\otimes a
\] has inverse \(\ell\otimes a\mapsto\ell\otimes(1\otimes a)\). Noetherianity of the base-changed algebra was proved above. The later exponent-one criterion mentioned at 47160--47163 is only a forward reference to More on Algebra, not a new theorem asserted or re-proved in this intake.

For the source's faithfully flat map \(A\to B\) with geometrically regular \(B\), first descend Noetherianity. If \(I\subset A\), finitely many original elements of \(I\) generate \(IB\); let their ideal be \(I_0\). Flatness identifies \((I/I_0)\otimes_AB\) with \(IB/I_0B=0\). Faithfulness gives \(I=I_0\). Thus every ideal of \(A\) is finite. After any finite purely inseparable \(L/k\), the map \(A\otimes_kL\to B\otimes_kL\) is faithfully flat by the same base-change module comparison, and the target is regular. Regular descent proves the source test. For a smooth \(A\to B\) with geometrically regular \(A\), finite presentation makes \(B\) Noetherian. Every finitely generated scalar extension of that smooth map is smooth, and regular smooth ascent proves the source assertion. These arguments also cover the zero algebras whenever allowed by the original maps.

\subsubsection{Finite purely inseparable descent to an actual subfield}\label{algebra-proof-071-finite-purely-inseparable-descent-to-an-actual-subfield}

Retain the original directed union \(k=\bigcup_i k_i\) and the finite purely inseparable extension \(k'/k\). The construction below proves the field and tensor identity used at 47213--47215. It does not assume that an arbitrary descended finite algebra is a field.

In characteristic \(p>0\), choose a \(k\)-basis \(e_1=1,e_2,\ldots,e_d\) of the original \(k'\). There is an exponent \(q=p^N\) such that every \(e_j^q\) belongs to \(k\). Retain all multiplication constants \[
 e_re_s=\sum_{j=1}^d c_{rsj}e_j,\qquad c_{rsj}\in k,
\] and all elements \(a_j=e_j^q\in k\). Directedness puts these finitely many original constants and elements in one \(k_i\). Define the actual subspace of \(k'\) \[
 k_i'=\left\{\sum_{j=1}^d b_je_j:b_j\in k_i\right\}.
\] It contains \(1\), is closed under addition and multiplication by the retained constants, and is a domain as a subring of the original field. Its displayed basis is linearly independent over \(k_i\), since it is independent over \(k\). For any nonzero \(x\in k_i'\), multiplication by \(x\) is an injective endomorphism of this finite-dimensional \(k_i\)-space, hence surjective. In particular \(xy=1\) for some \(y\in k_i'\). Thus \(k_i'\) is a field of degree \(d\) over \(k_i\).

For every original element \(x=\sum_jb_je_j\) of this field, Frobenius gives the full formula \[
 x^q=\sum_{j=1}^d b_j^q e_j^q=\sum_{j=1}^d b_j^q a_j\in k_i.
\] The mixed binomial coefficients vanish in characteristic \(p\). This proves that \(k_i'/k_i\) is finite purely inseparable, with the actual bound \(q\); none of its coefficients has been discarded.

The multiplication comparison \[
 \mu:k\otimes_{k_i}k_i'\xrightarrow{\sim}k',\qquad
 \lambda\otimes x\longmapsto\lambda x
\] sends the \(k\)-basis \(1\otimes e_j\) to the original basis \(e_j\). It is therefore bijective and multiplicative, and its inverse is explicitly \[
 \sum_j\lambda_je_j\longmapsto\sum_j\lambda_j\otimes e_j.
\] The retained multiplication constants verify multiplication in this inverse formula as well. In characteristic zero, a finite purely inseparable extension is \(k'=k\); choose any index, set \(k_i'=k_i\), and use the identity multiplication comparison. The degree-one case in positive characteristic is included too.

Since \(A\) is an algebra over the original \(k\), the needed algebra comparison is \[
 \Phi:A\otimes_{k_i}k_i'\xrightarrow{\sim}A\otimes_kk',
 \qquad a\otimes x\longmapsto a\otimes x.
\] For \(y=\sum_j\lambda_je_j\in k'\), its inverse sends \[
 a\otimes y\longmapsto\sum_j a\lambda_j\otimes e_j.
\] The unique basis expansion proves additivity; the \(k\)-balanced relation moves any coefficient from the second tensor factor to the first in precisely this formula. The multiplication constants prove multiplicativity and both composites fix each elementary tensor. This is exactly the source's scalar-extension identity with the corrected left-hand field \(k'\).

Noetherianity of the unchanged \(A\) follows from its geometric regularity over any one of the nonempty directed family of fields. Its geometric regularity over the selected \(k_i\) makes \(A\otimes_{k_i}k_i'\) regular. Transporting along the proved \(\Phi\) verifies every finite purely inseparable test over \(k\), and hence geometric regularity over \(k\).

The original printed \(k_i=k\otimes_{k_i}k_i'\) is not this statement. For example, take the constant directed system \(k_i=k=\mathbf F_p(t)\) and \(k'=k(t^{1/p})\). The right side of the corrected comparison has degree \(p\) over \(k\), whereas \(k_i=k\) has degree one; irreducibility follows from Eisenstein at \(t\). The two reports therefore identify one mathematical correction, already present as unit 0738, not two operations.

\subsubsection{Separable algebraic change of the ground field}\label{algebra-proof-071-separable-algebraic-change-of-the-ground-field}

Let \(k'/k\) be the original separable algebraic extension and \(A\) a \(k'\)-algebra. Noetherianity is a property of the same ring \(A\) in either statement. If \(L/k\) is finite purely inseparable, then \(L'=k'\otimes_kL\) is a field by the reducedness and Frobenius-unit proof retained in the geometric-normality note. It is finite of degree \([L:k]\) over \(k'\) and purely inseparable: for a common exponent \(q\), every \(\sum_jb_j\otimes\ell_j\) has \(q\)-th power \(\sum_jb_j^q\ell_j^q\) in the embedded \(k'\). The exact associativity isomorphism \[
 A\otimes_{k'}(k'\otimes_kL)\longrightarrow A\otimes_kL,
 \quad a\otimes(b\otimes\ell)\longmapsto ab\otimes\ell
\] has inverse \(a\otimes\ell\mapsto a\otimes(1\otimes\ell)\). Hence geometric regularity over \(k'\) implies the original finite test over \(k\).

For the converse, first suppose \(k'/k\) finite separable. The preceding finitely generated base-change argument makes \(A\otimes_kk'\) geometrically regular over its second \(k'\)-factor. The multiplication map \(A\otimes_kk'\to A\) is an étale localization. Here are its exact denominators. Choose a primitive element \(\alpha\) of \(k'/k\) and its original monic minimal polynomial \(P(T)\). In \(k'[T]\), write \[
 P(T)=(T-\alpha)Q(T),\qquad Q(\alpha)=P'(\alpha)\ne0.
\] The map \(A[T]/(P)\to A\otimes_kk'\), \(\sum_j a_jT^j\mapsto\sum_j a_j\otimes\alpha^j\), is an isomorphism by the original power basis. On the left, \(\alpha\) in a coefficient means its given image in \(A\). Inverting \(Q\) forces \(T=\alpha\), so evaluation gives \[
 (A[T]/(P))_Q\xrightarrow{\sim}A,
 \quad f(T)/Q(T)^n\longmapsto f(\alpha)/Q(\alpha)^n,
\] with inverse the coefficient inclusion. The element \(Q(\alpha)\) is already a unit in \(A\), including when \(A=0\), so all denominators are defined. This proves that the multiplication map is localization at the single indicated element, hence is étale by 39293 and smooth. Evaluation sends \(1\otimes r\) to the original \(r\in A\), so it preserves exactly the second \(k'\)-algebra structure. Geometric regular smooth ascent proves the desired assertion over \(k'\).

For arbitrary separable algebraic \(k'/k\), its finite intermediate extensions form a directed union. The argument just given proves that the same \(A\) is geometrically regular over each of them. The finite purely inseparable descent construction in the preceding section applies to this union and gives geometric regularity over the original \(k'\). It is the finite-test descent, not a claim that every infinite scalar extension of a Noetherian ring stays Noetherian.

\subsubsection{The original field tensor and its Cohen--Macaulay localizations}\label{algebra-proof-071-the-original-field-tensor-and-its-cohenmacaulay-localizations}

Let \(K/k,L/k\) be the original field extensions, one finitely generated. The tensor flip \(x\otimes y\mapsto y\otimes x\) is its own inverse, so suppose the source's \(K/k\) is finitely generated. Choose its original transcendence basis \(t_1,\ldots,t_r\); then \(K\) is finite of degree \(d\ge1\) over \(E=k(t_1,\ldots,t_r)\). Set \[
 U=k[t_1,\ldots,t_r]\setminus\{0\},\qquad
 R=E\otimes_kL\cong U^{-1}L[t_1,\ldots,t_r],
 \quad S=K\otimes_kL.
\] The displayed localization inverts the original nonzero polynomials with coefficients in \(k\), not every nonzero polynomial with coefficients in \(L\). The polynomial algebra over the field \(L\) is Cohen--Macaulay by the preceding source polynomial theorem, and so is each of its localizations. Thus \(R\) is Noetherian and Cohen--Macaulay, including the case \(r=0\), when \(R=L\).

The map \[
 K\otimes_E(E\otimes_kL)\xrightarrow{\sim}K\otimes_kL,
 \quad x\otimes(e\otimes\ell)\longmapsto xe\otimes\ell
\] has inverse \(x\otimes\ell\mapsto x\otimes(1\otimes\ell)\). An original \(E\)-basis of \(K\) therefore makes \(S\) finite free of rank \(d\) over \(R\). It is Noetherian, and this finite map is integral. For completeness, multiplying an element of \(S\) on its finite free \(R\)-module gives a matrix; its monic characteristic polynomial annihilates that matrix, and evaluation on the unit gives the integral equation for the original element.

Take any original prime \(\mathfrak q\subset S\), with contraction \(\mathfrak p\subset R\). Then \(R_{\mathfrak p}\to S_{\mathfrak q}\) is a flat local map of Noetherian local rings. We do not assume that localization at this one prime preserves finiteness as a module. Instead, any strict chain of primes in \(S\) below \(\mathfrak q\) contracts to a strict chain below \(\mathfrak p\), by incomparability for the original integral extension. Thus \[
 \dim S_{\mathfrak q}\le\dim R_{\mathfrak p}.
\] Let \(n=\dim R_{\mathfrak p}\). A parameter regular sequence of length \(n\) in the Cohen--Macaulay \(R_{\mathfrak p}\) remains regular on \(S_{\mathfrak q}\): flatness preserves every multiplication injection and quotient, and the local map is faithfully flat, so the final nonzero quotient remains nonzero. Consequently \[
 n\le\operatorname{depth}S_{\mathfrak q}
   \le\dim S_{\mathfrak q}\le n.
\] All inequalities are equalities. The target local ring is Cohen--Macaulay; since the prime was arbitrary, this proves the original tensor theorem. It uses the second alternative of 27492--27493 when needed, without an unjustified local finiteness assertion or any separability hypothesis on \(K/E\).

\subsubsection{The exact residue-fibre comparison and the earlier consequence}\label{algebra-proof-071-the-exact-residue-fibre-comparison-and-the-earlier-consequence}

Retain the final source notation \(S_K=K\otimes_kS\), \(\mathfrak q\subset S\) and \(\mathfrak q_K\subset S_K\) lying over it. Put \(A=S_{\mathfrak q}\), \(B=(S_K)_{\mathfrak q_K}\), and \(F=K\otimes_k\kappa(\mathfrak q)\). These are the original rings; \(A\to B\) is flat local, and \(A,B\) are Noetherian by the first section. Write \[
 D=K\otimes_kS_{\mathfrak q},\quad
 P=\mathfrak q_KD,\quad
 J=K\otimes_k\mathfrak q S_{\mathfrak q}.
\] The ideal \(P\) exists because every denominator from \(S\setminus\mathfrak q\) avoids \(\mathfrak q_K\); it is prime and contains \(J\). The exact quotient \(D/J\cong F\) sends \(u\otimes(a/v)\) to \(u\otimes\overline{a/v}\). Let \(\mathfrak q'\) be the image of \(P\) in \(F\). Quotient and localization then give \[
 B/\mathfrak q B\cong(D/J)_{P/J}=F_{\mathfrak q'}.
\] The forward map takes the residue class of \(z/w\in D_P\) to \(\bar z/\bar w\). Its inverse chooses lifts of numerator and denominator in \(D\); two choices differ by elements of \(J\), and the same fraction equality witnesses prove independence after quotienting. A denominator outside \(P/J\) has every lift outside \(P\), so both formulas are defined. Composing with the tensor flip identifies this with exactly the source's \((\kappa(\mathfrak q)\otimes_kK)_{\mathfrak q'}\), relabeling the prime by that proved isomorphism.

The field tensor just proved is Cohen--Macaulay; thus its localization \(B/\mathfrak q B\) is Cohen--Macaulay. The earlier complete flat local formulas, retained in the ascent note, are \[
 \dim B=\dim A+\dim(B/\mathfrak qB),\qquad
 \operatorname{depth}B=\operatorname{depth}A+
                      \operatorname{depth}(B/\mathfrak qB).
\] These finite local dimensions are not global bounds on the original rings. Subtracting the two full formulas gives \[
 \dim B-\operatorname{depth}B=
 \dim A-\operatorname{depth}A+
 \dim(B/\mathfrak qB)-\operatorname{depth}(B/\mathfrak qB).
\] The final summand is zero by the proved fibre theorem. Hence \(A\) is Cohen--Macaulay exactly when \(B\) is, proving the original final lemma at each specified pair of primes.

This review also resolves the scope of earlier group 650. Its proof at 27488--27503 extended the particular zero-dimensional-fibre assertion to a flat local map with any Cohen--Macaulay closed fibre. But the \href{https://github.com/stacks/stacks-project/blob/a04446e57ec1fbc252a871afcec7752fb2807b14/algebra.tex\#L46386}{later original lemma at 46386--46400} explicitly states the full equivalence, and the earlier dimension formula supplies the same dimension equality. Thus that broader implication is already in the original work. Preserve group 650 and its full regular-sequence proof for provenance, while refining its current ledger classification to a consequence already stated in later source content. It is a valid receiving proof, not an extension beyond the entire source. This comparison propagates to the final field-extension application just proved.

\subsubsection{Report accounting and continuation}\label{algebra-proof-071-report-accounting-and-continuation}

OCC-01032 and OCC-12449 both retain the single existing operation MC-STK-ERR-0738-OP1 at 47214. The basis construction proves the corrected field and both scalar-extension maps with all original constants. No new source operation, conceptual-extension group, external reading, build or publication is claimed. Full proofs and the refinement of group 650 stay editorial. Continue source review at 47319; the remaining thirteen reports have been read, but their proof intervals have not yet been adjudicated.

\hypertarget{algebra_fp_flat_notes_20260929}{}
\subsection{Algebra: finite-presentation models for flatness}\label{algebra-proof-072-algebra-finite-presentation-models-for-flatness}

Source: Stacks Project authors, \texttt{algebra.tex} 47319--47520, original authority \texttt{a04446e57ec1fbc252a871afcec7752fb2807b14}, SHA-256 \texttt{FA8BB92E58A4F78A2BD01B3B6A4A87DE0A0D279F5DD90641B574DD5FBFFFA4F3}. The \href{https://github.com/stacks/stacks-project/blob/a04446e57ec1fbc252a871afcec7752fb2807b14/algebra.tex\#L47319}{original source} remains identifiable. Seven physical reports retain four existing units 0739--0742. Two minimal unreported copyedits supply ``by'' in the inverse-image declaration and ``is'' in the next lemma's hypothesis. All complete proof details remain editorial.

Dependencies read in their original TeX: 9163--9181, 9707--9737, 32426--32481, 32592--32628, 32722--32762 and 33406--33506. The completed colimit-category, Noetherian-approximation, more-flatness and open-flatness notes retain the earlier full arguments, including their source corrections. Four indexed routing hits were returned: three are unread, and the Bhatt--Scholze Witt-vector Grassmannian source has only its previously recorded reading of lines 251--284. No additional contents were read in this query. No new mathematical extension or novelty is claimed here.

\subsubsection{Original presentations and their localized maps}\label{algebra-proof-072-original-presentations-and-their-localized-maps}

Retain the given finite-presentation ring map \(R\to S\), the finitely presented \(S\)-module \(M\), and its assumed flatness over \(R\). Fix the original presentations \[
 S=R[X_1,\ldots,X_n]/(g_1,\ldots,g_m),\qquad
 S^{\oplus a}\xrightarrow{H}S^{\oplus b}\longrightarrow M\longrightarrow0.
\] Each original polynomial coefficient and each entry of the full matrix \(H\) is retained. Choose polynomial representatives of the entries. The finite-type subrings \(R_\lambda\subset R\) containing their coefficients form a directed system; put \[
 S_\lambda=R_\lambda[X_1,\ldots,X_n]/(g_{1,\lambda},\ldots,g_{m,\lambda}),
 \qquad M_\lambda=\operatorname{coker}H_\lambda.
\] Every transition is the original coefficient map, fixing the variables and matrix positions. Right exactness gives the precise comparisons \[
 R_\mu\otimes_{R_\lambda}S_\lambda\cong S_\mu,
 \qquad S_\mu\otimes_{S_\lambda}M_\lambda\cong M_\mu.
\] The first sends \(r\otimes f(X)\) to \(rf_\mu(X)\). Its inverse sends a polynomial coefficient \(r\in R_\mu\) to \(r\otimes1\) and each original variable to \(1\otimes X_i\); the relations map to exactly the retained relations. The second sends \(s\otimes[v]\) to \([s v_\mu]\); its inverse sends each standard target generator to \(1\otimes[e_j]\). All matrix relations are preserved by right exactness. The same formulas identify the colimits with the original \(S,M\). An element or a defining equality involves finitely many coefficients, hence occurs at a finite stage. The base rings are Noetherian because they are finite type over \(\mathbf Z\); so are the \(S_\lambda\), and each finite \(M_\lambda\) is finitely presented. Empty variable, relation and matrix lists are allowed.

Take an original prime \(\mathfrak q\subset S\), its contraction \(\mathfrak p\subset R\), and the contractions \(\mathfrak q_\lambda,\mathfrak p_\lambda\). Every denominator outside a contracted prime stays outside the later contraction. Thus \[
 (R_\lambda)_{\mathfrak p_\lambda}\longrightarrow
 (S_\lambda)_{\mathfrak q_\lambda}
\] and their transition maps are local. Localization commutes with these colimits: each target fraction lifts with its original denominator, and a zero fraction has an annihilating denominator and equality that appear together at a later stage. Base change of \((S_\lambda)_{\mathfrak q_\lambda}\) first inverts the old denominator set; further localization at the contraction of \(\mathfrak q\) gives \((S_\mu)_{\mathfrak q_\mu}\). The map on fractions is \((s/u)\otimes(r/v)\mapsto sr/(uv)\), with the structural images understood; its further prime localization has the inverse induced by the two original factor maps. The module comparison follows by applying the same localizations to its actual presentation matrix. These are exactly the six conditions of 32426--32453, including localization rather than an unjustified unrestricted tensor isomorphism.

The complete earlier eventual-flatness proof now applies to \(M_{\mathfrak q}\), which is \(R_{\mathfrak p}\)-flat. Explicitly, at a local stage \(\lambda\) take the finite module \[
 T_\lambda=\operatorname{Tor}^{(R_\lambda)_{\mathfrak p_\lambda}}_1
 \big((M_\lambda)_{\mathfrak q_\lambda},\kappa(\mathfrak p_\lambda)\big)
 =\ker\big(\mathfrak p_\lambda(R_\lambda)_{\mathfrak p_\lambda}
       \otimes (M_\lambda)_{\mathfrak q_\lambda}
       \longrightarrow(M_\lambda)_{\mathfrak q_\lambda}\big).
\] It has finitely many original generators since the local stage rings are Noetherian and the module is finite. Their images in the tensors with the extended ideal vanish in the original flat limit module. The filtered-colimit equality witnesses therefore kill them all at one later stage. The change-of-base local flatness criterion already proved in the more-flatness note, with precisely these original ideals and localized modules, gives flatness at that stage. Later stages are localizations of its base changes, so remain flat. Thus each original \(\mathfrak q\) supplies a stage \(\lambda_{\mathfrak q}\) after which \((M_\lambda)_{\mathfrak q_\lambda}\) is \(R_\lambda\)-flat. A module over \((R_\lambda)_{\mathfrak p_\lambda}\) is flat over that ring exactly when it is flat over \(R_\lambda\): restriction uses localization flatness; the reverse follows by tensoring injections of localized modules, whose tensor products identify with the original ones.

\subsubsection{The entire flat locus and its direction under base change}\label{algebra-proof-072-the-entire-flat-locus-and-its-direction-under-base-change}

Set, exactly as retained correction 0740 states, \[
 U_\lambda=\{\mathfrak r\in\operatorname{Spec}(S_\lambda):
                    (M_\lambda)_{\mathfrak r}\text{ is flat over }R_\lambda\}.
\] This entire set is open by the source's openness theorem. The completed open-flatness note supplies its proof with the \(n=0\) and \(n=1\) resolution cases and the correct fibre global-dimension bound. Thus those earlier corrections are received here; we do not reuse the source's erroneous unrestricted syzygy formula.

For \(\mu\ge\lambda\), the inverse image of \(U_\lambda\) in \(\operatorname{Spec}(S_\mu)\) lies in \(U_\mu\). To prove this, let \(\mathfrak r_\mu\) contract to \(\mathfrak r_\lambda\in U_\lambda\). First base change the flat \(R_\lambda\)-module \((M_\lambda)_{\mathfrak r_\lambda}\) to \(R_\mu\). This is \(R_\mu\)-flat. Then localize its acting algebra at the prime induced by \(\mathfrak r_\mu\). The original presentation comparisons identify the result with \((M_\mu)_{\mathfrak r_\mu}\). Localization in that acting algebra preserves base-module flatness: it is a filtered colimit of copies of a flat base module with multiplication transition maps, and tensoring an injection commutes with that colimit. This proves the inclusion without assuming the transition \(R_\lambda\to R_\mu\) flat.

Let \(V_\lambda\subset\operatorname{Spec}(S)\) be the inverse image of \(U_\lambda\), with its original map. The inclusion just proved gives \(V_\lambda\subset V_\mu\). Eventual local flatness covers every original prime by some \(V_\lambda\). Quasi-compactness gives a finite subcover; choose a common upper index of those finitely many stages, not an unspecified maximum of the whole directed set. Monotonicity makes its \(V_{\lambda_0}\) the entire spectrum. If \(S=0\), the spectrum and this cover are empty; choose any stage of the nonempty directed set, and the equality is still valid.

An arbitrary open subset on which the module happens to be flat would not establish this cover or inclusion. For instance \(R=S=k\), \(M=k\) is flat at its sole prime, but the empty open subset also satisfies the old one-way property and misses that prime. Thus the retained change to the entire flat locus is mathematically necessary. The additional ``by'' in 47397 only makes the original inverse-image declaration grammatical; it does not change \(V_\lambda\).

\subsubsection{Every coefficient of the finite unit relation}\label{algebra-proof-072-every-coefficient-of-the-finite-unit-relation}

Write the complement of \(U_{\lambda_0}\) as \(V(I)\), for its original ideal \(I\subset S_{\lambda_0}\). Pulling back gives \(V(IS)=\varnothing\), so \(IS=S\): a proper ideal in a nonzero unital ring lies in a maximal ideal, while the zero ring already has the unit ideal as its only ideal. Consequently there are original elements \[
 f_1,\ldots,f_r\in I,\quad s_1,\ldots,s_r\in S,
 \qquad\sum_{i=1}^r f_i s_i=1.
\] The retained correction 0741 changes the inconsistent terminal index \(n\) to \(r\). Keep all \(f_i\), all \(s_i\), their products and coefficients. Lift the finitely many \(s_i\) to one later stage \(\mu\). At that stage the error \(1-\sum_i f_i s_{i,\mu}\) has image zero in the colimit. Enlarge once more to \(\nu\) where that same error is zero. Lifting the coefficients alone would not prove this equality; both finite steps are needed.

At this \(\nu\), \(IS_\nu=S_\nu\). The inverse image of \(U_{\lambda_0}=D(I)\) is therefore all of \(\operatorname{Spec}(S_\nu)\). The base-change inclusion puts it in \(U_\nu\), proving \(U_\nu=\operatorname{Spec}(S_\nu)\). This is the precise meaning of increasing the original stage; the later flat locus need not have the same defining ideal as the earlier one.

Flatness at all primes of the acting algebra implies that \(M_\nu\) is \(R_\nu\)-flat. Indeed, for any injection \(N_1\to N_2\) of \(R_\nu\)-modules, the kernel of its tensor map with \(M_\nu\) is an \(S_\nu\)-module. Its localizations at all primes are zero by the pointwise flatness. A nonzero module element has proper annihilator contained in a maximal ideal, and stays nonzero at a prime containing that annihilator; thus a module with all prime localizations zero is zero. The tensor map is injective. This proves part (1) with \(R_0=R_\nu,S_0=S_\nu,M_0=M_\nu\) and the original presentation comparisons. If the unit relation is empty because \(S=0\), the error is \(1\); a later stage kills it, and the same argument gives the zero stage algebra and module, which is flat.

\subsubsection{Arbitrary base systems and the two eventual inverse maps}\label{algebra-proof-072-arbitrary-base-systems-and-the-two-eventual-inverse-maps}

For part (2), let the original \(R\) be a directed colimit of the given \(R_\lambda\). The finite-type \(\mathbf Z\)-algebra \(R_0\) from part (1) is finitely presented, since the kernel of its polynomial presentation is a finitely generated ideal in a Noetherian polynomial ring. Lift the finite images of its generators to a common stage, then annihilate the finite defining relations at a later one. This gives a map \(R_0\to R_\lambda\) whose composite to \(R\) is the original one.

Define exactly \[
 S_\lambda=R_\lambda\otimes_{R_0}S_0,
 \qquad M_\lambda=S_\lambda\otimes_{S_0}M_0.
\] The algebra and module are finitely presented because their original finite relations and matrix base change to the indicated rings. The map \[
 (R_\lambda\otimes_{R_0}S_0)\otimes_{S_0}M_0
    \xrightarrow{\sim}R_\lambda\otimes_{R_0}M_0,
 \quad(r\otimes s)\otimes m\longleftrightarrow r\otimes sm
\] has inverse \(r\otimes m\mapsto(r\otimes1)\otimes m\). Thus \(M_\lambda\) is \(R_\lambda\)-flat, by base change of \(M_0\). The same associativity maps after tensoring with \(R\) recover exactly \(S,M\), proving part (2).

For part (3), retain the original prescribed system \((R_\lambda,S_\lambda,M_\lambda)\). First descend \(R_0\to R\) as above to a stage \(\lambda_0\). The finite presentation of \(S_0\) over \(R_0\) lifts its map to \(S=\operatorname{colim}S_\lambda\): lift its finite generator images and then all defining relation errors. Thus one obtains \(S_0\to S_{\lambda_1}\) over \(R_0\). For later indices the original map is \[
 \Psi_\lambda:R_\lambda\otimes_{R_0}S_0\longrightarrow S_\lambda,
 \qquad r\otimes s\longmapsto r\,s_\lambda.
\] Its colimit is the given presentation isomorphism. Both sides are base changes of finitely presented \(R_{\lambda_1}\)-algebras. Use their actual finite presentations to descend the inverse of \(\Psi_R\): lift each generator image, annihilate its finite relations, and then annihilate on each of the two finite generator lists the errors in both inverse composites. At a common later stage these maps are inverse algebra maps. This is the full finite-presentation argument behind the cited categorical lemma; it does not infer injectivity from a disappearing arbitrary kernel.

The finite \(S_0\)-module \(M_0\) is finitely presented since \(S_0\) is Noetherian. The same presentation argument lifts the map \(M_0\to M\) to an \(S_0\)-module map \(M_0\to M_{\lambda_2}\): lift the original generator images and kill each finite relation error. It induces \[
 \Theta_\lambda:S_\lambda\otimes_{S_0}M_0\longrightarrow M_\lambda,
 \qquad s\otimes m\longmapsto s\,m_\lambda.
\] Both sides, for \(\lambda\ge\lambda_2\), are base changes of finitely presented modules over \(S_{\lambda_2}\), and their limit map is the given isomorphism. Descend its inverse using the original presentation matrix; then kill the errors of both composites on the two finite generator lists. Thus \(\Theta_\lambda\) is an isomorphism for all sufficiently large indices. Combine this with the already established \(\Psi_\lambda\): the source of \(\Theta_\lambda\) is identified with \(R_\lambda\otimes_{R_0}M_0\), which is flat. This proves eventual flatness in the given system, including noninjective transitions.

\subsubsection{Faithful flatness, spectral images and a finite relation algebra}\label{algebra-proof-072-faithful-flatness-spectral-images-and-a-finite-relation-algebra}

Retain \(R\to A\to B\), with \(A\to B\) faithfully flat of finite presentation. The short added copula in 47472 makes this assumption grammatical and changes no property. Start with \(R=\mathbf Z\). Apply the preceding construction to \(A\to B\) and its rank-one \(B\)-module \(B\), using the presentation with one generator and no module relations. Every stage module is then its original stage algebra. The finite unit relation argument makes that stage algebra flat; an arbitrary finite module model alone would not suffice. This gives a finite-type \(\mathbf Z\)-algebra \(A_0\), a flat finitely presented \(A_0\)-algebra \(B_0\), and the actual isomorphism \[
 A\otimes_{A_0}B_0\xrightarrow{\sim}B.
\] The image \(W\) of \(\operatorname{Spec}(B_0)\to\operatorname{Spec}(A_0)\) is open by 9707--9737. That theorem applies exactly: finite presentation gives constructibility, flat going down makes the image stable under generalization, and this yields openness. Its standard-open argument retains the actual localizations. Write the complement as \(V(I_0)\).

Here is the required equality of images after any base change \(A_0\to D\). For a prime \(\mathfrak p\subset D\), let \(\mathfrak p_0\) be its contraction. The induced map \(\kappa(\mathfrak p_0)\to\kappa(\mathfrak p)\) is a field extension. The original fibre comparison is \[
 (D\otimes_{A_0}B_0)\otimes_D\kappa(\mathfrak p)
 \cong B_0\otimes_{A_0}\kappa(\mathfrak p)
 \cong (B_0\otimes_{A_0}\kappa(\mathfrak p_0))
            \otimes_{\kappa(\mathfrak p_0)}\kappa(\mathfrak p).
\] The first map sends \((d\otimes b)\otimes c\) to \(b\otimes dc\); the second sends \((b\otimes a)\otimes c\) to \(b\otimes ac\); their inverses insert the unit in the intermediate coefficient factor. Tensoring with a nonzero field extension reflects the zero vector space, by a field basis. Thus this fibre is nonzero exactly when \(B_0\otimes_{A_0}\kappa(\mathfrak p_0)\) is nonzero. A nonzero unital ring has a prime, so these conditions say precisely that \(\mathfrak p\) is in the base-changed image and that \(\mathfrak p_0\in W\). This proves the asserted equality of images, including zero fibres.

For \(D=A\), faithful flatness says the image is the whole spectrum. Therefore \(I_0A=A\). Keep a finite original relation \[
 \sum_{i=1}^r f_i r_i=1,\qquad f_i\in I_0,\quad r_i\in A.
\] The source's instruction to enlarge the model is made exact by the finite-presentation algebra \[
 A_1=A_0[T_1,\ldots,T_r]/\left(\sum_{i=1}^r f_iT_i-1\right),
 \qquad B_1=A_1\otimes_{A_0}B_0.
\] Send \(T_i\) to the actual \(r_i\) to define \(A_1\to A\); the relation verifies this map. No injection of \(A_0\) or \(A_1\) into \(A\) is assumed. The ring \(A_1\) is finite type and finitely presented over \(\mathbf Z\), and \(A_1\to B_1\) is flat of finite presentation by base change. In \(A_1\) the original extended ideal is the unit ideal, so the just-proved image equality makes \(\operatorname{Spec}(B_1)\to\operatorname{Spec}(A_1)\) surjective. The original flat/surjective-spectrum criterion then gives faithful flatness.

The receiving comparison is \[
 A\otimes_{A_1}(A_1\otimes_{A_0}B_0)
   \xrightarrow{\sim} A\otimes_{A_0}B_0\xrightarrow{\sim}B,
 \quad a\otimes(a_1\otimes b)\longmapsto aa_1\otimes b,
\] with inverse \(a\otimes b\mapsto a\otimes(1\otimes b)\). This proves exactly the original descent diagram after renaming \(A_1,B_1\) as the final model. If the finite relation is empty and \(A=0\), its relation algebra imposes \(-1=0\), giving the zero model; the resulting map and comparisons have the same required properties. Alternatively the nonempty relation with zero coefficient gives the identical conclusion. No nonzero-ring exclusion is introduced.

Finally retain arbitrary original \(R\). Let \(A_0'\to B_0'\) be the model over \(\mathbf Z\), with its maps to the original \(A,B\), and put exactly \[
 A_0=A_0'\otimes_{\mathbf Z}R,\qquad
 B_0=B_0'\otimes_{\mathbf Z}R.
\] The map \(A_0\to A\) sends \(a_0'\otimes r\) to the product of their original images. The map for \(B_0\) is defined the same way. Commutativity of \(R\to A\to B\) makes the original diagram commute. The isomorphism \[
 A_0\otimes_{A_0'}B_0'\cong B_0'\otimes_{\mathbf Z}R,
 \quad(a_0'\otimes r)\otimes b_0'\longmapsto a_0'b_0'\otimes r
\] has inverse \(b_0'\otimes r\mapsto(1\otimes r)\otimes b_0'\). Therefore \(A_0\to B_0\) is the actual base change of the faithfully flat finitely presented map, hence has those same properties. Also \(R\to A_0\) is finitely presented by the finite original \(\mathbf Z\)-presentation. Finally \[
 A\otimes_{A_0}B_0\xrightarrow{\sim}A\otimes_{A_0'}B_0',
 \quad a\otimes(b_0'\otimes r)\longmapsto ar\otimes b_0'
\] has inverse \(a\otimes b_0'\mapsto a\otimes(b_0'\otimes1)\). The balanced relations over \(A_0'\) and \(R\) verify both maps and composites. Composing with the original model isomorphism recovers \(B\), proving the requested general diagram.

\subsubsection{Dispositions and propagation}\label{algebra-proof-072-dispositions-and-propagation}

OCC-01033 and OCC-12450 retain 0739, the plural ``ring maps''. OCC-12451 retains 0740, which names the entire flat locus and supplies its copula. The adjacent unreported inverse-image declaration receives ``by'' at 47397. OCC-01034 and OCC-12452 retain 0741, the common terminal index \(r\). The unreported hypothesis at 47472 receives ``is''. OCC-01035 and OCC-12453 retain 0742, ``replaced by''. Thus seven reports retain four existing operations and add two source copyedits; expanded arguments are not substitutions into either translation branch.

This receives the original finite-presentation systems and matrices, eventual-flatness Tor argument, corrected openness theorem and its arbitrary-base-change inclusion, and the finite-stage inverse constructions from the earlier notes. No new conceptual-extension group is needed. Continue at 47521, with six remaining report payloads already read and their proof text still to be reviewed.

\hypertarget{algebra_fp_end_notes_20260929}{}
\subsection{Algebra: final colimit lemmas and a quasi-finite flat cover}\label{algebra-proof-073-algebra-final-colimit-lemmas-and-a-quasi-finite-flat-cover}

Source: Stacks Project authors, algebra.tex 47521--47885, original authority a04446e57ec1fbc252a871afcec7752fb2807b14, SHA-256 FA8BB92E58A4F78A2BD01B3B6A4A87DE0A0D279F5DD90641B574DD5FBFFFA4F3. The \href{https://github.com/stacks/stacks-project/blob/a04446e57ec1fbc252a871afcec7752fb2807b14/algebra.tex\#L47521}{original source}, including terminal apparatus, has been read. Six reports add seven missing copulas, retain two existing mathematical operations, and reject a proposed change to a valid redundant localization. This ends the source-reading denominator, not chapter admission, compilation or publication. Complete proofs remain separate editorial material.

Fresh original dependencies: 23491--23504, 30195--30230, 30361--30390, 30879--30903, 33622--33698, 33958--33975, 36638--36668, and 36726--36764. Earlier complete syntomic, smooth, étale, differential, flatness, Cohen--Macaulay, quasi-finite and colimit proofs are reused with their source versions and corrections retained. Of the two indexed hits, one is unread and the Bhatt--Lurie Absolute Prismatic Cohomology source has only its previously recorded reading of lines 11899--11920. No further contents were read, and no theorem from that bounded opening is adopted. No novelty claim is made.

\subsubsection{Original stage objects and the first four properties}\label{algebra-proof-073-original-stage-objects-and-the-first-four-properties}

Keep the directed set \(I\), its specified element \(0\), and the actual \(A_0\)-algebra map \(\varphi_0:B_0\to C_0\). The subsystem \(i\ge0\) is cofinal by directedness. Put \[
B_i=A_i\otimes_{A_0}B_0,\quad C_i=A_i\otimes_{A_0}C_0,\qquad
B=A\otimes_{A_0}B_0,\quad C=A\otimes_{A_0}C_0.
\] Every transition is the given coefficient map tensored with the identity. The comparison \(A_j\otimes_{A_i}B_i\cong B_j\) sends \(a\otimes(b\otimes x)\) to \(ab\otimes x\), with inverse \(a\otimes x\mapsto a\otimes(1\otimes x)\); the same formulas hold for \(C_i\). These are the original base-change systems with colimits \(B,C\). No transition is assumed injective or flat.

For finite descent, retain \(x_1,\ldots,x_m\) generating \(C_0\) over \(B_0\). Finiteness of \(C/B\) supplies a monic polynomial \(P_j\) for each original image of \(x_j\). Lift every lower coefficient to one \(B_i\), keeping leading coefficient one. Then kill the finitely many evaluation errors in a later \(C_i\). If \(d_j\) is the degree of the retained equation, repeated substitution expresses every monomial as a linear combination of \[
x_1^{e_1}\cdots x_m^{e_m},\qquad 0\le e_j<d_j.
\] Each substitution decreases the selected exponent and retains every lower coefficient and sign. This finite list therefore spans \(C_i\) over \(B_i\). The empty generator list gives the empty monomial and a surjective coefficient map. For a zero target ring, choose monic degree-one equations for the generators; the same proof applies.

For surjectivity, choose actual \(b_j\in B\) whose images are \(1\otimes x_j\). Lift them to one \(B_i\), then kill every error \(\varphi_i(b_{j,i})-1\otimes x_j\) at a common later stage. The image contains all target algebra generators and its coefficient ring, hence all of \(C_i\). Preimage lifting and the later equality are distinct necessary steps.

For unramifiedness, the original symbols \(dx_1,\ldots,dx_m\) generate each \(\Omega_{C_i/B_i}\). The derivative of every polynomial is its full partial-derivative sum times these symbols; coefficient derivatives vanish. The filtered-colimit comparison sends \(dc_i\) to the differential of its original image. Every element and every additive, Leibniz or coefficient relation involves finitely many elements and equalities, so it lifts to a common later stage. This proves both directions of \[
\operatorname{colim}_i\Omega_{C_i/B_i}\cong\Omega_{C/B}=0.
\] Kill the finite symbol list at one stage. The whole differential module vanishes there, and finite type persists from \(C_0/B_0\); the source's finite-type unramified criterion applies. No finite-presentation hypothesis is added.

For isomorphism descent, first use the proved surjective case. Write \(C_i=B_i/J\). The kernel \(J\) is finite by finite presentation. Explicitly choose \(C_i=B_i[Y]/(h_1,\ldots,h_s)\), and lifts \(b_j\in B_i\) of the images of \(Y_j\). Since \(Y_j-b_j\) belongs to the relation ideal, substitution gives \[
J=(h_1(b),\ldots,h_s(b)).
\] Indeed an element of \(B_i\) vanishing in \(C_i\) is a finite polynomial combination of the \(h_s\); evaluating that equality at \(b\) proves the assertion, with all coefficients retained. Each of these generators maps to zero in \(B\), since \(B\to C\) is injective. At a common later \(B_j\) all vanish. Right exactness and the original base-change comparisons give \(C_j=B_j/JB_j=B_j\), via exactly \(\varphi_j\). No tensor left-exactness or unchanged-kernel assumption is used.

\subsubsection{Conormal colimits and both inverse identities}\label{algebra-proof-073-conormal-colimits-and-both-inverse-identities}

Retain every original polynomial: \[
P_i=B_i[x_1,\ldots,x_n],\quad I_i=(f_{1,i},\ldots,f_{m,i}),\quad
C_i=P_i/I_i,\quad E_i=I_i/I_i^2,\quad F_i=\bigoplus_{k=1}^n C_i\,dx_k.
\] The \(f_{j,i}\) are images of \(f_{j,0}\), with all monomials, coefficients and signs. Let \(P,I,E,F\) be their actual limit objects. Every element of \(I\) is a finite polynomial combination of the retained generators and thus comes from some \(I_i\); a represented element becoming zero does so at a later \(P_i\). Therefore \(\operatorname{colim}I_i=I\). Every element of \(I^2\) is a finite sum of products of such expressions, proving similarly \(\operatorname{colim}I_i^2=I^2\). A class of \(I_i/I_i^2\) becomes zero precisely when its representative becomes such a finite sum of products: lift that sum and the equality to a later stage. Hence \[
\operatorname{colim}E_i\cong I/I^2,\qquad \operatorname{colim}F_i\cong F.
\] These arguments do not assume the conormal modules are finitely presented or that nonflat base change preserves their kernels.

The original differential is \[
d_i([f_{j,i}])=\sum_{k=1}^n\frac{\partial f_{j,i}}{\partial x_k}\,dx_k.
\] Differentiation commutes with coefficient maps. The \([f_{j,i}]\) generate \(E_i\), and the stated \(dx_k\) form the basis of \(F_i\). If \(B\to C\) is étale, its completed presentation criterion gives an isomorphism \(d:E\to F\). Set \(\xi_k=d^{-1}(dx_k)\). Lift this finite list to a stage and then annihilate at a common later stage every error in the two finite lists \[
d_i(\xi_{k,i})=dx_k,\qquad
\sum_{k=1}^n\frac{\partial f_{j,i}}{\partial x_k}\xi_{k,i}=[f_{j,i}].
\] Define \(h_i:F_i\to E_i\) by \(dx_k\mapsto\xi_{k,i}\). The first list gives \(d_i h_i=1\) on every basis vector; the second gives \(h_i d_i=1\) on every generator. Linearity proves both identities on the entire modules. Thus \(d_i\) is an isomorphism, and the unchanged finite presentation gives étaleness.

For smoothness, take a left inverse \(h:F\to E\) of the split injection \(d\), and lift \(\xi_k=h(dx_k)\). Kill just the second list of errors. Then \(h_i d_i=1\), so \(d_i\) is split injective and the source's finite-presentation smoothness criterion applies. No right inverse is required or asserted. Empty variable or relation lists obey the same equations: the remaining identities give precisely the stated splitting or inverse. This includes zero algebras.

\subsubsection{Relative global complete intersections and the finite cover}\label{algebra-proof-073-relative-global-complete-intersections-and-the-finite-cover}

For a relative global complete intersection \(B\to C\), retain the source's finite-type \(\mathbf Z\)-algebra \(R\), its map to \(B\), the actual relative global complete-intersection model \(R\to S\), and \(B\otimes_RS\cong C\). The earlier syntomic note proves this model theorem using all coefficients of the original equations and of the exact unit relation over dimension charts. Its base-change proof keeps all variables, equations and their number difference, including empty spectra.

Because \(R\) is finitely presented over \(\mathbf Z\), lift the finite generator images to one \(B_i\), then kill all defining relation errors. Put \(D_i=B_i\otimes_RS\), a finitely presented \(B_i\)-algebra. Lift the given map \(D_i\to C\) by its finite generator images and relations to a later \(C_i\), retaining its \(B_i\)-algebra structure.

The resulting map \(D_i\to C_i\) is itself finitely presented. Write \(D_i=B_i[Y]/(g)\), \(C_i=B_i[X]/(h)\), both finite presentations, and represent the images of \(Y_j\) by \(F_j(X)\). Then \[
C_i\cong D_i[X]/(h_1,\ldots,h_s,\ Y_1-F_1(X),\ldots,Y_t-F_t(X)).
\] The inverse maps send every original variable to its specified element. The original relations \(g(F(X))\) vanish modulo \((h)\) because the map already exists, so these formulas respect all relations. Apply the just-proved colimit isomorphism lemma to this map and the original base system. At a later stage it is an isomorphism. The map \(B_i\to D_i\) is the actual base change of the original \(R\to S\), so the relative global complete-intersection conclusion follows.

For syntomic descent retain \(g_1,\ldots,g_m\in C\), the original identity \(\sum_j a_jg_j=1\), and each relative global complete-intersection chart \(B\to C_{g_j}\). Lift all \(g_j,a_j\), then kill the unit-relation error. Keep the localization presentation \[
(C_i)_{g_{j,i}}=C_i[T]/(g_{j,i}T-1),
\] with inverse element \(T=g_{j,i}^{-1}\). Each map from \(B_i\) to this localization is finitely presented. The preceding result descends each of these finitely many charts; choose a common later stage. Their original unit relation still holds, so the established principal-cover syntomic criterion applies. Flatness and the fibre complete-intersection condition are checked on exactly those charts; finite presentation is already retained.

\subsubsection{The original algebra maps into the cover}\label{algebra-proof-073-the-original-algebra-maps-into-the-cover}

For the final theorem retain the original faithfully flat finitely presented \(R\to S\). If \(R=0\), unitality forces \(S=0\); choose \(S'=0\). Otherwise the preceding faithful-flat descent construction gives \(R_0\to S_0\), with \(R_0\) finite type over \(\mathbf Z\) and \(R\otimes_{R_0}S_0\cong S\). A cover \(S_0\to T_0\) of the required kind gives \[
S=R\otimes_{R_0}S_0\longrightarrow R\otimes_{R_0}T_0.
\] Finite presentation and flatness persist by the actual base-change presentations and tensor comparisons. For an \(R\)-module \(M\), \(M\otimes_R(R\otimes_{R_0}T_0)\cong M\otimes_{R_0}T_0\); original faithful flatness reflects its vanishing. Quasi-finiteness persists by the source residue-field fibre comparison. Thus it suffices to construct the cover over the Noetherian model. Use the source names \(R,S\) for that model below, but retain the map from the initial algebra throughout.

Let \(W\) be the full locus where the fibre local ring is Cohen--Macaulay. It is the union of the corrected fixed-relative-dimension opens from the earlier CM-locus proof and therefore open. It meets every nonempty fibre: a minimal prime of a finite-type fibre ring has zero-dimensional Noetherian local ring, hence Cohen--Macaulay. Faithfulness makes every fibre nonempty. Principal opens contained in \(W\) have open images in the base by flat finite-presentation openness, and these images cover the base. Quasi-compactness gives finitely many \(D(g_j)\). Put \[
\widetilde S=\prod_{j=1}^m S_{g_j},\qquad
S\longrightarrow\widetilde S,\quad s\longmapsto(s/1)_j.
\] This is the original map into a replacement algebra, not an identification with the product. Its base map is flat of finite presentation, surjective on spectra, and all its local fibres are Cohen--Macaulay. Finite products are finite direct sums as modules and have disjoint-union spectra, proving flatness and the local assertions; finite presentation is verified below. As in the source, now use \(S\) for this product while retaining the original map into it.

Fix \(\mathfrak p\subset R\). Choose a maximal ideal of the nonzero finite-type fibre \(S\otimes_R\kappa(\mathfrak p)\), with corresponding \(\mathfrak q\subset S\). Its residue field is finite over \(\kappa(\mathfrak p)\) by the finite-type field theorem. The nonzero local ring \[
\overline S_{\mathfrak q}=S_{\mathfrak q}/\mathfrak pS_{\mathfrak q}
\] is Cohen--Macaulay of finite dimension \(d\). Choose a parameter regular sequence \(\overline f_1,\ldots,\overline f_d\) and actual lifts \(f_j\in\mathfrak q S_{\mathfrak q}\). The full fibre ideal is primary for the maximal ideal. When \(d=0\), the sequence is empty and the fibre is already zero-dimensional.

Write \(f_j=a_j/u_j\) with \(u_j\notin\mathfrak q\), and take \(g=\prod_j u_j\), with \(g=1\) for the empty list. In \(S_g\) the same fractions are \(a_j\prod_{\ell\ne j}u_\ell/g\); these original unit factors are retained. Define \[
S'=S_g/(f_1,\ldots,f_d),\qquad
\mathfrak q'=\mathfrak qS_g/(f_1,\ldots,f_d).
\] Each \(f_j\) lies in \(\mathfrak qS_g\), because contraction of \(\mathfrak qS_{\mathfrak q}\) to \(S_g\) is \(\mathfrak qS_g\). Thus \(\mathfrak q'\) is a prime over \(\mathfrak p\), with the original point correspondence.

\subsubsection{The redundant localization and the complete denominator}\label{algebra-proof-073-the-redundant-localization-and-the-complete-denominator}

The printed map \(S_g\to S'_g\) at 47851 is valid. The image \(\bar g\in S'\) already has inverse given by the image of \(g^{-1}\). The map \(j:S'\to S'_g\), \(x\mapsto x/1\), has exact inverse \[
\rho:S'_g\longrightarrow S',\qquad x/\bar g^n\longmapsto x(\bar g^{-1})^n.
\] An equality of fractions is witnessed by a power of \(\bar g\) killing their cross-multiplication difference; invertibility proves their displayed images equal. Both composites are the identity by the same powers. Thus the printed map is the quotient followed by this isomorphism and gives precisely the intended prime correspondence. OCC-12454 proposes removal of a redundant localization, not repair of a defect; the source is retained.

The actual local fibre is \[
S'_{\mathfrak q'}/\mathfrak pS'_{\mathfrak q'}
\cong S_{\mathfrak q}/\big((f_1,\ldots,f_d)+\mathfrak pS_{\mathfrak q}\big)
\cong\overline S_{\mathfrak q}/(\overline f_1,\ldots,\overline f_d).
\] For the localization comparison, the image of \((a/g^n)/(b/g^t)\), \(b\notin\mathfrak q\), maps to the class of \(ag^t/(bg^n)\) in the original local quotient. Its inverse uses the same original fractions in the quotient localization. Quotienting by \(\mathfrak p\) gives exactly the sum of the two displayed ideals. The second map first reduces an original \(x\in S_{\mathfrak q}\) modulo \(\mathfrak p\), then modulo the fibre sequence; its kernel is the same sum ideal, proving the inverse quotient comparison. Unit 0744 restores the parentheses enclosing that entire denominator. Adding an ideal after forming the first quotient does not express this map.

The final quotient is nonzero and zero-dimensional because the sequence is a system of parameters in the original Cohen--Macaulay local ring. Its residue field remains the original \(\kappa(\mathfrak q)\), finite over \(\kappa(\mathfrak p)\). These two conditions give quasi-finiteness at \(\mathfrak q'\) by the isolated-fibre criterion. Zero local dimension alone would not suffice without the chosen closed fibre point. Unit 0743 therefore correctly changes the point at 47852 from \(\mathfrak q\) to \(\mathfrak q'\).

\subsubsection{Flatness, neighbourhoods, products and the final diagram}\label{algebra-proof-073-flatness-neighbourhoods-products-and-the-final-diagram}

The original map \(R_{\mathfrak p}\to S_{\mathfrak q}\) is flat local between Noetherian local rings. Apply the preceding completed regular-sequence lifting criterion to every actual \(f_j\): multiplication by its image is injective on the appropriate fibre quotient, the local injection criterion lifts the injection, and it proves that the next quotient is flat over \(R_{\mathfrak p}\). Each quotient remains nonzero because its fibre quotient is nonzero. Induction gives \(R_{\mathfrak p}\to S'_{\mathfrak q'}\) flat, hence flatness over \(R\) after the localization comparison. If \(d=0\), this is simply the original local flat map.

The algebra \(S'\) is finitely presented over \(R\): invert \(g\) by \(gZ-1\) and impose the finite list of \(f_j\), written as the exact polynomials \(a_j\prod_{\ell\ne j}u_\ell Z\). The flat and quasi-finite loci are both open and contain \(\mathfrak q'\). Choose a principal \(D(h)\) inside their intersection and containing that point. Lift \(h\) to \(a/g^n\in S_g\). Since \(h\notin\mathfrak q'\), \(a\notin\mathfrak q\). There is an exact isomorphism \[
(S')_h\cong S_{ga}/(f_1,\ldots,f_d).
\] One map sends each original fraction to the right-hand quotient; the inverse uses that \(g\) is already invertible and that inverting \(a/g^n\) is equivalent to inverting \(a\). Its inverse element for \(a\) is \(g^{-n}h^{-1}\). Both maps fix every original fraction and relation, giving the source's replacement \(g\mapsto gg'\) with \(g'=a\). The resulting flat quasi-finite finitely presented algebra has open image containing the original \(\mathfrak p\).

Perform this construction for every base prime, and choose finitely many resulting open images covering the quasi-compact base. Denote their actual algebras by \(T_1,\ldots,T_N\) and put \(T=\prod_jT_j\). All maps from the original \(S\), including its earlier CM product map and the initial model map, compose to the requested \(S\to T\). The map \(R\to T\) is flat as a finite direct sum of flat modules and is surjective on spectra by the chosen cover, hence faithfully flat. Every prime of a finite product lies in exactly one factor and has that factor's local ring, so quasi-finiteness holds at every prime.

Here are full equations for finite presentation of the product. Write \(T_j=R[X_{j,1},\ldots,X_{j,n_j}]/(h_{j,1},\ldots,h_{j,m_j})\). Introduce \(e_1,\ldots,e_N\) and these same variable lists, with the finite relations \[
e_j^2=e_j,\quad e_je_\ell=0\ (j\ne\ell),\quad
\sum_j e_j=1,\quad e_jX_{j,a}=X_{j,a},\quad
e_jh_{j,b}(X_{j,1},\ldots,X_{j,n_j})=0.
\] Map \(e_j\) to the actual coordinate idempotent and \(X_{j,a}\) to its original generator in that coordinate and zero elsewhere. Conversely these relations decompose the presented ring into its \(e_j\)-parts, each with unit \(e_j\). A coefficient \(r\) in that part means \(re_j\), and its polynomial relations are exactly the original \(h_{j,b}\); the remaining variable lists vanish there. This constructs the inverse maps to \(\prod_jT_j\), proves finite presentation and retains all factors. The already handled zero-base case covers an empty spectrum. The final diagram has exactly the three source properties.

\subsubsection{Report accounting and the chapter boundary}\label{algebra-proof-073-report-accounting-and-the-chapter-boundary}

OCC-01036 adds ``be'' before ``a map'' in the seven declarations at 47524, 47549, 47573, 47601, 47628, 47668 and 47708. The directed set, element \(0\), maps and hypotheses remain unchanged. OCC-12454 is rejected as a nondefect by the exact redundant-localization inverse. OCC-01037 and OCC-12455 retain 0743; OCC-01038 and OCC-12456 retain 0744. The chapter input, bibliography commands and terminator at 47880--47885 require no report disposition.

The supporting arguments receive the earlier finite-stage witness constructions, conormal criteria, complete-intersection models, corrected CM and flat loci, regular-sequence quotient flatness and quasi-finite fibre comparisons. No new extension group is created. The next numeric position is the verified one-past-end 47886. The cumulative replay must prove all 1,351 reports have exactly one disposition. That accounting alone does not constitute candidate admission, a successful build, visual review or publication.


\bibliographystyle{alpha}
\bibliography{my}
\end{document}
